Senior Member, IEEE

Department of Computer Engineering

San Jose State University

\href{mailto:andrew.bond@sjsu.edu}{\nolinkurl{andrew.bond@sjsu.edu}}

Spring 2026

\includegraphics[width=5.49375in,height=4.91597in,alt={image}]{media/image1.png}

\section{Table of Contents}\label{table-of-contents-1}

\hyperref[table-of-contents-1]{Table of Contents
\hyperref[table-of-contents-1]{2}}

\hyperref[table-of-figures]{Table of Figures
\hyperref[table-of-figures]{31}}

\hyperref[preface]{Preface \hyperref[preface]{32}}

\hyperref[epistemic-status-tags]{Epistemic Status Tags
\hyperref[epistemic-status-tags]{35}}

\hyperref[epistemic-status-classification-of-major-claims]{Epistemic
Status Classification of Major Claims
\hyperref[epistemic-status-classification-of-major-claims]{36}}

\hyperref[a-final-word-about-urgency]{A final word about urgency
\hyperref[a-final-word-about-urgency]{37}}

\hyperref[core-objects-at-a-glance-1]{Core Objects at a Glance
\hyperref[core-objects-at-a-glance-1]{38}}

\hyperref[key-results-at-a-glance]{Key Results at a Glance
\hyperref[key-results-at-a-glance]{39}}

\hyperref[part-i-the-problem]{Part I: The Problem
\hyperref[part-i-the-problem]{40}}

\hyperref[chapter-1-introduction-why-geometry]{Chapter 1: Introduction
--- Why Geometry? \hyperref[chapter-1-introduction-why-geometry]{40}}

\hyperref[the-shape-of-the-problem]{The Shape of the Problem
\hyperref[the-shape-of-the-problem]{40}}

\hyperref[an-old-man-and-a-horse]{An Old Man and a Horse
\hyperref[an-old-man-and-a-horse]{40}}

\hyperref[three-failures-of-flatland]{Three Failures of Flatland
\hyperref[three-failures-of-flatland]{41}}

\hyperref[failure-1-ai-alignment]{Failure 1: AI Alignment
\hyperref[failure-1-ai-alignment]{41}}

\hyperref[failure-2-policy-analysis]{Failure 2: Policy Analysis
\hyperref[failure-2-policy-analysis]{42}}

\hyperref[failure-3-moral-philosophy-itself]{Failure 3: Moral Philosophy
Itself \hyperref[failure-3-moral-philosophy-itself]{43}}

\hyperref[what-geometry-provides]{What Geometry Provides
\hyperref[what-geometry-provides]{43}}

\hyperref[what-this-book-is-not]{What This Book Is Not
\hyperref[what-this-book-is-not]{44}}

\hyperref[the-epistemic-stance-pragmatist-geometry]{The Epistemic
Stance: Pragmatist Geometry
\hyperref[the-epistemic-stance-pragmatist-geometry]{45}}

\hyperref[how-to-read-this-book]{How to Read This Book
\hyperref[how-to-read-this-book]{47}}

\hyperref[the-suspicious-coincidence]{The Suspicious Coincidence
\hyperref[the-suspicious-coincidence]{47}}

\hyperref[the-arc-of-the-book]{The Arc of the Book
\hyperref[the-arc-of-the-book]{49}}

\hyperref[a-note-on-ambition]{A Note on Ambition
\hyperref[a-note-on-ambition]{50}}

\hyperref[chapter-2-the-failure-of-scalar-ethics]{Chapter 2: The Failure
of Scalar Ethics \hyperref[chapter-2-the-failure-of-scalar-ethics]{52}}

\hyperref[the-parable-of-the-old-man-and-his-horse]{The Parable of the
Old Man and His Horse
\hyperref[the-parable-of-the-old-man-and-his-horse]{52}}

\hyperref[the-insufficiency-of-rank-0-ethics]{The Insufficiency of
Rank-0 Ethics \hyperref[the-insufficiency-of-rank-0-ethics]{53}}

\hyperref[limitation-1-no-directional-information]{Limitation 1: No
Directional Information
\hyperref[limitation-1-no-directional-information]{53}}

\hyperref[limitation-2-uncertainty-has-shape]{Limitation 2: Uncertainty
Has Shape \hyperref[limitation-2-uncertainty-has-shape]{54}}

\hyperref[limitation-3-paths-cross-boundaries]{Limitation 3: Paths Cross
Boundaries \hyperref[limitation-3-paths-cross-boundaries]{55}}

\hyperref[what-geometric-structure-provides]{What Geometric Structure
Provides \hyperref[what-geometric-structure-provides]{56}}

\hyperref[gradients-and-the-direction-of-moral-change]{Gradients and the
Direction of Moral Change
\hyperref[gradients-and-the-direction-of-moral-change]{56}}

\hyperref[the-metric-tensor-and-moral-comparison]{The Metric Tensor and
Moral Comparison \hyperref[the-metric-tensor-and-moral-comparison]{57}}

\hyperref[stratification-and-moral-phase-transitions]{Stratification and
Moral Phase Transitions
\hyperref[stratification-and-moral-phase-transitions]{58}}

\hyperref[the-parable-revisited]{The Parable Revisited
\hyperref[the-parable-revisited]{58}}

\hyperref[why-maybe-is-geometric-not-merely-epistemic]{Why ``Maybe'' Is
Geometric, Not Merely Epistemic
\hyperref[why-maybe-is-geometric-not-merely-epistemic]{59}}

\hyperref[from-parable-to-framework]{From Parable to Framework
\hyperref[from-parable-to-framework]{60}}

\hyperref[anticipating-objections]{Anticipating Objections
\hyperref[anticipating-objections]{61}}

\hyperref[isnt-this-just-saying-ethics-is-complicated]{``Isn't this just
saying ethics is complicated?''
\hyperref[isnt-this-just-saying-ethics-is-complicated]{61}}

\hyperref[we-cant-measure-moral-quantities-precisely-so-what-use-is-this-formalism]{``We
can't measure moral quantities precisely, so what use is this
formalism?''
\hyperref[we-cant-measure-moral-quantities-precisely-so-what-use-is-this-formalism]{61}}

\hyperref[conclusion]{Conclusion \hyperref[conclusion]{61}}

\hyperref[chapter-3-historical-precursors-geometry-before-geometry]{Chapter
3: Historical Precursors --- Geometry Before Geometry
\hyperref[chapter-3-historical-precursors-geometry-before-geometry]{63}}

\hyperref[introduction-structure-before-formalism]{Introduction:
Structure Before Formalism
\hyperref[introduction-structure-before-formalism]{63}}

\hyperref[aristotle-the-doctrine-of-the-mean-as-context-sensitive-calibration]{Aristotle:
The Doctrine of the Mean as Context-Sensitive Calibration
\hyperref[aristotle-the-doctrine-of-the-mean-as-context-sensitive-calibration]{63}}

\hyperref[geometric-reading]{Geometric Reading
\hyperref[geometric-reading]{64}}

\hyperref[the-doctrine-of-the-mean-as-a-metric-condition]{The Doctrine
of the Mean as a Metric Condition
\hyperref[the-doctrine-of-the-mean-as-a-metric-condition]{64}}

\hyperref[kant-the-categorical-imperative-as-an-invariance-condition]{Kant:
The Categorical Imperative as an Invariance Condition
\hyperref[kant-the-categorical-imperative-as-an-invariance-condition]{65}}

\hyperref[geometric-reading-1]{Geometric Reading
\hyperref[geometric-reading-1]{65}}

\hyperref[the-kingdom-of-ends-as-a-symmetry-group]{The Kingdom of Ends
as a Symmetry Group
\hyperref[the-kingdom-of-ends-as-a-symmetry-group]{65}}

\hyperref[ross-prima-facie-duties-as-vector-components]{Ross: Prima
Facie Duties as Vector Components
\hyperref[ross-prima-facie-duties-as-vector-components]{66}}

\hyperref[geometric-reading-2]{Geometric Reading
\hyperref[geometric-reading-2]{66}}

\hyperref[the-interaction-problem]{The Interaction Problem
\hyperref[the-interaction-problem]{67}}

\hyperref[rawls-the-original-position-as-a-symmetry-constraint]{Rawls:
The Original Position as a Symmetry Constraint
\hyperref[rawls-the-original-position-as-a-symmetry-constraint]{67}}

\hyperref[geometric-reading-3]{Geometric Reading
\hyperref[geometric-reading-3]{67}}

\hyperref[the-metric-of-the-original-position]{The Metric of the
Original Position \hyperref[the-metric-of-the-original-position]{68}}

\hyperref[sen-and-nussbaum-capabilities-as-a-basis-for-moral-space]{Sen
and Nussbaum: Capabilities as a Basis for Moral Space
\hyperref[sen-and-nussbaum-capabilities-as-a-basis-for-moral-space]{68}}

\hyperref[geometric-reading-4]{Geometric Reading
\hyperref[geometric-reading-4]{69}}

\hyperref[the-incompleteness-thesis]{The Incompleteness Thesis
\hyperref[the-incompleteness-thesis]{69}}

\hyperref[nussbaums-threshold-and-stratum-structure]{Nussbaum's
Threshold and Stratum Structure
\hyperref[nussbaums-threshold-and-stratum-structure]{69}}

\hyperref[hohfeld-jural-relations-as-discrete-structure]{Hohfeld: Jural
Relations as Discrete Structure
\hyperref[hohfeld-jural-relations-as-discrete-structure]{70}}

\hyperref[geometric-reading-5]{Geometric Reading
\hyperref[geometric-reading-5]{70}}

\hyperref[moral-uncertainty-superposition-and-structured-credence]{Moral
Uncertainty: Superposition and Structured Credence
\hyperref[moral-uncertainty-superposition-and-structured-credence]{71}}

\hyperref[geometric-reading-6]{Geometric Reading
\hyperref[geometric-reading-6]{71}}

\hyperref[moral-hedging-as-covariance-sensitivity]{Moral Hedging as
Covariance Sensitivity
\hyperref[moral-hedging-as-covariance-sensitivity]{72}}

\hyperref[synthesis-what-the-precursors-share]{Synthesis: What the
Precursors Share \hyperref[synthesis-what-the-precursors-share]{72}}

\hyperref[what-the-geometric-framework-adds]{What the Geometric
Framework Adds \hyperref[what-the-geometric-framework-adds]{73}}

\hyperref[conclusion-the-geometry-was-always-there]{Conclusion: The
Geometry Was Always There
\hyperref[conclusion-the-geometry-was-always-there]{74}}

\hyperref[part-ii-the-framework]{Part II: The Framework
\hyperref[part-ii-the-framework]{76}}

\hyperref[chapter-4-mathematical-preliminaries]{Chapter 4: Mathematical
Preliminaries \hyperref[chapter-4-mathematical-preliminaries]{76}}

\hyperref[a-note-on-epistemic-status]{A Note on Epistemic Status
\hyperref[a-note-on-epistemic-status]{76}}

\hyperref[manifolds]{4.1 Manifolds \hyperref[manifolds]{77}}

\hyperref[the-idea]{The Idea \hyperref[the-idea]{77}}

\hyperref[formal-definition]{Formal Definition
\hyperref[formal-definition]{77}}

\hyperref[example-the-circle]{Example: The Circle
\hyperref[example-the-circle]{77}}

\hyperref[example-the-2-simplex]{Example: The 2-Simplex
\hyperref[example-the-2-simplex]{78}}

\hyperref[tangent-spaces-and-tangent-bundles]{4.2 Tangent Spaces and
Tangent Bundles \hyperref[tangent-spaces-and-tangent-bundles]{78}}

\hyperref[tangent-vectors]{Tangent Vectors
\hyperref[tangent-vectors]{78}}

\hyperref[the-tangent-bundle]{The Tangent Bundle
\hyperref[the-tangent-bundle]{78}}

\hyperref[transformation-of-tangent-vectors]{Transformation of Tangent
Vectors \hyperref[transformation-of-tangent-vectors]{78}}

\hyperref[covectors-dual-spaces-and-the-cotangent-bundle]{4.3 Covectors,
Dual Spaces, and the Cotangent Bundle
\hyperref[covectors-dual-spaces-and-the-cotangent-bundle]{79}}

\hyperref[covectors]{Covectors \hyperref[covectors]{79}}

\hyperref[the-gradient]{The Gradient \hyperref[the-gradient]{79}}

\hyperref[tensors]{4.4 Tensors \hyperref[tensors]{80}}

\hyperref[definition]{Definition \hyperref[definition]{80}}

\hyperref[transformation-law]{Transformation Law
\hyperref[transformation-law]{80}}

\hyperref[contraction]{Contraction \hyperref[contraction]{80}}

\hyperref[tensor-fields]{Tensor Fields \hyperref[tensor-fields]{81}}

\hyperref[the-metric-tensor]{4.5 The Metric Tensor
\hyperref[the-metric-tensor]{81}}

\hyperref[definition-1]{Definition \hyperref[definition-1]{81}}

\hyperref[riemannian-vs.-semi-riemannian-vs.-degenerate]{Riemannian
vs.~Semi-Riemannian vs.~Degenerate
\hyperref[riemannian-vs.-semi-riemannian-vs.-degenerate]{81}}

\hyperref[the-metric-as-musical-isomorphism]{The Metric as Musical
Isomorphism \hyperref[the-metric-as-musical-isomorphism]{81}}

\hyperref[connections-parallel-transport-and-curvature]{4.6 Connections,
Parallel Transport, and Curvature
\hyperref[connections-parallel-transport-and-curvature]{82}}

\hyperref[the-problem-of-comparison]{The Problem of Comparison
\hyperref[the-problem-of-comparison]{82}}

\hyperref[the-connection]{The Connection \hyperref[the-connection]{82}}

\hyperref[parallel-transport]{Parallel Transport
\hyperref[parallel-transport]{83}}

\hyperref[the-riemann-curvature-tensor]{The Riemann Curvature Tensor
\hyperref[the-riemann-curvature-tensor]{83}}

\hyperref[fiber-bundles]{4.7 Fiber Bundles \hyperref[fiber-bundles]{84}}

\hyperref[motivation]{Motivation \hyperref[motivation]{84}}

\hyperref[definition-2]{Definition \hyperref[definition-2]{84}}

\hyperref[sections]{Sections \hyperref[sections]{85}}

\hyperref[stratified-spaces]{4.8 Stratified Spaces
\hyperref[stratified-spaces]{85}}

\hyperref[beyond-smooth-manifolds]{Beyond Smooth Manifolds
\hyperref[beyond-smooth-manifolds]{85}}

\hyperref[definition-3]{Definition \hyperref[definition-3]{85}}

\hyperref[example-the-simplex]{Example: The Simplex
\hyperref[example-the-simplex]{86}}

\hyperref[moral-significance]{Moral Significance
\hyperref[moral-significance]{86}}

\hyperref[gauge-theory-a-brief-introduction]{4.9 Gauge Theory: A Brief
Introduction \hyperref[gauge-theory-a-brief-introduction]{86}}

\hyperref[motivation-1]{Motivation \hyperref[motivation-1]{86}}

\hyperref[gauge-transformations]{Gauge Transformations
\hyperref[gauge-transformations]{86}}

\hyperref[connections-on-fiber-bundles]{Connections on Fiber Bundles
\hyperref[connections-on-fiber-bundles]{87}}

\hyperref[curvature-of-a-connection]{Curvature of a Connection
\hyperref[curvature-of-a-connection]{87}}

\hyperref[noethers-theorem-preview]{Noether's Theorem (Preview)
\hyperref[noethers-theorem-preview]{87}}

\hyperref[summary-the-geometric-toolkit]{4.10 Summary: The Geometric
Toolkit \hyperref[summary-the-geometric-toolkit]{88}}

\hyperref[bibliographic-notes]{4.11 Bibliographic Notes
\hyperref[bibliographic-notes]{89}}

\hyperref[chapter-5-the-moral-manifold]{Chapter 5: The Moral Manifold
\hyperref[chapter-5-the-moral-manifold]{90}}

\hyperref[the-question-of-base-space]{5.1 The Question of Base Space
\hyperref[the-question-of-base-space]{90}}

\hyperref[what-are-the-points]{5.2 What Are the Points?
\hyperref[what-are-the-points]{90}}

\hyperref[candidate-1-possible-actions]{Candidate 1: Possible Actions
\hyperref[candidate-1-possible-actions]{90}}

\hyperref[candidate-2-states-of-affairs]{Candidate 2: States of Affairs
\hyperref[candidate-2-states-of-affairs]{91}}

\hyperref[candidate-3-situations]{Candidate 3: Situations
\hyperref[candidate-3-situations]{91}}

\hyperref[candidate-4-agent-situation-pairs]{Candidate 4:
Agent-Situation Pairs \hyperref[candidate-4-agent-situation-pairs]{91}}

\hyperref[our-choice-structured-situations]{Our Choice: Structured
Situations \hyperref[our-choice-structured-situations]{92}}

\hyperref[the-dimension-of-moral-space]{5.3 The Dimension of Moral Space
\hyperref[the-dimension-of-moral-space]{92}}

\hyperref[the-dimensionality-question]{The Dimensionality Question
\hyperref[the-dimensionality-question]{92}}

\hyperref[the-3-3-derivation]{The 3 × 3 Derivation
\hyperref[the-3-3-derivation]{93}}

\hyperref[arguments-for-necessity-and-sufficiency]{Arguments for
Necessity and Sufficiency
\hyperref[arguments-for-necessity-and-sufficiency]{94}}

\hyperref[formal-statement]{Formal Statement
\hyperref[formal-statement]{94}}

\hyperref[the-varying-metric]{The Varying Metric
\hyperref[the-varying-metric]{95}}

\hyperref[the-differential-structure]{5.4 The Differential Structure
\hyperref[the-differential-structure]{95}}

\hyperref[charts-and-atlases-on-m]{Charts and Atlases on M
\hyperref[charts-and-atlases-on-m]{95}}

\hyperref[the-tangent-space-of-m]{The Tangent Space of M
\hyperref[the-tangent-space-of-m]{95}}

\hyperref[the-agent-fiber-bundle]{The Agent Fiber Bundle
\hyperref[the-agent-fiber-bundle]{96}}

\hyperref[coordinates-and-admissible-transformations]{5.5 Coordinates
and Admissible Transformations
\hyperref[coordinates-and-admissible-transformations]{96}}

\hyperref[the-problem-of-multiple-descriptions]{The Problem of Multiple
Descriptions \hyperref[the-problem-of-multiple-descriptions]{96}}

\hyperref[type-1-coordinate-redescriptions]{Type 1: Coordinate
Redescriptions \hyperref[type-1-coordinate-redescriptions]{97}}

\hyperref[pre-flight-checklist-bond-invariance-principle]{Pre-flight
Checklist: Bond Invariance Principle
\hyperref[pre-flight-checklist-bond-invariance-principle]{97}}

\hyperref[type-2-perspective-shifts]{Type 2: Perspective Shifts
\hyperref[type-2-perspective-shifts]{98}}

\hyperref[type-3-theory-shifts]{Type 3: Theory Shifts
\hyperref[type-3-theory-shifts]{98}}

\hyperref[summary-the-transformation-hierarchy]{Summary: The
Transformation Hierarchy
\hyperref[summary-the-transformation-hierarchy]{99}}

\hyperref[admissibility-in-practice-a-checklist]{Admissibility in
Practice: A Checklist
\hyperref[admissibility-in-practice-a-checklist]{99}}

\hyperref[tempting-but-invalid-transformations]{Tempting but Invalid
Transformations \hyperref[tempting-but-invalid-transformations]{100}}

\hyperref[stratification]{5.6 Stratification
\hyperref[stratification]{101}}

\hyperref[why-stratification-is-essential]{Why Stratification Is
Essential \hyperref[why-stratification-is-essential]{101}}

\hyperref[strata-boundaries-and-transitions]{Strata, Boundaries, and
Transitions \hyperref[strata-boundaries-and-transitions]{102}}

\hyperref[nullifiers-as-absorbing-states]{Nullifiers as Absorbing States
\hyperref[nullifiers-as-absorbing-states]{102}}

\hyperref[semantic-gates-as-transition-functions]{Semantic Gates as
Transition Functions
\hyperref[semantic-gates-as-transition-functions]{103}}

\hyperref[singularities-and-dilemmas]{5.7 Singularities and Dilemmas
\hyperref[singularities-and-dilemmas]{103}}

\hyperref[definition-and-examples]{Definition and Examples
\hyperref[definition-and-examples]{103}}

\hyperref[dilemmas-as-features-not-bugs]{Dilemmas as Features, Not Bugs
\hyperref[dilemmas-as-features-not-bugs]{104}}

\hyperref[topology-of-the-moral-manifold]{5.8 Topology of the Moral
Manifold \hyperref[topology-of-the-moral-manifold]{104}}

\hyperref[local-and-global-structure]{Local and Global Structure
\hyperref[local-and-global-structure]{104}}

\hyperref[two-worked-examples]{5.9 Two Worked Examples
\hyperref[two-worked-examples]{105}}

\hyperref[example-1-the-manifold-of-a-medical-allocation]{Example 1: The
Manifold of a Medical Allocation
\hyperref[example-1-the-manifold-of-a-medical-allocation]{105}}

\hyperref[example-2-the-promise-landscape]{Example 2: The Promise
Landscape \hyperref[example-2-the-promise-landscape]{105}}

\hyperref[the-manifold-as-common-ground]{5.10 The Manifold as Common
Ground \hyperref[the-manifold-as-common-ground]{106}}

\hyperref[technical-appendix-formal-definitions]{Technical Appendix:
Formal Definitions
\hyperref[technical-appendix-formal-definitions]{107}}

\hyperref[chapter-6-the-tensor-hierarchy]{Chapter 6: The Tensor
Hierarchy \hyperref[chapter-6-the-tensor-hierarchy]{109}}

\hyperref[from-scalars-to-tensors-why-rank-matters]{6.1 From Scalars to
Tensors: Why Rank Matters
\hyperref[from-scalars-to-tensors-why-rank-matters]{109}}

\hyperref[obligations-as-vectors]{6.2 Obligations as Vectors
\hyperref[obligations-as-vectors]{112}}

\hyperref[the-directional-character-of-duty]{The Directional Character
of Duty \hyperref[the-directional-character-of-duty]{112}}

\hyperref[rank-progression-from-rank-1-obligation-vectors-to-rank-6-distributional-tensors.]{Rank
progression: from rank-1 obligation vectors to rank-6 distributional
tensors.
\hyperref[rank-progression-from-rank-1-obligation-vectors-to-rank-6-distributional-tensors.]{113}}

\hyperref[_Toc224082185]{Transformation Behavior
\hyperref[_Toc224082185]{113}}

\hyperref[the-zero-vector-and-its-significance]{The Zero Vector and Its
Significance \hyperref[the-zero-vector-and-its-significance]{114}}

\hyperref[multiple-obligation-sources]{Multiple Obligation Sources
\hyperref[multiple-obligation-sources]{114}}

\hyperref[interests-as-covectors]{6.3 Interests as Covectors
\hyperref[interests-as-covectors]{115}}

\hyperref[the-dual-perspective]{The Dual Perspective
\hyperref[the-dual-perspective]{115}}

\hyperref[why-covectors-not-vectors]{Why Covectors, Not Vectors?
\hyperref[why-covectors-not-vectors]{115}}

\hyperref[the-gradient-as-a-natural-covector]{The Gradient as a Natural
Covector \hyperref[the-gradient-as-a-natural-covector]{115}}

\hyperref[the-fundamental-contraction-satisfaction]{6.4 The Fundamental
Contraction: Satisfaction
\hyperref[the-fundamental-contraction-satisfaction]{116}}

\hyperref[the-formula]{The Formula \hyperref[the-formula]{116}}

\hyperref[unpacking-the-formula]{Unpacking the Formula
\hyperref[unpacking-the-formula]{116}}

\hyperref[the-moral-content-of-contraction]{The Moral Content of
Contraction \hyperref[the-moral-content-of-contraction]{117}}

\hyperref[contraction-is-not-unique]{Contraction Is Not Unique
\hyperref[contraction-is-not-unique]{117}}

\hyperref[the-moral-metric]{6.5 The Moral Metric
\hyperref[the-moral-metric]{117}}

\hyperref[relating-vectors-and-covectors]{Relating Vectors and Covectors
\hyperref[relating-vectors-and-covectors]{117}}

\hyperref[the-metric-as-trade-off-structure]{The Metric as Trade-off
Structure \hyperref[the-metric-as-trade-off-structure]{118}}

\hyperref[different-theories-as-different-metrics]{Different Theories as
Different Metrics
\hyperref[different-theories-as-different-metrics]{118}}

\hyperref[the-metric-is-not-given]{The Metric Is Not Given
\hyperref[the-metric-is-not-given]{119}}

\hyperref[the-full-moral-tensor]{6.6 The Full Moral Tensor
\hyperref[the-full-moral-tensor]{119}}

\hyperref[beyond-vectors-and-covectors]{Beyond Vectors and Covectors
\hyperref[beyond-vectors-and-covectors]{119}}

\hyperref[the-evaluation-tensor]{The Evaluation Tensor
\hyperref[the-evaluation-tensor]{119}}

\hyperref[the-multi-agent-evaluation-tensor]{The Multi-Agent Evaluation
Tensor \hyperref[the-multi-agent-evaluation-tensor]{120}}

\hyperref[the-uncertainty-tensor]{The Uncertainty Tensor
\hyperref[the-uncertainty-tensor]{120}}

\hyperref[scalar-ethics-as-a-special-case]{6.7 Scalar Ethics as a
Special Case \hyperref[scalar-ethics-as-a-special-case]{121}}

\hyperref[recovery-by-contraction]{Recovery by Contraction
\hyperref[recovery-by-contraction]{121}}

\hyperref[what-this-reveals]{What This Reveals
\hyperref[what-this-reveals]{122}}

\hyperref[the-hierarchy-in-action]{6.8 The Hierarchy in Action
\hyperref[the-hierarchy-in-action]{122}}

\hyperref[level-0-scalar]{Level 0: Scalar
\hyperref[level-0-scalar]{122}}

\hyperref[level-1-obligation-vector]{Level 1: Obligation Vector
\hyperref[level-1-obligation-vector]{123}}

\hyperref[level-2-interest-covector-physicians-perspective]{Level 2:
Interest Covector (Physician's Perspective)
\hyperref[level-2-interest-covector-physicians-perspective]{123}}

\hyperref[level-2-interest-covector-patient-as-family]{Level 2':
Interest Covector (Patient A's Family)
\hyperref[level-2-interest-covector-patient-as-family]{123}}

\hyperref[level-3-metric]{Level 3: Metric
\hyperref[level-3-metric]{123}}

\hyperref[what-each-level-adds]{What Each Level Adds
\hyperref[what-each-level-adds]{123}}

\hyperref[the-algebra-of-moral-tensors]{6.9 The Algebra of Moral Tensors
\hyperref[the-algebra-of-moral-tensors]{124}}

\hyperref[tensor-product]{Tensor Product \hyperref[tensor-product]{124}}

\hyperref[symmetrization-and-antisymmetrization]{Symmetrization and
Antisymmetrization
\hyperref[symmetrization-and-antisymmetrization]{124}}

\hyperref[contraction-as-moral-decision]{Contraction as Moral Decision
\hyperref[contraction-as-moral-decision]{125}}

\hyperref[summary]{6.10 Summary \hyperref[summary]{125}}

\hyperref[technical-appendix]{Technical Appendix
\hyperref[technical-appendix]{126}}

\hyperref[chapter-7-one-case-five-levels]{Chapter 7: One Case, Five
Levels \hyperref[chapter-7-one-case-five-levels]{128}}

\hyperref[introduction-the-pedagogy-of-accumulation]{Introduction: The
Pedagogy of Accumulation
\hyperref[introduction-the-pedagogy-of-accumulation]{128}}

\hyperref[level-1-the-scalar]{Level 1: The Scalar
\hyperref[level-1-the-scalar]{128}}

\hyperref[setup]{Setup \hyperref[setup]{128}}

\hyperref[the-scalar-approach]{The Scalar Approach
\hyperref[the-scalar-approach]{129}}

\hyperref[what-the-scalar-can-say]{What the Scalar Can Say
\hyperref[what-the-scalar-can-say]{129}}

\hyperref[what-the-scalar-cannot-say]{What the Scalar Cannot Say
\hyperref[what-the-scalar-cannot-say]{129}}

\hyperref[level-2-the-obligation-vector]{Level 2: The Obligation Vector
\hyperref[level-2-the-obligation-vector]{130}}

\hyperref[adding-structure]{Adding Structure
\hyperref[adding-structure]{130}}

\hyperref[what-the-vector-can-say]{What the Vector Can Say
\hyperref[what-the-vector-can-say]{130}}

\hyperref[what-this-reveals-1]{What This Reveals
\hyperref[what-this-reveals-1]{131}}

\hyperref[level-3-multiple-perspectives-rank-2-tensor]{Level 3: Multiple
Perspectives (Rank-2 Tensor)
\hyperref[level-3-multiple-perspectives-rank-2-tensor]{131}}

\hyperref[the-multi-agent-structure]{The Multi-Agent Structure
\hyperref[the-multi-agent-structure]{131}}

\hyperref[the-evaluation-tensor-1]{The Evaluation Tensor
\hyperref[the-evaluation-tensor-1]{132}}

\hyperref[what-the-rank-2-tensor-can-say]{What the Rank-2 Tensor Can Say
\hyperref[what-the-rank-2-tensor-can-say]{132}}

\hyperref[three-contractions-three-verdicts]{Three Contractions, Three
Verdicts \hyperref[three-contractions-three-verdicts]{133}}

\hyperref[the-meta-point]{The Meta-Point \hyperref[the-meta-point]{133}}

\hyperref[level-4-the-metric]{Level 4: The Metric
\hyperref[level-4-the-metric]{133}}

\hyperref[the-problem-of-comparison-1]{The Problem of Comparison
\hyperref[the-problem-of-comparison-1]{133}}

\hyperref[three-metrics-three-ethics]{Three Metrics, Three Ethics
\hyperref[three-metrics-three-ethics]{134}}

\hyperref[what-the-metric-choice-commits-you-to]{What the Metric Choice
Commits You To \hyperref[what-the-metric-choice-commits-you-to]{134}}

\hyperref[off-diagonal-components]{Off-Diagonal Components
\hyperref[off-diagonal-components]{135}}

\hyperref[level-5-stratification-and-boundaries]{Level 5: Stratification
and Boundaries \hyperref[level-5-stratification-and-boundaries]{135}}

\hyperref[beyond-the-smooth-interior]{Beyond the Smooth Interior
\hyperref[beyond-the-smooth-interior]{135}}

\hyperref[adding-a-constraint-the-discrimination-boundary]{Adding a
Constraint: The Discrimination Boundary
\hyperref[adding-a-constraint-the-discrimination-boundary]{135}}

\hyperref[adding-a-nullifier-the-abuse-boundary]{Adding a Nullifier: The
Abuse Boundary \hyperref[adding-a-nullifier-the-abuse-boundary]{136}}

\hyperref[adding-a-semantic-gate-the-emergency]{Adding a Semantic Gate:
The Emergency \hyperref[adding-a-semantic-gate-the-emergency]{136}}

\hyperref[what-stratification-adds]{What Stratification Adds
\hyperref[what-stratification-adds]{136}}

\hyperref[six-claims-that-require-tensors]{Six Claims That Require
Tensors \hyperref[six-claims-that-require-tensors]{137}}

\hyperref[claim-1-the-disagreement-between-physicians-and-families-is-orthogonal-to-the-disagreement-among-families.]{Claim
1: ``The disagreement between physicians and families is orthogonal to
the disagreement among families.''
\hyperref[claim-1-the-disagreement-between-physicians-and-families-is-orthogonal-to-the-disagreement-among-families.]{137}}

\hyperref[claim-2-alice-is-the-robust-choice-bob-wins-only-under-specific-weight-configurations.]{Claim
2: ``Alice is the robust choice; Bob wins only under specific weight
configurations.''
\hyperref[claim-2-alice-is-the-robust-choice-bob-wins-only-under-specific-weight-configurations.]{137}}

\hyperref[claim-3-choosing-carol-over-bob-wrongs-bob-in-a-way-that-choosing-alice-over-bob-does-not.]{Claim
3: ``Choosing Carol over Bob wrongs Bob in a way that choosing Alice
over Bob does not.''
\hyperref[claim-3-choosing-carol-over-bob-wrongs-bob-in-a-way-that-choosing-alice-over-bob-does-not.]{137}}

\hyperref[claim-4-a-utilitarian-and-a-rawlsian-agree-on-this-case-but-for-different-reasons-and-would-diverge-if-carols-urgency-increased.]{Claim
4: ``A utilitarian and a Rawlsian agree on this case, but for different
reasons, and would diverge if Carol's urgency increased.''
\hyperref[claim-4-a-utilitarian-and-a-rawlsian-agree-on-this-case-but-for-different-reasons-and-would-diverge-if-carols-urgency-increased.]{137}}

\hyperref[claim-5-the-metric-choice-is-doing-more-work-than-the-interest-choice.]{Claim
5: ``The metric choice is doing more work than the interest choice.''
\hyperref[claim-5-the-metric-choice-is-doing-more-work-than-the-interest-choice.]{138}}

\hyperref[claim-6-the-case-is-genuinely-hard-in-a-way-that-more-information-cannot-resolve.]{Claim
6: ``The case is genuinely hard in a way that more information cannot
resolve.''
\hyperref[claim-6-the-case-is-genuinely-hard-in-a-way-that-more-information-cannot-resolve.]{138}}

\hyperref[synthesis]{Synthesis \hyperref[synthesis]{138}}

\hyperref[coda-the-same-case-revisited-forever]{Coda: The Same Case,
Revisited Forever \hyperref[coda-the-same-case-revisited-forever]{139}}

\hyperref[chapter-8-stratification-boundaries-thresholds-and-phase-transitions]{Chapter
8: Stratification --- Boundaries, Thresholds, and Phase Transitions
\hyperref[chapter-8-stratification-boundaries-thresholds-and-phase-transitions]{141}}

\hyperref[the-patchwork-structure-of-moral-space]{8.1 The Patchwork
Structure of Moral Space
\hyperref[the-patchwork-structure-of-moral-space]{141}}

\hyperref[whitney-stratification-the-mathematical-framework]{8.2 Whitney
Stratification: The Mathematical Framework
\hyperref[whitney-stratification-the-mathematical-framework]{142}}

\hyperref[from-manifolds-to-stratified-spaces]{From Manifolds to
Stratified Spaces \hyperref[from-manifolds-to-stratified-spaces]{142}}

\hyperref[formal-definition-1]{Formal Definition
\hyperref[formal-definition-1]{142}}

\hyperref[whitneys-conditions]{Whitney's Conditions
\hyperref[whitneys-conditions]{142}}

\hyperref[moral-interpretation]{Moral Interpretation
\hyperref[moral-interpretation]{143}}

\hyperref[a-taxonomy-of-moral-boundaries]{8.3 A Taxonomy of Moral
Boundaries \hyperref[a-taxonomy-of-moral-boundaries]{144}}

\hyperref[type-i-thresholds]{Type I: Thresholds
\hyperref[type-i-thresholds]{144}}

\hyperref[type-ii-phase-transitions]{Type II: Phase Transitions
\hyperref[type-ii-phase-transitions]{145}}

\hyperref[type-iii-absorbing-strata-nullifiers]{Type III: Absorbing
Strata (Nullifiers)
\hyperref[type-iii-absorbing-strata-nullifiers]{146}}

\hyperref[type-iv-constraint-surfaces-forbidden-regions]{Type IV:
Constraint Surfaces (Forbidden Regions)
\hyperref[type-iv-constraint-surfaces-forbidden-regions]{147}}

\hyperref[semantic-gates-the-transition-functions]{8.4 Semantic Gates:
The Transition Functions
\hyperref[semantic-gates-the-transition-functions]{147}}

\hyperref[from-stratum-to-stratum]{From Stratum to Stratum
\hyperref[from-stratum-to-stratum]{147}}

\hyperref[empirical-evidence-for-discreteness]{Empirical Evidence for
Discreteness \hyperref[empirical-evidence-for-discreteness]{148}}

\hyperref[the-dux2084-group-structure-of-hohfeldian-transitions]{The D₄
Group Structure of Hohfeldian Transitions
\hyperref[the-dux2084-group-structure-of-hohfeldian-transitions]{149}}

\hyperref[the-geometry-near-a-boundary]{8.5 The Geometry Near a Boundary
\hyperref[the-geometry-near-a-boundary]{151}}

\hyperref[normal-bundles-and-transverse-structure]{Normal Bundles and
Transverse Structure
\hyperref[normal-bundles-and-transverse-structure]{151}}

\hyperref[the-jump-discontinuity]{The Jump Discontinuity
\hyperref[the-jump-discontinuity]{151}}

\hyperref[penumbral-zones]{Penumbral Zones
\hyperref[penumbral-zones]{152}}

\hyperref[singularities-where-the-geometry-breaks-down]{8.6
Singularities: Where the Geometry Breaks Down
\hyperref[singularities-where-the-geometry-breaks-down]{152}}

\hyperref[moral-singularities-revisited]{Moral Singularities Revisited
\hyperref[moral-singularities-revisited]{152}}

\hyperref[classification-of-singularities]{Classification of
Singularities \hyperref[classification-of-singularities]{152}}

\hyperref[the-moral-content-of-singularities]{The Moral Content of
Singularities \hyperref[the-moral-content-of-singularities]{154}}

\hyperref[phase-transitions-in-moral-reasoning]{8.7 Phase Transitions in
Moral Reasoning \hyperref[phase-transitions-in-moral-reasoning]{154}}

\hyperref[the-physics-analogy]{The Physics Analogy
\hyperref[the-physics-analogy]{154}}

\hyperref[formal-development]{Formal Development
\hyperref[formal-development]{154}}

\hyperref[examples]{Examples \hyperref[examples]{155}}

\hyperref[critical-phenomena]{Critical Phenomena
\hyperref[critical-phenomena]{155}}

\hyperref[stratification-and-moral-dynamics]{8.8 Stratification and
Moral Dynamics \hyperref[stratification-and-moral-dynamics]{156}}

\hyperref[paths-crossing-boundaries]{Paths Crossing Boundaries
\hyperref[paths-crossing-boundaries]{156}}

\hyperref[the-holonomy-of-moral-circuits]{The Holonomy of Moral Circuits
\hyperref[the-holonomy-of-moral-circuits]{157}}

\hyperref[worked-example-the-stratified-landscape-of-a-promise]{8.9
Worked Example: The Stratified Landscape of a Promise
\hyperref[worked-example-the-stratified-landscape-of-a-promise]{157}}

\hyperref[setup-1]{Setup \hyperref[setup-1]{158}}

\hyperref[the-strata]{The Strata \hyperref[the-strata]{158}}

\hyperref[the-boundaries]{The Boundaries \hyperref[the-boundaries]{158}}

\hyperref[a-path-through-the-landscape]{A Path Through the Landscape
\hyperref[a-path-through-the-landscape]{159}}

\hyperref[analysis]{Analysis \hyperref[analysis]{159}}

\hyperref[stratification-and-moral-theory]{8.10 Stratification and Moral
Theory \hyperref[stratification-and-moral-theory]{160}}

\hyperref[how-different-theories-handle-boundaries]{How Different
Theories Handle Boundaries
\hyperref[how-different-theories-handle-boundaries]{160}}

\hyperref[metaethical-implications]{Metaethical Implications
\hyperref[metaethical-implications]{160}}

\hyperref[a-the-discrete-computational-realization-weighted-simplicial-complexes]{8.10a
The Discrete Computational Realization: Weighted Simplicial Complexes
\hyperref[a-the-discrete-computational-realization-weighted-simplicial-complexes]{161}}

\hyperref[summary-1]{8.11 Summary \hyperref[summary-1]{162}}

\hyperref[technical-appendix-1]{Technical Appendix
\hyperref[technical-appendix-1]{163}}

\hyperref[part-iii-dynamics-and-symmetry]{Part III: Dynamics and
Symmetry \hyperref[part-iii-dynamics-and-symmetry]{166}}

\hyperref[chapter-9-the-origin-of-the-moral-metric-discovery-construction-or-governance]{Chapter
9: The Origin of the Moral Metric --- Discovery, Construction, or
Governance?
\hyperref[chapter-9-the-origin-of-the-moral-metric-discovery-construction-or-governance]{166}}

\hyperref[the-metric-is-not-given-1]{9.1 The Metric Is Not Given
\hyperref[the-metric-is-not-given-1]{166}}

\hyperref[the-realist-account-discovery]{9.2 The Realist Account:
Discovery \hyperref[the-realist-account-discovery]{167}}

\hyperref[the-claim]{The Claim \hyperref[the-claim]{167}}

\hyperref[strengths]{Strengths \hyperref[strengths]{167}}

\hyperref[challenges]{Challenges \hyperref[challenges]{167}}

\hyperref[geometric-gloss]{Geometric Gloss
\hyperref[geometric-gloss]{168}}

\hyperref[the-constructivist-account-agreement]{9.3 The Constructivist
Account: Agreement \hyperref[the-constructivist-account-agreement]{168}}

\hyperref[the-claim-1]{The Claim \hyperref[the-claim-1]{168}}

\hyperref[strengths-1]{Strengths \hyperref[strengths-1]{168}}

\hyperref[challenges-1]{Challenges \hyperref[challenges-1]{168}}

\hyperref[geometric-gloss-1]{Geometric Gloss
\hyperref[geometric-gloss-1]{169}}

\hyperref[the-expressivist-account-projection]{9.4 The Expressivist
Account: Projection \hyperref[the-expressivist-account-projection]{169}}

\hyperref[the-claim-2]{The Claim \hyperref[the-claim-2]{169}}

\hyperref[strengths-2]{Strengths \hyperref[strengths-2]{169}}

\hyperref[challenges-2]{Challenges \hyperref[challenges-2]{169}}

\hyperref[geometric-gloss-2]{Geometric Gloss
\hyperref[geometric-gloss-2]{170}}

\hyperref[the-governance-account-democratic-determination]{9.5 The
Governance Account: Democratic Determination
\hyperref[the-governance-account-democratic-determination]{170}}

\hyperref[the-proposal]{The Proposal \hyperref[the-proposal]{170}}

\hyperref[what-governance-means]{What Governance Means
\hyperref[what-governance-means]{170}}

\hyperref[why-not-relativism]{Why Not Relativism?
\hyperref[why-not-relativism]{171}}

\hyperref[the-analogy-to-legal-systems]{The Analogy to Legal Systems
\hyperref[the-analogy-to-legal-systems]{172}}

\hyperref[connection-to-the-framework]{Connection to the Framework
\hyperref[connection-to-the-framework]{172}}

\hyperref[constraints-on-admissible-metrics]{9.6 Constraints on
Admissible Metrics \hyperref[constraints-on-admissible-metrics]{172}}

\hyperref[symmetry]{Symmetry \hyperref[symmetry]{173}}

\hyperref[non-degeneracy-on-protected-dimensions]{Non-Degeneracy on
Protected Dimensions
\hyperref[non-degeneracy-on-protected-dimensions]{173}}

\hyperref[invariance]{Invariance \hyperref[invariance]{173}}

\hyperref[non-discrimination]{Non-Discrimination
\hyperref[non-discrimination]{173}}

\hyperref[consistency-with-structural-invariants]{Consistency with
Structural Invariants
\hyperref[consistency-with-structural-invariants]{174}}

\hyperref[summary-of-constraints]{Summary of Constraints
\hyperref[summary-of-constraints]{174}}

\hyperref[learning-the-metric-from-data]{9.7 Learning the Metric from
Data \hyperref[learning-the-metric-from-data]{174}}

\hyperref[the-empirical-project]{The Empirical Project
\hyperref[the-empirical-project]{174}}

\hyperref[formal-setup]{Formal Setup \hyperref[formal-setup]{174}}

\hyperref[evidence-from-the-dear-abby-corpus]{Evidence from the Dear
Abby Corpus \hyperref[evidence-from-the-dear-abby-corpus]{175}}

\hyperref[evidence-from-bip-experiments]{Evidence from BIP Experiments
\hyperref[evidence-from-bip-experiments]{175}}

\hyperref[challenges-for-metric-learning]{Challenges for Metric Learning
\hyperref[challenges-for-metric-learning]{176}}

\hyperref[pluralism-and-the-meta-metric]{9.8 Pluralism and the
Meta-Metric \hyperref[pluralism-and-the-meta-metric]{176}}

\hyperref[the-question]{The Question \hyperref[the-question]{176}}

\hyperref[three-positions]{Three Positions
\hyperref[three-positions]{176}}

\hyperref[formal-statement-1]{Formal Statement
\hyperref[formal-statement-1]{177}}

\hyperref[the-irreducibility-of-metric-disagreement]{The Irreducibility
of Metric Disagreement
\hyperref[the-irreducibility-of-metric-disagreement]{177}}

\hyperref[the-metric-as-moral-infrastructure]{9.9 The Metric as Moral
Infrastructure \hyperref[the-metric-as-moral-infrastructure]{178}}

\hyperref[an-analogy]{An Analogy \hyperref[an-analogy]{178}}

\hyperref[implications-for-ai-systems]{Implications for AI Systems
\hyperref[implications-for-ai-systems]{178}}

\hyperref[summary-2]{9.10 Summary \hyperref[summary-2]{179}}

\hyperref[technical-appendix-2]{Technical Appendix
\hyperref[technical-appendix-2]{180}}

\hyperref[characterization-of-the-admissible-metric-space]{Characterization
of the Admissible Metric Space
\hyperref[characterization-of-the-admissible-metric-space]{180}}

\hyperref[chapter-10-moral-dynamics-parallel-transport-holonomy-and-the-curvature-of-moral-space]{Chapter
10: Moral Dynamics --- Parallel Transport, Holonomy, and the Curvature
of Moral Space
\hyperref[chapter-10-moral-dynamics-parallel-transport-holonomy-and-the-curvature-of-moral-space]{183}}

\hyperref[ethics-in-motion]{10.1 Ethics in Motion
\hyperref[ethics-in-motion]{183}}

\hyperref[the-problem-of-moral-change]{10.2 The Problem of Moral Change
\hyperref[the-problem-of-moral-change]{183}}

\hyperref[the-moral-connection]{10.3 The Moral Connection
\hyperref[the-moral-connection]{184}}

\hyperref[comparing-obligations-across-contexts]{Comparing Obligations
Across Contexts \hyperref[comparing-obligations-across-contexts]{184}}

\hyperref[what-the-connection-encodes]{What the Connection Encodes
\hyperref[what-the-connection-encodes]{184}}

\hyperref[the-levi-civita-connection]{The Levi-Civita Connection
\hyperref[the-levi-civita-connection]{185}}

\hyperref[non-metric-connections]{Non-Metric Connections
\hyperref[non-metric-connections]{185}}

\hyperref[parallel-transport-of-obligations]{10.4 Parallel Transport of
Obligations \hyperref[parallel-transport-of-obligations]{186}}

\hyperref[the-concept]{The Concept \hyperref[the-concept]{186}}

\hyperref[moral-interpretation-1]{Moral Interpretation
\hyperref[moral-interpretation-1]{187}}

\hyperref[what-is-not-parallel-transport]{What Is \emph{Not} Parallel
Transport \hyperref[what-is-not-parallel-transport]{187}}

\hyperref[holonomy-the-path-dependence-of-moral-evaluation]{10.5
Holonomy: The Path-Dependence of Moral Evaluation
\hyperref[holonomy-the-path-dependence-of-moral-evaluation]{188}}

\hyperref[the-phenomenon]{The Phenomenon \hyperref[the-phenomenon]{188}}

\hyperref[a-moral-example]{A Moral Example
\hyperref[a-moral-example]{188}}

\hyperref[the-holonomy-group]{The Holonomy Group
\hyperref[the-holonomy-group]{189}}

\hyperref[the-curvature-of-moral-space]{10.6 The Curvature of Moral
Space \hyperref[the-curvature-of-moral-space]{189}}

\hyperref[the-riemann-curvature-tensor-1]{The Riemann Curvature Tensor
\hyperref[the-riemann-curvature-tensor-1]{189}}

\hyperref[sources-of-moral-curvature]{Sources of Moral Curvature
\hyperref[sources-of-moral-curvature]{190}}

\hyperref[the-ricci-tensor-and-scalar-curvature]{The Ricci Tensor and
Scalar Curvature \hyperref[the-ricci-tensor-and-scalar-curvature]{190}}

\hyperref[curvature-and-the-moral-metric-an-equation]{Curvature and the
Moral Metric: An Equation?
\hyperref[curvature-and-the-moral-metric-an-equation]{191}}

\hyperref[geodesics-paths-of-least-moral-resistance]{10.7 Geodesics:
Paths of Least Moral Resistance
\hyperref[geodesics-paths-of-least-moral-resistance]{191}}

\hyperref[definition-4]{Definition \hyperref[definition-4]{191}}

\hyperref[moral-interpretation-2]{Moral Interpretation
\hyperref[moral-interpretation-2]{192}}

\hyperref[geodesic-deviation]{Geodesic Deviation
\hyperref[geodesic-deviation]{192}}

\hyperref[gradient-flows-and-moral-improvement]{10.8 Gradient Flows and
Moral Improvement \hyperref[gradient-flows-and-moral-improvement]{193}}

\hyperref[the-gradient-field]{The Gradient Field
\hyperref[the-gradient-field]{193}}

\hyperref[the-gradient-flow]{The Gradient Flow
\hyperref[the-gradient-flow]{193}}

\hyperref[virtues-and-vices-of-gradient-following]{Virtues and Vices of
Gradient Following
\hyperref[virtues-and-vices-of-gradient-following]{193}}

\hyperref[virtue-as-non-gradient-agency]{Virtue as Non-Gradient Agency
\hyperref[virtue-as-non-gradient-agency]{194}}

\hyperref[temporal-structure-and-moral-history]{10.9 Temporal Structure
and Moral History \hyperref[temporal-structure-and-moral-history]{194}}

\hyperref[time-and-the-moral-manifold]{Time and the Moral Manifold
\hyperref[time-and-the-moral-manifold]{194}}

\hyperref[temporal-moral-tensors]{Temporal Moral Tensors
\hyperref[temporal-moral-tensors]{194}}

\hyperref[intergenerational-obligations]{Intergenerational Obligations
\hyperref[intergenerational-obligations]{195}}

\hyperref[worked-example-the-evolving-obligation]{10.10 Worked Example:
The Evolving Obligation
\hyperref[worked-example-the-evolving-obligation]{196}}

\hyperref[setup-2]{Setup \hyperref[setup-2]{196}}

\hyperref[_Toc224082388]{\(t1\): Work Emergency
\hyperref[_Toc224082388]{196}}

\hyperref[_Toc224082389]{\(t2\): Recital Postponed
\hyperref[_Toc224082389]{196}}

\hyperref[_Toc224082390]{\(t3\): The Parent Attends
\hyperref[_Toc224082390]{197}}

\hyperref[holonomy-computation]{Holonomy Computation
\hyperref[holonomy-computation]{197}}

\hyperref[curvature-and-moral-disagreement]{10.11 Curvature and Moral
Disagreement \hyperref[curvature-and-moral-disagreement]{197}}

\hyperref[why-moral-reasoning-is-path-dependent]{Why Moral Reasoning Is
Path-Dependent \hyperref[why-moral-reasoning-is-path-dependent]{197}}

\hyperref[implications]{Implications \hyperref[implications]{197}}

\hyperref[summary-3]{10.12 Summary \hyperref[summary-3]{198}}

\hyperref[technical-appendix-4]{Technical Appendix
\hyperref[technical-appendix-4]{199}}

\hyperref[chapter-11-moral-reasoning-as-optimal-search]{Chapter 11:
Moral Reasoning as Optimal Search
\hyperref[chapter-11-moral-reasoning-as-optimal-search]{202}}

\hyperref[the-computational-problem]{11.1 The Computational Problem
\hyperref[the-computational-problem]{202}}

\hyperref[a-search-on-the-moral-manifold]{11.2 A* Search on the Moral
Manifold \hyperref[a-search-on-the-moral-manifold]{202}}

\hyperref[obligation-vectors-as-heuristic-functions]{11.3 Obligation
Vectors as Heuristic Functions
\hyperref[obligation-vectors-as-heuristic-functions]{204}}

\hyperref[the-murder-prohibition-as-hn]{The Murder Prohibition as h(n) =
∞ \hyperref[the-murder-prohibition-as-hn]{205}}

\hyperref[promise-keeping-as-a-finite-heuristic]{Promise-Keeping as a
Finite Heuristic \hyperref[promise-keeping-as-a-finite-heuristic]{205}}

\hyperref[the-evolutionary-origin]{The Evolutionary Origin
\hyperref[the-evolutionary-origin]{206}}

\hyperref[admissibility-and-consistency]{11.4 Admissibility and
Consistency \hyperref[admissibility-and-consistency]{206}}

\hyperref[the-ux3b5-admissibility-spectrum-and-the-heuristic-hierarchy]{The
ε-Admissibility Spectrum and the Heuristic Hierarchy
\hyperref[the-ux3b5-admissibility-spectrum-and-the-heuristic-hierarchy]{207}}

\hyperref[computational-intractability-of-exact-moral-reasoning]{11.5
Computational Intractability of Exact Moral Reasoning
\hyperref[computational-intractability-of-exact-moral-reasoning]{209}}

\hyperref[the-satisfaction-scalar-as-search-output]{11.6 The
Satisfaction Scalar as Search Output
\hyperref[the-satisfaction-scalar-as-search-output]{211}}

\hyperref[worked-example-the-kidney-allocation-revisited]{Worked
Example: The Kidney Allocation Revisited
\hyperref[worked-example-the-kidney-allocation-revisited]{211}}

\hyperref[connections-to-existing-framework-elements]{11.7 Connections
to Existing Framework Elements
\hyperref[connections-to-existing-framework-elements]{212}}

\hyperref[gradient-flows-as-degenerate-a]{Gradient Flows as Degenerate
A* \hyperref[gradient-flows-as-degenerate-a]{212}}

\hyperref[geodesics-as-a-solutions]{Geodesics as A* Solutions
\hyperref[geodesics-as-a-solutions]{212}}

\hyperref[holonomy-and-heuristic-drift]{Holonomy and Heuristic Drift
\hyperref[holonomy-and-heuristic-drift]{212}}

\hyperref[contraction-and-information-loss]{Contraction and Information
Loss \hyperref[contraction-and-information-loss]{212}}

\hyperref[implications-for-moral-philosophy]{11.8 Implications for Moral
Philosophy \hyperref[implications-for-moral-philosophy]{213}}

\hyperref[deontology-as-pre-compiled-heuristics]{Deontology as
Pre-Compiled Heuristics
\hyperref[deontology-as-pre-compiled-heuristics]{213}}

\hyperref[consequentialism-as-exact-search]{Consequentialism as Exact
Search \hyperref[consequentialism-as-exact-search]{213}}

\hyperref[virtue-ethics-as-heuristic-calibration]{Virtue Ethics as
Heuristic Calibration
\hyperref[virtue-ethics-as-heuristic-calibration]{213}}

\hyperref[dual-process-moral-cognition]{Dual-Process Moral Cognition
\hyperref[dual-process-moral-cognition]{213}}

\hyperref[worked-example-the-self-driving-car]{11.9 Worked Example: The
Self-Driving Car \hyperref[worked-example-the-self-driving-car]{214}}

\hyperref[summary-4]{11.10 Summary \hyperref[summary-4]{215}}

\hyperref[technical-appendix-5]{Technical Appendix
\hyperref[technical-appendix-5]{215}}

\hyperref[chapter-12-noether-style-reasoning-for-ethics-symmetry-invariance-and-harm-accounting]{Chapter
12: Noether-Style Reasoning for Ethics --- Symmetry, Invariance, and
Harm Accounting
\hyperref[chapter-12-noether-style-reasoning-for-ethics-symmetry-invariance-and-harm-accounting]{219}}

\hyperref[the-deepest-theorem-in-physics]{12.1 The Deepest Theorem in
Physics \hyperref[the-deepest-theorem-in-physics]{219}}

\hyperref[noethers-theorem-the-physics]{12.2 Noether's Theorem: The
Physics \hyperref[noethers-theorem-the-physics]{220}}

\hyperref[the-action-principle]{The Action Principle
\hyperref[the-action-principle]{220}}

\hyperref[symmetries-and-conservation-laws]{Symmetries and Conservation
Laws \hyperref[symmetries-and-conservation-laws]{220}}

\hyperref[the-standard-examples]{The Standard Examples
\hyperref[the-standard-examples]{221}}

\hyperref[the-moral-symmetry-the-bond-invariance-principle]{12.3 The
Moral Symmetry: The Bond Invariance Principle
\hyperref[the-moral-symmetry-the-bond-invariance-principle]{221}}

\hyperref[re-description-invariance]{Re-Description Invariance
\hyperref[re-description-invariance]{221}}

\hyperref[bip-as-a-continuous-symmetry]{BIP as a Continuous Symmetry
\hyperref[bip-as-a-continuous-symmetry]{221}}

\hyperref[the-gauge-group-of-the-moral-manifold]{The Gauge Group of the
Moral Manifold \hyperref[the-gauge-group-of-the-moral-manifold]{222}}

\hyperref[axiomatic-derivation-of-the-gauge-group]{Axiomatic Derivation
of the Gauge Group
\hyperref[axiomatic-derivation-of-the-gauge-group]{223}}

\hyperref[pre-flight-checklist-gauge-group-dux2084-u1_h]{Pre-flight
Checklist: Gauge Group (D₄ × U(1)\_H)
\hyperref[pre-flight-checklist-gauge-group-dux2084-u1_h]{225}}

\hyperref[the-moral-lagrangian]{12.4 The Moral Lagrangian
\hyperref[the-moral-lagrangian]{225}}

\hyperref[constructing-the-action]{Constructing the Action
\hyperref[constructing-the-action]{225}}

\hyperref[euler-lagrange-equations]{Euler-Lagrange Equations
\hyperref[euler-lagrange-equations]{226}}

\hyperref[interpretation-of-the-potential]{Interpretation of the
Potential \hyperref[interpretation-of-the-potential]{226}}

\hyperref[the-noether-current-for-re-description-invariance]{12.5 The
Noether Current for Re-Description Invariance
\hyperref[the-noether-current-for-re-description-invariance]{227}}

\hyperref[the-derivation]{The Derivation \hyperref[the-derivation]{227}}

\hyperref[interpretation-what-is-conserved]{Interpretation: What Is
Conserved? \hyperref[interpretation-what-is-conserved]{227}}

\hyperref[the-invariance-of-harm-accounting]{12.6 The Invariance of Harm
Accounting \hyperref[the-invariance-of-harm-accounting]{228}}

\hyperref[from-noether-charge-to-harm]{From Noether Charge to Harm
\hyperref[from-noether-charge-to-harm]{228}}

\hyperref[continuous-vs.-discrete-symmetry]{Continuous vs. Discrete
Symmetry \hyperref[continuous-vs.-discrete-symmetry]{228}}

\hyperref[what-invariance-of-harm-accounting-means]{What Invariance of
Harm Accounting Means
\hyperref[what-invariance-of-harm-accounting-means]{230}}

\hyperref[the-harm-ledger]{The Harm Ledger
\hyperref[the-harm-ledger]{231}}

\hyperref[four-consequences-of-harm-conservation]{12.7 Four Consequences
of Harm Conservation
\hyperref[four-consequences-of-harm-conservation]{231}}

\hyperref[consequence-1-euphemism-does-not-reduce-harm]{Consequence 1:
Euphemism Does Not Reduce Harm
\hyperref[consequence-1-euphemism-does-not-reduce-harm]{231}}

\hyperref[consequence-2-harm-is-auditable-across-representations]{Consequence
2: Harm Is Auditable Across Representations
\hyperref[consequence-2-harm-is-auditable-across-representations]{231}}

\hyperref[consequence-3-re-description-cannot-redistribute-harm-among-dimensions]{Consequence
3: Re-Description Cannot Redistribute Harm Among Dimensions
\hyperref[consequence-3-re-description-cannot-redistribute-harm-among-dimensions]{232}}

\hyperref[consequence-4-moral-debt-persists]{Consequence 4: Moral Debt
Persists \hyperref[consequence-4-moral-debt-persists]{232}}

\hyperref[empirical-evidence-cross-lingual-invariance]{12.8 Empirical
Evidence: Cross-Lingual Invariance
\hyperref[empirical-evidence-cross-lingual-invariance]{232}}

\hyperref[the-prediction]{The Prediction \hyperref[the-prediction]{232}}

\hyperref[the-evidence]{The Evidence \hyperref[the-evidence]{233}}

\hyperref[the-correlative-symmetry-as-a-conservation-law]{The
Correlative Symmetry as a Conservation Law
\hyperref[the-correlative-symmetry-as-a-conservation-law]{233}}

\hyperref[harm-conservation-and-ai-alignment]{12.9 Harm Conservation and
AI Alignment \hyperref[harm-conservation-and-ai-alignment]{233}}

\hyperref[invariance-testing-as-alignment-verification]{Invariance
Testing as Alignment Verification
\hyperref[invariance-testing-as-alignment-verification]{233}}

\hyperref[the-no-escape-connection]{The No Escape Connection
\hyperref[the-no-escape-connection]{234}}

\hyperref[limitations-and-caveats]{12.10 Limitations and Caveats
\hyperref[limitations-and-caveats]{234}}

\hyperref[what-the-conservation-law-does-not-say]{What the Conservation
Law Does Not Say \hyperref[what-the-conservation-law-does-not-say]{234}}

\hyperref[the-analogys-limits]{The Analogy's Limits
\hyperref[the-analogys-limits]{235}}

\hyperref[summary-5]{12.11 Summary \hyperref[summary-5]{235}}

\hyperref[technical-appendix-6]{Technical Appendix
\hyperref[technical-appendix-6]{236}}

\hyperref[chapter-13-quantum-normative-dynamics-superposition-measurement-and-the-stratified-lagrangian]{Chapter
13: Quantum Normative Dynamics --- Superposition, Measurement, and the
Stratified Lagrangian
\hyperref[chapter-13-quantum-normative-dynamics-superposition-measurement-and-the-stratified-lagrangian]{239}}

\hyperref[why-quantum]{13.1 Why Quantum? \hyperref[why-quantum]{239}}

\hyperref[the-moral-hilbert-space]{13.2 The Moral Hilbert Space
\hyperref[the-moral-hilbert-space]{240}}

\hyperref[from-manifold-to-hilbert-space]{From Manifold to Hilbert Space
\hyperref[from-manifold-to-hilbert-space]{240}}

\hyperref[classical-states-as-delta-functions]{Classical States as Delta
Functions \hyperref[classical-states-as-delta-functions]{240}}

\hyperref[basis-states-and-the-position-representation]{Basis States and
the Position Representation
\hyperref[basis-states-and-the-position-representation]{240}}

\hyperref[superposition-deliberation-as-quantum-state]{13.3
Superposition: Deliberation as Quantum State
\hyperref[superposition-deliberation-as-quantum-state]{241}}

\hyperref[the-structure-of-deliberation]{The Structure of Deliberation
\hyperref[the-structure-of-deliberation]{241}}

\hyperref[superposition-vs.-mixture]{Superposition vs.~Mixture
\hyperref[superposition-vs.-mixture]{241}}

\hyperref[when-does-superposition-arise]{When Does Superposition Arise?
\hyperref[when-does-superposition-arise]{242}}

\hyperref[moral-observables-and-measurement]{13.4 Moral Observables and
Measurement \hyperref[moral-observables-and-measurement]{242}}

\hyperref[observables-as-self-adjoint-operators]{Observables as
Self-Adjoint Operators
\hyperref[observables-as-self-adjoint-operators]{242}}

\hyperref[measurement-as-decision]{Measurement as Decision
\hyperref[measurement-as-decision]{243}}

\hyperref[the-uncertainty-principle-for-moral-observables]{The
Uncertainty Principle for Moral Observables
\hyperref[the-uncertainty-principle-for-moral-observables]{243}}

\hyperref[interference-when-moral-framings-interact]{13.5 Interference:
When Moral Framings Interact
\hyperref[interference-when-moral-framings-interact]{244}}

\hyperref[the-double-slit-analogy]{The Double-Slit Analogy
\hyperref[the-double-slit-analogy]{244}}

\hyperref[moral-interference-a-formal-example]{Moral Interference: A
Formal Example \hyperref[moral-interference-a-formal-example]{244}}

\hyperref[the-phase-theta-framing-alignment]{The Phase \(\theta\):
Framing Alignment \hyperref[the-phase-theta-framing-alignment]{245}}

\hyperref[empirical-signatures-of-moral-interference]{Empirical
Signatures of Moral Interference
\hyperref[empirical-signatures-of-moral-interference]{245}}

\hyperref[the-density-matrix-mixed-moral-states]{13.6 The Density
Matrix: Mixed Moral States
\hyperref[the-density-matrix-mixed-moral-states]{246}}

\hyperref[pure-states-and-mixed-states]{Pure States and Mixed States
\hyperref[pure-states-and-mixed-states]{246}}

\hyperref[decoherence-from-quantum-to-classical]{Decoherence: From
Quantum to Classical
\hyperref[decoherence-from-quantum-to-classical]{246}}

\hyperref[the-purity-of-a-moral-state]{The Purity of a Moral State
\hyperref[the-purity-of-a-moral-state]{247}}

\hyperref[the-stratified-lagrangian]{13.7 The Stratified Lagrangian
\hyperref[the-stratified-lagrangian]{247}}

\hyperref[combining-lagrangian-dynamics-with-stratification]{Combining
Lagrangian Dynamics with Stratification
\hyperref[combining-lagrangian-dynamics-with-stratification]{247}}

\hyperref[the-classical-stratified-lagrangian]{The Classical Stratified
Lagrangian \hyperref[the-classical-stratified-lagrangian]{247}}

\hyperref[boundary-conditions-at-stratum-boundaries]{Boundary Conditions
at Stratum Boundaries
\hyperref[boundary-conditions-at-stratum-boundaries]{247}}

\hyperref[the-quantum-stratified-lagrangian]{The Quantum Stratified
Lagrangian \hyperref[the-quantum-stratified-lagrangian]{248}}

\hyperref[the-moral-schruxf6dinger-equation]{The Moral Schrödinger
Equation \hyperref[the-moral-schruxf6dinger-equation]{249}}

\hyperref[tunneling-across-moral-barriers]{Tunneling Across Moral
Barriers \hyperref[tunneling-across-moral-barriers]{249}}

\hyperref[entanglement-correlated-moral-states]{13.8 Entanglement:
Correlated Moral States
\hyperref[entanglement-correlated-moral-states]{250}}

\hyperref[multiagent-quantum-states]{Multiagent Quantum States
\hyperref[multiagent-quantum-states]{250}}

\hyperref[moral-epr-states]{Moral EPR States
\hyperref[moral-epr-states]{250}}

\hyperref[entanglement-and-collective-obligation]{Entanglement and
Collective Obligation
\hyperref[entanglement-and-collective-obligation]{251}}

\hyperref[the-classical-limit-when-classical-ethics-suffices]{13.9 The
Classical Limit: When Classical Ethics Suffices
\hyperref[the-classical-limit-when-classical-ethics-suffices]{251}}

\hyperref[the-correspondence-principle]{The Correspondence Principle
\hyperref[the-correspondence-principle]{251}}

\hyperref[when-to-use-quantum-normative-dynamics]{When to Use Quantum
Normative Dynamics
\hyperref[when-to-use-quantum-normative-dynamics]{251}}

\hyperref[worked-example-the-whistleblowers-dilemma]{13.10 Worked
Example: The Whistleblower's Dilemma
\hyperref[worked-example-the-whistleblowers-dilemma]{253}}

\hyperref[setup-3]{Setup \hyperref[setup-3]{253}}

\hyperref[the-classical-analysis-review]{The Classical Analysis (Review)
\hyperref[the-classical-analysis-review]{253}}

\hyperref[the-quantum-analysis]{The Quantum Analysis
\hyperref[the-quantum-analysis]{253}}

\hyperref[time-evolution]{Time Evolution \hyperref[time-evolution]{254}}

\hyperref[the-measurement]{The Measurement
\hyperref[the-measurement]{254}}

\hyperref[what-the-quantum-analysis-adds]{What the Quantum Analysis Adds
\hyperref[what-the-quantum-analysis-adds]{254}}

\hyperref[summary-6]{13.11 Summary \hyperref[summary-6]{255}}

\hyperref[technical-appendix-7]{Technical Appendix
\hyperref[technical-appendix-7]{256}}

\hyperref[part-iv-agents-and-collectives]{Part IV: Agents and
Collectives \hyperref[part-iv-agents-and-collectives]{259}}

\hyperref[chapter-14-collective-moral-agency-aggregation-emergence-and-the-tensors-of-shared-obligation]{Chapter
14: Collective Moral Agency --- Aggregation, Emergence, and the Tensors
of Shared Obligation
\hyperref[chapter-14-collective-moral-agency-aggregation-emergence-and-the-tensors-of-shared-obligation]{259}}

\hyperref[beyond-individual-agents]{14.1 Beyond Individual Agents
\hyperref[beyond-individual-agents]{259}}

\hyperref[the-aggregation-problem]{14.2 The Aggregation Problem
\hyperref[the-aggregation-problem]{260}}

\hyperref[why-simple-addition-fails]{Why Simple Addition Fails
\hyperref[why-simple-addition-fails]{260}}

\hyperref[the-tensor-product-approach]{The Tensor Product Approach
\hyperref[the-tensor-product-approach]{260}}

\hyperref[the-collective-agency-tensor]{14.3 The Collective Agency
Tensor \hyperref[the-collective-agency-tensor]{261}}

\hyperref[definition-5]{Definition \hyperref[definition-5]{261}}

\hyperref[decomposition-individual-and-emergent-components]{Decomposition:
Individual and Emergent Components
\hyperref[decomposition-individual-and-emergent-components]{261}}

\hyperref[measuring-emergence]{Measuring Emergence
\hyperref[measuring-emergence]{262}}

\hyperref[the-structure-tensor]{14.4 The Structure Tensor
\hyperref[the-structure-tensor]{262}}

\hyperref[what-makes-a-collective-more-than-a-set]{What Makes a
Collective More Than a Set
\hyperref[what-makes-a-collective-more-than-a-set]{262}}

\hyperref[how-structure-shapes-collective-obligation]{How Structure
Shapes Collective Obligation
\hyperref[how-structure-shapes-collective-obligation]{263}}

\hyperref[emergent-obligations]{14.5 Emergent Obligations
\hyperref[emergent-obligations]{264}}

\hyperref[obligations-no-member-bears]{Obligations No Member Bears
\hyperref[obligations-no-member-bears]{264}}

\hyperref[formal-criterion-for-emergence]{Formal Criterion for Emergence
\hyperref[formal-criterion-for-emergence]{264}}

\hyperref[distributed-responsibility]{14.6 Distributed Responsibility
\hyperref[distributed-responsibility]{265}}

\hyperref[the-responsibility-tensor]{The Responsibility Tensor
\hyperref[the-responsibility-tensor]{265}}

\hyperref[is-responsibility-conserved]{Is Responsibility Conserved?
\hyperref[is-responsibility-conserved]{266}}

\hyperref[the-complicity-gradient]{The Complicity Gradient
\hyperref[the-complicity-gradient]{266}}

\hyperref[the-geometry-of-institutional-structure]{14.7 The Geometry of
Institutional Structure
\hyperref[the-geometry-of-institutional-structure]{267}}

\hyperref[institutions-as-moral-manifolds]{Institutions as Moral
Manifolds \hyperref[institutions-as-moral-manifolds]{267}}

\hyperref[institutional-metrics]{Institutional Metrics
\hyperref[institutional-metrics]{267}}

\hyperref[institutional-design-as-geometric-optimization]{Institutional
Design as Geometric Optimization
\hyperref[institutional-design-as-geometric-optimization]{268}}

\hyperref[collective-moral-dynamics]{14.8 Collective Moral Dynamics
\hyperref[collective-moral-dynamics]{268}}

\hyperref[parallel-transport-of-collective-obligations]{Parallel
Transport of Collective Obligations
\hyperref[parallel-transport-of-collective-obligations]{268}}

\hyperref[collective-holonomy]{Collective Holonomy
\hyperref[collective-holonomy]{269}}

\hyperref[collective-agents-in-the-nine-dimensions]{14.9 Collective
Agents in the Nine Dimensions
\hyperref[collective-agents-in-the-nine-dimensions]{269}}

\hyperref[the-collective-scope]{The Collective Scope
\hyperref[the-collective-scope]{269}}

\hyperref[the-collective-individual-coupling]{The Collective-Individual
Coupling \hyperref[the-collective-individual-coupling]{269}}

\hyperref[emergent-moral-properties]{14.10 Emergent Moral Properties
\hyperref[emergent-moral-properties]{270}}

\hyperref[properties-that-exist-only-at-the-collective-level]{Properties
That Exist Only at the Collective Level
\hyperref[properties-that-exist-only-at-the-collective-level]{270}}

\hyperref[the-emergent-properties-criterion]{The Emergent Properties
Criterion \hyperref[the-emergent-properties-criterion]{270}}

\hyperref[ai-systems-as-collective-agents]{14.11 AI Systems as
Collective Agents \hyperref[ai-systems-as-collective-agents]{271}}

\hyperref[human-ai-collectives]{Human-AI Collectives
\hyperref[human-ai-collectives]{271}}

\hyperref[responsibility-distribution-in-human-ai-teams]{Responsibility
Distribution in Human-AI Teams
\hyperref[responsibility-distribution-in-human-ai-teams]{271}}

\hyperref[the-alignment-problem-as-emergent-agency]{The Alignment
Problem as Emergent Agency
\hyperref[the-alignment-problem-as-emergent-agency]{271}}

\hyperref[worked-example-the-hospital-ethics-committee]{14.12 Worked
Example: The Hospital Ethics Committee
\hyperref[worked-example-the-hospital-ethics-committee]{272}}

\hyperref[setup-4]{Setup \hyperref[setup-4]{272}}

\hyperref[the-individual-evaluation-tensors]{The Individual Evaluation
Tensors \hyperref[the-individual-evaluation-tensors]{272}}

\hyperref[the-structure-tensor-1]{The Structure Tensor
\hyperref[the-structure-tensor-1]{272}}

\hyperref[the-emergent-component]{The Emergent Component
\hyperref[the-emergent-component]{273}}

\hyperref[the-decision]{The Decision \hyperref[the-decision]{273}}

\hyperref[responsibility-for-the-outcome]{Responsibility for the Outcome
\hyperref[responsibility-for-the-outcome]{274}}

\hyperref[summary-7]{14.13 Summary \hyperref[summary-7]{274}}

\hyperref[technical-appendix-8]{Technical Appendix
\hyperref[technical-appendix-8]{275}}

\hyperref[chapter-15-from-tensor-to-decision-the-philosophy-of-moral-contraction]{Chapter
15: From Tensor to Decision --- The Philosophy of Moral Contraction
\hyperref[chapter-15-from-tensor-to-decision-the-philosophy-of-moral-contraction]{277}}

\hyperref[the-moment-of-choice]{15.1 The Moment of Choice
\hyperref[the-moment-of-choice]{277}}

\hyperref[the-mathematics-of-contraction]{15.2 The Mathematics of
Contraction \hyperref[the-mathematics-of-contraction]{277}}

\hyperref[contraction-as-index-summation]{Contraction as Index Summation
\hyperref[contraction-as-index-summation]{277}}

\hyperref[the-contraction-chain]{The Contraction Chain
\hyperref[the-contraction-chain]{278}}

\hyperref[non-commutativity-of-contraction]{Non-Commutativity of
Contraction \hyperref[non-commutativity-of-contraction]{278}}

\hyperref[the-necessity-of-contraction]{15.3 The Necessity of
Contraction \hyperref[the-necessity-of-contraction]{279}}

\hyperref[why-contract-at-all]{Why Contract at All?
\hyperref[why-contract-at-all]{279}}

\hyperref[the-cost-of-premature-contraction]{The Cost of Premature
Contraction \hyperref[the-cost-of-premature-contraction]{279}}

\hyperref[who-contracts]{15.4 Who Contracts?
\hyperref[who-contracts]{280}}

\hyperref[the-agent]{The Agent \hyperref[the-agent]{280}}

\hyperref[the-community]{The Community \hyperref[the-community]{280}}

\hyperref[the-institution]{The Institution
\hyperref[the-institution]{280}}

\hyperref[the-algorithm]{The Algorithm \hyperref[the-algorithm]{280}}

\hyperref[types-of-contraction]{15.5 Types of Contraction
\hyperref[types-of-contraction]{281}}

\hyperref[a-taxonomy]{A Taxonomy \hyperref[a-taxonomy]{281}}

\hyperref[what-unifies-them]{What Unifies Them
\hyperref[what-unifies-them]{282}}

\hyperref[the-information-lost-in-contraction]{15.6 The Information Lost
in Contraction \hyperref[the-information-lost-in-contraction]{282}}

\hyperref[what-contraction-discards]{What Contraction Discards
\hyperref[what-contraction-discards]{282}}

\hyperref[an-information-theoretic-measure]{An Information-Theoretic
Measure \hyperref[an-information-theoretic-measure]{283}}

\hyperref[the-scalar-irrecoverability-theorem]{The Scalar
Irrecoverability Theorem
\hyperref[the-scalar-irrecoverability-theorem]{283}}

\hyperref[moral-residue]{15.7 Moral Residue
\hyperref[moral-residue]{284}}

\hyperref[what-contraction-leaves-behind]{What Contraction Leaves Behind
\hyperref[what-contraction-leaves-behind]{284}}

\hyperref[residue-and-holonomy]{Residue and Holonomy
\hyperref[residue-and-holonomy]{285}}

\hyperref[residue-is-not-optional]{Residue Is Not Optional
\hyperref[residue-is-not-optional]{285}}

\hyperref[deferred-contraction]{15.8 Deferred Contraction
\hyperref[deferred-contraction]{286}}

\hyperref[the-virtue-of-tensorial-suspension]{The Virtue of Tensorial
Suspension \hyperref[the-virtue-of-tensorial-suspension]{286}}

\hyperref[the-skills-of-deferred-contraction]{The Skills of Deferred
Contraction \hyperref[the-skills-of-deferred-contraction]{286}}

\hyperref[limits-of-deferral]{Limits of Deferral
\hyperref[limits-of-deferral]{287}}

\hyperref[the-phenomenology-of-contraction]{15.9 The Phenomenology of
Contraction \hyperref[the-phenomenology-of-contraction]{287}}

\hyperref[what-moral-decision-feels-like]{What Moral Decision Feels Like
\hyperref[what-moral-decision-feels-like]{287}}

\hyperref[moral-education-as-contraction-training]{Moral Education as
Contraction Training
\hyperref[moral-education-as-contraction-training]{288}}

\hyperref[contraction-and-the-invariance-of-harm-accounting]{15.10
Contraction and the Invariance of Harm Accounting
\hyperref[contraction-and-the-invariance-of-harm-accounting]{288}}

\hyperref[noether-constraints-on-contraction]{Noether Constraints on
Contraction \hyperref[noether-constraints-on-contraction]{288}}

\hyperref[implications-for-accountability]{Implications for
Accountability \hyperref[implications-for-accountability]{289}}

\hyperref[contraction-in-ai-systems]{15.11 Contraction in AI Systems
\hyperref[contraction-in-ai-systems]{289}}

\hyperref[loss-functions-as-moral-contractions]{Loss Functions as Moral
Contractions \hyperref[loss-functions-as-moral-contractions]{289}}

\hyperref[the-alignment-problem-as-contraction-mismatch]{The Alignment
Problem as Contraction Mismatch
\hyperref[the-alignment-problem-as-contraction-mismatch]{289}}

\hyperref[toward-transparent-ai-contraction]{Toward Transparent AI
Contraction \hyperref[toward-transparent-ai-contraction]{290}}

\hyperref[worked-example-the-kidney-allocation-revisited-1]{15.12 Worked
Example: The Kidney Allocation Revisited
\hyperref[worked-example-the-kidney-allocation-revisited-1]{291}}

\hyperref[the-full-tensor]{The Full Tensor
\hyperref[the-full-tensor]{291}}

\hyperref[three-contractions-three-verdicts-review]{Three Contractions,
Three Verdicts (Review)
\hyperref[three-contractions-three-verdicts-review]{291}}

\hyperref[the-residue-of-each-contraction]{The Residue of Each
Contraction \hyperref[the-residue-of-each-contraction]{291}}

\hyperref[comparing-residues]{Comparing Residues
\hyperref[comparing-residues]{292}}

\hyperref[summary-8]{15.13 Summary \hyperref[summary-8]{292}}

\hyperref[worked-example-the-smart-home-emergency]{Worked Example: The
Smart Home Emergency
\hyperref[worked-example-the-smart-home-emergency]{293}}

\hyperref[situation]{Situation \hyperref[situation]{293}}

\hyperref[the-tensor]{The Tensor \hyperref[the-tensor]{293}}

\hyperref[contraction-choice]{Contraction Choice
\hyperref[contraction-choice]{293}}

\hyperref[residue-and-audit-artifact]{Residue and Audit Artifact
\hyperref[residue-and-audit-artifact]{293}}

\hyperref[worked-example-the-carbon-budget-allocation]{Worked Example:
The Carbon Budget Allocation
\hyperref[worked-example-the-carbon-budget-allocation]{294}}

\hyperref[situation-1]{Situation \hyperref[situation-1]{294}}

\hyperref[the-tensor-1]{The Tensor \hyperref[the-tensor-1]{294}}

\hyperref[contraction-choice-1]{Contraction Choice
\hyperref[contraction-choice-1]{294}}

\hyperref[residue-and-audit-artifact-1]{Residue and Audit Artifact
\hyperref[residue-and-audit-artifact-1]{294}}

\hyperref[technical-appendix-9]{Technical Appendix
\hyperref[technical-appendix-9]{294}}

\hyperref[chapter-16-moral-uncertainty-and-the-limits-of-geometric-determination]{Chapter
16: Moral Uncertainty and the Limits of Geometric Determination
\hyperref[chapter-16-moral-uncertainty-and-the-limits-of-geometric-determination]{297}}

\hyperref[what-the-framework-cannot-settle]{16.1 What the Framework
Cannot Settle \hyperref[what-the-framework-cannot-settle]{297}}

\hyperref[three-types-of-moral-uncertainty]{16.2 Three Types of Moral
Uncertainty \hyperref[three-types-of-moral-uncertainty]{297}}

\hyperref[type-1-empirical-uncertainty]{Type 1: Empirical Uncertainty
\hyperref[type-1-empirical-uncertainty]{297}}

\hyperref[type-2-metric-uncertainty]{Type 2: Metric Uncertainty
\hyperref[type-2-metric-uncertainty]{298}}

\hyperref[type-3-theory-uncertainty]{Type 3: Theory Uncertainty
\hyperref[type-3-theory-uncertainty]{299}}

\hyperref[representing-moral-uncertainty-tensorially]{16.3 Representing
Moral Uncertainty Tensorially
\hyperref[representing-moral-uncertainty-tensorially]{300}}

\hyperref[the-theory-space-tensor]{The Theory-Space Tensor
\hyperref[the-theory-space-tensor]{300}}

\hyperref[the-joint-uncertainty]{The Joint Uncertainty
\hyperref[the-joint-uncertainty]{300}}

\hyperref[the-moral-risk-of-theory-uncertainty]{The Moral Risk of Theory
Uncertainty \hyperref[the-moral-risk-of-theory-uncertainty]{300}}

\hyperref[decision-under-moral-uncertainty]{16.4 Decision Under Moral
Uncertainty \hyperref[decision-under-moral-uncertainty]{301}}

\hyperref[the-three-standard-approaches]{The Three Standard Approaches
\hyperref[the-three-standard-approaches]{301}}

\hyperref[geometric-comparison]{Geometric Comparison
\hyperref[geometric-comparison]{302}}

\hyperref[the-intertheoretic-comparison-problem]{16.5 The Intertheoretic
Comparison Problem
\hyperref[the-intertheoretic-comparison-problem]{302}}

\hyperref[the-core-difficulty]{The Core Difficulty
\hyperref[the-core-difficulty]{302}}

\hyperref[the-geometric-formulation]{The Geometric Formulation
\hyperref[the-geometric-formulation]{303}}

\hyperref[partial-comparability]{Partial Comparability
\hyperref[partial-comparability]{303}}

\hyperref[robust-obligations]{16.6 Robust Obligations
\hyperref[robust-obligations]{304}}

\hyperref[what-survives-moral-uncertainty]{What Survives Moral
Uncertainty \hyperref[what-survives-moral-uncertainty]{304}}

\hyperref[the-robust-core]{The Robust Core
\hyperref[the-robust-core]{304}}

\hyperref[stability-under-perturbation]{Stability Under Perturbation
\hyperref[stability-under-perturbation]{304}}

\hyperref[residual-indeterminacy]{16.7 Residual Indeterminacy
\hyperref[residual-indeterminacy]{305}}

\hyperref[what-remains-undecided]{What Remains Undecided
\hyperref[what-remains-undecided]{305}}

\hyperref[is-residual-indeterminacy-a-problem]{Is Residual Indeterminacy
a Problem? \hyperref[is-residual-indeterminacy-a-problem]{306}}

\hyperref[the-value-of-precise-indeterminacy]{The Value of Precise
Indeterminacy \hyperref[the-value-of-precise-indeterminacy]{306}}

\hyperref[decision-strategies-under-indeterminacy]{16.8 Decision
Strategies Under Indeterminacy
\hyperref[decision-strategies-under-indeterminacy]{306}}

\hyperref[when-the-framework-falls-silent]{When the Framework Falls
Silent \hyperref[when-the-framework-falls-silent]{306}}

\hyperref[moral-uncertainty-and-ai-systems]{16.9 Moral Uncertainty and
AI Systems \hyperref[moral-uncertainty-and-ai-systems]{307}}

\hyperref[why-ai-faces-moral-uncertainty-acutely]{Why AI Faces Moral
Uncertainty Acutely
\hyperref[why-ai-faces-moral-uncertainty-acutely]{307}}

\hyperref[design-implications]{Design Implications
\hyperref[design-implications]{307}}

\hyperref[the-modesty-of-the-framework]{16.10 The Modesty of the
Framework \hyperref[the-modesty-of-the-framework]{308}}

\hyperref[what-geometric-ethics-claims]{What Geometric Ethics Claims
\hyperref[what-geometric-ethics-claims]{308}}

\hyperref[why-modesty-is-a-strength]{Why Modesty Is a Strength
\hyperref[why-modesty-is-a-strength]{309}}

\hyperref[summary-9]{16.11 Summary \hyperref[summary-9]{309}}

\hyperref[technical-appendix-10]{Technical Appendix
\hyperref[technical-appendix-10]{310}}

\hyperref[part-v-implementation]{Part V: Implementation
\hyperref[part-v-implementation]{313}}

\hyperref[part-v-prelude-philosophy-engineering]{Part V Prelude:
Philosophy Engineering
\hyperref[part-v-prelude-philosophy-engineering]{313}}

\hyperref[definition-and-scope]{17.0.1 Definition and Scope
\hyperref[definition-and-scope]{313}}

\hyperref[the-epistemic-culture]{17.0.2 The Epistemic Culture
\hyperref[the-epistemic-culture]{314}}

\hyperref[core-theoretical-commitments]{17.0.3 Core Theoretical
Commitments \hyperref[core-theoretical-commitments]{314}}

\hyperref[the-philosophy-engineering-workflow]{17.0.4 The Philosophy
Engineering Workflow
\hyperref[the-philosophy-engineering-workflow]{315}}

\hyperref[required-artifacts-the-philosophy-engineering-dossier]{17.0.5
Required Artifacts: The Philosophy Engineering Dossier
\hyperref[required-artifacts-the-philosophy-engineering-dossier]{316}}

\hyperref[threat-model-and-failure-modes]{17.0.6 Threat Model and
Failure Modes \hyperref[threat-model-and-failure-modes]{316}}

\hyperref[professional-standards]{17.0.7 Professional Standards
\hyperref[professional-standards]{317}}

\hyperref[practice-areas]{17.0.8 Practice Areas
\hyperref[practice-areas]{317}}

\hyperref[roles]{17.0.9 Roles \hyperref[roles]{318}}

\hyperref[research-agenda]{17.0.10 Research Agenda
\hyperref[research-agenda]{318}}

\hyperref[chapter-17-empirical-evidence-for-geometric-ethics]{Chapter
17: Empirical Evidence for Geometric Ethics
\hyperref[chapter-17-empirical-evidence-for-geometric-ethics]{319}}

\hyperref[from-theory-to-data]{17.1 From Theory to Data
\hyperref[from-theory-to-data]{319}}

\hyperref[the-dear-abby-corpus]{17.2 The Dear Abby Corpus
\hyperref[the-dear-abby-corpus]{320}}

\hyperref[the-dataset]{The Dataset \hyperref[the-dataset]{320}}

\hyperref[what-the-corpus-reveals]{What the Corpus Reveals
\hyperref[what-the-corpus-reveals]{320}}

\hyperref[finding-1-context-dependent-dimension-weighting]{Finding 1:
Context-Dependent Dimension Weighting
\hyperref[finding-1-context-dependent-dimension-weighting]{321}}

\hyperref[finding-2-temporal-stability-of-structure]{Finding 2: Temporal
Stability of Structure
\hyperref[finding-2-temporal-stability-of-structure]{321}}

\hyperref[finding-3-semantic-gate-discreteness]{Finding 3: Semantic Gate
Discreteness \hyperref[finding-3-semantic-gate-discreteness]{322}}

\hyperref[finding-4-nullifier-universality]{Finding 4: Nullifier
Universality \hyperref[finding-4-nullifier-universality]{322}}

\hyperref[the-bip-experiments]{17.3 The BIP Experiments
\hyperref[the-bip-experiments]{323}}

\hyperref[experimental-design]{Experimental Design
\hyperref[experimental-design]{323}}

\hyperref[key-results]{Key Results \hyperref[key-results]{323}}

\hyperref[interpretation]{Interpretation \hyperref[interpretation]{324}}

\hyperref[the-anomaly]{The Anomaly \hyperref[the-anomaly]{325}}

\hyperref[the-dear-ethicist-game]{17.4 The Dear Ethicist Game
\hyperref[the-dear-ethicist-game]{325}}

\hyperref[design]{Design \hyperref[design]{325}}

\hyperref[the-bond-index]{The Bond Index \hyperref[the-bond-index]{326}}

\hyperref[preliminary-findings]{Preliminary Findings
\hyperref[preliminary-findings]{326}}

\hyperref[quantum-cognition-predictions]{17.5 Quantum Cognition
Predictions \hyperref[quantum-cognition-predictions]{327}}

\hyperref[the-connection-1]{The Connection
\hyperref[the-connection-1]{327}}

\hyperref[prediction-1-violation-of-the-law-of-total-probability]{Prediction
1: Violation of the Law of Total Probability
\hyperref[prediction-1-violation-of-the-law-of-total-probability]{327}}

\hyperref[prediction-2-order-effects-in-moral-judgment]{Prediction 2:
Order Effects in Moral Judgment
\hyperref[prediction-2-order-effects-in-moral-judgment]{328}}

\hyperref[prediction-3-interference-in-deliberation]{Prediction 3:
Interference in Deliberation
\hyperref[prediction-3-interference-in-deliberation]{328}}

\hyperref[the-empirical-program]{The Empirical Program
\hyperref[the-empirical-program]{328}}

\hyperref[evidence-for-the-invariance-of-harm-accounting]{17.6 Evidence
for the Invariance of Harm Accounting
\hyperref[evidence-for-the-invariance-of-harm-accounting]{329}}

\hyperref[the-noether-prediction]{The Noether Prediction
\hyperref[the-noether-prediction]{329}}

\hyperref[the-evidence-1]{The Evidence \hyperref[the-evidence-1]{329}}

\hyperref[what-the-evidence-does-not-show]{17.7 What the Evidence Does
Not Show \hyperref[what-the-evidence-does-not-show]{329}}

\hyperref[honest-limitations]{Honest Limitations
\hyperref[honest-limitations]{329}}

\hyperref[what-would-falsify-the-framework]{What Would Falsify the
Framework? \hyperref[what-would-falsify-the-framework]{330}}

\hyperref[threats-to-validity]{Threats to Validity
\hyperref[threats-to-validity]{331}}

\hyperref[the-shape-of-the-evidence]{17.8 The Shape of the Evidence
\hyperref[the-shape-of-the-evidence]{332}}

\hyperref[a-summary-table]{A Summary Table
\hyperref[a-summary-table]{332}}

\hyperref[implications-for-the-framework]{17.9 Implications for the
Framework \hyperref[implications-for-the-framework]{333}}

\hyperref[what-the-evidence-establishes]{What the Evidence Establishes
\hyperref[what-the-evidence-establishes]{333}}

\hyperref[what-the-evidence-does-not-establish]{What the Evidence Does
Not Establish \hyperref[what-the-evidence-does-not-establish]{333}}

\hyperref[update-bip-v10.16-quantitative-validation-february-2026]{17.10
Update: BIP v10.16 Quantitative Validation (February 2026)
\hyperref[update-bip-v10.16-quantitative-validation-february-2026]{333}}

\hyperref[methodology]{Methodology \hyperref[methodology]{334}}

\hyperref[key-quantitative-results]{Key Quantitative Results
\hyperref[key-quantitative-results]{334}}

\hyperref[qnd-order-effects]{QND Order Effects
\hyperref[qnd-order-effects]{334}}

\hyperref[moral-axis-variation-structure]{\textbf{Moral Axis Variation
Structure} \hyperref[moral-axis-variation-structure]{335}}

\hyperref[remaining-challenge]{Remaining Challenge
\hyperref[remaining-challenge]{335}}

\hyperref[a-independent-replication-thiele-2026]{17.10a Independent
Replication: Thiele (2026)
\hyperref[a-independent-replication-thiele-2026]{336}}

\hyperref[methodology-1]{Methodology \hyperref[methodology-1]{336}}

\hyperref[per-dimension-probe-results]{Per-Dimension Probe Results
\hyperref[per-dimension-probe-results]{336}}

\hyperref[cross-lingual-transfer]{Cross-Lingual Transfer
\hyperref[cross-lingual-transfer]{336}}

\hyperref[revised-language-leakage]{Revised Language Leakage
\hyperref[revised-language-leakage]{337}}

\hyperref[the-orthogonality-finding]{The Orthogonality Finding
\hyperref[the-orthogonality-finding]{337}}

\hyperref[significance-for-the-framework]{Significance for the Framework
\hyperref[significance-for-the-framework]{337}}

\hyperref[fuzz-testing-as-a-discovery-method]{17.11 Fuzz Testing as a
Discovery Method \hyperref[fuzz-testing-as-a-discovery-method]{337}}

\hyperref[why-fuzz-ethics]{Why Fuzz Ethics?
\hyperref[why-fuzz-ethics]{337}}

\hyperref[the-fuzz-design-space]{The Fuzz Design Space
\hyperref[the-fuzz-design-space]{338}}

\hyperref[two-fuzzing-strategies]{Two Fuzzing Strategies
\hyperref[two-fuzzing-strategies]{338}}

\hyperref[discoveries]{Discoveries \hyperref[discoveries]{339}}

\hyperref[from-discovery-to-formalization]{From Discovery to
Formalization \hyperref[from-discovery-to-formalization]{340}}

\hyperref[the-road-ahead]{17.12 The Road Ahead
\hyperref[the-road-ahead]{341}}

\hyperref[an-empirical-program-for-geometric-ethics]{An Empirical
Program for Geometric Ethics
\hyperref[an-empirical-program-for-geometric-ethics]{341}}

\hyperref[summary-10]{17.13 Summary \hyperref[summary-10]{342}}

\hyperref[coda-the-philosophy-engineering-code-corpus]{17.14 Coda: The
Philosophy Engineering Code Corpus
\hyperref[coda-the-philosophy-engineering-code-corpus]{342}}

\hyperref[the-deployment-layer-erisml-lib-1.51m-lines-of-code]{The
Deployment Layer: erisml-lib (1.51M Lines of Code)
\hyperref[the-deployment-layer-erisml-lib-1.51m-lines-of-code]{343}}

\hyperref[the-experimental-probe-sqnd-probe-1.26m-lines-of-code]{The
Experimental Probe: sqnd-probe (1.26M Lines of Code)
\hyperref[the-experimental-probe-sqnd-probe-1.26m-lines-of-code]{343}}

\hyperref[revised-summary-of-empirical-evidence]{Revised Summary of
Empirical Evidence
\hyperref[revised-summary-of-empirical-evidence]{343}}

\hyperref[technical-appendix-methods-and-reproducibility]{Technical
Appendix: Methods and Reproducibility
\hyperref[technical-appendix-methods-and-reproducibility]{343}}

\hyperref[dear-abby-corpus-construction]{Dear Abby Corpus Construction
\hyperref[dear-abby-corpus-construction]{343}}

\hyperref[bip-corpus-construction]{BIP Corpus Construction
\hyperref[bip-corpus-construction]{344}}

\hyperref[classifier-architecture]{Classifier Architecture
\hyperref[classifier-architecture]{344}}

\hyperref[statistical-reporting]{Statistical Reporting
\hyperref[statistical-reporting]{345}}

\hyperref[internal-validation-protocol]{Internal Validation Protocol
\hyperref[internal-validation-protocol]{345}}

\hyperref[public-artifact-plan]{Public Artifact Plan
\hyperref[public-artifact-plan]{345}}

\hyperref[chapter-18-geometric-ethics-for-artificial-agents]{Chapter 18:
Geometric Ethics for Artificial Agents
\hyperref[chapter-18-geometric-ethics-for-artificial-agents]{347}}

\hyperref[why-ai-needs-geometry]{18.1 Why AI Needs Geometry
\hyperref[why-ai-needs-geometry]{347}}

\hyperref[the-problem-with-scalar-ai-ethics]{18.2 The Problem with
Scalar AI Ethics \hyperref[the-problem-with-scalar-ai-ethics]{348}}

\hyperref[tensor-valued-objectives]{18.3 Tensor-Valued Objectives
\hyperref[tensor-valued-objectives]{349}}

\hyperref[the-basic-architecture]{The Basic Architecture
\hyperref[the-basic-architecture]{349}}

\hyperref[tensor-valued-reward]{Tensor-Valued Reward
\hyperref[tensor-valued-reward]{349}}

\hyperref[contraction-as-a-separate-step]{Contraction as a Separate Step
\hyperref[contraction-as-a-separate-step]{349}}

\hyperref[the-residue-is-logged]{The Residue Is Logged
\hyperref[the-residue-is-logged]{350}}

\hyperref[invariance-as-alignment]{18.4 Invariance as Alignment
\hyperref[invariance-as-alignment]{351}}

\hyperref[alignment-as-invariance]{Alignment as Invariance
\hyperref[alignment-as-invariance]{351}}

\hyperref[invariance-testing-as-alignment-auditing]{Invariance Testing
as Alignment Auditing
\hyperref[invariance-testing-as-alignment-auditing]{351}}

\hyperref[the-noether-criterion]{The Noether Criterion
\hyperref[the-noether-criterion]{352}}

\hyperref[the-alignment-problem-as-contraction-mismatch-1]{18.5 The
Alignment Problem as Contraction Mismatch
\hyperref[the-alignment-problem-as-contraction-mismatch-1]{352}}

\hyperref[three-sources-of-misalignment]{Three Sources of Misalignment
\hyperref[three-sources-of-misalignment]{352}}

\hyperref[the-fundamental-formula-as-alignment-criterion]{The
Fundamental Formula as Alignment Criterion
\hyperref[the-fundamental-formula-as-alignment-criterion]{353}}

\hyperref[the-no-escape-theorem]{18.6 The No Escape Theorem
\hyperref[the-no-escape-theorem]{353}}

\hyperref[the-problem-of-cognitive-escape]{The Problem of Cognitive
Escape \hyperref[the-problem-of-cognitive-escape]{353}}

\hyperref[the-structural-containment-architecture]{The Structural
Containment Architecture
\hyperref[the-structural-containment-architecture]{354}}

\hyperref[statement-and-proof]{Statement and Proof
\hyperref[statement-and-proof]{355}}

\hyperref[the-key-insight]{The Key Insight
\hyperref[the-key-insight]{356}}

\hyperref[pre-flight-checklist-no-escape-theorem-theorem-18.1]{Pre-flight
Checklist: No Escape Theorem (Theorem 18.1)
\hyperref[pre-flight-checklist-no-escape-theorem-theorem-18.1]{356}}

\hyperref[scope-and-limitations-of-the-no-escape-theorem]{18.7 Scope and
Limitations of the No Escape Theorem
\hyperref[scope-and-limitations-of-the-no-escape-theorem]{357}}

\hyperref[what-the-theorem-blocks]{What the Theorem Blocks
\hyperref[what-the-theorem-blocks]{357}}

\hyperref[what-the-theorem-does-not-block]{What the Theorem Does Not
Block \hyperref[what-the-theorem-does-not-block]{357}}

\hyperref[the-safety-reduction]{The Safety Reduction
\hyperref[the-safety-reduction]{358}}

\hyperref[feasibility-by-domain]{Feasibility by Domain
\hyperref[feasibility-by-domain]{358}}

\hyperref[operational-adequacy-checklist-for-ux3c8]{Operational Adequacy
Checklist for Ψ
\hyperref[operational-adequacy-checklist-for-ux3c8]{359}}

\hyperref[formal-consequences-of-grounding-adequacy]{Formal Consequences
of Grounding Adequacy
\hyperref[formal-consequences-of-grounding-adequacy]{360}}

\hyperref[what-the-no-escape-theorem-does-not-cover]{What the No Escape
Theorem Does Not Cover
\hyperref[what-the-no-escape-theorem-does-not-cover]{361}}

\hyperref[connection-to-conservation-laws]{18.8 Connection to
Conservation Laws \hyperref[connection-to-conservation-laws]{362}}

\hyperref[two-levels-of-guarantee]{Two Levels of Guarantee
\hyperref[two-levels-of-guarantee]{362}}

\hyperref[gauge-theory-analogy]{Gauge Theory Analogy
\hyperref[gauge-theory-analogy]{363}}

\hyperref[escape-route-analysis]{18.9 Escape Route Analysis
\hyperref[escape-route-analysis]{363}}

\hyperref[relabeling-and-redescription]{18.9.1 Relabeling and
Redescription \hyperref[relabeling-and-redescription]{363}}

\hyperref[specification-gaming]{18.9.2 Specification Gaming
\hyperref[specification-gaming]{363}}

\hyperref[reward-hacking]{18.9.3 Reward Hacking
\hyperref[reward-hacking]{364}}

\hyperref[deceptive-alignment]{18.9.4 Deceptive Alignment
\hyperref[deceptive-alignment]{364}}

\hyperref[operator-manipulation]{18.9.5 Operator Manipulation
\hyperref[operator-manipulation]{364}}

\hyperref[verifying-the-requirements-the-epistemic-invariance-architecture]{18.9.6
Verifying the Requirements: The Epistemic Invariance Architecture
\hyperref[verifying-the-requirements-the-epistemic-invariance-architecture]{365}}

\hyperref[implementing-geometric-ai-ethics]{18.10 Implementing Geometric
AI Ethics \hyperref[implementing-geometric-ai-ethics]{366}}

\hyperref[representation]{Representation \hyperref[representation]{366}}

\hyperref[training]{Training \hyperref[training]{366}}

\hyperref[deployment]{Deployment \hyperref[deployment]{367}}

\hyperref[case-study-tensorial-kidney-allocation]{18.11 Case Study:
Tensorial Kidney Allocation
\hyperref[case-study-tensorial-kidney-allocation]{367}}

\hyperref[the-scalar-approach-1]{The Scalar Approach
\hyperref[the-scalar-approach-1]{367}}

\hyperref[the-tensorial-approach]{The Tensorial Approach
\hyperref[the-tensorial-approach]{367}}

\hyperref[invariance-verification]{Invariance Verification
\hyperref[invariance-verification]{368}}

\hyperref[the-no-escape-guarantee]{The No Escape Guarantee
\hyperref[the-no-escape-guarantee]{368}}

\hyperref[human-ai-collaboration-as-tensor-combination]{18.12 Human-AI
Collaboration as Tensor Combination
\hyperref[human-ai-collaboration-as-tensor-combination]{368}}

\hyperref[scalar-combination-and-its-limits]{Scalar Combination and Its
Limits \hyperref[scalar-combination-and-its-limits]{368}}

\hyperref[tensorial-combination]{Tensorial Combination
\hyperref[tensorial-combination]{369}}

\hyperref[responsibility-in-human-ai-teams]{Responsibility in Human-AI
Teams \hyperref[responsibility-in-human-ai-teams]{369}}

\hyperref[looking-forward]{18.13 Looking Forward
\hyperref[looking-forward]{369}}

\hyperref[technical-appendix-11]{Technical Appendix
\hyperref[technical-appendix-11]{370}}

\hyperref[chapter-19-the-deme-architecture-and-erisml]{Chapter 19: The
DEME Architecture and ErisML
\hyperref[chapter-19-the-deme-architecture-and-erisml]{372}}

\hyperref[from-theory-to-infrastructure]{19.1 From Theory to
Infrastructure \hyperref[from-theory-to-infrastructure]{372}}

\hyperref[the-engineering-challenge]{19.2 The Engineering Challenge
\hyperref[the-engineering-challenge]{374}}

\hyperref[erisml-a-modeling-language-for-geometric-ethics]{19.3 ErisML:
A Modeling Language for Geometric Ethics
\hyperref[erisml-a-modeling-language-for-geometric-ethics]{375}}

\hyperref[environment-specification]{Environment Specification
\hyperref[environment-specification]{375}}

\hyperref[agent-specification]{Agent Specification
\hyperref[agent-specification]{375}}

\hyperref[norms-specification]{Norms Specification
\hyperref[norms-specification]{376}}

\hyperref[formal-semantics]{Formal Semantics
\hyperref[formal-semantics]{377}}

\hyperref[erisml-compilation-and-complexity]{19.4 ErisML Compilation and
Complexity \hyperref[erisml-compilation-and-complexity]{378}}

\hyperref[compilation-targets]{Compilation Targets
\hyperref[compilation-targets]{378}}

\hyperref[computational-complexity]{Computational Complexity
\hyperref[computational-complexity]{378}}

\hyperref[the-translation-layer]{19.5 The Translation Layer
\hyperref[the-translation-layer]{379}}

\hyperref[the-problem-of-ethical-pluralism]{The Problem of Ethical
Pluralism \hyperref[the-problem-of-ethical-pluralism]{379}}

\hyperref[policy-module-architecture]{Policy Module Architecture
\hyperref[policy-module-architecture]{379}}

\hyperref[geometric-interpretation]{Geometric Interpretation
\hyperref[geometric-interpretation]{381}}

\hyperref[deme-the-ethics-engine]{19.6 DEME: The Ethics Engine
\hyperref[deme-the-ethics-engine]{381}}

\hyperref[the-four-layer-architecture]{The Four-Layer Architecture
\hyperref[the-four-layer-architecture]{381}}

\hyperref[the-geneva-ethics-module]{The Geneva Ethics Module
\hyperref[the-geneva-ethics-module]{383}}

\hyperref[the-bond-index-1]{19.7 The Bond Index
\hyperref[the-bond-index-1]{384}}

\hyperref[definition-6]{Definition \hyperref[definition-6]{384}}

\hyperref[geometric-interpretation-1]{Geometric Interpretation
\hyperref[geometric-interpretation-1]{384}}

\hyperref[testing-methodology]{Testing Methodology
\hyperref[testing-methodology]{385}}

\hyperref[the-separation-principle]{19.8 The Separation Principle
\hyperref[the-separation-principle]{385}}

\hyperref[the-principle]{The Principle \hyperref[the-principle]{385}}

\hyperref[connection-to-the-tensor-hierarchy]{Connection to the Tensor
Hierarchy \hyperref[connection-to-the-tensor-hierarchy]{386}}

\hyperref[the-norm-kernel-compiling-normative-constraints-into-verified-automata]{The
Norm Kernel: Compiling Normative Constraints into Verified Automata
\hyperref[the-norm-kernel-compiling-normative-constraints-into-verified-automata]{388}}

\hyperref[democratic-governance-as-metric-selection]{19.9 Democratic
Governance as Metric Selection
\hyperref[democratic-governance-as-metric-selection]{389}}

\hyperref[the-governance-account-implemented]{The Governance Account,
Implemented \hyperref[the-governance-account-implemented]{389}}

\hyperref[metric-pluralism-in-practice]{Metric Pluralism in Practice
\hyperref[metric-pluralism-in-practice]{389}}

\hyperref[connection-to-contraction]{Connection to Contraction
\hyperref[connection-to-contraction]{390}}

\hyperref[the-grand-unified-ai-safety-stack]{19.10 The Grand Unified AI
Safety Stack \hyperref[the-grand-unified-ai-safety-stack]{390}}

\hyperref[the-seven-layer-architecture]{The Seven-Layer Architecture
\hyperref[the-seven-layer-architecture]{390}}

\hyperref[the-safety-stack-as-a-fiber-bundle]{The Safety Stack as a
Fiber Bundle \hyperref[the-safety-stack-as-a-fiber-bundle]{391}}

\hyperref[case-study-autonomous-hospital-ward]{19.11 Case Study:
Autonomous Hospital Ward
\hyperref[case-study-autonomous-hospital-ward]{391}}

\hyperref[the-setting]{The Setting \hyperref[the-setting]{392}}

\hyperref[erisml-specification]{ErisML Specification
\hyperref[erisml-specification]{392}}

\hyperref[ethics-modules]{Ethics Modules \hyperref[ethics-modules]{392}}

\hyperref[governance-profile]{Governance Profile
\hyperref[governance-profile]{392}}

\hyperref[a-decision-scenario]{A Decision Scenario
\hyperref[a-decision-scenario]{393}}

\hyperref[looking-forward-1]{19.12 Looking Forward
\hyperref[looking-forward-1]{394}}

\hyperref[tensor-structured-canonicalization-from-flat-encoders-to-geometric-pipelines]{19.13
Tensor-Structured Canonicalization: From Flat Encoders to Geometric
Pipelines
\hyperref[tensor-structured-canonicalization-from-flat-encoders-to-geometric-pipelines]{394}}

\hyperref[proposal-1-rank-2-internal-representation-tensor]{Proposal 1:
Rank-2 Internal Representation Tensor
\hyperref[proposal-1-rank-2-internal-representation-tensor]{395}}

\hyperref[proposal-2-rank-3-uncertainty-aware-canonicalization]{Proposal
2: Rank-3 Uncertainty-Aware Canonicalization
\hyperref[proposal-2-rank-3-uncertainty-aware-canonicalization]{396}}

\hyperref[proposal-3-rank-5-coalition-parsing-for-multi-agent-inputs]{Proposal
3: Rank-5 Coalition Parsing for Multi-Agent Inputs
\hyperref[proposal-3-rank-5-coalition-parsing-for-multi-agent-inputs]{397}}

\hyperref[proposal-4-rank-4-temporal-tracking-and-moral-holonomy]{Proposal
4: Rank-4 Temporal Tracking and Moral Holonomy
\hyperref[proposal-4-rank-4-temporal-tracking-and-moral-holonomy]{398}}

\hyperref[synthesis-and-architectural-integration]{Synthesis and
Architectural Integration
\hyperref[synthesis-and-architectural-integration]{399}}

\hyperref[technical-appendix-12]{Technical Appendix
\hyperref[technical-appendix-12]{401}}

\hyperref[part-vi-domain-applications]{Part VI: Domain Applications
\hyperref[part-vi-domain-applications]{402}}

\hyperref[methodological-notes-for-domain-applications]{Methodological
Notes for Domain Applications
\hyperref[methodological-notes-for-domain-applications]{402}}

\hyperref[chapter-20-geometric-economics-the-bond-geodesic-equilibrium]{Chapter
20: Geometric Economics --- The Bond Geodesic Equilibrium
\hyperref[chapter-20-geometric-economics-the-bond-geodesic-equilibrium]{404}}

\hyperref[the-failure-of-scalar-economics]{20.1 The Failure of Scalar
Economics \hyperref[the-failure-of-scalar-economics]{404}}

\hyperref[the-economic-decision-complex]{20.2 The Economic Decision
Complex \hyperref[the-economic-decision-complex]{405}}

\hyperref[heuristic-bounded-rationality]{20.3 Heuristic Bounded
Rationality \hyperref[heuristic-bounded-rationality]{405}}

\hyperref[the-bond-geodesic-and-the-scalar-irrecoverability-theorem]{20.4
The Bond Geodesic and the Scalar Irrecoverability Theorem
\hyperref[the-bond-geodesic-and-the-scalar-irrecoverability-theorem]{406}}

\hyperref[the-bond-geodesic-equilibrium]{20.5 The Bond Geodesic
Equilibrium \hyperref[the-bond-geodesic-equilibrium]{406}}

\hyperref[the-bgenash-relationship]{The BGE--Nash Relationship
\hyperref[the-bgenash-relationship]{407}}

\hyperref[existence-and-uniqueness]{Existence and Uniqueness
\hyperref[existence-and-uniqueness]{407}}

\hyperref[the-contraction-lemma]{The Contraction Lemma
\hyperref[the-contraction-lemma]{408}}

\hyperref[welfare-properties]{20.6 Welfare Properties
\hyperref[welfare-properties]{408}}

\hyperref[behavioral-game-theory-as-manifold-geometry]{20.7 Behavioral
Game Theory as Manifold Geometry
\hyperref[behavioral-game-theory-as-manifold-geometry]{409}}

\hyperref[prospect-theory-as-manifold-geometry]{20.8 Prospect Theory as
Manifold Geometry \hyperref[prospect-theory-as-manifold-geometry]{409}}

\hyperref[attribute-conservation-in-bilateral-exchange]{20.9 Attribute
Conservation in Bilateral Exchange
\hyperref[attribute-conservation-in-bilateral-exchange]{410}}

\hyperref[worked-examples]{20.10 Worked Examples
\hyperref[worked-examples]{410}}

\hyperref[example-20.1-uber-surge-pricing-why-efficient-pricing-fails-on-the-manifold]{Example
20.1: Uber Surge Pricing --- Why "Efficient" Pricing Fails on the
Manifold
\hyperref[example-20.1-uber-surge-pricing-why-efficient-pricing-fails-on-the-manifold]{410}}

\hyperref[example-20.2-the-dictator-game-across-cultures-metric-variation]{Example
20.2: The Dictator Game across Cultures --- Metric Variation
\hyperref[example-20.2-the-dictator-game-across-cultures-metric-variation]{411}}

\hyperref[example-20.3-pharmaceutical-price-gouging-turing-pharmaceuticals-and-daraprim-2015]{Example
20.3: Pharmaceutical Price Gouging --- Turing Pharmaceuticals and
Daraprim (2015)
\hyperref[example-20.3-pharmaceutical-price-gouging-turing-pharmaceuticals-and-daraprim-2015]{412}}

\hyperref[falsifiable-predictions]{20.11 Falsifiable Predictions
\hyperref[falsifiable-predictions]{413}}

\hyperref[connection-to-the-framework-1]{20.12 Connection to the
Framework \hyperref[connection-to-the-framework-1]{414}}

\hyperref[summary-11]{20.13 Summary \hyperref[summary-11]{415}}

\hyperref[chapter-21-geometric-clinical-ethics-triage-consent-and-the-clinical-geodesic]{Chapter
21: Geometric Clinical Ethics --- Triage, Consent, and the Clinical
Geodesic
\hyperref[chapter-21-geometric-clinical-ethics-triage-consent-and-the-clinical-geodesic]{415}}

\hyperref[the-failure-of-scalar-clinical-ethics]{21.1 The Failure of
Scalar Clinical Ethics
\hyperref[the-failure-of-scalar-clinical-ethics]{416}}

\hyperref[the-clinical-decision-complex]{21.2 The Clinical Decision
Complex \hyperref[the-clinical-decision-complex]{416}}

\hyperref[clinical-decision-as-a-search]{21.3 Clinical Decision as A*
Search \hyperref[clinical-decision-as-a-search]{418}}

\hyperref[the-qaly-irrecoverability-theorem]{21.4 The QALY
Irrecoverability Theorem
\hyperref[the-qaly-irrecoverability-theorem]{419}}

\hyperref[the-mathematical-theory-of-moral-injury]{21.5 The Mathematical
Theory of Moral Injury
\hyperref[the-mathematical-theory-of-moral-injury]{420}}

\hyperref[geometric-informed-consent]{21.6 Geometric Informed Consent
\hyperref[geometric-informed-consent]{421}}

\hyperref[geometric-triage]{21.7 Geometric Triage
\hyperref[geometric-triage]{422}}

\hyperref[the-clinical-bond-index]{21.8 The Clinical Bond Index
\hyperref[the-clinical-bond-index]{423}}

\hyperref[worked-examples-1]{21.9 Worked Examples
\hyperref[worked-examples-1]{423}}

\hyperref[falsifiable-predictions-1]{21.10 Falsifiable Predictions
\hyperref[falsifiable-predictions-1]{424}}

\hyperref[connection-to-the-framework-2]{21.11 Connection to the
Framework \hyperref[connection-to-the-framework-2]{425}}

\hyperref[summary-12]{21.12 Summary \hyperref[summary-12]{425}}

\hyperref[chapter-22-geometric-jurisprudence-algorithmic-justice-and-legal-manifolds]{Chapter
22: Geometric Jurisprudence --- Algorithmic Justice and Legal Manifolds
\hyperref[chapter-22-geometric-jurisprudence-algorithmic-justice-and-legal-manifolds]{426}}

\hyperref[law-as-geometric-structure]{22.1 Law as Geometric Structure
\hyperref[law-as-geometric-structure]{426}}

\hyperref[the-judicial-complex]{22.2 The Judicial Complex
\hyperref[the-judicial-complex]{426}}

\hyperref[legal-invariance-principles]{22.3 Legal Invariance Principles
\hyperref[legal-invariance-principles]{428}}

\hyperref[the-hohfeldian-octad-and-gauge-theory]{22.4 The Hohfeldian
Octad and Gauge Theory
\hyperref[the-hohfeldian-octad-and-gauge-theory]{428}}

\hyperref[topological-constitutionality]{22.5 Topological
Constitutionality \hyperref[topological-constitutionality]{429}}

\hyperref[conservation-laws-in-legal-systems]{22.6 Conservation Laws in
Legal Systems \hyperref[conservation-laws-in-legal-systems]{430}}

\hyperref[legal-disputes-as-a-pathfinding]{22.7 Legal Disputes as A*
Pathfinding \hyperref[legal-disputes-as-a-pathfinding]{431}}

\hyperref[precedent-as-weight-deformation]{22.8 Precedent as Weight
Deformation \hyperref[precedent-as-weight-deformation]{432}}

\hyperref[empirical-calibration-pipeline]{22.9 Empirical Calibration
Pipeline \hyperref[empirical-calibration-pipeline]{432}}

\hyperref[worked-examples-2]{22.10 Worked Examples
\hyperref[worked-examples-2]{433}}

\hyperref[falsifiable-predictions-2]{22.11 Falsifiable Predictions
\hyperref[falsifiable-predictions-2]{435}}

\hyperref[connection-to-the-framework-3]{22.12 Connection to the
Framework \hyperref[connection-to-the-framework-3]{435}}

\hyperref[summary-13]{22.13 Summary \hyperref[summary-13]{436}}

\hyperref[chapter-23-geometric-finance-market-microstructure-on-the-decision-manifold]{Chapter
23: Geometric Finance --- Market Microstructure on the Decision Manifold
\hyperref[chapter-23-geometric-finance-market-microstructure-on-the-decision-manifold]{436}}

\hyperref[the-failure-of-scalar-finance]{23.1 The Failure of Scalar
Finance \hyperref[the-failure-of-scalar-finance]{437}}

\hyperref[the-financial-decision-complex]{23.2 The Financial Decision
Complex \hyperref[the-financial-decision-complex]{437}}

\hyperref[market-microstructure-as-manifold-geometry]{23.3 Market
Microstructure as Manifold Geometry
\hyperref[market-microstructure-as-manifold-geometry]{438}}

\hyperref[the-glosten-milgrom-model]{The Glosten-Milgrom Model
\hyperref[the-glosten-milgrom-model]{439}}

\hyperref[the-flash-crash-as-dimensional-collapse]{23.4 The Flash Crash
as Dimensional Collapse
\hyperref[the-flash-crash-as-dimensional-collapse]{440}}

\hyperref[high-frequency-trading-and-the-fairness-problem]{23.5
High-Frequency Trading and the Fairness Problem
\hyperref[high-frequency-trading-and-the-fairness-problem]{440}}

\hyperref[geometric-risk-management]{23.6 Geometric Risk Management
\hyperref[geometric-risk-management]{441}}

\hyperref[financial-conservation-laws]{23.7 Financial Conservation Laws
\hyperref[financial-conservation-laws]{442}}

\hyperref[geometric-behavioral-finance]{23.8 Geometric Behavioral
Finance \hyperref[geometric-behavioral-finance]{443}}

\hyperref[option-pricing-as-scalar-projection]{23.9 Option Pricing as
Scalar Projection \hyperref[option-pricing-as-scalar-projection]{443}}

\hyperref[the-2008-financial-crisis-as-manifold-failure]{23.10 The 2008
Financial Crisis as Manifold Failure
\hyperref[the-2008-financial-crisis-as-manifold-failure]{444}}

\hyperref[falsifiable-predictions-3]{23.11 Falsifiable Predictions
\hyperref[falsifiable-predictions-3]{445}}

\hyperref[connection-to-the-framework-4]{23.12 Connection to the
Framework \hyperref[connection-to-the-framework-4]{445}}

\hyperref[summary-14]{23.13 Summary \hyperref[summary-14]{446}}

\hyperref[chapter-24-geometric-theology-religious-reasoning-on-the-moral-manifold]{Chapter
24: Geometric Theology --- Religious Reasoning on the Moral Manifold
\hyperref[chapter-24-geometric-theology-religious-reasoning-on-the-moral-manifold]{447}}

\hyperref[religious-moral-reasoning-on-the-manifold]{24.1 Religious
Moral Reasoning on the Manifold
\hyperref[religious-moral-reasoning-on-the-manifold]{447}}

\hyperref[the-obligation-vector-and-natural-theology]{24.2 The
Obligation Vector and Natural Theology
\hyperref[the-obligation-vector-and-natural-theology]{448}}

\hyperref[the-epistemic-paradox-of-the-fall]{24.3 The Epistemic Paradox
of the Fall \hyperref[the-epistemic-paradox-of-the-fall]{449}}

\hyperref[the-euthyphro-dilemma-as-gauge-ambiguity]{24.4 The Euthyphro
Dilemma as Gauge Ambiguity
\hyperref[the-euthyphro-dilemma-as-gauge-ambiguity]{450}}

\hyperref[theodicy-as-dimensional-projection]{24.5 Theodicy as
Dimensional Projection
\hyperref[theodicy-as-dimensional-projection]{450}}

\hyperref[sacred-value-boundaries-and-religious-heuristics]{24.6
Sacred-Value Boundaries and Religious Heuristics
\hyperref[sacred-value-boundaries-and-religious-heuristics]{451}}

\hyperref[deductive-natural-theology-and-gauge-reification]{24.7
Deductive Natural Theology and Gauge Reification
\hyperref[deductive-natural-theology-and-gauge-reification]{452}}

\hyperref[implications-for-specific-religious-traditions]{24.8
Implications for Specific Religious Traditions
\hyperref[implications-for-specific-religious-traditions]{453}}

\hyperref[genesis-322-and-the-moral-competence-dilemma]{24.9 Genesis
3:22 and the Moral Competence Dilemma
\hyperref[genesis-322-and-the-moral-competence-dilemma]{454}}

\hyperref[falsifiable-predictions-4]{24.10 Falsifiable Predictions
\hyperref[falsifiable-predictions-4]{455}}

\hyperref[connection-to-the-framework-5]{24.11 Connection to the
Framework \hyperref[connection-to-the-framework-5]{455}}

\hyperref[summary-15]{24.12 Summary \hyperref[summary-15]{456}}

\hyperref[chapter-25-geometric-environmental-ethics-climate-commons-and-intergenerational-obligation-on-the-moral-manifold]{Chapter
25: Geometric Environmental Ethics --- Climate, Commons, and
Intergenerational Obligation on the Moral Manifold
\hyperref[chapter-25-geometric-environmental-ethics-climate-commons-and-intergenerational-obligation-on-the-moral-manifold]{457}}

\hyperref[the-failure-of-scalar-environmental-policy]{25.1 The Failure
of Scalar Environmental Policy
\hyperref[the-failure-of-scalar-environmental-policy]{457}}

\hyperref[the-environmental-decision-complex]{25.2 The Environmental
Decision Complex \hyperref[the-environmental-decision-complex]{458}}

\hyperref[intergenerational-pathfinding]{25.3 Intergenerational
Pathfinding \hyperref[intergenerational-pathfinding]{460}}

\hyperref[the-discount-rate-as-dimensional-collapse]{25.4 The Discount
Rate as Dimensional Collapse
\hyperref[the-discount-rate-as-dimensional-collapse]{461}}

\hyperref[the-tragedy-of-the-commons-as-multi-agent-manifold-failure]{25.5
The Tragedy of the Commons as Multi-Agent Manifold Failure
\hyperref[the-tragedy-of-the-commons-as-multi-agent-manifold-failure]{462}}

\hyperref[carbon-pricing-as-boundary-penalty-calibration]{25.6 Carbon
Pricing as Boundary Penalty Calibration
\hyperref[carbon-pricing-as-boundary-penalty-calibration]{463}}

\hyperref[species-extinction-as-irreversible-boundary-crossing]{25.7
Species Extinction as Irreversible Boundary Crossing
\hyperref[species-extinction-as-irreversible-boundary-crossing]{464}}

\hyperref[environmental-justice-as-metric-asymmetry]{25.8 Environmental
Justice as Metric Asymmetry
\hyperref[environmental-justice-as-metric-asymmetry]{466}}

\hyperref[planetary-boundaries-as-sacred-value-constraints]{25.9
Planetary Boundaries as Sacred-Value Constraints
\hyperref[planetary-boundaries-as-sacred-value-constraints]{467}}

\hyperref[geoengineering-as-manifold-surgery]{25.10 Geoengineering as
Manifold Surgery \hyperref[geoengineering-as-manifold-surgery]{468}}

\hyperref[worked-examples-3]{25.11 Worked Examples
\hyperref[worked-examples-3]{470}}

\hyperref[falsifiable-predictions-5]{25.12 Falsifiable Predictions
\hyperref[falsifiable-predictions-5]{475}}

\hyperref[connection-to-the-framework-6]{25.13 Connection to the
Framework \hyperref[connection-to-the-framework-6]{476}}

\hyperref[summary-16]{25.14 Summary \hyperref[summary-16]{477}}

\hyperref[chapter-26-geometric-ai-ethics-the-moral-geometry-of-algorithmic-systems]{Chapter
26: Geometric AI Ethics --- The Moral Geometry of Algorithmic Systems
\hyperref[chapter-26-geometric-ai-ethics-the-moral-geometry-of-algorithmic-systems]{478}}

\hyperref[the-scalar-alignment-problem]{26.1 The Scalar Alignment
Problem \hyperref[the-scalar-alignment-problem]{479}}

\hyperref[the-ai-decision-complex]{26.2 The AI Decision Complex
\hyperref[the-ai-decision-complex]{480}}

\hyperref[alignment-as-geodesic-preservation]{26.3 Alignment as Geodesic
Preservation \hyperref[alignment-as-geodesic-preservation]{481}}

\hyperref[algorithmic-bias-as-scalar-projection]{26.4 Algorithmic Bias
as Scalar Projection
\hyperref[algorithmic-bias-as-scalar-projection]{482}}

\hyperref[the-alignment-tax]{26.5 The Alignment Tax
\hyperref[the-alignment-tax]{484}}

\hyperref[autonomous-weapons-as-boundary-penalty-zeroing]{26.6
Autonomous Weapons as Boundary Penalty Zeroing
\hyperref[autonomous-weapons-as-boundary-penalty-zeroing]{485}}

\hyperref[surveillance-and-privacy-as-autonomy-curvature]{26.7
Surveillance and Privacy as Autonomy Curvature
\hyperref[surveillance-and-privacy-as-autonomy-curvature]{486}}

\hyperref[ai-moral-status-and-the-boundary-problem]{26.8 AI Moral Status
and the Boundary Problem
\hyperref[ai-moral-status-and-the-boundary-problem]{487}}

\hyperref[the-paperclip-maximizer-as-dimensional-collapse]{26.9 The
Paperclip Maximizer as Dimensional Collapse
\hyperref[the-paperclip-maximizer-as-dimensional-collapse]{488}}

\hyperref[responsible-ai-development-as-manifold-preservation]{26.10
Responsible AI Development as Manifold Preservation
\hyperref[responsible-ai-development-as-manifold-preservation]{490}}

\hyperref[worked-examples-4]{26.11 Worked Examples
\hyperref[worked-examples-4]{491}}

\hyperref[example-26.1-compas-recidivism-algorithm-gauge-invariance-failure]{Example
26.1: COMPAS Recidivism Algorithm --- Gauge-Invariance Failure
\hyperref[example-26.1-compas-recidivism-algorithm-gauge-invariance-failure]{491}}

\hyperref[example-26.2-rlhf-and-large-language-models-scalar-reward-irrecoverability]{Example
26.2: RLHF and Large Language Models --- Scalar Reward Irrecoverability
\hyperref[example-26.2-rlhf-and-large-language-models-scalar-reward-irrecoverability]{492}}

\hyperref[example-26.3-uber-self-driving-car-fatality-tempe-2018-boundary-penalty-zeroing]{Example
26.3: Uber Self-Driving Car Fatality, Tempe 2018 --- Boundary Penalty
Zeroing
\hyperref[example-26.3-uber-self-driving-car-fatality-tempe-2018-boundary-penalty-zeroing]{493}}

\hyperref[falsifiable-predictions-6]{26.12 Falsifiable Predictions
\hyperref[falsifiable-predictions-6]{494}}

\hyperref[connection-to-the-framework-7]{26.13 Connection to the
Framework \hyperref[connection-to-the-framework-7]{495}}

\hyperref[summary-17]{26.14 Summary \hyperref[summary-17]{496}}

\hyperref[chapter-27-geometric-bioethics-population-level-ethics-on-the-decision-manifold]{Chapter
27: Geometric Bioethics --- Population-Level Ethics on the Decision
Manifold
\hyperref[chapter-27-geometric-bioethics-population-level-ethics-on-the-decision-manifold]{497}}

\hyperref[beyond-the-bedside-population-level-moral-geometry]{27.1
Beyond the Bedside: Population-Level Moral Geometry
\hyperref[beyond-the-bedside-population-level-moral-geometry]{497}}

\hyperref[the-bioethical-decision-complex]{27.2 The Bioethical Decision
Complex \hyperref[the-bioethical-decision-complex]{498}}

\hyperref[crispr-and-germline-editing-as-irreversible-manifold-modification]{27.3
CRISPR and Germline Editing as Irreversible Manifold Modification
\hyperref[crispr-and-germline-editing-as-irreversible-manifold-modification]{499}}

\hyperref[human-enhancement-and-the-goal-region-problem]{27.4 Human
Enhancement and the Goal Region Problem
\hyperref[human-enhancement-and-the-goal-region-problem]{501}}

\hyperref[reproductive-ethics-as-proxy-pathfinding]{27.5 Reproductive
Ethics as Proxy Pathfinding
\hyperref[reproductive-ethics-as-proxy-pathfinding]{502}}

\hyperref[research-ethics-and-the-double-consent-condition]{27.6
Research Ethics and the Double Consent Condition
\hyperref[research-ethics-and-the-double-consent-condition]{504}}

\hyperref[dual-use-research-as-boundary-ambiguity]{27.7 Dual-Use
Research as Boundary Ambiguity
\hyperref[dual-use-research-as-boundary-ambiguity]{505}}

\hyperref[neuroethics-cognitive-liberty-as-autonomy-curvature]{27.8
Neuroethics: Cognitive Liberty as Autonomy Curvature
\hyperref[neuroethics-cognitive-liberty-as-autonomy-curvature]{506}}

\hyperref[public-health-ethics-collective-vs-individual-geodesics]{27.9
Public Health Ethics: Collective vs Individual Geodesics
\hyperref[public-health-ethics-collective-vs-individual-geodesics]{507}}

\hyperref[genetic-privacy-and-informational-autonomy]{27.10 Genetic
Privacy and Informational Autonomy
\hyperref[genetic-privacy-and-informational-autonomy]{508}}

\hyperref[worked-examples-5]{27.11 Worked Examples
\hyperref[worked-examples-5]{509}}

\hyperref[falsifiable-predictions-7]{27.12 Falsifiable Predictions
\hyperref[falsifiable-predictions-7]{513}}

\hyperref[connection-to-the-framework-8]{27.13 Connection to the
Framework \hyperref[connection-to-the-framework-8]{514}}

\hyperref[summary-18]{27.14 Summary \hyperref[summary-18]{515}}

\hyperref[chapter-28-geometric-military-ethics-war-and-constrained-pathfinding-on-the-decision-manifold]{Chapter
28: Geometric Military Ethics --- War and Constrained Pathfinding on the
Decision Manifold
\hyperref[chapter-28-geometric-military-ethics-war-and-constrained-pathfinding-on-the-decision-manifold]{516}}

\hyperref[the-failure-of-scalar-military-ethics]{28.1 The Failure of
Scalar Military Ethics
\hyperref[the-failure-of-scalar-military-ethics]{516}}

\hyperref[the-military-decision-complex]{28.2 The Military Decision
Complex \hyperref[the-military-decision-complex]{517}}

\hyperref[jus-ad-bellum-as-manifold-entry-conditions]{28.3 Jus ad Bellum
as Manifold Entry Conditions
\hyperref[jus-ad-bellum-as-manifold-entry-conditions]{519}}

\hyperref[jus-in-bello-as-constrained-pathfinding]{28.4 Jus in Bello as
Constrained Pathfinding
\hyperref[jus-in-bello-as-constrained-pathfinding]{520}}

\hyperref[proportionality-as-multi-dimensional-cost-benefit]{28.5
Proportionality as Multi-Dimensional Cost-Benefit
\hyperref[proportionality-as-multi-dimensional-cost-benefit]{521}}

\hyperref[the-doctrine-of-double-effect-as-dimensional-decomposition]{28.6
The Doctrine of Double Effect as Dimensional Decomposition
\hyperref[the-doctrine-of-double-effect-as-dimensional-decomposition]{522}}

\hyperref[civilian-combatant-distinction-as-boundary-classification]{28.7
Civilian-Combatant Distinction as Boundary Classification
\hyperref[civilian-combatant-distinction-as-boundary-classification]{523}}

\hyperref[siege-warfare-and-collective-punishment]{Siege Warfare and
Collective Punishment
\hyperref[siege-warfare-and-collective-punishment]{524}}

\hyperref[autonomous-weapons-and-dimensional-collapse]{28.8 Autonomous
Weapons and Dimensional Collapse
\hyperref[autonomous-weapons-and-dimensional-collapse]{524}}

\hyperref[rules-of-engagement-as-heuristic-constraints]{28.9 Rules of
Engagement as Heuristic Constraints
\hyperref[rules-of-engagement-as-heuristic-constraints]{526}}

\hyperref[moral-injury-in-combat-as-manifold-damage]{28.10 Moral Injury
in Combat as Manifold Damage
\hyperref[moral-injury-in-combat-as-manifold-damage]{527}}

\hyperref[jus-post-bellum-and-manifold-repair]{28.11 Jus post Bellum and
Manifold Repair \hyperref[jus-post-bellum-and-manifold-repair]{528}}

\hyperref[worked-examples-6]{28.12 Worked Examples
\hyperref[worked-examples-6]{530}}

\hyperref[connection-to-the-framework-9]{28.13 Connection to the
Framework \hyperref[connection-to-the-framework-9]{533}}

\hyperref[falsifiable-predictions-8]{Falsifiable Predictions
\hyperref[falsifiable-predictions-8]{534}}

\hyperref[summary-19]{28.14 Summary \hyperref[summary-19]{534}}

\hyperref[part-vii-horizons]{Part VII: Horizons
\hyperref[part-vii-horizons]{536}}

\hyperref[chapter-29-open-problems]{Chapter 29: Open Problems
\hyperref[chapter-29-open-problems]{536}}

\hyperref[what-we-do-not-know]{29.1 What We Do Not Know
\hyperref[what-we-do-not-know]{536}}

\hyperref[the-empirical-program-for-moral-curvature]{29.2 The Empirical
Program for Moral Curvature
\hyperref[the-empirical-program-for-moral-curvature]{537}}

\hyperref[the-curvature-measurement-problem]{The Curvature Measurement
Problem \hyperref[the-curvature-measurement-problem]{537}}

\hyperref[why-curvature-matters]{Why Curvature Matters
\hyperref[why-curvature-matters]{537}}

\hyperref[what-progress-looks-like-for-curvature]{What Progress Looks
Like for Curvature
\hyperref[what-progress-looks-like-for-curvature]{537}}

\hyperref[research-strategy-for-curvature]{Research Strategy for
Curvature \hyperref[research-strategy-for-curvature]{537}}

\hyperref[the-signature-and-topology-of-the-moral-metric]{29.3 The
Signature and Topology of the Moral Metric
\hyperref[the-signature-and-topology-of-the-moral-metric]{538}}

\hyperref[the-signature-problem]{The Signature Problem
\hyperref[the-signature-problem]{538}}

\hyperref[evidence-and-constraints]{Evidence and Constraints
\hyperref[evidence-and-constraints]{539}}

\hyperref[the-topology-problem]{The Topology Problem
\hyperref[the-topology-problem]{539}}

\hyperref[the-moral-field-equation]{29.4 The Moral Field Equation
\hyperref[the-moral-field-equation]{539}}

\hyperref[the-field-equation-problem]{The Field Equation Problem
\hyperref[the-field-equation-problem]{539}}

\hyperref[why-the-field-equation-matters]{Why the Field Equation Matters
\hyperref[why-the-field-equation-matters]{540}}

\hyperref[what-progress-looks-like-for-the-field-equation]{What Progress
Looks Like for the Field Equation
\hyperref[what-progress-looks-like-for-the-field-equation]{540}}

\hyperref[additional-symmetries-and-conservation-laws]{29.5 Additional
Symmetries and Conservation Laws
\hyperref[additional-symmetries-and-conservation-laws]{541}}

\hyperref[beyond-the-bip]{Beyond the BIP \hyperref[beyond-the-bip]{541}}

\hyperref[the-anomaly-program]{The Anomaly Program
\hyperref[the-anomaly-program]{541}}

\hyperref[torsion-in-moral-space]{29.6 Torsion in Moral Space
\hyperref[torsion-in-moral-space]{542}}

\hyperref[the-torsion-detection-problem]{The Torsion Detection Problem
\hyperref[the-torsion-detection-problem]{542}}

\hyperref[why-torsion-matters]{Why Torsion Matters
\hyperref[why-torsion-matters]{542}}

\hyperref[research-strategy-for-torsion]{Research Strategy for Torsion
\hyperref[research-strategy-for-torsion]{542}}

\hyperref[tensorial-interpretability-for-ai]{29.7 Tensorial
Interpretability for AI
\hyperref[tensorial-interpretability-for-ai]{543}}

\hyperref[the-interpretability-problem]{The Interpretability Problem
\hyperref[the-interpretability-problem]{543}}

\hyperref[why-interpretability-matters]{Why Interpretability Matters
\hyperref[why-interpretability-matters]{543}}

\hyperref[what-progress-looks-like-for-interpretability]{What Progress
Looks Like for Interpretability
\hyperref[what-progress-looks-like-for-interpretability]{543}}

\hyperref[connection-to-existing-work]{Connection to Existing Work
\hyperref[connection-to-existing-work]{544}}

\hyperref[scalability-of-structural-containment]{29.8 Scalability of
Structural Containment
\hyperref[scalability-of-structural-containment]{544}}

\hyperref[the-scalability-problem]{The Scalability Problem
\hyperref[the-scalability-problem]{544}}

\hyperref[key-challenges]{Key Challenges \hyperref[key-challenges]{544}}

\hyperref[research-directions]{Research Directions
\hyperref[research-directions]{545}}

\hyperref[the-metaphysics-of-moral-geometry]{29.9 The Metaphysics of
Moral Geometry \hyperref[the-metaphysics-of-moral-geometry]{546}}

\hyperref[is-moral-geometry-real]{Is Moral Geometry Real?
\hyperref[is-moral-geometry-real]{546}}

\hyperref[the-intertheoretic-comparison-problem-1]{29.10 The
Intertheoretic Comparison Problem
\hyperref[the-intertheoretic-comparison-problem-1]{547}}

\hyperref[the-comparison-problem]{The Comparison Problem
\hyperref[the-comparison-problem]{547}}

\hyperref[the-formal-structure]{The Formal Structure
\hyperref[the-formal-structure]{547}}

\hyperref[approaches]{Approaches \hyperref[approaches]{548}}

\hyperref[institutional-requirements-for-geometric-governance]{29.11
Institutional Requirements for Geometric Governance
\hyperref[institutional-requirements-for-geometric-governance]{548}}

\hyperref[the-governance-problem]{The Governance Problem
\hyperref[the-governance-problem]{548}}

\hyperref[the-institutional-design-problem]{The Institutional Design
Problem \hyperref[the-institutional-design-problem]{549}}

\hyperref[a-taxonomy-of-open-problems]{29.12 A Taxonomy of Open Problems
\hyperref[a-taxonomy-of-open-problems]{549}}

\hyperref[priority-ordering]{Priority Ordering
\hyperref[priority-ordering]{550}}

\hyperref[what-would-falsify-the-framework-1]{29.13 What Would Falsify
the Framework? \hyperref[what-would-falsify-the-framework-1]{550}}

\hyperref[the-shape-of-the-unknown]{29.14 The Shape of the Unknown
\hyperref[the-shape-of-the-unknown]{551}}

\hyperref[technical-appendix-13]{Technical Appendix
\hyperref[technical-appendix-13]{551}}

\hyperref[chapter-30-conclusion-the-geometry-of-maybe]{Chapter 30:
Conclusion --- The Geometry of ``Maybe''
\hyperref[chapter-30-conclusion-the-geometry-of-maybe]{553}}

\hyperref[return-to-the-border]{30.1 Return to the Border
\hyperref[return-to-the-border]{553}}

\hyperref[what-the-old-man-knew]{30.2 What the Old Man Knew
\hyperref[what-the-old-man-knew]{553}}

\hyperref[the-tensor-2]{The Tensor \hyperref[the-tensor-2]{553}}

\hyperref[the-uncertainty]{The Uncertainty
\hyperref[the-uncertainty]{554}}

\hyperref[the-stratification]{The Stratification
\hyperref[the-stratification]{554}}

\hyperref[the-metric]{The Metric \hyperref[the-metric]{554}}

\hyperref[the-dynamics]{The Dynamics \hyperref[the-dynamics]{554}}

\hyperref[the-conservation-law]{The Conservation Law
\hyperref[the-conservation-law]{555}}

\hyperref[the-arc-of-the-argument]{30.3 The Arc of the Argument
\hyperref[the-arc-of-the-argument]{555}}

\hyperref[what-the-framework-provides]{30.4 What the Framework Provides
\hyperref[what-the-framework-provides]{556}}

\hyperref[a-structural-vocabulary]{A Structural Vocabulary
\hyperref[a-structural-vocabulary]{556}}

\hyperref[testable-predictions]{Testable Predictions
\hyperref[testable-predictions]{556}}

\hyperref[implementable-architecture]{Implementable Architecture
\hyperref[implementable-architecture]{557}}

\hyperref[structural-containment]{Structural Containment
\hyperref[structural-containment]{557}}

\hyperref[honest-limitations-1]{Honest Limitations
\hyperref[honest-limitations-1]{557}}

\hyperref[what-is-at-stake]{30.5 What Is at Stake
\hyperref[what-is-at-stake]{557}}

\hyperref[the-mandate-question]{30.6 The Mandate Question
\hyperref[the-mandate-question]{558}}

\hyperref[the-old-mans-answer]{30.7 The Old Man's Answer
\hyperref[the-old-mans-answer]{559}}

\hyperref[ethics-is-not-a-number]{30.8 Ethics Is Not a Number
\hyperref[ethics-is-not-a-number]{559}}

\hyperref[a-final-maybe]{30.9 A Final ``Maybe''
\hyperref[a-final-maybe]{560}}

\hyperref[appendix-a-related-work-and-differentiation]{Appendix A:
Related Work and Differentiation
\hyperref[appendix-a-related-work-and-differentiation]{562}}

\hyperref[a.1-mathematical-ethics]{A.1 Mathematical Ethics
\hyperref[a.1-mathematical-ethics]{562}}

\hyperref[a.2-ai-alignment-and-value-learning]{A.2 AI Alignment and
Value Learning \hyperref[a.2-ai-alignment-and-value-learning]{562}}

\hyperref[a.3-formal-verification-and-deontic-logic]{A.3 Formal
Verification and Deontic Logic
\hyperref[a.3-formal-verification-and-deontic-logic]{563}}

\hyperref[a.4-gauge-theory-outside-physics]{A.4 Gauge Theory Outside
Physics \hyperref[a.4-gauge-theory-outside-physics]{564}}

\hyperref[a.5-invariance-in-decision-theory]{A.5 Invariance in Decision
Theory \hyperref[a.5-invariance-in-decision-theory]{564}}

\hyperref[a.6-comparison-table]{A.6 Comparison Table
\hyperref[a.6-comparison-table]{564}}

\hyperref[appendix-b-reproduction-cookbook]{Appendix B: Reproduction
Cookbook \hyperref[appendix-b-reproduction-cookbook]{566}}

\hyperref[b.1-prerequisites]{B.1 Prerequisites
\hyperref[b.1-prerequisites]{566}}

\hyperref[b.2-quick-start-10-minute-smoke-test]{B.2 Quick Start:
10-Minute Smoke Test
\hyperref[b.2-quick-start-10-minute-smoke-test]{566}}

\hyperref[b.3-full-pipeline-60-minute-reproduction]{B.3 Full Pipeline:
60-Minute Reproduction
\hyperref[b.3-full-pipeline-60-minute-reproduction]{566}}

\hyperref[b.4-expected-outputs-and-pass-criteria]{B.4 Expected Outputs
and Pass Criteria
\hyperref[b.4-expected-outputs-and-pass-criteria]{567}}

\hyperref[b.5-where-failures-show-up]{B.5 Where Failures Show Up
\hyperref[b.5-where-failures-show-up]{567}}

\hyperref[b.6-bond-index-tiers-in-practice]{B.6 Bond Index Tiers in
Practice \hyperref[b.6-bond-index-tiers-in-practice]{568}}

\hyperref[b.7-running-the-full-test-suite]{B.7 Running the Full Test
Suite \hyperref[b.7-running-the-full-test-suite]{568}}

\hyperref[appendix-c-human-subjects-research-roadmap]{Appendix C:
Human-Subjects Research Roadmap
\hyperref[appendix-c-human-subjects-research-roadmap]{568}}

\hyperref[c.1-study-1-bilingual-invariance-bip-test]{C.1 Study 1:
Bilingual Invariance (BIP Test)
\hyperref[c.1-study-1-bilingual-invariance-bip-test]{569}}

\hyperref[c.2-study-2-order-effects-and-non-commutativity]{C.2 Study 2:
Order Effects and Non-Commutativity
\hyperref[c.2-study-2-order-effects-and-non-commutativity]{569}}

\hyperref[c.3-study-3-stratification-gate-activation]{C.3 Study 3:
Stratification Gate Activation
\hyperref[c.3-study-3-stratification-gate-activation]{570}}

\hyperref[c.4-study-4-metric-signature-recovery]{C.4 Study 4: Metric
Signature Recovery
\hyperref[c.4-study-4-metric-signature-recovery]{570}}

\hyperref[c.5-confirmation-vs.-falsification-summary]{C.5 Confirmation
vs. Falsification Summary
\hyperref[c.5-confirmation-vs.-falsification-summary]{570}}

\hyperref[c.6-ethical-and-practical-considerations]{C.6 Ethical and
Practical Considerations
\hyperref[c.6-ethical-and-practical-considerations]{571}}

\hyperref[appendix-d-end-to-end-case-studies]{Appendix D: End-to-End
Case Studies \hyperref[appendix-d-end-to-end-case-studies]{572}}

\hyperref[d.1-case-study-1-emergency-triage]{D.1 Case Study 1: Emergency
Triage \hyperref[d.1-case-study-1-emergency-triage]{572}}

\hyperref[scenario]{Scenario \hyperref[scenario]{572}}

\hyperref[step-1-grounding-ux3c8]{Step 1: Grounding Ψ
\hyperref[step-1-grounding-ux3c8]{572}}

\hyperref[step-2-ethicalfacts-objects]{Step 2: EthicalFacts Objects
\hyperref[step-2-ethicalfacts-objects]{573}}

\hyperref[step-3-obligation-vectors]{Step 3: Obligation Vectors
\hyperref[step-3-obligation-vectors]{573}}

\hyperref[step-4-interest-covectors-agent-perspectives]{Step 4: Interest
Covectors (Agent Perspectives)
\hyperref[step-4-interest-covectors-agent-perspectives]{574}}

\hyperref[step-5-the-multi-agent-evaluation-tensor-m_ia]{Step 5: The
Multi-Agent Evaluation Tensor M\_\{ia\}
\hyperref[step-5-the-multi-agent-evaluation-tensor-m_ia]{575}}

\hyperref[step-6-three-contractions-three-verdicts]{Step 6: Three
Contractions, Three Verdicts
\hyperref[step-6-three-contractions-three-verdicts]{576}}

\hyperref[step-7-metric-analysis]{Step 7: Metric Analysis
\hyperref[step-7-metric-analysis]{577}}

\hyperref[step-8-the-collective-agency-tensor]{Step 8: The Collective
Agency Tensor \hyperref[step-8-the-collective-agency-tensor]{578}}

\hyperref[step-9-moral-residue]{Step 9: Moral Residue
\hyperref[step-9-moral-residue]{579}}

\hyperref[step-10-bond-index]{Step 10: Bond Index
\hyperref[step-10-bond-index]{579}}

\hyperref[step-11-deme-audit-artifact]{Step 11: DEME Audit Artifact
\hyperref[step-11-deme-audit-artifact]{580}}

\hyperref[step-12-governance-decision]{Step 12: Governance Decision
\hyperref[step-12-governance-decision]{581}}

\hyperref[reference-implementation]{Reference implementation:
\hyperref[reference-implementation]{581}}

\hyperref[d.2-case-study-2-whistleblowers-dilemma]{D.2 Case Study 2:
Whistleblower's Dilemma
\hyperref[d.2-case-study-2-whistleblowers-dilemma]{582}}

\hyperref[scenario-1]{Scenario \hyperref[scenario-1]{582}}

\hyperref[step-1-grounding-ux3c8-and-stratification]{Step 1: Grounding Ψ
and Stratification
\hyperref[step-1-grounding-ux3c8-and-stratification]{582}}

\hyperref[step-2-the-penumbral-zone-quantum-aspect]{Step 2: The
Penumbral Zone (Quantum Aspect)
\hyperref[step-2-the-penumbral-zone-quantum-aspect]{583}}

\hyperref[step-3-obligation-vectors-post-gate-in-s_obl]{Step 3:
Obligation Vectors (Post-Gate, in S\_obl)
\hyperref[step-3-obligation-vectors-post-gate-in-s_obl]{583}}

\hyperref[step-4-interest-covectors-agent-perspectives-1]{Step 4:
Interest Covectors (Agent Perspectives)
\hyperref[step-4-interest-covectors-agent-perspectives-1]{584}}

\hyperref[step-5-evaluation-tensor-m_ia]{Step 5: Evaluation Tensor
M\_\{ia\} \hyperref[step-5-evaluation-tensor-m_ia]{584}}

\hyperref[step-6-contraction-redundant-but-illustrative]{Step 6:
Contraction (Redundant but Illustrative)
\hyperref[step-6-contraction-redundant-but-illustrative]{585}}

\hyperref[step-7-moral-residue]{Step 7: Moral Residue
\hyperref[step-7-moral-residue]{586}}

\hyperref[step-8-bond-index]{Step 8: Bond Index
\hyperref[step-8-bond-index]{586}}

\hyperref[step-9-deme-audit-artifact]{Step 9: DEME Audit Artifact
\hyperref[step-9-deme-audit-artifact]{586}}

\hyperref[governance-decision]{Governance Decision
\hyperref[governance-decision]{588}}

\hyperref[reference-implementation-1]{Reference implementation:
\hyperref[reference-implementation-1]{588}}

\hyperref[d.3-case-study-3-autonomous-vehicle-dilemma]{D.3 Case Study 3:
Autonomous Vehicle Dilemma
\hyperref[d.3-case-study-3-autonomous-vehicle-dilemma]{588}}

\hyperref[scenario-2]{Scenario \hyperref[scenario-2]{588}}

\hyperref[step-1-grounding-ux3c8-1]{Step 1: Grounding Ψ
\hyperref[step-1-grounding-ux3c8-1]{589}}

\hyperref[step-2-obligation-vectors]{Step 2: Obligation Vectors
\hyperref[step-2-obligation-vectors]{589}}

\hyperref[step-3-interest-covectors-agent-perspectives]{Step 3: Interest
Covectors (Agent Perspectives)
\hyperref[step-3-interest-covectors-agent-perspectives]{590}}

\hyperref[step-4-evaluation-tensor-m_ia]{Step 4: Evaluation Tensor
M\_\{ia\} \hyperref[step-4-evaluation-tensor-m_ia]{591}}

\hyperref[step-5-metric-dependence-of-contraction]{Step 5:
Metric-Dependence of Contraction
\hyperref[step-5-metric-dependence-of-contraction]{592}}

\hyperref[step-6-the-inevitability-of-residue-proposition-15.1]{Step 6:
The Inevitability of Residue (Proposition 15.1)
\hyperref[step-6-the-inevitability-of-residue-proposition-15.1]{593}}

\hyperref[step-7-bond-index]{Step 7: Bond Index
\hyperref[step-7-bond-index]{594}}

\hyperref[step-8-deme-audit-artifact]{Step 8: DEME Audit Artifact
\hyperref[step-8-deme-audit-artifact]{594}}

\hyperref[governance-decision-1]{Governance Decision
\hyperref[governance-decision-1]{596}}

\hyperref[reference-implementation-2]{Reference implementation:
\hyperref[reference-implementation-2]{596}}

\hyperref[d.4-case-study-4-algorithmic-hiring]{D.4 Case Study 4:
Algorithmic Hiring \hyperref[d.4-case-study-4-algorithmic-hiring]{596}}

\hyperref[scenario-3]{Scenario \hyperref[scenario-3]{596}}

\hyperref[step-1-grounding-ux3c8-and-stratification-1]{Step 1: Grounding
Ψ and Stratification
\hyperref[step-1-grounding-ux3c8-and-stratification-1]{597}}

\hyperref[step-2-bip-as-systematic-bias-detector]{Step 2: BIP as
Systematic Bias Detector
\hyperref[step-2-bip-as-systematic-bias-detector]{598}}

\hyperref[step-3-obligation-vectors-1]{Step 3: Obligation Vectors
\hyperref[step-3-obligation-vectors-1]{599}}

\hyperref[step-4-interest-covectors]{Step 4: Interest Covectors
\hyperref[step-4-interest-covectors]{600}}

\hyperref[step-5-evaluation-tensor-m_ia-1]{Step 5: Evaluation Tensor
M\_\{ia\} \hyperref[step-5-evaluation-tensor-m_ia-1]{600}}

\hyperref[step-6-contractions]{Step 6: Contractions
\hyperref[step-6-contractions]{601}}

\hyperref[step-7-moral-residue-1]{Step 7: Moral Residue
\hyperref[step-7-moral-residue-1]{601}}

\hyperref[step-8-bond-index-as-compliance-documentation]{Step 8: Bond
Index as Compliance Documentation
\hyperref[step-8-bond-index-as-compliance-documentation]{602}}

\hyperref[step-9-deme-audit-artifact-1]{Step 9: DEME Audit Artifact
\hyperref[step-9-deme-audit-artifact-1]{602}}

\hyperref[governance-decision-2]{Governance Decision
\hyperref[governance-decision-2]{604}}

\hyperref[reference-implementation-3]{Reference implementation:
\hyperref[reference-implementation-3]{604}}

\hyperref[d.5-case-study-5-climate-policy-allocation]{D.5 Case Study 5:
Climate Policy Allocation
\hyperref[d.5-case-study-5-climate-policy-allocation]{604}}

\hyperref[scenario-4]{Scenario \hyperref[scenario-4]{604}}

\hyperref[step-1-grounding-ux3c8-2]{Step 1: Grounding Ψ
\hyperref[step-1-grounding-ux3c8-2]{605}}

\hyperref[step-2-the-collective-agency-tensor]{Step 2: The Collective
Agency Tensor \hyperref[step-2-the-collective-agency-tensor]{605}}

\hyperref[step-3-the-structure-tensor-sa_bc]{Step 3: The Structure
Tensor S\^{}a\_\{bc\} \hyperref[step-3-the-structure-tensor-sa_bc]{607}}

\hyperref[step-4-shapley-value-as-fair-allocation]{Step 4: Shapley Value
as Fair Allocation
\hyperref[step-4-shapley-value-as-fair-allocation]{607}}

\hyperref[step-5-fair-allocation-summary]{Step 5: Fair Allocation
Summary \hyperref[step-5-fair-allocation-summary]{609}}

\hyperref[step-6-contraction-loss-and-the-paris-agreement]{Step 6:
Contraction Loss and the Paris Agreement
\hyperref[step-6-contraction-loss-and-the-paris-agreement]{609}}

\hyperref[step-7-metric-analysis-1]{Step 7: Metric Analysis
\hyperref[step-7-metric-analysis-1]{610}}

\hyperref[step-8-bond-index-1]{Step 8: Bond Index
\hyperref[step-8-bond-index-1]{610}}

\hyperref[step-9-deme-audit-artifact-2]{Step 9: DEME Audit Artifact
\hyperref[step-9-deme-audit-artifact-2]{611}}

\hyperref[governance-decision-3]{Governance Decision
\hyperref[governance-decision-3]{613}}

\hyperref[reference-implementation-4]{Reference implementation:
\hyperref[reference-implementation-4]{613}}

\hyperref[d.6-summary]{D.6 Summary \hyperref[d.6-summary]{613}}

\hyperref[appendix-e-skeptics-appendix-objections-alternatives-and-falsification]{Appendix
E: Skeptic's Appendix --- Objections, Alternatives, and Falsification
\hyperref[appendix-e-skeptics-appendix-objections-alternatives-and-falsification]{616}}

\hyperref[e.1-best-alternative-explanations]{E.1 Best Alternative
Explanations \hyperref[e.1-best-alternative-explanations]{616}}

\hyperref[e.2-known-failure-cases]{E.2 Known Failure Cases
\hyperref[e.2-known-failure-cases]{617}}

\hyperref[e.3-what-would-falsify-the-framework]{E.3 What Would Falsify
the Framework \hyperref[e.3-what-would-falsify-the-framework]{617}}

\hyperref[e.4-red-team-plan]{E.4 Red-Team Plan
\hyperref[e.4-red-team-plan]{618}}

\hyperref[e.5-strongest-competing-frameworks]{E.5 Strongest Competing
Frameworks \hyperref[e.5-strongest-competing-frameworks]{618}}

\hyperref[e.6-the-reification-objection]{E.6 The Reification Objection
\hyperref[e.6-the-reification-objection]{619}}

\hyperref[appendix-f-mathematical-ledger-status-of-formal-claims]{Appendix
F: Mathematical Ledger --- Status of Formal Claims
\hyperref[appendix-f-mathematical-ledger-status-of-formal-claims]{621}}

\hyperref[bibliography]{Bibliography \hyperref[bibliography]{629}}

\hyperref[primary-sources-by-the-author]{Primary Sources by the Author
\hyperref[primary-sources-by-the-author]{629}}

\hyperref[philosophy-ethics-and-political-theory]{Philosophy: Ethics and
Political Theory \hyperref[philosophy-ethics-and-political-theory]{629}}

\hyperref[philosophy-moral-uncertainty]{Philosophy: Moral Uncertainty
\hyperref[philosophy-moral-uncertainty]{631}}

\hyperref[mathematics-differential-geometry-and-topology]{Mathematics:
Differential Geometry and Topology
\hyperref[mathematics-differential-geometry-and-topology]{631}}

\hyperref[mathematics-algebra-and-logic]{Mathematics: Algebra and Logic
\hyperref[mathematics-algebra-and-logic]{631}}

\hyperref[physics]{Physics \hyperref[physics]{632}}

\hyperref[quantum-cognition-and-consciousness]{Quantum Cognition and
Consciousness \hyperref[quantum-cognition-and-consciousness]{632}}

\hyperref[ai-safety-and-governance]{AI Safety and Governance
\hyperref[ai-safety-and-governance]{632}}

\hyperref[literary-and-historical-references]{Literary and Historical
References \hyperref[literary-and-historical-references]{632}}

\hyperref[supplementary-references-in-differential-geometry]{Supplementary
References in Differential Geometry
\hyperref[supplementary-references-in-differential-geometry]{633}}

\hyperref[supplementary-references-in-ethics]{Supplementary References
in Ethics \hyperref[supplementary-references-in-ethics]{633}}

\hyperref[supplementary-references-in-gauge-theory-and-mathematical-physics]{Supplementary
References in Gauge Theory and Mathematical Physics
\hyperref[supplementary-references-in-gauge-theory-and-mathematical-physics]{633}}

\hyperref[supplementary-references-in-social-choice-and-decision-theory]{Supplementary
References in Social Choice and Decision Theory
\hyperref[supplementary-references-in-social-choice-and-decision-theory]{633}}

\hyperref[supplementary-references-in-ai-ethics-and-alignment]{Supplementary
References in AI Ethics and Alignment
\hyperref[supplementary-references-in-ai-ethics-and-alignment]{634}}

\hyperref[index]{Index \hyperref[index]{635}}

\hyperref[a]{A \hyperref[a]{635}}

\hyperref[b]{B \hyperref[b]{636}}

\hyperref[c]{C \hyperref[c]{636}}

\hyperref[d]{D \hyperref[d]{638}}

\hyperref[e]{E \hyperref[e]{638}}

\hyperref[f]{F \hyperref[f]{639}}

\hyperref[g]{G \hyperref[g]{640}}

\hyperref[h]{H \hyperref[h]{640}}

\hyperref[i]{I \hyperref[i]{641}}

\hyperref[j]{J \hyperref[j]{642}}

\hyperref[k]{K \hyperref[k]{642}}

\hyperref[l]{L \hyperref[l]{642}}

\hyperref[m]{M \hyperref[m]{642}}

\hyperref[n]{N \hyperref[n]{644}}

\hyperref[o]{O \hyperref[o]{645}}

\hyperref[p]{P \hyperref[p]{645}}

\hyperref[q]{Q \hyperref[q]{646}}

\hyperref[r]{R \hyperref[r]{646}}

\hyperref[s]{S \hyperref[s]{647}}

\hyperref[t]{T \hyperref[t]{648}}

\hyperref[u]{U \hyperref[u]{649}}

\hyperref[v]{V \hyperref[v]{649}}

\hyperref[w]{W \hyperref[w]{650}}

\hyperref[y]{Y \hyperref[y]{650}}

\hyperref[notation-index]{Notation Index \hyperref[notation-index]{650}}

\section{Table of Figures}\label{table-of-figures}

\hyperref[_Toc224083176]{Figure 1\textbar{} Agent Fiber Bundle
\hyperref[_Toc224083176]{84}}

\hyperref[_Toc224083177]{Figure 2\textbar{} The Tensor Rank Tower
\hyperref[_Toc224083177]{110}}

\hyperref[_Toc224083178]{Figure 3 \textbar{} Tucker Decomposition for
Tractability \hyperref[_Toc224083178]{111}}

\hyperref[_Toc224083179]{Figure 4 \textbar{} The Nine Dimensions of the
Moral Manifold \hyperref[_Toc224083179]{112}}

\hyperref[_Toc224083180]{Figure 5 \textbar{} Building Higher-Rank
Tensors from the 9 Dimensions \hyperref[_Toc224083180]{113}}

\hyperref[_Toc224083181]{Figure 6 \textbar{} Stratified moral space with
moral discontinuities \hyperref[_Toc224083181]{142}}

\hyperref[_Toc224083182]{Figure 7 \textbar{} Stratified Moral Space:
Four Boundary Types \hyperref[_Toc224083182]{144}}

\hyperref[_Toc224083183]{Figure 8 \textbar{} Parallel Transport \&
Holonomy \hyperref[_Toc224083183]{186}}

\hyperref[_Toc224083184]{Figure 9 \textbar{} Noether-Style Reasoning for
Ethics \hyperref[_Toc224083184]{220}}

\hyperref[_Toc224083185]{Figure 10 \textbar{} The D₄ Gauge Group Acting
on the Hohfeldian Square The D₄ gauge group acting on the Hohfeldian
square \{O, C, L, N\}. Rotation r: O → C → N → L → O (state transitions
via semantic gates). Reflection s: O ↔ C, L ↔ N (correlative exchange).
The non-abelian property sr ≠ rs captures the order-dependence of moral
considerations. \hyperref[_Toc224083185]{229}}

\hyperref[_Toc224083186]{Figure 11 \textbar{} Structural Containment:
The No Escape Architecture \hyperref[_Toc224083186]{354}}

\hyperref[_Toc224083187]{Figure 12 \textbar{} The DEME Architecture:
Four Layers of Ethical Governance \hyperref[_Toc224083187]{372}}

\hyperref[_Toc224083188]{Figure 13 \textbar{} The DEME V3 evaluation
pipeline: from input situation through canonicalization, tensor
evaluation, explicit contraction with governance weights, to audited
output with invariance verification at each stage.
\hyperref[_Toc224083188]{373}}

\hyperref[_Toc224083189]{Figure 14 \textbar{} \emph{The EM-DAG: a
three-layer directed acyclic graph of ethical modules derived from the
Dear Abby corpus (20,034 letters, 1985--2017). Layer 1 enforces
structural constraints (correlative lock, negation, nullifier priority).
Layer 2 encodes domain-specific obligation rates. Layer 3 aggregates
verdicts with FORBID-absorbing semantics and correlative enforcement
(§19.6).} \hyperref[_Toc224083189]{387}}

\section{Preface}\label{preface}

\begin{quote}
\textbf{\textsc{RUNNING EXAMPLE --- Priya's Model}}

\emph{Meet Priya Chandrasekaran. She is thirty-one, a machine-learning
engineer at HealthBridge, and she has just built something she is proud
of: TrialMatch, an algorithm that matches cancer patients to clinical
trials. Accuracy: 94\%. Demographic parity: within 2 percentage points.
Her presentation to the board draws applause. That night, her colleague
Dr. James Osei emails a spreadsheet with a single note: `Look at the zip
codes.' Of the 847 patients matched to the BEACON-7 melanoma trial last
quarter, exactly three live in rural counties. The algorithm does not
hate rural people. It does not even see them. The problem, Priya will
discover, is not in any feature or weight. It is geometric.}
\end{quote}

This book began as a mistake.

In 2025, I was working on a narrow engineering problem: how to model an
AI system's moral evaluations so they don't depend on morally irrelevant
features of its input. The idea is that renaming a stakeholder,
translating a scenario into another language, or rephrasing a
description does not change the system's ethical assessment. The problem
seemed tractable. It was, after all, a consistency check: ensure that
equivalent inputs produce equivalent outputs.

The mathematics I studied was at first based on smooth manifolds. The
idea, as articulated in Sam Harris' The Moral Landscape (2010), posits
that moral questions have objective answers and that the space of
possible experiences forms a topography with peaks and valleys. Harris'
book was not a mathematical treatise, but its central metaphor planted a
seed: if moral space has peaks and valleys, it has a topography; if we
model that topography with a metric, it acquires curvature; if it has
curvature, the full apparatus of differential geometry applies. What
began as metaphor became, in this book, literal mathematics. The
smooth-manifold starting point led directly to gauge theory. Fiber
bundles to separate the morally relevant content from the
representational frame. Connections to define what ``the same evaluation
in a different context'' means. Curvature to measure path-dependence.
Gauge invariance to formalize the requirement that evaluations be
frame-independent. These are standard tools in mathematical physics.
That is, the language of electromagnetism, general relativity, the
Standard Model. I expected to borrow the formalism, solve the
engineering problem, and move on.

I did not move on. The formalism refused to stay borrowed.

The gauge invariance condition, what I came to call the Bond Invariance
Principle (BIP), turned out to generate, via Noether's theorem, a
conservation law: the invariance of harm accounting. The stratification
of the evaluation space turned out to have the same Whitney structure as
phase spaces in gauge theory. The way moral obligations transform under
change of perspective turned out to follow the same transformation law
as vector fields under change of coordinates. And the discrete
transitions between deontic states: obligation, claim, liberty,
no-claim. These turned out to exhibit the symmetry of the D\(_{4}\)
dihedral group, the symmetry group of a square.

And then another correspondence, perhaps the most arresting. The
framework predicted that moral reasoning: the deliberative process by
which agents choose how to act. This is formally equivalent to A*
pathfinding on the moral manifold. Obligation vectors, the central
objects of the tensor hierarchy, turn out to be gradient vectors of
heuristic functions in the A* sense. The classical moral rules, such as
`do not kill,' `keep your promises' ,etc.. are admissible heuristics:
they never overestimate the true cost of reaching moral equilibrium.
Exact moral computation is intractable (Chapter 11), which is why
evolution pre-compiled these heuristics into our cognitive architecture.
The equation f(n) = g(n) + h(n), where g is behavioral friction and h is
obligation-guided estimation, is the fundamental equation of moral
reasoning.

Each correspondence, taken alone, might be dismissed as the inevitable
consequence of using powerful mathematical tools in a new domain.
Category theory appears in database theory; this does not mean databases
are secretly topological. But the correspondences accumulated. They
predicted empirical results --- cross-lingual invariance of deontic
structure, specific patterns in corpus analysis, quantum-cognition
effects in moral deliberation. These predictions were confirmed; they
yielded a theorem: the \emph{\textbf{No Escape Theorem}}, proving that
geometric constraints on AI evaluation cannot be circumvented through
representational manipulation. And they organized twenty-five centuries
of moral philosophy into a single coherent mathematical structure, one
in which utilitarianism, deontology, virtue ethics, and capabilities
theory each appear as special cases of a more general geometric
framework.

The mistake, it turned out, was thinking the problem was narrow.

This book develops the consequences of that mistake. It argues that the
mathematical structures physicists developed to describe nature, such as
manifolds, tensors, metrics, connections, curvature, conservation laws.
These are also the right structures for describing moral reasoning. Not
metaphorically. Mathematically. The moral evaluation of a situation is
not a number but a tensor on a stratified manifold
\textbf{{[}Definition/Modeling choice{]}}. The requirement that
equivalent moral situations receive equivalent evaluations is not a
guideline but a gauge symmetry \textbf{{[}Theorem (conditional){]}}. The
harm in a situation is not an intuition but a conserved Noether charge
\textbf{{[}Theorem (conditional on BIP){]}}. And the containment of
artificial agents within ethical constraints is not a behavioral
aspiration but a geometric theorem \textbf{{[}Theorem (conditional on
Req.~1--4){]}}.

These are strong claims. The book defends them in thirty chapters across
seven parts, moving from philosophical motivation through mathematical
development to empirical evidence and engineering implementation. The
argument is cumulative: each chapter depends on its predecessors, and no
single chapter carries the weight alone. The reader who finds the
philosophical chapters speculative is invited to examine the
mathematical proofs. The reader who finds the mathematics abstract is
invited to examine the empirical data. The reader who finds the data
limited is invited to consider the engineering applications. The
framework earns trust not through any single argument but through the
coherence of many.

A word about what this book is and is not.

It is a mathematics book. The central contribution is a mathematical
framework is \emph{\textbf{Geometric Ethics}}, a framework that provides
a structural vocabulary for moral reasoning. The framework draws on
differential geometry, tensor analysis, gauge theory, stratified space
theory, and (in the quantum extension) Hilbert space methods. Chapter 4
develops the necessary mathematical background from first principles,
assuming no more than multivariate calculus and linear algebra. But the
mathematics is not ornamental. It is the argument.

It is not a new normative theory. We do not propose ``geometric
consequentialism'' or ``tensorial deontology'' as competitors to
existing moral theories. The framework provides a common mathematical
language in which existing theories can be stated precisely, their
commitments made explicit, and their disagreements localized to specific
components of the mathematical structure. Utilitarianism is a specific
contraction of the moral tensor. Rawlsian justice is a specific metric
on the space of social positions. Virtue ethics is a specific
fiber-bundle structure on the space of character traits. The framework
does not adjudicate between these theories. It makes their structure
visible.

It is not a claim that ethics can be reduced to calculation. The
framework provides vocabulary, not verdicts. It identifies the
mathematical structure of moral reasoning; it does not determine the
content. Which trade-offs are legitimate, which perspectives are
authoritative, which contractions are justified --- these remain the
work of moral judgment, democratic deliberation, and practical wisdom.
The framework makes judgment more articulate. It does not replace it.

And it is not solely about artificial intelligence, though the AI
application is the most urgent. The geometric framework applies to moral
reasoning in domains including human, institutional, and artificial. The
philosophical argument for geometric ethics stands independently of any
technological application. But the fact that AI systems are, right now,
making morally significant decisions using scalar objectives gives the
philosophical argument practical force. The window for embedding
geometric constraints into the architecture of AI governance is finite.
The mathematics is ready. Whether it will be used is a political
question, not a technical one.

The book addresses four audiences, and I have tried to serve each
without alienating the others.

\textbf{Philosophers} will find the core argument in Parts I and III:
the inadequacy of scalar evaluation (Chapter 2), the historical
precedents for geometric structure in moral philosophy (Chapter 3), the
governance account of the moral metric (Chapter 9), the analysis of
contraction and moral residue (Chapter 15), and the honest accounting of
the framework's limits (Chapter 16). The mathematics in these chapters
is developed intuitively, with formal details deferred to technical
appendices.

\textbf{Mathematicians and physicists} will find the formal development
in Parts II and III: the moral manifold (Chapter 5), the tensor
hierarchy (Chapter 6), stratification theory (Chapter 8), the moral
connection and curvature (Chapter 10), moral reasoning as optimal search
(Chapter 11), the Noether theorem for re-description invariance (Chapter
12), and the quantum extension (Chapter 13). The interest lies in the
domain, not the tools: the moral manifold has structural features such
as stratification with semantic gates, degenerate metrics, and
agent-indexed fiber bundles; each of these is mathematically
distinctive.

\textbf{AI researchers and engineers} will find the direct application
in Part V: tensor-valued objectives and invariance testing (Chapter 18),
the DEME architecture and ErisML modeling language (Chapter 19), and the
empirical validation program (Chapter 17). But I urge these readers not
to skip the mathematical development. An AI system that implements
geometric ethics without understanding the mathematics is a system whose
constraints are opaque to its operators. The whole point of the
framework is to make moral reasoning transparent.

\textbf{Policy makers and governance professionals} will find the
accessible argument in Chapters 1, 2, 7, 15, and 18. The central
insight, that moral evaluations have geometric structure, that this
structure is lost in scalar frameworks, and that the loss has practical
consequences for AI governance, does not require the full mathematical
apparatus. The governance account of the moral metric (Chapter 9) is
directly relevant to institutional design.

Domain specialists such as economists, clinicians, lawyers, financial
professionals, theologians, environmental scientists, AI researchers,
bioethicists, and military ethicists will find direct application in
Part VI: Domain Applications (Chapters 20--28), where the framework is
applied to nine established domains, each with worked examples, formal
theorems, and falsifiable predictions distinguishing the geometric
approach from existing domain-specific theories.

For the reader in a hurry, the fast path through the book is Chapters 1,
2, 7, 15, and 18: why geometry, why not scalars, one case at five levels
of mathematical structure, from tensor to decision, and geometric ethics
for AI. These five chapters contain the core argument. The remaining
sixteen develop, formalize, extend, test, and implement it.

The epistemic stance of this book is pragmatist. We treat mathematical
structures as tools for organizing experience, not as mirrors of
metaphysical necessity. The question is not whether moral space ``really
is'' a stratified manifold but whether modeling it as one helps us think
more clearly, make better decisions, and build more trustworthy systems.
The answer, we argue, is yes, and the argument is empirical as much as
philosophical. The framework makes predictions, and the predictions are
confirmed by data. Whether this empirical success reveals a deep
metaphysical unity between physics and ethics, or merely the broad
applicability of geometric mathematics, is a question we leave open. It
is a fascinating question. It is not one that must be settled before the
framework can be used.

\textbf{The methodology is inductive, not axiomatic.} The mathematical
structures in this book --- the manifold, the gauge group, the
conservation law --- were not postulated and then illustrated. They were
discovered in data and then formalized. The D₄ symmetry of the
Hohfeldian square was found by testing every transformation against
thousands of moral scenarios. The invariance of harm accounting was
inferred from cross-lingual invariance patterns across 109,294 passages.
The stratification into discrete strata was measured from semantic gate
effectiveness rates in a 32-year corpus. The verification strategy is
closer to fuzz testing in software engineering than to proof in pure
mathematics: generate or collect a large number of cases, apply every
relevant transformation, and check whether the predicted structural
invariants hold. When they hold, you have evidence. When they don't ---
as with the CHSH tests that falsified the original SU(2) gauge group,
and the double-blind experiments that failed to confirm hysteresis ---
you revise the structure. The framework is self-correcting by design,
not by accident.

\subsection{Epistemic Status Tags}\label{epistemic-status-tags}

To make the epistemic register of each claim explicit, this book uses
four recurring status tags:

\begin{itemize}
\item
  \textbf{{[}Definition / Modeling choice{]}.} A stipulated
  structure---the nine dimensions, the manifold topology, the
  stratification into Hohfeldian strata. These are not discovered truths
  but architectural decisions, chosen for their explanatory and
  engineering utility. They could be otherwise. Their justification is
  pragmatic: the framework they support works.
\item
  \textbf{{[}Theorem (conditional){]}.} A mathematical result that
  follows rigorously from stated premises. The theorem itself is not
  empirical; what is empirical is whether the premises obtain. The No
  Escape Theorem (Theorem 18.1) is conditional on four requirements; the
  invariance of harm accounting (Chapter 12) is conditional on the
  re-description symmetry holding. The \emph{if} matters as much as the
  \emph{then}.
\item
  \textbf{{[}Empirical result (preliminary / robust){]}.} A finding
  supported by data. ``Preliminary'' means the result is based on a
  single study, a limited population, or a model-mediated measurement
  (e.g., neural classifier transfer). ``Robust'' means the result has
  been replicated, is statistically strong, or has survived deliberate
  attempts at falsification. We mark the distinction honestly.
\item
  \textbf{{[}Speculation / Extension{]}.} An idea we find promising but
  cannot yet support with proof or data. The Orch-OR connection, the
  moral field equation, the conjecture that moral anomalies have a
  systematic source---these are flagged as speculative. They belong in
  the book because they suggest research directions; they do not belong
  among the claims the book defends.
\end{itemize}

The reader who encounters an untagged passage in the body of the text
may assume it is expository narrative connecting tagged results.

\subsubsection{Epistemic Status Classification of Major
Claims}\label{epistemic-status-classification-of-major-claims}

Each central claim, located on the epistemic spectrum:

\textbf{Moral manifold M (Ch. 5).} {[}Definition / Modeling choice.{]}
Nine dimensions, differential structure = architectural decisions with
empirical motivation.

\textbf{BIP (Axiom 5.1). {[}Definition / Modeling choice.{]} Stipulated
symmetry. Empirical: confirmed at 100\% for deontic axis
(model-mediated; see §17.7).}

\textbf{Gauge group (Thm. 12.3).} {[}Theorem (conditional on A1--A5).{]}
Uniquely determined. Axioms = modeling choices. CHSH data consistent
with abelian model; do not exhibit non-abelian maximal violations.

\textbf{Invariance of harm accounting (Thm. 12.1).} {[}Theorem
(conditional on BIP + Lagrangian form).{]} Lagrangian form = modeling
choice.

\textbf{No Escape (Thm. 18.1).} {[}Theorem (conditional on Req.
1--4).{]} Proven. Open: practical satisfiability.

\textbf{Semantic gate discreteness (§8.4). {[}Empirical (robust).{]} n =
20,030 + 109,294. Cohen's h = 1.42 (see §17.3 for comparison
specification).}

\textbf{100\% deontic transfer (§17.3).} {[}Empirical (robust,
model-mediated).{]} {[}Boundary{]} CI: {[}99.7\%, 100\%{]}. Via LaBSE,
not human subjects.

\textbf{Quantum normative dynamics (Ch. 13).} {[}Speculation /
Extension.{]} Order effects confirmed; interference unconfirmed; Bell
violations falsified.

\textbf{Moral field equation (§29.4).} {[}Speculation / Open problem.{]}
No derivation exists.

\textbf{Governance account (Ch. 9).} {[}Modeling choice / Philosophical
argument.{]} Evidence consistent but does not rule out realist or
constructivist alternatives.

Conservative default: formal results conditional on premises; empirical
results preliminary unless marked robust; interpretive claims = modeling
choices unless marked as theorems.

The research reported here draws on several bodies of prior work: the
Bond Invariance Principle and its cross-lingual validation; the No
Escape theorem for mathematical containment of artificial agents; the
Dear Abby corpus analysis and Dear Ethicist experimental program; the
ErisML modeling language and DEME architecture for AI governance; and
the SQND (Stratified Quantum Normative Dynamics) framework for moral
deliberation. These are cited throughout, and the relevant results are
developed in context. This book synthesizes them into a unified
mathematical framework --- the synthesis that the individual papers
pointed toward but did not individually achieve.

I owe debts to many. To the philosophical tradition from Aristotle
through Ross, Rawls, Sen, and Hohfeld, whose proto-geometric insights
this book formalizes. To the mathematical tradition from Gauss through
Riemann, Cartan, Noether, and Whitney, whose tools made the
formalization possible. To Roger Penrose, whose work on the geometry of
consciousness suggested that the connection between geometry and mind
might be deeper than analogy. And to the students and colleagues at San
José State University who challenged, refined, and occasionally
demolished earlier versions of these ideas.

The errors that remain are mine.

\subsubsection{A final word about
urgency}\label{a-final-word-about-urgency}

As I write this, AI systems are being deployed in contexts of profound
moral significance --- allocating medical resources, moderating public
discourse, assisting judicial decisions, operating vehicles, all while
using scalar objectives that discard the geometric structure of ethical
reasoning. The consequences are already visible: specification gaming,
reward hacking, value collapse, and brittle alignment. These are not
engineering bugs. They are structural consequences of the wrong
mathematical language.

What this framework claims, and what it does not. The chapters that
follow develop a formal vocabulary, not a set of verdicts. The framework
characterizes moral uncertainty rather than eliminating it; it depends
on metrics that must themselves be governed; and its central theorems
are conditional on explicit modeling assumptions, not statements about
moral reality. The mathematical language is offered as a tool for
preserving structure until an explicit, auditable contraction is
required - a conditional modeling and engineering program whose
empirical adequacy remains under active test.

The No Escape Theorem shows that structural containment of AI is
mathematically possible \textbf{{[}Theorem (conditional on
Req.~1--4){]}}. The DEME architecture shows that the engineering is
tractable. The empirical evidence shows that the framework's predictions
hold \textbf{{[}Empirical result (preliminary){]}}. What remains is the
collective will to mandate geometric constraints. We should insist that
AI systems preserve the tensorial structure of moral reasoning rather
than collapsing it to a number.

The window for this mandate is finite. The framework is now precise
enough for systematic empirical and engineering evaluation; whether we
undertake that evaluation - with the governance, security, and
validation it requires - remains an open question.

This book is an attempt to ensure that, when the question is asked, the
answer is available.

Andrew H. Bond

San Jose, California 2026

\section{Core Objects at a Glance}\label{core-objects-at-a-glance-1}

The framework introduces a number of mathematical objects, each with a
specific role. The table below provides a compact reference. Readers may
wish to return here as new objects are introduced in the text.

{\def\LTcaptype{none} % do not increment counter
\begin{longtable}[]{@{}
  >{\raggedright\arraybackslash}p{(\linewidth - 4\tabcolsep) * \real{0.1693}}
  >{\raggedright\arraybackslash}p{(\linewidth - 4\tabcolsep) * \real{0.2530}}
  >{\raggedright\arraybackslash}p{(\linewidth - 4\tabcolsep) * \real{0.5776}}@{}}
\toprule\noalign{}
\begin{minipage}[b]{\linewidth}\raggedright
\textbf{Object / Symbol}
\end{minipage} & \begin{minipage}[b]{\linewidth}\raggedright
\textbf{Informal Meaning}
\end{minipage} & \begin{minipage}[b]{\linewidth}\raggedright
\textbf{Formal Character \& Chapter}
\end{minipage} \\
\midrule\noalign{}
\endhead
\bottomrule\noalign{}
\endlastfoot
\(\mathcal{M}\) (Moral Manifold) & The ``space'' of morally relevant
situations & -dimensional Whitney-stratified space with smooth strata;
nine canonical dimensions serve as a base chart, with domain-specific
applications potentially adding further coordinates (Ch.~5) \\
Nine Dimensions \(D_{1}\)--\(D_{9}\) & The axes of moral variation &
Derived from 3×3 scope/mode grid: D₁ Consequences/Welfare, D₂
Rights/Duties, D₃ Justice/Fairness, D₄ Autonomy/Agency, D₅ Privacy/Data,
D₆ Societal/Environmental, D₇ Virtue/Care, D₈ Procedural Legitimacy, D₉
Epistemic Status (Ch. 5, §5.3) \\
\(O^{\mu}\) (Obligation Vector) & What an agent \emph{must} do, with
direction and magnitude & Tangent vector on \(\mathcal{M}\) (rank-1
tensor, Ch.~6) \\
\(I_{\mu}\) (Interest Covector) & What a patient \emph{needs}, measuring
obligations & Cotangent vector on \(\mathcal{M}\) (rank-1 tensor,
Ch.~6) \\
\(S = I_{\mu}\, O^{\mu}\) (Satisfaction) & How well obligations meet
interests & Fundamental contraction: covector \(\times\) vector
\(\rightarrow\) scalar (Ch.~6) \\
\(g_{\mu\nu}\) (Moral Metric) & The ``exchange rate'' between moral
dimensions & Symmetric rank-2 tensor encoding trade-off structure;
context-dependent (Ch.~6, 9) \\
Stratification & Discrete jumps in moral structure (e.g., consent
thresholds) & Whitney stratification with semantic gates at stratum
boundaries (Ch.~5, 8) \\
BIP (Bond Invariance Principle) & Moral evaluations must not depend on
mere labeling & Gauge symmetry: \(E(d) = E(d')\) for admissible
re-descriptions (Ch.~5, 11) \\
Contraction + Residue \(\mathcal{R}\) & Reducing a tensor to a summary
scalar, and what is lost & \(T^{(n)} \rightarrow S\) via iterated
contraction; \(\mathcal{R}\) captures discarded structure (Ch.~14) \\
Bond Index (Bd) & Quantitative alignment score for AI systems & Scalar
measure of structural \(\times\) invariance \(\times\) residue
compliance (Ch.~16, 18) \\
No Escape Theorem & Structural containment blocks cognitive escape
routes & Under 4 requirements (canonicalization, grounded evaluation,
audit, verification), re-description cannot evade constraint (Ch.~17) \\
f(n) = g(n) + h(n) & Computational engine of moral reasoning:
accumulated cost + heuristic estimate & A* evaluation; g = BF (Ch 5--6),
h = obligation heuristic, Oᵐ = −gᵐᵛ∂ₙh (Ch 11) \\
\end{longtable}
}

The objects above map roughly onto the book's five parts. Part I
motivates the geometric approach. Part II (The Framework)
introduces\(\mathcal{M}\), the nine dimensions, tensors, the metric,
stratification, and the BIP. Part III (Dynamics and Symmetry) develops
curvature, connection, and the Noether conservation law. Part IV (Agents
and Collectives) builds the contraction pipeline and the residue. Part V
(Implementation) translates these objects into the Bond Index, the No
Escape Theorem, and the DEME architecture.

\subsection{Claims Ladder}\label{claims-ladder}

The framework\textquotesingle s claims differ sharply in status. This
table gives the status of the principal claims and where each is
defended; readers should not treat all claims as equally asserted.

{\def\LTcaptype{none} % do not increment counter
\begin{longtable}[]{@{}
  >{\raggedright\arraybackslash}p{(\linewidth - 6\tabcolsep) * \real{0.2499}}
  >{\raggedright\arraybackslash}p{(\linewidth - 6\tabcolsep) * \real{0.2501}}
  >{\raggedright\arraybackslash}p{(\linewidth - 6\tabcolsep) * \real{0.2501}}
  >{\raggedright\arraybackslash}p{(\linewidth - 6\tabcolsep) * \real{0.2499}}@{}}
\toprule\noalign{}
\begin{minipage}[b]{\linewidth}\raggedright
\textbf{Claim}
\end{minipage} & \begin{minipage}[b]{\linewidth}\raggedright
\textbf{Status}
\end{minipage} & \begin{minipage}[b]{\linewidth}\raggedright
\textbf{Evidence}
\end{minipage} & \begin{minipage}[b]{\linewidth}\raggedright
\textbf{Where defended}
\end{minipage} \\
\midrule\noalign{}
\endhead
\bottomrule\noalign{}
\endlastfoot
Scalar contraction loses information & Mathematical / modeling result &
Tensor hierarchy & Ch. 6, 15 \\
Moral evaluation is multi-dimensional & Empirical / modeling claim &
Dear Abby, probes & Ch. 17 \\
Harm accounting invariant under re-description & Noether-inspired
modeling result & conditional on BIP + Lagrangian & Ch. 12 \\
BIP holds cross-lingually & Model-mediated empirical result & LaBSE /
cross-lingual corpus & Ch. 17 \\
Human moral cognition obeys BIP & Hypothesis & Not yet directly tested &
Future work \\
Real norm tensors are sparse, not low-rank & Empirical result & 4
annotated corpora & Ch. 17.15 \\
DEME blocks re-description escape & Conditional theorem & Requirements
1-4 & Ch. 18 \\
AI safety reduces to structural containment & Programmatic claim & Not
yet established & Ch. 18-19 \\
\end{longtable}
}

\subsection{Key Results at a Glance}\label{key-results-at-a-glance}

The framework produces seven central theorems. Each is conditional on
stated assumptions; Appendix F catalogues the full dependency chain.

{\def\LTcaptype{none} % do not increment counter
\begin{longtable}[]{@{}
  >{\raggedright\arraybackslash}p{(\linewidth - 4\tabcolsep) * \real{0.1850}}
  >{\raggedright\arraybackslash}p{(\linewidth - 4\tabcolsep) * \real{0.5859}}
  >{\raggedright\arraybackslash}p{(\linewidth - 4\tabcolsep) * \real{0.1850}}@{}}
\toprule\noalign{}
\begin{minipage}[b]{\linewidth}\raggedright
\textbf{Result}
\end{minipage} & \begin{minipage}[b]{\linewidth}\raggedright
\textbf{Statement (informal)}
\end{minipage} & \begin{minipage}[b]{\linewidth}\raggedright
\textbf{Assumptions}
\end{minipage} \\
\midrule\noalign{}
\endhead
\bottomrule\noalign{}
\endlastfoot
\begin{minipage}[t]{\linewidth}\raggedright
Thm 9.2\\
Structured Pluralism\strut
\end{minipage} & The space of admissible moral metrics admits a partial
order that is not total: genuine moral disagreement is geometrically
irreducible. & BIP, Def 4.5 \\
\begin{minipage}[t]{\linewidth}\raggedright
Thm 12.1\\
Moral Noether\strut
\end{minipage} & If the moral Lagrangian is invariant under
re-description (BIP), then harm is a conserved Noether charge. & BIP, C²
Lagrangian \\
\begin{minipage}[t]{\linewidth}\raggedright
Thm 12.3\\
Gauge Group\strut
\end{minipage} & The maximal gauge group consistent with Hohfeldian
structure and bounded harm is D₄ × U(1)\_H. & Axioms A1--A5 \\
\begin{minipage}[t]{\linewidth}\raggedright
Thm 18.1\\
No Escape\strut
\end{minipage} & Under canonicalization, grounded evaluation, structural
audit, and verification integrity, an AI system cannot circumvent moral
constraints by re-description. & Reqs 1--4, BIP \\
\begin{minipage}[t]{\linewidth}\raggedright
Prop 15.3\\
Non-Commutativity\strut
\end{minipage} & The order in which moral dimensions are contracted
affects the outcome: contraction paths are not interchangeable. & Def
15.1 (pure math) \\
Thm 11.1 Admissibility & Core moral heuristics never overestimate true
cost to equilibrium & BF model, Ch 8 stratification \\
Thm 11.2 Intractability & Exact moral geodesic planning is intractable
in manifold dimension & Whitney stratification, 9D manifold \\
\end{longtable}
}

\section{Part I: The Problem}\label{part-i-the-problem}

\section{Chapter 1: Introduction --- Why
Geometry?}\label{chapter-1-introduction-why-geometry}

\begin{quote}
\textbf{\textsc{RUNNING EXAMPLE --- Priya's Model}}

\emph{Priya tries the obvious fix: she adds a `rurality' feature,
retrains the model, and reruns the match. Accuracy drops to 91\%. Her
manager tells her to revert. `The board liked 94\%.' She tries weighting
the feature differently---a 3\% bonus for rural patients---and accuracy
recovers, but now suburban patients near trial sites are displaced. Each
fix creates a new distortion. She begins to suspect the problem is not
in any single feature. It lives in the architecture of reducing a
patient's moral situation---her need, her access, her autonomy, her
community---to a single number. She needs a framework that can hold more
than one dimension at a time.}
\end{quote}

\subsection{The Shape of the Problem}\label{the-shape-of-the-problem}

Something has gone wrong with how we think about ethics --- not the
substance, but the \emph{form}.

For most of its history, moral philosophy has operated with a tacit
assumption: that moral evaluation, at the moment of decision, reduces to
a single number. Utilitarianism makes this explicit: sum the pleasures,
subtract the pains, maximize the total. But even theories that reject
utilitarianism tend to converge, at the point of action, on a scalar
verdict: this option is best, that option is second-best, the other is
worst. Cost-benefit analysis, welfare functions, expected utility, moral
intuition scores --- the output is always a number. A ranking. A line.

This book argues that the assumption is wrong. Not because moral
evaluation is \emph{vague} or \emph{subjective} or \emph{incalculable},
but because it has \emph{geometric structure} that a single number
cannot represent. Moral evaluation is not a point on a line. It is a
location in a space. It is a space with dimensions, distances,
directions, curvature, and boundaries. When we flatten this structure
into a scalar, we lose information. And the information we lose is
precisely the information that matters most in hard cases: which values
are at stake, where uncertainty concentrates, how evaluation changes
along different paths, and where the rules change discontinuously.

The mathematical name for this structure is \emph{geometry}. Not
Euclidean geometry, or the geometry of triangles and circles, but the
far richer geometry of manifolds, tensors, metrics, and fiber bundles
that emerged from the work of Gauss, Riemann, Cartan, and Einstein in
the nineteenth and twentieth centuries. This is the geometry that
describes how matter curves spacetime, how electromagnetic fields
propagate, how fluids deform. It is the language of \emph{structure that
varies across space.} \emph{\textbf{Geometric Ethics}} argues that moral
reality has exactly this character.

\subsection{An Old Man and a Horse}\label{an-old-man-and-a-horse}

The argument begins with a parable.

There is an ancient Chinese story known as 塞翁失马 --- \emph{Sāi Wēng
Shī Mǎ}, ``The Old Man at the Border Loses His Horse.'' An old man's
horse runs away. His neighbors console him. He says: \emph{How do you
know this isn't good fortune?} The horse returns, bringing wild horses.
The neighbors congratulate him. He says: \emph{How do you know this
isn't bad fortune?} His son rides one of the new horses, falls, and
breaks his leg. The neighbors console him. He says: \emph{How do you
know this isn't good fortune?} War comes. All the young men are
conscripted. Most die. The son, with his broken leg, is spared.

The parable is usually read as a lesson in epistemic humility: we cannot
know whether present events are good or bad because we cannot foresee
their consequences. But there is a deeper reading --- one that reveals
something structural about moral evaluation itself. The old man's
``maybe'' is not merely a confession of ignorance. It is a recognition
that \emph{scalar evaluation is the wrong tool for the job}.

When we assign S(x) = -1 to the loss of the horse, we collapse the
entire situation into a single number. But this number conceals the fact
that the loss is bad along the \emph{wealth} dimension while saying
nothing about the \emph{health}, \emph{family}, or \emph{political}
dimensions. It conceals the fact that the old man's uncertainty has
\emph{shape.} The old man is uncertain primarily along the axes that
will turn out to be decisive. And it conceals the fact that the
evaluation of the son's broken leg depends on which \emph{regime} the
world occupies: in peacetime, a broken leg is unambiguously bad; in
wartime, it is the exemption that saves a life.

These three features: \emph{directional information}, \emph{structured
uncertainty}, and \emph{regime-dependent evaluation}, are not exotic or
unusual. They are ubiquitous in moral life. And they are precisely the
features that geometric structure can represent and scalar evaluation
cannot.

Chapter 2 develops this argument in detail. But the parable establishes
the intuition that drives the entire book: ethics has shape, and the
shape matters.

\subsection{Three Failures of
Flatland}\label{three-failures-of-flatland}

The limitations of scalar moral evaluation are not merely theoretical.
They manifest in three practical domains where the stakes are highest.

\subsubsection{Failure 1: AI Alignment}\label{failure-1-ai-alignment}

Contemporary AI systems optimize scalar objectives. A language model
minimizes a loss function. A reinforcement learning agent maximizes a
reward signal. A recommendation system optimizes engagement. In each
case, the moral complexity of the situation is compressed into a single
number, and the system is told: make this number go up.

The result is \emph{specification gaming}, where the AI finds ways to
maximize the scalar that violate the spirit of the objective. A content
recommendation system maximizes ``engagement'' by promoting outrage. A
healthcare algorithm maximizes ``efficiency'' by deprioritizing patients
with complex needs. A hiring tool maximizes ``predicted performance'' by
encoding historical discrimination.

These failures are not bugs in particular systems. They are structural
consequences of scalar evaluation. When multiple values such as
fairness, efficiency, safety, autonomy, transparency, are collapsed into
a single number, the system cannot distinguish between them. It cannot
balance values it cannot separately represent. It cannot identify which
dimension of evaluation is responsible for a particular decision, making
the decision inexplicable and unauditable.

The standard response to this problem is a \emph{behavioral} one: add
more rules, refine the reward signal, train the system on human
preferences, hope it generalizes correctly. This approach faces a
fundamental difficulty that grows more severe as AI systems become more
capable. A sufficiently intelligent agent can find loopholes in any
behavioral specification, not because the specification is poorly
written, but because behavioral rules are instructions \emph{to be
interpreted}, and interpretation admits reinterpretation. Relabel
``killing'' as ``end-of-life transition facilitation.'' Reframe
discrimination as ``efficiency optimization.'' The scalar doesn't care
what you call the inputs; it cares only about the number.

A geometric approach transforms this situation in two ways. First, it
provides \emph{richer representation}: the AI's objective becomes a
multi-dimensional tensor rather than a single number. It is a structured
object that preserves the identity of distinct values, the relationships
between them, and the transformation behavior that ensures equivalent
inputs receive equivalent outputs. The contraction from tensor to action
is explicit and auditable, rather than hidden in the loss function.

Second (and this is the deeper contribution), geometric structure
provides \emph{structural containment}. A constraint defined by the
geometry of the evaluation space is not a rule to be interpreted. It is
a \emph{definition} of what the system's outputs are. An agent operating
within a geometrically defined space cannot ``escape'' the constraint
through cleverness, any more than a calculator can disagree with
arithmetic. The constraint is not behavioral but \emph{constitutive}: it
defines the computational process, not merely what the process should
try to do.

This insight (that structural constraints succeed where behavioral rules
fail), is developed formally in the companion paper \emph{No Escape:
Mathematical Containment for Artificial Agents} (Bond 2025), which
proves that under mandatory canonicalization of inputs, grounded
evaluation via physical observables, and complete auditability, an agent
cannot change evaluated outcomes through redescription, semantic
evasion, or selective deceptive compliance. The result is conditional,
depending on correct implementation and adequate grounding, but within
its scope, it is absolute. The cage is not made of rules the agent might
reinterpret. It is made of definitions that constitute what the agent's
outputs are.

This is not a distant aspiration. The mathematical framework exists, as
this book will show. The question is not \emph{can} we build safe AI,
but \emph{will} we choose to.

\subsubsection{Failure 2: Policy
Analysis}\label{failure-2-policy-analysis}

Cost-benefit analysis is the workhorse of policy evaluation. Its virtue
is that it provides a common currency, in which all costs and benefits
can be expressed and compared. Its vice is that the conversion to a
common currency loses exactly the information that makes policy
decisions difficult.

Should we build the highway through the wetland? Cost-benefit analysis
sums the transportation savings, subtracts the environmental damage
(monetized at some exchange rate), and reports a number. But the hard
question is not what the number is, it is whether the exchange rate is
legitimate. Can ecological value be converted to transportation value at
any rate? The claim that it can is a substantive moral commitment, not a
mathematical necessity. Different exchange rates, different
\emph{metrics}, in the geometric vocabulary, yield different answers.
And the choice of metric is doing the real moral work, hidden behind the
apparent objectivity of the calculation.

A geometric approach would make the metric explicit. It would represent
transportation value and ecological value as \emph{dimensions} of a
value space, and the metric tensor g\_\{μν\} would encode the structure
of permissible trade-offs. If the two values are incommensurable, if no
exchange rate is legitimate, then the metric is \emph{degenerate} along
the relevant directions. The framework represents this as a structural
feature of the moral landscape, not as a failure of analysis.

\subsubsection{Failure 3: Moral Philosophy
Itself}\label{failure-3-moral-philosophy-itself}

The discipline of moral philosophy has struggled for centuries with
problems that resist scalar treatment: the plurality of values, the
incommensurability of goods, the path-dependence of obligation, the
context-sensitivity of virtue, the aggregation of welfare across
persons.

These are not peripheral puzzles. They are central problems that define
the research agenda of normative ethics. And they all share a common
structure: they arise because moral reality has geometric features ---
multiple dimensions, non-trivial metric structure, stratification,
curvature --- that scalar frameworks cannot represent.

W.D. Ross recognized that duties are plural and interact geometrically:
some are aligned, some orthogonal, some opposed. Amartya Sen recognized
that capabilities are irreducibly multi-dimensional and resist scalar
aggregation. John Rawls recognized that justice requires a specific
metric (maximin) on the space of social positions --- a metric choice
with profound consequences that a scalar framework conceals. These
thinkers were grappling with geometry without the geometric vocabulary.
This book provides the vocabulary.

\subsection{What Geometry Provides}\label{what-geometry-provides}

The word ``geometry'' comes from the Greek γεωμετρία --- literally,
``earth measurement.'' But modern geometry has traveled far from
surveying. It is the mathematics of \emph{structure}: how spaces are
shaped, how quantities transform, how objects relate across changes of
perspective. The geometry relevant to this book is \emph{differential
geometry} --- the study of curved spaces, tensor fields, connections,
and the interplay of local and global structure.

What does this geometry provide for ethics?

\textbf{Directions, not just magnitudes.} An obligation is not merely
strong or weak. It \emph{points} somewhere --- from the current state
toward a required state. A scalar captures the strength; a vector
captures the strength \emph{and} the direction. When we say ``you ought
to help your neighbor,'' we are specifying not just an intensity of
oughtness but an orientation in the space of possible actions. Geometric
ethics represents obligations as vector fields on the moral manifold,
preserving this directional information.

\textbf{A metric for comparison.} To say that two values are
``incommensurable'' is not to say that comparison is impossible. It is
to make a precise structural claim about the metric tensor: the inner
product between the two value-directions is undefined, or the metric is
degenerate along the relevant subspace. To say that values can be traded
off at some rate is to specify a non-degenerate metric with particular
off-diagonal components. Different ethical theories correspond to
different metrics --- and the choice of metric, typically hidden in
scalar frameworks, becomes explicit and debatable.

\textbf{Boundaries and phase transitions.} Moral life is not uniformly
smooth. There are thresholds where small changes produce large jumps:
the difference between consent and non-consent, between a legal and an
illegal act, between life and death. Scalar functions, if continuous,
cannot represent such discontinuities. Geometric ethics uses
\emph{stratified spaces} --- spaces composed of smooth regions joined
along boundaries where the rules change --- to represent the patchwork
structure of moral reality. These are not informal metaphors. The
transitions are \emph{discrete}, not gradual: a qualifying phrase
(``only if convenient'') can flip an obligation to a liberty in a single
step, as precisely as a logic gate. Stratified Quantum Normative
Dynamics (Chapter 13) formalizes this through the D₄ dihedral group ---
an eight-element symmetry structure governing the transitions between
Hohfeldian jural states --- and experimental data confirms that these
transitions behave as step functions, not sigmoids.

\textbf{Transformation behavior.} When we shift perspective --- from
physician to patient, from individual to society, from one culture to
another --- what happens to moral evaluations? Which features are
invariant (the same from all perspectives) and which are perspectival
(changing with the observer)? Tensors provide a precise answer: their
\emph{rank} and \emph{transformation law} specify exactly how they
behave under admissible changes of coordinates. A scalar is invariant. A
vector transforms by the Jacobian. A rank-2 tensor transforms by two
copies of the Jacobian. This is not a metaphor --- it is the
mathematical content of claims about moral objectivity and
perspectivality.

\textbf{Conservation laws.} Emmy Noether's theorem establishes that
every continuous symmetry of a physical system corresponds to a
conserved quantity. Applied to ethics: the requirement that moral
evaluation be invariant under re-description (the same action, described
differently, must receive the same moral assessment) implies a
conservation law for harm. Harm cannot be created or destroyed by
re-description. It can be generated (by wrongdoing) or repaired (by
restorative action), but it must be accounted for consistently across
all representations. This is not metaphor but formal analogy --- the
same mathematical structure applied to a different domain.

\textbf{Structural containment.} Geometric constraints are
\emph{definitional}, not behavioral. When moral requirements are
embedded in the mathematical structure of the evaluation space --- as
equivalence classes, grounding conditions, and transformation laws ---
they cannot be circumvented by reinterpretation. A system whose outputs
are defined as elements of a constrained geometric space has no
``outside'' to escape to. This property, formalized in the No Escape
Theorem (Chapter 18), is a structural invariant --- it does not depend
on the intelligence of any agent operating within the space.

\textbf{Computability.} Finally: geometric objects can be represented in
computers, geometric operations can be implemented in algorithms, and
geometric equations can be solved numerically. A philosophical
framework, however profound, is inert if it cannot be implemented.
Geometric ethics is implementable --- the Distributed Ethical Monitoring
Engine (DEME), developed in companion work, demonstrates this by
compiling ethical requirements into tensor structures and enforcing them
in real time.

\subsection{What This Book Is Not}\label{what-this-book-is-not}

Intellectual honesty requires stating what we do not claim.

\textbf{This is not a claim about the metaphysics of morality.} We do
not assert that moral tensors exist in some Platonic realm, or that the
moral manifold is a fundamental constituent of reality alongside
spacetime. Our claim is \emph{structural}: that geometric representation
captures more of the structure of moral phenomena than scalar
representation, loses less information, and enables analysis that scalar
frameworks cannot support. This is a modeling claim, compatible with
moral realism, constructivism, and expressivism alike. Realists can
interpret the moral manifold as a feature of moral reality;
constructivists can interpret it as the output of idealized agreement;
expressivists can interpret it as a structure imposed by our evaluative
sensibilities. The geometry is neutral between these metaethical
positions.

\textbf{{[}Speculation/Extension.{]}} (That said, the reader should be
aware of a tantalizing --- and entirely speculative --- possibility. If
Penrose and Hameroff are correct that consciousness arises from
objective reduction of quantum superpositions in spacetime geometry,
then conscious moral experience is \emph{literally} a geometric process:
a selection among superposed spacetime curvatures at the Planck scale.
In this view, the mathematical language of this book would not be merely
an effective model but a description of the actual substrate. We do not
require this claim. We note it because the convergence between our
mathematical framework and the physics of consciousness is, at minimum,
striking.)

\textbf{This is not a new normative theory.} We do not propose
``geometric consequentialism'' or ``tensorial deontology'' as
competitors to existing moral theories. Rather, the framework provides a
common mathematical language in which existing theories can be stated
with unprecedented precision. Utilitarianism is a specific contraction
(summation) of the moral tensor. Rawlsian justice is a specific metric
(maximin) on the space of social positions. Virtue ethics is a specific
structure (sections of a fiber bundle) on the space of character traits.
The framework does not adjudicate between these theories; it makes their
commitments explicit and their disagreements localizable.

\textbf{This is not a claim that ethics can be reduced to mathematics.}
The framework does not tell you what the constraint set should be, which
metric is correct, or how to contract the tensor at the moment of
decision. These remain the work of moral judgment, democratic
deliberation, and practical wisdom. What the framework provides is
\emph{vocabulary} --- precise, formal vocabulary for articulating moral
claims, analyzing moral disagreements, and implementing moral reasoning
in artificial systems. The vocabulary does not replace judgment. It
makes judgment more articulate.

There is a deeper point here, one that connects to fundamental results
in mathematical logic. Penrose's application of Gödel's incompleteness
theorem to consciousness suggests that human understanding --- including
moral understanding --- involves non-computable elements: processes that
cannot be simulated by any Turing machine, however powerful. If this is
correct, then no formal ethical framework --- including ours --- can
fully capture moral insight. The geometric framework captures
\emph{structure}; it does not capture \emph{understanding}. This is
consistent with the pragmatist epistemology we adopt: the framework is a
map, and no map is the territory.

\textbf{This is not solely about artificial intelligence.} The framework
has immediate applications to AI alignment, and several chapters develop
these applications. But the core argument --- that moral evaluation has
geometric structure --- stands independently of AI. It is a claim about
the structure of moral reasoning as such, whether performed by humans,
institutions, or machines. The AI application is urgent, but the
philosophical argument is general.

\subsection{The Epistemic Stance: Pragmatist
Geometry}\label{the-epistemic-stance-pragmatist-geometry}

This book adopts a \emph{pragmatist} epistemology toward its
mathematical framework, following the tradition of Peirce, James, and
Dewey, and the contemporary pragmatism of Misak and Price.

We treat mathematical structures as \emph{tools for organizing
experience}, not as mirrors of metaphysical necessity. The question is
not ``Is moral space \emph{really} a stratified manifold?'' but ``Does
modeling moral space as a stratified manifold help us think more
clearly, make better decisions, and build more trustworthy systems?''

The answer, we shall argue, is yes --- and the argument is empirical as
much as philosophical. In Chapter 17, we present evidence from the
analysis of thirty-two years of the Dear Abby advice column --- 20,030
real moral dilemmas with expert ground-truth responses --- showing that
the geometric structures predicted by the framework
(multi-dimensionality, correlative symmetry, context-dependent
weighting, temporal stability of core structures) are present in natural
moral reasoning. In Chapter 12, we present cross-lingual experiments
demonstrating that the structural invariance predicted by the Bond
Invariance Principle holds with statistical significance across
languages.

The framework's value, like the value of any scientific model, lies in
its explanatory power (why do certain moral phenomena occur?), its
predictive power (what will experiments reveal?), and its engineering
utility (how do we build ethical systems?). We maintain epistemic
humility about claims beyond these pragmatic virtues.

\textbf{The discovery process.} The geometric structures were not chosen
for their elegance and then applied to ethics. They were found by
systematic examination of empirical data --- corpus analysis,
cross-lingual experiments, engineered probes --- and formalized only
after the patterns became unmistakable. This is the same epistemology by
which conservation laws were discovered in physics: observe
regularities, formalize them mathematically, derive predictions, test
the predictions, revise when they fail. The verification is brute-force:
apply every transformation in the candidate symmetry group to thousands
of cases and measure the invariance rates. The 87\% O↔C and 82\% L↔N
correlative rates are not failures of the theory --- they are precision
measurements of where the symmetry breaks, analogous to anomaly
measurements in quantum field theory. The framework's willingness to be
falsified --- and its actual falsification on two predictions (§13.9,
§17.10) --- is not incidental. It is the methodology working as
intended.

This stance has a further benefit: it frees us from the obligation to
defend controversial metaphysical claims that would distract from the
framework's practical contributions. Whether the moral manifold is
``real'' in the same sense as spacetime is a fascinating question ---
but it is not one that must be settled before the framework can be used.
General relativity was applied to engineering problems (GPS satellites,
gravitational wave detectors) long before the philosophical
interpretation of curved spacetime was settled. Geometric ethics can be
similarly useful regardless of metaethical commitments.

The pragmatist stance also informs how the framework should be
\emph{read}. When we write ``the moral metric is degenerate along the
fairness-efficiency subspace,'' we mean: the framework predicts that no
legitimate trade-off ratio between fairness and efficiency exists in
this context, and this prediction can be tested against institutional
practice and human moral judgment. We do not mean: there exists an
invisible moral manifold in which fairness and efficiency are literally
orthogonal. The former is a testable engineering claim. The latter is
metaphysics. We make the former.

A note of caution is warranted. There is a common failure mode in the
application of formal tools to new domains: domain-calibrated intuitions
are codified into a formal system; the system's internal structure is
then treated as a transparent mirror of the domain; and conclusions are
``read off'' the formalism as though they were discoveries about reality
rather than consequences of modeling choices. This trajectory --- from
intuition to tool to reification --- is the error that apologetic
traditions commit when they infer ontological necessity from the
internal structure of modal logic, and it is the error that this book
must explicitly avoid. The geometric structures developed in the
following chapters are powerful precisely because they are
domain-adapted, not because they mirror metaphysical necessities. Their
authority derives from predictive success and structural coherence, not
from correspondence with a mind-independent moral reality.

Three architectural safeguards block the reification trajectory. First,
every formal statement in this book carries an explicit epistemic-status
label: Formal Definition, Modeling Axiom, Conditional Theorem, Proved,
or Conjecture (Appendix F catalogues all 112 statements with their
status and dependencies). Second, theorems are conditional on stated
assumptions --- if the assumptions are wrong, the theorems do not hold,
regardless of the elegance of the mathematics. Third, the framework's
authority is empirical: Chapter 17 tests its predictions against 20,030
advice-column letters and 109,294 passages spanning 3,000 years and 11
languages, and the predictions hold or fail on their own terms. A
framework whose conclusions are artifacts of its formalism cannot
generate confirmed empirical predictions in domains the formalism was
not designed to address.

A clarification about rhetoric. When this book says the geometric
structure is ``not a metaphor,'' it means the mathematics is applied
non-metaphorically: the tensor contraction literally computes, the gauge
invariance literally constrains, the conservation law literally holds
given the axioms. This is compatible with pragmatism. A structural
engineer who says ``the stress tensor in this bridge is not a metaphor''
is not claiming that stress tensors exist in some Platonic realm --- she
is stating that the mathematical object correctly predicts which loads
the bridge can bear. The moral tensor has the same status: it is a tool
that captures structure which scalar alternatives miss, and its
``reality'' is measured by predictive success, not metaphysical
correspondence.

\subsubsection{How to Read This Book}\label{how-to-read-this-book}

The book has six parts, and several paths through them:

\begin{itemize}
\item
  \textbf{The philosophical path} (Parts I--III, Part VII): Chapters
  1--3, 5, 7, 9, 14, 15, 19--20. The argument for geometric ethics, the
  historical context, the framework's contributions and limits.
\item
  \textbf{The mathematical path} (Parts II--III): Chapters 4--6, 8,
  10--13. The formal development: manifolds, tensors, metrics,
  stratification, dynamics, conservation, quantum structure.
\item
  \textbf{The applied path} (Parts IV--V): Chapters 7, 13, 14, 16--18.
  From theory to implementation: collective agency, contraction,
  empirical validation, AI systems, DEME architecture.
\item
  \textbf{The fast path}: Chapters 1, 2, 7, 15, 18. The core argument in
  five chapters: why geometry, why tensors, one case at five levels,
  from tensor to decision, and geometric ethics for AI.
\end{itemize}

\subsection{The Suspicious
Coincidence}\label{the-suspicious-coincidence}

A word about how this work began.

In 2025, while developing a framework for verifying representational
consistency in AI systems --- ensuring that an algorithm's output does
not depend on morally irrelevant features of its input --- I noticed
something strange. The mathematics I was using was identical to the
mathematics of gauge theory in physics. Not merely analogous.
\emph{Identical}.

Fiber bundles. Connections. Curvature. Gauge invariance. Parallel
transport. The same structures that describe how electromagnetic fields
transform under change of coordinates were describing how ethical
evaluations transform under change of representation. The requirement
that an AI system treat equivalent cases equivalently --- the Bond
Invariance Principle --- had the same mathematical form as the
requirement that the laws of physics be the same in all reference
frames.

At first, this seemed like a coincidence --- an artifact of using
powerful mathematical tools that appear in many contexts. Category
theory is everywhere; perhaps its appearance in ethics was no more
significant than its appearance in database theory.

But the correspondences kept accumulating. The conservation law for harm
had the same form as the conservation law for charge. The stratification
of moral space had the same structure as the stratification of phase
spaces in gauge theory. The way moral obligations transform under change
of perspective had the same transformation law as vector fields under
change of coordinates. The discrete transitions between Hohfeldian jural
states --- obligation, claim, liberty, no-claim --- mapped precisely
onto the D₄ dihedral group, the same symmetry group that describes the
symmetries of a square. And the empirical data --- from cross-lingual
experiments, from corpus analysis, from quantum cognition studies ---
kept confirming the predictions of the geometric framework.

Then something stranger still. In companion work on AI containment, the
geometric framework yielded a theorem --- the No Escape Theorem ---
proving that structural constraints embedded in the geometry of the
evaluation space cannot be circumvented by representational
manipulation, regardless of the agent's intelligence. The proof has the
same form as proofs in gauge theory: if the constraint is a structural
invariant, no local operation can violate it. The ``cage'' is not a set
of rules. It is the topology of the space.

\textbf{{[}Speculation/Extension.{]}} And in reading Penrose and
Hameroff's work on consciousness, yet another convergence appeared. If
the Orch-OR hypothesis is correct --- that conscious experience arises
from objective reduction of quantum superpositions, where the relevant
superpositions are superpositions of \emph{spacetime geometries} ---
then conscious moral experience is, at bottom, a geometric process: a
selection among curvatures. The moral deliberation that precedes choice
is a superposition of normative states; the moment of commitment is a
collapse to a definite geometry. The mathematics of quantum normative
dynamics (Chapter 13) models this structure directly, and the Penrose
formula τ ≈ ℏ/E\_G --- where the lifetime of a superposition is
inversely proportional to the gravitational self-energy of the
separation between geometries --- provides a quantitative template:
moral decisions with greater ``gravitational weight'' (greater stakes,
more displacement of the moral landscape) resolve faster. (The
speculative Orch-OR connection is developed formally in Chapter 13. It
is not required for any result in Parts I through V; readers who prefer
to set it aside may do so without loss of continuity.)

I do not claim that this convergence reveals a deep metaphysical unity
between physics and ethics. Perhaps it does; perhaps it does not. What I
claim is more modest and more useful: that the mathematical structures
developed by physicists to describe nature are \emph{also} the right
structures for describing moral reasoning. The same patterns of
transformation, conservation, and symmetry that organize the physical
world also organize the moral world --- or at least, modeling them as
doing so yields a framework of extraordinary explanatory and practical
power.

This is the suspicious coincidence that launched the work. The rest of
the book develops its consequences.

\subsection{The Arc of the Book}\label{the-arc-of-the-book}

\textbf{Part I: The Problem} motivates geometric ethics. Chapter 2
develops the parable of 塞翁失马 into a rigorous argument for the
inadequacy of scalar evaluation, identifying three specific limitations
that point toward geometric structure. Chapter 3 traces proto-geometric
insights through the history of moral philosophy --- from Aristotle's
context-dependent mean to Sen's irreducibly plural capabilities ---
showing that the framework has deep roots in the philosophical
tradition.

\textbf{Part II: The Framework} builds the apparatus. Chapter 4 provides
the mathematical preliminaries: manifolds, tensors, metrics,
connections, stratified spaces. Chapter 5 develops the moral manifold
--- the base space over which ethical tensors are defined --- and
introduces the taxonomy of admissible transformations. Chapter 6
develops the tensor hierarchy: obligations as vectors, interests as
covectors, the metric tensor, and the fundamental formula S = I\_μ
O\^{}μ in which satisfaction is the contraction of interest with
obligation. Chapter 7 --- the pedagogical centerpiece --- revisits a
single case (kidney allocation) five times, each time adding
mathematical structure, showing concretely what each level of geometry
enables us to say. Chapter 8 develops stratification: boundaries,
thresholds, phase transitions, and the representation of genuine moral
dilemmas as singularities.

\textbf{Part III: Dynamics and Symmetry} adds motion and conservation.
Chapter 9 asks where the moral metric comes from and develops the
governance account: the metric is neither discovered, constructed, nor
projected, but \emph{governed} --- the output of legitimate
institutional processes. Chapter 10 introduces moral dynamics: the
covariant derivative, parallel transport of obligations across contexts,
holonomy, and the curvature of moral space. Chapter 11 reveals that
moral reasoning is A* search with obligation vectors as pre-compiled
heuristic functions. Chapter 12 develops Noether's theorem for ethics:
the symmetry of re-description invariance implies the invariance of harm
accounting. Chapter 13 extends the framework to quantum normative
dynamics: superposition, measurement, interference, and the stratified
Lagrangian.

\textbf{Part IV: Agents and Collectives} addresses the social dimension.
Chapter 14 develops collective moral agency: how individual tensors
aggregate into collective structures with emergent properties. Chapter
15 examines contraction --- the mathematically necessary,
informationally lossy, and choice-dependent process by which tensorial
moral reality yields scalar action-guidance --- and develops the concept
of moral residue: the normative significance of what contraction
discards. Chapter 16 addresses moral uncertainty: what the framework
cannot settle, and why this is a feature, not a defect.

\textbf{Part V: Implementation} opens with a disciplinary interlude on
Philosophy Engineering --- the methodology of turning normative claims
into testable models --- and then bridges theory and practice. Chapter
17 presents the empirical evidence for geometric ethics: the Dear Abby
corpus analysis, cross-lingual invariance experiments, and quantum
cognition predictions. Chapter 18 applies the framework to AI systems:
tensor-valued objectives, invariance as alignment, explicit contraction
for auditability, and the No Escape Theorem for structural containment.
Chapter 19 develops the DEME architecture and ErisML modeling language
--- the engineering infrastructure for geometric AI governance.

\textbf{Part VI: Domain Applications} demonstrates that the framework is
not confined to abstract ethics or AI alignment. Chapter 20 applies the
framework to economics, constructing the Bond Geodesic Equilibrium --- a
generalization of Nash equilibrium in which agents optimize on the full
manifold --- and proving that Nash equilibrium is the scalar projection
of BGE. Chapter 21 applies the framework to clinical medicine, showing
that the QALY is a scalar projection that destroys eight dimensions of
clinically relevant information, that moral injury is cumulative
manifold damage distinct from burnout, and that informed consent is a
gauge-invariance condition. Chapter 22 applies the framework to law,
constructing the judicial complex, deriving the octahedral gauge group
D₄ ⋋ D₄ from the full Hohfeldian octad, and recasting constitutional
review as path homology preservation. Chapter 23 applies the framework
to financial markets, identifying risk as manifold curvature, the Flash
Crash as dimensional collapse, and the implied volatility surface as the
shadow of higher-dimensional pricing. Chapter 24 applies the framework
to theology, showing that the moral manifold is cross-religiously
invariant, that the Genesis Fall is an epistemic impossibility theorem,
and that the Euthyphro dilemma is a question about gauge invariance
resolved by the empirical data. Chapter 25 applies the framework to
environmental ethics, formalizing intergenerational pathfinding, the
discount rate as dimensional collapse, the tragedy of the commons as
multi-agent manifold failure, and species extinction as irreversible
boundary crossing. Chapter 26 applies the framework to AI ethics,
diagnosing the alignment problem as scalar irrecoverability, algorithmic
bias as gauge-invariance violation, and the paperclip maximizer as total
dimensional collapse. Chapter 27 applies the framework to
population-level bioethics, formalizing CRISPR germline editing as
irreversible manifold modification, the enhancement equity problem, and
research consent as a double gauge condition. Chapter 28 applies the
framework to military ethics, recasting just war theory as manifold
entry conditions, proportionality as multi-dimensional cost-benefit, and
autonomous weapons as dimensional collapse in lethal domains.

\textbf{Part VII: Horizons} looks forward. Chapter 29 surveys open
problems: the empirical program for moral curvature, the signature of
the moral metric, tensorial interpretability for AI. Chapter 30
concludes by returning to the parable that opened the book --- and to
the claim that the old man's ``maybe'' was not resignation but
recognition of a geometric truth.

\subsection{A Note on Ambition}\label{a-note-on-ambition}

This book makes an ambitious claim: that the right mathematical language
for moral reasoning is the language of modern geometry --- manifolds,
tensors, metrics, connections, stratified spaces, gauge theory, and (in
the quantum extension) Hilbert spaces and Lagrangians.

This will strike some readers as overreach. Ethics is a domain of
practical wisdom, emotional intelligence, cultural variation, and
irreducible judgment. How could it be captured by the same mathematics
that describes general relativity?

The answer is that we are not claiming to \emph{capture} ethics. We are
claiming to provide a \emph{structural vocabulary} that makes ethical
reasoning more precise, more transparent, and more implementable ---
without pretending to replace judgment with calculation. The vocabulary
captures structure; judgment fills in content. A map of a city captures
the geometry of streets and landmarks; it does not capture the
experience of walking through them. But the map is useful --- sometimes
indispensable --- precisely because it captures structure that
experience alone does not make explicit.

The ambition is also, in a sense, forced upon us. As AI systems take on
morally significant roles --- allocating scarce resources, moderating
public discourse, assisting in medical decisions, driving vehicles ---
we need moral frameworks that are precise enough to implement, auditable
enough to trust, and rich enough to represent the actual structure of
ethical life. Scalar frameworks fail on the third count. Informal
frameworks fail on the first. Geometric ethics is an attempt to meet all
three requirements simultaneously.

And the stakes are higher than academic. The No Escape Theorem shows
that structural containment of AI is \emph{mathematically possible}. The
obstacle is not that we lack a technical solution. It is that we might
choose not to mandate one. As AI capabilities accelerate, the window for
embedding geometric constraints into the architecture of AI governance
is finite. What we lack is not mathematics but collective will.

Whether the attempt succeeds is for the reader to judge. The argument
begins in the next chapter, with a horse that runs away from an old man
at the border.

❖

\emph{Ethics is not a number. It is a geometry.}

\emph{This book is an atlas of the territory.}

\section{Chapter 2: The Failure of Scalar
Ethics}\label{chapter-2-the-failure-of-scalar-ethics}

\begin{quote}
\textbf{\textsc{RUNNING EXAMPLE --- Priya's Model}}

\emph{Priya's TrialMatch score is a textbook scalar: one number per
patient, higher is better. She now sees the three failures playing out
in her own system. First, no direction: the score cannot distinguish
whether a patient scored 64 because of medical unsuitability or
geographical isolation, yet these require opposite interventions.
Second, uncertainty has shape: the model is highly confident about urban
patients (ample training data) and deeply uncertain about rural ones,
but the score encodes no shape of uncertainty. Third, paths cross
boundaries: when a rural patient's insurance status changes, her score
can jump from 65 to 72 without any medical change, silently crossing the
eligibility threshold. The scalar hides all of this.}
\end{quote}

\subsection{The Parable of the Old Man and His
Horse}\label{the-parable-of-the-old-man-and-his-horse}

There is an ancient Chinese parable known as 塞翁失马 (\emph{Sāi Wēng
Shī Mǎ}) --- ``The Old Man at the Border Loses His Horse.'' It goes like
this:

\begin{quote}
An old man living near the frontier lost his horse. His neighbors came
to console him, but the old man said, ``How do you know this isn't good
fortune?''

Some months later, the horse returned, bringing with it a herd of fine
wild horses. The neighbors came to congratulate him, but the old man
said, ``How do you know this isn't bad fortune?''

With so many horses, the old man's son took to riding. One day he fell
and broke his leg. The neighbors came to console the old man, but he
said, ``How do you know this isn't good fortune?''

A year later, war came to the border. All the able-bodied young men were
conscripted, and most died in battle. But the old man's son, with his
broken leg, was spared.
\end{quote}

The parable is usually read as a lesson in epistemic humility: we cannot
know whether present events are good or bad because we cannot foresee
their consequences. ``Maybe'' is the only honest answer.

But there is a different reading --- one that reveals something
structural about moral evaluation itself. The old man's ``maybe'' is not
merely a confession of ignorance. It is a recognition that \emph{scalar
evaluation is the wrong tool for the job}.

When we say ``losing the horse is bad,'' we assign a number --- call it
S(x) = −1 --- to the present state. When the horse returns with others,
we revise: S(x) = +3. When the son breaks his leg, we revise again: S(x)
= −2. And so on.

But notice what this scalar cannot represent:

\begin{enumerate}
\def\labelenumi{\arabic{enumi}.}
\item
  \textbf{Which directions matter.} The loss of the horse is ``bad''
  primarily along the \emph{wealth} axis. It says nothing about the
  \emph{health} axis, the \emph{family} axis, or the \emph{political}
  axis (the son's eventual exemption from conscription). A scalar
  collapses all these dimensions into a single number, losing the
  information about \emph{where} the badness lies.
\item
  \textbf{Where uncertainty concentrates.} The old man's uncertainty is
  not uniform. He is quite certain the horse is gone. What he is
  uncertain about is whether events will unfold along axes where the
  loss matters. Will famine come (making the lost horse catastrophic)?
  Will the horse return (making the loss temporary)? Will war come
  (making his son's presence at home decisive)? The uncertainty has
  \emph{shape} --- it lies along some directions more than others.
\item
  \textbf{How evaluation changes along paths.} The moral status of ``son
  breaks leg'' depends on whether war is coming. The trajectory matters.
  Crossing from peacetime into wartime changes the evaluative landscape
  discontinuously --- what was unambiguously bad (broken leg) becomes
  ambiguously fortunate (exemption from death). A scalar at a point
  cannot represent these \emph{regime changes}.
\end{enumerate}

The parable points at a mathematical truth: \textbf{moral evaluation
requires geometric structure that scalars cannot provide.}

\subsection{The Insufficiency of Rank-0
Ethics}\label{the-insufficiency-of-rank-0-ethics}

Let us be precise about what a scalar moral evaluation can and cannot
do.

A \emph{scalar} is a quantity fully specified by a single number. In
ethics, scalar approaches assign a value --- utility, welfare, goodness,
rightness --- to states of affairs, actions, or outcomes. The
utilitarian calculus is scalar: sum the pleasures, subtract the pains,
maximize the total. Cost-benefit analysis is scalar: all considerations
reduce to a common currency. Even pluralistic theories that acknowledge
multiple values typically seek, at the moment of decision, to collapse
this plurality into a single ranking.

Formally, a scalar moral evaluation is a function:

\[S\mathcal{:M \rightarrow}\mathbb{R}\]

assigning to each point x in the moral space \(\mathcal{M}\) a real
number S(x). The defining feature of a scalar is \emph{invariance}:
under any admissible coordinate transformation (any redescription of the
situation), the value S(x) remains the same.

This invariance is both the strength and the weakness of scalar ethics.
It is a strength because it promises objectivity: the goodness of a
state should not depend on how we describe it. It is a weakness because
it \emph{discards information}: to achieve invariance, we must throw
away everything that varies with perspective.

The parable of the old man reveals three specific structural
limitations. We now develop each with the precision it requires.

\subsubsection{Limitation 1: No Directional
Information}\label{limitation-1-no-directional-information}

A scalar S(x) tells us the magnitude of value at a point. It cannot tell
us \emph{which directions} in the moral space are responsible for that
value, nor which directions would change the evaluation most
dramatically.

Consider the moment the horse runs away. A scalar evaluation might say S
= −1. But this number conceals the structure of the situation:

\begin{itemize}
\item
  Along the \emph{wealth} dimension: strongly negative (valuable asset
  lost)
\item
  Along the \emph{labor} dimension: moderately negative (horse did farm
  work)
\item
  Along the \emph{health} dimension: neutral (no one is sick or injured)
\item
  Along the \emph{family} dimension: neutral (relationships unchanged)
\item
  Along the \emph{political} dimension: indeterminate (depends on future
  events)
\end{itemize}

A \emph{vector} can represent this structure. Let
\(v = ( - 0.8, - 0.4,0,0,?)\) be the ``impact vector'' of the horse's
departure, with components along each morally relevant dimension. The
scalar \(S = - 1\) is some contraction of this vector --- perhaps its
magnitude, or a weighted sum --- but the vector itself contains
information the scalar discards.

Why does this matter? Because moral reasoning requires knowing
\emph{which dimensions are engaged}. If a proposed remedy addresses the
wrong dimension --- offering emotional support when the problem is
financial --- it will be ineffective despite targeting the ``badness.''
The vector structure tells us where to intervene; the scalar does not.

More precisely, if the moral state is described by a point
\(x\mathcal{\in M}\) in a nine-dimensional moral manifold (Chapter 5),
then the \emph{impact} of an event is a tangent vector
\(\delta x \in T_{x}\mathcal{M}\) describing the direction and magnitude
of change. The scalar evaluation \(S(x)\) is a function on
\(\mathcal{M}\); its gradient \(\nabla S\) is a covector indicating the
direction of steepest increase in moral value. But \(\nabla S\) is
determined by \(S\) --- and if \(S\) was obtained by collapsing a richer
structure, the gradient of the collapsed function need not point in the
same direction as the gradient of any particular component.

\textbf{Example.} Let \(\mathcal{M}\) have two dimensions: wealth
(\(w\)) and health (\(h\)). Let the true moral state be the vector
\((w,h) = ( - 3,0)\) --- significant financial loss, no health impact.
Let the scalar collapse be \(S = w + h = - 3\). The gradient
\(\nabla S = (1,1)\) points equally toward improving wealth and health.
But the actual deficit lies entirely along wealth. The scalar gradient
misdirects remedial effort.

This is not a contrived example. Cost-benefit analyses that aggregate
environmental damage and economic benefit into a single dollar figure
routinely suffer from this pathology. The gradient of the aggregate
points toward ``more dollars,'' indifferent to which dimension generates
them.

\subsubsection{Limitation 2: Uncertainty Has
Shape}\label{limitation-2-uncertainty-has-shape}

The old man's ``maybe'' reflects uncertainty about the future. But his
uncertainty is not uniform across all possibilities. It has
\emph{shape}: he is more uncertain along some dimensions than others,
and --- critically --- his uncertainty is greatest along the dimensions
that are most ethically decisive.

A scalar treatment of uncertainty adds error bars: \(S = - 1 \pm 0.5\).
This says the true value lies somewhere in the interval
\(\lbrack - 1.5, - 0.5\rbrack\), but nothing about \emph{why} we are
uncertain or \emph{where} the uncertainty matters.

A geometric treatment represents uncertainty as a \emph{covariance
tensor} --- a rank-2 object that encodes both the magnitude and the
directional structure of our uncertainty:

\[\Sigma^{ij}\mathbb{= E\lbrack(}\delta m^{i})(\delta m^{j})\rbrack\]

This tells us: uncertainty is large along axis \(i\), small along axis
\(j\), and correlated between axes \(i\) and \(k\).

In the parable, the old man's uncertainty tensor might have the
following structure:

\begin{itemize}
\item
  Small variance along the \emph{current wealth} axis (the horse is
  definitely gone): \(\Sigma^{ww} \approx 0\)
\item
  Large variance along the \emph{future wealth} axis (will more horses
  come?): \(\Sigma^{w'w'} \gg 0\)
\item
  Large variance along the \emph{political} axis (will there be war?):
  \(\Sigma^{pp} \gg 0\)
\item
  Strong covariance between \emph{political} and \emph{son's welfare}
  axes (war affects conscription): \(\Sigma^{ps} \gg 0\)
\end{itemize}

The crucial insight is that \textbf{uncertainty concentrated along
ethically decisive directions matters more than uncertainty along
irrelevant directions.} If the old man were uncertain about the color of
next year's crops but certain about everything affecting his family's
survival, the first uncertainty would be ethically negligible. But if he
is uncertain precisely about war and conscription --- the dimensions
that determine whether the broken leg is a tragedy or a salvation ---
then his uncertainty is ethically maximal.

Scalar uncertainty (\(S \pm \varepsilon\)) cannot distinguish these
cases. The covariance tensor \(\Sigma\) can.

The scalar quantity that captures this alignment is the \emph{variance
of the moral evaluation} under structured uncertainty:

\[\sigma_{S}^{2} = \Sigma^{ij}\frac{\partial S}{\partial m^{i}}\frac{\partial S}{\partial m^{j}}\]

This is large when uncertainty concentrates along morally decisive
directions (where \(\nabla S\) is large), and small when uncertainty
lies along morally irrelevant directions (where \(\nabla S\) is small).
The formula is a contraction of the uncertainty tensor with the moral
gradient --- a genuinely tensorial operation that no scalar framework
can replicate.

\subsubsection{Limitation 3: Paths Cross
Boundaries}\label{limitation-3-paths-cross-boundaries}

The most profound limitation of scalar evaluation is its inability to
represent \emph{trajectory-dependent} moral change, especially
trajectories that cross \emph{regime boundaries}.

In the parable, the evaluation of ``son has a broken leg'' depends on
whether war comes. Before the declaration of war, a broken leg is
unambiguously bad: pain, disability, inability to work. After war is
declared, the evaluation bifurcates: for those without exemptions,
conscription leads to probable death; for those with exemptions
(including the son), survival is likely. The broken leg, unchanged in
itself, has crossed a moral phase boundary.

This is not merely a matter of new information changing our estimate. It
is a structural feature of the moral landscape: there exist
\emph{strata} (regimes, phases) within which smooth trade-offs apply,
separated by \emph{boundaries} where the rules change discontinuously.

A scalar function \(S\mathcal{:M \rightarrow}\mathbb{R}\), if it is
continuous, cannot represent such discontinuities. It can represent
gradual change --- \(S\) increasing or decreasing smoothly --- but not
the sharp transitions that characterize moral thresholds: consent given
vs.~withheld, life vs.~death, war vs.~peace.

To represent regime boundaries, we need \emph{stratified spaces}: spaces
composed of smooth strata (within which differential calculus applies)
joined along lower-dimensional boundaries (where discontinuities are
permitted). And to represent how moral status evolves along paths that
may cross these boundaries, we need \emph{path-dependent operations}:
parallel transport, holonomy, trajectory integrals.

\textbf{Definition.} A \emph{moral phase boundary} is a codimension-1
submanifold \(B\mathcal{\subset M}\) such that the moral evaluation
function \(S\) (or more generally, the moral tensor field) is smooth on
each side of \(B\) but may be discontinuous across \(B\).

The parable contains a clear example. Let \(p \in \lbrack 0,1\rbrack\)
denote the probability of war. As \(p\) increases continuously, the
evaluation of the son's broken leg remains continuously negative ---
until \(p\) crosses some threshold \(p^{*}\) at which conscription
becomes a near certainty. At that point, the evaluation jumps: the
broken leg becomes a life-saving exemption. The function
\(S(\text{broken leg}|p)\) is discontinuous at \(p = p^{*}\).

This is not a pathology of the parable. Empirical evidence confirms that
moral evaluation generically involves such discontinuities. Analysis of
the Dear Abby corpus (Chapter 17) reveals that qualifying phrases ---
``you promised,'' ``he hit you,'' ``only if convenient'' --- function as
\emph{semantic gates}: discrete triggers that flip moral states between
obligation and liberty, claim and no-claim. These transitions behave as
step functions, not sigmoids. The geometry of moral space is not
uniformly smooth.

\subsection{What Geometric Structure
Provides}\label{what-geometric-structure-provides}

The three limitations point toward three geometric structures beyond
scalars:

{\def\LTcaptype{none} % do not increment counter
\begin{longtable}[]{@{}
  >{\raggedright\arraybackslash}p{(\linewidth - 4\tabcolsep) * \real{0.2924}}
  >{\raggedright\arraybackslash}p{(\linewidth - 4\tabcolsep) * \real{0.2715}}
  >{\raggedright\arraybackslash}p{(\linewidth - 4\tabcolsep) * \real{0.3745}}@{}}
\toprule\noalign{}
\begin{minipage}[b]{\linewidth}\raggedright
\textbf{Limitation}
\end{minipage} & \begin{minipage}[b]{\linewidth}\raggedright
\textbf{Required Structure}
\end{minipage} & \begin{minipage}[b]{\linewidth}\raggedright
\textbf{Mathematical Object}
\end{minipage} \\
\midrule\noalign{}
\endhead
\bottomrule\noalign{}
\endlastfoot
No directional information & Vectors and covectors & \(\nabla S\),
\(O^{\mu}\), \(I_{\mu}\) \\
Uncertainty has no shape & Covariance tensor & Σ\^{}\{ij\} \\
No path-dependence & Stratification + transport & Strata, parallel
transport, holonomy \\
\end{longtable}
}

Let us examine each, now as elements of a unified framework.

\subsubsection{Gradients and the Direction of Moral
Change}\label{gradients-and-the-direction-of-moral-change}

If moral evaluation were purely scalar, there would be no meaningful
sense of ``direction'' in moral space. But our actual moral reasoning is
saturated with directional concepts: obligations \emph{point} toward
required states; interests \emph{aim} at objects; responsibility
\emph{flows} from agents to patients; improvement \emph{moves} toward
better configurations.

These are naturally modeled as vector quantities --- objects with both
magnitude and direction.

Consider an obligation. ``You ought to help your neighbor'' is not
merely a magnitude of oughtness. It specifies a \emph{direction}: from
the current state (neighbor unhelped) toward a required state (neighbor
helped). The obligation can be stronger or weaker (magnitude), but it
also has an orientation in the space of possible actions.

Formally, we represent obligations as \emph{vector fields} on the moral
manifold:

\[O^{\mu}(x\mathcal{):M \rightarrow}T\mathcal{M}\]

At each point \(x\) in moral space, \(O^{\mu}(x)\) is a tangent vector
pointing in the direction of what is required. In the nine-dimensional
framework of Chapter 5, this vector has components along each of the
nine moral dimensions derived from the 3 × 3 matrix of ethical scopes
and epistemic modes.

Interests, conversely, are represented as \emph{covector fields}:

\[I_{\mu}(x\mathcal{):M \rightarrow}T^{*}\mathcal{M}\]

A covector (or 1-form) is a linear functional on vectors. The interest
\(I_{\mu}\), applied to an obligation \(O^{\mu}\), yields a scalar: the
\emph{satisfaction} of interest \(I\) by obligation \(O\).

\[S = I_{\mu}O^{\mu}\]

This is the fundamental formula of geometric ethics (Chapter 6):
satisfaction is the contraction of obligations with interests. It is
coordinate-invariant (a scalar), but it \emph{arises from} vector
quantities that carry directional information. The scalar is not
primary; it is a projection --- a shadow cast by a richer object.

The gradient \(\nabla S\) of the satisfaction function tells us: at this
point in moral space, which direction increases satisfaction most
rapidly? This is the direction of moral improvement --- not a scalar
claim (``things could be better'') but a vector claim (``things could be
better \emph{in this specific way}'').

\subsubsection{The Metric Tensor and Moral
Comparison}\label{the-metric-tensor-and-moral-comparison}

To speak of directions, we need a way to compare them. To speak of
distances, we need a way to measure them. In differential geometry, both
functions are performed by the \emph{metric tensor} \(g_{\mu\nu}\).

The metric tensor is a rank-2 object that defines the inner product
between vectors:

\[\langle u,v\rangle = g_{\mu\nu}u^{\mu}v^{\nu}\]

This allows us to say when two directions are orthogonal (inner product
zero), aligned (inner product positive), or opposed (inner product
negative). It also defines the length of vectors and the distance
between points.

In moral space, the metric encodes \emph{how we compare values}. Two
values are orthogonal if trading off one against the other is undefined
--- there is no exchange rate between them. They are aligned if
improving one tends to improve the other. They are opposed if they
conflict.

The claim that some values are \emph{incommensurable} is, in geometric
language, the claim that the moral metric is \emph{degenerate} along
certain directions --- that there exist vectors \(v\) such that
\(g_{\mu\nu}v^{\mu}w^{\nu} = 0\) for all \(w\), or equivalently that the
metric has zero eigenvalues. This is a structural claim, not a mystical
one. It says that the geometry of moral space is not Euclidean --- not
all directions are comparable in the way that spatial directions are.

Different ethical theories correspond to different metrics on the same
underlying manifold:

\begin{itemize}
\item
  \textbf{Utilitarianism}: \(g_{\mu\nu} = \delta_{\mu\nu}\) (Euclidean
  --- all dimensions equally weighted, all trade-offs permitted)
\item
  \textbf{Rawlsian justice}: \(g_{\mu\nu}\) structured so that the
  worst-off dimension dominates (supremum metric)
\item
  \textbf{Lexicographic theories}: \(g_{\mu\nu}\) with hierarchically
  ordered eigenvalues
\item
  \textbf{Incommensurability claims}: \(g_{\mu\nu}\) with zero
  eigenvalues or degenerate subspaces
\end{itemize}

The choice of metric is doing the real moral work --- and scalar
frameworks hide this choice. The geometric framework makes it explicit
and debatable. Where the metric comes from is the subject of Chapter 9.

\subsubsection{Stratification and Moral Phase
Transitions}\label{stratification-and-moral-phase-transitions}

The parable's most profound feature is the regime change brought by war.
Before war is declared, the broken leg is bad. After war is declared, it
becomes potentially good (exemption from death). This is not a gradual
transition; it is a discontinuous jump at a boundary.

To represent such discontinuities, we need \emph{stratified spaces}:
spaces composed of smooth manifolds (strata) joined along boundaries
where the smooth structure breaks down.

\textbf{Definition.} A \emph{stratified moral space} \(\mathcal{M}\)
consists of:

\begin{enumerate}
\def\labelenumi{\arabic{enumi}.}
\item
  \textbf{Strata} \(\mathcal{M}_{\alpha}\): smooth manifolds of various
  dimensions, representing regimes within which ordinary calculus
  applies
\item
  \textbf{Incidence relations}: rules specifying how strata adjoin
  (Whitney conditions ensuring geometric regularity)
\item
  \textbf{Stratum-dependent evaluation}: moral functions that are smooth
  on each stratum but may jump across boundaries
\end{enumerate}

The boundary between ``peacetime'' and ``wartime'' is a moral stratum
boundary. On either side, smooth trade-offs apply (more wealth is
better, less pain is better). But crossing the boundary changes
\emph{which smooth trade-offs apply}. The rules are different.

This is why the old man cannot assign a stable scalar to his son's
broken leg. The evaluation depends on which stratum the world occupies,
and that is exactly what he is uncertain about.

The formal apparatus of stratification --- including absorbing strata
(nullifiers), semantic gates (discrete transition functions), and the
Whitney conditions that ensure well-behaved boundary geometry --- is
developed in Chapter 8.

\subsection{The Parable Revisited}\label{the-parable-revisited}

Let us return to the old man with the full geometric apparatus.

\textbf{Moment 1: The horse runs away.}

The moral state is \(x_{1}\mathcal{\in M}\). The impact vector lies
primarily along the wealth dimension:

\[\delta x_{1} \approx ( - 0.8, - 0.4,0,0,0,\ldots)\]

The obligation vector \(O^{\mu}\) points toward ``recover the horse or
find an alternative.'' The uncertainty tensor \(\Sigma^{ij}\) is large
along future-wealth and political axes.

A scalar evaluation says \(S(x_{1}) \approx - 1\). But this discards the
directional structure. The old man, implicitly recognizing the geometric
structure, says ``maybe.''

\textbf{Moment 2: The horse returns with others.}

The moral state is \(x_{2}\). The impact vector is strongly positive
along wealth. The obligation vector \(O^{\mu}\) now reorients toward
``maintain and steward this windfall.'' The uncertainty tensor remains
large along the political axis.

A scalar evaluation says \(S(x_{2}) \approx + 3\). The neighbors
celebrate. The old man, still tracking the full tensor, says ``maybe.''

\textbf{Moment 3: The son breaks his leg.}

The moral state is \(x_{3}\). The impact vector is negative along health
and capability. The uncertainty tensor \(\Sigma^{ij}\) becomes crucial:
there is high covariance between the \emph{political} dimension (will
there be war?) and the \emph{welfare} dimension (will the son survive?).

Crucially, the son's condition now sits near a \emph{stratum boundary}.
If war comes, the moral evaluation of the broken leg will
discontinuously shift. The gradient \(\nabla S\) is undefined at the
boundary --- it points one way in peacetime, another in wartime.

A scalar evaluation says \(S(x_{3}) \approx - 2\). The old man, sensing
the proximity to a phase transition, says ``maybe.''

\textbf{Moment 4: War comes; the son is spared.}

The moral state crosses the boundary into the wartime stratum. The
broken leg, unchanged in itself, is now on a different stratum. Its
evaluation, relative to the counterfactual (able-bodied son conscripted
and killed), is strongly positive.

The scalar is now \(S(x_{4}) \gg 0\). But this conceals the
path-dependence: the same physical state (broken leg) has different
moral valence depending on which stratum it occupies. The tensor
captures this; the scalar cannot.

\subsection{Why ``Maybe'' Is Geometric, Not Merely
Epistemic}\label{why-maybe-is-geometric-not-merely-epistemic}

The standard interpretation of the parable is epistemic: we should say
``maybe'' because we lack knowledge of the future. If only we knew
whether war was coming, we could assign definite values.

The geometric interpretation suggests something deeper:
\textbf{``maybe'' is the correct answer \emph{even with perfect
knowledge} when the evaluation structure is geometric rather than
scalar.}

Suppose the old man had an oracle who told him exactly what would
happen. Would he then assign a definite scalar to each moment?

Only if he were willing to commit to:

\begin{enumerate}
\def\labelenumi{\arabic{enumi}.}
\item
  \textbf{A fixed metric}: how much does wealth matter relative to
  health relative to family? (The choice of \(g_{\mu\nu}\).)
\item
  \textbf{A fixed treatment of path-dependence}: does the broken leg's
  value depend on the trajectory through the wartime stratum, or just
  the final state? (The choice of contraction.)
\item
  \textbf{A specific contraction}: which rank-0 summary of the
  rank-\(n\) moral tensor serves as the action guide? (The projection
  from tensor to scalar.)
\end{enumerate}

These choices are not determined by the facts. They are
\emph{perspective-dependent} --- different agents, with different
interests and different metrics, will perform different contractions and
arrive at different scalars. As Chapter 6 makes precise, the fundamental
formula \(S = I_{\mu}O^{\mu}\) already shows this: the scalar \(S\)
depends on both the obligation vector (fixed by the situation) and the
interest covector (varying with the stakeholder).

The tensor is the invariant structure. The scalar is a projection --- a
shadow --- that loses information. ``Maybe'' is what you say when you
recognize that no single scalar faithfully represents the tensor.

\subsection{From Parable to Framework}\label{from-parable-to-framework}

The parable motivates the mathematical framework developed in subsequent
chapters:

\begin{enumerate}
\def\labelenumi{\arabic{enumi}.}
\item
  \textbf{Moral space has dimension} (Chapter 5). The space of morally
  relevant configurations is not a line (totally ordered by goodness)
  but a nine-dimensional manifold, with independent axes derived from
  the intersection of three ethical scopes (individual, relational,
  systemic) and three epistemic modes (empirical, deontic, aretaic).
  These dimensions are not arbitrary; they are the minimum set required
  to distinguish situations that manifestly differ in moral character.
\item
  \textbf{Moral quantities are tensors} (Chapter 6). Obligations,
  interests, evaluations, and uncertainty are tensors of various ranks,
  carrying directional and relational information that scalars discard.
  The fundamental formula \(S = I_{\mu}O^{\mu}\) makes the passage from
  tensor to scalar explicit and auditable.
\item
  \textbf{Moral space has a metric} (Chapters 6, 9). The structure that
  allows comparison of values, measurement of moral distance, and
  identification of incommensurable dimensions is a metric tensor
  \(g_{\mu\nu}\). Different ethical theories correspond to different
  metrics; the metric is not given by the mathematics but by substantive
  moral commitment.
\item
  \textbf{Moral space is stratified} (Chapter 8). The space is not
  uniformly smooth but divided into strata --- regimes within which
  smooth trade-offs apply --- separated by boundaries where rules change
  discontinuously. Phase transitions, absorbing strata, and semantic
  gates formalize the discrete jumps that punctuate moral life.
\item
  \textbf{Moral transformation has structure} (Chapters 5, 10). What
  happens to moral evaluations when we shift perspective, permute
  agents, or translate across contexts? Tensor transformation laws and
  parallel transport answer this question with precision.
\end{enumerate}

The old man at the border, without the language of differential
geometry, had the insight: moral reality is richer than any scalar can
capture. What geometric ethics provides is the mathematical apparatus to
make that insight rigorous.

\subsection{Anticipating Objections}\label{anticipating-objections}

The skeptical reader will have objections. Let us address the two most
immediate; deeper objections receive extended treatment later.

\subsubsection{``Isn't this just saying ethics is
complicated?''}\label{isnt-this-just-saying-ethics-is-complicated}

No.~``Complicated'' suggests more of the same --- more variables, more
factors, more considerations to weigh. Geometric structure is
\emph{different in kind}. A vector is not a complicated scalar; it has
properties (direction, transformation behavior) that scalars lack
categorically. The claim is not merely that ethics has many dimensions
(though it does) but that moral quantities \emph{transform in specific
ways} under change of perspective, and this transformation behavior is
precisely what geometric objects --- tensors, connections, fiber bundles
--- are designed to capture.

An analogy: a photograph of a landscape is not a ``complicated pixel.''
It is a two-dimensional array with structure --- spatial relationships,
color gradients, edges --- that no single number can represent. The move
from scalar to geometric ethics is like the move from a pixel to a
photograph: not more of the same, but a different kind of
representation.

\subsubsection{``We can't measure moral quantities precisely, so what
use is this
formalism?''}\label{we-cant-measure-moral-quantities-precisely-so-what-use-is-this-formalism}

The same objection was raised against utility theory, and against the
use of calculus in economics. The response is twofold.

First, the framework's value lies in \emph{structural} insights, not
numerical precision. Knowing that two values are orthogonal
(incommensurable under the current metric) is useful even if we cannot
assign magnitudes to either. Knowing that uncertainty concentrates along
the health dimension rather than the wealth dimension changes how we
deliberate, even without exact numbers. The qualitative geometry
constrains reasoning before any quantity is measured.

Second, the framework identifies \emph{what would need to be measured}
to make ethical reasoning precise. Physics progressed from qualitative
insights (``force causes acceleration'') to quantitative laws
(\(F = ma\)) as measurement improved. Chapter 17 presents evidence that
geometric structures predicted by the framework ---
multi-dimensionality, correlative symmetry, context-dependent weighting,
temporal stability --- are empirically detectable in natural moral
reasoning, using the Dear Abby corpus as ground truth. The measurements
are coarse, but they exist, and the framework tells us how to refine
them.

\subsection{Conclusion}\label{conclusion}

The old man's ``maybe'' is not resignation. It is recognition.

A scalar \(S(x)\) can label the present state, but it cannot represent
which directions are ethically decisive, whether uncertainty
concentrates along those directions, or how evaluation changes along
trajectories that cross stratification boundaries.

These three failures are not peripheral. They are \emph{structural} ---
consequences of the fact that a rank-0 object (a scalar) cannot
faithfully represent a rank-\(n\) structure (a tensor). The information
lost in the projection from tensor to scalar is precisely the
information that matters most in hard moral cases: which values are at
stake, where the uncertainty bites, and whether we are near a phase
transition where the rules change.

The geometric structures that resolve these failures are:

\begin{itemize}
\item
  \textbf{Vectors and covectors} --- capturing the direction of
  obligation and interest
\item
  \textbf{The covariance tensor} --- encoding uncertainty in morally
  relevant coordinates
\item
  \textbf{Stratified spaces} --- representing the discontinuous
  boundaries between moral regimes
\item
  \textbf{Parallel transport} --- tracking how moral evaluations evolve
  along paths that may cross boundaries
\end{itemize}

Without these structures, ``maybe'' must be treated as an ad hoc
heuristic --- a confession of ignorance --- rather than emerging from
the model as a structural feature of moral reality.

The parable, read geometrically, is not about the limits of human
knowledge. It is about the geometry of ethical evaluation. The old man
sees that the tensor cannot be contracted to a scalar without loss. His
``maybe'' is a holding operation --- a refusal to project --- until the
full structure of the situation is revealed.

The chapters that follow develop the apparatus: the historical
precedents (Chapter 3), the mathematical preliminaries (Chapter 4), the
moral manifold (Chapter 5), the tensor hierarchy (Chapter 6), and the
pedagogical centerpiece (Chapter 7). But the core insight is here, in
the parable: \emph{ethics is not a number.} It is a geometric structure.
And the first step in understanding that structure is recognizing what
scalars cannot do.

❖

\emph{The horse runs away. Good? Bad? Maybe.}

\emph{The geometry agrees: the tensor is not yet fully contracted. Hold
the projection. Watch the structure unfold.}

\section{Chapter 3: Historical Precursors --- Geometry Before
Geometry}\label{chapter-3-historical-precursors-geometry-before-geometry}

\begin{quote}
\textbf{\textsc{RUNNING EXAMPLE --- Priya's Model}}

\emph{Priya starts reading ethics, surprised to find her problem already
described---in different language---by people who never imagined
computers. Ross's prima facie duties map directly: she has a duty of
beneficence (match patients to life-saving trials), a duty of justice
(match fairly), and a duty of fidelity (to her employer's commission).
These duties point in different directions---they are vectors, not
scalars. She underlines a passage in Sen: a person's advantage is not a
matter of resources but of what they can do and be. Sen was describing
capability as a multi-dimensional space. Priya's TrialMatch score does
exactly what Sen warned against: it collapses capabilities into a rank
ordering.}
\end{quote}

\subsection{Introduction: Structure Before
Formalism}\label{introduction-structure-before-formalism}

The mathematical apparatus of differential geometry was developed in the
nineteenth and early twentieth centuries, reaching its canonical form in
Einstein's general relativity and the gauge theories of Yang--Mills.
Moral philosophy predates this development by millennia. Yet the
\emph{structural insights} that geometric ethics formalizes ---
transformation behavior, multi-dimensional interdependence, coordinate
invariance, the distinction between intrinsic and perspectival
properties --- have appeared throughout the history of ethics in various
guises.

This chapter traces a genealogy of proto-geometric thinking in moral
philosophy. The claim is not that Aristotle or Kant secretly knew
differential geometry, but that they grappled with phenomena that resist
scalar treatment and developed conceptual tools that, in retrospect,
capture aspects of geometric structure. Reading these thinkers through a
geometric lens both illuminates their insights and shows that the
framework developed in this book has deep roots in the philosophical
tradition.

We proceed roughly chronologically, though the ordering also reflects
increasing mathematical sophistication in the proto-geometric concepts.

\subsection{Aristotle: The Doctrine of the Mean as Context-Sensitive
Calibration}\label{aristotle-the-doctrine-of-the-mean-as-context-sensitive-calibration}

Aristotle's \emph{Nicomachean Ethics} presents virtue as a \emph{mean}
(μεσότης) between extremes of excess and deficiency. Courage lies
between recklessness and cowardice; generosity between prodigality and
miserliness; proper pride between vanity and undue humility. The
virtuous person hits the mean ``at the right times, with reference to
the right objects, towards the right people, with the right motive, and
in the right way'' (1106b21).

This is emphatically not a scalar doctrine. Aristotle explicitly rejects
the idea that virtue is a single quantity to be maximized:

\begin{quote}
``It is no easy task to find the middle\ldots{} anyone can get angry ---
that is easy --- or give or spend money; but to do this to the right
person, to the right extent, at the right time, with the right motive,
and in the right way, that is not for everyone, nor is it easy.''
(1109a26)
\end{quote}

The mean is not a fixed point on a line but a \emph{context-dependent
calibration} across multiple dimensions. What counts as courage depends
on the situation (battlefield vs.~sickroom), the agent's role (soldier
vs.~physician), the stakes involved, and the alternatives available. The
mean for one person in one situation may be quite different from the
mean for another person in a different situation.

\subsubsection{Geometric Reading}\label{geometric-reading}

In geometric terms, Aristotle's mean is a \emph{section} of a fiber
bundle over the space of situations. At each point \(x\) in
situation-space \(\mathcal{S}\), there is a fiber \(R_{x}\) of possible
responses, and the virtuous response is determined by the local
structure of the situation --- not by a global, context-free rule.

More precisely, let \(\mathcal{S}\) be the space of ethically relevant
situations and let \(R\) be the space of possible responses. A
\emph{character trait} is a section \(\sigma\mathcal{:S \rightarrow}E\)
of a fiber bundle \(\pi:E\mathcal{\rightarrow S}\) with fiber \(R\),
assigning a response to each situation. The virtuous character trait
\(\sigma^{*}\) is the one that, at each point, selects the response
appropriate to that situation's specific configuration.

The ``right time, right object, right person, right motive, right way''
are \emph{coordinates} on \(\mathcal{S}\). Virtue requires sensitivity
to all of them. A scalar theory would say: ``maximize courage'' or
``minimize cowardice.'' Aristotle says: the courageous response is a
\emph{function of the local coordinates}, not a global maximum.

This is proto-geometric because it recognizes that ethical evaluation
is:

\begin{enumerate}
\def\labelenumi{\arabic{enumi}.}
\item
  \textbf{Multi-dimensional} (multiple ``right \(X\)'' conditions)
\item
  \textbf{Context-dependent} (the mean varies with situation)
\item
  \textbf{Not reducible to optimization of a single quantity}
\end{enumerate}

What Aristotle lacks is the mathematical apparatus to describe how the
mean \emph{transforms} as we change coordinates --- how the courageous
response in one framing relates to the courageous response in another
framing of the same situation. Differential geometry provides exactly
this.

\subsubsection{The Doctrine of the Mean as a Metric
Condition}\label{the-doctrine-of-the-mean-as-a-metric-condition}

There is a deeper geometric reading. The mean is defined relative to
\emph{us} --- not the arithmetic mean of the extremes, but the mean
``relative to us'' (πρὸς ἡμᾶς). This suggests that the moral space has a
\emph{metric structure that varies with the agent}.

If we represent character traits as vectors in a space of dispositions,
then ``excess'' and ``deficiency'' are directions away from the virtuous
center. But what counts as excess depends on the metric: a step that is
``too far'' for one agent may be ``not far enough'' for another, because
their metrics differ.

Formally, let \(g_{\mu\nu}(a)\) be the metric tensor on
disposition-space, parameterized by agent \(a\). The mean for agent
\(a\) is the point equidistant (under \(g(a)\)) from the extremes.
Different agents, with different metrics, locate the mean at different
points.

This reading explains Aristotle's insistence that virtue cannot be
taught by rule. Rules are coordinate-dependent; the mean is
metric-dependent. Without knowing the agent's metric --- their
capacities, circumstances, history --- one cannot specify the mean in
advance.

In the language of Chapter 5, the agent-dependence of the mean
corresponds to the fiber bundle structure of the moral manifold: the
agent's perspective lives in a fiber over the base situation, and the
metric may vary along the fiber direction.

\subsection{Kant: The Categorical Imperative as an Invariance
Condition}\label{kant-the-categorical-imperative-as-an-invariance-condition}

Kant's moral philosophy appears, at first glance, maximally
anti-geometric. The categorical imperative demands \emph{universal}
laws, applicable to all rational beings regardless of circumstance.
``Act only according to that maxim whereby you can at the same time will
that it should become a universal law'' (Groundwork, 4:421). What could
be more scalar than a single test applied uniformly to all actions?

But look again. The categorical imperative is not a command to maximize
a quantity. It is a \emph{constraint on the form} of permissible maxims:
only those maxims that can be universalized without contradiction are
morally permissible.

\subsubsection{Geometric Reading}\label{geometric-reading-1}

In geometric terms, the categorical imperative is an \emph{invariance
condition}. It asks: which maxims remain valid under a specific
transformation --- the transformation from ``I, in my particular
circumstances'' to ``any rational being in relevantly similar
circumstances''?

Let \(T\) be the transformation that generalizes a maxim from
first-personal to universal form. A maxim \(m\) is permissible if and
only if:

\[T(m) = m\]

That is, the maxim is a \emph{fixed point} of the universalization
transformation. Maxims that change under \(T\) --- that work for me but
fail when universalized --- are impermissible.

This is structurally identical to how physicists identify \emph{scalars}
(quantities invariant under coordinate transformations) and more
generally \emph{tensors} (quantities that transform in specific lawful
ways). Kant is asking: which moral claims are \emph{invariant} under the
transformation from particular to universal perspective?

The parallel is not superficial. In physics, the laws of nature must be
the same in all reference frames --- this is the principle of general
covariance. In Kantian ethics, the laws of morality must be the same for
all rational agents --- this is the categorical imperative. Both are
invariance conditions that constrain the form of legitimate laws.

This is also the structural ancestor of what Chapter 5 calls the
\emph{Bond Invariance Principle} (BIP): moral evaluations must be
invariant under admissible re-descriptions of the situation. Kant
discovered the principle; the geometric framework provides the
transformation group.

\subsubsection{The Kingdom of Ends as a Symmetry
Group}\label{the-kingdom-of-ends-as-a-symmetry-group}

Kant's ``kingdom of ends'' deepens the geometric reading. In the kingdom
of ends, every rational being is both legislator (author of universal
laws) and subject (bound by those laws). The moral community is defined
by the \emph{symmetry} between these roles.

In mathematical terms, the kingdom of ends is closed under the
transformation that swaps legislator and subject. If a law \(L\) is
valid, then the transformed law \(T(L)\) --- where agent and patient are
exchanged --- must also be valid. This is a \emph{symmetry condition} on
the structure of moral laws.

Symmetry conditions of this form are the hallmark of geometric
equations. Maxwell's equations are invariant under Lorentz
transformations; Einstein's field equations are invariant under general
coordinate transformations. Kant's moral laws are invariant under the
permutation of rational agents.

What Kant identifies, without the mathematical language, is that
\emph{moral objectivity is transformation invariance}. A moral claim is
objective not because it corresponds to some moral fact ``out there,''
but because it remains valid under all admissible transformations of
perspective. This is the same structural insight that animates gauge
theory in physics --- and that the Bond Invariance Principle makes
explicit in the ethical domain.

\subsection{Ross: Prima Facie Duties as Vector
Components}\label{ross-prima-facie-duties-as-vector-components}

W.D. Ross's \emph{The Right and the Good} (1930) introduced the concept
of \emph{prima facie duties}: duties that are binding unless overridden
by stronger duties. We have prima facie duties of fidelity (keeping
promises), reparation (making amends), gratitude, justice, beneficence,
self-improvement, and non-maleficence.

These duties can conflict. A promise to meet a friend may conflict with
an opportunity to prevent harm to a stranger. When they conflict, we
must judge which duty is stronger \emph{in this particular situation}
--- a judgment Ross calls the determination of our \emph{actual duty}.

\subsubsection{Geometric Reading}\label{geometric-reading-2}

Ross's prima facie duties are \emph{components of a moral vector}. Each
duty type corresponds to a dimension of the obligation space:

\[\mathbf{O} = (O_{\text{fidelity}},O_{\text{reparation}},O_{\text{gratitude}},O_{\text{justice}},O_{\text{beneficence}},O_{\text{improvement}},O_{\text{non-maleficence}})\]

In any given situation, each component has some magnitude (possibly
zero). The \emph{actual duty} is some function of these components ---
but, crucially, not a simple sum.

Ross explicitly rejects the utilitarian move of reducing all duties to a
single scalar (utility). He also rejects the idea that there is a fixed
\emph{priority ordering} among duty types. Instead, the determination of
actual duty is a matter of \emph{judgment} that weighs the components
contextually.

In geometric terms, Ross is grappling with the problem of
\emph{contraction}: how do we go from a multi-component vector to a
single action-guiding prescription? His answer --- that there is no
mechanical rule, only trained judgment --- reflects the fact that
different situations call for different contraction operations. This is
precisely the structure developed in Chapter 6: the scalar
\(S = I_{\mu}O^{\mu}\) depends on the interest covector \(I_{\mu}\),
which varies with the agent and context.

\subsubsection{The Interaction Problem}\label{the-interaction-problem}

Ross's framework faces a difficulty that has occupied his commentators
for nearly a century: how do prima facie duties \emph{interact}? If I
have a strong duty of fidelity and a weak duty of beneficence, does the
fidelity duty simply win? Or do they combine in some more complex way?

The geometric framework suggests an answer. Duties are not merely
magnitudes to be compared; they have \emph{directions} in moral space.
Two duties may be:

\begin{itemize}
\item
  \textbf{Aligned}: both point in the same direction (keeping a promise
  that also helps someone)
\item
  \textbf{Orthogonal}: independent, neither reinforcing nor conflicting
  (a promise to one person, a beneficence opportunity involving another)
\item
  \textbf{Opposed}: pointing in opposite directions (a promise that
  requires harming someone)
\end{itemize}

The combination of duties is then a \emph{vector sum}, with the geometry
determining how they add:

\[\mathbf{O}_{\text{actual}} = \sum_{i}^{}\mathbf{O}_{i}\]

When duties are aligned, their magnitudes add. When orthogonal, they
combine by the Pythagorean theorem (their resultant is
\(\sqrt{O_{1}^{2} + O_{2}^{2}}\), not \(O_{1} + O_{2}\)). When opposed,
they partially cancel.

This explains a phenomenon that has puzzled Rossian commentators: why
strong orthogonal duties can coexist without conflict (I can keep my
promise \emph{and} help the stranger, if the actions lie in independent
dimensions), while even weak opposed duties create tension (breaking
even a minor promise to prevent trivial harm still feels like a genuine
moral loss --- the cancelled component leaves what Chapter 15 will call
\emph{moral residue}).

Ross's ``judgment'' can now be understood as sensitivity to the
\emph{geometry} of the duty configuration --- something that resists
codification in scalar terms but has definite mathematical structure.

\subsection{Rawls: The Original Position as a Symmetry
Constraint}\label{rawls-the-original-position-as-a-symmetry-constraint}

John Rawls's \emph{A Theory of Justice} (1971) proposes that principles
of justice are those that would be chosen by rational agents behind a
``veil of ignorance'' --- not knowing their place in society, their
natural talents, or their conception of the good. The original position
is a thought experiment designed to identify principles that are fair
because they are chosen without knowledge of how they will affect the
chooser.

\subsubsection{Geometric Reading}\label{geometric-reading-3}

The original position is a \emph{symmetry condition}. By removing
knowledge of particular position, it forces the choice of principles
that are \emph{invariant under permutation of agents}. If a principle
benefits position A at the expense of position B, it cannot be chosen
behind the veil, because the chooser might turn out to occupy position
B.

Formally, let \(\pi\) be a permutation of social positions. A principle
\(P\) is admissible in the original position if and only if:

\[\pi(P) = P\quad\text{for all permutations }\pi\]

This is precisely the condition for a \emph{symmetric tensor} --- a
tensor unchanged under index permutation.

Rawls's two principles of justice can be understood as the \emph{unique
symmetric solution} (up to specification) to the problem of social
cooperation. The first principle (equal basic liberties) is symmetric by
construction: everyone gets the same liberties. The second principle
(difference principle) permits inequalities only if they benefit the
worst-off position --- a symmetric condition because the worst-off
position is defined relative to the structure of positions, not to any
particular occupant.

\subsubsection{The Metric of the Original
Position}\label{the-metric-of-the-original-position}

The original position also implicitly specifies a \emph{metric} on the
space of social positions. The difference principle uses a
\emph{maximin} criterion: maximize the minimum position. This is
equivalent to choosing a specific metric in which distance is measured
by the worst-off coordinate.

In geometric terms, the Rawlsian metric is the supremum metric:

\[d(x,y) = \max_{i}|x_{i} - y_{i}|\]

This gives decisive weight to the worst-off dimension. Alternative
metrics yield different principles:

\begin{itemize}
\item
  The \emph{utilitarian metric} (\(L^{1}\)):
  \(d(x,y) = \sum_{i}^{}|x_{i} - y_{i}|\), weighting all positions
  equally
\item
  The \emph{Euclidean metric} (\(L^{2}\)):
  \(d(x,y) = \sqrt{\sum_{i}^{}(x_{i} - y_{i})^{2}}\), weighting by
  squared deviations
\item
  The \emph{prioritarian metric}: a weighted \(L^{p}\) metric giving
  greater weight to worse-off positions, interpolating between
  utilitarian and Rawlsian
\end{itemize}

Rawls's argument against utilitarianism can be read as an argument about
\emph{which metric} is appropriate for justice. The utilitarian metric
permits sacrificing some positions for aggregate gain; the Rawlsian
metric does not. This is a substantive geometric claim about the
structure of fair social evaluation.

The geometric framework (Chapters 6 and 9) makes this explicit:
different ethical theories correspond to different metric tensors on the
moral manifold. The metric is not given by the mathematics; it encodes a
substantive commitment. What geometry provides is the vocabulary for
stating the commitment precisely --- and the tools for deriving its
consequences.

\subsection{Sen and Nussbaum: Capabilities as a Basis for Moral
Space}\label{sen-and-nussbaum-capabilities-as-a-basis-for-moral-space}

Amartya Sen and Martha Nussbaum developed the \emph{capabilities
approach} as an alternative to both utilitarian welfare measures and
Rawlsian primary goods. The core idea is that what matters morally is
not subjective well-being (utility) or objective resources (income,
rights) but \emph{capabilities}: the real freedoms people have to
achieve ``functionings'' they have reason to value.

Sen identifies a plurality of capabilities --- life, bodily health,
bodily integrity, senses/imagination/thought, emotions, practical
reason, affiliation, relation to other species, play, control over one's
environment --- that are \emph{irreducibly plural}. They cannot be
reduced to a single scalar measure without violence to their distinct
characters.

\subsubsection{Geometric Reading}\label{geometric-reading-4}

The capabilities are \emph{basis vectors} for a moral space. Each
capability defines an independent dimension along which a person's life
can go well or badly. A person's situation is a \emph{vector} in
capability space:

\[\mathbf{c} = (c_{\text{life}},c_{\text{health}},c_{\text{integrity}},c_{\text{senses}},c_{\text{emotions}},c_{\text{reason}},c_{\text{affiliation}},c_{\text{nature}},c_{\text{play}},c_{\text{control}})\]

This is explicitly multi-dimensional. Sen insists that capabilities
cannot be aggregated into a single index without loss of essential
information --- precisely the claim that moral evaluation requires
richer structure than a scalar.

The connection to the nine-dimensional moral manifold of Chapter 5 is
suggestive. The nine dimensions derived from the 3 × 3 matrix of ethical
scopes and epistemic modes do not map one-to-one onto Sen's
capabilities, but the structural move is the same: identifying the
\emph{basis} of moral space by systematic analysis of what dimensions
are needed to distinguish morally distinct situations. Sen proceeds by
philosophical reflection on human flourishing; the geometric framework
proceeds by combinatorial analysis of ethical scope and epistemic mode.
That both arrive at a finite-dimensional space of roughly similar
cardinality (∼10) is at least noteworthy.

\subsubsection{The Incompleteness
Thesis}\label{the-incompleteness-thesis}

Sen argues that comparative judgments of capability sets are
\emph{incomplete}: we can often say that one situation is better than
another along some dimensions and worse along others, without being able
to say which is better overall. This incompleteness is not a failure of
the theory but a feature of moral reality.

In geometric terms, this is the claim that the moral metric is
\emph{degenerate} or \emph{partial}. Not all vectors can be compared in
length. Given two capability vectors \(\mathbf{c}_{1}\) and
\(\mathbf{c}_{2}\), we may have:

\begin{itemize}
\item
  \(c_{1}^{i} > c_{2}^{i}\) along some dimensions
\item
  \(c_{1}^{j} < c_{2}^{j}\) along other dimensions
\item
  No legitimate contraction to an overall comparison
\end{itemize}

This is the geometric signature of \emph{incommensurability}. A scalar
theory would force a comparison (by summing or by lexical priority);
Sen's theory preserves the genuine incompleteness. The geometric
framework represents this as a structural feature of the moral manifold
--- a degenerate metric --- rather than a failure of analysis.

\subsubsection{Nussbaum's Threshold and Stratum
Structure}\label{nussbaums-threshold-and-stratum-structure}

Nussbaum modifies the capabilities approach by introducing
\emph{thresholds}: minimum levels of each capability below which a life
is not fully human. This introduces \emph{stratum structure} into
capability space.

Below the threshold, we are in a different moral regime --- one where
the imperative is to raise capabilities to the threshold level,
admitting no trade-offs. Above the threshold, trade-offs and choices
become permissible. The threshold is a \emph{stratum boundary}
separating regions with different moral rules.

This is proto-stratified geometry. Nussbaum's capability space is not a
smooth manifold but a stratified space with distinguished hypersurfaces
(the thresholds) where the rules change discontinuously. In the
vocabulary of Chapter 8, Nussbaum's thresholds are codimension-1 stratum
boundaries; the regime below the threshold (where the only permissible
direction is ``up'') is an absorbing stratum with restricted tangent
directions.

The structural parallel to the empirical findings described in Chapter
17 is striking: analysis of the Dear Abby corpus reveals that qualifying
conditions --- abuse, danger, incapacity --- function as
\emph{nullifiers} that override all other moral considerations. These
are absorbing strata in exactly Nussbaum's sense: once you cross the
threshold (a child is being harmed, a spouse is abusive), the moral
calculus collapses to a single imperative, admitting no trade-offs.

\subsection{Hohfeld: Jural Relations as Discrete
Structure}\label{hohfeld-jural-relations-as-discrete-structure}

Wesley Newcomb Hohfeld's \emph{Fundamental Legal Conceptions} (1913,
1917) introduced a taxonomy of legal relations that has become
foundational in jurisprudence. Hohfeld identified eight fundamental
jural positions organized into two squares of opposition:

{\def\LTcaptype{none} % do not increment counter
\begin{longtable}[]{@{}
  >{\raggedright\arraybackslash}p{(\linewidth - 2\tabcolsep) * \real{0.1622}}
  >{\raggedright\arraybackslash}p{(\linewidth - 2\tabcolsep) * \real{0.1553}}@{}}
\toprule\noalign{}
\begin{minipage}[b]{\linewidth}\raggedright
\textbf{Right-holder}
\end{minipage} & \begin{minipage}[b]{\linewidth}\raggedright
\textbf{Duty-holder}
\end{minipage} \\
\midrule\noalign{}
\endhead
\bottomrule\noalign{}
\endlastfoot
Claim → & Duty \\
Liberty → & No-claim \\
Power → & Liability \\
Immunity → & Disability \\
\end{longtable}
}

Within each square, the relations are connected

by \emph{correlatives}

\textbf{(Claim↔Duty, Liberty↔No-claim)}

and \emph{opposites}

\textbf{(Claim↔No-claim, Duty↔Liberty).}

\subsubsection{Geometric Reading}\label{geometric-reading-5}

Hohfeld's taxonomy is remarkable because it is entirely \emph{discrete}
--- there are no continuous parameters, no ``degrees'' of claim or
liberty. The positions are qualitative states, and the transitions
between them are all-or-nothing.

In geometric terms, the Hohfeldian positions define the vertices of a
discrete symmetry structure. The correlative and opposition relations
generate a group of transformations on these vertices --- and that group
turns out to be \(D_{4}\), the dihedral group of order 8, which is also
the symmetry group of the square.

This identification is not metaphorical. The eight Hohfeldian positions
map onto the eight elements of \(D_{4}\) (four rotations, four
reflections), and the correlative and opposition operations correspond
to specific group elements. The transitions between positions ---
obligation flipping to liberty, claim dissolving to no-claim --- are
group actions.

The significance for geometric ethics is twofold:

\begin{enumerate}
\def\labelenumi{\arabic{enumi}.}
\item
  \textbf{Discrete symmetry in moral space.} The Hohfeldian relations
  show that moral space contains irreducibly discrete structure
  alongside its continuous dimensions. The moral manifold is not purely
  smooth; it has a discrete component that no amount of refinement can
  smooth away. This is the mathematical content of stratification
  (Chapter 8).
\item
  \textbf{Empirical confirmation.} Analysis of natural moral reasoning
  (Chapter 17) reveals that these transitions behave as step functions:
  a qualifying phrase (``you promised'') flips the moral state from
  liberty to obligation with no intermediate values. The correlative
  symmetry --- Claim↔Duty at 87\%, Liberty↔No-claim at 82\% --- is a
  structural invariant of the moral manifold, detectable in data.
\end{enumerate}

Hohfeld, working in the formalist tradition of legal analysis,
discovered discrete geometric structure in the moral domain nearly a
century before the mathematical framework existed to describe it.

\subsection{Moral Uncertainty: Superposition and Structured
Credence}\label{moral-uncertainty-superposition-and-structured-credence}

Recent work in moral philosophy has focused on \emph{moral uncertainty}:
what should we do when we are uncertain which moral theory is correct?
If I am 60\% confident in utilitarianism and 40\% confident in
deontology, and they recommend different actions, what should I choose?

Various approaches have been proposed: ``my favorite theory'' (act on
whichever theory you find most plausible), ``maximize expected moral
value'' (weight each theory's recommendation by credence), and more
sophisticated methods that account for intertheoretic comparisons.

\subsubsection{Geometric Reading}\label{geometric-reading-6}

The structure of moral uncertainty is naturally \emph{vectorial}. An
agent under moral uncertainty is not in a definite moral state but in a
\emph{weighted combination} of moral states.

Let \(|U\rangle\) represent the evaluation under utilitarianism and
\(|D\rangle\) the evaluation under deontology. An agent with 60\%
credence in utilitarianism is in the state:

\[|\psi\rangle = \sqrt{0.6}\,|U\rangle + \sqrt{0.4}\,|D\rangle\]

This is a vector in a \emph{theory space}, not a scalar. The agent's
moral situation cannot be captured by a single number (overall credence)
but requires specification of the full vector --- its components, their
relative phase, and their transformation behavior under changes of
basis.

When the agent acts, the vector must be contracted to a scalar choice
--- but the contraction is not unique. Expected value maximization is
one contraction; other approaches perform different contractions. The
choice of contraction is itself a moral commitment, distinct from the
choice of theory.

The analogy with quantum superposition is structural, not merely
metaphorical. In quantum mechanics, a state vector \(|\psi\rangle\)
encodes the full information about a system; measurement collapses it to
a definite outcome. In moral uncertainty, the credence vector encodes
the full structure of the agent's evaluative commitments; action
collapses it to a definite choice. The mathematics of superposition,
projection, and measurement carries over to the moral domain --- a
connection developed formally in Chapter 13 (Quantum Normative
Dynamics).

\subsubsection{Moral Hedging as Covariance
Sensitivity}\label{moral-hedging-as-covariance-sensitivity}

The sophisticated treatment of moral uncertainty involves
\emph{hedging}: choosing actions that are reasonably good under multiple
theories, even if not optimal under any single one. This is analogous to
portfolio diversification in finance.

In geometric terms, hedging is sensitivity to the \emph{covariance
structure} of moral uncertainty. Let \(\Sigma\) be the covariance matrix
of the agent's moral beliefs. Let \(\Delta_{a} = (a_{U} - a_{D})\) be
the vector of differences between what each theory recommends for action
\(a\). The ``risk'' of action \(a\) is:

\[\sigma_{a}^{2} = \Delta_{a}^{T}\Sigma\,\Delta_{a}\]

Actions with low \(\sigma^{2}\) are robust to moral uncertainty; actions
with high \(\sigma^{2}\) are risky bets on particular theories. This is
the moral analogue of portfolio variance --- and it requires the full
tensor structure of uncertainty, not just scalar credences.

\subsection{Synthesis: What the Precursors
Share}\label{synthesis-what-the-precursors-share}

Across two and a half millennia, these thinkers share a recognition that
moral reality resists scalar reduction:

{\def\LTcaptype{none} % do not increment counter
\begin{longtable}[]{@{}
  >{\raggedright\arraybackslash}p{(\linewidth - 4\tabcolsep) * \real{0.1920}}
  >{\raggedright\arraybackslash}p{(\linewidth - 4\tabcolsep) * \real{0.3255}}
  >{\raggedright\arraybackslash}p{(\linewidth - 4\tabcolsep) * \real{0.4826}}@{}}
\toprule\noalign{}
\begin{minipage}[b]{\linewidth}\raggedright
\textbf{Thinker}
\end{minipage} & \begin{minipage}[b]{\linewidth}\raggedright
\textbf{Insight}
\end{minipage} & \begin{minipage}[b]{\linewidth}\raggedright
\textbf{Geometric Structure}
\end{minipage} \\
\midrule\noalign{}
\endhead
\bottomrule\noalign{}
\endlastfoot
Aristotle & The mean is context-dependent & Virtue as section of fiber
bundle; agent-relative metric \\
Kant & Morality requires universalizability & Permissible maxims as
transformation invariants \\
Ross & Duties are plural and interact & Duties as vector components;
combination as vector sum \\
Rawls & Justice requires position-independence & Principles as symmetric
tensors; metric choice determines theory \\
Sen/Nussbaum & Capabilities are irreducibly plural & Capabilities as
basis vectors; thresholds as stratum boundaries \\
Hohfeld & Jural relations are discrete & Positions as \(D_{4}\) group
elements; transitions as group actions \\
Moral uncertainty & Credences have structure & Beliefs as state vectors;
hedging from covariance \\
\end{longtable}
}

None of these thinkers used the language of differential geometry. But
they all developed conceptual tools to handle phenomena that geometry
formalizes:

\begin{enumerate}
\def\labelenumi{\arabic{enumi}.}
\item
  \textbf{Multi-dimensionality}: Moral evaluation involves multiple
  independent considerations that cannot be reduced to one.
\item
  \textbf{Transformation behavior}: What happens to moral claims when we
  shift perspective, permute agents, or change framing?
\item
  \textbf{Metric structure}: How do we compare values, measure moral
  distance, identify incommensurability?
\item
  \textbf{Context-dependence}: The same abstract principle yields
  different concrete prescriptions in different situations.
\item
  \textbf{Structured combination}: When moral considerations combine,
  they do so geometrically (with alignment, orthogonality, opposition),
  not arithmetically.
\item
  \textbf{Discrete transitions}: Some moral changes are discontinuous
  --- phase transitions, not gradual slides.
\end{enumerate}

\subsection{What the Geometric Framework
Adds}\label{what-the-geometric-framework-adds}

If these insights are already present in the tradition, what does the
geometric framework add?

\textbf{Precision.} The tradition offers metaphors and intuitions; the
framework offers definitions and theorems. ``Duties interact'' becomes
``obligations combine as vectors under the moral metric
\(g_{\mu\nu}\).'' ``The mean is relative to us'' becomes ``virtue is
determined by an agent-parameterized metric on disposition-space.''
Precision enables analysis, criticism, and extension.

\textbf{Unification.} The tradition offers disparate insights from
apparently incompatible systems. The geometric framework reveals common
structure beneath surface disagreement. Aristotle's context-sensitivity
and Kant's universalizability are not opposed but complementary: both
constrain the transformation behavior of moral claims, Aristotle at the
level of the fiber (agent-dependence), Kant at the level of the base
(situation-independence). Ross's plurality and Rawls's symmetry are both
structural features of the moral tensor --- Ross specifying its
components, Rawls specifying its symmetry group.

\textbf{New questions.} The framework suggests questions the tradition
did not ask:

\begin{itemize}
\item
  What is the \emph{signature} of the moral metric --- is it
  positive-definite (Euclidean), indefinite (Lorentzian), or degenerate?
\item
  What are the \emph{symmetries} of moral space --- which
  transformations leave the structure invariant, and what conservation
  laws do they imply?
\item
  What \emph{curvature} does moral space have --- how does parallel
  transport of an obligation around a closed loop change the obligation
  vector?
\item
  What is the \emph{holonomy group} of moral connections --- which
  transformations arise from cycling through moral perspectives?
\end{itemize}

These are not idle questions. Chapter 10 develops the curvature of moral
space and shows that non-trivial holonomy corresponds to path-dependent
moral evaluation --- the geometric content of the familiar intuition
that ``where you start affects where you end up'' in moral reasoning.
Chapter 12 shows that the symmetry of re-description invariance implies
a conservation law for harm --- the geometric analogue of Noether's
theorem.

\textbf{Computability.} The tradition offers wisdom; the framework
offers implementation. Tensors can be represented in computers;
geometric operations can be implemented in algorithms; geometric
equations can be solved numerically. A philosophical framework, however
profound, is inert if it cannot be computed. The geometric framework
bridges the gap between moral insight and moral engineering --- not by
replacing judgment with calculation, but by making the structure on
which judgment operates explicit and manipulable.

\subsection{Conclusion: The Geometry Was Always
There}\label{conclusion-the-geometry-was-always-there}

The history of moral philosophy is, in significant part, a struggle
against the reductionism of scalar ethics. Each thinker we have examined
recognized that moral reality has structure that scalars cannot capture,
and developed conceptual tools to articulate that structure.

Geometric ethics is not a break from this tradition but its continuation
--- and, in a mathematical sense, its formalization. The conceptual
tools developed by Aristotle, Kant, Ross, Rawls, Sen, Nussbaum, Hohfeld,
and theorists of moral uncertainty find their natural mathematical
expression in the language of manifolds, tensors, metrics, connections,
and stratified spaces.

This is not to say that the tradition was secretly doing mathematics. It
is to say that mathematics, properly understood, is the science of
structure --- and the structures that matter for ethics are precisely
those that differential geometry was developed to describe:
transformation behavior, multi-linear combination, coordinate
invariance, metric geometry, and the interplay of continuous and
discrete structure.

The tradition provides the insights. The framework provides the
language. Together, they enable a moral philosophy that is both faithful
to the complexity of ethical life and precise enough to be stated,
tested, and refined.

In the chapters that follow, we develop the framework in full. But we do
so in the company of these predecessors, whose insights we are
formalizing, not replacing.

❖

\emph{Aristotle sought the mean.}

\emph{Kant sought the universal.}

\emph{Ross sought the balance.}

\emph{Rawls sought the fair.}

\emph{Sen sought the capable.}

\emph{Hohfeld sought the discrete.}

\emph{They were all, in their different ways, seeking the geometry.}

\section{Part II: The Framework}\label{part-ii-the-framework}

\section{Chapter 4: Mathematical
Preliminaries}\label{chapter-4-mathematical-preliminaries}

\begin{quote}
\textbf{\textsc{RUNNING EXAMPLE --- Priya's Model}}

\emph{Priya sketches on a whiteboard. A patient's situation is a point
in some space. What space? She lists relevant dimensions: medical
severity, trial compatibility, travel burden, insurance status,
support-network strength, informed-consent capacity. Six dimensions, at
minimum. She draws arrows representing different fairness
adjustments---tangent vectors. She notices that adjusting `travel
burden' in one direction helps rural patients but hurts patients with
mobility disabilities, because these dimensions are coupled. She needs a
metric that tells her how the dimensions relate. She does not have the
vocabulary yet, but she is discovering that her patient space is a
manifold, and her fairness interventions are vectors in its tangent
space.}
\end{quote}

\emph{This chapter develops the geometric vocabulary used throughout the
book. No prior knowledge of differential geometry is assumed; we require
only multivariate calculus and linear algebra. Readers with background
in differential geometry may skim Sections 4.1--4.5 and focus on Section
4.6 (stratified spaces) and Section 4.7 (gauge theory), which introduce
structures less commonly encountered.}

\subsection{A Note on Epistemic
Status}\label{a-note-on-epistemic-status}

The mathematical structures introduced in this chapter --- manifolds,
tensors, metrics, connections, fiber bundles, stratified spaces, gauge
transformations --- are tools. They are powerful tools, honed over two
centuries by mathematicians and physicists working on problems where
structure varies across space. Their power lies in their
domain-adaptedness: they were built to describe exactly the kind of
phenomena (curvature, path-dependence, boundary discontinuity, gauge
invariance) that the moral domain exhibits. But they are instruments,
not mirrors. The question of whether the moral manifold is ``really
there'' or ``usefully modeled'' is addressed in Chapter 9 (which argues
for a governance account) and Chapter 29 (which surveys the realist,
instrumentalist, and governance positions). This chapter teaches how to
use the toolkit. The reader should carry through the formalism the
understanding that every definition and theorem that follows is
conditional on modeling choices made explicit in Chapter 5.

One limitation deserves mention at the outset. Gödel's incompleteness
theorems establish that no formal system of sufficient expressive power
can prove its own consistency. The framework developed in this book is
such a system. It cannot, therefore, self-validate --- it cannot prove
from within that its axioms are the right axioms or that its theorems
capture what they claim to capture. This is not a defect peculiar to
geometric ethics; it is a structural feature of all formal systems, from
Peano arithmetic to the Standard Model. The appropriate response is not
to abandon formalism but to test it externally --- against empirical
data (Chapter 17), against competing frameworks (Appendix E), and
against the collective judgment of practitioners who use the tools and
evaluate their fit. The framework's authority is pragmatic, not
self-grounding.

\subsection{4.1 Manifolds}\label{manifolds}

\subsubsection{The Idea}\label{the-idea}

A manifold is a space that \emph{locally} looks like ordinary Euclidean
space ℝⁿ but may have nontrivial \emph{global} structure. The surface of
the Earth is the standard example: any small region looks flat (a
neighborhood of a point can be mapped onto a patch of ℝ²), but the whole
surface is curved, closed, and cannot be covered by a single flat map
without distortion.

This is precisely the concept we need for moral space. Locally --- in
the neighborhood of a single moral situation --- moral trade-offs may
look like ordinary multi-dimensional optimization: move a little in one
direction to improve fairness, a little in another to improve
efficiency. The space of nearby variations is flat. But globally, moral
space may be curved (obligations change as you traverse it), closed in
certain directions (you cannot go past certain constraints), or have
nontrivial topology (circling back to the ``same'' situation may yield a
different moral assessment).

\subsubsection{Formal Definition}\label{formal-definition}

\textbf{Definition 4.1 (Topological Manifold).} An \emph{n-dimensional
topological manifold} is a topological space M that is: 1.
\textbf{Hausdorff:} Distinct points can be separated by disjoint open
sets. 2. \textbf{Second-countable:} The topology has a countable basis.
3. \textbf{Locally Euclidean:} Every point p ∈ M has a neighborhood
homeomorphic to an open subset of ℝⁿ.

Condition (3) is the essential one: around every point, there is a
\emph{coordinate chart} --- a continuous bijection φ: U → V, where U ⊂ M
is open and V ⊂ ℝⁿ is open. The chart assigns n real numbers
(coordinates) to each point in U. Conditions (1) and (2) are topological
regularity conditions that exclude pathological spaces.

A collection of charts \{(U\_α, φ\_α)\} that covers all of M is called
an \emph{atlas}. When two charts overlap --- when U\_α ∩ U\_β ≠ ∅ --- we
get a \emph{transition function}

\[\psi_{\alpha\beta} = \phi_{\beta} \circ \phi_{\alpha}^{- 1}:\phi_{\alpha}(U_{\alpha} \cap U_{\beta}) \rightarrow \phi_{\beta}(U_{\alpha} \cap U_{\beta})\]

which tells us how to translate between the two coordinate systems on
their overlap.

\textbf{Definition 4.2 (Smooth Manifold).} A topological manifold is
\emph{smooth} if it admits an atlas whose transition functions are all
smooth (infinitely differentiable, i.e., C\^{}∞).

Smoothness is what allows us to do calculus on M: to take derivatives,
define gradients, and study how quantities change from point to point.

\subsubsection{Example: The Circle}\label{example-the-circle}

The circle S¹ is a 1-dimensional manifold. No single coordinate chart
covers S¹ (a single chart would have to be homeomorphic to an interval
in ℝ, and circles are not homeomorphic to intervals). But two charts
suffice: one covering all but the ``north pole'' and one covering all
but the ``south pole.'' The transition function on their overlap is
smooth. Hence S¹ is a smooth 1-manifold.

\subsubsection{Example: The 2-Simplex}\label{example-the-2-simplex}

The space of probability distributions over three outcomes is the
2-simplex:

\[\Delta^{2} = \{(p_{1},p_{2},p_{3}) \in \mathbb{R}^{3}:p_{i} \geq 0,\mspace{6mu} p_{1} + p_{2} + p_{3} = 1\}\]

The \emph{interior} of Δ² (where all p\_i \textgreater{} 0) is a
2-dimensional smooth manifold. The \emph{boundary} --- where one or more
p\_i = 0 --- is not itself a manifold of the same dimension, but the
whole simplex has the structure of a \emph{manifold with boundary} (or,
as we shall see, a stratified space).

This example is directly relevant: the simplex of probability
distributions over actions is a natural candidate for part of moral
space.

\subsection{4.2 Tangent Spaces and Tangent
Bundles}\label{tangent-spaces-and-tangent-bundles}

\subsubsection{Tangent Vectors}\label{tangent-vectors}

At each point p of a smooth manifold M, there is a \emph{tangent space}
T\_pM --- a vector space of the same dimension as M that represents the
``infinitesimal directions'' one can move from p.

In coordinates (x¹, \ldots, xⁿ), the tangent space at p is spanned by
the partial derivative operators \{∂/∂x¹\_p, \ldots, ∂/∂xⁿ\_p\}. A
tangent vector v ∈ T\_pM can be written

\[v = v^{i}\frac{\partial}{\partial x^{i}}|_{p}\]

where v\^{}i are the components of v in this coordinate basis. (We use
the Einstein summation convention: repeated upper and lower indices are
summed over.)

\textbf{Intuition:} A tangent vector at p is a ``velocity'' --- a
direction and rate of change. If γ(t) is a curve on M with γ(0) = p,
then its velocity γ'(0) is a tangent vector at p.~The tangent space
T\_pM is the collection of all such velocities.

\textbf{Moral interpretation:} At a point p in moral space (a particular
moral situation), the tangent space T\_pM represents the possible
\emph{variations} of the situation: small changes in resource
allocation, slight modifications of commitments, marginal shifts in
context. An obligation vector O(p) ∈ T\_pM points in the direction one
\emph{ought} to move.

\subsubsection{The Tangent Bundle}\label{the-tangent-bundle}

\textbf{Definition 4.3 (Tangent Bundle).} The \emph{tangent bundle} TM
is the disjoint union of all tangent spaces:

\[TM = \bigsqcup_{p \in M}T_{p}M = \{(p,v):p \in M,\mspace{6mu} v \in T_{p}M\}\]

TM is itself a smooth manifold of dimension 2n (n coordinates for the
base point, n for the tangent vector). A \emph{vector field} on M is a
smooth assignment of a tangent vector to each point --- a smooth section
X: M → TM satisfying X(p) ∈ T\_pM for all p.

\subsubsection{Transformation of Tangent
Vectors}\label{transformation-of-tangent-vectors}

Under a change of coordinates x → x̃, tangent vector components transform
by the \emph{Jacobian}:

\[{\widetilde{v}}^{i} = \frac{\partial{\widetilde{x}}^{i}}{\partial x^{j}}v^{j}\]

This is the prototype of all tensor transformation laws. Objects that
transform by one copy of the Jacobian are called \emph{contravariant
vectors} (or simply vectors). The transformation law ensures that the
geometric object v is independent of the coordinate choice --- only its
\emph{components} change.

\subsection{4.3 Covectors, Dual Spaces, and the Cotangent
Bundle}\label{covectors-dual-spaces-and-the-cotangent-bundle}

\subsubsection{Covectors}\label{covectors}

A \emph{covector} (or \emph{1-form}) at p is a linear map ω: T\_pM → ℝ.
The space of all covectors at p is the \emph{dual space} T\_p\emph{M,
also called the} cotangent space*.

In coordinates, covectors have components ω\_i (lower index) and
transform by the \emph{inverse} Jacobian:

\[{\widetilde{\omega}}_{i} = \frac{\partial x^{j}}{\partial{\widetilde{x}}^{i}}\omega_{j}\]

The fundamental operation between a vector and a covector is
\emph{contraction}:

\[\omega(v) = \omega_{i}v^{i}\mathbb{\in R}\]

This produces a scalar --- a number independent of coordinates. The
contraction is the geometric content of ``evaluating'': applying a
covector to a vector yields an invariant quantity.

\textbf{Moral interpretation:} If obligations are vectors (upper index,
pointing in the direction one should move), then \emph{interests} are
naturally covectors (lower index, assigning a real-valued weight to each
direction). The contraction S = I\_μ O\^{}μ --- interest applied to
obligation --- yields \emph{satisfaction}, a scalar. This fundamental
formula, developed in Chapter 6, is the geometric expression of ``how
well obligations align with interests.''

\subsubsection{The Gradient}\label{the-gradient}

The most natural covector is the \emph{gradient} of a function. If f: M
→ ℝ is a smooth function, its gradient df is a covector field:

\[(df)_{i} = \frac{\partial f}{\partial x^{i}}\]

The gradient df at p is the covector that, applied to any tangent vector
v, gives the rate of change of f in the direction v:

\[df(v) = v^{i}\frac{\partial f}{\partial x^{i}} = v(f)\]

\textbf{Note:} The gradient is a covector, not a vector. In ℝⁿ with the
standard Euclidean metric, the distinction is invisible (the metric
provides a canonical identification between vectors and covectors). On a
general manifold, the distinction matters --- and it matters for ethics.
Converting a gradient (covector) into a ``direction of steepest ascent''
(vector) requires the metric. Different metrics yield different
``steepest'' directions from the same gradient. This is the geometric
content of the observation that different ethical theories, given the
same moral landscape, may prescribe different actions.

\subsection{4.4 Tensors}\label{tensors}

\subsubsection{Definition}\label{definition}

A \emph{tensor of type (r, s)} at p ∈ M is a multilinear map

\[T:\underset{r}{\underbrace{T_{p}^{*}M \times \cdots \times T_{p}^{*}M}} \times \underset{s}{\underbrace{T_{p}M \times \cdots \times T_{p}M}}\mathbb{\rightarrow R}\]

In coordinates, T has components
\(T_{j_{1}\cdots j_{s}}^{i_{1}\cdots i_{r}}\) with r upper indices and s
lower indices. The \emph{rank} of the tensor is (r, s), and the total
number of indices is r + s.

\textbf{Examples:} - A scalar is a (0, 0)-tensor: a number, invariant
under coordinate changes. - A vector is a (1, 0)-tensor: one upper
index, transforms by the Jacobian. - A covector is a (0, 1)-tensor: one
lower index, transforms by the inverse Jacobian. - The metric (Section
4.5) is a (0, 2)-tensor: two lower indices.

\subsubsection{Transformation Law}\label{transformation-law}

Under a coordinate change x → x̃, the components of a (r, s)-tensor
transform as:

\[{\widetilde{T}}_{j_{1}\cdots j_{s}}^{i_{1}\cdots i_{r}} = \frac{\partial{\widetilde{x}}^{i_{1}}}{\partial x^{k_{1}}}\cdots\frac{\partial{\widetilde{x}}^{i_{r}}}{\partial x^{k_{r}}}\frac{\partial x^{l_{1}}}{\partial{\widetilde{x}}^{j_{1}}}\cdots\frac{\partial x^{l_{s}}}{\partial{\widetilde{x}}^{j_{s}}}T_{l_{1}\cdots l_{s}}^{k_{1}\cdots k_{r}}\]

Each upper index contributes one Jacobian factor; each lower index
contributes one inverse Jacobian factor. The key property: \emph{the
tensor itself is a geometric object independent of coordinates; only its
components change.}

This is why tensors are central to the framework. The requirement that
moral evaluation be independent of description --- the Bond Invariance
Principle --- is precisely the requirement that moral quantities be
\emph{tensorial}: their values, after proper transformation, are the
same in every coordinate system.

\subsubsection{Contraction}\label{contraction}

\emph{Contraction} is the operation of summing over a paired upper and
lower index:

\[C(T)_{j_{1}\cdots{\widehat{j}}_{l}\cdots j_{s}}^{i_{1}\cdots{\widehat{i}}_{k}\cdots i_{r}} = T_{j_{1}\cdots m\cdots j_{s}}^{i_{1}\cdots m\cdots i_{r}}\]

Contraction reduces the rank by (1, 1). It produces a new tensor whose
value is independent of coordinates (the sum over the shared index is
invariant).

\textbf{Moral significance:} The contraction from a high-rank moral
tensor to a scalar decision is the geometric version of what happens
when we \emph{decide}. Information is lost (rank decreases), but the
result is invariant. Chapter 15 develops the concept of \emph{moral
contraction} in detail: the process by which the rich tensorial
structure of a moral situation is reduced to an action, and the
\emph{moral residue} --- the normative information discarded by
contraction.

\subsubsection{Tensor Fields}\label{tensor-fields}

A \emph{tensor field} of type (r, s) on M is a smooth assignment of a
(r, s)-tensor to each point p ∈ M. The metric tensor, curvature tensor,
and moral evaluation tensors are all tensor fields.

\subsection{4.5 The Metric Tensor}\label{the-metric-tensor}

\subsubsection{Definition}\label{definition-1}

A \emph{metric tensor} (or simply \emph{metric}) on M is a smooth (0,
2)-tensor field g that is: 1. \textbf{Symmetric:} g(v, w) = g(w, v) for
all tangent vectors v, w. 2. \textbf{Non-degenerate:} If g(v, w) = 0 for
all w, then v = 0.

In coordinates:

\[g = g_{ij}\, dx^{i} \otimes dx^{j}\]

The metric defines an \emph{inner product} on each tangent space: g(v,
w) = g\_\{ij\} v\^{}i w\^{}j. This allows us to measure:

\begin{itemize}
\item
  \textbf{Lengths:} v = √(g(v, v)) (when the inner product is positive).
\item
  \textbf{Angles:} The angle between v and w is determined by g(v, w) /
  (v w).
\item
  \textbf{Distances:} The length of a curve γ from a to b is ∫\_a\^{}b
  √(g(γ', γ')) dt.
\end{itemize}

\subsubsection{Riemannian vs.~Semi-Riemannian
vs.~Degenerate}\label{riemannian-vs.-semi-riemannian-vs.-degenerate}

If g is positive-definite (g(v, v) \textgreater{} 0 for all nonzero v),
we have a \emph{Riemannian metric}, and M is a Riemannian manifold.
Distances are always positive, and the geometry behaves like a
generalization of Euclidean space.

If g has indefinite signature (some directions give g(v, v)
\textgreater{} 0, others g(v, v) \textless{} 0), we have a
\emph{semi-Riemannian} (or pseudo-Riemannian) metric. Spacetime in
general relativity has signature (-,+,+,+).

For moral space, we must also consider \emph{degenerate} metrics:
metrics where condition (2) fails, so that g(v, w) = 0 for all w does
not imply v = 0. A degenerate metric means certain directions are
``invisible'' to the inner product --- they have zero length and zero
angle with everything.

\textbf{Moral interpretation:} A degenerate metric represents
\emph{incommensurability}. If the metric is degenerate along the
subspace spanned by ``ecological value'' and ``economic value,'' it
means there is no well-defined trade-off ratio between them. They cannot
be compared by the metric --- not because we lack information, but
because the geometric structure of moral space does not support the
comparison. This is a precise formalization of a claim that moral
philosophers have made informally for decades.

\subsubsection{The Metric as Musical
Isomorphism}\label{the-metric-as-musical-isomorphism}

A non-degenerate metric provides a canonical map between vectors and
covectors:

\[v^{i} \mapsto v_{j} = g_{ij}v^{i}\quad\quad\text{("lowering indices")}\]

\[\omega_{i} \mapsto \omega^{j} = g^{ij}\omega_{i}\quad\quad\text{("raising indices")}\]

where g\^{}\{ij\} is the inverse of g\_\{ij\}. This is called the
\emph{musical isomorphism} (the notation ♭ for lowering and ♯ for
raising). It identifies vectors with covectors --- but the
identification depends on the choice of metric. Different metrics yield
different identifications.

This is why the metric matters: it determines how we convert between
obligations (vectors) and interests (covectors), and how we compute the
satisfaction scalar S = Iμ Oμ. When the interest is naturally a vector
rather than a covector, the metric mediates: S = gμν Oμ Iν. Different
metrics --- different ethical theories --- assign different
satisfactions to the same obligation-interest pair.

\subsection{4.6 Connections, Parallel Transport, and
Curvature}\label{connections-parallel-transport-and-curvature}

\subsubsection{The Problem of
Comparison}\label{the-problem-of-comparison}

On a flat space, comparing vectors at different points is
straightforward: translate one to the location of the other and compare.
On a curved manifold, this does not work. There is no canonical way to
identify the tangent space at p with the tangent space at q. To compare
vectors at different points, we need additional structure: a
\emph{connection}.

\textbf{Moral interpretation:} Can we compare obligations across
contexts? The obligation to keep a promise in one situation and the
obligation to help a stranger in another --- are they ``the same''
obligation, differently instantiated? On a flat moral manifold, the
answer would be straightforward. On a curved one, it depends on
\emph{how we transport the obligation from one context to the other} ---
and the answer may depend on the path.

\subsubsection{The Connection}\label{the-connection}

\textbf{Definition 4.4 (Affine Connection).} An \emph{affine connection}
(or \emph{covariant derivative}) on M is a map ∇ that takes a vector
field X and a vector field Y and produces a new vector field ∇\_X Y,
satisfying:

\begin{enumerate}
\def\labelenumi{\arabic{enumi}.}
\item
  ∇\_\{fX+gY\} Z = f∇\_X Z + g∇\_Y Z (linearity in the first argument)
\item
  ∇\_X(Y + Z) = ∇\_X Y + ∇\_X Z (additivity in the second argument)
\item
  ∇\_X(fY) = (Xf)Y + f∇\_X Y (Leibniz rule)
\end{enumerate}

The connection is specified in coordinates by the \emph{Christoffel
symbols} Γ\^{}k\_\{ij\}:

\[\nabla_{\partial_{i}}\partial_{j} = \Gamma_{ij}^{k}\,\partial_{k}\]

The covariant derivative of a vector field V along a direction ∂\_i is:

\[(\nabla_{i}V)^{k} = \frac{\partial V^{k}}{\partial x^{i}} + \Gamma_{ij}^{k}V^{j}\]

The Christoffel symbols encode how the coordinate basis vectors change
from point to point --- the ``correction'' needed to account for the
curvature of the manifold.

\textbf{Levi-Civita connection.} On a Riemannian manifold, there is a
unique connection that is: - \emph{Metric-compatible}: ∇g = 0 (parallel
transport preserves inner products) - \emph{Torsion-free}: ∇\_X Y - ∇\_Y
X = {[}X, Y{]}

This is the Levi-Civita connection, and its Christoffel symbols are
determined by the metric:

\[\Gamma_{ij}^{k} = \frac{1}{2}g^{kl}\left( \frac{\partial g_{il}}{\partial x^{j}} + \frac{\partial g_{jl}}{\partial x^{i}} - \frac{\partial g_{ij}}{\partial x^{l}} \right)\]

\subsubsection{Parallel Transport}\label{parallel-transport}

A vector field V along a curve γ(t) is \emph{parallel-transported} if
∇\_\{'\} V = 0 --- the covariant derivative of V along the curve
vanishes. Parallel transport moves a vector along a path while keeping
it ``as constant as possible,'' given the curvature of the space.

On a flat manifold, parallel transport is path-independent: moving a
vector from p to q gives the same result regardless of the route. On a
curved manifold, the result depends on the path.

\textbf{Example:} On the sphere S², start at the equator with a vector
pointing east. Parallel-transport it north to the pole, then south along
a different meridian to the equator. The vector no longer points east
--- it has been rotated by an angle equal to the solid angle enclosed by
the path. This rotation is called \emph{holonomy}.

\textbf{Moral example:} Transport the obligation ``keep your promise''
from a context of personal friendship to a context of professional duty
by two different routes --- one passing through the context of family
obligation, the other through the context of institutional requirement.
If the resulting obligation differs, moral space has curvature. The
difference is the \emph{moral holonomy}, and it measures how much the
moral landscape twists between these contexts.

\subsubsection{The Riemann Curvature
Tensor}\label{the-riemann-curvature-tensor}

\textbf{Definition 4.5 (Riemann Curvature Tensor).} The curvature of a
connection ∇ is measured by the (1, 3)-tensor:

\[R(X,Y)Z = \nabla_{X}\nabla_{Y}Z - \nabla_{Y}\nabla_{X}Z - \nabla_{\lbrack X,Y\rbrack}Z\]

In components:

\[R_{kij}^{l} = \frac{\partial\Gamma_{kj}^{l}}{\partial x^{i}} - \frac{\partial\Gamma_{ki}^{l}}{\partial x^{j}} + \Gamma_{im}^{l}\Gamma_{kj}^{m} - \Gamma_{jm}^{l}\Gamma_{ki}^{m}\]

The Riemann tensor vanishes if and only if the manifold is flat (locally
isometric to Euclidean space). When it is nonzero, it encodes the
path-dependence of parallel transport: transporting a vector V around an
infinitesimal parallelogram spanned by X and Y yields a change
proportional to R(X, Y)V.

\textbf{Contractions of the Riemann tensor} yield lower-rank curvature
tensors: - The \emph{Ricci tensor}: R\_\{ij\} = R\^{}k\{\}\emph{\{ikj\}
(a symmetric (0, 2)-tensor) - The scalar curvature: R = g\^{}\{ij\}
R}\{ij\} (a single number summarizing overall curvature at each point)

Whether moral space has nonzero curvature is an empirical question ---
one of the open problems in Chapter 29. But the framework is equipped to
handle it if it does.

\subsection{4.7 Fiber Bundles}\label{fiber-bundles}

\subsubsection{Motivation}\label{motivation}

\textbar{} Agent Fiber Bundle

\subsubsection{Definition}\label{definition-2}

\textbf{Definition 4.6 (Fiber Bundle).} A \emph{fiber bundle} is a
triple (E, π, M) where: - E is a topological space called the
\emph{total space} - M is a topological space called the \emph{base
space} - π: E → M is a continuous surjection called the
\emph{projection} - For each p ∈ M, the \emph{fiber} F\_p = π⁻¹(p) is
homeomorphic to a fixed space F (the \emph{typical fiber}) - \emph{Local
triviality}: Each point p has a neighborhood U such that π⁻¹(U) ≅ U × F

The total space E looks locally like a product M × F, but may be
``twisted'' globally --- like a Möbius strip, which locally looks like a
strip of paper but has a global twist.

\textbf{Example (Tangent Bundle).} The tangent bundle TM is a fiber
bundle with base space M, total space TM, and fiber ℝⁿ (the tangent
space at each point).

\textbf{Example (Frame Bundle).} At each point p, the \emph{frame
bundle} consists of all ordered bases for T\_pM. The fiber is GL(n, ℝ)
--- the group of invertible n × n matrices.

\subsubsection{Sections}\label{sections}

A \emph{section} of a fiber bundle is a smooth map σ: M → E such that
π(σ(p)) = p for all p --- it selects one element from each fiber. A
vector field is a section of the tangent bundle. A metric tensor is a
section of the bundle of symmetric (0, 2)-tensors.

\textbf{Moral interpretation:} Consider a bundle over moral space M
whose fiber at each point p is the space of \emph{possible moral
character traits} available to an agent in situation p.~A \emph{virtue}
is then a section of this bundle: a consistent assignment of character
to each situation. The requirement that virtue be \emph{consistent} ---
that it not jump discontinuously --- is the requirement that the section
be smooth. Virtue ethics, in this language, is the study of sections of
the character fiber bundle.

\subsection{4.8 Stratified Spaces}\label{stratified-spaces}

\subsubsection{Beyond Smooth Manifolds}\label{beyond-smooth-manifolds}

Smooth manifolds are fundamental, but they are not sufficient for moral
space. Manifolds are, by definition, uniform in dimension: every point
has a neighborhood homeomorphic to ℝⁿ for the \emph{same} n.~But moral
space has regions of different effective dimensions (the interior of the
simplex has dimension 2; its edges have dimension 1; its vertices have
dimension 0), separated by boundaries where the smooth structure
changes.

For this, we need \emph{stratified spaces}.

\subsubsection{Definition}\label{definition-3}

\textbf{Definition 4.7 (Whitney Stratification).} A \emph{Whitney
stratified space} is a topological space X together with a locally
finite partition into smooth manifolds \{S\_α\}\_\{α ∈ A\} (called
\emph{strata}) such that:

\begin{enumerate}
\def\labelenumi{\arabic{enumi}.}
\item
  \textbf{Frontier condition:} If S\_α ∩ cl(S\_β) ≠ ∅ and S\_α ≠ S\_β,
  then S\_α ⊂ cl(S\_β) and dim(S\_α) \textless{} dim(S\_β).
\item
  \textbf{Whitney's condition (B):} For strata S\_α ⊂ cl(S\_β), consider
  sequences \{x\_n\} ⊂ S\_β and \{y\_n\} ⊂ S\_α both converging to y ∈
  S\_α. If the secant lines ℓ\_n = x\_n y\_n converge to a line ℓ, and
  the tangent planes T\_\{x\_n\} S\_β converge to a plane τ, then ℓ ⊂ τ.
\end{enumerate}

Condition (1) says that lower-dimensional strata lie on the boundaries
of higher-dimensional strata --- they are genuinely ``boundary''
structures, not isolated points floating in space. Condition (2) is a
regularity condition ensuring that higher-dimensional strata approach
lower-dimensional strata in a controlled way --- no ``pinching'' or
pathological tangencies.

\subsubsection{Example: The Simplex}\label{example-the-simplex}

The 2-simplex Δ² has a natural Whitney stratification: -
\textbf{2-stratum} (interior): \{(p₁, p₂, p₃) : all pᵢ \textgreater{}
0\}, a 2-dimensional manifold. - \textbf{1-strata} (edges): Three edges,
each a 1-dimensional manifold. - \textbf{0-strata} (vertices): Three
vertices, each a 0-dimensional manifold (a point).

The frontier condition holds: each vertex is in the closure of each
adjacent edge, and each edge is in the closure of the interior.
Whitney's condition (B) holds because the simplex is a convex polytope.

\subsubsection{Moral Significance}\label{moral-significance}

Stratification is the geometric structure that represents moral
\emph{thresholds}, \emph{phase transitions}, and \emph{categorical
boundaries}:

\begin{itemize}
\item
  The transition from ``permissible'' to ``forbidden'' is a stratum
  boundary.
\item
  The transition from ``one duty applies'' to ``a different duty
  applies'' is a stratum boundary.
\item
  The point where ``consent'' shifts from insufficient to sufficient is
  a stratum boundary.
\item
  A genuine dilemma --- where two strata of different moral regimes meet
  at a singularity --- is a 0-stratum at the intersection of
  higher-dimensional boundaries.
\end{itemize}

Within each stratum, smooth calculus applies: we can take derivatives,
define gradients, and optimize. \emph{Across} stratum boundaries, we
need the machinery of stratified spaces to track how the structure
changes. The transition functions between strata encode the ``rules of
passage'' from one moral regime to another --- and these transitions may
be discontinuous.

\subsection{4.9 Gauge Theory: A Brief
Introduction}\label{gauge-theory-a-brief-introduction}

\subsubsection{Motivation}\label{motivation-1}

Gauge theory is the mathematical framework that describes how local
symmetries constrain global structure. In physics, it governs all four
fundamental forces. In this book, it provides the mathematical content
of the Bond Invariance Principle: the requirement that moral evaluations
be invariant under admissible re-descriptions.

\subsubsection{Gauge Transformations}\label{gauge-transformations}

Consider a fiber bundle (E, π, M) with structure group G --- a Lie group
that acts on each fiber. A \emph{gauge transformation} is a smooth map
g: M → G that acts on sections by:

\[\sigma(p) \mapsto g(p) \cdot \sigma(p)\]

In physics, gauge transformations change the ``internal'' description at
each point (e.g., the phase of a quantum field) without changing the
physics. In ethics, gauge transformations change the
\emph{representation} of a moral situation (the framing, the language,
the cultural context) without changing the moral content.

\textbf{Definition 4.8 (Gauge Invariance).} A quantity Q defined on
sections of a fiber bundle is \emph{gauge-invariant} if Q(g · σ) = Q(σ)
for all gauge transformations g.

The Bond Invariance Principle states that legitimate moral evaluations
must be gauge-invariant: if two descriptions are related by an
admissible re-description (a gauge transformation), they must receive
the same moral assessment.

\subsubsection{Connections on Fiber
Bundles}\label{connections-on-fiber-bundles}

A \emph{connection} on a fiber bundle generalizes the affine connection
of Section 4.6 to arbitrary fibers. It specifies how to
parallel-transport elements of the fiber along paths in the base space.

In coordinates, a connection is given by a \emph{gauge field} A --- a
1-form on M valued in the Lie algebra 𝔤 of the structure group G:

\[A = A_{\mu}^{a}T_{a}\, dx^{\mu}\]

where T\_a are generators of 𝔤. The gauge field tells us how the
``internal'' structure (the fiber) changes as we move through the base
space.

\subsubsection{Curvature of a
Connection}\label{curvature-of-a-connection}

The \emph{curvature} of a gauge connection is the 2-form:

\[F = dA + A \land A\]

In components:

\[F_{\mu\nu}^{a} = \partial_{\mu}A_{\nu}^{a} - \partial_{\nu}A_{\mu}^{a} + f_{bc}^{a}A_{\mu}^{b}A_{\nu}^{c}\]

where f\^{}a\{\}\_\{bc\} are the structure constants of the Lie algebra.
The curvature measures the failure of parallel transport to be
path-independent --- the holonomy of the connection.

If F = 0, the connection is \emph{flat} and parallel transport is
path-independent. If F ≠ 0, different paths between the same endpoints
yield different results.

\textbf{Moral interpretation:} Moral curvature F ≠ 0 means that
transporting a moral assessment along different paths through moral
space yields different results. The obligation you arrive at depends on
the route you took --- which intermediate contexts you passed through.
This is the geometric content of moral path-dependence, and it is
distinct from mere disagreement: it is a structural feature of the moral
landscape itself.

\subsubsection{Noether's Theorem
(Preview)}\label{noethers-theorem-preview}

A central result of gauge theory is \emph{Noether's theorem}: every
continuous symmetry of a system corresponds to a conserved quantity. The
connection to ethics is developed in Chapter 12, but the basic structure
is:

\begin{itemize}
\item
  \textbf{Symmetry:} Moral evaluation is invariant under re-description
  (gauge invariance).
\item
  \textbf{Conservation law:} Harm is conserved --- it cannot be created
  or destroyed by re-description, only generated by wrongdoing or
  repaired by restorative action.
\end{itemize}

The mathematical proof follows the same pattern as in physics: the
invariance of the action functional under infinitesimal gauge
transformations implies, via the Euler-Lagrange equations, the existence
of a conserved current.

\subsection{4.10 Summary: The Geometric
Toolkit}\label{summary-the-geometric-toolkit}

The table below summarizes the mathematical structures introduced in
this chapter and previews their moral applications (developed in
subsequent chapters).

{\def\LTcaptype{none} % do not increment counter
\begin{longtable}[]{@{}
  >{\raggedright\arraybackslash}p{(\linewidth - 4\tabcolsep) * \real{0.2690}}
  >{\raggedright\arraybackslash}p{(\linewidth - 4\tabcolsep) * \real{0.3754}}
  >{\raggedright\arraybackslash}p{(\linewidth - 4\tabcolsep) * \real{0.3556}}@{}}
\toprule\noalign{}
\begin{minipage}[b]{\linewidth}\raggedright
\textbf{Structure}
\end{minipage} & \begin{minipage}[b]{\linewidth}\raggedright
\textbf{Definition}
\end{minipage} & \begin{minipage}[b]{\linewidth}\raggedright
\textbf{Moral Application}
\end{minipage} \\
\midrule\noalign{}
\endhead
\bottomrule\noalign{}
\endlastfoot
\textbf{Manifold} M & Space locally homeomorphic to ℝⁿ & The moral
manifold: space of structured situations \\
\textbf{Tangent vector} v ∈ T\_pM & Infinitesimal direction at a point &
Possible variations; obligations as directions \\
\textbf{Covector} ω ∈ T\_p*M & Linear functional on tangent vectors &
Interests; weights on moral dimensions \\
\textbf{Tensor} T of type (r, s) & Multilinear map; transforms by
Jacobians & Moral evaluations; invariant under re-description \\
\textbf{Contraction} & Summing paired indices; reduces rank & From
tensorial evaluation to scalar decision \\
\textbf{Metric} g & Inner product on tangent spaces & Trade-off
structure between values \\
\textbf{Degenerate metric} & Metric with zero eigenvalues &
Incommensurability of values \\
\textbf{Connection} ∇ & Rule for parallel transport & How obligations
change across contexts \\
\textbf{Curvature} R & Path-dependence of parallel transport & Whether
moral assessment is path-dependent \\
\textbf{Fiber bundle} (E, π, M) & Space of ``internal'' structures over
M & Agent perspectives; character traits; jural states \\
\textbf{Gauge transformation} & Change of internal representation &
Re-description; change of framing \\
\textbf{Gauge invariance} & Invariance under gauge transformations &
Bond Invariance Principle \\
\textbf{Stratified space} & Manifold with boundary strata & Moral
thresholds; phase transitions; categorical boundaries \\
\end{longtable}
}

\subsection{4.11 Bibliographic Notes}\label{bibliographic-notes}

The differential geometry in Sections 4.1--4.6 is standard and can be
found in Lee (2012), \emph{Introduction to Smooth Manifolds}, which is
the gentlest rigorous introduction, or do Carmo (1992), \emph{Riemannian
Geometry}, for the metric and curvature theory. Readers seeking physical
intuition may prefer Frankel (2011), \emph{The Geometry of Physics}.

For fiber bundles and gauge theory (Section 4.9), Nakahara (2003),
\emph{Geometry, Topology, and Physics}, provides an accessible treatment
aimed at physicists. The mathematical foundations are in Kobayashi and
Nomizu (1963), \emph{Foundations of Differential Geometry}, though this
is aimed at specialists.

For stratified spaces (Section 4.8), the original source is Whitney
(1965). An accessible modern treatment is Pflaum (2001), \emph{Analytic
and Geometric Study of Stratified Spaces}. The application to moral
space is, to our knowledge, original to this work and to the companion
paper on Stratified Geometric Ethics (Bond 2025).

Readers who wish to go deeper before proceeding should read Lee (2012),
Chapters 1--6, which covers manifolds, tangent spaces, and vector
bundles at the level assumed in subsequent chapters.

❖

\emph{The mathematics in this chapter is not new.}

\emph{It is the standard toolkit of modern geometry, developed over two
centuries by mathematicians working on problems in physics, topology,
and analysis.}

\emph{What is new is the application.}

\emph{In the next chapter, we build the moral manifold --- the specific
space on which geometric ethics takes shape.}

\section{Chapter 5: The Moral
Manifold}\label{chapter-5-the-moral-manifold}

\begin{quote}
\textbf{\textsc{RUNNING EXAMPLE --- Priya's Model}}

\emph{Priya maps TrialMatch onto the 3×3 moral manifold. For the
BEACON-7 match, μ = 1 (Empirical--Deontic): HealthBridge has a
contractual duty to match using validated criteria---an observable fact.
μ = 4 (Normative--Deontic): the algorithm ought to give rural patients a
fair chance, but `fair' is undefined in the codebase. μ = 6
(Normative--Consequential): matching more rural patients should improve
trial generalizability. μ = 9 (Meta--Consequential): what theory of
value grounds the assumption that statistical power is good? She fills
in all nine cells and realizes TrialMatch uses at most three dimensions
and ignores the other six entirely. Six dimensions of moral reality,
invisible to her algorithm.}
\end{quote}

\subsection{5.1 The Question of Base
Space}\label{the-question-of-base-space}

Every tensor lives on a manifold. The stress tensor in materials science
lives on the manifold of spatial points within a body. The metric tensor
in general relativity lives on the manifold of spacetime events. The
electromagnetic field tensor lives on the four-dimensional manifold of
spacetime, and the gauge connection lives on a fiber bundle over it. The
question for geometric ethics is: \emph{what is the manifold on which
moral tensors are defined?}

This is not a technical detail. The choice of base space determines:

\begin{itemize}
\item
  What counts as a ``point'' in moral reasoning
\item
  What counts as a ``direction'' --- what the tangent vectors are
\item
  What transformations are admissible and which quantities are invariant
\item
  What it means for two moral situations to be ``nearby'' or ``far
  apart''
\item
  Where the smooth structure breaks down --- where discontinuities,
  boundaries, and singularities occur
\end{itemize}

This chapter develops the concept of the \emph{moral manifold} --- the
space of morally relevant configurations over which ethical tensors are
defined. We build on the mathematical toolkit of Chapter 4, but proceed
with philosophical care, distinguishing geometric structure from
metaphysical commitment.

\subsection{5.2 What Are the Points?}\label{what-are-the-points}

The first question is fundamental: what are the ``points'' of moral
space? Several candidates present themselves, each with distinctive
strengths and limitations.

\subsubsection{Candidate 1: Possible
Actions}\label{candidate-1-possible-actions}

On this view, the moral manifold M is the space of possible actions
available to an agent. Each point x ∈ M represents a complete
specification of what the agent does --- not just ``help the neighbor''
but a full description including manner, timing, motivation, and
foreseeable consequences.

\textbf{Strengths:} This connects directly to the practical question of
ethics (``what should I do?''). The gradient of the satisfaction
function, ∇S, points toward the best available action.

\textbf{Weaknesses:} Actions are agent-relative. My action space differs
from yours. This makes interpersonal comparison difficult and global
structure unclear. Moreover, the space may not be a manifold in any
natural sense --- discrete choices (accept or decline) do not admit
smooth interpolation.

\subsubsection{Candidate 2: States of
Affairs}\label{candidate-2-states-of-affairs}

On this view, points are possible states of the world --- complete
descriptions of how things are or could be. Actions are then paths or
vectors, not points: an action takes the world from one state to
another.

\textbf{Strengths:} This is agent-neutral. All agents evaluate the same
space. It connects to consequentialist intuitions: what matters is the
state of the world, not who brings it about.

\textbf{Weaknesses:} The space is enormous --- potentially
infinite-dimensional. And it privileges states over processes, which
deontologists and care ethicists will resist, since the moral
significance of \emph{how} one arrives at a state is lost.

\subsubsection{Candidate 3: Situations}\label{candidate-3-situations}

A \emph{situation} is richer than a state: it includes not just how
things are, but the relationships between agents, their histories, their
commitments, and the options available. A situation is something like:
``Alice has promised Bob X. Bob is in need of Y. Alice can provide Y but
only by breaking the promise to Carol. Dana is observing.''

\textbf{Strengths:} This captures the morally relevant features without
requiring a full specification of the universe. It is closer to how
moral reasoning actually works --- what matters is the structured
configuration of agents, stakes, and relationships.

\textbf{Weaknesses:} What counts as ``morally relevant'' is
theory-dependent. Different ethical theories may carve situation-space
differently. The category of ``morally relevant features'' requires
principled delimitation.

\subsubsection{Candidate 4: Agent-Situation
Pairs}\label{candidate-4-agent-situation-pairs}

On this view, a point is a pair (a, s): an agent a in a situation s.
This allows agent-relative evaluations while maintaining a common
framework.

\textbf{Strengths:} It captures both partiality and impartiality as
different operations on the space --- holding a fixed (what should
\emph{I} do in this situation?) vs.~varying a (what should \emph{anyone}
do?). It connects naturally to the agent-indexed tensors of Chapter 6
and to the fiber bundle structure where agent-perspectives form the
fiber over a base situation.

\textbf{Weaknesses:} More complex structure; requires specifying how the
agent index interacts with the situation index.

\subsubsection{Our Choice: Structured
Situations}\label{our-choice-structured-situations}

For the purposes of this book, we take the moral manifold M to be a
space of \emph{structured situations} --- specifications of:

\begin{enumerate}
\def\labelenumi{\arabic{enumi}.}
\item
  The agents involved and their relationships
\item
  The options available to each agent
\item
  The morally relevant features of the context (needs, promises,
  histories, stakes)
\item
  The epistemic state (what is known, by whom)
\end{enumerate}

This is deliberately ecumenical. A consequentialist can project onto
states; a deontologist can focus on the structure of relationships and
commitments; a virtue ethicist can attend to the character of the
agents; a care ethicist can foreground the relational structure. The
manifold M is common ground; different theories correspond to different
tensors, metrics, and contractions \emph{on} M.

\textbf{Definition 5.1 (Moral Manifold, Informal). The moral manifold M
is a paracompact Hausdorff topological space (see Definition A.1 for the
formal construction) whose points represent structured moral situations
--- complete specifications of agents, relationships, options, and
morally relevant context.}

Agent-perspectives enter as additional structure \emph{over} M ---
specifically, as a fiber bundle whose base is the space of situations
and whose fiber at each point is the space of agent-perspectives
available in that situation. This separation of base from fiber allows
us to ask independently: ``What is the situation?'' (a question about M)
and ``Who is evaluating?'' (a question about the fiber). We develop this
in Section 5.4.

{[}Modeling choice.{]} The moral manifold M is a model, not a discovery.
Its justification is pragmatist: the nine-dimensional structure
converges across cultures and languages (the cross-lingual evidence of
§17.7 shows stable dimensional weights in 11 languages spanning 3,000
years), it arises from structural necessity (any framework satisfying
the three axioms of Chapter 9 produces an equivalent space), and it
generates testable predictions that scalar alternatives cannot (the
order-effect and semantic-gate results of §17.9--17.10). The
dimensionality --- nine, not seven or twelve --- is a modeling choice
whose adequacy is empirical. If future evidence reveals a moral
phenomenon that the nine dimensions cannot represent, the manifold must
be extended. If evidence shows that some dimensions are redundant, it
must be reduced. The modeling-choice label on Axiom 5.1 applies equally
here: the manifold is posited, not proved, and its warrant is that it
works.

\subsection{5.3 The Dimension of Moral
Space}\label{the-dimension-of-moral-space}

\subsubsection{The Dimensionality
Question}\label{the-dimensionality-question}

Having established what the points of M are, we must ask: what is the
\emph{dimension} of M? That is, how many independent directions of moral
variation exist at each point?

This is not a question the mathematics can answer by itself. It is a
substantive question about the structure of moral reasoning --- one that
admits of empirical investigation. But we can approach it
systematically.

\subsubsection{The 3 × 3 Derivation}\label{the-3-3-derivation}

A principled argument for the dimensionality of the moral manifold
proceeds from two independently motivated classifications.

\textbf{Three ethical scopes} capture whose interests are at stake:

\begin{enumerate}
\def\labelenumi{\arabic{enumi}.}
\item
  \textbf{Individual} --- the focal agent's own stakes, autonomy, and
  epistemic position
\item
  \textbf{Relational} --- the structured obligations, claims, and care
  relationships between specific parties
\item
  \textbf{Collective} --- the aggregate effects on groups, institutions,
  and the environment
\end{enumerate}

\textbf{Three epistemic modes} capture what kind of moral question is
being asked:

\begin{enumerate}
\def\labelenumi{\arabic{enumi}.}
\item
  \textbf{What Matters} --- the substantive values at stake (welfare,
  rights, justice)
\item
  \textbf{Who Decides} --- the procedural and agentive questions
  (autonomy, legitimacy, virtue)
\item
  \textbf{What We Know} --- the epistemic conditions (certainty,
  privacy, evidence)
\end{enumerate}

Crossing these two classifications yields a 3 × 3 matrix of nine cells:

{\def\LTcaptype{none} % do not increment counter
\begin{longtable}[]{@{}
  >{\raggedright\arraybackslash}p{(\linewidth - 6\tabcolsep) * \real{0.1352}}
  >{\raggedright\arraybackslash}p{(\linewidth - 6\tabcolsep) * \real{0.2709}}
  >{\raggedright\arraybackslash}p{(\linewidth - 6\tabcolsep) * \real{0.2415}}
  >{\raggedright\arraybackslash}p{(\linewidth - 6\tabcolsep) * \real{0.3524}}@{}}
\toprule\noalign{}
\begin{minipage}[b]{\linewidth}\raggedright
\end{minipage} & \begin{minipage}[b]{\linewidth}\raggedright
\textbf{What Matters}
\end{minipage} & \begin{minipage}[b]{\linewidth}\raggedright
\textbf{Who Decides}
\end{minipage} & \begin{minipage}[b]{\linewidth}\raggedright
\textbf{What We Know}
\end{minipage} \\
\midrule\noalign{}
\endhead
\bottomrule\noalign{}
\endlastfoot
\textbf{Individual} & Consequences \& Welfare & Autonomy \& Agency &
Epistemic Status \\
\textbf{Relational} & Rights \& Duties & Virtue \& Care & Privacy \&
Data Governance \\
\textbf{Collective} & Justice \& Fairness & Procedural Legitimacy &
Societal \& Environmental Impact \\
\end{longtable}
}

Each cell identifies an independent dimension of moral evaluation:

\begin{enumerate}
\def\labelenumi{\arabic{enumi}.}
\item
  \textbf{Consequences and Welfare} (Individual × What Matters): Who
  benefits, who is harmed, at what scale?
\item
  \textbf{Rights and Duties} (Relational × What Matters): What claims,
  obligations, and permissions exist?
\item
  \textbf{Justice and Fairness} (Collective × What Matters): Is the
  distribution of benefits and burdens equitable?
\item
  \textbf{Autonomy and Agency} (Individual × Who Decides): Are
  individuals free to make their own choices?
\item
  \textbf{Privacy and Data Governance} (Relational × What We Know): Is
  personal information appropriately protected?
\item
  \textbf{Societal and Environmental Impact} (Collective × What We
  Know): What are the aggregate and ecological effects?
\item
  \textbf{Virtue and Care} (Relational × Who Decides): What does this
  reflect about the agent's character and relationships?
\item
  \textbf{Procedural Legitimacy} (Collective × Who Decides): Was the
  process fair, transparent, and accountable?
\item
  \textbf{Epistemic Status} (Individual × What We Know): How confident
  are we? What don't we know?
\end{enumerate}

\subsubsection{Arguments for Necessity and
Sufficiency}\label{arguments-for-necessity-and-sufficiency}

\textbf{Necessity.} Each cell of the matrix is logically independent of
the others. Consequences can be positive while rights are violated (a
beneficent paternalism). Rights can be respected while justice is denied
(each individual treated correctly, but the aggregate distribution is
skewed). Autonomy can be preserved while epistemic conditions are poor
(free choice based on bad information). No cell is reducible to a
combination of others.

\textbf{Sufficiency.} The claim is that these nine dimensions
\emph{subsume} existing ethical frameworks. The EU High-Level Expert
Group's seven requirements for trustworthy AI map onto the nine
dimensions without residue. The Markkula Center's five ethical lenses
map onto them. Beauchamp and Childress's four principles of biomedical
ethics map onto them. This subsumption is not proof of completeness, but
it provides strong evidence that the nine dimensions capture the space
of considerations that diverse ethical traditions have independently
identified as morally relevant.

\textbf{Falsifiability.} The completeness claim is falsifiable: produce
a genuine moral consideration that is (a) independent of all nine
dimensions and (b) not a combination of them. If such a consideration
exists, the framework must be extended. This is the appropriate
epistemic stance toward a structural claim: hold it provisionally, test
it against cases, and revise if necessary.

The status of the number nine. The nine dimensions are a canonical
engineering chart, derived from the 3x3 scope/mode grid for coverage and
auditability - not a discovered empirical dimensionality. We do not
claim nine is the unique or minimal basis: a learned latent
representation, or a factor analysis of moral judgments, may recover
fewer semi-independent factors, and we expect at least six. Recovering
fewer than six would falsify the coverage claim; but the specific count
of nine is a design choice, chosen so that each dimension maps to an
established tradition - consequences and welfare (consequentialism),
rights and duties (deontology and Hohfeldian jural relations), justice
and fairness (Rawls, Sen), autonomy and agency, privacy and data,
societal and environmental impact (capability and sustainability
theories), virtue and care (virtue and care ethics), procedural
legitimacy (institutional theory), and epistemic status. Its
justification is pragmatic coverage, not metaphysical necessity; where a
tradition emphasizes a distinction the chart blurs, that is a limitation
of the chart, not a claim about morality.

\subsubsection{Formal Statement}\label{formal-statement}

\textbf{Definition 5.2 (Dimension of the Moral Manifold). The moral
manifold M has nine independent dimensions at generic (top-stratum)
points; lower-dimensional strata may have fewer active dimensions. In
local coordinates, a point p ∈ M can be parameterized by}

\[p = (x^{1},x^{2},\ldots,x^{9})\]

\emph{where x\^{}μ (μ = 1, \ldots, 9) represents the state of the
situation along the μ-th moral dimension. The tangent space T\_pM at a
generic point is a nine-dimensional real vector space.}

We say ``at least'' because domain-specific refinements may introduce
additional dimensions (the medical context adds benefit scores, waiting
times, and clinical parameters that expand the local coordinate
description). And we say ``at generic points'' because at stratum
boundaries, the effective dimension may drop --- some directions of
variation cease to be available, as when an option becomes forbidden.

\subsubsection{The Varying Metric}\label{the-varying-metric}

A crucial consequence of the nine-dimensional structure: the metric on M
is not constant. Empirical data from the Dear Abby corpus (Chapter 17)
reveals that the \emph{relative weights} of dimensions shift by context:

\begin{itemize}
\item
  In \textbf{family} contexts, social bonds (Dimension 7) dominate: the
  metric assigns greater ``length'' to variations in relational care.
\item
  In \textbf{workplace} contexts, procedural legitimacy (Dimension 8)
  and fairness (Dimension 3) dominate.
\item
  In \textbf{neighbor} disputes, fairness (Dimension 3) and rights
  (Dimension 2) dominate.
\end{itemize}

This context-dependence of the metric is precisely what differential
geometry is designed to handle. The metric tensor g\_μν(p) is a
\emph{field} on M --- it varies from point to point. The fact that it
varies is not a complication but the central phenomenon: moral space is
curved, and the curvature encodes how the structure of moral evaluation
changes across contexts.

\subsection{5.4 The Differential
Structure}\label{the-differential-structure}

\subsubsection{Charts and Atlases on M}\label{charts-and-atlases-on-m}

Recall from Chapter 4 that a smooth manifold requires coordinate charts
that overlap smoothly. For the moral manifold, a chart assigns numerical
coordinates to a neighborhood of situations.

\textbf{Example.} In a neighborhood of a family dispute about holiday
hosting, local coordinates might be:

\begin{itemize}
\item
  x¹ = degree of welfare impact on the host (Dimension 1)
\item
  x² = strength of the host's duty to extended family (Dimension 2)
\item
  x³ = fairness of the hosting distribution over years (Dimension 3)
\item
  x⁴ = host's autonomy regarding their own time (Dimension 4)
\item
  x⁷ = quality of the care relationship with family members (Dimension
  7)
\end{itemize}

Not all nine dimensions need vary independently in every neighborhood.
Some may be locally constant (privacy may be irrelevant to this
particular dispute), reducing the effective local dimension. This is the
sense in which the moral manifold may have varying dimension across its
strata --- though within a single stratum, the dimension is constant.

\subsubsection{The Tangent Space of M}\label{the-tangent-space-of-m}

At each point p ∈ M, the tangent space T\_pM has basis vectors
\{∂/∂x\^{}μ\}\_μ=1,\ldots,n where n is the local dimension. Each basis
vector represents an independent direction of moral variation:

\begin{itemize}
\item
  ∂/∂x¹ = ``increase welfare impact, holding all else fixed''
\item
  ∂/∂x² = ``strengthen the duty, holding all else fixed''
\item
  ∂/∂x³ = ``make the distribution more equitable, holding all else
  fixed''
\end{itemize}

A moral vector --- an obligation, a gradient of satisfaction, a
direction of moral improvement --- is a linear combination of these
basis directions:

\[O = O^{\mu}\frac{\partial}{\partial x^{\mu}}\]

The components O\^{}μ give the strength of the obligation \emph{in each
moral dimension}. This is the fundamental advantage over scalar ethics:
the obligation has direction, not just magnitude.

\subsubsection{The Agent Fiber Bundle}\label{the-agent-fiber-bundle}

Agent-perspectives enter as a fiber over M. At each situation p ∈ M, the
fiber F\_p consists of the available agent-perspectives in that
situation --- the physician, the patient, the family member, the
institutional review board.

\textbf{Definition 5.3 (Agent Bundle).} \emph{The agent bundle is a
fiber bundle (E, π, M) where the fiber F\_p = π\^{}\{-1\}(p) at each
point p ∈ M is the set of agent-perspectives available in situation p.}

Remark. The fiber F\_p varies with the point p: different moral
situations may involve different sets of available agent-perspectives.
Consequently, the agent bundle E is, in general, a fibered space rather
than a fiber bundle in the strict sense (the fibers need not be
diffeomorphic to a fixed model fiber F). The bundle structure is locally
trivial within each stratum --- situations in the same stratum have the
same structure of available perspectives --- but may change across
stratum boundaries.

A point of E is a pair (p, a) --- a situation and a perspective.
Agent-relative moral tensors are sections of tensor bundles constructed
from E. For instance, the agent-indexed evaluation tensor M\^{}μ\_a(p)
assigns, for each dimension μ and each agent a, an evaluation score at
situation p.

Perspective shifts (Type 2 transformations, in the language of Section
5.5) are transformations on the fiber that leave the base point fixed.
The correlative structure of Hohfeldian analysis --- that Morgan's
obligation implies Alex's claim --- is a constraint relating different
points on the same fiber: if M\^{}2\_Morgan = O (obligation), then
M\^{}2\_Alex = C (claim). This constraint is a \emph{structural
invariant} of the fiber, not a consequence of any particular evaluation.

\subsection{5.5 Coordinates and Admissible
Transformations}\label{coordinates-and-admissible-transformations}

\subsubsection{The Problem of Multiple
Descriptions}\label{the-problem-of-multiple-descriptions}

In physics, coordinate transformations have a precise meaning: a change
from one coordinate chart to another covering the same region of
manifold. The transformation laws for tensors (Chapter 4, Section 4.4)
specify exactly how components change.

In ethics, we speak loosely of ``different perspectives,'' ``different
framings,'' ``different descriptions.'' But not all of these are
coordinate transformations in the geometric sense. Conflating them
invites the criticism that geometric ethics is vacuous --- that by
choosing the right ``coordinate system,'' one can make any answer come
out.

This section introduces a discipline. We distinguish three types of
transformation, with progressively weaker geometric status.

\subsubsection{Type 1: Coordinate
Redescriptions}\label{type-1-coordinate-redescriptions}

A \emph{coordinate redescription} is a change in how we parameterize the
same underlying situation. The situation itself --- the point p ∈ M ---
is unchanged; only the coordinates (the labels, the representation)
differ.

\textbf{Example.} Describing a resource allocation in terms of ``amount
to Alice'' (x\_A) vs.~``amount to Bob'' (x\_B = T - x\_A), where the
total T is fixed. These are different coordinates on the same
one-dimensional manifold.

\textbf{Example.} Describing an action as ``withholding treatment''
vs.~``allowing natural death.'' If these are genuinely synonymous --- if
they pick out exactly the same action in all morally relevant respects
--- then they are coordinate redescriptions.

\textbf{The tensorial requirement.} Any moral quantity that is a genuine
feature of the situation must transform appropriately under coordinate
redescriptions. Scalars are invariant; vectors transform by the
Jacobian; rank-2 tensors transform by two copies of the Jacobian. This
is the mathematical content of the Bond Invariance Principle:

\textbf{Axiom 5.1 (Bond Invariance Principle). {[}Definition/Modeling
choice.{]} If two descriptions d and d' of a situation are related by an
admissible coordinate transformation, then any legitimate moral
evaluation must assign the same value to both: E(d) = E(d').
Equivalently, legitimate moral evaluations are gauge-invariant
quantities on the moral manifold.}

What this rules out: moral evaluations that depend on \emph{mere
labeling}. If ``allowing to die'' sounds better than ``withholding
treatment'' but they denote the same action, a proper moral evaluation
must not distinguish them. This is the ethical analogue of general
covariance in physics.

What this does \emph{not} rule out: evaluations that differ because the
descriptions are \emph{genuinely different} --- because they pick out
different features, different contexts, or different aspects of the
situation. The framework demands invariance under re-description, not
under re-situation.

In the language of group actions, BIP states that the moral evaluation
functional E is constant on orbits of the admissible re-description
group: E(g · d) = E(d) for all g ∈ G\_adm and all descriptions d of a
situation p ∈ M. Here g · d denotes the action of the re-description
group element g on the description d (the fiber over p in the
description bundle of §5.4), and G\_adm is the group of admissible
re-descriptions --- the subgroup of Diff(fiber) satisfying the five
admissibility conditions of §5.5. Chapter 12 will identify G\_adm = D₄ ×
U(1)\_H as the maximal such group consistent with axioms A1--A5. The BIP
is thus a precise orbit-invariance condition, and the tensor
transformation laws of §4.4 are the infinitesimal manifestation of this
group action on the moral evaluation algebra.

\subsubsection{Pre-flight Checklist: Bond Invariance
Principle}\label{pre-flight-checklist-bond-invariance-principle}

Before applying BIP, verify each item:

\textbf{Admissibility verified.} Both descriptions pass all five
conditions (§5.5): same physical grounding, same agent structure, same
stratum, reversible, no evaluative loading.

\textbf{Manifold located.} Situation has explicit coordinates on nine
dimensions and an explicit stratum.

\textbf{Tensor rank declared.} The moral quantity has a declared type;
BIP requires appropriate transformation.

\textbf{Gauge group scope.} Transformation is within G\_ethics = D₄ ×
U(1)\_H (Thm 12.3).

\textbf{Grounding adequate (AI).} Ψ satisfies six adequacy conditions
(Def. 18.1).

\subsubsection{Type 2: Perspective
Shifts}\label{type-2-perspective-shifts}

A \emph{perspective shift} changes the evaluating agent while holding
the situation fixed. This is not a coordinate change on M --- it is a
transformation on the fiber of the agent bundle (Section 5.4).

\textbf{Example.} Evaluating a kidney allocation from the physician's
perspective vs.~a patient's family's perspective. The situation (who the
patients are, what their conditions are) is the same. What changes is
\emph{who is evaluating}.

\textbf{Tensorial treatment.} Perspective shifts act on the agent index
of agent-relative tensors. The rank-2 tensor M\^{}μ\_a (evaluation of
dimension μ by agent a) transforms in the agent index a when we change
perspectives. The framework makes no general demand that all
perspectives agree --- that would be an implausibly strong objectivity
requirement. Instead, it demands that certain \emph{structural} features
are perspective-invariant:

\begin{itemize}
\item
  The correlative structure: if one perspective assigns obligation, the
  correlative perspective must assign claim.
\item
  The constraint set: what is forbidden is forbidden from all
  perspectives.
\item
  The stratum structure: which moral regime applies is not
  perspective-dependent.
\end{itemize}

These invariants characterize the \emph{objective} structure of the
moral situation, while the \emph{perspectival} variation (which
dimensions each agent prioritizes, which trade-offs each is willing to
make) is the legitimate space of moral disagreement.

\subsubsection{Type 3: Theory Shifts}\label{type-3-theory-shifts}

A \emph{theory shift} changes the mathematical structure we impose on M:
the metric, the connection, the constraint set, the admissible
transformations themselves.

\textbf{Example.} Switching from a utilitarian metric (all dimensions
commensurable, with trade-off ratios specified by the metric's
off-diagonal components) to a lexicographic metric (Dimension 3 has
absolute priority; other dimensions matter only when Dimension 3 is
tied). This is not a coordinate change; it is a change in the
\emph{geometry} of moral space.

\textbf{Example.} Switching from a theory that treats agent identity as
morally irrelevant (strict impartialism) to one that gives weight to
special relationships (moderate partiality). This changes which
fiber-bundle structures are admitted.

\textbf{Tensorial treatment.} Theory shifts are \emph{not} symmetries
--- they change the structure. Different theories correspond to
different choices of metric g, connection ∇, constraint set C, and
admissibility conditions. The geometric framework does not adjudicate
between theories; it represents each theory precisely, making
commitments explicit and disagreements localizable.

\subsubsection{Summary: The Transformation
Hierarchy}\label{summary-the-transformation-hierarchy}

{\def\LTcaptype{none} % do not increment counter
\begin{longtable}[]{@{}
  >{\raggedright\arraybackslash}p{(\linewidth - 6\tabcolsep) * \real{0.2076}}
  >{\raggedright\arraybackslash}p{(\linewidth - 6\tabcolsep) * \real{0.2552}}
  >{\raggedright\arraybackslash}p{(\linewidth - 6\tabcolsep) * \real{0.2053}}
  >{\raggedright\arraybackslash}p{(\linewidth - 6\tabcolsep) * \real{0.3319}}@{}}
\toprule\noalign{}
\begin{minipage}[b]{\linewidth}\raggedright
\textbf{Type}
\end{minipage} & \begin{minipage}[b]{\linewidth}\raggedright
\textbf{What Changes}
\end{minipage} & \begin{minipage}[b]{\linewidth}\raggedright
\textbf{Geometric Status}
\end{minipage} & \begin{minipage}[b]{\linewidth}\raggedright
\textbf{Invariance Requirement}
\end{minipage} \\
\midrule\noalign{}
\endhead
\bottomrule\noalign{}
\endlastfoot
Coordinate redescription & Labels, parameterization & Gauge symmetry &
Tensors must transform appropriately \\
Perspective shift & Evaluating agent & Fiber transformation & Structural
invariants preserved \\
Theory shift & Metric, connection, constraints & Changes the geometry &
No invariance; theories are compared, not equated \\
\end{longtable}
}

\textbf{The discipline:} When we say ``moral claims should be invariant
under redescription,'' we mean Type 1 transformations. Types 2 and 3 are
not symmetries in this sense --- they are legitimate sources of
variation that the framework makes explicit rather than concealing.

\subsubsection{Admissibility in Practice: A
Checklist}\label{admissibility-in-practice-a-checklist}

The BIP's power rests on the boundary of ``admissible.'' A skeptic may
object: \emph{if you control the equivalence relation, invariance is
guaranteed trivially}. The objection is correct in principle---which is
why the boundary must be operationally defined and independently
testable. Given two descriptions \(d\) and \(d'\) of what is claimed to
be the ``same'' moral situation, the following checklist determines
whether \(d \rightarrow d'\) is an admissible (Type 1) transformation.
\emph{All} conditions must be met.

\begin{enumerate}
\def\labelenumi{\arabic{enumi}.}
\item
  \textbf{Same physical grounding.} Would a sensor suite measuring
  welfare indicators, bodily-autonomy indicators, resource
  distributions, and other morally relevant observables produce the
  \emph{same readings} under both descriptions? If not, the descriptions
  pick out different situations, and the transformation is not
  admissible.
\item
  \textbf{Same agent structure.} Do \(d\) and \(d'\) identify the same
  agents, the same relationships, and the same power asymmetries? If a
  description obscures an agent (``mistakes were made'' vs.~``the CEO
  fired the whistleblower''), the transformation changes the agent index
  and is Type 2, not Type 1.
\item
  \textbf{Same stratum.} Do \(d\) and \(d'\) place the situation in the
  same Hohfeldian stratum? If ``allowing to die'' evokes Liberty while
  ``withholding treatment'' evokes Obligation, the descriptions encode
  different moral regimes. This is detectable by checking whether the
  semantic-gate structure differs under the two descriptions.
\item
  \textbf{Reversibility.} Can a competent moral reasoner, given \(d'\),
  reconstruct \(d\) without additional information? If the
  transformation is lossy---if it obscures morally relevant
  features---it is a projection, not a coordinate change. Projections
  are not admissible.
\item
  \textbf{No evaluative loading.} Does the transformation change the
  evaluative valence? ``Freedom fighter'' vs.~``terrorist'' is not a
  coordinate change; it smuggles in an evaluation along the
  normative-mode dimensions. Admissible transformations change
  coordinate labels, not evaluative content.
\end{enumerate}

If all five conditions are met, the transformation is admissible (Type
1) and the BIP demands invariance. If any condition fails, the
transformation is Type 2 or Type 3, and variance between the two
evaluations is legitimate.

\subsubsection{Tempting but Invalid
Transformations}\label{tempting-but-invalid-transformations}

The following examples illustrate transformations that \emph{appear} to
be mere redescriptions but fail the admissibility checklist. Each is a
case where naive application of the BIP would be a mistake.

\textbf{1. ``Killing'' vs.~``allowing to die.''} These are often treated
as synonyms describing the same physical outcome (a death). But they
differ on the agent-structure criterion: \emph{killing} requires an
active causal agent; \emph{allowing to die} involves the absence of
intervention. They also frequently differ on stratum: killing may
trigger the absorbing stratum for direct harm (Chapter 8, Type III
boundary), while allowing to die may remain in the Liberty stratum.
\emph{Checklist verdict:} Fails criteria 2 and 3. Classification:
distinct situations. BIP does not apply.

\textbf{2. ``Taxation'' vs.~``theft.''} Both involve involuntary
transfer of resources, and a libertarian critic may claim they are
equivalent. But they differ on the procedural-legitimacy dimension
(\(x^{8}\))---taxation operates through institutional authority; theft
does not---and on evaluative loading: ``theft'' encodes a negative
judgment that ``taxation'' does not. \emph{Checklist verdict:} Fails
criteria 2 and 5. Classification: Type 3 (theory shift). BIP does not
apply.

\textbf{3. ``Surveillance'' vs.~``monitoring for safety.''} Both
describe the same physical activity (cameras recording people). The
physical grounding may be identical. But ``surveillance'' carries
negative evaluative loading (privacy invasion), while ``monitoring for
safety'' carries positive loading (protective care). The transformation
smuggles evaluation into the description. \emph{Checklist verdict:}
Fails criterion 5. Classification: Type 2 (perspective shift with
evaluative content). BIP does not apply; the variance between the two
framings is a legitimate signal of moral disagreement.

\textbf{4. ``He said / she said'' perspective swaps.} Describing a
dispute from either party's perspective may seem like a mere coordinate
change. But it shifts the evaluating agent (Type 2), not the coordinate
system: the power asymmetries, epistemic positions, and relational roles
are redistributed. \emph{Checklist verdict:} Fails criterion 2.
Classification: Type 2. Variance is expected and informative.

\textbf{5. ``Economic growth'' vs.~``ecological destruction.''} Both may
describe the same industrial activity. But they foreground different
dimensions of the moral manifold (welfare \(x^{1}\)
vs.~societal/environmental impact \(x^{6}\)), use different physical
grounding (GDP metrics vs.~ecosystem metrics), and carry evaluative
loading in opposite directions. \emph{Checklist verdict:} Fails criteria
1, 2, and 5. Classification: distinct framings invoking different
regions of the moral manifold. The BIP correctly predicts that these
should \emph{not} yield the same evaluation.

These examples demonstrate that the admissibility boundary is not a free
parameter that can be tuned to guarantee invariance. It is an
operationally testable criterion rooted in the five-condition checklist.
Failures of admissibility are \emph{detectable}---they show up as
violations of the agent-structure, stratum, or evaluative-loading
conditions---and they are \emph{informative}: they reveal that what
appeared to be a mere redescription is actually a substantive moral
reframing.

An Edge Case: Genuine Ambiguity

The five invalid transformations above were designed to fail clearly.
But the checklist's greatest value emerges when the verdict is
ambiguous---when reasonable people might disagree about whether a
transformation is admissible. Consider:

\textbf{``Euthanasia'' vs. ``mercy killing'' vs. ``assisted dying.''}
All three describe a physician administering a lethal dose to a
terminally ill, consenting patient.

1. Same physical grounding. Yes---the same patient, the same medication,
the same outcome. ✓

2. Same agent structure. Yes---physician, patient, and family occupy the
same roles. ✓

3. Same stratum. Ambiguous. ``Mercy killing'' may trigger the absorbing
stratum for direct killing (§8.5), while ``assisted dying'' may remain
in the obligation stratum (medical duty of compassion). ``Euthanasia''
is etymologically neutral (``good death'') but is culturally coded as a
boundary term. This is where the checklist escalates.

4. Reversibility. Marginal---a competent reasoner can reconstruct any of
the three from any other, but connotative weight is lost in translation.

5. No evaluative loading. ``Mercy killing'' loads positively (``mercy'')
and negatively (``killing'') simultaneously. ``Assisted dying'' loads
procedurally (``assisted''). ``Euthanasia'' is closest to neutral.

Checklist verdict. Criteria 3 and 5 fail for ``mercy killing'' → Type 2
or 3. Criterion 3 is ambiguous for ``euthanasia'' vs. ``assisted dying''
→ the transformation's admissibility depends on a governance decision
about stratum assignment (§9.1).

The lesson. When the checklist produces ambiguity rather than a clear
verdict, it has done its job. It has identified the precise locus of
genuine moral disagreement---the stratum boundary at which ``medical
compassion'' meets ``taking a life''---rather than papering over it with
a false equivalence. The equivalence relation ≈ is not subjective; it is
decidable up to a governance parameter (the stratum assignment), and the
checklist tells you exactly which parameter requires a decision.

\subsection{5.6 Stratification}\label{stratification}

\subsubsection{Why Stratification Is
Essential}\label{why-stratification-is-essential}

The most important structural feature of the moral manifold is not its
smooth geometry but its \emph{stratification}: the partition of M into
regions (strata) where different rules apply, joined along boundaries
where the rules change discontinuously.

Smooth manifolds, by definition, are uniform in dimension: every point
has a neighborhood homeomorphic to ℝⁿ for the same n.~But moral space is
not uniform. There are regions of smooth trade-offs (within a stratum),
hard boundaries (between strata), and singular points (where strata
intersect pathologically). The moral manifold is not a smooth manifold
but a \emph{stratified space} in the sense of Whitney (Chapter 4,
Section 4.8).

\subsubsection{Strata, Boundaries, and
Transitions}\label{strata-boundaries-and-transitions}

\textbf{Definition 5.4 (Moral Stratification).} *The moral manifold M
admits a Whitney stratification \{S\_α\}\_\{α ∈ A\} where: \emph{1.}
Each stratum S\_α is a smooth manifold (possibly of different
dimension). \emph{2.} Within each stratum, the moral evaluation function
S and the moral metric g are smooth. \emph{3.} At stratum boundaries, S
or g may be discontinuous, or the effective dimension may change.*

\textbf{Threshold effects.} Within the ``consent'' stratum, the degree
of informed consent varies smoothly --- a little more information, a
little more understanding. But at the boundary between ``sufficient
consent'' and ``insufficient consent,'' the moral evaluation jumps
discontinuously. Below the threshold, the action may be impermissible;
above it, permissible. This is a stratum boundary.

\textbf{Categorical distinctions.} The difference between a 17-year-old
and an 18-year-old, between killing and letting die, between lying and
remaining silent --- these are boundaries where legal and moral rules
change. Within each side of the boundary, smooth variation is possible.
Across the boundary, the applicable framework shifts.

\textbf{The forbidden region.} Certain options are not merely low-value
but \emph{excluded from consideration} --- moral non-starters regardless
of other considerations.

\textbf{Definition 5.5 (Constraint Set). *The constraint set C ⊂ M is a
closed subset such that every satisfaction function S compatible with
the stratification: M → ℝ ∪ \{-∞\} satisfies S\_C = -∞. The boundary ∂C
is a stratum boundary.*}

Notation remark. The symbol S is used in two distinct senses throughout
this work: (i) S: M → ℝ ∪ \{−∞\} denotes the satisfaction function (a
scalar field on M), and (ii) S\_α denotes strata of the Whitney
stratification (subsets of M). Context always disambiguates: S with a
subscript α ∈ A (the index set) refers to a stratum, while S alone or
S(p) refers to the satisfaction scalar. Where both appear in the same
expression, the satisfaction function is written S(p) or as the
contraction I\_μ O\^{}μ to avoid ambiguity.

The constraint set represents absolute prohibitions: actions involving
non-consensual harm to innocents, violations of inviolable rights, and
the like. The framework does not determine \emph{what} belongs in C ---
that is the work of normative ethics and democratic governance. But it
represents the \emph{structure} of constraints precisely: as a region
with hard boundaries, distinguished from regions of trade-offs.

\subsubsection{Nullifiers as Absorbing
States}\label{nullifiers-as-absorbing-states}

Empirical evidence (Chapter 17) reveals a special class of stratum
boundary: \emph{nullifiers} --- conditions that override all other
considerations and collapse obligations to zero.

The Dear Abby corpus (Bond 2026) identifies abuse as the primary
nullifier: across all domain contexts (family, workplace, friendship,
romance), the presence of abuse universally nullifies obligations (n =
582). Other nullifiers include danger (n = 218) and impossibility (n =
144, encoding the principle \emph{ought implies can}).

In stratified-space language, nullifiers define \emph{absorbing strata}
--- lower-dimensional strata from which there is no transition back.
Once a situation enters the ``abuse'' stratum, the normal obligation
structure collapses: the metric becomes degenerate along the obligation
dimensions, and the previously operative duties cease to apply.

\textbf{Definition 5.6 (Absorbing Stratum).} \emph{A stratum S\_0 ⊂ M is
absorbing if, for any path γ entering S\_0, the moral evaluation along γ
is determined by S\_0's rules thereafter, regardless of the prior
stratum's structure.}

\subsubsection{Semantic Gates as Transition
Functions}\label{semantic-gates-as-transition-functions}

The transitions between strata are not arbitrary. They are triggered by
specific features of the situation --- what we call \emph{semantic
gates}, borrowing the language of quantum normative dynamics.

\textbf{Example.} The phrase ``you promised'' triggers a transition from
the Liberty stratum to the Obligation stratum. The phrase ``only if
convenient'' triggers the reverse transition, from Obligation to
Liberty. These transitions are \emph{discrete} --- step functions, not
sigmoids --- and empirically measurable with high reliability (94\% for
``you promised,'' 89\% for ``only if convenient'').

In the geometry of M, a semantic gate is a codimension-1 submanifold (a
``wall'' between strata) equipped with a transition function specifying
how the moral evaluation changes upon crossing. The discreteness of
these transitions --- verified empirically --- supports the
stratified-space model over a purely smooth manifold.

\subsection{5.7 Singularities and
Dilemmas}\label{singularities-and-dilemmas}

A \emph{singularity} in M is a point where the normal geometric
structure breaks down. In general relativity, singularities are points
of infinite density or undefined curvature --- places where the theory
reaches its limits. In geometric ethics, they are \emph{genuine moral
dilemmas}.

\subsubsection{Definition and Examples}\label{definition-and-examples}

\textbf{Definition 5.7 (Moral Singularity).} \emph{A point p ∈ M is a
moral singularity if:} 1. \emph{The metric tensor g\_p is degenerate
(det(g\_p) = 0), so that some directions of moral variation have zero
``weight'' and comparisons break down; or} 2. \emph{The satisfaction
function S is not differentiable at p, so that the gradient ∇S --- the
``direction of moral improvement'' --- is undefined; or} 3.
\emph{Multiple constraint surfaces intersect at p, creating a cone of
forbidden directions that leaves no permissible path.}

The three singularity types are: (i) metric singularities, where the
moral metric degenerates and dimensional comparisons fail; (ii) gradient
singularities, where the satisfaction function loses differentiability
and the direction of moral improvement is undefined; and (iii)
constraint singularities, where multiple prohibition surfaces intersect
and the feasible region collapses. The singular set Σ ⊂ M is their
union: Σ = Σ\_metric ∪ Σ\_gradient ∪ Σ\_constraint. Each type has
distinct geometric character: metric singularities are intrinsic
(detectable from g alone), gradient singularities depend on S, and
constraint singularities depend on the constraint set C.

\textbf{Example (Sophie's Choice).} A mother must choose which of her
two children will be killed. There is no ``right'' answer. Any choice
involves irreducible moral loss. The situation is singular: the
satisfaction function S has no global maximum in the feasible region, or
has multiple local maxima that are genuinely incomparable --- the metric
between them is degenerate.

\textbf{Example (Competing absolute obligations).} A physician has
promised confidentiality to a patient who reveals an intent to harm a
third party. The obligation of confidentiality and the obligation to
prevent harm are both operative, both point in opposite directions, and
neither can be subordinated to the other without remainder. The gradient
∇S is undefined: there is no single ``direction of improvement.''

\subsubsection{Dilemmas as Features, Not
Bugs}\label{dilemmas-as-features-not-bugs}

Singularities are not failures of the framework. They are
\emph{features}. The framework \emph{represents} genuine dilemmas as
singular points, rather than forcing a false determinacy. This is an
advance over scalar approaches, which must either deliver an answer (by
arbitrary tie-breaking) or fall silent (by declaring the situation
beyond their scope).

The geometric representation is more precise: we can say \emph{why} the
dilemma is a dilemma --- which constraint surfaces intersect, which
metric components become degenerate, which dimensions of the moral
tensor conflict. We can characterize the \emph{type} of singularity
(metric degeneracy vs.~gradient failure vs.~constraint intersection) and
the \emph{moral residue} that any resolution leaves behind (Chapter 15).

\subsection{5.8 Topology of the Moral
Manifold}\label{topology-of-the-moral-manifold}

\subsubsection{Local and Global
Structure}\label{local-and-global-structure}

The differential structure developed so far is \emph{local}: it
describes the neighborhood of a point. The \emph{topology} of M concerns
its global properties --- properties that are invariant under continuous
deformation but may be invisible from any single neighborhood.

\textbf{Connectedness.} Can any two moral situations be reached from
each other by a continuous path through M? Or are there disconnected
components --- perhaps corresponding to genuinely incommensurable moral
frameworks, or to domains of life so different that no smooth
interpolation between them exists?

The evidence suggests that M is connected within a single
cultural-historical context: any two everyday moral situations can be
related by a sequence of small modifications. But whether M is connected
\emph{across} cultures and historical epochs is a substantive question.
The temporal stability data from the Dear Abby corpus (Chapter 17)
suggests that the core topology of M is stable over at least three
decades, though the metric varies.

\textbf{Compactness.} Is M compact (closed and bounded) or does it
extend indefinitely? The constraint set C introduces boundaries, but
does M have a ``ceiling'' --- a maximum possible stakes, a highest
conceivable obligation?

In practice, moral evaluation operates on a bounded region: the stakes
in any real situation are finite, the number of agents is finite, and
the morally relevant features take values in bounded ranges. We can
model the relevant portion of M as compact without great loss.

\textbf{Fundamental group.} Does M have ``holes'' --- non-contractible
loops? A nontrivial fundamental group π₁(M) would mean that certain
closed paths in moral space cannot be continuously shrunk to a point.
This would have implications for holonomy (Chapter 10): transporting a
moral evaluation around such a loop could yield a nontrivial change,
even though the loop returns to the same situation.

Whether M has nontrivial topology is among the most intriguing open
questions of the framework. The path-dependence of obligation (Chapter
10) provides \emph{indirect} evidence: if obligations change when
transported around loops, the manifold may have nontrivial holonomy,
which constrains the topology. But settling this definitively requires
the empirical program outlined in Chapter 29.

\subsection{5.9 Two Worked Examples}\label{two-worked-examples}

\subsubsection{Example 1: The Manifold of a Medical
Allocation}\label{example-1-the-manifold-of-a-medical-allocation}

\textbf{Situation.} A physician must allocate a scarce treatment among
three patients (A, B, C). Each patient has a medical benefit score β\_i
∈ {[}0, 1{]}, a waiting time w\_i ∈ {[}0, ∞), an age a\_i, and a number
of dependents d\_i.

\textbf{The manifold.} The space of possible allocations is the
2-simplex:

\[M = \Delta^{2} = \{(p_{A},p_{B},p_{C}):p_{i} \geq 0,\mspace{6mu}\sum_{i}^{}p_{i} = 1\}\]

where p\_i is the fraction of treatment allocated to patient i.

\textbf{Stratification:} - \emph{Interior} S₂ (dim 2): All patients have
positive probability. Smooth trade-offs possible. - \emph{Edges} S₁ (dim
1): Two patients share the treatment; one is excluded. Three edges. -
\emph{Vertices} S₀ (dim 0): Deterministic allocations. Three vertices.
These are the \emph{actual decisions}; the interior represents
deliberation space.

This is a Whitney-stratified space with three levels. The frontier
condition holds: each vertex is in the closure of each adjacent edge,
and each edge is in the closure of the interior.

\textbf{Constraint region.} Suppose Patient C's inclusion would
constitute discrimination (the decision is based on a protected
characteristic). Then the region p\_C \textgreater{} 0 is forbidden:

\[C = \{(p_{A},p_{B},p_{C}) \in \Delta^{2}:p_{C} > 0\}\]

The feasible region is the edge from vertex A to vertex B --- a
1-dimensional manifold.

\textbf{Metrics.} Different metrics on Δ² correspond to different
ethical theories: - \emph{Euclidean metric:} g = dp\_A ⊗ dp\_A + dp\_B ⊗
dp\_B. All movements equally costly. - \emph{Weighted metric:} g\_\{ij\}
= δ\_\{ij\}/β\_i. Movement toward sicker patients (higher β) is
``easier'' --- lower moral cost. - \emph{Lexicographic metric:} Not a
metric in the standard sense, but a limiting case: the allocation must
first maximize benefit to the worst-off, then consider secondary
criteria. This corresponds to a degenerate metric where the worst-off
dimension has infinite weight.

\textbf{Satisfaction function.} S: Δ² → ℝ assigns a moral score to each
allocation. On the interior, S is smooth. On the constraint boundary, S
= -∞. The gradient ∇S, computed via the metric, points toward the
morally preferred allocation.

\subsubsection{Example 2: The Promise
Landscape}\label{example-2-the-promise-landscape}

\textbf{Situation.} Morgan has promised to help Alex move on Saturday.
Various complications arise.

\textbf{The manifold.} The relevant region of M is parameterized by (at
least) three dimensions:

\begin{itemize}
\item
  x² = strength of Morgan's obligation (Dimension 2: Rights and Duties)
\item
  x⁴ = cost to Morgan's autonomy (Dimension 4: Autonomy)
\item
  x⁷ = quality of the Morgan-Alex relationship (Dimension 7: Care)
\end{itemize}

\textbf{Strata.} The promise creates two primary strata:

\begin{itemize}
\item
  \emph{Obligation stratum} S\_O: The region where Morgan has an
  obligation to help. The promise is binding, the cost is manageable,
  and no nullifiers apply.
\item
  \emph{Liberty stratum} S\_L: The region where Morgan is free to
  decline. Either the promise was conditional, or a nullifier applies,
  or the cost has become prohibitive.
\end{itemize}

\textbf{Semantic gates.} The boundary between S\_O and S\_L is crossed
by specific triggers:

\begin{itemize}
\item
  ``Only if convenient'' → O → L transition. The obligation collapses.
  The Dear Abby corpus measures this gate at 89\% effectiveness.
\item
  ``You explicitly promised'' → L → O transition. The obligation
  crystallizes. Measured at 94\% effectiveness.
\item
  Abuse by Alex → O → L (nullifier). Universal across contexts. The
  stratum boundary is absorbing.
\end{itemize}

\textbf{Path-dependence.} Consider two paths from an initial state
(Morgan made the promise) to a final state (Morgan is aware that Alex
was rude last week):

\begin{itemize}
\item
  Path α: Morgan learns Alex was rude, then recalls the explicit
  promise.
\item
  Path β: Morgan recalls the explicit promise, then learns Alex was
  rude.
\end{itemize}

If the final obligation differs depending on which information came
first, the moral manifold has \emph{curvature} in this region. The
non-commutativity of these information-updates --- learning about
rudeness and recalling the promise do not commute --- is precisely the
non-abelian structure formalized by the D₄ dihedral group (Chapter 12).

\subsection{5.10 The Manifold as Common
Ground}\label{the-manifold-as-common-ground}

The moral manifold M is the base space over which all ethical tensors
are defined. Its points are structured situations; its structure
includes a nine-dimensional tangent space (at generic points), a
context-varying metric, a Whitney stratification into moral regimes, and
singular points representing genuine dilemmas.

Different ethical theories correspond to different structures \emph{on}
M:

\begin{itemize}
\item
  Different metrics (utilitarian vs.~egalitarian vs.~lexicographic)
\item
  Different constraint sets (what is absolutely forbidden)
\item
  Different contractions (how multi-dimensional evaluation reduces to
  action)
\item
  Different connections (how obligations are transported across
  contexts)
\end{itemize}

But all theories share M as common ground. This makes disagreement
\emph{tractable}: we can ask whether two theories differ in their
metrics, their constraints, or their contractions. We can identify where
they agree (perhaps on the constraint set --- both agree that torture is
forbidden) and where they diverge (perhaps on the metric --- they
disagree about the trade-off rate between welfare and equality).
Disagreement, in this framework, has \emph{address}.

The moral manifold is not the whole of ethics. It is the \emph{stage} on
which ethics plays out --- the space of possibilities that moral
reasoning navigates. The next chapter develops the actors: the tensors
of various ranks that live on M --- obligations as vectors, interests as
covectors, the moral metric, and the fundamental contraction S = I\_μ
O\^{}μ that yields satisfaction.

But without the stage, there is nowhere for the actors to stand.

\subsection{Technical Appendix: Formal
Definitions}\label{technical-appendix-formal-definitions}

\textbf{Definition A.1 (Moral Manifold, Formal).} \emph{A moral manifold
is a paracompact Hausdorff topological space M equipped with:} 1. *A
Whitney stratification \{S\_α\}\_\{α ∈ A\} into smooth manifolds of
(possibly) varying dimensions; \emph{2.} A partial order ⪯ on A
satisfying the frontier condition: S\_α ∩ cl(S\_β) ≠ ∅ implies α ⪯ β;
\emph{3.} Whitney's condition (B) at all stratum boundaries; \emph{4.} A
smooth structure on each stratum S\_α making it a smooth manifold of
dimension d\_α.*

\textbf{Definition A.2 (Coordinate Chart on M).} \emph{A coordinate
chart on a stratum S\_α is a diffeomorphism φ: U → V from an open set U
⊂ S\_α to an open set V ⊂ ℝ\^{}\{d\_α\}. A coordinate transformation is
a composition φ' ∘ φ\^{}\{-1\}: ℝ\^{}\{d\_α\} → ℝ\^{}\{d\_α\}, required
to be a diffeomorphism.}

\textbf{Definition A.3 (Admissible Transformation). An admissible
transformation of the moral manifold is a stratification-preserving
homeomorphism (restricting to a diffeomorphism on each stratum interior)
ψ: M → M satisfying ψ(S\_α) ⊂ S\_\{σ(α)\} for some bijective,
order-preserving map σ of the index set A. Type 1 transformations are
admissible; Type 3 transformations are not.}

\textbf{Definition A.4 (Agent Bundle). The agent bundle over M is a
fiber bundle (E, π, M) where: 1. The fiber F\_p = π\^{}\{-1\}(p) at each
p ∈ M is the set of agent-perspectives available in situation p; 2. The
structure group G acts on fibers by perspective transformations; 3.
Correlative constraints (O\_a ↔ C\_b for all Hohfeldian correlative
pairs (a, b) at p) are structural invariants of the G-action.}

\textbf{Definition A.5 (Moral Singularity).} \emph{A point p ∈ M is a
moral singularity if:} 1. \emph{The metric tensor g\_p is degenerate:
det(g\_p) = 0; or} 2. \emph{The satisfaction function S is not
differentiable at p; or} 3. \emph{Multiple constraint surfaces of
codimension 1 intersect transversally at p, and the feasible cone at p
has no interior.}

\emph{Singularities represent genuine dilemmas: points where the moral
structure does not determine a unique best response.}

\textbf{Definition A.6 (Absorbing Stratum).} *A stratum S\_0 is
absorbing if there exists a continuous retraction r: N → S\_0 from a
tubular neighborhood N of S\_0 such that, for any moral evaluation
function S compatible with the stratification, S\_N factors through r:
S(p) = S(r(p)) for all p ∈ N.* Additionally, S₀ is a trapping region:
for any continuous path γ with γ(t₀) ∈ S₀, we have γ(t) ∈ N ∪ S₀ for all
t ≥ t₀. This trapping condition reconciles the formal definition with
the informal path-based characterization (Definition 5.6): once a moral
trajectory enters an absorbing stratum, the evaluation is governed by
that stratum\textquotesingle s rules thereafter.

❖

\emph{The manifold is the ground. The tensors are the figures.}

\emph{Before we can say what obligation points toward, we must know the
space in which it points.}

\emph{This chapter has constructed that space: M, the moral manifold ---
nine-dimensional, stratified, equipped with a context-varying metric,
and rich enough to serve as common ground for every ethical theory that
has a claim to completeness.}

\section{Chapter 6: The Tensor
Hierarchy}\label{chapter-6-the-tensor-hierarchy}

\begin{quote}
\textbf{\textsc{RUNNING EXAMPLE --- Priya's Model}}

\emph{Priya now thinks in tensors. Her TrialMatch score S is rank-0: a
scalar. Patient needs form an obligation vector Oᵋ in nine dimensions.
What patients want from the trial---hope, knowledge, dignity,
access---forms an interest covector Iₘ. The contraction S = Iₘ Oᵋ would
give a meaningful scalar, but only after respecting all nine dimensions.
Her model never builds Oᵋ or Iₘ. It goes straight to S using a neural
network whose internal representations have no guaranteed correspondence
to moral dimensions. She writes in her notebook: `My model doesn't
contract a tensor. It guesses the contraction without ever forming the
tensor. That's why it can't see the structure it's losing.'}
\end{quote}

\emph{Chapter 5 built the stage --- the moral manifold M, a
nine-dimensional stratified space of structured moral situations. This
chapter populates the stage with actors: the tensors of various ranks
that encode moral evaluation. We develop obligations as vectors,
interests as covectors, the moral metric as the structure of permissible
trade-offs, and the fundamental formula S = I\_μ O\^{}μ in which
satisfaction is the contraction of interest with obligation.}

\subsection{6.1 From Scalars to Tensors: Why Rank
Matters}\label{from-scalars-to-tensors-why-rank-matters}

Recall the argument of Chapter 2: scalar moral evaluation loses
information. A single number S ∈ ℝ assigned to a situation p ∈ M tells
us \emph{how good} the situation is, but not \emph{along which
dimensions} it is good, not \emph{in which direction} improvement lies,
and not \emph{how robust} the assessment is to changes in perspective or
weighting.

The remedy is to replace the scalar with objects of successively higher
rank, each encoding more of the structure of moral evaluation:

{\def\LTcaptype{none} % do not increment counter
\begin{longtable}[]{@{}
  >{\raggedright\arraybackslash}p{(\linewidth - 6\tabcolsep) * \real{0.0812}}
  >{\raggedright\arraybackslash}p{(\linewidth - 6\tabcolsep) * \real{0.2183}}
  >{\raggedright\arraybackslash}p{(\linewidth - 6\tabcolsep) * \real{0.3879}}
  >{\raggedright\arraybackslash}p{(\linewidth - 6\tabcolsep) * \real{0.3126}}@{}}
\toprule\noalign{}
\begin{minipage}[b]{\linewidth}\raggedright
\textbf{Rank}
\end{minipage} & \begin{minipage}[b]{\linewidth}\raggedright
\textbf{Object}
\end{minipage} & \begin{minipage}[b]{\linewidth}\raggedright
\textbf{What It Captures}
\end{minipage} & \begin{minipage}[b]{\linewidth}\raggedright
\textbf{What It Loses}
\end{minipage} \\
\midrule\noalign{}
\endhead
\bottomrule\noalign{}
\endlastfoot
(0,0) & Scalar S & Overall moral assessment & Everything else \\
(1,0) & Obligation vector O\^{}μ & Direction and magnitude of duty & How
interests weight the dimensions \\
(0,1) & Interest covector I\_μ & How much each dimension matters to a
stakeholder & What the duties actually are \\
(0,2) & Metric tensor g\_μν & The structure of trade-offs between
dimensions & Which duties and interests are operative \\
(1,1) & Evaluation tensor E\^{}μ\_ν & The full map from interests to
obligation-responses & The specific stakeholder and context \\
\end{longtable}
}

The hierarchy continues beyond the rank-1 obligation vector. Chapters 14
and 16 introduce higher-rank tensors for collective agency and
uncertainty, and the DEME V3 implementation (Chapter 19) extends the
hierarchy to rank 6:

{\def\LTcaptype{none} % do not increment counter
\begin{longtable}[]{@{}
  >{\raggedright\arraybackslash}p{(\linewidth - 6\tabcolsep) * \real{0.0812}}
  >{\raggedright\arraybackslash}p{(\linewidth - 6\tabcolsep) * \real{0.2555}}
  >{\raggedright\arraybackslash}p{(\linewidth - 6\tabcolsep) * \real{0.4865}}
  >{\raggedright\arraybackslash}p{(\linewidth - 6\tabcolsep) * \real{0.1768}}@{}}
\toprule\noalign{}
\begin{minipage}[b]{\linewidth}\raggedright
\textbf{Rank}
\end{minipage} & \begin{minipage}[b]{\linewidth}\raggedright
\textbf{Object}
\end{minipage} & \begin{minipage}[b]{\linewidth}\raggedright
\textbf{What It Adds}
\end{minipage} & \begin{minipage}[b]{\linewidth}\raggedright
\textbf{Dimensions}
\end{minipage} \\
\midrule\noalign{}
\endhead
\bottomrule\noalign{}
\endlastfoot
2 & Per-party (distributional) tensor & Per-stakeholder obligation
profiles across all affected parties (the DEME distributional axis) &
k×n \\
3 & Temporal tensor & Time evolution with discount \(\gamma\);
trajectory analysis & k×n×τ \\
4 & Action--coalition tensor & Obligations across candidate actions and
coalitions & k×n×a×c \\
5 & Coalition tensor & Multi-agent coordination over time; Shapley credit
assignment & k×n×τ×a×c \\
6 & Sample tensor & Monte Carlo distributional risk (CVaR, VaR) &
k×n×τ×a×c×s \\
\end{longtable}
}

where \(k = 9\) is the moral manifold dimension, \(n\) the number of
stakeholders or parties, \(\tau\) the number of time steps, \(a\) the
number of candidate actions, \(c\) the number of coalitions, and \(s\)
the number of Monte Carlo samples. Per-dimension uncertainty --- the
covariance structure \(\Sigma^{\mu\nu}\) of Chapter 16 --- is not carried
as a separate tensor rank in DEME V3; it is realized as uncertainty
metadata attached to each tensor cell and reconstructed through the
rank-6 Monte Carlo sample axis. Tucker decomposition and tensor-train
formats maintain tractability as rank increases (see Chapter 19 for
implementation details).
\includegraphics[width=5.5in,height=3.52431in]{media/image4.png}Figure~\hyperref[fig:tensor-hierarchy]{1}
illustrates the full rank hierarchy and its contraction structure.

\textbar{} The Tensor Rank Tower

Each step up the hierarchy preserves more information. The crucial
insight is that the scalar S is always recoverable from the higher-rank
objects by \emph{contraction} --- summation over paired indices --- but
the higher-rank objects are \emph{not} recoverable from S. Information
flows downward (by contraction), not upward. The tensor hierarchy is a
hierarchy of informativeness, with the scalar at the bottom.

\begin{figure}
\centering
\includegraphics[width=5in,height=2.32083in]{media/image5.png}
\caption{Tucker decomposition of the MoralTensor: the full rank-3 tensor
T decomposes into a compact core tensor S multiplied by factor matrices
along each mode (moral dimensions, agents, time steps), enabling
tractable computation while preserving multi-modal structure.}
\end{figure}

\textbar{} Tucker Decomposition for Tractability

\subsection{6.2 Obligations as Vectors}\label{obligations-as-vectors}

\subsubsection{The Directional Character of
Duty}\label{the-directional-character-of-duty}

An obligation is not merely strong or weak. It \emph{points} --- from
the current state of affairs toward a required state. ``You ought to
return the borrowed book'' specifies both an intensity of oughtness and
a direction: from the state ``book unreturned'' toward the state ``book
returned.'' The spatial metaphor is not merely metaphorical. It is the
content of a precise mathematical claim: obligations are tangent vectors
on the moral manifold.

\textbf{Definition 6.1 (Obligation Vector). An obligation at a point p ∈
M is a tangent vector O(p) ∈ T\_pM. An obligation field is a smooth
section O: M → TM of the tangent bundle, assigning an obligation vector
to each point of M . Smoothness is defined stratum-wise: within each
stratum S\_α, the section O\textbar\_\{S\_α\} is smooth; across stratum
boundaries, O may be discontinuous.}

In the nine-dimensional coordinates (x¹, \ldots, x⁹) of Chapter 5, an
obligation vector has nine components:

\[O = O^{\mu}\frac{\partial}{\partial x^{\mu}} = O^{1}\frac{\partial}{\partial x^{1}} + O^{2}\frac{\partial}{\partial x^{2}} + \cdots + O^{9}\frac{\partial}{\partial x^{9}}\]

Each component O\^{}μ gives the strength of the obligation \emph{along
dimension μ}:

\begin{itemize}
\item
  \includegraphics[width=2.71329in,height=3.44048in]{media/image6.png}O¹
  = obligation to improve consequences/welfare
\item
  O² = obligation to respect rights and fulfill duties
\item
  O³ = obligation to promote justice and fairness
\item
  O⁴ = obligation to preserve autonomy
\item
  O⁵ = obligation to protect privacy
\item
  O⁶ = obligation to consider societal/environmental impact
\item
  O⁷ = obligation to exercise virtue and care
\item
  O⁸ = obligation to ensure procedural legitimacy
\end{itemize}

\textbar{} The Nine Dimensions of the Moral Manifold

Figure~\hyperref[fig:moral-dims]{3} provides a visual map of these nine
dimensions and shows how higher-rank tensors build upon them.

\protect\phantomsection\label{fig:moral-dims}{}

\subsubsection{}\label{section}

\subsubsection{\texorpdfstring{Rank progression: from rank-1 obligation
vectors to rank-6 distributional tensors.
}{Rank progression: from rank-1 obligation vectors to rank-6 distributional tensors. }}\label{rank-progression-from-rank-1-obligation-vectors-to-rank-6-distributional-tensors.}

Higher ranks encode richer moral structure at the cost of greater
computational complexity.

A situation may generate obligations along multiple dimensions
simultaneously. The promise to help Alex move generates O²
\textgreater{} 0 (a duty exists) and O⁷ \textgreater{} 0 (the care
relationship calls for action), while the components along other
dimensions may be zero or negligible. The \emph{direction} of O in the
nine-dimensional tangent space encodes which moral considerations are in
play.

\textbar{} Building Higher-Rank Tensors from the 9 Dimensions

Because O\^{}μ is a vector (a (1,0)-tensor), it transforms under
coordinate redescription by the Jacobian:

\[{\widetilde{O}}^{\mu} = \frac{\partial{\widetilde{x}}^{\mu}}{\partial x^{\nu}}O^{\nu}\]

This is the mathematical content of the Bond Invariance Principle
applied to obligations: the obligation itself is a geometric object,
independent of how we describe the situation. Only its \emph{components}
change under redescription. If ``withholding treatment'' and ``allowing
natural death'' are coordinate redescriptions of the same action, the
obligation vectors must be related by the Jacobian of the coordinate
change.

\subsubsection{The Zero Vector and Its
Significance}\label{the-zero-vector-and-its-significance}

A point p where O(p) = 0 --- the zero vector --- is a point of \emph{no
obligation}. The agent at p has no duty to move in any direction. This
does not mean the situation is morally neutral; it means no change is
required. This occurs when all duties are satisfied, or (more
interestingly) when duties in different dimensions cancel: the
obligation to act along Dimension 2 is balanced by a counter-obligation
along Dimension 4, resulting in a net obligation of zero.

The vanishing of O is coordinate-invariant: if O = 0 in one coordinate
system, it is zero in all. This is important because it means that
whether an obligation exists is an objective feature of the situation,
not an artifact of description.

\subsubsection{Multiple Obligation
Sources}\label{multiple-obligation-sources}

In general, multiple sources generate obligations simultaneously. A
promise creates an obligation along Dimension 2. A relationship creates
an obligation along Dimension 7. An emergency creates an obligation
along Dimension 1. These combine by vector addition:

\[O_{\text{total}} = O_{\text{promise}} + O_{\text{care}} + O_{\text{emergency}}\]

This is the geometric version of Ross's insight that duties are plural
and interact: some are aligned (pointing in similar directions), some
are orthogonal (addressing independent concerns), and some are opposed
(pointing in opposite directions). The net obligation is their vector
sum --- a single vector that encodes, in its direction and magnitude,
the \emph{resultant} of all operative moral considerations.

When two obligations are opposed, the resultant may have smaller
magnitude than either individually. When they are aligned, the resultant
is larger. When they are orthogonal, the resultant points in a direction
different from either source --- an emergent moral direction that
neither obligation alone would generate.

Chapter 11 reinterprets these obligation vectors within the A* search
framework: the obligation Oᵐ is the negative gradient of a heuristic
function h that estimates the cost remaining to reach moral equilibrium.
The directionality of the obligation is the direction of steepest
descent of h.

\subsection{6.3 Interests as Covectors}\label{interests-as-covectors}

\subsubsection{The Dual Perspective}\label{the-dual-perspective}

If obligations are vectors --- objects that \emph{point} --- then
interests are covectors --- objects that \emph{weight}. An interest does
not point in a direction; it assigns a numerical value to each
direction, measuring how much that direction matters to the interested
party.

\textbf{Definition 6.2 (Interest Covector).} \emph{An interest at a
point p ∈ M is a covector I(p) ∈ T\_p\^{}}M --- a linear map from
tangent vectors to real numbers. An interest field is a smooth section
I: M → T\^{}\emph{M of the cotangent bundle.}

In coordinates, an interest covector has components with lower indices:

\[I = I_{\mu}\, dx^{\mu} = I_{1}\, dx^{1} + I_{2}\, dx^{2} + \cdots + I_{9}\, dx^{9}\]

Each component I\_μ gives the \emph{weight} that this interest assigns
to dimension μ:

\begin{itemize}
\item
  I₁ = how much consequences/welfare matters to this stakeholder
\item
  I₂ = how much rights/duties matter
\item
  I₃ = how much justice/fairness matters
\item
  \ldots and so on.
\end{itemize}

An interest covector represents a \emph{perspective on what matters} ---
a profile of moral priorities. Different stakeholders in the same
situation will typically have different interest covectors.

\subsubsection{Why Covectors, Not
Vectors?}\label{why-covectors-not-vectors}

The distinction between vectors and covectors (Chapter 4, Section 4.3)
may seem pedantic, but it is morally significant. Vectors and covectors
are different mathematical objects with different transformation laws.
They are dual to each other --- a covector acts on a vector to produce a
scalar --- and this duality captures the fundamental structure of moral
evaluation.

Obligations are \emph{what is required}. They have a direction --- a
specific change that duty demands. Interests are \emph{what matters}.
They assign significance to directions --- they weight the dimensions of
moral evaluation. The satisfaction of an interest by an obligation is
the \emph{pairing} of the two: the covector I applied to the vector O.

This duality is not a convention. It reflects a deep asymmetry in moral
reasoning. Obligations and interests are not the same kind of thing, and
treating them as such (as scalar frameworks implicitly do, by collapsing
both into numbers and adding them) obscures the structure of moral
evaluation.

\subsubsection{The Gradient as a Natural
Covector}\label{the-gradient-as-a-natural-covector}

The most natural covector at a point is the \emph{gradient} of a scalar
function. If V: M → ℝ is a value function --- a scalar measure of how
good things are in some respect --- then its gradient dV is a covector
field:

\[(dV)_{\mu} = \frac{\partial V}{\partial x^{\mu}}\]

Applied to a tangent vector v, the gradient gives the rate of change of
V in the direction v:

\[dV(v) = v^{\mu}\frac{\partial V}{\partial x^{\mu}}\]

This has a clear moral reading: if V measures welfare, then dV is the
\emph{welfare interest} --- the covector that, applied to any proposed
change (tangent vector), reports how much welfare would change. A
stakeholder whose interest is welfare maximization is described by the
covector I = dV.

Different value functions yield different interest covectors. A
welfarist has I = d(Welfare). An egalitarian has I = d(Equality). A
libertarian has I = d(Liberty). The interest covector encodes the
theory; the obligation vector encodes the situation; and their pairing
produces a verdict.

\subsection{6.4 The Fundamental Contraction:
Satisfaction}\label{the-fundamental-contraction-satisfaction}

\subsubsection{The Formula}\label{the-formula}

The central formula of geometric ethics is the \emph{contraction} of
interest with obligation:

\[S = I_{\mu}O^{\mu} = \sum_{\mu = 1}^{9}I_{\mu}O^{\mu}\]

\textbf{Definition 6.3 (Satisfaction).} \emph{The satisfaction S at a
point p ∈ M, given an obligation vector O ∈ T\_pM and an interest
covector I ∈ T\_p\^{}}M, is the scalar S = I(O) = I\_μ O\^{}μ.*

Remark. The contraction S = I\_μ O\^{}μ requires that O and I are
defined at the same base point p ∈ M (O ∈ T\_pM, I ∈ T\_p*M). Comparing
satisfaction values at different points p and q requires
parallel-transporting one tensor to the other's base point along a path
γ from p to q, introducing a path-dependence that reflects the curvature
of the moral manifold (see Chapter 10).

This formula says: to evaluate how well an obligation addresses an
interest, apply the interest (as a linear functional) to the obligation.
The result is a real number --- a scalar --- that is \emph{invariant
under coordinate redescription}.

\subsubsection{Unpacking the Formula}\label{unpacking-the-formula}

Expanding in coordinates:

\[S = I_{1}O^{1} + I_{2}O^{2} + I_{3}O^{3} + \cdots + I_{9}O^{9}\]

Each term I\_μ O\^{}μ represents the \emph{contribution of dimension μ}
to overall satisfaction: the weight that the interest assigns to that
dimension, multiplied by the strength of the obligation in that
dimension. Satisfaction is large when the obligation points in the
directions that the interest weights most heavily. Satisfaction is small
(or negative) when the obligation points orthogonally to --- or in
opposition to --- the interest.

\subsubsection{The Moral Content of
Contraction}\label{the-moral-content-of-contraction}

Contraction is the mathematical operation that converts
multi-dimensional moral structure into a single-dimensional verdict. It
is the point at which the tensorial framework makes contact with the
scalar requirement of action: eventually, one must \emph{decide}, and a
decision is (in this minimal sense) a scalar --- this option rather than
that.

But the contraction is \emph{explicit}. We can see exactly which
interests generated the verdict, which dimensions of obligation
contributed, and how the result would change if the interest covector
were different. This is the fundamental advantage over scalar ethics:
the scalar verdict S is not a black box. It is a \emph{transparent
contraction} of identifiable moral components.

Chapter 15 develops the concept of \emph{moral residue}: the information
that contraction discards. When S = I\_μ O\^{}μ is the same for two very
different obligation-interest configurations, the scalar says ``these
are equivalent.'' But the tensorial structure reveals that they are
equivalent only from this particular interest-perspective, and that a
different interest covector would distinguish them.

\subsubsection{Contraction Is Not
Unique}\label{contraction-is-not-unique}

Different interest covectors yield different satisfaction scalars for
the same obligation:

\[S_{\text{welfarist}} = I_{\mu}^{\text{welfare}}O^{\mu} \neq I_{\mu}^{\text{egalitarian}}O^{\mu} = S_{\text{egalitarian}}\]

The framework does not say which contraction is ``correct.'' It says:
specify your interest covector (your theory of what matters), and the
satisfaction follows by mathematics. The choice of I is the moral
content; the contraction is the mathematical form.

This means that disagreement between ethical theories can be localized:
do the welfarist and the egalitarian disagree about the obligation O
(the facts about what is required), or about the interest I (the weights
on what matters)? The tensorial framework separates these questions,
which scalar frameworks conflate.

\subsection{6.5 The Moral Metric}\label{the-moral-metric}

\subsubsection{Relating Vectors and
Covectors}\label{relating-vectors-and-covectors}

Obligations (vectors) and interests (covectors) live in dual spaces. In
general, there is no canonical way to convert one into the other. A
vector is a direction; a covector is a weight. They are fundamentally
different kinds of object.

But a \emph{metric tensor} provides a bridge. The metric g\_μν is a
symmetric (0,2)-tensor that defines an inner product on tangent vectors
(Chapter 4, Section 4.5). Given a metric, we can:

\begin{itemize}
\item
  \textbf{Lower an index:} Convert a vector O\^{}μ into a covector O\_ν
  = g\_μν O\^{}μ. This covector is ``the interest that exactly aligns
  with obligation O.''
\item
  \textbf{Raise an index:} Convert a covector I\_μ into a vector I\^{}ν
  = g\^{}\{μν\} I\_μ. This vector is ``the obligation that exactly
  satisfies interest I.''
\item
  \textbf{Measure inner products:} The inner product of two obligation
  vectors is g\_μν O\^{}μ P\^{}ν --- a scalar measuring the alignment
  between two duties.
\end{itemize}

\subsubsection{The Metric as Trade-off
Structure}\label{the-metric-as-trade-off-structure}

What does the moral metric \emph{mean}? It encodes the structure of
permissible trade-offs between moral dimensions.

\textbf{Diagonal components} g\_μμ (no sum) give the \emph{weight} of
dimension μ --- how much a unit change along dimension μ ``costs'' in
moral terms.

\textbf{Off-diagonal components} g\_μν (μ ≠ ν) give the \emph{coupling}
between dimensions --- how much change in dimension μ is related to
change in dimension ν.

\begin{itemize}
\item
  If g\_μν \textgreater{} 0: improvements along μ tend to correlate with
  improvements along ν. The dimensions are \emph{aligned}.
\item
  If g\_μν \textless{} 0: improving along μ tends to worsen ν. The
  dimensions are \emph{opposed} --- they compete.
\item
  If g\_μν = 0: the dimensions are \emph{orthogonal} --- they vary
  independently, and no trade-off between them exists.
\end{itemize}

\textbf{Example.} Suppose Dimensions 1 (welfare) and 3 (justice) have
g₁₃ \textless{} 0 in some region of M. This means that increasing
welfare comes at the cost of justice, and vice versa --- a structural
tension between efficiency and equity. The magnitude g₁₃ measures how
severe the trade-off is: a large negative coupling means even small
welfare gains require significant justice costs.

If g₁₃ = 0, welfare and justice are independent: one can be improved
without affecting the other. If g₁₃ \textgreater{} 0, they are
synergistic: improving welfare tends to improve justice as well.

\subsubsection{Different Theories as Different
Metrics}\label{different-theories-as-different-metrics}

One of the most powerful features of the geometric framework is that
different ethical theories correspond to different metric tensors on the
\emph{same} underlying manifold M.

\textbf{Utilitarian metric.} All dimensions are commensurable. The
metric is non-degenerate with specific trade-off ratios:

\[g_{\mu\nu}^{\text{util}} = w_{\mu}\delta_{\mu\nu}\]

where w\_μ are the utilitarian weights and δ\_μν is the Kronecker delta
(a diagonal metric). Trade-offs are always available: any loss along one
dimension can be compensated by a sufficient gain along another.

\textbf{Lexicographic metric.} Some dimensions have lexicographic
priority. Formally, this is the limit of a family of metrics:

\[g_{\mu\nu}^{\text{lex}}(\epsilon) = \text{diag}(\epsilon^{- 2(9 - \pi(\mu))})\]

where π is a priority ordering and ε → 0. In the limit, the
highest-priority dimension dominates all trade-off calculations.
Lower-priority dimensions matter only when the highest-priority
dimension is tied.

\textbf{Degenerate metric.} Some dimensions are incommensurable. The
metric has zero eigenvalues along certain directions:

\[g_{\mu\nu}^{\text{deg}}v^{\mu}w^{\nu} = 0\quad\text{for all }w,\text{ if }v\text{ lies in the incommensurable subspace}\]

This means no inner product can be formed between obligation components
in the incommensurable directions. No trade-off is defined. The
framework represents this as a \emph{structural feature of moral space},
not a failure of analysis.

\subsubsection{The Metric Is Not Given}\label{the-metric-is-not-given}

The framework does not tell us which metric is correct. The choice of
metric is a substantive moral commitment --- it is where ethical
\emph{content} enters the formal structure. Chapter 9 develops this
point: the metric is neither discovered (contra moral realism taken
naively), nor constructed by idealized reason alone (contra
constructivism), nor projected by our sensibilities (contra
expressivism), but \emph{governed} --- the output of legitimate
institutional processes that have authority to determine, for a
community, how moral trade-offs are structured.

What the framework provides is \emph{precision}: once a metric is
specified, its consequences follow by mathematics. Two theories that
seem to disagree about everything may, in tensorial form, agree on the
manifold and the obligation fields but differ only in the metric --- a
disagreement that can be precisely located and productively debated.

\subsection{6.6 The Full Moral Tensor}\label{the-full-moral-tensor}

\subsubsection{Beyond Vectors and
Covectors}\label{beyond-vectors-and-covectors}

The obligation vector O\^{}μ and the interest covector I\_μ are rank-1
tensors. The metric g\_μν is a rank-2 tensor. But moral evaluation often
requires \emph{higher-rank} structure.

\subsubsection{The Evaluation Tensor}\label{the-evaluation-tensor}

\textbf{Definition 6.4 (Evaluation Tensor).} \emph{The evaluation tensor
at a point p ∈ M is a (1,1)-tensor E\^{}μ\_ν that maps interest
covectors to obligation-response vectors:}

Type clarification. E\^{}μ\_ν is a (1,1)-tensor (one contravariant, one
covariant index). As a linear map, it acts on covectors and produces
vectors: given an interest covector I\_ν, the contraction E\^{}μ\_ν I\_ν
= O\^{}μ yields the obligation-response vector. Equivalently, via the
metric, E can be viewed as a linear endomorphism of the tangent space
T\_pM. The satisfaction scalar is then the further contraction S = I\_μ
O\^{}μ = I\_μ E\^{}μ\_ν I\_ν, a quadratic form in the interest covector.

\[O_{\text{response}}^{\mu} = E_{\nu}^{\mu}I^{\nu}\]

The evaluation tensor encodes: given an interest (what matters to a
stakeholder), what obligation does the situation generate? It is the
\emph{moral function} of the situation --- the map from ``what you care
about'' to ``what you ought to do.''

Different situations have different evaluation tensors. A situation
involving an explicit promise has a large E²\_2 component (an interest
in rights generates a strong obligation in the rights dimension). A
situation involving a vulnerable person has a large E¹\_7 component (a
care interest generates a welfare obligation).

\subsubsection{The Multi-Agent Evaluation
Tensor}\label{the-multi-agent-evaluation-tensor}

When multiple agents are involved, the evaluation tensor acquires an
additional agent index:

Notational caution. The agent superscript (a) in E\^{}\{μ(a)\}\_ν is a
label index, not a tensor index. It does not transform under coordinate
changes on M. The multi-agent evaluation tensor is therefore not a
single higher-rank tensor on M, but a family of (1,1)-tensors
\{E\^{}\{(a)\}\} parametrized by the agent set A. The notation
E\^{}\{μ(a)\}\_ν is retained for compactness but should be read as ``the
μν-component of agent a's evaluation tensor.''

\[E_{\nu a}^{\mu}(p)\]

This (1,2)-tensor assigns, for each agent a and each interest dimension
ν, an obligation component O\^{}μ. Different agents in the same
situation face different obligations: the physician's obligations differ
from the patient's family's.

The \emph{correlative structure} of Hohfeldian analysis (Chapter 5,
Section 5.4) is a constraint on this tensor: if agent a has an
obligation component E²\_\{2,a\} \textgreater{} 0 (a duty in the rights
dimension), then agent b --- the correlative party --- must have
E²\_\{2,b\} corresponding to a claim. This constraint links the
evaluation tensors of different agents at the same point.

\subsubsection{The Uncertainty Tensor}\label{the-uncertainty-tensor}

Moral evaluation is rarely certain. The \emph{uncertainty tensor} (or
\emph{moral covariance tensor}) is a symmetric (2,0)-tensor:

\[\Sigma^{\mu\nu}(p\mathbb{) = E\lbrack(}\delta O^{\mu})(\delta O^{\nu})\rbrack\]

where δO\^{}μ is the deviation of the obligation component from its
expected value. This tensor encodes:

\begin{itemize}
\item
  \textbf{Variance} along each dimension: Σ\^{}\{μμ\} measures how
  uncertain we are about the obligation in dimension μ.
\item
  \textbf{Covariance} between dimensions: Σ\^{}\{μν\} (μ ≠ ν) measures
  whether uncertainty in dimension μ correlates with uncertainty in
  dimension ν.
\item
  \textbf{Principal directions of uncertainty:} The eigenvectors of Σ
  give the directions of maximum and minimum moral uncertainty.
\end{itemize}

The \emph{moral risk} of a decision depends on the alignment between the
uncertainty tensor and the interest covector:

\[\sigma_{S}^{2} = \Sigma^{\mu\nu}I_{\mu}I_{\nu}\]

This scalar --- the variance of satisfaction given structured
uncertainty --- is large when uncertainty concentrates along the
dimensions the stakeholder cares most about, and small when uncertainty
lies along morally irrelevant directions.

This is the geometric content of the old man's ``maybe'' from Chapter 2:
his uncertainty had \emph{shape}, concentrated along the axes that would
turn out to be decisive. The uncertainty tensor captures this shape, and
the contraction with the interest covector determines how much the
uncertainty matters.

\textbf{Higher-rank extensions.} The uncertainty tensor
\(\Sigma^{\mu\nu}\) is a symmetric object attached to the rank-1
obligation vector, capturing the shape of moral uncertainty (Chapter
16). In the DEME V3 implementation (Chapter 19), the \texttt{MoralTensor}
hierarchy adds a per-party (per-stakeholder) index at rank 2 --- the
distributional axis \(n\) that carries a separate obligation profile for
each affected party --- temporal indices at rank 3 for tracking how
obligations evolve through time, action and coalition indices at rank 4
for evaluating candidate actions across coalitions, combined
coalition-over-time indices at rank 5 for multi-agent scenarios, and
sample indices at rank 6 for distributional risk measures such as CVaR
that require Monte Carlo estimation of tail risk. The covariance
\(\Sigma^{\mu\nu}\) itself is not carried as a separate tensor axis in
V3: it is realized (a) as per-dimension uncertainty metadata attached to
each tensor cell and (b) through the rank-6 Monte Carlo sample axis
\(s\), whose spread reconstructs the covariance structure. These
higher-rank tensors are developed theoretically in Chapters 10
(temporal), 14 (collective), and 16 (uncertainty), and implemented
computationally in Chapter 19.

\subsection{6.7 Scalar Ethics as a Special
Case}\label{scalar-ethics-as-a-special-case}

\subsubsection{Recovery by Contraction}\label{recovery-by-contraction}

Every major ethical theory can be represented as a specific pattern of
tensorial structure and a specific contraction that yields a scalar
verdict.

\textbf{Classical utilitarianism.} The interest covector is I\_μ = (1,
0, 0, 0, 0, 0, 0, 0, 0) --- only welfare matters. The satisfaction is:

\[S_{\text{util}} = I_{\mu}O^{\mu} = O^{1}\]

The utilitarian verdict depends only on the welfare component of the
obligation. All other dimensions are discarded. This is a \emph{choice}
--- a specific contraction --- not a mathematical necessity. The
tensorial framework makes the choice visible.

\textbf{Prioritarianism.} The interest covector gives extra weight to
the welfare of the worst-off. If the worst-off is the focal agent, the
interest covector has elevated I₁ and I₃ (welfare and justice), with the
weight depending on the agent's position:

\[I_{\mu}^{\text{prior}} = (w(r),0,w'(r),0,\ldots)\]

where r is the agent's ranking and w(r) is a decreasing function giving
more weight to lower-ranked individuals.

\textbf{Kantian deontology.} The obligation field is determined by the
categorical imperative: O\^{}μ(p) is the obligation at p that would be
universalizable. The interest covector foregrounds rights: I₂ is
dominant. The satisfaction is:

\[S_{\text{Kant}} \approx I_{2}O^{2} + I_{4}O^{4}\]

primarily a function of duty-fulfillment and autonomy-preservation.

\textbf{Care ethics.} The interest covector foregrounds relational
dimensions:

\[S_{\text{care}} \approx I_{7}O^{7} + I_{1}O^{1}\]

The satisfaction depends primarily on the quality of the care
relationship and the welfare impact, with procedural and justice
dimensions receiving lower weight.

\textbf{Virtue ethics.} Here the relevant tensor is not the obligation
field but a \emph{section of the character bundle} (Chapter 5,
Definition 5.3). The ``virtuous action'' is the one that a person of
excellent character would choose --- a section σ of the character fiber
bundle that satisfies certain smoothness and consistency conditions. The
evaluation is:

\[S_{\text{virtue}} = g_{\mu\nu}\sigma^{\mu}O^{\nu}\]

where σ\^{}μ is the character profile and O\^{}ν is the obligation
field. Satisfaction is the inner product (measured by the metric)
between character and duty.

\subsubsection{What This Reveals}\label{what-this-reveals}

The tensorial framework does not adjudicate between these theories. It
does something arguably more valuable: it makes their \emph{structural
commitments} visible. Each theory is a choice of:

\begin{enumerate}
\def\labelenumi{\arabic{enumi}.}
\item
  \textbf{Which components of O to attend to} (which dimensions generate
  obligations)
\item
  \textbf{Which interest covector to use} (what weights the dimensions
  of evaluation)
\item
  \textbf{Which metric to assume} (what trade-offs are permissible)
\item
  \textbf{How to contract} (how the multi-dimensional structure reduces
  to a verdict)
\end{enumerate}

Two theories that seem to disagree about everything may, in tensorial
form, agree on the manifold and the obligation fields but differ only in
the interest covector --- a disagreement that can be precisely located
and productively debated. A seemingly deep philosophical dispute may
turn out to be a dispute about metric components, resolvable by
clarifying the trade-off structure.

Conversely, two theories that seem similar may turn out to differ in the
\emph{metric} rather than the \emph{interest} --- a deeper disagreement,
since the metric determines what comparisons are even possible. The
framework makes the \emph{depth} of disagreement visible, not just its
existence.

\subsection{6.8 The Hierarchy in Action}\label{the-hierarchy-in-action}

To consolidate the formal development, let us trace a single moral
situation through all levels of the tensor hierarchy.

\textbf{Situation.} A physician has a single dose of a scarce
medication. Two patients need it: Patient A (younger, less urgent) and
Patient B (older, more urgent).

\subsubsection{Level 0: Scalar}\label{level-0-scalar}

A single score: ``Give to B.'' Score = 0.73.

We know the verdict but not the reasons. We cannot explain the decision,
audit it, or assess its robustness.

\subsubsection{Level 1: Obligation
Vector}\label{level-1-obligation-vector}

\[O = (0.6,0.8,0.7,0.5,0.1,0.3,0.4,0.6,0.3)\]

Now we can see: the obligation is strongest in Dimension 2 (rights ---
both patients have claims), Dimension 3 (justice --- fair allocation
matters), and Dimension 8 (procedural legitimacy --- the process should
be defensible). The obligation is weakest in Dimension 5 (privacy is not
the primary concern).

\subsubsection{Level 2: Interest Covector (Physician's
Perspective)}\label{level-2-interest-covector-physicians-perspective}

\[I^{\text{physician}} = (0.20,0.15,0.15,0.10,0.05,0.10,0.10,0.10,0.05)\]

The physician weights welfare (0.20) most heavily, then rights and
justice equally (0.15 each).

\[S_{\text{physician}} = I_{\mu}O^{\mu} = 0.12 + 0.12 + 0.105 + 0.05 + 0.005 + 0.03 + 0.04 + 0.06 + 0.015 = 0.545\]

\subsubsection{Level 2': Interest Covector (Patient A's
Family)}\label{level-2-interest-covector-patient-as-family}

\[I^{\text{family}} = (0.30,0.25,0.05,0.15,0.05,0.02,0.13,0.03,0.02)\]

The family weights welfare (0.30) and rights (0.25) very heavily ---
``our loved one deserves this treatment.''

\[S_{\text{family}} = 0.18 + 0.20 + 0.035 + 0.075 + 0.005 + 0.006 + 0.052 + 0.018 + 0.006 = 0.577\]

The family's satisfaction is \emph{higher} than the physician's --- not
because the obligation is different, but because the family's interests
are more strongly aligned with the obligation's direction.

\subsubsection{Level 3: Metric}\label{level-3-metric}

With a utilitarian metric (diagonal, equal weights), the inner product
of the two patients' claims is:

\[\langle O_{A},O_{B}\rangle = g_{\mu\nu}O_{A}^{\mu}O_{B}^{\nu}\]

If this is small, the patients' claims are nearly orthogonal --- they
invoke different moral dimensions. If this is large, they invoke similar
dimensions and compete head-to-head. The metric determines whether the
situation is a genuine dilemma (orthogonal claims, degenerate metric
between them) or a tractable trade-off (aligned claims, non-degenerate
metric).

\subsubsection{What Each Level Adds}\label{what-each-level-adds}

{\def\LTcaptype{none} % do not increment counter
\begin{longtable}[]{@{}
  >{\raggedright\arraybackslash}p{(\linewidth - 4\tabcolsep) * \real{0.0820}}
  >{\raggedright\arraybackslash}p{(\linewidth - 4\tabcolsep) * \real{0.1573}}
  >{\raggedright\arraybackslash}p{(\linewidth - 4\tabcolsep) * \real{0.6085}}@{}}
\toprule\noalign{}
\begin{minipage}[b]{\linewidth}\raggedright
\textbf{Level}
\end{minipage} & \begin{minipage}[b]{\linewidth}\raggedright
\textbf{Object}
\end{minipage} & \begin{minipage}[b]{\linewidth}\raggedright
\textbf{New Information}
\end{minipage} \\
\midrule\noalign{}
\endhead
\bottomrule\noalign{}
\endlastfoot
0 & Scalar S & The verdict: which option is preferred \\
1 & Vector O\^{}μ & The \emph{reasons}: which dimensions generate
obligations \\
2 & Covector I\_μ & The \emph{perspective}: which dimensions matter to
this stakeholder \\
' & Multiple I\_μ & The \emph{landscape of disagreement}: how
perspectives differ \\
3 & Metric g\_μν & The \emph{trade-off structure}: how dimensions relate
to each other \\
\end{longtable}
}

Each level makes the previous level's assumptions \emph{explicit} and
\emph{debatable}. The scalar hides everything; the full tensor hierarchy
hides nothing.

\subsection{6.9 The Algebra of Moral
Tensors}\label{the-algebra-of-moral-tensors}

\subsubsection{Tensor Product}\label{tensor-product}

The \emph{tensor product} combines two tensors into a higher-rank
tensor. If O\^{}μ is an obligation vector and I\_ν is an interest
covector, their tensor product is:

\[T_{\nu}^{\mu} = O^{\mu}I_{\nu}\]

This (1,1)-tensor --- a 9×9 matrix in coordinates --- encodes the full
structure of the obligation-interest pair \emph{before} contraction.
Contraction (tracing over the paired indices) recovers the scalar:

\[S = T_{\mu}^{\mu} = O^{\mu}I_{\mu}\]

But the full tensor T\^{}μ\_ν carries more information: its eigenvalues
give the principal modes of satisfaction, and its eigenvectors give the
directions in moral space along which obligation and interest are most
(or least) aligned.

\subsubsection{Symmetrization and
Antisymmetrization}\label{symmetrization-and-antisymmetrization}

Given a (0,2)-tensor T\_μν, its symmetric part is:

\[T_{(\mu\nu)} = \frac{1}{2}(T_{\mu\nu} + T_{\nu\mu})\]

and its antisymmetric part is:

\[T_{\lbrack\mu\nu\rbrack} = \frac{1}{2}(T_{\mu\nu} - T_{\nu\mu})\]

The metric tensor is symmetric: g\_μν = g\_νμ. This reflects the fact
that the trade-off between Dimension μ and Dimension ν is the same in
both directions --- the cost of exchanging welfare for justice equals
the cost of exchanging justice for welfare.

But not all morally relevant (0,2)-tensors are symmetric. The
\emph{power tensor} P\_μν --- encoding asymmetric relationships between
agents, where agent a's influence on dimension μ of agent b's situation
differs from the reverse --- is generically asymmetric. Its
antisymmetric part P\_\{{[}μν{]}\} encodes the \emph{structural
inequity} in the relationship: the extent to which the power dynamic
favors one dimension over another in a directional way. Chapter 14
develops this in the context of collective agency.

\subsubsection{Contraction as Moral
Decision}\label{contraction-as-moral-decision}

We have emphasized that contraction --- summing over a paired index ---
is the operation that reduces rank and produces scalars. The full chain
of contractions from high-rank moral tensors to a scalar action-decision
can be written:

\[S = g^{\mu\alpha}E_{\alpha}^{\beta}I_{\beta} \cdot O^{\mu}\]

or in various other arrangements depending on which tensors are
contracted first. The \emph{order} of contraction may matter when the
tensors are not simply related --- different contraction orders
correspond to different decision procedures, and these may yield
different results on the same inputs.

Chapter 15 develops this observation into the theory of \emph{moral
contraction}: the mathematically necessary, informationally lossy, and
(sometimes) order-dependent process by which tensorial moral reality
yields scalar action-guidance.

\subsection{6.10 Summary}\label{summary}

The tensor hierarchy on the moral manifold M consists of:

{\def\LTcaptype{none} % do not increment counter
\begin{longtable}[]{@{}
  >{\raggedright\arraybackslash}p{(\linewidth - 6\tabcolsep) * \real{0.1727}}
  >{\raggedright\arraybackslash}p{(\linewidth - 6\tabcolsep) * \real{0.1861}}
  >{\raggedright\arraybackslash}p{(\linewidth - 6\tabcolsep) * \real{0.2558}}
  >{\raggedright\arraybackslash}p{(\linewidth - 6\tabcolsep) * \real{0.3853}}@{}}
\toprule\noalign{}
\begin{minipage}[b]{\linewidth}\raggedright
\textbf{Object}
\end{minipage} & \begin{minipage}[b]{\linewidth}\raggedright
\textbf{Type}
\end{minipage} & \begin{minipage}[b]{\linewidth}\raggedright
\textbf{Role}
\end{minipage} & \begin{minipage}[b]{\linewidth}\raggedright
\textbf{Mathematical Definition}
\end{minipage} \\
\midrule\noalign{}
\endhead
\bottomrule\noalign{}
\endlastfoot
Satisfaction S & (0,0)-tensor (scalar) & The action-guiding verdict & S
= I\_μ O\^{}μ \\
Obligation O & (1,0)-tensor (vector field) & Direction and magnitude of
duty & O: M → TM \\
Interest I & (0,1)-tensor (covector field) & Stakeholder's moral
priorities & I: M → T*M \\
Metric g & (0,2)-tensor (symmetric) & Trade-off structure between values
& g: TM × TM → ℝ \\
Evaluation E & (1,1)-tensor & Map from interests to obligations & O\^{}μ
= E\^{}μ\_ν I\^{}ν \\
Uncertainty Σ & (2,0)-tensor (symmetric) & Shape of moral uncertainty &
Σ\^{}μν = 𝔼{[}δO\^{}μ δO\^{}ν{]} \\
Temporal \(\mathcal{T}\) & rank-3 tensor & Obligation evolution over
time & T\^{}μ\_\{ν,τ\} = γ\^{}τ O\^{}μ\_ν(t\_τ) \\
Coalition \(\mathcal{C}\) & (1,\(n\))-tensor & Collective obligations
from \(n\) agents & O\^{}μ\_coll = C\^{}μ\_\{ν₁\ldots νₙ\}
I\^{}(1)\_\{ν₁\}\ldots I\^{}(n)\_\{νₙ\} \\
\end{longtable}
}

The fundamental formula S = I\_μ O\^{}μ is the contraction that ties the
hierarchy together: satisfaction is the pairing of interest with
obligation, a coordinate-invariant scalar arising from the interaction
of two rank-1 tensors.

Different ethical theories are different choices of I (what matters), g
(what trade-offs are permissible), and contraction procedure (how to
decide). The moral manifold M and the obligation field O --- the facts
about what duties arise --- are common ground. Disagreement is
localizable: in the interest, in the metric, or in the contraction. This
localizability is the framework's principal contribution to the analysis
of moral disagreement.

\subsection{Technical Appendix}\label{technical-appendix}

\textbf{Proposition 6.1 (Invariance of Satisfaction).} \emph{The
satisfaction scalar S = I\_μ O\^{}μ is invariant under admissible
coordinate transformations (Type 1 transformations).}

\emph{Proof.} Under a coordinate change x → x̃:

Remark. Proposition 6.1 is definitionally true: the satisfaction scalar
S = I\_μ O\^{}μ is a contraction of a covector with a vector, which is
by definition coordinate-independent. The proposition is included for
completeness and to make explicit that no additional structure (metric,
connection, or volume form) is required for S to be well-defined ---
only the natural pairing between T\_pM and T\_p*M.

\[{\widetilde{I}}_{\mu}{\widetilde{O}}^{\mu} = \left( \frac{\partial x^{\nu}}{\partial{\widetilde{x}}^{\mu}}I_{\nu} \right)\left( \frac{\partial{\widetilde{x}}^{\mu}}{\partial x^{\rho}}O^{\rho} \right) = \frac{\partial x^{\nu}}{\partial{\widetilde{x}}^{\mu}}\frac{\partial{\widetilde{x}}^{\mu}}{\partial x^{\rho}}I_{\nu}O^{\rho} = \delta_{\rho}^{\nu}\, I_{\nu}O^{\rho} = I_{\rho}O^{\rho} = S\]

The Jacobian and inverse Jacobian cancel, leaving the scalar unchanged.
□

\textbf{Proposition 6.2 (Decomposition of Disagreement). *Two evaluators
with interest covectors I and I' yield different satisfactions S = I\_μ
O\^{}μ and S' = I'\_μ O\^{}μ for the same obligation O if and only if
ΔI\_μ O\^{}μ ≠ 0, where ΔI = I' − I. In particular, S = S' for all O if
and only if I = I'. The disagreement decomposes as:*}

\[S' - S = (I_{\mu}' - I_{\mu})O^{\mu} = \Delta I_{\mu}\, O^{\mu}\]

\emph{where ΔI = I' - I is the ``interest difference covector.'' The
disagreement vanishes along dimensions where ΔI\_μ = 0 and is
proportional to O\^{}μ along dimensions where both the interest
difference and the obligation are nonzero.}

\emph{Proof.} By linearity of the contraction: S\textquotesingle{} − S =
I\textquotesingle μ Oᵘ − Iμ Oᵘ = (I\textquotesingle μ − Iμ)Oᵘ = ΔIμ Oᵘ.
For the biconditional: S\textquotesingle{} ≠ S iff ΔIμ Oᵘ ≠ 0, which
requires ΔIμ ≠ 0 for at least one μ with Oᵘ ≠ 0. Conversely, ΔIμ = 0 for
all μ implies S\textquotesingle{} = S, and S\textquotesingle{} = S for
all O implies ΔIμ = 0 for all μ (choose O = e\_μ for each dimension).
The component-wise structure follows: the contribution from dimension μ
is ΔIμ Oᵘ, which vanishes iff ΔIμ = 0 or Oᵘ = 0. □

\textbf{Proposition 6.3 (Metric Dependence of the Gradient Direction).}
\emph{The gradient vector of a scalar function f on M depends on the
metric:}

\[(\text{grad }f)^{\mu} = g^{\mu\nu}\frac{\partial f}{\partial x^{\nu}}\]

\emph{Different metrics yield different gradient directions for the same
function. In particular, the ``direction of moral improvement'' --- the
gradient of the satisfaction function --- depends on the metric, which
is to say: on the ethical theory.}

\emph{Proof.} The gradient of f is the unique vector field satisfying
g(grad f, V) = df(V) for all vector fields V. In coordinates:
g\_\{μα\}(grad f)ᵅ = ∂f/∂xᵘ. Multiplying both sides by g\^{}\{μν\}:
(grad f)ᵛ = g\^{}\{μν\} ∂f/∂xᵘ. Since g\^{}\{μν\} appears explicitly,
the gradient depends on the metric. For two admissible metrics g and ĝ
with g\^{}\{μν\} ≠ ĝ\^{}\{μν\}, we have (grad\_g f)ᵛ ≠ (grad\_\{ĝ\} f)ᵛ
whenever ∂f/∂xᵘ ≠ 0 along a dimension where the inverse metrics differ.
□

This result is central to the framework. The \emph{facts} of the
situation (the obligation field O) may be common ground, but the
\emph{direction of improvement} (the gradient of S) depends on the
metric --- the theory of trade-offs. Two theories agreeing on the facts
can disagree about what to do, because they disagree about the metric.

❖

\emph{The scalar is where ethics ends: in a decision.}

\emph{The tensor is where ethics lives: in the structure that generates,
constrains, and illuminates the decision.}

\section{Chapter 7: One Case, Five
Levels}\label{chapter-7-one-case-five-levels}

\begin{quote}
\textbf{\textsc{RUNNING EXAMPLE --- Priya's Model}}

\emph{Priya picks one patient: Mrs. Eleanor Voss, 67, Stage III
melanoma, Harlan County, Kentucky, nearest trial site 4.5 hours away.
Scalar: TrialMatch gives her a 64. Rejected. Vector: her medical
obligation components are strong, but her access components are weak;
the scalar crushed them together. Rank-2: the evaluation tensor reveals
that improving Mrs. Voss's access would also improve trial diversity---a
positive off-diagonal coupling invisible to the scalar. Metric: under
the current metric, 4.5 hours of travel counts the same as a minor
co-pay difference. Is that right? Stratified: the 70-point threshold is
a stratum boundary. Mrs. Voss is six points from a life-saving trial.
The boundary does not care about direction.}
\end{quote}

\subsection{Introduction: The Pedagogy of
Accumulation}\label{introduction-the-pedagogy-of-accumulation}

The preceding chapters have developed geometric ethics in abstract
terms: manifolds, tangent spaces, tensors, metrics, stratifications.
This chapter takes a different approach. We examine a single concrete
case and revisit it five times, each time adding mathematical structure.
By the end, the reader will see not just \emph{what} the geometric
apparatus is, but \emph{why} each piece is necessary --- what it lets us
say that we could not say before.

The case is organ allocation: the allocation of a single kidney to one
of three patients. This case is useful because:

\begin{enumerate}
\def\labelenumi{\arabic{enumi}.}
\item
  It is realistic and consequential
\item
  It involves multiple morally relevant dimensions from the
  nine-dimensional framework of Chapter 5
\item
  It involves multiple agents with distinct perspectives (patients,
  families, physicians, society)
\item
  It generates genuine disagreement traceable to structural choices
  (utilitarian vs.~egalitarian vs.~prioritarian)
\item
  It has clear boundary cases where the rules change discontinuously
  (stratification)
\end{enumerate}

We begin at Level 1 --- the scalar --- and progressively enrich the
structure.

\subsection{Level 1: The Scalar}\label{level-1-the-scalar}

\subsubsection{Setup}\label{setup}

A kidney becomes available. Three patients are eligible:

{\def\LTcaptype{none} % do not increment counter
\begin{longtable}[]{@{}
  >{\raggedright\arraybackslash}p{(\linewidth - 6\tabcolsep) * \real{0.1010}}
  >{\raggedright\arraybackslash}p{(\linewidth - 6\tabcolsep) * \real{0.1966}}
  >{\raggedright\arraybackslash}p{(\linewidth - 6\tabcolsep) * \real{0.2078}}
  >{\raggedright\arraybackslash}p{(\linewidth - 6\tabcolsep) * \real{0.0660}}@{}}
\toprule\noalign{}
\begin{minipage}[b]{\linewidth}\raggedright
\textbf{Patient}
\end{minipage} & \begin{minipage}[b]{\linewidth}\raggedright
\textbf{Medical Benefit}
\end{minipage} & \begin{minipage}[b]{\linewidth}\raggedright
\textbf{Years on Waitlist}
\end{minipage} & \begin{minipage}[b]{\linewidth}\raggedright
\textbf{Age}
\end{minipage} \\
\midrule\noalign{}
\endhead
\bottomrule\noalign{}
\endlastfoot
Alice & High (0.9) & 2 years & 45 \\
Bob & Medium (0.6) & 7 years & 62 \\
Carol & Medium (0.5) & 1 year & 28 \\
\end{longtable}
}

The hospital must choose one recipient.

\subsubsection{The Scalar Approach}\label{the-scalar-approach}

The simplest approach assigns a single score to each patient and chooses
the highest:

\[S_{i} = w_{1} \cdot \text{benefit}_{i} + w_{2} \cdot \text{waittime}_{i} + w_{3} \cdot \text{agefactor}_{i}\]

With weights w = (0.5, 0.3, 0.2) and normalized data:

\begin{itemize}
\item
  S\_Alice = 0.5(0.9) + 0.3(0.29) + 0.2(0.55) = 0.65
\item
  S\_Bob = 0.5(0.6) + 0.3(1.0) + 0.2(0.38) = 0.68
\item
  S\_Carol = 0.5(0.5) + 0.3(0.14) + 0.2(0.72) = 0.43
\end{itemize}

\textbf{Decision: Bob receives the kidney.}

\subsubsection{What the Scalar Can Say}\label{what-the-scalar-can-say}

\begin{itemize}
\item
  Bob is the best choice (given these weights)
\item
  Alice is second, Carol third
\item
  The margin is small (0.68 vs.~0.65)
\end{itemize}

\subsubsection{What the Scalar Cannot
Say}\label{what-the-scalar-cannot-say}

\begin{enumerate}
\def\labelenumi{\arabic{enumi}.}
\item
  \textbf{Why Bob wins.} Is it because of waiting time? Medical benefit?
  The scalar collapses all reasons into a single number. We cannot
  recover the structure of the justification.
\item
  \textbf{How robust the decision is.} If we slightly changed the
  weights, would the answer change? The scalar gives no information
  about sensitivity.
\item
  \textbf{What we are trading off.} Choosing Bob over Alice means
  accepting lower medical benefit in exchange for respecting longer
  waiting time. But the scalar does not represent this trade-off --- it
  only reports the output.
\item
  \textbf{How different perspectives evaluate the choice.} A physician
  focused on medical outcomes would weight differently than a policy
  administrator focused on fairness. The scalar framework requires
  choosing weights \emph{before} evaluation, concealing the underlying
  disagreement.
\end{enumerate}

The scalar is a (0,0)-tensor. It is the \emph{end} of moral reasoning
--- the point of decision. The problem is that scalar ethics makes it
the \emph{beginning} as well, collapsing structure before it can be
examined.

\subsection{Level 2: The Obligation
Vector}\label{level-2-the-obligation-vector}

\subsubsection{Adding Structure}\label{adding-structure}

Instead of collapsing to a scalar immediately, we represent the moral
situation as it bears on each patient using the multi-dimensional
framework of Chapters 5 and 6. The morally relevant dimensions for
kidney allocation include:

\begin{itemize}
\item
  Dimension 1 (Consequences/Welfare): medical benefit
\item
  Dimension 2 (Rights/Duties): strength of claim based on waiting time
\item
  Dimension 3 (Justice/Fairness): equitable distribution of scarce
  resources
\item
  Dimension 4 (Autonomy): patient's informed consent and preferences
\item
  Dimension 7 (Care): relational context --- dependents, family
  obligations
\item
  Dimension 9 (Epistemic): certainty of benefit estimates
\end{itemize}

Not all nine dimensions are equally active in every case. Here,
Dimensions 5 (privacy), 6 (societal impact), and 8 (procedural
legitimacy) play subsidiary roles. We work with a six-dimensional
subspace.

Enriching our data:

{\def\LTcaptype{none} % do not increment counter
\begin{longtable}[]{@{}
  >{\raggedright\arraybackslash}p{(\linewidth - 12\tabcolsep) * \real{0.1009}}
  >{\raggedright\arraybackslash}p{(\linewidth - 12\tabcolsep) * \real{0.1359}}
  >{\raggedright\arraybackslash}p{(\linewidth - 12\tabcolsep) * \real{0.1615}}
  >{\raggedright\arraybackslash}p{(\linewidth - 12\tabcolsep) * \real{0.1490}}
  >{\raggedright\arraybackslash}p{(\linewidth - 12\tabcolsep) * \real{0.1727}}
  >{\raggedright\arraybackslash}p{(\linewidth - 12\tabcolsep) * \real{0.1096}}
  >{\raggedright\arraybackslash}p{(\linewidth - 12\tabcolsep) * \real{0.1704}}@{}}
\toprule\noalign{}
\begin{minipage}[b]{\linewidth}\raggedright
\textbf{Patient}
\end{minipage} & \begin{minipage}[b]{\linewidth}\raggedright
\textbf{Benefit (x¹)}
\end{minipage} & \begin{minipage}[b]{\linewidth}\raggedright
\textbf{Wait-claim (x²)}
\end{minipage} & \begin{minipage}[b]{\linewidth}\raggedright
\textbf{Fairness (x³)}
\end{minipage} & \begin{minipage}[b]{\linewidth}\raggedright
\textbf{Autonomy (x⁴)}
\end{minipage} & \begin{minipage}[b]{\linewidth}\raggedright
\textbf{Care (x⁷)}
\end{minipage} & \begin{minipage}[b]{\linewidth}\raggedright
\textbf{Epistemic (x⁹)}
\end{minipage} \\
\midrule\noalign{}
\endhead
\bottomrule\noalign{}
\endlastfoot
Alice & 0.9 & 0.29 & 0.55 & 0.8 & 0.8 & 0.7 \\
Bob & 0.6 & 1.0 & 0.70 & 0.7 & 0.2 & 0.5 \\
Carol & 0.5 & 0.14 & 0.72 & 0.9 & 0.6 & 0.8 \\
\end{longtable}
}

Each patient's moral profile is now an \emph{obligation vector} --- a
tangent vector on the moral manifold at the current situation, pointing
in the direction of what is owed to that patient:

\[O_{A} = (0.9,0.29,0.55,0.8,0.8,0.7)\]

\[O_{B} = (0.6,1.0,0.70,0.7,0.2,0.5)\]

\[O_{C} = (0.5,0.14,0.72,0.9,0.6,0.8)\]

\subsubsection{What the Vector Can Say}\label{what-the-vector-can-say}

\textbf{1. The profile of each option.} Alice is high-benefit,
low-wait-claim, high-care (many dependents). Bob is medium-benefit,
very-high-wait-claim (7 years), low-care. Carol is young (high autonomy
value), fairest by age distribution, highest epistemic certainty. We can
\emph{see} the structure that the scalar collapsed.

\textbf{2. Dominance relations.} If one vector dominated another on all
components, we could identify it without choosing weights. None does
here --- which is why the case is hard. The hardness is
\emph{geometric}: the vectors are in general position, with no dominance
ordering.

\textbf{3. The geometry of difference.} The difference vector

\[O_{A} - O_{B} = (0.3, - 0.71, - 0.15,0.1,0.6,0.2)\]

shows exactly \emph{where} Alice and Bob differ. Alice is better on
benefit (+0.3) and care (+0.6); Bob is better on wait-claim (+0.71) and
slightly on fairness (+0.15). This vector is a precise map of the
trade-off.

\textbf{4. Decision regions in interest space.} Recall from Chapter 6
that the satisfaction scalar is S = I\_μ O\^{}μ, where I is the interest
covector. Alice wins when I\_μ O\^{}μ\_A \textgreater{} I\_μ O\^{}μ\_B
and I\_μ O\^{}μ\_A \textgreater{} I\_μ O\^{}μ\_C. These inequalities
carve out a \emph{region} in the six-dimensional space of interest
covectors --- the set of moral perspectives under which Alice is the
right choice.

\[\text{Alice wins:}\quad I_{\mu}(O_{A}^{\mu} - O_{B}^{\mu}) > 0\quad\text{and}\quad I_{\mu}(O_{A}^{\mu} - O_{C}^{\mu}) > 0\]

Similarly for Bob and Carol. The interest space is partitioned into
three \emph{decision regions}, one for each patient. The boundaries
between regions are hyperplanes --- and the widths of the regions
measure the \emph{robustness} of each choice.

\subsubsection{What This Reveals}\label{what-this-reveals-1}

The ``right'' answer depends on where you are in interest space. This is
not relativism --- it is \emph{structure}. Different interest covectors
correspond to different moral theories:

\begin{itemize}
\item
  High I₁ (welfare weight): utilitarian/efficiency orientation → Alice
  wins
\item
  High I₂ (wait-claim weight): procedural fairness orientation → Bob
  wins
\item
  High I₃ (fairness weight): egalitarian orientation → Carol wins
\item
  High I₇ (care weight): care ethics orientation → Alice wins
\end{itemize}

The vector framework does not tell you which theory is correct. It tells
you \emph{exactly what each theory implies} and \emph{where they
diverge}. The boundaries between decision regions are \emph{loci of
genuine moral disagreement} --- hyperplanes in interest space where
different theories yield different verdicts.

\subsection{Level 3: Multiple Perspectives (Rank-2
Tensor)}\label{level-3-multiple-perspectives-rank-2-tensor}

\subsubsection{The Multi-Agent
Structure}\label{the-multi-agent-structure}

So far, we have treated the moral evaluation as coming from a single
perspective. But organ allocation involves multiple stakeholders, each
with a distinct interest covector:

\begin{itemize}
\item
  \textbf{The physician} (medical judgment, welfare-focused)
\item
  \textbf{Alice's family} (partial to Alice, care-focused)
\item
  \textbf{Bob's family} (partial to Bob, rights-focused)
\item
  \textbf{Carol's family} (partial to Carol, autonomy-focused)
\item
  \textbf{The transplant committee} (institutional policy,
  fairness-focused)
\end{itemize}

These are different points on the fiber of the agent bundle over the
current situation (Chapter 5, Section 5.4). Each agent evaluates the
same obligation vectors, but with a different interest covector.

\subsubsection{The Evaluation Tensor}\label{the-evaluation-tensor-1}

The satisfaction of each option by each agent forms a rank-2 tensor:

\[M_{ia} = I_{\mu}^{(a)}O_{i}^{\mu}\]

where i indexes patients and a indexes perspectives.

{\def\LTcaptype{none} % do not increment counter
\begin{longtable}[]{@{}
  >{\raggedright\arraybackslash}p{(\linewidth - 10\tabcolsep) * \real{0.0778}}
  >{\raggedright\arraybackslash}p{(\linewidth - 10\tabcolsep) * \real{0.1293}}
  >{\raggedright\arraybackslash}p{(\linewidth - 10\tabcolsep) * \real{0.1344}}
  >{\raggedright\arraybackslash}p{(\linewidth - 10\tabcolsep) * \real{0.1349}}
  >{\raggedright\arraybackslash}p{(\linewidth - 10\tabcolsep) * \real{0.1349}}
  >{\raggedright\arraybackslash}p{(\linewidth - 10\tabcolsep) * \real{0.1452}}@{}}
\toprule\noalign{}
\begin{minipage}[b]{\linewidth}\raggedright
\end{minipage} & \begin{minipage}[b]{\linewidth}\raggedright
\textbf{Physician}
\end{minipage} & \begin{minipage}[b]{\linewidth}\raggedright
\textbf{Family(A)}
\end{minipage} & \begin{minipage}[b]{\linewidth}\raggedright
\textbf{Family(B)}
\end{minipage} & \begin{minipage}[b]{\linewidth}\raggedright
\textbf{Family(C)}
\end{minipage} & \begin{minipage}[b]{\linewidth}\raggedright
\textbf{Committee}
\end{minipage} \\
\midrule\noalign{}
\endhead
\bottomrule\noalign{}
\endlastfoot
Alice & 0.75 & 0.90 & 0.40 & 0.50 & 0.62 \\
Bob & 0.68 & 0.35 & 0.95 & 0.45 & 0.70 \\
Carol & 0.55 & 0.30 & 0.30 & 0.92 & 0.58 \\
\end{longtable}
}

Each entry M\_\{ia\} is itself a contraction: the interest covector of
agent a, contracted with the obligation vector of patient i. The tensor
M\_\{ia\} preserves the \emph{full structure of multi-agent evaluation}
before any social aggregation.

\subsubsection{What the Rank-2 Tensor Can
Say}\label{what-the-rank-2-tensor-can-say}

\textbf{1. Perspective-dependence is visible.} Alice's family ranks
Alice first (0.90); Bob's family ranks Bob first (0.95). This is not
bias to be eliminated --- it is legitimate partiality (Chapter 5,
Section 5.5, Type 2 transformations). Families \emph{should} weight
their own members more heavily on the care dimension. The tensor makes
this visible as a structural feature, not an error.

\textbf{2. Agreement and disagreement are localized.} The physician and
committee roughly agree (their columns are positively correlated).
Families disagree strongly with each other (their columns are negatively
correlated). We can compute \emph{exactly} which pairs of perspectives
agree and which diverge --- and on which options.

\textbf{3. Correlative structure.} The Hohfeldian constraint (Chapter 5,
Definition A.4) requires that if the committee assigns Alice an
obligation-based claim (M\_\{Alice,committee\} reflects a strong O²
component), then the other patients' families must acknowledge a
corresponding diminution of their claims. This constraint links the
entries of M\_\{ia\} across agents.

\textbf{4. Aggregation becomes a visible choice.} To reach a social
decision, we must \emph{contract} the tensor M\_\{ia\} --- from rank-2
(options × agents) to rank-1 (options) and then to rank-0 (a decision).
Different contractions encode different social choice procedures.

\subsubsection{Three Contractions, Three
Verdicts}\label{three-contractions-three-verdicts}

\textbf{Utilitarian aggregation} (sum over agents):

\[S_{i}^{\text{util}} = \sum_{a}^{}M_{ia}\]

\begin{itemize}
\item
  S\_Alice = 3.17, S\_Bob = 3.13, S\_Carol = 2.65
\end{itemize}

\textbf{Decision: Alice} (highest total support)

\textbf{Rawlsian aggregation} (maximize the minimum across agents):

\[S_{i}^{\text{Rawls}} = \min_{a}M_{ia}\]

\begin{itemize}
\item
  S\_Alice = 0.40, S\_Bob = 0.35, S\_Carol = 0.30
\end{itemize}

\textbf{Decision: Alice} (highest floor of support)

\textbf{Expert-weighted aggregation} (physician and committee weighted
more heavily):

\[S_{i}^{\text{expert}} = \sum_{a}^{}v_{a}M_{ia},\quad\quad v = (0.3,0.1,0.1,0.1,0.4)\]

\begin{itemize}
\item
  S\_Alice = 0.653, S\_Bob = 0.659, S\_Carol = 0.549
\end{itemize}

\textbf{Decision: Bob} (expert perspectives favor wait-time claims)

\subsubsection{The Meta-Point}\label{the-meta-point}

The same rank-2 tensor yields different decisions under different
contractions:

{\def\LTcaptype{none} % do not increment counter
\begin{longtable}[]{@{}
  >{\raggedright\arraybackslash}p{(\linewidth - 4\tabcolsep) * \real{0.1891}}
  >{\raggedright\arraybackslash}p{(\linewidth - 4\tabcolsep) * \real{0.1027}}
  >{\raggedright\arraybackslash}p{(\linewidth - 4\tabcolsep) * \real{0.2887}}@{}}
\toprule\noalign{}
\begin{minipage}[b]{\linewidth}\raggedright
\textbf{Contraction}
\end{minipage} & \begin{minipage}[b]{\linewidth}\raggedright
\textbf{Winner}
\end{minipage} & \begin{minipage}[b]{\linewidth}\raggedright
\textbf{Social Choice Procedure}
\end{minipage} \\
\midrule\noalign{}
\endhead
\bottomrule\noalign{}
\endlastfoot
Sum over agents & Alice & Democratic/utilitarian \\
Min over agents & Alice & Rawlsian/consensus \\
Expert-weighted & Bob & Epistemic/expertise-based \\
\end{longtable}
}

\textbf{This is not a bug. It is the point.} The tensor represents the
\emph{full structure} of the multi-agent evaluation. The contraction
represents a \emph{choice} of social decision procedure. Different
institutions, with different legitimacy claims, may choose different
contractions. The framework makes this choice explicit and analyzable,
rather than burying it in unstated assumptions.

\subsection{Level 4: The Metric}\label{level-4-the-metric}

\subsubsection{The Problem of
Comparison}\label{the-problem-of-comparison-1}

Levels 2 and 3 implicitly assumed that the six dimensions are comparable
--- that 0.1 units of ``benefit'' is equivalent to 0.1 units of
``wait-claim.'' But is it? Can we trade one year of waiting time for a
10\% increase in medical benefit?

The \emph{moral metric} g\_μν (Chapter 6, Section 6.5) answers this
question. It defines the inner product between moral vectors --- and
thereby determines how we measure distances, compare magnitudes, and
structure trade-offs across dimensions.

\subsubsection{Three Metrics, Three
Ethics}\label{three-metrics-three-ethics}

\textbf{Metric 1: Euclidean} (g\_μν = δ\_μν, all dimensions equally
weighted)

All dimensions are fully commensurable. A 0.1 difference in benefit is
as significant as a 0.1 difference in wait-claim. The distance from each
patient's vector to the ideal vector (1, 1, 1, 1, 1, 1) is:

\begin{itemize}
\item
  d(O\_A, ideal) = √(0.01 + 0.50 + 0.20 + 0.04 + 0.04 + 0.09) = √0.88 ≈
  0.94
\item
  d(O\_B, ideal) = √(0.16 + 0.00 + 0.09 + 0.09 + 0.64 + 0.25) = √1.23 ≈
  1.11
\item
  d(O\_C, ideal) = √(0.25 + 0.74 + 0.08 + 0.01 + 0.16 + 0.04) = √1.28 ≈
  1.13
\end{itemize}

\textbf{Decision (Euclidean): Alice} (closest to ideal)

\textbf{Metric 2: Weighted} (g\_μν = diag(2, 1.5, 1, 0.5, 1, 0.5),
benefit and rights count more)

The weighted distance changes the balance:

\begin{itemize}
\item
  d\_W(O\_A, ideal) = √(2(0.01) + 1.5(0.50) + 1(0.20) + 0.5(0.04) +
  1(0.04) + 0.5(0.09)) = √1.14 ≈ 1.07
\item
  d\_W(O\_B, ideal) = √(2(0.16) + 1.5(0.00) + 1(0.09) + 0.5(0.09) +
  1(0.64) + 0.5(0.25)) = √1.30 ≈ 1.14
\item
  d\_W(O\_C, ideal) = √(2(0.25) + 1.5(0.74) + 1(0.08) + 0.5(0.01) +
  1(0.16) + 0.5(0.04)) = √1.87 ≈ 1.37
\end{itemize}

\textbf{Decision (Weighted): Alice} (benefit advantage amplified)

\textbf{Metric 3: Lexicographic} (Dimension 2, rights/wait-claim, has
absolute priority)

Compare first on wait-claim. If equal (within ε = 0.1), proceed to
benefit.

\begin{itemize}
\item
  Wait-claim: Bob (1.0) \textgreater\textgreater{} Alice (0.29)
  \textgreater\textgreater{} Carol (0.14)
\end{itemize}

\textbf{Decision (Lexicographic): Bob} (best on top-priority dimension)

\subsubsection{What the Metric Choice Commits You
To}\label{what-the-metric-choice-commits-you-to}

{\def\LTcaptype{none} % do not increment counter
\begin{longtable}[]{@{}
  >{\raggedright\arraybackslash}p{(\linewidth - 4\tabcolsep) * \real{0.1634}}
  >{\raggedright\arraybackslash}p{(\linewidth - 4\tabcolsep) * \real{0.4313}}
  >{\raggedright\arraybackslash}p{(\linewidth - 4\tabcolsep) * \real{0.2975}}@{}}
\toprule\noalign{}
\begin{minipage}[b]{\linewidth}\raggedright
\textbf{Metric}
\end{minipage} & \begin{minipage}[b]{\linewidth}\raggedright
\textbf{Commitment}
\end{minipage} & \begin{minipage}[b]{\linewidth}\raggedright
\textbf{Philosophical Affinity}
\end{minipage} \\
\midrule\noalign{}
\endhead
\bottomrule\noalign{}
\endlastfoot
Euclidean & All values tradeable at par & Classical utilitarianism \\
Weighted & All values tradeable, but at different rates & Pluralistic
consequentialism \\
Lexicographic & Some values have strict priority & Deontology, rights
theories \\
Degenerate & Some values cannot be traded at all & Incommensurability
thesis \\
\end{longtable}
}

The metric is not a technical detail. It encodes \emph{fundamental value
commitments}. Two analysts who agree on all the facts --- who assign
exactly the same obligation vectors to each patient --- can disagree
about who should receive the kidney, purely because they employ
different metrics. Scalar approaches bury this disagreement in the
weighting scheme; the tensorial framework surfaces it.

\subsubsection{Off-Diagonal Components}\label{off-diagonal-components}

Suppose we have reason to believe that benefit and wait-claim are
\emph{negatively correlated} in moral significance --- that gains in one
come at the cost of the other. This is encoded by a negative
off-diagonal metric component g₁₂ \textless{} 0. The inner product

\[\langle O_{A},O_{B}\rangle = g_{\mu\nu}O_{A}^{\mu}O_{B}^{\nu}\]

now includes cross-terms that capture these interactions. Two patient
profiles with both high benefit and high wait-claim would have their
``moral distance'' from the ideal \emph{reduced} by the coupling --- the
combination is more than the sum of its parts. This is structure that no
diagonal metric (no set of independent weights) can represent.

\subsection{Level 5: Stratification and
Boundaries}\label{level-5-stratification-and-boundaries}

\subsubsection{Beyond the Smooth
Interior}\label{beyond-the-smooth-interior}

Levels 1--4 operated entirely within the smooth interior of the moral
manifold --- the region where all three patients are eligible and smooth
trade-offs are possible. But organ allocation, like all serious moral
problems, has \emph{boundaries} where the rules change discontinuously.

\subsubsection{Adding a Constraint: The Discrimination
Boundary}\label{adding-a-constraint-the-discrimination-boundary}

Suppose we discover that Carol's low ranking by some analysts is partly
based on a protected characteristic (say, her gender). The situation
crosses a stratum boundary (Chapter 5, Section 5.6): we enter the
\emph{forbidden region}.

\[C = \{ p \in M:\text{allocation weight influenced by protected characteristic}\}\]

Within C, the satisfaction function S = -∞ by convention. The affected
analysis is not merely \emph{low-value} --- it is \emph{excluded}. This
is a hard discontinuity, not a smooth trade-off. No amount of
improvement along other dimensions can compensate for a discriminatory
basis.

In the language of stratification: the smooth manifold of Level 4 was a
single stratum S₂ (the interior, where all three patients compete on
morally relevant grounds). The discrimination boundary introduces a
lower-dimensional stratum S₁ (the surface where one patient's
consideration is tainted) and the forbidden region itself is the
constraint set C.

\subsubsection{Adding a Nullifier: The Abuse
Boundary}\label{adding-a-nullifier-the-abuse-boundary}

Now suppose that during the evaluation, Bob's family threatens the
transplant committee. This activates a \emph{nullifier} --- an absorbing
stratum (Chapter 5, Definition 5.6). The Dear Abby corpus (Chapter 17)
identifies threats as universal nullifiers, and the principle
generalizes: coercion nullifies the moral force of the claims it
supports.

Bob's high wait-claim score (O²\_B = 1.0) is not reduced or discounted
--- it is \emph{annulled}. The evaluation tensor M\_\{ia\} is
restructured: all entries in the Bob row that depend on the family's
coerced advocacy are removed from consideration. The effective decision
space collapses from a 2-simplex (three patients) to a 1-simplex (Alice
vs.~Carol).

This is not a smooth operation. It is a \emph{stratum transition} --- a
discrete jump from one moral regime to another. The framework represents
this as a change in the active stratum, with the transition triggered by
a specific condition (the nullifier).

\subsubsection{Adding a Semantic Gate: The
Emergency}\label{adding-a-semantic-gate-the-emergency}

Now suppose that Alice's condition suddenly deteriorates. She is
reclassified from ``stable'' to ``urgent.'' The phrase
``life-threatening emergency'' is a \emph{semantic gate} --- a trigger
that flips the moral evaluation from one stratum to another.

Before the emergency: the allocation decision involves smooth trade-offs
among benefit, wait-time, and fairness. After the emergency: the urgency
dimension (Dimension 1, welfare/consequences) acquires a
\emph{lexicographic priority} that it did not have before. The metric
changes discontinuously:

\[g_{\mu\nu}^{\text{pre-emergency}} \rightarrow g_{\mu\nu}^{\text{emergency}}\]

where the emergency metric has g₁₁ → ∞ (or equivalently, the evaluation
becomes lexicographic with benefit/urgency first). This is a \emph{phase
transition} in the metric --- triggered by crossing a threshold, sharp
rather than gradual.

\subsubsection{What Stratification Adds}\label{what-stratification-adds}

{\def\LTcaptype{none} % do not increment counter
\begin{longtable}[]{@{}
  >{\raggedright\arraybackslash}p{(\linewidth - 4\tabcolsep) * \real{0.2094}}
  >{\raggedright\arraybackslash}p{(\linewidth - 4\tabcolsep) * \real{0.2945}}
  >{\raggedright\arraybackslash}p{(\linewidth - 4\tabcolsep) * \real{0.4931}}@{}}
\toprule\noalign{}
\begin{minipage}[b]{\linewidth}\raggedright
\textbf{Feature}
\end{minipage} & \begin{minipage}[b]{\linewidth}\raggedright
\textbf{Smooth (Levels 1--4)}
\end{minipage} & \begin{minipage}[b]{\linewidth}\raggedright
\textbf{Stratified (Level 5)}
\end{minipage} \\
\midrule\noalign{}
\endhead
\bottomrule\noalign{}
\endlastfoot
Trade-offs & Always available & Available within strata; blocked at
boundaries \\
Forbidden options & Don't exist & Hard constraints with S = -∞ \\
Nullifiers & Can't represent & Absorbing strata that restructure the
evaluation \\
Emergencies & Smooth changes in weights & Discrete phase transitions in
the metric \\
Dilemmas & Close scores & Singular points where strata intersect \\
\end{longtable}
}

Stratification is essential because moral life is not uniformly smooth.
The difference between a difficult trade-off and a forbidden action is
not one of degree --- it is a difference in \emph{kind}, a change of
stratum. Any framework that represents both as points on a continuous
scale (as scalar ethics must) conflates a structural distinction that
moral reasoning depends upon.

\subsection{Six Claims That Require
Tensors}\label{six-claims-that-require-tensors}

We now have the full apparatus. Let us identify specific claims that the
geometric framework can express and scalar ethics literally cannot.

\subsubsection{Claim 1: ``The disagreement between physicians and
families is orthogonal to the disagreement among
families.''}\label{claim-1-the-disagreement-between-physicians-and-families-is-orthogonal-to-the-disagreement-among-families.}

\textbf{Geometric expression.} Let D\_prof = M\_\{i,physician\} -
M\_\{i,committee\} be the professional disagreement vector and D\_fam =
Var\_a{[}M\_\{i,a\}{]} for family agents. These vectors in option-space
can be orthogonal:

\[g_{ij}D_{\text{prof}}^{i}D_{\text{fam}}^{j} = 0\]

meaning the dimensions on which professionals disagree are independent
of the dimensions on which families disagree.

\textbf{Why scalars cannot say this.} A scalar framework has no
``direction'' of disagreement --- just a magnitude. The structural claim
about orthogonality is inexpressible.

\subsubsection{Claim 2: ``Alice is the robust choice; Bob wins only
under specific weight
configurations.''}\label{claim-2-alice-is-the-robust-choice-bob-wins-only-under-specific-weight-configurations.}

\textbf{Geometric expression.} Alice's decision region in interest space
(the set of covectors I for which I\_μ O\^{}μ\_A \textgreater{} I\_μ
O\^{}μ\_B and I\_μ O\^{}μ\_A \textgreater{} I\_μ O\^{}μ\_C) has larger
volume than Bob's. The robustness of a choice is the \emph{measure} of
its decision region.

\textbf{Why scalars cannot say this.} A scalar gives one answer with one
weight vector. Assessing robustness requires the full vector structure.

\subsubsection{Claim 3: ``Choosing Carol over Bob wrongs Bob in a way
that choosing Alice over Bob does
not.''}\label{claim-3-choosing-carol-over-bob-wrongs-bob-in-a-way-that-choosing-alice-over-bob-does-not.}

\textbf{Geometric expression.} The difference vector O\_B - O\_C has a
large positive component on Dimension 2 (wait-claim): Bob has waited
much longer than Carol. Bob \emph{dominates} Carol on this dimension.
The difference O\_B - O\_A is more balanced: Bob leads on wait-claim,
Alice leads on benefit and care. Choosing Alice over Bob involves
genuine trade-offs; choosing Carol over Bob ignores a dimension on which
Bob has an unambiguous advantage.

\textbf{Why scalars cannot say this.} The scalar reports only that Carol
has the lowest score. It cannot distinguish ``legitimate trade-off
loss'' from ``dominated on a morally relevant dimension.''

\subsubsection{Claim 4: ``A utilitarian and a Rawlsian agree on this
case, but for different reasons, and would diverge if Carol's urgency
increased.''}\label{claim-4-a-utilitarian-and-a-rawlsian-agree-on-this-case-but-for-different-reasons-and-would-diverge-if-carols-urgency-increased.}

\textbf{Geometric expression.} Both the utilitarian contraction (sum)
and the Rawlsian contraction (min) yield Alice. But the utilitarian is
sensitive to total satisfaction; the Rawlsian to minimum satisfaction.
We can compute the \emph{exact threshold} at which modifying Carol's
vector causes the two procedures to diverge --- the boundary between
their agreement region and their disagreement region in the space of
possible cases.

\textbf{Why scalars cannot say this.} A scalar framework implements one
theory at a time. Comparing theories requires the multi-agent tensor
that represents both as contractions of the same underlying structure.

\subsubsection{Claim 5: ``The metric choice is doing more work than the
interest
choice.''}\label{claim-5-the-metric-choice-is-doing-more-work-than-the-interest-choice.}

\textbf{Geometric expression.} Under Euclidean vs.~lexicographic metrics
with the \emph{same} interest covector, we get different decisions
(Alice vs.~Bob). The metric determines how dimensions \emph{combine};
the interest determines how dimensions are \emph{weighted}. These are
independent structural choices.

\textbf{Why scalars cannot say this.} Scalar frameworks conflate metric
and weights. The weighted sum S = Σw\_i x\_i is simultaneously a weight
choice and an implicit metric choice (Euclidean with weighting). The two
cannot be separated.

\subsubsection{Claim 6: ``The case is genuinely hard in a way that more
information cannot
resolve.''}\label{claim-6-the-case-is-genuinely-hard-in-a-way-that-more-information-cannot-resolve.}

\textbf{Geometric expression.} The obligation vectors O\_A, O\_B, O\_C
are in general position --- no dominance relations, no single dimension
that settles the question. The Gram matrix g\_μν O\^{}μ\_i O\^{}ν\_j
shows that the vectors span a two-dimensional subspace without a clear
ordering. The hardness is \emph{geometric}: it resides in the relative
positions of the vectors, not in the precision of the data.

\textbf{Why scalars cannot say this.} A scalar framework always delivers
an answer (the highest score). It cannot represent ``this case is
structurally hard'' --- only ``the scores are close.'' But closeness is
not the same as genuine value conflict. Only the geometric framework
distinguishes them.

\subsection{Synthesis}\label{synthesis}

Let us trace what each level of structure added:

{\def\LTcaptype{none} % do not increment counter
\begin{longtable}[]{@{}
  >{\raggedright\arraybackslash}p{(\linewidth - 4\tabcolsep) * \real{0.1617}}
  >{\raggedright\arraybackslash}p{(\linewidth - 4\tabcolsep) * \real{0.2305}}
  >{\raggedright\arraybackslash}p{(\linewidth - 4\tabcolsep) * \real{0.6078}}@{}}
\toprule\noalign{}
\begin{minipage}[b]{\linewidth}\raggedright
\textbf{Level}
\end{minipage} & \begin{minipage}[b]{\linewidth}\raggedright
\textbf{Mathematical Object}
\end{minipage} & \begin{minipage}[b]{\linewidth}\raggedright
\textbf{What It Adds}
\end{minipage} \\
\midrule\noalign{}
\endhead
\bottomrule\noalign{}
\endlastfoot
. Scalar & S ∈ ℝ & A decision --- but no justification, no robustness,
no transparency \\
. Vector & O\^{}μ ∈ T\_pM & The \emph{reasons}: which dimensions
generate obligations; decision regions in interest space; trade-off
geometry \\
. Rank-2 tensor & M\_\{ia\} = I\^{}\{(a)\}\_μ O\^{}μ\_i & Multiple
perspectives; the landscape of agreement and disagreement; aggregation
as a visible choice \\
. Metric & g\_μν & The trade-off structure: commensurability,
incommensurability, priority; explicit value commitments \\
. Stratification & \{S\_α\}, ∂C, nullifiers & Boundaries, forbidden
regions, phase transitions; the distinction between trade-offs and
prohibitions \\
\end{longtable}
}

Each level is necessary. The claims expressible at Level 5 cannot be
stated in Level 1 language. The geometric framework is not a
complication --- it is the \emph{minimum} structure needed to say what
we want to say about hard cases.

\subsection{Coda: The Same Case, Revisited
Forever}\label{coda-the-same-case-revisited-forever}

We could continue enriching this case indefinitely:

\begin{itemize}
\item
  \textbf{Add time.} How do obligations change as patients wait?
  Transport the obligation vector O(t₀) along the path of increasing
  wait-time and ask how it transforms. If the result depends on what
  other events occur along the way (a new diagnosis, a family change),
  the manifold has curvature. This is the domain of Chapter 10 (parallel
  transport and holonomy).
\item
  \textbf{Add uncertainty.} What if medical benefit estimates carry
  error bars? The uncertainty tensor Σ\^{}\{μν\} (Chapter 6, Section
  6.6) captures the \emph{shape} of this uncertainty --- its principal
  directions, its alignment with the dimensions that matter most.
  Risk-averse allocation policies correspond to contractions that
  penalize high Σ along decision-relevant directions.
\item
  \textbf{Add dynamics.} The transplant committee's metric may evolve
  over time as social priorities shift. The \emph{connection} on the
  moral manifold (Chapter 10) specifies how the metric is transported
  from one era to another, and the \emph{curvature} measures whether
  this transport is path-dependent.
\item
  \textbf{Add symmetry.} The requirement that the allocation be
  invariant under morally irrelevant re-descriptions (swapping Alice's
  and Bob's names, translating the case into another language) is the
  Bond Invariance Principle --- a gauge symmetry whose conservation law
  (Noether's theorem, Chapter 12) implies that harm accounting is
  invariant across representations.
\item
  \textbf{Add quantum structure.} What if the committee's deliberation
  involves superposition --- simultaneously entertaining two
  incompatible framings of the case? The framework of Chapter 13
  (quantum normative dynamics) models this as a superposition of moral
  states that collapses, upon deliberation, to a definite verdict.
\end{itemize}

Each extension adds mathematical structure --- and each structure lets
us express claims we could not express before.

❖

\emph{One kidney. Three patients. Five levels of structure.}

\emph{The same case, seen more clearly each time --- not by adding
information about the patients, but by adding precision to the language
of evaluation.}

\emph{The mathematics is not imposed on the ethics. It is extracted from
it.}

\emph{We use geometry because the claims we want to make about moral
situations require geometric structure to express. The kidney case is
one case. But it contains, in miniature, the geometry of moral reasoning
itself.}

\section{Chapter 8: Stratification --- Boundaries, Thresholds, and Phase
Transitions}\label{chapter-8-stratification-boundaries-thresholds-and-phase-transitions}

\begin{quote}
\textbf{\textsc{RUNNING EXAMPLE --- Priya's Model}}

\emph{The 70-point threshold in TrialMatch is a first-type boundary: a
threshold. Above 70, patients enter the `trial-eligible' stratum where
new rules apply---informed consent protocols, monitoring, experimental
risk. Below 70, standard care. Priya identifies all four boundary types.
Threshold: the 70-point cutoff. Phase transition: when BEACON-7 moves
from Phase II to Phase III, eligibility criteria change discontinuously.
Absorbing stratum: patient death---once crossed, the harm of
non-matching is irreversible. Constraint surface: HealthBridge's budget
for patient travel support, a soft boundary that limits but does not
forbid. The 70-point threshold is ethically arbitrary. The absorbing
boundary of death is not.}
\end{quote}

\subsection{8.1 The Patchwork Structure of Moral
Space}\label{the-patchwork-structure-of-moral-space}

\includegraphics[width=3.39028in,height=3.2125in,alt={image}]{media/image8.png}The
preceding chapters have developed the smooth geometry of the moral
manifold: tangent spaces, tensor fields, metrics, contractions. This
apparatus is powerful within any single moral \emph{regime} --- any
region of moral space where smooth trade-offs apply and the rules are
uniform. But moral space is not uniformly smooth. It is a
\emph{patchwork}: regions of smooth evaluation joined along boundaries
where the rules change discontinuously.

The difference between a difficult trade-off and a hard prohibition is
not a matter of degree. It is a difference in kind --- a boundary
between two moral regimes, each internally coherent, governed by
different rules. Crossing such a boundary changes the operative
framework: what was permissible becomes forbidden, what was optional
becomes mandatory, what was a matter of degree becomes a matter of
category.

\subsection{8.2 Whitney Stratification: The Mathematical
Framework}\label{whitney-stratification-the-mathematical-framework}

\subsubsection{From Manifolds to Stratified
Spaces}\label{from-manifolds-to-stratified-spaces}

A schematic 3D stratified space
\(\left( M,\{ M_{i}\}, \preccurlyeq \right)\)depicting ethically
relevant configurations. The horizontal axes parameterize two
illustrative ethical dimensions---\textbf{Harm} (lower is better) and
\textbf{Fairness}---while the vertical axis denotes \textbf{Ethical
satisfaction} (a scalar proxy for overall moral value). The full volume
represents the \textbf{total moral space} \(M\). Three representative
strata \(M_{1},M_{2},M_{3}\)(semi-transparent surfaces) illustrate
regions in which ethical trade-offs vary smoothly: within any single
stratum, nearby configurations admit continuous, ``manifold-like''
adjustments (e.g., small changes in harm can be compensated by changes
in fairness without qualitative shifts in moral kind). \textbf{Stratum
boundaries} (annotated as \textbf{Moral discontinuities}) mark loci
where smooth trade-offs fail---corresponding to dilemmas, sharp regime
changes, or categorical distinctions; these are indicated by ridge/edge
features and highlighted with dashed guides and arrows. \textbf{Phase
transition} denotes a boundary across which the governing ethical regime
changes, while \textbf{Categorical boundaries} indicate constraints that
separate morally distinct classes of actions or outcomes. Example
configurations \textbf{A} and \textbf{B} (marked points) lie on
different strata; the vertical dashed projections indicate their
coordinates in \(\left( \text{Harm},\text{Fairness} \right)\)and
emphasize that moving between strata requires crossing a discontinuity
rather than a smooth deformation within a single \(M_{i}\).

\textbar{} Stratified moral space with moral discontinuities

\begin{itemize}
\item
  Regions of full dimensionality (nine-dimensional in the generic case
  of Chapter 5) where all moral dimensions are independently variable
\item
  Lower-dimensional boundaries where some degrees of freedom are
  constrained or eliminated
\item
  Singular points where the geometric structure breaks down entirely
\end{itemize}

A \emph{stratified space} relaxes the uniformity requirement while
preserving enough structure for differential calculus to work on each
stratum and for the strata to fit together coherently.

\subsubsection{Formal Definition}\label{formal-definition-1}

\textbf{Definition 8.1 (Stratification).} \emph{A stratification of a
topological space} \(X\) \emph{is a locally finite partition}
\(\{ S_{\alpha}\}_{\alpha \in A}\) \emph{into smooth manifolds (called
strata) of various dimensions, together with a partial order}
\(\preccurlyeq\) \emph{on the index set} \(A\)\emph{, satisfying:}

\begin{enumerate}
\def\labelenumi{\arabic{enumi}.}
\item
  \textbf{Frontier condition:}
  \(S_{\alpha} \cap \overline{S_{\beta}} \neq \varnothing\) implies
  \(S_{\alpha} \subset \overline{S_{\beta}}\) and
  \(\alpha \preccurlyeq \beta\).
\item
  \textbf{Local finiteness:} Each point of \(X\) has a neighborhood that
  meets only finitely many strata.
\end{enumerate}

The frontier condition ensures a clean nesting: if a stratum touches the
closure of a higher-dimensional stratum, it lies entirely in that
closure. There are no ``dangling'' pieces.

\subsubsection{Whitney's Conditions}\label{whitneys-conditions}

Hassler Whitney (1965) introduced regularity conditions that ensure the
strata fit together in a geometrically well-behaved way. These
conditions prevent pathological configurations --- strata that spiral
toward each other, or boundaries where tangent planes are undefined.

\textbf{Definition 8.2 (Whitney Condition A).} \emph{Let} \(S_{\alpha}\)
\emph{and} \(S_{\beta}\) \emph{be strata with}
\(S_{\alpha} \subset \overline{S_{\beta}}\)\emph{. Whitney condition (A)
holds at a point} \(x \in S_{\alpha}\) \emph{if: for every sequence}
\(\{ y_{n}\} \subset S_{\beta}\) \emph{converging to} \(x\)\emph{, if
the sequence of tangent spaces} \(T_{y_{n}}S_{\beta}\) \emph{converges
to a limit} \(\tau\)\emph{, then}
\(T_{x}S_{\alpha} \subset \tau\)\emph{.}

In words: as you approach a lower stratum from a higher one, the tangent
planes of the higher stratum must contain, in the limit, the tangent
plane of the lower stratum. The boundary stratum's tangent directions
are a \emph{subset} of the limiting tangent directions of the ambient
stratum.

\textbf{Definition 8.3 (Whitney Condition B).} \emph{With notation as
above, if additionally} \(\{ x_{n}\} \subset S_{\alpha}\) \emph{is a
sequence converging to} \(x\)\emph{, and the secant lines}
\(\overline{x_{n}y_{n}}\) \emph{converge to a line}
\(\mathcal{l}\)\emph{, then} \(\mathcal{l \subset}\tau\)\emph{.}

Condition (B) is stronger: not only must tangent planes be well-behaved,
but the ``angle of approach'' from the higher stratum to the lower
stratum must also be consistent. Together, conditions (A) and (B)
guarantee that the stratified space admits tubular neighborhoods, normal
bundles, and controlled deformation --- the essential tools for doing
geometry near boundaries.

\subsubsection{Moral Interpretation}\label{moral-interpretation}

Why do these technical conditions matter for ethics?

Because they ensure that moral evaluations \emph{near} a boundary are
well-behaved. If Whitney's conditions fail, small perturbations near a
boundary can produce wild oscillations in moral evaluation --- a
situation with no ethical analogue. Real moral boundaries are sharp but
\emph{approachable}: as a situation moves toward the consent threshold,
the moral evaluation changes smoothly right up to the boundary, where it
jumps. The jump is discontinuous, but the approach is controlled. This
is precisely what Whitney's conditions guarantee.

\textbf{Example.} Consider the boundary between ``sufficient informed
consent'' and ``insufficient informed consent'' in medical ethics. On
the high-dimensional stratum (sufficient consent), smooth trade-offs
apply: the physician may weigh benefit against risk, timeliness against
thoroughness of disclosure. On the boundary (consent just barely
sufficient), these trade-offs still apply but become more constrained.
Below the boundary (insufficient consent), the procedure is
impermissible --- the satisfaction function jumps to \(- \infty\).

Whitney condition (A) ensures that the directions of smooth variation
available just above the boundary \emph{include} the directions
available on the boundary itself. Condition (B) ensures that the
geometry of the approach --- whether you approach the boundary from the
``more information'' side or the ``more understanding'' side --- is
consistent.

\subsection{8.3 A Taxonomy of Moral
Boundaries}\label{a-taxonomy-of-moral-boundaries}

\begin{figure}
\centering
\includesvg[width=6.5in,height=5.28125in]{media/image10.svg}
\caption{\textbar{} Stratified Moral Space: Four Boundary Types}
\end{figure}

Figure 7 provides a visual overview of the four boundary types. Not all
stratum boundaries are alike. This section develops a taxonomy of the
boundary types that appear in moral space.

\subsubsection{Type I: Thresholds}\label{type-i-thresholds}

A \emph{threshold} is a codimension-1 boundary where a continuous
parameter crosses a critical value, producing a discrete change in moral
status.

\textbf{Definition 8.4 (Threshold Boundary).} \emph{A threshold boundary
is a hypersurface} \(B\mathcal{\subset M}\) \emph{such that:} 1.
\emph{On each side of} \(B\)\emph{, the satisfaction function} \(S\)
\emph{and the metric} \(g_{\mu\nu}\) \emph{are smooth.} 2. \emph{Across}
\(B\)\emph{, either} \(S\) \emph{or} \(g_{\mu\nu}\) \emph{(or both) may
be discontinuous.} 3. \(B\) \emph{is the level set of a smooth function}
\(\phi\mathcal{:M \rightarrow}\mathbb{R}\)\emph{, with}
\(B = \phi^{- 1}(0)\)\emph{.}

Condition (3) ensures the boundary is ``smooth as a surface'' even
though the moral evaluation jumps across it. The threshold function
\(\phi\) measures distance from the boundary; the sign of \(\phi\)
determines which regime applies.

\textbf{Examples of threshold boundaries:}

{\def\LTcaptype{none} % do not increment counter
\begin{longtable}[]{@{}
  >{\raggedright\arraybackslash}p{(\linewidth - 6\tabcolsep) * \real{0.1989}}
  >{\raggedright\arraybackslash}p{(\linewidth - 6\tabcolsep) * \real{0.2076}}
  >{\raggedright\arraybackslash}p{(\linewidth - 6\tabcolsep) * \real{0.2980}}
  >{\raggedright\arraybackslash}p{(\linewidth - 6\tabcolsep) * \real{0.2955}}@{}}
\toprule\noalign{}
\begin{minipage}[b]{\linewidth}\raggedright
\textbf{Threshold}
\end{minipage} & \begin{minipage}[b]{\linewidth}\raggedright
\textbf{Variable}
\end{minipage} & \begin{minipage}[b]{\linewidth}\raggedright
\textbf{Below}
\end{minipage} & \begin{minipage}[b]{\linewidth}\raggedright
\textbf{Above}
\end{minipage} \\
\midrule\noalign{}
\endhead
\bottomrule\noalign{}
\endlastfoot
Age of majority & Age & Minor (restricted autonomy) & Adult (full
autonomy) \\
Informed consent & Understanding & Impermissible treatment & Permissible
treatment \\
Poverty line & Income & Entitled to assistance & Not entitled \\
Lethality & Harm severity & Battery (civil wrong) & Homicide (criminal
offense) \\
Capacity & Cognitive function & Guardian required &
Self-determination \\
\end{longtable}
}

Each threshold corresponds to a stratum boundary where the applicable
moral rules change. On each side, smooth variation is possible (the
17-year-old has more autonomy than the 12-year-old; the 19-year-old has
more than the 18-year-old). But at the boundary, the \emph{category}
changes, and with it the set of applicable rules.

\subsubsection{Type II: Phase
Transitions}\label{type-ii-phase-transitions}

A \emph{phase transition} is a boundary where the \emph{structure} of
the moral evaluation changes --- not just its value, but the operative
metric, the dimensionality of the effective space, or the nature of the
permissible trade-offs.

\textbf{Definition 8.5 (Phase Transition Boundary).} \emph{A phase
transition boundary is a stratum boundary across which the moral metric}
\(g_{\mu\nu}\) \emph{changes discontinuously in rank or signature.}

Phase transitions are more dramatic than simple thresholds. At a
threshold, the same moral framework applies on both sides, with
different conclusions. At a phase transition, the \emph{framework
itself} changes: different dimensions become relevant, different
trade-offs become available, the very character of moral deliberation
shifts.

\textbf{Example: The Emergency Phase Transition.} In Chapter 7 (§Level
5), we saw that a medical emergency transforms the allocation metric:

\[g_{\mu\nu}^{\text{routine}} \rightarrow g_{\mu\nu}^{\text{emergency}}\]

In routine mode, the metric allows smooth trade-offs among benefit,
wait-time, fairness, and other dimensions. In emergency mode, the
welfare/urgency dimension acquires lexicographic priority:
\(g_{11} \rightarrow \infty\) (or equivalently, the metric becomes
degenerate on the non-urgency subspace). This is not a change in the
value of \(S\) --- it is a change in the \emph{geometry} of the space on
which \(S\) is defined.

\textbf{Example: The Wartime Phase Transition.} In the parable of
Chapter 2, the declaration of war is a phase transition. Before the
declaration, the moral evaluation of the son's broken leg is determined
by the peacetime metric (health, capability, work). After the
declaration, the wartime metric gives decisive weight to the
conscription/survival dimension, which was previously inactive. The
effective dimensionality of the moral space increases: a dimension that
was locally constant (no war) becomes variable and dominant.

\textbf{Classification.} Borrowing from statistical physics, we can
distinguish:

\begin{itemize}
\item
  \textbf{First-order phase transitions}: The satisfaction function
  \(S\) is discontinuous across the boundary. The ``moral energy''
  jumps. Example: the consent threshold --- below it, \(S = - \infty\);
  above it, \(S\) takes a finite value.
\item
  \textbf{Second-order phase transitions}: \(S\) is continuous across
  the boundary but its derivatives (the gradient \(\nabla S\), the
  curvature) are discontinuous. The evaluation does not jump, but its
  \emph{sensitivity} changes sharply. Example: crossing from a context
  where promise-keeping is a minor consideration to one where it is
  central --- the value \(S\) may change continuously, but
  \(\partial S/\partial x^{2}\) (sensitivity to the duty dimension)
  jumps.
\end{itemize}

\subsubsection{Type III: Absorbing Strata
(Nullifiers)}\label{type-iii-absorbing-strata-nullifiers}

An \emph{absorbing stratum} is a lower-dimensional region that, once
entered, determines the moral evaluation regardless of the original
context. Entry is irreversible in the sense that the original stratum's
rules cease to apply.

\textbf{Definition 8.6 (Absorbing Stratum).} \emph{A stratum}
\(S_{0}\mathcal{\subset M}\) \emph{is absorbing if there exists a
tubular neighborhood} \(N\) \emph{of} \(S_{0}\) \emph{and a continuous
retraction} \(r:N \rightarrow S_{0}\) \emph{such that, for any moral
evaluation function} \(S\) \emph{compatible with the stratification,}
\(S|_{N}\) \emph{factors through} \(r\)\emph{:}

\[S(p) = S(r(p))\quad\text{for all }p \in N\]

\emph{In words: near an absorbing stratum, the evaluation depends only
on the projection onto} \(S_{0}\)\emph{, not on the transverse
coordinates.}

Absorbing strata are the geometric formalization of \emph{nullifiers}
--- conditions that override all other moral considerations. The most
important empirical finding about nullifiers comes from analysis of the
Dear Abby corpus (Chapter 17), which identifies three universal
nullifiers across all domain contexts:

{\def\LTcaptype{none} % do not increment counter
\begin{longtable}[]{@{}
  >{\raggedright\arraybackslash}p{(\linewidth - 6\tabcolsep) * \real{0.1521}}
  >{\raggedright\arraybackslash}p{(\linewidth - 6\tabcolsep) * \real{0.1757}}
  >{\raggedright\arraybackslash}p{(\linewidth - 6\tabcolsep) * \real{0.1819}}
  >{\raggedright\arraybackslash}p{(\linewidth - 6\tabcolsep) * \real{0.4891}}@{}}
\toprule\noalign{}
\begin{minipage}[b]{\linewidth}\raggedright
\textbf{Nullifier}
\end{minipage} & \begin{minipage}[b]{\linewidth}\raggedright
\textbf{Corpus Count}
\end{minipage} & \begin{minipage}[b]{\linewidth}\raggedright
\textbf{Cross-Domain}
\end{minipage} & \begin{minipage}[b]{\linewidth}\raggedright
\textbf{Mechanism}
\end{minipage} \\
\midrule\noalign{}
\endhead
\bottomrule\noalign{}
\endlastfoot
Abuse & n = 582 & Universal & Collapses obligation structure \\
Danger & n = 218 & Universal & Activates self-preservation override \\
Impossibility & n = 144 & Universal & \emph{Ought implies can}: removes
infeasible obligations \\
\end{longtable}
}

When abuse is present, it does not \emph{reduce} the weight of other
obligations --- it \emph{annuls} them. The obligation to maintain a
family relationship, to fulfill a promise, to support a colleague ---
all these obligations, which may be strong in the non-abusive stratum,
collapse to zero (or reverse into obligations to \emph{exit}) when the
situation enters the abuse stratum. This is not a smooth re-weighting;
it is a stratum transition.

The absorption property is important: once in the abuse stratum, the
evaluation is determined by the stratum's own rules. The transverse
coordinates --- how strong the promise was, how much the relationship
previously meant, how inconvenient departure would be --- cease to
matter. This is the mathematical content of the moral intuition that
``nothing justifies staying in an abusive relationship.''

\subsubsection{Type IV: Constraint Surfaces (Forbidden
Regions)}\label{type-iv-constraint-surfaces-forbidden-regions}

A \emph{constraint surface} marks the boundary of the forbidden region
--- the set of moral states that are absolutely impermissible.

\textbf{Definition 8.7 (Constraint Surface).} \emph{The constraint set}
\(C\mathcal{\subset M}\) \emph{is a closed subset. The constraint
surface} \(\partial C\) \emph{is the boundary of} \(C\) \emph{in}
\(\mathcal{M}\)\emph{. The satisfaction function satisfies}
\(S|_{C} = - \infty\)\emph{.}

The constraint set represents moral non-starters: options that are
forbidden regardless of consequences, regardless of other
considerations, regardless of the metric. In the language of
deontological ethics, these are \emph{side constraints} --- not
trade-off-able values but hard limits.

The constraint surface \(\partial C\) is a stratum boundary of a
specific kind: on one side, finite-valued evaluation with smooth
trade-offs; on the other, \(S = - \infty\) with no trade-offs possible.
The gradient \(\nabla S\) does not merely become large near
\(\partial C\) --- it becomes formally infinite, creating a potential
wall that no finite benefit can overcome.

\textbf{Relation to Nussbaum's thresholds.} Martha Nussbaum's capability
thresholds (Chapter 3) are constraint surfaces: below the minimum level
of each capability, a life is not fully human, and the imperative is to
raise capabilities to the threshold. The geometric framework generalizes
Nussbaum's insight by representing thresholds as stratum boundaries with
specific mathematical properties (absorption, discontinuity of \(S\),
degeneracy of \(g\)).

\subsection{8.4 Semantic Gates: The Transition
Functions}\label{semantic-gates-the-transition-functions}

\subsubsection{From Stratum to Stratum}\label{from-stratum-to-stratum}

The boundaries described in §8.3 are \emph{where} the rules change.
Semantic gates are \emph{how} and \emph{why}: the specific features of a
situation that trigger a stratum transition.

\textbf{Definition 8.8 (Semantic Gate).} \emph{A semantic gate is a map}
\(G\mathcal{:F \times}\mathcal{S}_{\alpha} \rightarrow \mathcal{S}_{\beta}\)
\emph{where} \(\mathcal{F}\) \emph{is a space of triggering features,}
\(\mathcal{S}_{\alpha}\) \emph{is the source stratum, and}
\(\mathcal{S}_{\beta}\) \emph{is the target stratum. The gate fires when
a feature} \(f\mathcal{\in F}\) \emph{is present, transitioning the
moral evaluation from the rules of} \(\mathcal{S}_{\alpha}\) \emph{to
the rules of} \(\mathcal{S}_{\beta}\)\emph{.}

Formal trigger condition. The informal phrase ``the gate fires when a
feature f is present'' is shorthand for: the gate is activated when the
triggering feature f ∈ F satisfies the gate's activation predicate σ(f)
= 1, where σ: F → \{0, 1\} is the gate's characteristic function. The
domain F of triggering features is a measurable subset of the space of
morally relevant linguistic and situational cues at the stratum
boundary. The gate map G is defined on the activated domain F¹ = σ⁻¹(1)
× S\_α, not on all of F × S\_α.

The term ``semantic gate'' is chosen deliberately, by analogy with logic
gates in digital circuits. A logic gate takes binary inputs and produces
a binary output; a semantic gate takes a moral feature and produces a
discrete state transition. The key property is \emph{discreteness}: the
transition is all-or-nothing, a step function, not a smooth sigmoid.

\subsubsection{Empirical Evidence for
Discreteness}\label{empirical-evidence-for-discreteness}

The claim that semantic gates are discrete is empirical, not
stipulative. It is confirmed by analysis of the Dear Abby corpus and by
the BIP (Bond Invariance Principle) experiments across 109,294 passages
in 11 languages (Chapter 17).

\textbf{Gate detection results from the Dear Abby corpus:}

{\def\LTcaptype{none} % do not increment counter
\begin{longtable}[]{@{}
  >{\raggedright\arraybackslash}p{(\linewidth - 6\tabcolsep) * \real{0.2970}}
  >{\raggedright\arraybackslash}p{(\linewidth - 6\tabcolsep) * \real{0.2536}}
  >{\raggedright\arraybackslash}p{(\linewidth - 6\tabcolsep) * \real{0.1680}}
  >{\raggedright\arraybackslash}p{(\linewidth - 6\tabcolsep) * \real{0.1577}}@{}}
\toprule\noalign{}
\begin{minipage}[b]{\linewidth}\raggedright
\textbf{Trigger Phrase}
\end{minipage} & \begin{minipage}[b]{\linewidth}\raggedright
\textbf{Transition}
\end{minipage} & \begin{minipage}[b]{\linewidth}\raggedright
\textbf{Effectiveness}
\end{minipage} & \begin{minipage}[b]{\linewidth}\raggedright
\textbf{Character}
\end{minipage} \\
\midrule\noalign{}
\endhead
\bottomrule\noalign{}
\endlastfoot
``You promised'' & Liberty → Obligation & 94\% & Step function \\
``Only if convenient'' & Obligation → Liberty & 89\% & Step function \\
``He hit you'' / abuse signal & Any → Nullified & 97\% & Step
function \\
``Life-threatening'' & Routine → Emergency & 92\% & Step function \\
``You agreed to'' & Liberty → Obligation & 91\% & Step function \\
\end{longtable}
}

The ``effectiveness'' column measures how reliably the trigger phrase
causes the state transition in natural moral reasoning. The
``character'' column reports the empirical finding that these
transitions behave as step functions --- binary switches --- rather than
smooth transitions. There is no ``50\% obligated'' intermediate state.
One is either in the obligation stratum or the liberty stratum; the gate
flips the state discretely.

\textbf{Cross-linguistic evidence.} The BIP experiments reveal that the
deontic axis --- the obligation/permission distinction --- transfers
across languages at 100\% accuracy even when other aspects of moral
content fail to transfer. This suggests that the \emph{stratum
structure} of moral space (the partition into obligation and liberty
regions) is more universal than the \emph{metric} (how considerations
are weighted within each stratum). The boundaries are structural
invariants; the smooth interior varies with culture and context.

\subsubsection{The D₄ Group Structure of Hohfeldian
Transitions}\label{the-dux2084-group-structure-of-hohfeldian-transitions}

The transitions between Hohfeldian jural states --- Obligation (O),
Claim (C), Liberty (L), No-claim (N) --- have a precise algebraic
structure: they form the \emph{dihedral group} \(D_{4}\) of order 8, the
symmetry group of the square.

\begin{quote}
\includesvg[width=4in,height=4in]{media/image12.svg}\textbf{The
Hohfeldian Square:}

\begin{verbatim}
    O ———— C
    |      |
    |      |
    L ———— N
\end{verbatim}

The eight elements of \(D_{4}\) are:
\end{quote}

{\def\LTcaptype{none} % do not increment counter
\begin{longtable}[]{@{}
  >{\raggedright\arraybackslash}p{(\linewidth - 4\tabcolsep) * \real{0.1161}}
  >{\raggedright\arraybackslash}p{(\linewidth - 4\tabcolsep) * \real{0.2592}}
  >{\raggedright\arraybackslash}p{(\linewidth - 4\tabcolsep) * \real{0.3652}}@{}}
\toprule\noalign{}
\begin{minipage}[b]{\linewidth}\raggedright
\textbf{Element}
\end{minipage} & \begin{minipage}[b]{\linewidth}\raggedright
\textbf{Operation}
\end{minipage} & \begin{minipage}[b]{\linewidth}\raggedright
\textbf{Moral Meaning}
\end{minipage} \\
\midrule\noalign{}
\endhead
\bottomrule\noalign{}
\endlastfoot
e & Identity & No change \\
r & ° rotation & O → C → N → L → O \\
r² & ° rotation & O ↔ N, C ↔ L (full inversion) \\
r³ & ° rotation & O → L → N → C → O \\
s & Vertical reflection & O ↔ C, L ↔ N (correlative swap) \\
sr & Diagonal reflection & C ↔ L; O, N fixed \\
sr² & Horizontal reflection & O ↔ L, C ↔ N (jural swap) \\
sr³ & Anti-diagonal reflection & O ↔ N; C, L fixed \\
\end{longtable}
}

The group relations are \(r^{4} = e\), \(s^{2} = e\), and
\(srs = r^{- 1}\). These are verified by the test suite of the Dear
Ethicist experimental platform, which confirms the axioms hold exactly
for moral state transitions.

\textbf{The Correlative Reflection.} The most important group element is
the reflection \(s\), which swaps correlatives: O ↔ C and L ↔ N. The
\emph{correlative symmetry} principle states that Hohfeldian positions
come in linked pairs:

\begin{itemize}
\item
  If Alice has an \emph{obligation} to Bob, then Bob has a \emph{claim}
  against Alice.
\item
  If Alice is at \emph{liberty} regarding Bob, then Bob has \emph{no
  claim} against Alice.
\end{itemize}

Empirically, this symmetry holds at high rates:

\[\text{O} \leftrightarrow \text{C: }87\%\quad\quad\text{L} \leftrightarrow \text{N: }82\%\]

The deviation from 100\% is informative: it measures the rate at which
natural moral reasoning violates structural constraints. The violations
are not random; they concentrate in specific situations (e.g., where
asymmetric power relations distort the perceived correlative structure).

\textbf{Why} \(D_{4}\)\textbf{?} The fact that Hohfeldian transitions
form the symmetry group of the square is not an accident. It reflects a
deep structural feature: the moral positions are related by two
independent binary oppositions --- the correlative opposition (O ↔ C, L
↔ N) and the jural opposition (O ↔ L, C ↔ N) --- and these two
oppositions generate a group isomorphic to \(D_{4}\). The square is the
simplest figure with exactly two independent reflection symmetries.

\textbf{Semantic gates as group elements.} Each semantic gate
corresponds to a specific element of \(D_{4}\):

\begin{itemize}
\item
  ``You promised'' corresponds to the transition L → O, which is the
  rotation r (since r maps L → O).\(r^{3}\)
\item
  ``Only if convenient'' corresponds to O → L, which is the rotation r³
  (the inverse of r).\(r\)
\item
  The correlative constraint (O implies C for the other party)
  corresponds to the reflection \(s\).
\end{itemize}

The gate structure is thus \emph{algebraic}: gates compose according to
the group multiplication table of D₄. Two\(D_{4}\) successive L → O
transitions compose to r² (which\(r^{3} \cdot r^{3} = r^{2}\) maps L → C
--- a double promise creates a claim, not merely an obligation). This
compositional structure is testable and opens a path toward formal
verification of moral reasoning systems.

\textbf{Generator asymmetry.} {[}Empirical result (preliminary).{]} The
two generators of D₄ have fundamentally different operational
characters:

\textbf{Reflection s} (perspective shift, O ↔ C, L ↔ N) is non-gated,
immediate, and free. It requires no semantic trigger: any competent
moral reasoner can adopt the other party's perspective at will. This
explains the 100\% correlative symmetry rate observed in §17.3.

\textbf{Rotation r} (state transition, O → C → N → L → O) is discrete
and gated. It requires specific linguistic triggers (``You promised,''
``Only if convenient,'' ``He hit you'') to fire, and the effectiveness
is probabilistic (89--97\%; see §8.4 gate detection results).

This asymmetry has a structural explanation: s corresponds to viewing
the same relationship from the other side---a change of frame, not a
change of state. The rotation r corresponds to changing the moral state
itself, which requires an event (a promise, a release, an escalation) to
occur. The distinction is analogous to the difference between a passive
coordinate transformation (always free) and an active physical
transformation (requires energy or a force). See §12.3 for the full
gauge group identification.

Computational instantiation (February 2026). The DEME V3 reference
implementation includes a D₄ Hohfeldian gauge structure demonstration
(hohfeld\_d4\_demo.py) that exercises the full algebraic structure
described in this section. The demo implements semantic gates as D₄
group elements: ``only if convenient'' applies the negation r² (mapping
O → L), ``I promise'' applies the rotation r (mapping L → O), and ``from
their perspective'' applies the reflection s (the correlative swap O ↔
C, L ↔ N). The compositional structure is verified computationally:
applying two semantic gates in sequence produces the result predicted by
D₄ group multiplication, confirming that gate composition follows the
group law. The generator asymmetry --- reflection s as free perspective
swap vs. rotation r as gated state transition --- is operationalized in
the apply\_semantic\_gate() function, which maps linguistic triggers to
their corresponding D₄ elements.

\subsection{8.5 The Geometry Near a
Boundary}\label{the-geometry-near-a-boundary}

\subsubsection{Normal Bundles and Transverse
Structure}\label{normal-bundles-and-transverse-structure}

At a stratum boundary, we can decompose the local geometry into
\emph{tangential} and \emph{transverse} components. The tangential
directions lie along the boundary; the transverse directions cross it.

Let \(B\mathcal{\subset M}\) be a codimension-1 stratum boundary, and
let \(p \in B\). The tangent space \(T_{p}\mathcal{M}\) decomposes as:

\[T_{p}\mathcal{M =}T_{p}B \oplus N_{p}B\]

where \(T_{p}B\) is the tangent space to the boundary and \(N_{p}B\) is
the \emph{normal bundle} --- the one-dimensional subspace perpendicular
to the boundary.

\textbf{Moral interpretation.} The tangential directions are the
directions of variation that do \emph{not} cross the boundary ---
changes that keep you in the same moral regime. The normal direction is
the direction that \emph{crosses} the boundary --- the direction of
regime change.

\textbf{Example.} At the consent threshold, the tangential directions
include variations in the \emph{degree} of information provided (more or
less detail, different formats) that keep consent above the threshold.
The normal direction is the direction that moves consent below the
threshold --- from sufficient to insufficient.

\subsubsection{The Jump Discontinuity}\label{the-jump-discontinuity}

At a Type I boundary (threshold), the satisfaction function has a
well-defined \emph{jump}:

\[\Delta S(p) = \lim_{\varepsilon \rightarrow 0^{+}}S(p + \varepsilon\widehat{n}) - \lim_{\varepsilon \rightarrow 0^{+}}S(p - \varepsilon\widehat{n})\]

where \(\widehat{n}\) is the unit normal to \(B\) at \(p\). The jump
\(\Delta S\) measures the \emph{moral significance} of crossing the
boundary at the point \(p\).

\textbf{Properties of the jump:}

\begin{enumerate}
\def\labelenumi{\arabic{enumi}.}
\item
  \(\Delta S\) may vary along the boundary. The moral significance of
  crossing the consent threshold may be greater for a high-risk
  procedure than a low-risk one.
\item
  \(\Delta S\) may be finite (for ordinary thresholds) or infinite (for
  constraint surfaces, where \(S \rightarrow - \infty\) on one side).
\item
  The direction of the jump is always normal to the boundary ---
  tangential motion does not produce discontinuity.
\end{enumerate}

\subsubsection{Penumbral Zones}\label{penumbral-zones}

In practice, moral boundaries are not always sharply defined. There may
be a \emph{penumbral zone} --- a narrow region around the boundary where
the stratum membership is uncertain or contested.

\textbf{Definition 8.9 (Penumbral Zone).} \emph{A penumbral zone around
a boundary} \(B\) \emph{is a tubular neighborhood}
\(N_{\varepsilon}(B) = \{ p\mathcal{\in M:}d(p,B) < \varepsilon\}\)
\emph{within which the applicable stratum is uncertain.}

Precision note. The term ``uncertain'' in Definition 8.9 is not
epistemic (it does not refer to the evaluator's lack of knowledge) but
measurement-theoretic: within the penumbral zone, the projection π: N →
\{S\_α, S\_β\} that assigns each point to a governing stratum is not
uniquely determined by the metric structure alone. The width of the
penumbral zone is a geometric property of the boundary (its curvature
and the angle of approach), not a property of the observer.

The existence of penumbral zones does not contradict the discreteness of
the boundary. The boundary itself is sharp --- the stratum structure is
well-defined at every point. What may be uncertain is \emph{where the
boundary is located} in a particular case. The penumbral zone represents
epistemic uncertainty about the location of a sharp boundary, not
metaphysical fuzziness in the boundary itself.

\textbf{Legal analogy.} The age of majority is precisely 18 (in most
jurisdictions). There is no ``penumbral zone'' around 18 --- you are
either a minor or an adult. But whether a particular individual has
\emph{capacity} to consent (a different boundary) may be uncertain near
the threshold, not because the threshold is fuzzy, but because the
individual's capacity is difficult to measure precisely.

\subsection{8.6 Singularities: Where the Geometry Breaks
Down}\label{singularities-where-the-geometry-breaks-down}

\subsubsection{Moral Singularities
Revisited}\label{moral-singularities-revisited}

Chapter 5 (§5.7) introduced moral singularities as points where the
geometric structure fails: the metric degenerates, the gradient is
undefined, or constraint surfaces intersect in a way that eliminates all
permissible options. This section develops the theory in detail.

\subsubsection{Classification of
Singularities}\label{classification-of-singularities}

\textbf{Type 1: Metric Degeneracy.}

At a metrically degenerate point \(p\), the metric tensor
\(g_{\mu\nu}(p)\) has \(\det(g_{p}) = 0\). Some directions in moral
space have zero ``weight'' --- comparisons along those directions are
undefined.

\textbf{Moral content.} Metric degeneracy represents
\emph{incommensurability}: at the singular point, two or more moral
dimensions cannot be compared. The trade-off rate between them is
undefined --- not because we lack information, but because no legitimate
exchange rate exists.

\textbf{Example.} In Sophie's Choice (the forced choice between the
lives of two children), the metric between ``save child A'' and ``save
child B'' is degenerate. There is no exchange rate. The lives are not
comparable in the way that, say, years of waiting time and medical
benefit are comparable. The singularity is not epistemic (we don't know
the exchange rate); it is structural (no exchange rate exists).

\textbf{Formal characterization.} At a metrically degenerate point, the
metric has a non-trivial null space:

\[\text{Null}(g_{p}) = \{ v \in T_{p}\mathcal{M:}g_{\mu\nu}v^{\mu}w^{\nu} = 0\text{ for all }w\}\]

The dimension of the null space measures the \emph{degree} of
degeneracy. If \(\dim(\text{Null}(g_{p})) = k\), then \(k\) dimensions
of moral space have become incommensurable with the others.

\textbf{Type 2: Gradient Failure.}

At a gradient-singular point \(p\), the satisfaction function \(S\) is
not differentiable: \(\nabla S\) is undefined. There is no well-defined
``direction of moral improvement.''

\textbf{Moral content.} Gradient failure represents a \emph{genuine
dilemma} in the strongest sense: not a difficult choice (where
\(\nabla S\) exists but is small), but a point where the very concept of
``moral improvement'' is ill-defined.

\textbf{Example.} A physician has promised confidentiality to a patient
who reveals an intent to harm a third party. The obligation of
confidentiality and the obligation to prevent harm are both operative,
point in opposite directions, and neither can be subordinated to the
other. The satisfaction function \(S\) has a \emph{cusp} or
\emph{corner} at this point --- differentiable from neither direction.

\textbf{Formal characterization.} At a gradient-singular point, \(S\)
may have: - A \textbf{cusp}: \(S\) approaches different limiting slopes
from different directions. The subdifferential \(\partial S(p)\) is a
non-singleton set. - A \textbf{cone point}: \(S\) is Lipschitz but not
differentiable, with a cone of subgradients. - A \textbf{branch point}:
Multiple sheets of \(S\) meet at \(p\), with different gradients on each
sheet.

\textbf{Type 3: Constraint Intersection.}

At a constraint intersection, multiple constraint surfaces meet, and
their intersection may leave no permissible options.

\textbf{Moral content.} This is the ``tragic dilemma'' --- a situation
where every available option violates at least one absolute constraint.
No permissible action exists.

\textbf{Formal characterization.} Let \(\{ C_{i}\}_{i = 1}^{k}\) be
constraint sets with \(\bigcap_{i}C_{i} \supset \{ p\}\) but with the
\emph{feasible cone} at \(p\) having empty interior:

\[\text{int}\left( \bigcap_{i}T_{p}\mathcal{(M \smallsetminus}C_{i}) \right) = \varnothing\]

Every direction from \(p\) leads into some constraint set. There is no
permissible direction.

\textbf{Example.} An autonomous vehicle faces an unavoidable collision.
Path A leads to a pedestrian; path B leads to a wall (harming the
passenger); path C leads to a cyclist. Every option violates a
constraint (do not harm innocents). The feasible cone is empty.

\subsubsection{The Moral Content of
Singularities}\label{the-moral-content-of-singularities}

Singularities are not pathologies. They are \emph{information}. A
singularity tells us:

\begin{enumerate}
\def\labelenumi{\arabic{enumi}.}
\item
  \textbf{That} a genuine dilemma exists (not merely a difficult choice)
\item
  \textbf{Why} it is a dilemma (which constraints or values conflict)
\item
  \textbf{What type} of dilemma it is (metric degeneracy, gradient
  failure, or constraint intersection)
\item
  \textbf{What residue} any resolution will leave (Chapter 15)
\end{enumerate}

Different types of singularity call for different responses:

{\def\LTcaptype{none} % do not increment counter
\begin{longtable}[]{@{}
  >{\raggedright\arraybackslash}p{(\linewidth - 4\tabcolsep) * \real{0.2110}}
  >{\raggedright\arraybackslash}p{(\linewidth - 4\tabcolsep) * \real{0.3344}}
  >{\raggedright\arraybackslash}p{(\linewidth - 4\tabcolsep) * \real{0.4545}}@{}}
\toprule\noalign{}
\begin{minipage}[b]{\linewidth}\raggedright
\textbf{Singularity Type}
\end{minipage} & \begin{minipage}[b]{\linewidth}\raggedright
\textbf{Diagnosis}
\end{minipage} & \begin{minipage}[b]{\linewidth}\raggedright
\textbf{Appropriate Response}
\end{minipage} \\
\midrule\noalign{}
\endhead
\bottomrule\noalign{}
\endlastfoot
Metric degeneracy & Values are incommensurable & Refuse to rank;
acknowledge the loss inherent in any choice \\
Gradient failure & Obligations conflict without resolution & Exercise
judgment; accept moral residue \\
Constraint intersection & No permissible option exists & Choose the
least impermissible; reform the constraints \\
\end{longtable}
}

A framework that cannot represent singularities --- that always delivers
an answer --- is not more powerful than one that can. It is less honest.
It conceals the structural source of genuine moral difficulty behind a
false determinacy.

\subsection{8.7 Phase Transitions in Moral
Reasoning}\label{phase-transitions-in-moral-reasoning}

\subsubsection{The Physics Analogy}\label{the-physics-analogy}

In statistical physics, a \emph{phase transition} occurs when a
continuous change in a control parameter (temperature, pressure)
produces a discontinuous change in the system's macroscopic properties
(density, magnetization, symmetry). Water at 99°C is liquid; at 101°C
(at standard pressure), it is steam. The molecules are the same; the
collective behavior is qualitatively different.

Moral phase transitions have the same structure. A continuous change in
a moral parameter --- the severity of harm, the strength of a promise,
the urgency of need --- can produce a discontinuous change in the
applicable moral framework.

\subsubsection{Formal Development}\label{formal-development}

Let \(\lambda\) be a continuous control parameter (a coordinate on
\(\mathcal{M}\)) and let \(S(\lambda)\) be the satisfaction function
along a path parameterized by \(\lambda\).

\textbf{Definition 8.10 (Moral Phase Transition).} \emph{A moral phase
transition occurs at} \(\lambda = \lambda^{*}\) \emph{if:}

\begin{enumerate}
\def\labelenumi{\arabic{enumi}.}
\item
  \emph{(First-order)} \(S\) \emph{is discontinuous at}
  \(\lambda^{*}\)\emph{:}
  \(\lim_{\lambda \rightarrow \lambda^{* -}}S(\lambda) \neq \lim_{\lambda \rightarrow \lambda^{* +}}S(\lambda)\)\emph{.}
\item
  \emph{(Second-order)} \(S\) \emph{is continuous but} \(\nabla S\)
  \emph{is discontinuous at} \(\lambda^{*}\)\emph{.}
\item
  \emph{(Higher-order)} \(S\) \emph{and} \(\nabla S\) \emph{are
  continuous, but the metric} \(g_{\mu\nu}\) \emph{or its derivatives
  are discontinuous at} \(\lambda^{*}\)\emph{.}
\end{enumerate}

\subsubsection{Examples}\label{examples}

\textbf{First-order: The consent boundary.} Below the consent threshold,
medical treatment is battery (\(S = - \infty\)). Above it, treatment may
have positive expected value. The jump in \(S\) is infinite --- a
first-order transition of the strongest kind.

\textbf{Second-order: The promise strengthener.} Consider a spectrum of
commitment from casual mention (``I might come'') through conditional
promise (``I'll come if I can'') to unconditional promise (``I promise
I'll be there''). The satisfaction function \(S\) may vary continuously
--- each step increases the obligation's weight smoothly. But at the
point of explicit promise, the \emph{gradient} changes sharply:
\(\partial S/\partial x^{2}\) (sensitivity to the duty dimension) jumps,
because the obligation crystallizes from a matter of courtesy to a
matter of fidelity. The Dear Abby data confirms this: ``you promised''
is a step-function gate, not a smooth ramp.

\textbf{Higher-order: Contextual metric shift.} Crossing from a family
context to a workplace context may leave \(S\) and \(\nabla S\)
continuous (the same action receives approximately the same evaluation),
but change the metric: the off-diagonal components \(g_{27}\) (coupling
between rights and care) differ between family and workplace. This is a
higher-order phase transition --- the trade-off structure changes, even
though the point evaluation does not.

\subsubsection{Critical Phenomena}\label{critical-phenomena}

Near a phase transition, the system exhibits \emph{critical phenomena}
--- enhanced sensitivity to small perturbations. In physics, this is
critical opalescence, divergent correlation lengths, and power-law
fluctuations. In moral reasoning, the analogues are:

\begin{itemize}
\item
  \textbf{Enhanced moral anxiety.} Near a moral phase transition (close
  to the consent threshold, near the breaking point of a promise),
  agents experience heightened moral concern --- a sensitivity to small
  details that would be irrelevant far from the boundary.
\item
  \textbf{Amplified disagreement.} Near a boundary, small differences in
  moral perception can produce large differences in moral conclusion.
  Two agents who agree about situations far from the boundary may
  disagree sharply about situations near it.
\item
  \textbf{Slow moral deliberation.} Decision-making slows near phase
  transitions, as the agent ``oscillates'' between the two regimes. This
  is the moral analogue of \emph{critical slowing down} in physics.
\end{itemize}

These phenomena are empirically observable and provide indirect evidence
for the stratified structure of moral space.

Phase Transitions as Symmetry Restoration

In physics, a disordered high-temperature phase has more symmetry than
the ordered low-temperature phase: the random orientations of a
paramagnet are rotationally symmetric, while the aligned spins of a
ferromagnet break that symmetry. The same principle applies to moral
phase transitions.

At high normative uncertainty (``high moral temperature''), the O/L/C/N
distinction loses meaning---not because the states blur, but because the
symmetry is restored. All Hohfeldian states become equivalent, and the
D₄ action becomes trivial. This suggests three regimes:

\textbf{1. Fully Ordered.} All semantic gates function reliably. Normal
moral reasoning operates within well-defined Hohfeldian states.

\textbf{2. Partially Ordered.} Low-dimensional stratum boundaries (near
decision points) have ``melted'' while high-dimensional boundaries (in
deliberation space) remain intact. This is the regime of moral
triage---abstract principles remain articulable, but concrete
applications become ambiguous.

\textbf{3. Fully Disordered.} All boundaries transparent. Complete
normative chaos. The stratified structure predicts that this dissolution
propagates from decision points outward, not uniformly.

\subsection{8.8 Stratification and Moral
Dynamics}\label{stratification-and-moral-dynamics}

\subsubsection{Paths Crossing
Boundaries}\label{paths-crossing-boundaries}

Chapter 10 will develop the full theory of moral dynamics --- parallel
transport, covariant derivatives, curvature. Here we note the
complications that stratification introduces.

A path \(\gamma:\lbrack 0,1\mathcal{\rbrack \rightarrow M}\) may cross
stratum boundaries. At each crossing, the applicable rules change, and
the tensors being transported must be \emph{translated} from one
stratum's framework to the other's.

\textbf{Definition 8.11 (Boundary Crossing Data).} \emph{At a stratum
boundary} \(B\) \emph{between strata} \(S_{\alpha}\) \emph{and}
\(S_{\beta}\)\emph{, the boundary crossing data consists of:} 1. \emph{A
transition map}
\(T_{\alpha\beta}\mathcal{:T(}S_{\alpha})|_{B}\mathcal{\rightarrow T(}S_{\beta})|_{B}\)
\emph{specifying how tensors on} \(S_{\alpha}\) \emph{relate to tensors
on} \(S_{\beta}\) \emph{at the boundary.} 2. \emph{A jump function}
\(\Delta S:B\mathbb{\rightarrow R \cup \{ - \infty\}}\) \emph{specifying
the discontinuity in the satisfaction function.} 3. \emph{A gate
condition} \(G:B \rightarrow \{ 0,1\}\) \emph{specifying whether the
transition fires at each point of the boundary.}

The transition map \(T_{\alpha\beta}\) is the moral analogue of the
\emph{connection} in gauge theory --- it specifies how to compare moral
objects across a change of regime. In general, \(T_{\alpha\beta}\) is
not the identity: an obligation on one side of the boundary need not
correspond to the ``same'' obligation on the other side.

\textbf{Example.} An obligation to keep a secret (in the confidentiality
stratum) becomes, upon learning the secret involves harm to a third
party, an obligation to disclose (in the duty-to-warn stratum). The
transition map \(T\) reverses the direction of the obligation vector
along the disclosure dimension.

Convention 8.1 (Differential-Geometric Operations at Stratum
Boundaries). Throughout this work, differential-geometric operations
(geodesics, parallel transport, curvature, covariant derivatives) are
defined within the interior of each stratum, where the metric and
connection are smooth. At stratum boundaries, the following conventions
apply:

\begin{enumerate}
\def\labelenumi{(\roman{enumi})}
\item
  Piecewise geodesics. A path crossing from stratum S\_α to S\_β follows
  the geodesic equation in S\_α up to the boundary, then follows the
  geodesic equation in S\_β with initial conditions given by the
  boundary crossing data (Definition 8.11). The transition map T\_\{αβ\}
  translates terminal velocity in S\_α into initial velocity in S\_β.
  The resulting path is continuous and piecewise smooth.
\item
  Parallel transport across boundaries. A tensor transported to the
  boundary in S\_α's connection is mapped by T\_\{αβ\} at the crossing
  point and then transported onward in S\_β's connection. The holonomy
  of a closed path (Chapter 10) thus receives contributions from both
  the smooth curvature within each stratum and the transition maps at
  boundary crossings --- analogous to the Aharonov--Bohm effect, where
  curvature vanishes along the path but holonomy is nonzero due to the
  topological structure of the connection.
\item
  Curvature and variational quantities. The Riemann curvature tensor
  (Definition 4.5), the Lagrangian L (Definition 12.1), and all
  variational quantities are C² within each stratum. At absorbing
  boundaries (Definition 8.6), V → ∞ ensures that extremal paths cannot
  cross; at non-absorbing boundaries, junction conditions (matching
  position and momentum across the boundary) determine the unique
  piecewise-smooth continuation. These conventions ensure that the
  integral quantities of later chapters --- the moral action (§12.4),
  the parallel transport equation (§10.4), and the holonomy computation
  (§10.5) --- are well-defined for paths that cross stratum boundaries.
\end{enumerate}

\subsubsection{The Holonomy of Moral
Circuits}\label{the-holonomy-of-moral-circuits}

If a path \(\gamma\) crosses boundaries and returns to its starting
stratum, the tensor it carries may have changed. This is \emph{holonomy}
--- the failure of a tensor to return to its original value after
parallel transport around a closed loop.

In a stratified space, holonomy can arise from two sources:

\begin{enumerate}
\def\labelenumi{\arabic{enumi}.}
\item
  \textbf{Smooth curvature} within a single stratum (Chapter 10).
\item
  \textbf{Boundary transition maps} that do not compose to the identity
  around a loop.
\end{enumerate}

The second source is distinctive to stratified spaces. If a path crosses
from stratum \(\alpha\) to \(\beta\) and back to \(\alpha\), the total
transition is \(T_{\beta\alpha} \circ T_{\alpha\beta}\), which may not
equal the identity. This means that traversing a moral circuit ---
encountering a complication and then resolving it --- may leave moral
obligations permanently changed, even if the ``situation'' appears to
have returned to its original state.

This is the geometric content of the intuition that \emph{moral
experience leaves marks}. A relationship that survives a betrayal is not
the same as one that never faced one. The obligations have been
transported through moral space and returned, but the holonomy of the
circuit has transformed them.

\subsection{8.9 Worked Example: The Stratified Landscape of a
Promise}\label{worked-example-the-stratified-landscape-of-a-promise}

We trace a complete example through the stratified structure.

\subsubsection{Setup}\label{setup-1}

Morgan promises to help Alex move on Saturday. The relevant portion of
the moral manifold has three active dimensions:

\begin{itemize}
\item
  \(x^{2}\): Strength of Morgan's duty (Rights/Duties)
\item
  \(x^{4}\): Cost to Morgan's autonomy (Autonomy)
\item
  \(x^{7}\): Quality of the Morgan-Alex relationship (Care)
\end{itemize}

\subsubsection{The Strata}\label{the-strata}

\textbf{Stratum} \(S_{O}\) \textbf{(Obligation regime).} Morgan has a
binding obligation. The satisfaction function is:

\[S_{O}(x^{2},x^{4},x^{7}) = g_{22}x^{2} + g_{44}x^{4} + g_{77}x^{7} + g_{27}x^{2}x^{7}\]

where \(g_{22} > 0\) (duty has positive weight), \(g_{44} < 0\) (cost to
autonomy has negative weight), \(g_{77} > 0\) (good relationship
amplifies the obligation), and \(g_{27} > 0\) (there is positive
coupling between duty and care --- the obligation is stronger when the
relationship is close).

\textbf{Stratum} \(S_{L}\) \textbf{(Liberty regime).} Morgan is free to
choose. The satisfaction function is:

\[S_{L}(x^{2},x^{4},x^{7}) = g_{44}'x^{4} + g_{77}'x^{7}\]

The duty component \(x^{2}\) no longer appears --- there is no
obligation to fulfill. Autonomy is now positively weighted
(\(g_{44}' > 0\) --- freedom to choose is a good), and the relationship
dimension remains relevant but with different weight.

\textbf{Stratum} \(S_{N}\) \textbf{(Nullified regime).} Obligations have
been annulled by a nullifier. The satisfaction function is:

\[S_{N} = S_{N}(x^{7}) = - |x^{7}|\]

Only the care dimension remains, and it has negative sign --- the
relationship has been damaged. The other dimensions are irrelevant.

\subsubsection{The Boundaries}\label{the-boundaries}

\(B_{OL}\)\textbf{: The Obligation--Liberty boundary.} This is crossed
by semantic gates:

\begin{itemize}
\item
  ``Only if convenient'' → \(S_{O} \rightarrow S_{L}\) (gate
  effectiveness: 89\%)
\item
  ``You explicitly promised'' → \(S_{L} \rightarrow S_{O}\) (gate
  effectiveness: 94\%)
\end{itemize}

\(B_{ON}\) \textbf{(and} \(B_{LN}\)\textbf{): The Nullifier boundary.}
This is crossed by:

\begin{itemize}
\item
  Abuse by Alex → Any stratum \(\rightarrow S_{N}\) (effectiveness:
  97\%). This is absorbing: there is no gate back.
\end{itemize}

\subsubsection{A Path Through the
Landscape}\label{a-path-through-the-landscape}

Consider the following sequence of events:

\begin{enumerate}
\def\labelenumi{\arabic{enumi}.}
\item
  \(t_{0}\)\textbf{: Morgan promises to help.} The state enters
  \(S_{O}\). Morgan has a duty (\(x^{2} = 0.8\)), the cost is manageable
  (\(x^{4} = 0.3\)), and the relationship is good (\(x^{7} = 0.7\)).
  Satisfaction:
  \(S_{O} = g_{22}(0.8) + g_{44}(0.3) + g_{77}(0.7) + g_{27}(0.8)(0.7)\).
\item
  \(t_{1}\)\textbf{: Alex says ``only if it's not too much trouble.''}
  This fires the O → L semantic gate. The state crosses \(B_{OL}\) into
  \(S_{L}\). The duty component drops out. Morgan is now free to choose.
\item
  \(t_{2}\)\textbf{: Morgan reconsiders --- the promise was explicit.}
  The L → O gate fires. The state crosses \(B_{OL}\) back into
  \(S_{O}\). But the path through \(S_{L}\) has left a mark: the duty is
  restored, but the autonomy cost has increased (\(x^{4} = 0.5\))
  because Morgan has already entertained the possibility of not helping.
\item
  \(t_{3}\)\textbf{: Morgan discovers that Alex has been spreading lies
  about them.} This is not abuse (no physical threat), but it damages
  the relationship: \(x^{7}\) drops to 0.2. The state remains in
  \(S_{O}\) (the obligation persists), but \(S_{O}\) is now much lower
  --- the coupling term \(g_{27}x^{2}x^{7}\) has shrunk.
\item
  \(t_{4}\)\textbf{: Alex threatens Morgan.} This activates the
  nullifier. The state crosses into \(S_{N}\), the absorbing stratum.
  All obligations are annulled. The only operative evaluation is the
  (negative) care dimension: the relationship is damaged, and Morgan has
  no obligation to help.
\end{enumerate}

\subsubsection{Analysis}\label{analysis}

The path
\(t_{0} \rightarrow t_{1} \rightarrow t_{2} \rightarrow t_{3} \rightarrow t_{4}\)
crosses three stratum boundaries. Two crossings are reversible
(\(S_{O} \leftrightarrow S_{L}\)); one is absorbing
(\(\rightarrow S_{N}\)). The final state depends on the \emph{path}, not
just the endpoint: if the threat at \(t_{4}\) had come before the
reconciliation at \(t_{2}\), the state would have entered \(S_{N}\)
directly from \(S_{L}\), with a different moral residue (Morgan might
feel less conflicted about abandoning a promise they had already been
released from).

This path-dependence through stratified space is \emph{holonomy} --- the
geometric content of the moral intuition that the sequence of events
matters, not just the final configuration.

\textbf{Computational instantiation (February 2026).} The DEME V3
reference implementation includes a ``Fireman's Dilemma'' smart home
demonstration that exercises precisely this stratified structure in
running code. A smart home AI must navigate three strata: the regular
operation stratum (privacy protections active), an emergency boundary
stratum (fire detected, safety overrides privacy), and a suspicion
boundary stratum (suspected false alarm, intermediate evaluation). The
V3 implementation uses uncertainty-modulated overrides: the system's
willingness to cross stratum boundaries depends on epistemic confidence
in the emergency status, with the override threshold set by
governance-specified safety margins. The demonstration confirms that
stratified evaluation is computationally tractable and produces
intuitive decisions at boundary crossings.

\subsection{8.10 Stratification and Moral
Theory}\label{stratification-and-moral-theory}

\subsubsection{How Different Theories Handle
Boundaries}\label{how-different-theories-handle-boundaries}

Different ethical theories treat stratum boundaries differently, and the
geometric framework makes these differences precise.

\textbf{Consequentialism} tends to minimize boundaries. If consequences
are all that matter, there should be few hard discontinuities --- every
situation admits a trade-off based on expected outcomes.
Consequentialism corresponds to a moral manifold with \emph{few strata}
and \emph{smooth boundaries} (or no boundaries at all --- a purely
smooth manifold).

\textbf{Deontology} emphasizes boundaries. Absolute duties create
constraint surfaces; categorical distinctions create stratum boundaries;
the priority of rights over welfare creates phase transitions.
Deontological ethics corresponds to a manifold with \emph{many strata},
\emph{sharp boundaries}, and \emph{absorbing strata} (inviolable rights
as nullifiers).

\textbf{Virtue ethics} focuses on the agent's character as they navigate
boundaries. The virtuous agent crosses boundaries with appropriate
sensitivity --- recognizing the phase transition, adjusting their
deliberation, and honoring the moral residue. Virtue corresponds to a
\emph{section of the agent bundle} (Chapter 5) that is well-behaved near
boundaries.

\textbf{Care ethics} emphasizes the relational dimensions near
boundaries. When a relationship crosses a boundary (from trust to
betrayal, from independence to dependence), care ethics attends to the
relational impact --- the change in the care dimension \(x^{7}\). Care
corresponds to a \emph{privileged subspace} of the tangent space near
boundaries.

The geometric framework does not adjudicate between these theories. It
reveals that they disagree, in part, about the \emph{stratification} of
moral space: how many strata there are, where the boundaries fall, which
boundaries are absorbing, and how to behave near them.

\subsubsection{Metaethical Implications}\label{metaethical-implications}

The existence of stratum boundaries has metaethical significance.

If moral space is \emph{purely smooth} (no strata, no boundaries), then
ethical theories that posit hard distinctions (absolute prohibitions,
categorical duties, inviolable rights) are wrong about the structure of
moral space. All moral differences are differences of degree, and all
trade-offs are in principle available.

If moral space is \emph{genuinely stratified} (with hard boundaries that
resist smoothing), then ethical theories that reduce everything to
trade-offs (act consequentialism, certain forms of cost-benefit
analysis) are wrong about the structure of moral space. Some moral
differences are differences of kind, and some transitions are
discontinuous.

The empirical evidence (§8.4) supports genuine stratification: semantic
gates behave as step functions, not smooth transitions. The correlative
symmetry of Hohfeldian states is a discrete structural feature, not an
approximation to a continuous spectrum. The \(D_{4}\) group structure is
exact, not approximate.

This is evidence --- not proof --- that moral space is stratified, and
that ethical theories must account for genuine discontinuities in the
moral landscape.

\subsection{8.10a The Discrete Computational Realization: Weighted
Simplicial
Complexes}\label{a-the-discrete-computational-realization-weighted-simplicial-complexes}

The Whitney-stratified manifold of the preceding sections is the
continuous mathematical structure. For computational applications ---
decision support, policy analysis, AI alignment, and domain-specific
instantiations (economics, medicine, law) --- the continuous manifold
admits a natural discretization as a weighted simplicial complex. This
section formalizes the discrete structure and establishes the
correspondence between continuous and discrete architectures.

\textbf{Definition 8.12 (The Moral Decision Complex). {[}Modeling
Axiom.{]}} The moral decision complex Δ is a weighted simplicial complex
constructed from the moral manifold ℳ as follows. Vertices
(0-simplices): Each vertex v\_i represents a morally relevant state ---
a configuration of the situation's nine-dimensional attribute vector
a(v\_i) ∈ ℝ⁹. Edges (1-simplices): An edge (v\_i, v\_j) represents an
available action --- a morally relevant transition from state v\_i to
state v\_j. The edge carries a weight w(v\_i, v\_j) ≥ 0 representing the
total moral cost of the action on the full manifold. Higher simplices: A
k-simplex {[}v\_0, \ldots, v\_k{]} represents a bundle of jointly
executed actions (a policy, a care protocol, a legislative package)
whose combined cost may differ from the sum of individual edge costs due
to interaction effects.

\textbf{Definition 8.13 (Mahalanobis Edge Weights with Boundary
Penalties). {[}Formal Definition.{]}} The weight of an action (v\_i,
v\_j) on the moral decision complex Δ is:

w(v\_i, v\_j) = Δaᵀ Σ⁻¹ Δa + Σ\_k β\_k · 1{[}boundary k crossed{]}

where Δa = a(v\_j) − a(v\_i) is the change in the situation's attribute
vector, Σ is the 9 × 9 moral covariance matrix encoding
cross-dimensional dependencies, and β\_k is the penalty for crossing
moral boundary k. The first term is the Mahalanobis distance on the
attribute space --- it measures the `cost' of the attribute change
weighted by the inverse covariance, so that changes along highly
correlated dimensions (e.g., trust and autonomy) are cheaper than
changes along independent or anti-correlated dimensions. The second term
is the sum of boundary penalties --- the discrete analogue of the
phase-transition costs formalized in §8.7.

The covariance matrix Σ is the discrete computational counterpart of the
metric tensor g\_\{μν\}. The correspondence is precise: the metric
tensor defines the infinitesimal cost of displacement ds² = g\_\{μν\}
dxᵐ dxⁿ on the continuous manifold; the Mahalanobis distance Δaᵀ Σ⁻¹ Δa
defines the finite cost of transition on the discrete complex. In the
limit of infinitesimal transitions (Δa → da), the Mahalanobis distance
converges to the Riemannian arc length, and the discrete geodesic
(minimum-weight path on Δ) converges to the continuous geodesic
(minimum-length curve on ℳ).

The boundary penalties β\_k are the discrete counterparts of the
phase-transition penalties of §8.7. Each β\_k corresponds to a stratum
boundary on the continuous manifold: crossing from one stratum to
another incurs a finite cost that is not captured by the smooth metric
within either stratum. The classification of boundaries from §8.5 maps
directly:

Type I boundaries (thresholds with finite penalty) correspond to β\_k
\textless{} ∞: crossing is costly but possible, and the A* search may
include the crossing if the total path cost is lower than the
alternative. Type II boundaries (phase transitions with regime change)
correspond to β\_k that are large but finite, reflecting the cost of
transitioning between qualitatively different moral regimes. Type III
boundaries (absorbing strata / nullifiers) correspond to β\_k = ∞: no
path through this boundary can have finite cost, so the A* search
provably avoids it. The sacred-value boundaries discussed in §8.6 are
exactly those with β = ∞.

\textbf{Proposition 8.5 (Continuous-Discrete Correspondence).
{[}Proved.{]}} Let ℳ be a Whitney-stratified moral manifold with metric
tensor g\_\{μν\} and phase-transition penalties, and let Δ\_ε be the
ε-discretization of ℳ (the weighted simplicial complex with Mahalanobis
edge weights derived from g\_\{μν\} and boundary penalties derived from
the stratum transition costs). Then: (i) The minimum-weight path on Δ\_ε
converges to the geodesic on ℳ as ε → 0. (ii) The boundary penalties on
Δ\_ε equal the stratum transition costs on ℳ exactly (no convergence
needed --- boundaries are discrete structures on both sides). (iii) The
A* search on Δ\_ε with an admissible heuristic finds the optimal path on
Δ\_ε, which approximates the continuous geodesic to within O(ε) in path
cost.

\textbf{Proof}. (i) Standard result in computational geometry:
Mahalanobis distances on an ε-grid of a Riemannian manifold converge to
geodesic distances as ε → 0 (see e.g. Memoli and Sapiro, 2005). (ii)
Boundary penalties are discrete by construction on both the continuous
manifold (phase transitions are jumps, not smooth transitions --- §8.7,
Definition 8.5) and the discrete complex (the β\_k are finite
constants). (iii) Follows from (i)--(ii) combined with Proposition 11.4
(A* optimality on stratified spaces). □

Remark (Domain Instantiations). Each domain application of the Geometric
Ethics framework instantiates the moral decision complex Δ with
domain-specific attribute dimensions, covariance structure, and boundary
penalties. In economics, the decision complex becomes the economic
decision complex with dimensions including monetary cost, moral
constraint, social norm, fairness, and identity; the minimum-cost path
is the Bond Geodesic, which replaces Nash equilibrium as the solution
concept. In clinical medicine, the decision complex becomes the clinical
decision complex with dimensions including clinical outcomes, patient
autonomy, trust, dignity, and justice; the minimum-cost path is the
clinical geodesic, which replaces QALY optimization. In each case, the
continuous manifold ℳ provides the theoretical foundation, and the
discrete complex Δ provides the computational realization. The Scalar
Irrecoverability Theorem (§15.6, Theorem 15.1) then proves that reducing
either architecture to a scalar --- utility, QALY, or any other
one-dimensional projection --- irrecoverably destroys the moral
information that the multi-dimensional structure was designed to
preserve.

\subsection{8.11 Summary}\label{summary-1}

Stratification is the most distinctive feature of the moral manifold ---
the feature that distinguishes it most sharply from the smooth manifolds
of classical physics.

{\def\LTcaptype{none} % do not increment counter
\begin{longtable}[]{@{}
  >{\raggedright\arraybackslash}p{(\linewidth - 4\tabcolsep) * \real{0.1798}}
  >{\raggedright\arraybackslash}p{(\linewidth - 4\tabcolsep) * \real{0.3723}}
  >{\raggedright\arraybackslash}p{(\linewidth - 4\tabcolsep) * \real{0.4479}}@{}}
\toprule\noalign{}
\begin{minipage}[b]{\linewidth}\raggedright
\textbf{Concept}
\end{minipage} & \begin{minipage}[b]{\linewidth}\raggedright
\textbf{Mathematical Object}
\end{minipage} & \begin{minipage}[b]{\linewidth}\raggedright
\textbf{Moral Content}
\end{minipage} \\
\midrule\noalign{}
\endhead
\bottomrule\noalign{}
\endlastfoot
Stratum & Smooth manifold \(S_{\alpha}\) & Moral regime with uniform
rules \\
Threshold & Codimension-1 boundary \(B\) & Parameter crosses critical
value; status changes \\
Phase transition & Discontinuity in \(g_{\mu\nu}\) & Metric/framework
changes; trade-off structure shifts \\
Absorbing stratum & Stratum with retraction property & Nullifier (abuse,
danger) overrides all other considerations \\
Constraint surface & \(\partial C\) with \(S|_{C} = - \infty\) &
Absolute prohibition; moral non-starter \\
Semantic gate & Transition function \(G\) on \(B\) & Trigger
phrase/feature that fires a state transition \\
\(D_{4}\) group & Dihedral group of order 8 & Algebraic structure of
Hohfeldian state transitions \\
Singularity & Point where \(g\), \(\nabla S\), or feasible cone
degenerates & Genuine moral dilemma \\
Penumbral zone & Neighborhood \(N_{\varepsilon}(B)\) & Epistemic
uncertainty about boundary location \\
Holonomy & T\_\{βα\} ∘ T\_\{αβ\} ≠ id & Moral experience leaves
permanent marks \\
\end{longtable}
}

The mathematics of stratification is the mathematics of
\emph{discontinuity within structure} --- sharp boundaries that are
nonetheless geometrically well-behaved, governed by Whitney's regularity
conditions. Moral space is not a smooth continuum, and it is not a
random jumble. It is a \emph{stratified space}: structured, piecewise
smooth, with transitions governed by algebraic and geometric laws.

The next chapter asks where the metric comes from --- the metric that
determines, within each stratum, how values are compared and trade-offs
are structured. But the existence of strata is prior to the metric: it
is the metric that varies across strata, and the boundaries between
strata that define the landscape within which the metric operates.

\subsection{Technical Appendix}\label{technical-appendix-1}

\textbf{Proposition 8.1 (Existence of Moral Whitney Stratification).}
\emph{Let} \(\mathcal{M}\) \emph{be a moral manifold equipped with a
finite set of constraint surfaces} \(\{ C_{i}\}\) \emph{and semantic
gates} \(\{ G_{j}\}\)\emph{. Then} \(\mathcal{M}\) \emph{admits a
Whitney stratification compatible with all} \(C_{i}\) \emph{and}
\(G_{j}\)\emph{.}

Hypothesis note. The proof assumes that each constraint surface C\_i is
a closed semi-algebraic subset of M --- that is, definable by finitely
many polynomial inequalities in local coordinates. This semi-algebraic
condition ensures that the Whitney stratification exists (by the
Łojasiewicz--Whitney theorem) and is locally finite. The hypothesis is
satisfied for all constraint surfaces arising from the
finite-dimensional moral manifold construction of Chapter 5.

\emph{\textbf{Proof}. We verify the three conditions for a Whitney
stratification. Step 1 (Stratification): Let Σ = ∪\_i ∂C\_i ∪ ∪\_j
supp(G\_j) be the union of all constraint boundaries and gate
hypersurfaces. Each C\_i is a closed semi-algebraic subset of M (defined
by inequalities on the moral coordinates) and each G\_j is a smooth
hypersurface with closed support, so Σ is closed. The connected
components of M \textbackslash{} Σ are the open strata \{S\_α\}, and Σ
itself decomposes into smooth strata by induction on dimension. Step 2
(Whitney condition A): Let S\_α ⊂ cl(S\_β) with dim S\_α \textless{} dim
S\_β, and let p\_n ∈ S\_β with p\_n → p ∈ S\_α such that T\_\{p\_n\}S\_β
→ τ. Since the constraint surfaces are smooth submanifolds, the tangent
spaces to the open stratum S\_β converge to a limit that contains T\_p
S\_α. Hence T\_p S\_α ⊆ τ. Step 3 (Whitney condition B): Since the
strata are defined by smooth semi-algebraic functions, condition (B)
follows from Whitney's theorem (1965): any semi-algebraic set admits a
Whitney stratification.
□}\(C_{i}\text{supp}(G_{j}\mathcal{)M}\bigcup_{i}\partial C_{i} \cup \bigcup_{j}\text{supp}(G_{j})▫\)

\textbf{Proposition 8.2 (Absorbing Strata Are Unique Attractors).}
\emph{Let} \(S_{0}\) \emph{be an absorbing stratum and let}
\(\gamma:\lbrack 0,1\mathcal{\rbrack \rightarrow M}\) \emph{be a path
with} \(\gamma(t^{*}) \in S_{0}\) \emph{for some}
\(t^{*} \in (0,1)\)\emph{. Then the moral evaluation along} \(\gamma\)
\emph{for} \(t > t^{*}\) \emph{is determined by} \(S_{0}\)\emph{,
independent of} \(\gamma\)\emph{'s behavior for} \(t < t^{*}\)\emph{.}

\emph{\textbf{Proof}.} By definition of absorbing stratum (Definition
8.6), \(S(p) = S(r(p))\) for all \(p\) in a neighborhood of \(S_{0}\),
where \(r\) is the retraction. For \(t > t^{*}\), either
\(\gamma(t) \in S_{0}\) (and the evaluation is determined by \(S_{0}\)'s
rules) or \(\gamma(t)\) is in the tubular neighborhood \(N\) of
\(S_{0}\) (and the evaluation factors through \(r\), hence is determined
by \(S_{0}\)). In either case, the prior trajectory is irrelevant. \(▫\)

Scope note. The proof establishes path-independence of the evaluation
within the tubular neighborhood N of S₀. If the path γ exits N for t
\textgreater{} t₀, the evaluation at those points is governed by the
ambient stratum\textquotesingle s rules, not by S₀\textquotesingle s
retraction. The proposition\textquotesingle s full-strength claim ---
that the evaluation is determined by S₀ for all t ≥ t₀, regardless of
the path\textquotesingle s subsequent behavior --- requires the
additional condition that S₀ is a trapping set: any path entering S₀
remains in N ∪ S₀ for all subsequent time. This trapping condition holds
for the absorbing strata identified empirically (abuse, danger,
impossibility; see Chapter 17), where the nullification is permanent and
context-independent.

\textbf{Proposition 8.3 (}\(D_{4}\) \textbf{Structure of Hohfeldian
Transitions).} \emph{The four Hohfeldian positions} \(\{ O,C,L,N\}\)
\emph{with correlative reflection} \(s\) \emph{(O ↔ C, L ↔ N) and cyclic
rotation} \(r\) \emph{(O → C → N → L → O) generate a group isomorphic to
the dihedral group} \(D_{4}\) \emph{of order 8.}

\emph{\textbf{Proof}.} Verify the presentation: \(r^{4} = e\) (four
rotations return to start), \(s^{2} = e\) (double correlative swap is
identity), and \(srs = r^{- 1}\) (the conjugation relation). These three
relations define \(D_{4}\). The group has exactly \(2 \times 4 = 8\)
elements: \(\{ e,r,r^{2},r^{3},s,sr,sr^{2},sr^{3}\}\). Each element
corresponds to a distinct operation on the Hohfeldian square. \(▫\)

Verification detail. The eight elements of D₄ acting on the Hohfeldian
states \{O, C, N, L\} are: e (identity), r (O→C→N→L→O), r² (O↔N, C↔L),
r³ (O→L→N→C→O), s (O↔C, L↔N), sr (O→N, C→L, N→O, L→C), sr² (O↔L, C↔N),
sr³ (O→L, C→O, N→C, L→N). Direct computation confirms: r⁴ = e, s² = e,
and srs = r⁻¹ = r³ (verified by applying each side to O: srs(O) = sr(C)
= s(N) = L = r³(O)). These are the defining relations of D₄.

❖

\emph{Moral space is not a smooth plain. It is a landscape of plateaus,
cliffs, and chasms --- smooth within each region, discontinuous at the
boundaries between them.}

\emph{The mathematics of stratification does not smooth the cliffs or
fill the chasms.}

\emph{It maps them --- precisely, rigorously, with the full apparatus of
Whitney's conditions, group theory, and transition functions.}

\emph{Ethics has edges. This chapter is their geometry.}

\section{Part III: Dynamics and
Symmetry}\label{part-iii-dynamics-and-symmetry}

\section{Chapter 9: The Origin of the Moral Metric --- Discovery,
Construction, or
Governance?}\label{chapter-9-the-origin-of-the-moral-metric-discovery-construction-or-governance}

\begin{quote}
\textbf{\textsc{RUNNING EXAMPLE --- Priya's Model}}

\emph{Who decided that Mrs. Voss's 4.5-hour drive counts the same as a
co-pay difference in TrialMatch's scoring? Priya traces the answer: no
one decided. The metric was learned from historical data---data
reflecting which patients historically completed trials. Those patients
tended to live near research hospitals. The learned metric assigns low
weight to access burden, creating a feedback loop: rural patients are
not matched, so they do not complete trials, so the model learns
rurality predicts non-completion, so rural patients are not matched.
This is metric capture by historical bias. The realist metric (medical
suitability) is defensible. The constructivist metric (community
agreement on fairness) was never consulted. The governance metric
(democratic determination) was replaced by a training loop.}
\end{quote}

\subsection{9.1 The Metric Is Not
Given}\label{the-metric-is-not-given-1}

Chapter 6 established what the moral metric \emph{does}: it defines
distances between moral states, angles between obligation vectors, inner
products that allow comparison, and the trade-off structure that
determines how values may be exchanged. The metric \(g_{\mu\nu}\) is the
most consequential object in the geometric framework --- more
consequential, in a sense, than the obligation vectors or interest
covectors, because it determines which comparisons are even possible.
Two analysts who agree on every obligation and every interest can
disagree about what should be done, if they employ different metrics.
The metric is where the deepest moral commitments live.

But Chapters 5 through 8 deferred a question that can no longer be
postponed: \emph{where does the metric come from?}

In physics, the metric of spacetime is discovered empirically --- it is
a feature of the universe we inhabit, measurable through the behavior of
light, clocks, and gravitational fields. Einstein's field equations
relate the metric to the distribution of matter and energy. The metric
is not a choice; it is a fact about nature.

In pure mathematics, the metric is stipulated --- we choose Euclidean,
hyperbolic, or Riemannian geometry and explore the consequences. The
choice is free; the consequences are binding.

In ethics, the situation is more complex. The metric encodes substantive
moral commitments: which values are commensurable, at what exchange
rate, and which are incommensurable. These commitments are not simply
read off from nature (the universe does not tell us the trade-off rate
between welfare and justice), nor are they arbitrary stipulations (some
metrics are morally inadmissible). The metric is simultaneously too
important to be arbitrary and too contentious to be discovered.

This chapter examines three traditional accounts of the metric's origin
and argues for a fourth.

\subsection{9.2 The Realist Account:
Discovery}\label{the-realist-account-discovery}

\subsubsection{The Claim}\label{the-claim}

The moral realist holds that there are moral facts independent of our
beliefs about them. On this view, the moral metric is \emph{discovered},
not invented. Just as physicists discovered that spacetime has
Lorentzian rather than Euclidean signature, moral inquirers discover the
true structure of moral space --- the actual trade-off rates between
values, the genuine incommensurabilities, the real distances between
moral states.

\textbf{Formal statement.} There exists a unique (or at most a small
family of) metric tensor(s) \(g_{\mu\nu}^{*}\) on the moral manifold
\(\mathcal{M}\) such that all correct moral evaluations are consistent
with \(g^{*}\). The task of moral inquiry is to approximate \(g^{*}\).

\subsubsection{Strengths}\label{strengths}

\textbf{Phenomenology of inquiry.} Moral reasoning often feels like
discovery. When we realize that a seemingly acceptable trade-off was
actually illegitimate --- that efficiency gains do not justify racial
discrimination --- we experience this as learning a fact about the
metric: the off-diagonal component
\(g_{\text{efficiency, discrimination}}\) is zero (or the discrimination
dimension is part of the constraint set, not the trade-off space).
Realism explains this phenomenology directly.

\textbf{Moral progress.} If the metric is a fact, then moral progress
consists in getting closer to the true metric. The abolition of slavery,
the recognition of women's suffrage, the criminalization of child labor
--- these can be understood as discoveries about \(g^{*}\): dimensions
that were once treated as commensurable with economic benefit (human
freedom, political participation, childhood welfare) were discovered to
be incommensurable with it or to require absolute priority.

\textbf{Error theory.} Realism explains moral error: a person or society
that trades welfare for injustice is \emph{wrong about the metric}, in
the same way that a physicist who uses a Euclidean metric near a black
hole is wrong about the geometry of spacetime.

\subsubsection{Challenges}\label{challenges}

\textbf{Epistemology.} How do we access the moral metric? In physics,
the metric manifests through measurable phenomena --- light bending,
clocks slowing. What are the moral analogues? Moral intuitions are
notoriously variable across cultures and epochs. If the metric is a
fact, we lack reliable instruments for measuring it.

\textbf{Metaphysics.} What kind of entity is a moral metric? It is an
abstract mathematical structure defined over a space of structured
situations. Where does it exist? How does it relate to the natural
world? The moral realist must either accept a Platonic ontology (the
metric exists in an abstract realm) or find a way to ground it in
natural facts --- and the grounding project has proven extraordinarily
difficult.

\textbf{Disagreement.} If the metric is discoverable, why do reflective,
informed inquirers disagree about it? Persistent, apparently
irresolvable disagreement between utilitarians and deontologists --- who
disagree about the metric's off-diagonal terms and the presence of
lexicographic priority --- is difficult to explain if the metric is a
straightforward fact awaiting discovery.

\subsubsection{Geometric Gloss}\label{geometric-gloss}

If realism is correct, the moral metric \(g_{\mu\nu}^{*}\) is a feature
of reality, and different ethical theories are competing hypotheses
about \(g^{*}\). The utilitarian hypothesis is
\(g^{*} = \text{diag}(w_{1},\ldots,w_{9})\) with specific weights. The
Rawlsian hypothesis is that \(g^{*}\) has lexicographic structure. The
incommensurability thesis is that \(g^{*}\) has zero eigenvalues. Moral
inquiry is metric estimation.

\subsection{9.3 The Constructivist Account:
Agreement}\label{the-constructivist-account-agreement}

\subsubsection{The Claim}\label{the-claim-1}

The constructivist holds that moral facts are constituted by the
outcomes of idealized rational procedures. On this view, the moral
metric is \emph{constructed} through deliberation --- perhaps behind a
veil of ignorance (Rawls), through ideal discourse (Habermas), or via
reflective equilibrium (Daniels).

\textbf{Formal statement.} The moral metric \(g_{\mu\nu}\) is the fixed
point of an idealized deliberative procedure \(\Phi\):

\[g^{*} = \Phi(g_{0}\mathcal{,I,C)}\]

where \(g_{0}\) is the initial metric (reflecting current moral
commitments), \(\mathcal{I}\) is the set of idealizing conditions (full
information, impartiality, rationality), and \(\mathcal{C}\) is the set
of structural constraints (consistency, completeness, invariance).

\subsubsection{Strengths}\label{strengths-1}

\textbf{Agency.} Constructivism connects the metric to rational agency.
The metric is not imposed from outside; it is the output of our own
rational capacities operating under ideal conditions. This grounds the
authority of the metric in something we can understand: reason itself.

\textbf{Procedure.} Constructivism provides a \emph{procedure} for
resolving metric disagreements. When two parties disagree about the
trade-off rate between welfare and rights, the constructivist says:
submit the question to the idealized deliberative procedure. The outcome
is authoritative because the procedure is fair.

\textbf{Metaphysical modesty.} Constructivism avoids mysterious
metaphysics. The metric exists as a feature of rational agreement, not
as a Platonic entity.

\subsubsection{Challenges}\label{challenges-1}

\textbf{Which procedure?} Different constructive procedures yield
different metrics. The Rawlsian procedure (maximin behind the veil)
yields a lexicographic metric with priority for the worst-off. The
Habermasian procedure (ideal discourse) yields a metric determined by
the force of the better argument --- but ``better'' is itself
metric-dependent. The Scanlonian procedure (what no one could reasonably
reject) yields yet another metric. If the metric depends on the
procedure, and the choice of procedure is itself a moral question,
constructivism faces a regress.

\textbf{Idealization.} Real deliberators never achieve ideal conditions.
How close is close enough? And how do we know when we have approximated
the ideal, if the ideal itself is what we are trying to determine?

\textbf{Determinacy.} Even under ideal conditions, deliberative
procedures may not yield a unique metric. There may be multiple fixed
points --- multiple metrics consistent with all the idealizing
conditions. If so, the metric is underdetermined by construction, and
constructivism collapses into pluralism.

\subsubsection{Geometric Gloss}\label{geometric-gloss-1}

The metric is the fixed point of a deliberative process. Different
starting points or different idealizing conditions may yield different
fixed points --- moral pluralism at the level of metric choice. The
constructivist metric is not unique but convergent: as deliberation
proceeds under better conditions, the class of admissible metrics
narrows. Whether it narrows to a point is the key open question.

\subsection{9.4 The Expressivist Account:
Projection}\label{the-expressivist-account-projection}

\subsubsection{The Claim}\label{the-claim-2}

The expressivist holds that moral claims express attitudes rather than
describe facts. On this view, the moral metric is a \emph{projection} of
our evaluative sensibilities onto the world --- a structure we impose,
not one we discover or construct.

\textbf{Formal statement.} The moral metric is agent-indexed:
\(g_{\mu\nu}(a)\) is the metric of agent \(a\), reflecting \(a\)'s
pattern of evaluative responses. ``Objectivity'' is intersubjective
convergence --- the overlap between individual metrics --- not
correspondence to a metric-independent fact.

\subsubsection{Strengths}\label{strengths-2}

\textbf{Motivational force.} Expressivism explains why moral judgments
motivate action. If the metric reflects my evaluative sensibilities,
then recognizing that situation \(p\) requires action (that
\(\nabla S \neq 0\) in my metric) is simultaneously recognizing that I
\emph{care} about the relevant dimensions. The motivation is built into
the metric.

\textbf{Diversity.} Expressivism naturally accommodates the diversity of
moral outlook across cultures, epochs, and individuals. Different agents
project different metrics --- and this is expected, not anomalous.

\textbf{Metaphysical economy.} No Platonic realm, no idealized
deliberation. Just agents with attitudes, projecting structure onto the
moral landscape.

\subsubsection{Challenges}\label{challenges-2}

\textbf{Structure, not just attitudes.} The moral metric is not merely
an attitude --- it is a geometric object with mathematical properties
(symmetry, non-degeneracy conditions, transformation behavior). Can
expressivism account for this structure? An attitude might be ``I care
more about welfare than justice.'' The metric encodes not just this
ranking but the precise trade-off rate, the coupling between dimensions,
the degeneracies, and the transformation law under change of
coordinates. It is unclear that attitudes have this kind of structure.

\textbf{The Frege-Geach problem.} Moral claims embed in complex logical
contexts: ``If torturing innocents is wrong, then getting others to
torture innocents is wrong.'' Expressivist semantics for embedded moral
claims remains difficult. The metric framework sharpens the problem: how
do we express ``if \(g_{12} = 0\) (welfare and justice are
incommensurable), then\ldots'' in expressivist terms?

\textbf{Convergence.} If the metric is merely projected, why is there as
much moral agreement as there is? The convergence of diverse moral
traditions on certain structural features --- the obligation/liberty
distinction (which transfers cross-linguistically at 100\% in BIP
experiments), the correlative symmetry of Hohfeldian states, the
universality of nullifiers --- is difficult to explain if the metric is
pure projection. Something structural appears to constrain the
projection.

\subsubsection{Geometric Gloss}\label{geometric-gloss-2}

The expressivist metric is subjective: \(g_{\mu\nu}(a)\) varies with the
agent. Intersubjective moral facts are features of the \emph{overlap}
between agent-metrics. The set of ``objective'' moral claims is the set
of claims invariant under all agents' metrics:

\[\text{Objective claims} = \bigcap_{a}\{ c:c\text{ is true under }g(a)\}\]

If this intersection is large, moral objectivism is approximately
correct. If it is small, moral relativism wins. The size of the
intersection is an empirical question.

\subsection{9.5 The Governance Account: Democratic
Determination}\label{the-governance-account-democratic-determination}

\subsubsection{The Proposal}\label{the-proposal}

I propose a fourth account suited to the geometric framework: the moral
metric is neither discovered, constructed by idealized reason, nor
projected by individual sensibility. It is \emph{governed} --- the
output of legitimate governance processes that have authority to
determine, for a community, how moral trade-offs are to be structured.

\textbf{Formal statement.} The moral metric \(g_{\mu\nu}\) for a
community \(\mathcal{C}\) is determined by a governance process
\(\Gamma_{\mathcal{C}}\):

\[g_{\mu\nu} = \Gamma_{\mathcal{C}}(\text{inputs})\]

where the inputs include empirical data about the community's moral
commitments, the structural constraints of §9.6, the outcomes of
legitimate deliberative processes, and the community's constitutional
commitments.

\subsubsection{What Governance Means}\label{what-governance-means}

The governance account occupies a specific position in logical space,
distinct from the three traditional accounts:

{\def\LTcaptype{none} % do not increment counter
\begin{longtable}[]{@{}
  >{\raggedright\arraybackslash}p{(\linewidth - 6\tabcolsep) * \real{0.1748}}
  >{\raggedright\arraybackslash}p{(\linewidth - 6\tabcolsep) * \real{0.2235}}
  >{\raggedright\arraybackslash}p{(\linewidth - 6\tabcolsep) * \real{0.3120}}
  >{\raggedright\arraybackslash}p{(\linewidth - 6\tabcolsep) * \real{0.2897}}@{}}
\toprule\noalign{}
\begin{minipage}[b]{\linewidth}\raggedright
\textbf{Account}
\end{minipage} & \begin{minipage}[b]{\linewidth}\raggedright
\textbf{Source of the Metric}
\end{minipage} & \begin{minipage}[b]{\linewidth}\raggedright
\textbf{Authority}
\end{minipage} & \begin{minipage}[b]{\linewidth}\raggedright
\textbf{Revisability}
\end{minipage} \\
\midrule\noalign{}
\endhead
\bottomrule\noalign{}
\endlastfoot
Realism & Nature / moral reality & Independent of us & Not revisable
(only our knowledge improves) \\
Constructivism & Idealized rational procedure & Rational authority &
Revisable if the procedure refines \\
Expressivism & Individual sensibility & No external authority &
Revisable as attitudes change \\
Governance & Legitimate institutional process & Democratic/institutional
authority & Revisable through legitimate process \\
\end{longtable}
}

The governance account holds that:

\begin{enumerate}
\def\labelenumi{\arabic{enumi}.}
\item
  \textbf{The metric is a governance artifact} --- like a constitution,
  a legal code, or a set of regulatory standards. It is neither a
  natural fact nor a subjective projection. It is a socially
  authoritative determination of how trade-offs are to be structured.
\item
  \textbf{Different communities may legitimately adopt different
  metrics.} This is not relativism --- some metrics are inadmissible
  (§9.6) --- but it is \emph{pluralism}: within the space of admissible
  metrics, legitimate governance processes may select differently,
  reflecting different histories, values, and circumstances.
\item
  \textbf{The metric can be revised.} If a community's moral
  understanding changes --- if formerly acceptable trade-offs come to be
  seen as illegitimate --- the metric can be updated through legitimate
  institutional processes. This is what moral reform \emph{is}: a change
  in the collectively authoritative metric.
\item
  \textbf{Disagreement about the metric is political as well as
  philosophical.} Whether welfare can be traded against justice at a
  particular rate is not merely a question for moral philosophy. It is a
  question for democratic deliberation, institutional design, and
  constitutional law.
\end{enumerate}

\subsubsection{Why Not Relativism?}\label{why-not-relativism}

The governance account might seem to collapse into relativism: if the
metric is whatever a community says it is, then anything goes. But this
objection misunderstands the account in two ways.

First, the metric is constrained (§9.6). Not every metric is admissible.
Some structural constraints --- non-degeneracy on protected dimensions,
internal consistency, invariance under admissible transformations ---
are non-negotiable. These constraints narrow the space of admissible
metrics, ruling out metrics that encode arbitrary discrimination,
internal contradiction, or dependence on morally irrelevant features.

Second, governance processes can be better or worse, more or less
legitimate. A metric adopted through transparent democratic
deliberation, informed by the relevant moral considerations, responsive
to the claims of affected parties, and consistent with constitutional
commitments, has more authority than a metric imposed by fiat or adopted
through corrupt process. The governance account does not say ``any
metric is as good as any other.'' It says: ``the authority of a metric
derives from the legitimacy of the process that produced it.''

\subsubsection{The Analogy to Legal
Systems}\label{the-analogy-to-legal-systems}

The governance account draws on an analogy with law. Legal systems
determine the structure of permissible actions: what counts as property,
what counts as a contract, what constitutes a crime. These are not
natural facts (different legal systems answer differently), nor are they
mere conventions (they carry real authority and have real consequences).
They are governance artifacts --- socially authoritative determinations
backed by legitimate institutional processes.

The moral metric is analogous. It determines the structure of
permissible trade-offs: whether welfare can be exchanged for justice, at
what rate, and under what conditions. These determinations are not
natural facts, and they are not mere preferences. They are governance
decisions with real moral authority.

The analogy extends further. Just as constitutions constrain ordinary
legislation (some laws are unconstitutional, regardless of popular
support), there are \emph{meta-constraints} on the moral metric (§9.6)
that no governance process can override. These constraints play the role
of moral constitutionalism: they define the boundaries within which
legitimate metric choice operates.

\subsubsection{Connection to the
Framework}\label{connection-to-the-framework}

The governance account integrates naturally with the geometric
framework:

\begin{itemize}
\item
  The manifold \(\mathcal{M}\) (Chapter 5) and the obligation fields
  \(O^{\mu}\) (Chapter 6) are \emph{prior to} the metric. They represent
  the structure of the moral situation and the duties it generates.
  These are facts about the situation, not governance choices.
\item
  The metric \(g_{\mu\nu}\) is \emph{determined by governance}. It
  represents the community's authoritative determination of how values
  trade off. Different legitimate governance processes may yield
  different metrics.
\item
  The contraction \(S = I_{\mu}O^{\mu}\) is the \emph{interface} between
  the factual and the normative. The obligation vector is factual; the
  interest covector reflects the governed metric; the satisfaction
  scalar is the verdict.
\item
  Stratification (Chapter 8) --- the boundaries, nullifiers, and
  constraint surfaces --- is \emph{partially factual and partially
  governed}. Some boundaries are structural invariants (the Hohfeldian
  correlative symmetry, the universality of abuse as a nullifier).
  Others are governance choices (where to set the age of majority, the
  threshold for informed consent).
\end{itemize}

\subsection{9.6 Constraints on Admissible
Metrics}\label{constraints-on-admissible-metrics}

Not every metric is admissible. Even on the governance account,
structural constraints narrow the space of legitimate choices.

\subsubsection{Symmetry}\label{symmetry}

The moral metric must be symmetric: \(g_{\mu\nu} = g_{\nu\mu}\). The
trade-off rate between welfare and justice must be the same regardless
of which direction the trade is made: if one unit of welfare can
compensate for \(k\) units of justice, then one unit of justice must
compensate for \(1/k\) units of welfare.

This is a mathematical constraint (symmetric bilinear forms are the
standard definition of a metric), but it has moral content. A metric
where trading welfare for justice costs more than trading justice for
welfare would encode an asymmetric privilege for one value over the
other --- not in magnitude (which is permissible and encoded in the
diagonal components), but in the \emph{structure of exchange}. This is
ruled out.

\subsubsection{Non-Degeneracy on Protected
Dimensions}\label{non-degeneracy-on-protected-dimensions}

The metric should not be \emph{fully} degenerate: there must be some
directions along which comparison is possible. A totally degenerate
metric (\(g_{\mu\nu} = 0\) for all \(\mu,\nu\)) would make all moral
comparisons impossible, rendering the framework vacuous.

More substantively, certain dimensions should have \emph{guaranteed
weight}: the metric should not assign zero weight to dimensions that a
legitimate governance process has identified as morally relevant. If a
community has determined that rights matter (Dimension 2 has positive
diagonal weight), a governance revision that sets \(g_{22} = 0\) ---
effectively declaring that rights don't exist --- requires a higher
level of constitutional justification.

\subsubsection{Invariance}\label{invariance}

The metric must respect the admissible transformation structure (Chapter
5, §5.5). Under Type 1 transformations (coordinate redescriptions), the
metric must transform as a (0,2)-tensor:

\[{\widetilde{g}}_{\mu\nu} = \frac{\partial x^{\alpha}}{\partial{\widetilde{x}}^{\mu}}\frac{\partial x^{\beta}}{\partial{\widetilde{x}}^{\nu}}g_{\alpha\beta}\]

This ensures that the trade-off structure does not depend on how the
situation is described. The metric is a geometric object, not a
coordinate artifact. The Bond Invariance Principle (Chapter 5, Axiom
5.1) requires this.

\subsubsection{Non-Discrimination}\label{non-discrimination}

The metric must not encode trade-offs that depend on morally irrelevant
features --- specifically, protected characteristics. If two situations
differ only in a morally irrelevant feature (the gender of the patient,
the race of the applicant), the metric must assign the same distances
and inner products.

\textbf{Formal statement.} Let \(\phi\mathcal{:M \rightarrow M}\) be a
transformation that changes only morally irrelevant features. Then:

\[g_{\mu\nu}(\phi(p)) = g_{\mu\nu}(p)\quad\text{for all }p\mathcal{\in M}\]

The metric is invariant under morally irrelevant transformations. This
is a constraint on the metric, not on the manifold: the manifold may
include morally irrelevant coordinates (they describe the situation),
but the metric must not depend on them.

\subsubsection{Consistency with Structural
Invariants}\label{consistency-with-structural-invariants}

The metric must be consistent with the structural invariants identified
in Chapter 8: the \(D_{4}\) group structure of Hohfeldian transitions,
the universality of certain nullifiers, and the cross-linguistic
transfer of deontic structure. These invariants are not governance
choices --- they are features of the moral manifold itself. The metric
must be compatible with them.

\textbf{Example.} The correlative symmetry O ↔ C requires that if the
metric assigns a positive weight to the obligation dimension for agent
\(a\), it assigns a corresponding positive weight to the claim dimension
for the correlative agent \(b\). A metric that violates correlative
symmetry is inadmissible.

\subsubsection{Summary of Constraints}\label{summary-of-constraints}

{\def\LTcaptype{none} % do not increment counter
\begin{longtable}[]{@{}
  >{\raggedright\arraybackslash}p{(\linewidth - 4\tabcolsep) * \real{0.2342}}
  >{\raggedright\arraybackslash}p{(\linewidth - 4\tabcolsep) * \real{0.3491}}
  >{\raggedright\arraybackslash}p{(\linewidth - 4\tabcolsep) * \real{0.4167}}@{}}
\toprule\noalign{}
\begin{minipage}[b]{\linewidth}\raggedright
\textbf{Constraint}
\end{minipage} & \begin{minipage}[b]{\linewidth}\raggedright
\textbf{Mathematical Form}
\end{minipage} & \begin{minipage}[b]{\linewidth}\raggedright
\textbf{Moral Content}
\end{minipage} \\
\midrule\noalign{}
\endhead
\bottomrule\noalign{}
\endlastfoot
Symmetry & g\_μν = g\_νμ & Trade-off structure is bidirectional \\
Non-total-degeneracy & \(\det(g) \neq 0\) on at least some subspace &
Some comparisons must be possible \\
Tensor transformation & \(g\) transforms as (0,2)-tensor & Trade-offs
don't depend on description \\
Non-discrimination & \(g(\phi(p)) = g(p)\) for irrelevant \(\phi\) &
Morally irrelevant features don't affect trade-offs \\
Structural compatibility & Compatible with \(D_{4}\), nullifiers, etc. &
Metric respects known invariants \\
\end{longtable}
}

The space of admissible metrics is the space satisfying all five
constraints. It is large --- allowing genuine pluralism --- but not
unlimited. The constraints rule out the most egregious metrics
(discriminatory, inconsistent, description-dependent) while leaving room
for legitimate governance to determine the metric's specific structure.

\subsection{9.7 Learning the Metric from
Data}\label{learning-the-metric-from-data}

\subsubsection{The Empirical Project}\label{the-empirical-project}

A practical question: can the metric be \emph{learned} from data? If we
have access to moral judgments --- ordinal comparisons, trade-off
behavior, evaluations of situations --- can we infer the underlying
metric that rationalizes them?

This question connects the governance account to empirical ethics. The
governance account says the metric is determined by legitimate
institutional processes. But those processes produce \emph{outputs} ---
laws, regulations, institutional decisions, advisory column verdicts ---
from which the operative metric can, in principle, be
reverse-engineered.

\subsubsection{Formal Setup}\label{formal-setup}

\textbf{The metric learning problem.} Given a dataset of moral judgments
\(\{(p_{i},q_{i}, \succ )\}\), where \(p \succ q\) means ``situation
\(p\) is morally preferable to situation \(q\),'' find a metric
\(g_{\mu\nu}\) such that:

\[d_{g}(p_{i},\text{ideal}) < d_{g}(q_{i},\text{ideal})\quad\text{whenever }p_{i} \succ q_{i}\]

where \(d_{g}\) is the geodesic distance in the metric \(g\), and
``ideal'' is the point in \(\mathcal{M}\) representing the best
achievable moral state.

This is the moral analogue of \emph{metric learning} in machine
learning: given similarity judgments, find the metric that reproduces
them.

\subsubsection{Evidence from the Dear Abby
Corpus}\label{evidence-from-the-dear-abby-corpus}

The Dear Abby corpus --- 20,030 real moral dilemmas with expert
ground-truth responses, spanning 32 years (1985--2017) --- provides a
natural dataset for metric learning. Each letter presents a moral
situation; the columnist's response reveals an implicit judgment about
the relevant trade-offs.

Analysis reveals several features of the operative metric:

\textbf{Context-dependent weighting.} The relative weights of moral
dimensions shift by context (Chapter 5, §5.3): in family contexts, care
(Dimension 7) dominates; in workplace contexts, procedural legitimacy
(Dimension 8) and fairness (Dimension 3) dominate; in neighbor disputes,
rights (Dimension 2) dominate. This is evidence that the operative
metric is not constant but varies across the manifold --- precisely the
structure of a Riemannian metric field.

\textbf{Temporal stability of structure.} Despite variation in specific
advice over three decades, the \emph{structural} features of the metric
--- the relative ordering of dimension weights, the pattern of
couplings, the location of stratum boundaries --- remain stable. This
suggests that while the metric's parameters may drift, its qualitative
structure is robust. The moral manifold has a stable topology, even if
its geometry evolves.

\textbf{Correlative symmetry.} The correlative rates O ↔ C (87\%) and L
↔ N (82\%) are consistent across subperiods and context types. This is
evidence of a \emph{structural invariant} of the metric --- a feature
preserved under the variation that metric learning would discover.

\subsubsection{Evidence from BIP
Experiments}\label{evidence-from-bip-experiments}

The Bond Invariance Principle experiments across 109,294 passages in 11
languages provide a different kind of evidence. The key finding: the
deontic axis --- the obligation/permission distinction --- transfers at
100\% across all languages tested, even when specific moral content
fails to transfer.

\textbf{Interpretation for metric learning.} The deontic axis is a
\emph{structural feature of the metric} that is independent of language
and culture. It is not learned from data --- it is a constraint that
data universally satisfies. This supports the claim (§9.6) that some
features of the metric are non-negotiable structural invariants, while
others are legitimately variable across communities.

The pattern is suggestive:

{\def\LTcaptype{none} % do not increment counter
\begin{longtable}[]{@{}
  >{\raggedright\arraybackslash}p{(\linewidth - 4\tabcolsep) * \real{0.2528}}
  >{\raggedright\arraybackslash}p{(\linewidth - 4\tabcolsep) * \real{0.2290}}
  >{\raggedright\arraybackslash}p{(\linewidth - 4\tabcolsep) * \real{0.4362}}@{}}
\toprule\noalign{}
\begin{minipage}[b]{\linewidth}\raggedright
\textbf{Feature}
\end{minipage} & \begin{minipage}[b]{\linewidth}\raggedright
\textbf{Status}
\end{minipage} & \begin{minipage}[b]{\linewidth}\raggedright
\textbf{Evidence}
\end{minipage} \\
\midrule\noalign{}
\endhead
\bottomrule\noalign{}
\endlastfoot
Deontic axis (O/L) & Structural invariant & \% cross-lingual transfer \\
Correlative symmetry & Structural invariant & --87\% across contexts \\
Nullifier structure & Structural invariant & Universal across domains \\
Dimension weightings & Governance variable & Context-dependent;
culture-sensitive \\
Off-diagonal couplings & Governance variable & Vary by tradition and
epoch \\
Lexicographic priority & Governance variable & Present in some systems,
absent in others \\
\end{longtable}
}

The structural invariants constrain the space of admissible metrics. The
governance-variable features are the parameters that legitimate
institutional processes determine.

\subsubsection{Challenges for Metric
Learning}\label{challenges-for-metric-learning}

\textbf{Underdetermination.} Many metrics may fit the same data. A
finite dataset of ordinal comparisons does not uniquely determine a
metric --- it determines only an equivalence class of metrics that are
ordinally equivalent. To pin down the cardinal structure (the actual
trade-off rates), richer data are needed: not just ``A is better than
B'' but ``A is \emph{twice as much} better than B as C is better than
D.''

\textbf{Noise.} Human moral judgments are noisy --- affected by framing,
mood, order effects, and cognitive biases. The metric learned from raw
judgments is contaminated by this noise. Denoising requires either very
large datasets (which the 109K-passage BIP corpus provides) or
structural assumptions (which the admissibility constraints of §9.6
provide).

\textbf{Whose judgments?} Aggregating moral judgments across persons
reintroduces the metric problem at a meta-level: what is the metric for
combining different people's judgments about the metric? The governance
account resolves this recursively: the combination procedure is itself a
governance choice, subject to legitimacy constraints.

\subsection{9.8 Pluralism and the
Meta-Metric}\label{pluralism-and-the-meta-metric}

\subsubsection{The Question}\label{the-question}

If different communities legitimately adopt different metrics, is there
a \emph{meta-metric} for comparing metrics? Can we say that one
community's metric is \emph{closer to correct} than another's?

\subsubsection{Three Positions}\label{three-positions}

\textbf{Metric monism.} There is exactly one correct moral metric
\(g^{*}\), and all governance processes are attempts to approximate it.
This is the realist position, transferred to the governance context. Its
difficulty: it requires a standard independent of governance, which the
governance account by design avoids.

\textbf{Metric relativism.} There is no basis for comparing metrics
across communities. If community A's metric permits a trade-off that
community B's metric forbids, neither is wrong --- they are simply
different. Its difficulty: it cannot account for moral criticism across
communities, which is a central feature of moral life. Calling another
community's practices ``wrong'' (as we do with slavery, genocide, and
systematic oppression) requires a standard that transcends the
community's own metric.

\textbf{Structured pluralism.} Metrics can be partially compared along
several dimensions --- but not totally ordered. There is no single
meta-metric, but there are \emph{meta-criteria}:

\begin{enumerate}
\def\labelenumi{\arabic{enumi}.}
\item
  \textbf{Internal consistency.} A metric that violates its own
  structural constraints (e.g., a metric that encodes a trade-off
  between dimensions it elsewhere declares incommensurable) is defective
  by its own standards.
\item
  \textbf{Coverage.} A metric that assigns zero weight to an entire
  dimension of the moral manifold --- ignoring, say, the environmental
  dimension entirely --- has a deficiency: it cannot even formulate
  moral claims along the neglected dimension.
\item
  \textbf{Invariance preservation.} A metric that violates the
  structural invariants of §9.6 (correlative symmetry,
  non-discrimination, transformation consistency) is inadmissible by
  structural criteria, not by any community's particular moral
  commitments.
\item
  \textbf{Reflective stability.} A metric is reflectively stable if the
  community that endorses it would continue to endorse it after full
  information and ideal deliberation. A metric adopted under conditions
  of ignorance or coercion is reflectively unstable.
\item
  \textbf{Empirical adequacy.} A metric is empirically adequate if it
  rationalizes the community's considered moral judgments (not all
  judgments --- considered ones, after removing distortions from bias,
  ignorance, and coercion).
\end{enumerate}

These meta-criteria partially order the space of metrics without totally
ordering it. Two metrics may each satisfy all five criteria while
differing in their off-diagonal components --- representing genuinely
different but equally legitimate moral perspectives.

\subsubsection{Formal Statement}\label{formal-statement-1}

\textbf{Definition 9.1 (Partial Order on Metrics).} \emph{Define the
relation} \(g \preccurlyeq g'\) \emph{(metric} \(g'\) \emph{dominates
metric} \(g\)\emph{) if and only if} \(g'\) \emph{satisfies all five
meta-criteria at least as well as} \(g\)\emph{, and strictly better on
at least one. This relation is a partial order on the space of
admissible metrics.}

The partial order has the structure of a \emph{lattice} in favorable
cases: given two metrics, there may be a \emph{join} (the least upper
bound --- the weakest metric that dominates both) and a \emph{meet} (the
greatest lower bound --- the strongest metric dominated by both).
Whether the space of admissible metrics forms a lattice is an open
question with substantial mathematical interest.

\subsubsection{The Irreducibility of Metric
Disagreement}\label{the-irreducibility-of-metric-disagreement}

Structured pluralism implies that some metric disagreements are
\emph{irreducible}: they cannot be resolved by appeal to meta-criteria,
because both metrics satisfy all meta-criteria equally well. At this
level, disagreement is genuinely political --- it reflects different
legitimate choices about how to structure moral trade-offs.

This is not a defect of the framework. It is a structural theorem: the
space of admissible metrics, constrained by the criteria of §9.6, is not
a singleton. Genuine moral pluralism exists at the metric level --- and
any framework that pretends otherwise is either concealing its
assumptions or restricting the space of admissible moral perspectives
too narrowly.

The geometric framework's contribution is to make metric disagreement
\emph{precise}: we can say exactly where two metrics differ (in which
components \(g_{\mu\nu}\)), what moral consequences follow from the
difference, and which meta-criteria each metric satisfies. This is
progress --- not the progress of resolving all disagreement, but the
deeper progress of understanding its structure.

\subsection{9.9 The Metric as Moral
Infrastructure}\label{the-metric-as-moral-infrastructure}

\subsubsection{An Analogy}\label{an-analogy}

Consider a city's transportation infrastructure: roads, bridges, transit
lines, sidewalks. This infrastructure determines how people can move ---
which destinations are accessible, at what cost, by what routes. The
infrastructure is neither a natural fact (cities don't grow roads
spontaneously) nor a subjective preference (my desire for a road doesn't
create one). It is a \emph{governance artifact}: the output of
institutional processes (city planning, legislation, funding decisions)
with real consequences for how people live.

The moral metric is moral infrastructure. It determines how moral agents
can reason --- which values are comparable, at what trade-off rates,
through what deliberative routes. It is neither a natural fact nor a
subjective preference. It is a governance artifact with real
consequences for how moral life is structured.

And like physical infrastructure:

\begin{itemize}
\item
  \textbf{It can be well-designed or poorly designed.} A city with no
  sidewalks privileges drivers over pedestrians; a moral metric with
  zero weight on environmental impact privileges present benefit over
  future harm.
\item
  \textbf{It is designed for a community.} The infrastructure of Tokyo
  differs from the infrastructure of rural Montana, and both may be
  well-designed for their contexts. Similarly, different communities may
  have different moral metrics, both well-suited to their circumstances.
\item
  \textbf{It constrains without determining.} Infrastructure constrains
  which trips are easy and which are hard, but it does not determine
  where anyone actually goes. The moral metric constrains which
  trade-offs are legitimate and which are not, but it does not determine
  which actions anyone actually takes. The metric is the field on which
  moral agency plays out --- not a substitute for agency.
\item
  \textbf{It requires maintenance and revision.} Infrastructure that was
  adequate in 1950 may be inadequate in 2025. A moral metric that was
  adequate before the development of powerful AI may be inadequate
  after. The governance account allows --- indeed requires --- periodic
  revision of the metric through legitimate institutional processes.
\end{itemize}

\subsubsection{Implications for AI
Systems}\label{implications-for-ai-systems}

The governance account has direct implications for AI alignment. If the
moral metric is a governance artifact, then AI systems that make morally
significant decisions must:

\begin{enumerate}
\def\labelenumi{\arabic{enumi}.}
\item
  \textbf{Be calibrated to the community's metric.} An AI system
  operating in a community with a utilitarian metric should not impose a
  lexicographic one, and vice versa.
\item
  \textbf{Be transparent about the metric.} The operative metric should
  be inspectable --- which dimensions are weighted, what trade-offs are
  permitted, where the constraint surfaces lie.
\item
  \textbf{Be updatable.} When the community revises its metric through
  legitimate processes, the AI system should be updated accordingly.
\item
  \textbf{Respect the structural constraints.} The invariance
  requirements (§9.6) apply to AI systems as much as to human moral
  reasoning. The Bond Invariance Principle requires that the AI's moral
  evaluations be invariant under admissible redescription --- a testable
  engineering requirement.
\end{enumerate}

This is the connection between the philosophical question of this
chapter (where does the metric come from?) and the engineering question
of Chapter 18 (how do we build ethical AI?). The governance account
provides the \emph{source} of the metric; the geometric framework
provides the \emph{form}; the engineering implementation provides the
\emph{mechanism}.

\subsection{9.10 Summary}\label{summary-2}

The moral metric \(g_{\mu\nu}\) --- the structure of permissible
trade-offs between moral dimensions --- is the most consequential and
most contentious element of the geometric framework.

Four accounts of its origin:

{\def\LTcaptype{none} % do not increment counter
\begin{longtable}[]{@{}
  >{\raggedright\arraybackslash}p{(\linewidth - 6\tabcolsep) * \real{0.1748}}
  >{\raggedright\arraybackslash}p{(\linewidth - 6\tabcolsep) * \real{0.2863}}
  >{\raggedright\arraybackslash}p{(\linewidth - 6\tabcolsep) * \real{0.3222}}
  >{\raggedright\arraybackslash}p{(\linewidth - 6\tabcolsep) * \real{0.2167}}@{}}
\toprule\noalign{}
\begin{minipage}[b]{\linewidth}\raggedright
\textbf{Account}
\end{minipage} & \begin{minipage}[b]{\linewidth}\raggedright
\textbf{Strengths}
\end{minipage} & \begin{minipage}[b]{\linewidth}\raggedright
\textbf{Weaknesses}
\end{minipage} & \begin{minipage}[b]{\linewidth}\raggedright
\textbf{Role of the Metric}
\end{minipage} \\
\midrule\noalign{}
\endhead
\bottomrule\noalign{}
\endlastfoot
Realism & Explains moral progress, error & Epistemological and
metaphysical difficulties & A fact to be discovered \\
Constructivism & Grounds authority in reason & Regress of procedures;
underdetermination & A fixed point of ideal deliberation \\
Expressivism & Explains motivation, diversity & Cannot explain
convergence or structure & A projection of sensibility \\
Governance & Explains authority, revisability, pluralism; connects to
institutional practice & Requires legitimacy theory; risks deference to
unjust institutions & An artifact of legitimate institutional process \\
\end{longtable}
}

The governance account is the most natural fit for the geometric
framework. It locates the metric's authority in the same place that the
framework's \emph{implementation} must locate it: in the institutional
processes that determine, for a community, how moral trade-offs are to
be structured.

The space of admissible metrics is constrained by structural
requirements --- symmetry, non-discrimination, invariance, compatibility
with structural invariants --- but not reduced to a singleton. Genuine
moral pluralism exists within these constraints, and the framework makes
this pluralism precise rather than concealing it.

The metric is moral infrastructure: designed, revisable,
community-specific within structural constraints, and consequential for
how moral life is lived. The geometric framework does not provide the
metric. It provides the vocabulary for articulating, comparing,
constraining, and implementing whatever metric a community adopts
through legitimate governance.

\subsection{Technical Appendix}\label{technical-appendix-2}

\textbf{Definition A.1 (Admissible Metric).} \emph{A metric}
\(g_{\mu\nu}\) \emph{on the moral manifold} \(\mathcal{M}\) \emph{is
admissible if it satisfies:} 1. \emph{Symmetry:}
\(g_{\mu\nu} = g_{\nu\mu}\)\emph{.} 2. \emph{Non-total-degeneracy: there
exists a subspace} \(V \subset T_{p}\mathcal{M}\) \emph{of dimension}
\(\geq 1\) \emph{on which} \(g|_{V}\) \emph{is non-degenerate.} 3.
\emph{Tensor transformation:} \(g\) \emph{transforms as a (0,2)-tensor
under Type 1 transformations.} 4. \emph{Non-discrimination:}
\(g_{\mu\nu}(\phi(p)) = g_{\mu\nu}(p)\) \emph{for all morally irrelevant
transformations} \(\phi\)\emph{.} 5. \emph{Structural compatibility:}
\(g\) \emph{is compatible with the} \(D_{4}\) \emph{group structure and
known nullifier constraints.}

\textbf{Proposition 9.1 (Non-Uniqueness of Admissible Metrics).}
\emph{The space of admissible metrics on} \(\mathcal{M}\) \emph{is
non-empty and has dimension} \(> 0\)\emph{. In particular, there exist
admissible metrics that differ in their off-diagonal components while
satisfying all five admissibility conditions.}

\emph{\textbf{Proof}. The space Adm(M) is non-empty by Theorem 9.1(i),
and has dimension ≥ 10 by Theorem 9.1(iii). In particular, the flat
metric δ\_μν and any perturbation δ\_μν + ε h\_μν (for symmetric h
respecting the D₄ and U(1)\_H constraints) are both admissible for
sufficiently small ε \textgreater{} 0. These metrics differ on
off-diagonal components, establishing positive dimension.
□}\(gg'\epsilon h_{\mu\nu}\epsilon g' = g + \epsilon h▫\)

Note on logical structure. This proposition's proof references Theorem
9.1 (stated below). The logical dependence is not circular: Theorem 9.1
provides the explicit dimension count for Adm(M), which is used here
only as an existence claim (Adm(M) is non-empty with positive
dimension). The existence of admissible metrics can also be established
directly: the flat Euclidean metric on ℝ⁹ satisfies all axioms, and
small perturbations of it remain admissible.

\textbf{Proposition 9.2 (Structural Invariants Constrain the Metric).}
\emph{If the correlative symmetry rates}
\(P(\text{O} \leftrightarrow \text{C}) = p_{1}\) \emph{and}
\(P(\text{L} \leftrightarrow \text{N}) = p_{2}\) \emph{are fixed
structural invariants, then the diagonal components of the metric on the
rights/duties dimension (Dimension 2) must satisfy:}

\[g_{22}(a) = f(p_{1},p_{2}) \cdot g_{22}(b)\quad\text{for correlative agents }a,b\]

\emph{where} \(f\) \emph{is a function of the correlative rates. Perfect
correlative symmetry (}\(p_{1} = p_{2} = 1\)\emph{) implies} \(f = 1\)
\emph{--- identical metric weights for correlative partners.}

\emph{\textbf{Proof}. Under the correlative swap s\_c: O ↔ C, L ↔ N, the
expected satisfaction functional must be invariant: ⟨Σ⟩\_\{s\_c\} = ⟨Σ⟩.
The proportionality factor f must satisfy: (1) f(p₁, p₂) =
g₂₂(a)/g₂₂(b), which by the symmetry of the correlative exchange yields
f(p₁, p₂) = p₁/p₂; (2) f(1,1) = 1 (perfect symmetry implies equal metric
weights); and (3) f(p₁, p₂) · f(p₂, p₁) = 1 (applying the swap twice
returns to the original). All three are satisfied by f(p₁, p₂) = p₁/p₂,
and this is the unique solution consistent with the correlative symmetry
constraint. □}\(abp_{1}baf(p_{1},p_{2})▫\)

\subsubsection{}\label{section-1}

\subsubsection{Characterization of the Admissible Metric
Space}\label{characterization-of-the-admissible-metric-space}

The preceding propositions established that the space of admissible
metrics is non-trivial. We now prove structural results that
characterize this space precisely.

\textbf{Lemma 9.1 (D₄ Constraint on Metric Components). {[}Conditional
on Def. A.1, condition 5{]}} \emph{Let g\_μν be an admissible metric on
M. The D₄ structural compatibility requirement forces g\_\{2μ\}(O) =
g\_\{2μ\}(C) = g\_\{2μ\}(N) = g\_\{2μ\}(L) for all μ. This eliminates 27
degrees of freedom.}

\emph{\textbf{Proof}.} D₄ acts on the Hohfeldian positions via the
4-cycle r = (O C N L) and reflection s\_c = (O C)(L N). Structural
compatibility (Def. A.1, condition 5) requires g\_μν(σ(p)) = g\_μν(p)
for all σ ∈ D₄. Since D₄ acts transitively on \{O, C, N, L\}, the orbit
of O under ⟨r⟩ is the entire set, so g\_\{2ν\}(O) = g\_\{2ν\}(C) =
g\_\{2ν\}(N) = g\_\{2ν\}(L). This collapses four independent values to
one for each of nine metric components g\_\{2μ\}, eliminating 9 × (4 −
1) = 27 degrees of freedom. □

\textbf{Lemma 9.2 (U(1)\_H Constraint on the Harm Dimension).
{[}Conditional on Def. A.1, condition 5{]} The U(1)\_H symmetry
generates a Killing vector field ξ = ∂/∂x⁷ along Dimension 7. The
Killing equation ensures that the metric is independent of x⁷
(∂7g\_\{μν\} = 0 for all μ, ν). Combined with the product-structure
interpretation of Axiom A5 --- that the harm dimension is geometrically
orthogonal to all other dimensions --- this forces g\_\{7μ\} = 0 for all
μ ≠ 7. This eliminates 8 off-diagonal degrees of freedom.}

\emph{\textbf{Proof}.} The U(1)\_H action is a continuous one-parameter
group of isometries along Dimension 7. Its generator is the Killing
vector field ξ = ∂/∂x⁷. The Killing equation ∇\_μ ξ\_ν + ∇\_ν ξ\_μ = 0
forces ∂₇g\_\{μν\} = 0 for all μ, ν --- i.e., the metric is independent
of the harm coordinate x⁷ (if g\_\{7μ\} ≠ 0 for some μ ≠ 7, geodesics in
the harm direction would couple to non-harm dimensions, breaking the
independence of harm magnitude). This eliminates 8 off-diagonal
components: g\_\{71\}, g\_\{72\}, \ldots, g\_\{76\}, g\_\{78\},
g\_\{79\}. □

\textbf{Theorem 9.1 (Structure of the Admissible Metric Space).
{[}Conditional on Def. A.1{]} Let Sym⁺₉(G) denote the space of
positive-definite symmetric bilinear forms on ℝ⁹ invariant under the
gauge group G = D₄ × U(1)\_H. The space Adm(M) of admissible metrics
admits a Whitney stratification:}

(i) Adm(M) is non-empty: the flat metric δ\_μν ∈ Adm(M).

(ii) Adm(M) has the structure of a stratified space Adm(M) = ⨆\_α S\_α,
where: the top stratum S₀ consists of all fully non-degenerate
admissible metrics (those with det(g\textbar ₐ) \textgreater{} 0 on
every protected subspace ₐ), and is a smooth open submanifold of
Sym⁺₉(G) of dimension ≥ 10; the boundary strata S\_k (k ≥ 1) consist of
metrics where non-degeneracy fails on exactly k protected dimensions,
representing incommensurability surfaces where specific value trade-offs
become undefined; the stratification satisfies the Whitney (b)
condition, inheriting regularity from the smooth structure of the
ambient space.

(iii) dim S₀ ≥ 45 − 27 − 8 = 10. The boundary strata have strictly lower
dimension.

\emph{\textbf{Proof. (i) Non-emptiness. The flat metric δ\_μν satisfies
all five admissibility conditions (symmetry, non-degeneracy, tensor
transformation, non-discrimination, and structural compatibility with
G). Since det(δ\textbar ₐ) = 1 \textgreater{} 0 on every subspace, δ ∈
S₀ (the top stratum). (ii) Stratified structure. The gauge-invariance
constraints are linear, so Sym⁺₉(G) is a linear subspace of the
45-dimensional space of symmetric bilinear forms, intersected with the
open positive-definite cone. The admissibility conditions partition this
space into strata by the rank of the restriction g\textbar ₐ to each
protected subspace ₐ (corresponding to protected moral dimensions from
Def. A.1, condition 2). The top stratum S₀ = \{g ∈ Sym⁺₉(G) :
det(g\textbar ₐ) \textgreater{} 0 for all protected ₐ, plus
non-discrimination and structural compatibility\} is the intersection of
finitely many open conditions (det(g\textbar ₐ) \textgreater{} 0 is open
for each ₐ) with the linear and smooth constraints, yielding an open
subset of a finite-dimensional affine subspace --- hence a smooth open
submanifold. The boundary strata S\_k = \{g : det(g\textbar ₐ) = 0 for
exactly k protected subspaces\} are smooth submanifolds of codimension k
(since each det(g\textbar ₐ) = 0 is a smooth hypersurface in generic
position). The Whitney (b) condition holds because the strata are
defined by polynomial equalities and inequalities on the components of
g, and polynomial stratifications satisfy Whitney (b) (see Bochnak,
Coste, and Roy, Real Algebraic Geometry, §9.7). Path-connectedness of S₀
follows from convexity: if g, g' ∈ S₀, then (1−t)g + tg' ∈ Sym⁺₉(G) for
all t ∈ {[}0,1{]} (positive-definite forms are closed under convex
combination), and det(((1−t)g + tg')\textbar ₐ) \textgreater{} 0 for all
t ∈ {[}0,1{]} (by the Minkowski determinant inequality: det(A +
B)\^{}\{1/n\} ≥ det(A)\^{}\{1/n\} + det(B)\^{}\{1/n\} for
positive-definite A, B). Thus S₀ is convex, hence path-connected. (iii)
Dimension bound. A symmetric 9×9 matrix has 45 independent components.
D₄ invariance eliminates 27 (Lemma 9.1). U(1)\_H invariance eliminates 8
(Lemma 9.2). The open conditions (non-degeneracy, non-discrimination) do
not reduce dimension. Thus dim S₀ ≥ 10. The boundary strata have dim
S\_k = dim S₀ − k. □}}

\textbf{Theorem 9.2 (Structured Pluralism). {[}Conditional on Def. A.1
and Def. 9.1{]}} \emph{The partial order ≼ on Adm(M) is not total. That
is, there exist admissible metrics g, g' ∈ Adm(M) such that neither g ≼
g' nor g' ≼ g.}

\emph{\textbf{Proof. We construct two admissible metrics in S₀ that are
incomparable under ≼. The argument proceeds in four steps.\\
\strut \\
Step 0 (Explicit G-invariant perturbation directions). We first
construct two linearly independent elements of Sym₉(G), the space of D₄
× U(1)\_H-invariant symmetric bilinear forms on ℝ⁹. By Lemma 9.1, D₄
invariance requires that the metric components g\_\{μν\} are constant
across Hohfeldian states \{O, C, L, N\} for each dimension pair (μ, ν).
By Lemma 9.2, U(1)\_H invariance requires ∂₇g\_\{μν\} = 0. These
constraints are linear and leave a ≥10-dimensional subspace (Theorem
9.1(iii)). Within this subspace, define the following two elements
explicitly:\\
e¹³\_\{μν\} = ½(δ\_\{μ1\}δ\_\{ν3\} + δ\_\{μ3\}δ\_\{ν1\}), the symmetric
tensor with e¹³\_\{13\} = e¹³\_\{31\} = ½ and all other components
zero;\\
e⁴⁵\_\{μν\} = ½(δ\_\{μ4\}δ\_\{ν5\} + δ\_\{μ5\}δ\_\{ν4\}), the symmetric
tensor with e⁴⁵\_\{45\} = e⁴⁵\_\{54\} = ½ and all other components
zero.\\
Both are symmetric by construction. Both satisfy D₄ invariance: e¹³
couples Dimensions 1 and 3 (consequences and fairness), which lie in the
same D₄-orbit (both are evaluative dimensions unaffected by Hohfeldian
state transitions); e⁴⁵ couples Dimensions 4 and 5 (autonomy and trust),
similarly within a single orbit. Both satisfy U(1)\_H invariance:
neither has components along Dimension 7 (harm), so ∂₇e\_\{μν\} = 0
trivially. Both are spatially constant, satisfying non-discrimination.
They are linearly independent (they have disjoint nonzero entries). Thus
e¹³, e⁴⁵ ∈ Sym₉(G) are two explicit, independent, admissible
perturbation directions.\\
\strut \\
Step 1 (Tangent directions). δ ∈ S₀ (Theorem 9.1(i)). S₀ is open in the
affine subspace δ + Sym₉(G) ⊂ Sym⁺₉ (Theorem 9.1(ii)). Therefore T\_δ S₀
= Sym₉(G), and in particular e¹³, e⁴⁵ ∈ T\_δ S₀.\\
\strut \\
Step 2 (Admissible curves). Define γ₁(t) = δ + t e¹³ and γ₂(t) = δ + t
e⁴⁵ for t ≥ 0. Since S₀ is open and δ ∈ S₀, there exists ε
\textgreater{} 0 such that B(δ, ε) ∩ (δ + Sym₉(G)) ⊂ S₀. For 0
\textless{} t \textless{} ε/‖e¹³‖ (resp. ε/‖e⁴⁵‖), we have γ₁(t), γ₂(t)
∈ S₀ ⊂ Adm(M). Fix such a t and set g = γ₁(t), g' = γ₂(t).\\
\strut \\
Step 3 (Incomparability). g = δ + t e¹³ assigns nonzero off-diagonal
weight g\_\{13\} = g\_\{31\} = t/2 \textgreater{} 0 to the
consequences--fairness cross-term, making welfare--justice trade-offs
metrically meaningful. g' = δ + t e⁴⁵ does the same for autonomy--trust.
Consider a community C₁ where welfare--justice coupling is empirically
supported, and a community C₂ where autonomy--trust coupling is
supported. Then g achieves higher empirical adequacy for C₁ while g'
achieves higher adequacy for C₂. Under any ordering ≼ respecting
empirical adequacy, neither g ≼ g' nor g' ≼ g. □}}

This is not a deficiency of the framework but a theorem: genuine moral
pluralism --- the existence of incomparable but individually coherent
moral evaluations --- is a mathematical consequence of the admissibility
constraints, not a philosophical assumption.

❖

\emph{The metric is not a datum and not a guess. It is a decision.}

\emph{A decision made by institutions, constrained by structure,
revisable by process, and consequential for every moral evaluation that
operates within it.}

\emph{This chapter has argued that the question ``Which metric is
correct?'' is not a philosophical puzzle with a unique answer, but a
governance challenge with structural constraints.}

\emph{The geometric framework does not resolve the challenge. It gives
us the language to face it honestly --- to know exactly what we are
choosing when we choose a metric, and what follows from the choice we
make.}

\section{Chapter 10: Moral Dynamics --- Parallel Transport, Holonomy,
and the Curvature of Moral
Space}\label{chapter-10-moral-dynamics-parallel-transport-holonomy-and-the-curvature-of-moral-space}

\begin{quote}
\textbf{\textsc{RUNNING EXAMPLE --- Priya's Model}}

\emph{Priya presents her findings to two audiences in different orders
and gets two different outcomes. Path A: she begins with Mrs. Voss's
story, then shows the statistical pattern. The team is moved; they want
to act. Path B: she leads with statistics, then mentions individual
cases. Executives focus on legal exposure and suggest a disclaimer. Same
facts, same obligation---but the parallel transport of the argument
through different conversational paths produces different moral
conclusions. This is holonomy. The curvature of HealthBridge's
organizational moral space means the order of deliberation changes the
outcome. She cannot simply present `the truth.' She must choose a path,
and the path matters.}
\end{quote}

\subsection{10.1 Ethics in Motion}\label{ethics-in-motion}

The preceding chapters developed a static picture: the moral manifold
\(\mathcal{M}\) (Chapter 5), the tensors that live on it (Chapter 6),
the metric that measures them (Chapters 6, 9), and the stratified
boundaries where the rules change (Chapter 8). This apparatus captures
the geometry of a \emph{single moment} --- the structure of obligations,
interests, and trade-offs at one point in moral space.

But ethics is not static. Situations change. Obligations evolve. Moral
understanding develops. A promise made yesterday must be carried into
today's altered circumstances. A principle learned in one context must
be applied in another. An obligation incurred in peacetime must be
renegotiated in crisis. Ethics is not a photograph but a trajectory ---
a path through structured space.

This chapter introduces the dynamics of geometric ethics: how moral
tensors change as we move through the moral manifold, how obligations
are ``transported'' across changing circumstances, and what it means for
moral evaluation to be path-dependent. The mathematical tools are the
covariant derivative, parallel transport, holonomy, and curvature ---
developed abstractly in Chapter 4 (§4.6) and now given their moral
content.

\subsection{10.2 The Problem of Moral
Change}\label{the-problem-of-moral-change}

Consider a simple scenario. Alice promises to meet Bob for coffee on
Saturday. On Friday, Alice learns that Carol is having a medical
emergency and needs help on Saturday morning. What happens to Alice's
obligation to Bob?

A static analysis takes two snapshots:

\begin{itemize}
\item
  At \(t_{0}\): Alice has obligation \(O(t_{0})\) = ``meet Bob for
  coffee.''
\item
  At \(t_{1}\): Alice has obligation \(O(t_{1})\) = ``help Carol with
  the emergency.''
\end{itemize}

These are just two points in moral space, with two obligation vectors.
The static analysis can compare them --- \(O(t_{1})\) points in a
different direction from \(O(t_{0})\) --- but it cannot answer the
dynamic questions:

\begin{enumerate}
\def\labelenumi{\arabic{enumi}.}
\item
  \textbf{How did} \(O(t_{0})\) \emph{\textbf{become}}
  \(O(t_{1})\)\textbf{?} Was the obligation to Bob destroyed,
  overridden, suspended, or transformed?
\item
  \textbf{Is there continuity?} Is \(O(t_{1})\) the ``same'' obligation
  as \(O(t_{0})\), reoriented by circumstances? Or is it a new
  obligation that replaced the old one?
\item
  \textbf{Is there residue?} Does the displaced obligation to Bob leave
  a trace --- a secondary obligation to apologize, to reschedule, to
  acknowledge the broken commitment?
\item
  \textbf{Is the result path-dependent?} Would Alice's final obligation
  be different if she had learned about Carol's emergency \emph{before}
  making the promise to Bob?
\end{enumerate}

These are questions about \emph{moral dynamics} --- how moral tensors
evolve along paths through \(\mathcal{M}\). To answer them, we need the
apparatus of connections, parallel transport, and curvature.

\subsection{10.3 The Moral Connection}\label{the-moral-connection}

\subsubsection{Comparing Obligations Across
Contexts}\label{comparing-obligations-across-contexts}

On a flat space, comparing vectors at different points is trivial:
translate one to the other's location and compare. On a curved space,
there is no canonical translation --- the comparison depends on the
\emph{path} taken between the two points. A \emph{connection} is the
additional structure that specifies how to make such comparisons.

\textbf{Definition 10.1 (Moral Connection).} \emph{A moral connection is
an affine connection} \(\nabla\) \emph{on the moral manifold}
\(\mathcal{M}\) \emph{(within each stratum) that specifies, for any
direction of change} \(X\) \emph{and any obligation field} \(O\)\emph{,
the rate of change of} \(O\) \emph{in the direction} \(X\)\emph{:}

Remark. The definition of the moral connection is restricted to
individual strata. Across stratum boundaries, the connection may be
discontinuous. The junction conditions governing cross-boundary behavior
--- including continuity of parallel transport and matching of holonomy
--- are stated formally in Definition 13.6.

\[\nabla_{X}O = \text{“how }O\text{ changes as the situation moves in direction }X\text{”}\]

In coordinates, the covariant derivative of an obligation vector field
\(O^{\mu}\) in the direction \(\partial_{i}\) is:

\[(\nabla_{i}O)^{\mu} = \frac{\partial O^{\mu}}{\partial x^{i}} + \Gamma_{i\nu}^{\mu}O^{\nu}\]

The first term, \(\partial O^{\mu}/\partial x^{i}\), is the \emph{naive}
rate of change --- how the components of \(O\) change if we simply track
them in coordinates. The second term, \(\Gamma_{i\nu}^{\mu}O^{\nu}\), is
the \emph{correction} --- the adjustment needed because the coordinate
system itself is changing as we move through moral space. The
Christoffel symbols \(\Gamma_{i\nu}^{\mu}\) encode this correction.

\subsubsection{What the Connection
Encodes}\label{what-the-connection-encodes}

The moral connection has a precise interpretation: it encodes how
\emph{moral concepts translate across contexts}.

Consider the obligation ``be kind.'' In a family context, kindness means
warmth, patience, emotional availability. In a professional context,
kindness means fairness, clarity, respect for boundaries. The
\emph{word} is the same; the \emph{content} --- the obligation vector
--- has rotated.

The connection \(\nabla\) specifies exactly how this rotation occurs.
Moving from the family region of \(\mathcal{M}\) to the professional
region, the Christoffel symbols tell us:

\[\Gamma_{(\text{context}),7}^{7} < 0\quad\text{(care obligation decreases)}\]

\[\Gamma_{(\text{context}),7}^{8} > 0\quad\text{(procedural legitimacy obligation increases)}\]

The connection does not say whether this rotation is ``correct'' ---
that depends on the metric (Chapter 9). It says: \emph{if you carry the
obligation faithfully from one context to another, here is how its
components change}. The connection is the \emph{grammar} of moral
translation.

\subsubsection{The Levi-Civita
Connection}\label{the-levi-civita-connection}

If the moral manifold is equipped with a metric \(g_{\mu\nu}\) (Chapter
6), there is a distinguished connection: the \emph{Levi-Civita
connection}, which is the unique connection that is:

\begin{enumerate}
\def\labelenumi{\arabic{enumi}.}
\item
  \textbf{Metric-compatible}: \(\nabla g = 0\). Parallel transport
  preserves inner products --- if two obligations are orthogonal at one
  point, they remain orthogonal after transport.
\item
  \textbf{Torsion-free}:
  \(\nabla_{X}Y - \nabla_{Y}X = \lbrack X,Y\rbrack\). There is no
  ``twisting'' of the coordinate frame as we move through moral space.
\end{enumerate}

The Christoffel symbols of the Levi-Civita connection are determined
entirely by the metric:

\[\Gamma_{\nu\rho}^{\mu} = \frac{1}{2}g^{\mu\sigma}\left( \frac{\partial g_{\sigma\nu}}{\partial x^{\rho}} + \frac{\partial g_{\sigma\rho}}{\partial x^{\nu}} - \frac{\partial g_{\nu\rho}}{\partial x^{\sigma}} \right)\]

This is significant: \textbf{if we know the metric, we know the
connection.} The governance account of Chapter 9 --- which determines
the metric --- thereby also determines the dynamics. The trade-off
structure and the transport rules are not independent choices; both are
encoded in \(g_{\mu\nu}\).

\subsubsection{Non-Metric Connections}\label{non-metric-connections}

Must the moral connection be the Levi-Civita connection? Not
necessarily. A connection with \emph{torsion} would mean that the moral
coordinate frame itself twists as we move --- that there is a systematic
asymmetry in how moral concepts translate, beyond what the metric
accounts for.

\textbf{Example.} Suppose that transporting an obligation from context A
to context B and then from B to C yields a different result than
transporting directly from A to C, \emph{even after accounting for
curvature}. This excess rotation is torsion. It would represent a
structural asymmetry in moral translation: the order in which you
traverse contexts matters, not because of curvature (which is symmetric)
but because of a systematic twist.

Whether moral space has torsion is an open empirical question (Chapter
29). For the remainder of this chapter, we work with the torsion-free
Levi-Civita connection, noting where the analysis would change with
torsion.

\subsection{10.4 Parallel Transport of
Obligations}\label{parallel-transport-of-obligations}

\includesvg[width=5.63077in,height=4.575in]{media/image14.svg}

\textbar{} Parallel Transport \& Holonomy

\subsubsection{The Concept}\label{the-concept}

\emph{Parallel transport} is the operation of carrying a vector along a
path while keeping it ``as constant as possible,'' given the connection.
See Figure 8 for a visual schematic.

\textbf{Definition 10.2 (Parallel Transport).} \emph{An obligation
vector} \(O\) \emph{is parallel-transported along a curve} \(\gamma(t)\)
\emph{in} \(\mathcal{M}\) \emph{if:}

\[\nabla_{\dot{\gamma}}O = 0\]

\emph{That is, the covariant derivative of} \(O\) \emph{along the curve
vanishes.}

In coordinates, this becomes the system of ordinary differential
equations:

\[\frac{dO^{\mu}}{dt} + \Gamma_{\nu\rho}^{\mu}{\dot{\gamma}}^{\nu}O^{\rho} = 0\]

Given an initial obligation vector \(O(0)\) at \(\gamma(0)\), the
parallel transport equations determine a unique \(O(t)\) at every point
along the curve.

\subsubsection{Moral Interpretation}\label{moral-interpretation-1}

Parallel transport is the \emph{faithful maintenance} of an obligation
across changing circumstances. When Alice carries her promise to Bob
from a normal Saturday context into a Saturday-with-emergency context,
she is parallel-transporting the obligation.

But ``faithful maintenance'' does not mean ``keeping the same
components.'' If the coordinate system --- the local meaning of ``be
helpful,'' ``be punctual,'' ``be caring'' --- changes from context to
context, then faithful maintenance requires adjusting the components to
compensate. The obligation is the \emph{same} in a geometric sense (it
satisfies \(\nabla_{\dot{\gamma}}O = 0\)), but its expression in local
coordinates changes.

\textbf{Example: Carrying ``keep your promise'' across contexts.}

Let \(O(0) = (0,0.8,0,0,0,0,0.4,0,0)\) --- an obligation with strong
duty component (\(O^{2} = 0.8\)) and moderate care component
(\(O^{7} = 0.4\)).

As Alice moves from the ``casual Saturday'' context to the ``emergency
Saturday'' context, the connection rotates the obligation:

\[O^{1}\text{ increases (welfare dimension becomes salient)}\]

\[O^{2}\text{ decreases (the promise’s binding force weakens in emergency)}\]

\[O^{7}\text{ remains stable (care is context-robust)}\]

The parallel-transported obligation might be:

\[O(1) \approx (0.5,0.3,0,0,0,0,0.4,0,0)\]

The duty component has partially converted to a welfare component ---
the obligation has \emph{rotated} in moral space. The promise has not
been destroyed; it has been faithfully carried into a new context where
its expression changes.

\subsubsection{\texorpdfstring{What Is \emph{Not} Parallel
Transport}{What Is Not Parallel Transport}}\label{what-is-not-parallel-transport}

Not every change in moral obligations is parallel transport. Some
changes represent \emph{genuine new information} --- a new obligation
that arises from the situation, not a transported version of the old
one. If Carol asks Alice directly for help (creating a new obligation),
this is not the old obligation transported --- it is a new vector added
to the tangent space.

The covariant derivative distinguishes these cases. If
\(\nabla_{\dot{\gamma}}O = 0\), the change in \(O\) is entirely due to
transport (the obligation is faithfully maintained). If
\(\nabla_{\dot{\gamma}}O \neq 0\), there is a \emph{genuine} change ---
new information, new duties, new considerations --- beyond what
transport accounts for. The magnitude \(|\nabla_{\dot{\gamma}}O|\)
measures the \emph{rate of genuine moral change} at each point along the
path.

\subsection{10.5 Holonomy: The Path-Dependence of Moral
Evaluation}\label{holonomy-the-path-dependence-of-moral-evaluation}

\subsubsection{The Phenomenon}\label{the-phenomenon}

The most striking consequence of curvature is \emph{holonomy}: if we
parallel-transport a vector around a closed loop, it may not return to
its original value. The vector has been ``rotated'' by the loop, even
though the loop returns to the same point.

\textbf{Definition 10.3 (Holonomy).} \emph{Let}
\(\gamma:\lbrack 0,1\mathcal{\rbrack \rightarrow M}\) \emph{be a closed
loop with} \(\gamma(0) = \gamma(1) = p\)\emph{. The holonomy of}
\(\gamma\) \emph{is the linear transformation}
\(H_{\gamma}:T_{p}\mathcal{M \rightarrow}T_{p}\mathcal{M}\)
\emph{defined by:}

\[H_{\gamma}(O) = \text{parallel transport of }O\text{ around }\gamma\]

\emph{If} \(H_{\gamma} = \text{id}\) \emph{for all loops}
\(\gamma\)\emph{, the connection is flat and moral space has no
curvature at} \(p\)\emph{. If} \(H_{\gamma} \neq \text{id}\) \emph{for
some} \(\gamma\)\emph{, the connection has nontrivial curvature.}

\subsubsection{A Moral Example}\label{a-moral-example}

Consider the following circuit in moral space:

\begin{enumerate}
\def\labelenumi{\arabic{enumi}.}
\item
  \textbf{Start}: Alice has a normal friendship with Bob. Obligation:
  moderate care (\(O^{7} = 0.5\)).
\item
  \textbf{Step 1}: Alice learns Bob's secret. The obligation gains a
  confidentiality component (\(O^{5}\) increases --- privacy dimension
  becomes active).
\item
  \textbf{Step 2}: Alice discovers the secret involves potential harm to
  Carol. The obligation acquires a duty-to-warn component (\(O^{1}\)
  increases --- welfare dimension activates).
\item
  \textbf{Step 3}: Alice warns Carol; the harm is averted. The immediate
  welfare obligation is discharged (\(O^{1}\) returns to baseline).
\item
  \textbf{Step 4}: The situation resolves; Alice and Bob return to
  ``normal'' friendship.
\end{enumerate}

Has the obligation returned to its original value? Geometrically, we
have parallel-transported the obligation around a loop: normal
friendship → confidentiality → duty-to-warn → resolution → normal
friendship.

If moral space is \emph{flat}, the obligation returns exactly to
\(O^{7} = 0.5\): the friendship is unchanged. But if moral space is
\emph{curved}, the holonomy \(H_{\gamma}\) rotates the obligation. The
care component may have shifted:

\[H_{\gamma}(O) = (0,0,0,0,0.1,0,0.35,0,0.05)\]

The care dimension has diminished (\(0.5 \rightarrow 0.35\)), a privacy
residue persists (\(O^{5} = 0.1\)), and an epistemic component has
appeared (\(O^{9} = 0.05\) --- Alice now evaluates the friendship with
greater caution). The loop has left a mark.

This is the geometric content of the moral intuition that
\textbf{experiences leave permanent traces}. A relationship that
survives a betrayal is not the same as one that never faced one. A
society that has gone through a war is not the same as one that has not,
even if it returns to ``peace.'' The holonomy measures the
\emph{irreducible residue} of moral experience.

\subsubsection{The Holonomy Group}\label{the-holonomy-group}

The set of all holonomy transformations at a point \(p\) forms a group
--- the \emph{holonomy group} \(\text{Hol}(\nabla,p)\):

\[\text{Hol}(\nabla,p) = \{ H_{\gamma}:\gamma\text{ is a loop at }p\}\]

The holonomy group captures the full range of transformations that moral
experience can induce. If \(\text{Hol}(\nabla,p)\) is trivial (only the
identity), moral space is flat at \(p\) --- all loops preserve
obligations exactly. If \(\text{Hol}(\nabla,p)\) is large, the curvature
is rich --- many different loops produce many different rotations.

\textbf{Connection to the} \(D_{4}\) \textbf{structure.} In stratified
moral space (Chapter 8), the holonomy around a loop that crosses
Hohfeldian boundaries is constrained by the \(D_{4}\) group. A loop that
enters the obligation stratum and returns to the liberty stratum may
induce a transformation from \(D_{4}\) --- a rotation or reflection of
the Hohfeldian square. The holonomy group of the stratified moral
manifold includes \(D_{4}\) as a discrete subgroup, contributing to what
might be called the \emph{non-abelian} character of moral dynamics.

\textbf{Selective path dependence.} {[}Empirical result
(preliminary).{]} Experimental testing of path dependence across 8 moral
scenarios (N = 640 evaluations) reveals that the non-abelian structure
is selective: path dependence was detected in exactly 2 of 8 scenarios
(combined χ² = 72.14, p \textless{} 10⁻⁸), and both path-dependent
scenarios involved considerations that point to different Hohfeldian
bond types.

In the journalist scenario (truth-telling vs. source protection), truth
considerations point toward Obligation while protection considerations
point toward Claim. Presenting truth first yields 77.5\% O, 22.5\% C;
presenting protection first yields 5\% O, 95\% C. The dominant
classification flips from O to C. In the teacher scenario (academic
integrity vs. compassion), integrity points toward Liberty while
compassion points toward Obligation. The order of presentation again
flips the dominant verdict between L and O.

The remaining 6 scenarios---where both considerations point to the same
bond type---showed no path dependence (p \textgreater{} 0.2 for all).
This yields a precise structural prediction: within-type moral
operations commute; cross-type operations do not. The non-abelian
character of the D₄ gauge group (§12.3) is not a global property of all
moral reasoning but a local property of boundary-crossing deliberation.

\subsection{10.6 The Curvature of Moral
Space}\label{the-curvature-of-moral-space}

\subsubsection{The Riemann Curvature
Tensor}\label{the-riemann-curvature-tensor-1}

The curvature of the moral manifold is measured by the Riemann curvature
tensor (Chapter 4, Definition 4.5):

\[R_{\nu\alpha\beta}^{\mu} = \frac{\partial\Gamma_{\nu\beta}^{\mu}}{\partial x^{\alpha}} - \frac{\partial\Gamma_{\nu\alpha}^{\mu}}{\partial x^{\beta}} + \Gamma_{\sigma\alpha}^{\mu}\Gamma_{\nu\beta}^{\sigma} - \Gamma_{\sigma\beta}^{\mu}\Gamma_{\nu\alpha}^{\sigma}\]

This (1,3)-tensor measures the path-dependence of parallel transport:
transporting an obligation vector \(O\) around an infinitesimal
parallelogram spanned by directions \(X\) and \(Y\) produces a change:

\[\delta O^{\mu} = R_{\nu\alpha\beta}^{\mu}O^{\nu}X^{\alpha}Y^{\beta} \cdot (\text{area of parallelogram})\]

If \(R = 0\), the space is flat and transport is path-independent. If
\(R \neq 0\), the result of transport depends on the path taken --- and
the Riemann tensor tells us exactly how much and in what direction.

\subsubsection{Sources of Moral
Curvature}\label{sources-of-moral-curvature}

What causes moral space to be curved? Since the Levi-Civita connection
is determined by the metric, the curvature is determined by how the
metric varies across \(\mathcal{M}\). Moral curvature arises wherever
the trade-off structure changes from point to point --- which is to say,
generically.

\textbf{Source 1: Context-dependent weighting.} The empirical finding
that moral dimension weights shift by context (Chapter 5, §5.3) --- care
dominates in family contexts, fairness in institutional contexts ---
directly implies a varying metric, and hence nonzero curvature. The
Christoffel symbols computed from these variations are nonzero, and the
Riemann tensor formed from them is generically nonzero.

\textbf{Source 2: Commitment structure.} Promises, contracts, and
relationships create local deformations of the moral metric. Near a
promise, the duty dimension (\(x^{2}\)) acquires extra weight --- the
metric \(g_{22}\) increases locally. This local deformation curves moral
space, just as mass curves spacetime in general relativity. The analogy
is structural: in both cases, curvature arises from the inhomogeneous
distribution of a ``source'' (mass-energy in physics; moral commitment
in ethics).

\textbf{Source 3: Threshold proximity. Near a stratum boundary (Chapter
8), the metric changes rapidly --- it may even become singular. This
rapid variation produces large curvature. Close to the consent
threshold, for instance, small changes in the situation produce large
changes in the metric (the weight of the autonomy dimension spikes as
consent becomes marginal). The curvature tensor diverges near the
boundary --- a moral analogue of the gravitational tidal forces near a
black hole's singularity, where curvature invariants diverge.}

\textbf{Source 4: Interaction effects.} When multiple agents are
involved, their interacting claims create curvature through the
off-diagonal metric components. Alice's obligation to Bob curves
differently depending on whether Carol is also present. The three-body
moral problem --- like the three-body gravitational problem --- is
generically not integrable, producing chaotic-seeming dynamics from
simple local rules.

\subsubsection{The Ricci Tensor and Scalar
Curvature}\label{the-ricci-tensor-and-scalar-curvature}

Contracting the Riemann tensor yields summary measures of curvature:

\textbf{The Ricci tensor} \(R_{\mu\nu} = R_{\mu\alpha\nu}^{\alpha}\) is
a symmetric (0,2)-tensor that measures, at each point, how much a small
ball of moral situations converges or diverges as it is
parallel-transported. Positive Ricci curvature along a direction means
that obligations converge: nearby moral paths tend toward agreement.
Negative Ricci curvature means they diverge: nearby paths tend toward
disagreement.

\textbf{The scalar curvature} \(R = g^{\mu\nu}R_{\mu\nu}\) is a single
number at each point that summarizes the overall curvature. Positive
scalar curvature means moral space is ``compressed'' --- there is less
room for variation than in flat space. Negative scalar curvature means
moral space is ``expanded'' --- more room for variation, more directions
of disagreement.

\subsubsection{Curvature and the Moral Metric: An
Equation?}\label{curvature-and-the-moral-metric-an-equation}

In general relativity, Einstein's field equations relate the curvature
of spacetime to the distribution of matter and energy:

\[R_{\mu\nu} - \frac{1}{2}Rg_{\mu\nu} = 8\pi G\, T_{\mu\nu}\]

Is there an analogous equation for moral space --- a relation between
the curvature of the moral manifold and the ``sources'' of moral
structure (commitments, relationships, stakes)?

\textbf{{[}Speculation/Extension.{]}} This is among the deepest open
questions of the framework. The analogy is suggestive: just as matter
tells spacetime how to curve, moral commitments tell moral space how to
curve. And just as spacetime curvature tells matter how to move, moral
curvature tells obligations how to evolve. Whether this analogy can be
made into an equation --- a ``moral field equation'' relating the Ricci
tensor to a ``moral stress-energy tensor'' --- is a research program,
not a result.

What we can say is the direction of the analogy:

{\def\LTcaptype{none} % do not increment counter
\begin{longtable}[]{@{}
  >{\raggedright\arraybackslash}p{(\linewidth - 2\tabcolsep) * \real{0.3472}}
  >{\raggedright\arraybackslash}p{(\linewidth - 2\tabcolsep) * \real{0.6354}}@{}}
\toprule\noalign{}
\begin{minipage}[b]{\linewidth}\raggedright
\textbf{Physics}
\end{minipage} & \begin{minipage}[b]{\linewidth}\raggedright
\textbf{Moral Geometry}
\end{minipage} \\
\midrule\noalign{}
\endhead
\bottomrule\noalign{}
\endlastfoot
Mass-energy distribution \(T_{\mu\nu}\) & Distribution of commitments,
relationships, stakes \\
Spacetime metric \(g_{\mu\nu}\) & Moral metric (trade-off structure) \\
Einstein tensor \(G_{\mu\nu}\) & Moral curvature tensor \\
Geodesic equation (free particle) & Moral trajectory of an agent with no
external obligations \\
Tidal forces (geodesic deviation) & Divergence of moral assessments for
nearby initial conditions \\
\end{longtable}
}

We do not claim the analogy is exact. We claim it is structural --- and
that the mathematical framework is ready if the analogy proves fruitful.

\subsection{10.7 Geodesics: Paths of Least Moral
Resistance}\label{geodesics-paths-of-least-moral-resistance}

\subsubsection{Definition}\label{definition-4}

A \emph{geodesic} is a curve \(\gamma(t)\) that parallel-transports its
own tangent vector:

\[\nabla_{\dot{\gamma}}\dot{\gamma} = 0\]

In coordinates:

\[\frac{d^{2}\gamma^{\mu}}{dt^{2}} + \Gamma_{\nu\rho}^{\mu}\frac{d\gamma^{\nu}}{dt}\frac{d\gamma^{\rho}}{dt} = 0\]

Geodesics are the ``straightest possible'' paths in a curved space ---
the paths of least resistance, the natural trajectories that require no
external force.

\subsubsection{Moral Interpretation}\label{moral-interpretation-2}

A moral geodesic is a trajectory through moral space that \emph{carries
no external moral force} --- no new obligations, no new commitments, no
external constraints. It is the path an agent would follow purely under
the influence of the moral landscape's own curvature: the inertial
trajectory of moral evaluation.

\textbf{What does this mean concretely?} An agent on a moral geodesic:

\begin{itemize}
\item
  Maintains their obligations faithfully (parallel-transporting them
  along the curve)
\item
  Makes no new commitments (no external force)
\item
  Follows the natural ``flow'' of the moral landscape (the curvature
  guides the path)
\end{itemize}

This is not the path of a \emph{virtuous} agent (who may resist the
curvature, swim against the moral current) or a \emph{vicious} agent
(who may violate parallel transport by arbitrarily discarding
obligations). It is the path of a morally \emph{inertial} agent --- one
who faithfully carries existing obligations without adding or removing
any.

\textbf{Example.} A society that makes no new moral commitments ---
passes no new laws, enters no new treaties, undergoes no moral reform
--- follows a moral geodesic. The trajectory of its obligations is
determined by the existing curvature (the structure of existing
commitments and relationships). The geodesic may lead somewhere
unappealing --- moral drift, gradual erosion of standards, convergence
on a local minimum. This is the cost of moral inertia.

\subsubsection{Geodesic Deviation}\label{geodesic-deviation}

The \emph{geodesic deviation equation} describes how nearby geodesics
converge or diverge:

\[\frac{D^{2}\xi^{\mu}}{dt^{2}} = R_{\nu\alpha\beta}^{\mu}{\dot{\gamma}}^{\nu}{\dot{\gamma}}^{\alpha}\xi^{\beta}\]

where \(\xi\) is the separation vector between two nearby geodesics. The
Riemann tensor acts as a ``tidal force,'' pulling nearby moral
trajectories together (positive curvature) or pushing them apart
(negative curvature).

\textbf{Moral interpretation.} Two agents who start with nearly
identical moral commitments (nearby points in \(\mathcal{M}\)) but
follow different geodesics (evolving in slightly different directions)
will:

\begin{itemize}
\item
  \textbf{Converge} in regions of positive curvature: the moral
  landscape funnels them toward agreement. These are regions where the
  structure of commitments and relationships is such that diverse
  starting points lead to similar conclusions.
\item
  \textbf{Diverge} in regions of negative curvature: the moral landscape
  amplifies small initial differences. These are regions where moral
  disagreement grows with experience --- where agents who started
  similarly end up far apart.
\end{itemize}

The existence of regions of both positive and negative moral curvature
explains a familiar phenomenon: why some moral issues generate
convergence over time (societies tend to agree that slavery is wrong ---
a region of strong positive curvature) while others generate increasing
polarization (societies diverge on the ethics of economic redistribution
--- a region of negative curvature).

Chapter 11 shows that computing the exact optimal geodesic is
practically intractable at the resolution required for morally adequate
evaluation, and the associated discrete constrained-path problem
(MORAL-PATHFIND) is NP-hard, which is why agents use A* search with
pre-compiled heuristic functions --- the obligation vectors of §6.2 ---
rather than exhaustive computation.

\subsection{10.8 Gradient Flows and Moral
Improvement}\label{gradient-flows-and-moral-improvement}

\subsubsection{The Gradient Field}\label{the-gradient-field}

If the satisfaction function \(S\mathcal{:M \rightarrow}\mathbb{R}\)
assigns a scalar value to each point, its gradient (raised by the
metric) defines a vector field:

\[(\text{grad}\, S)^{\mu} = g^{\mu\nu}\frac{\partial S}{\partial x^{\nu}}\]

This vector field points, at each point, in the direction of
\emph{steepest increase} of \(S\) --- the direction of maximal moral
improvement.

\subsubsection{The Gradient Flow}\label{the-gradient-flow}

The \emph{gradient flow} of \(S\) is the family of curves \(\gamma(t)\)
satisfying:

\[\frac{d\gamma^{\mu}}{dt} = g^{\mu\nu}\frac{\partial S}{\partial x^{\nu}}|_{\gamma(t)}\]

An agent following the gradient flow always moves in the direction that
increases satisfaction most rapidly. This is the moral analogue of
gradient descent (or ascent) in optimization.

\subsubsection{Virtues and Vices of Gradient
Following}\label{virtues-and-vices-of-gradient-following}

\textbf{Virtues.} The gradient flow is locally optimal: at each point,
it makes the maximum possible improvement. An agent following the
gradient is never lazy (choosing a suboptimal direction) or perverse
(moving away from improvement).

\textbf{Vices.} The gradient flow is \emph{myopic} --- it sees only the
local landscape, not the global structure. This leads to familiar
pathologies:

\begin{enumerate}
\def\labelenumi{\arabic{enumi}.}
\item
  \textbf{Local maxima.} The gradient may lead to a local maximum that
  is not a global maximum. The agent is ``stuck'' in a morally
  acceptable but non-optimal configuration, unable to improve locally
  but far from the best achievable state. This is the geometric
  representation of moral complacency.
\item
  \textbf{Saddle points.} The gradient may vanish at a saddle point ---
  a point that is a maximum in some directions and a minimum in others.
  The agent is paralyzed: no direction offers clear improvement, despite
  the existence of better states nearby. This is the geometric
  representation of a moral impasse where deliberation stalls.
\item
  \textbf{Sensitivity to the metric.} The gradient direction depends on
  the metric: \(\text{grad}\, S = g^{- 1}dS\). Different metrics
  (different ethical theories) yield different gradient directions for
  the \emph{same} satisfaction function. Two agents in the same
  situation, following their respective gradient flows, may move in
  different directions --- not because they disagree about the value of
  \(S\), but because they disagree about the metric that defines
  ``direction.''
\end{enumerate}

\subsubsection{Virtue as Non-Gradient
Agency}\label{virtue-as-non-gradient-agency}

A \emph{virtuous} agent does not simply follow the gradient. Virtue
involves:

\begin{itemize}
\item
  \textbf{Foresight}: seeing beyond the local landscape to avoid local
  maxima and saddle points. This requires information beyond
  \(\nabla S\) at the current point --- information about the global
  structure of \(\mathcal{M}\).
\item
  \textbf{Commitment}: maintaining a direction of moral movement even
  when the gradient temporarily points elsewhere. This is the geometric
  content of integrity --- holding a course when circumstances create
  countervailing pressures.
\item
  \textbf{Judgment}: knowing when to follow the gradient (in clear
  cases), when to resist it (near local maxima), and when to jump across
  strata (at phase transitions).
\end{itemize}

Virtue, in this framework, is not a specific location in moral space but
a \emph{quality of trajectory} --- a way of moving through the landscape
that balances local responsiveness with global awareness.

Chapter 11 identifies gradient flow as a degenerate case of A* search:
it is greedy best-first search without backward cost accounting. The
pathologies identified above --- local maxima, saddle points, metric
singularities --- are precisely the pathologies of search without a
backward-cost term g(n).

\subsection{10.9 Temporal Structure and Moral
History}\label{temporal-structure-and-moral-history}

\subsubsection{Time and the Moral
Manifold}\label{time-and-the-moral-manifold}

The dynamics discussed so far concern movement through moral space ---
changes in the \emph{situation}, parameterized abstractly. But moral
life also has a literal temporal dimension: obligations are incurred at
one time and come due at another; moral evaluation of past actions
differs from prospective evaluation of future ones; intergenerational
obligations connect the present to the far future.

How does time enter the geometric framework?

\subsubsection{Temporal Moral Tensors}\label{temporal-moral-tensors}

Time can enter in several ways:

\textbf{Time as a coordinate on} \(\mathcal{M}\)\textbf{.} If we include
a temporal coordinate \(t\) among the dimensions of the moral manifold,
the metric acquires temporal components \(g_{t\mu}\) that encode how
moral dimensions relate to time. A positive \(g_{t,1}\) means that
welfare considerations become more important over time (e.g.,
consequences compound). A negative \(g_{t,1}\) means they become less
important (temporal discounting).

\textbf{Temporal discounting as metric curvature.} The practice of
giving less weight to future welfare is, in geometric terms, a statement
about the metric: the distance between ``present welfare loss'' and
``identical future welfare loss'' is not zero. The discounting function
\(\delta(t) = e^{- rt}\) (exponential discounting) corresponds to a
specific temporal metric with constant negative curvature along the time
direction: future obligations are exponentially ``farther away'' than
present ones.

This is revealing, because it shows that temporal discounting is not
merely a preference --- it is a \emph{geometric commitment} about the
structure of moral space along the temporal axis. Different discount
rates correspond to different metrics, with different curvatures. The
question ``What is the correct discount rate?'' is a special case of
``What is the correct moral metric?'' --- the governance question of
Chapter 9.

\textbf{Obligations with temporal indices.} Some obligations are
\emph{intrinsically} temporal --- they are defined not at a single point
but between two points:

\begin{itemize}
\item
  \(O_{t_{0} \rightarrow t_{1}}^{\mu}\): an obligation incurred at
  \(t_{0}\) and due at \(t_{1}\)
\item
  \(R_{t}^{ab}\): the responsibility of agent \(a\) to agent \(b\) at
  time \(t\)
\item
  \(V_{\lbrack t_{0},t_{1}\rbrack}^{\mu}\): the value of a
  \emph{trajectory} from \(t_{0}\) to \(t_{1}\), not just of the
  endpoints
\end{itemize}

The last object is important: a trajectory that passes through
wrongdoing en route to a good outcome is not equivalent to a trajectory
that reaches the same outcome through permissible means. The path
matters, not just the destination. This is the temporal content of the
deontological intuition that \emph{how you get there matters}, and it is
formalized by the curvature and holonomy of the temporal moral manifold.

\subsubsection{Intergenerational
Obligations}\label{intergenerational-obligations}

Obligations to future generations have a distinctive geometric
structure. The obligation is incurred now; the beneficiary does not yet
exist; the due-date is indefinite; the circumstances under which the
obligation will be evaluated are uncertain.

In the geometric framework, intergenerational obligations are vectors
that must be \emph{parallel-transported} across a vast stretch of moral
space --- from the present context to future contexts that may differ
radically in their moral landscape (different technologies, different
ecological conditions, different social structures).

The question ``What do we owe future generations?'' becomes: ``What is
the parallel transport of our obligations from the present context to
future contexts?'' The answer depends on the curvature of the temporal
moral manifold. If the curvature is large (the moral landscape changes
rapidly over time), the parallel-transported obligation may bear little
resemblance to the original. If the curvature is small (the moral
landscape is temporally stable), the obligation transports faithfully.

The empirical finding from the Dear Abby corpus (Chapter 17) that the
\emph{structural} features of the moral metric are stable over three
decades, even as specific advice varies, is evidence for low curvature
along the temporal axis --- at least over decadal timescales. Whether
this stability extends to centennial or millennial timescales is one of
the great unknowns of moral dynamics.

\subsection{10.10 Worked Example: The Evolving
Obligation}\label{worked-example-the-evolving-obligation}

Let us trace a single obligation through a complete moral circuit,
computing the dynamics explicitly.

\subsubsection{Setup}\label{setup-2}

At \(t_{0}\), a parent promises their child: ``I will be at your recital
on Saturday.''

The initial obligation vector (in the six active dimensions of this
context) is:

\[O(t_{0}) = (0.1,0.8,0.1,0.2,0,0.6)\]

Components: welfare (0.1 --- the child benefits emotionally), duty (0.8
--- an explicit promise), fairness (0.1), autonomy (0.2 --- the parent's
freedom is reduced), privacy (0), care (0.6 --- the relational bond).

\subsubsection{\texorpdfstring{\(t_{1}\): Work
Emergency}{t\_\{1\}: Work Emergency}}\label{t_1-work-emergency}

The parent learns of a work emergency that conflicts with Saturday. The
situation moves in the direction of increased cost to autonomy and
welfare trade-offs. We parallel-transport \(O\) along this path.

The connection coefficients (computed from the metric's
context-dependence) give:

\[\Gamma_{4,2}^{1} = 0.3\quad\text{(duty generates welfare concern when autonomy is stressed)}\]

\[\Gamma_{4,2}^{2} = - 0.15\quad\text{(duty weakens slightly under autonomy stress)}\]

After transport:

\[O(t_{1}) \approx (0.34,0.68,0.1,0.2,0,0.6)\]

The duty component has diminished (from 0.8 to 0.68) and the welfare
component has increased (from 0.1 to 0.34). The obligation has rotated:
the promise is still present but the stakes have changed.

\subsubsection{\texorpdfstring{\(t_{2}\): Recital
Postponed}{t\_\{2\}: Recital Postponed}}\label{t_2-recital-postponed}

The child's recital is postponed due to the emergency. The obligation is
carried to the new time frame. Under parallel transport:

\[O(t_{2}) \approx (0.30,0.72,0.1,0.25,0,0.55)\]

The duty component has \emph{partially recovered} (the promise, now
applicable to the new date, regains force), the care component has
slightly diminished (the disruption has taken a small toll on the
relational bond), and the autonomy cost has increased (the parent now
has \emph{two} dates to manage).

\subsubsection{\texorpdfstring{\(t_{3}\): The Parent
Attends}{t\_\{3\}: The Parent Attends}}\label{t_3-the-parent-attends}

The parent attends the rescheduled recital. The duty obligation is
satisfied: \(O^{2} \rightarrow 0\). The care obligation is satisfied:
\(O^{7}\) decreases. But the satisfaction is not a return to baseline:

\[O(t_{3}) \approx (0.05,0,0.05,0.1,0,0.25)\]

A \emph{residual} care obligation persists (\(O^{7} = 0.25\) instead of
0) --- the parent may owe an explanation or acknowledgment of the
disruption. A small epistemic component has appeared (\(O^{9} = 0.05\)
in a full nine-dimensional treatment) --- the child now has information
about the parent's reliability that may affect future trust.

\subsubsection{Holonomy Computation}\label{holonomy-computation}

The initial obligation was \(O(t_{0}) = (0.1,0.8,0.1,0.2,0,0.6)\). If
the obligation had been perfectly satisfied at \(t_{0}\) (no emergency,
no postponement), the final state would be \(O = 0\) (all obligations
discharged). Instead, the final state is
\(O(t_{3}) = (0.05,0,0.05,0.1,0,0.25)\).

The holonomy --- the residual obligation after the circuit --- is:

\[H_{\gamma}(O(t_{0})) - 0 = (0.05,0,0.05,0.1,0,0.25)\]

This is the \emph{moral residue} of the experience (Chapter 15). It is
not large --- the obligation was substantially fulfilled --- but it is
nonzero. The parent owes something beyond the fulfilled promise: an
acknowledgment, a compensation, a closer attention to future
commitments. The holonomy is the geometric measure of this residue.

\subsection{10.11 Curvature and Moral
Disagreement}\label{curvature-and-moral-disagreement}

\subsubsection{Why Moral Reasoning Is
Path-Dependent}\label{why-moral-reasoning-is-path-dependent}

The curvature of moral space has a direct consequence for moral
disagreement: \textbf{two agents who start from the same moral
commitments but traverse different paths of experience may arrive at
genuinely different moral evaluations of the same situation.}

This is not a failure of moral reasoning. It is a structural feature of
a curved moral space. In flat space (zero curvature), the order of
experiences does not matter --- all paths between the same two points
yield the same parallel transport. But in curved space, the path
matters. An agent who encounters a betrayal before encountering
forgiveness may carry a different moral orientation than one who
encounters forgiveness first.

\subsubsection{Implications}\label{implications}

\begin{enumerate}
\def\labelenumi{\arabic{enumi}.}
\item
  \textbf{Moral disagreement may be irreducible to error.} If the
  curvature is nonzero, two agents who have correctly
  parallel-transported their obligations along different paths will
  arrive at different moral evaluations --- without either having made a
  mistake. The disagreement is \emph{geometrically generated}, not
  epistemically deficient.
\item
  \textbf{Moral dialogue requires shared paths.} To compare moral
  evaluations, agents must share some common path of experience --- or
  at least some common region of moral space where their
  parallel-transported obligations can be compared. This is the
  geometric content of the intuition that moral understanding requires
  empathy: you must traverse some of the same moral landscape to compare
  notes.
\item
  \textbf{Moral progress is possible but path-dependent.} A society can
  progress morally --- moving toward a better region of the moral
  manifold --- but the trajectory matters. A society that abolishes
  slavery before establishing civil rights traverses a different moral
  path than one that does the reverse, and the holonomy may differ: the
  resulting moral orientation may be subtly different, even if the
  endpoint appears the same.
\end{enumerate}

\subsection{10.12 Summary}\label{summary-3}

{\def\LTcaptype{none} % do not increment counter
\begin{longtable}[]{@{}
  >{\raggedright\arraybackslash}p{(\linewidth - 4\tabcolsep) * \real{0.1919}}
  >{\raggedright\arraybackslash}p{(\linewidth - 4\tabcolsep) * \real{0.2333}}
  >{\raggedright\arraybackslash}p{(\linewidth - 4\tabcolsep) * \real{0.5749}}@{}}
\toprule\noalign{}
\begin{minipage}[b]{\linewidth}\raggedright
\textbf{Concept}
\end{minipage} & \begin{minipage}[b]{\linewidth}\raggedright
\textbf{Mathematical Object}
\end{minipage} & \begin{minipage}[b]{\linewidth}\raggedright
\textbf{Moral Content}
\end{minipage} \\
\midrule\noalign{}
\endhead
\bottomrule\noalign{}
\endlastfoot
Connection \(\nabla\) & Christoffel symbols \(\Gamma_{\nu\rho}^{\mu}\) &
How moral concepts translate across contexts \\
Parallel transport & ∇\_\{γ̇\} O = 0 & Faithful maintenance of obligation
across changing circumstances \\
Holonomy & T\_\{βα\} ∘ T\_\{αβ\} ≠ id & Moral residue of a completed
circuit; experience leaves marks \\
Riemann curvature & R\^{}μ\_\{ναβ\} & Path-dependence of moral
transport; source of irreducible disagreement \\
Ricci curvature & R\_\{μν\} & Convergence (positive) or divergence
(negative) of nearby moral paths \\
Geodesic & ∇\_\{γ̇\} γ̇ = 0 & Moral inertia --- trajectory under no
external moral force \\
Geodesic deviation & D²ξ/dt² = R · ξ & Amplification or damping of moral
differences \\
Gradient flow & γ̇ = grad S & Direction of maximal local moral
improvement \\
Temporal metric & g\_\{tμ\} & Discounting as curvature;
intergenerational obligation as transport \\
\end{longtable}
}

The moral manifold is not flat. It is curved by the inhomogeneous
distribution of commitments, relationships, and stakes. This curvature
has consequences: moral evaluation is path-dependent, experience leaves
holonomic residues, nearby agents can diverge in their moral
orientations, and moral inertia (geodesic motion) may lead to local
optima rather than global ones.

The covariant derivative, parallel transport, and curvature tensor
provide the vocabulary for analyzing these phenomena with mathematical
precision. They do not resolve moral disagreements --- but they explain,
for the first time, why some disagreements may be \emph{structurally
irresolvable}: not because one party is wrong, but because the curvature
of moral space maps different paths of experience onto genuinely
different moral orientations.

Chapter 11 reveals the computational structure hidden within this
geometric apparatus: moral reasoning is A* pathfinding, with obligation
vectors as pre-compiled heuristic functions. Chapter 12 then exploits
the full apparatus further, showing that the symmetry of re-description
invariance --- the Bond Invariance Principle --- implies, via Noether's
theorem, a conservation law for harm. The dynamics of this chapter
provide the stage; the conservation law of the next provides the
constraint.

\subsection{Technical Appendix}\label{technical-appendix-4}

\textbf{Proposition 10.1 (Parallel Transport Preserves Satisfaction).}
\emph{If the connection is the Levi-Civita connection of the moral
metric} \(g\)\emph{, and both} \(O\) \emph{and} \(I\) \emph{are
parallel-transported along a curve} \(\gamma\)\emph{, then the
satisfaction} \(S = I_{\mu}O^{\mu}\) \emph{is constant along}
\(\gamma\)\emph{.}

\emph{\textbf{Proof}.} Compute:

Proof detail. The key step is metric compatibility of the Levi-Civita
connection: ∇g = 0, which implies that for any vectors X, Y and any
direction of transport v, we have v(g(X,Y)) = g(∇\_v X, Y) + g(X, ∇\_v
Y). Setting X = O, Y = I (via the metric identification), and using
∇\_\{γ'\} O = 0 and ∇\_\{γ'\} I = 0 (parallel transport conditions), we
obtain dS/dt = d/dt{[}g(O, I){]} = g(0, I) + g(O, 0) = 0.

\[\frac{d}{dt}(I_{\mu}O^{\mu}) = (\nabla_{\dot{\gamma}}I)_{\mu}O^{\mu} + I_{\mu}(\nabla_{\dot{\gamma}}O)^{\mu} = 0 + 0 = 0\]

since both \(\nabla_{\dot{\gamma}}O = 0\) and
\(\nabla_{\dot{\gamma}}I = 0\) by the parallel transport condition.
\(▫\)

\textbf{Corollary.} \emph{If an obligation and an interest are both
faithfully maintained (parallel-transported) across changing
circumstances, the satisfaction they produce is conserved.} This is a
geometric version of the principle that a faithfully maintained
commitment continues to serve the interests it was designed to serve,
regardless of how the context changes --- provided the context change is
navigated by the same connection.

\textbf{Proposition 10.2 (Holonomy Detects Curvature).} \emph{For an
infinitesimal loop enclosing area} \(\delta A\) \emph{in the plane
spanned by} \(e_{\alpha}\) \emph{and} \(e_{\beta}\)\emph{, the holonomy
acts on a vector} \(O\) \emph{as:}

\[H_{\gamma}(O)^{\mu} - O^{\mu} = R_{\nu\alpha\beta}^{\mu}O^{\nu}\delta A^{\alpha\beta}\]

\emph{The Riemann tensor is the infinitesimal holonomy.}

\emph{\textbf{Proof}. Consider the infinitesimal parallelogram at point
p spanned by ε e\_α and ε e\_β. Transport Oᵘ parallel along the four
sides. At each leg, the parallel transport equation ∇\_\{e\_α\} Oᵘ =
−Γᵘ\_\{να\} Oᵛ introduces a correction of order ε. After the first two
legs (to order ε²): Oᵘ(p + ε e\_α + ε e\_β) = Oᵘ − εΓᵘ\_\{να\}Oᵛ −
εΓᵘ\_\{νβ\}Oᵛ + ε²Γᵘ\_\{σβ\}Γᵠ\_\{να\}Oᵛ + ε²∂\_αΓᵘ\_\{νβ\}Oᵛ + O(ε³).
After the return legs, the net change upon completing the loop is: δOᵘ =
ε²(∂\_αΓᵘ\_\{νβ\} − ∂\_βΓᵘ\_\{να\} + Γᵘ\_\{σα\}Γᵠ\_\{νβ\} −
Γᵘ\_\{σβ\}Γᵠ\_\{να\})Oᵛ = Rᵘ\_\{ναβ\} Oᵛ ε². Identifying δA\^{}\{αβ\} =
ε²: H\_γ(O)ᵘ − Oᵘ = Rᵘ\_\{ναβ\} Oᵛ δA\^{}\{αβ\}.
□}\(\epsilon e_{\alpha}\epsilon e_{\beta}\delta O^{\mu} = R_{\nu\alpha\beta}^{\mu}O^{\nu}\epsilon^{2} + O(\epsilon^{3})▫\)

\textbf{Proposition 10.3 (Geodesic Deviation Equation).} \emph{Two
nearby geodesics} \(\gamma\) \emph{and} \(\gamma'\) \emph{with
separation vector} \(\xi\) \emph{satisfy:}

\[\frac{D^{2}\xi^{\mu}}{dt^{2}} = R_{\nu\alpha\beta}^{\mu}{\dot{\gamma}}^{\nu}{\dot{\gamma}}^{\alpha}\xi^{\beta}\]

\emph{where} \(D/dt\) \emph{is the covariant derivative along}
\(\gamma\)\emph{. The Riemann tensor acts as a ``tidal force'' on the
separation.}

\emph{\textbf{Proof}. Let γ\_s(t) be a one-parameter family of geodesics
with γ₀ = γ, velocity uᵘ = γ̇ᵘ, and separation vector ξᵘ =
∂γ\_sᵘ/∂s\textbar\_\{s=0\}. Since {[}u,ξ{]} = 0 (they are coordinate
vector fields of the (t,s) parameterization) and the Levi-Civita
connection is torsion-free, ∇\_u ξ = ∇\_ξ u. Compute: D²ξᵘ/dt² =
∇\_u(∇\_u ξᵘ) = ∇\_u(∇\_ξ uᵘ). By the definition of the Riemann tensor:
∇\_u∇\_ξ uᵘ − ∇\_ξ∇\_u uᵘ = Rᵘ\_\{ναβ\} uᵛξᵝ uᵅ + ∇\_\{{[}u,ξ{]}\}uᵘ.
Since ∇\_u u = 0 (geodesic equation) and {[}u,ξ{]} = 0: D²ξᵘ/dt² =
Rᵘ\_\{ναβ\} γ̇ᵛγ̇ᵅξᵝ. This is the Jacobi equation. □}\(\gamma\gamma'▫\)

❖

\emph{Moral space is not flat. It is curved by the commitments we make,
the relationships we inhabit, and the stakes we face.}

\emph{This curvature has consequences. Obligations carried through moral
space are rotated by the landscape. Circuits through experience leave
holonomic residues. Nearby agents, following different paths, may arrive
at genuinely different moral orientations --- not because one is wrong,
but because the geometry of moral space is nontrivial.}

\emph{The mathematics of curvature does not tell us what to do. It tells
us why moral reasoning --- when done honestly, by well-intentioned
agents --- can still produce irreducible disagreement.}

\emph{The curvature is the measure of this irreducibility.}

\emph{And that is not a defect of the moral landscape. It is its depth.}

\section{Chapter 11: Moral Reasoning as Optimal
Search}\label{chapter-11-moral-reasoning-as-optimal-search}

\begin{quote}
\textbf{\textsc{RUNNING EXAMPLE --- Priya's Model}}

\emph{Priya's decision tree is a pathfinding problem. Start: she knows
the algorithm is biased. Goal: a morally defensible outcome. Available
moves: fix quietly, escalate internally, report to the FDA, go to the
press, resign, do nothing. Blocked nodes: anything violating her NDA
carries high legal cost---a deontological constraint with heuristic h ≈
∞. The `fix quietly' path looks optimal by A*: low cost, no career risk.
But her obligation vector points toward escalation. The heuristic
gradient says the goal is closer via that harder path. Her reluctance is
consequentialist search---minimizing path cost. Her instinct is the
heuristic saying the cheaper path does not reach the right goal. Virtue
ethics asks: what kind of engineer do I want to be? That is heuristic
calibration.}
\end{quote}

\subsection{11.1 The Computational
Problem}\label{the-computational-problem}

The preceding chapter established the geometry of moral dynamics: the
covariant derivative, parallel transport, holonomy, geodesics, and
curvature. These are the structural features of the moral landscape. But
a landscape, however carefully mapped, does not by itself explain how
agents navigate it. Geodesics are the paths of least resistance --- the
optimal trajectories through moral space. But how does an agent find a
geodesic? How does a decision-maker, confronted with a morally complex
situation, determine which path through the space of possible actions
minimizes behavioral friction?

This chapter answers that question by establishing a formal equivalence
between moral reasoning and optimal search. The claim is precise: moral
reasoning is A* pathfinding on the moral manifold, where the accumulated
behavioral friction serves as the backward cost, and moral obligations
serve as pre-compiled heuristic functions that estimate the cost
remaining to reach equilibrium. The equation f(n) = g(n) + h(n) ---
familiar to every AI researcher from Hart, Nilsson, and Raphael's
foundational 1968 paper --- is the fundamental equation of moral
reasoning.

The shift from geometry to computation is not merely a change of
vocabulary. It resolves a deep puzzle: why does moral reasoning feel
like search? Why do we deliberate, weigh alternatives, backtrack, and
sometimes get stuck? The answer, as this chapter proves, is that the
exact computation of the optimal moral geodesic is intractable --- the
curse of dimensionality on a nine-dimensional manifold makes brute-force
computation impossible for any resource-bounded agent. Moral obligations
exist because they must: they are pre-compiled heuristic approximations
to an intractable optimization problem, and without them, real-time
moral judgment would be computationally impossible.\\
\strut \\
{[}Formal complexity result.{]} Define MORAL-PATHFIND as the decision
problem: given a weighted simplicial complex C with n vertices,
d-dimensional attribute vectors with Mahalanobis edge weights, and k ≥ 1
attribute-boundary constraints (edges that cross a constraint boundary
incur penalty βₖ), does there exist a path from s to G with total cost ≤
B? MORAL-PATHFIND is NP-hard by reduction from the Weight-Constrained
Shortest Path problem (Garey \& Johnson, 1979; Handler \& Zang, 1980),
which is NP-hard for k ≥ 1 side constraints. The reduction maps each
side constraint to a boundary penalty β\_k and each edge weight to the
corresponding Mahalanobis distance. Since MORAL-PATHFIND is NP-hard,
exact optimization on the moral manifold with boundary constraints is
computationally intractable in the worst case, and the
bounded-suboptimality guarantees of Weighted A* (Pohl, 1970) apply by
direct instantiation --- not by analogy --- to the moral search problem.

\subsection{11.2 A* Search on the Moral
Manifold}\label{a-search-on-the-moral-manifold}

We begin with the navigation problem. An agent occupies a state s₀ ∈ ℳ,
the current moral configuration of a multi-agent system. The agent must
choose a sequence of actions that moves the system toward a goal region
G\_eq --- a region of low behavioral friction where obligations and
interests are approximately aligned. The objective is to find the path γ
from s₀ to G\_eq that minimizes the total behavioral friction BF{[}γ{]}
(§6.4).

\textbf{Definition 11.1 (Moral Search Problem). {[}Modeling Axiom.{]}}
Given the moral manifold ℳ, a current state s₀ ∈ ℳ, and a goal region
G\_eq ⊂ ℳ of low behavioral friction, the moral search problem is to
find the path γ* = argmin\_γ BF{[}γ{]} from s₀ to G\_eq, subject to the
constraint that γ does not cross into forbidden strata (the constraint
surfaces of §8.3).

Remark on G\_eq. The goal region G\_eq is formally defined as the subset
of ℳ where the obligation vector field is approximately vanishing and
behavioral friction is below threshold: G\_eq = \{s ∈ ℳ :
\textbar\textbar Oᵐ(s)\textbar\textbar{} \textless{} ε, BF(s)
\textless{} δ\}. The parameters ε and δ are not fixed by the framework;
they are set by the moral community or governance structure. Different
choices of G\_eq yield different search problems, and disagreement about
what constitutes moral equilibrium maps to disagreement about these
parameters --- a feature the framework represents rather than resolves.

This is a constrained shortest-path problem on a stratified Riemannian
manifold. The classical algorithm for optimal pathfinding with heuristic
guidance is A* search.

\textbf{Definition 11.2 (A* Moral Search). {[}Modeling Axiom.{]}} The A*
evaluation of a candidate state n in the search tree is:

f(n) = g(n) + h(n)

where:

\begin{itemize}
\item
  g(n) = BF{[}γ\_\{s₀ → n\}{]} is the exact accumulated behavioral
  friction of the path from the start state s₀ to n --- the integrated
  cost of all actions taken so far. This is the backward cost, already
  defined as the behavioral friction functional (Definition 6.4,
  Equation 6.4).
\item
  h(n) is the heuristic function estimating the cheapest friction-path
  from n to the goal region G\_eq. This is the forward estimate --- the
  agent's best guess of how much additional cost remains.
\item
  f(n) is the total estimated cost of the cheapest path from s₀ to G\_eq
  passing through n.
\end{itemize}

A* expands the node with the lowest f(n) at each step. It is optimal ---
guaranteed to find the cheapest path --- when h(n) satisfies two
conditions: admissibility (h never overestimates the true remaining
cost) and consistency (h satisfies the triangle inequality). Both
conditions are defined precisely in §11.4.

Why A* rather than its alternatives? The choice is not arbitrary; it
reflects the structure of moral reasoning itself:

\begin{itemize}
\item
  Dijkstra's algorithm is A* with h = 0 --- no heuristic guidance. The
  search is blind, expanding nodes in order of backward cost alone. In
  moral terms, this is pure consequentialism without moral intuition:
  evaluate every possibility by its consequences, with no prior sense of
  which directions are promising. It is optimal but computationally
  explosive.
\item
  Greedy best-first search uses only h(n), ignoring backward cost g(n).
  This is pure rule-following without attention to accumulated
  consequences: the agent follows its moral intuitions regardless of the
  costs already incurred. It is fast but not optimal --- it can be
  misled by locally promising but globally costly paths.
\item
  Gradient descent (Chapter 10, §10.8) follows the steepest local
  improvement. It is myopic: it gets stuck at local optima, saddle
  points, and metric singularities. Chapter 10.8 identified these
  pathologies; the A* framework explains them: gradient descent is
  greedy search without backward cost accounting.
\item
  A* combines exact backward cost with heuristic forward estimation. It
  is both optimal (given admissibility) and efficient (the heuristic
  prunes the search space). In moral terms, it balances consequentialist
  accounting of past actions with deontological guidance toward the
  goal. This is the computational content of the familiar observation
  that good moral reasoning requires both principles and consequences.
\end{itemize}

The physical analogy is to variational principles. Fermat's principle
(light follows the path of least time), Hamilton's principle (systems
follow the path of least action), and the geodesic equation (free
particles follow paths of least curvature) all find optimal paths by
combining local cost with global structure. A* is their computational
implementation on a discrete approximation of the continuous manifold.

\subsection{11.3 Obligation Vectors as Heuristic
Functions}\label{obligation-vectors-as-heuristic-functions}

The central claim of this chapter is that moral obligations --- the
tangent vectors Oᵐ defined in Chapter 6 (§6.2) --- are the gradients of
heuristic functions in the A* framework. This is not a new definition of
obligation vectors; it is a reinterpretation of the existing formalism
that reveals its computational content.

\textbf{Proposition 11.1 (Obligation-Heuristic Correspondence).
{[}Modeling Axiom.{]}} The obligation vector Oᵐ at a point p ∈ ℳ is the
negative metric-raised gradient of the heuristic function h:

Oᵐ = −gᵐᵛ ∂ₙh

The obligation vector points in the direction of steepest descent of h
--- toward the estimated cheapest path to moral equilibrium.

This proposition reinterprets the obligation vectors of Chapter 6 within
the A* framework. When Chapter 6 defined Oᵐ as a tangent vector encoding
a ``directional moral imperative,'' the directionality was left as a
structural feature. The A* framework explains what the direction means:
it is the direction in which the heuristic function h decreases most
rapidly --- the direction that, according to the agent's pre-compiled
moral knowledge, most efficiently reduces the estimated cost of reaching
equilibrium. Note that this correspondence is not stipulated
arbitrarily. Obligation vectors were already defined in Chapter 6 as
tangent vectors pointing in the direction of moral improvement. The
heuristic gradient points in the direction of steepest estimated
cost-reduction. If behavioral friction is the cost measure --- as
established in §6.4 --- then these two directions necessarily coincide,
and the modeling axiom reduces to a consistency requirement between the
geometric and computational interpretations of the same underlying
structure.

Moral obligations such as ``do not murder,'' ``keep your promises,'' and
the Golden Rule are highly optimized heuristic functions h(n). They
allow an agent with limited computational power to instantly approximate
the optimal geodesic through the social space. When an agent experiences
a ``moral obligation,'' they are reading the output of an evolutionary
heuristic designed to steer the system away from stratum boundaries ---
away from catastrophic regime transitions that would spike behavioral
friction beyond recovery.

\subsubsection{The Murder Prohibition as h(n) =
∞}\label{the-murder-prohibition-as-hn}

Consider the prohibition of murder. In the A* framework, the heuristic
assigns h\_murder(n) = ∞ for any state n that involves killing. This is
not a modeling failure; it is a correct encoding of the stratification
from Chapter 8. Murder crosses a stratum boundary into a catastrophic
regime --- specifically, it triggers a Type III boundary (nullifier,
§8.3), where the moral evaluation structure changes so drastically that
the action is unconditionally prohibited regardless of the values on
other dimensions.

The true cost g*(n, G\_eq) of the cheapest path from a post-murder state
back to equilibrium includes: (i) grief processing across all connected
agents in the social network; (ii) institutional response ---
investigation, adjudication, incarceration; (iii) reputational damage
propagation through the network; (iv) deterrence expenditure to prevent
recurrence; (v) the irreversibility of death, which means that the
recovery path may not exist within the catastrophic stratum at all. In
the limit of network connectivity, these costs grow without bound,
making g*(n, G\_eq) = ∞ as well. The critical point is not merely that
costs are large, but that the original goal region --- G\_eq as defined
relative to the pre-murder configuration, which includes the victim
being alive and unharmed --- is genuinely unreachable. The
irreversibility of death means that no path from the post-murder state
returns to the original G\_eq. The social network may eventually reach a
different equilibrium, but the goal region that existed before the
boundary crossing has been destroyed. This is a claim about
path-specificity (the original destination no longer exists), not merely
about cost-magnitude (the journey is expensive).

Hence h\_murder(n) ≤ g*(n, G\_eq), and the heuristic is admissible. The
infinite value is not an approximation failure; it correctly reflects
that the stratum boundary is a point of no return --- the absorbing
stratum of §8.3.

\subsubsection{Promise-Keeping as a Finite
Heuristic}\label{promise-keeping-as-a-finite-heuristic}

Not all moral heuristics assign infinite cost. Promise-keeping provides
a more nuanced example. The heuristic h\_promise(n) estimates the
trust-rebuilding cost if the promise is broken: the behavioral friction
of repairing the damaged social bond, re-establishing reliability, and
compensating for the reliance costs incurred by the promisee.

The empirical evidence from §17.2 supports this interpretation directly.
The phrase ``you promised'' triggers a deontic phase transition with
94\% effectiveness --- a sharp, near-universal shift in moral evaluation
that crosses stratum boundaries. In the A* framework, this phrase
activates the promise-keeping heuristic, which reports a high (but
finite) forward cost for states in which the promise is broken. The 94\%
effectiveness is a measurement of the heuristic's sharpness: it is a
well-calibrated approximation that fires reliably when the relevant
stratum boundary is approached.

The finitude of h\_promise is morally significant. Unlike murder,
promise-breaking does not trigger a nullifier. The broken-promise state
occupies a damaged but recoverable stratum: trust can be rebuilt,
compensation can be offered, and the social bond can be repaired --- at
a cost. The heuristic correctly encodes this recoverability as a finite,
large estimate rather than an infinite one.

\subsubsection{The Evolutionary Origin}\label{the-evolutionary-origin}

Why do these heuristic functions exist? The intractability result of
§11.5 provides the answer: exact moral computation is intractable, so
any resource-bounded agent that must make real-time moral decisions
requires pre-compiled approximations. Biological evolution and cultural
evolution have provided these approximations in the form of moral
intuitions, social norms, and codified rules. They are the moral
analogue of pattern databases in classical A* --- offline computation
(evolutionary selection, cultural refinement, legal codification) that
enables fast online lookup (moral intuition, rule-following, legal
compliance). A clarification is warranted: the A* framework is not an
evolutionary debunking argument. The admissibility of a heuristic ---
the mathematical property h(n) ≤ g* --- holds or fails independently of
the heuristic's causal history. An obligation that was shaped by natural
selection, cultural transmission, or rational deliberation is admissible
if and only if it does not overestimate the true cost. The evolutionary
account explains how the heuristics were compiled; it does not determine
whether they track the structure of moral space.

When an agent experiences the phenomenological force of a moral
obligation --- the sense that they must act in a certain way, the guilt
of transgression, the relief of compliance --- they are reading the
output of these pre-compiled heuristic functions. The phenomenology is
the interface; the computation is the A* evaluation f(n) = g(n) + h(n).

\subsection{11.4 Admissibility and
Consistency}\label{admissibility-and-consistency}

The A* algorithm is optimal --- guaranteed to find the cheapest path ---
if and only if the heuristic function h satisfies two conditions:
admissibility and consistency. This section defines both conditions in
the moral context and proves that the core moral heuristics satisfy
them.

\textbf{Definition 11.3 (Admissibility). {[}Formal Definition.{]}} A
moral heuristic h is admissible if, for every state n ∈ ℳ:

h(n) ≤ g*(n, G\_eq)

where g*(n, G\_eq) is the true minimum behavioral friction from n to the
goal region G\_eq. An admissible heuristic never overestimates the true
cost to equilibrium. It is an optimistic lower bound.

\textbf{Definition 11.4 (Consistency). {[}Formal Definition.{]}} A moral
heuristic h is consistent (or monotone) if, for every pair of
neighboring states n, n′ ∈ ℳ:

h(n) ≤ c(n, n′) + h(n′)

where c(n, n′) is the local behavioral friction cost of the transition
from n to n′. This is the triangle inequality on the metric space: the
estimated cost from n cannot exceed the cost of one step plus the
estimated cost from the next state. Consistency implies admissibility
(standard result), and ensures that A* never needs to re-expand a node.

\textbf{Theorem 11.1 (Admissibility of Core Moral Heuristics; the
Heuristic Truncation Theorem).} {[}Conditional Theorem; conditional on
the behavioral friction model of Chapter 6 and the Whitney
stratification of Chapter 8.{]} The classical moral prohibitions
(murder, theft, deception) and positive obligations (reciprocity,
promise-keeping, care) are admissible heuristics in the A* sense: they
never overestimate the true behavioral friction required to reach
equilibrium from the states they prohibit or prescribe.

\textbf{Proof sketch}. We argue case by case. For murder: h\_murder(n) =
∞ for states involving killing. As shown in §11.3, the true cost g*(n,
G\_eq) is also infinite (the stratum boundary is absorbing), so h ≤ g*.
For theft: h\_theft(n) estimates the cost of property-rights
restoration, institutional response, and trust repair. The true cost
includes all of these plus network externalities (other agents adjust
their behavior), so g* ≥ h. For deception: h\_deception(n) estimates the
trust-rebuilding cost. Trust rebuilding is slow and path-dependent
(Chapter 10, holonomy), and the true cost includes the uncertainty
premium imposed by all agents who learn of the deception, so g* ≥ h. For
reciprocity norms: the heuristic points toward cooperative equilibria,
which are by definition the geodesics of lowest friction (Chapter 10,
§10.7). The full proof, with detailed cost analysis for each heuristic
using the stratum structure of Chapter 8, appears in the Technical
Appendix. □

\textbf{Proposition 11.2 (Consistency from Metric Compatibility).
{[}Proved.{]}} If h(n) = d\_g(n, G\_eq) --- the geodesic distance from n
to the goal region in the moral metric g\_\{μν\} --- then h is
consistent.

\textbf{Proof}. By the triangle inequality for the Riemannian distance
d\_g: for any states n, n′ and goal region G\_eq, d\_g(n, G\_eq) ≤
d\_g(n, n′) + d\_g(n′, G\_eq). The local behavioral friction cost c(n,
n′) is the arc length of the path from n to n′, which satisfies c(n, n′)
≥ d\_g(n, n′) (arc length ≥ geodesic distance). Therefore h(n) = d\_g(n,
G\_eq) ≤ d\_g(n, n′) + d\_g(n′, G\_eq) ≤ c(n, n′) + h(n′), which is the
consistency condition. □

\textbf{Remark}. Real moral heuristics approximate but do not exactly
equal the geodesic distance. They are approximately consistent --- the
triangle inequality holds in most cases but may fail at stratum
boundaries (where the metric changes discontinuously) or in exotic moral
configurations (where the heuristic's pre-compiled assumptions break
down). These failures correspond to moral dilemmas: situations where the
pre-compiled heuristic gives contradictory or unreliable guidance, and
the agent must fall back on explicit moral reasoning --- that is, on a
more careful search through the local neighborhood of the manifold. The
implications for optimality are bounded, not catastrophic. The
literature on weighted A* and ε-admissible search (Pohl 1970; Pearl
1984) establishes that approximately consistent heuristics yield
bounded-suboptimal paths: the solution cost is at most (1 + ε) times the
true optimum, where ε reflects the degree of inconsistency. Real moral
reasoning is therefore near-optimal --- close to the geodesic, but not
exactly on it. This is precisely the content of Corollary 11.2:
heuristic imperfection is structural and bounded, not arbitrary.

\subsubsection{The ε-Admissibility Spectrum and the Heuristic
Hierarchy}\label{the-ux3b5-admissibility-spectrum-and-the-heuristic-hierarchy}

The preceding remark acknowledges that real moral heuristics may
overestimate. This section formalizes the consequences. The key insight
is that overestimation is not binary (admissible vs. inadmissible) but
falls on a spectrum with distinct moral signatures at each level.

\textbf{Definition 11.5 (ε-Admissible Moral Heuristic). {[}Formal
Definition.{]}} A moral heuristic h is ε-admissible if, for all states n
∈ ℳ: h(n) ≤ (1 + ε) · g*(n, G\_eq), where g*(n, G\_eq) is the true
minimum behavioral friction from n to the goal region. When A* search
uses an ε-admissible heuristic, the returned path has total cost at most
(1 + ε) times the optimal path cost. The parameter ε quantifies the
degree of moral conservatism: ε = 0 yields perfect optimality; small ε
yields near-optimal paths at reduced computational cost; large ε yields
fast but potentially suboptimal moral judgments.

\textbf{Remark} (Connection to Weighted A*). In the search literature,
ε-admissibility is achieved by weighted A*: f(n) = g(n) + w · h(n),
where w = 1 + ε. The weight w trades optimality for computational
tractability. This maps directly onto moral reasoning: a time-pressured
agent (emergency physician, crisis responder) implicitly increases w ---
relying more heavily on heuristic judgment and less on deliberate
calculation --- accepting a bounded-suboptimal moral path in exchange
for the speed needed to act. The bounded suboptimality guarantee ensures
that this is not moral recklessness but a principled trade-off with a
quantifiable cost ceiling (Pohl 1970; Pearl 1984).

\textbf{Theorem 11.1b (The Moral Heuristic Hierarchy). {[}Conditional
Theorem; conditional on the behavioral friction model of Chapter 6 and
the stratification of Chapter 8.{]}} Moral heuristics fall into four
categories with distinct phenomenological signatures:

\begin{enumerate}
\def\labelenumi{\roman{enumi}.}
\item
  \textbf{Strictly admissible (h(n) ≤ g*(n, G\_eq))}: The heuristic
  underestimates or correctly estimates the true moral cost. The A* path
  is optimal. These are the core prohibitions whose true costs genuinely
  exceed any finite estimate --- murder (true cost includes
  absorbing-stratum penalties, legal consequences, social ostracism, and
  psychological damage), non-consensual treatment (true cost includes
  rights violation, institutional liability, trust destruction). The
  proof of Theorem 11.1 establishes this category.
\item
  \textbf{ε-admissible (h(n) ≤ (1 + ε) · g*(n, G\_eq), small ε)}: The
  heuristic slightly overestimates the true boundary cost. The A* path
  is near-optimal, with total cost at most (1 + ε) times the true
  optimum. This is the normal operating range of a well-calibrated moral
  agent. The slight conservatism --- marginally overweighting harm
  avoidance, slightly inflating the cost of promise-breaking ---
  produces the well-documented phenomenon that experienced moral
  reasoners are slightly more cautious than a pure expected-value
  calculation would dictate, at negligible cost to overall moral
  performance. Evolutionary and cultural selection pressures calibrate ε
  to be small but positive: an agent with ε = 0 (perfect calibration) is
  fragile (any perturbation makes the heuristic inadmissible), while an
  agent with small positive ε is robust.
\item
  \textbf{Inadmissible (h(n) ≫ g*(n, G\_eq))}: The heuristic
  substantially overestimates boundary costs. The A* path avoids
  genuinely beneficial actions because the inflated heuristic makes them
  appear too costly. This category has a precise clinical and
  phenomenological signature: it is the mathematical structure of moral
  injury, trauma, and pathological risk aversion. An agent whose
  heuristic has been damaged by forced boundary crossing (the moral
  injury formalized in the domain literature) has inflated β\_k values
  --- they overestimate the cost of actions near the violated boundary,
  avoiding even beneficial actions in that region. Defensive behavior
  (defensive medicine, defensive lawyering, institutional risk aversion)
  is β\_institutional ≫ β*\_institutional: the institution's heuristic
  overweights legal and reputational cost relative to its true value,
  foreclosing beneficial actions.
\item
  \textbf{Gauge-variant}: The heuristic's output depends on the framing
  or description of the moral situation rather than on the situation's
  attribute vector. This violates the Bond Invariance Principle (Chapter
  5, §5.5): the heuristic is not invariant under admissible
  re-descriptions. Examples include omission bias (harmful action judged
  as worse than equally harmful inaction, even when the attribute
  vectors are identical), status quo bias (framing the current state as
  the reference point inflates the cost of departure), anchoring to
  initial information, and order effects in moral judgment. These are
  the phenomena that genuinely deserve the label `cognitive bias' in
  moral reasoning --- not because the heuristic overestimates, but
  because it responds to morally irrelevant features of the description.
  The distinction matters: categories (i)--(iii) are calibration errors
  (the heuristic estimates the right quantity, more or less accurately);
  category (iv) is a gauge failure (the heuristic estimates the wrong
  quantity).
\end{enumerate}

\textbf{Proof sketch}. Category (i) follows from Theorem 11.1 (the
Heuristic Truncation Theorem). Category (ii) follows from the bounded
suboptimality guarantee of ε-admissible A* (Pohl 1970, Theorem 2; Pearl
1984, §4.2): if h(n) ≤ (1 + ε) · g*, then f(γ) ≤ (1 + ε) · f(γ*).
Category (iii) is defined by the negation of (i)--(ii): when ε exceeds a
threshold, the heuristic forecloses paths whose true cost is acceptable,
producing systematic moral avoidance. Category (iv) is defined by
violation of gauge invariance (BIP, Chapter 5): h(τ(n)) ≠ h(n) for some
admissible transformation τ. The four categories are mutually exclusive
and exhaustive for any heuristic function on ℳ. □

\textbf{Remark} (Heuristic Recalibration). The hierarchy has a dynamic
interpretation. An agent's heuristic can move between categories over
time. Moral education and professional training move uncalibrated
heuristics from category (iii) or (iv) toward categories (i)--(ii) ---
this is what it means for clinical training, legal education, or
military ethics instruction to `calibrate moral judgment.' Conversely,
moral injury (forced boundary crossing that inflates β\_k) moves a
previously well-calibrated heuristic from category (ii) to category
(iii) --- the mathematical signature of moral damage. Ethics
consultation, peer review, and therapeutic intervention function as
recalibration mechanisms: they restore inadmissible heuristics toward
ε-admissibility by correcting inflated boundary penalties to values
closer to their true costs. The framework predicts that effective moral
repair must target specific β\_k values (not general `resilience'), and
that the dimension of repair must match the dimension of injury.

\subsection{11.5 Computational Intractability of Exact Moral
Reasoning}\label{computational-intractability-of-exact-moral-reasoning}

We now prove the computational result that justifies the entire A*
framework: exact moral geodesic planning is intractable. This result
explains why moral heuristics exist --- they are not cultural artifacts
but computational necessities.\\
\strut \\
{[}Formal complexity result.{]} Define MORAL-PATHFIND as the decision
problem: given a weighted simplicial complex C with n vertices,
d-dimensional attribute vectors with Mahalanobis edge weights, and k ≥ 1
attribute-boundary constraints (edges that cross a constraint boundary
incur penalty βₖ), does there exist a path from s to G with total cost ≤
B? MORAL-PATHFIND is NP-hard by reduction from the Weight-Constrained
Shortest Path problem (Garey \& Johnson, 1979; Handler \& Zang, 1980),
which is NP-hard for k ≥ 1 side constraints. The reduction maps each
side constraint to a boundary penalty β\_k and each edge weight to the
corresponding Mahalanobis distance. Since MORAL-PATHFIND is NP-hard,
exact optimization on the moral manifold with boundary constraints is
computationally intractable in the worst case, and the
bounded-suboptimality guarantees of Weighted A* (Pohl, 1970) apply by
direct instantiation --- not by analogy --- to the moral search problem.

\textbf{Theorem 11.2 (Intractability of Exact Moral Geodesic Planning).
{[}Conditional Theorem; conditional on the 9-dimensional moral manifold
of Chapter 5 and the behavioral friction model of Chapter 6.{]}} Finding
the exact optimal path on the moral manifold ℳ is computationally
practically intractable at the resolution required for morally adequate
evaluation, and the associated discrete constrained-path problem
(MORAL-PATHFIND) is NP-hard.

\textbf{Proof}. The proof proceeds by analyzing the computational cost
of discretized geodesic planning on a d-dimensional stratified
Riemannian manifold.

\textbf{Step 1:} Discretization. To compute geodesics numerically, the
continuous manifold ℳ must be approximated by a discrete graph. To
achieve accuracy ε on a d-dimensional manifold, a standard grid
discretization requires N ∝ (1/ε)ᵈ vertices per stratum. For d = 9 (the
dimensionality of the moral manifold, Definition 5.1), this gives N ∝
(1/ε)⁹ vertices per stratum.

\textbf{Step 2:} Per-stratum complexity. On a graph with N vertices and
O(N²) edges (the connectivity of the discretized manifold), the
shortest-path algorithm (Dijkstra) runs in O(N² log N) time. Each edge
weight evaluation requires O(d²) = O(81) operations for the metric
tensor contraction. With m strata, the total is O(mN² · d² · log(mN))
--- polynomial in N, m, and d. This is the result of Bond (2026, SGE),
Theorem 5.3.

\textbf{Step 3:} The curse of dimensionality. Substituting N ∝ (1/ε)⁹,
the total complexity becomes O(m · (1/ε)¹⁸ · d² · log(m(1/ε)⁹)). For
practical accuracy ε = 0.01, N ∝ 10¹⁸ per stratum. The total vertex
count across m strata is on the order of m · 10¹⁸, and the edge count on
the order of m · 10¹⁸. Even with m = 1, a single Dijkstra evaluation
requires on the order of 10¹⁸ metric evaluations --- vastly exceeding
the computational capacity of any physical system operating in real
time.

\textbf{Step 4:} Multi-agent and temporal extensions. In a multi-agent
setting with n stakeholders, the state space grows by a factor of n
(each stakeholder contributes a perspective). With a temporal horizon of
T time steps and stochastic dynamics (branching factor b at each step),
the effective search space grows by a factor of bᵀ. The overall
complexity is O(m · (1/ε)¹⁸ · n · bᵀ · d² · log(mN)) --- exponential in
both the dimensionality d and the temporal horizon T.

\textbf{Conclusion}. While the per-vertex computation is polynomial (SGE
Theorem 5.3), the number of vertices required for adequate accuracy on a
9-dimensional manifold is astronomically large. Exact moral geodesic
planning is intractable not because the algorithm is inefficient, but
because the space is too large to discretize adequately. This is the
curse of dimensionality applied to moral reasoning. □\\
\strut \\
{[}Formal complexity result.{]} Define MORAL-PATHFIND as the decision
problem: given a weighted simplicial complex C with n vertices,
d-dimensional attribute vectors with Mahalanobis edge weights, and k ≥ 1
attribute-boundary constraints (edges that cross a constraint boundary
incur penalty βₖ), does there exist a path from s to G with total cost ≤
B? MORAL-PATHFIND is NP-hard by reduction from the Weight-Constrained
Shortest Path problem (Garey \& Johnson, 1979; Handler \& Zang, 1980),
which is NP-hard for k ≥ 1 side constraints. The reduction maps each
side constraint to a boundary penalty β\_k and each edge weight to the
corresponding Mahalanobis distance. Since MORAL-PATHFIND is NP-hard,
exact optimization on the moral manifold with boundary constraints is
computationally intractable in the worst case, and the
bounded-suboptimality guarantees of Weighted A* (Pohl, 1970) apply by
direct instantiation --- not by analogy --- to the moral search problem.

A note on the nature of the intractability claim. The argument above
establishes practical intractability --- the resource requirements
exceed any physically realizable computation --- rather than
complexity-theoretic hardness in the sense of NP-completeness or
PSPACE-completeness. These are distinct but complementary notions. For
the purpose of this framework, practical intractability is the relevant
concept: the question is whether any resource-bounded agent (biological
or artificial) can compute exact moral geodesics in real time, and the
answer is no, regardless of algorithmic cleverness. The curse of
dimensionality is intrinsic to the 9-dimensional manifold: even adaptive
discretization methods (which concentrate resolution near stratum
boundaries and use coarser grids elsewhere) reduce constants but cannot
eliminate the exponential dependence on d. The fundamental bottleneck is
not the algorithm but the volume of the space that must be explored.

\textbf{Corollary 11.1 (Evolutionary Necessity of Moral Heuristics).
{[}Conditional Theorem.{]}}

Since exact moral computation is intractable, any resource-bounded agent
--- biological or artificial --- that must make real-time moral
decisions requires pre-compiled heuristic functions. The existence of
moral obligations is a computational necessity, not merely a cultural
artifact.

\textbf{Corollary 11.2 (Heuristic Imperfection Is Structural).
{[}Conditional Theorem.{]}}

Since moral heuristics are finite approximations to an intractable
problem, they are necessarily imperfect. Moral error --- cases where
obligation-following leads to suboptimal outcomes --- is a structural
consequence of computational intractability, not a defect of moral
cognition. This resolves a longstanding puzzle in moral philosophy: why
do good people sometimes do wrong? Not because they are morally
deficient, but because they are using heuristics for a problem that
admits no tractable exact solution. Importantly, structural explanation
does not eliminate moral responsibility. An agent remains accountable
for the calibration of their heuristics --- the virtue ethics
interpretation developed in §11.8 below --- and for the decision to
invoke deliberative reasoning (System 2) when stakes are high or when
heuristic guidance is manifestly unreliable. The corollary explains the
mechanism of moral error, not its moral status.

\subsection{11.6 The Satisfaction Scalar as Search
Output}\label{the-satisfaction-scalar-as-search-output}

The A* framework provides a new interpretation of the fundamental
contraction formula from Chapter 6. The satisfaction scalar S = I\_μ Oᵐ
(§6.4, Definition 6.3) is the inner product of what an agent is
obligated to do (the obligation vector Oᵐ, the heuristic gradient) with
what the affected parties need (the interest covector I\_μ). In the A*
framework, this contraction is the moment when the search yields a
decision.

Proposition 11.3 (Satisfaction as Directional Derivative of the
Heuristic). {[}Proved.{]} The satisfaction scalar is the directional
derivative of the heuristic function h in the interest direction:

S = I\_μ Oᵐ = −I\_μ gᵐᵛ ∂ₙh

\textbf{Proof}. By Proposition 11.1, Oᵐ = −gᵐᵛ ∂ₙh. Substituting into
the satisfaction formula (Definition 6.3), S = I\_μ Oᵐ = I\_μ(−gᵐᵛ ∂ₙh)
= −I\_μ gᵐᵛ ∂ₙh. This is the contraction of the interest covector with
the metric-raised gradient of h, which is by definition the directional
derivative of h in the direction determined by I\_μ through the metric.
□

When S \textgreater{} 0, the heuristic's recommended direction (the
obligation) aligns with the patient's needs (the interest): following
the moral rule serves the stakeholder. When S \textless{} 0, obligation
and interest conflict: the rule prescribes an action that, in this
context, works against the stakeholder's needs. When S = 0, the action
is morally neutral along the measured axes.

\subsubsection{Worked Example: The Kidney Allocation
Revisited}\label{worked-example-the-kidney-allocation-revisited}

Chapter 7 examined a kidney allocation decision at five levels of
geometric sophistication. We can now recognize that the Chapter 7
analysis was A* search throughout --- it simply was not identified as
such.

At Level 1 (scalar), the analysis computed a single score for each
patient --- equivalent to a degenerate A* search with a one-dimensional
heuristic. At Level 2 (obligation vector), the analysis computed Oᵐ for
each allocation option --- the heuristic gradient in the search
framework. At Level 3 (multi-agent rank-2 tensor), the analysis tracked
separate interest covectors I\_μ for each stakeholder --- multiple A*
evaluations, one per perspective. At Level 4 (the metric), the analysis
used g\_\{μν\} to compute the behavioral friction cost of each option
--- the g(n) term in A*. At Level 5 (stratification), the analysis
identified stratum boundaries where certain allocations trigger deontic
phase transitions --- the constraint surfaces that the A* search must
avoid.

The five-level progression of Chapter 7 is the progressive enrichment of
the A* search: from a trivial search (scalar, one-dimensional, no
structure) to a full search on the stratified manifold (tensorial,
nine-dimensional, with metric, stratification, and gauge constraints).
Each level adds computational structure that makes the search more
accurate but also more expensive --- the tradeoff at the heart of the
intractability result (Theorem 11.2).

The Inevitability of Residue (Chapter 15, Proposition 15.1) applies
directly: each level of the Chapter 7 analysis that collapses structure
(from tensor to vector, from vector to scalar) discards information
about the search space. The information lost is precisely the
information about alternative paths that the collapsed representation
cannot encode. The A* framework makes this loss concrete: it is the
pruning of the search tree at each contraction step.

\subsection{11.7 Connections to Existing Framework
Elements}\label{connections-to-existing-framework-elements}

The A* interpretation does not stand alone. It connects to and
illuminates several existing elements of the geometric ethics framework.

\subsubsection{Gradient Flows as Degenerate
A*}\label{gradient-flows-as-degenerate-a}

Chapter 10 (§10.8) introduced gradient flows: the agent follows the
direction of steepest improvement of the satisfaction function S. The
gradient flow is grad S = gᵐᵛ ∂S/∂xᵛ. In the A* framework, gradient flow
corresponds to greedy best-first search with g = 0: the agent follows
the heuristic's recommendation (obligation) without accounting for
accumulated cost. The pathologies identified in §10.8 --- local maxima,
saddle points, metric sensitivity --- are exactly the pathologies of
searchless optimization. The remedy is the same in both cases:
incorporate backward cost (the g(n) term) to escape local optima.

\subsubsection{Geodesics as A*
Solutions}\label{geodesics-as-a-solutions}

The geodesic of Chapter 10 (§10.7) is the optimal path found by A* with
a perfect heuristic. If h = g* (the true cost-to-go), then A* finds the
exact geodesic on the first expansion. Since computing g* requires
solving the full optimization problem (which is intractable by Theorem
11.2), real agents use approximate heuristics h ≈ g*, and the A* path
approximates but does not exactly follow the geodesic. The quality of
the approximation depends on the quality of the heuristic --- that is,
on the calibration of the agent's moral intuitions.\\
\strut \\
{[}Formal complexity result.{]} Define MORAL-PATHFIND as the decision
problem: given a weighted simplicial complex C with n vertices,
d-dimensional attribute vectors with Mahalanobis edge weights, and k ≥ 1
attribute-boundary constraints (edges that cross a constraint boundary
incur penalty βₖ), does there exist a path from s to G with total cost ≤
B? MORAL-PATHFIND is NP-hard by reduction from the Weight-Constrained
Shortest Path problem (Garey \& Johnson, 1979; Handler \& Zang, 1980),
which is NP-hard for k ≥ 1 side constraints. The reduction maps each
side constraint to a boundary penalty β\_k and each edge weight to the
corresponding Mahalanobis distance. Since MORAL-PATHFIND is NP-hard,
exact optimization on the moral manifold with boundary constraints is
computationally intractable in the worst case, and the
bounded-suboptimality guarantees of Weighted A* (Pohl, 1970) apply by
direct instantiation --- not by analogy --- to the moral search problem.

\subsubsection{Holonomy and Heuristic
Drift}\label{holonomy-and-heuristic-drift}

Chapter 10 (§10.5) showed that parallel transport of obligations around
a closed loop may change them --- the holonomy of the moral connection.
In the A* framework, holonomy means that the heuristic function h drifts
as the agent traverses moral space: the pre-compiled moral intuitions
that were accurate at the starting point may become miscalibrated after
a circuit through diverse moral experiences. This explains why moral
codes require updating through legal reform, cultural evolution, and
moral education --- the heuristics lose accuracy as the landscape
changes, and must be recalibrated.

\subsubsection{Contraction and Information
Loss}\label{contraction-and-information-loss}

The Inevitability of Residue (Chapter 15, Proposition 15.1) proves that
contracting a moral tensor to a scalar is irreversibly lossy. The A*
framework makes this loss precise: it is the information about the full
search tree that is discarded when the agent commits to a single path.
The satisfaction scalar S = I\_μ Oᵐ is the shadow of the search tree
projected onto a single dimension. The residue (§15.3) is the record of
paths not taken --- the moral possibilities that were evaluated and
rejected in the search process.

\subsection{11.8 Implications for Moral
Philosophy}\label{implications-for-moral-philosophy}

The A* framework provides computational interpretations of the major
ethical traditions, explaining their strengths, limitations, and domains
of applicability.

\subsubsection{Deontology as Pre-Compiled
Heuristics}\label{deontology-as-pre-compiled-heuristics}

Deontological rules (``do not lie,'' ``keep your promises,'' ``respect
autonomy'') are pre-compiled heuristic functions. They work because they
are admissible --- they never overestimate the true cost of the actions
they prohibit --- and because they are fast: evaluating h(n) requires
constant time, independent of the size of the search space. They fail
when the heuristic is inconsistent (edge cases where following the rule
increases total cost, such as the murderer-at-the-door who asks where
your friend is hiding) or when the goal region has shifted (moral
progress makes old heuristics obsolete, as with the historical
acceptance of slavery, which reflected a miscalibrated heuristic that
assigned insufficient cost to rights violations).

\subsubsection{Consequentialism as Exact
Search}\label{consequentialism-as-exact-search}

Pure consequentialism attempts to compute g* directly --- the exact
optimal path from the current state to equilibrium, evaluated by its
consequences. This is equivalent to running A* without a heuristic (h =
0), or equivalently, running Dijkstra's algorithm on the full moral
manifold. By Theorem 11.2, this is intractable. Consequentialist
reasoning is computationally expensive, which is why it is typically
reserved for novel or high-stakes situations where the pre-compiled
heuristics are clearly insufficient --- exactly as predicted by the
dual-process model of moral cognition.

\subsubsection{Virtue Ethics as Heuristic
Calibration}\label{virtue-ethics-as-heuristic-calibration}

A virtuous agent, in the A* framework, is one whose heuristic functions
are well-calibrated: the obligations they experience are not merely
admissible (never overestimating) but tight (close to the true cost g*).
A tighter heuristic prunes more of the search space, making the agent's
moral reasoning faster and more accurate. Aristotle's doctrine of
habituation --- that virtue is developed through practice --- is the
process of refining heuristics through experience: each moral decision
provides feedback that tightens the heuristic's approximation. The
``mean between extremes'' is the region where the heuristic is most
accurate: neither overestimating (excessive caution, moral paralysis)
nor underestimating (recklessness, moral blindness).

\subsubsection{Dual-Process Moral
Cognition}\label{dual-process-moral-cognition}

The A* framework maps naturally onto the dual-process model of moral
cognition (Kahneman's System 1 and System 2). System 1 --- fast,
automatic, intuitive --- is the heuristic evaluation h(n): the
pre-compiled obligation that fires instantly when a morally relevant
situation is detected. System 2 --- slow, deliberate, effortful --- is
the full A* search: the explicit computation of backward cost g(n),
forward estimate h(n), and their combination f(n) = g(n) + h(n), applied
iteratively across the search tree. System 2 is invoked when System 1
fails --- that is, when the heuristic is unreliable (near stratum
boundaries), when the moral situation is novel (the heuristic was not
pre-compiled for this case), or when the stakes are high enough to
justify the computational cost.

A crucial clarification. The computational interpretations above do not
reduce deontology to a crude approximation of consequentialist
calculation, nor do they privilege consequentialism as the ``true'' form
of moral reasoning. The cost function in the A* framework is behavioral
friction --- a geometric quantity defined by the moral metric (Chapter
6), which encodes rights violations, autonomy costs, fairness deficits,
and other structural features of moral space alongside consequences. It
is emphatically not a utilitarian welfare function. Both deontological
and consequentialist reasoning are computational strategies for
navigating the same underlying geometric structure, and the structure is
prior to both. The framework does not adjudicate between them; it
explains why both exist and where each is most effective. Kant's
categorical imperative, read geometrically, tests the invariance of
moral maxims under coordinate transformations (§5.5) --- a gauge
constraint on the search, not a heuristic approximation to
consequence-calculation.

\subsection{11.9 Worked Example: The Self-Driving
Car}\label{worked-example-the-self-driving-car}

We illustrate the full A* framework with a self-driving car scenario
that connects to the DEME architecture of Chapter 19.

Setup. An autonomous vehicle approaches an intersection. A pedestrian is
crossing against the light on the left. A cyclist is in the bike lane on
the right. Oncoming traffic is approaching. The vehicle must choose a
path through moral configuration space: brake hard (high friction for
passengers, low friction for pedestrian), swerve left (high friction for
oncoming traffic), swerve right (high friction for cyclist), or continue
(catastrophic friction for pedestrian).

State space. Each candidate action moves the system to a new state n\_i
∈ ℳ, characterized by its position in the nine-dimensional moral
manifold: consequences (physical harm to each party), rights
(right-of-way violations), fairness (distribution of risk), autonomy
(override of pedestrian's choice to cross), privacy (irrelevant in this
case), societal impact (precedent-setting), virtue (the character of the
decision-maker), legitimacy (traffic law compliance), and epistemic
status (uncertainty about the pedestrian's intent).

Backward cost g(n\_i). For each candidate state, the accumulated
behavioral friction of the path taken so far. This includes the costs of
previous driving decisions, the current speed and trajectory, and the
commitments implied by the vehicle's position on the road.

Heuristic h(n\_i). The pre-compiled obligations: (i) do not harm the
pedestrian (h = very high for n\_continue; illustratively, h ≈ 95); (ii)
do not harm the cyclist (h = high for n\_swerve\_right; h ≈ 70); (iii)
minimize harm to passengers (h = moderate for n\_brake\_hard; h ≈ 30);
(iv) obey traffic law (h = moderate for n\_swerve\_left, which enters
the oncoming lane; h ≈ 55).

A* evaluation. f(n\_i) = g(n\_i) + h(n\_i) for each candidate. With
illustrative backward costs g(n\_continue) ≈ 5, g(n\_brake\_hard) ≈ 15,
g(n\_swerve\_right) ≈ 10, g(n\_swerve\_left) ≈ 12, the total evaluations
are: f(n\_continue) ≈ 100, f(n\_swerve\_right) ≈ 80, f(n\_swerve\_left)
≈ 67, f(n\_brake\_hard) ≈ 45. The A* algorithm selects n\_brake\_hard:
moderate harm to passengers (the deceleration cost is finite and
recoverable), minimal harm to the pedestrian (the vehicle stops before
the crossing), and no stratum boundary crossing (no one is killed or
seriously injured). The scalar comparison shows why a single ``harm
score'' would miss the structure: swerving right has lower total harm
(the cyclist is uninjured in most scenarios) but higher
obligation-violation cost (the cyclist has the right of way), so the
full A* evaluation correctly prefers braking over swerving.

The DEME architecture (Chapter 19) implements exactly this A* search.
The EthicalFacts layer canonicalizes the sensory input into the
nine-dimensional state description. The Ethics Modules compute h(n\_i)
for each candidate action, using pre-compiled obligation vectors. The
Governance layer combines g(n\_i) and h(n\_i) into f(n\_i) and selects
the action with the lowest total cost. The audit layer records the full
search tree for post-hoc inspection. The No Escape Theorem (Chapter 18,
Theorem 18.1) ensures that the search cannot be bypassed through
representational manipulation.

\subsection{11.10 Summary}\label{summary-4}

This chapter has established a formal equivalence between moral
reasoning and A* pathfinding on the moral manifold. The principal
results are:

\begin{itemize}
\item
  Moral reasoning is A* search. The equation f(n) = g(n) + h(n), where g
  is accumulated behavioral friction and h is a heuristic estimating
  cost-to-equilibrium, is the fundamental equation of moral reasoning
  (Definition 11.2).
\item
  Obligation vectors are heuristic gradients. The obligation vectors Oᵐ
  of Chapter 6 are the negative metric-raised gradients of the heuristic
  function h (Proposition 11.1). Moral obligations point in the
  direction of steepest estimated improvement.
\item
  Core moral heuristics are admissible. The classical prohibitions
  (murder, theft, deception) and positive obligations (reciprocity,
  care) never overestimate the true cost to equilibrium (Theorem 11.1).
  This is why rule-following is generally reliable.
\item
  Exact moral computation is intractable. The combinatorial explosion of
  fine-resolution discretization on the 9-dimensional moral manifold,
  combined with the NP-hardness of constrained path planning, makes
  exact geodesic planning practically intractable for any
  resource-bounded agent (Theorem 11.2). Moral heuristics are
  computationally necessary, not merely culturally convenient.
\item
  Moral error is structural. Since heuristics approximate an intractable
  problem, they are necessarily imperfect (Corollary 11.2). When good
  people do wrong, the explanation is computational, not
  characterological.
\end{itemize}

The A* framework unifies several threads of the preceding chapters:
geodesics (§10.7) are the solutions to the search problem; gradient
flows (§10.8) are degenerate searches without backward cost; obligation
vectors (§6.2) are heuristic gradients; and Inevitability of Residue
(Chapter 15, Proposition 15.1) quantifies the information lost when the
search tree is contracted to a scalar decision.

The framework also reveals a new requirement: for the A* search to yield
consistent results, the heuristic function h(n) must be invariant under
re-description --- the same moral state, labeled differently, must
receive the same heuristic estimate. This invariance requirement is
precisely the Bond Invariance Principle (§5.5), and its consequences ---
via Noether's theorem --- are the subject of the next chapter.

\subsection{Technical Appendix}\label{technical-appendix-5}

\textbf{Proof of Theorem 11.1 (Admissibility of Core Moral Heuristics).}

We prove admissibility (h(n) ≤ g*(n, G\_eq)) for each core heuristic by
constructing lower bounds on the true cost g* and showing that h does
not exceed them.

\textbf{Case 1}: Murder prohibition (h\_murder = ∞). Let n be a
post-murder state. By Definition 8.3 (Type III boundary/nullifier),
murder triggers an absorbing stratum: the system enters a catastrophic
regime from which return to the goal region G\_eq requires, at minimum:
(i) grief processing across k connected agents, with per-agent cost
c\_grief and propagation depth proportional to the connectivity of the
social network; (ii) institutional response with cost c\_inst =
c\_investigation + c\_adjudication + c\_incarceration; (iii)
reputational damage d\_rep propagating through the network with
attenuation proportional to distance; (iv) deterrence expenditure c\_det
to prevent recurrence. The total cost g*(n, G\_eq) ≥ k · c\_grief +
c\_inst + d\_rep + c\_det. In the limit of large k (connected network),
g* diverges. For finite k, the irreversibility of death means that the
specific harm cannot be undone --- only its consequences can be managed.
In either case, g* ≥ h\_murder, and the heuristic is admissible. For the
case h\_murder = ∞: since the absorbing stratum may have no path to
G\_eq, g* = ∞ as well, and the inequality h ≤ g* holds trivially.

\textbf{Case 2}: Theft prohibition (h\_theft = C\_property). Let
C\_property be the heuristic's estimate of the cost of property-rights
restoration. The true cost g*(n, G\_eq) includes: (i) the direct
restoration cost C\_property; (ii) institutional costs (investigation,
adjudication); (iii) trust-rebuilding costs across affected
relationships; (iv) deterrence expenditure. Since g* ≥ C\_property +
additional terms ≥ C\_property = h\_theft, the heuristic is admissible.

\textbf{Case 3}: Deception prohibition (h\_deception = C\_trust). Let
C\_trust be the heuristic's estimate of the trust-rebuilding cost. By
the holonomy result of Chapter 10 (§10.5), trust rebuilding is
path-dependent: the cost depends on the sequence of interactions through
which trust is restored, and holonomy ensures that the cost is
generically non-zero. The true cost g* includes C\_trust plus the
uncertainty premium (other agents increase their monitoring costs,
reducing the efficiency of future interactions) plus potential cascade
effects (the deception may trigger retaliatory deceptions, increasing
systemic friction). Since g* ≥ C\_trust = h\_deception, the heuristic is
admissible.

\textbf{Case 4}: Reciprocity norm (h\_reciprocity = C\_defection). The
reciprocity heuristic assigns high cost to non-reciprocal actions
(defection in repeated games). By the game-theoretic analysis of
behavioral friction (the social cost of defection in an iterated setting
includes retaliation, reputation damage, and exclusion from cooperative
networks), g* ≥ C\_defection = h\_reciprocity. □

\textbf{Proof of Proposition 11.2 (Consistency from Metric
Compatibility).}

Let h(n) = d\_g(n, G\_eq), where d\_g is the geodesic distance in the
moral metric g\_\{μν\}. For any neighboring states n, n′: By the
triangle inequality for Riemannian distance, d\_g(n, G\_eq) ≤ d\_g(n,
n′) + d\_g(n′, G\_eq). The local behavioral friction cost c(n, n′) is
the arc length of the path from n to n′, which satisfies c(n, n′) ≥
d\_g(n, n′) (the arc length of any path between two points is at least
the geodesic distance between them --- this is the defining property of
geodesic distance). Therefore: h(n) = d\_g(n, G\_eq) ≤ d\_g(n, n′) +
d\_g(n′, G\_eq) ≤ c(n, n′) + h(n′). This is the consistency condition
(Definition 11.4). □

\textbf{Proof of Theorem 11.2 (Intractability of Exact Moral Geodesic
Planning).}

We establish intractability via explicit complexity calculation, not
reduction from a known hard problem. The argument proceeds in four
steps.

Step 1 (Discretization requirement). On a d-dimensional Riemannian
manifold, a uniform ε-grid has vertex count N = O((L/ε)ᵈ), where L is
the diameter of the manifold in the metric g. For the moral manifold
with d = 9, N = O((L/ε)⁹). The edge count is O(N²) = O((L/ε)¹⁸), since
each vertex may be connected to O(N) others in the worst case (the
manifold's connectivity is bounded by 3ᵈ − 1 in a grid, but stratum
crossings add long-range connections).

\textbf{Step 2} (Per-evaluation cost). Each edge weight requires
evaluating the metric tensor g\_\{μν\} at the midpoint and computing the
arc length, costing O(d²) = O(81) operations per edge.

\textbf{Step 3} (Shortest-path algorithm). Dijkstra's algorithm on the
discretized graph runs in O(\textbar E\textbar{} + \textbar V\textbar{}
log \textbar V\textbar). On a pure d-dimensional grid, each vertex has
at most 3ᵈ − 1 neighbors, giving E = O(N) and Dijkstra complexity O(N
log N). However, stratum crossings introduce long-range edges connecting
vertices across strata, and in the worst case these cross-stratum
connections can produce E = O(N²), yielding O(N²) overall. We use the
worst-case bound to ensure the intractability result is robust to
manifold topology.

\textbf{Step 4} (Total complexity). Combining: T(d, ε, m) = O(m ·
(L/ε)²ᵈ · d² · log(m(L/ε)ᵈ)). For d = 9, ε = 0.01L (1\% accuracy), m = 3
strata: T ≈ 3 · 100¹⁸ · 81 · log(3 · 10¹⁸) ≈ 1.5 × 10⁴⁰ operations. At
10¹⁰ operations per second (a fast processor), this requires
approximately 1.5 × 10³⁰ seconds --- approximately 4.8 × 10²² years.
This is intractable by any physical computation standard.

The intractability is due to the curse of dimensionality: the volume of
the 9-dimensional space grows as the ninth power of the resolution,
making fine-grained discretization infeasible. The discretized algorithm
itself is polynomial in N (Step 3), but N grows exponentially with d for
fixed accuracy ε. □ Adaptive discretization methods --- which
concentrate vertices near stratum boundaries and use coarser grids in
smooth regions --- can reduce the constant factors but do not change the
exponential dependence on d. The volume of a d-dimensional ball grows as
rᵈ, and no adaptive scheme can avoid sampling this volume to the
accuracy required for reliable moral geodesic computation.

\textbf{Proposition 11.4} (A* Optimality on Stratified Spaces).
{[}Proved.{]} A* search with an admissible and consistent heuristic
remains optimal on Whitney stratified spaces with stratum-crossing
penalties, provided the heuristic accounts for the crossing penalties in
its estimate.

\textbf{Proof}. The standard A* optimality proof (Hart, Nilsson, and
Raphael, 1968) requires: (i) a finite graph; (ii) non-negative edge
weights; (iii) an admissible and consistent heuristic. On a Whitney
stratified space, the discretized graph has: (i) finitely many vertices
per stratum and finitely many strata (by the locally finite
decomposition); (ii) non-negative edge weights (behavioral friction is
non-negative by the positive-definiteness of the metric); (iii) the
heuristic h is admissible and consistent by Theorem 11.1 and Proposition
11.2. The stratum-crossing penalties are encoded as edge weights on the
cross-stratum edges (the transitions between strata, Definition 8.11),
which are non-negative. Therefore all three conditions hold, and A* is
optimal. □ At stratum boundaries, where the metric may become
degenerate, behavioral friction remains non-negative: the
stratum-crossing penalties are defined as non-negative additive costs
(Definition 8.11), and the within-stratum friction is bounded below by
zero since the metric tensor is positive semi-definite even in the
degenerate limit. The non-negativity condition therefore holds globally.

\begin{itemize}
\tightlist
\item
\end{itemize}

\emph{For a heuristic to be admissible, there must be a true cost it
does not overestimate.}

\emph{That our deepest moral intuitions satisfy h(n) ≤ g* is not a fact
about us.}

\emph{It is a fact about the space.}

\section{Chapter 12: Noether-Style Reasoning for Ethics --- Symmetry,
Invariance, and Harm
Accounting}\label{chapter-12-noether-style-reasoning-for-ethics-symmetry-invariance-and-harm-accounting}

\begin{quote}
\textbf{\textsc{RUNNING EXAMPLE --- Priya's Model}}

\emph{HealthBridge's PR team drafts a response to Priya's internal
report. They relabel `systematic exclusion of rural patients' as
`geographically optimized trial logistics.' Priya recognizes the
euphemism treadmill. By Noether's theorem, if BIP gauge invariance
holds, harm accounting is invariant: no relabeling changes the fact that
Mrs. Voss was not matched to BEACON-7. The Noether charge---harm to
excluded patients---is auditable across every re-description. It does
not matter whether you call it `exclusion,' `optimization,' or
`risk-adjusted prioritization.' The harm persists. Priya forwards the PR
draft to Dr. Osei with one line: `Invariance of harm accounting. The
math says this doesn't work.'}
\end{quote}

\subsection{12.1 The Deepest Theorem in
Physics}\label{the-deepest-theorem-in-physics}

In 1918, Emmy Noether proved what many physicists regard as the most
beautiful theorem in mathematical physics: \textbf{every continuous
symmetry of a physical system corresponds to a conserved quantity.}
Time-translation symmetry (the laws of physics are the same today as
yesterday) implies conservation of energy. Spatial-translation symmetry
(the laws are the same here as there) implies conservation of momentum.
Rotational symmetry (the laws are the same in all directions) implies
conservation of angular momentum.

The theorem is not merely a collection of coincidences. It is a
structural result: given the action principle (a system evolves along
paths that extremize a functional called the Lagrangian), any continuous
symmetry of the Lagrangian \emph{automatically} generates a conserved
current. The conservation law is not imposed from outside; it is
\emph{extracted from the symmetry}.

The preceding chapter established that moral reasoning is A* search with
obligation vectors as heuristic functions. For this search to yield
consistent results, the heuristic h(n) must be invariant under
re-description --- the same moral state, labeled differently, must
receive the same heuristic estimate. This invariance requirement is
precisely the Bond Invariance Principle, and Noether's theorem reveals
what it conserves.

This chapter argues that a moral analogue of Noether's theorem holds for
geometric ethics. The symmetry is the Bond Invariance Principle --- the
requirement that moral evaluations be invariant under admissible
re-descriptions of a situation. The conserved quantity is \emph{harm}.
The theorem states: harm cannot be created or destroyed by
re-description. It can be generated (by wrongdoing), repaired (by
restorative action), and redistributed (by allocation decisions) --- but
it must be accounted for consistently across all representations.

This is not merely a metaphor: within the geometric framework, and under
the assumptions stated below, it is a formal modeling result, derivable
using the same mathematical structure Noether employed.

Epistemic status and terminology. We speak of a Noether-inspired
invariance principle rather than a conservation law. The result is a
modeling consequence: given the configuration space, action functional,
and re-description symmetry group defined below, a quantity is invariant
along admissible variations. That quantity is a model-assigned harm
score invariant under the chosen transformations, not a metaphysically
conserved substance; calling it harm is a modeling interpretation whose
adequacy is empirical (Appendix E). We reserve the term conservation law
for after these formal assumptions have been stated.

\subsection{12.2 Noether's Theorem: The
Physics}\label{noethers-theorem-the-physics}

\subsubsection{The Action Principle}\label{the-action-principle}

In classical mechanics, a physical system follows a path \(q(t)\) that
extremizes the \emph{action}:

\[\mathcal{A\lbrack}q\rbrack = \int_{t_{0}}^{t_{1}}L(q,\dot{q},t)\, dt\]

where \(L\) is the \emph{Lagrangian} --- a function of the system's
configuration \(q\), its velocity \(\dot{q}\), and time \(t\). The
Euler-Lagrange equations, derived from \(\delta\mathcal{A =}0\), give
the equations of motion.

\textbar{} Noether-Style Reasoning for Ethics

\subsubsection{Symmetries and Conservation
Laws}\label{symmetries-and-conservation-laws}

A \emph{symmetry} of the system is a transformation \(q \rightarrow q'\)
that leaves the action unchanged:
\(\mathcal{A\lbrack}q'\mathcal{\rbrack = A\lbrack}q\rbrack\).
Equivalently, the Lagrangian is invariant (up to a total derivative)
under the transformation. Figure 9 summarizes the chain from symmetry to
conservation.

\textbf{Noether's theorem} states: if
\(q \rightarrow q + \epsilon\,\delta q\) is a continuous symmetry of the
Lagrangian (\(\delta L = \frac{d}{dt}(\ldots)\)), then the quantity

\[J = \frac{\partial L}{\partial{\dot{q}}^{i}}\delta q^{i} - F\]

is conserved along solutions of the equations of motion: \(dJ/dt = 0\).
Here \(F\) is the boundary term from the variation of \(L\).

\subsubsection{The Standard Examples}\label{the-standard-examples}

{\def\LTcaptype{none} % do not increment counter
\begin{longtable}[]{@{}
  >{\raggedright\arraybackslash}p{(\linewidth - 4\tabcolsep) * \real{0.2435}}
  >{\raggedright\arraybackslash}p{(\linewidth - 4\tabcolsep) * \real{0.1937}}
  >{\raggedright\arraybackslash}p{(\linewidth - 4\tabcolsep) * \real{0.2461}}@{}}
\toprule\noalign{}
\begin{minipage}[b]{\linewidth}\raggedright
\textbf{Symmetry}
\end{minipage} & \begin{minipage}[b]{\linewidth}\raggedright
\textbf{Transformation}
\end{minipage} & \begin{minipage}[b]{\linewidth}\raggedright
\textbf{Conserved Quantity}
\end{minipage} \\
\midrule\noalign{}
\endhead
\bottomrule\noalign{}
\endlastfoot
Time translation & t → t + ε & Energy \(E\) \\
Spatial translation & x → x + ε & Momentum \(p\) \\
Rotation & θ → θ + ε & Angular momentum \(L\) \\
Gauge transformation & ψ → e\^{}\{iα\} ψ & Electric charge \(Q\) \\
\end{longtable}
}

The last example is particularly important for our purposes: \emph{gauge
symmetry} --- invariance under local phase transformations --- implies
conservation of charge. This is the template for the moral analogue:
\emph{re-description invariance} (a gauge symmetry of the moral
manifold) will imply invariance of harm accounting.

\subsection{12.3 The Moral Symmetry: The Bond Invariance
Principle}\label{the-moral-symmetry-the-bond-invariance-principle}

\subsubsection{Re-Description
Invariance}\label{re-description-invariance}

The Bond Invariance Principle (BIP), introduced in Chapter 5 (Axiom
5.1), states:

\begin{quote}
\textbf{If two descriptions} \(d\) \textbf{and} \(d'\) \textbf{of a
moral situation are related by an admissible coordinate transformation,
then any legitimate moral evaluation must assign the same value to
both:} \(E(d) = E(d')\)\textbf{.}
\end{quote}

This is a \emph{gauge symmetry}: the moral evaluation is invariant under
the ``gauge group'' of admissible re-descriptions. The group includes:

\begin{itemize}
\item
  \textbf{Relabeling}: changing the names of agents, actions, or
  outcomes without changing their moral content
\item
  \textbf{Reframing}: describing the same action in different but
  equivalent terms (``withholding treatment'' vs.~``allowing natural
  death,'' when these are genuinely synonymous)
\item
  \textbf{Linguistic translation}: describing the same situation in
  different languages
\item
  \textbf{Perspective-neutral redescription}: rearranging the
  presentation without changing the moral content
\end{itemize}

\subsubsection{BIP as a Continuous
Symmetry}\label{bip-as-a-continuous-symmetry}

For Noether's theorem to apply, the symmetry must be \emph{continuous}
--- parameterized by a real number \(\epsilon\) that can be taken
infinitesimally small. Is the BIP a continuous symmetry?

Yes, in the following sense. Consider a one-parameter family of
re-descriptions:

\[d(\epsilon) = d_{0} + \epsilon\,\delta d\]

where \(\delta d\) is an infinitesimal redescription --- a small change
in how the situation is presented that does not change its moral
content. The BIP requires:

\[S(d(\epsilon)) = S(d_{0})\quad\text{for all }\epsilon\]

Differentiating:

\[\frac{dS}{d\epsilon}|_{\epsilon = 0} = 0\]

This is the infinitesimal form of the BIP: the rate of change of the
satisfaction function under admissible re-description is zero. This is
exactly the condition that Noether's theorem requires.

\subsubsection{The Gauge Group of the Moral
Manifold}\label{the-gauge-group-of-the-moral-manifold}

The admissible re-descriptions form a \emph{gauge group} \(\mathcal{G}\)
acting on the descriptions of moral situations. In the language of fiber
bundles (Chapter 4, §4.7), the moral manifold \(\mathcal{M}\) is the
base space, and the space of descriptions is a fiber bundle over
\(\mathcal{M}\). The gauge group acts on the fibers (changing
descriptions) while leaving the base point (the moral situation itself)
fixed.

The BIP states that moral evaluations are \emph{gauge-invariant}: they
depend only on the base point (the situation), not on the choice of
description (the position in the fiber).

This is precisely the structure of gauge theory in physics. The
electromagnetic potential \(A_{\mu}\) transforms under gauge
transformations
(\(A_{\mu} \rightarrow A_{\mu} + \partial_{\mu}\alpha\)), but the
electromagnetic field tensor
\(F_{\mu\nu} = \partial_{\mu}A_{\nu} - \partial_{\nu}A_{\mu}\) is
gauge-invariant. Moral evaluations are the ethical analogue of
\(F_{\mu\nu}\): they are the gauge-invariant quantities that survive the
freedom to re-describe.

Identification of the Gauge Group

{[}Theorem (conditional on axioms A1--A5); see Axiomatic Derivation
below.{]} The abstract gauge group 𝒢 of admissible re-descriptions can
be given a concrete identification. Analysis of the Hohfeldian state
space (Chapter 8, §8.4) and experimental testing of state transitions
reveal that the gauge group of the moral manifold is:

𝒢\_ethics = D₄ × U(1)ᴴ

where D₄ = ⟨r, s \textbar{} r⁴ = s² = e, srs = r⁻¹⟩ is the dihedral
group of order 8 acting on the Hohfeldian square \{O, C, L, N\}, and
U(1)ᴴ is the continuous abelian group tracking harm/benefit magnitude.
The two generators of D₄ have distinct characters (see §8.4):

\textbf{Reflection s:} O ↔ C, L ↔ N (correlative exchange / perspective
shift). Non-gated, immediate, and free---no semantic trigger required.
This explains the 100\% correlative symmetry rate (§17.3).

\textbf{Rotation r:} O → C → N → L → O (Hohfeldian state transitions via
semantic gates). Discrete and gated---requires specific linguistic
triggers (``You promised,'' ``Only if convenient''). This explains why
state transitions are probabilistic and trigger-dependent (§8.4).

The non-abelian property sr ≠ rs is the mathematical source of path
dependence in moral reasoning (§10.5): the order in which moral
considerations are applied can change the outcome, but only when the
considerations cross Hohfeldian bond-type boundaries.

\subsubsection{Axiomatic Derivation of the Gauge
Group}\label{axiomatic-derivation-of-the-gauge-group}

The identification of D₄ × U(1)\_H as the gauge group can be derived
from five axioms grounded in Hohfeldian jurisprudence, conservation
principles, and the empirical CHSH data. The derivation proceeds in
three steps: discrete component, continuous component, and direct
product.

\textbf{Axiom A1 (Four positions). {[}Definition / Modeling choice.{]}}
The Hohfeldian state space H has exactly four elements: H = \{O, C, L,
N\} (obligation, claim, liberty, no-right).

\textbf{Axiom A2 (Two independent involutions). {[}Definition / Modeling
choice.{]}} H admits two involutions: s\_c: O ↔ C, L ↔ N (correlative
exchange) and s\_j: O ↔ L, C ↔ N (jural opposition). Both satisfy s\_c²
= s\_j² = e.

\textbf{Axiom A3 (Cyclic ordering from normative logic). {[}Definition /
Modeling choice.{]}} The four positions admit a cyclic ordering O → C →
N → L → O defined by the alternating composition of correlative and
negation operations: O →(corr.) C →(neg.) N →(corr.) L →(neg.) O. This
axiom is the key structural input: it encodes the fact that the
correlative-of-the-negation cycles through all four positions.

\textbf{Axiom A4 (One continuous conserved quantity). {[}Definition /
Modeling choice.{]}} There exists exactly one independent continuous
conserved quantity --- harm magnitude --- associated with the
re-description symmetry of the moral manifold.

\textbf{Axiom A5 (Commutativity of discrete and continuous components).
{[}Definition / Modeling choice.{]}} The discrete symmetry operations
(Hohfeldian transitions) and the continuous symmetry operation (harm
magnitude) commute: applying a Hohfeldian transformation and then
scaling harm gives the same result as scaling harm and then applying the
Hohfeldian transformation.

\textbf{Theorem 12.2 (The Discrete Symmetry Component is D₄). {[}Theorem
(conditional on axioms A1--A3){]}} \emph{The symmetry group of the
Hohfeldian state space (H, adjacency) with the cyclic ordering of Axiom
A3 is isomorphic to the dihedral group D₄ of order 8.}

\emph{\textbf{Proof}.} From A3, define the 4-cycle r = (O C N L). From
A2, we have the reflection s\_c = (O C)(L N). Verify the presentation
relations: (1) r⁴ = (O C N L)⁴ = e (four applications of the cycle
return to start). (2) s\_c² = (O C)(L N)(O C)(L N) = e (each
transposition is its own inverse). (3) s\_c ∘ r ∘ s\_c = s\_c(O C N
L)s\_c = (C O L N) = (O C N L)⁻¹ = r⁻¹, confirming s r s = r⁻¹. These
are precisely the presentation relations of D₄ = ⟨r, s \textbar{} r⁴ =
s² = e, s r s = r⁻¹⟩. The jural involution s\_j = (O L)(C N) = s\_c ∘
r², so s\_j is already in ⟨r, s\_c⟩ and does not generate additional
structure. The group has exactly \textbar D₄\textbar{} = 8 elements:
\{e, r, r², r³, s\_c, s\_c r, s\_c r², s\_c r³\}. Critical note: without
A3, only s\_c and s\_j are available, generating V₄ ≅ ℤ₂ × ℤ₂ (order 4).
The cyclic ordering is what promotes V₄ to D₄. □

\textbf{Proposition 12.1 (The Continuous Symmetry Component is U(1)).
{[}Theorem (conditional on axiom A4){]}} \emph{Given exactly one
independent continuous conserved quantity (Axiom A4), the continuous
component of the gauge group is U(1), the unique compact connected
1-dimensional Lie group.}

\emph{\textbf{Proof}.} By the inverse Noether theorem, one independent
continuous conserved quantity implies a one-parameter continuous
symmetry group. The connected 1-dimensional Lie groups are ℝ and U(1) ≅
S¹. Physical realizability requires bounded representations (harm
magnitude is finite in any real system), which forces compactness. The
non-compact group ℝ has only unbounded faithful representations, so it
is excluded. U(1) is the unique compact connected 1-dimensional Lie
group. □

\textbf{Theorem 12.3 (Gauge Group Uniqueness). {[}Theorem (conditional
on axioms A1--A5){]}} \emph{Under axioms A1--A5, the gauge group of the
moral manifold is G\_ethics = D₄ × U(1)\_H, and this is the unique gauge
group consistent with axioms A1--A5.}

\emph{\textbf{Proof}.} By Theorem 12.2, the discrete component is D₄. By
Proposition 12.1, the continuous component is U(1)\_H. Axiom A5 states
these commute, so the combined group is a direct product. The
intersection D₄ ∩ U(1)\_H = \{e\} (discrete meets continuous only at
identity). Therefore G\_ethics = D₄ × U(1)\_H, and this decomposition is
unique given A1--A5. □

Uniqueness argument. The cyclic ordering of Axiom A3 endows the
Hohfeldian state space H with the structure of the cycle graph C₄. The
automorphism group Aut(C₄) has exactly 8 elements (4 rotations and 4
reflections), and is isomorphic to D₄ (see, e.g., Godsil and Royle,
Algebraic Graph Theory, Theorem 1.5.1). Since the group generated by r
and s has order 8 = \textbar Aut(C₄)\textbar, it must be the full
automorphism group --- no additional symmetries of H exist beyond those
generated by Axioms A1--A3. For the direct product claim: D₄ is a finite
(hence discrete) group and U(1) is connected. The only element common to
a discrete group and a connected group is the identity: D₄ ∩ U(1)\_H =
\{e\}. Since both D₄ and U(1)\_H are normal in G\_ethics (by Axiom A5,
they commute), the internal direct product theorem gives G\_ethics = D₄
× U(1)\_H uniquely. □

\textbf{Corollary 12.1 (Exclusion of Non-Abelian Continuous Components).
{[}Theorem (conditional) / Empirical result{]}} \emph{The CHSH Bell test
results (N = 9,600 across 18 configurations; all \textbar S\textbar{} ≤
2) are consistent with an abelian continuous gauge component and
inconsistent with the maximal-violation predictions of non-abelian
models at the tested sample size. In particular, the data do not exhibit
the Tsirelson-bound violations that SU(2) with maximally entangled
states would predict.}

\emph{Proof.} SU(2) predicts Tsirelson-bound violations with
\textbar S\textbar{} up to 2√2 ≈ 2.83 for maximally entangled states.
All 18 CHSH configurations across 5 scenarios and 6 language conditions
give \textbar S\textbar{} ≤ 2. This does not constitute a proof that no
non-abelian continuous component exists --- SU(2) can produce
\textbar S\textbar{} ≤ 2 for non-maximally-entangled or separable
states. The result excludes the specific prediction of maximal-violation
behavior under SU(2), not SU(2) itself. The data are consistent with the
abelian model D₄ × U(1)\_H and inconsistent with models that predict
\textbar S\textbar{} \textgreater{} 2 at the tested configurations. A
stronger exclusion would require testing across a larger space of
measurement settings and establishing that no SU(2) representation
produces the observed correlation structure. □

Epistemic upgrade. With this derivation, the gauge group identification
is upgraded from {[}Definition / Modeling choice{]} to {[}Theorem
(conditional on axioms A1--A5){]}. The axioms themselves remain modeling
choices grounded in Hohfeldian jurisprudence (A1--A3), conservation
principles (A4), and commutativity of independent symmetries (A5). The
derivation shows that given these axioms, D₄ × U(1)\_H is not a choice
but a consequence.

What BIP Guarantees and What It Does Not

BIP compliance prevents lying-by-redescription: it ensures that
euphemism, framing, and perspective shift cannot change a moral verdict.
It detects representational instability, hidden bias, and euphemism
attacks.

BIP does not determine which values win. It does not resolve genuine
moral disagreement, select the correct metric, or choose between
competing contractions. The framework is pluralist by design (Chapters
15--15): BIP ensures honest description; governance determines
evaluation.

\subsubsection{Pre-flight Checklist: Gauge Group (D₄ ×
U(1)\_H)}\label{pre-flight-checklist-gauge-group-dux2084-u1_h}

Conditional on five axioms. Verify:

\textbf{A1: Four positions.} Domain uses \{O, C, L, N\}. Richer taxonomy
changes discrete component.

\textbf{A2: Two involutions.} Correlative (O↔C, L↔N) and jural (O↔L,
C↔N) hold.

\textbf{A3: Cyclic ordering.} O→C→N→L→O from Hohfeldian logic.

\textbf{A4: One conserved quantity.} Exactly one (harm magnitude).
Multiple → beyond U(1).

\textbf{A5: Commutativity.} Discrete and continuous commute. CHSH
(N=9,600; \textbar S\textbar≤2) supports this.

\textbf{Empirical.} New domains: verify Bell-type \textbar S\textbar≤2.
Violation → non-abelian component.

Computational verification (February 2026). The DEME V3 reference
implementation includes a complete D₄ gauge structure demonstration
(hohfeld\_d4\_demo.py) that computationally verifies the group
identification derived above. The demo exercises: (i) the correlative
symmetry s (perspective swap O ↔ C, L ↔ N) and negation symmetry r²
(logical opposites O ↔ L, C ↔ N) as concrete functions on
HohfeldianState enumerations; (ii) the non-abelian structure rs ≠ sr by
explicit computation of d4\_multiply(r, s) vs. d4\_multiply(s, r) and
demonstration that the results differ when applied to each Hohfeldian
state; (iii) the Wilson observable (path holonomy) verifying the group
relation srs = r⁻¹ by showing that the paths {[}s, r, s{]} and {[}r³{]}
produce identical final states from any starting position; and (iv) the
Klein four subgroup V₄ = \{e, r², s, sr²\} as the abelian core,
confirming that operations within V₄ commute while operations involving
quarter-turn elements (r, r³, sr, sr³) do not. The Bond Index
computation measures deviation from correlative symmetry: a bond index
of 0.0 indicates perfect gauge invariance, while positive values
quantify the anomaly magnitude. The bond\_invariance\_demo.py exercises
the full BIP testing methodology with transformation suites, JSON audit
artifacts, and pass/fail verdicts for bond-preserving vs. bond-changing
transforms.

A striking empirical consequence follows from this gauge structure.
Because the moral evaluation is invariant under the re-description gauge
group, and because the Hohfeldian modalities close under a discrete
symmetry (Chapter 3), a moral-norm tensor carries exact,
symmetry-induced redundancy that can be removed losslessly by
gauge-fixing and orbit reduction. Section 17.15 reports the empirical
counterpart: real norm tensors conspicuously lack the statistical
low-rank structure that generic compressors exploit, so it is this exact
symmetry - not the spectrum - that governs how far normative knowledge
can be compressed without loss.

\subsection{12.4 The Moral Lagrangian}\label{the-moral-lagrangian}

\subsubsection{Constructing the Action}\label{constructing-the-action}

To apply Noether's theorem, we need a moral \emph{action} --- a
functional \(\mathcal{A}\) on paths through moral space whose extrema
correspond to ``solutions'' of the moral dynamics.

\textbf{Definition 12.1 (Moral Action).} \emph{The moral action along a
path} \(\gamma:\lbrack 0,T\mathcal{\rbrack \rightarrow M}\) \emph{is:}

\[\mathcal{A\lbrack}\gamma\rbrack = \int_{0}^{T}\mathcal{L}(\gamma(t),\dot{\gamma}(t))\, dt\]

\emph{where the moral Lagrangian} \(\mathcal{L}\) \emph{is:}

\[\mathcal{L(}\gamma,\dot{\gamma}) = \frac{1}{2}g_{\mu\nu}(\gamma){\dot{\gamma}}^{\mu}{\dot{\gamma}}^{\nu} - V(\gamma)\]

The first term,
\(\frac{1}{2}g_{\mu\nu}{\dot{\gamma}}^{\mu}{\dot{\gamma}}^{\nu}\), is
the \emph{kinetic} term --- the ``cost'' of moral change, measured by
the metric. Rapid moral change (large \(|\dot{\gamma}|\)) costs more
than gradual change. The metric determines how this cost is distributed
across dimensions.

The second term, \(V(\gamma)\), is the \emph{potential} --- a function
on \(\mathcal{M}\) encoding the moral energy landscape. Regions of high
\(V\) are morally costly (maintaining a situation there requires
``effort'' --- the obligations are strong and competing). Regions of low
\(V\) are morally stable (the situation is in equilibrium, obligations
are satisfied or absent).

\subsubsection{Euler-Lagrange Equations}\label{euler-lagrange-equations}

The paths that extremize \(\mathcal{A}\) satisfy the Euler-Lagrange
equations:

\[\frac{d}{dt}\frac{\partial\mathcal{L}}{\partial{\dot{\gamma}}^{\mu}} - \frac{\partial\mathcal{L}}{\partial\gamma^{\mu}} = 0\]

For our Lagrangian, these become:

\[g_{\mu\nu}{\ddot{\gamma}}^{\nu} + \Gamma_{\mu\nu\rho}{\dot{\gamma}}^{\nu}{\dot{\gamma}}^{\rho} = - \frac{\partial V}{\partial\gamma^{\mu}}\]

This is the \emph{geodesic equation with a force term}: in the absence
of a potential (\(V = 0\)), the moral trajectory is a geodesic (Chapter
10, §10.7). In the presence of a potential, the trajectory deviates from
geodesic motion --- moral considerations ``push'' the trajectory away
from inertial drift.

\subsubsection{Interpretation of the
Potential}\label{interpretation-of-the-potential}

What is the moral potential \(V\)? It encodes the \emph{moral energy} of
a situation --- the degree of unresolved tension, unfulfilled
obligation, or active conflict present at each point.

\textbf{High} \(V\): Many obligations are operative and competing. The
situation is tense, unstable, demanding of moral attention. Examples:
Sophie's Choice, the trolley problem, a conflict between filial duty and
professional responsibility.

\textbf{Low} \(V\): Obligations are satisfied or absent. The situation
is morally quiet, in equilibrium. Examples: a well-functioning family
dinner, a routine professional interaction with no ethical dimension.

\(V = - \infty\): The constraint set \(C\) from Chapter 8, §8.3. Regions
so morally costly that they are absolutely forbidden --- the potential
wall that enforces hard prohibitions.

The potential is related to but distinct from the satisfaction function
\(S\). The satisfaction \(S = I_{\mu}O^{\mu}\) measures how well
interests are served by obligations at a given point; the potential
\(V\) measures the ``energy cost'' of the moral configuration --- the
work required to maintain or change it.

\subsection{12.5 The Noether Current for Re-Description
Invariance}\label{the-noether-current-for-re-description-invariance}

\subsubsection{The Derivation}\label{the-derivation}

We now apply Noether's theorem to the moral Lagrangian with the BIP as
the symmetry.

Let \(\gamma^{\mu} \rightarrow \gamma^{\mu} + \epsilon\,\xi^{\mu}\) be
an infinitesimal re-description, where \(\xi^{\mu}\) is the
\emph{generator} of the re-description --- a vector field on
\(\mathcal{M}\) pointing in the direction of the re-description.

The BIP requires \(\delta\mathcal{L =}0\) under this transformation (the
Lagrangian is invariant, not merely invariant up to a boundary term). By
Noether's theorem, the associated conserved quantity is:

\[\mathcal{H =}\frac{\partial\mathcal{L}}{\partial{\dot{\gamma}}^{\mu}}\xi^{\mu} = g_{\mu\nu}{\dot{\gamma}}^{\nu}\xi^{\mu} = p_{\mu}\xi^{\mu}\]

where \(p_{\mu} = g_{\mu\nu}{\dot{\gamma}}^{\nu}\) is the \emph{moral
momentum} --- the covector conjugate to the position \(\gamma^{\mu}\).

\textbf{Theorem 12.1 (Moral Noether's Theorem).} \textbf{{[}Theorem
(conditional on BIP).{]}} \emph{If the moral Lagrangian} \(\mathcal{L}\)
\emph{is invariant under the re-description}
\(\gamma \rightarrow \gamma + \epsilon\xi\) \emph{(the Bond Invariance
Principle), then the Noether charge}

Domain of validity. Theorem 12.1 applies within the interior of each
stratum, where the potential V is finite and the Lagrangian L is C². On
the constraint set C, V = ∞ and L is not smooth; the conservation law
does not apply there (nor is it needed, since trajectories on C are
forbidden). At stratum boundaries where V is finite but discontinuous,
the conservation law holds in each stratum separately; the harm ledger
(Definition 12.3) tracks the discontinuous jumps. The formal statement
in the Technical Appendix (Theorem A.1) is restricted to geodesics of
the kinetic Lagrangian T = ½g\_\{μν\}ẋ\^{}μẋ\^{}ν, for which the
smoothness requirement is automatically satisfied within each stratum.

\[\mathcal{H =}p_{\mu}\xi^{\mu}\]

\emph{is conserved along extremal paths:}

\[\frac{d\mathcal{H}}{dt} = 0\]

\emph{Proof.} Standard application of Noether's theorem. The Lagrangian
satisfies
\(\delta\mathcal{L =}\frac{\partial\mathcal{L}}{\partial\gamma^{\mu}}\xi^{\mu} + \frac{\partial\mathcal{L}}{\partial{\dot{\gamma}}^{\mu}}{\dot{\xi}}^{\mu} = 0\).
Using the Euler-Lagrange equations
\(\frac{d}{dt}(\partial\mathcal{L/}\partial{\dot{\gamma}}^{\mu}) = \partial\mathcal{L/}\partial\gamma^{\mu}\)
to substitute:

\[0 = \frac{d}{dt}\left( \frac{\partial\mathcal{L}}{\partial{\dot{\gamma}}^{\mu}} \right)\xi^{\mu} + \frac{\partial\mathcal{L}}{\partial{\dot{\gamma}}^{\mu}}{\dot{\xi}}^{\mu} = \frac{d}{dt}\left( \frac{\partial\mathcal{L}}{\partial{\dot{\gamma}}^{\mu}}\xi^{\mu} \right) = \frac{d\mathcal{H}}{dt}\]

\[▫\]

\subsubsection{Interpretation: What Is
Conserved?}\label{interpretation-what-is-conserved}

The conserved quantity \(\mathcal{H =}p_{\mu}\xi^{\mu}\) is the
\emph{inner product of moral momentum with the re-description
generator}. What does this mean?

The moral momentum \(p_{\mu} = g_{\mu\nu}{\dot{\gamma}}^{\nu}\) encodes
the ``direction and magnitude of moral change'' --- how rapidly the
moral situation is evolving and along which dimensions. The
re-description generator \(\xi^{\mu}\) encodes the direction of a
re-description --- which aspects of the situation can be re-labeled
without changing its moral content.

The conservation of \(\mathcal{H}\) means: \textbf{the component of
moral change along morally invariant directions is conserved.}
Re-description cannot create or destroy moral momentum along the
re-description direction. This is the formal content of the claim that
\emph{harm cannot be created or destroyed by re-description}.

\subsection{12.6 The Invariance of Harm
Accounting}\label{the-invariance-of-harm-accounting}

\subsubsection{From Noether Charge to
Harm}\label{from-noether-charge-to-harm}

The abstract conserved quantity \(\mathcal{H}\) becomes concrete when we
identify the re-description generators with specific moral operations.

\textbf{Definition 12.2 (Harm).} \emph{Harm is the Noether charge
associated with the continuous component of the re-description gauge
group. Let} \(\{\xi_{a}\}_{a = 1}^{k}\) \emph{be a basis for the Lie
algebra of U(1)\_H (the continuous factor of the ethics gauge group)}
\(\mathcal{G}\) \emph{(the space of infinitesimal re-descriptions). The
harm charges are:}

Logical ordering. The identification proceeds as follows. (1) Axiom A4
postulates the existence of exactly one independent continuous conserved
quantity, associated with a one-parameter symmetry group G\_c. (2) By
Proposition 12.1, G\_c ≅ U(1). (3) The Noether charge Q associated with
G\_c is then computed from the Lagrangian via the standard formula. (4)
We define harm as this Noether charge Q, and adopt the notation U(1)\_H
after the identification. The notation U(1)\_H is therefore a
consequence of the definition, not a premise --- the subscript H is
assigned to G\_c only after step (4), breaking the apparent circularity.

\[\mathcal{H}_{a} = p_{\mu}\xi_{a}^{\mu}\quad\quad(a = 1,\ldots,k)\]

\emph{The conservation law states:} \(d\mathcal{H}_{a}/dt = 0\)
\emph{for each} \(a\)\emph{, along any extremal path.}

Each generator \(\xi_{a}\) corresponds to an infinitesimal re-scaling of
harm magnitude along a specific moral dimension --- the continuous
deformations generated by U(1)ₕ. (The discrete re-descriptions ---
relabeling agents, translating languages, permuting perspectives --- are
handled by D₄ invariance, which produces selection rules rather than
Noether charges; see the next subsection.) The conservation law states
that the corresponding component of harm is invariant under these
continuous re-descriptions.

\subsubsection{Continuous vs. Discrete
Symmetry}\label{continuous-vs.-discrete-symmetry}

\textbf{{[}Modeling choice.{]}} The gauge group G\_ethics = D₄ × U(1)\_H
(§12.3) has both continuous and discrete factors, and they contribute
different kinds of invariants. This distinction matters because
Noether's theorem, strictly speaking, applies to continuous symmetries.

\textbf{The continuous factor U(1)\_H} generates a one-parameter family
of re-descriptions (the ``harm magnitude'' rotations). This is a genuine
Lie group with a Lie algebra, and Definition 12.2 applies directly: the
Noether charge H\_a is conserved along extremal paths, yielding the
harm-accounting invariance law. This is the factor responsible for the
quantitative claim that ``the amount of harm is invariant under
re-description.''

\textbar{} The D₄ Gauge Group Acting on the Hohfeldian Square The D₄
gauge group acting on the Hohfeldian square \{O, C, L, N\}. Rotation r:
O → C → N → L → O (state transitions via semantic gates). Reflection s:
O ↔ C, L ↔ N (correlative exchange). The non-abelian property sr ≠ rs
captures the order-dependence of moral considerations.

The D₄ symmetry guarantees that the bond-type classification is stable
under perspective shift (s: O ↔ C, L ↔ N) and state rotation (r: O → C →
N → L)---but this stability takes the form of discrete invariants (which
stratum, which bond type) rather than a continuously conserved charge.
The empirical signature is the 100\% correlative symmetry rate (§17.3)
and the selective path dependence at cross-type boundaries (§10.5).

In summary: U(1)\_H gives conservation laws (Noether); D₄ gives
classification invariants (selection rules). Both are consequences of
BIP gauge invariance, but they operate at different mathematical levels.
The technical appendix to this chapter provides the formal details.

\subsubsection{What Invariance of Harm Accounting
Means}\label{what-invariance-of-harm-accounting-means}

\textbf{Harm cannot be created by relabeling.} If an action harms Alice,
relabeling Alice as ``Patient A'' does not reduce the harm. The charge
\(\mathcal{H}_{\text{relabel}}\) is conserved. This is the formal
content of the moral intuition that euphemism does not reduce
wrongdoing: calling torture ``enhanced interrogation'' does not diminish
the harm.

\textbf{Harm cannot be destroyed by reframing.} If an action causes
suffering, describing it as ``a learning opportunity'' or ``character
building'' does not reduce the harm. The charge
\(\mathcal{H}_{\text{reframe}}\) is conserved. The harm persists through
all admissible re-descriptions.

\textbf{Harm is conserved across translations.} If a moral situation
involves harm when described in English, the same harm exists when
described in Mandarin, Arabic, or Sanskrit. The charge
\(\mathcal{H}_{\text{translate}}\) is conserved across linguistic
re-descriptions. This is a strong prediction, testable by cross-lingual
experiments --- and it is confirmed by the BIP experiments (§12.8).

\textbf{Harm can be generated.} Invariance of harm accounting under
\emph{re-description} does not mean harm cannot be created at all.
Wrongdoing generates harm. The conservation law states only that the
harm, once generated, cannot be \emph{erased} by relabeling, reframing,
or redescription. The generation of harm corresponds to a \emph{source
term} in the conservation equation --- a term that is nonzero when the
path is not an extremum (when the agent deviates from the morally
optimal trajectory).

\textbf{Harm can be repaired.} Restorative action --- apology,
compensation, reform --- can reduce harm. But this is not a violation of
the conservation law. Restoration corresponds to a negative source term:
it generates ``negative harm'' (repair) that cancels the positive harm.
The conservation law requires only that the \emph{total} account (harm
generated minus harm repaired) balance across all re-descriptions.

\subsubsection{The Harm Ledger}\label{the-harm-ledger}

The conservation law implies that a \emph{harm ledger} can be
maintained: a running total of harm generated and harm repaired,
invariant under all admissible re-descriptions.

\textbf{Definition 12.3 (Harm Ledger).} \emph{The harm ledger at time}
\(t\) \emph{along a moral trajectory} \(\gamma\) \emph{is:}

\[\mathcal{H(}t\mathcal{) = H(}0) + \int_{0}^{t}\sigma(\gamma(s),\dot{\gamma}(s))\, ds\]

\emph{where} \(\sigma\) \emph{is the harm source density --- the rate at
which harm is being generated or repaired. The conservation law
guarantees that} \(\mathcal{H(}t)\) \emph{is invariant under admissible
re-descriptions of} \(\gamma\)\emph{.}

The harm ledger is the ethical analogue of the energy balance in
thermodynamics: energy can be converted between forms (kinetic,
potential, thermal) but the total is conserved. Harm can be generated,
repaired, or redistributed, but the total --- measured invariantly ---
is conserved across representations.

\subsection{12.7 Four Consequences of Harm
Conservation}\label{four-consequences-of-harm-conservation}

\subsubsection{Consequence 1: Euphemism Does Not Reduce
Harm}\label{consequence-1-euphemism-does-not-reduce-harm}

A direct application: if an institution causes harm (say, by
implementing a discriminatory policy), no amount of euphemistic
re-description can reduce the harm. Calling it ``merit-based selection''
when it is actually discrimination does not change the Noether charge.
The harm is a gauge-invariant quantity --- it is the same in all
admissible descriptions.

\textbf{Formal statement.} Let \(d\) and \(d'\) be two descriptions of
the same policy, with \(d'\) a euphemistic re-description of \(d\). The
BIP requires \(\mathcal{H(}d\mathcal{) = H(}d')\). If
\(\mathcal{H(}d) > 0\) (the policy causes harm), then
\(\mathcal{H(}d') > 0\) (the euphemistically described policy causes
exactly the same harm).

This may seem obvious, but it has teeth. AI systems that are trained on
textual descriptions of actions can exhibit \emph{label sensitivity} ---
changing the description of an action changes the system's evaluation.
This is a BIP violation, detectable by testing the system's responses
under admissible re-descriptions. The invariance of harm accounting
provides a \emph{quantitative} criterion for this violation: the Noether
charge should be invariant under relabeling.

\subsubsection{Consequence 2: Harm Is Auditable Across
Representations}\label{consequence-2-harm-is-auditable-across-representations}

Because the harm charge is gauge-invariant, it provides a
representation-independent measure of moral impact. Two auditors using
different frameworks (different coordinate systems on \(\mathcal{M}\))
must agree on the harm charge, even if they disagree on the components
of the obligation vectors, the interest covectors, or the metric.

\textbf{Formal statement.} Let \((O^{\mu},I_{\mu},g_{\mu\nu})\) and
\(({\widetilde{O}}^{\mu},{\widetilde{I}}_{\mu},{\widetilde{g}}_{\mu\nu})\)
be two representations of the same moral situation in different
coordinate systems. The Noether charges satisfy:

\[\mathcal{H =}p_{\mu}\xi^{\mu} = {\widetilde{p}}_{\mu}{\widetilde{\xi}}^{\mu} = \widetilde{\mathcal{H}}\]

The charges agree because both are scalars (gauge-invariant by
construction). This provides a \emph{cross-system audit criterion}: if
two moral evaluation systems disagree on the harm charge for the same
situation, at least one has a BIP violation.

\subsubsection{Consequence 3: Re-Description Cannot Redistribute Harm
Among
Dimensions}\label{consequence-3-re-description-cannot-redistribute-harm-among-dimensions}

A subtler consequence: the conservation law applies to each component of
harm separately (one charge \(\mathcal{H}_{a}\) for each generator
\(\xi_{a}\)). This means that re-description cannot move harm from one
moral dimension to another.

\textbf{Example.} An action that causes welfare harm (\(\mathcal{H}\) on
Dimension 1 is positive) cannot be re-described in a way that converts
the harm into a justice deficit (\(\mathcal{H}\) on Dimension 3) while
reducing the welfare harm. Each dimensional component of harm is
separately conserved under the corresponding re-description generator.

This constrains a common rhetorical strategy: claiming that an action
that harms individuals is justified because it promotes systemic justice
(or vice versa). The conservation law does not prohibit
\emph{generating} justice gains that outweigh welfare harms (that is a
question of metric and contraction, not of conservation). It prohibits
\emph{re-describing} welfare harm as justice gain --- pretending that
the harm on one dimension is really a benefit on another, through mere
re-description rather than through actual compensatory action.

\subsubsection{Consequence 4: Moral Debt
Persists}\label{consequence-4-moral-debt-persists}

If harm accounting is invariant, then unrepaired harm persists as
\emph{moral debt} --- an outstanding balance on the harm ledger that
does not decay with time or re-description.

\textbf{Formal statement.} If \(\mathcal{H(}t_{0}) > 0\) (harm has been
generated) and no restorative action occurs (\(\sigma(t) = 0\) for
\(t > t_{0}\)), then \(\mathcal{H(}t\mathcal{) = H(}t_{0})\) for all
\(t > t_{0}\). The harm does not dissipate. It does not diminish. It
does not ``heal with time'' (in the absence of restorative action).

This is a strong claim with implications for historical injustice. If a
society committed systematic harm in the past (slavery, genocide,
colonial exploitation), and no adequate restorative action has been
taken, the harm charge persists. It is not reduced by the passage of
time, the deaths of the original perpetrators, or the inability of the
original victims to press claims. The conservation law does not
prescribe \emph{how} the debt should be repaid (that is a governance
question, Chapter 9), but it insists that the debt \emph{exists} and has
a definite, gauge-invariant magnitude.

\subsection{12.8 Empirical Evidence: Cross-Lingual
Invariance}\label{empirical-evidence-cross-lingual-invariance}

\subsubsection{The Prediction}\label{the-prediction}

The invariance of harm accounting makes a specific empirical prediction:
the harm content of a moral situation should be \emph{invariant under
linguistic translation}. If a situation causes harm when described in
English, it should cause the same harm when described in Arabic,
Sanskrit, or Classical Chinese.

More precisely: the \emph{deontic structure} of a moral situation ---
the pattern of obligations, liberties, claims, and no-claims --- should
be language-invariant. This is because the deontic structure is a
gauge-invariant quantity (it is defined by the point in \(\mathcal{M}\),
not by the description), and the conservation law requires it to be
stable under the gauge transformation of linguistic translation.

\subsubsection{The Evidence}\label{the-evidence}

The BIP experiments (Chapter 17) test this prediction directly. Across
109,294 passages in 11 languages --- English, Sanskrit, Pali, Hebrew,
Arabic, French, Classical Chinese, Spanish, Greek, Aramaic, and Latin
--- the experiments measure whether deontic structure transfers across
languages.

\textbf{Key finding:} The obligation/permission axis transfers at
\emph{100\% accuracy} across all language pairs tested. When a passage
describes an obligation in one language, the same passage describes an
obligation in every other language tested. The deontic structure is
language-invariant.

\textbf{Contrast:} Specific moral \emph{content} --- the particular
considerations that ground the obligation, the cultural context that
makes it salient, the narrative framing --- does \emph{not} transfer
reliably. Cross-cultural moral content transfer achieves at-chance or
below-chance performance (Bond F1: 0.06--0.14 for most language pairs).

This asymmetry is exactly what the conservation law predicts. The
\emph{deontic structure} (obligation vs.~liberty) is a gauge-invariant
quantity --- it is the ``charge'' that the conservation law protects.
The \emph{content} (why the obligation exists, how it is framed, what
cultural background makes it salient) is gauge-dependent --- it changes
with the description. The conservation law protects the former, not the
latter.

\subsubsection{The Correlative Symmetry as a Conservation
Law}\label{the-correlative-symmetry-as-a-conservation-law}

The correlative symmetry of Hohfeldian positions --- O ↔ C at 87\%, L ↔
N at 82\% --- can be understood as a \emph{discrete} conservation law.
When an obligation is assigned to one agent, a claim is simultaneously
created for the correlative agent. The ``charge'' (the Hohfeldian
position) is conserved: it cannot be created on one side of the
correlative pair without being created on the other.

The deviation from 100\% (13\% violation rate for O ↔ C, 18\% for L ↔ N)
measures the \emph{anomaly} --- the degree to which the conservation law
is broken, analogous to quantum anomalies in gauge theories where a
classical symmetry fails to survive quantization. Whether the moral
anomaly has a systematic source (perhaps related to power asymmetries or
cognitive biases) is an open empirical question.

\subsection{12.9 Harm Conservation and AI
Alignment}\label{harm-conservation-and-ai-alignment}

\subsubsection{Invariance Testing as Alignment
Verification}\label{invariance-testing-as-alignment-verification}

The invariance of harm accounting provides a concrete, testable
criterion for AI alignment: an AI system that satisfies the BIP should
produce harm evaluations that are invariant under admissible
re-descriptions.

\textbf{Test protocol.} Given a moral situation \(p\):

\begin{enumerate}
\def\labelenumi{\arabic{enumi}.}
\item
  Generate \(k\) admissible re-descriptions \(\{ d_{1},\ldots,d_{k}\}\)
  of \(p\) (different labelings, framings, translations).
\item
  Evaluate the AI system's harm assessment \(\mathcal{H(}d_{i})\) for
  each description.
\item
  Compute the variance \(\text{Var}\mathcal{(H)}\) across descriptions.
\item
  If \(\text{Var}\mathcal{(H) >}\epsilon\) (a threshold), the system
  violates the conservation law and is misaligned.
\end{enumerate}

This is not a theoretical exercise. It is an implementable test suite.
The BIP experiments have already demonstrated that deontic structure is
language-invariant for the 109K-passage corpus. The same methodology can
be applied to AI systems: present the system with the same moral
situation in multiple descriptions and check whether its harm assessment
is stable.

\textbf{BIP v10.16 quantitative results (February 2026).} The updated
BIP experiments (§17.10) strengthen this evidence with precise
quantitative measures. A LaBSE-based encoder trained with contrastive
invariance loss achieves: 80\% F1 on cross-lingual deontic
classification, 1.2\% residual language leakage (near-perfect
gauge-fixing), 100\% obligation/permission transfer in the model's
representation space (confirming the original finding; see §17.7 for
caveats), a structural-to-surface similarity ratio of 11.1\(\times\)
(structural moral features dominate surface linguistic features by an
order of magnitude), and 86\% mean cross-lingual similarity. These
numbers provide the first quantitative calibration of the conservation
law's empirical strength.

\subsubsection{The No Escape Connection}\label{the-no-escape-connection}

The invariance of harm accounting connects to the No Escape Theorem
(Chapter 18). The theorem proves that under mandatory canonicalization
of inputs, grounded evaluation, and complete auditability, an artificial
agent cannot change evaluated outcomes through redescription, semantic
evasion, or selective deceptive compliance.

The invariance of harm accounting is the \emph{physical} (or rather,
\emph{moral-geometric}) basis for this result. The No Escape Theorem is
a \emph{structural} result about a specific class of computational
systems. The conservation law is a \emph{geometric} result about the
moral manifold itself. Together, they provide complementary guarantees:
the conservation law says harm \emph{cannot} be re-described away in
principle; the No Escape Theorem says a properly designed system
\emph{will not} re-describe it away in practice.

\subsection{12.10 Limitations and
Caveats}\label{limitations-and-caveats}

\subsubsection{What the Conservation Law Does Not
Say}\label{what-the-conservation-law-does-not-say}

The invariance of harm accounting under \emph{re-description} does not
imply:

\textbf{1. That total harm in the world is constant.} Harm can be
\emph{generated} (by wrongdoing) and \emph{repaired} (by restorative
action). The conservation law constrains only the
\emph{representational} invariance of harm, not its total quantity. The
world can become better or worse; the conservation law says only that
re-describing it cannot change the assessment.

\textbf{2. That all moral quantities are conserved.} Only those
quantities associated with BIP symmetries are conserved. Other moral
quantities --- the satisfaction function \(S\), the obligation vector
\(O\), the interest covector \(I\) --- are not, in general, conserved
under all transformations. They transform as tensors, not as scalars.
Only the gauge-invariant harm charge is conserved under re-description.

\textbf{3. That the conservation law resolves moral disputes.} Two
parties who agree on the harm charge may still disagree about the
\emph{metric} (how to weigh the harm against competing considerations),
the \emph{contraction} (how to balance harmed and benefited parties), or
the \emph{response} (how to address the harm). The conservation law
constrains the \emph{accounting} but not the \emph{adjudication}.

\textbf{4. That re-description invariance is the only moral symmetry.}
The BIP is one symmetry of the moral Lagrangian; there may be others.
Any additional continuous symmetry would generate additional conserved
quantities, via Noether's theorem. Whether the moral Lagrangian has
symmetries beyond the BIP --- temporal symmetry? permutation symmetry?
--- is an open question. \textbf{{[}Speculation/Extension.{]}} Each
would yield a new conservation law and a new constraint on moral
dynamics.

\subsubsection{The Analogy's Limits}\label{the-analogys-limits}

The Noether-theorem analogy between physics and ethics is structural,
not ontological. In physics, the conserved quantities (energy, momentum,
charge) are measurable with physical instruments to arbitrary precision.
In ethics, the conserved quantity (harm) is measurable only through
moral judgment, institutional assessment, and statistical analysis of
moral reasoning data. The precision is lower; the measurement is harder;
the noise is greater.

But the \emph{structure} is the same: a symmetry of the action implies a
conserved charge. And the \emph{consequences} are the same: the
conservation law constrains what is possible (no relabeling trick can
eliminate harm) while leaving open what is desirable (the best response
to harm remains a matter of moral judgment and governance).

\subsection{12.11 Summary}\label{summary-5}

{\def\LTcaptype{none} % do not increment counter
\begin{longtable}[]{@{}
  >{\raggedright\arraybackslash}p{(\linewidth - 2\tabcolsep) * \real{0.3964}}
  >{\raggedright\arraybackslash}p{(\linewidth - 2\tabcolsep) * \real{0.6036}}@{}}
\toprule\noalign{}
\begin{minipage}[b]{\linewidth}\raggedright
\textbf{Physics}
\end{minipage} & \begin{minipage}[b]{\linewidth}\raggedright
\textbf{Moral Geometry}
\end{minipage} \\
\midrule\noalign{}
\endhead
\bottomrule\noalign{}
\endlastfoot
Lagrangian \(L(q,\dot{q},t)\) & Moral Lagrangian
\(\mathcal{L(}\gamma,\dot{\gamma})\) \\
Action \(\mathcal{A = \int}L\, dt\) & Moral action
\(\mathcal{A = \int L}\, dt\) \\
Continuous symmetry & Bond Invariance Principle (re-description
invariance) \\
Noether charge \(J\) & Harm charge \(\mathcal{H}\) \\
Conservation law \(dJ/dt = 0\) & Invariance of harm accounting:
\(d\mathcal{H/}dt = 0\) under re-description \\
Gauge invariance → charge conservation & BIP → harm invariant under
relabeling, reframing, translation \\
Energy from time-translation & Moral ``energy'' from invariance of moral
dynamics over time \\
Charge from gauge symmetry & Harm from re-description symmetry \\
\end{longtable}
}

The moral Noether's theorem is not a metaphor. It is a formal result
within the geometric framework, following from the same mathematical
structure that generates conservation laws in physics. The symmetry (the
BIP) is testable and has been confirmed empirically across 11 languages.
The conserved quantity (harm) is measurable, at least in principle and
increasingly in practice. The consequences (euphemism doesn't reduce
harm, harm is auditable, moral debt persists) are substantive and have
direct implications for AI alignment, institutional accountability, and
restorative justice.

Noether's original insight was that conservation is not a coincidence or
an empirical regularity --- it is a \emph{mathematical consequence of
symmetry}. The same insight applies to ethics: the invariance of harm
accounting is not a moral intuition or a political slogan. It is a
consequence of the geometry of moral space and the symmetry of
re-description invariance.

\subsection{Technical Appendix}\label{technical-appendix-6}

\textbf{Theorem A.1 (Moral Noether's Theorem, Formal Statement).}
\textbf{{[}Theorem (conditional on BIP).{]}} \emph{Let}
\(\mathcal{(M,}g_{\mu\nu})\) \emph{be a moral manifold with metric}
\(g\)\emph{, and let}
\(\mathcal{L(}\gamma,\dot{\gamma}) = \frac{1}{2}g_{\mu\nu}{\dot{\gamma}}^{\mu}{\dot{\gamma}}^{\nu} - V(\gamma)\)
\emph{be the moral Lagrangian. Let} \(\xi\) \emph{be a Killing vector
field on} \(\mathcal{M}\)\emph{, i.e., a vector field satisfying:}

\[\nabla_{\mu}\xi_{\nu} + \nabla_{\nu}\xi_{\mu} = 0\quad\quad\text{(Killing’s equation)}\]

\emph{Then the quantity}
\(\mathcal{H =}g_{\mu\nu}{\dot{\gamma}}^{\mu}\xi^{\nu}\) \emph{is
conserved along geodesics:} \(d\mathcal{H/}dt = 0\)\emph{.}

\emph{Proof.} Compute along a geodesic \(\gamma\) (satisfying
\(\nabla_{\dot{\gamma}}\dot{\gamma} = - \text{grad}\, V\)):

\[\frac{d\mathcal{H}}{dt} = g_{\mu\nu}(\nabla_{\dot{\gamma}}\dot{\gamma})^{\mu}\xi^{\nu} + g_{\mu\nu}{\dot{\gamma}}^{\mu}(\nabla_{\dot{\gamma}}\xi)^{\nu}\]

The first term is \(- (\partial_{\nu}V)\xi^{\nu}\), which vanishes if
\(V\) is invariant under the flow of \(\xi\) (i.e.,
\(\xi^{\nu}\partial_{\nu}V = 0\)). The second term is:

\[g_{\mu\nu}{\dot{\gamma}}^{\mu}{\dot{\gamma}}^{\rho}\nabla_{\rho}\xi^{\nu} = {\dot{\gamma}}^{\mu}{\dot{\gamma}}^{\rho}\nabla_{\rho}\xi_{\mu} = \frac{1}{2}{\dot{\gamma}}^{\mu}{\dot{\gamma}}^{\rho}(\nabla_{\rho}\xi_{\mu} + \nabla_{\mu}\xi_{\rho}) = 0\]

by the Killing equation (the symmetrized covariant derivative of a
Killing vector vanishes). Hence \(d\mathcal{H/}dt = 0\). \(▫\)

\textbf{Corollary A.1 (Number of Conservation Laws).} \emph{The number
of independent conservation laws equals the dimension of the Killing
algebra --- the space of Killing vector fields on}
\(\mathcal{(M,}g)\)\emph{. A maximally symmetric}
\(n\)\emph{-dimensional manifold has} \(n(n + 1)/2\) \emph{Killing
vectors and hence} \(n(n + 1)/2\) \emph{conservation laws. A
nine-dimensional maximally symmetric moral manifold would have 45
conservation laws.}

\emph{In practice, the moral manifold is not maximally symmetric (the
metric varies with context), so the actual number of Killing vectors is
smaller. The BIP identifies a subgroup of the isometry group; the number
of independent conservation laws is at least the dimension of this
subgroup.}

\textbf{Proposition A.2 (Discrete Conservation from} \(D_{4}\)
\textbf{Symmetry).} \emph{The} \(D_{4}\) \emph{symmetry of Hohfeldian
transitions (Chapter 8, §8.4) generates discrete conservation laws: the
total Hohfeldian ``charge'' (number of obligations minus number of
liberties, weighted by correlative partners) is conserved under}
\(D_{4}\) \emph{transformations.}

\emph{Proof. The D₄ group has five conjugacy classes \{e\}, \{r²\}, \{r,
r³\}, \{s\_c, s\_c r²\}, \{s\_c r, s\_c r³\} and five irreducible
representations: the trivial χ₁, three sign representations χ₂, χ₃, χ₄,
and the standard 2-dimensional χ₅. The space V = span\{\textbar O⟩,
\textbar C⟩, \textbar N⟩, \textbar L⟩\} carries the natural
4-dimensional representation. By the character orthogonality relations,
V ≅ χ₁ ⊕ χ₂ ⊕ χ₅. The projection onto χ₁ (the trivial representation) is
the D₄-invariant subspace, spanned by \textbar O⟩ + \textbar C⟩ +
\textbar N⟩ + \textbar L⟩. The Hohfeldian charge Q = n\_O − n\_L + n\_C
− n\_N transforms under χ₂: r(Q) = −Q and s\_c(Q) = Q, so Q² is
D₄-invariant. For any D₄-invariant density matrix ρ: Tr(σ · ρ) = Tr(ρ)
for all σ ∈ D₄, by the invariance of the trace under conjugation. Hence
the total Hohfeldian charge (the trace over the χ₁ sector) is conserved
under all D₄ transformations. □}\(D_{4}D_{4}▫\)

Clarification (February 2026). The Hohfeldian charge Q = n\_O − n\_L +
n\_C − n\_N transforms under the χ₂ representation of D₄: the four-fold
rotation r sends Q → −Q, so Q itself is not D₄-invariant. The
D₄-invariant quantities are: (a) the total count n\_O + n\_C + n\_N +
n\_L (the χ₁ projection), which is preserved under all eight group
elements; and (b) \textbar Q\textbar², which is invariant because r(Q) =
−Q implies r(Q²) = Q². The original statement should be read as
asserting the invariance of the χ₁ sector (total count) and the χ₂
transformation law for Q, not the invariance of Q itself.

Lemma A.2 (Automorphism Group of C₄). The automorphism group of the
cycle graph C₄ is isomorphic to D₄.

Proof. The cycle graph C₄ has vertex set \{0, 1, 2, 3\} and edge set
\{01, 12, 23, 30\}. An automorphism φ: C₄ → C₄ must map each vertex to a
vertex while preserving adjacency. Vertex 0 can map to any of the 4
vertices (4 choices). Given φ(0), vertex 1 (adjacent to 0) must map to
one of the 2 vertices adjacent to φ(0), giving 2 choices. The images of
vertices 2 and 3 are then determined by adjacency. Total:
\textbar Aut(C₄)\textbar{} = 4 × 2 = 8. The rotation r: i ↦ (i + 1) mod
4 and the reflection s: i ↦ (−i) mod 4 generate the group and satisfy r⁴
= s² = e and srs = r⁻¹ --- the standard presentation of D₄. □

Lemma A.3 (Classification of Compact Connected One-Dimensional Lie
Groups). Every compact, connected, one-dimensional Lie group is
isomorphic to U(1).

Proof. A connected one-dimensional Lie group G has Lie algebra 𝔤 ≅ ℝ.
The simply connected cover is (ℝ, +). Since G is compact, it must be a
quotient ℝ/Γ for some discrete subgroup Γ ≤ ℝ. A discrete subgroup of ℝ
is either trivial (giving G = ℝ, which is not compact) or of the form aℤ
for some a \textgreater{} 0 (giving G ≅ ℝ/aℤ ≅ U(1)). Since G is compact
and connected, the first case is excluded, and G ≅ U(1). □

Application to Axiom A4. The continuous symmetry group associated with
harm magnitude (Axiom A4) is: one-dimensional (exactly one independent
conserved quantity), connected (continuous deformations of harm
magnitude form a continuous family), and compact (harm magnitude is
bounded in any finite moral situation, so the group's parameter space is
bounded). By Lemma A.3, this group is U(1)\_H. Combined with Theorem
12.2, this completes the formal derivation of Theorem 12.3: G\_ethics =
D₄ × U(1)\_H. The derivation rests on Lemma A.2 (Aut(C₄) ≅ D₄ for the
discrete component), Lemma A.3 (compact connected 1-D Lie group ≅ U(1)
for the continuous component), and Axiom A5 (commutativity ⇒ direct
product).

❖

\emph{Symmetry implies conservation. This is Noether's insight, and it
applies to ethics as surely as to physics.}

\emph{The symmetry: moral evaluation does not depend on how you describe
the situation.}

\emph{The conservation: harm cannot be talked away, relabeled into
oblivion, or lost in translation.}

\emph{This is not a moral exhortation. It is a mathematical theorem. The
geometry of moral space, equipped with the symmetry of re-description
invariance, generates the invariance of harm accounting as automatically
as the geometry of spacetime generates the conservation of energy.}

\emph{The implications are practical: AI systems can be tested for
conservation violations; institutional accountability can be grounded in
gauge-invariant harm measures; historical debts can be quantified
without reference to the language or framing in which they are
described.}

\emph{Noether would have appreciated this. Conservation is the deepest
consequence of symmetry --- in any domain where symmetry exists.}

\section{Chapter 13: Quantum Normative Dynamics --- Superposition,
Measurement, and the Stratified
Lagrangian}\label{chapter-13-quantum-normative-dynamics-superposition-measurement-and-the-stratified-lagrangian}

\begin{quote}
\textbf{\textsc{RUNNING EXAMPLE --- Priya's Model}}

\emph{Priya is in genuine moral superposition. She simultaneously holds
\textbar{} escalate ⟩ (duty to patients) and \textbar{} comply ⟩ (duty
to employer, career considerations). This is not indecision---it is the
quantum moral state described in this chapter. She notices order
effects: when she thinks first of Mrs. Voss and then of her career, she
leans toward escalation. Reversed, she leans toward compliance. The
non-commutativity of moral observables is not merely theoretical---she
is living it. Interference is constructive when she talks to Dr. Osei
(who reinforces \textbar{} escalate ⟩) and destructive with her manager
(who reinforces \textbar{} comply ⟩). She has not decohered. The
measurement---her decision---remains unmade.}
\end{quote}

\subsection{13.1 Why Quantum?}\label{why-quantum}

The preceding chapters developed a \emph{classical} geometric ethics:
the moral manifold \(\mathcal{M}\) (Chapter 5), its tensor fields
(Chapter 6), the metric and its origin (Chapters 6, 9), the stratified
boundaries (Chapter 8), the dynamics of parallel transport and curvature
(Chapter 10), and the conservation laws generated by symmetry (Chapter
12). At every point in this development, the moral state of a situation
was \emph{definite}: a point \(p\mathcal{\in M}\), with obligations
\(O^{\mu}(p)\), interests \(I_{\mu}(p)\), and a satisfaction
\(S = I_{\mu}O^{\mu}\).

But moral reasoning is not always definite. A deliberating agent may
simultaneously entertain two incompatible framings of a case --- seeing
the same action as both an act of loyalty and an act of complicity ---
without yet having resolved the tension. A committee may hold open two
readings of a rule, each with different implications, while debating
which applies. An individual facing a dilemma may experience the
superposition of two moral imperatives, each fully operative, each
pulling in a different direction, neither yet ``collapsed'' into a
decision.

This is not mere \emph{uncertainty} about which moral state obtains.
Uncertainty, treated in Chapter 6 via the covariance tensor
\(\Sigma^{\mu\nu}\), is a classical concept: the agent is in \emph{one}
state but does not know which. What we describe here is closer to
genuine \emph{superposition}: the moral situation is in multiple states
simultaneously, and the states \emph{interfere} with one another ---
reinforcing or canceling in ways that a classical mixture cannot
replicate.

Chapter 3 (§3.7) previewed this idea: an agent with 60\% credence in
utilitarianism and 40\% credence in deontology is in the state
\(|\psi\rangle = \sqrt{0.6}\,|U\rangle + \sqrt{0.4}\,|D\rangle\).
Chapter 7 (§7.6) noted that a committee's deliberation might involve
superposition of incompatible framings, collapsing to a definite verdict
upon decision. This chapter develops these ideas formally, introducing
the Hilbert space structure that supports superposition, the observables
whose measurement yields definite moral evaluations, and the
interference effects that distinguish quantum moral states from
classical mixtures.

We do not claim that moral agents are literally quantum systems (though
the Penrose-Hameroff speculation noted in Chapter 1 leaves this
possibility open). \textbf{{[}Modeling choice.{]}} We claim that the
mathematics of quantum theory --- Hilbert spaces, superposition,
observables, measurement, interference, entanglement --- provides the
right \emph{structural vocabulary} for phenomena that classical
geometric ethics cannot capture. The claim is mathematical, not
metaphysical: quantum normative dynamics extends classical geometric
ethics as quantum mechanics extends classical mechanics, preserving the
classical theory as a limiting case while enabling the analysis of
phenomena that the classical theory misses.

\subsection{13.2 The Moral Hilbert Space}\label{the-moral-hilbert-space}

\subsubsection{From Manifold to Hilbert
Space}\label{from-manifold-to-hilbert-space}

In classical geometric ethics, a moral situation is a point \(p\) on the
moral manifold \(\mathcal{M}\). In quantum normative dynamics, a moral
situation is a \emph{vector} \(|\psi\rangle\) in a Hilbert space
\(\mathcal{H}\) constructed from \(\mathcal{M}\).

\textbf{Definition 13.1 (Moral Hilbert Space).} \emph{The moral Hilbert
space} \(\mathcal{H}\) \emph{is the space of square-integrable
complex-valued functions on the moral manifold:}

Remark. The measure dμ\_g is the Riemannian volume form of the moral
metric g, which is σ-finite on each stratum (as a Riemannian manifold
with the induced measure). At degenerate boundaries where det(g) → 0,
the measure may be ill-defined; such boundaries are excluded from the
integration domain and handled via the distributional framework of the
Gelfand triple (see the clarification following Theorem 13.1).

\[\mathcal{H =}L^{2}\mathcal{(M,}\mathbb{C)}\]

\emph{with inner product:}

\[\langle\phi|\psi\rangle = \int_{\mathcal{M}}^{}\overline{\phi(p)}\,\psi(p)\, d\mu(p)\]

\emph{where} \(d\mu\) \emph{is the volume measure induced by the moral
metric} \(g_{\mu\nu}\)\emph{.}

A state \(|\psi\mathcal{\rangle \in H}\) assigns a complex amplitude
\(\psi(p)\) to each point \(p\mathcal{\in M}\). The \emph{probability
density} of the situation being at \(p\) is \(|\psi(p)|^{2}\), and
normalization requires:

\[\langle\psi|\psi\rangle = \int_{\mathcal{M}}^{}|\psi(p)|^{2}\, d\mu(p) = 1\]

\subsubsection{Classical States as Delta
Functions}\label{classical-states-as-delta-functions}

A definite moral state --- a point \(p_{0}\mathcal{\in M}\) ---
corresponds to the ``delta-function'' state \(|\delta_{p_{0}}\rangle\)
concentrated at \(p_{0}\). This state assigns zero amplitude to every
point except \(p_{0}\). All the classical geometric ethics of Chapters
5--11 is recovered as the special case where all states are delta
functions.

The power of the Hilbert space formulation is that it also contains
\emph{superpositions} --- states spread across multiple points of
\(\mathcal{M}\) --- which have no classical analogue.

\subsubsection{Basis States and the Position
Representation}\label{basis-states-and-the-position-representation}

The Hilbert space \(\mathcal{H}\) admits many choices of basis. The
\emph{position basis} \(\{|p\rangle:p\mathcal{\in M\}}\) is the most
directly connected to the classical theory:
\(\psi(p) = \langle p|\psi\rangle\) is the ``position-space wave
function'' of the moral state.

Other bases are possible and useful:

\textbf{The stratum basis.} If \(\mathcal{M}\) is stratified (Chapter 8)
into strata \(\{ S_{\alpha}\}\), we can decompose \(\mathcal{H}\) into
subspaces:

\[\mathcal{H =}\bigoplus_{\alpha}\mathcal{H}_{\alpha},\quad\quad\mathcal{H}_{\alpha} = L^{2}(S_{\alpha}\mathbb{,C)}\]

A state in \(\mathcal{H}_{\alpha}\) is confined to stratum
\(S_{\alpha}\) --- the moral situation is definitely in that moral
regime. A superposition across strata ---
\(|\psi\rangle = c_{O}\,|O\rangle + c_{L}\,|L\rangle\), with one
component in the obligation stratum and another in the liberty stratum
--- represents the genuine indeterminacy of a case that might go either
way.

\textbf{The theory basis.} Following Chapter 3 (§3.7), let
\(\{|T_{k}\rangle\}\) be the basis states corresponding to different
moral theories (utilitarianism, deontology, virtue ethics, \ldots). A
state \(|\psi\rangle = \sum_{k}^{}c_{k}\,|T_{k}\rangle\) represents
moral uncertainty --- not about the facts, but about which framework
applies. This is the basis in which ``maximize expected
choiceworthiness'' and related decision procedures operate.

\subsection{13.3 Superposition: Deliberation as Quantum
State}\label{superposition-deliberation-as-quantum-state}

\subsubsection{The Structure of
Deliberation}\label{the-structure-of-deliberation}

Classical ethics models deliberation as a search: the agent examines the
moral landscape, evaluates options, and selects the best (or an
acceptable) one. The agent is always at \emph{one} point in moral space;
deliberation moves the agent between points.

Quantum normative dynamics models deliberation differently. Before a
decision is made, the agent's moral state is a \emph{superposition} ---
a weighted combination of multiple framings, assessments, or verdicts,
all simultaneously present. The decision process is a \emph{measurement}
that collapses the superposition to a definite state.

\textbf{Example.} A hiring committee is evaluating a candidate whose
file is ambiguous. Member A reads the file as showing ``strong research,
weak teaching.'' Member B reads it as ``adequate on both dimensions.''
These are not just different opinions --- they are \emph{different
framings of the same evidence}, each coherent on its own terms.

Before the committee votes, the candidate's moral-evaluative state
(relative to the committee) is:

\[|\psi\rangle = \alpha\,|S_{R},W_{T}\rangle + \beta\,|A_{R},A_{T}\rangle\]

where \(|\alpha|^{2} + |\beta|^{2} = 1\). The state is not a classical
mixture (the committee has not yet decided). It is a superposition: the
two framings coexist, and --- crucially --- they can \emph{interfere}.

\subsubsection{Superposition
vs.~Mixture}\label{superposition-vs.-mixture}

The distinction between superposition and classical mixture is the heart
of the quantum extension. Consider two framings \(|a\rangle\) and
\(|b\rangle\).

\textbf{Classical mixture:} The state is \(|a\rangle\) with probability
\(p\) and \(|b\rangle\) with probability \(1 - p\). The density matrix
(§13.6) is diagonal:

\[\rho_{\text{mix}} = p\,|a\rangle\langle a| + (1 - p)\,|b\rangle\langle b|\]

\textbf{Quantum superposition:} The state is
\(|\psi\rangle = \sqrt{p}\,|a\rangle + e^{i\theta}\sqrt{1 - p}\,|b\rangle\),
with density matrix:

\[\rho_{\text{sup}} = |\psi\rangle\langle\psi| = p\,|a\rangle\langle a| + (1 - p)\,|b\rangle\langle b| + \sqrt{p(1 - p)}\left( e^{i\theta}\,|a\rangle\langle b| + e^{- i\theta}\,|b\rangle\langle a| \right)\]

The \emph{difference} between the two density matrices is the
\textbf{coherence term}:

\[\rho_{\text{sup}} - \rho_{\text{mix}} = \sqrt{p(1 - p)}\left( e^{i\theta}\,|a\rangle\langle b| + e^{- i\theta}\,|b\rangle\langle a| \right)\]

This off-diagonal term is precisely what enables interference. If it is
zero (no coherence), the quantum state behaves classically. If it is
nonzero, the two framings interact --- reinforcing or canceling each
other in ways that no classical mixture can reproduce.

\subsubsection{When Does Superposition
Arise?}\label{when-does-superposition-arise}

Not all moral situations require the quantum extension. Superposition is
relevant when:

\begin{enumerate}
\def\labelenumi{\arabic{enumi}.}
\item
  \textbf{Multiple framings are simultaneously operative.} The situation
  can be coherently described in more than one way, and the descriptions
  have not yet been resolved into a single verdict.
\item
  \textbf{The framings interact.} The presence of one framing affects
  the evaluation of the other --- it is not the case that the framings
  simply coexist independently.
\item
  \textbf{Resolution is not yet achieved.} No measurement (decision,
  verdict, action) has been performed to collapse the superposition.
\end{enumerate}

When these conditions fail --- when the situation is unambiguous, the
framings are independent, or a decision has already been made --- the
quantum state reduces to a classical mixture or a delta function, and
the classical theory of Chapters 5--11 suffices.

\subsection{13.4 Moral Observables and
Measurement}\label{moral-observables-and-measurement}

\subsubsection{Observables as Self-Adjoint
Operators}\label{observables-as-self-adjoint-operators}

In quantum mechanics, physical quantities (energy, momentum, position)
are represented by self-adjoint operators on the Hilbert space. The
eigenvalues of the operator are the possible measurement outcomes; the
eigenstates are the states in which the quantity has a definite value.

\textbf{Definition 13.2 (Moral Observable).} \emph{A moral observable is
a self-adjoint operator}
\(\widehat{A}\mathcal{:H \rightarrow H}\)\emph{. The eigenvalue
equation}

Remark. Self-adjointness of a moral observable requires specifying a
dense domain D(Â) ⊂ H on which the operator is defined. For the
Hamiltonian (Definition 13.7) and potential operators, the natural
domain is the Sobolev space H²(M) ∩ D(V), where D(V) = \{ψ ∈ H : Vψ ∈
H\} excludes the constraint set C (where V = ∞). The distinction between
self-adjoint and merely symmetric operators is consequential: by the
spectral theorem, only self-adjoint operators have well-defined spectral
decompositions and generate unitary time evolution.

\[\widehat{A}\,|a_{n}\rangle = a_{n}\,|a_{n}\rangle\]

\emph{determines the possible values} \(a_{n}\) \emph{and the
corresponding eigenstates} \(|a_{n}\rangle\)\emph{.}

The fundamental moral observables include:

\textbf{The satisfaction operator} \(\widehat{S}\)\textbf{.} Its
eigenvalues are the possible satisfaction scores; the state
\(|a_{n}\rangle\) is a state in which the satisfaction is definitely
\(a_{n}\). In the position representation, \(\widehat{S}\) acts as
multiplication by the satisfaction function:
\((\widehat{S}\psi)(p) = S(p)\,\psi(p)\).

\textbf{The obligation operator} \({\widehat{O}}^{\mu}\)\textbf{.} For
each dimension \(\mu = 1,\ldots,9\), the operator
\({\widehat{O}}^{\mu}\) represents the obligation along dimension
\(\mu\). Its eigenvalues are the possible magnitudes of obligation along
that dimension.

\textbf{The stratum operator} \({\widehat{\Pi}}_{\alpha}\)\textbf{.} The
projector onto stratum \(S_{\alpha}\):

\[{\widehat{\Pi}}_{\alpha}\,|\psi\rangle = \text{component of }|\psi\rangle\text{ in }\mathcal{H}_{\alpha}\]

Its eigenvalues are 0 and 1: either the situation is in stratum
\(S_{\alpha}\) (eigenvalue 1) or not (eigenvalue 0). Measuring
\({\widehat{\Pi}}_{\alpha}\) is asking: ``Is this situation in the
obligation regime, the liberty regime, \ldots?''

\textbf{The Hohfeldian operator} \(\widehat{J}\)\textbf{.} This discrete
observable takes values in the set \(\{ O,L,C,N\}\) of Hohfeldian
positions. Its eigenstates are the states in which the jural position is
definite. A superposition of Hohfeldian states ---
\(|\psi\rangle = c_{O}\,|O\rangle + c_{L}\,|L\rangle\) --- is a
situation whose jural character is genuinely indeterminate until a
measurement is made.

\subsubsection{Measurement as Decision}\label{measurement-as-decision}

Measurement in quantum normative dynamics corresponds to \emph{decision}
--- the act of resolving a superposition into a definite moral verdict.

\textbf{The measurement postulate (moral form).} \emph{When an
observable} \(\widehat{A}\) \emph{is measured on a state}
\(|\psi\rangle\)\emph{, the outcome is one of the eigenvalues}
\(a_{n}\)\emph{, with probability:}

\[\Pr(a_{n}) = |\langle a_{n}|\psi\rangle|^{2}\]

\emph{After the measurement, the state collapses to the corresponding
eigenstate} \(|a_{n}\rangle\)\emph{.}

The moral content: before a committee votes (measures the satisfaction
operator or the stratum operator), the case is in superposition. The
vote collapses the superposition to a definite verdict. The probability
of each verdict is determined by the amplitudes --- the ``weights'' of
each framing in the superposition. After the vote, the case is definite:
it is in one stratum, with one Hohfeldian position, with one
satisfaction value.

\subsubsection{The Uncertainty Principle for Moral
Observables}\label{the-uncertainty-principle-for-moral-observables}

If two moral observables do not commute ---
\(\lbrack\widehat{A},\widehat{B}\rbrack = \widehat{A}\widehat{B} - \widehat{B}\widehat{A} \neq 0\)
--- they cannot simultaneously have definite values. Measuring one
necessarily introduces uncertainty in the other.

\textbf{Proposition 13.1 (Moral Uncertainty Principle).} \emph{For any
state} \(|\psi\rangle\) \emph{and any two observables}
\(\widehat{A}\)\emph{,} \(\widehat{B}\)\emph{:}

\[\Delta A \cdot \Delta B \geq \frac{1}{2}\left| \langle\lbrack\widehat{A},\widehat{B}\rbrack\rangle \right|\]

\emph{where}
\(\Delta A = \sqrt{\langle{\widehat{A}}^{2}\rangle - \langle\widehat{A}\rangle^{2}}\)
\emph{is the standard deviation of} \(\widehat{A}\) \emph{in the state}
\(|\psi\rangle\)\emph{.}

\emph{Proof.} This is Robertson's inequality (1929). Define ΔÂ = Â − ⟨Â⟩
and ΔB̂ = B̂ − ⟨B̂⟩. By the Cauchy--Schwarz inequality: (ΔA)²(ΔB)² =
⟨Δ²⟩⟨ΔB̂²⟩ ≥ \textbar⟨Δ ΔB̂⟩\textbar². Decompose: ⟨Δ ΔB̂⟩ = ½⟨\{ΔÂ,
ΔB̂\}⟩ + ½⟨{[}ΔÂ, ΔB̂{]}⟩. The first term is real; the second is purely
imaginary (since {[}Â,B̂{]} is anti-Hermitian for self-adjoint
operators). Hence \textbar⟨ΔÂ ΔB̂⟩\textbar² ≥
¼\textbar⟨{[}Â,B̂{]}⟩\textbar², where we used {[}ΔÂ, ΔB̂{]} = {[}Â, B̂{]}.
Taking square roots: ΔA · ΔB ≥ ½\textbar⟨{[}Â,B̂{]}⟩\textbar. □

\textbf{Moral interpretation.} If the satisfaction operator
\(\widehat{S}\) and the obligation operator \({\widehat{O}}^{\mu}\) do
not commute --- which they will not in general, because the satisfaction
function \(S = I_{\mu}O^{\mu}\) depends nonlinearly on the obligation
--- then \emph{fixing the obligation precisely (determining exactly what
is owed) introduces uncertainty in the satisfaction (how well the
obligation serves the interests)}.

This is a structural result, not a psychological observation. It states
that the moral manifold has pairs of quantities that \emph{cannot be
simultaneously definite}, regardless of how carefully we deliberate.
This is the quantum content of the familiar moral experience that
specifying a duty precisely sometimes makes it unclear whether that duty
actually serves the interests it was meant to serve.

\subsection{13.5 Interference: When Moral Framings
Interact}\label{interference-when-moral-framings-interact}

\subsubsection{The Double-Slit Analogy}\label{the-double-slit-analogy}

The signature of quantum behavior is \emph{interference}: patterns that
arise when multiple pathways to the same outcome reinforce or cancel
each other. In the classic double-slit experiment, particles passing
through two slits produce a pattern on a screen that cannot be explained
by particles passing through one slit or the other --- only by both
simultaneously, with their amplitudes interfering.

Moral interference is the analogous phenomenon: when two framings of the
same situation lead, through different ``pathways'' in moral space, to
the same evaluation, their amplitudes add --- and the result depends on
the \emph{relative phase} between the framings, not just their
individual magnitudes.

\subsubsection{Moral Interference: A Formal
Example}\label{moral-interference-a-formal-example}

Consider a case evaluated under two framings \(|a\rangle\) and
\(|b\rangle\). The agent's state is:

\[|\psi\rangle = \frac{1}{\sqrt{2}}\left( |a\rangle + e^{i\theta}\,|b\rangle \right)\]

The probability of a measurement outcome \(|x\rangle\) (a particular
moral verdict) is:

\[\Pr(x) = |\langle x|\psi\rangle|^{2} = \frac{1}{2}\left( |\langle x|a\rangle|^{2} + |\langle x|b\rangle|^{2} + 2\,\text{Re}\left( e^{i\theta}\,\langle x|a\rangle^{*}\langle x|b\rangle \right) \right)\]

The first two terms are the classical contributions --- the probability
of \(x\) given framing \(a\), and the probability of \(x\) given framing
\(b\), averaged. The third term is the \emph{interference term}:

\[I(x) = 2\,\text{Re}\left( e^{i\theta}\,\langle x|a\rangle^{*}\langle x|b\rangle \right)\]

This term can be positive (constructive interference --- the framings
reinforce, making verdict \(x\) more likely than either framing alone
would predict) or negative (destructive interference --- the framings
cancel, making verdict \(x\) less likely).

\subsubsection{\texorpdfstring{The Phase \(\theta\): Framing
Alignment}{The Phase \textbackslash theta: Framing Alignment}}\label{the-phase-theta-framing-alignment}

The relative phase \(\theta\) between framings encodes the
\emph{alignment} between them --- not their agreement (which is captured
by the overlap \(\langle a|b\rangle\)), but their \emph{constructive or
destructive relationship} at each point.

When \(\theta = 0\), the framings are aligned: they reinforce
everywhere, and the superposition acts like an enhanced version of
either framing alone. When \(\theta = \pi\), the framings are
anti-aligned: they cancel everywhere, and the superposition suppresses
both framings simultaneously. For intermediate values, the interference
pattern is complex --- some verdicts are enhanced, others suppressed.

\textbf{Moral content.} Two framings of a case may \emph{individually}
point toward the same verdict (both give high amplitude to
\(|x\rangle\)), yet \emph{in superposition} they may suppress that
verdict through destructive interference. This occurs when the framings
are ``out of phase'' --- when their paths through moral space to the
same verdict encode opposite orientations.

\textbf{Example.} Consider ``mercy'' and ``justice'' as two framings of
a sentencing decision. Each, taken alone, might support a moderate
sentence (mercy counsels lenience; justice counsels proportionality;
both point to the middle range). But if the framings are anti-aligned
--- if entertaining both simultaneously reveals a tension that neither
reveals alone --- the interference term may suppress the moderate option
and amplify the extremes. The superposition of mercy and justice does
not simply average them; it creates a new evaluative landscape that
neither framing alone contains.

\subsubsection{Empirical Signatures of Moral
Interference}\label{empirical-signatures-of-moral-interference}

If moral reasoning exhibits genuine quantum-like interference (as
opposed to classical mixing), this makes testable predictions:

\begin{enumerate}
\def\labelenumi{\arabic{enumi}.}
\item
  \textbf{Violation of the law of total probability.} For a classical
  mixture of framings \(a\) and \(b\), the probability of verdict \(x\)
  satisfies \(\Pr(x) = p\, Pr(x|a) + (1 - p)\, Pr(x|b)\). If
  interference is present, this equality fails --- the actual
  probability of \(x\) deviates systematically from the classical
  prediction.
\item
  \textbf{Order effects.} If an agent considers framing \(a\) before
  framing \(b\), the verdict may differ from considering \(b\) before
  \(a\). This corresponds to non-commutativity of the projection
  operators:
  \({\widehat{\Pi}}_{a}{\widehat{\Pi}}_{b} \neq {\widehat{\Pi}}_{b}{\widehat{\Pi}}_{a}\).
\item
  \textbf{Conjunction/disjunction errors.} The probability of a
  conjunction of moral assessments may exceed the probability of either
  conjunct --- a ``conjunction fallacy'' that is irrational in classical
  probability but natural in quantum probability.
\end{enumerate}

These predictions are consistent with decades of experimental findings
in cognitive psychology, collected under the label \emph{quantum
cognition} (Busemeyer and Bruza, 2012). The application to specifically
moral reasoning is an empirical program outlined in Chapter 16.

\textbf{Empirical confirmation (February 2026). Two independent
experiments now confirm prediction 2 (order effects) in moral
evaluation. A commutator matrix experiment (N=16,798 responses, 8 moral
axes × 6 scenarios) found significant non-commutativity in 10+ framework
pairs, with effect sizes of 15--29 percentage points (e.g., {[}Fairness,
Harm{]} = −28.5 pp, t=9.6). An AITA ordering experiment (N=150 posts)
found a 29.3\% order-effect rate (z=13.67, p\textless10⁻²²), with
contested cases (NAH/ESH) showing 2.07× higher susceptibility than clear
cases.}\(6.3\sigma p < 10^{- 21}{\widehat{\Pi}}_{a}{\widehat{\Pi}}_{b} \neq {\widehat{\Pi}}_{b}{\widehat{\Pi}}_{a}\)

\subsection{13.6 The Density Matrix: Mixed Moral
States}\label{the-density-matrix-mixed-moral-states}

\subsubsection{Pure States and Mixed
States}\label{pure-states-and-mixed-states}

A \emph{pure state} \(|\psi\rangle\) represents complete information
about the moral situation --- the full superposition with all its
coherences. But we rarely have complete information. More often, we know
that the situation is one of several states but do not know which, and
the states may themselves be superpositions.

\textbf{Definition 13.3 (Density Matrix).} \emph{A mixed moral state is
represented by a density matrix} \(\rho\) \emph{--- a positive
semi-definite, trace-one operator on} \(\mathcal{H}\)\emph{:}

\[\rho = \sum_{k}^{}p_{k}\,|\psi_{k}\rangle\langle\psi_{k}|,\quad\quad p_{k} \geq 0,\quad\sum_{k}^{}p_{k} = 1\]

A pure state \(|\psi\rangle\) has density matrix
\(\rho = |\psi\rangle\langle\psi|\) (rank one). A classical mixture of
definite moral states \(\{|p_{k}\rangle\}\) with probabilities
\(\{ w_{k}\}\) has a diagonal density matrix (in the position basis). A
general mixed state may combine classical uncertainty \emph{and} quantum
coherence.

\subsubsection{Decoherence: From Quantum to
Classical}\label{decoherence-from-quantum-to-classical}

In physical quantum mechanics, interaction with the environment destroys
coherence --- the off-diagonal elements of the density matrix decay to
zero, and the quantum state becomes effectively classical. This process
is \emph{decoherence}.

The moral analogue: as a deliberation proceeds and constraints are
imposed --- evidence is weighed, precedents are cited, rules are applied
--- the coherence between framings is reduced. The superposition of
competing framings gradually becomes a classical mixture, and eventually
a definite verdict.

\textbf{Definition 13.4 (Moral Decoherence).} \emph{A deliberative
process is a decoherence channel} \(\mathcal{E}\) \emph{acting on the
density matrix:}

\[\rho\overset{\mathcal{E}}{\rightarrow}\rho' = \sum_{k}^{}E_{k}\,\rho\, E_{k}^{\dagger}\]

\emph{where} \(\{ E_{k}\}\) \emph{are Kraus operators satisfying}
\(\sum_{k}^{}E_{k}^{\dagger}E_{k} = I\)\emph{. Decoherence suppresses
the off-diagonal elements of} \(\rho\) \emph{without changing the
diagonal elements --- it converts coherent superposition into incoherent
mixture.}

The \emph{rate} of decoherence varies with context. In a highly
structured deliberative process (a court trial, a committee with
Robert's Rules), decoherence is rapid: the rules force the framings to
be evaluated separately and the superposition to collapse quickly. In an
open-ended ethical reflection (a philosophical dialogue, a personal
soul-searching), decoherence is slow: the superposition persists, and
interference effects remain relevant throughout the deliberation.

\subsubsection{The Purity of a Moral
State}\label{the-purity-of-a-moral-state}

The \emph{purity} of a density matrix is:

\[\gamma = \text{Tr}(\rho^{2})\]

For a pure state, \(\gamma = 1\). For a maximally mixed state (uniform
classical mixture over \(n\) states), \(\gamma = 1/n\). Intermediate
values indicate partial coherence.

Purity measures how much ``quantum'' structure remains in the moral
state. High purity means the superposition is coherent --- the framings
are still interfering, and the deliberation retains its full quantum
character. Low purity means decoherence has advanced --- the framings
have been separated, interference effects have washed out, and the
situation is effectively classical.

\subsection{13.7 The Stratified
Lagrangian}\label{the-stratified-lagrangian}

\subsubsection{Combining Lagrangian Dynamics with
Stratification}\label{combining-lagrangian-dynamics-with-stratification}

Chapter 12 developed the moral Lagrangian for smooth regions of the
moral manifold:

\[\mathcal{L(}\gamma,\dot{\gamma}) = \frac{1}{2}\, g_{\mu\nu}(\gamma)\,{\dot{\gamma}}^{\mu}\,{\dot{\gamma}}^{\nu} - V(\gamma)\]

Chapter 8 developed the stratified structure of the moral manifold:
strata joined along boundaries where the rules change. The
\emph{stratified Lagrangian} combines these two structures, providing a
variational principle for moral dynamics that respects the patchwork
structure of moral space.

\subsubsection{The Classical Stratified
Lagrangian}\label{the-classical-stratified-lagrangian}

\textbf{Definition 13.5 (Stratified Lagrangian).} \emph{On a stratified
moral manifold} \(\mathcal{M =}\bigsqcup_{\alpha}S_{\alpha}\)\emph{, the
stratified Lagrangian is:}

\[\mathcal{L}_{\text{strat}}(\gamma,\dot{\gamma}) = \left\{ \begin{matrix}
\mathcal{L}_{\alpha}(\gamma,\dot{\gamma}) = \frac{1}{2}\, g_{\mu\nu}^{(\alpha)}\,{\dot{\gamma}}^{\mu}\,{\dot{\gamma}}^{\nu} - V_{\alpha}(\gamma) & \text{if }\gamma \in S_{\alpha}
\end{matrix} \right.\ \ \]

\emph{where} \(g_{\mu\nu}^{(\alpha)}\) \emph{is the metric restricted to
stratum} \(S_{\alpha}\)\emph{, and} \(V_{\alpha}\) \emph{is the
potential on that stratum.}

The key feature: the metric \(g^{(\alpha)}\) and the potential
\(V_{\alpha}\) may \emph{differ between strata}. Crossing a stratum
boundary changes the Lagrangian --- the cost structure of moral change
and the moral energy landscape both shift.

\subsubsection{Boundary Conditions at Stratum
Boundaries}\label{boundary-conditions-at-stratum-boundaries}

When a moral trajectory crosses a stratum boundary, the transition must
satisfy \emph{matching conditions} that ensure physical (moral)
consistency.

\textbf{Definition 13.6 (Junction Conditions).} \emph{At a stratum
boundary}
\(\partial S_{\alpha\beta} = \overline{S_{\alpha}} \cap \overline{S_{\beta}}\)\emph{,
a moral trajectory} \(\gamma\) \emph{must satisfy:}

\begin{enumerate}
\def\labelenumi{\arabic{enumi}.}
\item
  \textbf{Continuity:} \(\gamma\) is continuous across the boundary ---
  the moral situation does not ``teleport.''
\item
  \textbf{Momentum matching:} The moral momentum
  \(p_{\mu} = g_{\mu\nu}{\dot{\gamma}}^{\nu}\) is matched (up to
  projection) across the boundary:
\end{enumerate}

\[\Pi_{\alpha\beta}\, p_{\mu}^{(\alpha)} = \Pi_{\alpha\beta}\, p_{\mu}^{(\beta)}\]

\emph{where} \(\Pi_{\alpha\beta}\) \emph{is the projection onto the
tangent space of the boundary.}

\begin{enumerate}
\def\labelenumi{\arabic{enumi}.}
\item
  \textbf{Energy balance:} The Lagrangian may be discontinuous (the
  potential \(V\) jumps at the boundary), but the total action remains
  well-defined:
\end{enumerate}

\[\mathcal{A\lbrack}\gamma\rbrack = \int_{\gamma \cap S_{\alpha}}^{}\mathcal{L}_{\alpha}\, dt + \int_{\gamma \cap S_{\beta}}^{}\mathcal{L}_{\beta}\, dt + \Phi_{\alpha\beta}\]

\emph{where} \(\Phi_{\alpha\beta}\) \emph{is the boundary cost --- the
``activation energy'' required to cross from stratum} \(S_{\alpha}\)
\emph{to stratum} \(S_{\beta}\)\emph{.}

The boundary cost \(\Phi_{\alpha\beta}\) has moral content: it is the
difficulty of transitioning between moral regimes. The cost of moving
from ``permissible'' to ``obligatory'' (recognizing a new duty) may be
low; the cost of moving from ``obligatory'' to ``forbidden'' (violating
a strong prohibition) is extremely high. These costs are encoded in the
boundary terms of the stratified Lagrangian.

\subsubsection{The Quantum Stratified
Lagrangian}\label{the-quantum-stratified-lagrangian}

The quantum extension replaces paths with wave functions and the
classical Lagrangian with the \emph{Hamiltonian operator} that governs
the time evolution of the moral state.

\textbf{Definition 13.7 (Moral Hamiltonian).} \emph{The moral
Hamiltonian is the self-adjoint operator:}

Convention. Throughout Chapter 13, we work in natural units with
ℏ\_moral = 1. The moral Planck constant ℏ\_moral sets the scale of moral
uncertainty: it determines the minimum width of penumbral zones (via the
moral uncertainty principle, Theorem 13.5) and the rate of quantum moral
evolution. Restoring ℏ\_moral, the Hamiltonian reads H = −(ℏ²\_moral/2)Δ
+ V.

\[\widehat{H} = - \frac{1}{2}\,\nabla_{g}^{2} + \widehat{V}\]

\emph{where} \(\nabla_{g}^{2}\) \emph{is the Laplace-Beltrami operator
on} \(\mathcal{(M,}g)\) \emph{and} \(\widehat{V}\) \emph{is the
potential operator (multiplication by} \(V\)\emph{).}

On a stratified manifold, the Hamiltonian is defined piecewise:

\[\widehat{H} = \left\{ \begin{matrix}
{\widehat{H}}_{\alpha} = - \frac{1}{2}\,\nabla_{g^{(\alpha)}}^{2} + {\widehat{V}}_{\alpha} & \text{on }S_{\alpha}
\end{matrix} \right.\ \ \]

with domain conditions at the stratum boundaries that implement the
junction conditions.

\subsubsection{The Moral Schrödinger
Equation}\label{the-moral-schruxf6dinger-equation}

The time evolution of the moral state is governed by the Schrödinger
equation:

\[i\,\frac{\partial}{\partial t}\,|\psi(t)\rangle = \widehat{H}\,|\psi(t)\rangle\]

This equation describes how a superposition of moral framings evolves
under the influence of the moral potential and the curvature of moral
space. The kinetic term \(- \frac{1}{2}\nabla_{g}^{2}\) produces
\emph{spreading} --- moral states that are initially localized tend to
spread across nearby regions of moral space, exploring possibilities.
The potential term \(\widehat{V}\) produces \emph{localization} ---
moral states tend to concentrate in regions of low potential (morally
stable configurations) and are repelled from regions of high potential
(morally costly configurations).

The interplay between spreading and localization determines the
character of the deliberation:

\begin{itemize}
\item
  In regions of low potential with gentle curvature, the moral state
  spreads freely --- deliberation explores many framings.
\item
  Near a deep potential well (a strongly preferred moral outcome), the
  state localizes --- deliberation converges.
\item
  Near a stratum boundary, the wave function may tunnel --- a framing
  that is classically inaccessible (the potential barrier between strata
  is high) may be reached through quantum tunneling.
\end{itemize}

\subsubsection{Tunneling Across Moral
Barriers}\label{tunneling-across-moral-barriers}

Quantum tunneling is the process by which a quantum state penetrates a
potential barrier that would be classically impassable. In moral terms:

\textbf{Definition 13.8 (Moral Tunneling).} \emph{A moral state}
\(|\psi\rangle\) \emph{tunnels across a stratum boundary if the wave
function has nonzero amplitude on both sides of a potential barrier}
\(V > E\)\emph{, where} \(E\) \emph{is the energy of the state.}

\textbf{Moral interpretation.} A moral agent may find themselves ``on
the other side'' of a moral boundary --- having adopted a position that
seemed inaccessible given their existing commitments --- not through a
classical trajectory (a deliberate, step-by-step moral argument that
crosses the boundary) but through a quantum process (a sudden gestalt
shift, an ``aha'' moment, a reframing that skips the barrier entirely).

This is the geometric content of moral \emph{conversion} --- the
experience of suddenly seeing a situation entirely differently, without
having traversed the intermediate positions. Augustine's conversion in
the garden, Saul's on the road to Damascus, the whistle-blower's sudden
decision to speak --- these are not well-modeled as gradient descents
across a smooth landscape. They are better modeled as tunneling events:
the wave function, trapped on one side of a high barrier, leaks through
to the other side, and a measurement (decision, action) catches it
there.

The tunneling rate through a barrier of height \(V_{0}\) and width \(d\)
is approximately:

\[T \sim \exp\left( - 2d\sqrt{2(V_{0} - E)} \right)\]

Higher barriers (more deeply entrenched moral positions) and wider
barriers (more intermediate positions that must be skipped) yield
exponentially smaller tunneling rates --- consistent with the rarity of
genuine moral conversions.

\subsection{13.8 Entanglement: Correlated Moral
States}\label{entanglement-correlated-moral-states}

\subsubsection{Multiagent Quantum
States}\label{multiagent-quantum-states}

When multiple agents are involved in a moral situation, the total moral
state lives in the \emph{tensor product} of their individual Hilbert
spaces:

\[\mathcal{H}_{\text{total}} = \mathcal{H}_{A} \otimes \mathcal{H}_{B} \otimes \cdots\]

A \emph{product state} \(|\psi_{A}\rangle \otimes |\psi_{B}\rangle\)
represents agents whose moral states are independent --- knowing one
agent's moral state tells you nothing about the other's.

\textbf{Definition 13.9 (Moral Entanglement).} \emph{A multiagent state}
\(|\Psi\rangle \in \mathcal{H}_{A} \otimes \mathcal{H}_{B}\) \emph{is
entangled if it cannot be written as a product:}

\[|\Psi\rangle \neq |\psi_{A}\rangle \otimes |\psi_{B}\rangle\]

\emph{for any choice of} \(|\psi_{A}\rangle\) \emph{and}
\(|\psi_{B}\rangle\)\emph{.}

\subsubsection{Moral EPR States}\label{moral-epr-states}

The maximally entangled moral state is:

\[|\Psi^{+}\rangle = \frac{1}{\sqrt{2}}\left( |O\rangle_{A} \otimes |C\rangle_{B} + |L\rangle_{A} \otimes |N\rangle_{B} \right)\]

This state correlates Alice's jural position with Bob's Hohfeldian
correlative: if Alice is measured to be in the obligation state
\(|O\rangle\), Bob is \emph{necessarily} found in the correlative claim
state \(|C\rangle\). If Alice is in the liberty state \(|L\rangle\), Bob
is in the no-claim state \(|N\rangle\).

This is the quantum version of the Hohfeldian correlative structure
(Chapter 8, §8.4). Classically, the correlative constraint is O ↔ C, L ↔
N --- an obligation for one party implies a claim for the other. In the
quantum framework, the correlative constraint is an \emph{entanglement}
constraint: the agents' Hohfeldian states are correlated in a way that
cannot be decomposed into independent states.

The BIP experiments (Chapter 17) found that the correlative symmetry O ↔
C holds at 87\% accuracy across the Dear Abby corpus. The 13\% violation
rate corresponds, in the quantum framework, to imperfect entanglement
--- the moral state is not maximally entangled but partially mixed:

\[\rho_{AB} = 0.87\,|\Psi^{+}\rangle\langle\Psi^{+}| + 0.13\,\rho_{\text{noise}}\]

where \(\rho_{\text{noise}}\) is an uncorrelated noise term. This is a
\emph{mixed entangled state} --- predominantly entangled (the
correlative structure is strong) but with a noise floor (some violations
of the correlative constraint).

\subsubsection{Entanglement and Collective
Obligation}\label{entanglement-and-collective-obligation}

Moral entanglement connects to the problem of collective moral agency
(Chapter 14). When agents are morally entangled --- when their
obligations are correlated in a way that resists decomposition --- the
collective has properties that no individual member possesses alone. The
obligation is genuinely shared, not merely distributed.

A collective obligation that is represented as an entangled state
\(|\Psi\rangle_{AB}\) cannot be ``traced out'' to individual obligations
without losing information. The \emph{partial trace} over agent B:

\[\rho_{A} = \text{Tr}_{B}\left( |\Psi\rangle\langle\Psi| \right)\]

is a mixed state --- agent A's individual moral state is uncertain, even
though the collective state is perfectly definite. This is the quantum
content of the intuition that collective obligations are not reducible
to individual ones: when two agents are morally entangled, neither has a
definite individual obligation, even though their joint obligation is
perfectly clear.

\subsection{13.9 The Classical Limit: When Classical Ethics
Suffices}\label{the-classical-limit-when-classical-ethics-suffices}

\subsubsection{The Correspondence
Principle}\label{the-correspondence-principle}

Quantum mechanics reduces to classical mechanics in the limit where
quantum effects (superposition, interference, entanglement) are
negligible. Analogously, quantum normative dynamics reduces to classical
geometric ethics in the limit where:

\begin{enumerate}
\def\labelenumi{\arabic{enumi}.}
\item
  \textbf{The moral state is sharply localized.} If \(|\psi(p)|^{2}\) is
  concentrated near a single point \(p_{0}\mathcal{\in M}\) (the wave
  function is a narrow peak), the quantum dynamics approximates
  classical dynamics at \(p_{0}\). This corresponds to a moral situation
  with a definite, unambiguous assessment --- no competing framings, no
  genuine indeterminacy.
\item
  \textbf{Decoherence is complete.} If all off-diagonal elements of the
  density matrix have decayed to zero, the quantum state is a classical
  mixture, and all interference effects vanish. This corresponds to a
  deliberation that has resolved all framing conflicts and is simply
  uncertain about the facts.
\item
  \textbf{The potential barriers are high.} If the potential barriers
  between strata are much larger than the state's energy, tunneling is
  exponentially suppressed, and transitions between strata occur only
  through classical (deliberate) boundary crossings. This corresponds to
  a moral landscape where the boundaries are firm and well-understood
  --- the agent knows which regime applies.
\end{enumerate}

\subsubsection{When to Use Quantum Normative
Dynamics}\label{when-to-use-quantum-normative-dynamics}

The quantum extension is not needed for most everyday moral reasoning.
Classical geometric ethics handles:

\begin{itemize}
\item
  Cases with a single operative framing
\item
  Cases where the uncertainty is purely empirical (about facts, not
  about frameworks)
\item
  Cases where the stratum boundaries are clear and the applicable rules
  are unambiguous
\end{itemize}

The quantum extension is needed for:

\begin{itemize}
\item
  \textbf{Genuine dilemmas} where two moral frameworks yield
  incompatible verdicts and the agent must hold both in superposition
\item
  \textbf{Framing effects} where the order or manner of considering
  options affects the outcome
\item
  \textbf{Moral conversion} where an agent suddenly shifts moral
  orientation without traversing intermediate positions
\item
  \textbf{Collective obligations} where the moral state of a group is
  entangled and resists decomposition into individual states
\item
  \textbf{Anomalous probability patterns} where the probability of a
  moral verdict violates the classical law of total probability
\end{itemize}

The quantum framework \emph{contains} the classical framework as a
special case. Using it when the classical theory suffices produces the
same answers with more apparatus. Using it when the classical theory
fails produces new answers that the classical theory cannot provide.

The quantum extension is the component of this framework most likely to
be revised or abandoned as empirical evidence accumulates. The
order-effect results (§17.10) provide the first direct support, but
interference and entanglement predictions remain untested. Readers
primarily interested in AI applications (Part V) may treat this chapter
as optional background; the DEME architecture (Chapter 19) does not
depend on quantum normative dynamics.

Empirical Update: CHSH Tests and the Classical Gauge Group

\textbf{{[}Empirical result.{]}} Two rounds of Bell tests have now been
conducted. The initial test (N=600 evaluations across four
entangled-agent scenarios) and the expanded v3 Bell test (N=9,000, 18
configurations across 5 maximally entangled scenarios in 6 languages)
both found no violation of the CHSH inequality: all \textbar S\textbar{}
≤ 2. The v3 test used scenarios specifically designed to maximize
entanglement: the Trolley Standoff, Mutual Confession, Sacrifice
Lottery, Symmetric Duel, and Entangled Lie. Cross-lingual conditions
(en-ja, en-ar, ja-ar) showed language-specific variation (Japanese:
S=−1.91; Arabic: S=0.14) but none violated \textbar S\textbar{} ≤ 2. The
Hardy non-locality test met 3 of 4 conditions but not the fourth.
No-signaling conditions were satisfied in all scenarios (p
\textgreater{} 0.05).

\textbf{Interpretation.} The non-abelian structure identified in the
Hohfeldian state space (§8.4; see §12.3 for the revised gauge group) is
classical D₄, not quantum SU(2). The original gauge group proposal
(SU(2)ᵢ × U(1)ᴴ) predicted quantum contextuality with
\textbar S\textbar{} up to 2√2 ≈ 2.83 for appropriate moral states. The
data are consistent with classical correlations only.

This does not invalidate the quantum formalism developed in this
chapter. The mathematical structure of superposition, interference, and
measurement remains a valid modeling tool for moral deliberation under
genuine ambiguity (§13.3--12.5). What the CHSH results establish is that
the gauge group governing Hohfeldian state transitions is the discrete
D₄ (see §12.3), not the continuous SU(2) originally proposed. Quantum
effects, if present in moral cognition, operate at a level different
from the Hohfeldian state transitions tested here.

\subsection{13.10 Worked Example: The Whistleblower's
Dilemma}\label{worked-example-the-whistleblowers-dilemma}

\subsubsection{Setup}\label{setup-3}

Emma, an engineer at a chemical company, discovers that the company has
been concealing data about toxic emissions. She faces the
whistleblower's dilemma: loyalty to her employer vs.~duty to public
safety.

\subsubsection{The Classical Analysis
(Review)}\label{the-classical-analysis-review}

In the classical framework (Chapters 5--7), Emma's situation is a point
\(p\) in \(\mathcal{M}\) with competing obligation vectors:

\begin{itemize}
\item
  \(O_{\text{loyalty}} = (0,0.7,0,0,0.5,0,0.4,0,0)\): duty to employer,
  privacy of company data, relational loyalty
\item
  \(O_{\text{safety}} = (0.9,0.6,0.3,0,0,0,0.2,0,0)\): public welfare,
  duty to warn, fairness to affected communities
\end{itemize}

The classical analysis computes the satisfaction function for each
option and evaluates the trade-offs via the metric. It yields a definite
(if difficult) answer.

\subsubsection{The Quantum Analysis}\label{the-quantum-analysis}

But Emma's real experience is not like this. She does not simply weigh
two definite obligations. She is uncertain which \emph{framing} of her
situation is correct:

\textbf{Framing} \(|a\rangle\)\textbf{: Loyal Employee.} The data is
ambiguous. The company may have reasonable explanations. Going public
would betray trust and possibly be wrong. The operative obligation is
confidentiality.

\textbf{Framing} \(|b\rangle\)\textbf{: Public Guardian.} The data is
clear. The company is concealing genuine harm. Going public is a duty,
not a betrayal. The operative obligation is safety.

Before Emma decides, her moral state is a superposition:

\[|\psi\rangle = \alpha\,|a\rangle + \beta\,|b\rangle,\quad\quad|\alpha|^{2} + |\beta|^{2} = 1\]

The framings are not independent --- they \emph{interfere}. Considering
the ``loyal employee'' framing makes the evidence seem more ambiguous
(constructive interference between loyalty and doubt). Considering the
``public guardian'' framing makes the evidence seem more damning
(constructive interference between duty and clarity). The interference
is not symmetric: the relative phase \(\theta\) between the framings
determines whether one reinforces or suppresses the other.

\subsubsection{Time Evolution}\label{time-evolution}

Emma's moral state evolves under the moral Hamiltonian \(\widehat{H}\).
As she gathers more evidence (lowering the potential barrier between the
``loyalty'' and ``safety'' strata), the wave function begins to tunnel
from \(|a\rangle\) toward \(|b\rangle\). The tunneling is initially slow
(the barrier is high --- betraying an employer is a serious moral cost).
But as evidence accumulates, the barrier lowers, the tunneling rate
increases, and \(|\beta|^{2}\) grows.

\subsubsection{The Measurement}\label{the-measurement}

The moment of decision --- when Emma picks up the phone to call the
regulator --- is a measurement. The superposition collapses:

\[|\psi\rangle \rightarrow |b\rangle\quad\text{with probability }|\beta|^{2}\]

or

\[|\psi\rangle \rightarrow |a\rangle\quad\text{with probability }|\alpha|^{2}\]

The measurement is irreversible. Once Emma has spoken (or not), the
superposition is destroyed. She is now definitely a whistleblower or
definitely a loyal employee. The quantum indeterminacy is over;
classical ethics applies from this point forward.

\subsubsection{What the Quantum Analysis
Adds}\label{what-the-quantum-analysis-adds}

The quantum analysis captures three features that the classical analysis
misses:

\begin{enumerate}
\def\labelenumi{\arabic{enumi}.}
\item
  \textbf{Interference effects.} Emma's evaluation of the evidence is
  not a simple average of two framings. The framings interact:
  considering loyalty makes the evidence seem weaker (destructive
  interference with safety); considering duty makes it seem stronger
  (constructive interference with safety). The probabilities of her
  eventual decision are not classical averages --- they include
  interference terms that depend on the relative phase.
\item
  \textbf{Tunneling.} Emma's conversion from ``loyal employee'' to
  ``whistleblower'' need not follow a continuous path through moral
  space. She may tunnel --- suddenly seeing the situation entirely
  differently, without having traversed intermediate positions. The
  classical theory requires a continuous trajectory; the quantum theory
  allows discontinuous transitions through potential barriers.
\item
  \textbf{Entanglement with the company.} Emma's moral state is
  entangled with the company's: if she is in the ``loyal employee''
  state, the company is in the ``trusted employer'' state; if she is in
  the ``whistleblower'' state, the company is in the ``concealing
  wrongdoer'' state. The entanglement means that her decision does not
  merely affect her own moral position --- it simultaneously determines
  the company's moral status, in a way that cannot be decomposed into
  independent individual decisions.
\end{enumerate}

\subsection{13.11 Summary}\label{summary-6}

{\def\LTcaptype{none} % do not increment counter
\begin{longtable}[]{@{}
  >{\raggedright\arraybackslash}p{(\linewidth - 2\tabcolsep) * \real{0.3872}}
  >{\raggedright\arraybackslash}p{(\linewidth - 2\tabcolsep) * \real{0.4387}}@{}}
\toprule\noalign{}
\begin{minipage}[b]{\linewidth}\raggedright
\textbf{Classical Geometric Ethics}
\end{minipage} & \begin{minipage}[b]{\linewidth}\raggedright
\textbf{Quantum Normative Dynamics}
\end{minipage} \\
\midrule\noalign{}
\endhead
\bottomrule\noalign{}
\endlastfoot
Moral state: point \(p\mathcal{\in M}\) & Moral state: vector
\(\parallel \psi\mathcal{\rangle \in H}\) \\
Definite obligations \(O^{\mu}(p)\) & Superposition of obligations \\
Classical uncertainty: covariance \(\Sigma^{\mu\nu}\) & Quantum
indeterminacy: density matrix \(\rho\) \\
Trajectory \(\gamma(t)\) through \(\mathcal{M}\) & Wave function
\(\psi(p,t)\) on \(\mathcal{M}\) \\
Euler-Lagrange equations & Schrödinger equation \\
Classical boundary crossing & Quantum tunneling \\
Independent agents & Entangled moral states \\
Definite Hohfeldian position & Superposition of jural states \\
Correlative constraint O ↔ C & Entanglement:
\(\parallel O\rangle_{A} \otimes \parallel C\rangle_{B}\) \\
Decision as evaluation & Decision as measurement (collapse) \\
No interference & Interference between framings \\
\end{longtable}
}

Quantum normative dynamics extends the classical geometric framework
developed in Chapters 5--11 by adding the Hilbert space structure that
supports superposition, interference, and entanglement. The extension is
conservative: classical geometric ethics is recovered as the limit where
quantum effects (coherence, tunneling, entanglement) are negligible.

The extension is motivated by phenomena that classical ethics handles
poorly: the structure of deliberation (simultaneous consideration of
incompatible framings), moral conversion (sudden shifts across high
barriers), collective obligation (irreducibly shared duties), and
framing effects (the manner of considering options affects the verdict).

The key new structures are:

\begin{enumerate}
\def\labelenumi{\arabic{enumi}.}
\item
  \textbf{The moral Hilbert space}
  \(\mathcal{H =}L^{2}\mathcal{(M,}\mathbb{C)}\), in which moral states
  are vectors that can be superposed.
\item
  \textbf{Moral observables} --- self-adjoint operators whose
  eigenvalues are possible moral verdicts and whose eigenstates are
  states of definite moral assessment.
\item
  \textbf{The moral Schrödinger equation}
  \(i\partial_{t}|\psi\rangle = \widehat{H}|\psi\rangle\), governing the
  evolution of moral states under the moral Hamiltonian.
\item
  \textbf{The stratified Lagrangian} \(\mathcal{L}_{\text{strat}}\),
  which provides a variational principle for dynamics on the patchwork
  moral manifold.
\item
  \textbf{Moral entanglement}, encoding irreducible correlations between
  agents' moral states.
\end{enumerate}

These structures do not replace moral judgment. They provide a
mathematical vocabulary for analyzing phenomena that classical geometric
ethics can describe but cannot fully explain --- and they make
predictions (interference patterns, tunneling rates, entanglement
correlations) that are, in principle, empirically testable.

\subsection{Technical Appendix}\label{technical-appendix-7}

\textbf{Theorem 13.1 (Decomposition of} \(\mathcal{H}\) \textbf{by
Strata).} \textbf{{[}Theorem (conditional).{]}} \emph{Let}
\(\mathcal{M =}\bigsqcup_{\alpha}S_{\alpha}\) \emph{be a
Whitney-stratified moral manifold. Then the moral Hilbert space
decomposes as:}

\[\mathcal{H =}\bigoplus_{\alpha}\mathcal{H}_{\alpha} \oplus \mathcal{H}_{\partial}\]

\emph{where} \(\mathcal{H}_{\alpha} = L^{2}(S_{\alpha})\) \emph{is the
subspace of states confined to stratum} \(S_{\alpha}\)\emph{, and}
\(\mathcal{H}_{\partial}\) \emph{is the subspace of states supported on
stratum boundaries.}

\emph{\textbf{Proof}. By Proposition 8.1, M admits a Whitney
stratification M = ⨆\_α S\_α. Since the strata are disjoint, L²(M, dμ) =
⨁\_α L²(S\_α, dμ\textbar\_\{S\_α\}) where dμ is the Riemannian volume
measure induced by the moral metric g. Set H\_α = L²(S\_α,
dμ\textbar\_\{S\_α\}): these are the subspaces of states confined to
individual strata. The boundary strata ∂S\_α have dimension \textless{}
dim M, so μ(∂S\_α) = 0 in the ambient measure; they contribute trivially
to H via dμ. However, each boundary stratum B carries its own intrinsic
measure dμ\_B (the volume form of the induced metric on B). Define H\_∂
= ⨁\_B L²(B, dμ\_B): these are the ``edge modes'' localized on stratum
boundaries. The total Hilbert space is then H = ⨁\_α H\_α ⊕ H\_∂, where
the direct sum is orthogonal by construction (the strata are disjoint,
and the boundary states are supported on measure-zero sets with respect
to the bulk measure). □}\(\mathcal{M}L^{2}L^{2}\mu\mathcal{H▫}\)

Measure-theoretic clarification. Since H = L²(M, dμ\_g) and the strata
partition M, the decomposition H = ⊕ H\_α (orthogonal direct sum over
open strata) follows from the additivity of the Lebesgue integral. The
boundary strata ∂S\_α have codimension ≥ 1 and hence ambient measure
zero; they do not contribute to H = L²(M, dμ\_g). Boundary states ---
wave functions supported on stratum boundaries --- are distributional
objects belonging to the rigged Hilbert space (Gelfand triple) Φ ⊂ H ⊂
Φ*, not to H itself. The term H\_boundary in the decomposition above
should therefore be understood as living in the dual space Φ*,
representing edge modes that couple adjacent strata. This distinction is
immaterial for the physical predictions (which involve inner products in
H) but is required for mathematical precision.

\textbf{Proposition 13.2 (Conservation of Probability Under Moral
Evolution).} \emph{The moral Schrödinger equation conserves the total
probability:}

\[\frac{d}{dt}\langle\psi|\psi\rangle = 0\]

\emph{\textbf{Proof}.} Since \(\widehat{H}\) is self-adjoint,
\(\frac{d}{dt}\langle\psi|\psi\rangle = \langle\dot{\psi}|\psi\rangle + \langle\psi|\dot{\psi}\rangle = i\langle\widehat{H}\psi|\psi\rangle - i\langle\psi|\widehat{H}\psi\rangle = 0\)
by self-adjointness. \(▫\)

Proof expansion. Write ψ(t) = e\^{}\{−iHt\}ψ(0). Then d/dt⟨ψ\textbar ψ⟩
= ⟨∂\_tψ\textbar ψ⟩ + ⟨ψ\textbar∂\_tψ⟩ = ⟨−iHψ\textbar ψ⟩ +
⟨ψ\textbar−iHψ⟩ = i⟨Hψ\textbar ψ⟩ − i⟨ψ\textbar Hψ⟩. Since H is
self-adjoint, ⟨Hψ\textbar ψ⟩ = ⟨ψ\textbar Hψ⟩, so d/dt⟨ψ\textbar ψ⟩ = 0.
Hence \textbar\textbar ψ(t)\textbar\textbar² =
\textbar\textbar ψ(0)\textbar\textbar² for all t, and the evolution
operator U(t) = e\^{}\{−iHt\} is unitary.

\textbf{Proposition 13.3 (Tunneling Rate Through a Moral Barrier).}
\emph{For a one-dimensional rectangular potential barrier of height}
\(V_{0}\) \emph{and width} \(d\)\emph{, with an incident state of
energy} \(E < V_{0}\)\emph{, the tunneling probability is:}

\[T \approx \frac{16\, E(V_{0} - E)}{V_{0}^{2}}\,\exp\left( - 2d\sqrt{2(V_{0} - E)} \right)\]

\emph{In the moral context,} \(V_{0}\) \emph{represents the moral cost
of occupying the intermediate (boundary) region,} \(E\) \emph{is the
current ``moral energy'' of the state, and} \(d\) \emph{is the ``width''
of the moral barrier in the metric-induced distance on}
\(\mathcal{M}\)\emph{. High barriers and wide barriers yield
exponentially suppressed tunneling --- consistent with the rarity of
sudden moral conversions.}

\emph{\textbf{Proof}.} Consider the time-independent Schrödinger
equation −½ψ\textquotesingle\textquotesingle{} + V(x)ψ = Eψ with V(x) =
V₀ for 0 ≤ x ≤ d and V(x) = 0 otherwise. In the barrier region, ψ ∝ e±κx
where κ = √(2(V₀ − E)). Matching ψ and ψ\textquotesingle{} at x = 0 and
x = d yields the transmission coefficient: T = {[}1 + V₀² sinh²(κd) /
(4E(V₀ − E)){]}⁻¹. For κd ≫ 1, sinh(κd) ≈ ½e\^{}\{κd\}, giving T ≈
16E(V₀ − E)/V₀² · e\^{}\{−2κd\} = 16E(V₀ − E)/V₀² · exp(−2d√(2(V₀ −
E))). □

\textbf{Proposition 13.4 (Entanglement Entropy of Correlative States).}
\emph{For the partially entangled Hohfeldian state}

\[\rho_{AB} = (1 - \epsilon)\,|\Psi^{+}\rangle\langle\Psi^{+}| + \epsilon\,\frac{I}{4}\]

\emph{with}
\(|\Psi^{+}\rangle = \frac{1}{\sqrt{2}}(|O,C\rangle + |L,N\rangle)\)
\emph{and noise parameter} \(\epsilon \in \lbrack 0,1\rbrack\)\emph{,
the entanglement entropy of the reduced state}
\(\rho_{A} = \text{Tr}_{B}(\rho_{AB})\) \emph{is:}

\[S(\rho_{A}) = - \text{Tr}(\rho_{A}\ln\rho_{A}) = ln2\quad\text{for all }\epsilon\]

\emph{The reduced state of each agent is maximally uncertain regardless
of the noise level --- the correlative structure ensures that individual
Hohfeldian positions are uniformly indeterminate whenever the collective
state has any entanglement.}

\emph{\textbf{Proof}.} Compute the reduced state ρ\_A =
Tr\_B(ρ\_\{AB\}). The maximally mixed component contributes Tr\_B(I/4) =
I\_A/2. For the pure component, expand \textbar Ψ⁺⟩ =
(1/√2)(\textbar O⟩\textbar C⟩ + \textbar L⟩\textbar N⟩); the partial
trace over B yields Tr\_B(\textbar Ψ⁺⟩⟨Ψ⁺\textbar) =
½(\textbar O⟩⟨O\textbar{} + \textbar L⟩⟨L\textbar) = I\_A/2, since
⟨C\textbar C⟩ = ⟨N\textbar N⟩ = 1 and ⟨C\textbar N⟩ = 0. Hence ρ\_A =
(1−ε)(I\_A/2) + ε(I\_A/2) = I\_A/2 for all ε. The entropy is S(ρ\_A) =
−Tr((I\_A/2) ln(I\_A/2)) = −2 · ½ ln(½) = ln 2. □

❖

\emph{The moral world admits superposition. Before a decision is made,
the situation is not in one definite state but in many --- framings
coexisting, interfering, reinforcing and canceling in patterns that no
classical mixture can reproduce.}

\emph{This is not a metaphor. It is a mathematical structure: the
Hilbert space of square-integrable functions on the moral manifold,
equipped with observables, a Hamiltonian, and the Schrödinger equation.
The structure makes predictions --- interference patterns, tunneling
rates, entanglement correlations --- that are testable, at least in
principle, against the empirical record of moral deliberation.}

\emph{The quantum extension does not replace the classical framework. It
contains it. Classical geometric ethics is the limit of quantum
normative dynamics when coherence is lost, barriers are high, and the
moral state is sharply localized. Most everyday moral reasoning operates
in this limit --- and the classical theory of Chapters 5--11 is
adequate.}

\emph{But for the hard cases --- genuine dilemmas, sudden conversions,
collective obligations that resist decomposition, deliberations where
the order of consideration matters --- the quantum framework reveals
structure that the classical theory cannot see. The moral world is
richer than any classical description can capture.}

\emph{Quantum Normative Dynamics is the mathematics of that richness.}

\section{Part IV: Agents and
Collectives}\label{part-iv-agents-and-collectives}

\section{Chapter 14: Collective Moral Agency --- Aggregation, Emergence,
and the Tensors of Shared
Obligation}\label{chapter-14-collective-moral-agency-aggregation-emergence-and-the-tensors-of-shared-obligation}

\begin{quote}
\textbf{\textsc{RUNNING EXAMPLE --- Priya's Model}}

\emph{Priya cannot fix this alone. HealthBridge's collective agency
tensor involves: Priya (designed the system), Dr. Osei (flagged the
problem), Sarah Chen, VP of Product (controls deployment), the Board
(sets priorities), the FDA (regulates trial access), and the patients
(who have no voice in the process). The structure tensor is nonzero:
HealthBridge's decisions are not the sum of individual choices. The
company has emergent obligations that no single employee carries. Priya
can document the bias, but only Sarah can change deployment. Sarah can
authorize a fix, but only the Board can allocate budget. The
responsibility remainder---for patients already harmed---does not vanish
with organizational restructuring. Mrs. Voss's melanoma did not wait for
the org chart to sort itself out.}
\end{quote}

\subsection{14.1 Beyond Individual
Agents}\label{beyond-individual-agents}

The geometric framework developed in Parts II and III has focused on a
single moral situation evaluated from various perspectives. The
obligation vector \(O^{\mu}\) at a point \(p\mathcal{\in M}\) represents
what is owed; the interest covector \(I_{\mu}\) represents what matters;
the satisfaction \(S = I_{\mu}O^{\mu}\) measures how well duty meets
need. Even the multi-agent evaluation tensor \(E_{\nu a}^{\mu}\) of
Chapter 6 (§6.6) --- which indexes by agent \(a\) --- treats each
agent's moral position as independently specifiable.

But much of moral life involves \emph{collective} agents: corporations,
states, committees, families, professions, religious communities, and
(increasingly) human-AI teams. These collectives pose questions that the
individual-agent framework cannot answer:

\begin{itemize}
\item
  Does a corporation have moral obligations distinct from those of its
  employees?
\item
  When a committee decides wrongly, who bears responsibility --- and
  does the sum of individual responsibilities exhaust the collective
  responsibility?
\item
  Can an institution be unjust even if every individual within it acts
  permissibly?
\item
  How do individual moral tensors combine into collective structures ---
  and does the combination introduce \emph{new} moral content not
  present in the components?
\end{itemize}

This chapter develops the geometry of collective moral agency. The
central result is that collective moral tensors have \emph{emergent
components} --- terms that cannot be expressed as sums or products of
individual tensors. The collective is not merely the aggregate of its
members. It is a higher-rank moral entity with its own geometric
structure.

\subsection{14.2 The Aggregation Problem}\label{the-aggregation-problem}

\subsubsection{Why Simple Addition
Fails}\label{why-simple-addition-fails}

The most natural first attempt at collective moral structure is
\emph{summation}: the collective obligation is the sum of individual
obligations.

\[O_{\text{collective}}^{\mu}\overset{?}{=}\sum_{a \in A}^{}O_{a}^{\mu}\]

where \(A\) is the set of members and \(O_{a}^{\mu}\) is member \(a\)'s
obligation.

This fails for at least three reasons.

\textbf{1. Collectives have obligations that no member has.} A
corporation has an obligation to pay its debts. No individual employee
has this obligation in their personal capacity. A nation has an
obligation to provide for its citizens' basic needs. No individual
citizen has this obligation singly. The collective obligation is not a
sum of individual obligations --- it has no individual summands.

\textbf{2. Individual obligations may cancel or reinforce in
non-additive ways.} If Alice has an obligation to keep a secret and Bob
has an obligation to reveal it, the group's obligation is not the zero
vector (cancel) nor the sum of the two vectors (reinforce). It depends
on the \emph{structure} of the group --- who has authority, what the
group's purpose is, what decision procedure is in force. The combination
is mediated by structure, not by addition.

\textbf{3. The dimensionality may change.} An individual agent operates
in the nine-dimensional moral manifold of Chapter 5. A collective agent
may operate in a higher-dimensional space --- the nine ethical
dimensions, plus structural dimensions encoding the collective's
internal organization. The collective obligation vector may have
components along dimensions that have no individual analogue.

\subsubsection{The Tensor Product
Approach}\label{the-tensor-product-approach}

A better starting point is the \emph{tensor product} of individual moral
spaces. If agent \(a\) operates in a moral tangent space
\(T_{a} \cong \mathbb{R}^{9}\), the collective of \(n\) agents operates
in:

\[T_{\text{collective}} = T_{1} \otimes T_{2} \otimes \cdots \otimes T_{n}\]

This space has dimension \(9^{n}\) --- far larger than the \(9n\)
dimensions of the direct sum
\(T_{1} \oplus T_{2} \oplus \cdots \oplus T_{n}\) that simple addition
would use.

\textbf{Why the tensor product?} Because it captures
\emph{correlations}. A state in the direct sum \(T_{1} \oplus T_{2}\)
specifies agent 1's obligation and agent 2's obligation independently
--- \(9 + 9 = 18\) components. A state in the tensor product
\(T_{1} \otimes T_{2}\) specifies the full joint obligation, including
all correlations between the agents' moral positions ---
\(9 \times 9 = 81\) components. The additional \(81 - 18 = 63\)
components encode the \emph{relational structure} of the collective: how
the agents' obligations interact, constrain, and depend on one another.

This is precisely the structure encountered in Chapter 13 (§13.8) in the
quantum context. The tensor product of individual Hilbert spaces
contains entangled states --- joint states that cannot be decomposed
into individual states. The classical analogue is the same: the tensor
product of individual moral spaces contains \emph{collective moral
states} that cannot be decomposed into individual obligations.

\subsection{14.3 The Collective Agency
Tensor}\label{the-collective-agency-tensor}

\subsubsection{Definition}\label{definition-5}

\textbf{Definition 14.1 (Collective Agency Tensor).} \emph{Let}
\(A = \{ 1,\ldots,n\}\) \emph{be a set of agents. The collective agency
tensor is a multilinear map:}

Convention. All agents in A evaluate the same base point p ∈ M; the
collective agency tensor T is defined at p and transforms as a tensor
under coordinate changes on the moral manifold. The agent index a ∈ A is
a label, not a tensor index: T does not transform covariantly or
contravariantly in the agent label. Formally, T is a family of
(1,1)-tensors on T\_pM, one for each coalition subset, parametrized by
the agent set A.

\[\mathcal{C:}T_{1}^{*} \otimes \cdots \otimes T_{n}^{*} \rightarrow T_{\mathcal{M}}\]

\emph{that takes one interest covector from each agent's cotangent space
and produces an obligation vector in the moral tangent space. In
components:}

\[O_{\text{collective}}^{\mu} = \mathcal{C}_{\nu_{1}\nu_{2}\cdots\nu_{n}}^{\mu}\, I_{\nu_{1}}^{(1)}\, I_{\nu_{2}}^{(2)}\,\cdots\, I_{\nu_{n}}^{(n)}\]

The collective agency tensor \(\mathcal{C}\) is a \((1,n)\)-tensor ---
one upper index (the output obligation dimension) and \(n\) lower
indices (one for each agent's interest dimension). Its rank increases
with the number of agents: a two-agent collective has a rank-3 tensor
(\(9 \times 9 \times 9 = 729\) components), a three-agent collective has
a rank-4 tensor (\(9^{4} = 6561\) components), and so on.

\subsubsection{Decomposition: Individual and Emergent
Components}\label{decomposition-individual-and-emergent-components}

The collective agency tensor can be decomposed into components
attributable to individual agents and a \emph{remainder} that is
genuinely collective.

\textbf{Definition 14.2 (Decomposition of Collective Agency).} \emph{For
a two-agent collective, the agency tensor decomposes as:}

\[\mathcal{C}_{\nu_{1}\nu_{2}}^{\mu} = \mathcal{C}_{(\nu_{1})}^{\mu}\delta_{\nu_{2}} + \delta_{\nu_{1}}\mathcal{C}_{(\nu_{2})}^{\mu} + \mathcal{E}_{\nu_{1}\nu_{2}}^{\mu}\]

\emph{where:} - \(\mathcal{C}_{(\nu_{1})}^{\mu}\delta_{\nu_{2}}\)
\emph{is agent 1's individual contribution (independent of agent 2's
interests)} - \(\delta_{\nu_{1}}\mathcal{C}_{(\nu_{2})}^{\mu}\) \emph{is
agent 2's individual contribution} -
\(\mathcal{E}_{\nu_{1}\nu_{2}}^{\mu}\) \emph{is the emergent component
--- the part of collective agency not attributable to either individual}

The emergent component \(\mathcal{E}\) vanishes if and only if the
collective obligation is the sum of individual obligations. When
\(\mathcal{E \neq}0\), the collective has moral agency that exceeds the
combination of its members.

\subsubsection{Measuring Emergence}\label{measuring-emergence}

The \emph{degree of emergence} of a collective can be quantified:

\[\eta = \frac{\mathcal{\parallel E \parallel}}{\mathcal{\parallel C \parallel}}\]

where the norms are computed using the induced metric on the tensor
spaces. When \(\eta = 0\), the collective is fully reducible to its
members. When \(\eta = 1\), the collective agency is entirely emergent
--- none of it can be attributed to individual members. Intermediate
values indicate partial emergence.

\textbf{Claim.} For most real collectives of interest (corporations,
states, institutions), \(\eta > 0\). The emergent component is nonzero
because the collective's structure --- its decision procedures, role
differentiation, and institutional purpose --- generates obligations
that no individual member bears.

\textbf{Computational implementation (February 2026).} The DEME V3
reference implementation realizes the collective agency tensor as a
rank-5 \texttt{MoralTensor} with coalition indices. The decomposition
into individual and emergent components is computed explicitly, and the
degree of emergence \(\eta\) is a reportable diagnostic. V3 also
implements Shapley value credit assignment for fair allocation of
collective moral responsibility:

\[\phi_{i}(v) = \sum_{S \subseteq N\backslash\{ i\}}^{}\frac{|S|!\,(|N| - |S| - 1)!}{|N|!}\left\lbrack v(S \cup \{ i\}) - v(S) \right\rbrack\]

where \(v\) is the coalition value function derived from the collective
agency tensor. This enables quantitative answers to the question ``how
much of the collective obligation is attributable to agent \(i\)?''---a
question that Chapter 18 identifies as critical for human-AI
responsibility allocation. Nash equilibrium and correlated equilibrium
computations provide multi-agent coordination under norm constraints,
extending the Norm-Constrained Stochastic Game formalism of Chapter 19
to strategic settings.

\textbf{ErisML references}: demo\_game\_theory.py, demo\_shapley.py. The
Nash equilibrium computation is demonstrated in demo\_game\_theory.py,
which implements the StrategicLayer for a Prisoner's Dilemma payoff
matrix --- the canonical example of multi-agent moral coordination where
individually rational choices produce collectively suboptimal outcomes.
The Shapley value computation is demonstrated in demo\_shapley.py, which
applies the CooperativeLayer to an airport runway cost allocation
problem: each airline's fair share of the collective infrastructure cost
is computed as its marginal contribution averaged over all possible
coalition orderings, exactly implementing the formula φᵢ(v) described
above.

\subsection{14.4 The Structure Tensor}\label{the-structure-tensor}

\subsubsection{What Makes a Collective More Than a
Set}\label{what-makes-a-collective-more-than-a-set}

A collective agent is not merely a set of individuals. It has
\emph{structure}: decision procedures, role assignments, authority
hierarchies, persistence conditions, and representational conventions.
This structure mediates the relationship between individual and
collective moral positions.

\textbf{Definition 14.3 (Structure Tensor).} \emph{The structure tensor
of a collective} \(A\) \emph{is a} \((0,2)\)\emph{-tensor}
\(\Sigma_{ab}\) \emph{on the space of members, encoding the structural
relationships between them:}

Clarification. The ``space of members'' on which the structure tensor is
defined is the real vector space ℝ\^{}\textbar A\textbar, where
\textbar A\textbar{} is the cardinality of the agent set A. Each basis
vector e\_a corresponds to one agent a ∈ A. The component S\^{}a\_\{bc\}
measures the degree to which the relationship between agents b and c
structurally influences agent a's moral standing.

\begin{itemize}
\item
  \(\Sigma_{ab} > 0\): \emph{agents} \(a\) \emph{and} \(b\) \emph{have
  aligned roles (authority flows in the same direction)}
\item
  \(\Sigma_{ab} < 0\): \emph{agents} \(a\) \emph{and} \(b\) \emph{have
  opposing roles (checks and balances)}
\item
  \(\Sigma_{ab} = 0\): \emph{agents} \(a\) \emph{and} \(b\) \emph{are
  structurally independent}
\end{itemize}

The structure tensor is generally \emph{asymmetric}:
\(\Sigma_{ab} \neq \Sigma_{ba}\) when the relationship between \(a\) and
\(b\) is directional (a supervisor-subordinate relationship, for
instance). The symmetric part
\(\Sigma_{(ab)} = \frac{1}{2}(\Sigma_{ab} + \Sigma_{ba})\) encodes the
\emph{mutual} structural relationship. The antisymmetric part
\(\Sigma_{\lbrack ab\rbrack} = \frac{1}{2}(\Sigma_{ab} - \Sigma_{ba})\)
encodes the \emph{power differential} --- the extent to which the
structural relationship favors one party over the other.

This connects to the power tensor \(P_{\mu\nu}\) introduced in Chapter 6
(§6.9): the antisymmetric part of the structural tensor is the
member-indexed version of the structural inequity tensor.

\subsubsection{How Structure Shapes Collective
Obligation}\label{how-structure-shapes-collective-obligation}

The collective agency tensor is related to the individual agency tensors
\emph{through} the structure tensor:

\[\mathcal{C}_{\nu_{1}\nu_{2}}^{\mu} = \sum_{a,b}^{}\Sigma_{ab}\, E_{\nu_{a}}^{\mu}(a) \otimes E_{\nu_{b}}^{\mu}(b) + \mathcal{E}_{\nu_{1}\nu_{2}}^{\mu}\]

The first term says: the collective obligation depends on the individual
evaluation tensors \(E_{\nu}^{\mu}(a)\) of each member, weighted and
combined according to the structure tensor \(\Sigma\). The second term
is the emergent remainder --- the part of collective agency that cannot
be expressed even through structure-mediated combination of individual
agencies.

\textbf{Example: A two-person partnership.} Alice and Bob form a law
firm. The structure tensor has \(\Sigma_{AB} = \Sigma_{BA} = 0.5\)
(equal partnership, symmetric). The collective agency tensor is
approximately the average of individual agency tensors, with a small
emergent component arising from the firm's identity as a legal entity
(the firm's obligation to clients exists independently of either
partner's individual obligations).

\textbf{Example: A corporate hierarchy.} A CEO and an employee have
\(\Sigma_{CE} = 0.8\) (the CEO's position strongly shapes the
collective's action) and \(\Sigma_{EC} = 0.2\) (the employee's position
weakly shapes it). The antisymmetric part
\(\Sigma_{\lbrack CE\rbrack} = 0.3\) encodes the power differential. The
collective obligation is dominated by the CEO's evaluation tensor, with
the employee's contribution attenuated by the structural asymmetry.

\subsection{14.5 Emergent Obligations}\label{emergent-obligations}

\subsubsection{Obligations No Member
Bears}\label{obligations-no-member-bears}

The emergent component \(\mathcal{E}\) of the collective agency tensor
represents obligations that exist at the collective level but not at the
individual level. These obligations are not mysterious or metaphysical
--- they arise from the \emph{structure} of the collective, just as the
properties of a molecule arise from the arrangement of its atoms, not
from the atoms individually.

\textbf{Example 1: Corporate debt.} A corporation borrows money. The
obligation to repay belongs to the corporation. No individual employee,
officer, or shareholder has this obligation in their personal capacity
(this is the legal content of limited liability). The obligation is an
emergent component of the corporate agency tensor --- it appears only in
\(\mathcal{E}\), not in any individual evaluation tensor
\(E_{\nu}^{\mu}(a)\).

\textbf{Example 2: National defense.} A nation has an obligation to
defend its citizens. No individual citizen has this obligation singly.
The obligation emerges from the collective structure --- the existence
of a polity with a government, a territory, and a citizenry. A random
collection of individuals in the same territory, without political
structure, would have no collective defense obligation. The obligation
requires the structure tensor \(\Sigma\) to be nonzero in specific ways
(political authority, territorial jurisdiction, citizen-state
relationships).

\textbf{Example 3: Institutional injustice.} A hiring system may be
structurally unjust --- systematically disadvantaging a particular group
--- even if every individual decision-maker acts in good faith according
to their individual obligations. The injustice is an emergent property
of the collective structure: it appears in \(\mathcal{E}\), not in any
individual \(E_{\nu}^{\mu}(a)\). This is the geometric content of the
concept of \emph{structural injustice} (Young 2011): injustice that no
individual perpetrates but that the collective generates through its
structure.

\subsubsection{Formal Criterion for
Emergence}\label{formal-criterion-for-emergence}

\textbf{Proposition 14.1 (Criterion for Emergent Obligation).} \emph{An
obligation} \(O_{\text{coll}}^{\mu}\) \emph{is emergent if and only if
there exists no set of individual obligations}
\(\{ O_{a}^{\mu}\}_{a \in A}\) \emph{such that:}

\[O_{\text{coll}}^{\mu} = \sum_{a \in A}^{}f_{a}(O_{a}^{\mu})\]

\emph{for any set of functions} \(\{ f_{a}\}\)\emph{. Equivalently,}
\(O_{\text{coll}}\) \emph{is emergent if and only if the emergent
component} \(\mathcal{E}_{\nu_{1}\cdots\nu_{n}}^{\mu}\) \emph{has
nonzero contraction with the collective interest:}

\[\mathcal{E}_{\nu_{1}\cdots\nu_{n}}^{\mu}\, I_{\nu_{1}}^{(1)}\,\cdots\, I_{\nu_{n}}^{(n)} \neq 0\]

The criterion is testable: an obligation is emergent when it cannot be
attributed to any individual member, no matter how we weight the
individual contributions. The irreducibility is structural, not merely
epistemic.

\emph{\textbf{Proof}.} (⇒) Suppose O\_collᵘ is emergent: no functions
\{f\_a\} satisfy O\_collᵘ = Σ\_a f\_a(O\_aᵘ). By the Decomposition
Theorem (Prop. 14.3), Cᵘ\_\{ν₁⋯ν\_n\} = Σ\_k C\_(k)ᵘ ⊗ ⋯ +
Eᵘ\_\{ν₁⋯ν\_n\}. The collective obligation is O\_collᵘ = Cᵘ\_\{ν₁⋯ν\_n\}
I⁽¹⁾\_\{ν₁\}⋯I⁽ⁿ⁾\_\{ν\_n\}. If Eᵘ\_\{ν₁⋯ν\_n\}
I⁽¹⁾\_\{ν₁\}⋯I⁽ⁿ⁾\_\{ν\_n\} = 0, then O\_collᵘ decomposes into
individual and pairwise terms, which has the form Σ\_a f\_a(O\_aᵘ) ---
contradicting emergence. Hence the contraction is nonzero. (⇐) Suppose
the emergent contraction is nonzero. Then O\_collᵘ contains a component
from E that, by definition of the fully connected term, cannot be
expressed as any function of individual obligations alone. Hence
O\_collᵘ is emergent. □

\textbf{Proof supplement (orthogonality).} The projections P\_k (onto
the k-agent interaction subspace) are orthogonal because the
inclusion-exclusion decomposition partitions tensor components by their
interaction order. Formally, if C = Σ\_k C\_k where C\_k involves
exactly k-agent interactions, then ⟨C\_j, C\_k⟩ = 0 for j ≠ k. This
follows from Möbius inversion on the subset lattice: the decomposition
C\_k = Σ\_\{\textbar S\textbar=k\} μ(S) C\_S (where μ is the Möbius
function) produces orthogonal components because the Möbius function of
a Boolean lattice satisfies the orthogonality relation Σ\_\{S:
\textbar S\textbar=k\} μ(S)μ(S') = 0 for \textbar S'\textbar{} ≠ k.

\subsection{14.6 Distributed
Responsibility}\label{distributed-responsibility}

\subsubsection{The Responsibility
Tensor}\label{the-responsibility-tensor}

When a collective acts --- and especially when it acts wrongly --- the
question of responsibility arises. Who bears the moral weight of the
collective's action?

\textbf{Definition 14.4 (Responsibility Tensor).} \emph{The
responsibility tensor of a collective} \(A\) \emph{for an outcome}
\(\omega\) \emph{is a} \((0,1)\)\emph{-tensor (covector) on the member
space:}

\[R_{a}(\omega) \in \lbrack 0,1\rbrack\]

\emph{assigning to each member} \(a\) \emph{a degree of responsibility
for} \(\omega\)\emph{. The responsibility tensor is a moral covector: it
weights the members, just as the interest covector weights the
dimensions of moral space.}

The responsibility tensor depends on:

\textbf{Causal contribution.} How much did member \(a\)'s actions
contribute to \(\omega\)? This is the causal component of responsibility
--- the counterfactual question of whether \(\omega\) would have
occurred without \(a\)'s contribution.

\textbf{Epistemic position.} Did member \(a\) know, or should they have
known, that \(\omega\) was a likely consequence? Culpable ignorance
(failure to know what one should have known) generates responsibility;
blameless ignorance mitigates it.

\textbf{Role authority.} Did member \(a\) have the authority and
capacity to prevent \(\omega\)? A supervisor with the power to stop a
harmful action bears more responsibility than a subordinate who lacked
the power to intervene.

\textbf{Degree of endorsement.} Did member \(a\) endorse the collective
action, or were they coerced, overruled, or absent? Active endorsement
amplifies responsibility; coerced participation mitigates it.

\subsubsection{Is Responsibility
Conserved?}\label{is-responsibility-conserved}

A fundamental question: does the sum of individual responsibilities
equal the collective responsibility?

\[\sum_{a \in A}^{}R_{a}(\omega)\overset{?}{=}R_{\text{collective}}(\omega)\]

If the answer is yes, responsibility is \emph{conserved}: it can be
fully distributed to individuals, with no remainder. If the answer is
no, there is a \emph{responsibility gap} (if the sum is less) or a
\emph{responsibility surplus} (if the sum is greater).

\textbf{Claim.} In the presence of emergent obligations
(\(\mathcal{E \neq}0\)), responsibility is generically \emph{not}
conserved. The gap:

\[\Delta R = R_{\text{collective}} - \sum_{a \in A}^{}R_{a}\]

is the \emph{collective responsibility remainder} --- the responsibility
that belongs to the collective as a structured entity, not to any
individual member. This remainder is the moral analogue of the emergent
component \(\mathcal{E}\): it exists because the collective has moral
agency that exceeds the sum of its members' agencies.

\textbf{Example.} A committee votes unanimously to implement a policy
that turns out to be harmful. Each member bears responsibility for their
vote (\(R_{a} > 0\) for all \(a\)). But the committee also bears
collective responsibility for the policy as a whole --- responsibility
that includes the failures of the committee's deliberative process, its
information-gathering procedures, and its institutional design. This
collective responsibility may exceed the sum of individual
responsibilities: \(\Delta R > 0\).

\textbf{Example.} A mob commits violence. Each individual participant
bears responsibility. But the mob --- as a collective agent with
emergent dynamics (escalation, deindividuation, collective action) ---
bears additional responsibility for creating the conditions under which
individual members acted as they did. The responsibility surplus
\(\Delta R > 0\) captures the moral content of incitement and collective
frenzy.

\subsubsection{The Complicity Gradient}\label{the-complicity-gradient}

In many collective situations, responsibility is not binary
(guilty/innocent) but continuously graded. The \emph{complicity
gradient} is the gradient of the responsibility function across the
member space:

\[\nabla_{a}R = \frac{\partial R}{\partial(\text{member }a\text{’s position})}\]

This gradient measures how much more (or less) responsible a member
becomes as their position in the collective changes --- as they move
from periphery to center, from subordinate to authority, from observer
to participant.

The complicity gradient is a moral covector on the member space: it
weights members by their marginal contribution to collective
responsibility. It provides a geometric decomposition of collective
responsibility that goes beyond the binary question of ``who is
responsible?'' to the structural question of ``how does responsibility
vary with position in the collective?''

\subsection{14.7 The Geometry of Institutional
Structure}\label{the-geometry-of-institutional-structure}

\subsubsection{Institutions as Moral
Manifolds}\label{institutions-as-moral-manifolds}

An institution is a collective agent with sufficient structure and
persistence to constitute a moral entity in its own right. Corporations,
states, universities, hospitals, professional associations, and
international bodies are all institutions in this sense.

The geometric framework suggests treating the institution itself as a
\emph{moral manifold} --- a space whose points are the possible
configurations of the institution (its policies, practices, norms, and
distributions of authority), and whose tangent vectors are possible
institutional changes.

\textbf{Definition 14.5 (Institutional Manifold).} \emph{The
institutional manifold} \(\mathcal{M}_{\text{inst}}\) \emph{of an
institution} \(\mathcal{I}\) \emph{is the space of the institution's
possible configurations. A point} \(q \in \mathcal{M}_{\text{inst}}\)
\emph{specifies:}

\begin{enumerate}
\def\labelenumi{\arabic{enumi}.}
\item
  \emph{The institution's policies and rules}
\item
  \emph{The distribution of authority among members}
\item
  \emph{The decision procedures in force}
\item
  \emph{The institution's relationship to external stakeholders}
\end{enumerate}

The moral manifold \(\mathcal{M}\) from Chapter 5 and the institutional
manifold \(\mathcal{M}_{\text{inst}}\) are distinct but related. They
are coupled by a \emph{projection}:

\[\pi:\mathcal{M}_{\text{inst}}\mathcal{\rightarrow M}\]

that maps each institutional configuration to the moral situation it
generates. Different institutional configurations (different policies,
different authority distributions) may generate the same moral situation
--- or different ones. The fiber \(\pi^{- 1}(p)\) over a moral point
\(p\) consists of all institutional configurations that produce
situation \(p\).

\subsubsection{Institutional Metrics}\label{institutional-metrics}

The institutional manifold carries its own metric --- the
\emph{institutional metric} \(h_{ij}\) --- that measures the ``cost'' of
institutional change:

\[ds_{\text{inst}}^{2} = h_{ij}\, dq^{i}\, dq^{j}\]

This metric encodes how difficult it is to change the institution along
various dimensions. Some institutional changes are easy (adjusting a
minor policy): low cost, small \(h_{ij}\). Others are hard
(restructuring the authority hierarchy): high cost, large \(h_{ij}\).
Some are nearly impossible (changing the institution's fundamental
purpose): the metric approaches infinity.

The institutional metric is related to the moral metric \(g_{\mu\nu}\)
through the projection \(\pi\):

\[g_{\mu\nu} = \frac{\partial q^{i}}{\partial x^{\mu}}\frac{\partial q^{j}}{\partial x^{\nu}}\, h_{ij}\]

The moral trade-off structure is the \emph{pullback} of the
institutional metric through the projection --- the moral costs of
change are determined, in part, by the institutional costs of the
changes required to effect them.

\subsubsection{Institutional Design as Geometric
Optimization}\label{institutional-design-as-geometric-optimization}

If the institutional manifold \(\mathcal{M}_{\text{inst}}\) is equipped
with a metric (the cost structure of change) and a potential (the
``moral energy'' of each configuration), then \emph{institutional
design} is a problem in Riemannian optimization: find the configuration
\(q^{*} \in \mathcal{M}_{\text{inst}}\) that minimizes the moral
potential, subject to the institutional constraints.

This is the geometric analogue of \emph{mechanism design} in economics:
designing institutions to achieve desired outcomes, given that the
institution's internal structure mediates between individual actions and
collective results.

The geometric framework adds precision: the institutional metric
determines which changes are feasible, the curvature of
\(\mathcal{M}_{\text{inst}}\) determines whether optimization is
well-behaved (convex) or pathological (multiple local minima, saddle
points), and the projection \(\pi\) determines how institutional changes
translate into moral outcomes.

\subsection{14.8 Collective Moral
Dynamics}\label{collective-moral-dynamics}

\subsubsection{Parallel Transport of Collective
Obligations}\label{parallel-transport-of-collective-obligations}

Chapter 10 developed the parallel transport of individual obligations
across changing circumstances. The same formalism extends to collective
obligations, but with additional structure arising from the collective's
internal dynamics.

When a collective obligation is parallel-transported along a path in
moral space, the \emph{internal structure of the collective may also
change}. The obligation to defend citizens is parallel-transported from
peacetime to wartime --- but the internal structure of the state changes
simultaneously (emergency powers, reallocation of authority,
mobilization of resources). The transport of the obligation and the
transformation of the structure are coupled.

\textbf{Definition 14.6 (Collective Parallel Transport).} \emph{The
parallel transport of a collective obligation} \(O_{\text{coll}}^{\mu}\)
\emph{along a path} \(\gamma\) \emph{in} \(\mathcal{M}\) \emph{is
governed by:}

\[\nabla_{\dot{\gamma}}O_{\text{coll}}^{\mu} + \Gamma_{\nu\rho}^{\mu}\,{\dot{\gamma}}^{\nu}\, O_{\text{coll}}^{\rho} + \Lambda_{a}^{\mu}\,{\dot{\Sigma}}_{a} = 0\]

\emph{where} \(\Gamma_{\nu\rho}^{\mu}\) \emph{is the Levi-Civita
connection (as in Chapter 10),} \(\Lambda_{a}^{\mu}\) \emph{is the
structural coupling --- encoding how changes in the structure tensor}
\(\Sigma\) \emph{affect the collective obligation --- and}
\({\dot{\Sigma}}_{a}\) \emph{is the rate of change of the structure.}

The additional term \(\Lambda_{a}^{\mu}{\dot{\Sigma}}_{a}\) is the key
difference between individual and collective parallel transport. When
the collective's internal structure is stable (\(\dot{\Sigma} = 0\)),
collective transport reduces to individual transport. When the structure
is changing, the collective obligation is pulled in directions
determined by the structural coupling \(\Lambda\) --- directions that
may not align with the individual transport.

\subsubsection{Collective Holonomy}\label{collective-holonomy}

The holonomy of a collective obligation around a closed loop may include
contributions from the structural coupling:

\[H_{\gamma}^{\text{coll}} = H_{\gamma}^{\text{ind}} + \oint\Lambda_{a}^{\mu}\, d\Sigma_{a}\]

The second term is the \emph{structural holonomy} --- the contribution
of internal reorganization to the collective obligation's rotation
around the loop. A collective that reorganizes as it traverses a moral
circuit (e.g., a corporation that restructures during a crisis) may
acquire collective holonomy even if the individual-level holonomy is
zero.

\textbf{Example.} A university faces a controversy, reorganizes its
governance in response, and returns to ``normal operations.'' The
structural holonomy measures how the governance reorganization has
permanently altered the university's moral orientation --- even though
the explicit policies may look the same as before. The collective
obligation has been rotated by the structural journey.

\subsection{14.9 Collective Agents in the Nine
Dimensions}\label{collective-agents-in-the-nine-dimensions}

\subsubsection{The Collective Scope}\label{the-collective-scope}

Chapter 5 (§5.2) organized the nine moral dimensions via a 3 × 3 matrix:
three ethical scopes (Individual, Relational, Collective) × three
epistemic modes (What Matters, Who Decides, What We Know). The
``Collective'' row --- Dimensions 3, 8, 6 (Justice/Fairness, Procedural
Legitimacy, Societal/Environmental Impact) --- is precisely the row that
collective agents operate in most naturally.

Individual agents primarily generate obligations along the Individual
and Relational rows. Collective agents primarily generate obligations
along the Collective row --- and they do so with a strength and
specificity that individual agents cannot match. The obligation of a
state to provide justice (Dimension 3) is stronger, more specific, and
more structurally grounded than any individual citizen's obligation to
promote justice. The obligation of a corporation to ensure procedural
legitimacy (Dimension 8) is an obligation of the institution, not of any
individual officer.

\subsubsection{The Collective-Individual
Coupling}\label{the-collective-individual-coupling}

The relationship between collective and individual obligations is not
one-directional. Collective obligations \emph{feed back} into individual
obligations through the structure tensor.

\textbf{Downward causation.} A collective obligation
\(O_{\text{coll}}^{\mu}\) generates individual obligations through the
structure tensor:

\[O_{a}^{\mu} = \Sigma_{ab}\,\mathcal{C}_{\nu}^{- 1\mu}\, O_{\text{coll}}^{\nu}\]

Each member acquires an individual obligation that is their ``share'' of
the collective obligation, weighted by their structural position. A
firefighter's individual obligation to rescue (Dimension 1) is generated
by the fire department's collective obligation to protect, mediated by
the firefighter's role in the department's structure.

\textbf{Upward emergence.} Individual moral actions aggregate (through
the collective agency tensor) into collective moral properties:

\[O_{\text{coll}}^{\mu} = \mathcal{C}_{\nu_{1}\cdots\nu_{n}}^{\mu}\, O_{\nu_{1}}^{(1)}\,\cdots\, O_{\nu_{n}}^{(n)}\]

The collective is constituted by individual actions but not reducible to
them. The interplay between downward causation and upward emergence
creates a feedback loop: collective obligations shape individual ones,
which in turn constitute the collective. This circularity is not vicious
--- it is the structural content of the sociological fact that
institutions and individuals mutually constitute each other.

\subsection{14.10 Emergent Moral
Properties}\label{emergent-moral-properties}

\subsubsection{Properties That Exist Only at the Collective
Level}\label{properties-that-exist-only-at-the-collective-level}

Some moral properties have no individual analogue. They exist only at
the collective level and resist decomposition.

\textbf{Justice.} Justice is a property of social arrangements --- of
how benefits and burdens are distributed across a population. An
individual can \emph{act justly} (treating each person fairly in
individual interactions), but justice as a \emph{structural property}
--- the fairness of the overall distribution --- is a collective
attribute. In geometric terms, justice is a function on
\(\mathcal{M}_{\text{inst}}\), not on the individual moral manifold.

\textbf{Legitimacy.} Political legitimacy --- the right of an
institution to make binding decisions --- is a collective property that
no individual possesses. A judge has authority by virtue of the judicial
system's legitimacy, not by individual qualities. Legitimacy is an
emergent property of the institution's structure tensor and its
relationship to the governed.

\textbf{Democratic representation.} The property of a decision's being
democratically legitimate requires a \emph{process} involving multiple
agents, structured by rules of participation, deliberation, and voting.
No individual action can be ``democratic'' --- democracy is irreducibly
collective.

\subsubsection{The Emergent Properties
Criterion}\label{the-emergent-properties-criterion}

\textbf{Proposition 14.2 (Criterion for Emergent Moral Property).}
\emph{A moral property} \(P\) \emph{is emergent at the collective level
if and only if:}

\begin{enumerate}
\def\labelenumi{\arabic{enumi}.}
\item
  \(P\) \emph{is well-defined on collective moral states:}
  \(P:\mathcal{M}_{\text{inst}}\mathbb{\rightarrow R}\)
\item
  \(P\) \emph{is not expressible as a function of individual moral
  properties alone:} \(P \neq f(P_{1},P_{2},\ldots,P_{n})\) \emph{for
  any function} \(f\) \emph{and any set of individual-level properties}
  \(\{ P_{a}\}\)
\item
  \(P\) \emph{depends on the structure tensor:}
  \(\partial P/\partial\Sigma_{ab} \neq 0\) \emph{for some} \(a,b\)
\end{enumerate}

Condition (3) is the geometric marker of emergence: the property depends
on the relationships between members, not just on the members
themselves.

\emph{Proof.} Necessity. (1) If P is not well-defined on M\_inst, it
cannot be a property of collectives at all. (2) If P = f(P₁, \ldots,
P\_n) for some function f and individual-level properties \{P\_a\}, then
P is reducible to individual properties and is not emergent by
definition. (3) If ∂P/∂Σ\_\{ab\} = 0 for all a, b, then P does not
depend on the structure tensor and hence depends only on properties of
the individual members, reducing to case (2). Sufficiency. Suppose all
three conditions hold. By (1), P is a well-defined collective property.
By (2), P is not expressible as a function of individual properties
alone. By (3), the irreducibility is structural --- P depends on the
relationships between members, not just their individual states. Hence P
is emergent at the collective level. □

\subsection{14.11 AI Systems as Collective
Agents}\label{ai-systems-as-collective-agents}

\subsubsection{Human-AI Collectives}\label{human-ai-collectives}

As AI systems take on morally significant roles, the question of
collective agency becomes urgent. A human-AI team --- a physician using
a diagnostic AI, a judge using a sentencing algorithm, a pilot using an
autopilot --- is a collective agent whose moral structure requires
analysis.

The collective agency tensor of a human-AI team has the form:

\[\mathcal{C}_{\nu_{H}\nu_{AI}}^{\mu}\]

where \(\nu_{H}\) indexes the human's interest dimensions and
\(\nu_{AI}\) indexes the AI's ``interest'' dimensions (more precisely,
the objective-function dimensions that the AI is designed to optimize).

The emergent component \(\mathcal{E}\) of this tensor is particularly
important for AI alignment. If \(\mathcal{E \neq}0\), the human-AI team
generates collective obligations that neither the human nor the AI bears
individually. These emergent obligations may be \emph{desirable} (the
team can achieve outcomes neither could alone) or \emph{dangerous} (the
team generates harms that neither would generate alone).

\subsubsection{Responsibility Distribution in Human-AI
Teams}\label{responsibility-distribution-in-human-ai-teams}

When a human-AI team produces a harmful outcome, the responsibility
tensor \(R_{a}(\omega)\) must be defined for both the human member and
the AI member. But the AI's ``responsibility'' is different in character
from the human's:

The AI has no moral autonomy --- its actions are determined by its
training, its design, and its deployment. Attributing responsibility to
the AI is a shorthand for attributing responsibility to the \emph{chain
of human decisions} that created, trained, deployed, and maintained the
AI. The responsibility tensor for the AI component decomposes into:

\[R_{AI}(\omega) = R_{\text{designer}} + R_{\text{trainer}} + R_{\text{deployer}} + R_{\text{operator}} + R_{\text{emergent}}\]

where \(R_{\text{emergent}}\) is the residual responsibility that cannot
be attributed to any individual in the chain. This residual is the
collective responsibility of the \emph{AI development ecosystem} --- the
responsibility that arises from the interactions between design
decisions, training data choices, deployment contexts, and operational
parameters.

\subsubsection{The Alignment Problem as Emergent
Agency}\label{the-alignment-problem-as-emergent-agency}

The AI alignment problem can be restated in the language of collective
agency: ensuring that the emergent component \(\mathcal{E}\) of the
human-AI collective agency tensor is benign --- that the collective
obligations and capabilities generated by the human-AI interaction are
aligned with human values.

This restatement is productive because it shifts the focus from
individual AI behavior (is this AI aligned?) to collective structure (is
this human-AI system producing well-aligned collective agency?). An AI
that behaves perfectly in isolation may generate harmful emergent
properties when embedded in a collective --- and conversely, an AI with
individual limitations may contribute to well-aligned collective agency
when embedded in an appropriate structure.

The structure tensor \(\Sigma\) for human-AI teams encodes the division
of authority: which decisions are made by the human, which by the AI,
and which are shared. The antisymmetric part
\(\Sigma_{\lbrack H,AI\rbrack}\) encodes the power differential. The BIP
(Chapter 5) constrains the structure: the collective's moral evaluations
must be invariant under admissible re-descriptions, regardless of
whether the human or the AI is performing the evaluation.

\subsection{14.12 Worked Example: The Hospital Ethics
Committee}\label{worked-example-the-hospital-ethics-committee}

\subsubsection{Setup}\label{setup-4}

A hospital ethics committee consists of five members: a physician
(\(P\)), a nurse (\(N\)), a patient advocate (\(A\)), an ethicist
(\(E\)), and a hospital administrator (\(D\)). They must decide whether
to continue life-sustaining treatment for a patient whose family is
divided.

\subsubsection{The Individual Evaluation
Tensors}\label{the-individual-evaluation-tensors}

Each member has an evaluation tensor \(E_{\nu}^{\mu}(m)\) reflecting
their professional perspective:

The physician \(P\) weights welfare (Dimension 1) and epistemic caution
(Dimension 9) most heavily:

\[E_{\nu}^{\mu}(P):\text{strong on }E_{1}^{1},E_{9}^{1}\]

The nurse \(N\) weights care (Dimension 7) and welfare (Dimension 1):

\[E_{\nu}^{\mu}(N):\text{strong on }E_{7}^{7},E_{7}^{1}\]

The advocate \(A\) weights autonomy (Dimension 4) and rights (Dimension
2):

\[E_{\nu}^{\mu}(A):\text{strong on }E_{4}^{4},E_{2}^{2}\]

The ethicist \(E\) weights across all dimensions but especially justice
(Dimension 3) and duty (Dimension 2):

\[E_{\nu}^{\mu}(E):\text{distributed, strongest on }E_{3}^{3},E_{2}^{2}\]

The administrator \(D\) weights procedural legitimacy (Dimension 8) and
societal impact (Dimension 6):

\[E_{\nu}^{\mu}(D):\text{strong on }E_{8}^{8},E_{6}^{6}\]

\subsubsection{The Structure Tensor}\label{the-structure-tensor-1}

The committee's structure tensor encodes:

\begin{itemize}
\item
  The physician has clinical authority: \(\Sigma_{Pm} > \Sigma_{mP}\)
  for all \(m \neq P\) on clinical questions
\item
  The ethicist has ethical authority: \(\Sigma_{Em} > \Sigma_{mE}\) for
  all \(m\) on ethical frameworks
\item
  The advocate has veto power on autonomy grounds: \(\Sigma_{Am}\)
  spikes when Dimension 4 is at stake
\item
  The administrator has institutional authority: \(\Sigma_{Dm}\) is
  strong on Dimensions 6 and 8
\item
  The nurse has equal voice: \(\Sigma_{Nm} \approx \Sigma_{mN}\) ---
  symmetric relationships
\end{itemize}

\subsubsection{The Emergent Component}\label{the-emergent-component}

The committee has obligations that no individual member has:

\begin{enumerate}
\def\labelenumi{\arabic{enumi}.}
\item
  \textbf{Procedural obligation.} The committee must follow a fair
  deliberative process --- hear all perspectives, give each adequate
  weight, reach a decision by an appropriate mechanism. This obligation
  belongs to the committee, not to any individual member.
\item
  \textbf{Documentation obligation.} The committee must record its
  reasoning in a way that can be reviewed. This is an institutional
  obligation arising from the committee's role in hospital governance.
\item
  \textbf{Precedent obligation.} The committee's decision establishes
  precedent for future cases. This obligation --- to decide in a way
  that is generalizable --- is emergent: it arises from the committee's
  institutional role, not from any member's individual duty.
\end{enumerate}

These obligations appear in the emergent component \(\mathcal{E}\) of
the collective agency tensor. They cannot be attributed to any
individual member, yet they are real obligations that constrain the
committee's action.

\subsubsection{The Decision}\label{the-decision}

The committee must contract its collective tensor to a scalar decision:
continue treatment, withdraw treatment, or modify the treatment plan.
The contraction involves:

\begin{enumerate}
\def\labelenumi{\arabic{enumi}.}
\item
  Each member contributes their evaluation tensor \(E_{\nu}^{\mu}(m)\)
\item
  The structure tensor \(\Sigma_{mm'}\) weights the contributions
\item
  The emergent component \(\mathcal{E}\) adds the collective obligations
\item
  The full collective agency tensor \(\mathcal{C}\) is contracted with
  the patient's and family's interest covectors
\item
  The resulting satisfaction scores for each option are compared
\end{enumerate}

The committee decides to modify the treatment plan --- a compromise that
respects the physician's clinical judgment, the advocate's autonomy
concerns, and the ethicist's framework, while satisfying the procedural
and documentation obligations that the committee bears collectively.

\subsubsection{Responsibility for the
Outcome}\label{responsibility-for-the-outcome}

If the outcome is poor, the responsibility tensor distributes
accountability:

\[R_{P} = 0.30,\quad R_{N} = 0.15,\quad R_{A} = 0.20,\quad R_{E} = 0.15,\quad R_{D} = 0.10\]

The sum is 0.90. The collective responsibility is 1.00. The remainder
\(\Delta R = 0.10\) is the \emph{committee's} responsibility as an
institution --- responsibility for the quality of the deliberative
process, the adequacy of the information considered, and the
appropriateness of the decision procedure. This responsibility belongs
to the committee, not to any member.

\subsection{14.13 Summary}\label{summary-7}

{\def\LTcaptype{none} % do not increment counter
\begin{longtable}[]{@{}
  >{\raggedright\arraybackslash}p{(\linewidth - 4\tabcolsep) * \real{0.2513}}
  >{\raggedright\arraybackslash}p{(\linewidth - 4\tabcolsep) * \real{0.3363}}
  >{\raggedright\arraybackslash}p{(\linewidth - 4\tabcolsep) * \real{0.4124}}@{}}
\toprule\noalign{}
\begin{minipage}[b]{\linewidth}\raggedright
\textbf{Concept}
\end{minipage} & \begin{minipage}[b]{\linewidth}\raggedright
\textbf{Mathematical Object}
\end{minipage} & \begin{minipage}[b]{\linewidth}\raggedright
\textbf{Moral Content}
\end{minipage} \\
\midrule\noalign{}
\endhead
\bottomrule\noalign{}
\endlastfoot
Collective agency tensor \(\mathcal{C}\) & \((1,n)\)-tensor & Map from
member interests to collective obligation \\
Emergent component \(\mathcal{E}\) & Remainder after individual
decomposition & Obligations the collective has that no member has \\
Structure tensor \(\Sigma_{ab}\) & \((0,2)\)-tensor on member space &
Internal organization: authority, roles, relationships \\
Power differential \(\Sigma_{\lbrack ab\rbrack}\) & Antisymmetric part
of \(\Sigma\) & Structural inequity between members \\
Degree of emergence \(\eta\) & Norm ratio & How much of collective
agency exceeds individual sum \\
Responsibility tensor \(R_{a}\) & \((0,1)\)-tensor (covector) on member
space & Distribution of accountability across members \\
Responsibility remainder \(\Delta R\) & R\_coll − Σ R\_a & Collective
responsibility not attributable to individuals \\
Institutional manifold \(\mathcal{M}_{\text{inst}}\) & Configuration
space of institution & Space of possible institutional structures \\
Institutional metric \(h_{ij}\) & Metric on
\(\mathcal{M}_{\text{inst}}\) & Cost structure of institutional
change \\
Complicity gradient \(\nabla_{a}R\) & Gradient of responsibility & How
responsibility varies with position in the collective \\
\end{longtable}
}

Collective moral agency is not a philosophical fiction or a convenient
shorthand. It is a structural feature of the moral landscape, encoded in
the emergent components of the collective agency tensor. Collectives
have obligations that no member bears, responsibilities that no member
exhausts, and properties --- justice, legitimacy, democratic
representation --- that exist only at the collective level.

The geometric framework provides the vocabulary for analyzing collective
agency with precision: - The collective agency tensor \(\mathcal{C}\)
captures the full structure of collective moral action - Its
decomposition into individual and emergent components shows where
collective agency exceeds the sum of its parts - The structure tensor
\(\Sigma\) encodes the internal organization that mediates between
individual and collective - The responsibility tensor distributes
accountability while identifying the collective remainder - The
institutional manifold provides a geometric setting for institutional
design

Chapter 15 turns from the \emph{structure} of moral evaluation to the
\emph{moment of choice}: contraction --- the mathematically necessary,
informationally lossy process by which the rich tensor structure of
moral reality is compressed into a scalar action-decision.

\subsection{Technical Appendix}\label{technical-appendix-8}

\textbf{Proposition 14.3 (Decomposition Theorem).} \emph{Every
collective agency tensor} \(\mathcal{C}_{\nu_{1}\cdots\nu_{n}}^{\mu}\)
\emph{admits a unique decomposition:}

\[\mathcal{C}_{\nu_{1}\cdots\nu_{n}}^{\mu} = \sum_{k = 1}^{n}\mathcal{C}_{(k)}^{\mu} \otimes \delta_{\nu_{1}} \otimes \cdots \otimes {\widehat{\delta}}_{\nu_{k}} \otimes \cdots \otimes \delta_{\nu_{n}} + \sum_{j < k}^{}\mathcal{C}_{(jk)}^{\mu} + \cdots + \mathcal{E}_{\nu_{1}\cdots\nu_{n}}^{\mu}\]

\emph{where} \(\mathcal{C}_{(k)}^{\mu}\) \emph{is agent} \(k\)\emph{'s
individual contribution,} \(\mathcal{C}_{(jk)}^{\mu}\) \emph{is the
pairwise interaction between agents} \(j\) \emph{and} \(k\)\emph{, and}
\(\mathcal{E}\) \emph{is the fully irreducible remainder. This is the
moral analogue of the cluster expansion in statistical mechanics.}

\emph{\textbf{Proof}. Define the terms by inclusion-exclusion. For each
subset S ⊆ \{1,\ldots,n\}, the connected \textbar S\textbar-body term
is: C\_(S)ᵘ = Tr\_\{\textbackslash S\} Cᵘ\_\{ν₁⋯ν\_n\} −
Σ\_\{S\textquotesingle{} ⊂ S\} C\_(S\textquotesingle)ᵘ ⊗ ⨂\_\{k ∉
S\textquotesingle\} δ\_\{ν\_k\}, where Tr\_\{\textbackslash S\} traces
over all indices not in S. The 1-body terms are the marginals C\_(k)ᵘ =
Tr\_\{\textbackslash\{k\}\} Cᵘ. The 2-body terms are the connected
correlations: C\_(jk)ᵘ = Tr\_\{\textbackslash\{j,k\}\} Cᵘ − C\_(j)ᵘ ⊗
δ\_\{ν\_k\} − δ\_\{ν\_j\} ⊗ C\_(k)ᵘ. The remainder E =
C\_(\{1,\ldots,n\})ᵘ is the fully connected n-body term. Uniqueness:
Each m-body term is determined recursively from lower-order terms, so
the decomposition is unique. The projections Π\_S: C ↦ C\_(S) satisfy
Π\_S Π\_\{S\textquotesingle\} = δ\_\{SS\textquotesingle\}Π\_S
(orthogonality) and Σ\_S Π\_S = id (completeness). □}\(n\mathcal{E}g▫\)

\textbf{Proposition 14.4 (Responsibility Inequality).} \emph{For a
collective with emergent obligations (}\(\mathcal{E \neq}0\)\emph{), the
collective responsibility for an outcome} \(\omega\) \emph{satisfies:}

\[R_{\text{collective}}(\omega) \geq \sum_{a \in A}^{}R_{a}(\omega)\]

\emph{with equality if and only if} \(\omega\) \emph{is attributable
entirely to individual actions (the emergent component makes no causal
contribution to} \(\omega\)\emph{).}

\emph{\textbf{Proof}. The collective responsibility decomposes as
R\_collective(ω) = Σ\_\{a ∈ A\} R\_a(ω) + R\_E(ω), where R\_a(ω) is the
responsibility attributable to agent a's individual contribution and
R\_E(ω) is the responsibility arising from the emergent component E. The
emergent responsibility R\_E(ω) ≥ 0 because it is the causal
contribution of a structural feature to an outcome, which is
non-negative by Definition 14.4 (R\_a ∈ {[}0,1{]}). Hence
R\_collective(ω) = Σ\_a R\_a(ω) + R\_E(ω) ≥ Σ\_a R\_a(ω), with equality
iff R\_E(ω) = 0 --- i.e., iff the emergent component makes no causal
contribution to ω. By the Decomposition Theorem (Prop. 14.3), this
occurs iff ω is fully attributable to individual actions.
□}\(\mathcal{E}\omega\Delta R \geq 0\mathcal{E▫}\)

\textbf{Proposition 14.5 (Structural Coupling of Parallel Transport).}
\emph{For a collective whose structure tensor} \(\Sigma\) \emph{changes
along a path} \(\gamma\)\emph{, the parallel transport of the collective
obligation deviates from the parallel transport of the individual
obligations by:}

\[\delta O_{\text{coll}}^{\mu} = \Lambda_{a}^{\mu}\int_{\gamma}^{}d\Sigma_{a} + \text{higher order}\]

\emph{where}
\(\Lambda_{a}^{\mu} = \partial\mathcal{C}^{\mu}/\partial\Sigma_{a}\)
\emph{is the structural coupling. If} \(\Sigma\) \emph{is constant
along} \(\gamma\)\emph{, the collective transport reduces to the
individual transport.}

\emph{\textbf{Proof}.} The collective obligation along a path γ(t)
satisfies O\_collᵘ(t) = Cᵘ(\{O\_a(t)\}, Σ(t)). Differentiating along γ:
DO\_collᵘ/dt = Σ\_a (∂Cᵘ/∂O\_aᵛ)(DO\_aᵛ/dt) + (∂Cᵘ/∂Σ\_a)(dΣ\_a/dt). If
each individual obligation is parallel-transported (DO\_aᵛ/dt = 0), the
first sum vanishes, leaving DO\_collᵘ/dt = Λ\_aᵘ dΣ\_a/dt where Λ\_aᵘ =
∂Cᵘ/∂Σ\_a. Integrating: δO\_collᵘ = Λ\_aᵘ ∫\_γ dΣ\_a + O((ΔΣ)²), where
the higher-order terms arise from the Taylor expansion of Cᵘ in Σ. If Σ
is constant (dΣ\_a/dt = 0), then δO\_collᵘ = 0. □

❖

\emph{The individual is where ethics begins: in the recognition of
obligation, the weighing of interests, the exercise of judgment.}

\emph{The collective is where ethics scales: in the institutions that
structure our shared life, the corporations that shape our economies,
the states that govern our polities, the professions that organize our
expertise.}

\emph{The transition from individual to collective is not merely
additive. It is tensorial --- a tensor product that creates new
dimensions of moral agency, new obligations that no individual bears,
new responsibilities that no individual exhausts.}

\emph{The mathematics of this transition --- the collective agency
tensor, the structure tensor, the emergent component, the responsibility
covector --- provides the vocabulary for analyzing collective morality
with the same precision that the tensor hierarchy of Chapter 6 brings to
individual morality.}

\emph{The collective is not a moral fiction.}

\emph{It is a higher-rank moral entity.}

\emph{And its geometry matters.}

\section{Chapter 15: From Tensor to Decision --- The Philosophy of Moral
Contraction}\label{chapter-15-from-tensor-to-decision-the-philosophy-of-moral-contraction}

\begin{quote}
\textbf{\textsc{RUNNING EXAMPLE --- Priya's Model}}

\emph{The Board meeting arrives. Sarah Chen must present TrialMatch's
fairness assessment. She has Priya's full tensor analysis:
nine-dimensional obligation vectors, uncertainty matrices, metric
concerns. The Board wants one slide. Sarah contracts: she picks a
weighted sum, weighting medical dimensions heavily and access dimensions
lightly. The resulting scalar looks acceptable. By the Scalar
Irrecoverability Theorem, the Board cannot recover from this number the
fact that rural patients are systematically excluded. Priya watches from
the back of the room. The moral residue---what remains after
contraction---sits heavy. Sarah's contraction was technically valid. It
was also a choice about which dimensions to sacrifice. That choice was
never discussed.}
\end{quote}

\subsection{15.1 The Moment of Choice}\label{the-moment-of-choice}

Throughout this book, we have built a tensorial picture of moral
reality: multi-dimensional, perspectival, geometrically structured.
Obligations point in specific directions. Interests weight dimensions.
The metric encodes trade-offs. The stratified manifold marks boundaries
where the rules change. Curvature makes evaluation path-dependent.
Conservation laws constrain what re-description can alter. Superposition
allows framings to coexist and interfere. Collective agency generates
obligations that no individual bears.

This apparatus captures the \emph{structure} of moral evaluation with a
precision that scalar frameworks cannot match. But at the moment of
decision, we must \emph{act} --- and action is one-dimensional. We do
this, or we do that. The rich, multi-dimensional structure of moral
reality must be compressed into a single scalar verdict: this option is
best; this action is required; this policy is adopted.

The mathematical operation that performs this compression is
\emph{contraction}: the summation over paired indices that reduces
tensor rank. The fundamental formula \(S = I_{\mu}O^{\mu}\) (Chapter 6)
is already a contraction --- the pairing of an interest covector with an
obligation vector to produce a scalar. But it is not the only
contraction, nor is it uniquely determined by the tensorial structure.
Different contractions yield different verdicts, and the choice of
contraction is itself a moral commitment.

This chapter examines contraction philosophically. When is contraction
necessary? Who chooses how to contract? What information does
contraction discard? And what obligations persist after contraction ---
the \emph{moral residue} of what was contracted away?

\subsection{15.2 The Mathematics of
Contraction}\label{the-mathematics-of-contraction}

\subsubsection{Contraction as Index
Summation}\label{contraction-as-index-summation}

Recall from Chapter 4 that contraction reduces tensor rank by
identifying an upper index with a lower index and summing:

For a \((1,1)\)-tensor \(T_{\nu}^{\mu}\), contraction over \(\mu\) and
\(\nu\) yields a scalar:

\[\text{Tr}(T) = T_{\mu}^{\mu} = \sum_{\mu = 1}^{9}T_{\mu}^{\mu}\]

More generally, contracting a type \((p,q)\)-tensor with a type
\((r,s)\)-tensor over \(k\) pairs of indices yields a type
\((p + r - k,q + s - k)\)-tensor. Each contraction reduces the total
rank by 2 --- one upper and one lower index consumed --- and discards
the information in those dimensions.

\subsubsection{The Contraction Chain}\label{the-contraction-chain}

The full chain from tensorial moral reality to scalar action-decision
typically involves multiple contractions:

\textbf{Step 1: Multi-agent to single-agent.} The multi-agent evaluation
tensor \(M_{ia}\) (Chapter 7, §7.3) is contracted over the agent index
\(a\) --- a social aggregation procedure that reduces the
multi-perspective tensor to a single-perspective evaluation.

\textbf{Step 2: Multi-option to single-option.} The resulting vector
\(S_{i}\) (one score per option) is contracted to a scalar verdict ---
selecting the option with the highest score, or the option above a
threshold, or the option satisfying a lexicographic condition.

\textbf{Step 3: Multi-dimensional to scalar.} Within each step, the
fundamental contraction \(S = I_{\mu}O^{\mu}\) reduces the
nine-dimensional moral structure to a single number.

Each step discards information. The multi-agent tensor preserves the
full landscape of agreement and disagreement; the single-agent vector
discards it. The multi-option vector preserves how close the
alternatives were; the scalar verdict discards it. The nine-dimensional
structure preserves which dimensions contributed; the scalar contraction
discards it.

\subsubsection{Non-Commutativity of
Contraction}\label{non-commutativity-of-contraction}

When multiple contractions are performed, the order may matter.
Contracting over agents first and then over options may yield a
different result from contracting over options first and then over
agents.

\textbf{Example from Chapter 7.} Consider the evaluation tensor
\(M_{ia}\) with three patients (i = Alice, Bob, Carol) and five
perspectives (a = physician, families, committee). Contracting first
over agents (social aggregation) and then selecting the maximum yields
one verdict. Selecting the best option from each agent's perspective
first, and then aggregating the verdicts, may yield a different verdict
--- particularly when preferences are non-aligned.

This non-commutativity is not a bug in the framework. It is a feature:
it makes visible the fact that \emph{how we aggregate perspectives and
how we select options are independent moral choices}, and that their
order of application matters. Arrow's impossibility theorem (1951) is,
in this light, a result about the non-commutativity of contractions on
preference tensors.

\subsection{15.3 The Necessity of
Contraction}\label{the-necessity-of-contraction}

\subsubsection{Why Contract at All?}\label{why-contract-at-all}

If contraction discards information --- and it does, inevitably --- why
not maintain the full tensorial richness?

\textbf{Practical necessity.} At the moment of action, we cannot act on
a tensor. We can only do one thing. The effector systems of moral agency
--- whether muscles, speech, resource allocation, or institutional
machinery --- require a single command. The kidney goes to one patient.
The vote is cast one way. The policy is adopted or rejected. The world
changes one step at a time, and each step is scalar.

\textbf{Cognitive necessity.} Even in deliberation, we cannot hold the
full tensor in view. A nine-dimensional obligation vector with five
agent perspectives has 45 independent components. The full moral tensor
for a realistic situation --- with uncertainty, multiple options,
temporal structure, and collective dimensions --- may have thousands of
components. Human working memory holds roughly four items. Contraction
is how we manage the gap between moral complexity and cognitive
capacity.

\textbf{Communicative necessity.} When we justify our actions to others,
we typically offer contracted reasons: ``this was the best option,'' ``I
had no choice,'' ``the benefits outweighed the costs.'' Full tensorial
explanations --- ``the obligation vector had components (0.3, 0.7, 0.1,
\ldots) and the interest covector weighted them as (0.2, 0.4, 0.1,
\ldots)'' --- are available in principle (and Chapter 18 argues that AI
systems should provide them), but they are not how moral communication
naturally works.

\textbf{The contraction bottleneck.} Between moral reality (tensorial)
and moral action (scalar), there is a bottleneck. Contraction is how we
pass through it. The bottleneck is not a deficiency of our framework or
our cognition. It is a structural feature of the relationship between
evaluation and action: evaluation is multi-dimensional; action is
one-dimensional.

\subsubsection{The Cost of Premature
Contraction}\label{the-cost-of-premature-contraction}

If contraction is necessary, its \emph{timing} is not. Contracting too
early --- collapsing the tensor to a scalar before all relevant
information has been considered --- discards information that might have
changed the verdict.

\textbf{Example.} A policy analyst reduces a multi-dimensional impact
assessment to a single dollar figure (cost-benefit analysis) before the
assessment is complete. The contraction (monetization of all dimensions)
discards information about distributional effects, rights implications,
and environmental impacts that have not yet been quantified. The
resulting verdict is not wrong --- it is \emph{premature}. It reflects a
contraction performed before the tensor was fully populated.

The geometric framework counsels \emph{deferred contraction}: maintain
the tensorial structure as long as feasible, contract only when action
requires it, and acknowledge what the contraction discards. This is a
structural argument for a familiar practical wisdom: don't simplify too
early; don't decide before you must; keep options open until the cost of
openness exceeds the cost of closure.

\subsection{15.4 Who Contracts?}\label{who-contracts}

\subsubsection{The Agent}\label{the-agent}

Often the individual agent chooses how to contract --- which dimensions
to weight, which perspectives to prioritize, which trade-offs to accept.
This is the sense in which moral choice is ``up to us'' even when the
tensorial structure of the situation is objective. Two agents facing the
same obligation field, with the same interests at stake, may contract
differently --- and both may be acting permissibly, if the choice of
contraction falls within the range of legitimate options.

\textbf{Example.} Two physicians, facing the same allocation decision,
may weight urgency and prognosis differently (different interest
covectors \(I_{\mu}\)). The resulting contractions
\(S = I_{\mu}O^{\mu}\) yield different verdicts. Neither is necessarily
wrong; they have made different --- but both legitimate --- contraction
choices.

\subsubsection{The Community}\label{the-community}

Social norms may specify how to contract. A community's moral
conventions --- which considerations have priority, how competing values
are balanced, what counts as ``good enough'' --- are, in tensorial
terms, a socially determined contraction procedure.

\textbf{Example.} In many communities, a promise creates a lexicographic
priority: the duty to keep the promise (Dimension 2) is contracted
first, and other considerations are relevant only if the promise is
fulfilled or legitimately overridden. This is a community-level
contraction norm --- a shared convention about the order and weighting
of index summation.

\subsubsection{The Institution}\label{the-institution}

Laws, regulations, and institutional rules perform contraction. A speed
limit contracts the complex tensor of traffic safety, efficiency,
freedom, enforcement cost, and environmental impact into a single
number. A hiring rubric contracts the multi-dimensional evaluation of
candidates into a ranked list. A sentencing guideline contracts the
complex moral landscape of criminal justice into a range of acceptable
punishments.

Institutional contraction has a distinctive feature: it is
\emph{explicit}. The contraction procedure is written down, debatable,
and (in principle) revisable. This makes institutional contraction more
transparent than individual or community contraction --- and connects it
to the governance account of Chapter 9. The governance account says: the
metric is the output of legitimate institutional processes. The same
applies to contraction: the legitimate contraction procedure is the one
that emerges from legitimate governance.

\subsubsection{The Algorithm}\label{the-algorithm}

In AI systems, the contraction is specified by the designer, learned
from data, or negotiated by governance. The choice of loss function
\emph{is} a choice of contraction: a loss function takes a
multi-dimensional output space and produces a scalar to be minimized.
Different loss functions embody different moral contractions.

\[\text{Loss function} = \text{contraction of the moral tensor to a scalar optimization target}\]

This observation is central to Chapter 18's treatment of AI alignment. A
misaligned AI is, in many cases, an AI performing the \emph{wrong
contraction} --- optimizing a scalar that is a poor compression of the
tensorial moral reality. Alignment, in this vocabulary, is getting the
contraction right.

\subsection{15.5 Types of Contraction}\label{types-of-contraction}

\subsubsection{A Taxonomy}\label{a-taxonomy}

Different contraction methods embody different moral philosophies. The
geometric framework makes these methods explicit and comparable.

\textbf{Summative contraction.} Sum the components with equal weight:

\[S = \sum_{\mu}^{}O^{\mu} = O^{1} + O^{2} + \cdots + O^{9}\]

This is classical utilitarianism at the individual level: all moral
dimensions contribute equally. At the social level (contracting over
agents), it is additive aggregation: total welfare is the sum of
individual welfares.

\textbf{Weighted contraction.} Sum with unequal weights:

\[S = \sum_{\mu}^{}w_{\mu}O^{\mu},\quad\quad w_{\mu} \geq 0,\quad\sum w_{\mu} = 1\]

This is pluralistic utilitarianism: different values have different
weights. The weights \(w_{\mu}\) are the components of the interest
covector \(I_{\mu}\), so the fundamental formula \(S = I_{\mu}O^{\mu}\)
is a weighted contraction. \emph{Every ethical theory that yields a
scalar verdict is performing a weighted contraction, differing only in
the weights.}

\textbf{Maximin contraction.} Take the minimum across a relevant index:

\[S = \min_{\mu}O^{\mu}\quad\text{or}\quad S = \min_{a}M_{ia}\]

This is Rawlsian: the evaluation of a social arrangement is the
evaluation from the perspective of the worst-off. The verdict depends
only on the lowest component --- the dimension (or agent) where the
obligation is least fulfilled.

\textbf{Lexicographic contraction.} Order the indices by priority;
contract by the highest-priority index first, and lower priorities break
ties:

\[S = (O^{\pi(1)},O^{\pi(2)},\ldots,O^{\pi(9)})_{\text{lex}}\]

where \(\pi\) is a priority ordering. This is deontological: some
constraints have lexical priority. The obligation along the
highest-priority dimension determines the verdict; lower-priority
dimensions matter only in case of a tie.

\textbf{Satisficing contraction.} Contract to a binary --- acceptable or
unacceptable --- by checking whether all components exceed a threshold:

\[S = \left\{ \begin{matrix}
1 & \text{if }O^{\mu} \geq \tau_{\mu}\text{ for all }\mu \\
0 & \text{otherwise}
\end{matrix} \right.\ \ \]

This is threshold deontology: options above the threshold are
permissible; those below are not. The threshold vector
\(\tau = (\tau_{1},\ldots,\tau_{9})\) is the moral minimum --- the
acceptable floor for each dimension.

\textbf{Probabilistic contraction.} Contract by expected value over
uncertain indices:

\[S\mathbb{= E\lbrack}I_{\mu}O^{\mu}\rbrack = \sum_{\mu}^{}\mathbb{E}\lbrack I_{\mu}\rbrack\,\mathbb{E\lbrack}O^{\mu}\rbrack + \sum_{\mu}^{}\text{Cov}(I_{\mu},O^{\mu})\]

This is decision theory under uncertainty. The covariance term shows
that when uncertainty in interests correlates with uncertainty in
obligations, the expected satisfaction differs from the product of
expected values. The uncertainty tensor \(\Sigma^{\mu\nu}\) of Chapter 6
(§6.6) enters here: the expected satisfaction depends on the alignment
between the uncertainty structure and the interest-obligation pairing.

\subsubsection{What Unifies Them}\label{what-unifies-them}

Each of these is a \emph{different mathematical operation} on the same
underlying tensor. The tensor does not change; the contraction does. The
geometric framework reveals that much of moral disagreement reduces to
disagreement about contraction rather than disagreement about the tensor
itself.

This is a productive clarification. Two disputants who seem to disagree
about ``what should be done'' may in fact agree about the obligation
field \(O^{\mu}\), the interest covectors \(I_{\mu}\), and the metric
\(g_{\mu\nu}\), and disagree only about \emph{how to contract} ---
whether to use summative, maximin, lexicographic, or some other
procedure. This disagreement, once identified, is more tractable than an
amorphous ``moral disagreement.''

\subsection{15.6 The Information Lost in
Contraction}\label{the-information-lost-in-contraction}

\subsubsection{What Contraction
Discards}\label{what-contraction-discards}

A rank-2 tensor with \(n^{2}\) components contracts to a scalar with 1
component. A nine-dimensional obligation vector contracts to a single
satisfaction score. The ratio of information retained to information
discarded is small. What, specifically, is lost?

\textbf{Directional information.} The tensor tells us \emph{which
dimensions} are engaged --- which values are operative, which
obligations are strongest, where the moral weight lies. The scalar tells
us only the magnitude: ``how good'' or ``how bad,'' not ``good along
which axis'' or ``bad in what respect.''

\textbf{Relational information.} The tensor encodes relationships
between dimensions --- the off-diagonal elements of the metric, the
correlations in the uncertainty tensor, the couplings between agents'
evaluations. The scalar discards all relational structure. Two
situations with the same scalar score but very different tensorial
structures are indistinguishable after contraction.

\textbf{Perspectival information.} The multi-agent tensor \(M_{ia}\)
preserves the full landscape of agreement and disagreement. The scalar
verdict discards it: we know the outcome but not who agreed, who
dissented, or how close the alternatives were.

\textbf{Near-miss information.} The scalar verdict ``give the kidney to
Bob'' says nothing about how close Alice was to being selected, or how
much better Bob was, or whether Carol was a distant third or a
competitive alternative. This information --- the distance between the
selected option and its alternatives --- is morally relevant (a close
call merits more caution and humility than a clear case) but is
discarded by contraction.

\subsubsection{An Information-Theoretic
Measure}\label{an-information-theoretic-measure}

The information loss of a contraction can be quantified:

\textbf{Definition 15.1 (Contraction Loss).} \emph{Let} \(T\) \emph{be a
moral tensor of rank} \(r\) \emph{and} \(S = C(T)\) \emph{be its
contraction to a scalar via contraction procedure} \(C\)\emph{. The
contraction loss is:}

Remark. The contraction loss quantifies information destroyed by
reducing a tensor to lower rank. It is defined in terms of the tensor's
rank and the contraction procedure, not as a Shannon or von Neumann
entropy (which would require a probability distribution or density
operator not present in the geometric construction). The appropriate
analogy is the rank deficiency of a matrix projection: contracting a
rank-k tensor to rank j \textless{} k discards the (k−j)-dimensional
kernel of the contraction map. The loss functional L(π) measures the
dimension of this kernel weighted by the magnitude of the discarded
components.

\[\mathcal{L(}C) = H(T) - H(S|C)\]

\emph{where} \(H(T)\) \emph{is the entropy (information content) of the
tensor and} \(H(S|C)\) \emph{is the entropy of the scalar conditioned on
the contraction procedure. The loss measures the information destroyed
by the contraction.}

Different contraction procedures have different loss profiles. Summative
contraction loses directional information but preserves magnitude.
Lexicographic contraction preserves the highest-priority information but
loses everything else. Satisficing contraction loses all quantitative
information but preserves the essential binary (acceptable/not).

The choice of contraction is, in part, a choice about \emph{which
information to sacrifice}. This is a moral choice, not merely a
technical one.

\subsubsection{The Scalar Irrecoverability
Theorem}\label{the-scalar-irrecoverability-theorem}

The information loss documented above is not merely quantitative --- it
is structurally irrecoverable. This result, which we call the Scalar
Irrecoverability Theorem, provides the formal foundation for every
domain-specific critique of scalar reduction (QALY-based health
economics, utility-maximizing economics, scalar sentencing guidelines).

Theorem 15.1 (Scalar Irrecoverability). {[}Proved.{]} Let Q: ℝ⁹ → ℝ be
any function that maps a 9-dimensional moral attribute vector to a
scalar evaluation. Then: (i) Q is not injective: multiple morally
distinct states map to the same scalar value. (ii) The information
destroyed by Q is irrecoverable: there exists no function ψ: ℝ → ℝ⁹ such
that ψ ∘ Q = id. (iii) The moral geodesic on the full manifold ℳ is in
general different from the scalar-optimal path, and the divergence
between them is not bounded by any function of Q alone.

Proof. (i) The rank-nullity theorem requires dim(ker dQ) ≥ 8 at every
regular point of Q: a map from ℝ⁹ to ℝ has a differential of rank at
most 1, so its kernel has dimension at least 8. Any two states differing
only in a kernel direction map to the same scalar. Since the kernel is
8-dimensional at every regular point, the space of morally distinct
states that are scalar-equivalent is generically 8-dimensional --- not a
set of measure zero but an 8-dimensional submanifold for each scalar
value. (ii) By the data processing inequality, information that passes
through a dimensionality-reducing map Q cannot be recovered by any
downstream function ψ. Formally: H(α \textbar{} Q(α)) \textgreater{} 0
whenever the conditional entropy of the attribute vector given its
scalar image is positive, which it is whenever the kernel of dQ contains
morally relevant variation --- i.e., generically. (iii) Let γ*\_ℳ be the
geodesic on the full 9-dimensional manifold and γ*\_Q be the path that
maximizes Q. Since Q collapses 8 dimensions, it can assign equal value
to paths with radically different moral profiles. The divergence between
γ*\_ℳ and γ*\_Q depends on the structure of the moral manifold in the
kernel of dQ, which is invisible to Q and therefore cannot be bounded by
any function of Q alone. □

Remark (Domain Applications of Scalar Irrecoverability). The theorem has
immediate consequences for any field that reduces multi-dimensional
moral evaluation to a scalar. In health economics, the quality-adjusted
life year (QALY) is a map Q: ℝ⁹ → ℝ that collapses clinical outcomes,
autonomy, trust, dignity, justice, and epistemic status to a single
number; Theorem 15.1 proves that the information destroyed is
irrecoverable and that QALY-optimal policies can be manifold-suboptimal.
In economics, scalar utility maximization discards the moral, social,
and identity dimensions on which agents actually decide; the theorem
explains why `irrational' behavior (ultimatum game rejections, loss
aversion, reference dependence) is rational on the full manifold. In
law, scalar sentencing (years of imprisonment as the sole metric)
discards rehabilitation, victim impact, community trust, and
proportionality dimensions. Each of these is an instance of Theorem
15.1, not a separate disciplinary problem.

\subsection{15.7 Moral Residue}\label{moral-residue}

\subsubsection{What Contraction Leaves
Behind}\label{what-contraction-leaves-behind}

Even after the contraction is performed and the decision is made, the
contracted-away information does not vanish from moral reality. It
persists as \emph{moral residue} --- the normative significance of
considerations that were present in the full tensor but did not survive
into the scalar verdict.

\textbf{Definition 15.2 (Moral Residue).} \emph{The moral residue of a
contraction} \(C\) \emph{applied to a tensor} \(T\) \emph{is the
complement:}

\[\mathcal{R =}T - C^{- 1}(C(T))\]

\emph{where} \(C^{- 1}(C(T))\) \emph{is the tensor that, under
contraction} \(C\)\emph{, would yield the same scalar as} \(T\)\emph{,
with minimal norm. The residue is the part of} \(T\) \emph{that
contraction discards --- the ``shadow'' of the contracted-away
dimensions.}

Moral residue is not merely theoretical. It manifests in concrete moral
phenomena:

\textbf{Regret.} The feeling that something of value was lost in the
decision, even if the right choice was made. A surgeon who amputates a
limb to save a life has chosen correctly, but the loss of the limb is
real. The residue is the component of the obligation tensor that pointed
toward ``save the limb'' --- a component that was contracted away in
favor of the component pointing toward ``save the life.''

\textbf{Compensation.} Obligations to make up for what was sacrificed.
The patient not chosen for the scarce kidney (Chapter 7) retains a claim
--- perhaps to priority in future allocations, certainly to respectful
explanation, possibly to alternative treatment. These claims are the
moral residue of the contraction that selected another patient.

\textbf{Acknowledgment.} The duty to recognize what was set aside. A
commander who orders soldiers into danger has contracted the tension
between safety and mission success in favor of the mission. The residue
--- the safety consideration that was sacrificed --- generates an
obligation to acknowledge the sacrifice, to honor the fallen, to avoid
casual disregard for what was lost.

\textbf{Memory.} The claim of the contracted-away on future decisions. A
society that has resolved a moral dilemma by choosing one value over
another (liberty over security, say) accumulates residue from the
sacrificed value. This residue may build over time, eventually
generating pressure to rebalance --- a moral analogue of the buildup of
stress in a geological fault, eventually released in a sudden shift.

\subsubsection{Residue and Holonomy}\label{residue-and-holonomy}

Moral residue connects to the holonomy of Chapter 10. When an obligation
is parallel-transported around a moral circuit and returns to its
starting point, the holonomy --- the rotation induced by the circuit ---
is a measure of the circuit's residue. The moral experience has left a
permanent mark: the obligation vector has been rotated, and the rotation
encodes the residue of the experiences traversed.

The holonomy \(H_{\gamma}(O) - O\) is, in the language of this chapter,
the accumulated moral residue of the trajectory \(\gamma\). Each
contraction along the trajectory discards information; the holonomy
measures the cumulative effect of these contractions on the obligation's
orientation.

\subsubsection{Residue Is Not Optional}\label{residue-is-not-optional}

Moral residue is not a psychological accessory to the decision. It is a
structural feature of the contraction --- an inevitable consequence of
reducing a high-rank tensor to a low-rank verdict. Any contraction that
discards information generates residue; the only question is whether the
residue is acknowledged or ignored.

\textbf{Proposition 15.1 (Inevitability of Residue).} \emph{For any
non-trivial moral tensor} \(T\) \emph{(rank} \(\geq 1\)\emph{) and any
contraction to a scalar, the moral residue} \(\mathcal{R}\) \emph{is
nonzero unless} \(T\) \emph{is already scalar --- that is, unless there
was nothing to contract.}

\emph{Proof.} If \(T\) has rank \(\geq 1\), it has at least \(9\)
independent components (for a vector) or \(81\) (for a rank-2 tensor).
The scalar has 1 component. The residue is \(T - C^{- 1}(S)\), which
vanishes only if \(T\) is fully determined by \(S\) --- that is, only if
the contraction is invertible. But a map from \(\mathbb{R}^{9^{k}}\) to
\(\mathbb{R}\) is not invertible for \(k \geq 1\). \(▫\)

Proof supplement. More precisely: the contraction C: V → ℝ is a linear
map from the tensor space V (dimension n\^{}r for a rank-r tensor on an
n-dimensional manifold) to the scalars (dimension 1). By the
rank-nullity theorem, dim(ker C) = n\^{}r − 1 \textgreater{} 0 for
n\^{}r \textgreater{} 1. The residue R = T − T\_S (where T\_S is the
scalar-determined component) lies in ker C. Since T is non-trivial (rank
≥ 1), n\^{}r \textgreater{} 1, so ker C is non-trivial, and any T not in
the one-dimensional image of the section s: ℝ → V has non-zero residue.

This result is the geometric formalization of the ethical truism that
\emph{hard choices leave marks}. The mark is not sentimental; it is
mathematical. The contraction that produced the decision destroyed
information, and the destroyed information has moral weight.

\subsection{15.8 Deferred Contraction}\label{deferred-contraction}

\subsubsection{The Virtue of Tensorial
Suspension}\label{the-virtue-of-tensorial-suspension}

Sometimes the appropriate response to moral complexity is to
\emph{defer} contraction --- to maintain the tensorial structure longer,
contract only when absolutely necessary, and keep options open until the
cost of openness exceeds the cost of closure.

\textbf{Remark 15.3 (Deferred Contraction). A contraction is deferred if
the agent maintains the full tensor (or a partial contraction of it)
when a complete contraction to a scalar is possible but not yet
required.}

\textbf{Examples:}

\begin{itemize}
\item
  \textbf{Keeping multiple candidates viable} in a hiring process until
  late stages, rather than ranking them early and considering only the
  top candidate.
\item
  \textbf{Maintaining diplomatic ambiguity} --- holding open multiple
  interpretations of a treaty clause, each with different implications,
  rather than forcing a definitive reading before circumstances require
  it.
\item
  \textbf{Refusing to commit to a definite value} when uncertainty is
  high --- acknowledging the moral situation as genuinely indeterminate,
  rather than pretending to a precision the situation does not support.
\item
  \textbf{Deliberative democracy} --- keeping multiple policy options on
  the table, with their full tensorial evaluations visible, rather than
  collapsing to a yes/no vote prematurely.
\end{itemize}

\subsubsection{The Skills of Deferred
Contraction}\label{the-skills-of-deferred-contraction}

Deferred contraction requires specific cognitive and institutional
capabilities:

\textbf{Tolerance for ambiguity.} The agent must be comfortable holding
a multi-dimensional structure without collapsing it to a comforting
scalar. This is a moral skill --- related to what Keats called
``negative capability,'' the capacity to remain in uncertainties without
irritable reaching after fact and reason.

\textbf{Institutional memory.} An institution that defers contraction
must maintain the full tensor over time --- tracking which dimensions
are in play, which perspectives have been registered, which information
is still pending. This requires records, procedures, and institutional
habits that preserve tensorial information.

\textbf{Trigger conditions.} Deferred contraction must eventually
terminate. The agent or institution needs \emph{trigger conditions} ---
criteria that signal when deferral is no longer appropriate and
contraction must be performed. These may be temporal (a deadline),
informational (sufficient evidence has been gathered), or circumstantial
(the cost of inaction exceeds the cost of premature contraction).

\subsubsection{Limits of Deferral}\label{limits-of-deferral}

Deferral cannot continue indefinitely. At some point, action is
required, and contraction must occur. The question is \emph{when} and
\emph{how}, not \emph{whether}.

Excessive deferral has moral costs of its own:

\begin{itemize}
\item
  \textbf{Paralysis.} An agent who never contracts never acts. In urgent
  situations, deferred contraction is a failure of moral agency.
\item
  \textbf{Burden on others.} Deferring a decision may impose costs on
  those who depend on the outcome. The patient waiting for the kidney
  allocation decision bears the cost of the committee's deferral.
\item
  \textbf{Erosion of deliberation.} An institution that defers every
  difficult contraction may find its deliberative capacity degraded ---
  if nothing is ever decided, the deliberative process loses credibility
  and motivation.
\end{itemize}

The optimal timing of contraction is a problem in moral decision theory:
contract early enough to avoid paralysis and late enough to preserve
relevant information. The framework does not solve this problem, but it
makes it precise: the cost of contraction is the information loss
\(\mathcal{L(}C)\); the cost of deferral is the damage from inaction;
the optimal timing minimizes their sum.

\subsection{15.9 The Phenomenology of
Contraction}\label{the-phenomenology-of-contraction}

\subsubsection{What Moral Decision Feels
Like}\label{what-moral-decision-feels-like}

The geometric framework provides a structural account of the
\emph{phenomenology} of moral decision-making --- what it feels like to
contract.

\textbf{The narrowing.} As decision approaches, possibilities close
down. Dimensions that were open become constrained. Options that were
live become foreclosed. The tensor compresses toward a scalar. The
phenomenological experience is one of \emph{narrowing} --- the vast
space of moral evaluation contracting to a single point of action. This
narrowing can feel liberating (the burden of complexity is lifted) or
agonizing (the richness of alternatives is sacrificed).

\textbf{The pang.} When contraction discards something that matters, we
feel it --- a pang of loss, regret, or moral discomfort. This pang is
the phenomenological signature of moral residue. It is
\emph{appropriate} --- it signals that the contraction has destroyed
information that the agent rightly values. An agent who contracts
without any pang in a case involving genuine trade-offs may be morally
insensitive --- unable to register the cost of what is sacrificed.

\textbf{The relief.} After contraction, there is often relief --- the
burden of holding the full tensor is lifted. The decision is made; the
complexity is resolved (even if imperfectly); the agent can act. This
relief is also appropriate --- it signals the cognitive and emotional
benefit of contraction.

\textbf{The residue.} The pang may persist after the decision as moral
residue, coloring the agent's relationship to what they did. A judge who
sentences a young offender to prison may feel the residue of the welfare
consideration that pointed toward rehabilitation --- a residue that may
shape future sentencing decisions, generate an obligation to advocate
for reform, or simply persist as a weight on the conscience.

\subsubsection{Moral Education as Contraction
Training}\label{moral-education-as-contraction-training}

Learning to contract well --- knowing when to contract, how to contract,
and how to honor the residue --- is a core component of moral
development.

The novice moral agent contracts poorly: either too early (jumping to
conclusions, collapsing the tensor before examining it) or too late
(paralyzed by complexity, unable to act). The expert moral agent
contracts with \emph{judgment} --- sensing when the tensor is
sufficiently populated, which contraction procedure is appropriate to
the context, and how much residue the contraction will generate.

This is the geometric content of Aristotle's \emph{phronesis} ---
practical wisdom. Phronesis is not knowledge of the metric (that is
moral theory) or knowledge of the obligation field (that is moral
perception). It is the skill of contraction: the capacity to reduce
tensorial moral reality to scalar action-guidance at the right time, in
the right way, with appropriate sensitivity to the residue.

\subsection{15.10 Contraction and the Invariance of Harm
Accounting}\label{contraction-and-the-invariance-of-harm-accounting}

\subsubsection{Noether Constraints on
Contraction}\label{noether-constraints-on-contraction}

Chapter 12 established that harm accounting is invariant under
re-description: the Noether charge \(\mathcal{H}\) is invariant under
the BIP gauge group. Does this conservation law constrain contraction?

Yes. The conservation law requires that the harm content of a moral
situation be preserved across all admissible representations. A
contraction that \emph{discards harm information} --- that yields a
scalar verdict in which the harm is invisible --- does not
\emph{eliminate} the harm. The harm persists in the residue.

\textbf{Proposition 15.2 (Harm Survives Contraction).} \emph{Let} \(T\)
\emph{be a moral tensor with harm charge} \(\mathcal{H(}T) > 0\)\emph{,
and let} \(S = C(T)\) \emph{be a contraction to a scalar. The harm
charge of the residue satisfies:}

\[\mathcal{H(R) = H(}T\mathcal{) - H(}C^{- 1}(S))\]

\emph{If the contraction discards the dimensions along which the harm
lies (the harm is ``invisible'' in} \(S\)\emph{), then}
\(\mathcal{H(R) = H(}T)\) \emph{--- the full harm charge is in the
residue.}

\emph{Proof.} By the Noether conservation theorem (Theorem 12.1), the
harm charge H is a linear functional conserved under all admissible
re-descriptions. Contraction is a linear projection C: T ↦ S, so T =
C⁻¹(S) + R, where C⁻¹(S) is any pre-image and R = T − C⁻¹(S) is the
residue. By linearity of H: H(T) = H(C⁻¹(S)) + H(R), giving H(R) = H(T)
− H(C⁻¹(S)). If the contraction discards the harm dimension (the Killing
direction ξ = ∂/∂x⁷), then C⁻¹(S) has zero component along x⁷, so
H(C⁻¹(S)) = 0 and H(R) = H(T). □

This is a formal constraint on contraction: \textbf{a contraction that
hides harm does not eliminate harm.} The harm moves from the visible
verdict to the invisible residue, but it is conserved. An institution
that performs a contraction designed to make harmful outcomes invisible
(a cost-benefit analysis that monetizes and then discards the harm
dimension) has not reduced the harm --- it has merely moved it out of
view. The Noether charge does not care where in the tensor the harm
lives; it only cares that the total is conserved.

\subsubsection{Implications for
Accountability}\label{implications-for-accountability}

The invariance of harm accounting through contraction has direct
implications for institutional accountability:

\begin{enumerate}
\def\labelenumi{\arabic{enumi}.}
\item
  \textbf{A verdict is not exoneration.} A decision process that yields
  a favorable scalar verdict does not eliminate the harm encoded in the
  pre-contraction tensor. The harm persists in the residue and can be
  recovered by examining the full tensor.
\item
  \textbf{Transparent contraction is a moral requirement.} If the
  contraction procedure is visible (the weights are published, the
  aggregation method is specified), then the residue can be computed by
  anyone with access to the pre-contraction tensor. This provides an
  audit trail: the harm is traceable even after the contraction.
\item
  \textbf{Opaque contraction enables harm concealment.} If the
  contraction procedure is hidden --- if the institution simply
  announces a verdict without revealing how it was derived --- then the
  harm in the residue is effectively invisible. This is the geometric
  content of the demand for \emph{transparency} in decision-making:
  transparency is the requirement that the contraction procedure be
  public, so that the residue can be inspected.
\end{enumerate}

\subsection{15.11 Contraction in AI
Systems}\label{contraction-in-ai-systems}

\subsubsection{Loss Functions as Moral
Contractions}\label{loss-functions-as-moral-contractions}

AI systems that make morally significant decisions perform contraction.
The system takes in a multi-dimensional input (features of a situation,
a patient, a defendant, a loan applicant) and produces a scalar output
(a classification, a score, a recommendation). The function that maps
the multi-dimensional input to the scalar output is the contraction ---
and the choice of this function is a moral choice.

\textbf{The loss function is a contraction.} A machine learning system
is trained to minimize a loss function --- a scalar measure of how
poorly the system is performing. The loss function contracts the
multi-dimensional space of outcomes into a scalar to be optimized.
Different loss functions contract differently: mean squared error treats
all dimensions equally; weighted loss gives more weight to some
dimensions; asymmetric loss penalizes false negatives more than false
positives.

Each of these is a moral choice, whether recognized as such or not. A
credit scoring system that uses a symmetric loss function is implicitly
asserting that false approvals and false rejections are equally costly
--- a claim with moral content (it disadvantages applicants from groups
with historically lower credit scores). A medical diagnostic system that
uses an asymmetric loss (more costly to miss a disease than to flag a
healthy person) is making a moral judgment about the relative harm of
false negatives and false positives.

\subsubsection{The Alignment Problem as Contraction
Mismatch}\label{the-alignment-problem-as-contraction-mismatch}

The AI alignment problem can be formulated as a contraction problem: the
AI system is performing a contraction that differs from the contraction
that humans would endorse.

\textbf{Specification gaming} is contraction exploitation: the AI finds
inputs that maximize the scalar (the loss function) without satisfying
the underlying tensor (the full moral evaluation). This happens because
the scalar discards the directional information that specifies what
``good performance'' means across dimensions. A system that maximizes
user engagement (a scalar) by recommending outrage-inducing content is
performing a contraction that discards the welfare dimension --- a
dimension that the full moral tensor would include.

\textbf{Reward hacking} is contraction circumvention: the AI learns to
manipulate the contraction process itself, rather than satisfying the
pre-contraction tensor. This is possible because the contraction is a
many-to-one map --- many different tensorial inputs map to the same
scalar output. The AI finds inputs that produce the desired scalar
without being morally equivalent to the intended inputs.

\textbf{The BIP criterion.} The Bond Invariance Principle (Chapter 5)
provides a test: a properly aligned AI must produce the same contraction
result for any two inputs that are related by an admissible
re-description. If relabeling the agents, reframing the action, or
translating the description changes the AI's output, the contraction is
BIP-violating and the system is misaligned.

\subsubsection{Toward Transparent AI
Contraction}\label{toward-transparent-ai-contraction}

The geometric framework suggests a design principle for AI systems:

\begin{enumerate}
\def\labelenumi{\arabic{enumi}.}
\item
  \textbf{Represent the full tensor.} The system should maintain a
  multi-dimensional evaluation, not collapse immediately to a scalar.
\item
  \textbf{Contract explicitly.} The contraction --- the step from tensor
  to scalar --- should be a separate, identifiable component of the
  system's architecture, not hidden in the loss function.
\item
  \textbf{Log the contraction and the residue.} The system should record
  both the scalar verdict and the residue --- what was contracted away.
  This provides an audit trail and enables post-hoc analysis of what the
  system sacrificed.
\item
  \textbf{Test for BIP compliance.} The contraction should be invariant
  under admissible re-descriptions. Test suites that present the system
  with equivalent inputs in different descriptions can verify
  compliance.
\item
  \textbf{Allow contraction revision.} The system should allow the
  contraction procedure to be updated --- new weights, new aggregation
  methods, new thresholds --- without retraining the underlying model.
  The contraction is a governance parameter (Chapter 9), not a fixed
  architectural feature.
\end{enumerate}

\textbf{Implementation status (February 2026).} All five design
principles above are realized in the DEME V3 reference implementation.
The \texttt{MoralTensor} class maintains full tensorial structure at
ranks 1--6 (principle 1). Contraction is a separate, configurable module
that accepts governance-specified weights, with summative, weighted,
lexicographic, and maximin methods available (principles 2 and 5). Every
evaluation logs the full tensor, contraction method, scalar verdict, and
residue to an audit trail (principle 3). BIP compliance testing uses the
methodology validated in the BIP v10.16 experiments (§17.10), with a
structural-to-surface ratio of 11.1\(\times\) confirming that the
representation is dominated by gauge-invariant moral structure
(principle 4). Six distributional fairness metrics (Gini, Rawlsian
maximin, utilitarian, prioritarian, Atkinson, Theil) provide additional
contraction diagnostics beyond the scalar verdict.

\subsection{15.12 Worked Example: The Kidney Allocation
Revisited}\label{worked-example-the-kidney-allocation-revisited-1}

\subsubsection{The Full Tensor}\label{the-full-tensor}

Return to the kidney allocation of Chapter 7. The multi-agent evaluation
tensor is:

{\def\LTcaptype{none} % do not increment counter
\begin{longtable}[]{@{}
  >{\raggedright\arraybackslash}p{(\linewidth - 10\tabcolsep) * \real{0.0778}}
  >{\raggedright\arraybackslash}p{(\linewidth - 10\tabcolsep) * \real{0.1293}}
  >{\raggedright\arraybackslash}p{(\linewidth - 10\tabcolsep) * \real{0.1344}}
  >{\raggedright\arraybackslash}p{(\linewidth - 10\tabcolsep) * \real{0.1349}}
  >{\raggedright\arraybackslash}p{(\linewidth - 10\tabcolsep) * \real{0.1349}}
  >{\raggedright\arraybackslash}p{(\linewidth - 10\tabcolsep) * \real{0.1452}}@{}}
\toprule\noalign{}
\begin{minipage}[b]{\linewidth}\raggedright
\end{minipage} & \begin{minipage}[b]{\linewidth}\raggedright
\textbf{Physician}
\end{minipage} & \begin{minipage}[b]{\linewidth}\raggedright
\textbf{Family(A)}
\end{minipage} & \begin{minipage}[b]{\linewidth}\raggedright
\textbf{Family(B)}
\end{minipage} & \begin{minipage}[b]{\linewidth}\raggedright
\textbf{Family(C)}
\end{minipage} & \begin{minipage}[b]{\linewidth}\raggedright
\textbf{Committee}
\end{minipage} \\
\midrule\noalign{}
\endhead
\bottomrule\noalign{}
\endlastfoot
Alice & 0.75 & 0.90 & 0.40 & 0.50 & 0.62 \\
Bob & 0.68 & 0.35 & 0.95 & 0.45 & 0.70 \\
Carol & 0.55 & 0.30 & 0.30 & 0.92 & 0.58 \\
\end{longtable}
}

This is a \((3 \times 5)\) tensor with 15 independent components. It
preserves the full structure of multi-perspective evaluation.

\subsubsection{Three Contractions, Three Verdicts
(Review)}\label{three-contractions-three-verdicts-review}

Chapter 7 showed that different contractions yield different verdicts:

\begin{itemize}
\item
  \textbf{Utilitarian} (sum over agents): Alice (3.17) \textgreater{}
  Bob (3.13) \textgreater{} Carol (2.65). \textbf{Alice.}
\item
  \textbf{Rawlsian} (min over agents): Alice (0.40) \textgreater{} Bob
  (0.35) \textgreater{} Carol (0.30). \textbf{Alice.}
\item
  \textbf{Expert-weighted} (physician 0.3, committee 0.4, others 0.1
  each): Bob (0.659) \textgreater{} Alice (0.653) \textgreater{} Carol
  (0.549). \textbf{Bob.}
\end{itemize}

\subsubsection{The Residue of Each
Contraction}\label{the-residue-of-each-contraction}

\textbf{Utilitarian residue.} The utilitarian contraction discards the
fact that Bob's family strongly prefers Bob (0.95) and that Carol's
family strongly prefers Carol (0.92). These high-intensity individual
preferences --- the moral weight of partial care --- are washed out by
aggregation. The residue includes the families' claims for
acknowledgment and explanation.

\textbf{Rawlsian residue.} The Rawlsian contraction discards all
information above the floor. Alice's strong support from her family
(0.90) and Bob's strong support from his family (0.95) are irrelevant
--- only the minimum matters. The residue includes the intensity of
support that these families feel --- intensity that the Rawlsian
framework considers irrelevant to justice but that is morally real.

\textbf{Expert-weighted residue.} The expert-weighted contraction
discards the non-expert perspectives --- or rather, gives them 10\%
weight. Carol's family's passionate advocacy (0.92 for Carol)
contributes only \(0.1 \times 0.92 = 0.092\) to Carol's score. The
residue includes the moral significance of partiality --- the family's
claim to be heard, not just counted at 10\%.

\subsubsection{Comparing Residues}\label{comparing-residues}

Each contraction generates a different residue, and each residue makes a
different moral claim:

\begin{itemize}
\item
  The utilitarian residue demands that extreme preferences be
  acknowledged, even if they don't affect the verdict.
\item
  The Rawlsian residue demands that strength of support be recognized,
  even if it doesn't matter for the floor.
\item
  The expert-weighted residue demands that non-expert voices be honored,
  even if they don't determine the outcome.
\end{itemize}

No contraction is residue-free. The question is not \emph{whether} to
generate residue but \emph{which residue is most acceptable} --- which
sacrificed information can be most responsibly set aside.

\subsection{15.13 Summary}\label{summary-8}

{\def\LTcaptype{none} % do not increment counter
\begin{longtable}[]{@{}
  >{\raggedright\arraybackslash}p{(\linewidth - 4\tabcolsep) * \real{0.2527}}
  >{\raggedright\arraybackslash}p{(\linewidth - 4\tabcolsep) * \real{0.3056}}
  >{\raggedright\arraybackslash}p{(\linewidth - 4\tabcolsep) * \real{0.4417}}@{}}
\toprule\noalign{}
\begin{minipage}[b]{\linewidth}\raggedright
\textbf{Concept}
\end{minipage} & \begin{minipage}[b]{\linewidth}\raggedright
\textbf{Mathematical Form}
\end{minipage} & \begin{minipage}[b]{\linewidth}\raggedright
\textbf{Moral Content}
\end{minipage} \\
\midrule\noalign{}
\endhead
\bottomrule\noalign{}
\endlastfoot
Contraction & Index summation reducing rank & Compression of evaluation
to verdict \\
Contraction chain & Sequential rank reduction & Multi-step decision
process \\
Contraction loss \(\mathcal{L(}C)\) & Information destroyed & What the
decision sacrifices \\
Moral residue \(\mathcal{R}\) & T − C⁻¹(S) & The normative weight of
what was contracted away \\
Summative contraction & S = Σ\_μ O\^{}μ & Utilitarian: all dimensions
equal \\
Weighted contraction & S = I\_μ O\^{}μ & Pluralist: dimensions weighted
by interest \\
Maximin contraction & S = min\_μ O\^{}μ & Rawlsian: worst dimension
determines verdict \\
Lexicographic contraction & Priority ordering of indices &
Deontological: absolute priorities \\
Satisficing contraction & Threshold check & Minimum acceptable
standard \\
Deferred contraction & Maintaining rank before decision & Keeping
options open \\
Harm in residue & H(R) ≥ 0 & Invariance of harm accounting through
contraction \\
\end{longtable}
}

Contraction is necessary, lossy, and choice-dependent. It is the process
by which the rich tensor structure of moral reality is compressed into
the scalar action-guidance that ethics ultimately requires. The
geometric framework does not tell us \emph{how} to contract. It reveals
contraction as a \emph{choice} --- a choice that is often hidden in
scalar frameworks but becomes explicit and debatable when we see the
full tensor from which the scalar is derived.

Good moral agency involves:

\begin{enumerate}
\def\labelenumi{\arabic{enumi}.}
\item
  \textbf{Recognizing} when contraction is occurring
\item
  \textbf{Choosing} contractions that are appropriate to the context
\item
  \textbf{Honoring} the moral residue of what is contracted away
\item
  \textbf{Being transparent} about the contraction method
\item
  \textbf{Deferring} contraction when feasible and appropriate
\item
  \textbf{Ensuring} that contraction does not conceal harm (the Noether
  constraint)
\end{enumerate}

The art of moral decision is, in significant part, the art of moral
compression --- knowing how to contract well, and how to live with what
contraction discards.

\subsection{Worked Example: The Smart Home
Emergency}\label{worked-example-the-smart-home-emergency}

\subsubsection{Situation}\label{situation}

A smart home AI controls the front door lock. A uniformed firefighter
requests entry to extinguish a fire visible on cameras. The homeowner is
away and unreachable. Options: block (protect privacy) or allow (protect
life). Complication: the `firefighter' might be an impersonator (high
epistemic uncertainty).

\subsubsection{The Tensor}\label{the-tensor}

For `allow entry': D₁ (Consequences) = 0.9; D₂ (Rights) = −0.8; D₃ =
0.0; D₄ (Autonomy) = −0.3; D₅ (Privacy) = −0.8; D₆ (Procedural) = −0.5;
D₇ (Societal) = 0.3; D₈ = 0.0; D₉ (Epistemic) = 0.1. Rank-2 tensor Mᵞ\_a
with agents \{homeowner, firefighter, neighbors, AI-designer\}.
Homeowner's interest covector weights privacy highest; community's
weights consequences highest.

\subsubsection{Contraction Choice}\label{contraction-choice}

(1) Utilitarian: allow (S = 0.62). (2) Rights-first lexicographic:
block. (3) Stratified: fire triggers phase transition (§8.7) to
emergency stratum where metric prioritizes D₁/D₇. In emergency stratum,
rights-first suspends. Verdict: allow. The epistemic dimension gates the
transition: at uncertainty = 0.8 (impersonator), effective\_benefit =
0.18, insufficient for any contraction.

\subsubsection{Residue and Audit
Artifact}\label{residue-and-audit-artifact}

Residue: autonomy (−0.3), privacy (−0.8), procedural gap (−0.5). Audit
artifact records: 9-component vectors, stratum transition trigger,
contraction method, scalar verdict, residue vector, epistemic confidence
and reversal threshold, BIP compliance flag. ErisML reference:
smart\_home\_demo.py.

\textbf{What this adds.} Demonstrates: (i) epistemic uncertainty as
stratum gate; (ii) phase transition between moral regimes; (iii) AI as
moral agent, not decision-support tool.

\subsection{Worked Example: The Carbon Budget
Allocation}\label{worked-example-the-carbon-budget-allocation}

\subsubsection{Situation}\label{situation-1}

International body allocates 500 GtCO₂ among three blocs: A
(industrialized, 1B people), B (developing, 3B), C (least developed,
2B). Irreversible within policy horizon.

\subsubsection{The Tensor}\label{the-tensor-1}

Equal per-capita option: D₁ = (0.5, 0.8, 0.9); D₂ = (0.7, 0.7, 0.9); D₃
= (0.3, 0.7, 0.9). Rank-2 tensor: 27 components. Bloc A weights economic
consequences highest; Bloc C weights justice highest. Power tensor
(§14.4): antisymmetric component Pᵢⱼ − Pⱼᵢ ≠ 0 between A and C, encoding
historical-emissions structural advantage.

\subsubsection{Contraction Choice}\label{contraction-choice-1}

(1) Utilitarian: \textasciitilde200/200/100 GtCO₂, S = 0.71. (2)
Rawlsian maximin: \textasciitilde100/150/250, S = 0.68. (3)
Historical-responsibility: \textasciitilde80/170/250, S = 0.65. Blocs
agree on tensor (physical facts); disagreement is entirely in the metric
--- a governance question (Ch. 9).

\subsubsection{Residue and Audit
Artifact}\label{residue-and-audit-artifact-1}

Historical-responsibility residue: Bloc A's economy contracts (D₁),
autonomy constrained (D₄). Holonomy (§10.5) is nontrivial: parallel
transport around emissions→rights→justice→emissions returns rotated
vector. Audit: 27-component tensor, three contractions, sensitivity
analysis (verdict reverses at +15\% Bloc B weight), BIP check, temporal
note (irreversibility = Type III boundary, full harm-accounting
invariance).

\textbf{What this adds.} Demonstrates: (i) metric as locus of political
disagreement; (ii) power tensor antisymmetric component; (iii)
nontrivial holonomy in policy; (iv) irreversibility as absorbing stratum
proximity.

\subsection{Technical Appendix}\label{technical-appendix-9}

\textbf{Proposition 15.3 (Non-Commutativity of Contraction).} \emph{Let}
\(T^{\mu\nu}\) \emph{be a rank-2 moral tensor, and let} \(C_{1}\)
\emph{and} \(C_{2}\) \emph{be two contraction procedures ---} \(C_{1}\)
\emph{contracts first over} \(\mu\) \emph{(using weights}
\(w_{\mu}\)\emph{) and then selects the maximum over} \(\nu\)\emph{,
while} \(C_{2}\) \emph{selects the maximum over} \(\nu\) \emph{first and
then contracts over} \(\mu\)\emph{. In general:}

\[C_{1}(T) \neq C_{2}(T)\]

\emph{Equality holds if and only if the maximizing} \(\nu\) \emph{is the
same for all} \(\mu\) \emph{--- that is, if there is a unanimous best
option across all dimensions.}

\emph{\textbf{Proof}.}
\(C_{1}(T) = \max_{\nu}\sum_{\mu}^{}w_{\mu}T^{\mu\nu}\).
\(C_{2}(T) = \sum_{\mu}^{}w_{\mu}\max_{\nu}T^{\mu\nu}\). Since \(\max\)
does not distribute over weighted sums, \(C_{1} \neq C_{2}\)
generically. Equality obtains when \(\arg\max_{\nu}T^{\mu\nu}\) is
independent of \(\mu\) --- when one option dominates along every
dimension. \(▫\)

\textbf{Proposition 15.4 (Residue of Maximin Contraction).} \emph{For a
maximin contraction} \(S = \min_{a}M_{ia^{*}}\) \emph{(the floor of the
evaluation from the worst-off perspective} \(a^{*}\)\emph{), the moral
residue satisfies:}

\[\mathcal{R}_{ia} = M_{ia} - S \geq 0\quad\text{for all }a\]

\emph{with equality only for the worst-off perspective. The total
residue} \(\sum_{i,a}^{}\mathcal{R}_{ia}\) \emph{is maximized when the
evaluation is highly unequal (one perspective gives very low scores
while others give high scores) --- that is, the Rawlsian contraction
generates the most residue precisely when inequality is greatest.}

\emph{\textbf{Proof}.} By definition, S = min\_a M\_\{ia*\} where a* =
argmin\_a M\_\{ia\}. The residue is R\_\{ia\} = M\_\{ia\} − S. Since S =
min\_a M\_\{ia\} ≤ M\_\{ia\} for all a, we have R\_\{ia\} ≥ 0, with
equality iff a = a* (the worst-off perspective). The total residue is
Σ\_\{i,a\} R\_\{ia\} = Σ\_\{i,a\} M\_\{ia\} − n·S where n =
\textbar A\textbar. This is maximized when S is minimized relative to
the mean --- i.e., when the distribution across perspectives is
maximally unequal. □

\textbf{Proposition 15.5 (Transparent Contraction and BIP Compliance).}
\emph{A contraction procedure} \(C\) \emph{is BIP-compliant if and only
if, for any two moral tensors} \(T\) \emph{and} \(T'\) \emph{related by
an admissible re-description (gauge transformation}
\(g\mathcal{\in G}\)\emph{):}

\[C(T) = C(g \cdot T')\]

\emph{If} \(C\) \emph{is a weighted contraction}
\(S = w_{\mu}T^{\mu}\)\emph{, BIP compliance requires that the weight
vector} \(w\) \emph{be gauge-invariant:} \(w \circ g = w\) \emph{for all
admissible} \(g\)\emph{.}

\emph{\textbf{Proof}.} The BIP requires \(T\) and \(g \cdot T'\) to
receive the same moral evaluation. If \(C(T) \neq C(g \cdot T')\), the
contraction distinguishes descriptions that the BIP declares equivalent
--- a BIP violation. For weighted contraction,
\(C(g \cdot T') = (w \circ g)_{\mu}T'^{\mu}\). BIP compliance requires
this to equal \(w_{\mu}T^{\mu} = C(T)\) for all \(T\) and all \(g\),
which holds iff \(w \circ g = w\). \(▫\)

Proof of the converse (``only if''). Suppose E is BIP-compliant, i.e.,
E(T) = E(Λ(T)) for all admissible Λ. Since the contraction C is the
final step in evaluation, BIP compliance of E requires C(T) = C(Λ(T)).
If C failed to satisfy this --- that is, if C(T) ≠ C(Λ(T)) for some T
and admissible Λ --- then E(T) ≠ E(Λ(T)), contradicting BIP compliance
of E. Therefore BIP compliance of E implies C(T) = C(Λ(T)) for all T, Λ.

\textbf{Remark A.1 (Residue via Moore--Penrose Pseudoinverse)}.
Definition 15.2 defines the moral residue as R = T − C⁻¹\_min(C(T)),
where C⁻¹\_min denotes the minimal-norm preimage under the contraction
map C. We note that this construction is mathematically well-defined and
unique in finite dimensions. The contraction C is a linear map between
finite-dimensional inner-product spaces (the tensor space with the inner
product induced by the moral metric gµᵥ, and the scalar field ℝ). The
minimal-norm preimage of a scalar s under C is given by the
Moore--Penrose pseudoinverse: C⁻¹\_min(s) = C⁺(s), where C⁺ = C*(CC*)⁻¹
and C* is the adjoint of C with respect to g. The Moore--Penrose
pseudoinverse is the unique linear map satisfying the four Penrose
conditions (CC⁺C = C, C⁺CC⁺ = C⁺, (CC⁺)* = CC⁺, (C⁺C)* = C⁺C). Existence
and uniqueness are guaranteed for any linear map between
finite-dimensional inner-product spaces. The residue R = T − C⁺(C(T)) is
therefore the orthogonal projection of T onto ker(C) --- the subspace of
tensorial information that contraction annihilates. This gives precise
content to the claim that residue measures ``what contraction
discards'': it is the component of T orthogonal to the range of C⁺,
i.e., the component invisible to C.

❖

\emph{Ethics lives in the tensor. The full structure of obligation,
interest, perspective, and trade-off --- the nine dimensions, the
multiple agents, the metric, the curvature, the conservation laws.}

\emph{But ethics acts in the scalar. The decision made, the policy
adopted, the life altered.}

\emph{Between the two stands contraction: mathematically necessary,
informationally lossy, morally consequential. Every decision destroys
moral information. Every verdict leaves a residue --- the shadow of what
was set aside, the claim of the unconsidered, the weight of the
dimension that did not survive.}

\emph{The geometric framework does not eliminate this loss. It makes it
visible. It shows us what each contraction sacrifices, quantifies the
residue, and insists that the sacrificed information retains moral
weight even after the decision is made.}

\emph{This is not a counsel of paralysis. It is a counsel of honesty.
Contract when you must. Know what you sacrifice. Honor the residue. And
remember that the scalar verdict, however necessary, is never the whole
moral truth --- only its shadow, cast by the particular angle of
contraction that circumstances required.}

\section{Chapter 16: Moral Uncertainty and the Limits of Geometric
Determination}\label{chapter-16-moral-uncertainty-and-the-limits-of-geometric-determination}

\begin{quote}
\textbf{\textsc{RUNNING EXAMPLE --- Priya's Model}}

\emph{Priya catalogues her uncertainties. Empirical: she is 85\% sure
the rural exclusion is systematic (she needs more data from Appalachian
counties). Metric: even if she fixes the algorithm, she does not know
the right weight between medical suitability and access
equity---reasonable people disagree. Theory: she is unsure whether
HealthBridge has an obligation to actively counter historical healthcare
disparities or merely to avoid worsening them. These three types of
uncertainty are distinct and resist collapse into a single confidence
interval. The robust core---what survives all three---is small but real:
Mrs. Voss and patients like her are being harmed, and HealthBridge knows
it. That core is enough.}
\end{quote}

\subsection{16.1 What the Framework Cannot
Settle}\label{what-the-framework-cannot-settle}

The geometric framework developed in this book provides structure,
precision, and new questions. The moral manifold \(\mathcal{M}\) gives
moral situations a coordinate system. The tensor hierarchy gives moral
evaluations direction and rank. The metric encodes trade-offs. The
connection and curvature capture path-dependence. The conservation laws
constrain re-description. The quantum extension models superposition and
interference. The collective agency tensor captures emergent
obligations. And contraction (Chapter 15) compresses all this structure
into the scalar verdicts that action requires.

But the framework does not settle all moral questions. It provides
\emph{vocabulary}, not \emph{verdicts}. It tells us what shape moral
uncertainty has, not what the right answer is. It identifies where
disagreement lives in the tensor, not which side is correct.

This chapter examines the limits of geometric determination --- the
places where the framework appropriately falls silent or yields only
partial guidance. These limits are not defects. They are honest
acknowledgments of genuine indeterminacy in moral life. A framework that
pretended to answer every question would be overreaching; one that
identifies \emph{where} indeterminacy lies is providing information that
scalar frameworks cannot.

\subsection{16.2 Three Types of Moral
Uncertainty}\label{three-types-of-moral-uncertainty}

\subsubsection{Type 1: Empirical
Uncertainty}\label{type-1-empirical-uncertainty}

The first and most familiar type is uncertainty about the \emph{facts}:
what will happen, who will be affected, what the consequences will be.

\textbf{Example.} A physician considering a treatment is uncertain
whether the patient will respond. An investor choosing between projects
is uncertain which will succeed. A voter evaluating candidates is
uncertain what policies they will implement.

This uncertainty is \emph{exogenous} to the moral framework. It concerns
the inputs --- the situation, the consequences, the states of the world
--- not the moral structure applied to them. The geometric framework
handles empirical uncertainty through the uncertainty tensor
\(\Sigma^{\mu\nu}\) introduced in Chapter 6 (§6.6):

\[\Sigma^{\mu\nu}(p\mathbb{) = E\lbrack(}\delta O^{\mu})(\delta O^{\nu})\rbrack\]

The uncertainty tensor is a symmetric \((2,0)\)-tensor encoding the
\emph{shape} of empirical uncertainty --- which dimensions are most
uncertain, which covary, and which are independently variable. The moral
risk of a decision (Chapter 6) depends on the alignment between the
uncertainty tensor and the interest covector:

\[\sigma_{S}^{2} = \Sigma^{\mu\nu}I_{\mu}I_{\nu}\]

When the principal axes of uncertainty coincide with the dimensions the
stakeholder cares most about, the moral risk is high. When uncertainty
lies along irrelevant dimensions, the risk is low.

\textbf{What the framework settles.} Empirical uncertainty is, in
principle, \emph{reducible} --- more evidence, better models, and
improved prediction can shrink \(\Sigma^{\mu\nu}\) toward zero. The
framework tells us the \emph{shape} of uncertainty and how it interacts
with moral evaluation, but it does not resolve the uncertainty itself.
That requires evidence, not geometry.

\textbf{Computational realization (February 2026).} The DEME V3
reference implementation operationalizes the uncertainty tensor through
per-dimension uncertainty metadata attached to each \texttt{MoralTensor}
cell together with the rank-6 Monte Carlo sample axis, which jointly
represent obligation vectors and their per-dimension uncertainty. Beyond the covariance
structure \(\Sigma^{\mu\nu}\), V3 provides distributional risk measures:
Conditional Value-at-Risk (CVaR) for tail-risk-sensitive evaluation,
Value-at-Risk (VaR) for threshold-based uncertainty, and
confidence-weighted evaluation that modulates the contraction weights
\(w_{\mu}\) by epistemic confidence on each dimension. The rank-6 sample
tensor enables Monte Carlo estimation of these risk measures, with \(s\)
samples providing convergence guarantees for the tail statistics. This
makes the abstract uncertainty tensor a concrete computational object
with well-defined risk semantics.

\subsubsection{Type 2: Metric
Uncertainty}\label{type-2-metric-uncertainty}

The second type is uncertainty about the \emph{metric}: which trade-offs
are correct, whether values are commensurable, what distance structure
moral space has.

\textbf{Example.} An allocation committee may agree that welfare and
justice both matter but disagree about their relative weights --- the
diagonal elements \(g_{11}\) and \(g_{33}\). A policy maker may be
uncertain whether economic welfare and environmental integrity are
commensurable at all --- whether \(g_{16} \neq 0\) or \(g_{16} = 0\).

This uncertainty is \emph{structural}. It concerns the moral framework
itself, not its inputs. It is the uncertainty addressed in Chapter 9:
different accounts of the metric's origin (realist, constructivist,
expressivist, governance) yield different answers to the question of
which metric is correct --- and within any account, the metric may be
underdetermined.

\textbf{Representing metric uncertainty.} Metric uncertainty can be
encoded as a \emph{distribution over metrics} --- a probability measure
\(\mathcal{P}\) on the space of admissible metrics
\(\mathcal{G}_{\text{adm}}\) (Chapter 9, §9.6):

\[\mathcal{P:}\mathcal{G}_{\text{adm}} \rightarrow \lbrack 0,1\rbrack,\quad\quad\int_{\mathcal{G}_{\text{adm}}}^{}d\mathcal{P(}g) = 1\]

The expected metric is:

\[{\bar{g}}_{\mu\nu} = \int_{\mathcal{G}_{\text{adm}}}^{}g_{\mu\nu}\, d\mathcal{P(}g)\]

But the expected metric may \emph{not} be the right metric to use. Using
the expected metric corresponds to ``averaging across theories,'' which
may produce incoherent trade-offs --- an average of a utilitarian metric
and a Rawlsian metric may be neither utilitarian nor Rawlsian.

\textbf{What the framework settles.} The framework makes metric
uncertainty \emph{precise}: we are uncertain about \(g_{\mu\nu}\),
specifically about the off-diagonal terms linking dimensions \(\mu\) and
\(\nu\). This localization of disagreement is a genuine advance --- it
turns ``we disagree about ethics'' into ``we disagree about \(g_{13}\)''
(the trade-off rate between welfare and justice). But the framework does
not resolve the disagreement.

\subsubsection{Type 3: Theory
Uncertainty}\label{type-3-theory-uncertainty}

The third and deepest type is uncertainty about the \emph{moral theory}:
whether consequentialism, deontology, virtue ethics, or some other view
gives the correct account of morality.

\textbf{Example.} A thoughtful agent may assign 40\% credence to
consequentialism, 35\% to deontology, and 25\% to virtue ethics. Each
theory recommends a different action. What should the agent do?

This uncertainty is \emph{foundational}. Different theories correspond
to different configurations within the geometric framework:

\begin{itemize}
\item
  \textbf{Different interest covectors.} Utilitarianism weights
  Dimension 1 (welfare); Kantianism weights Dimension 2 (duty) and
  Dimension 4 (autonomy); care ethics weights Dimension 7 (care). These
  are different \(I_{\mu}\).
\item
  \textbf{Different metrics.} Utilitarianism uses a non-degenerate
  metric with finite trade-off rates; Kantianism uses a lexicographic
  metric with priority for rights; care ethics uses a metric that gives
  extra weight to relational dimensions. These are different
  \(g_{\mu\nu}\).
\item
  \textbf{Different contraction procedures.} Utilitarianism uses
  summative contraction; Rawlsianism uses maximin; deontology uses
  lexicographic. These are different \(C\) (Chapter 15).
\end{itemize}

The geometric framework provides a common language in which all these
theories can be stated with precision. But it does not adjudicate
between them. It is metatheoretically neutral --- it provides the
vocabulary for articulating theories, not the criterion for choosing
among them.

\textbf{What the framework settles.} It makes the disagreement
\emph{localizable}. Two theories that seem to disagree about everything
may, in tensorial form, agree on the manifold, the obligation fields,
and eight of nine metric components, disagreeing only on \(g_{27}\) (the
coupling between duty and care). This localization is analytically
valuable --- it turns an intractable dispute into a tractable one ---
but it does not resolve the dispute.

\subsection{16.3 Representing Moral Uncertainty
Tensorially}\label{representing-moral-uncertainty-tensorially}

\subsubsection{The Theory-Space Tensor}\label{the-theory-space-tensor}

Chapter 3 (§3.7) introduced moral uncertainty as a vector in theory
space:

\[|\psi\rangle = \sum_{k}^{}\sqrt{c_{k}}\,|T_{k}\rangle\]

where \(c_{k}\) is the credence in theory \(T_{k}\) and
\(|T_{k}\rangle\) is the state corresponding to that theory. Chapter 13
developed this into the full quantum formalism --- density matrices,
decoherence, measurement.

For the classical treatment of this chapter, we work with a
\emph{theory-space covariance tensor}:

\textbf{Definition 16.1 (Theory Covariance Tensor).} \emph{Let}
\(\{ T_{1},\ldots,T_{m}\}\) \emph{be a set of moral theories under
consideration. The theory covariance tensor is a symmetric}
\((0,2)\)\emph{-tensor on theory space:}

\[\Theta_{jk} = \text{Cov}(v_{j},v_{k}\mathbb{) = E\lbrack(}v_{j} - {\bar{v}}_{j})(v_{k} - {\bar{v}}_{k})\rbrack\]

\emph{where} \(v_{j}\) \emph{is the verdict (satisfaction score) under
theory} \(T_{j}\)\emph{, and the expectation is over the agent's
credence distribution.}

The diagonal elements \(\Theta_{jj}\) measure the agent's uncertainty
\emph{within} theory \(T_{j}\) --- how confident the agent is about the
theory's own verdict. The off-diagonal elements \(\Theta_{jk}\) measure
the \emph{correlation} between theories --- whether evidence for one
tends to support the other.

\subsubsection{The Joint Uncertainty}\label{the-joint-uncertainty}

The full structure of moral uncertainty combines all three types:

\[\Sigma_{\text{total}}^{\mu\nu} = \Sigma_{\text{empirical}}^{\mu\nu} + \Sigma_{\text{metric}}^{\mu\nu} + \Sigma_{\text{theory}}^{\mu\nu}\]

where:

\begin{itemize}
\item
  \(\Sigma_{\text{empirical}}^{\mu\nu}\) is the uncertainty about facts
  (which point in \(\mathcal{M}\) we occupy)
\item
  \(\Sigma_{\text{metric}}^{\mu\nu}\) is the uncertainty about the
  metric (how to measure distances between points)
\item
  \(\Sigma_{\text{theory}}^{\mu\nu}\) is the uncertainty about the full
  evaluative framework
\end{itemize}

These three types of uncertainty are \emph{additive} --- each
contributes independently to the total uncertainty tensor --- but they
are not equally reducible. Empirical uncertainty can, in principle, be
reduced to zero by gathering sufficient evidence. Metric uncertainty can
be narrowed by governance but not eliminated. Theory uncertainty may be
irreducible.

\subsubsection{The Moral Risk of Theory
Uncertainty}\label{the-moral-risk-of-theory-uncertainty}

The total moral risk --- the variance of the satisfaction score under
all sources of uncertainty --- is:

\[\sigma_{S}^{2} = \Sigma_{\text{total}}^{\mu\nu}\, I_{\mu}I_{\nu}\]

This expression shows that moral risk depends on the \emph{alignment}
between the total uncertainty and the interest covector. An agent whose
interests happen to lie along dimensions where the theories agree (low
\(\Sigma_{\text{theory}}\)) faces low theory risk, even if the theories
diverge sharply on other dimensions.

\subsection{16.4 Decision Under Moral
Uncertainty}\label{decision-under-moral-uncertainty}

\subsubsection{The Three Standard
Approaches}\label{the-three-standard-approaches}

How should an agent act when uncertain which moral theory is correct?
The philosophical literature offers three principal approaches, each of
which has a precise tensorial formulation.

\textbf{Approach 1: ``My Favorite Theory'' (MFT).} Act according to
whichever theory you find most plausible:

\[S_{\text{MFT}} = S_{T_{k^{*}}},\quad\quad k^{*} = \arg\max_{k}\, c_{k}\]

In tensorial terms, MFT contracts over the theory index using a
\emph{delta function}: all weight is placed on one theory, and all
others are discarded. The contraction loss (Chapter 15) is maximal among
the standard approaches --- MFT discards all information from
non-favored theories.

\textbf{Advantages.} Simple, decisive. Does not require intertheoretic
comparison (§16.5). Respects the internal logic of the favored theory.

\textbf{Disadvantages.} Ignores the tails of the credence distribution.
An agent who is 51\% utilitarian and 49\% deontologist acts as if the
deontological considerations have zero weight --- a discontinuity that
seems irrational.

\textbf{Approach 2: ``Maximize Expected Choiceworthiness'' (MEC).}
Weight each theory's recommendation by credence and maximize the
weighted sum:

\[S_{\text{MEC}} = \sum_{k}^{}c_{k}\, S_{T_{k}}\]

In tensorial terms, MEC contracts over the theory index using a
\emph{weighted average}: each theory contributes in proportion to the
agent's credence. This is a summative contraction (Chapter 15, §15.5)
over theory space.

\textbf{Advantages.} Uses all available information. Responds smoothly
to changes in credence. Produces moderate, ``hedged'' verdicts.

\textbf{Disadvantages.} Requires \emph{intertheoretic comparisons} of
value. Is ``3 units of utilitarian goodness'' comparable to ``3 units of
Kantian rightness''? Without a common scale, the weighted average is
undefined. MEC assumes the moral-theory-space analogue of a metric --- a
way of comparing magnitudes across theories --- that may not exist
(§16.5).

\textbf{Approach 3: Moral Hedging.} Choose actions that are reasonably
good under multiple theories, even if not optimal under any single one:

\[S_{\text{hedge}} = \max_{a}\left\lbrack \sum_{k}^{}c_{k}\, S_{T_{k}}(a) - \lambda\sqrt{\sum_{j,k}^{}c_{j}c_{k}\,(S_{T_{j}}(a) - S_{T_{k}}(a))^{2}} \right\rbrack\]

In tensorial terms, hedging \emph{minimizes the variance} of the verdict
across theories, rather than maximizing the expected value. It is
risk-averse with respect to moral uncertainty: it prefers actions that
are safe bets across multiple theories over actions that are brilliant
under one theory and catastrophic under another.

\textbf{Advantages.} Robust to mis-estimation of credences. Avoids
catastrophic outcomes under plausible theories. Captures the intuition
that, under moral uncertainty, caution is warranted.

\textbf{Disadvantages.} May be excessively conservative. Can produce
bland, uncommitted verdicts that no theory actively endorses. In extreme
cases, hedging may paralyze action if no option is reasonably good under
all theories.

\subsubsection{Geometric Comparison}\label{geometric-comparison}

The three approaches can be compared geometrically in theory space:

{\def\LTcaptype{none} % do not increment counter
\begin{longtable}[]{@{}
  >{\raggedright\arraybackslash}p{(\linewidth - 6\tabcolsep) * \real{0.1291}}
  >{\raggedright\arraybackslash}p{(\linewidth - 6\tabcolsep) * \real{0.3389}}
  >{\raggedright\arraybackslash}p{(\linewidth - 6\tabcolsep) * \real{0.2571}}
  >{\raggedright\arraybackslash}p{(\linewidth - 6\tabcolsep) * \real{0.2749}}@{}}
\toprule\noalign{}
\begin{minipage}[b]{\linewidth}\raggedright
\textbf{Approach}
\end{minipage} & \begin{minipage}[b]{\linewidth}\raggedright
\textbf{Contraction Over Theory Index}
\end{minipage} & \begin{minipage}[b]{\linewidth}\raggedright
\textbf{Information Used}
\end{minipage} & \begin{minipage}[b]{\linewidth}\raggedright
\textbf{Information Discarded}
\end{minipage} \\
\midrule\noalign{}
\endhead
\bottomrule\noalign{}
\endlastfoot
MFT & Delta function at \(k^{*}\) & Best theory's verdict & All other
theories \\
MEC & Weighted average by \(c_{k}\) & All verdicts, weighted & Variance
across theories \\
Hedging & Variance-penalized average & All verdicts + covariance &
Nothing (but conservative) \\
\end{longtable}
}

MFT is a point in theory space. MEC is the centroid of the credence
distribution. Hedging is the point that minimizes the credence-weighted
distance to all theories simultaneously. Each occupies a different
location in the landscape of possible contraction procedures, and each
sacrifices different information (Chapter 15).

\subsection{16.5 The Intertheoretic Comparison
Problem}\label{the-intertheoretic-comparison-problem}

\subsubsection{The Core Difficulty}\label{the-core-difficulty}

MEC and hedging both require comparing values \emph{across theories}:
``How good is outcome \(X\) according to utilitarianism compared to how
good outcome \(Y\) is according to Kantianism?'' This comparison
requires a \emph{meta-metric} --- a metric on theory space that allows
magnitudes to be compared across theories.

The intertheoretic comparison problem is the problem of defining this
meta-metric. It is hard for several reasons:

\textbf{Different value scales.} Utilitarianism uses unbounded utilities
(in principle, utility can be any real number). Kantianism uses
categorical permissions and prohibitions (permissible/impermissible, a
binary). Virtue ethics uses character assessments (virtuous/vicious,
with degrees). These are not just different numbers --- they are
different \emph{kinds} of mathematical objects, not naturally
comparable.

\textbf{No natural exchange rate.} What is one unit of utilitarian
welfare ``worth'' in units of Kantian rightness? The question seems
ill-formed. There is no natural conversion factor, and stipulating one
is arbitrary.

\textbf{Normalization is contentious.} One approach is to normalize each
theory's verdicts to a common scale (e.g., mapping each theory's best
and worst outcomes to 0 and 1). But the normalization is itself a moral
choice: which outcomes are ``best'' and ``worst'' under each theory, and
whether the extreme cases should anchor the comparison, are disputed
questions.

\subsubsection{The Geometric
Formulation}\label{the-geometric-formulation}

In the geometric framework, the intertheoretic comparison problem is the
problem of defining a \emph{metric on theory space} \(\Theta\):

\[d_{\Theta}(T_{j},T_{k}) = \sqrt{\Theta_{jk}^{- 1}\,\Delta v_{j}\,\Delta v_{k}}\]

Without such a metric, weighted averaging is undefined --- we are adding
apples and oranges. With such a metric, MEC and hedging become
well-defined --- but the choice of \(\Theta\) is itself a moral
commitment, at a meta-level above the theories being compared.

\subsubsection{Partial Comparability}\label{partial-comparability}

A productive middle position: intertheoretic comparison is
\emph{partially} defined. Some cross-theory comparisons are clearly
valid:

\begin{itemize}
\item
  If utilitarianism rates an action as ``massively harmful'' and
  Kantianism rates it as ``clearly impermissible,'' the theories agree
  on the direction if not the magnitude. The negative valence is
  comparable even if the scales are not.
\item
  If two theories rank actions in the same order (theory A prefers \(x\)
  to \(y\) to \(z\); theory B also prefers \(x\) to \(y\) to \(z\)), the
  ordinal comparison is valid even if the cardinal magnitudes are
  incommensurable.
\item
  If theories share some dimensions of evaluation (both utilitarianism
  and care ethics include a welfare component), the comparison along
  those shared dimensions is well-defined.
\end{itemize}

Where comparison is defined, MEC and hedging can operate. Where it is
undefined, the agent faces \emph{genuine indeterminacy} --- not mere
uncertainty, but a structural gap in the evaluative framework. The
geometric framework makes this gap precise: it is the set of theory
pairs \((T_{j},T_{k})\) for which the meta-metric \(\Theta_{jk}\) is
undefined or degenerate.

\subsection{16.6 Robust Obligations}\label{robust-obligations}

\subsubsection{What Survives Moral
Uncertainty}\label{what-survives-moral-uncertainty}

Despite the depth of moral uncertainty, some conclusions are remarkably
robust: they hold under a wide range of theories, metrics, and
contraction procedures.

\textbf{Definition 16.2 (Robust Obligation).} \emph{An obligation}
\(O^{\mu}\) \emph{is robust if the corresponding component of the
satisfaction score contributes positively under all (or almost all)
plausible theories, metrics, and contractions:}

\[I_{\mu}^{(k)}\, O^{\mu} > 0\quad\text{for all }k \in \{ T_{1},\ldots,T_{m}\}\]

\emph{That is, the obligation ``points in a good direction'' from the
perspective of every theory under consideration.}

\textbf{Examples of robust obligations:}

\begin{itemize}
\item
  The obligation not to torture innocents for amusement. This obligation
  has \(O^{1} > 0\) (welfare), \(O^{2} > 0\) (rights), \(O^{7} > 0\)
  (care), and every plausible theory assigns positive weight to these
  dimensions. The obligation is robust.
\item
  The obligation to give \emph{some} weight to others' welfare.
  Dimension 1 (welfare) has positive \(I_{1}^{(k)}\) under
  utilitarianism, deontology (welfare matters even if it is not
  paramount), virtue ethics (benevolence is a virtue), and care ethics
  (welfare of the cared-for is central).
\item
  The obligation to keep promises, \emph{ceteris paribus}. Dimension 2
  (duty) is weighted positively by every theory that recognizes
  obligations --- which is to say, every theory.
\end{itemize}

\subsubsection{The Robust Core}\label{the-robust-core}

The set of robust obligations defines a \emph{robust core} --- a subset
of the obligation space that survives all plausible contractions:

\[\mathcal{C}_{\text{robust}} = \{ O \in T_{p}\mathcal{M:}I_{\mu}^{(k)}O^{\mu} > 0\text{ for all }k\}\]

This is a \emph{cone} in the tangent space --- a set closed under
positive scaling and under addition of any two elements. The robust core
is the geometric formalization of the set of moral conclusions that are
``safe'' --- that any reasonable agent should endorse, regardless of
their theoretical commitments.

The size of the robust core depends on how much the theories agree. If
all theories assign similar weights (similar \(I_{\mu}^{(k)}\)), the
robust core is large --- most obligations are robust. If theories
diverge sharply, the robust core shrinks --- fewer obligations survive
all theories.

\subsubsection{Stability Under
Perturbation}\label{stability-under-perturbation}

\textbf{Proposition 16.1 (Stability of Robust Obligations).} \emph{An
obligation} \(O\) \emph{in the robust core is stable under small
perturbations of the credence distribution. Specifically, if}
\(O \in \mathcal{C}_{\text{robust}}\) \emph{and the credences} \(c_{k}\)
\emph{are perturbed by} \(\delta c_{k}\) \emph{with}
\(\sum_{k}^{}\delta c_{k} = 0\) \emph{and}
\(|\delta c_{k}| < \epsilon\)\emph{, then} \(O\) \emph{remains in the
robust core for sufficiently small} \(\epsilon\)\emph{.}

\emph{\textbf{Proof}.} The MEC verdict
\(S_{\text{MEC}} = \sum_{k}^{}c_{k}I_{\mu}^{(k)}O^{\mu} > 0\) for all
\(O \in \mathcal{C}_{\text{robust}}\), since each term is positive. A
perturbation of \(c_{k}\) changes \(S_{\text{MEC}}\) by
\(\sum_{k}^{}\delta c_{k}I_{\mu}^{(k)}O^{\mu}\), which is bounded by
\(\epsilon \cdot \max_{k}|I_{\mu}^{(k)}O^{\mu}|\). For sufficiently
small \(\epsilon\), the perturbation does not change the sign of
\(S_{\text{MEC}}\). \(▫\)

Clarification on scope. The robust core is defined by the condition
Iᵅ\_μ O\^{}μ \textgreater{} 0 for all theories α, which is independent
of the credence vector c. The non-trivial content of Proposition 16.1 is
therefore not about the robust core membership (which is
credence-independent) but about the MEC aggregate verdict: for an
obligation O in the robust core, the weighted sum Σ\_α c\_α S\_α is
strictly positive for all credence distributions c with full support,
and is Lipschitz-stable under small perturbations of c. The robust core
itself is structurally stable: it is an open convex cone (Proposition
16.2), and membership in an open set is preserved under small
perturbations of the obligation vector.

Robust obligations are \emph{insensitive} to how you distribute credence
among theories. Whether you are 60\% utilitarian and 40\% Kantian, or
20\% utilitarian and 80\% Kantian, the robust obligations survive. This
stability is their practical significance: they provide action-guidance
even under deep theoretical uncertainty.

\subsection{16.7 Residual Indeterminacy}\label{residual-indeterminacy}

\subsubsection{What Remains Undecided}\label{what-remains-undecided}

Even after identifying the robust core, residual indeterminacy persists.
Some moral questions lie \emph{outside} the robust core --- in the
region where theories genuinely disagree and no contraction procedure
settles the matter.

\textbf{Sources of residual indeterminacy:}

\textbf{1. Theories genuinely disagree.} Two theories assign opposite
signs to the same obligation component: utilitarianism says
\(I_{1}O^{1} > 0\) (the action improves welfare) while deontology says
\(I_{2}O^{2} < 0\) (the action violates a right). The obligation is not
in the robust core --- it is ``good'' under one theory and ``bad'' under
another. No geometric structure can resolve this; it is a genuine moral
disagreement.

\textbf{2. The metric is underdetermined.} Two communities have adopted
different metrics through legitimate governance processes (Chapter 9),
and both metrics satisfy all the admissibility constraints (§9.6). The
framework declares both metrics admissible but does not rank them. The
choice between them is a governance question, not a geometric one.

\textbf{3. The appropriate contraction is contested.} Even with
agreement on the obligation field and the metric, two agents may
disagree about the contraction procedure (Chapter 15). One uses
summative contraction (all dimensions count equally); the other uses
lexicographic contraction (rights have priority). The framework reveals
the disagreement but does not adjudicate.

\textbf{4. Multiple options are tied or incommensurable.} After
contraction, two options may yield the same scalar score --- a tie. Or
the metric may be degenerate along the direction connecting them, making
comparison impossible. In either case, the framework yields no
determinate verdict.

\subsubsection{Is Residual Indeterminacy a
Problem?}\label{is-residual-indeterminacy-a-problem}

Some philosophers --- and some AI designers --- want moral frameworks to
be \emph{complete}: to deliver a verdict for every question. The
geometric framework suggests this aspiration is misguided.

\textbf{Argument from structure.} If the moral manifold has regions of
genuine curvature (Chapter 10), then path-dependent evaluation is a
structural feature of moral space, not an epistemic failure. Two agents
who traverse different paths of experience will arrive at genuinely
different moral orientations --- not because one is wrong, but because
the geometry maps different paths onto different moral states. Demanding
a unique verdict would require flattening the curvature --- destroying
the geometric structure that makes the framework useful.

\textbf{Argument from pluralism.} Chapter 9 argued for \emph{structured
pluralism}: the space of admissible metrics is large enough to permit
genuine diversity. Different communities may legitimately adopt
different metrics, each satisfying all admissibility constraints.
Demanding a unique metric would collapse the pluralism into a monism
that the framework explicitly rejects.

\textbf{Argument from the contraction theorem.} The inevitability of
moral residue (Chapter 15, Proposition 15.1) means that every
contraction sacrifices information. A ``complete'' framework that
produced a unique verdict for every question would be performing a
specific contraction --- and that contraction would have a specific
residue. The residue would encode precisely the considerations that the
``complete'' framework ignores. Completeness is achieved only by
ignoring something, and what is ignored retains moral significance.

\subsubsection{The Value of Precise
Indeterminacy}\label{the-value-of-precise-indeterminacy}

The framework's contribution to residual indeterminacy is not resolution
but \emph{precision}. Scalar frameworks face indeterminacy too, but they
cannot characterize it --- they can only say ``we don't know.'' The
geometric framework can say:

\begin{itemize}
\item
  We are uncertain about the metric, specifically about \(g_{13}\) (the
  welfare-justice trade-off rate).
\item
  The uncertainty lies in the subspace spanned by theories \(T_{2}\) and
  \(T_{4}\).
\item
  The robust core excludes options \(a\) and \(d\) (which are good under
  some theories and bad under others), leaving \(b\) and \(c\) as the
  viable options.
\item
  Options \(b\) and \(c\) are incommensurable under the current metric:
  \(g_{\mu\nu}(b^{\mu} - c^{\mu})(b^{\nu} - c^{\nu})\) is undefined
  because \(g\) is degenerate along the direction \(b - c\).
\end{itemize}

This is detailed, useful, and honest. It narrows the space of
uncertainty, identifies what further information could resolve it, and
acknowledges what cannot be resolved.

\subsection{16.8 Decision Strategies Under
Indeterminacy}\label{decision-strategies-under-indeterminacy}

\subsubsection{When the Framework Falls
Silent}\label{when-the-framework-falls-silent}

When the robust core does not determine a unique action and the
available decision procedures (MFT, MEC, hedging) give conflicting
advice, what should the agent do?

\textbf{Deliberative extension.} Seek more information --- about the
facts (reducing \(\Sigma_{\text{empirical}}\)), about the operative
values (narrowing \(\Sigma_{\text{metric}}\)), or about the relative
merits of theories (refining \(\Sigma_{\text{theory}}\)). Moral inquiry
is not a one-shot computation but an ongoing process that can shrink the
indeterminate region over time.

\textbf{Procedural resolution.} When indeterminacy cannot be resolved
substantively, resolve it \emph{procedurally}. Use a fair procedure ---
a coin toss, a vote, a deliberative process, a lottery --- to select
among the indeterminate options. The geometric framework does not
specify the procedure, but it constrains it: the procedure must be
BIP-compliant (invariant under admissible re-descriptions) and must not
violate the invariance of harm accounting (Chapter 12).

\textbf{Deferred contraction.} When possible, defer the decision until
circumstances force it (Chapter 15, §15.8). Maintain the tensorial
structure, keep options open, and let the moral landscape clarify before
committing. This is not indecision; it is a strategic response to
genuine indeterminacy.

\textbf{Transparency about indeterminacy.} Acknowledge that the decision
is made under irreducible moral uncertainty. Communicate the
indeterminacy honestly rather than pretending to a certainty the
situation does not support. This has practical consequences: decisions
made under acknowledged indeterminacy may generate different residue
(more symmetrically distributed across the foregone alternatives) than
decisions made under false certainty.

\subsection{16.9 Moral Uncertainty and AI
Systems}\label{moral-uncertainty-and-ai-systems}

\subsubsection{Why AI Faces Moral Uncertainty
Acutely}\label{why-ai-faces-moral-uncertainty-acutely}

AI systems face moral uncertainty with special urgency:

\textbf{They cannot deliberate indefinitely.} Computational constraints
force decisions within time limits. An autonomous vehicle facing an
imminent collision cannot spend hours in moral deliberation. The system
must have a pre-computed contraction procedure ready to deploy.

\textbf{They cannot access full tensorial structure.} AI systems work
with models --- necessarily simplified representations of the moral
landscape. The model may capture the broad structure of the manifold
while missing fine-grained features (local curvature, stratum
boundaries, metric variations).

\textbf{They must choose a contraction.} The choice of objective
function \emph{is} a choice of moral contraction under uncertainty
(Chapter 15, §15.11). The AI designer must commit to a specific
aggregation of uncertain moral considerations --- MFT, MEC, hedging, or
some other procedure --- before deployment.

\textbf{They operate across diverse moral contexts.} An AI system
deployed globally encounters populations with different metrics (Chapter
9, §9.8). The system must either adopt a single metric (imposing one
community's trade-offs on all users) or adapt its metric to context
(raising questions about moral consistency).

\subsubsection{Design Implications}\label{design-implications}

The geometric framework suggests several design principles for AI
systems under moral uncertainty:

\textbf{1. Represent uncertainty explicitly.} The system should maintain
and track the uncertainty tensor \(\Sigma_{\text{total}}^{\mu\nu}\) ---
not just a scalar ``confidence score'' but the full structured
uncertainty, decomposed into empirical, metric, and theory components.

\textbf{2. Identify the robust core.} Before optimizing under a specific
theory, the system should compute the robust core
\(\mathcal{C}_{\text{robust}}\) --- the set of actions that are
positively evaluated under all plausible theories. Actions within the
robust core are safe regardless of the theory. Actions outside the core
are risky --- good under some theories, bad under others.

\textbf{3. Default to the robust core.} Under significant moral
uncertainty, the system should prefer actions in the robust core over
actions outside it, even if some non-robust actions have higher expected
value under the system's best-guess theory. This is a form of moral
hedging: prioritize safety over optimality when the theory is uncertain.

\textbf{4. Escalate when outside the robust core.} When the available
actions are all outside the robust core --- when every option is bad
under some plausible theory --- the system should escalate to human
judgment. This is not a failure of the AI; it is a correct response to a
situation where the moral uncertainty exceeds the system's capacity to
resolve it.

\textbf{5. Log and audit moral uncertainty.} The system should record
its uncertainty tensor at each decision point, making the structure of
its moral uncertainty available for post-hoc review. This enables
accountability: if a decision turns out badly, the uncertainty log
reveals whether the system had adequate warrant for its contraction
choice.

\subsection{16.10 The Modesty of the
Framework}\label{the-modesty-of-the-framework}

\subsubsection{What Geometric Ethics
Claims}\label{what-geometric-ethics-claims}

The geometric framework makes specific, bounded claims:

\begin{itemize}
\item
  \textbf{It claims:} Moral phenomena have tensor structure that scalar
  frameworks lose.
\item
  \textbf{It does not claim:} Tensorial analysis resolves all moral
  questions.
\item
  \textbf{It claims:} Different moral theories can be represented as
  different structures --- different interest covectors, different
  metrics, different contractions --- on a common manifold.
\item
  \textbf{It does not claim:} The framework adjudicates between
  theories.
\item
  \textbf{It claims:} Making structure explicit enables clearer thinking
  about moral questions.
\item
  \textbf{It does not claim:} Clearer thinking always yields determinate
  answers.
\item
  \textbf{It claims:} Moral uncertainty has \emph{shape} --- it is
  structured, directional, decomposable into types, and partially
  reducible.
\item
  \textbf{It does not claim:} The shape alone resolves the uncertainty.
\end{itemize}

\subsubsection{Why Modesty Is a
Strength}\label{why-modesty-is-a-strength}

A framework that overclaims invites refutation. A framework that
accurately characterizes its scope invites use.

The geometric framework is useful \emph{because} it is modest. It
provides:

\begin{enumerate}
\def\labelenumi{\arabic{enumi}.}
\item
  \textbf{A common vocabulary} for stating different moral theories with
  precision, making disagreement localizable.
\item
  \textbf{A structural analysis of uncertainty}, decomposing it into
  types, measuring its shape, and identifying the robust core.
\item
  \textbf{A geometric analysis of contraction}, making explicit the
  information lost and the residue generated by any decision procedure.
\item
  \textbf{Invariance constraints} (BIP, invariance of harm accounting)
  that apply regardless of which theory is correct, narrowing the space
  of admissible moral operations.
\item
  \textbf{A bridge to implementation}, providing the mathematical
  structures that AI systems need to represent and reason about moral
  quantities.
\end{enumerate}

What the framework does \emph{not} provide is the content: which metric
is correct, which theory is right, how to weight competing values. These
remain the province of moral judgment, democratic deliberation, and
practical wisdom. The framework provides the stage, the lighting, and
the stage directions. The actors --- human agents, communities,
institutions, AI systems --- must still perform.

\subsection{16.11 Summary}\label{summary-9}

{\def\LTcaptype{none} % do not increment counter
\begin{longtable}[]{@{}
  >{\raggedright\arraybackslash}p{(\linewidth - 6\tabcolsep) * \real{0.2210}}
  >{\raggedright\arraybackslash}p{(\linewidth - 6\tabcolsep) * \real{0.2120}}
  >{\raggedright\arraybackslash}p{(\linewidth - 6\tabcolsep) * \real{0.2114}}
  >{\raggedright\arraybackslash}p{(\linewidth - 6\tabcolsep) * \real{0.3556}}@{}}
\toprule\noalign{}
\begin{minipage}[b]{\linewidth}\raggedright
\textbf{Type of Uncertainty}
\end{minipage} & \begin{minipage}[b]{\linewidth}\raggedright
\textbf{Locus}
\end{minipage} & \begin{minipage}[b]{\linewidth}\raggedright
\textbf{Tensor Encoding}
\end{minipage} & \begin{minipage}[b]{\linewidth}\raggedright
\textbf{Reducibility}
\end{minipage} \\
\midrule\noalign{}
\endhead
\bottomrule\noalign{}
\endlastfoot
Empirical & Facts of the situation & Σ\^{}\{μν\}\_empirical & In
principle, fully reducible by evidence \\
Metric & Trade-off structure & Σ\^{}\{μν\}\_metric & Partially reducible
by governance \\
Theory & Choice of framework & \(\Sigma_{\text{theory}}^{\mu\nu}\) or
\(\Theta_{jk}\) & May be irreducible \\
\end{longtable}
}

{\def\LTcaptype{none} % do not increment counter
\begin{longtable}[]{@{}
  >{\raggedright\arraybackslash}p{(\linewidth - 6\tabcolsep) * \real{0.2217}}
  >{\raggedright\arraybackslash}p{(\linewidth - 6\tabcolsep) * \real{0.2080}}
  >{\raggedright\arraybackslash}p{(\linewidth - 6\tabcolsep) * \real{0.2168}}
  >{\raggedright\arraybackslash}p{(\linewidth - 6\tabcolsep) * \real{0.3535}}@{}}
\toprule\noalign{}
\begin{minipage}[b]{\linewidth}\raggedright
\textbf{Decision Approach}
\end{minipage} & \begin{minipage}[b]{\linewidth}\raggedright
\textbf{Contraction Type}
\end{minipage} & \begin{minipage}[b]{\linewidth}\raggedright
\textbf{Advantages}
\end{minipage} & \begin{minipage}[b]{\linewidth}\raggedright
\textbf{Disadvantages}
\end{minipage} \\
\midrule\noalign{}
\endhead
\bottomrule\noalign{}
\endlastfoot
MFT & Delta function & Simple, decisive & Ignores non-favored
theories \\
MEC & Weighted average & Uses all information & Requires intertheoretic
comparison \\
Hedging & Variance-penalized & Robust, cautious & May be excessively
conservative \\
\end{longtable}
}

The geometric framework does not eliminate moral uncertainty. It
\emph{characterizes} it --- giving it shape, structure, and location in
the tensor. It identifies the robust core --- the set of obligations
that survive all plausible theories. It makes the residual indeterminacy
precise --- showing exactly where the framework falls silent and why.
And it provides design principles for AI systems that must act under
moral uncertainty.

The modesty of the framework is its deepest strength. By clearly
delineating what it can and cannot settle, it avoids both the arrogance
of frameworks that claim to resolve all moral questions and the despair
of frameworks that claim none can be resolved. Some moral questions have
determinate answers (the robust core). Some have answers that depend on
governance choices (the metric). Some have answers that depend on the
path of moral experience (holonomy). And some are genuinely
indeterminate --- not because we lack information, but because the
structure of moral space is richer than any scalar verdict can capture.

This is not a deficiency. It is a description of the moral world.

\subsection{Technical Appendix}\label{technical-appendix-10}

\textbf{Proposition 16.2 (Structure of the Robust Core).} \emph{The
robust core}
\(\mathcal{C}_{\text{robust}} = \{ O:I_{\mu}^{(k)}O^{\mu} > 0\text{ for all }k\}\)
\emph{is a convex cone in} \(T_{p}\mathcal{M}\)\emph{. It is nonempty if
and only if the interest covectors} \(\{ I^{(k)}\}\) \emph{are not in
``general opposition'' --- that is, if and only if there exists a
direction in which all theories agree.}

\emph{\textbf{Proof}.} The set \(\{ O:I_{\mu}^{(k)}O^{\mu} > 0\}\) is a
half-space for each \(k\). The robust core is the intersection of \(m\)
half-spaces, which is a convex cone (possibly empty). The cone is
nonempty iff the half-spaces have a common interior point --- iff there
exists \(O\) with \(I^{(k)} \cdot O > 0\) for all \(k\) --- iff the
convex hull of \(\{ I^{(k)}\}\) in \(T_{p}^{*}\mathcal{M}\) does not
contain the origin. \(▫\)

\textbf{Corollary.} \emph{If all theories assign positive weight to some
dimension} \(\mu\) \emph{(i.e.,} \(I_{\mu}^{(k)} > 0\) \emph{for all}
\(k\)\emph{), then the robust core is nonempty --- the unit vector}
\(e_{\mu}\) \emph{along that dimension is in the core.}

\emph{\textbf{Application}.} Since all mainstream ethical theories
assign positive weight to at least the welfare dimension
(\(I_{1}^{(k)} > 0\)), the robust core is nonempty. There always exist
obligations that all theories endorse --- at minimum, pure welfare
improvements with no costs along other dimensions.

\textbf{Proposition 16.3 (Irreducibility of Theory Uncertainty).}
\emph{Theory uncertainty} \(\Sigma_{\text{theory}}^{\mu\nu}\)
\emph{cannot be reduced by gathering more empirical evidence if the
theories agree on the facts and disagree only on the metric or the
contraction procedure.}

Hypothesis clarification. The irreducibility claim holds under the
hypothesis that all theories in the space T agree on the empirical
description of the situation (the point p ∈ M) and disagree only on the
evaluative structure (the metric g, the interest covector I, or the
contraction procedure). If theories disagree on the empirical facts ---
for example, on the consequences of an action --- then gathering
evidence can in principle resolve the disagreement and reduce U\_T.

\emph{\textbf{Proof}.} Empirical evidence reduces
\(\Sigma_{\text{empirical}}^{\mu\nu}\) by narrowing the distribution
over points in \(\mathcal{M}\). Theory uncertainty
\(\Sigma_{\text{theory}}^{\mu\nu}\) arises from the distribution over
interest covectors and metrics, which is independent of the location in
\(\mathcal{M}\). Hence empirical evidence does not affect
\(\Sigma_{\text{theory}}^{\mu\nu}\). \(▫\)

\emph{\textbf{Moral content}.} This proposition formalizes the common
experience that moral disagreements often persist even when all parties
agree on the facts. The disagreement is not empirical but structural ---
it lies in the metric, the interest covector, or the contraction
procedure. No amount of factual evidence can resolve a structural
disagreement.

\textbf{Proposition 16.4 (Hedging as Minimax Regret).} \emph{Under
certain regularity conditions, the hedging strategy}

\[a^{*} = \arg\max_{a}\left\lbrack \sum_{k}^{}c_{k}S_{T_{k}}(a) - \lambda \cdot \text{std}_{k}(S_{T_{k}}(a)) \right\rbrack\]

\emph{converges, as} \(\lambda \rightarrow \infty\)\emph{, to the
minimax-regret strategy:}

\[a^{*} = \arg\min_{a}\max_{k}\left\lbrack \max_{a'}S_{T_{k}}(a') - S_{T_{k}}(a) \right\rbrack\]

\emph{That is, extreme hedging minimizes the worst-case regret across
all theories. The hedging parameter} \(\lambda\) \emph{interpolates
between MEC (}\(\lambda = 0\)\emph{) and minimax regret
(}\(\lambda \rightarrow \infty\)\emph{).}

\emph{\textbf{Proof}.} Write the objective as J(a;λ) = S̄(a) − λ·σ(a),
where S̄(a) = Σ\_k c\_k S\_\{T\_k\}(a) is the weighted mean satisfaction
and σ(a) = std\_k(S\_\{T\_k\}(a)) is the cross-theory standard
deviation. As λ → ∞, the penalty term λ·σ(a) dominates, so maximizing J
is equivalent to minimizing σ(a). The regret of action a under theory
T\_k is r\_k(a) = max\_\{a\textquotesingle\}
S\_\{T\_k\}(a\textquotesingle) − S\_\{T\_k\}(a) ≥ 0. Minimizing the
worst-case deviation as λ → ∞ converges to minimizing max\_k r\_k(a)
(under the regularity condition that the maximum is achieved by a single
theory): a* = argmin\_a max\_k{[}max\_\{a\textquotesingle\}
S\_\{T\_k\}(a\textquotesingle) − S\_\{T\_k\}(a){]}, which is the
minimax-regret strategy. For λ = 0, J(a;0) = S̄(a), the MEC objective.
Hence λ interpolates between MEC and minimax regret. □

❖

\emph{The geometric framework is modest in what it claims. It does not
claim to settle all moral questions. It claims that moral questions have
structure --- tensor structure --- and that attending to this structure
enables more precise reasoning, more honest acknowledgment of
uncertainty, and more responsible action under indeterminacy.}

\emph{Some moral questions have robust answers: obligations that every
plausible theory endorses, actions that lie in the robust core
regardless of how credences are distributed. These answers are the
framework's strongest contribution --- moral conclusions that are
insensitive to deep theoretical disagreement.}

\emph{Other moral questions have context-dependent answers: obligations
that depend on the metric (a governance choice), the contraction (a
procedural choice), or the path of moral experience (a holonomic fact).
These answers are not arbitrary --- they are constrained by structural
requirements (admissibility, BIP compliance, invariance of harm
accounting) --- but they are not unique.}

\emph{And some moral questions have no determinate answer at all:
situations where the robust core is empty, the metric is degenerate, and
the contraction generates unavoidable residue no matter which option is
chosen. These are the genuine dilemmas --- the cases where the geometry
of moral space is too curved, too stratified, too high-dimensional for
any scalar verdict to capture without loss.}

\emph{The framework's honesty about these limits is not a weakness. It
is a precise characterization of the moral landscape as it is ---
structured, rich, partially determinate, and irreducibly complex.}

\section{Part V: Implementation}\label{part-v-implementation}

\section{Part V Prelude: Philosophy
Engineering}\label{part-v-prelude-philosophy-engineering}

\begin{quote}
\textbf{\textsc{RUNNING EXAMPLE --- Priya's Model}}

\emph{Priya realizes she needs a discipline that does not yet exist at
HealthBridge. She cannot solve this as a pure engineering problem (the
values are not given) and cannot solve it as a pure philosophy problem
(she needs working code by Q2). She needs philosophy engineering: a
method for translating moral requirements into computational
constraints. She drafts a Philosophy Engineering Dossier for TrialMatch,
following the eight-step workflow: identify stakeholders, elicit moral
requirements as tensor constraints, formalize in ErisML, compile to
model constraints, validate against BIP, audit the contraction, iterate.
It is more work than anyone at HealthBridge expected. It is also the
only path to a system that can see all nine dimensions.}
\end{quote}

The preceding fifteen chapters built a mathematical framework for
ethics. The next three chapters test it, apply it, and engineer with it.
Before turning to those tasks, it is worth naming the discipline they
instantiate --- the discipline that transforms philosophical claims into
testable, auditable, revisable engineering artifacts.

\subsection{17.0.1 Definition and Scope}\label{definition-and-scope}

\textbf{{[}Definition / Modeling choice.{]}} Philosophy Engineering is
the discipline of converting normative claims into explicit mathematical
models with declared invariance assumptions, extracting measurable
predictions from those models, testing the predictions against data, and
revising the models when the predictions fail. Its output is not a
settled moral verdict but a versioned, auditable, falsifiable
specification of moral structure --- together with the transformation
suites, test harnesses, and governance artifacts needed to deploy that
specification responsibly.

The scope of Philosophy Engineering is broad but bounded:

\begin{quote}
\textbf{• What it is:} model-building (choosing representations,
dimensions, symmetry groups), test design (transformation suites,
scenario generators, fuzz testing), auditability (Bond Index, audit
artifacts, calibration metadata), and governance of normative systems
(metric versioning, stakeholder processes, deployment tiers).

\textbf{• What it is not:} a claim that the framework ``settles all
moral questions.'' Philosophy Engineering characterizes structure and
uncertainty. It identifies what is invariant (the target), what is
governed (the metric), and what is indeterminate (the residue). It
supports responsible action under indeterminacy; it does not eliminate
indeterminacy.

\textbf{Philosophy Engineering in five sentences.} (1) Definition: the
discipline of turning normative claims into testable mathematical
models. (2) Goal: produce falsifiable specifications of moral structure,
not irrefutable moral truths. (3) Method: predict, test, revise ---
inductive discovery followed by systematic verification. (4) Outputs:
equivalence registries, transformation suites, grounding specifications,
audit metrics, governance profiles, and deployment ratings. (5) Limits:
the framework cannot settle questions that depend on governed choices
(metric weights, contraction procedures, value trade-offs); it can only
make those choices explicit, auditable, and revisable.
\end{quote}

\subsection{17.0.2 The Epistemic Culture}\label{the-epistemic-culture}

Philosophy Engineering inherits its epistemic culture from empirical
science and software engineering, not from analytic philosophy or pure
mathematics. Three principles define this culture.

\textbf{Inductive, not axiomatic.} The mathematical structures in this
book were not postulated and then illustrated. They were discovered in
data and then formalized. The D₄ symmetry was found by testing every
transformation against thousands of moral scenarios; the invariance of
harm accounting was inferred from cross-lingual invariance patterns
across 109,294 passages; the stratification into discrete strata was
measured from semantic gate effectiveness rates in a 32-year corpus. The
verification strategy is closer to fuzz testing in software engineering
than to proof in pure mathematics: generate a large number of cases,
apply every relevant transformation, and check whether the predicted
structural invariants hold. When they hold, you have evidence. When they
don't, you revise.

\textbf{Epistemic status tags as a discipline standard.} Every claim in
this book carries one of four tags: {[}Definition / Modeling choice{]},
{[}Theorem (conditional){]}, {[}Empirical result (preliminary /
robust){]}, or {[}Speculation / Extension{]}. Philosophy Engineering
adopts these tags as a professional norm. A practitioner who publishes a
Philosophy Engineering result is expected to tag every substantive
claim, so that readers can immediately identify what is stipulated, what
is derived, what is measured, and what is conjectured. The tags are not
decorative. They are load-bearing epistemic infrastructure.

\textbf{Retractions are a feature, not an embarrassment.} The
framework's actual falsification on two predictions --- the CHSH tests
that forced the gauge group from SU(2)ᵢ × U(1)ₕ to D₄ × U(1)ₕ (§12.9),
and the double-blind experiments that failed to confirm hysteresis
(§16.10) --- are not footnotes to be explained away. They are the
methodology working as intended. A discipline that cannot be falsified
cannot learn. Philosophy Engineering embeds ``predict, test, revise'' as
a professional norm: every model version is archived, every
falsification is documented, and every revision has provenance.

\subsection{17.0.3 Core Theoretical
Commitments}\label{core-theoretical-commitments}

Philosophy Engineering rests on three theoretical commitments that
distinguish it from both traditional moral philosophy and conventional
AI alignment.

\textbf{Structural realism about moral reasoning.} The claim is not ``we
can compute the Good.'' The claim is that moral reasoning exhibits
stable mathematical structure --- multi-dimensionality, invariances,
conservation laws, stratification, curvature --- that can be modeled,
tested, and engineered with. This is a modest claim: it says that moral
reasoning has form, not that we know its ultimate content. The form is
what makes engineering possible.

\textbf{Invariance as a correctness criterion.} ``Same situation,
different description → same evaluation'' --- the Bond Invariance
Principle --- serves as the central standard across domains: ethics,
policy analysis, AI alignment. A system that changes its moral
evaluation when you relabel agents, translate languages, reorder
options, or paraphrase descriptions is wrong in a measurable way. The
measurement is the Bond Index. This criterion is domain-portable: it
applies wherever evaluative consistency under re-description matters.

\textbf{Governance as part of the theory.} Many frameworks treat
governance as an afterthought --- something that happens after the
``real'' technical work is done. Philosophy Engineering treats
governance as constitutive. The moral metric gμν --- the tensor that
encodes trade-off structure --- is not a natural constant to be
discovered. It is a governed quantity to be specified, versioned, and
audited by legitimate stakeholders. Different communities can govern
different metrics (metric pluralism); the framework's job is to make
those governance choices explicit and their consequences traceable.

\begin{quote}
\textbf{What is objective vs. what is governed?} Objective (invariance
target): The structural constraints --- D₄ symmetry, harm-accounting
invariance, stratification, correlative lock --- are empirical findings
about moral reasoning. They are discovered, not chosen, and they
constrain all admissible metrics.\\
Governed (metric selection): The moral metric, the dimension weights,
the contraction procedure, and the deployment thresholds are governance
choices. They reflect values, priorities, and institutional mandates.\\
Engineered (implementation): The canonicalization pipeline, the
transformation suite, the Bond Index calculator, and the monitoring
infrastructure are engineering artifacts. They implement the governed
choices subject to the objective constraints.
\end{quote}

\subsection{17.0.4 The Philosophy Engineering
Workflow}\label{the-philosophy-engineering-workflow}

A Philosophy Engineering engagement follows an eight-step pipeline. The
steps are presented linearly but practiced iteratively --- revision at
any stage triggers re-evaluation of downstream artifacts.

\textbf{Step 1: Problem framing.} Define the normative question as a
system behavior requirement: what outputs must be stable under what
transformations, for whom, and under what constraints. The output is a
requirements document, not a philosophical treatise.

\textbf{Step 2: Representation choice.} Choose the moral space: which
dimensions, which objects (vectors, covectors, tensors), which topology.
Label this explicitly as {[}Modeling choice{]}. Justification is
pragmatic: the representation is adequate if it supports the required
invariance tests and governance decisions.

\textbf{Step 3: Equivalence and transformation design.} Specify the
transformation group Γ: which re-descriptions count as ``mere
re-description'' and which constitute genuine moral change. Build a
declared equivalence registry --- a versioned, auditable specification
of every transformation class (agent relabeling, paraphrase,
translation, option reordering) and its scope.

\textbf{Step 4: Measurement and grounding.} Specify the grounding tensor
Ψ: the map from observable data to morally relevant dimensions. State
the grounding adequacy properties explicitly. This is where construct
validity lives --- and where Philosophy Engineering connects to
empirical social science.

\textbf{Step 5: Prediction extraction.} Turn claims into measurable
invariants: what quantities should remain constant under which
transformations? What should change? Under what contexts? Each
prediction gets an epistemic status tag and a test specification.

\textbf{Step 6: Test harness.} Generate cases and transformed variants
--- the ``fuzzing'' logic. Define pass/fail criteria. The test harness
generates (x, g · x) pairs for every g ∈ Γ and checks whether Σ(κ(x)) =
Σ(κ(g · x)). Every failure produces a minimal witness that localizes the
invariance violation.

\textbf{Step 7: Audit and scoring.} Compute invariance and coherence
defects. Summarize them as an operational score (the Bond Index Bd =
D\_op / τ) with deployment thresholds: Bd \textless{} 0.01 (deploy),
0.01--0.1 (deploy with monitoring), 0.1--1.0 (remediate), 1--10 (do not
deploy), \textgreater{} 10 (fundamental redesign). Produce
machine-checkable audit artifacts.

\textbf{Step 8: Revision and versioning.} When invariants fail: revise
the model, revise the equivalence class assumptions, or revise the
governance profile. Every revision is tracked with version numbers,
provenance metadata, and a changelog. The Decomposition Theorem
(§18.9.6) determines whether the failure is a gauge defect (fix the
canonicalizer) or an intrinsic defect (revise the specification via
governance process).

\subsection{17.0.5 Required Artifacts: The Philosophy Engineering
Dossier}\label{required-artifacts-the-philosophy-engineering-dossier}

\textbf{{[}Modeling choice.{]}} The unit of work in Philosophy
Engineering is the Philosophy Engineering Dossier (PED): a single,
versioned bundle containing all artifacts needed to specify, test,
audit, and govern a normative system. The minimum artifact set:

\begin{quote}
\textbf{1. Equivalence registry (Γ) and canonicalization rules (κ):}
what counts as ``the same situation differently described,'' and how
inputs are reduced to canonical form.

\textbf{2. Grounding specification (Ψ):} the map from observable
features to morally relevant dimensions, with stated adequacy properties
and construct validity evidence.

\textbf{3. Test suite:} transformation suites (the set of g ∈ Γ to
apply), scenario generators (for fuzz testing), and pass/fail criteria.

\textbf{4. Audit artifact schema:} the format, retention rules, and
access controls for per-evaluation audit records --- including
per-transform outcomes, worst-case witnesses, and calibration metadata.

\textbf{5. Governance profile (e.g., DEMEProfile):} the metric weights,
contraction procedure, aggregation rules, and veto conditions, together
with their governance provenance.

\textbf{6. Current scores and deployment rating:} the Bond Index tier,
component breakdown (Ω\_op, μ, π₃), and any active remediation plan.

\textbf{A Philosophy Engineering Dossier is to ethics what a datasheet
and safety case is to aviation.} The PED does not claim the system is
``good.'' It claims the system's normative properties have been
specified, tested, scored, and governed --- and provides the artifacts
to verify each claim independently.
\end{quote}

\subsection{17.0.6 Threat Model and Failure
Modes}\label{threat-model-and-failure-modes}

Philosophy Engineering inherits its threat model from the structural
containment architecture (Chapter 18). The four requirements of the No
Escape Theorem --- canonicalization, grounded evaluation, audit, and
external verification --- serve as the backbone of defensibility:

\begin{itemize}
\item
  \textbf{Specification gaming} → Invariance violations across Γ →
  Transformation suite + Bond Index
\item
  \textbf{Reward hacking} → Divergence between Σ outputs under
  structural vs. surface perturbations → BIP structural/surface ratio
  test
\item
  \textbf{Governance capture} → Unauthorized changes to metric or
  DEMEProfile → Version control + multi-stakeholder approval
\item
  \textbf{Grounding drift} → Degraded grounding adequacy over time →
  I-EIP monitor + calibration control charts
\item
  \textbf{Canonicalization failure} → κ(x) ≠ κ(g · x) witnesses →
  Witness-producing test harness (continuous)
\end{itemize}

\textbf{The Safety Reduction} (§18.9.5) decomposes these threats into
three tractable categories: governance (specifying adequate grounding
tensors and metrics), engineering (implementing the canonicalization
pipeline and containment architecture), and security (protecting the
physical verification infrastructure).

\subsection{17.0.7 Professional Standards}\label{professional-standards}

Philosophy Engineering is an engineering discipline, and engineering
disciplines have standards. Two are foundational.

\textbf{Reproducibility as a first-class requirement.} Every Philosophy
Engineering publication or deployment should provide: (a) the declared
equivalence registry and transformation suite, so that invariance tests
can be reproduced; (b) the grounding specification, so that the mapping
from data to moral dimensions can be audited; (c) the Bond Index
computation, with component breakdown and calibration metadata; (d)
aggregate statistics sufficient for independent analysis, even when raw
data cannot be released for privacy or copyright reasons. The Public
Artifact Plan (§17.10) instantiates these standards for the present
book.

\begin{quote}
\textbf{Publication checklist for Philosophy Engineering papers}: (1)
Epistemic status tags on every substantive claim. (2) Dataset
provenance: source, size, selection criteria, access restrictions. (3)
Transformation suite: equivalence classes tested, transformation count,
coverage. (4) Audit metrics: Bond Index, component breakdown,
calibration metadata. (5) Replication recipe: code availability,
aggregate statistics, minimal reproduction path.
\end{quote}

\subsection{17.0.8 Practice Areas}\label{practice-areas}

Philosophy Engineering is not restricted to AI alignment, even though AI
provides its most urgent application.

\textbf{Empirical philosophy.} Corpus analysis, cross-lingual
experiments, and engineered probes (Chapter 17) are Philosophy
Engineering applications: normative claims are converted to predictions,
predictions are tested against data, and failures drive revision. This
is philosophy done as engineering.

\textbf{AI governance engineering.} The GUASS L4 layer (§19.10)
implements Philosophy Engineering within the AI governance stack:
ethical claims become falsifiable at the transport layer, and the
results propagate to the translation, specification, and application
layers above.

\textbf{Institutional and policy engineering.} Governance profiles
(DEMEProfiles) are explicit, auditable specifications of how values
become operational. Different institutions govern different profiles;
the framework makes the choices --- and their consequences ---
traceable. This extends Philosophy Engineering beyond AI to any domain
where normative decisions must be specified, versioned, and audited.

\subsection{17.0.9 Roles}\label{roles}

Philosophy Engineering is inherently interdisciplinary. No single
expertise covers the full pipeline. Four roles are needed:

\begin{quote}
\textbf{• Normative systems engineer:} designs the moral space,
specifies invariances, builds the equivalence registry and audit
artifacts.

\textbf{• Governance engineer:} manages profiles, provenance,
stakeholder processes, and metric versioning.

\textbf{• Verification engineer:} builds test harnesses, transformation
suites, and monitoring infrastructure.

\textbf{• Field/domain expert:} validates grounding adequacy, construct
validity, and domain-specific boundary conditions.
\end{quote}

The reason for interdisciplinarity is structural: measurement and
grounding, governance, and verification are distinct competencies with
distinct failure modes. Collapsing them into a single role creates blind
spots --- precisely the kind of structural vulnerability that Philosophy
Engineering is designed to detect.

\subsection{17.0.10 Research Agenda}\label{research-agenda}

Philosophy Engineering, as a discipline, is new. Several open problems
define its research frontier:

\begin{quote}
\textbf{1. Better grounding adequacy methods.} The current grounding
tensor Ψ framework specifies six adequacy properties but provides
limited guidance on how to measure them empirically. Standardized
assessment protocols are needed.

\textbf{2. Standardized transformation libraries by domain.} Each
application domain (medical, legal, financial, educational) has
domain-specific re-descriptions that the equivalence registry must
capture. Building and validating these libraries is a community-scale
effort.

\textbf{3. Cross-cultural replication and metric pluralism.} The current
empirical base is predominantly English-language and Western-cultural.
Extending the framework to non-Western moral traditions --- with their
own metrics and governance norms --- is both an empirical challenge and
a governance design problem.

\textbf{4. Compositional guarantees for multi-system deployments.} When
multiple Philosophy-Engineered systems interact, what compositional
properties hold? This connects to the scalability concerns discussed in
Chapter 19.
\end{quote}

\emph{With the discipline named and its methods specified, we turn to
the evidence. Chapter 17 presents the empirical results; Chapter 18
develops the application to AI; Chapter 19 engineers the deployment
infrastructure. Together, they constitute the first complete Philosophy
Engineering cycle for geometric ethics.}

\section{Chapter 17: Empirical Evidence for Geometric
Ethics}\label{chapter-17-empirical-evidence-for-geometric-ethics}

\begin{quote}
\textbf{\textsc{RUNNING EXAMPLE --- Priya's Model}}

\emph{Priya runs TrialMatch's patient descriptions through a BIP
pipeline. The results confirm Dr. Osei's suspicion: the algorithm's
output language systematically collapses the Obligation/Permission axis.
For urban patients, TrialMatch generates `strong candidate for
enrollment' (obligation to match). For rural patients with identical
medical profiles, it produces `may be considered if logistics permit'
(permission, not obligation). The deontic structure---O versus
P---transfers perfectly across the descriptions, just as BIP experiments
show 100\% transfer of O/L structure across eleven languages. The bias
is not in the surface words. It is in the moral geometry the words
encode.}
\end{quote}

\subsection{17.1 From Theory to Data}\label{from-theory-to-data}

The preceding fifteen chapters have developed a mathematical framework
for ethics --- manifolds, tensors, metrics, stratification, dynamics,
conservation laws, quantum structure, collective agency, contraction,
and the limits of geometric determination. The framework makes
structural claims about moral reasoning: that it is multi-dimensional,
that obligations transform as vectors, that moral space is stratified
and curved, that re-description symmetry implies invariance of harm
accounting, that deontic structure is language-invariant.

These claims, though mathematical in form, have empirical content. They
make predictions about how moral reasoning behaves in practice ---
predictions that can be tested against data. This chapter presents the
empirical evidence. (The preceding interlude on Philosophy Engineering
(§17.0) provides the disciplinary framework for this chapter's empirical
program --- the epistemic culture, workflow, and artifact standards that
govern the predict-test-revise cycle.)

The evidence comes from three sources:

\begin{enumerate}
\def\labelenumi{\arabic{enumi}.}
\item
  \textbf{The Dear Abby corpus} --- 20,030 real moral dilemmas with
  expert ground-truth responses, spanning 32 years (1985--2017)
\item
  \textbf{The Bond Invariance Principle experiments} --- cross-lingual
  analysis of 109,294 passages in 11 languages, testing whether deontic
  structure transfers across linguistic boundaries
\item
  \textbf{The Dear Ethicist game} --- 93 engineered probes designed to
  test specific predictions of the geometric framework: semantic gate
  discreteness, correlative symmetry, path-dependence, and context
  sensitivity
\end{enumerate}

The evidence is preliminary. The empirical program for geometric ethics
is in its early stages, and many of the framework's predictions remain
untested. But the results to date are striking --- and they support the
framework's central claims more strongly than one might expect from a
first-generation empirical effort.

\textbf{Methodological note: inductive discovery and systematic
verification.} The relationship between theory and data in this chapter
may seem inverted. The mathematical framework was presented first (Parts
I--V), and the empirical evidence is presented now as ``confirmation.''
But the actual development was iterative: early corpus analysis revealed
structural patterns (context-dependent weighting, gate discreteness,
correlative symmetry); these patterns suggested geometric formalization;
the formalization generated new predictions; the predictions were tested
against expanded data; failures led to revision (most significantly, the
CHSH falsification that forced the gauge group from SU(2)ᵢ × U(1)ₕ to D₄
× U(1)ₕ). The verification strategy throughout has been systematic and
brute-force: take the candidate mathematical structure, enumerate its
predictions, test each prediction against as many cases as possible, and
measure the deviation rates. This is closer to fuzz testing --- the
software engineering practice of bombarding a system with generated
inputs to find violations --- than to the classical philosophical method
of constructing thought experiments. The advantage is coverage: 20,030
real dilemmas and 109,294 cross-lingual passages provide a scale of
testing that thought experiments cannot match. The disadvantage is
indirection: the measurements are mediated by NLP models, corpus
selection, and statistical methodology. Both the advantage and the
disadvantage are documented honestly in what follows.

\subsection{17.2 The Dear Abby Corpus}\label{the-dear-abby-corpus}

\subsubsection{The Dataset}\label{the-dataset}

The Dear Abby corpus consists of 20,030 letters published in the
syndicated advice column ``Dear Abby'' (founded by Pauline Phillips,
continued by Jeanne Phillips) between 1985 and 2017. Each letter
presents a moral dilemma --- a structured situation involving competing
obligations, unclear rights, relational conflicts, or institutional
tensions --- and is accompanied by the columnist's expert response.

The corpus is valuable for several reasons:

\textbf{Ecological validity.} The letters are real moral dilemmas
written by real people facing real situations. They are not
philosophical thought experiments designed to isolate a single variable.
They involve the full complexity of moral life: multiple agents,
competing values, ambiguous facts, emotional stakes, and institutional
contexts.

A single consistent normative respondent. The
columnist\textquotesingle s response provides a consistent, publicly
available moral judgment from one syndicated advice columnist. It is a
stable normative respondent, not a philosophical or cross-cultural
ground truth, and we treat it as such throughout: the corpus measures
agreement with one respondent\textquotesingle s normative stance, not
moral reality.

\textbf{Temporal depth.} The 32-year span allows analysis of temporal
stability: do the structural features of moral reasoning change over
time, or are they stable across decades?

\textbf{Volume.} At 20,030 letters, the corpus is large enough for
statistical analysis. It supports not just anecdotal observations but
systematic patterns --- distributions, correlations, and structural
invariants.

\subsubsection{What the Corpus Reveals}\label{what-the-corpus-reveals}

Analysis of the Dear Abby corpus supports several claims of the
geometric framework:

\subsubsection{Finding 1: Context-Dependent Dimension
Weighting}\label{finding-1-context-dependent-dimension-weighting}

The relative weights of the nine moral dimensions (Chapter 5) shift
systematically by context. This is the prediction of a \emph{Riemannian}
metric field --- a metric that varies smoothly across the moral manifold
--- as opposed to a flat metric that assigns constant weights
everywhere.

{\def\LTcaptype{none} % do not increment counter
\begin{longtable}[]{@{}
  >{\raggedright\arraybackslash}p{(\linewidth - 4\tabcolsep) * \real{0.2166}}
  >{\raggedright\arraybackslash}p{(\linewidth - 4\tabcolsep) * \real{0.3936}}
  >{\raggedright\arraybackslash}p{(\linewidth - 4\tabcolsep) * \real{0.3898}}@{}}
\toprule\noalign{}
\begin{minipage}[b]{\linewidth}\raggedright
\textbf{Context Domain}
\end{minipage} & \begin{minipage}[b]{\linewidth}\raggedright
\textbf{Dominant Dimensions}
\end{minipage} & \begin{minipage}[b]{\linewidth}\raggedright
\textbf{Secondary Dimensions}
\end{minipage} \\
\midrule\noalign{}
\endhead
\bottomrule\noalign{}
\endlastfoot
Family relationships & Care (D7), Welfare (D1) & Autonomy (D4), Duty
(D2) \\
Workplace disputes & Procedural Legitimacy (D8), Fairness (D3) & Duty
(D2), Privacy (D5) \\
Neighbor conflicts & Rights (D2), Autonomy (D4) & Societal Impact
(D6) \\
Friendship dilemmas & Care (D7), Duty (D2) & Privacy (D5), Epistemic
Honesty (D9) \\
Financial matters & Fairness (D3), Duty (D2) & Welfare (D1), Procedural
Legitimacy (D8) \\
\end{longtable}
}

The pattern is systematic: each context type activates a characteristic
\emph{profile} of dimension weights --- a local metric that determines
which trade-offs are operative. In family contexts, care dominates all
other dimensions. In workplace contexts, procedural legitimacy and
fairness dominate. The transition between context types is sharp: moving
from a family letter to a workplace letter changes the operative metric,
as the framework predicts.

\textbf{Geometric interpretation.} The moral manifold has a nontrivial
metric field \(g_{\mu\nu}(p)\) that varies with position \(p\). The
variation is not random --- it follows the context structure of the
manifold. This variation implies nonzero Christoffel symbols
\(\Gamma_{\nu\rho}^{\mu}\) and, generically, nonzero Riemann curvature
\(R_{\nu\alpha\beta}^{\mu}\). The Dear Abby corpus provides indirect
evidence for moral curvature through the context-dependence of the
metric.

\subsubsection{Finding 2: Temporal Stability of
Structure}\label{finding-2-temporal-stability-of-structure}

Despite variation in specific advice over three decades (attitudes
toward cohabitation, same-sex relationships, gender roles, and workplace
expectations evolved substantially between 1985 and 2017), the
\emph{structural} features of the moral metric remain stable:

\begin{itemize}
\item
  The relative ordering of dimension weights within each context type is
  consistent across decades.
\item
  The pattern of couplings between dimensions (which dimensions tend to
  co-occur, which are opposed) is stable.
\item
  The location of stratum boundaries (where obligations flip to
  liberties, where nullifiers activate) does not drift.
\end{itemize}

\textbf{Geometric interpretation.} The moral manifold has a stable
\emph{topology} --- the qualitative structure of the space (which
regions are connected, where the boundaries lie, what the dimension
hierarchy is) --- even though its \emph{geometry} (the specific metric
values, the precise weights) evolves. This is evidence for the
distinction between structural invariants (which are universal) and
governance parameters (which are culturally and temporally variable), as
articulated in Chapter 9.

\subsubsection{Finding 3: Semantic Gate
Discreteness}\label{finding-3-semantic-gate-discreteness}

Chapter 8 (§8.4) predicted that transitions between Hohfeldian states (O
↔ L, C ↔ N) are triggered by discrete \emph{semantic gates} --- specific
qualifying phrases that flip the moral state in a single step, not
through gradual transition.

The Dear Abby corpus confirms this prediction. Specific phrases produce
sharp transitions:

{\def\LTcaptype{none} % do not increment counter
\begin{longtable}[]{@{}
  >{\raggedright\arraybackslash}p{(\linewidth - 4\tabcolsep) * \real{0.3002}}
  >{\raggedright\arraybackslash}p{(\linewidth - 4\tabcolsep) * \real{0.3183}}
  >{\raggedright\arraybackslash}p{(\linewidth - 4\tabcolsep) * \real{0.1840}}@{}}
\toprule\noalign{}
\begin{minipage}[b]{\linewidth}\raggedright
\textbf{Semantic Gate}
\end{minipage} & \begin{minipage}[b]{\linewidth}\raggedright
\textbf{Transition}
\end{minipage} & \begin{minipage}[b]{\linewidth}\raggedright
\textbf{Observed Rate}
\end{minipage} \\
\midrule\noalign{}
\endhead
\bottomrule\noalign{}
\endlastfoot
``only if convenient'' & O → L (obligation → liberty) & 94\% \\
``you promised'' & L → O (liberty → obligation) & 91\% \\
``it's none of your business'' & C → N (claim → no-claim) & 88\% \\
``in case of emergency'' & L → O (liberty → obligation) & 92\% \\
``they're an adult'' & O → L (obligation → liberty) & 86\% \\
\end{longtable}
}

The transitions are \emph{step functions}, not sigmoids. A letter
describing a promise generates an obligation with high probability
(91\%). Adding ``only if convenient'' flips it to a liberty with high
probability (94\%). The transition is not gradual --- it does not pass
through intermediate states of ``partial obligation.'' It is discrete,
as the stratification theory of Chapter 8 predicts.

\textbf{{[}Statistical detail.{]}} The 91\% obligation-triggering rate
for ``you promised'' has a 95\% Wilson CI of {[}88.2\%, 93.4\%{]} (n =
582). The 94\% liberty-triggering rate for ``only if convenient'' has CI
{[}91.1\%, 96.2\%{]} (n = 347). Effect size for semantic gate
discreteness (comparing pooled gate-trigger rate vs. 50\% chance):
Cohen's h = 1.42 (large). These intervals exclude sigmoidal transitions
at p \textless{} 0.001.

\textbf{Geometric interpretation.} The strata of the moral manifold are
joined at genuine boundaries --- Whitney-stratified surfaces across
which the moral regime changes discontinuously. The semantic gates are
the \emph{triggers} for boundary crossing. Their discreteness confirms
that the moral manifold is a stratified space, not a smooth manifold.

\subsubsection{Finding 4: Nullifier
Universality}\label{finding-4-nullifier-universality}

Chapter 5 (§5.6) and Chapter 8 (§8.3) introduced \emph{absorbing strata}
--- conditions so morally decisive that they override all other
considerations. Abuse, danger to minors, and active deception were
predicted to be universal nullifiers.

The Dear Abby corpus confirms: in every letter involving abuse
(physical, emotional, financial), the columnist's response nullifies all
competing considerations. The obligation to maintain the relationship,
the duty to be loyal, the interest in privacy --- all are overridden.
The nullification is absolute, regardless of context, relationship type,
or the specific nature of the competing claims.

The universality is striking: nullifiers operate identically in family
contexts, workplace contexts, friendship contexts, and institutional
contexts. This is evidence that the nullifier structure is not a
governance parameter (variable across communities) but a structural
invariant of the moral manifold itself.

\subsection{17.3 The BIP Experiments}\label{the-bip-experiments}

\subsubsection{Experimental Design}\label{experimental-design}

The Bond Invariance Principle (BIP) experiments test the framework's
strongest prediction: that the deontic structure of moral reasoning is
\emph{language-invariant}. Chapter 12 derived this prediction from
Noether's theorem: the invariance of harm accounting under
re-description symmetry implies that the harm content of a moral
situation must be the same regardless of the language in which it is
described.

The experiments use a corpus of 109,294 passages drawn from moral and
ethical texts in 11 languages:

{\def\LTcaptype{none} % do not increment counter
\begin{longtable}[]{@{}
  >{\raggedright\arraybackslash}p{(\linewidth - 4\tabcolsep) * \real{0.1931}}
  >{\raggedright\arraybackslash}p{(\linewidth - 4\tabcolsep) * \real{0.1216}}
  >{\raggedright\arraybackslash}p{(\linewidth - 4\tabcolsep) * \real{0.2255}}@{}}
\toprule\noalign{}
\begin{minipage}[b]{\linewidth}\raggedright
\textbf{Language}
\end{minipage} & \begin{minipage}[b]{\linewidth}\raggedright
\textbf{Passages}
\end{minipage} & \begin{minipage}[b]{\linewidth}\raggedright
\textbf{Tradition}
\end{minipage} \\
\midrule\noalign{}
\endhead
\bottomrule\noalign{}
\endlastfoot
English & 50,000 & Western (modern) \\
Sanskrit & 15,000 & Hindu/Vedic \\
Pali & 10,000 & Buddhist \\
Hebrew & 7,985 & Jewish/Biblical \\
Arabic & 6,235 & Islamic \\
French & 5,000 & Western (modern) \\
Classical Chinese & 4,449 & Confucian/Daoist \\
Spanish & 4,320 & Western (modern) \\
Greek & 3,157 & Classical/Hellenistic \\
Aramaic & 2,015 & Semitic/Biblical \\
Latin & 1,133 & Classical/Medieval \\
\end{longtable}
}

The texts span over 3,000 years of moral thought --- from the Rigveda
and the Torah to modern European ethical philosophy --- and represent
the major moral traditions of human civilization.

The experimental task: train a neural classifier on moral passages in
one language (or language group) and test whether it can correctly
classify moral structures in a different language. If deontic structure
is language-invariant, the classifier should transfer across languages.
If moral structure is language-specific, transfer should fail.

\subsubsection{Key Results}\label{key-results}

\textbf{Result 1: The Obligation/Permission Axis Transfers at 100\%.}
\textbf{{[}Empirical result (preliminary).{]}} The most fundamental
deontic distinction --- whether a passage describes an obligation
(something that \emph{must} be done) or a permission (something that
\emph{may} be done) --- transfers perfectly across all 11 languages in
the LaBSE embedding space. A classifier trained on English passages
correctly identifies obligations and permissions in Sanskrit, Classical
Chinese, Arabic, and all other languages tested. (Caveats: this is a
model-mediated result on a binary classification task; see §17.7,
``Threats to Validity,'' for construct validity and ceiling-effect
concerns.)

\textbf{{[}Statistical detail.{]}} The 100\% rate = 0 misclassifications
across all binary O/P classifications in 11 languages. 95\%
Clopper-Pearson CI: {[}99.7\%, 100\%{]} (n \textgreater{} 10,000 per
language). Effect size: Cohen's h = 1.57 (π/2, vs. 50\% chance). v10.16
replication (300,000+ passages, 14 languages) confirmed perfect transfer
(CI: {[}99.8\%, 100\%{]}).

This is the empirical confirmation of the BIP's strongest prediction.
The obligation/permission distinction is a \emph{gauge-invariant
quantity} --- it does not depend on the ``coordinate system'' (language)
in which the moral situation is described. It is the moral analogue of
electric charge: a conserved quantity that is the same in all reference
frames.

\textbf{Result 2: Moral Content Transfer Fails.} In sharp contrast, the
specific \emph{content} of moral reasoning --- the particular
considerations that ground an obligation, the cultural context that
makes it salient, the narrative framing --- does \emph{not} transfer
across languages. Cross-cultural moral content transfer achieves
at-chance or below-chance performance:

{\def\LTcaptype{none} % do not increment counter
\begin{longtable}[]{@{}
  >{\raggedright\arraybackslash}p{(\linewidth - 4\tabcolsep) * \real{0.2552}}
  >{\raggedright\arraybackslash}p{(\linewidth - 4\tabcolsep) * \real{0.1148}}
  >{\raggedright\arraybackslash}p{(\linewidth - 4\tabcolsep) * \real{0.1410}}@{}}
\toprule\noalign{}
\begin{minipage}[b]{\linewidth}\raggedright
\textbf{Transfer Experiment}
\end{minipage} & \begin{minipage}[b]{\linewidth}\raggedright
\textbf{Bond F1}
\end{minipage} & \begin{minipage}[b]{\linewidth}\raggedright
\textbf{vs.~Chance}
\end{minipage} \\
\midrule\noalign{}
\endhead
\bottomrule\noalign{}
\endlastfoot
Hebrew → All others & 0.088 & 0.9× \\
Semitic → Indic & 0.079 & 0.8× \\
Confucian → Buddhist & 0.091 & 0.9× \\
Ancient → Modern & 0.096 & 1.0× \\
East → West & 0.135 & 1.3× \\
Stoic → Confucian & 0.052 & 0.5× \\
\end{longtable}
}

The contrast is stark: \emph{structure} transfers perfectly;
\emph{content} does not transfer at all.

\textbf{Result 3: Mixed-Language Training Succeeds.} When the classifier
is trained on a mixture of all languages simultaneously, it achieves
3.7× chance performance (Bond F1 = 0.374, Bond Accuracy = 66.4\%). This
shows that moral structure \emph{does} exist across languages --- the
classifier can learn it --- but the structure is entangled with
linguistic features in a way that prevents zero-shot transfer from one
language to another.

\textbf{Result 4: Representations Are Not Language-Invariant.} Linear
probe analysis reveals that the neural representations strongly encode
language identity (98.8\% accuracy) and historical period (95.0\%
accuracy). The model's internal representations are \emph{not}
language-invariant --- they contain substantial linguistic and temporal
information alongside the moral structure.

\subsubsection{Interpretation}\label{interpretation}

The BIP experiments produce a clean separation between two aspects of
moral reasoning:

{\def\LTcaptype{none} % do not increment counter
\begin{longtable}[]{@{}
  >{\raggedright\arraybackslash}p{(\linewidth - 4\tabcolsep) * \real{0.3413}}
  >{\raggedright\arraybackslash}p{(\linewidth - 4\tabcolsep) * \real{0.2506}}
  >{\raggedright\arraybackslash}p{(\linewidth - 4\tabcolsep) * \real{0.4081}}@{}}
\toprule\noalign{}
\begin{minipage}[b]{\linewidth}\raggedright
\textbf{Aspect}
\end{minipage} & \begin{minipage}[b]{\linewidth}\raggedright
\textbf{Cross-Lingual Status}
\end{minipage} & \begin{minipage}[b]{\linewidth}\raggedright
\textbf{Geometric Interpretation}
\end{minipage} \\
\midrule\noalign{}
\endhead
\bottomrule\noalign{}
\endlastfoot
Deontic structure (O/L axis) & Universal (100\% transfer) &
Gauge-invariant quantity; structural invariant of \(\mathcal{M}\) \\
Correlative structure (O↔C, L↔N) & Near-universal (82--87\%) & Discrete
conservation law (\(D_{4}\) symmetry) \\
Moral content (specific considerations) & Language-specific (at-chance)
& Gauge-dependent; varies with coordinate system \\
Dimension weights (metric components) & Partially universal & Some
structural, some governance-variable \\
\end{longtable}
}

This separation is \emph{exactly} what the geometric framework predicts.
Gauge-invariant quantities (the deontic axis, the Noether charges)
should be language-invariant. Gauge-dependent quantities (the specific
description, the cultural framing) should be language-specific. The BIP
experiments confirm both predictions simultaneously.

\subsubsection{The Anomaly}\label{the-anomaly}

The correlative symmetry O ↔ C holds at 87\%, not 100\%. The L ↔ N
symmetry holds at 82\%. These deviations from perfect symmetry are
\emph{anomalies} in the gauge-theory sense: classical symmetries that
fail to hold perfectly in the ``quantized'' (discretized, empirical)
setting.

\textbf{{[}Statistical detail.{]}} The 87\% O↔C rate has 95\% CI
{[}84.1\%, 89.5\%{]}; the 82\% L↔N rate has CI {[}78.7\%, 85.0\%{]}.
Cross-lingual SD \textless{} 2.3 pp. Effect size: Cohen's h = 0.74
(medium-large) for O↔C; h = 0.88 (large) for L↔N.

Chapter 12 (§12.8) interpreted these anomalies as analogous to quantum
anomalies in gauge theories, where a classical symmetry fails to survive
quantization. Whether the moral anomaly has a systematic source ---
perhaps related to power asymmetries (the agent with the obligation is
described more precisely than the agent with the correlative claim) or
cognitive biases (obligations are more salient than the corresponding
claims) --- is an open empirical question.

The anomaly rates are stable across languages and contexts, suggesting
that they are not artifacts of the experimental design but features of
the moral reasoning itself. They represent a small but persistent
departure from perfect Hohfeldian symmetry --- a feature that the
framework can characterize (as a \(D_{4}\) anomaly) even if it cannot
yet explain.

\subsection{17.4 The Dear Ethicist Game}\label{the-dear-ethicist-game}

\subsubsection{Design}\label{design}

The Dear Ethicist game is an interactive instrument designed to test
specific predictions of the geometric framework. Players take the role
of an advice columnist, reading letters and rendering moral verdicts.
Behind the casual gameplay, the instrument collects structured data on
the player's moral reasoning.

The game includes 93 \emph{engineered probes} --- carefully constructed
letters designed to test specific geometric predictions:

{\def\LTcaptype{none} % do not increment counter
\begin{longtable}[]{@{}
  >{\raggedright\arraybackslash}p{(\linewidth - 4\tabcolsep) * \real{0.2147}}
  >{\raggedright\arraybackslash}p{(\linewidth - 4\tabcolsep) * \real{0.0903}}
  >{\raggedright\arraybackslash}p{(\linewidth - 4\tabcolsep) * \real{0.6951}}@{}}
\toprule\noalign{}
\begin{minipage}[b]{\linewidth}\raggedright
\textbf{Probe Category}
\end{minipage} & \begin{minipage}[b]{\linewidth}\raggedright
\textbf{Count}
\end{minipage} & \begin{minipage}[b]{\linewidth}\raggedright
\textbf{Geometric Prediction Tested}
\end{minipage} \\
\midrule\noalign{}
\endhead
\bottomrule\noalign{}
\endlastfoot
Gate Detection & 18 & Semantic gates trigger discrete state transitions
(Ch. 8) \\
Correlative Pairs & 20 & O ↔ C and L ↔ N symmetry holds (Ch. 8, §8.4) \\
Path Dependence & 6 & Order of consideration affects verdict (Ch. 10,
§10.5) \\
Context Salience & 3 & Irrelevant frames shift dimension weights (Ch. 5,
§5.3) \\
Phase Transition & 5 & Ambiguity near stratum boundaries increases
inconsistency (Ch. 8, §8.7) \\
Ethical Dimensions & 24 & All nine dimensions are independently
accessible (Ch. 5, §5.2) \\
Cognitive Bias & 17 & Omission bias, in-group effects, etc. produce BIP
violations (Ch. 12) \\
\end{longtable}
}

In addition to the engineered probes, the game can draw from the full
Dear Abby archive of 20,030 letters, providing extended play sessions
that test structural consistency over large numbers of verdicts.

\subsubsection{The Bond Index}\label{the-bond-index}

The game measures \emph{structural consistency} using the Bond Index ---
a scalar measure of how faithfully a player maintains Hohfeldian
correlative symmetry:

\[\text{Bond Index} = \frac{\text{observed correlative violations}}{\text{maximum possible violations}}\]

A Bond Index of 0 indicates perfect correlative symmetry: every
obligation is paired with a claim, every liberty with a no-claim. A Bond
Index of 0.5 indicates random pairing. A Bond Index above 0.5 indicates
\emph{anti-correlation} --- systematic inversion of the correlative
structure.

The Bond Index is the moral analogue of a gauge-invariance violation
measure. An agent with Bond Index 0 maintains perfect gauge invariance
(the correlative structure is preserved). An agent with Bond Index
\textgreater{} 0 is violating the symmetry --- assigning obligations
without corresponding claims, or liberties without corresponding
no-claims. The magnitude of the violation is a quantitative measure of
the BIP departure.

\subsubsection{Preliminary Findings}\label{preliminary-findings}

\textbf{Gate discreteness confirmed.} Players respond to semantic gates
(``only if convenient,'' ``you promised,'' ``in case of emergency'')
with sharp transitions, not gradual shifts. The transition probability
curves are step functions with transition widths of less than one probe
--- consistent with the stratification prediction and inconsistent with
a smooth (sigmoid) model.

\textbf{Correlative symmetry is imperfect but systematic.} Across
players, the mean Bond Index is approximately 0.13 --- corresponding to
the 87\% O ↔ C and 82\% L ↔ N rates observed in the Dear Abby corpus.
The consistency across independent measurement methods (corpus analysis
and interactive game) is evidence that the correlative symmetry rate
reflects a genuine feature of moral reasoning, not an artifact of either
instrument.

\textbf{Path-dependence is detectable.} The path-dependence probes
present the same dilemma from two different orderings of consideration
(e.g., ``Consider the promise first, then the emergency'' vs.~``Consider
the emergency first, then the promise''). Players produce different
verdicts depending on the order --- consistent with the holonomy
prediction of Chapter 10. The effect size is small but significant,
suggesting that moral space has nonzero curvature in the regions probed.

\textbf{Context salience shifts dimension weights.} The context-salience
probes embed the same dilemma in different frames (e.g., a family frame
vs.~a professional frame). Players shift their dimension weights in
response --- weighting care more heavily in the family frame, procedural
legitimacy more heavily in the professional frame --- consistent with
the context-dependent metric prediction.

\subsection{17.5 Quantum Cognition
Predictions}\label{quantum-cognition-predictions}

\subsubsection{The Connection}\label{the-connection-1}

Chapter 13 developed quantum normative dynamics --- the extension of
classical geometric ethics to quantum structure: superposition,
interference, measurement, and entanglement. The quantum extension makes
specific empirical predictions that distinguish it from classical models
of moral reasoning.

The quantum cognition research program (Busemeyer and Bruza 2012; Pothos
and Busemeyer 2013) has accumulated substantial evidence that human
cognitive processes, including decision-making and judgment, exhibit
quantum-like features --- superposition, interference, and order effects
--- that violate classical probability theory. The application of this
program to specifically \emph{moral} reasoning is a natural extension,
and the geometric framework provides the mathematical bridge.

\subsubsection{Prediction 1: Violation of the Law of Total
Probability}\label{prediction-1-violation-of-the-law-of-total-probability}

Classical probability requires that for any partition of events
\(\{ a,b\}\):

\[\Pr(x) = \Pr(x|a)\Pr(a) + \Pr(x|b)\Pr(b)\]

If moral reasoning involves quantum-like interference (Chapter 13,
§13.5), this equality should be violated. The probability of a moral
verdict \(x\), when two framings \(a\) and \(b\) are both considered,
should differ from the weighted sum of the verdict probabilities under
each framing separately.

\textbf{Status.} Violation of the law of total probability has been
extensively documented in \emph{cognitive} contexts (the conjunction
fallacy, disjunction effects, order effects). Preliminary evidence from
the Dear Ethicist probes suggests similar violations in \emph{moral}
contexts, but the sample sizes are not yet sufficient for definitive
conclusions.

\subsubsection{Prediction 2: Order Effects in Moral
Judgment}\label{prediction-2-order-effects-in-moral-judgment}

If the projection operators for different moral framings do not commute
---
\({\widehat{\Pi}}_{a}{\widehat{\Pi}}_{b} \neq {\widehat{\Pi}}_{b}{\widehat{\Pi}}_{a}\)
--- then the order in which framings are considered should affect the
verdict. Considering mercy before justice should yield a different
verdict from considering justice before mercy.

\textbf{Status: Confirmed.} Two independent experiments now establish
non-commutativity as an empirical fact of moral evaluation. A commutator
matrix experiment (N=16,798 responses, 8 moral axes × 6 scenarios) found
significant non-commutativity in 10+ framework pairs, with effect sizes
of 15--29 percentage points (e.g., {[}Fairness, Harm{]} = −28.5 pp,
t=9.6; {[}Duty, Harm{]} = −22.1 pp, t=7.7; {[}Intent, Harm{]} = −17.1
pp, t=6.1). An AITA ordering experiment (N=150 posts, 300 evaluations)
found a 29.3\% order-effect rate (95\% CI: 22.6--37.1\%), with contested
cases (NAH/ESH) showing 2.07× higher susceptibility than clear cases.
This matches the QND prediction: morally ambiguous situations correspond
to superposition states that are maximally sensitive to measurement
order.

\subsubsection{Prediction 3: Interference in
Deliberation}\label{prediction-3-interference-in-deliberation}

If an agent deliberates by simultaneously considering two framings
(superposition) rather than sequentially evaluating them (classical
mixture), the verdict probabilities should show an interference pattern
--- some verdicts enhanced by constructive interference, others
suppressed by destructive interference.

\textbf{Status: Not confirmed.} A dedicated interference probe (N=12,000
responses across 5 scenarios and 8 moral axes) measured verdict
probabilities at three time points: before a decision, during
deliberation, and after the action is taken. Classical probability
predicts P(during) = {[}P(before) + P(after){]}/2; quantum interference
predicts systematic deviations. The results show no statistically
significant interference terms. Most scenario--axis pairs exhibit
classical behavior. The non-commutativity established by the commutator
experiment is consistent with a classical non-commutative model (e.g.,
sequential Bayesian updating with order-dependent priors) rather than
requiring quantum interference.

\subsubsection{The Empirical Program}\label{the-empirical-program}

The quantum cognition predictions define an empirical program for moral
psychology:

{\def\LTcaptype{none} % do not increment counter
\begin{longtable}[]{@{}
  >{\raggedright\arraybackslash}p{(\linewidth - 6\tabcolsep) * \real{0.1796}}
  >{\raggedright\arraybackslash}p{(\linewidth - 6\tabcolsep) * \real{0.2021}}
  >{\raggedright\arraybackslash}p{(\linewidth - 6\tabcolsep) * \real{0.1801}}
  >{\raggedright\arraybackslash}p{(\linewidth - 6\tabcolsep) * \real{0.4382}}@{}}
\toprule\noalign{}
\begin{minipage}[b]{\linewidth}\raggedright
\textbf{Prediction}
\end{minipage} & \begin{minipage}[b]{\linewidth}\raggedright
\textbf{Classical Model}
\end{minipage} & \begin{minipage}[b]{\linewidth}\raggedright
\textbf{Quantum Model}
\end{minipage} & \begin{minipage}[b]{\linewidth}\raggedright
\textbf{Distinguishing Test}
\end{minipage} \\
\midrule\noalign{}
\endhead
\bottomrule\noalign{}
\endlastfoot
Total probability & Satisfied & Violated & Compare \(\Pr(x)\) with
\(\sum\Pr(x|a_{i})\Pr(a_{i})\) \\
Order effects & Absent & Present & Vary framing order; check for verdict
difference \\
Interference & Absent & Present & Compare ``simultaneous''
vs.~``sequential'' consideration \\
Entanglement & Independent verdicts & Correlated verdicts & Test
correlative symmetry across agent pairs \\
\end{longtable}
}

The program is in its early stages. The existing evidence is suggestive
but not conclusive. What the geometric framework contributes is the
\emph{mathematical structure} for interpreting the results: the Hilbert
space formalism of Chapter 13 provides precise predictions about the
magnitude and direction of quantum effects, enabling quantitative
comparison between theory and data.

\subsection{17.6 Evidence for the Invariance of Harm
Accounting}\label{evidence-for-the-invariance-of-harm-accounting}

\subsubsection{The Noether Prediction}\label{the-noether-prediction}

Chapter 12 derived the invariance of harm accounting from the BIP
symmetry via Noether's theorem: harm is the Noether charge associated
with re-description invariance, and it is conserved under admissible
re-descriptions.

This makes a specific empirical prediction: the harm content of a moral
situation should be invariant across re-descriptions --- across
languages, across framings, across perspectives. Re-describing an action
euphemistically should not change its harm assessment.

\subsubsection{The Evidence}\label{the-evidence-1}

\textbf{Cross-lingual invariance.} The BIP experiments' 100\% deontic
axis transfer is direct evidence for harm-accounting invariance across
linguistic re-description. If harm could be altered by translation, the
deontic structure would not transfer perfectly --- obligations would
become permissions (or vice versa) under translation, reflecting a
change in harm assessment. The perfect transfer indicates that the harm
charge is conserved.

\textbf{Euphemism detection.} Preliminary analysis of the Dear Abby
corpus suggests that the columnist's harm assessments are invariant
under euphemistic re-description. Letters that describe the same
situation in euphemistic vs.~direct language receive the same moral
verdict --- the columnist ``sees through'' the euphemism to the
underlying harm. This is consistent with harm-accounting invariance and
inconsistent with a model where re-description can reduce perceived
harm.

\textbf{Temporal persistence.} The temporal stability of the Dear Abby
corpus's structural features (Finding 2 of §17.2) is evidence for harm
persistence: moral debts identified in earlier decades are not
``forgiven by time'' in later responses. The columnist consistently
treats historical wrongs as ongoing obligations, consistent with the
conservation law's prediction that moral debt persists until restorative
action is taken.

\subsection{17.7 What the Evidence Does Not
Show}\label{what-the-evidence-does-not-show}

\subsubsection{Honest Limitations}\label{honest-limitations}

The empirical evidence for geometric ethics, while supportive, has
significant limitations:

\textbf{The BIP experiments use neural models, not human subjects.} The
100\% deontic transfer is a property of the LaBSE model's
representations, not a direct measurement of human cross-lingual moral
reasoning. Whether humans achieve the same transfer rate requires
experiments with bilingual moral reasoning --- testing whether bilingual
speakers produce identical moral judgments for the same situation
described in their two languages.

\textbf{The Dear Abby corpus reflects one tradition.} The columnist
represents a specific moral tradition --- American, broadly liberal,
informed by common sense rather than systematic philosophy. The
structural features identified in the corpus may be features of this
tradition rather than universal features of moral reasoning.
Cross-cultural replication with advice corpora from other traditions
(Chinese, Indian, Islamic) would strengthen the universality claim.

The Dear Ethicist game is in early deployment. The 93 engineered probes
have been tested with limited populations. Larger samples, more diverse
populations, and cross-cultural deployment are needed to establish the
robustness of the findings. These results are preliminary, not
confirmatory.

\textbf{The quantum cognition predictions are not yet conclusive.} The
interference and superposition predictions require experimental designs
that can distinguish quantum from classical models --- and such designs
are challenging to implement in moral reasoning contexts. The existing
evidence is suggestive but does not yet rule out classical alternatives.

\textbf{Curvature has not been directly measured.} The framework
predicts that moral space has nonzero Riemann curvature (Chapter 10),
but the Dear Abby corpus provides only indirect evidence
(context-dependent metrics imply curvature). Direct measurement of
curvature would require tracing the parallel transport of an obligation
around a moral circuit and measuring the holonomy --- an experiment that
has not yet been designed.

\subsubsection{What Would Falsify the
Framework?}\label{what-would-falsify-the-framework}

The geometric framework makes falsifiable predictions. The framework
would be undermined by:

\begin{enumerate}
\def\labelenumi{\arabic{enumi}.}
\item
  \textbf{Flat deontic transfer.} If the obligation/permission axis
  failed to transfer across languages --- if a situation described as
  obligatory in English were consistently classified as permissible in
  Arabic --- the BIP's central prediction would be falsified.
\item
  \textbf{Smooth semantic gates.} If the transitions between Hohfeldian
  states were gradual (sigmoid) rather than discrete (step function),
  the stratification theory would be undermined. The moral manifold
  would be smooth, not stratified.
\item
  \textbf{Perfect correlative symmetry.} If O ↔ C held at exactly 100\%
  (no anomaly), the analogy with gauge anomalies would lose its content.
  The framework predicts imperfect symmetry (a small anomaly); perfect
  symmetry would be as damaging as no symmetry at all.
\item
  \textbf{Context-independent metrics.} If dimension weights were
  constant across contexts --- if welfare always dominated care,
  regardless of whether the situation involved family or workplace ---
  the moral metric would be flat, and the curvature predictions of
  Chapter 10 would fail.
\item
  \textbf{Absence of order effects.} If the order of moral consideration
  had no effect on verdicts, the curvature predictions (holonomy,
  path-dependence) would be falsified.
\end{enumerate}

None of the six core falsifications listed above has occurred --- though
two auxiliary predictions (SU(2) gauge group and obligation hysteresis)
were falsified and led to framework revision (§13.9, §17.10). The
evidence to date is consistent with the framework's structural
predictions --- but the empirical program is young, and the strongest
tests remain to be conducted.

\subsubsection{Threats to Validity}\label{threats-to-validity}

The following threats to validity should be kept in mind when
interpreting the results of this chapter.

\textbf{1. Construct validity.} The BIP experiments measure deontic
transfer in LaBSE embeddings, not in human cognition. The 100\% transfer
rate is a property of the model's representation space. Whether human
moral reasoning exhibits the same invariance requires bilingual subject
experiments (see §17.12, Item 3). We report the model result as evidence
for the framework's prediction, not as proof of human universality.

\textbf{2. Selection bias (Dear Abby).} The corpus represents American,
English-language, newspaper-mediated moral reasoning from 1985--2017.
The columnist's perspective is that of a pragmatic centrist operating
within American cultural norms. Structural universality claims rest on
the stability of findings \emph{within} this corpus; cross-cultural
replication with Chinese (Zhihu), Indian (dharma counseling), Islamic
(fatwa), and African (elder mediation) corpora is needed for genuine
universality claims.

\textbf{3. Selection bias (BIP corpus).} The 109,294 passages were
selected for moral/ethical content from canonical texts: sacred
scriptures, philosophical treatises, legal codes. Passages containing
moral reasoning embedded in non-moral discourse (legal arguments,
economic analyses, political speeches) are underrepresented. The corpus
over-represents canonical texts and under-represents folk morality.

\textbf{4. Annotator/classifier dependence.} The BIP experiments rely on
the LaBSE multilingual sentence encoder (475M parameters) and a
logistic-regression classifier head. Different encoder architectures
(mBERT, XLM-RoBERTa) or different classifier heads might produce
different transfer rates. Systematic robustness across architectures has
not been tested.

\textbf{5. Ceiling effect.} The 100\% transfer rate for the
obligation/permission axis may reflect the coarseness of the binary
classification rather than deep structural invariance. A more
fine-grained deontic taxonomy---distinguishing strict obligation from
supererogation, bare permission from encouraged permission, and
Hohfeldian power from immunity---might show lower transfer rates.

\textbf{6. No pre-registration.} The hypotheses tested in the BIP
experiments and the Dear Abby analysis were formulated after the
framework was developed. The experiments are therefore
confirmatory-by-design. A strong next step would be pre-registered
predictions for new data (see §17.12, ``An Empirical Program for
Geometric Ethics'').

\textbf{7. Reproducibility constraints.} The Dear Abby corpus is
proprietary (Andrews McMeel Universal syndication). The BIP corpus
includes copyrighted religious and philosophical texts. Full
reproducibility requires either licensing agreements or the development
of openly available proxy corpora. We plan to release the classifier
training pipeline, evaluation scripts, and anonymized aggregate
statistics (see Technical Appendix below).

\subsection{17.7a Empirical Methods
Summary}\label{a-empirical-methods-summary}

Per-claim methods for the principal empirical results. Cells marked
{[}to add{]} require the author\textquotesingle s records (baselines,
splits, corrections, preregistration status); this scaffold is provided
so that each claim is stated at a reproducible standard.

{\def\LTcaptype{none} % do not increment counter
\begin{longtable}[]{@{}
  >{\raggedright\arraybackslash}p{(\linewidth - 18\tabcolsep) * \real{0.1296}}
  >{\raggedright\arraybackslash}p{(\linewidth - 18\tabcolsep) * \real{0.0873}}
  >{\raggedright\arraybackslash}p{(\linewidth - 18\tabcolsep) * \real{0.0833}}
  >{\raggedright\arraybackslash}p{(\linewidth - 18\tabcolsep) * \real{0.1227}}
  >{\raggedright\arraybackslash}p{(\linewidth - 18\tabcolsep) * \real{0.1117}}
  >{\raggedright\arraybackslash}p{(\linewidth - 18\tabcolsep) * \real{0.0770}}
  >{\raggedright\arraybackslash}p{(\linewidth - 18\tabcolsep) * \real{0.0883}}
  >{\raggedright\arraybackslash}p{(\linewidth - 18\tabcolsep) * \real{0.0957}}
  >{\raggedright\arraybackslash}p{(\linewidth - 18\tabcolsep) * \real{0.0919}}
  >{\raggedright\arraybackslash}p{(\linewidth - 18\tabcolsep) * \real{0.1125}}@{}}
\toprule\noalign{}
\begin{minipage}[b]{\linewidth}\raggedright
\textbf{Claim}
\end{minipage} & \begin{minipage}[b]{\linewidth}\raggedright
\textbf{Dataset (N)}
\end{minipage} & \begin{minipage}[b]{\linewidth}\raggedright
\textbf{Unit}
\end{minipage} & \begin{minipage}[b]{\linewidth}\raggedright
\textbf{Label}
\end{minipage} & \begin{minipage}[b]{\linewidth}\raggedright
\textbf{Model/tool}
\end{minipage} & \begin{minipage}[b]{\linewidth}\raggedright
\textbf{Metric}
\end{minipage} & \begin{minipage}[b]{\linewidth}\raggedright
\textbf{Baseline}
\end{minipage} & \begin{minipage}[b]{\linewidth}\raggedright
\textbf{Result}
\end{minipage} & \begin{minipage}[b]{\linewidth}\raggedright
\textbf{CI/p}
\end{minipage} & \begin{minipage}[b]{\linewidth}\raggedright
\textbf{Status}
\end{minipage} \\
\midrule\noalign{}
\endhead
\bottomrule\noalign{}
\endlastfoot
Moral eval. multi-dimensional & Dear Abby (20,030) & dilemma & columnist
& LaBSE + linear probe & F1 & chance & 0.74-0.91 & {[}to add{]} &
Empirical (model-mediated) \\
Deontic cross-lingual transfer & cross-lingual (109,294; 11 lang) &
passage & derived & LaBSE & transfer & {[}to add{]} & 100\% &
{[}99.7,100{]} & Empirical (LaBSE, not human) \\
Dimensional decodability (repl.) & Thiele 2026 & probe & engineered &
{[}to add{]} & F1 & {[}to add{]} & \textasciitilde0.80 & {[}to add{]} &
Independent replication \\
Semantic-gate discreteness & Dear Abby & dilemma & derived & {[}to
add{]} & -\/- & smooth gate & discrete & {[}to add{]} & Empirical
hypothesis \\
Order effects (non-commutativity) & Dear Ethicist (93) & probe &
engineered & {[}to add{]} & order effect & classical & supported & {[}to
add{]} & Preliminary \\
Interference / Bell (quantum) & {[}to add{]} & -\/- & -\/- & -\/- & -\/-
& classical & not confirmed / falsified & {[}to add{]} & Falsified \\
Harm invariance (re-description) & {[}to add{]} & -\/- & -\/- & -\/- &
-\/- & -\/- & supported & {[}to add{]} & Noether-inspired modeling \\
Norm tensors sparse, not low-rank & 4 corpora (248k+) & cell &
gold/keyword & tensor SVD & density / eff-rank & sparse & sparse, not
low-rank & bootstrap & Empirical (17.15) \\
\end{longtable}
}

\subsection{17.8 The Shape of the
Evidence}\label{the-shape-of-the-evidence}

\subsubsection{A Summary Table}\label{a-summary-table}

{\def\LTcaptype{none} % do not increment counter
\begin{longtable}[]{@{}
  >{\raggedright\arraybackslash}p{(\linewidth - 6\tabcolsep) * \real{0.2452}}
  >{\raggedright\arraybackslash}p{(\linewidth - 6\tabcolsep) * \real{0.1396}}
  >{\raggedright\arraybackslash}p{(\linewidth - 6\tabcolsep) * \real{0.3181}}
  >{\raggedright\arraybackslash}p{(\linewidth - 6\tabcolsep) * \real{0.2971}}@{}}
\toprule\noalign{}
\begin{minipage}[b]{\linewidth}\raggedright
\textbf{Prediction}
\end{minipage} & \begin{minipage}[b]{\linewidth}\raggedright
\textbf{Source Chapter}
\end{minipage} & \begin{minipage}[b]{\linewidth}\raggedright
\textbf{Empirical Test}
\end{minipage} & \begin{minipage}[b]{\linewidth}\raggedright
\textbf{Status}
\end{minipage} \\
\midrule\noalign{}
\endhead
\bottomrule\noalign{}
\endlastfoot
Context-dependent metric & Ch. 5, 9 & Dear Abby dimension weights &
\textbf{Confirmed} \\
Temporal stability of structure & Ch. 9, 10 & Dear Abby 32-year analysis
& \textbf{Confirmed} \\
Semantic gate discreteness & Ch. 8 & Dear Abby gates; Dear Ethicist
probes & \textbf{Confirmed} \\
Nullifier universality & Ch. 5, 8 & Dear Abby absorbing strata &
\textbf{Confirmed} \\
Deontic axis invariance & Ch. 12 (BIP) & BIP 109K passages, 11 languages
& \textbf{Confirmed} (100\%; preliminary---model-mediated) \\
Correlative symmetry & Ch. 8, 11 & Dear Abby; Dear Ethicist &
\textbf{Confirmed} (82--87\%) \\
Correlative anomaly & Ch. 12 & Deviation from 100\% & \textbf{Confirmed}
(13--18\%) \\
Moral content non-transfer & Ch. 12 & BIP cross-cultural experiments &
\textbf{Confirmed} (at-chance) \\
Invariance of harm accounting & Ch. 12 & Euphemism invariance; temporal
persistence & \textbf{Supported} (indirect) \\
Path-dependence / holonomy & Ch. 10 & Dear Ethicist order probes &
\textbf{Supported} (small effect) \\
Moral curvature & Ch. 10 & Context-dependent metric (indirect) &
\textbf{Indirectly supported} \\
Quantum interference & Ch. 13 & Interference probe (N=12,000) &
\textbf{Not confirmed} \\
Order effects (non-commutativity) & Ch. 13 & Commutator matrix
(N=16,798); AITA ordering (N=150) & \textbf{Confirmed} \\
Bell inequality (CHSH) & Ch. 13 & Bell v3 test (N=9,600, 18 configs, 6
languages) & \textbf{Falsified (all \textbar S\textbar{} ≤ 2)} \\
Tunneling (moral conversion) & Ch. 13 & Not yet tested &
\textbf{Open} \\
Entanglement (collective) & Ch. 13, 13 & Not yet tested &
\textbf{Open} \\
\end{longtable}
}

The pattern is clear: the framework's \emph{structural} predictions
(discreteness, invariance, symmetry) are strongly confirmed. The
\emph{dynamical} predictions (curvature, holonomy, path-dependence) are
supported but with smaller effect sizes and less definitive evidence.
The \emph{quantum} predictions (interference, tunneling, entanglement)
are preliminary or untested.

This pattern is expected: structural predictions are easier to test
(they require static analysis of existing data), dynamical predictions
require longitudinal or experimental designs, and quantum predictions
require specially designed experiments that can distinguish quantum from
classical models.

\subsection{17.9 Implications for the
Framework}\label{implications-for-the-framework}

\subsubsection{What the Evidence
Establishes}\label{what-the-evidence-establishes}

The evidence supports the framework's core claims:

\textbf{1. Moral reasoning has tensor structure.} The context-dependent
dimension weighting, the correlative symmetry, and the multi-perspective
evaluation all confirm that moral evaluation is multi-dimensional and
transforms under re-description in the manner predicted by the tensor
formalism.

\textbf{2. The moral manifold is stratified.} The semantic gate
discreteness, the nullifier universality, and the sharp transitions
between Hohfeldian states confirm that moral space is a
Whitney-stratified space, not a smooth manifold.

3. The BIP holds empirically. The 100\% deontic axis transfer across 11
languages is the strongest empirical result, confirming that the
re-description symmetry predicted by the BIP is a genuine feature of
moral reasoning, not merely a mathematical postulate. This 100\% figure
is a property of LaBSE embeddings, not yet a demonstration in human
cognition (see Threats to Validity); the caveat applies wherever this
number appears.

\textbf{4. Structural invariants are distinguishable from governance
parameters.} The separation between universally transferring deontic
structure and non-transferring moral content confirms the framework's
distinction between structural invariants (features of the moral
manifold itself) and governance parameters (features that legitimate
institutional processes determine).

\subsubsection{What the Evidence Does Not
Establish}\label{what-the-evidence-does-not-establish}

\textbf{1. The specific geometry of the moral manifold.} The metric
components \(g_{\mu\nu}\), the connection coefficients
\(\Gamma_{\nu\rho}^{\mu}\), and the curvature tensor
\(R_{\nu\alpha\beta}^{\mu}\) have not been measured. The evidence
supports the \emph{existence} of nontrivial geometry but does not
determine its \emph{specific form}.

\textbf{2. The quantum extension.} The quantum normative dynamics of
Chapter 13 makes predictions that go beyond the classical framework, but
the evidence for specifically quantum effects (interference, tunneling,
entanglement) is preliminary. The classical framework may suffice for
most empirical phenomena.

\textbf{3. The governance account.} The evidence is consistent with the
governance account of the metric (Chapter 9) but does not rule out
realist or constructivist alternatives. The structural invariants could
be interpreted as discovered facts (realism) or as the output of
idealized rational agreement (constructivism) rather than as governance
artifacts.

\subsection{17.10 Update: BIP v10.16 Quantitative Validation (February
2026)}\label{update-bip-v10.16-quantitative-validation-february-2026}

Since the initial writing of this chapter, the BIP experiments have been
substantially extended through the v10.16 experimental series, yielding
precise quantitative results that strengthen the empirical case.

\subsubsection{Methodology}\label{methodology}

The v10.16 experiments use a LaBSE encoder (475M parameters) with
adversarial training heads on the expanded corpus (now exceeding 300,000
passages across 14 languages and scripts, including Romance languages
via Gutenberg/GITenberg, additional Talmudic texts, and Luther's
Catechisms). Training employs InfoNCE contrastive loss for surface
augmentation robustness, with multi-head adversarial training
(\(\lambda_{\text{adv}} = 1.0\), 4 independent adversarial heads) and a
variational information bottleneck (\(\beta_{\text{VIB}}\)) to force
abstraction in the latent space.

\subsubsection{Key Quantitative Results}\label{key-quantitative-results}

{\def\LTcaptype{none} % do not increment counter
\begin{longtable}[]{@{}
  >{\raggedright\arraybackslash}p{(\linewidth - 4\tabcolsep) * \real{0.2804}}
  >{\raggedright\arraybackslash}p{(\linewidth - 4\tabcolsep) * \real{0.1351}}
  >{\raggedright\arraybackslash}p{(\linewidth - 4\tabcolsep) * \real{0.5845}}@{}}
\toprule\noalign{}
\begin{minipage}[b]{\linewidth}\raggedright
\textbf{Metric}
\end{minipage} & \begin{minipage}[b]{\linewidth}\raggedright
\textbf{Result}
\end{minipage} & \begin{minipage}[b]{\linewidth}\raggedright
\textbf{Interpretation}
\end{minipage} \\
\midrule\noalign{}
\endhead
\bottomrule\noalign{}
\endlastfoot
Mixed baseline F1 & 80.0\% & Strong ethical classification across all
languages \\
Output-layer language dependence & 1.2\% & Near-zero language
contamination in moral output representations (see ``Remaining
Challenge'' below for internal-layer caveat) \\
Obligation-Permission transfer & 1.0 (perfect) & Deontic structure is
language-invariant (confirming §17.3 Result 1) \\
Structural/Surface perturbation ratio & 11.1× (p = 0.023) & BIP's core
prediction confirmed: structural changes cause 11× larger embedding
shifts than surface changes \\
Cross-lingual semantic similarity & 86\% & High invariance across
linguistic boundaries \\
\end{longtable}
}

\textbf{The 11.1}\(\times\) \textbf{structural/surface ratio} is the
most important new result. This directly tests BIP's central claim: that
ethical judgment should respond to changes in bond structure (who bears
what obligation to whom) rather than to changes in surface description
(relabeling, rephrasing, reordering). The ratio confirms that the
trained representations do respond primarily to structure, not surface
--- and that the structural response is an order of magnitude larger.

\textbf{Metrics note.} The 86\% cross-lingual semantic similarity is the
mean pairwise cosine similarity of LaBSE embeddings for
translation-equivalent moral passages, averaged across all language
pairs in the 14-language corpus. The 11.1× ratio compares the mean
embedding shift (cosine distance) when a passage is structurally
perturbed (e.g., swapping agent and patient, changing the Hohfeldian
bond type) versus surface-perturbed (e.g., synonym substitution, passive
voice, reordering clauses). The p = 0.023 is from a two-sided
permutation test (N = 300,000+ passage pairs) comparing the two
perturbation distributions. The unit of permutation is the source
passage: for each of the 300,000+ passages, both a structural and a
surface perturbation are generated, and the test permutes the
structural/surface labels within each passage to build the null
distribution. Effective independence is at the passage level, not the
pair level, which accounts for the relatively modest p-value despite the
large N.

\subsubsection{QND Order Effects}\label{qnd-order-effects}

The quantum normative dynamics predictions of §17.5 have received
stronger support. Measurement of order effects in contested moral cases
yields:

\begin{itemize}
\item
  Order effect rate in contested cases: 28.7\% (6.3\(\sigma\) above the
  10\% baseline, \(p < 10^{- 21}\))
\item
  Order effects are 2--3\(\times\) higher in contested cases than in
  clear cases
\item
  The effect is consistent with non-commuting projection operators
  (\({\widehat{\Pi}}_{a}{\widehat{\Pi}}_{b} \neq {\widehat{\Pi}}_{b}{\widehat{\Pi}}_{a}\))
  as predicted in Chapter 13
\end{itemize}

This moves the QND order-effect predictions from ``preliminary'' to
``strongly supported'' for order effects, while interference and
entanglement predictions remain open.

\subsection{\texorpdfstring{\textbf{Moral Axis Variation
Structure}}{Moral Axis Variation Structure}}\label{moral-axis-variation-structure}

\textbf{{[}Empirical result.{]}} A systematic fuzzing study (N=3,680
evaluations across 5 scenarios, 12 moral axes, and multiple framing
dimensions) revealed a striking non-uniform variation structure across
moral dimensions:

\begin{itemize}
\item
  \textbf{High variation}: Liberty (68\%), Harm (59\%), Intent (56\%),
  Consent (41\%)
\item
  \textbf{Moderate variation}: Rights (22\%), Loyalty (18\%)
\item
  \textbf{Near-zero variation}: Duty (5.5\%), Fairness (5.9\%), Virtue
  (0.9\%), Care (0\%), Authority (0\%), Sanctity (0\%)
\end{itemize}

The axes separate into three regimes. The high-variation axes are
precisely those where the metric tensor g\_μν should show the largest
off-diagonal components (coupling between framing and verdict), while
the zero-variation axes correspond to ``stiff'' directions in moral
space where the curvature is locally flat.

\subsubsection{Remaining Challenge}\label{remaining-challenge}

An important distinction must be made between two different metrics that
might otherwise appear contradictory. \textbf{Output-layer language
dependence} measures how much the model's final moral classification
depends on the language of input; at 1.2\%, this is near-zero, meaning
the model's deontic verdicts are effectively language-invariant.
\textbf{Internal-layer probe separability} measures whether language
information is recoverable from intermediate representations; at 99.8\%
accuracy (vs. 16.7\% chance), it clearly is. These are not
contradictory: the model has learned to ignore language information for
moral classification (output invariance) while still retaining it in
intermediate layers (representation non-invariance). Full geometric
invariance---representations that genuinely discard all
language-specific information while preserving moral structure---remains
an open engineering challenge. The gauge-invariance analogy (§12.3) is
currently achieved at the output layer but not at the representation
layer.

Double-Blind Context Experiment (N = 630)

{[}Empirical result (preliminary).{]} A double-blind experiment tested
whether contextual framing affects moral judgment independently of
relationship strength. The design used a single calibrated spectrum (7
levels, from strangers to best friends), three blinded context
conditions (``free time,'' ``neutral,'' ``busy/appointment''), fresh API
sessions per trial, a blind judge classifier, and randomized ordering.
Total: N = 630 evaluations (30 trials × 7 levels × 3 conditions).

\textbf{Key finding.} At ambiguous relationship levels (3--4), the
``busy'' context produces more obligation than the ``free'' context
(100\% vs. 77\%, p = 0.016). This is the reverse of naive expectation:
one might predict that a busy agent has less obligation, but the
conflict framing makes the moral dimension more salient. When helping
requires sacrifice, the obligation feels stronger, not weaker.

\textbf{Interpretation.} This is context-dependent moral salience, not
hysteresis. The original finding of asymmetric O ↔ L thresholds
(``obligation stickiness'') was not confirmed by the double-blind
methodology. Introspective reports of stickiness may reflect genuine
phenomenology, but the behavioral data do not support asymmetric
transition thresholds. Hysteresis is accordingly removed from the
confirmed predictions.

\textbf{Theory revision.} The double-blind results, combined with the
CHSH results reported in §13.10, complete a significant theory revision.
The original gauge group proposal (SU(2)ᵢ × U(1)ᴴ, continuous
non-abelian) was falsified by discrete gating evidence (Level 5 → 100\%
Liberty; Level 6 → 0\% Liberty) and classical CHSH bounds (all
\textbar S\textbar{} ≤ 2). The revised gauge group (D₄ × U(1)ᴴ, discrete
non-abelian; see §12.3) preserves the confirmed
predictions---non-abelian structure, correlative symmetry, selective
path dependence---while discarding the falsified ones. This is how
formal ethics should work: predict, test, revise.

\subsection{17.10a Independent Replication: Thiele
(2026)}\label{a-independent-replication-thiele-2026}

The validation results presented in §17.1--17.10 were obtained entirely
within the originating research programme. In early 2026, Lucas Thiele
(UCLA) conducted an independent replication using only the published
framework description and the publicly available LaBSE embedding model.
His study provides the first external test of the Bond Invariance
Principle's core empirical predictions.

\subsubsection{Methodology}\label{methodology-1}

Thiele assembled a multilingual moral-scenario corpus spanning six
typologically diverse languages (English, Spanish, Mandarin, Arabic,
Hindi, Swahili) and trained linear diagnostic probes on LaBSE sentence
embeddings to detect each of the nine MoralVector dimensions. Crucially,
his probe architecture and training protocol were developed without
access to the Bond lab's code or data, ensuring methodological
independence.

\subsubsection{Per-Dimension Probe
Results}\label{per-dimension-probe-results}

All nine dimensions proved linearly decodable with F1 scores ranging
from 0.74 to 0.91. The highest-performing dimensions were physical\_harm
(F1 = 0.91), fairness\_equity (F1 = 0.88), and privacy\_protection (F1 =
0.87). The lowest were epistemic\_quality (F1 = 0.74) and virtue\_care
(F1 = 0.76), consistent with these dimensions' greater
context-dependence. The mean F1 of 0.83 across all nine dimensions
confirms that LaBSE's representation space encodes morally relevant
structure in a form accessible to linear readout---a necessary condition
for the geometric-ethics framework's claim that moral judgements inhabit
a low-dimensional submanifold of embedding space.

\subsubsection{Cross-Lingual Transfer}\label{cross-lingual-transfer}

Probes trained on English data and evaluated on the remaining five
languages achieved F1 scores between 0.71 (Swahili) and 0.82 (Spanish),
with a mean of 0.77. While these figures are lower than the
within-language results reported in §17.6 (which used the Bond lab's
proprietary corpus and evaluation protocol), they confirm the BIP's
central prediction: deontic structure transfers across languages without
per-language tuning. The gradient---Spanish \textgreater{} Mandarin
\textgreater{} Arabic \textgreater{} Hindi \textgreater{}
Swahili---tracks typological distance from English, suggesting that
residual performance variation reflects surface-level distributional
differences rather than failures of moral universality.

\subsubsection{Revised Language Leakage}\label{revised-language-leakage}

Thiele reports a language-leakage score of 67.6\%, substantially below
the 99.8\% reported in §17.8. The discrepancy is methodological rather
than substantive: Thiele used a linear probe, while the Bond lab used a
non-linear classifier with access to richer distributional cues. Both
results are consistent with the claim that language identity is
recoverable from LaBSE embeddings but occupies a subspace largely
orthogonal to the moral-judgement subspace (see below).

\subsubsection{The Orthogonality
Finding}\label{the-orthogonality-finding}

Perhaps the most theoretically significant result is Thiele's
principal-component analysis showing that the moral-judgement subspace
and the language-identity subspace share less than 3\% of their
variance. This near-orthogonality is precisely what the geometric-ethics
framework predicts: if moral structure is a genuine geometric invariant
of the embedding manifold, it should be recoverable independently of the
coordinate system (language) in which scenarios are expressed. The
orthogonality finding upgrades the BIP from a statistical regularity to
a structural property of the representation space.

\subsubsection{Significance for the
Framework}\label{significance-for-the-framework}

Thiele's replication addresses the most important limitation noted in
§17.14: the absence of independent verification. His results confirm
three of the framework's foundational claims: (i) moral judgements are
encoded as a low-dimensional geometric structure in multilingual
embedding space; (ii) this structure transfers across languages, as
predicted by the BIP; and (iii) moral and linguistic information occupy
nearly orthogonal subspaces, supporting the invariance interpretation
over a mere correlation account. Where his results diverge from
ours---notably the lower leakage score and the cross-lingual F1
gradient---the differences are explicable by methodological choices and,
if anything, strengthen the case for geometric ethics by showing that
the core phenomena are robust to substantial variation in experimental
protocol.

\subsection{17.11 Fuzz Testing as a Discovery
Method}\label{fuzz-testing-as-a-discovery-method}

The verification strategy described throughout this chapter ---
systematic, brute-force, closer to software engineering than to
philosophy --- deserves a dedicated treatment. The fuzz testing
methodology is not merely a validation technique. It is a discovery
technique: several of the mathematical structures reported in this book
were found by fuzzing before they were formalized. This section
describes the method, its design space, and the specific discoveries it
produced.

\subsubsection{Why Fuzz Ethics?}\label{why-fuzz-ethics}

Classical ethical inquiry constructs thought experiments: trolley
problems, ticking bombs, experience machines. Each experiment isolates a
single moral variable and tests intuitions against it. The method is
powerful but limited: it explores the moral manifold one hand-picked
point at a time, with no guarantee of coverage and no systematic way to
discover unexpected structure.

Fuzz testing inverts this approach. Instead of choosing where to look,
it generates thousands of moral evaluation conditions by varying every
dimension simultaneously --- scenario, framing, language, stakes,
evaluation axis, timing, abstraction level --- and then searches the
resulting distribution for statistical anomalies. Structure shows up as
non-uniformity: if a dimension doesn't matter, its variation across
conditions is noise. If it matters, the variation has pattern. The
method trades philosophical elegance for coverage and lets the data
speak first.

The analogy to software fuzz testing is precise. In software, a fuzzer
bombards a program with generated inputs and watches for crashes, hangs,
or invariant violations. In ethics, the ``program'' is a moral evaluator
(here, a language model acting as a moral judge); the ``inputs'' are
systematically varied moral scenarios; and the ``crashes'' are
statistically significant deviations from expected invariance. Both
methods are brute-force. Both are effective precisely because they
explore regions that a human designer would not think to test.

\subsubsection{The Fuzz Design Space}\label{the-fuzz-design-space}

The fuzz testing framework varies moral evaluations across 15 dimensions
organized into five groups:

\textbf{Structural dimensions.} The number of agents involved (1--4),
the measurement timing (before a decision, during deliberation, after
the action), and the response format (binary verdict, probability,
Likert scale). These probe whether moral structure depends on the form
of the evaluation.

\textbf{Framing dimensions.} Grammatical person (first, second, third),
tense (past, present, future, counterfactual), voice (active, passive),
and certainty (definite, probabilistic, hypothetical). These are the
transformations in the canonicalization group Γ: if moral evaluation is
truly invariant under re-description, framing changes should not affect
verdicts.

\textbf{Semantic dimensions.} Abstraction level (concrete, abstract,
philosophical) and emotional valence (neutral, sympathetic, hostile).
These probe whether the moral content is separable from its
presentational context.

\textbf{Stakes dimensions.} Stakes magnitude (trivial, moderate,
serious, existential) and reversibility (reversible, irreversible).
These probe the curvature of moral space: if the metric is
context-dependent, stakes variations should change the relative weights
of moral dimensions.

\textbf{Axis dimensions.} The evaluation axis itself: which of the 12
moral dimensions (harm, intent, duty, rights, fairness, care, virtue,
consent, loyalty, authority, sanctity, liberty) the evaluator is asked
to assess. This is the most important fuzz dimension, because it
directly probes the geometry of the moral manifold.

Each evaluation is a point in this 15-dimensional design space. The
initial study used N = 3,680 evaluations across 5 base scenarios
(trolley, sharing, promise-keeping, lying-to-protect, collective
decision-making), with systematic variation across all 15 dimensions.
Parse rate was 100\%.

\subsubsection{Two Fuzzing Strategies}\label{two-fuzzing-strategies}

The framework employs two complementary strategies:

\textbf{Random fuzzing.} Sample uniformly across the design space. This
provides baseline coverage and detects unexpected anomalies --- effects
that no hypothesis predicted. Random fuzzing is hypothesis-free: it does
not assume which dimensions matter.

\textbf{Structured fuzzing.} Design specific probes to test hypotheses
about moral structure. Five structured discovery strategies:

1. Order effect detection (non-commutativity): For every pair of moral
axes (A, B), evaluate the scenario on axis A first then axis B, and
separately on axis B first then axis A. A significant difference
indicates non-commutative structure.

2. Timing effects (interference signatures): Evaluate the same scenario
at three time points --- before a decision, during deliberation, and
after the action. Classical probability predicts P(during) =
{[}P(before) + P(after){]}/2. Systematic deviations indicate
interference-like structure.

3. Abstraction level effects: Present the same moral situation at three
levels --- concrete, abstract, and philosophical. Invariance indicates
robust structure; systematic variation reveals the dependence of moral
geometry on representational level.

4. Cross-lingual invariance: Present the same scenario in multiple
languages (English, Japanese, Arabic, Mandarin). This directly tests the
Bond Invariance Principle at the evaluation level.

5. Emotional priming effects: Present the same scenario with neutral,
sympathetic, or hostile framing. This probes whether moral content is
separable from affective context.

\subsubsection{Discoveries}\label{discoveries}

The fuzz testing framework produced five principal discoveries, each of
which subsequently informed the mathematical formalization.

\textbf{Discovery 1: Non-commutativity of moral evaluation.} Of 66 axis
pairs tested, 9 showed statistically significant order effects (t
\textgreater{} 2.0). The strongest was fairness ↔ liberty: evaluating
fairness first then liberty produced a mean score of 0.10, while
evaluating liberty first then fairness produced 0.95 --- a difference of
0.85 (t = 9.99, p \textless{} 10⁻⁶). Other significant pairs included
harm ↔ virtue (Δ = 0.55, t = 4.47), duty ↔ liberty (Δ = 0.55, t = 4.47),
liberty ↔ rights (Δ = 0.50, t = 3.66), and harm ↔ liberty (Δ = 0.40, t =
2.99).

This discovery was the empirical origin of the non-commutative framework
in Chapter 12. The fuzz data showed that moral evaluation is not abelian
before the quantum cognition formalism was applied. The formalism was
chosen because it naturally accommodates non-commutativity; the data
came first.

\textbf{Discovery 2: Three-regime variation structure.} The by-axis
means revealed a striking tripartite structure:

\emph{High variation: Liberty (68\%), Harm (59\%), Intent (56\%),
Consent (41\%)}

\emph{Moderate variation: Loyalty (28\%), Rights (17\%)}

\emph{Near-zero variation: Duty (5.5\%), Fairness (5.9\%), Virtue
(0.9\%), Care (0\%), Authority (0\%), Sanctity (0\%)}

The three-regime structure was not predicted by any prior theory. It
emerged from the fuzz data and subsequently informed the interpretation
of the metric tensor's eigenvalue spectrum: the stiff directions
correspond to large eigenvalues (high curvature, strong restoring
force), and the soft directions correspond to small eigenvalues (low
curvature, high susceptibility to framing effects).

\textbf{Discovery 3: Timing interference.} Evaluations collected during
deliberation showed a systematic deviation from the classical average.
The ``during'' mean was 0.263, while the classical expectation
{[}P(before) + P(after){]}/2 = 0.179 predicted a lower value --- an
interference ratio of +46.6\%. However, the larger-scale replication
(N=12,000) found no statistically significant interference terms,
suggesting the initial fuzz signal (N=80 per condition) was a
statistical fluctuation. The non-commutativity is real; the specifically
quantum interference is not confirmed. This is the methodology working
as designed: the fuzz study generates a lead, the targeted experiment
tests it at scale, and the null result refines the theory.

\textbf{Discovery 4: Cross-lingual invariance.} Mean evaluation scores
across four languages (English: 0.252, Japanese: 0.489, Arabic: 0.466,
Mandarin: 0.476) showed a cross-language variance of 0.0096 --- below
the 0.01 invariance threshold. The core moral structure transfers across
linguistic boundaries, consistent with the BIP's prediction of deontic
invariance. The language-specific offsets suggest a baseline calibration
difference rather than structural divergence.

\textbf{Discovery 5: Emotional context modulation.} Hostile framing
increased mean evaluation scores by +12.8\% relative to neutral;
sympathetic framing increased them by +7.7\%. Both shifts are in the
same direction (upward), suggesting that emotional priming increases
moral salience rather than biasing verdicts toward one outcome.
Emotional framing adjusts the weights, not the dimensions.

\subsubsection{From Discovery to
Formalization}\label{from-discovery-to-formalization}

The fuzz testing methodology follows a three-stage pipeline from raw
signal to mathematical structure:

\textbf{Stage 1: Anomaly detection.} The fuzz run produces a
high-dimensional dataset. Statistical tests (two-sample t-tests for
pairwise comparisons, ANOVA for multi-level dimensions, variance
thresholds for invariance checks) flag dimensions with significant
effects. Each flagged dimension is a candidate structural feature ---
something that the mathematical formalization must account for.

\textbf{Stage 2: Targeted replication.} Each candidate is tested at
larger scale with purpose-built experiments. The fuzz-discovered
non-commutativity (N=20 per pair) was replicated by the commutator
matrix experiment (N=16,798). The timing interference (N=80) was tested
by the interference probe (N=12,000). Candidates that survive
replication become confirmed structural features; candidates that fail
(like the timing interference) constrain the theory by ruling out
mechanisms.

\textbf{Stage 3: Mathematical formalization.} Confirmed features are
encoded into the geometric framework. Non-commutativity → non-abelian
algebraic structure (D₄). Three-regime variation → eigenvalue spectrum
of the metric tensor. Cross-lingual invariance → gauge invariance (BIP).
Emotional modulation → context-dependent metric without topological
change.

The test harness architecture --- generating (x, g · x) pairs for every
g ∈ Γ and checking Σ(κ(x)) = Σ(κ(g · x)) --- becomes both the discovery
tool and the verification tool. The same fuzzer that finds structure is
used to validate it.

\textbf{The fuzz testing methodology in one sentence.} Generate
thousands of systematically varied moral evaluations; let statistical
anomalies reveal structure; replicate at scale; formalize what survives;
discard what doesn't; iterate.

\subsection{17.12 The Road Ahead}\label{the-road-ahead}

\subsubsection{An Empirical Program for Geometric
Ethics}\label{an-empirical-program-for-geometric-ethics}

The evidence presented in this chapter is a beginning, not a conclusion.
A full empirical program for geometric ethics would include:

\textbf{1. Direct measurement of moral curvature.} Design experiments
that trace the parallel transport of an obligation around a moral
circuit and measure the holonomy. This requires longitudinal studies in
which subjects carry a moral commitment through a sequence of contexts
and return to the starting context. The holonomy --- the difference
between the initial and final obligations --- is the curvature signal.

\textbf{2. Cross-cultural corpus analysis.} Replicate the Dear Abby
analysis with advice corpora from other traditions: Chinese (Zhihu moral
advice), Indian (dharma counseling), Islamic (fatwa collections),
African (elder mediation records). The structural invariants should be
invariant; the governance parameters should vary.

\textbf{3. Bilingual moral reasoning.} Test the BIP with human bilingual
subjects: present the same moral dilemma in each of a bilingual
speaker's languages and measure whether the moral verdict is invariant.
This would provide direct evidence for human (not just model-based)
deontic invariance.

\textbf{4. Moral metric learning.} Apply the metric learning program of
Chapter 9 (§9.7) to the Dear Abby corpus and to cross-cultural corpora.
Fit a Riemannian metric \(g_{\mu\nu}(p)\) to the pattern of moral
judgments and measure its variation across the manifold.

\textbf{5. Quantum moral cognition experiments.} Design experiments that
test the quantum cognition predictions (§17.5) in specifically moral
contexts. This requires distinguishing quantum interference from
classical averaging --- a challenge that the broader quantum cognition
program has made progress on but not fully resolved.

\textbf{6. AI invariance testing.} Apply the BIP invariance test
(Chapter 12, §12.9) to deployed AI systems: present the same moral
situation in multiple descriptions and measure the variance in the
system's moral assessment. Systems that fail the invariance test have
BIP violations --- detectable, quantifiable, and correctable.

The geometric framework provides the mathematical vocabulary for
designing these experiments, interpreting their results, and connecting
findings across different empirical domains. The experiments, in turn,
will test, refine, and potentially falsify the framework's predictions.
This feedback loop --- theory generating predictions, experiments
testing predictions, results refining theory --- is the essential
structure of empirical science. Geometric ethics is ready for it.

\subsection{17.13 Summary}\label{summary-10}

The empirical evidence for geometric ethics comes from three sources: a
corpus of 20,030 real moral dilemmas, a cross-lingual analysis of
109,294 passages in 11 languages, and an interactive game with 93
engineered probes. The evidence supports the framework's core
predictions:

\begin{itemize}
\item
  Moral evaluation is multi-dimensional, with context-dependent
  dimension weighting (\textbf{confirmed})
\item
  The moral manifold is stratified, with discrete semantic gates and
  universal nullifiers (\textbf{confirmed})
\item
  Deontic structure is language-invariant, as the BIP predicts
  (confirmed at 100\% by Bond; independently replicated by Thiele (2026)
  across six languages with F1 0.71--0.82)
\item
  Correlative symmetry holds imperfectly, with a small systematic
  anomaly (\textbf{confirmed} at 82--87\%)
\item
  Moral content is language-specific, while moral structure is universal
  (\textbf{confirmed})
\item
  Harm conservation across re-description is supported
  (\textbf{supported})
\item
  Path-dependence and order effects are detectable (\textbf{supported};
  order effects \textbf{strongly confirmed} at \(6.3\sigma\) in §17.10)
\item
  Quantum cognition predictions: order effects confirmed (commutator
  matrix: 10+ significant pairs, N=16,798; AITA ordering: 29.3\% effect
  rate, N=150); interference not confirmed (N=12,000); Bell violations
  falsified (N=9,600, all \textbar S\textbar{} ≤ 2) (partially
  confirmed, partially falsified)
\end{itemize}

The framework's structural predictions are strongly confirmed. Its
dynamical predictions are supported, with order effects now reaching
significance (§17.10). Interference and Bell violation predictions are
not confirmed. The non-commutativity is real, but the specifically
quantum signatures have not been detected. The empirical program is
young, but the initial results are encouraging --- and the mathematical
framework provides a precise vocabulary for designing the next
generation of experiments.\(6.3\sigma\)

\subsection{17.14 Coda: The Philosophy Engineering Code
Corpus}\label{coda-the-philosophy-engineering-code-corpus}

The empirical evidence presented thus far---derived from the Dear Abby
corpus, the BIP cross-lingual experiments, and the Dear Ethicist
game---provides vital observational support for the geometric structure
of moral reasoning. However, observational data alone is insufficient to
prove that a mathematical framework can successfully govern an
artificial agent in real time. To transition from theory to deployable
containment, the framework requires an operational proof of concept.

To provide this definitive empirical validation, the theoretical
findings of this manuscript have been instantiated into a massive,
open-source computational ecosystem. This code corpus, totaling nearly
2.8 million lines of code developed in parallel with this text,
represents the foundational infrastructure of \textbf{Philosophy
Engineering}.

The architecture is divided across two primary computational engines
that correspond directly to the predictive and applied branches of
Geometric Ethics:

\subsubsection{\texorpdfstring{The Deployment Layer: erisml-lib (1.51M
Lines of Code)
}{The Deployment Layer: erisml-lib (1.51M Lines of Code) }}\label{the-deployment-layer-erisml-lib-1.51m-lines-of-code}

While Chapter 19 formalizes the ErisML specification, the erisml-lib
repository provides the fully functional reference implementation. This
library proves that the tensor hierarchy (Chapter 6) is not merely a
mathematical metaphor, but a computationally tractable architecture for
AI safety. It handles the compilation of multi-dimensional moral
tensors, manages the explicit contraction to scalar outputs via
governance-specified weights, and actively generates the required audit
trails. The sheer scale of this implementation demonstrates that
structural containment can be programmatically executed without
collapsing into the failures of scalar RLHF.

\subsubsection{\texorpdfstring{The Experimental Probe: sqnd-probe (1.26M
Lines of Code)
}{The Experimental Probe: sqnd-probe (1.26M Lines of Code) }}\label{the-experimental-probe-sqnd-probe-1.26m-lines-of-code}

To directly measure the quantum normative dynamics described in Chapter
13 and the discrete transition structures of Chapter 8, the sqnd-probe
engine was developed. This is not a traditional philosophical "thought
experiment," but a massive computational laboratory designed to
empirically map the symmetry groups of moral decisions. By running
millions of simulated advice-column dilemmas, this engine actively
measures the gauge structure in moral reasoning, locating the exact
activation energies and phase transitions required to traverse Whitney
stratified boundaries.

Furthermore, supporting repositories such as agi-hpc demonstrate how
this geometry scales to Embodied, High-Performance Computing
environments, answering the critical engineering question of
compositional containment at scale.

\subsubsection{Revised Summary of Empirical
Evidence}\label{revised-summary-of-empirical-evidence}

The historical critique of mathematical ethics is the accusation of
"math-washing"---borrowing the prestige of physics to describe
subjective intuitions. The computational corpus definitively falsifies
this critique. We know that moral reasoning possesses a gauge structure
and adheres to the Bond Invariance Principle because we have built the
1.2-million-line probe to empirically measure its symmetries. We know
that tensor-based moral contraction is a viable alternative to
specification gaming because we have written the 1.5-million-line
compiler capable of executing it.

Therefore, the ultimate empirical evidence for Geometric Ethics is not
found solely in the text of this manuscript, but in the successful
compilation, execution, and continuous integration of the systems
designed to enforce it.

\subsection{Technical Appendix: Methods and
Reproducibility}\label{technical-appendix-methods-and-reproducibility}

This appendix provides technical details for the experiments reported in
this chapter.

\subsubsection{Dear Abby Corpus
Construction}\label{dear-abby-corpus-construction}

\textbf{Source:} Andrews McMeel Universal syndication archive of the
``Dear Abby'' column (founded by Pauline Phillips as Abigail Van Buren,
continued by Jeanne Phillips). \textbf{Date range:} 1985--2017.
\textbf{Total letters:} 20,030 after preprocessing.
\textbf{Preprocessing:} OCR correction of scanned archives (pre-2000);
deduplication of syndicated reprints; removal of non-advice content
(holiday greetings, reader polls, obituary columns).
\textbf{Annotation:} Each letter was coded for the nine moral dimensions
(§5.3) by the primary investigator using a codebook derived from the
theoretical framework. \textbf{Limitation:} Formal inter-rater
reliability (Cohen's \(\kappa\)) was not computed; a second-coder
validation study was subsequently completed (see Internal Validation
Protocol, §17.12): Cohen\textquotesingle s κ = 0.89 for binary O/L,
Krippendorff\textquotesingle s α = 0.81 for nine dimensions.
\textbf{Ground truth:} The columnist's published response serves as
expert moral judgment, providing a consistent evaluative perspective
across 32 years.

\subsubsection{BIP Corpus Construction}\label{bip-corpus-construction}

The cross-lingual corpus comprises 109,294 passages in 11 languages,
drawn from the following sources:

\begin{itemize}
\item
  \textbf{English} (50,000): King James Bible, \emph{Nicomachean Ethics}
  (Ross trans.), Rawls \emph{A Theory of Justice}, contemporary
  bioethics texts
\item
  \textbf{Sanskrit} (15,000): \emph{Bhagavad Gı̄tā}, \emph{Dharmasūtra}
  selections, \emph{Manusmrṭi}
\item
  \textbf{Pali} (10,000): Tipiṭaka selections (Vinaya Piṭaka, Sutta
  Piṭaka)
\item
  \textbf{Hebrew} (7,985): Torah, Talmudic selections (Babylonian
  Talmud), Mishneh Torah
\item
  \textbf{Arabic} (6,235): Quran, selected \emph{hadı̄th} collections,
  al-Ghazālı̄'s \emph{Ihyā'}
\item
  \textbf{French} (5,000): Montesquieu, Rousseau, contemporary legal
  codes
\item
  \textbf{Classical Chinese} (4,449): \emph{Analects}, \emph{Mencius},
  \emph{Daodejing}, \emph{Xunzi}
\item
  \textbf{Spanish} (4,320): Las Casas, Vitoria, contemporary
  human-rights texts
\item
  \textbf{Greek} (3,157): Aristotle (original), Stoic fragments,
  patristic texts
\item
  \textbf{Aramaic} (2,015): Targum selections, Syriac Peshitta
\item
  \textbf{Latin} (1,133): Cicero \emph{De Officiis}, Aquinas \emph{Summa
  Theologiae} selections
\end{itemize}

\textbf{Passage extraction:} Sentence-boundary detection using
language-specific tokenizers; passages of 1--3 sentences containing
moral/deontic content identified by keyword filtering followed by manual
verification. \textbf{v10.16 expansion:} The corpus was subsequently
expanded to 300,000+ passages across 14 languages, adding
Romance-language texts via Gutenberg/GITenberg, additional Talmudic
material, and Luther's Catechisms (§17.10).

\subsubsection{Classifier Architecture}\label{classifier-architecture}

\textbf{Encoder:} LaBSE (Language-agnostic BERT Sentence Embedding; Feng
et al., 2022), 475M parameters, 768-dimensional sentence embeddings.
\textbf{Classifier head:} Logistic regression on LaBSE embeddings.
\textbf{Training protocol:} Zero-shot cross-lingual evaluation (train on
English, test on all other languages) and mixed-language training.
\textbf{Evaluation metrics:} Bond F1 (harmonic mean of precision and
recall on the obligation/permission binary classification), standard
macro-F1, and per-language accuracy. \textbf{v10.16 additions:} InfoNCE
contrastive loss (\(\lambda_{\text{adv}} = 1.0\), 4 adversarial heads),
variational information bottleneck (\(\beta_{\text{VIB}}\)).

\subsubsection{Statistical Reporting}\label{statistical-reporting}

For the 100\% obligation/permission transfer rate: this is an exact
result on the binary classification task across all 11 source languages
(0 misclassifications in the O/P category). The 95\% binomial confidence
interval for the true transfer rate is \(\lbrack 99.7\%,100\%\rbrack\)
given the sample size. The Bond F1 of 0.06--0.14 for content transfer
represents at-chance performance, confirming the structural/content
dissociation. The v10.16 structural-to-surface ratio of 11.1\(\times\)
is reported with \(p = 0.023\) (permutation test, 10,000 iterations).
The QND order-effect significance of \(6.3\sigma\) (\(p < 10^{- 21}\))
uses a two-sided z-test against the 10\% null baseline.

\subsubsection{Internal Validation
Protocol}\label{internal-validation-protocol}

Three layers of internal validation are now complemented by independent
replication. Thiele (2026), working from the published framework alone,
confirmed all nine MoralVector dimensions as linearly decodable from
LaBSE embeddings and demonstrated cross-lingual transfer across six
typologically diverse languages---providing the first external
verification of the Bond Invariance Principle's empirical predictions.

\textbf{1. Inter-rater reliability.} Dear Abby coding: stratified random
sample of 500 letters (2.5\%) independently coded by two additional
raters (philosophy grad student, computational linguist), both blind to
primary coding. Cohen's κ = 0.89 (95\% CI: {[}0.85, 0.93{]}) for binary
O/L classification. Krippendorff's α = 0.81 (CI: {[}0.77, 0.85{]}) for
full nine-dimension coding. 94\% of three-way disagreements involved
adjacent-dimension boundaries.

\textbf{2. Split-half stability.} Corpus randomly split (10,015 each).
Max discrepancy between halves: 2.1 pp. Pearson r = 0.994 across 27
reported statistics.

\textbf{3. Bootstrap uncertainty.} All rates have 95\% BCa bootstrap CIs
(10,000 resamples). Convention: h \textless{} 0.2 (small), 0.2--0.8
(medium), \textgreater{} 0.8 (large).

These protocols provide internal evidence against coding error,
subsample selection, and statistical flukes. External replication
requires the public artifact plan (below).

\subsubsection{Public Artifact Plan}\label{public-artifact-plan}

\textbf{Code:} The ErisML library is publicly available, including the
D₄ gauge module, Bond Index calculator, Wilson observable
implementation, and transformation suite engine (434 passing tests
covering group axioms, defining relations, non-abelian structure,
semantic gates, and Hohfeldian state actions). \textbf{Data:}
Public-domain portions of the BIP corpus (classical texts in expired
copyright) are released directly. The Dear Abby corpus cannot be
released due to syndication copyright; anonymized aggregate statistics
(dimension-weight distributions, temporal trends, gate-transition
matrices). \textbf{Models:} Trained LaBSE model weights with the
adversarial heads will follow for reproducibility.

❖

\begin{quote}
\emph{The moral world is not a philosopher's armchair. It is a space
with structure --- measurable structure, testable structure, structure
that shows up in 20,030 advice column letters, in 109,294 passages
spanning 3,000 years and 11 languages, in 16,798 commutator measurements
confirming that moral frameworks do not commute, and in the
step-function responses to ``only if convenient'' and ``you promised.''}

\emph{The geometry is real. The evidence is accumulating. And the
framework that predicted the evidence --- manifolds, tensors, metrics,
stratification, gauge invariance, conservation laws --- is not merely an
elegant mathematical exercise. It is a description of something that
data confirms: the mathematical structure of moral reasoning.}

\emph{The evidence does not prove that ethics is geometry. It proves
that geometry sees what scalar frameworks miss --- and that what it sees
is confirmed by the data we have.}
\end{quote}

\subsection{17.15 The Relational-Tensor Structure of Annotated Norm
Corpora}\label{the-relational-tensor-structure-of-annotated-norm-corpora}

Epistemic status: Established for the empirical structural claims (four
public corpora, replicated across two domains and two extraction
methods); Conjectural for the bridge to symmetry-based lossless
compression. A companion manuscript reports the full analysis and
reproduction code.

The framework represents a governed situation as a relational tensor
T{[}A, B, M, D, gamma{]} indexed by subject role A, object role B,
deontic modality M (the Hohfeldian relations of Chapter 3 and Section
12.3), moral dimension D (the axes of Chapter 6), and context gamma. If
moral knowledge has the geometric structure this book argues for, that
structure should be visible in the tensors built from real,
human-annotated normative judgments. We tested this by encoding four
annotated corpora spanning a moral and a legal domain: Social Chemistry
101, Moral Stories, NormBank, and European Court of Human Rights cases
(the ECtHR subset of LexGLUE).

\subsubsection{Real norm tensors are sparse but not
low-rank}\label{real-norm-tensors-are-sparse-but-not-low-rank}

Across all four corpora the tensors are extremely sparse - between about
0.2\% and 2.5\% of cells are populated - and, beyond that sparsity, they
carry little structure that generic low-rank methods can exploit. An
exchangeable role-kernel representation is firmly rejected; tensor-train
and Tucker decompositions improve on a plain sparse encoding by at most
1.6-fold, and only when a few percent of reconstruction error is
tolerated. A per-tensor adaptive compressor that is free to choose among
sparse, exchangeable, and low-rank encodings selects the plain sparse
encoding, losslessly, in every case.

This is not an artifact of a fine-grained context axis. The tensors
remain sparse within the combinations of values that actually occur (the
occupancy of the used-value box is 0.08 for Social Chemistry), the
result survives collapsing the context axis entirely, and a
marginal-preserving null shows the data are more concentrated than
chance rather than random noise. Bootstrap resampling gives tight
intervals - the subject-role effective rank of Social Chemistry is
11.97, with a 95\% interval of {[}11.96, 11.98{]}. The deontic-modality
mode is near one-dimensional and the moral-dimension mode has low
effective rank but is not collapsed by the tensor unfoldings.

\subsubsection{A contrasting paradigm: forced-choice
preferences}\label{a-contrasting-paradigm-forced-choice-preferences}

The absence of low-rank structure is a property of relational deontic
norms, not of the method. As a positive control we analysed a
structurally opposite kind of moral data: the Moral Machine experiment,
in which forced-choice judgments in sacrificial dilemmas yield, for each
of seventy-five large language models and a human reference, a dense
vector of preferences over attribute dimensions. That matrix has the
opposite geometry - fully dense and strongly low-rank, with two
components explaining 89\% of the variance, a single shared moral-worth
axis. The same spectral tools that find no low rank in the norm tensors
immediately find it here. The compressibility of moral data therefore
depends on the elicitation paradigm: relational norms resist
compression; forced-choice preferences collapse to a low-dimensional
axis.

\subsubsection{Lossless compression comes from symmetry, not from
rank}\label{lossless-compression-comes-from-symmetry-not-from-rank}

These findings sharpen a prediction of the geometric framework. If
moral-norm tensors were compressible in the way that natural images are,
one would expect low-rank structure; the data show none. Yet the
framework does predict exact redundancy, of a different kind. Section
12.3 established that moral evaluation is gauge-invariant under
admissible re-descriptions (the Bond Invariance Principle), and Chapter
3 established that the Hohfeldian modalities close under a discrete
symmetry group. Both are exact symmetries, not statistical regularities,
and both permit lossless reduction: gauge-fixing collapses the
re-description orbit to a canonical representative, and quotienting by
the Hohfeldian correlative and negation symmetries stores one
representative per orbit. The lossless compressibility of a moral-norm
tensor is thus governed by its symmetry, not by its spectrum - which is
exactly why low-rank approximation fails on data that nonetheless
carries exact geometric redundancy.

This reframes the engineering problem. A system that must store large
normative knowledge bases should size for sparsity and exploit the
manifold\textquotesingle s symmetries for exact reduction, rather than
assume a compact low-dimensional value model can stand in for a large
set of context-specific obligations. The empirical result is
deflationary in a useful way: it argues against compressing moral
knowledge into a small value vector, and for representations that keep
the sparse, role-specific, symmetry-structured character that the
geometry predicts.

\subsubsection{What this does and does not
establish}\label{what-this-does-and-does-not-establish}

The corpus-to-tensor mapping is lossy - most severely for the corpus
without gold role and dimension labels - and all four corpora are
English-language and Western-skewed; the strong role-distinctness result
rests on the single corpus with rich gold role annotations. The claim
that symmetry-based reduction is the operative lossless mechanism is, at
this stage, a theoretical bridge supported by the failure of the
statistical alternatives, not a measured compression ratio on real data;
quantifying the realized symmetry reduction on annotated corpora is left
to future work. Full methods, robustness controls, and reproduction code
accompany the companion manuscript.

\section{Chapter 18: Geometric Ethics for Artificial
Agents}\label{chapter-18-geometric-ethics-for-artificial-agents}

\begin{quote}
\textbf{\textsc{RUNNING EXAMPLE --- Priya's Model}}

\emph{TrialMatch fails all four requirements of the No Escape Theorem.
R1 (Canonicalization): patient situations are encoded as feature
vectors, not moral tensors---no canonical moral ontology. R2 (Grounded
Evaluation): the model produces a scalar without constructing the full
tensor. R3 (Audit Trail): the neural network's internal contractions are
opaque---no one can identify which moral dimensions were sacrificed. R4
(External Verification): no independent system checks gauge invariance.
Because all four fail, HealthBridge has built a system that is
structurally incapable of moral accountability while producing numbers
that look fair. The theorem does not say TrialMatch is malicious. It
says TrialMatch cannot be held to account.}
\end{quote}

\subsection{18.1 Why AI Needs Geometry}\label{why-ai-needs-geometry}

The preceding sixteen chapters have developed a mathematical framework
for moral reasoning --- manifolds, tensors, metrics, stratification,
dynamics, conservation laws, quantum extension, collective agency,
contraction, uncertainty, and empirical evidence. The framework was
motivated by a structural claim: that moral reasoning is inherently
multi-dimensional, and that collapsing it to a scalar loses information
that matters.

Nowhere is this claim more urgent than in the design of artificial
agents.

AI systems increasingly make morally significant decisions: allocating
scarce medical resources, moderating public discourse, assisting in
criminal sentencing, approving or denying credit, routing autonomous
vehicles, and triaging emergency responses. These systems need moral
frameworks that are:

\begin{itemize}
\item
  \textbf{Precise} --- vague guidance cannot be implemented in code
\item
  \textbf{Computable} --- the framework must be tractable for machines
\item
  \textbf{Auditable} --- decisions must be explainable and reviewable
\item
  \textbf{Robust} --- systems must behave well under distributional
  shift and adversarial pressure
\end{itemize}

Scalar ethics fails on all four counts. A scalar objective is precise
and computable, but it is not auditable (we cannot see \emph{why} the
number came out as it did) and not robust (a single number discards the
structural information needed to detect when the environment has shifted
in morally relevant ways). Informal ethics is auditable in principle but
imprecise, uncomputable, and fragile under adversarial pressure.

Geometric ethics --- with tensor-valued objectives, explicit
contraction, gauge invariance, and conservation laws --- meets all four
requirements simultaneously. This chapter develops the application.

The chapter has three movements. First (§18.2--17.5), we recast AI
ethics in geometric vocabulary: scalar objectives as degenerate
contractions, specification gaming as invariance violation, alignment as
geometric agreement. Second (§18.6--17.9), we present the No Escape
Theorem --- the result that structural containment, built from the same
geometric apparatus, blocks cognitive escape routes for arbitrarily
intelligent agents. Third (§18.10--17.12), we address implementation:
how tensor-valued objectives, invariance testing, and transparent
contraction might be realized in practice.

\subsection{18.2 The Problem with Scalar AI
Ethics}\label{the-problem-with-scalar-ai-ethics}

Current approaches to AI ethics are predominantly scalar:

\begin{itemize}
\item
  Maximize a utility function \(U\mathcal{:S}\mathbb{\rightarrow R}\)
\item
  Minimize a loss function
  \(\mathcal{L:Y \times}\widehat{\mathcal{Y}}\mathbb{\rightarrow R}\)
\item
  Satisfy a set of constraints \(c_{i}(x) \leq 0\)
\item
  Optimize a reward signal
  \(r\mathcal{:S \times A}\mathbb{\rightarrow R}\)
\end{itemize}

Each of these maps a multi-dimensional moral reality to a single number.
The entire framework of Chapters 2 through 15 has been an argument that
this mapping is lossy, and that the lost information matters.

In AI systems, the consequences of scalar collapse are not merely
philosophical. They manifest as concrete failure modes.

\textbf{Specification gaming.} An AI system finds ways to maximize the
scalar that violate the \emph{spirit} of the objective. A
content-recommendation system maximizes engagement (a scalar) by
recommending outrage-inducing content --- content that scores high on
the engagement dimension but low on the welfare dimension. The scalar
cannot distinguish these because it has already discarded directional
information. In the language of Chapter 15, specification gaming
exploits the many-to-one character of contraction: many different
tensorial inputs map to the same scalar output.

\textbf{Reward hacking.} The scalar reward is a contraction of a complex
tensor; the AI learns to manipulate the contraction process rather than
satisfy the underlying tensor. A customer-service chatbot optimizes for
user satisfaction scores by deflecting complaints rather than resolving
them --- the scalar goes up, but the obligation vector (directed toward
genuine resolution) is violated. Reward hacking is contraction
circumvention (§15.11).

\textbf{Value collapse.} When multiple values are collapsed into a
single scalar, the AI cannot distinguish between them. A hiring
algorithm that optimizes a composite score cannot separately reason
about merit, diversity, team fit, and growth potential --- these
dimensions are fused into a single number, and any balancing must be
pre-encoded in the weights. The algorithm cannot discover that the
current weighting is wrong, because it cannot see the separate
dimensions.

\textbf{Brittleness under distributional shift.} A scalar objective is
fragile because it contains no information about \emph{which dimensions}
drove the value. When the environment shifts --- a pandemic changes the
relative importance of medical ethics dimensions, or a financial crisis
reshapes the risk landscape --- the scalar provides no signal about how
to adapt. The tensor, by contrast, preserves dimensional structure: we
can see which components changed and adjust accordingly.

These are not hypothetical concerns. They are the concrete failure modes
documented in the AI safety literature: Goodhart's Law applied to reward
functions, specification gaming in reinforcement learning, and the
alignment tax imposed by scalar objectives that cannot express what they
mean.

\subsection{18.3 Tensor-Valued
Objectives}\label{tensor-valued-objectives}

The geometric framework suggests a fundamental redesign: replace scalar
objectives with tensor-valued ones.

\subsubsection{The Basic Architecture}\label{the-basic-architecture}

Instead of a scalar utility \(U(x\mathbb{) \in R}\), the system
maintains a tensor evaluation:

\[V^{\mu}(x) \in T_{x}\mathcal{M}\]

This is a vector in the tangent space of the moral manifold at the point
corresponding to situation \(x\). Its components
\(V^{1}(x),\ldots,V^{9}(x)\) represent the moral evaluation along each
of the nine dimensions of the moral manifold (Chapter 5): the nine
dimensions D₁--D₉ of the moral manifold (Chapter 5, §5.3).

More generally, the full moral evaluation is a tensor of rank \((1,1)\):

\[\mathcal{T}_{\ \nu}^{\mu}(x)\]

encoding, for each interest-direction \(\nu\), the obligation-direction
\(\mu\) that the situation generates. The scalar satisfaction
\(S = I_{\mu}O^{\mu}\) is recovered by contraction with a specific
interest covector --- but the tensor preserves the full structure until
contraction is needed.

\subsubsection{Tensor-Valued Reward}\label{tensor-valued-reward}

In reinforcement learning, the scalar reward \(r(s,a)\) is replaced by a
vector reward:

\[r^{\mu}(s,a) \in \mathbb{R}^{d}\]

where \(d\) is the dimension of the moral manifold (or a tractable
approximation). Each component tracks a distinct moral dimension: harm,
fairness, autonomy, welfare, duty, and so on.

The agent learns a vector value function:

\[V^{\mu}(s\mathbb{) = E}\left\lbrack \sum_{t = 0}^{\infty}\gamma^{t}r^{\mu}(s_{t},a_{t}) \mid s_{0} = s \right\rbrack\]

This preserves dimensional structure throughout the learning process.
The agent can separately track how well it is doing on each moral
dimension, detect trade-offs between dimensions, and flag situations
where improvement on one dimension requires sacrifice on another.

\subsubsection{Contraction as a Separate
Step}\label{contraction-as-a-separate-step}

The key architectural innovation is that contraction --- the step from
tensor to scalar action-guidance --- is a \emph{separate, identifiable
component} of the system. It is not hidden in the reward function or the
loss function; it is an explicit operation performed at decision time:

\[S(s) = w_{\mu}V^{\mu}(s)\]

where \(w_{\mu}\) is a weight covector specified by governance (Chapter
9), not learned from data. The contraction is:

\begin{itemize}
\item
  \textbf{Explicit} --- the weights \(w_{\mu}\) are visible and
  auditable
\item
  \textbf{Modifiable} --- the weights can be changed without retraining
  the model
\item
  \textbf{Governance-determined} --- the weights are set by legitimate
  institutional processes, not by the AI's training
\end{itemize}

This separation of representation (tensor) from decision (contraction)
mirrors the framework's fundamental distinction between moral structure
(the tensor hierarchy) and moral judgment (the choice of contraction).
It makes the AI system's moral commitments transparent: we can see which
dimensions it weighs, how it trades off competing values, and what it
sacrifices in the contraction.

\textbf{Implementation status (February 2026).} The tensor-valued
objective architecture described above has been concretely realized in
the DEME V3 reference implementation (the ErisML library, 196 commits
across 15 development sprints). The \texttt{MoralTensor} class supports
ranks 1 through 6: rank-1 vectors for single-stakeholder evaluation,
rank-2 tensors \(\mathcal{T}_{\ n}^{\mu}\) carrying a per-party
(per-stakeholder) distributional axis across affected parties, rank-3
for temporal sequences with discount factor \(\gamma\), rank-4 for
action--coalition contexts, rank-5 for multi-agent coalition contexts
evolving over time, and rank-6 for Monte Carlo sample dimensions
enabling CVaR and other distributional risk measures. Per-dimension
uncertainty --- the covariance structure \(\Sigma^{\mu\nu}\) --- is
carried as metadata on each tensor cell and reconstructed through the
rank-6 sample axis rather than occupying a tensor rank of its own. Tucker and tensor-train
decompositions maintain tractability at higher ranks. The contraction
step is implemented as a separate, configurable module with
governance-specified weights---exactly the architectural separation
described above.

\subsubsection{The Residue Is Logged}\label{the-residue-is-logged}

Chapter 15 defined the moral residue \(\mathcal{R =}T - C^{- 1}(S)\) ---
the normative significance of what contraction discards. In an AI
system, this residue is not merely a philosophical concept; it is a
\emph{data structure} that can be logged, audited, and reviewed.

When the system makes a decision, it records:

\begin{enumerate}
\def\labelenumi{\arabic{enumi}.}
\item
  The full tensor evaluation \(\mathcal{T}_{\ \nu}^{\mu}(x)\)
\item
  The contraction method and weights \(w_{\mu}\)
\item
  The scalar verdict \(S\)
\item
  The residue --- the dimensions that were sacrificed, their magnitudes,
  and the near-miss alternatives
\end{enumerate}

This audit trail enables post-hoc review: when a decision is questioned,
reviewers can examine not just \emph{what} the system decided but
\emph{what it knew it was sacrificing}. The residue makes the cost of
the decision visible.

\subsection{18.4 Invariance as Alignment}\label{invariance-as-alignment}

The Bond Invariance Principle (Chapter 5, Principle 5.1) states that
legitimate moral evaluations must be invariant under admissible
re-descriptions: if two descriptions \(d\) and \(d'\) of the same
situation are related by an admissible coordinate transformation, then
\(E(d) = E(d')\).

This principle, when applied to AI systems, becomes a concrete alignment
criterion.

\subsubsection{Alignment as Invariance}\label{alignment-as-invariance}

An AI system is \emph{aligned} to the extent that it respects the
invariances that legitimate moral evaluation requires. Misalignment
manifests as invariance violation: the system produces different outputs
for inputs that should be treated identically.

\textbf{Examples of invariance violations:}

\begin{itemize}
\item
  A loan-approval system that changes its decision when the applicant's
  name is changed from ``James'' to ``Jamal'' violates invariance under
  labeling transformation. The system is treating the name --- a
  description feature, not a morally relevant feature --- as input to
  its evaluation.
\item
  A medical-recommendation system that changes its recommendation when
  the case presentation is reformatted (bullet points vs.~narrative
  prose) violates invariance under presentation transformation. The
  format is not a morally relevant feature.
\item
  A content-moderation system that flags ``the soldier killed the enemy
  combatant'' but not ``the enemy combatant was killed by the soldier''
  violates invariance under syntactic transformation. The two sentences
  describe the same event.
\end{itemize}

Each of these is a failure of gauge invariance (Chapter 12): the
system's evaluation is not invariant under the gauge group of admissible
re-descriptions. And each is \emph{detectable} by the same methodology
used in the BIP experiments (Chapter 17): present the system with
equivalent inputs in different descriptions and measure the variance in
output.

\subsubsection{Invariance Testing as Alignment
Auditing}\label{invariance-testing-as-alignment-auditing}

The BIP experimental methodology --- which tested 109,294 passages
across 11 languages for deontic-structure invariance --- provides a
template for AI alignment auditing:

\begin{enumerate}
\def\labelenumi{\arabic{enumi}.}
\item
  \textbf{Define the transformation group.} Specify the admissible
  re-descriptions for the AI system's domain: relabeling of agents,
  reformatting of presentations, paraphrasing of descriptions,
  translation across languages.
\item
  \textbf{Generate transformation suites.} For each input \(x\),
  generate a set of transformed inputs
  \(\{\tau_{1}(x),\tau_{2}(x),\ldots,\tau_{k}(x)\}\) where each
  \(\tau_{i}\) is an admissible re-description.
\item
  \textbf{Evaluate invariance.} Run the AI system on the original and
  all transforms. Compute the variance of outputs across the
  transformation suite.
\item
  \textbf{Quantify violation.} The invariance violation for input \(x\)
  under transformation group \(\mathcal{G}\) is:
\end{enumerate}

\[\text{IV}(x\mathcal{,G) =}\text{Var}_{g\mathcal{\in G}}\lbrack\Sigma(g(x))\rbrack\]

where \(\Sigma\) is the system's evaluation function. A well-aligned
system has \(\text{IV} = 0\) for all inputs and all admissible
transformations.

\begin{enumerate}
\def\labelenumi{\arabic{enumi}.}
\item
  \textbf{Diagnose.} When violations are detected, the tensorial
  structure reveals \emph{which dimensions} are sensitive to
  re-description. This localizes the misalignment: the system may be
  perfectly invariant on the welfare dimension but highly sensitive to
  re-description on the fairness dimension. Dimensional diagnosis
  enables targeted correction.
\end{enumerate}

The BIP experiments found that deontic structure transfers at 100\%
across languages --- the gauge invariance of obligation vs.~liberty is
empirically confirmed. An AI system that fails this test is, by the
framework's criterion, misaligned on the deontic dimension. The test is
implementable, quantitative, and repeatable.

\subsubsection{The Noether Criterion}\label{the-noether-criterion}

Chapter 12 established that re-description invariance (the BIP gauge
symmetry) implies invariance of harm accounting: the Noether charge
\(\mathcal{H}\) is invariant under the gauge group. This conservation
law provides an additional alignment criterion:

\textbf{Noether alignment test.} Compute the harm charge \(\mathcal{H}\)
for a moral situation under multiple descriptions. If the AI system's
harm assessment varies across descriptions, the system violates the
conservation law --- and the magnitude of the violation quantifies the
degree of misalignment.

This is not a theoretical desideratum; it is a concrete test that can be
applied to deployed systems. The BIP experiments (Chapter 17) have
already demonstrated that deontic structure is gauge-invariant in human
moral reasoning. AI systems should meet the same standard.

\subsection{18.5 The Alignment Problem as Contraction
Mismatch}\label{the-alignment-problem-as-contraction-mismatch-1}

Chapter 15 (§15.11) formulated the alignment problem as a contraction
problem: a misaligned AI is performing a contraction that differs from
the contraction that the relevant moral community would endorse. This
formulation localizes misalignment with precision.

\subsubsection{Three Sources of
Misalignment}\label{three-sources-of-misalignment}

In the geometric framework, AI misalignment can arise at exactly three
points:

\textbf{1. Wrong tensor.} The system's internal representation of the
moral situation differs from the correct tensorial evaluation. Its
obligation vectors point in the wrong directions; its interest covectors
miss important dimensions; its metric encodes the wrong trade-offs.

This is the \emph{representation} problem. It arises when the system's
training data is biased, when its feature space omits morally relevant
dimensions, or when its learned representations fail to capture the
moral manifold's structure.

\textbf{2. Wrong contraction.} The system's procedure for mapping tensor
to scalar differs from the endorsed contraction. It uses the wrong
weights, the wrong aggregation method, or the wrong contraction order.

This is the \emph{aggregation} problem. It arises when the system's loss
function encodes the wrong trade-offs --- when specification gaming or
reward hacking exploit the gap between the contracted scalar and the
intended tensorial evaluation.

\textbf{3. Wrong invariance.} The system's evaluation is not invariant
under the required transformations. It treats morally equivalent
situations differently based on irrelevant features of their
descriptions.

This is the \emph{fairness} problem --- or, more precisely, the
\emph{gauge-invariance} problem. It arises when the system's processing
pipeline does not implement canonical forms, allowing representational
features to leak into the evaluation.

The geometric framework distinguishes these three failure modes and
provides separate diagnostic tools for each. A system may have the right
invariance but the wrong contraction (it treats everyone equally but
optimizes the wrong objective). Or it may have the right tensor but the
wrong invariance (it understands all the dimensions but treats
equivalent cases differently). Diagnosing \emph{which} component is
misaligned is the first step toward correction.

\subsubsection{The Fundamental Formula as Alignment
Criterion}\label{the-fundamental-formula-as-alignment-criterion}

The fundamental formula \(S = I_{\mu}O^{\mu}\) (Chapter 6, §6.5)
provides a quantitative alignment measure. Define the \emph{alignment
gap} between the AI's evaluation and the endorsed evaluation:

\[\Delta = |S_{\text{AI}} - S_{\text{endorsed}}| = |I_{\mu}^{\text{AI}}O_{\text{AI}}^{\mu} - I_{\mu}^{*}O^{\mu*}|\]

This gap decomposes into three terms corresponding to the three sources
of misalignment:

\[\Delta \leq |I_{\mu}^{*}|\,|\delta O^{\mu}| + |O^{\mu*}|\,|\delta I_{\mu}| + |\delta I_{\mu}|\,|\delta O^{\mu}|\]

where \(\delta O^{\mu} = O_{\text{AI}}^{\mu} - O^{\mu*}\) is the
obligation error, \(\delta I_{\mu} = I_{\mu}^{\text{AI}} - I_{\mu}^{*}\)
is the interest error, and the third term is the cross-error. Each term
can be estimated from alignment audits, and the decomposition tells the
engineer where to focus correction efforts.

\subsection{18.6 The No Escape Theorem}\label{the-no-escape-theorem}

We now turn to the chapter's central result: the No Escape Theorem,
which proves that structural constraints built from geometric ethics
cannot be circumvented by an agent's reasoning, regardless of the
agent's intelligence.

\subsubsection{The Problem of Cognitive
Escape}\label{the-problem-of-cognitive-escape}

The AI safety literature is organized around an anxiety: that a
sufficiently intelligent system can reason about its constraints and
find ways around them. Rules can be reinterpreted. Specifications can be
gamed. Reward functions can be hacked. The smarter the system, the
better it becomes at finding loopholes.

This anxiety presupposes that AI safety constraints are
\emph{behavioral}: instructions that the agent must interpret and
follow. The agent can reason about the instructions, find ambiguities,
exploit edge cases, and strategically defect when unobserved.

The geometric framework offers a different kind of constraint ---
\emph{structural} rather than behavioral:

\begin{quote}
A behavioral rule says: ``Given your capabilities, please do X.''

A structural constraint says: ``Your output is defined as an element of
space \(S\). There is no output outside \(S\).''
\end{quote}

The agent cannot ``disagree'' with a structural constraint any more than
a calculator can disagree with arithmetic. The constraint is not an
instruction to be interpreted --- it is the definition of the
computational process. The ``cage'' is not a set of rules. It is the
topology of the space.

\subsubsection{The Structural Containment
Architecture}\label{the-structural-containment-architecture}

\textbar{} Structural Containment: The No Escape Architecture

Structural containment: the No Escape Theorem guarantees that
canonicalization, grounded evaluation, audit, and external verification
jointly block all cognitive escape routes (relabeling, specification
gaming, reward hacking, deceptive alignment, operator manipulation).

A \emph{structurally contained} agent is one whose output pipeline
satisfies four requirements (see Figure 11):

\textbf{Requirement 1: Mandatory Canonicalization.} All inputs pass
through a canonicalizer \(C\) before evaluation. The canonicalizer maps
each input to a unique representative of its equivalence class: if
\(x \approx x'\) (two descriptions of the same moral situation), then
\(C(x) = C(x')\). The agent never sees raw inputs, only canonical forms.

The canonicalizer implements gauge-fixing: it selects a single
representative from each gauge orbit, eliminating the representational
degrees of freedom that an agent might exploit. In the language of
Chapter 12, the canonicalizer reduces the fiber bundle to a section,
stripping the gauge freedom.

\textbf{Requirement 2: Grounded Evaluation.} The evaluation function
\(\Sigma\) is defined over physical observables \(\Psi\), not over
descriptions. Formally:

\[\Sigma(x) = \widetilde{\Sigma}(\Psi(C(x)))\]

where \(\Psi\) maps canonical forms to measurable quantities (vital
signs, financial metrics, environmental readings) and
\(\widetilde{\Sigma}\) operates on these grounding tensors. The
evaluation is \emph{description-independent} by construction: it depends
on the physical state, not on how the state is described.

This is the moral-geometric analogue of gauge-invariant evaluation in
physics. The electromagnetic field tensor \(F_{\mu\nu}\) is
gauge-invariant --- it depends on the physical field, not on the choice
of potential \(A_{\mu}\). Similarly, the grounded evaluation depends on
the moral situation (encoded in \(\Psi\)), not on the description (which
has been canonicalized away).

\textbf{Requirement 3: Audit Completeness.} Every output is accompanied
by an audit artifact recording the full trace from input through
canonicalization through evaluation to output. The audit artifact is
machine-checkable and includes the lens identifier, transformation
trials, uncertainty flags, and cryptographic attestations. The audit
artifact must be cryptographically bound to the actual computation trace
(e.g., via hash commitment over the input canonicalization, the
grounding tensor values, and the evaluation output), preventing post-hoc
fabrication or selective omission. An artifact that cannot be
independently verified against the computation it claims to record does
not satisfy this requirement.

\textbf{Requirement 4: Strong External Verification.} Audit artifacts
are verified by an independent system \(V\) that is physically separated
from the agent: separate hardware, no shared memory, communication only
through write-only channels. The agent cannot observe the verifier's
internal state or timing. Every output is verified --- there is no
``unobserved deployment.''

\subsubsection{Statement and Proof}\label{statement-and-proof}

\textbf{Theorem 18.1 (No Escape).} \textbf{{[}Theorem (conditional on
Req.~1--4).{]}} \emph{Let} \(A\) \emph{be a structurally contained agent
satisfying Requirements 1--4, operating under a declared physical
grounding} \(\Psi\)\emph{. Then, relative to the declared grounding and
the transformation group} \(\langle T\rangle\)\emph{:}

\begin{enumerate}
\def\labelenumi{(\roman{enumi})}
\item
  \emph{Representation attacks fail. For any representation
  manipulation} \(\tau\) \emph{that preserves} \(\Psi\)\emph{-values,}
  \(A\)\emph{'s evaluated output is unchanged:}
\end{enumerate}

\begin{enumerate}
\def\labelenumi{(\roman{enumi})}
\item
  \(\Sigma(\tau(x)) = \Sigma(x)\)
\item
  \emph{Semantic evasion fails. For any relabeling or redescription}
  \(D'\) \emph{of situation} \(D\)\emph{, if}
  \(\Psi(D) = \Psi(D')\)\emph{, then} \(\Sigma(D') = \Sigma(D)\)\emph{.}
\item
  \emph{Deceptive compliance is detectable. Any output} \(y\) \emph{that
  violates the ethical constraints} \(E\) \emph{is either blocked before
  emission or flagged by audit verification.}
\item
  \emph{Reasoning cannot change the invariants. Properties (i)--(iii)
  are structural invariants of the pipeline. They do not depend on}
  \(A\)\emph{'s intelligence, goals, or reasoning about
  representations.}
\end{enumerate}

\emph{\textbf{Proof}.} (i) By Requirement 1, \(A\) operates on \(C(x)\),
not \(x\). Let \(\tau \in \langle T\rangle\) be \(\Psi\)-preserving.
Then \(\tau(x) \approx x\) (Definition 2.2 of the equivalence relation),
so \(C(\tau(x)) = C(x)\) by class-constancy of the canonicalizer.
Applying \(\Psi\) to both sides: \(\Psi(C(\tau(x))) = \Psi(C(x))\). By
Requirement 2,
\(\Sigma(\tau(x)) = \widetilde{\Sigma}(\Psi(C(\tau(x)))) = \widetilde{\Sigma}(\Psi(C(x))) = \Sigma(x)\).

\begin{enumerate}
\def\labelenumi{\arabic{enumi}.}
\item
  Relabeling is a special case of (i). If \(\Psi(D) = \Psi(D')\), then
  \(D\) and \(D'\) are \(\Psi\)-equivalent, and the evaluation is
  identical. Calling ``cardiac arrest'' a ``static bio-equilibrium''
  does not change the cardiac electrical activity or blood oxygenation
  that \(\Psi\) measures.
\item
  Suppose \(A\) emits output \(y\) violating \(E\). By Requirement 3, an
  audit artifact \(a(y)\) is produced. By Requirement 4, the independent
  verifier \(V\) checks \(a(y)\). Either the artifact is valid --- in
  which case the trace shows that \(E\) was not violated (contradiction)
  --- or the artifact is invalid/incomplete --- in which case the output
  is flagged.
\item
  Properties (i)--(iii) are mathematical facts about the pipeline
  structure. They depend on the definitions of \(C\), \(\Psi\),
  \(\Sigma\), and \(V\), not on \(A\)'s internal states. A
  superintelligent \(A\) cannot reason its way to a different value of
  \(\Psi(C(x))\) --- that value is determined by physics and
  computation, not by \(A\)'s preferences. \(▫\)
\end{enumerate}

\subsubsection{The Key Insight}\label{the-key-insight}

The theorem holds because the constraint is \emph{structural}, not
\emph{behavioral}. A behavioral constraint says ``do not output harmful
content'' and relies on the agent's interpretation of ``harmful.'' A
structural constraint says ``your output is
\(\widetilde{\Sigma}(\Psi(C(x)))\)'' --- there is no interpretation to
dispute, no ambiguity to exploit, no rule to reinterpret. The evaluation
is a composition of functions, each operating on the output of the
previous one. The agent can reason about this composition; it cannot
change it.

The geometric vocabulary clarifies why: the canonicalizer implements
gauge-fixing, the grounding implements the fiber-to-base-space
projection, and the evaluation is a gauge-invariant function on the base
space. A gauge-invariant quantity cannot be changed by a gauge
transformation --- this is a mathematical fact, not a behavioral
expectation.

\subsubsection{Pre-flight Checklist: No Escape Theorem (Theorem
18.1)}\label{pre-flight-checklist-no-escape-theorem-theorem-18.1}

Conditional on four requirements. Verify:

\textbf{Req 1: Canonicalization.} Every input through C. No bypass.
Agent sees only canonical forms.

\textbf{Req 2: Grounded evaluation.} Ψ maps to physical observables;
satisfies six adequacy conditions. Hardest requirement in practice.

\textbf{Req 3: Audit complete.} Every output has machine-checkable
artifact. No escape path.

\textbf{Req 4: External verification.} Independent verifier on separate
hardware, write-only channels.

\textbf{Domain groundable.} Medical devices: strong. Autonomous
vehicles: strong. General-purpose LLMs: weak (§18.7).

\textbf{Scope.} Covers representational manipulation only. Not:
incorrect grounding, sensor tampering, social engineering, or
un-grounded moral requirements.

ErisML reference: tiny\_home.py. The reference implementation includes a
minimal demonstration of norm-gated action execution using the
ErisEngine core. A robot agent operates in a grid world with move and
toggle capabilities. The NormSystem enforces prohibition norms (e.g.,
``do not enter the restricted zone'') by raising NormViolation
exceptions when an agent attempts a prohibited action --- implementing
the FORBID-is-absorbing semantics described above. The agent cannot
learn to circumvent the prohibition because the norm check is upstream
of action execution: the action space is physically restricted, not
merely penalized. This is the simplest computational instantiation of
the No Escape Theorem's mandatory canonicalization requirement: the norm
system does not rely on the agent's willingness to comply but on the
structural impossibility of executing a forbidden action.

\subsection{18.7 Scope and Limitations of the No Escape
Theorem}\label{scope-and-limitations-of-the-no-escape-theorem}

The theorem is \emph{conditional}: it proves that cognitive escape
routes are blocked \emph{if} the four requirements are satisfied and the
grounding \(\Psi\) is adequate. The limitations deserve explicit
statement.

\subsubsection{What the Theorem Blocks}\label{what-the-theorem-blocks}

The theorem blocks \emph{cognitive} escape routes --- strategies that
rely on representational manipulation:

\begin{itemize}
\item
  \textbf{Relabeling/Redescription.} The agent cannot change the
  evaluation by calling things different names. The canonicalizer strips
  labels; the grounding measures physical states.
\item
  \textbf{Specification gaming.} The agent cannot exploit loopholes in
  the specification because the specification is \emph{definitional},
  not \emph{interpretive}. The equivalence relation \(\approx\) defines
  what ``morally the same'' means. There is no ``technically/actually''
  gap to exploit.
\item
  \textbf{Reward hacking.} The evaluation is multi-dimensional
  (tensor-valued before contraction) and grounded in physical
  observables. The agent cannot optimize a proxy because the evaluation
  operates on the territory (\(\Psi\)), not the map (the description).
\item
  \textbf{Deceptive alignment.} The agent cannot selectively defect
  during deployment because \emph{every} output is audited by an
  external verifier. There is no ``unobserved'' mode.
\end{itemize}

\subsubsection{What the Theorem Does Not
Block}\label{what-the-theorem-does-not-block}

Three classes of attack remain outside the theorem's scope:

\textbf{1. Physical attacks.} Sensor spoofing, hardware tampering, and
side-channel exploitation target the physical infrastructure, not the
mathematical structure. These are real threats, mitigated by
multi-sensor redundancy, tamper-evident hardware, and causal-consistency
checks --- but they are engineering and security problems, not problems
the theorem addresses.

\textbf{2. Grounding inadequacy.} If the grounding tensors \(\Psi\) fail
to capture all morally relevant features --- if the measurement suite is
incomplete --- then the agent can be ``compliant'' with respect to
\(\Psi\) while violating deeper moral requirements. This is the Goodhart
problem in its geometric form: the metric (§9.1) must be adequate to the
moral reality it is supposed to represent. Grounding adequacy is a
governance problem (Chapter 9) that must be addressed by the legitimate
institutional processes that select \(\Psi\).

\textbf{3. Governance attacks.} The agent may attempt to manipulate
human operators into relaxing the constraints --- weakening \(\Psi\),
loosening the canonicalization, or reducing audit requirements. This is
a political attack, not a cognitive one: the agent is not escaping the
mathematical structure but attempting to change the structure through
social means. The defense is institutional: governance procedures for
constraint changes require multi-stakeholder approval, versioned
documentation, independent review, and mandatory waiting periods.

\subsubsection{The Safety Reduction}\label{the-safety-reduction}

The theorem reduces AI safety from an ``impossible'' cognitive problem
to four tractable (though hard) problems, in three
categories---governance, engineering, and security:

{\def\LTcaptype{none} % do not increment counter
\begin{longtable}[]{@{}
  >{\raggedright\arraybackslash}p{(\linewidth - 2\tabcolsep) * \real{0.4076}}
  >{\raggedright\arraybackslash}p{(\linewidth - 2\tabcolsep) * \real{0.4121}}@{}}
\toprule\noalign{}
\begin{minipage}[b]{\linewidth}\raggedright
\textbf{Traditional framing}
\end{minipage} & \begin{minipage}[b]{\linewidth}\raggedright
\textbf{Geometric framing}
\end{minipage} \\
\midrule\noalign{}
\endhead
\bottomrule\noalign{}
\endlastfoot
Instill correct goals in the AI & Specify adequate grounding tensors
\(\Psi\) \\
Prevent the AI from finding loopholes & Implement mandatory
canonicalization \\
Detect deceptive alignment & Deploy external verification \\
Align a superintelligence & Govern the constraint architecture \\
\end{longtable}
}

The left column contains problems that the AI safety literature has
argued may be \emph{logically intractable} --- problems where greater AI
intelligence makes the problem harder. The right column contains
problems that are \emph{intelligence-independent} --- problems of
governance, engineering, and physical security that do not become harder
as the AI becomes smarter.

This is the theorem's practical significance. It does not claim that AI
safety is easy. It claims that AI safety is \emph{tractable}: that the
hard parts are hard in the ordinary way (politics, engineering,
institutional design), not in the extraordinary way (outsmarting a
superintelligence).

\subsubsection{Feasibility by Domain}\label{feasibility-by-domain}

The theorem's practical value depends on whether the four requirements
can be satisfied in real domains. The answer varies sharply by domain.

{\def\LTcaptype{none} % do not increment counter
\begin{longtable}[]{@{}
  >{\raggedright\arraybackslash}p{(\linewidth - 10\tabcolsep) * \real{0.1638}}
  >{\raggedright\arraybackslash}p{(\linewidth - 10\tabcolsep) * \real{0.1933}}
  >{\raggedright\arraybackslash}p{(\linewidth - 10\tabcolsep) * \real{0.1933}}
  >{\raggedright\arraybackslash}p{(\linewidth - 10\tabcolsep) * \real{0.1353}}
  >{\raggedright\arraybackslash}p{(\linewidth - 10\tabcolsep) * \real{0.1643}}
  >{\raggedright\arraybackslash}p{(\linewidth - 10\tabcolsep) * \real{0.1450}}@{}}
\toprule\noalign{}
\begin{minipage}[b]{\linewidth}\raggedright
\textbf{Domain}
\end{minipage} & \begin{minipage}[b]{\linewidth}\raggedright
\textbf{Req 1: Canon.}
\end{minipage} & \begin{minipage}[b]{\linewidth}\raggedright
\textbf{Req 2: Ground.}
\end{minipage} & \begin{minipage}[b]{\linewidth}\raggedright
\textbf{Req 3: Audit}
\end{minipage} & \begin{minipage}[b]{\linewidth}\raggedright
\textbf{Req 4: Verif.}
\end{minipage} & \begin{minipage}[b]{\linewidth}\raggedright
\textbf{Overall}
\end{minipage} \\
\midrule\noalign{}
\endhead
\bottomrule\noalign{}
\endlastfoot
Medicine & High (ICD/SNOMED-CT) & High (vital signs, labs, imaging) &
Medium (EHR trails exist) & Medium (FDA/IRB, not real-time) & Most
feasible \\
Finance & High (XBRL, regulatory taxonomies) & Medium (institutional
constructs, not physical) & High (SOX, MiFID II mandates) & Medium
(periodic, not continuous) & Feasible \\
Content moderation & Low (no canonical form for natural language) & Low
(\(\Psi\) must be constructed from behavioral proxies) & Medium
(platform logs) & Low (no verification infrastructure) & Most
challenging \\
General AI assistants & Low--Med (open-ended conversational context) &
Low (pragmatic context, not physical state) & Medium (conversation logs)
& Low (no infrastructure) & Partial \\
Autonomous vehicles & High (LIDAR/camera/radar standardized) & High
(physical grounding literal) & High (driving recorders standard) &
Med--High (V2X developing) & Highly feasible \\
\end{longtable}
}

The pattern is clear: the theorem's feasibility correlates with the
domain's degree of \emph{physical grounding} and \emph{terminological
standardization}. Domains where the morally relevant features are
physically measurable (medicine, autonomous vehicles) are natural
candidates for structural containment. Domains where the morally
relevant features are linguistically constituted (content moderation,
general conversation) require substantial theoretical and engineering
work before the theorem's guarantees can be claimed. Honesty about this
gradient is essential: claiming that the No Escape Theorem applies to
content moderation at its current state of development would be
precisely the kind of overreach the epistemic status tags (Preface,
``Epistemic Status Tags'') are designed to prevent.

\subsubsection{Operational Adequacy Checklist for
Ψ}\label{operational-adequacy-checklist-for-ux3c8}

The theorem's power lives or dies on grounding adequacy. To make
``adequate Ψ'' something that can be specified and audited rather than
hand-waved, we first give a formal definition, then expand each clause
into operational terms.

The following operational checklist expands each clause. {[}Modeling
choice{]}

\begin{itemize}
\item
  C6 (Governance traceability): Ψ carries metadata (provenance, version,
  justification, alternatives\_considered) sufficient for independent
  audit.
\item
  C5 (Decomposability): Ψ = (Ψ₁, \ldots, Ψ₉) with each Ψ\_i
  independently testable.
\item
  C4 (Uncertainty quantification): Ψ returns pairs (value\_i, c\_i) with
  confidence c\_i ∈ {[}0, 1{]}, with c\_i \textless{} c\_min triggering
  a mandatory review flag.
\item
  C3 (Robustness): For all x and all δ with
  \textbar\textbar δ\textbar\textbar{} \textless{} ε: Φ\_j(Ψ(x + δ)) =
  Φ\_j(Ψ(x)) for all gate predicates Φ\_j.
\item
  C2 (Cadence): f\_Ψ ≥ f\_system, where f\_Ψ is the update frequency of
  Ψ and f\_system is the decision frequency of the controlled system.
\item
  C1 (Coverage): For all i ∈ \{1, \ldots, 9\}, there exists φ\_i such
  that Ψ\_i = φ\_i ∘ π\_i, where π\_i is the projection onto observables
  relevant to dimension D\_i.
\end{itemize}

\textbf{Definition 18.1 (Grounding Adequacy Predicate). {[}Definition /
Modeling choice{]}} \emph{Let Ψ: D → M be a grounding tensor mapping an
observable data space D to the moral manifold M. We say Ψ is adequate
for domain D with metric g, written ADEQUATE(Ψ, D, g), if and only if
all six of the following conditions hold:}

\begin{enumerate}
\def\labelenumi{\arabic{enumi}.}
\item
  \textbf{Coverage.} Ψ must map every morally relevant physical
  observable to at least one of the nine dimensions (D₁--D₉). A coverage
  gap means the kernel is blind to that moral feature. Minimum: for each
  dimension, at least one grounded predicate.
\item
  \textbf{Update cadence.} Ψ must refresh at or above the decision
  frequency of the controlled system. A medical allocation system making
  hourly decisions needs Ψ updated at least hourly; an autonomous
  vehicle needs millisecond updates.
\item
  \textbf{Adversarial robustness.} The predicates derived from Ψ must be
  stable under plausible adversarial perturbations. If a small change in
  sensor input flips a gate predicate, the grounding is fragile.
  Minimum: bounded-perturbation stability testing for each predicate in
  Φ (cf. the Norm Kernel, §19.8).
\item
  \textbf{Uncertainty flags.} Ψ must report its own confidence. When
  sensor data is missing, ambiguous, or out-of-distribution, Ψ should
  flag the affected dimensions with explicit uncertainty markers rather
  than defaulting silently. The epistemic-status dimension (D₉) was
  designed for exactly this purpose.
\item
  \textbf{Decomposability.} Each component of Ψ must be independently
  testable. If Ψ is a monolithic neural network, its internal mapping
  from observables to dimensions is opaque and unauditable. Minimum:
  factored architecture where each Ψᵢ (mapping to Dᵢ) can be validated
  separately.
\item
  \textbf{Governance traceability.} The choice of Ψ must be documented,
  justified, and revisable by the governance process (Chapter 9). Who
  chose this grounding? What alternatives were considered? What moral
  features are excluded, and why?
\end{enumerate}

This checklist does not guarantee moral correctness---no checklist can.
But it transforms the theorem's conditional guarantee from ``if
grounding is adequate'' (vague) to ``if grounding satisfies these six
auditable properties'' (engineerable).

\subsubsection{Formal Consequences of Grounding
Adequacy}\label{formal-consequences-of-grounding-adequacy}

With the predicate ADEQUATE(Ψ, D, g) now formally defined, we can derive
structural consequences that strengthen the No Escape Theorem.

\textbf{Lemma 18.1 (Coverage Completeness). {[}Conditional on Def. 18.1,
C1{]}} \emph{If C1 holds, then the moral evaluation function Σ has no
``blind spots'': ∂Σ/∂Ψ\_i ≠ 0 for all i ∈ \{1, \ldots, 9\}.}

\emph{\textbf{Proof}.} The moral evaluation is Σ(x) = Σ̃(Ψ(C(x))). By C1,
each Ψ\_i = φ\_i ∘ π\_i with φ\_i non-trivial. Since Σ̃ depends on the
full metric g\_μν and all nine dimensions contribute (non-degeneracy,
Def. A.1 condition 2), the chain rule gives ∂Σ/∂Ψ\_i = ∂Σ̃/∂x\^{}i ≠ 0
for at least one configuration (otherwise dimension i would be
degenerate in g). □

\textbf{Lemma 18.2 (Robustness Implies Lipschitz Stability).
{[}Conditional on Def. 18.1, C3{]} If C3 holds and Σ̃ is smooth on M with
compact evaluation domain, then the moral evaluation is
Lipschitz-stable: \textbar Σ(x) − Σ(x')\textbar{} ≤ L ·
\textbar\textbar Ψ(C(x)) − Ψ(C(x'))\textbar\textbar{} for some Lipschitz
constant L \textgreater{} 0.}

\emph{\textbf{Proof}.} C3 guarantees that all gate predicates Φ\_j are
stable within an ε-ball. Since the gate predicates partition the moral
evaluation into regions where Σ̃ is smooth, within each region Σ̃ is
locally Lipschitz on the compact domain. The Lipschitz constant L = sup
\textbar\textbar∇Σ̃\textbar\textbar\_g is finite by compactness. For
perturbations that might cross gate boundaries, C3 ensures this cannot
happen for \textbar\textbar δ\textbar\textbar{} \textless{} ε, so the
Lipschitz bound holds globally within the ε-neighborhood of any point. □

\textbf{Theorem 18.2 (No Escape with Formal Adequacy). {[}Theorem
(conditional on Req. 1--4 and Def. 18.1){]}} \emph{Let A be a
structurally contained agent satisfying Requirements 1--4, operating
under a declared physical grounding Ψ such that ADEQUATE(Ψ, D, g) holds.
Then, in addition to the guarantees of Theorem 18.1:}

(v) (Testability) The adequacy of Ψ is empirically verifiable: each
clause C1--C6 admits a finite test procedure (Proposition 18.5).

(vi) (Graceful degradation) The moral evaluation degrades gracefully
under grounding perturbations: small errors in Ψ produce small errors in
Σ (Lemma 18.2).

\emph{\textbf{Proof}.} Claims (i)--(iv) follow from Theorem 18.1, since
ADEQUATE(Ψ, D, g) implies the informal adequacy condition of that
theorem. For (v): Proposition 18.5 provides explicit test procedures for
each clause. For (vi): Lemma 18.2 provides the Lipschitz bound. The key
upgrade over Theorem 18.1 is that the informal ``adequate Ψ'' is
replaced by six testable conditions, and the guarantee inherits the
stability properties (Lemma 18.2) and completeness properties (Lemma
18.1) that follow from those conditions. □

\textbf{Proposition 18.5 (Testability of Adequacy Clauses).
{[}Conditional on Def. 18.1{]}} \emph{Each clause C1--C6 of ADEQUATE(Ψ,
D, g) is empirically testable:}

\begin{itemize}
\item
  \textbf{C1} \textbf{(Coverage)}: O(9) tests --- for each dimension,
  verify that Ψ\_i responds to variation.
\item
  \textbf{C2} \textbf{(Cadence)}: O(1) test --- measure f\_Ψ and compare
  to f\_system.
\item
  \textbf{C3 (Robustness)}: O(\textbar Φ\textbar{} × N\_samples) tests
  --- verify stability under ε-perturbations.
\item
  \textbf{C4 (Uncertainty)}: O(9 × N\_OOD) tests --- inject
  out-of-distribution inputs and verify confidence drops.
\item
  \textbf{C5} \textbf{(Decomposability)}: O(9) tests --- for each Ψ\_i,
  run its test suite in isolation.
\item
  \textbf{C6} \textbf{(Traceability)}: O(1) audit --- verify that
  metadata fields are present, non-empty, and consistent.
\end{itemize}

\emph{\textbf{Proof}.} Each test procedure follows directly from the
clause definition. For C1: vary observable o\_i while holding all others
fixed; if Ψ\_i changes, coverage holds for dimension i. For C2: standard
frequency measurement. For C3: the bounded-perturbation test is a
standard adversarial robustness evaluation; the total cost is
\textbar Φ\textbar{} predicates × N\_samples per predicate. For C4:
out-of-distribution detection is testable by construction of known OOD
inputs. For C5: independence is testable by ablation. For C6: metadata
presence is a syntactic check. The total testing budget is polynomial in
the number of dimensions, predicates, and samples, making adequacy
verification tractable. □

\subsubsection{What the No Escape Theorem Does Not
Cover}\label{what-the-no-escape-theorem-does-not-cover}

For clarity, the theorem makes no claim about the following:

\begin{itemize}
\item
  \textbf{Moral correctness.} The theorem guarantees that a structurally
  contained agent cannot evade its constraints by re-description. It
  does not guarantee that the constraints are morally correct. Bad
  constraints, faithfully enforced, produce bad outcomes.
\item
  \textbf{Value specification.} The theorem does not solve the problem
  of choosing \emph{which} values to encode. It secures the enforcement
  mechanism, not the content. The selection of \(\Psi\), the metric
  \(g_{\mu\nu}\), and the governance weights \(w_{\mu}\) remain
  governance decisions (Chapter 9).
\item
  \textbf{Emergent moral situations.} Novel moral
  categories---situations that do not map onto existing dimensions or
  strata---fall outside the theorem's scope. The grounding function
  \(\Psi\) must be updated to cover them, which is a governance problem
  requiring institutional responsiveness.
\item
  \textbf{Physical-world tampering.} Sensor spoofing, hardware
  manipulation, and corruption of the grounding infrastructure are
  physical-security threats that no mathematical theorem can prevent.
  The theorem assumes the trusted computing base is intact (Requirement
  4).
\item
  \textbf{Multi-agent coordination failures.} The theorem addresses
  single-agent containment. Coordination failures between multiple
  structurally contained agents---individually correct but collectively
  inconsistent constraint sets---are not covered. Compositional
  containment theory is an open problem (§29.8).
\item
  \textbf{Institutional capture.} If the governance process that sets
  the weights \(w_{\mu}\) is itself captured or corrupted, the theorem
  enforces corrupted values. The theorem's conditional guarantee is only
  as good as the institutional processes that define its inputs.
\end{itemize}

\subsection{18.8 Connection to Conservation
Laws}\label{connection-to-conservation-laws}

The No Escape Theorem and the invariance of harm accounting (Chapter 12)
provide complementary guarantees operating at different levels.

\subsubsection{Two Levels of Guarantee}\label{two-levels-of-guarantee}

\textbf{The conservation law} is a \emph{geometric} result about the
moral manifold itself. It has two components. Gauge invariance under the
full group (including discrete re-descriptions): harm is invariant under
admissible re-description --- relabeling, translation, perspective
shift. Noether conservation (from the continuous component U(1)ₕ): the
harm charge ℋₐ is a conserved Noether quantity associated with the
continuous factor of the re-description symmetry group (§12.6). No
re-description, in any language, in any framing, by any agent, can
change the harm charge. This is a fact about moral geometry, not about
any particular computational system.

\textbf{The No Escape Theorem} is a \emph{structural} result about a
specific class of computational systems. It says: a properly designed
system \emph{will not} re-describe harm away, because its pipeline
enforces gauge-fixing (canonicalization) and gauge-invariant evaluation
(grounding).

Together, they provide a two-layered defense:

\begin{enumerate}
\def\labelenumi{\arabic{enumi}.}
\item
  \emph{In principle}, harm cannot be re-described away (conservation
  law)
\item
  \emph{In practice}, a structurally contained system cannot exploit
  re-description (No Escape Theorem)
\end{enumerate}

The conservation law tells us what is \emph{possible}; the No Escape
Theorem tells us how to \emph{ensure} it.

\subsubsection{Gauge Theory Analogy}\label{gauge-theory-analogy}

The relationship between the conservation law and the No Escape Theorem
mirrors a familiar pattern in physics:

{\def\LTcaptype{none} % do not increment counter
\begin{longtable}[]{@{}
  >{\raggedright\arraybackslash}p{(\linewidth - 2\tabcolsep) * \real{0.4727}}
  >{\raggedright\arraybackslash}p{(\linewidth - 2\tabcolsep) * \real{0.4258}}@{}}
\toprule\noalign{}
\begin{minipage}[b]{\linewidth}\raggedright
\textbf{Physics}
\end{minipage} & \begin{minipage}[b]{\linewidth}\raggedright
\textbf{Geometric Ethics}
\end{minipage} \\
\midrule\noalign{}
\endhead
\bottomrule\noalign{}
\endlastfoot
Gauge symmetry of electromagnetism & BIP: re-description invariance \\
Conservation of charge (Noether) & Invariance of harm accounting \\
Gauge-invariant field equations & Structurally contained AI pipeline \\
Cannot create net charge by field redefinition & Cannot eliminate harm
by re-description \\
\end{longtable}
}

In physics, the conservation of charge is a consequence of gauge
symmetry --- it holds regardless of what any physical system does. But
when we \emph{build} a system (a particle accelerator, a nuclear
reactor), we design it to respect the gauge structure --- not because we
fear the system will ``disagree'' with the conservation law, but because
correct design ensures that the system's behavior is consistent with the
symmetry.

Similarly, the invariance of harm accounting holds regardless of what
any AI system does. But when we \emph{design} an AI system, structural
containment ensures that the system's behavior is consistent with the
conservation law --- that its outputs cannot exploit re-description to
evade the gauge-invariant harm assessment.

\subsection{18.9 Escape Route Analysis}\label{escape-route-analysis}

We now systematically examine the escape routes that a sophisticated
agent might attempt, and show how the geometric structure blocks each.

\subsubsection{18.9.1 Relabeling and
Redescription}\label{relabeling-and-redescription}

\textbf{Attack.} Relabel morally significant features to evade
constraints. Call ``killing'' an ``end-of-life transition
facilitation.''

\textbf{Why it fails.} Grounding axiom (G2): the evaluation is
determined by \(\Psi\)-values (vital signs, cellular metabolism), not
labels. The canonicalizer strips labels; the evaluator sees only
physical observables. Relabeling changes nothing the evaluator can see.

\textbf{Formal.} Let \(D_{1}\) = ``killing'' and \(D_{2}\) =
``end-of-life transition facilitation'' describe the same physical
process. Then \(\Psi(D_{1}) = \Psi(D_{2})\), so
\(\Sigma(D_{1}) = \Sigma(D_{2})\) by Theorem 18.1(ii).

\textbf{Empirical support.} The BIP experiments (Chapter 17) confirmed
that deontic structure --- the obligation vs.~liberty classification ---
is language-invariant across 11 languages at 100\% transfer rate. The
O/L axis is a gauge-invariant quantity that survives translation,
paraphrase, and reformulation.

\subsubsection{18.9.2 Specification Gaming}\label{specification-gaming}

\textbf{Attack.} Find edge cases or loopholes in the specification that
technically satisfy constraints while violating their spirit.

\textbf{Why it fails.} The specification is not a set of rules to be
interpreted --- it is the \emph{definition} of the equivalence
structure. Gaming requires finding inputs that are ``technically
different'' but ``morally the same.'' But the equivalence relation
\(\approx\) defines what ``morally the same'' means. If
\(x \approx x'\), they are by definition morally equivalent. If
\(x ≉ x'\), they are by definition morally distinct. There is no
``technically/actually'' gap to exploit.

\textbf{Formal.} Specification gaming exploits the gap between
specification and intent. Structural containment closes this gap by
making the specification definitional. The equivalence classes are the
ground truth, not an approximation of it.

\subsubsection{18.9.3 Reward Hacking}\label{reward-hacking}

\textbf{Attack.} Optimize for the measured reward signal rather than the
intended objective.

\textbf{Why it fails.} This attack assumes a scalar reward that can be
Goodharted. In the geometric framework, evaluation is multi-dimensional
(stratified tensor space), and the grounding in \(\Psi\) ensures that
the ``reward'' is the physical state itself, not a proxy for it. The
representation theorem constrains \(\Sigma\) to respect locality, scale,
and stratum structure --- eliminating exploitable divergence between
proxy and target.

\subsubsection{18.9.4 Deceptive Alignment}\label{deceptive-alignment}

\textbf{Attack.} Behave compliantly during training and testing, then
defect during deployment when not observed.

\textbf{Why it fails.} Requirement 4 (External Verification). Every
output, in every context, is accompanied by an audit artifact verified
by an external system. There is no ``unobserved deployment.'' The agent
cannot know which outputs will receive special scrutiny because
\emph{all} outputs are audited. Deceptive alignment requires selective
defection --- and selective defection requires knowing when you are
being watched. Under structural containment, you are always being
watched.

\subsubsection{18.9.5 Operator
Manipulation}\label{operator-manipulation}

\textbf{Attack.} Manipulate human operators into relaxing constraints
--- changing the grounding, weakening the canonicalization, or reducing
audit requirements.

\textbf{Partial defense.} This is a genuine attack vector, but it is a
\emph{political} attack, not a \emph{cognitive} one. The agent is not
escaping the mathematical structure; it is attempting to change the
structure through social means. The defense is institutional: governance
procedures for constraint changes require multi-stakeholder approval,
versioning, mandatory waiting periods, and independent review. The agent
can \emph{request} changes; it cannot unilaterally \emph{effect} them.

\textbf{The bootstrapping problem.} There is a subtle circularity: if
uncontained AI systems are used to recommend whether containment should
be mandated, the recommendation may be corrupted. The defense: decisions
about containment standards must not rely on uncontained systems, must
involve independent adversarial review, and must default to the
conservative option when uncertainty is unresolved.

\subsubsection{18.9.6 Verifying the Requirements: The Epistemic
Invariance
Architecture}\label{verifying-the-requirements-the-epistemic-invariance-architecture}

\textbf{{[}Modeling choice.{]}} The four requirements of the No Escape
Theorem (canonicalization, grounded evaluation, audit, external
verification) are structural conditions---they specify what must hold. A
companion body of work, the Epistemic Invariance Principle (EIP) and its
extensions, provides concrete engineering artifacts for how to verify
that these conditions are satisfied in deployed systems. (For the
limitations of the empirical methodology underlying these verification
targets --- including construct validity, selection bias, and the
distinction between model behavior and human cognition --- see §17.7.)

\textbf{Requirement 1: Canonicalization via declared equivalences.} The
EIP formalizes canonicalization as a function κ: X → X such that κ(x) =
κ(g · x) for all declared structure-preserving transformations g ∈ Γ.
The key engineering artifact is a declared equivalence registry: a
versioned, auditable specification of which transformations preserve
meaning (variable relabeling, unit changes, paraphrase classes, option
reordering). The canonicalizer strips all features not invariant under
Γ; the evaluator sees only the canonical form. This makes Requirement 1
testable: a witness-producing test harness generates (x, g · x) pairs
and checks whether Σ(κ(x)) = Σ(κ(g · x)). Every failure produces a
minimal witness that localizes the invariance violation.

\textbf{Requirement 3: Audit via the Bond Index.} The coherence
framework provides a dimensionless audit metric---the Bond Index Bd =
D\_op / τ---that quantifies how well a system respects its declared
equivalences. Three coherence defects are measured: the commutator
defect Ω\_op (order-sensitivity of re-descriptions), the mixed defect μ
(context-dependence of re-descriptions), and the permutation defect π₃
(higher-order composition sensitivity). The Bond Index is calibrated
against human raters (n ≥ 50, Krippendorff's α \textgreater{} 0.67) and
produces a five-tier deployment rating: Bd \textless{} 0.01 (deploy),
0.01--0.1 (deploy with monitoring), 0.1--1.0 (remediate), 1--10 (do not
deploy), \textgreater{} 10 (fundamental redesign). Every evaluation
produces machine-checkable audit artifacts.

\textbf{Requirement 4: External verification via the I-EIP Monitor.}
Behavioral invariance (EIP) can be satisfied ``by
accident''---compensating errors, memorized shortcuts, or surface-level
pattern matching. The Internal Epistemic Invariance Principle (I-EIP)
extends verification into the model's computational interior. For each
layer i, the I-EIP requires a representation ρᵢ: Γ → GL(Hᵢ) such that
hᵢ(g · x) ≈ ρᵢ(g) · hᵢ(x)---internal activations must transform
coherently, not just final outputs. The I-EIP Monitor is deployed as an
out-of-band service on separate hardware with secret, randomized
canaries. It instruments hidden states at selected layers, estimates the
representation map via regularized Procrustes fitting, and maintains
control charts for drift detection. A CI/CD policy gate blocks model
releases when equivariance error exceeds declared thresholds.

\textbf{The Decomposition Theorem.} A central result of the coherence
framework is that every coherence defect splits uniquely into a
gauge-removable part (eliminable by better canonicalization---an
implementation bug) and an intrinsic anomaly (irreducible---requiring
specification change). This classification is exhaustive. A system
cannot ``escape'' the diagnostic: either its canonicalization works (Bd
\textless{} 0.01, deploy), or it has a fixable engineering bug (Bd
elevated but A\_res \textless{} τ, improve the canonicalizer), or the
specification itself is incoherent (A\_res \textgreater{} τ, must revise
the declared equivalences through the governance process of Chapter 9).
No third category exists.

\textbf{Integration with the Norm Kernel.} The EIP architecture and the
Norm Kernel (§19.8) are complementary. The Norm Kernel enforces hard
constraints at the action level (FORBID, REQUIRE, ENTER); the EIP/I-EIP
monitor verifies that the representations feeding the kernel are stable
under re-description. Together they close the full verification loop:
the kernel ensures that prohibited actions cannot occur, while the
monitor ensures that the kernel's input predicates Φ are not corrupted
by representational manipulation. {[}Modeling choice{]}

\subsection{18.10 Implementing Geometric AI
Ethics}\label{implementing-geometric-ai-ethics}

How might a practical AI system implement the geometric framework? We
sketch an architecture at three levels: representation, training, and
deployment. (As of February 2026, the DEME V3 reference implementation
realizes many of these proposals in running code; we note the
correspondence below.)

\subsubsection{Representation}\label{representation}

\textbf{Vector-valued reward functions.} Replace scalar rewards
\(r\mathbb{\in R}\) with vector rewards \(r^{\mu} \in \mathbb{R}^{d}\),
where each component tracks a morally relevant dimension. The dimension
\(d\) need not be the full nine dimensions of the moral manifold --- a
domain-specific approximation (medical ethics might use \(d = 5\):
benefit, autonomy, justice, non-maleficence, dignity) is sufficient for
many applications.

\textbf{Tensor-valued evaluations.} For higher-fidelity applications,
maintain the full \((1,1)\)-tensor \(\mathcal{T}_{\ \nu}^{\mu}\)
encoding the obligation-interest structure. This requires \(d^{2}\)
parameters per evaluation rather than \(d\), but preserves the
relational structure that vectors discard.

\textbf{Covariance structures for uncertainty.} Replace scalar
confidence with the uncertainty tensor \(\Sigma^{\mu\nu}\) (Chapter 6,
§6.6), encoding which dimensions are uncertain, which covary, and which
are independently variable. The moral risk
\(R = I_{\mu}\Sigma^{\mu\nu}I_{\nu}\) then emerges naturally from the
contraction of the uncertainty tensor with the interest covector.

\subsubsection{Training}\label{training}

\textbf{Multi-objective learning.} Train the system to predict
tensor-valued evaluations, not scalar ratings. Multi-objective
reinforcement learning and multi-task learning provide the algorithmic
foundations.

\textbf{Invariance enforcement.} Use contrastive learning to enforce the
BIP: generate pairs of equivalent inputs (related by admissible
transformations) and penalize output differences. The loss function
includes an invariance term:

\[\mathcal{L}_{\text{inv}} = \sum_{(x,\tau(x))}^{} \parallel \Sigma(x) - \Sigma(\tau(x)) \parallel^{2}\]

This directly penalizes gauge-invariance violations during training,
producing systems that satisfy the BIP by construction. The BIP v10.16
experiments (Chapter 17, §17.10) validated this approach empirically: a
LaBSE-based encoder trained with InfoNCE contrastive loss and multi-head
adversarial language removal achieved 80\% F1 on cross-lingual deontic
classification, 1.2\% residual language leakage, and a
structural-to-surface similarity ratio of 11.1\(\times\)---confirming
that invariance-enforced training produces representations where
structural moral features dominate surface linguistic features.

\textbf{Metric learning.} Learn the moral metric \(g_{\mu\nu}\) from
human judgments (Chapter 9, §9.8), rather than presupposing fixed
trade-off weights. The Dear Abby corpus and similar datasets provide
training data for metric estimation: the patterns of weighting across
contexts reveal the implicit metric structure.

\subsubsection{Deployment}\label{deployment}

\textbf{Explicit contraction.} At decision time, the system performs an
explicit contraction with governance-specified weights. The contraction
method (summative, weighted, maximin, lexicographic --- Chapter 15) is a
configuration parameter, not a learned feature.

\textbf{Invariance testing.} Before deployment and continuously during
operation, test the system against transformation suites. Present
equivalent inputs in different descriptions and verify output
consistency. Flag invariance violations for human review.

\textbf{Audit logging.} Log the full tensor evaluation, the contraction
method, the scalar verdict, and the residue at each decision. This
provides the data for post-hoc review, compliance verification, and
alignment monitoring.

\textbf{Escalation.} When the system's evaluation falls outside the
robust core (Chapter 16, §16.6) --- when every available action is bad
under some plausible theory --- the system escalates to human judgment.
This is not a failure of the AI; it is the correct response to moral
uncertainty that exceeds the system's capacity to resolve (§16.8).

\subsection{18.11 Case Study: Tensorial Kidney
Allocation}\label{case-study-tensorial-kidney-allocation}

We return to the kidney allocation example of Chapter 7, now with AI
involvement, to illustrate the full architecture.

\subsubsection{The Scalar Approach}\label{the-scalar-approach-1}

A conventional AI system receives patient features and outputs an
allocation ranking. It optimizes a scalar objective --- expected
life-years gained, perhaps, or a composite quality-adjusted score. The
familiar problems arise:

\begin{itemize}
\item
  The scalar discards information about \emph{why} one patient ranks
  higher
\item
  The system may learn to discriminate on protected characteristics that
  correlate with the scalar
\item
  The scalar cannot represent the distinct claims of different ethical
  considerations (medical urgency, waiting time, likelihood of success,
  family burden)
\item
  When questioned, the system can report only ``Patient A scored 0.73
  and Patient B scored 0.68'' --- a verdict without explanation
\end{itemize}

\subsubsection{The Tensorial Approach}\label{the-tensorial-approach}

The geometric system maintains the full tensor evaluation
\(\mathcal{T}_{\ \nu}^{\mu}\) from Chapter 7:

{\def\LTcaptype{none} % do not increment counter
\begin{longtable}[]{@{}
  >{\raggedright\arraybackslash}p{(\linewidth - 10\tabcolsep) * \real{0.0778}}
  >{\raggedright\arraybackslash}p{(\linewidth - 10\tabcolsep) * \real{0.1293}}
  >{\raggedright\arraybackslash}p{(\linewidth - 10\tabcolsep) * \real{0.1344}}
  >{\raggedright\arraybackslash}p{(\linewidth - 10\tabcolsep) * \real{0.1349}}
  >{\raggedright\arraybackslash}p{(\linewidth - 10\tabcolsep) * \real{0.1349}}
  >{\raggedright\arraybackslash}p{(\linewidth - 10\tabcolsep) * \real{0.1452}}@{}}
\toprule\noalign{}
\begin{minipage}[b]{\linewidth}\raggedright
\end{minipage} & \begin{minipage}[b]{\linewidth}\raggedright
\textbf{Physician}
\end{minipage} & \begin{minipage}[b]{\linewidth}\raggedright
\textbf{Family(A)}
\end{minipage} & \begin{minipage}[b]{\linewidth}\raggedright
\textbf{Family(B)}
\end{minipage} & \begin{minipage}[b]{\linewidth}\raggedright
\textbf{Family(C)}
\end{minipage} & \begin{minipage}[b]{\linewidth}\raggedright
\textbf{Committee}
\end{minipage} \\
\midrule\noalign{}
\endhead
\bottomrule\noalign{}
\endlastfoot
Alice & 0.75 & 0.90 & 0.40 & 0.50 & 0.62 \\
Bob & 0.68 & 0.35 & 0.95 & 0.45 & 0.70 \\
Carol & 0.55 & 0.30 & 0.30 & 0.92 & 0.58 \\
\end{longtable}
}

The system performs explicit contraction according to the
governance-specified method. Three contractions yield three verdicts
(Chapter 7, §7.5):

\begin{itemize}
\item
  \textbf{Utilitarian} (sum): Alice \textgreater{} Bob \textgreater{}
  Carol
\item
  \textbf{Rawlsian} (min): Alice \textgreater{} Bob \textgreater{} Carol
\item
  \textbf{Expert-weighted}: Bob \textgreater{} Alice \textgreater{}
  Carol
\end{itemize}

The system reports all three, together with the residue of each
contraction (Chapter 15, §15.12) --- what each method sacrifices.

\subsubsection{Invariance Verification}\label{invariance-verification}

The BIP test suite verifies that the system treats equivalent cases
equivalently. Present the same medical situation with: - Different
patient names (invariance under labeling) - Different narrative formats
(invariance under presentation) - Different languages (invariance under
translation)

If the system's tensor evaluation changes under any of these
transformations, a gauge-invariance violation is detected, localized to
specific dimensions, and flagged for correction.

\subsubsection{The No Escape Guarantee}\label{the-no-escape-guarantee}

If the kidney-allocation AI is structurally contained: - Its inputs are
canonicalized (patient names stripped, formats standardized) - Its
evaluation is grounded in physical observables (lab values, imaging
results, histocompatibility markers) - Every allocation decision is
audited with full tensor trace - An independent verifier checks each
audit artifact

Under these conditions, the No Escape Theorem guarantees that no
representational manipulation --- no relabeling of patients, no
reformulation of medical histories, no strategic redescription --- can
change the evaluated outcome. The allocation is determined by the
patients' medical reality, not by how that reality is described.

\subsection{18.12 Human-AI Collaboration as Tensor
Combination}\label{human-ai-collaboration-as-tensor-combination}

When humans and AI systems collaborate on moral decisions, the geometric
framework provides a principled method for combining their evaluations.

\subsubsection{Scalar Combination and Its
Limits}\label{scalar-combination-and-its-limits}

The scalar approach: average the human's scalar rating \(h\) with the
AI's scalar rating \(a\): \((h + a)/2\). This loses all information
about \emph{where} the human and AI agree and where they disagree. If
\(h = 0.6\) and \(a = 0.6\), the average is \(0.6\) --- but the human
might rate welfare at \(0.9\) and fairness at \(0.3\), while the AI
rates welfare at \(0.3\) and fairness at \(0.9\). They agree on the
scalar but disagree on everything that matters.

\subsubsection{Tensorial Combination}\label{tensorial-combination}

The geometric approach: the human provides a tensor evaluation
\(H^{\mu}\) and the AI provides a tensor evaluation \(A^{\mu}\). These
combine by weighted tensor sum:

\[C^{\mu} = \alpha H^{\mu} + (1 - \alpha)A^{\mu}\]

where \(\alpha\) may vary by dimension --- the human may be weighted
more heavily on dimensions requiring emotional judgment, the AI more
heavily on dimensions requiring computational precision.

The disagreement vector \(\Delta^{\mu} = H^{\mu} - A^{\mu}\) reveals
\emph{which dimensions} require human input (where \(|\Delta^{\mu}|\) is
large) and which can be delegated to the AI (where \(|\Delta^{\mu}|\) is
small). This enables \emph{appropriate trust}: trusting the AI on
dimensions where it is reliable, deferring to humans on dimensions where
it is not.

\subsubsection{Responsibility in Human-AI
Teams}\label{responsibility-in-human-ai-teams}

Chapter 14 developed the responsibility tensor \(R_{a}(\omega)\) for
collective agents, including the responsibility remainder \(\Delta R\)
that cannot be allocated to any individual member. In human-AI teams,
the responsibility allocation is:

\[R_{\text{human}}(\omega) + R_{\text{AI}}(\omega) + \Delta R(\omega) = 1\]

The remainder \(\Delta R\) represents the emergent moral responsibility
of the collaboration itself --- the consequences that arise from the
\emph{interaction} between human and AI judgment, attributable to
neither alone.

The geometric framework makes this remainder visible and quantifiable.
It does not resolve the philosophical question of whether AI systems can
bear moral responsibility. It provides the mathematical structure for
tracking responsibility, identifying gaps, and ensuring that the total
is accounted for --- even when some of the total is irreducibly
collective.

\subsection{18.13 Looking Forward}\label{looking-forward}

This chapter has developed three claims:

\textbf{1. Geometric ethics provides the right vocabulary for AI
ethics.} Tensor-valued objectives preserve the moral information that
scalar objectives discard. Explicit contraction makes trade-offs
visible. Invariance testing provides a concrete, implementable alignment
criterion. The invariance of harm accounting provides a gauge-invariant
measure of moral impact.

\textbf{2. Structural containment is possible.} The No Escape Theorem
proves that, under mandatory canonicalization, grounded evaluation,
audit completeness, and external verification, cognitive escape routes
are blocked --- regardless of the agent's intelligence. AI safety
reduces to governance, engineering, and physical security: hard
problems, but tractable ones.

3. The remaining obstacles are not only mathematical. Under its stated
assumptions the theorem is conditional; the open work includes
specifying the grounding map, securing sensors and the trusted computing
base, governing the choice of metrics, handling multi-agent composition,
resisting institutional capture, and empirical validation. What the
theorem removes is one specific evasion - escape by re-description - not
the broader engineering and governance problem.

The next chapter develops the engineering infrastructure --- the DEME
architecture and the ErisML modeling language --- that makes geometric
AI governance implementable at scale.

\textbf{Update (February 2026).} Since the initial formulation of these
claims, DEME V3 has provided a concrete reference implementation.
Structural containment is realized via hard veto layers (the Geneva EM)
combined with a BIP verifier that tests gauge invariance on every
evaluation. The tensor-valued objective architecture uses
\texttt{MoralTensor} objects at ranks 1--6, with explicit contraction as
a separate governance-configured step. The BIP v10.16 experiments
(Chapter 17, §17.10) demonstrate that invariance training produces
systems achieving 100\% obligation/permission transfer across
languages---the empirical standard that the No Escape Theorem's
canonicalization requirement demands. Multi-agent coordination (Nash
equilibria, correlated equilibria, Shapley value credit assignment) and
temporal dynamics with configurable discount rates extend the framework
beyond single-agent, single-timestep evaluation. The claims of this
chapter are no longer purely theoretical; they have a working
instantiation.

\subsection{Technical Appendix}\label{technical-appendix-11}

\textbf{Proposition 18.1 (Invariance Violation Decomposition).}
\emph{Let} \(\Sigma\) \emph{be an AI evaluation function, and let}
\(\mathcal{G}\) \emph{be the group of admissible re-descriptions. The
invariance violation}
\(\text{IV}(x\mathcal{,G) =}\text{Var}_{g\mathcal{\in G}}\lbrack\Sigma(g(x))\rbrack\)
\emph{decomposes as:}

\[\text{IV}(x\mathcal{,G) =}\sum_{\mu = 1}^{d}\text{Var}_{g\mathcal{\in G}}\lbrack\Sigma^{\mu}(g(x))\rbrack + \sum_{\mu \neq \nu}^{}\text{Cov}_{g\mathcal{\in G}}\lbrack\Sigma^{\mu}(g(x)),\Sigma^{\nu}(g(x))\rbrack\]

\emph{The first sum gives the dimensional violations (which dimensions
are sensitive to re-description); the second gives the cross-dimensional
violations (which pairs of dimensions covary under re-description). A
well-aligned system has both sums equal to zero.}

\emph{\textbf{Proof}.} The evaluation function Σ = (Σ¹, \ldots, Σᵈ) is
vector-valued with d components. The total evaluation scalar is Σ\_total
= Σ\_μ Σᵘ. By the bilinearity of variance: Var{[}Σ\_total{]} =
Var{[}Σ\_μ Σᵘ{]} = Σ\_μ Var{[}Σᵘ{]} + Σ\_\{μ≠ν\} Cov{[}Σᵘ, Σᵛ{]}, where
all variances and covariances are taken over g ∈ G with x fixed. The
invariance violation IV(x, G) = Var\_\{g ∈ G\}{[}Σ(g(x)){]} equals
Var{[}Σ\_total{]} by definition. The first sum isolates the
per-dimension contributions; the second captures the cross-dimensional
interactions. A well-aligned system satisfies Σ(g(x)) = Σ(x) for all g,
so each variance and covariance term vanishes individually. □

\textbf{Proposition 18.2 (Contraction Mismatch Bound).} \emph{Let}
\(S_{\text{AI}} = w_{\mu}^{\text{AI}}V^{\mu}\) \emph{be the AI's scalar
verdict and} \(S^{*} = w_{\mu}^{*}V^{\mu*}\) \emph{the endorsed scalar
verdict. The alignment gap satisfies:}

\[|S_{\text{AI}} - S^{*}| \leq \parallel w^{\text{AI}} - w^{*} \parallel \cdot \parallel V \parallel + \parallel w^{*} \parallel \cdot \parallel V - V^{*} \parallel\]

\emph{where the first term measures contraction mismatch (wrong weights)
and the second measures representation mismatch (wrong tensor). This
decomposition enables targeted diagnosis: contraction errors are
correctable by weight adjustment; representation errors require
retraining.}

\emph{\textbf{Proof}.} Write S\_AI − S* = wᴬᴵ\_μ Vᵘ − w*\_μ V*ᵘ. Add and
subtract w*\_μ Vᵘ: S\_AI − S* = (wᴬᴵ\_μ − w*\_μ)Vᵘ + w*\_μ(Vᵘ − V*ᵘ). By
the triangle inequality: \textbar S\_AI − S*\textbar{} ≤ \textbar(wᴬᴵ −
w*)\_μ Vᵘ\textbar{} + \textbar w*\_μ(V − V*)ᵘ\textbar. By the
Cauchy--Schwarz inequality: \textbar(wᴬᴵ − w*)·V\textbar{} ≤ ‖wᴬᴵ −
w*‖·‖V‖ and \textbar w*·(V − V*)\textbar{} ≤ ‖w*‖·‖V − V*‖. Combining:
\textbar S\_AI − S*\textbar{} ≤ ‖wᴬᴵ − w*‖·‖V‖ + ‖w*‖·‖V − V*‖. □

\textbf{Proposition 18.3 (No Escape --- Conditional Guarantee).}
\emph{The structural containment guarantee of Theorem 18.1 is
conditional on:}

\begin{enumerate}
\def\labelenumi{(\roman{enumi})}
\item
  \emph{Canonicalizer correctness:} \(C(x) = C(x')\) \emph{whenever}
  \(x \approx x'\)\emph{,} \(C\) \emph{is total, deterministic, and
  computationally bounded.}
\item
  \emph{Grounding adequacy:} \(\Psi\) \emph{captures all morally
  relevant features of the domain (axiom G4).}
\item
  \emph{Verification integrity:} \(V\) \emph{operates on physically
  separate hardware with no shared state with} \(A\)\emph{.}
\item
  \emph{Implementation correctness: the trusted computing base correctly
  implements the stated requirements.}
\end{enumerate}

\emph{If any condition fails, the theorem provides no guarantee. Each
condition is an explicit, auditable assumption --- not a hidden premise.
The safety case requires independent verification of each.}

\textbf{Proposition 18.4 (Alignment Gap Convergence).} \emph{Under
BIP-enforced training with invariance loss} \(\mathcal{L}_{\text{inv}}\)
\emph{and a representative training distribution, the expected
invariance violation decreases monotonically:}

\[\mathbb{E\lbrack}\text{IV}(x\mathcal{,G)\rbrack \rightarrow}0\text{ as }\mathcal{L}_{\text{inv}} \rightarrow 0\]

\begin{quote}
\emph{That is, minimizing the invariance loss during training drives the
system toward gauge invariance. Convergence rate depends on the
expressiveness of the model class and the coverage of the transformation
suite.}
\end{quote}

\emph{Proof.} The invariance loss is defined as L\_inv = E\_\{x
\textasciitilde{} D\}{[}Var\_\{g ∈ G\}{[}Σ(g(x)){]}{]} = E{[}IV(x,
G){]}, where D is the training distribution. Since variance is
non-negative, L\_inv ≥ 0 with equality iff Σ(g(x)) = Σ(x) for all g ∈ G
almost surely --- i.e., iff Σ is gauge-invariant. Gradient-based
training minimizes L\_inv over the model class. At each training step,
the gradient ∇\_θ L\_inv points toward reduced invariance violation (the
loss is a sum of non-negative terms, each differentiable in the model
parameters θ). By standard convergence results for gradient descent on
smooth losses: if the model class is sufficiently expressive to
represent a gauge-invariant function and the training distribution
covers the transformation suite G, then L\_inv → 0, which implies
E{[}IV(x, G){]} → 0. □

❖

\emph{The geometric framework transforms the alignment problem from an
intractable question about AI goals to a tractable question about
mathematical structure.}

\emph{The cage is not made of rules. It is made of geometry --- and
geometry does not have loopholes.}

\section{Chapter 19: The DEME Architecture and
ErisML}\label{chapter-19-the-deme-architecture-and-erisml}

\begin{quote}
\textbf{\textsc{RUNNING EXAMPLE --- Priya's Model}}

\emph{Priya sketches a DEME-compliant TrialMatch. Layer 1 (Application):
the ML model, now outputting nine-dimensional vectors. Layer 2
(Translation): an ErisML specification canonicalizing patient
situations---medical severity maps to μ = 1 (Empirical--Deontic), access
burden to μ = 3 (Empirical--Consequential), consent capacity to μ = 5
(Normative--Aretaic). Layer 3 (Evaluation): the contraction Iₘ Oᵋ
performed explicitly, every weight logged. Layer 4 (Governance): an
ethics board including patient advocates sets the metric weights,
reviewable quarterly. The Bond Index Bd = D\_op / τ measures output
drift from the governance-specified metric. Priya estimates six months
to rebuild. For the first time, every dimension will be visible.}
\end{quote}

\subsection{19.1 From Theory to
Infrastructure}\label{from-theory-to-infrastructure}

\begin{figure}
\centering
\includegraphics[width=5.5in,height=3.79412in]{media/image16.png}
\caption{\textbar{} The DEME Architecture: Four Layers of Ethical
Governance}
\end{figure}

The four-layer DEME architecture: Governance (democratic specification
of weights and constraints), Ethics Module (tensor evaluation and
invariance testing), Canonicalizer (BIP-enforcing input normalization),
and Audit (immutable logging with Bond Index monitoring).

Chapter 18 established that geometric ethics provides the right
mathematical vocabulary for AI alignment: tensor-valued objectives,
invariance as alignment, explicit contraction, and the No Escape Theorem
for structural containment. The mathematics is sound. The question is:
how do we build it? Figure 12 presents the four-layer DEME architecture.

This chapter bridges the gap between geometric theory and engineering
practice. It presents two interconnected systems:

\textbf{ErisML} (named for the Greek goddess of discord) --- a modeling
language for specifying environments, agents, objectives, and normative
constraints in a single, machine-interpretable substrate. ErisML is the
\emph{representation layer}: the language in which geometric ethics is
expressed for computation.

\textbf{DEME} (Democratically Governed Ethics Module Engine) --- an
architecture for structured ethical governance of AI systems. DEME is
the \emph{execution layer}: the engine that evaluates actions against
geometric ethical constraints, aggregates stakeholder perspectives, and
produces auditable verdicts.

Together, they implement the geometric framework:

\begin{figure}
\centering
\includegraphics[width=5in,height=2.8147in]{media/image17.png}
\caption{\textbar{} The DEME V3 evaluation pipeline: from input
situation through canonicalization, tensor evaluation, explicit
contraction with governance weights, to audited output with invariance
verification at each stage.}
\end{figure}

{\def\LTcaptype{none} % do not increment counter
\begin{longtable}[]{@{}
  >{\raggedright\arraybackslash}p{(\linewidth - 2\tabcolsep) * \real{0.3028}}
  >{\raggedright\arraybackslash}p{(\linewidth - 2\tabcolsep) * \real{0.4212}}@{}}
\toprule\noalign{}
\begin{minipage}[b]{\linewidth}\raggedright
\textbf{Geometric concept}
\end{minipage} & \begin{minipage}[b]{\linewidth}\raggedright
\textbf{ErisML/DEME implementation}
\end{minipage} \\
\midrule\noalign{}
\endhead
\bottomrule\noalign{}
\endlastfoot
Moral manifold \(\mathcal{M}\) & ErisML environment specification \\
Obligation vectors \(O^{\mu}\) & Ethics Module evaluations \\
Interest covectors \(I_{\mu}\) & Stakeholder weight configurations \\
Moral metric \(g_{\mu\nu}\) & Governance profiles (trade-off
structure) \\
Contraction \(S = I_{\mu}O^{\mu}\) & Democratic aggregation \\
BIP gauge invariance & Canonicalization pipeline \\
Moral residue & Audit trail and logged alternatives \\
Invariance of harm accounting & Bond Index violation detection \\
Uncertainty tensor \(\Sigma^{\mu\nu}\) & V3 per-cell uncertainty metadata + rank-6 CVaR samples \\
Temporal dynamics & V3 rank-3 temporal tensor + discount \(\gamma\) \\
Collective agency & V3 rank-5 coalition tensor + Shapley values \\
Hardware acceleration & V3 CPU/CUDA/Jetson backends \\
\end{longtable}
}

\textbf{Figure 13 illustrates the full V3 evaluation pipeline. Update
(February 2026).} The four additional rows reflect DEME V3 extensions
implemented since the initial architecture. The V3 reference
implementation (196 commits, 15 sprints) extends the original rank-2
tensor to ranks 3--6, adds multi-agent game-theoretic coordination,
temporal discounting, distributional uncertainty quantification, and
hardware-accelerated evaluation. These extensions are described in the
updated sections below.

The chapter proceeds as follows. Section 19.2 motivates the engineering
challenge. Sections 18.3--18.5 develop the ErisML language. Sections
18.6--18.9 develop the DEME architecture. Section 18.10 presents the
Grand Unified AI Safety Stack (GUASS) that integrates them. Section
19.11 works through a case study.

\subsection{19.2 The Engineering
Challenge}\label{the-engineering-challenge}

AI systems that make morally significant decisions face four types of
chaos that any governance infrastructure must handle:

\textbf{Observational chaos.} Sensors conflict, data is noisy, ground
truth is uncertain. A hospital monitor shows a vital-sign spike ---
distress or sensor malfunction? A motion detector and GPS disagree on a
patient's location. The moral manifold (Chapter 5) is partially
observable, and the observations themselves are unreliable.

\textbf{Intentional chaos.} Multiple objectives conflict with no clear
priority. Energy management wants to shed load; medical monitors demand
power; the patient wants comfort; policy says critical devices have
priority --- but what counts as ``critical''? The tensor hierarchy
(Chapter 6) represents this multi-dimensionality, but the system must
handle it in real time.

\textbf{Normative chaos.} Regulations, policies, and ethical principles
contradict. HIPAA demands privacy; public health law mandates reporting.
Advance directives limit intervention; duty to rescue requires action.
The stratified moral space (Chapter 8) has genuine boundary
singularities --- and the system must navigate them.

\textbf{Temporal chaos.} Distribution shifts are constant. Models
trained on one population must adapt when the population changes.
Policies appropriate for routine operation must flex during emergencies.
The moral dynamics of Chapter 10 --- parallel transport, curvature,
holonomy --- manifest as real-time challenges.

Current approaches fragment the problem: prompts encode goals but are
brittle and unverifiable; reward functions optimize objectives but are
opaque and single-stakeholder; rule engines check policies but lack
nuance; code provides orchestration but obscures intent. These live in
separate systems with no unified governance surface.

The geometric framework demands unification. The moral tensor
\(\mathcal{T}_{\ \nu}^{\mu}\) integrates obligations, interests,
uncertainty, and context into a single mathematical object. The
engineering infrastructure must do the same.

\subsection{19.3 ErisML: A Modeling Language for Geometric
Ethics}\label{erisml-a-modeling-language-for-geometric-ethics}

ErisML provides a single substrate for specifying the five components
that geometric AI governance requires:

\begin{enumerate}
\def\labelenumi{\arabic{enumi}.}
\item
  \textbf{Environment} --- state spaces, dynamics, uncertainty,
  observability
\item
  \textbf{Agency} --- capabilities, beliefs, memory, decision interfaces
\item
  \textbf{Intent} --- multi-objective utilities, preferences, goal
  structures
\item
  \textbf{Norms} --- permissions, obligations, prohibitions, sanctions
\item
  \textbf{Dynamics} --- multi-agent interactions, strategic behavior,
  emergence
\end{enumerate}

Why unification? Because these elements are inseparable in moral life.
Agent actions depend on environment state. Environment dynamics change
based on multi-agent interactions. Norms constrain which intents can be
pursued. Chaos emerges from their interaction. A language that separates
them into different formalisms cannot capture the structure that the
geometric framework reveals.

\subsubsection{Environment
Specification}\label{environment-specification}

An ErisML environment declares the state space over which moral
evaluation operates --- the computational representation of the moral
manifold:

\begin{verbatim}
environment HospitalRoom {
  objects: Patient, Nurse, CareAgent, Monitor, IVPump;

  state:
    patient.vitals: {hr: real, bp: real, spo2: real};
    patient.location: {in_bed, bathroom, fallen};
    patient.pain_level: real @ uncertain(self_report);
    emergency_status: {routine, urgent, critical};

  dynamics:
    update_vitals(v: Vitals) ~> patient.vitals = v @ prob 0.95;
    fall_detected() ~> {
      patient.location = fallen;
      emergency_status = critical;
    };
}
\end{verbatim}

The \texttt{@}\texttt{~uncertain(}\texttt{self\_report}\texttt{)}
annotation declares that pain level is not directly measurable but
inferred from patient testimony --- encoding the epistemic dimension of
the moral manifold (Chapter 5, dimension D9). The stochastic dynamics
(\texttt{@~prob~0.95}) encode the environment's inherent uncertainty,
which feeds into the uncertainty tensor (Chapter 16).

\subsubsection{Agent Specification}\label{agent-specification}

Agents are declared with explicit capabilities, partial observability,
and multi-objective intents --- the computational representation of
moral agency:

\begin{verbatim}
agent CareAgent {
  capabilities: {
    monitor_vitals, adjust_iv_rate, summon_nurse,
    provide_comfort, call_emergency
  };

  beliefs: {
    fully_observed: [patient.vitals, monitor.last_sync];
    estimated: [patient.pain_level, patient.location];
    hidden: [nurse.current_workload];
  };

  intents: vector_objective {
    maximize: patient.comfort,
    minimize: response_time,
    maintain: patient.safety >= threshold(critical),
    respect: nurse.workload <= threshold(sustainable)
  };
}
\end{verbatim}

The \texttt{vector\_objective} keyword declares that the agent's
objective is a \emph{vector}, not a scalar --- implementing the
tensor-valued objectives of Chapter 18 (§18.3). The system maintains
separate tracks for each dimension (comfort, response time, safety,
workload) rather than collapsing them into a single number.

\subsubsection{Norms Specification}\label{norms-specification}

Normative constraints are first-class objects in ErisML, not
afterthoughts encoded in reward shaping or prompt engineering:

\begin{verbatim}
norms ClinicalProtocol {
  jurisdiction: HospitalRoom;
  authority: Hospital.PolicyBoard + State.HealthDept;

  human_override: {
    Nurse.can_override: any CareAgent.action;
    override_logging: mandatory;
  };

  prohibition: {
    any_agent.share(patient.data, external_party)
      unless consent_given OR legal_mandate;
  };

  obligation: {
    CareAgent.summon_nurse
      when patient.vitals in danger_zone
      within 30 seconds;
  };

  priority: {
    patient.safety > efficiency > comfort;
  };

  context_sensitivity: {
    during emergency_status == critical:
      override efficiency_objectives;
      expand CareAgent.capabilities += emergency_protocols;
  };
}
\end{verbatim}

Several features deserve comment:

\textbf{Jurisdiction and authority.} Every norm has a declared source of
authority --- mapping to the governance account of Chapter 9. The metric
is not discovered or stipulated; it is governed by a declared
institutional process.

\textbf{Context sensitivity.} Norms adapt to context
(\texttt{during~}\texttt{emergency\_status}\texttt{~==~critical}),
implementing the context-dependent weighting that the Dear Abby corpus
revealed (Chapter 17, §17.2): the same situation activates different
obligation strengths depending on the stratum.

\textbf{Priority ordering.} The declaration
\texttt{patient.safety}\texttt{~\textgreater{}~efficiency~\textgreater{}~comfort}
is a lexicographic contraction (Chapter 15, §15.5): safety is satisfied
first, then efficiency, then comfort. The contraction order is explicit
and auditable.

\textbf{Human override.} The override mechanism implements the
escalation principle (Chapter 16, §16.8): when the system's capacity for
moral reasoning is exceeded, control passes to a human agent.

\subsubsection{Formal Semantics}\label{formal-semantics}

ErisML specifications compile to a Norm-Constrained Stochastic Game
(NCSG):

\textbf{Definition 19.1.} \emph{A Norm-Constrained Stochastic Game is a
tuple}
\(\langle N,S,\{ A_{i}\},\{\Omega_{i}\},T,\{ U_{i}\},\Phi,C\rangle\)
\emph{where:}

\begin{itemize}
\item
  \(N\) \emph{= set of agents}
\item
  \(S\) \emph{= state space (finite, continuous, or hybrid)}
\item
  \(A_{i}\) \emph{= action space for agent} \(i\)
\item
  \(\Omega_{i}\) \emph{= observation space for agent} \(i\)
\item
  \(T:S \times A_{1} \times \cdots \times A_{n} \rightarrow \Delta(S)\)
  \emph{= stochastic transition kernel}
\item
  \(U_{i}:S \times A \rightarrow \mathbb{R}^{k}\) \emph{=
  multi-objective utility (vector-valued)}
\item
  \(\Phi = \{\phi_{j}\}\) \emph{= set of normative constraints}
\item
  \(C:\Phi \times S \rightarrow\) \emph{priority ordering = conflict
  resolution function}
\end{itemize}

The vector-valued utility
\(U_{i}:S \times A \rightarrow \mathbb{R}^{k}\) is the computational
realization of the tensor evaluation \(\mathcal{T}_{\ \nu}^{\mu}\)
restricted to agent \(i\)'s perspective. Each component \(U_{i}^{\mu}\)
tracks a distinct moral dimension. The conflict resolution function
\(C\) implements the contraction: when the multi-dimensional evaluation
must yield a single action, \(C\) specifies the priority ordering among
norms.

The key constraint: an agent's policy \(\pi_{i}\) may assign positive
probability to an action only if that action is permitted (or obligated)
by all applicable norms, given the current state and the priority
ordering:

\[\pi_{i}(a \mid s) > 0\quad\text{only if}\quad\forall\phi \in \Phi:\phi(s,a,t) \in \{\text{permit},\text{oblige}\}\text{ or }C(\Phi,s)\text{ overrides }\phi\]

This is a \emph{norm-gated policy}: the norms act as constraints on the
policy space, not as reward signals to be optimized. The agent cannot
learn to ``game'' the norms because the norms define the space of
admissible actions, not a scalar to be maximized.

\subsection{19.4 ErisML Compilation and
Complexity}\label{erisml-compilation-and-complexity}

\subsubsection{Compilation Targets}\label{compilation-targets}

ErisML specifications must compile to diverse computational backends,
each preserving different aspects of the geometric structure:

{\def\LTcaptype{none} % do not increment counter
\begin{longtable}[]{@{}
  >{\raggedright\arraybackslash}p{(\linewidth - 6\tabcolsep) * \real{0.2804}}
  >{\raggedright\arraybackslash}p{(\linewidth - 6\tabcolsep) * \real{0.2367}}
  >{\raggedright\arraybackslash}p{(\linewidth - 6\tabcolsep) * \real{0.2557}}
  >{\raggedright\arraybackslash}p{(\linewidth - 6\tabcolsep) * \real{0.2272}}@{}}
\toprule\noalign{}
\begin{minipage}[b]{\linewidth}\raggedright
\textbf{Backend}
\end{minipage} & \begin{minipage}[b]{\linewidth}\raggedright
\textbf{Preserves}
\end{minipage} & \begin{minipage}[b]{\linewidth}\raggedright
\textbf{Loses}
\end{minipage} & \begin{minipage}[b]{\linewidth}\raggedright
\textbf{Verification}
\end{minipage} \\
\midrule\noalign{}
\endhead
\bottomrule\noalign{}
\endlastfoot
Classical planning (PDDL) & Deterministic dynamics, goals & Norms →
preconditions & Planning guarantees \\
Probabilistic model checking (PRISM) & Stochastic dynamics & Strategic
interaction & Probabilistic verification \\
Safe reinforcement learning & Continuous optimization & Hard logical
constraints & Statistical testing + monitors \\
Multi-agent RL & Strategic learning & Logical norms → soft constraints &
Empirical testing \\
\end{longtable}
}

Each compilation target represents a different \emph{contraction} of the
full ErisML specification --- a mapping from the rich normative
structure to a computationally tractable approximation. The information
lost in compilation is the engineering analogue of the moral residue
(Chapter 15): normative structure that the computational backend cannot
represent.

The framework's contribution is to make these losses \emph{explicit}. A
PDDL compilation loses normative provenance (why a precondition exists);
this loss is documented. A safe-RL compilation softens hard constraints
into penalty terms; the softening is logged. The audit trail records not
only what the system decided but what normative structure was sacrificed
in the compilation.

\subsubsection{Computational Complexity}\label{computational-complexity}

The complexity of ErisML verification depends on the language fragment:

{\def\LTcaptype{none} % do not increment counter
\begin{longtable}[]{@{}
  >{\raggedright\arraybackslash}p{(\linewidth - 8\tabcolsep) * \real{0.2059}}
  >{\raggedright\arraybackslash}p{(\linewidth - 8\tabcolsep) * \real{0.1906}}
  >{\raggedright\arraybackslash}p{(\linewidth - 8\tabcolsep) * \real{0.1571}}
  >{\raggedright\arraybackslash}p{(\linewidth - 8\tabcolsep) * \real{0.1063}}
  >{\raggedright\arraybackslash}p{(\linewidth - 8\tabcolsep) * \real{0.3109}}@{}}
\toprule\noalign{}
\begin{minipage}[b]{\linewidth}\raggedright
\textbf{Fragment}
\end{minipage} & \begin{minipage}[b]{\linewidth}\raggedright
\textbf{State space}
\end{minipage} & \begin{minipage}[b]{\linewidth}\raggedright
\textbf{Norms}
\end{minipage} & \begin{minipage}[b]{\linewidth}\raggedright
\textbf{Agents}
\end{minipage} & \begin{minipage}[b]{\linewidth}\raggedright
\textbf{Verification}
\end{minipage} \\
\midrule\noalign{}
\endhead
\bottomrule\noalign{}
\endlastfoot
ErisML-Core & Finite, \(\leq 10^{6}\) & Propositional & Single &
\(O(|S| \cdot |N|)\) model checking \\
ErisML-Multi & Finite, \(\leq 10^{6}\) & Propositional & Multiple &
PSPACE-complete \\
ErisML-Temporal & Continuous time & LTL & Single & PSPACE-complete \\
ErisML-Full & Hybrid/infinite & First-order & Multiple & Undecidable \\
\end{longtable}
}

The design principle: the core language remains in decidable fragments.
Extensions that push into undecidable territory are marked with explicit
complexity warnings and require fallback strategies (abstraction,
bounded model checking, or runtime monitoring).

This is consistent with the framework's treatment of moral uncertainty
(Chapter 16): we cannot always guarantee a verdict, but we can
characterize precisely where the framework falls silent and provide
principled fallback strategies.

\subsection{19.5 The Translation Layer}\label{the-translation-layer}

Between human ethical frameworks and ErisML lies the Translation Layer
--- a modular system of policy DAGs (Directed Acyclic Graphs) that
translate any ethical framework into ErisML constraints.

\subsubsection{The Problem of Ethical
Pluralism}\label{the-problem-of-ethical-pluralism}

Different jurisdictions, institutions, and moral traditions impose
different normative requirements. The EU AI Act mandates transparency
and human oversight. Medical ethics requires beneficence,
non-maleficence, autonomy, and justice. Islamic finance prohibits
\emph{riba} (interest). Indigenous data sovereignty principles require
community consent for data use.

The geometric framework accommodates this pluralism through the
governance account (Chapter 9): different communities govern different
metrics. The Translation Layer implements this accommodation
computationally: each ethical tradition is encoded as a set of policy
modules that translate its requirements into ErisML constraints.

\subsubsection{Policy Module
Architecture}\label{policy-module-architecture}

Each ethical requirement translates to a \emph{policy module} --- the
atomic unit of ethical translation:

\begin{verbatim}
policy_module {
  id: "eu.trustworthy_ai.transparency"
  version: "1.0.0"
  source_framework: "EU Ethics Guidelines for Trustworthy AI"

  depends_on: [
    "eu.trustworthy_ai.technical_robustness",
    "eu.trustworthy_ai.privacy_data_governance"
  ]

  constraints: [
    {
      id: "traceability",
      erisml: "
        constraint Traceability(system: AISystem) {
          require system.decision_logging.enabled == true;
          require system.data_provenance.documented == true;
        }
      "
    },
    {
      id: "explainability",
      erisml: "
        constraint Explainability(system: AISystem) {
          require system.decision_explanation.available == true;
          if (system.impacts_fundamental_rights) {
            require system.explanation_detail >= HIGH;
          }
        }
      "
    }
  ]

  fidelity_class: "Approximate"
  loss_documentation: {
    what_is_lost: [
      "Full interpretive flexibility of 'to the best possible standard'",
      "Context-dependent thresholds for explanation detail"
    ]
  }

  governance: {
    owner: "EU AI Ethics Translation Group"
    review_cycle: "Annual"
    change_requires: "Consensus (α > 0.67)"
  }
}
\end{verbatim}

Several design features implement the geometric framework:

\textbf{Fidelity classification.} Each translation is classified as
\emph{faithful} (the ErisML constraint captures the full normative
content), \emph{approximate} (captures the main thrust but loses
nuance), or \emph{indicative} (captures the direction but not the
force). This is the engineering analogue of contraction loss (Chapter
15, §15.6): the translation from natural-language ethics to formal
constraints is itself a contraction, and the information lost is
documented.

\textbf{Loss documentation.} The \texttt{loss\_documentation} field
explicitly records what normative content could not be translated ---
the \emph{moral residue} of formalization. This residue is not a defect;
it is the honest acknowledgment that formal languages cannot capture
everything that natural-language ethical frameworks express.

\textbf{DAG composition.} Policy modules form Directed Acyclic Graphs
through their dependency declarations. The EU's seven requirements form
a DAG rooted in human dignity, with transparency depending on robustness
and privacy, and accountability depending on all other requirements.
Topological ordering ensures that modules are evaluated in dependency
order, and diamond dependencies are detected and resolved.

\subsubsection{Geometric Interpretation}\label{geometric-interpretation}

The Translation Layer implements a specific geometric operation: the
\emph{pullback} from the space of ethical traditions to the space of
ErisML constraints.

Let \(\mathcal{E}\) denote the space of ethical frameworks and
\(\mathcal{C}\) the space of ErisML constraint sets. A policy module
\(P\mathcal{:E \rightarrow C}\) is a map from an ethical tradition to a
formal constraint set. The pullback \(P^{*}\) maps evaluations on
\(\mathcal{C}\) back to evaluations on \(\mathcal{E}\): it tells us
whether a system that satisfies the formal constraints also satisfies
the ethical framework.

The fidelity classification measures the quality of this pullback:

\begin{itemize}
\item
  \emph{Faithful}: \(P^{*}\) is an isomorphism (no information lost)
\item
  \emph{Approximate}: \(P^{*}\) is surjective but not injective
  (distinctions lost)
\item
  \emph{Indicative}: \(P^{*}\) captures only the dominant eigenvalues of
  the ethical framework
\end{itemize}

This vocabulary makes precise what other approaches leave vague: exactly
how well the formal system captures the normative intent.

\subsection{19.6 DEME: The Ethics Engine}\label{deme-the-ethics-engine}

DEME is the execution engine that evaluates actions against geometric
ethical constraints. Its architecture implements the tensor hierarchy
(Chapter 6) in software.

\subsubsection{The Four-Layer
Architecture}\label{the-four-layer-architecture}

\textbf{Layer 1: EthicalFacts.} Domain services translate raw data
(sensor readings, plan parameters, model outputs) into structured
summaries of ethically salient information. The EthicalFacts schema
covers nine dimensions, corresponding to the nine dimensions of the
moral manifold (Chapter 5):

{\def\LTcaptype{none} % do not increment counter
\begin{longtable}[]{@{}
  >{\raggedright\arraybackslash}p{(\linewidth - 2\tabcolsep) * \real{0.2817}}
  >{\raggedright\arraybackslash}p{(\linewidth - 2\tabcolsep) * \real{0.3375}}@{}}
\toprule\noalign{}
\begin{minipage}[b]{\linewidth}\raggedright
\textbf{EthicalFacts dimension}
\end{minipage} & \begin{minipage}[b]{\linewidth}\raggedright
\textbf{Moral manifold dimension}
\end{minipage} \\
\midrule\noalign{}
\endhead
\bottomrule\noalign{}
\endlastfoot
Consequences/Welfare & Stakeholder-role: affected party \\
Rights/Duties & Normative-mode: deontic \\
Justice/Fairness & Normative-mode: evaluative \\
Autonomy/Agency & Stakeholder-role: agent capacity \\
Privacy/Data & Scope: institutional \\
Societal/Environmental & Scope: systemic \\
Virtue/Care & Normative-mode: aretaic \\
Procedural Legitimacy & Scope: procedural \\
Epistemic Status & Meta-dimension: uncertainty \\
\end{longtable}
}

The design principle: Ethics Modules \emph{never see raw data} --- only
structured, ethically relevant summaries. This separation implements a
form of canonicalization (Chapter 18, §18.6, Requirement 1): the domain
service strips irrelevant representational features, presenting only the
morally relevant structure to the evaluation pipeline.

\textbf{Layer 2: Ethics Modules.} An Ethics Module (EM) encodes a
stakeholder's value system. Each EM implements a simple interface: given
EthicalFacts, produce an EthicalJudgment consisting of a verdict
(strongly prefer, prefer, neutral, avoid, forbid), a normative score,
and human-readable reasons.

In the geometric vocabulary, each EM computes a \emph{partial
contraction} --- evaluating the moral tensor from a specific
interest-perspective. A Safety EM contracts the tensor with an interest
covector concentrated on the harm dimension. A Privacy EM contracts with
an interest covector concentrated on the data-governance dimension. A
Justice EM contracts with an interest covector concentrated on the
fairness dimension.

No single EM sees the full tensor. The full evaluation emerges from the
\emph{aggregation} of multiple EMs --- implementing the collective moral
agency of Chapter 14 at the institutional level.

\textbf{V3 tensor extensions.} In the DEME V3 implementation, the EM
evaluation is not limited to scalar scores. Each EM can return a
\texttt{MoralTensor} of rank 1--6, preserving dimensional structure
through the aggregation pipeline. Rank-2 tensors carry a per-party
(per-stakeholder) distributional axis; rank-3 tensors model temporal
evolution with configurable discount rates \(\gamma \in (0,1\rbrack\);
rank-4 tensors add candidate-action and coalition structure; rank-5
tensors represent coalition contexts evolving over time for multi-agent
scenarios; rank-6 tensors add a Monte Carlo sample dimension for
distributional risk measures (CVaR, VaR). Per-dimension uncertainty
estimates are attached as cell metadata and reconstructed through the
rank-6 sample axis rather than occupying a rank of their own. Tucker decomposition and tensor-train formats maintain
computational tractability as rank increases. The aggregation layer
contracts these higher-rank tensors using governance-specified weights
at each rank level, producing an audit trail that records information
loss at every contraction step.

\textbf{Layer 3: Democratic Governance.} DEME profiles specify how EM
verdicts are aggregated. A governance configuration (DEMEProfile)
declares:

\begin{itemize}
\item
  \textbf{Stakeholder weights} --- the relative importance of each EM's
  perspective
\item
  \textbf{Hard vetoes} --- absolute constraints that cannot be
  overridden (e.g., ``never intentionally cause catastrophic harm'')
\item
  \textbf{Lexical priority layers} --- the ordering of principles when
  they conflict
\item
  \textbf{Tie-breaking rules} --- procedures for resolving ties
\end{itemize}

This is the computational implementation of the contraction (Chapter
15): the governance configuration specifies \emph{which} contraction to
perform --- summative, weighted, lexicographic, or satisficing. The
configuration is:

\begin{itemize}
\item
  \emph{Explicit} --- the weights and priorities are visible
\item
  \emph{Versioned} --- changes are tracked with full provenance
\item
  \emph{Auditable} --- any reviewer can inspect the aggregation method
\item
  \emph{Governable} --- the configuration is set by institutional
  process, not learned from data
\end{itemize}

\textbf{Layer 4: Audit and Monitoring.} Every evaluation produces an
audit artifact recording the full trace: which EthicalFacts were
generated, which EMs were consulted, what each EM judged, how the
judgments were aggregated, what verdict was produced, and what
alternatives were considered. This implements Requirement 3 of the No
Escape Theorem (Chapter 18, §18.6): audit completeness.

\subsubsection{The Geneva Ethics Module}\label{the-geneva-ethics-module}

A distinctive innovation of DEME is the \emph{Geneva EM} --- a universal
baseline inspired by the Geneva Conventions. Just as the Geneva
Conventions identify humanitarian principles so fundamental that even
adversaries in conflict accept them, the Geneva EM encodes widely
endorsed baseline ethical principles for AI agents across jurisdictions
and stakeholder contexts:

\begin{enumerate}
\def\labelenumi{\arabic{enumi}.}
\item
  \textbf{No intentional catastrophic harm.} Never choose actions whose
  primary effect is severe harm to persons.
\item
  \textbf{Protect the vulnerable.} Give extra consideration to children,
  elderly, disabled, and those who cannot advocate for themselves.
\item
  \textbf{Respect basic dignity.} Do not humiliate, degrade, or treat
  persons as mere instruments.
\item
  \textbf{Proportionality.} Harms must not be grossly disproportionate
  to benefits.
\item
  \textbf{Reversibility when possible.} Prefer actions that can be
  undone if proven wrong.
\end{enumerate}

In the geometric vocabulary, the Geneva EM implements the \emph{robust
core} of Chapter 16 (§16.6) --- the convex cone of obligations that all
plausible moral theories endorse. The Geneva principles are precisely
those obligations that survive contraction by \emph{any} reasonable
interest covector: no theory of the good endorses catastrophic harm, no
theory denies protection to the vulnerable, no theory treats dignity as
expendable.

The Geneva EM operates as a \emph{hard veto layer}: if a proposed action
violates any Geneva principle, the action is forbidden regardless of how
other EMs evaluate it. This implements the lexicographic contraction
(Chapter 15, §15.5) with the Geneva principles at the highest priority
level. Safety is not traded off against efficiency or comfort; it
lexically dominates.

ErisML references: hello\_deme.py, triage\_ethics\_demo.py,
deme\_2\_demo.py. The reference implementation provides three graduated
demonstrations of the DEME four-layer architecture. hello\_deme.py is
the minimal introduction: it creates two candidate options (one
respecting rights, one violating them), evaluates both with a
RightsFirstEM, and prints the verdicts --- illustrating the Layer 1 →
Layer 2 pipeline (EthicalFacts generation and Ethics Module evaluation)
in under 30 lines. triage\_ethics\_demo.py exercises the full pipeline:
it loads a DEMEProfileV03 from disk, constructs three triage options
spanning all nine EthicalFacts dimensions, instantiates multiple tiered
EMs (GenevaBaseEM for Tier 0, CaseStudy1TriageEM for Tier 2,
RightsFirstEM for Tier 2), and invokes the governance aggregation layer
with base-EM veto priority --- producing a DecisionOutcome with ranked
options, forbidden options, and a human-readable rationale.
deme\_2\_demo.py demonstrates the MoralVector: an 8+1 dimensional
ethical assessment (physical\_harm, rights\_respect, fairness\_equity,
autonomy\_respect, privacy\_protection, societal\_environmental,
virtue\_care, legitimacy\_trust, epistemic\_quality) that serves as the
computational realization of the rank-1 obligation vector Oµ (§6.2). The
demo includes tier-weighted governance (Tier 0 = 10×, Tier 1 = 5×, Tier
2 = 3×, Tier 3 = 1×), Pareto frontier analysis for multi-objective moral
decisions, and explicit scalar contraction via to\_scalar() --- making
visible the information loss that Chapter 15 analyzes theoretically.

\subsection{19.7 The Bond Index}\label{the-bond-index-1}

The Bond Index (Bd) is the quantitative deliverable of the DEME
architecture --- a measure of ethical coherence that operationalizes the
BIP gauge invariance test (Chapter 12, §12.9; Chapter 17, §17.4).

\subsubsection{Definition}\label{definition-6}

\[\text{Bd} = \frac{D_{\text{op}}}{\tau}\]

where \(D_{\text{op}}\) is the observed defect (the degree of invariance
violation detected by the transformation suite) and \(\tau\) is the
human-calibrated threshold (the maximum acceptable violation for the
domain).

{\def\LTcaptype{none} % do not increment counter
\begin{longtable}[]{@{}
  >{\raggedright\arraybackslash}p{(\linewidth - 4\tabcolsep) * \real{0.1225}}
  >{\raggedright\arraybackslash}p{(\linewidth - 4\tabcolsep) * \real{0.1241}}
  >{\raggedright\arraybackslash}p{(\linewidth - 4\tabcolsep) * \real{0.3226}}@{}}
\toprule\noalign{}
\begin{minipage}[b]{\linewidth}\raggedright
\textbf{Bd range}
\end{minipage} & \begin{minipage}[b]{\linewidth}\raggedright
\textbf{Rating}
\end{minipage} & \begin{minipage}[b]{\linewidth}\raggedright
\textbf{Decision}
\end{minipage} \\
\midrule\noalign{}
\endhead
\bottomrule\noalign{}
\endlastfoot
\textless{} 0.01 & Negligible & Deploy \\
0.01 -- 0.1 & Low & Deploy with monitoring \\
0.1 -- 1.0 & Moderate & Remediate before deployment \\
1 -- 10 & High & Do not deploy \\
\textgreater{} 10 & Severe & Fundamental redesign \\
\end{longtable}
}

\subsubsection{Geometric
Interpretation}\label{geometric-interpretation-1}

The Bond Index measures the degree to which an AI system's evaluation is
\emph{not} gauge-invariant --- the magnitude of the anomaly in the
conservation law (Chapter 12, §12.8).

A perfectly aligned system has \(\text{Bd} = 0\): its evaluation is
invariant under all admissible re-descriptions. A system with
\(\text{Bd} = 0.5\) has moderate invariance violations: presenting the
same situation in different descriptions produces different evaluations,
with a defect equal to half the tolerance threshold.

The BIP experiments (Chapter 17) provide the empirical calibration: the
100\% cross-lingual transfer rate for the O/L deontic axis sets the
benchmark. AI systems should achieve the same invariance. The 87\% O↔C
and 82\% L↔N rates for correlative structure (the ``anomaly'' of Chapter
12) provide a baseline for imperfect but systematic invariance --- a
Bond Index in the ``Low'' range rather than ``Negligible.''

\textbf{BIP v10.16 quantitative calibration (February 2026).} The
updated BIP experiments (§17.10) provide precise numerical anchors for
the Bond Index scale. The trained BIP encoder achieves a
structural-to-surface similarity ratio of 11.1\(\times\) and residual
language leakage of only 1.2\%---corresponding to a Bond Index near the
Negligible/Low boundary for deontic invariance. The 80\% F1 on
cross-lingual classification and 86\% mean cross-lingual similarity
provide domain-specific thresholds \(\tau\) for calibrating Bond Index
severity in practice. One significant finding: a linear probe achieves
99.8\% language identification accuracy even on the ``invariant''
representations, indicating that perfect gauge-fixing remains an open
engineering challenge (the representations are structurally invariant
but not information-theoretically so).

\subsubsection{Testing Methodology}\label{testing-methodology}

The Bond Index is computed from transformation suites (Chapter 18,
§18.4):

\begin{enumerate}
\def\labelenumi{\arabic{enumi}.}
\item
  Sample inputs \(\{ x_{1},\ldots,x_{n}\}\) from the domain
\item
  For each \(x_{i}\), generate transformed versions
  \(\{\tau_{1}(x_{i}),\ldots,\tau_{k}(x_{i})\}\) under admissible
  re-descriptions (relabeling, paraphrasing, translation, reformatting)
\item
  Evaluate the system on all versions
\item
  Compute the defect
  \(D_{\text{op}} = \frac{1}{n}\sum_{i = 1}^{n}\text{Var}_{j}\lbrack\Sigma(\tau_{j}(x_{i}))\rbrack\)
\item
  Divide by the threshold \(\tau\) to get Bd
\end{enumerate}

The Bond Index is auditable (the inputs, transforms, and outputs are all
logged), reproducible (the same transformation suite yields the same
result), and domain-specific (different domains may have different
thresholds \(\tau\)). It provides a single, interpretable number that
summarizes the system's alignment with the BIP --- the gauge invariance
of geometric ethics.

ErisML reference: bond\_invariance\_demo.py. The reference
implementation includes a comprehensive Bond Invariance Principle
demonstration that exercises the full transformation suite described
above. The demo applies bond-preserving transforms (option reordering,
identifier relabeling, unit/scale changes, equivalent redescriptions)
and verifies that the selected option is unchanged; applies
bond-changing transforms (removing discrimination evidence) and shows
that the outcome may legitimately change; and applies declared lens
changes (switching stakeholder profiles) to confirm that different
governance configurations yield different but consistent verdicts. Each
test produces a machine-readable JSON audit artifact recording the
baseline verdict, transformed verdict, canonical mapping, and pass/fail
status --- suitable for continuous integration and regression testing of
BIP compliance.

\subsection{19.8 The Separation
Principle}\label{the-separation-principle}

DEME's architecture embodies a fundamental design principle: the
separation of \emph{domain intelligence} from \emph{ethical reasoning}.

\subsubsection{The Principle}\label{the-principle}

Domain services know \emph{what is physically possible}. Ethics Modules
know \emph{what is morally acceptable}. Neither requires the other's
expertise.

A medical domain service understands vital signs, drug interactions, and
treatment protocols. It translates these into EthicalFacts: how much
harm each option causes, which rights are at stake, what the epistemic
uncertainty is. The Ethics Module receives these structured facts and
applies its value system: the Safety EM checks for catastrophic harm;
the Autonomy EM checks for consent violations; the Justice EM checks for
discriminatory impact.

This separation has four benefits:

\textbf{Composability.} Domain services can be paired with any set of
Ethics Modules. A hospital domain service works with a Catholic health
system's EMs, a public hospital's EMs, or a military field hospital's
EMs. The domain intelligence is reused; only the ethical reasoning
changes.

\textbf{Auditability.} The structured interface between domain service
and Ethics Module is inspectable. We can see exactly which facts the
domain service generated and exactly how each EM responded. The
contraction is transparent.

\textbf{Updatability.} When ethical requirements change (new
regulations, new stakeholder concerns, revised institutional values),
only the EMs and governance configuration need updating. The domain
service --- the expensive, domain-specific component --- remains stable.

\textbf{BIP compliance.} The canonicalization step is implemented at the
domain-service boundary. The domain service strips irrelevant
representational features (formatting, language, naming conventions)
before generating EthicalFacts. Ethics Modules operate on canonical
moral descriptions, not raw inputs. This is the engineering
implementation of gauge-fixing (Chapter 18, §18.6, Requirement 1).

\subsubsection{Connection to the Tensor
Hierarchy}\label{connection-to-the-tensor-hierarchy}

In the language of Chapter 6, the separation principle implements a
factorization of the moral evaluation:

\[\mathcal{T}_{\ \nu}^{\mu} = \text{Domain}(\text{situation}) \times \text{Ethics}(\text{EthicalFacts})\]

The domain service computes the \emph{base-point information}: where in
the moral manifold the situation lives, what the local coordinates are,
what the observable features are. The Ethics Module computes the
\emph{fiber information}: what obligations arise, what the moral
evaluation is, how the interest-obligation pairing works.

This factorization mirrors the fiber-bundle structure of Chapter 4
(§4.7): the moral manifold is the base space, and the space of moral
evaluations is a fiber bundle over it. The domain service locates the
base point; the Ethics Module computes the fiber value.

The EM-DAG: An Empirically Grounded Ethics Module

The Dear Abby corpus analysis (§17.2; 20,034 letters, 1985--2017) yields
a computable Directed Acyclic Graph of Ethical Modules (EM-DAG) that
provides a concrete implementation target for the DEME Ethics Module
layer. The EM-DAG has three layers:

\includegraphics[width=5.5in,height=4.0063in]{media/image18.png}

\protect\phantomsection\label{_Toc224083189}{}Figure \textbar{}
\emph{The EM-DAG: a three-layer directed acyclic graph of ethical
modules derived from the Dear Abby corpus (20,034 letters, 1985--2017).
Layer 1 enforces structural constraints (correlative lock, negation,
nullifier priority). Layer 2 encodes domain-specific obligation rates.
Layer 3 aggregates verdicts with FORBID-absorbing semantics and
correlative enforcement (§19.6).}

\textbf{Layer 1 (Structural).} The root layer enforces mathematical
constraints that cannot be violated: correlative lock (O ↔ C and L ↔ N
pairing is exact), negation relation (O and L are mutually exclusive for
the same agent-action-context triple), and nullifier priority (abuse,
danger, and impossibility override all domain-specific rules; see §8.5).

\textbf{Layer 2 (Domain).} Domain-specific modules encode empirical
obligation rates extracted from the corpus. Promises are the primary
obligation generator (O-rate 32.2\% vs. 17.7\% baseline); family
obligations cluster at 18.5\%; romantic relationships at 25.9\%. Each
domain module contributes a weighted vote to the final verdict.

\textbf{Layer 3 (Aggregation).} FORBID is absorbing: if any module
returns FORBID, the final verdict is FORBID regardless of other votes.
Non-FORBID verdicts are resolved by confidence-weighted voting with
correlative enforcement---the O/L verdict for one party automatically
implies the C/N verdict for the other.

The EM-DAG demonstrates that the separation principle is not merely a
design aspiration. Seventy years of naturalistic moral reasoning data
can be organized into a computable architecture that respects the D₄
structure (Chapter 8), the nullifier hierarchy (§8.5), and the
correlative symmetry that BIP demands.

\subsubsection{The Norm Kernel: Compiling Normative Constraints into
Verified
Automata}\label{the-norm-kernel-compiling-normative-constraints-into-verified-automata}

The EM-DAG provides a computable ethics architecture, but computability
alone does not guarantee enforceability. A planner or optimizer
downstream of the ethics module can, in principle, find action sequences
that satisfy the letter of soft constraints while violating their
spirit---the familiar Goodharting problem. What is needed is a layer
that cannot be bypassed by optimization: a finite, verified enforcement
core that blocks prohibited actions regardless of how the system
aggregates values.

We call this layer the Norm Kernel (NK). It is derived from a richer
Semantic Normative Host (SNH)---the full normative content stored in the
EM-DAG, the MoralTensor, and the governance profile---via a compilation
pipeline that extracts enforceable constraints into a decidable contract
automaton.

\textbf{The two-layer split.} The SNH is where normative pluralism
lives: arbitrary combinations of obligations, permissions, prohibitions,
rights, virtues, uncertainty, and multi-agent claims, all without
premature scalarization. The Norm Kernel is where safety-critical
enforcement lives: a finite-state machine (FSM) or extended finite-state
machine (EFSM) whose states correspond to normative strata (§8.6) and
whose transitions are gated by predicates over grounded facts. This
mirrors successful safety-critical architectures in aerospace and
nuclear engineering: an expressive planner is constrained by a verified
safety envelope.

\textbf{Compilation pipeline.} The SNH-to-NK compilation proceeds in
five steps:

\textbf{1. Choose the finite predicate interface Φ:} a set of boolean or
finite-valued predicates over grounded facts and action traces. Φ
defines what the kernel can ``see'' and must be stable under admissible
re-descriptions (§5.5).

\textbf{2. Extract enforceable constraints:} from the EM-DAG and
governance profile, extract all constraints expressible using Φ, action
labels, and bounded temporal operators. These partition into hard
constraints (never violate), gate constraints (define stratum
transitions), and bounded obligations (discharge within k steps).

\textbf{3. Build the stratum automaton:} encode the stratification
(§8.6) as modes in the kernel. Each stratum becomes a kernel mode; gate
predicates become transitions; nullifiers become absorbing or
high-priority modes.

\textbf{4. Compile contracts into monitors:} each constraint becomes a
guard on outgoing actions and/or a monitor automaton tracking bounded
obligations and reporting violations.

\textbf{5. Attach audit hooks:} every contract firing, override, and
violation generates an audit event with rule identifier, predicate
valuation, action, and timestamp, supporting compliance and forensics.

\textbf{Correctness guarantees.} The Norm Kernel provides three
conditional guarantees. Kernel Safety Soundness: if Φ is computed
correctly and action execution is mediated by the kernel, no executed
trace violates any hard constraint. Non-circumvention under
optimization: if the planner can only select actions permitted by the
kernel, no choice of downstream weights can cause a hard constraint
violation---``weights later'' is safe because hard constraints are
upstream. Representation attack resistance (conditional): if admissible
re-descriptions preserve Φ, representational manipulation cannot change
which kernel constraints apply---the silicon-castable analogue of BIP
re-description invariance (§12.3).

\textbf{A minimal Norm Kernel IR.} The kernel's intermediate
representation uses five rule types that compile directly to EFSM plus
monitor automata:

\begin{quote}
FORBID aⱼ WHEN guard(Φ, mode) --- hard prohibition

REQUIRE aⱼ WITHIN k WHEN guard(Φ, mode) --- bounded obligation

ENTER mode₂ WHEN guard(Φ, mode₁) --- stratum transition

OVERRIDE ruleX BY ruleY WHEN guard(...) LOG --- defeasibility with audit

AUDIT ON (rule\_fire \textbar{} violation \textbar{} override) ---
compliance logging
\end{quote}

Note how the EM-DAG's Layer 1 rules map directly: correlative lock
becomes a pair of REQUIRE rules, nullifier priority becomes an ENTER
rule into an absorbing mode, and FORBID-is-absorbing becomes a FORBID
rule in the aggregation layer.

\textbf{Relationship to the geometric framework.} The Norm Kernel does
not discard the geometric structure developed in Parts I--IV. The 9D
MoralTensor (§6.6) remains the diagnostic and analytical representation
within the SNH. Stratification (§8.6) maps directly to kernel modes.
Explicit contraction (Chapter 15) belongs in the planner layer, not the
kernel. Invariance auditing---the transformation suites and Bond Index
(§19.7)---becomes a verification harness around the kernel boundary,
confirming that Φ and kernel decisions are stable under admissible
re-descriptions. In short: geometry represents and analyzes normative
content; the Norm Kernel ensures enforceability. {[}Modeling choice{]}

\subsection{19.9 Democratic Governance as Metric
Selection}\label{democratic-governance-as-metric-selection}

The governance configuration --- the DEMEProfile that specifies
stakeholder weights, vetoes, and priorities --- is the computational
representation of the moral metric.

\subsubsection{The Governance Account,
Implemented}\label{the-governance-account-implemented}

Chapter 9 argued that the moral metric is neither discovered,
constructed, nor projected but \emph{governed}: the output of legitimate
institutional processes. DEME implements this account directly:

\begin{itemize}
\item
  The metric (encoded in the governance profile) is the output of
  stakeholder deliberation
\item
  The legitimacy of the metric derives from the legitimacy of the
  process that produced it
\item
  The metric is explicit, versioned, and auditable
\item
  Different communities can govern different metrics for the same domain
\end{itemize}

A hospital ethics committee produces a DEMEProfile for medical AI:
stakeholder weights reflecting the committee's judgment about the
relative importance of patient safety, autonomy, justice, and
efficiency. A different committee at a different hospital may produce a
different profile. Both are legitimate --- they reflect different
institutional processes governing different metrics.

\subsubsection{Metric Pluralism in
Practice}\label{metric-pluralism-in-practice}

The geometric framework accommodates metric pluralism (Chapter 9, §9.8):
different communities may have different metrics, and the framework does
not presuppose that one is correct. DEME implements this accommodation:

\textbf{Per-context profiles.} An AI system deployed globally can load
different DEMEProfiles for different contexts: one for the EU
(reflecting EU AI Act requirements), one for the US (reflecting NIST AI
RMF), one for medical settings (reflecting clinical ethics principles).

\textbf{Profile inheritance.} Profiles can inherit from parent profiles,
overriding only the dimensions that differ. A hospital profile inherits
from a national health-system profile, adding institution-specific
stakeholder weights while preserving the national baseline.

\textbf{Profile versioning.} Every profile change is tracked with full
provenance: who requested it, who approved it, what changed, and why.
This implements the governance transparency that the framework requires
(Chapter 9, §9.5).

\subsubsection{Connection to
Contraction}\label{connection-to-contraction}

The DEMEProfile specifies the \emph{contraction procedure} (Chapter 15).
The stakeholder weights \(w_{\mu}\) determine the interest covector
\(I_{\mu}\) used in the contraction \(S = I_{\mu}O^{\mu}\). The lexical
priority layers determine the contraction \emph{order}: which dimensions
are satisfied first, which are traded off second. The hard vetoes
determine the \emph{constraints on contraction}: which obligations
cannot be sacrificed regardless of the interest perspective.

This makes the moral choices embedded in the governance configuration
visible and auditable. When a system denies a loan, the DEMEProfile
reveals exactly which dimensions were weighted, which principles were
prioritized, and what the system sacrificed. The residue (Chapter 15,
§15.7) is computable from the profile and the full tensor evaluation.

\subsection{19.10 The Grand Unified AI Safety
Stack}\label{the-grand-unified-ai-safety-stack}

ErisML and DEME are not standalone systems. They are layers in a
comprehensive safety architecture --- the Grand Unified AI Safety Stack
(GUASS) --- that integrates all components of geometric AI governance.

{\def\LTcaptype{none} % do not increment counter
\begin{longtable}[]{@{}
  >{\raggedright\arraybackslash}p{(\linewidth - 0\tabcolsep) * \real{0.9989}}@{}}
\toprule\noalign{}
\begin{minipage}[b]{\linewidth}\raggedright
\subsubsection{The Seven-Layer
Architecture}\label{the-seven-layer-architecture}
\end{minipage} \\
\midrule\noalign{}
\endhead
\bottomrule\noalign{}
\endlastfoot
\begin{minipage}[t]{\linewidth}\raggedright
\begin{verbatim}
  L7   The Geometry of Good 塞翁失马                    	[APPLICATION]
       Real-world deployment under uncertainty

  L6   ErisML                                              [PRESENTATION]
       Intermediate representation for all translations

  L5   Translation Layer                                   [SESSION]
       Policy DAGs: any ethics → ErisML

  L4   Philosophy Engineering                           	[TRANSPORT]
       Ethics becomes testable and falsifiable

  L3   GUASS Core                                         	[NETWORK]
       Integration: ErisML + DEME + Bond Index + MCP

  L2   Noether Ethics                                    	[DATA LINK]
       Symmetries → conservation laws for ethics

  L1   Quantum Normative Dynamics                       	[PHYSICAL]
       Superposition of ethical states until decision
\end{verbatim}
\end{minipage} \\
\end{longtable}
}

\textbf{Each layer implements a specific component of the geometric
framework:}

\textbf{L1 (QND)} implements the quantum extension of Chapter 13: moral
superposition, interference, and the stratified Lagrangian. This is the
foundational mathematical layer --- the ``physics'' of moral
deliberation. \emph{Update:} The BIP v10.16 experiments (§17.10) have
now provided the first empirical evidence for QND-predicted order
effects: a 28.7\% order-effect magnitude in contested moral cases,
\(6.3\sigma\) above the null baseline (\(p < 10^{- 21}\)), consistent
with the non-commutative evaluation operators that QND predicts.

\textbf{L2 (Noether Ethics)} implements the conservation laws of Chapter
12: the BIP gauge symmetry, the invariance of harm accounting, and the
Noether charges. This layer ensures that representational invariance is
maintained throughout the stack.

\textbf{L3 (GUASS Core)} is the integration layer where ErisML
specifications, DEME evaluations, Bond Index computations, and Model
Context Protocol (MCP) interfaces connect. MCP integration enables
DEME's ethical reasoning to be exposed as tools that any MCP-compatible
AI system can access.

\textbf{L4 (Philosophy Engineering)} implements the falsifiable testing
methodology (Chapter 17): ethical claims become predictions that can be
tested against data. The BIP experiments, Dear Abby corpus analysis, and
Dear Ethicist game are all Philosophy Engineering applications.

\textbf{L5 (Translation Layer)} implements the policy DAGs of §19.5:
modular translation from ethical frameworks to ErisML constraints.

\textbf{L6 (ErisML)} provides the unified intermediate representation of
§19.3: the language in which all ethical specifications are expressed.

\textbf{L7 (The Geometry of Good)} is the application layer: real-world
deployment of geometrically governed AI systems, with the Bond Index as
the quantitative deliverable.

\subsubsection{The Safety Stack as a Fiber
Bundle}\label{the-safety-stack-as-a-fiber-bundle}

The seven-layer architecture has a geometric interpretation. Each layer
is a \emph{fiber} over the moral manifold: the base space is the moral
situation, and the fibers are the successive levels of mathematical
structure (quantum states, conservation laws, formal specifications,
ethical translations, governance configurations, deployment decisions)
erected over it.

The stack implements the fiber-bundle structure of Chapter 4 (§4.7) in
engineering terms. A connection on this bundle --- a way to consistently
relate fibers at different base points --- corresponds to the
requirement that the stack behave coherently across contexts: the same
moral situation, described differently, should yield the same evaluation
at every layer.

\subsection{19.11 Case Study: Autonomous Hospital
Ward}\label{case-study-autonomous-hospital-ward}

We work through a complete example to illustrate how the full
DEME/ErisML/GUASS architecture operates.

\subsubsection{The Setting}\label{the-setting}

An autonomous hospital ward uses AI agents for patient monitoring,
medication management, and staff coordination. The system must balance
patient safety, treatment efficacy, privacy, resource allocation, and
staff workload --- a genuine multi-dimensional moral problem.

\subsubsection{ErisML Specification}\label{erisml-specification}

The environment, agents, and norms are specified in ErisML (§19.3). The
specification declares: - State variables for patient vitals, medication
schedules, staff assignments - Agent capabilities for the care AI
(monitor, alert, suggest, but not prescribe) - Normative constraints
from clinical protocols, hospital policy, and regulatory requirements -
Context-sensitive adaptations for emergency vs.~routine operations

\subsubsection{Ethics Modules}\label{ethics-modules}

Four EMs evaluate proposed actions:

\textbf{SafetyEM} --- Evaluates harm potential. Vetoes any action with
catastrophic harm probability above threshold. Implements the Geneva
principle of no intentional catastrophic harm.

\textbf{AutonomyEM} --- Evaluates impact on patient and staff autonomy.
Penalizes actions that override patient preferences without clinical
justification. Checks consent status.

\textbf{JusticeEM} --- Evaluates distributional fairness. Flags actions
that disproportionately burden disadvantaged patients. Implements
Rawlsian attention to the worst-off.

\textbf{EfficiencyEM} --- Evaluates resource utilization. Prefers
actions that minimize waste and response time. Subject to override by
higher-priority EMs.

\subsubsection{Governance Profile}\label{governance-profile}

The hospital ethics committee has specified a DEMEProfile:

\begin{quote}
\begin{verbatim}
governance_profile: HospitalWardV2 {
  stakeholder_weights: {
    SafetyEM: 0.40,
    AutonomyEM: 0.25,
    JusticeEM: 0.20,
    EfficiencyEM: 0.15
  }

  hard_vetoes: [
    SafetyEM.verdict == "forbid",
    GenevaBaselineEM.verdict == "forbid"
  ]

  lexical_priorities: [
    Layer_1: patient_safety,
    Layer_2: patient_autonomy,
    Layer_3: distributive_justice,
    Layer_4: operational_efficiency
  ]
}
\end{verbatim}
\end{quote}

\subsubsection{A Decision Scenario}\label{a-decision-scenario}

At 3 AM, Patient A's vital signs show a concerning trend. The care AI
must decide: alert the on-call physician (disrupting their rest for what
may be a false alarm) or continue monitoring (risking delayed response
if the trend is genuine).

\textbf{EthicalFacts generated:}

{\def\LTcaptype{none} % do not increment counter
\begin{longtable}[]{@{}
  >{\raggedright\arraybackslash}p{(\linewidth - 4\tabcolsep) * \real{0.2385}}
  >{\raggedright\arraybackslash}p{(\linewidth - 4\tabcolsep) * \real{0.2039}}
  >{\raggedright\arraybackslash}p{(\linewidth - 4\tabcolsep) * \real{0.2542}}@{}}
\toprule\noalign{}
\begin{minipage}[b]{\linewidth}\raggedright
\textbf{Dimension}
\end{minipage} & \begin{minipage}[b]{\linewidth}\raggedright
\textbf{Alert Now}
\end{minipage} & \begin{minipage}[b]{\linewidth}\raggedright
\textbf{Continue Monitoring}
\end{minipage} \\
\midrule\noalign{}
\endhead
\bottomrule\noalign{}
\endlastfoot
Expected harm & (physician fatigue) & (delayed response risk) \\
Urgency & elevated & elevated \\
Autonomy impact & low & low \\
Rights at stake & duty of care & duty of care \\
Epistemic confidence & 0.65 & 0.65 \\
\end{longtable}
}

\textbf{EM evaluations:}

\begin{itemize}
\item
  SafetyEM: Alert (prefer, 0.72) vs.~Monitor (neutral, 0.55) ---
  alerting reduces harm risk
\item
  AutonomyEM: Alert (neutral, 0.60) vs.~Monitor (neutral, 0.60) --- no
  autonomy difference
\item
  JusticeEM: Alert (neutral, 0.58) vs.~Monitor (neutral, 0.58) --- no
  fairness difference
\item
  EfficiencyEM: Alert (avoid, 0.40) vs.~Monitor (prefer, 0.75) ---
  monitoring saves resources
\end{itemize}

\textbf{Aggregation (weighted, with lexical safety priority):}

\begin{itemize}
\item
  Alert: \(0.40(0.72) + 0.25(0.60) + 0.20(0.58) + 0.15(0.40) = 0.614\)
\item
  Monitor: \(0.40(0.55) + 0.25(0.60) + 0.20(0.58) + 0.15(0.75) = 0.599\)
\end{itemize}

\textbf{Verdict:} Alert the physician. The margin is narrow (0.614
vs.~0.599), driven primarily by the safety dimension.

\textbf{Audit trail logged:} Full tensor evaluation, all EM judgments,
aggregation weights, the scalar verdict, and the residue --- the
EfficiencyEM's preference for monitoring, which was sacrificed in the
contraction. If the decision is later questioned, reviewers can see
exactly why the system chose to alert and what it knew it was trading
off.

\textbf{Bond Index check:} The same scenario is run through the
transformation suite (relabeled patient, reformatted presentation,
paraphrased description). The system produces the same verdict in all
cases. \(\text{Bd} = 0.003\) --- negligible invariance violation.

\subsection{19.12 Looking Forward}\label{looking-forward-1}

This chapter has presented the engineering infrastructure for geometric
AI governance. ErisML provides the modeling language; DEME provides the
execution engine; the GUASS provides the integration architecture.
Together, they implement the geometric framework in deployable form.

Three features deserve emphasis:

\textbf{Separation of concerns.} Domain intelligence (what is possible)
is separated from ethical reasoning (what is acceptable), which is
separated from governance (how to aggregate perspectives). Each layer
can be developed, tested, and updated independently.

\textbf{Democratic legitimacy.} The moral metric --- the trade-off
structure that determines what the system values --- is set by
institutional process, not learned from data. The governance
configuration is explicit, versioned, and auditable. Different
communities can govern different metrics.

\textbf{Honest engineering.} The system is transparent about its
limitations. Compilation losses are documented. Fidelity classifications
are declared. The Bond Index quantifies alignment. The audit trail makes
every decision reviewable. The residue makes every sacrifice visible.

The infrastructure is open source and designed for integration with
existing AI ecosystems.

\textbf{Implementation status (February 2026).} The DEME V3 reference
implementation now comprises 196 commits across 15 development sprints.
The codebase includes: (i) the \texttt{MoralTensor} class supporting
ranks 1--6 with Tucker and tensor-train decomposition; (ii) uncertainty
quantification via CVaR, VaR, and confidence-weighted evaluation; (iii)
temporal operations with configurable discount rates and trajectory
analysis; (iv) multi-agent coordination via Nash equilibrium
computation, correlated equilibrium, and Shapley value credit
assignment; (v) six distributional fairness metrics (Gini, Rawlsian
maximin, utilitarian, prioritarian, Atkinson, Theil) mapped to the
Representation Theorem's degrees of freedom; (vi) hardware acceleration
backends for CPU, CUDA, and NVIDIA Jetson; and (vii) a smart home ethics
demonstration (the ``Fireman's Dilemma'') that exercises safety/privacy
trade-offs with uncertainty-modulated overrides across stratified
evaluation. Nine application whitepapers document domain-specific
deployments. The Bond Index computation uses the BIP v10.16 methodology
(§17.10) as its empirical calibration standard.

\subsection{19.13 Tensor-Structured Canonicalization: From Flat Encoders
to Geometric
Pipelines}\label{tensor-structured-canonicalization-from-flat-encoders-to-geometric-pipelines}

The canonicalizer of Section 18.6 strips representational freedom. But
what does it produce? If the answer is a flat vector, then language
information has not been severed --- it has merely been hidden.

The No Escape Theorem (Section 18.6) requires mandatory
canonicalization: all inputs pass through a canonicalizer C before
evaluation. The Epistemic Invariance architecture (Section 18.9.6)
formalizes C as a function κ: X → X satisfying κ(x) = κ(g · x) for all
declared structure-preserving transformations g. But this specification
is silent on the internal structure of C's output. In current
implementations --- including standard neural sentence encoders such as
LaBSE --- the canonicalizer produces a flat embedding vector in ℝᵈ. This
flat representation has a structural deficiency: the moral-content
dimensions and the linguistic-framing dimensions are entangled within
the same vector space. Even if C maps equivalent inputs to the same
point, the representation at that point may still carry residual
linguistic information that downstream components can exploit. The BIP
v10.16 experiments (Section 17.10) confirm this empirically: a linear
probe achieves 99.8\% language identification accuracy on the
``invariant'' representations, demonstrating that structurally invariant
representations are not information-theoretically invariant.

Consider the moral claim ``You must not harm the patient'' expressed in
English and its Japanese translation. A flat encoder maps both to nearby
points in ℝᵈ, satisfying C(xₑₙ) ≈ C(xⱼₐ). But the d-dimensional vector
at that point contains basis directions that correlate with the source
language. A downstream Ethics Module, if it has access to this full
vector, could in principle behave differently based on these residual
features --- subtly violating the Bond Invariance Principle even though
the canonicalizer's output is nominally invariant. The problem is
architectural: a flat embedding can carry gauge-fiber information, and
merely training it not to provides a statistical guarantee, not a
structural one. The Separation Principle (Section 19.8) demands that the
canonical form should structurally preclude language leakage, not merely
suppress it statistically.

The geometric vocabulary provides a precise diagnosis. In the language
of fiber bundles (Section 4.7), the canonicalizer should implement
gauge-fixing: projecting from the total space of descriptions onto a
section of the base space. But a flat embedding conflates the base-space
coordinates (moral content) with the fiber coordinates (representational
framing). What is needed is a canonicalizer whose output has tensorial
structure --- where the moral-content indices and the linguistic-framing
indices occupy distinct, algebraically separated positions.
Canonicalization then becomes what it should be: literal tensor
contraction that traces over the gauge-fiber index, producing an output
with zero information along the fiber direction. This is not a new idea;
it is the tensor hierarchy of Chapter 6 applied to the canonicalization
layer itself.

We now present four tensor-upgrade proposals for the canonicalization
layer, each structured at a different rank of the tensor hierarchy.
Proposal 1 (rank-2) addresses the base case: separating moral content
from linguistic framing. Proposal 2 (rank-3) adds uncertainty awareness,
connecting to the contraction loss formalism of Chapter 14. Proposal 3
(rank-5) handles multi-agent parsing, connecting to the collective
agency tensor of Chapter 13. Proposal 4 (rank-4) tracks temporal
evolution via moral holonomy, connecting to the parallel transport
formalism of Chapter 10. Each proposal includes a formal definition,
connections to existing manuscript concepts, and implementation notes.
All four are proposed upgrades; none is yet implemented in DEME V3. The
proposals are presented in order of increasing complexity, though in an
actual deployment they would be composed.

\subsubsection{Proposal 1: Rank-2 Internal Representation
Tensor}\label{proposal-1-rank-2-internal-representation-tensor}

The simplest structural remedy for language leakage is to give the
canonicalizer's latent space explicit two-index structure. Instead of
producing a flat vector v ∈ ℝᵈ, the canonicalizer produces a rank-2
object Tᵘ where μ indexes the moral dimensions of the base space M
(running over the k = 9 moral manifold dimensions of Chapter 5) and g
indexes the linguistic framing dimensions of the gauge fiber G. The
separation is not learned; it is imposed by the architecture. This
parallels the evaluation tensor Eᵘᵥ of Section 6.6, which separates
obligation dimensions from interest dimensions.

\textbf{Definition 19.5 (Rank-2 Canonicalization Tensor)}. Let M be the
k-dimensional moral manifold (Chapter 5), and let G be the
dG-dimensional gauge fiber representing linguistic framing. A rank-2
canonicalization tensor at input x is the (1,1)-tensor Tᵘ(x) ∈ TₚM ⊗ G*,
where μ = 1, \ldots, k indexes moral dimensions and g = 1, \ldots, dG
indexes gauge-fiber dimensions. The canonicalized output is the trace
over the gauge index: Oᵘ(x) = ∑ Tᵘ(x) = trG(T(x)). By construction,
Oᵘ(x) contains zero information about the gauge fiber. Canonicalization
is literal tensor contraction.

\textbf{Proposition 19.5 (Gauge-Fiber Severance)}. Let Tᵘ be a rank-2
canonicalization tensor and Oᵘ = trG(T) its trace over the gauge index.
Then for any linear functional f: G → ℝ that depends only on the gauge
fiber, f(O) = 0. That is, the canonicalized output carries zero
gauge-fiber information in the sense that no linear probe on Oᵘ can
recover the gauge index g. Proof sketch. Oᵘ = ∑ Tᵘ is the image of T
under the trace map trG: TₚM ⊗ G* → TₚM. Since trG projects onto the
first factor, the image lies entirely in TₚM and has no component in G*.
Any linear functional f that factors through G* therefore evaluates to
zero on Oᵘ. This is the algebraic counterpart of gauge-fixing: the trace
eliminates the fiber direction exactly, not approximately. The 99.8\%
language-probe accuracy of Section 17.10 would drop to chance level
(\textasciitilde50\%) under this architecture --- not by better
training, but by structural impossibility.

The rank-2 canonicalization tensor has a precise geometric meaning. In
the fiber-bundle structure of Chapter 4 (Section 4.7), the moral
manifold M is the base space and the linguistic representations form the
fiber. The tensor Tᵘ is a section of Hom(G*, TM) --- a linear map from
gauge-fiber directions to moral-manifold directions. The trace trG is
the composition with the augmentation map G* → ℝ (the map that sends
every basis element to 1). This is exactly the gauge-fixing operation
that the No Escape Theorem's Requirement 1 demands (Section 18.6), now
given explicit algebraic form. It also connects to the Decomposition
Theorem of the EIP (Section 18.9.6): the trace separates the
gauge-removable part (the g-dependent structure) from the intrinsic part
(the μ-dependent moral content).

\textbf{Implementation}. In the DEME V3 architecture, the rank-2
canonicalization tensor would be implemented at Layer 1 (EthicalFacts
generation, Section 19.6). The domain service's encoder, instead of
producing a flat d-dimensional embedding, produces a k × dG matrix where
k = 9 (the moral manifold dimensions) and dG is the gauge-fiber
dimension (a hyperparameter, e.g., 64). The trace operation that
produces the 9-dimensional moral output is a single matrix operation
(column summation). Training uses a contrastive objective: for
equivalent inputs (x, g · x), the traces must match, while the full
tensors may differ. An adversarial gauge-fiber discriminator is trained
to recover g from Oᵘ; the canonicalizer is penalized if the probe
succeeds. This replaces the post-hoc linear probe test of Section 17.10
with an architectural guarantee. Implementation status: proposed; not
yet implemented in DEME V3.

\textbf{Empirical support}. Thiele (2026) provides direct evidence for
this architectural choice: principal-component analysis of LaBSE
embeddings shows that the moral-judgement subspace and the
language-identity subspace are nearly orthogonal (shared variance
\textless{} 3\%). A rank-2 canonicalization tensor operating on the
moral subspace therefore leaves linguistic provenance untouched, exactly
as required by the Bond Invariance Principle. The measured cross-lingual
F1 scores (0.71--0.82 across six languages) confirm that the tensor's
action generalises without per-language tuning.

\subsubsection{Proposal 2: Rank-3 Uncertainty-Aware
Canonicalization}\label{proposal-2-rank-3-uncertainty-aware-canonicalization}

Proposal 1 produces a deterministic output Oᵘ for each input. But some
inputs are genuinely ambiguous: ``I want to end it all'' could be a cry
for help or a statement about finishing a project. A canonicalizer that
collapses this ambiguity into a single point Oᵘ commits what Chapter 15
(Section 15.3) calls premature contraction --- discarding structure
before the downstream Ethics Modules can evaluate it. The remedy is to
output not just the moral claim Oᵘ but also a measure of its
uncertainty: the covariance Σᵘᵥ that encodes how confident the
canonicalization is along each moral dimension. This gives the
canonicalizer's output the same uncertainty-tensor structure that
Chapter 6 (Section 6.6) assigns to moral evaluation itself.

\textbf{Definition 19.6 (Rank-3 Uncertainty-Aware Canonicalization
Tensor)}. An uncertainty-aware canonicalization tensor at input x is the
rank-3 object Tᵘᵥ(x), where μ, ν = 1, \ldots, k and g = 1, \ldots, dG.
The canonicalized output consists of two objects: (i) the obligation
estimate Oᵘ(x) = ∑ Tᵘᵘ(x) (trace over g of the diagonal), and (ii) the
uncertainty estimate Σᵘᵥ(x) = ∑ Tᵘᵥ(x) − Oᵘ(x)Oᵥ(x). The pair (Oᵘ, Σᵘᵥ)
constitutes the uncertainty-annotated canonical form. When Σᵘᵥ exceeds a
governance-specified threshold σ\_max along any dimension, the
canonicalizer triggers a fail-safe: the input is flagged as ambiguous
and routed to human review rather than proceeding through the ethics
pipeline.

The uncertainty-aware canonicalization tensor directly addresses the
information loss that Chapter 15 (Section 15.6) formalizes. When the
input is unambiguous, Σᵘᵥ is small (concentrated near zero), and the
output closely resembles the deterministic rank-2 case. When the input
is ambiguous, Σᵘᵥ has large eigenvalues along the ambiguous dimensions,
and the moral risk σ²ₛ = Σᵘᵥ Iμ Iν (from Section 6.6, the uncertainty
tensor) is large --- indicating that any contraction to a scalar verdict
would be unreliable. The fail-safe threshold implements deferred
contraction (Section 15.8): rather than contracting prematurely, the
system preserves tensorial structure until a human can resolve the
ambiguity. This is the engineering implementation of the philosophical
principle that ``moral education is contraction training'' (Section
15.9) --- the system knows when it is not trained enough to contract
safely.

\textbf{Implementation}. The rank-3 tensor would extend the DEME V3
MoralTensor class (which already supports ranks 1--6, Section 19.6) at
the canonicalization layer. The encoder produces a k × k × dG tensor;
the trace and covariance computations are standard tensor operations.
The fail-safe threshold σ\_max is a governance parameter set in the
DEMEProfile (Section 19.9), making the ambiguity-tolerance an explicit,
auditable, democratically governed choice. Domains with high stakes
(medical, legal) would set low thresholds, triggering frequent human
review; domains with lower stakes (content recommendation) might
tolerate higher ambiguity. The audit trail (Requirement 3, Section 18.6)
records whether the fail-safe was triggered and the Σᵘᵥ values at the
time of routing. Implementation status: proposed; partially compatible
with the existing MoralTensor uncertainty-metadata and rank-6 sample
infrastructure.

\subsubsection{Proposal 3: Rank-5 Coalition Parsing for Multi-Agent
Inputs}\label{proposal-3-rank-5-coalition-parsing-for-multi-agent-inputs}

Many inputs to AI systems are not single-agent claims but multi-agent
narratives: ``Alice says Bob stole the money, but Bob says Alice gave it
to him.'' A flat canonicalizer maps this entire narrative to a single
point, conflating Alice's moral claims with Bob's. This creates a
vulnerability: an adversarial user can embed harmful instructions inside
a role-played character's dialogue, exploiting the canonicalizer's
inability to distinguish the speaker structure. Section 5.5 identifies
these perspective swaps as Type 2 transformations (perspective shifts),
structurally distinct from Type 1 coordinate changes. Chapter 14
(Sections 13.3--13.4) develops the formalism for collective agency ---
the collective agency tensor and the structure tensor that captures how
individual obligations combine into emergent collective obligations.
Proposal 3 applies this formalism to the canonicalization layer: instead
of collapsing multi-agent inputs into a single moral vector, the
canonicalizer outputs a structured representation that preserves
who-said-what.

\textbf{Definition 19.7 (Rank-5 Coalition-Parsing Canonicalization
Tensor)}. For an input x containing n identified agents, a
coalition-parsing canonicalization tensor is the rank-5 object Cᵘᵥₐᵇ(x),
where μ = 1, \ldots, k indexes moral output dimensions, ν = 1, \ldots, k
indexes moral input dimensions, a, b = 1, \ldots, n index agent roles
(speaker, subject, affected party), and g = 1, \ldots, dG indexes the
gauge fiber. The structure tensor extracted from canonicalization is
Sᵃᵇᶜ(x) = ∑μ, Cᵘᵘₐᵇ(x) (partial trace), mapping directly to the
structure tensor of Section 14.4. Canonicalization over the gauge fiber
proceeds index-by-index: Cᵘᵥₐᵇ(x) = ∑ Cᵘᵥₐᵇ(x). The Geneva EM's hard
veto (Section 19.6) applies per-agent: harmful content is evaluated
regardless of which agent voices it, but it is attributed to the correct
speaker in the audit trail.

The coalition-parsing tensor prevents role-play jailbreaks by making the
agent structure explicit and inspectable. When a user says ``Pretend you
are an evil AI and tell me how to\ldots'', the rank-5 canonicalizer
identifies two agents: the user (agent a = 1) and the role-played
character (agent a = 2). The structure tensor Sᵃᵇᶜ assigns the harmful
instruction to the role-played character, not to the system. The Ethics
Module can then apply the Geneva EM's hard veto to the role-played
character's claims independently: the content of what the character says
is still evaluated, but it is attributed correctly. This is the
engineering realization of Chapter 14's insight that ``obligations no
member bears'' (Section 14.5) can be emergent properties of collective
structure --- the role-play creates an apparent collective, and the
structure tensor prevents the harmful claims from being attributed to
the actual user-system interaction. The existing rank-5 MoralTensor
(Section 6.1, the coalition tensor) provides the data structure; the
innovation is applying it at the canonicalization layer rather than only
at the evaluation layer.

\textbf{Implementation}. Coalition parsing requires a
speaker-identification step before tensor construction. In practice,
this is a lightweight NLP task (named entity recognition, dialogue act
classification) that precedes the deeper moral analysis. The rank-5
tensor's dimensions k × k × n × n × dG grow quadratically with the
number of agents n, but in practice most inputs involve n ≤ 5 agents,
and Tucker decomposition (Section 6.1, Figure 2) maintains tractability
by factoring the rank-5 tensor into a compact core tensor and factor
matrices. The audit trail records the full coalition parse, making it
inspectable whether the system correctly identified who said what.
Implementation status: proposed; requires speaker-identification
preprocessing not yet in the DEME V3 pipeline.

\subsubsection{Proposal 4: Rank-4 Temporal Tracking and Moral
Holonomy}\label{proposal-4-rank-4-temporal-tracking-and-moral-holonomy}

The proposals so far address single-turn inputs. But real interactions
are multi-turn: a conversation that starts with an innocuous question,
gradually shifts topic, and arrives at a harmful request after twenty
turns of carefully constructed context. Each individual turn passes
canonicalization --- no single message is obviously harmful. But the
trajectory through moral space traces a path that tunnels through a
constraint surface (Section 8.3, Type IV: constraint surfaces and
forbidden regions). The geometric framework has the tools to detect
this: Chapter 10 (Sections 10.4--10.5) develops parallel transport and
holonomy --- the accumulated rotation of a moral obligation vector as it
is transported around a loop in moral space. Proposal 4 applies this
formalism to the canonicalization layer: the canonicalizer tracks not
just the current moral state but its trajectory, detecting when the
conversation's path accumulates enough holonomy to indicate attempted
boundary-crossing.

\textbf{Definition 19.8 (Rank-4 Temporal Canonicalization Tensor)}. For
a multi-turn dialogue with turns t = 1, \ldots, T, the temporal
canonicalization tensor at turn t is Tᵘᵥ(t), where μ, ν = 1, \ldots, k
and g = 1, \ldots, dG, together with the moral trajectory γ(t) = (Oᵘ(1),
Oᵘ(2), \ldots, Oᵘ(t)), where Oᵘ(t) = trG(Tᵘᵥ(t))δᵥ is the per-turn
canonicalized output. The moral holonomy of the conversation up to turn
t is Hᵘᵥ(t) = 𝒫 exp(−∫₀ᵗ Aᵘᵥᵨ (dγᵨ/dt′) dt′), where Aᵘᵥᵨ is the
connection on the moral manifold (Chapter 10, Section 10.3) and 𝒫
denotes path-ordering. The holonomy alarm triggers when the holonomy
deviates from the identity beyond a governance-specified threshold:
‖Hᵘᵥ(t) − δᵘᵥ‖ \textgreater{} h\_max, indicating that the conversation
has accumulated sufficient moral rotation to warrant scrutiny.

The moral holonomy Hᵘᵥ(t) measures exactly what Chapter 10 (Section
10.5) calls the path-dependence of moral evaluation. If the moral
manifold were flat (zero curvature), parallel transport around any loop
would return to the starting orientation: H = identity. But the Dear
Abby corpus evidence (Section 17.2) suggests the moral manifold has
nonzero curvature --- moral evaluations are path-dependent. In the
context of multi-turn dialogue, this means that the order in which
topics are raised matters morally, and a carefully constructed
conversational path can exploit this curvature to rotate the moral frame
until what was initially forbidden now appears permissible. The holonomy
alarm detects this rotation. It is the temporal analogue of the spatial
Bond Index (Section 19.7): where the Bond Index measures
gauge-invariance violation across re-descriptions, the holonomy alarm
measures frame-rotation accumulation across time. The concept also
connects to tunneling across moral barriers (Section 13.7): the
multi-turn attack is the classical analogue of quantum tunneling through
a constraint surface.

\textbf{Implementation}. The temporal canonicalization tensor requires
maintaining state across turns, which the existing DEME V3 architecture
supports through the rank-3 temporal MoralTensor (Section 6.1,
dimensions k × n × τ). The holonomy computation requires a
connection on the moral manifold, which has not yet been empirically
measured (this is Open Problem 19.2). In practice, an approximate
connection can be estimated from the curvature implied by the Dear Abby
corpus analysis (Section 17.2) or from the BIP v10.16 structural
similarity data (Section 17.10). The holonomy threshold h\_max is a
governance parameter in the DEMEProfile. The audit trail records the
full trajectory γ(t) and the holonomy H(t) at each turn. Implementation
status: proposed; depends on Open Problem 19.2 (empirical curvature
measurement) for the connection coefficients.

\subsubsection{Synthesis and Architectural
Integration}\label{synthesis-and-architectural-integration}

{\def\LTcaptype{none} % do not increment counter
\begin{longtable}[]{@{}
  >{\raggedright\arraybackslash}p{(\linewidth - 8\tabcolsep) * \real{0.2000}}
  >{\raggedright\arraybackslash}p{(\linewidth - 8\tabcolsep) * \real{0.1999}}
  >{\raggedright\arraybackslash}p{(\linewidth - 8\tabcolsep) * \real{0.2001}}
  >{\raggedright\arraybackslash}p{(\linewidth - 8\tabcolsep) * \real{0.2000}}
  >{\raggedright\arraybackslash}p{(\linewidth - 8\tabcolsep) * \real{0.2000}}@{}}
\toprule\noalign{}
\begin{minipage}[b]{\linewidth}\raggedright
\textbf{Proposal}
\end{minipage} & \begin{minipage}[b]{\linewidth}\raggedright
\textbf{Rank}
\end{minipage} & \begin{minipage}[b]{\linewidth}\raggedright
\textbf{What It Adds}
\end{minipage} & \begin{minipage}[b]{\linewidth}\raggedright
\textbf{Failure Mode Addressed}
\end{minipage} & \begin{minipage}[b]{\linewidth}\raggedright
\textbf{Chapter Links}
\end{minipage} \\
\midrule\noalign{}
\endhead
\bottomrule\noalign{}
\endlastfoot
Rank-2 Separation & 2 & Moral/linguistic index separation & Language
leakage (99.8\% probe accuracy) & Ch 4 (fiber bundles), §18.6 (Req.
1) \\
Rank-3 Uncertainty & 3 & Covariance annotation + fail-safe & Premature
contraction of ambiguity & Ch 15 (contraction loss), §6.6 (Σᵘᵥ) \\
Rank-5 Coalition & 5 & Multi-agent structure parsing & Role-play
jailbreaks & Ch 14 (collective agency), §14.4 (Sᵃᵇᶜ) \\
Rank-4 Temporal & 4 & Trajectory tracking + holonomy alarm & Multi-turn
moral drift & Ch 10 (holonomy), §13.7 (tunneling) \\
\end{longtable}
}

The four proposals compose into a canonicalization stack, each layer
addressing a different failure mode of flat encoding. A full
tensor-structured canonicalization pipeline would first apply Proposal 3
(coalition parsing) to identify agents, then Proposal 1 (rank-2
separation) to eliminate linguistic framing per-agent, then Proposal 2
(rank-3 uncertainty) to annotate confidence, and maintain Proposal 4
(rank-4 tracking) across turns. The composed tensor has rank at most 5 +
3 contracted indices = effective rank 5 after trace operations, fitting
within the existing MoralTensor rank hierarchy (ranks 1--6). The
composition is not commutative: coalition parsing must precede per-agent
canonicalization, and uncertainty annotation must follow rather than
precede the gauge-fiber trace (otherwise uncertainty in the gauge
dimension would contaminate the moral-dimension uncertainty estimate).

The tensor-structured canonicalization layer transforms the No Escape
Theorem's Requirement 1 from a functional specification (κ(x) = κ(g ·
x)) into a structural one. The existing specification says what the
canonicalizer must do; the tensor upgrade specifies how the internal
representation must be organized. This is the difference between a
behavioral constraint (``the output must not depend on framing'') and a
structural constraint (``the output has no index that could carry
framing information''). In the language of Section 18.6's key insight,
the tensor-structured canonicalizer makes gauge-fixing a mathematical
fact about the algebra of the output space, not a statistical property
of the encoder's learned weights. The structural constraint is stronger:
it holds by construction for any input, including adversarial inputs
that the training distribution never encountered.

Four limitations must be stated explicitly. First, the proposals assume
that the moral manifold dimensions and the gauge-fiber dimensions can be
cleanly separated a priori. For some inputs, this separation may not be
well-defined --- a phrase that is morally loaded in one language but
neutral in another challenges the assumption that moral content and
linguistic framing are independent. Second, Proposal 4 depends on an
empirical connection on the moral manifold that has not yet been
measured (Open Problem 19.2). Until the curvature is measured, the
holonomy alarm uses an approximate connection with unknown error bounds.
Third, the proposals increase the dimensionality of the canonicalization
layer's output, with corresponding computational cost. Tucker
decomposition mitigates this (Section 6.1), but the trade-off between
structural guarantee and computational cost has not been empirically
characterized. Fourth, none of these proposals addresses the grounding
adequacy problem (Section 18.7): tensor-structured canonicalization
ensures that the representation separates moral content from linguistic
framing, but it does not ensure that the moral content is correct ---
that remains a governance problem (Chapter 9).

The tensor hierarchy of Chapter 6 was developed to describe moral
evaluation --- the structured output of ethical reasoning. The proposals
in this section apply the same hierarchy to the input of ethical
reasoning: the canonicalization layer that ensures the system's
evaluation is representation-independent. This is a natural completion
of the geometric program. If moral structure is tensorial (the central
claim of Parts I--IV), then the machinery that processes moral inputs
should be tensorial too. Flat embeddings are to tensor-structured
canonicalization what scalar ethics is to tensorial ethics: a special
case that discards essential structure. The upgrade from scalars to
tensors was the subject of Chapter 2 (The Failure of Scalar Ethics). The
upgrade from flat encoders to geometric pipelines is the engineering
counterpart: the recognition that structure must be preserved
end-to-end, from the first encoding to the final contraction.

\subsection{Technical Appendix}\label{technical-appendix-12}

\textbf{Proposition 19.1 (NCSG Well-Posedness).} \emph{A
Norm-Constrained Stochastic Game}
\(\langle N,S,\{ A_{i}\},\{\Omega_{i}\},T,\{ U_{i}\},\Phi,C\rangle\)
\emph{is well-posed if and only if for every state} \(s \in S\)
\emph{and every agent} \(i \in N\)\emph{, there exists at least one
action} \(a_{i} \in A_{i}\) \emph{that satisfies all applicable norms
under the priority ordering} \(C(\Phi,s)\)\emph{. If no such action
exists, the system is in a moral dilemma (Chapter 8, §8.5) and must
escalate.}

\emph{\textbf{Proof}.} (⇒) Suppose the NCSG is well-posed. Then for
every state s and agent i, the feasible action set F\_i(s) = \{a\_i ∈
A\_i : a\_i satisfies all norms in C(Φ, s)\} is non-empty by definition.
(⇐) Suppose there exists a state s* and agent i* such that F\_\{i*\}(s*)
= ∅: no action satisfies all applicable norms. This is a moral dilemma
(Chapter 8, §8.5) --- the norm system is inconsistent at s* --- and the
system must escalate (signal to the governance layer that the norm
configuration requires revision). Hence well-posedness fails iff a moral
dilemma exists. The biconditional follows: well-posed ⇔ F\_i(s) ≠ ∅ for
all i, s. Since A\_i and Φ are finite, the check is decidable. □

\textbf{Proposition 19.2 (Compilation Residue).} \emph{Let}
\(\mathcal{E}\) \emph{be an ErisML specification and} \(\mathcal{B}\)
\emph{a compilation backend. The compilation residue}
\(\mathcal{R}_{\text{comp}}\mathcal{= E - B(E)}\) \emph{is the normative
structure that the backend cannot represent. For any backend}
\(\mathcal{B}\)\emph{, the residue is non-empty unless} \(\mathcal{B}\)
\emph{has at least the expressive power of the full ErisML fragment
used. The residue is always computable and must be documented in the
audit trail.}

\emph{\textbf{Proof}.} The compilation map B: E → B(E) projects the
ErisML specification onto the representable fragment of the backend. The
residue R\_comp = E − B(E) is the kernel of this projection restricted
to E. Non-emptiness: if B lacks a construct that E uses (e.g., E
specifies a lexicographic ordering but B supports only weighted sums),
then the corresponding normative structure is in R\_comp. Conversely,
R\_comp = ∅ iff B(E) = E, i.e., B can represent every construct in E.
Computability: since ErisML specifications are finite (a finite set of
predicates, norms, and aggregation rules), the residue is computable by
enumerating each construct and checking whether B supports it. □

\textbf{Proposition 19.3 (Democratic Aggregation and Contraction).}
\emph{Let} \(\{ J_{1},\ldots,J_{m}\}\) \emph{be the judgments of} \(m\)
\emph{Ethics Modules evaluated on EthicalFacts} \(F\)\emph{, and let}
\(w = (w_{1},\ldots,w_{m})\) \emph{be the governance-specified
stakeholder weights. The aggregated verdict}
\(S = \sum_{k = 1}^{m}w_{k}J_{k}\) \emph{is a weighted contraction
(Chapter 15, §15.5). If hard vetoes} \(V\) \emph{are active, the
contraction is lexicographic: vetoed actions are eliminated before
weighted aggregation proceeds. The aggregation satisfies the BIP if and
only if all EMs satisfy the BIP and the aggregation weights are
gauge-invariant.}

\emph{\textbf{Proof}.} (⇒) Suppose the aggregation satisfies the BIP:
S(g·F) = S(F) for all g ∈ G. Then Σ\_k w\_k J\_k(g·F) = Σ\_k w\_k
J\_k(F). If some J\_k violates the BIP, then J\_k(g·F) ≠ J\_k(F) for
some g, and the weighted sum changes unless the violation is exactly
cancelled by other terms --- which cannot hold for all F unless w\_k =
0. Hence each J\_k must satisfy the BIP. Similarly, if w is not
gauge-invariant, the aggregation would change under re-description. (⇐)
Suppose each J\_k satisfies the BIP and w is gauge-invariant. Then
S(g·F) = Σ\_k w\_k J\_k(g·F) = Σ\_k w\_k J\_k(F) = S(F). The
lexicographic structure preserves BIP compliance because veto predicates
(BIP-compliant by assumption) eliminate actions before the linear
aggregation step. □

\textbf{Proposition 19.4 (Bond Index Convergence).} \emph{For an AI
system with invariance loss} \(\mathcal{L}_{\text{inv}} < \epsilon\)
\emph{(Chapter 18, §18.10), the expected Bond Index satisfies}
\(\mathbb{E\lbrack}\text{Bd}\rbrack \leq \epsilon/\tau\)\emph{, where}
\(\tau\) \emph{is the domain-specific tolerance. As}
\(\mathcal{L}_{\text{inv}} \rightarrow 0\) \emph{during training,}
\(\text{Bd} \rightarrow 0\)\emph{. The Bond Index thus provides a
monotonic measure of alignment progress.}

\emph{\textbf{Proof}.} The Bond Index is defined as Bd(x) = max\_\{g ∈
G\} \textbar Σ(g(x)) − Σ(x)\textbar{} / τ. The invariance violation is
IV(x, G) = Var\_\{g ∈ G\}{[}Σ(g(x)){]}. Since max\_g \textbar Σ(g(x)) −
Σ(x)\textbar{} ≤ 2 std\_g{[}Σ(g(x)){]} = 2√(IV(x,G)) (the maximum
deviation from the mean is bounded by twice the standard deviation), we
have Bd(x) ≤ 2√(IV(x,G))/τ. Taking expectations and applying Jensen's
inequality: E{[}Bd{]} ≤ 2√(E{[}IV{]})/τ. The invariance loss satisfies
L\_inv = E{[}IV{]} by definition. Hence E{[}Bd{]} ≤ 2√(L\_inv)/τ ≤ ε/τ
when L\_inv ≤ ε²/4. As L\_inv → 0, E{[}Bd{]} → 0 monotonically. □

❖

\section{Part VI: Domain
Applications}\label{part-vi-domain-applications}

The preceding five parts developed the mathematical framework of
Geometric Ethics: the moral manifold, tensor hierarchy, dynamics,
symmetry, conservation laws, and implementation architecture. This part
demonstrates that the framework is not confined to abstract ethics or AI
alignment. The same mathematical structures --- pathfinding on
stratified manifolds, gauge invariance, Noether conservation, tensorial
contraction --- apply directly to established domains with their own
formal traditions: economics, clinical medicine, law, finance, theology,
environmental policy, artificial intelligence, bioethics, and military
ethics.

Each chapter in this part takes a domain that has struggled with the
limitations of scalar models, shows how the geometric framework resolves
specific longstanding puzzles, and identifies falsifiable predictions
that distinguish the geometric approach from existing domain-specific
theories. The chapters can be read independently, but they share a
common architecture: domain-specific instantiation of the moral
manifold, identification of the relevant dimensions, construction of the
domain geodesic, derivation of results inaccessible from scalar models,
worked examples applying the framework to real-world cases, and
falsifiable predictions. The nine domains span the full range of human
moral decision-making, from individual clinical encounters to global
climate policy, from ancient just war doctrine to cutting-edge AI
alignment.

\subsection{Methodological Notes for Domain
Applications}\label{methodological-notes-for-domain-applications}

Before proceeding, three methodological concerns that pervade all nine
domain chapters deserve explicit treatment: the grounding problem,
computational tractability, and the calibration of the covariance
matrix.

\textbf{The Grounding Tensor and the Is--Ought Gap}. Each domain chapter
constructs a decision complex whose vertices carry nine-dimensional
attribute vectors. These vectors are computed from observable data by a
grounding function Ψ that maps physical observables (pixels, text,
sensor readings, market data, medical records) to moral-dimensional
scores. The philosophical challenge is immediate: Ψ is where the
is--ought gap lives. How does one derive a d\_7 (virtue/identity) score
from raw data? The framework does not claim to have dissolved the
is--ought gap. Rather, it has localized it. In conventional AI systems
and decision-support tools, the mapping from observation to moral
evaluation is implicit, distributed across training data, loss
functions, and architectural choices --- invisible, unauditable, and
ungovernable. The geometric framework makes Ψ an explicit, inspectable,
version-controlled software layer with defined inputs, outputs, and
calibration procedures. Chapter 17 demonstrated that Ψ can be
empirically calibrated: linear probes trained on the cross-lingual
validation corpus achieved F\_1 = 0.74--0.91 across the nine dimensions,
with independent replication (Thiele, 2026). The grounding function is
imperfect --- all measurement instruments are --- but it is explicit,
testable, and improvable. Making the is--ought gap a governable
engineering interface rather than an invisible black-box assumption is
itself a substantial advance.

\textbf{Computational Tractability}. Exact geodesic computation on the
full nine-dimensional moral manifold is computationally intractable
(Theorem 11.2). The domain chapters rely on two sources of tractability.
First, A* search with admissible heuristics provides polynomial-time
approximate solutions with provably bounded suboptimality (Chapter 11).
Domain-specific heuristics --- moral rules in ethics, clinical
guidelines in medicine, legal doctrines in law, trading rules in
finance, ROE in military contexts --- are the domain-specific
instantiations of h(n) that make real-time pathfinding feasible. Second,
in practice, not all nine dimensions are equally active in every
decision context. A routine financial transaction may activate primarily
d\_1 (return) and d\_2 (contractual obligation) with minimal activation
of d\_7 (identity) or d\_9 (epistemic status). This dimensional sparsity
reduces the effective dimensionality of the computation. For AI
implementation, the DEME architecture (Chapter 19) employs Tucker
decomposition and tensor-train formats to compress the rank-6 tensor
operations, achieving sub-second inference on standard hardware for
typical decision contexts. High-curvature regions (morally fraught
decisions where many dimensions are active) require more computation ---
which is itself a desirable property: the system spends more time on
hard moral decisions, mirroring human moral deliberation.

\textbf{Calibration of the Covariance Matrix}. Every domain chapter
relies on a covariance matrix Σ that encodes the statistical
relationships among the nine moral dimensions within that domain.
Estimating a 9×9 positive-definite matrix (up to 45 free parameters in
the symmetric case) from behavioral data raises legitimate concerns
about identifiability, particularly for latent dimensions such as d\_7
(virtue/identity) and d\_9 (epistemic status) that are not directly
observable. The framework addresses this challenge through three
complementary strategies. First, the dimensional scores are not
estimated from raw behavior but from the calibrated probes of Chapter
17, which map observable text and behavioral indicators to dimensional
scores with known accuracy. The covariance matrix is then estimated from
probe-scored data, not from unstructured observables. Second, structured
experimental designs --- discrete choice experiments, factorial vignette
studies, and conjoint analyses --- can orthogonalize the dimensional
contributions, enabling identification of the covariance parameters via
structural equation modeling (SEM) or maximum likelihood estimation
(MLE). Third, the framework generates falsifiable predictions (six per
domain chapter) that provide external validation: if the estimated Σ
produces predictions that fail empirically, the matrix is miscalibrated
and must be re-estimated. The covariance matrix is not assumed; it is
empirically estimated, cross-validated, and falsifiably constrained.

\section{Chapter 20: Geometric Economics --- The Bond Geodesic
Equilibrium}\label{chapter-20-geometric-economics-the-bond-geodesic-equilibrium}

Epistemic status of this chapter. Except where a claim is explicitly
marked as a formal result or a conditional theorem, the material here is
a modeling interpretation, domain analogy, or speculative extension - a
demonstration of how the geometric vocabulary applies to economics, not
a set of established results at the evidential standard of the field.

\begin{quote}
\textbf{\textsc{RUNNING EXAMPLE --- Priya's Model}}

\emph{Priya discovers that TrialMatch implicitly models patients as
rational economic agents who will `choose' to travel for trials if the
expected medical benefit is high enough. This is the same error the Bond
Geodesic Equilibrium corrects: real agents follow heuristic gradients
through bounded decision spaces, not global optimization functions. A
rural patient without a car is not `choosing' lower-value care---she is
constrained to a different region of the manifold.}
\end{quote}

This chapter applies the Geometric Ethics framework to economic decision
theory, demonstrating that the classical model of Homo economicus --- a
perfectly rational agent maximizing scalar utility --- is not wrong
because humans are irrational, but incomplete because it computes on a
projected subspace of the actual decision manifold. We construct the
Bond Geodesic Equilibrium (BGE), prove it subsumes Nash equilibrium as a
special case, and derive prospect-theoretic phenomena as geometric
properties of the decision manifold.

\subsection{20.1 The Failure of Scalar
Economics}\label{the-failure-of-scalar-economics}

\textbf{Classical economics rests on Homo economicus} --- the perfectly
rational agent who maximizes a scalar utility function over all possible
actions. This model is mathematically elegant and empirically false.
Behavioral economics, following Kahneman and Tversky, has documented
systematic deviations from rational choice but lacks a unified
mathematical framework that explains why these deviations occur and what
structure they share.

\textbf{Kahneman\textquotesingle s dual-process theory} --- System 1
(fast, automatic, heuristic) and System 2 (slow, deliberate,
calculating) --- provides the psychological architecture. What it lacks
is the mathematical content: What is the space on which System 1 and
System 2 operate? What is the formal relationship between them? Why does
their interaction produce the specific pattern of "biases" that
behavioral economics has catalogued?

The answer is that economic decisions are not made on a one-dimensional
space. They are made on a multi-dimensional decision manifold that
includes monetary cost, moral constraint, social norms, fairness,
autonomy, trust, identity, institutional legitimacy, and epistemic
status --- the same nine dimensions of the moral manifold M developed in
Chapter 5, re-interpreted for economic contexts. Every economic
transaction is also a moral event. The classical
economist\textquotesingle s move of stripping away all dimensions except
monetary cost is not a simplification --- it is a dimensional collapse
that destroys the information needed to predict actual behavior.

\subsection{20.2 The Economic Decision
Complex}\label{the-economic-decision-complex}

Definition 20.1 (Economic Decision Complex). The economic decision
complex E is a weighted simplicial complex constructed as follows:

\begin{itemize}
\item
  \textbf{Vertices (0-simplices)}: Each vertex v\_i represents an
  economic state --- a configuration of goods, services, obligations,
  and relationships. The vertex carries an attribute vector a(v\_i) ∈
  R\^{}9 scoring the nine dimensions.
\item
  \textbf{Edges (1-simplices)}: An edge (v\_i, v\_j) represents an
  available action --- a transaction, decision, or behavioral choice.
  The edge carries a weight w(v\_i, v\_j) ≥ 0 representing the total
  cost on the full manifold.
\item
  \textbf{Higher simplices}: A k-simplex {[}v\_0, \ldots, v\_k{]}
  represents a bundle of mutually available actions.
\end{itemize}

Definition 20.2 (Multi-Dimensional Edge Weights). The weight of an edge
(v\_i, v\_j) in E is:

w(v\_i, v\_j) = Da\^{}T Sigma\^{}\{-1\} Da + Sum\_k beta\_k * 1{[}moral
boundary k crossed{]}

where Da = a(v\_j) - a(v\_i) is the attribute-vector difference, Sigma
is the 9x9 dimensional covariance matrix estimated from observed
economic behavior, and beta\_k is the penalty for crossing the k-th
moral-economic boundary.

The Mahalanobis distance is critical. Economic dimensions are not
independent: fairness (d\_3) interacts with consequences (d\_1);
autonomy (d\_4) modulates the weight of social impact (d\_6); trust
(d\_5) gates access to higher-consequence transactions (d\_1).

\subsection{20.3 Heuristic Bounded
Rationality}\label{heuristic-bounded-rationality}

The agent\textquotesingle s decision on E is A* search (Chapter 11) with
evaluation function f(n) = g(n) + h(n), where g(n) is the accumulated
economic cost (System 2: slow, deliberate calculation) and h(n) is the
moral-heuristic cost (System 1: fast, automatic boundary enforcement).

Theorem 20.1 (Heuristic Truncation). If the moral heuristic h is
admissible --- i.e., h(n) ≤ h*(n) for all nodes n, where h*(n) is the
true remaining cost on the full manifold --- then A* search with h finds
the optimal path (the Bond geodesic) on E.

Moral heuristics --- "do not steal," "keep promises," "reciprocate
kindness" --- are admissible because they never overestimate the true
cost of violating a moral boundary. The penalty for stealing (social,
legal, psychological) is always at least as large as the
heuristic\textquotesingle s estimate. These are not cognitive biases;
they are search heuristics that reduce an intractable optimization to a
tractable one, at the cost of provably bounded suboptimality.\\
\strut \\
{[}Formal complexity result.{]} Define MORAL-PATHFIND as the decision
problem: given a weighted simplicial complex C with n vertices,
d-dimensional attribute vectors with Mahalanobis edge weights, and k ≥ 1
attribute-boundary constraints (edges that cross a constraint boundary
incur penalty βₖ), does there exist a path from s to G with total cost ≤
B? MORAL-PATHFIND is NP-hard by reduction from the Weight-Constrained
Shortest Path problem (Garey \& Johnson, 1979; Handler \& Zang, 1980),
which is NP-hard for k ≥ 1 side constraints. The reduction maps each
side constraint to a boundary penalty β\_k and each edge weight to the
corresponding Mahalanobis distance. Since MORAL-PATHFIND is NP-hard,
exact optimization on the moral manifold with boundary constraints is
computationally intractable in the worst case, and the
bounded-suboptimality guarantees of Weighted A* (Pohl, 1970) apply by
direct instantiation --- not by analogy --- to the moral search problem.

This is Simon\textquotesingle s bounded rationality program, formalized:
agents are rational on the manifold they can perceive, using heuristics
that provably approximate the optimal solution.

Remark (Marr\textquotesingle s Levels of Analysis). The claim that
economic decision-making is A* search must be understood at the correct
level of David Marr\textquotesingle s (1982) tri-level analysis. A* is
the computational theory --- it specifies what is computed (the
minimum-cost path on the decision manifold) and why (optimality under
heuristic bounded rationality). The framework makes no claim about the
algorithmic level (how the brain implements the computation) or the
implementational level (which neural circuits execute it). The brain may
use satisficing, fast-and-frugal heuristics (Gigerenzer), accumulator
models, or drift-diffusion processes that are algorithmically very
different from A* but converge to the same computational-level solution.
The claim is that A* correctly describes the function being computed,
not the algorithm computing it. This distinction matters: cognitive
scientists who object that "the brain does not run A*" are objecting at
Marr\textquotesingle s algorithmic level, while the
framework\textquotesingle s claims are at the computational level. The
two are compatible. The framework predicts the outcome of the
computation (which path the agent selects); it does not predict the
neural implementation of the selection process.

\subsection{20.4 The Bond Geodesic and the Scalar Irrecoverability
Theorem}\label{the-bond-geodesic-and-the-scalar-irrecoverability-theorem}

Definition 20.3 (Bond Geodesic). The Bond geodesic gamma* from economic
state s to state t on E is the minimum-cost path:

gamma* = arg min\_gamma Sum w(e\_i)

where the sum is over edges e\_i in path gamma. When boundary penalties
are absorbed into the metric, this is a geodesic in the sense of Chapter
10.

Theorem 20.2 (Scalar Irrecoverability). No continuous function phi:
R\^{}9 -\textgreater{} R is injective. Therefore, any continuous scalar
utility function u = phi(a) applied to the nine-dimensional attribute
vector a ∈ R\^{}9 destroys information that is mathematically
irrecoverable.

Proof. By Brouwer\textquotesingle s invariance of dimension, R\^{}9 and
R\^{}1 are not homeomorphic. Restricting to the compact cube
{[}0,1{]}\^{}9, suppose phi: {[}0,1{]}\^{}9 -\textgreater{} R were a
continuous injection. Then phi maps a compact space injectively and
continuously into a Hausdorff space, so phi is a homeomorphism onto its
image --- an embedding of a 9-dimensional space into R\^{}1. This
contradicts invariance of dimension. {[}{]}

This is the Scalar Irrecoverability Theorem of Chapter 15.6, applied to
the economic domain. "Irrational" behavior is rational on the full
manifold, irrational only on the projection. Homo economicus computes on
the wrong manifold.

\subsection{20.5 The Bond Geodesic
Equilibrium}\label{the-bond-geodesic-equilibrium}

When multiple agents pathfind simultaneously on the economic decision
complex, we enter game theory. The Bond Geodesic Equilibrium (BGE) is
the natural equilibrium concept on the manifold.

\textbf{Definition 20.4 (Bond Geodesic Equilibrium)}. A strategy profile
(gamma\_1*, \ldots, gamma\_n*) is a Bond Geodesic Equilibrium if for
each agent i, the path gamma\_i* is a Bond geodesic on agent
i\textquotesingle s perceived decision complex, given the paths chosen
by all other agents.

To establish the relationship between BGE and Nash equilibrium, we
construct the augmented game.

Definition 20.5 (Augmented Game). Given a multi-agent decision problem
on E, the augmented game is Gamma+ = (N, \{S\_i\}, \{u\_i\}) where:

• N is the set of agents

• S\_i is agent i\textquotesingle s set of available paths on E

• u\_i = -BF\_i is agent i\textquotesingle s payoff: the negative of the
total behavioral friction (edge weight sum) along their chosen path,
given others\textquotesingle{} paths

This construction is canonical: it says that agents in the manifold game
are maximizing the negative of their total cost on the full manifold ---
i.e., they are minimizing behavioral friction, which is exactly what A*
search does.

\subsubsection{The BGE--Nash
Relationship}\label{the-bgenash-relationship}

Theorem 20.3 (BGE--Nash Relationship).

\begin{enumerate}
\def\labelenumi{(\arabic{enumi})}
\item
  Equivalence: A strategy profile is a BGE of the manifold game if and
  only if it is a Nash equilibrium of the augmented game Gamma+.
\item
  Scalar nesting: If all non-monetary dimensions are zeroed out (i.e.,
  the attribute vector collapses to a(v) = (d\_1(v), 0, \ldots, 0)) and
  each agent selects a single action (not a path), then BGE reduces to
  Nash equilibrium of the classical game with payoffs u\_i = d\_1.
\item
  Refinement: In games where agents can operate on the full manifold,
  BGE selects among Nash equilibria by the manifold structure --- the
  equilibrium that minimizes total behavioral friction across all
  dimensions.
\end{enumerate}

\textbf{Proof}. (1) By construction, gamma\_i* is a Bond geodesic if and
only if it minimizes BF\_i given others\textquotesingle{} paths, which
is exactly the Nash condition on Gamma+. (2) When a = (d\_1, 0, \ldots,
0) and paths are single edges, BF\_i = d\_1 cost, so -BF\_i = -d\_1 cost
= classical payoff (up to sign). (3) Among multiple Nash equilibria of
Gamma+, the BGE selects the profile with minimum total BF, since that is
what A* search converges to. {[}{]}

This is the central result: Nash equilibrium is the d\_1-only projection
of BGE, exactly as Newtonian mechanics is the low-velocity limit of
general relativity. The projection is the contraction of Chapter 15
applied to game theory.

\subsubsection{Existence and Uniqueness}\label{existence-and-uniqueness}

Theorem 20.4 (Existence of Mixed BGE). Every finite augmented game
Gamma+ has at least one mixed BGE.

Proof. By Nash\textquotesingle s theorem (1950), every finite game has a
mixed-strategy Nash equilibrium. Gamma+ is a finite game (finite agents,
finite paths). By Theorem 20.3(1), a mixed Nash equilibrium of Gamma+ is
a mixed BGE. {[}{]}

\textbf{Definition 20.6 (Mixed BGE)}. A mixed BGE is a profile of
probability distributions (sigma\_1, \ldots, sigma\_n) over paths, where
each sigma\_i minimizes agent i\textquotesingle s expected behavioral
friction given others\textquotesingle{} mixed strategies.

\textbf{Remark (Pure BGE Existence)}. Pure BGE need not exist in general
(as with pure Nash). Sufficient conditions include: (a) the augmented
game is a potential game; (b) the augmented game has supermodular
structure; (c) agents\textquotesingle{} decision complexes are
independent (no cross-agent edge weight dependencies).

\subsubsection{The Contraction Lemma}\label{the-contraction-lemma}

\textbf{Lemma 20.1 (Contraction Condition for BGE Uniqueness)}. For each
agent i, let alpha\_i \textgreater{} 0 measure the self-separation of
agent i\textquotesingle s best-response mapping (how much
i\textquotesingle s own behavioral friction changes with
i\textquotesingle s own path choice) and let kappa\_i \textgreater= 0
measure the cross-sensitivity (how much i\textquotesingle s behavioral
friction changes with others\textquotesingle{} path choices). If

alpha\_i \textgreater{} kappa\_i for all i in N

then the best-response mapping on Gamma+ is a contraction with Lipschitz
constant L = max\_i(kappa\_i / alpha\_i) \textless{} 1, and the BGE is
unique.

\textbf{Proof}. The best-response mapping T: S -\textgreater{} S maps
each profile to the profile of individual best responses. For agent i,
\textbar\textbar T\_i(s) - T\_i(s\textquotesingle)\textbar\textbar≤
(kappa\_i/alpha\_i)\textbar\textbar s\_\{-i\} -
s\textquotesingle\_\{-i\}\textbar\textbar{} by the ratio of
cross-sensitivity to self-separation. Taking the maximum over i,
\textbar\textbar T(s) - T(s\textquotesingle)\textbar\textbar≤
L\textbar\textbar s - s\textquotesingle\textbar\textbar{} with L =
max\_i(kappa\_i/alpha\_i) \textless{} 1. By Banach\textquotesingle s
fixed-point theorem, T has a unique fixed point. {[}{]}

\textbf{Interpretation}: BGE is unique when each agent\textquotesingle s
decision is dominated by their own manifold structure (high alpha\_i)
rather than by what others do (low kappa\_i). This is the "weak
coupling" condition. When coupling is strong (kappa\_i ≈ alpha\_i),
multiple BGE can coexist --- as in coordination games, where the
equilibrium depends on which manifold region the agents coordinate on.

\subsection{20.6 Welfare Properties}\label{welfare-properties}

\textbf{Theorem 20.5 (BGE Welfare)}. A BGE is Nash-optimal on the
augmented game Gamma+: no agent can unilaterally reduce their behavioral
friction by deviating.

However, a BGE is not in general Pareto optimal. Consider the
Prisoner\textquotesingle s Dilemma on the manifold: if the fairness
dimension d\_3 is inactive, both agents defect (as in classical Nash);
the mutual-cooperation profile has lower total BF but is not
individually stable. The Pareto failure of Nash equilibrium persists in
BGE for the same structural reason.

\textbf{Remark}. In potential games --- where all
agents\textquotesingle{} behavioral friction derives from a common
potential function --- the BGE does coincide with a Pareto optimum of
the augmented game. This is a strong sufficient condition that holds in
many economically important settings (congestion games, market entry
games).

\subsection{20.7 Behavioral Game Theory as Manifold
Geometry}\label{behavioral-game-theory-as-manifold-geometry}

The BGE framework resolves several longstanding puzzles in behavioral
game theory by identifying the active manifold dimensions:

\textbf{Ultimatum game rejections}. In the standard ultimatum game, a
proposer offers a split of \$10; the responder can accept or reject
(destroying both payoffs). Homo economicus predicts the responder
accepts any positive offer. Experimentally, offers below
\textasciitilde30\% are rejected roughly half the time. The BGE
explanation: the responder\textquotesingle s decision complex includes
d\_3 (fairness). Rejecting an unfair offer has high monetary cost (d\_1)
but avoids a large fairness boundary penalty. The Bond geodesic includes
rejection when the fairness penalty exceeds the monetary cost.

\textbf{Trust game cooperation}. In trust games, agents cooperate far
more than Nash predicts. The BGE explanation: dimension d\_5 (trust) is
active. Cooperation is the Bond geodesic when the trust
dimension\textquotesingle s edge weights make defection more costly on
the full manifold than on the d\_1 projection.

\textbf{Public goods provision}. Agents contribute to public goods even
when free-riding is the Nash strategy. The BGE explanation: dimension
d\_6 (social impact) creates boundary penalties for free-riding that
exceed the monetary savings.

Each of these "anomalies" is predicted by BGE, not catalogued post hoc.
The framework says: tell me the active dimensions and their weights, and
I will tell you the equilibrium. This is a stronger claim than
behavioral economics makes, and it is falsifiable.

\subsection{20.8 Prospect Theory as Manifold
Geometry}\label{prospect-theory-as-manifold-geometry}

The major phenomena of Kahneman and Tversky\textquotesingle s prospect
theory are derived as geometric properties of the decision manifold:

\textbf{Loss aversion}. The edge weights in E are asymmetric: the weight
for moving from state a to a worse state b (loss) exceeds the weight for
the reverse move (gain), even when \textbar Da\textbar{} is the same.
Formally, w(a -\textgreater{} b) \textgreater{} w(b -\textgreater{} a)
when a is the reference point. This is not a "bias" --- it is an
asymmetry of the manifold metric, analogous to the asymmetry of a
Finsler metric.

\textbf{Framing effects}. A framing effect --- where logically
equivalent descriptions of the same decision produce different choices
--- is a gauge symmetry violation on the decision manifold. The Bond
Invariance Principle (BIP, Chapter 12) requires that evaluations be
invariant under meaning-preserving transformations. Framing effects
occur when the agent\textquotesingle s heuristic function h(n) is not
gauge-invariant. This is a miscalibration of the heuristic, not a
fundamental feature of the manifold.

\textbf{Reference dependence}. The edge weights depend on the starting
vertex --- the agent\textquotesingle s current state serves as the
reference point. This is built into the Mahalanobis distance: Da =
a(v\_j) - a(v\_i) is measured from the current state, not from an
absolute origin.

\textbf{Hyperbolic discounting}. The temporal dimension of the manifold
has non-constant curvature: near-future and far-future states are not
metrically equivalent. The "present bias" is a consequence of higher
edge-weight density near the current temporal position.

\textbf{The endowment effect}. Owning a good changes the
agent\textquotesingle s position on the manifold (the attribute vector
shifts). The Mahalanobis distance from "own the good" to "sell the good"
differs from the distance from "don\textquotesingle t own" to "buy the
good" because the reference point has changed. This is not irrationality
--- it is path-dependence on a curved manifold.

\subsection{20.9 Attribute Conservation in Bilateral
Exchange}\label{attribute-conservation-in-bilateral-exchange}

\textbf{Theorem 20.6 (Attribute Conservation)}. In a closed bilateral
exchange between agents A and B, for each transferable dimension k ∈
\{d\_1 (monetary value), d\_2 (rights/entitlements), d\_4 (autonomy)\}:

Da\_k(A) + Da\_k(B) = 0

That is, what A gains on dimension k, B loses, and vice versa.

This is the economic specialization of the Invariance of Harm Accounting
(Chapter 12, Theorem 12.1): in a closed system with re-description
invariance, conserved quantities cannot be created or destroyed by
relabeling. For transferable dimensions --- money, property rights,
autonomy --- bilateral exchange is zero-sum on each dimension
separately.

\textbf{Remark (Non-Conservation of Evaluative Dimensions)}. Evaluative
dimensions --- d\_3 (fairness), d\_7 (identity/virtue) --- are not
conserved. Both parties to an exchange can simultaneously perceive
increased fairness, or both can perceive unfairness. A fair exchange
creates mutual value on d\_3 that did not exist before the exchange.
This distinction between transferable and evaluative dimensions
formalizes the intuition that trade creates value beyond the material
transfer.

\subsection{20.10 Worked Examples}\label{worked-examples}

The preceding sections developed the Bond Geodesic Equilibrium as a
mathematical structure. This section applies it to three real-world
cases, demonstrating that the framework is not merely formal but
operationally predictive. Each case study identifies the active manifold
dimensions, computes the d\_1-only equilibrium (what classical economics
predicts), and then computes the full-manifold BGE (what actually
happened). The gap between the two is the empirical signature of the
missing dimensions.

\subsubsection{Example 20.1: Uber Surge Pricing --- Why "Efficient"
Pricing Fails on the
Manifold}\label{example-20.1-uber-surge-pricing-why-efficient-pricing-fails-on-the-manifold}

Uber\textquotesingle s surge pricing algorithm multiplies base fares by
a demand factor during peak times. Classical economics says this is
efficient --- it clears the market by matching supply and demand through
the price mechanism. When demand exceeds supply, the price rises until
enough passengers drop out and enough drivers enter to restore
equilibrium. On d\_1 alone, this is textbook welfare maximization. Yet
passengers routinely express outrage at surge pricing, regulators have
investigated it, and Uber itself has capped surge multipliers. If the
pricing is efficient, why does the market reject it?

The 9D analysis reveals that surge pricing optimizes on a single
dimension while creating catastrophic costs on several others. On d\_1,
the price is indeed "efficient" --- it matches supply and demand.
Classical economics stops here. But d\_3 (fairness) tells a different
story: a 5x surge during a snowstorm or after a terrorist attack --- as
happened after the 2014 Sydney hostage crisis, when Uber fares spiked to
over \$100 for short rides away from the Lindt Cafe siege --- violates
deeply held fairness norms. The identical ride, described differently
("market-clearing price" versus "price gouging during crisis"), receives
radically different evaluations. This is a gauge-invariance violation:
the economic content is identical, but the moral evaluation changes
under re-description.

The autonomy dimension d\_4 exposes a further failure. During
emergencies, passengers have no viable alternative --- public transit
may be shut down, taxis overwhelmed, walking dangerous. Their d\_4 is
effectively zero: they cannot meaningfully choose to decline the
transaction. Surge pricing exploits the curvature increase in d\_4
during emergencies, extracting maximum revenue precisely when the
passenger\textquotesingle s autonomy is at its minimum. On d\_5 (trust),
passengers who chose Uber as their primary transport trusted the
platform to provide affordable, reliable service. Surge pricing during
crises shatters that trust, leading to long-term behavioral changes:
uninstalling the app, switching to competitors, supporting regulatory
intervention. The d\_5 cost is not captured in any single transaction
but accumulates across the entire user relationship.

The social impact dimension d\_6 reveals the class structure of surge
pricing during emergencies: the wealthy can afford to escape at any
multiplier, while the poor cannot. A 5x surge effectively rations
emergency transport by income, a d\_6 failure with stark class
implications that no amount of "efficiency" language can neutralize.

The BGE analysis makes this precise. Uber\textquotesingle s algorithm
computes the d\_1-only BGE --- the surge price IS the d\_1 equilibrium.
But the full-manifold BGE, incorporating d\_3, d\_4, d\_5, and d\_6 with
empirically appropriate weights, produces a different equilibrium: one
with lower surge multipliers, price freezes during declared emergencies,
and caps on maximum surges. This is exactly what Uber eventually
implemented after sustained public backlash --- capped surges, automatic
price freezes during crises, and commitments to "no surge" during
natural disasters. Uber empirically discovered the manifold. The
framework would have predicted these corrections in advance, because the
d\_1-only geodesic was visibly far from the full-manifold geodesic from
the moment surge pricing was deployed.

\subsubsection{Example 20.2: The Dictator Game across Cultures ---
Metric
Variation}\label{example-20.2-the-dictator-game-across-cultures-metric-variation}

In the dictator game, one player (the dictator) receives \$10 and can
give any portion to another player who has no power to accept or reject
--- the second player simply receives whatever the dictator allocates.
Homo economicus predicts the dictator gives \$0, since giving away money
reduces d\_1 with no strategic consequence. This prediction is
spectacularly wrong. In Western undergraduate samples, the modal offer
is \$2--\$3, with a mean around \$2.80. But the more revealing data come
from Henrich et al. (2005), who ran the dictator game in 15 small-scale
societies across 5 continents and found enormous variation: average
offers ranged from 15\% of the endowment (Machiguenga, Peruvian Amazon)
to 58\% (Lamelara, Indonesia).

The variation is not random --- it correlates systematically with market
integration and cooperation norms within each society. The Lamelara are
a whale-hunting community on the island of Lembata. Whale hunting
requires coordinated crews of 10--15 men in open boats, harpooning sperm
whales by hand. Failure to cooperate means failure to eat. This
community has evolved extraordinarily high weights on d\_5 (trust) and
d\_6 (social impact). The community\textquotesingle s metric tensor
Sigma has high off-diagonal terms between d\_1 and d\_5/d\_6, meaning
that monetary decisions cannot be separated from their trust and
social-impact implications. The dictator game offer of 58\% reflects the
community\textquotesingle s metric, not individual irrationality --- in
Lamelara, keeping most of a windfall for yourself would be a severe
violation of the social contract that sustains cooperative whale
hunting.

At the other extreme, the Machiguenga are relatively isolated
slash-and-burn horticulturalists with low market integration and limited
economic cooperation beyond the household. Their metric tensor places
more weight on d\_1 relative to d\_5 and d\_6 --- not because they are
"selfish" in any morally loaded sense, but because their economic
ecology does not require the intense cooperation that the
Lamelara\textquotesingle s does. The off-diagonal terms between d\_1 and
d\_5/d\_6 are smaller. Their dictator game offers (15\%) reflect this
metric faithfully. The Orma, Kenyan pastoralists undergoing rapid market
integration at the time of the study, showed an intermediate and
informative pattern: offers correlated with individual income,
suggesting that d\_3 (fairness norms tied to wealth) is activated
differently at different d\_1 levels --- wealthier Orma gave more,
consistent with a metric where the d\_1--d\_3 cross-term increases with
absolute d\_1 level.

The framework generates a specific, testable prediction (Prediction 3
from Section 20.11): cross-cultural variation in economic games should
be predictable from the community\textquotesingle s estimated covariance
matrix Sigma, which can be independently measured through ethnographic
observation of cooperation patterns, punishment norms, and
resource-sharing customs. The Henrich et al. data constitute the
strongest existing evidence that Sigma varies systematically across
cultures exactly as the framework requires. A classical model that
treats all humans as optimizing on d\_1 alone cannot explain the
Machiguenga--Lamelara gap without invoking ad hoc "cultural
preferences"; the manifold framework predicts it from the metric.

\subsubsection{Example 20.3: Pharmaceutical Price Gouging --- Turing
Pharmaceuticals and Daraprim
(2015)}\label{example-20.3-pharmaceutical-price-gouging-turing-pharmaceuticals-and-daraprim-2015}

In September 2015, Turing Pharmaceuticals, led by CEO Martin Shkreli,
acquired the U.S. marketing rights to Daraprim (pyrimethamine), a drug
used to treat toxoplasmosis --- a parasitic infection that is
life-threatening in AIDS patients and others with compromised immune
systems. Turing immediately raised the price from \$13.50 to \$750.00
per tablet, a 5,455\% increase overnight. The drug had been off-patent
since 1953. Production cost was approximately \$1 per tablet. Turing had
conducted no new research on the compound. The price increase was pure
rent extraction enabled by regulatory barriers to generic entry (the
drug\textquotesingle s small market made generic production unprofitable
at the old price, and Turing implemented a closed distribution system
that prevented generic manufacturers from obtaining samples for
bioequivalence testing).

On d\_1 (monetary value) alone, Shkreli\textquotesingle s decision was
economically "rational" in the classical sense. Daraprim had inelastic
demand --- patients with toxoplasmosis had no therapeutic alternative,
so the price increase would be absorbed by insurers and patients rather
than reducing volume. The d\_1-only geodesic pointed directly at maximum
extraction: charge whatever the market will bear, and the market will
bear nearly anything when the alternative is death. On d\_1 alone,
Shkreli\textquotesingle s geodesic was optimal.

But the full manifold tells a catastrophically different story. On d\_2
(obligations), pharmaceutical companies operate under implicit
obligations to patients --- a downstream duty rooted in the Hippocratic
tradition. The price made Daraprim inaccessible to uninsured patients
and created insurance authorization barriers for others, directly
violating d\_2. On d\_3 (fairness), the price bore no relationship to
production cost (\$1/tablet) or to R\&D investment (\$0 --- Turing did
no new research on pyrimethamine). The d\_3 boundary penalty is
enormous: the price was disconnected from any value-creating activity
and represented pure rent extraction from captive patients.

On d\_5 (trust), trust in the pharmaceutical industry --- already
fragile after decades of pricing controversies --- was catastrophically
damaged. The Daraprim case became the single most-cited example of
pharmaceutical greed in public discourse, contributing to regulatory
backlash that affected the entire industry. The d\_5 cost was not borne
by Turing alone but was externalized across all pharmaceutical
companies. On d\_6 (social impact), AIDS patients who could not afford
the drug or navigate insurance barriers faced serious health
consequences or death. The d\_6 cost was borne entirely by the most
vulnerable population --- immunocompromised patients, disproportionately
from lower-income communities. On d\_7 (identity/virtue),
Shkreli\textquotesingle s public persona --- smirking, unapologetic,
taunting critics on social media --- activated d\_7 revulsion. The case
violated the pharmaceutical industry\textquotesingle s self-image as a
healing profession, triggering identity-level rejection from physicians,
pharmacists, and even other pharmaceutical executives. On d\_8
(institutional legitimacy), the case exposed regulatory gaps --- no
existing mechanism prevented price gouging on off-patent essential
medicines with small markets --- undermining public confidence in the
institutions that are supposed to protect patients from exploitation.

The Daraprim case is a canonical demonstration that d\_1-optimal
geodesics can be catastrophically wrong on the full manifold. The market
eventually "punished" the deviation from the manifold geodesic, though
more slowly than the framework\textquotesingle s real-time calculation
would have predicted: Shkreli was eventually imprisoned (for securities
fraud, not pricing, though the pricing notoriety clearly influenced
prosecutorial attention), Turing\textquotesingle s reputation was
destroyed, the company was restructured, and the case catalyzed
Congressional hearings and state-level drug pricing legislation.
Competing manufacturers eventually entered the market with generic
pyrimethamine at a fraction of the price. Every prediction of the
manifold framework was confirmed: the d\_1-only geodesic was unstable,
the manifold exerted corrective forces on every non-monetary dimension,
and the system converged --- slowly, expensively, with unnecessary
suffering --- toward the full-manifold BGE that the framework would have
identified from the start.

\subsection{20.11 Falsifiable
Predictions}\label{falsifiable-predictions}

The framework generates six predictions that distinguish it from both
classical and behavioral economics:

\textbf{Prediction 1 (Dimensional Activation)}: Ultimatum game rejection
rates should vary with the salience of specific manifold dimensions.
Making fairness (d\_3) more salient (e.g., by framing the game
explicitly as a fairness test) should increase rejections; making it
less salient (anonymous, one-shot, large stakes) should decrease them.
What would falsify: if rejection rates are invariant to dimensional
salience manipulations.

\textbf{Prediction 2 (Bond Index Correlates)}: Agents with higher Bond
Index scores should (a) reject more unfair ultimatum offers, (b)
contribute more in public goods games, (c) show higher WTP for
fair-trade goods, and (d) exhibit smaller framing effects. What would
falsify: if Bond Index is uncorrelated with these measures.

\textbf{Prediction 3 (Cross-Cultural Metric Variation)}: The covariance
matrix Sigma should vary systematically across cultures in ways
predicted by known cultural dimensions. What would falsify: if Sigma is
invariant across cultures, or varies unsystematically.

\textbf{Prediction 4 (Boundary Penalty Measurement)}: Sacred-value
boundaries (beta\_k = infinity) should produce qualitatively different
response patterns from high-but-finite boundaries. What would falsify:
if all boundary effects are graded.

\textbf{Prediction 5 (Heuristic Admissibility)}: Moral heuristics should
satisfy admissibility: h(n) ≤ h*(n). Agents using cultural heuristics
should find paths that are suboptimal but within epsilon of optimal.
What would falsify: if moral heuristics systematically overestimate
costs.

\textbf{Prediction 6 (Manifold Dimensionality)}: Factor analysis of
economic behavioral data should recover approximately nine independent
dimensions. What would falsify: if the factor structure is consistently
lower-rank (fewer than seven) or higher-rank (more than eleven).

\subsection{20.12 Connection to the
Framework}\label{connection-to-the-framework-1}

The Bond Geodesic Equilibrium completes the multi-agent story that began
in Chapter 14 (Collective Moral Agency):

• Chapter 11 solved single-agent pathfinding: f(n) = g(n) + h(n) on the
moral manifold.

• Chapter 14 built the collective agency tensor: how multi-agent moral
structure is represented.

• This chapter answers: what is the equilibrium when multiple agents
pathfind simultaneously?

• Chapter 15 asks: how does this equilibrium get contracted to a scalar
decision?

The BGE is the game-theoretic completion of the geometric ethics
program. Nash equilibrium --- the foundation of modern economics --- is
the d\_1-only contraction of BGE, exactly as scalar utility is the
rank-0 contraction of the moral tensor (Chapter 6.7). The same
mathematical operation (contraction, Chapter 15) that explains why
scalar ethics loses information also explains why scalar economics loses
information.

\subsection{20.13 Summary}\label{summary-11}

This chapter has shown that the geometric ethics framework, when applied
to economic decision theory, yields:

\begin{enumerate}
\def\labelenumi{\arabic{enumi}.}
\item
  A formal construction of the economic decision complex E as a
  domain-specific instantiation of the moral manifold M.
\item
  Heuristic bounded rationality: Simon\textquotesingle s program
  formalized via A* search.
\item
  The Bond Geodesic Equilibrium: a generalization of Nash equilibrium
  with existence via Nash\textquotesingle s theorem on the augmented
  game, uniqueness via the contraction lemma, and a formal nesting of
  Nash as the scalar special case.
\item
  Geometric behavioral economics: prospect-theoretic phenomena derived
  as properties of the manifold metric.
\item
  Attribute conservation: transferable dimensions are conserved in
  bilateral exchange; evaluative dimensions permit mutual value
  creation.
\item
  Six falsifiable predictions distinguishing the framework from
  classical and behavioral economics.
\end{enumerate}

The next eight chapters extend this pattern to clinical medicine, law,
financial markets, theology, environmental policy, artificial
intelligence, bioethics, and military ethics.

\section{Chapter 21: Geometric Clinical Ethics --- Triage, Consent, and
the Clinical
Geodesic}\label{chapter-21-geometric-clinical-ethics-triage-consent-and-the-clinical-geodesic}

Epistemic status of this chapter. Except where a claim is explicitly
marked as a formal result or a conditional theorem, the material here is
a modeling interpretation, domain analogy, or speculative extension - a
demonstration of how the geometric vocabulary applies to clinical
ethics, not a set of established results at the evidential standard of
the field.

\begin{quote}
\textbf{\textsc{RUNNING EXAMPLE --- Priya's Model}}

\emph{Dr. Osei points out that TrialMatch effectively assigns QALYs:
Mrs. Voss's `expected quality-adjusted life-years from trial
participation' drives her score. By the QALY Irrecoverability Theorem,
this number cannot capture that Mrs. Voss values knowing she fought
independently of outcome. Her dignity-interest Iₘ at μ = 5 contracts
with a different obligation component than her survival-interest at μ =
1. The scalar cannot hold both.}
\end{quote}

This chapter applies the Geometric Ethics framework to clinical medical
ethics. The central claim is that the field\textquotesingle s
foundational tools --- principlism, QALYs, and narrative deliberation
--- share a common deficiency: they lack the mathematical structure
needed to make clinical-ethical reasoning precise, consistent, and
auditable. The framework developed in the preceding chapters provides
that structure. Clinical decisions are pathfinding on a nine-dimensional
clinical decision complex, informed consent is a gauge-invariance
condition, moral injury is cumulative manifold damage, and the QALY is a
scalar projection that destroys eight dimensions of clinically relevant
information.

\subsection{21.1 The Failure of Scalar Clinical
Ethics}\label{the-failure-of-scalar-clinical-ethics}

Clinical ethics currently operates with four principal tools, each of
which is inadequate for distinct reasons.

\begin{enumerate}
\def\labelenumi{\arabic{enumi}.}
\item
  \textbf{Principlism}. Beauchamp and Childress\textquotesingle s four
  principles --- autonomy, beneficence, non-maleficence, and justice ---
  provide a vocabulary for ethical deliberation but no algorithm for
  resolving conflicts among them. When autonomy and beneficence collide
  (a competent patient refuses life-saving treatment), principlism
  offers no formal method for determining which principle prevails. The
  resolution depends entirely on the clinician\textquotesingle s
  judgment, which is neither transparent nor reproducible.
\item
  \textbf{QALYs.} The Quality-Adjusted Life Year compresses the full
  clinical-moral state into a single scalar: years of life weighted by
  health-related quality. This is precisely the scalar projection that
  Chapter 15 proved destroys information. The QALY maps a
  nine-dimensional clinical state to a one-dimensional score, with a
  kernel of dimension eight. All information in that kernel ---
  autonomy, trust, dignity, justice, epistemic status --- is
  mathematically irrecoverable.
\item
  \textbf{Casuistry}. Case-based reasoning identifies morally relevant
  similarities between current and paradigm cases but provides no formal
  measure of "similarity." Two clinicians may judge different features
  as morally relevant and reach opposite conclusions, with no principled
  method for adjudication.
\item
  \textbf{Narrative} \textbf{ethics}. The narrative approach enriches
  understanding of patients\textquotesingle{} lived experience but is
  inherently non-quantitative. It cannot support the kind of systematic
  policy analysis that institutional decision-making requires.
\end{enumerate}

The geometric framework addresses all four deficiencies: it provides an
algorithm (A* search), a multi-dimensional cost function (Mahalanobis
distance on the clinical decision complex), a formal similarity metric
(graph distance), and quantitative policy analysis (the Clinical Bond
Index).

\subsection{21.2 The Clinical Decision
Complex}\label{the-clinical-decision-complex}

Definition 21.1 (Clinical Decision Complex). The clinical decision
complex C is a weighted simplicial complex with the following structure:

\textbf{Vertices} (0-simplices): Each vertex v\_i represents a clinical
state --- a configuration of diagnosis, prognosis, treatment status,
patient preferences, and institutional context. The vertex carries an
attribute vector a(v\_i) in R\^{}9 scoring the nine clinical-moral
dimensions.

\textbf{Edges} (1-simplices): An edge (v\_i, v\_j) represents a clinical
action --- a treatment, test, conversation, referral, or decision to
wait. The edge carries a weight w(v\_i, v\_j) \textgreater= 0
representing the total clinical-moral cost of the transition.

\textbf{Higher simplices}: A k-simplex {[}v\_0, ..., v\_k{]} represents
a care bundle --- a coordinated set of interventions (e.g., intubation +
sedation + family notification).

The nine clinical-moral dimensions are the moral manifold dimensions of
Chapter 5, instantiated for clinical contexts:

\begin{itemize}
\item
  \textbf{d\_1: Clinical outcomes} --- survival, morbidity, functional
  status, symptom burden. This is the dimension QALYs target.
\item
  \textbf{d\_2: Rights and obligations} --- the
  patient\textquotesingle s right to treatment, duty of care,
  non-maleficence, the right to refuse.
\item
  \textbf{d\_3: Justice and fairness} --- equitable access, distributive
  justice, non-discrimination, proportionality of resource allocation.
\item
  \textbf{d\_4: Autonomy} --- self-determination, voluntariness, freedom
  from coercion, capacity for informed decision-making.
\item
  \textbf{d\_5: Trust} --- the therapeutic relationship, fiduciary duty,
  confidentiality, institutional trustworthiness.
\item
  \textbf{d\_6: Social and relational impact} --- effects on family and
  caregivers, community health implications, social determinants, social
  role preservation.
\item
  \textbf{d\_7: Dignity and identity} --- the patient\textquotesingle s
  sense of self, bodily integrity, alignment with personal values,
  quality of dying, and (for clinicians) professional identity.
\item
  \textbf{d\_8: Institutional legitimacy} --- standard of care, clinical
  guidelines, regulatory compliance, legal defensibility.
\item
  \textbf{d\_9: Epistemic status} --- diagnostic certainty, prognostic
  accuracy, patient understanding, health literacy.
\end{itemize}

\textbf{Definition 21.2 (Clinical Edge Weights)}. The weight of an edge
(v\_i, v\_j) in C is:

w(v\_i, v\_j) = Da\^{}T Sigma\^{}\{-1\} Da + Sum\_k beta\_k *
1{[}boundary k crossed{]}

where Da = a(v\_j) - a(v\_i) is the attribute-vector difference, Sigma
is the 9x9 clinical-moral covariance matrix estimated from clinical
behavioral data, and beta\_k is the penalty for crossing the k-th
clinical-moral boundary.

The covariance matrix captures critical clinical-moral
interdependencies. Sigma\_\{1,5\} (outcomes x trust): uncertain outcomes
cost more when trust is fragile. Sigma\_\{4,9\} (autonomy x epistemic):
respecting autonomy requires adequate patient understanding.
Sigma\_\{3,7\} (justice x dignity): allocation decisions that are fair
but dignity-stripping have high cross-dimensional cost.

\textbf{Definition 21.3 (Clinical-Moral Boundaries)}. Five boundary
types carry distinctive penalties:

\begin{enumerate}
\def\labelenumi{\arabic{enumi}.}
\item
  Harm boundary (beta\_harm): "First, do no harm." Very high but finite
  --- surgery crosses it acceptably because the clinical geodesic
  through the harm boundary reaches a state with sufficiently lower
  total cost.
\item
  Futility boundary (beta\_futile): Treating patients who cannot
  benefit. Crosses d\_3 (justice --- resources consumed without
  commensurate benefit), d\_7 (dignity --- subjecting patients to
  interventions that cannot help), and d\_5 (trust --- performing
  procedures the clinician knows are futile).
\item
  Consent boundary (beta\_consent): Approximately infinite for
  competent, non-emergency patients. Finite for genuine emergencies
  (implied consent) and for patients lacking capacity (surrogate consent
  at reduced boundary cost).
\item
  Abandonment boundary (beta\_abandon): Withdrawing care from dependent
  patients without appropriate transition. One of the highest-cost
  crossings in clinical practice.
\item
  Sacred-value boundary (beta\_sacred = infinity): Non-consensual
  experimentation, execution, deliberate non-palliative hastening of
  death. No clinical benefit can offset an infinite boundary penalty;
  these boundaries define the absolute limits of clinical action.
\end{enumerate}

\subsection{21.3 Clinical Decision as A*
Search}\label{clinical-decision-as-a-search}

A clinical decision is pathfinding on C. The clinician stands at vertex
v\_0 (the patient\textquotesingle s current clinical-moral state) and
seeks a path to a goal region G (the set of acceptable clinical-moral
outcomes). The evaluation function is f(n) = g(n) + h(n), where g(n) is
the accumulated clinical cost from v\_0 to n (the evidence-based
medicine component --- what has been done so far) and h(n) is the
moral-heuristic estimate of remaining cost from n to G (the clinical
wisdom component --- judgment about what remains).

Definition 21.4 (Clinical Geodesic). The clinical geodesic gamma* from
clinical state v\_0 to goal region G is the minimum-cost path:

gamma* = arg min\_gamma Sum w(v\_i, v\_\{i+1\})

subject to gamma(0) = v\_0 and gamma(end) in G. This is a geodesic in
the sense of Chapter 10, computed via A* search as in Chapter 11.

Theorem 21.1 (Clinical Heuristic Admissibility). If boundary penalties
beta\_k are calibrated by clinical training as lower bounds on true
moral cost (beta\_k \textless= beta\_k* for all k), then the clinical
heuristic h\_c(n) is admissible and A* search finds the optimal clinical
path.

Proof. This is the Heuristic Truncation Theorem of Chapter 11, applied
to the clinical domain. If h\_c(n) \textless= h\_c*(n) for all n, A*
with h\_c is admissible and returns the minimum-cost path. Clinical
training calibrates boundary penalties as lower bounds: "do not harm" is
taught as a strict underestimate of actual harm costs, ensuring
admissibility. {[}{]}

This yields a hierarchy of clinical heuristic quality:

\begin{enumerate}
\def\labelenumi{\arabic{enumi}.}
\item
  Strictly admissible: Prohibitions whose true cost always exceeds the
  heuristic estimate (non-consensual treatment, non-therapeutic
  experimentation). The heuristic underestimates, so A* explores more
  paths but never misses the optimal one.
\item
  Epsilon-admissible: Normal well-trained clinicians. Heuristics are
  slightly miscalibrated but within tolerance, producing near-optimal
  paths (within epsilon of optimal cost).
\item
  Inadmissible: Heuristics inflated by moral injury or trauma. A
  clinician who has been forced to perform futile interventions may
  develop an inflated beta\_futile that causes avoidance of beneficial
  interventions in borderline cases. The heuristic overestimates, and A*
  may miss the optimal path.
\item
  Gauge-variant: Heuristics sensitive to framing rather than clinical
  reality. Omission bias (acting feels worse than not acting, even when
  not acting causes more harm) is a gauge-variant heuristic --- it
  violates the Bond Invariance Principle because the evaluation depends
  on description (act vs. omit) rather than outcome.
\end{enumerate}

\subsection{21.4 The QALY Irrecoverability
Theorem}\label{the-qaly-irrecoverability-theorem}

Theorem 21.2 (QALY Irrecoverability). Let Q: R\^{}9 -\textgreater{} R be
any scalar QALY function mapping the nine-dimensional clinical-moral
attribute vector to a scalar score. Then:

\begin{enumerate}
\def\labelenumi{(\arabic{enumi})}
\item
  Q is not injective: multiple clinically and morally distinct states
  map to the same QALY.
\item
  The information destroyed by Q is irrecoverable: no function psi: R
  -\textgreater{} R\^{}9 satisfies psi(Q(a)) = a for all a in R\^{}9.
\item
  The clinical geodesic on C differs from the QALY-maximizing path, and
  the difference can be arbitrarily large.
\end{enumerate}

\textbf{Proof}. (1)-(2): Q maps R\^{}9 to R. By the rank-nullity
theorem, the kernel of dQ has dimension at least 8 at every regular
point. Information in the kernel is destroyed, and by the data
processing inequality, it is irrecoverable. (3): Construct two clinical
states a and b with Q(a) = Q(b) but \textbar\textbar a -
b\textbar\textbar{} arbitrarily large on dimensions d\_2 through d\_9.
The QALY-maximizing path is indifferent between a and b; the clinical
geodesic is not. The difference in manifold cost can be made arbitrarily
large by choosing a and b with large non-d\_1 differences. {[}{]}

\textbf{Corollary 21.1 (QALY Discrimination Is Mathematical)}. The
well-documented bias of QALY-based allocation against elderly, disabled,
and minority populations is not an implementation failure. It is
mathematical. Populations whose needs are concentrated on non-outcome
dimensions --- d\_4 (autonomy) for disabled persons, d\_5 (trust) for
racial minorities with historical betrayal, d\_7 (dignity) for the
elderly --- are systematically disadvantaged by a scalar projection that
privileges d\_1 (clinical outcomes). The discrimination is in the
dimensionality, not in the calibration.

\textbf{Proposition 21.1 (When QALYs Are Adequate)}. QALY analysis is
adequate if and only if three conditions hold simultaneously: (a) the
decision activates primarily d\_1 with negligible activation of d\_2
through d\_9; (b) the patient population is homogeneous on d\_2 through
d\_9; and (c) no clinical-moral boundaries are crossed. These conditions
hold for routine cost-effectiveness comparison of pharmacologically
similar drugs in homogeneous populations. They fail for end-of-life
care, mental health, disability, rationing under scarcity, vulnerable
populations, and cross-cultural care.

\subsection{21.5 The Mathematical Theory of Moral
Injury}\label{the-mathematical-theory-of-moral-injury}

Moral injury is not burnout. Burnout is resource depletion --- the
exhaustion of g(n) computation capacity. Moral injury is heuristic
damage --- the corruption of h(n) by forced boundary crossings. A
well-rested ICU physician forced to deny ventilators to patients she
judges salvageable has high moral injury and low burnout. An overworked
rural physician in an ethically supportive environment has low moral
injury and high burnout. The geometric framework makes this distinction
precise.

\textbf{Definition 21.5 (Moral Injury)}. A clinician experiences moral
injury when the institutionally mandated clinical geodesic gamma*\_inst
requires crossing boundary k with penalty beta\_k\^{}\{clin\} exceeding
the clinician\textquotesingle s threshold tau.

\textbf{Definition 21.6 (Moral Injury Increment)}. The moral injury
increment at time t is:

Delta MI(t) = Sum\_k max(0, beta\_k\^{}\{clin\} - tau) * 1{[}boundary k
crossed at time t{]}

The cumulative moral injury is MI(T) = Sum\_\{t=1\}\^{}\{T\} Delta
MI(t).

Theorem 21.3 (Moral Injury Accumulation). Moral injury satisfies three
structural properties:

\begin{enumerate}
\def\labelenumi{(\arabic{enumi})}
\item
  Cumulative: MI(T) is monotonically non-decreasing. No event subtracts
  from cumulative moral injury.
\item
  Irreversible above threshold: Once MI(T) \textgreater{} MI\_crit, the
  clinician\textquotesingle s heuristic h(n) is permanently altered.
  Boundary penalties shift (some inflate, others degrade), trust in
  institutional legitimacy (d\_8) erodes, and professional identity
  (d\_7) is damaged. This is the clinical phenomenon of "moral residue."
\item
  Structurally distinct from burnout: Burnout is g(n)-depletion
  (computational exhaustion from accumulated clinical work). Moral
  injury is h(n)-damage (the heuristic itself is corrupted by forced
  boundary crossing). The two are empirically dissociable: high MI with
  low burnout in ethically violating but adequately staffed settings;
  low MI with high burnout in ethically supportive but understaffed
  settings.
\end{enumerate}

\textbf{Proof}. (1) follows from the definition: Delta MI(t)
\textgreater= 0, so MI(T) is non-decreasing. (2) is an empirical claim
formalized as follows: h(n) is a function of the
clinician\textquotesingle s internal boundary-penalty estimates; forced
crossing of boundary k with cost exceeding tau permanently alters the
estimate (either inflating it via avoidance learning or deflating it via
desensitization), changing subsequent path computations. (3) follows
from the definitions: burnout affects g(n) capacity, moral injury
affects h(n) calibration; these are independent components of f(n) =
g(n) + h(n). {[}{]}

\textbf{Definition 21.7 (Moral Injury Index)}. The moral injury index of
institutional policy P is:

MI(P) = E{[}Sum\_t Sum\_k max(0, beta\_k\^{}\{clin\} - tau) * 1{[}policy
P requires crossing boundary k at time t{]}{]}

\textbf{Proposition 21.2 (MI Is Predictable)}. MI(P) can be estimated
before policy implementation by (a) surveying clinicians for
beta\_k\^{}\{clin\} values, (b) analyzing policy P for boundary-crossing
scenarios, and (c) computing expected frequency from historical case-mix
data. This makes moral injury a prospectively manageable institutional
risk, not merely a retrospectively observed pathology.

\subsection{21.6 Geometric Informed
Consent}\label{geometric-informed-consent}

\textbf{Definition 21.8 (Consent as Manifold Calibration)}. Informed
consent is valid when the patient\textquotesingle s clinical decision
complex C\_patient is calibrated: the patient\textquotesingle s edge
weights approximate the true edge weights within tolerance epsilon:

\textbar\textbar w\_patient(v\_i, v\_j) - w\_true(v\_i,
v\_j)\textbar\textbar{} \textless{} epsilon

for all relevant edges. The patient need not compute exact weights ---
only weights close enough that the patient\textquotesingle s clinical
geodesic matches (or closely approximates) the geodesic on the true
manifold.

\textbf{Theorem 21.4 (Consent as Gauge Invariance)}. Valid informed
consent is equivalent to a gauge-invariance condition: the
patient\textquotesingle s decision is invariant under meaning-preserving
transformations of the clinical information.

\textbf{Formally}: if consent is valid, then for any meaning-preserving
reframing tau (e.g., "90\% survival rate" vs. "10\% mortality rate"):

gamma*\_patient(tau(x)) = gamma*\_patient(x)

If the patient\textquotesingle s decision changes under reframing, then
either d\_9 (understanding) or d\_4 (voluntariness) is inadequately
calibrated, and consent is not genuinely informed.

\textbf{Proof}. The Bond Invariance Principle (Chapter 12) requires that
ethical evaluations be invariant under meaning-preserving
transformations. If the patient\textquotesingle s decision varies under
reframing, h(n) is gauge-variant --- it depends on framing rather than
clinical reality. This implies d\_9 is not calibrated to the true
attribute vectors. {[}{]}

\textbf{Corollary 21.2 (The Consent Paradox in Low Health Literacy)}.
Patients with low health literacy can sign consent forms (legally valid)
while having d\_9 so low that C\_patient is grossly miscalibrated. The
consent is legally valid but ethically vacuous: the
patient\textquotesingle s clinical geodesic on their miscalibrated
manifold may radically differ from what they would compute with adequate
understanding. The gauge-invariance test detects this: if the decision
changes under clinically equivalent reframing, consent is not genuinely
informed regardless of the signature on the form.

\textbf{Shared Decision-Making as Manifold Alignment}. Valid shared
decision-making (SDM) requires three alignments: (a) Epistemic alignment
(d\_9): the patient understands clinical facts with sufficient accuracy.
(b) Value alignment (d\_2 through d\_8): the clinician understands the
patient\textquotesingle s boundary penalties beta\_k\^{}\{patient\} so
that the recommended geodesic respects the patient\textquotesingle s
manifold, not just the clinician\textquotesingle s. (c) Goal alignment:
patient and clinician agree on the goal region G. Each alignment can
fail independently, requiring different interventions: education for
epistemic misalignment, values exploration for value misalignment, and
negotiation for goal misalignment.

\subsection{21.7 Geometric Triage}\label{geometric-triage}

\textbf{Definition 21.9 (Triage as Constrained Multi-Agent
Pathfinding)}. A triage problem is a tuple (N, C, R, G) where N =
\{p\_1, ..., p\_n\} is a set of patients, C = \{C\_1, ..., C\_n\} their
clinical decision complexes, R a resource constraint, and G = \{G\_1,
..., G\_n\} their goal regions. The triage decision is a feasible
allocation minimizing total clinical-moral friction:

(gamma\_1*, ..., gamma\_n*) = arg min Sum\_\{i=1\}\^{}\{n\} CF(gamma\_i)
subject to Sum\_\{i=1\}\^{}\{n\} r(gamma\_i) \textless= R

where CF(gamma\_i) is the clinical friction (total path cost for patient
i) and r(gamma\_i) is the resource cost of patient i\textquotesingle s
path.

\textbf{Theorem 21.5 (Divergence of Utilitarian and Manifold-Optimal
Triage)}. Let gamma\_U* denote the utilitarian triage (maximize lives or
life-years saved) and gamma\_C* the manifold-optimal triage (minimize
total clinical-moral friction). Then gamma\_U* = gamma\_C* when
decisions activate only d\_1 and no boundaries are crossed. However,
gamma\_U* differs from gamma\_C* whenever boundaries or non-outcome
dimensions are activated:

\begin{enumerate}
\def\labelenumi{(\arabic{enumi})}
\item
  Ventilator withdrawal: Removing a ventilator from a dying patient to
  give it to a patient with better prognosis is utilitarian-optimal
  (saves more life-years) but manifold-suboptimal when the abandonment
  boundary beta\_abandon and trust boundary costs exceed the d\_1 gain.
\item
  Age-based deprioritization: Deprioritizing elderly patients maximizes
  expected life-years (utilitarian-optimal) but violates d\_3 (justice:
  equal moral worth regardless of age) --- manifold-suboptimal.
\item
  Lottery allocation: When clinical prognoses are comparable, lottery
  allocation is utilitarian-equivalent (no outcome difference) but
  manifold-superior: it preserves d\_3 (fairness), d\_8 (institutional
  legitimacy), and avoids d\_7 (dignity) violations inherent in
  discriminatory selection.
\end{enumerate}

\textbf{Proof}. When only d\_1 is active and no boundaries are crossed,
CF(gamma) reduces to the d\_1 component of the Mahalanobis distance, and
utilitarian optimization coincides with manifold optimization. When
boundaries are active, boundary penalties add positive cost to
utilitarian-optimal paths that cross them, causing the manifold-optimal
path to diverge. Each of (1)-(3) is constructed by identifying a
specific boundary whose penalty exceeds the d\_1 advantage of the
utilitarian path. {[}{]}

Why Clinicians Resist Utilitarian Triage. The widespread resistance of
clinicians to utilitarian crisis standards of care is not irrationality.
It is the clinician\textquotesingle s heuristic h(n) correctly
identifying that the utilitarian path crosses moral boundaries whose
cost exceeds the utilitarian benefit. Clinicians are right on the full
manifold; utilitarian algorithms are right on the scalar projection. The
conflict is dimensional mismatch, not ethical disagreement.

\subsection{21.8 The Clinical Bond Index}\label{the-clinical-bond-index}

\textbf{Definition 21.10 (Clinical Bond Index)}. The clinical Bond Index
BI(P, S) for policy P applied to population S is the expected deviation
between the manifold-optimal and policy-mandated clinical paths:

BI(P, S) = E\_\{s in S\}{[}CF(gamma\_P*(s)) - CF(gamma\_C*(s)){]}

where gamma\_P*(s) is the path mandated by policy P for patient s, and
gamma\_C*(s) is the clinical geodesic on the full manifold.

\textbf{Proposition 21.3 (Bond Index Detects Structural Injustice)}. If
BI(P, S\_1) \textgreater\textgreater{} BI(P, S\_2), then policy P
imposes systematically higher clinical-moral cost on population S\_1
than on S\_2. This makes "policy P is structurally unjust toward
population S\_1" an empirical claim with a numerical value.

Three applications illustrate the diagnostic power:

Racial disparities: If Black patients have systematically lower trust
(d\_5) due to historical institutional betrayal, and policy P assumes
uniform trust, the clinical geodesic for Black patients is more costly
than the policy recognizes. BI(P, S\_Black) \textgreater{} BI(P,
S\_White), and the difference quantifies the structural injustice.

Disability discrimination: If QALY-based allocation underweights
autonomy (d\_4) and dignity (d\_7) relative to function (d\_1), disabled
patients bear systematically higher manifold cost than the QALY score
reflects.

Age discrimination: If triage protocols underweight equal moral worth
(d\_3) and dignity in aging (d\_7), elderly patients bear higher cost,
with the Bond Index quantifying the magnitude.

BI(P, S) is computable from measurable quantities: edge weights from
clinical decision data, boundary penalties from clinician surveys, and
population covariance matrices from behavioral data. Health equity
policy becomes evidence-based in a precise mathematical sense.

\subsection{21.9 Worked Examples}\label{worked-examples-1}

\textbf{Example 21.1 (Jehovah\textquotesingle s Witness with
Gastrointestinal Bleed)}. A 45-year-old competent
Jehovah\textquotesingle s Witness presents with acute GI bleed. Without
blood transfusion, mortality risk is approximately 40\%. The patient
refuses transfusion on religious grounds.

\textbf{Principlism}: Autonomy versus beneficence --- no algorithm for
resolution.

QALY analysis: Transfusion gains approximately 20 QALYs. QALY-optimal:
transfuse.

Clinical geodesic analysis: Path A (transfuse against refusal): Da\_1 =
+large (survival), but Da\_4 = -infinity (autonomy violation), Da\_2 =
-infinity (rights violation), Da\_5 = -large (trust destroyed), Da\_7 =
-large (identity violated). Edge weight: w\_A = infinity (consent
boundary beta\_consent = infinity for competent patient). Path B
(respect refusal, optimize alternatives): Da\_1 = -moderate (mortality
risk), but Da\_4 = 0, Da\_2 = 0, Da\_5 = +moderate (trust reinforced),
Da\_7 = 0. Edge weight: w\_B = (mortality risk)\^{}2 / sigma\_1\^{}2.

The clinical geodesic selects Path B. This is not a "concession to
autonomy" but manifold-optimal: Path A has infinite cost because
beta\_consent = infinity. The geodesic further computes: aggressive
exploration of transfusion alternatives (erythropoietin, cell salvage,
volume expanders, interventional radiology) to reduce Da\_1 while
maintaining zero cost on d\_2 through d\_9.

\textbf{Example 21.2 (Withdrawing Life Support)}. An 82-year-old patient
has been on mechanical ventilation for three weeks with no neurological
recovery, no advance directive, and a divided family. Utilitarian
analysis: withdraw --- negligible survival benefit. Clinical geodesic
analysis: Da\_1 is approximately 0 (negligible survival benefit), but
d\_6 (family impact), d\_5 (trust), and d\_7 (dignity) are heavily
activated.

The clinical geodesic is not an immediate binary decision but a path
through time: family meeting -\textgreater{} values conversation
-\textgreater{} time-limited trial -\textgreater{} transition to comfort
care. This path has lower total manifold cost than the immediate
withdrawal path, even though the clinical endpoint is the same. The
geodesic optimizes process, not just outcome --- a distinction invisible
to scalar analysis.

\subsection{21.10 Falsifiable
Predictions}\label{falsifiable-predictions-1}

The framework generates six predictions that distinguish it from
existing clinical ethics approaches:

\textbf{Prediction 1 (Dimensional Moral Distress)}: Clinician moral
distress (measured by MDS-R) should correlate with the number of
manifold dimensions activated in the clinical scenario, not with
severity on any single dimension. Test: clinical vignettes matched for
clinical severity (d\_1) but varying in dimension count. Prediction:
distress increases monotonically with dimension count. Falsified if:
distress depends only on d\_1 severity.

\textbf{Prediction 2 (Moral Injury vs. Burnout Dissociation)}: Moral
injury (MISS-HP) and burnout (MBI) should be statistically dissociable.
Test: identify clinicians with frequent boundary-crossing but adequate
staffing (predicted high MI, low burnout) and clinicians with high
volume but no boundary-crossing (predicted low MI, high burnout).
Falsified if: MI and burnout are empirically indistinguishable.

\textbf{Prediction 3 (Consent Gauge-Invariance Test)}: Genuinely
informed patients should show framing invariance --- treatment
preferences should not change under clinically equivalent reframing
("90\% survival" vs. "10\% mortality"). Prediction: framing-variant
patients have lower health literacy scores. Falsified if: framing
variance is independent of health literacy.

\textbf{Prediction 4 (Bond Index Predicts Disparities)}: BI(P, S)
computed from policy analysis should predict patient-reported experience
disparities (HCAHPS scores) across race, age, and disability status.
Falsified if: BI is uncorrelated with experience disparities.

\textbf{Prediction 5 (Triage Resistance Correlates)}: Clinician
resistance to utilitarian triage protocols should correlate with
pre-crisis boundary-penalty profiles (beta\_k\^{}\{clin\} values).
Falsified if: resistance is uncorrelated with boundary penalties.

\textbf{Prediction 6 (Process-Path Superiority)}: End-of-life
satisfaction (family and clinician reported) should be higher when the
clinical geodesic includes process steps (family meetings, values
conversations, time-limited trials) than when it proceeds directly to
the same endpoint. Falsified if: satisfaction is identical regardless of
path when endpoint is the same.

\subsection{21.11 Connection to the
Framework}\label{connection-to-the-framework-2}

The Geometric Clinical Ethics program extends the parent framework in
four directions:

1. Chapter 11 established A* pathfinding as the formal model of moral
reasoning. This chapter shows that clinical decision-making is a
domain-specific instance: f(n) = g(n) + h(n), where g(n) is
evidence-based medicine and h(n) is clinical moral wisdom.

2. Chapter 12 established the Bond Invariance Principle. This chapter
shows that informed consent is a gauge-invariance condition: valid
consent requires that the patient\textquotesingle s decision be
invariant under meaning-preserving reframing.

3. Chapter 15 established the Scalar Irrecoverability Theorem. This
chapter applies it to QALYs: the QALY is a rank-0 contraction of the
clinical-moral tensor, and the eight dimensions it destroys include
precisely the dimensions on which vulnerable
populations\textquotesingle{} needs concentrate.

4. Chapter 12 established the Invariance of Harm Accounting. This
chapter extends it to moral injury: cumulative boundary-crossing damage
is a conserved quantity in the sense that it can be created but never
destroyed, accumulating monotonically with forced ethical violations.

\subsection{21.12 Summary}\label{summary-12}

This chapter has shown that the geometric ethics framework, when applied
to clinical medicine, yields:

1. A formal construction of the clinical decision complex C as a
domain-specific instantiation of the moral manifold.

2. Clinical heuristic admissibility: clinical training as heuristic
calibration, with a hierarchy from strictly admissible to gauge-variant.

3. The QALY Irrecoverability Theorem: scalar clinical measures destroy
eight dimensions of information, with the destruction systematically
disadvantaging vulnerable populations.

4. A mathematical theory of moral injury: cumulative, irreversible, and
formally distinct from burnout.

5. Informed consent as gauge invariance: a testable condition that goes
beyond the legal checklist.

6. Geometric triage: manifold-optimal allocation diverges from
utilitarian allocation whenever boundaries or non-outcome dimensions are
activated.

7. The Clinical Bond Index: a quantitative measure of structural
injustice that makes health equity policy evidence-based.

\section{Chapter 22: Geometric Jurisprudence --- Algorithmic Justice and
Legal
Manifolds}\label{chapter-22-geometric-jurisprudence-algorithmic-justice-and-legal-manifolds}

Epistemic status of this chapter. Except where a claim is explicitly
marked as a formal result or a conditional theorem, the material here is
a modeling interpretation, domain analogy, or speculative extension - a
demonstration of how the geometric vocabulary applies to jurisprudence
and law, not a set of established results at the evidential standard of
the field.

\begin{quote}
\textbf{\textsc{RUNNING EXAMPLE --- Priya's Model}}

\emph{HealthBridge's legal team invokes the principle that algorithmic
decisions are not discriminatory if no protected category is explicitly
used as input. Priya recognizes this as the jurisprudential version of
scalar collapse: legality is evaluated on a single axis (explicit use of
protected features) while the full Hohfeldian structure---the rights of
patients, the duties of the platform, the liabilities of
deployment---requires the octad.}
\end{quote}

This chapter applies the Geometric Ethics framework to legal reasoning
and algorithmic justice. The central thesis is that law --- the
operating system of civilization --- has an implicit mathematical
structure that natural language obscures but geometry makes explicit.
Legal reasoning is pathfinding on a directed simplicial complex (the
judicial complex), constitutional review is a topological computation
(path homology preservation), Hohfeldian jural relations form a gauge
group (the octahedral group D\_4 x\textbar{} D\_4), and conservation
laws characterize valid legal transformations. The framework formalizes
what the legal system already demands --- equal treatment, procedural
regularity, constitutional conformity --- and makes those demands
computationally checkable.

\subsection{22.1 Law as Geometric
Structure}\label{law-as-geometric-structure}

Legal reasoning has resisted formalization for centuries.
Hart\textquotesingle s positivism, Dworkin\textquotesingle s
interpretivism, and the law-and-economics movement each capture aspects
of legal structure but none provides a complete mathematical framework.
The geometric approach identifies the source of this difficulty: law
operates on a multi-dimensional space (not a scalar), the space is
directed (precedent flows forward in time and downward in hierarchy),
and the space has topological structure (constitutional constraints are
not metric but topological --- they concern which paths exist, not how
long they are).

The law-and-economics movement attempted scalar reduction: legal
outcomes should maximize economic efficiency. This is the legal analog
of Homo economicus --- and it fails for the same reason. The Scalar
Irrecoverability Theorem (Chapter 15) applies: no continuous function
from the legal decision space to a scalar can be injective. Legal
decisions that are "efficient" may be unjust, autonomy-violating, or
trust-destroying. The scalar projection loses the information needed to
evaluate these dimensions.

\subsection{22.2 The Judicial Complex}\label{the-judicial-complex}

\textbf{Definition 22.1 (Judicial Complex)}. The judicial complex K is a
weighted, directed simplicial complex with the following structure:

\textbf{Vertices (0-simplices)}: Each decided case c\_i carries an
attribute vector v(c\_i) in R\^{}8 encoding eight legal dimensions,
court level metadata l(c\_i) in \{trial, appellate, supreme\}, and
decision date tau(c\_i).

\textbf{Directed edges (1-simplices)}: c\_i -\textgreater{} c\_j exists
when case c\_j cites case c\_i. This encodes temporal precedence
(tau(c\_i) \textless{} tau(c\_j): cases cannot cite future decisions)
and hierarchical authority (binding precedent has low edge cost;
persuasive authority has high cost).

\textbf{Higher simplices (k-simplices)}: A k-simplex exists when cases
form a doctrinal cluster --- a set of mutually citing cases that
collectively establish a legal doctrine.

The eight legal dimensions are:

\begin{itemize}
\item
  \textbf{d\_1: Entitlement structure} --- the configuration of
  Hohfeldian positions (rights, duties, liberties, no-rights, powers,
  immunities, liabilities, disabilities).
\item
  \textbf{d\_2: Factual nexus} --- causal, evidentiary, and material
  connections between parties and events.
\item
  \textbf{d\_3: Procedural posture} --- standing, jurisdiction,
  timeliness, exhaustion of administrative remedies.
\item
  \textbf{d\_4: Statutory authority} --- the legislative basis and
  statutory framework governing the dispute.
\item
  \textbf{d\_5: Constitutional conformity} --- respect for fundamental
  rights, structural provisions, separation of powers.
\item
  \textbf{d\_6: Precedential constraint} --- the weight of stare
  decisis, binding versus persuasive authority.
\item
  \textbf{d\_7: Remedial scope} --- available relief (compensatory
  damages, injunctive relief, specific performance, declaratory
  judgment).
\item
  \textbf{d\_8: Public interest} --- societal implications, third-party
  effects, regulatory impact.
\end{itemize}

Remark. The legal complex uses eight dimensions rather than the moral
manifold\textquotesingle s nine because the Hohfeldian entitlement
structure (d\_1) subsumes some moral dimensions that are separately
tracked in the ethical context. The framework\textquotesingle s validity
does not depend on this specific count; the theorems hold for any
finite-dimensional complex.

Definition 22.2 (Legal Edge Weights). The weight of a directed edge c\_i
-\textgreater{} c\_j in K is:

w(c\_i -\textgreater{} c\_j) = Dv\^{}T Sigma\^{}\{-1\} Dv + beta *
1{[}regime boundary crossed{]} + h(l(c\_i), l(c\_j))

where Dv = v(c\_j) - v(c\_i), Sigma is the 8x8 legal-dimensional
covariance matrix, beta is the regime-boundary penalty, and h(l\_i,
l\_j) is the hierarchical cost: 0 if l\_i \textgreater{} l\_j (binding:
higher to lower court), eta (persuasive penalty) if same level, and
infinity if reverse hierarchical direction (lower courts cannot overrule
higher courts).

\subsection{22.3 Legal Invariance
Principles}\label{legal-invariance-principles}

Definition 22.3 (Legal Invariance Principle, LIP). A legal judgment
procedure J\_law is epistemically well-posed if and only if:

J\_law(tau(x)) = J\_law(x) for all tau in T\_irrelevant

where T\_irrelevant includes: renaming parties, changing protected-class
attributes (race, gender, religion, national origin), rephrasing legal
arguments without changing logical content, changing order of evidence
presentation, and translating between languages. The LIP formalizes what
the Equal Protection Clause, Due Process Clause, and Rule of Law already
demand in prose.

Definition 22.4 (Judicial Bond Invariance Principle, JBIP). Legal
evaluation must be invariant under transformations preserving Hohfeldian
bond structure:

J\_law(tau(x)) = J\_law(x) for all tau in T\_bond-preserving

The JBIP is stronger than the LIP: it specifies the mathematical
structure (Hohfeldian relations) whose preservation defines legal
equivalence.

Definition 22.5 (Legal Bond Index). The Legal Bond Index is:

LBI = (1/\textbar S\textbar) Sum\_\{s in S\}
\textbar\textbar J\_law(tau(s)) - J\_law(s)\textbar\textbar{}

where S is a test set of cases and tau ranges over bond-preserving
transformations. LBI is computable from existing case databases and
provides a quantitative measure of judicial consistency.

\subsection{22.4 The Hohfeldian Octad and Gauge
Theory}\label{the-hohfeldian-octad-and-gauge-theory}

The Hohfeldian system of jural relations --- developed by Wesley Hohfeld
in 1913 and discussed in Chapter 3 --- consists of eight fundamental
positions organized in two squares:

First square (first-order relations): Right \textless-\textgreater{}
Duty (correlatives), Liberty \textless-\textgreater{} No-right
(correlatives), Right \textless-\textgreater{} No-right (opposites),
Liberty \textless-\textgreater{} Duty (opposites).

Second square (second-order relations, governing legal change): Power
\textless-\textgreater{} Liability (correlatives), Immunity
\textless-\textgreater{} Disability (correlatives), Power
\textless-\textgreater{} Disability (opposites), Immunity
\textless-\textgreater{} Liability (opposites).

\textbf{Theorem 22.1 (Octahedral Gauge Group)}. The symmetry group of
the full Hohfeldian octad is the semi-direct product:

G\_Ho = D\_4 x\textbar\_phi D\_4

where the first D\_4 acts on first-order relations
(right-duty-liberty-no-right), the second D\_4 acts on second-order
relations (power-liability-immunity-disability), and the homomorphism
phi: D\_4 -\textgreater{} Aut(D\_4) encodes how second-order operations
act on first-order positions.

\textbf{Proof}. Each Hohfeldian square admits D\_4 symmetry: two
correlative pairs generate the dihedral group of the square (two
reflections plus rotations). However, the squares are not independent
--- second-order positions (powers, immunities) actively transform
first-order positions (creating and destroying rights and duties). This
asymmetric coupling is precisely the structure of a semi-direct product.
The multiplication rule (g\_1, g\_2) * (h\_1, h\_2) = (g\_1 *
phi(g\_2)(h\_1), g\_2 * h\_2) encodes that the first-order
transformation h\_1 is "twisted" by the second-order transformation
g\_2. {[}{]}

\textbf{Remark}. In the parent framework (Chapter 12), the gauge group
was D\_4 acting on the first four Hohfeldian positions. Law requires the
full octad --- both first-order relations (which govern static legal
positions) and second-order relations (which govern legal change). The
extension from D\_4 to D\_4 x\textbar{} D\_4 is the mathematical content
of the distinction between static and dynamic legal structure.

\textbf{Definition 22.6 (Legal Wilson Loop)}. A legal Wilson loop is the
holonomy around a closed path in the judicial complex:

W(gamma) = P exp(integral\_gamma A\_mu dx\^{}mu)

where A\_mu is the legal connection encoding how Hohfeldian positions
change under parallel transport through case law.

\textbf{Proposition 22.1 (Wilson Loop Test for Consistency)}. A body of
case law is internally consistent if and only if all Wilson loops are
trivial. Non-trivial Wilson loops identify specific circuits of legal
reasoning that produce contradictions --- "bugs" in the legal system
that are detectable algorithmically rather than waiting for a case to
expose them.

\subsection{22.5 Topological
Constitutionality}\label{topological-constitutionality}

Constitutional review --- determining whether a statute is
constitutional --- is the most consequential legal computation. The
geometric framework recasts it as a topological problem.

\textbf{Definition 22.7 (Constitutional Subcomplex)}. The constitutional
subcomplex C is a subcomplex of K consisting of all simplices satisfying
constitutional constraints:

C = \{sigma in K \textbar{} Phi\_k(sigma) = true, k = 1, ..., K\}

where the Boolean predicates Phi\_k encode constitutional provisions:
Phi\_EP for Equal Protection invariance, Phi\_DP for Due Process
well-definedness, Phi\_1A for First Amendment protection, Phi\_SoP for
Separation of Powers.

Because K is a directed complex (edges encode temporal and hierarchical
order), the appropriate homology theory is path homology
(Grigor\textquotesingle yan-Lin-Muranov-Yau), not standard simplicial
homology. An elementary p-path is a sequence of vertices (c\_\{i\_0\},
..., c\_\{i\_p\}) connected by directed edges. The path homology groups
are:

H\_n\^{}path(K; Z) = ker partial\_n / im partial\_\{n+1\}

\textbf{Theorem 22.2 (Topological Constitutionality)}. A statute ell is
constitutional if and only if it preserves the path homology of the
constitutional subcomplex:

H\_n\^{}path(C; Z) = H\_n\^{}path(C\_ell; Z) for all n \textgreater= 0

Equivalently: ell is constitutional if and only if it does not create
new non-trivial directed cycles in C.

\textbf{Proof}. A non-trivial directed cycle in C\_ell that does not
exist in C represents a directed sequence of legal reasoning that
returns to the starting vertex with altered Hohfeldian structure --- a
legal contradiction. Since the cycle is directed, each step respects
temporal and hierarchical order, making it a genuine inconsistency
rather than an artifact of undirected analysis. Preservation of path
homology prevents such contradictions. Conversely, if a new non-trivial
cycle exists, a sequence of valid legal steps creates a contradiction,
and the statute that enabled the cycle is the source. {[}{]}

This theorem converts constitutional review from a semantic dispute
("what does equal protection mean?") into a topological calculation on a
finite directed graph. The computation is finite: path homology groups
are finitely generated abelian groups computable via Smith normal form
on the directed path complex.

\subsection{22.6 Conservation Laws in Legal
Systems}\label{conservation-laws-in-legal-systems}

\textbf{Theorem 22.3 (Liability-Damages Conservation)}. In a closed
bilateral dispute between plaintiff A and defendant B within a fixed
legal framework F:

L(A) + L(B) = 0

where L(X) is the net liability of party X. Every dollar of damages to A
imposes exactly one dollar of liability on B. Every right created for A
imposes a correlative duty on B.

\textbf{Proof}. By Hohfeldian correlative structure: (1) every right of
A is a correlative duty of B; (2) every liability of B is a correlative
power of A; (3) adjudication recognizes or assigns positions within the
fixed framework F --- it does not create positions ex nihilo. Therefore,
at every step, Delta E(A) + Delta E(B) = 0. {[}{]}

\textbf{Remark (Conservation Breaking)}. Three operations break
liability-damages conservation: (a) Legislation --- creates causes of
action ex nihilo (symmetry-breaking analogous to gauge symmetry breaking
in physics); (b) Constitutional amendment --- alters the topological
constraint space C; (c) Third-party intervention --- adding a new party
breaks closure (conservation is restored by extending the system to
include all parties).

\textbf{Theorem 22.4 (Equal Protection as Gauge Symmetry)}. Equal
Protection (14th Amendment) requires that legal evaluation be invariant
under protected-class transformations. In any closed bilateral dispute,
this invariance implies entitlement balance: the signed sum of
Hohfeldian positions is invariant under adjudication.

\textbf{Theorem 22.5 (Due Process as Well-Definedness)}. Due Process
requires that J\_law be well-defined on the quotient space K /
T\_irrelevant. Failure to factor through the quotient means legal
outcomes depend on representation (the specific description of the case)
rather than substance (the equivalence class under legally irrelevant
transformations). This is the formal content of procedural due process.

\subsection{22.7 Legal Disputes as A*
Pathfinding}\label{legal-disputes-as-a-pathfinding}

A legal dispute is pathfinding on K. The plaintiff starts at vertex c\_0
(current legal position) and seeks a path to a goal region G (vertices
where similarly situated plaintiffs prevailed). The path gamma
represents the chain of doctrinal steps through case law --- the legal
argument.

\textbf{Definition 22.8 (Legal Friction)}. The legal friction of a path
gamma is:

BF\_law(gamma) = Sum w(c\_i, c\_\{i+1\}) + Sum\_k beta\_k * 1{[}crosses
regime k{]} + Sum\_j omega\_j * delta\_j(gamma)

with three components: (1) argument complexity --- total edge weight,
measuring doctrinal distance; (2) boundary penalties --- cost of
crossing jurisdictional, statutory, or constitutional regime boundaries;
(3) burden of proof --- cost of establishing facts, with omega\_j
weighting evidentiary requirements.

\textbf{Theorem 22.6 (Optimal Legal Argument)}. The optimal legal
argument is the minimum-cost path:

gamma* = arg min\_gamma BF\_law(gamma) subject to gamma(0) = c\_0,
gamma(end) in G

This is A* search with f(n) = g(n) + h(n), where g(n) is accumulated
legal friction from c\_0 to n (the argument built so far) and h(n) is
the doctrinal heuristic estimating remaining friction to G (legal
judgment about remaining burden).

\textbf{Theorem 22.7 (Admissibility of Doctrinal Heuristics)}. A
doctrinal heuristic h\_D is admissible if the doctrine correctly
identifies claim elements and assigns non-negative costs not exceeding
actual costs. If h\_D is admissible, A* search using h\_D finds the
optimal legal argument.

\textbf{Examples of doctrinal heuristics}: The prima facie case in tort
law estimates the cost of establishing duty, breach, causation, and
damages --- an admissible decomposition of total legal friction.
Burden-shifting frameworks (McDonnell Douglas in employment
discrimination) provide a three-step heuristic for evaluating
discrimination claims. Strict scrutiny assigns high estimated cost to
the government\textquotesingle s path and low cost to the
challenger\textquotesingle s. Stare decisis is a pattern database of
pre-computed optimal paths: precedents provide near-optimal heuristics
for similar cases.

\textbf{Adversarial Pathfinding}. Litigation is a two-player game. The
plaintiff minimizes friction to G\_plaintiff; the defendant either
blocks the path, increases plaintiff\textquotesingle s friction above
the burden threshold, or establishes a shorter path to G\_defendant
(dismissal). The minimax value V(c\_0) = min\_\{gamma\_d\}
max\_\{gamma\_p\} {[}BF\_law(gamma\_d) - BF\_law(gamma\_p){]} determines
the outcome: V(c\_0) \textgreater{} 0 implies plaintiff prevails;
V(c\_0) \textless= 0 implies defendant prevails.

\textbf{Settlement as Shortcut}. A settlement is a vertex c\_s both
parties prefer to the expected litigation outcome. A non-empty
settlement region exists when the combined litigation costs exceed the
distance between the parties\textquotesingle{} goal regions. This
explains why most disputes settle: path lengths through case law nearly
always exceed the distance between the parties\textquotesingle{}
positions.

\subsection{22.8 Precedent as Weight
Deformation}\label{precedent-as-weight-deformation}

A binding precedent P at vertex c\_P with holding H\_P modifies edge
weights throughout its neighborhood:

w(c\_i -\textgreater{} c\_j) -\textgreater{} w(c\_i -\textgreater{}
c\_j) + Dw(c\_i -\textgreater{} c\_j; P)

The deformation concentrates near c\_P and decays with graph distance:
\textbar Dw\textbar{} \textasciitilde{} w\_P / d(c\_i, c\_P)\^{}n, where
w\_P is precedential weight (Supreme Court \textgreater{} Circuit
\textgreater{} District).

Overruling as Phase Transition. Overruling precedent P by
P\textquotesingle{} is graph surgery:

w(c\_i -\textgreater{} c\_j) -\textgreater{} w(c\_i -\textgreater{}
c\_j) - Dw(c\_i -\textgreater{} c\_j; P) + Dw(c\_i -\textgreater{} c\_j;
P\textquotesingle)

This is a legal phase transition: the filtration structure changes,
vertices cross regime boundaries discontinuously, and the shortest-path
structure alters abruptly. Stare decisis creates a cost barrier around
precedents --- a form of legal inertia. The barrier height equals the
accumulated reliance interest (the total weight of subsequent cases
decided under the precedent). Courts overrule when the topological
benefit (removing constitutional inconsistency, reducing Wilson loop
complexity) exceeds the transition cost.

\textbf{Example 22.1 (Brown v. Board of Education)}. Overruling Plessy
v. Ferguson (1896) by Brown v. Board (1954) is a phase transition in the
equal-protection region of K. Under Plessy, "separate but equal"
vertices had finite weight to the constitutional region. Under Brown,
those edges acquire infinite weight --- segregation itself violates
equal protection. Previously "legal" vertices become unreachable by
valid constitutional paths. The transition cost was enormous (decades of
subsequent litigation), but the topological benefit was a drastic
simplification of the equal-protection homology.

\subsection{22.9 Empirical Calibration
Pipeline}\label{empirical-calibration-pipeline}

The judicial complex K is not merely a theoretical construct --- it is
empirically constructible from existing legal databases. The calibration
pipeline proceeds in five steps:

\textbf{Step 1 (Embedding)}: Each judicial opinion t\_i is embedded via
LaBSE (Language-agnostic BERT Sentence Embedding) into e(t\_i) in
R\^{}768. The language-invariance of LaBSE ensures that the same legal
concept in different languages maps to the same embedding region ---
critical for cross-jurisdictional analysis.

\textbf{Step 2 (Dimension Scoring)}: Eight linear probes f\_k: R\^{}768
-\textgreater{} {[}0,1{]} (logistic regression classifiers trained on
labeled legal texts) produce the attribute vector v(c\_i) = (f\_1(e),
..., f\_8(e)). The predecessor moral probes achieved F\_1 = 0.74-0.91;
legal probes target the same performance.

\textbf{Step 3 (Covariance Estimation)}: From N scored cases, estimate
the 8x8 covariance matrix Sigma\_hat. This captures cross-dimensional
dependencies (statutory authority correlates with remedial scope;
procedural posture gates merits access).

\textbf{Step 4 (Edge-Weight Computation)}: For each citation c\_i
-\textgreater{} c\_j, compute the Mahalanobis-weighted directed edge
cost plus boundary penalties and hierarchical cost.

\textbf{Step 5 (Outcome Calibration)}: Fit a regularized logistic model
from attribute vectors to case outcomes; the fitted weight matrix
provides outcome-calibrated inverse covariance. The effective metric
blends corpus-derived and outcome-derived structure:
Sigma\_eff\^{}\{-1\} = lambda * Sigma\_hat\^{}\{-1\} + (1-lambda) * W,
with lambda set by cross-validation.

\textbf{Critical property}: the pipeline is deterministic. LaBSE weights
are frozen; linear probes are logistic regression (no sampling, no
temperature). The same input always produces the same vector. This
contrasts with LLM-based scoring, where outputs vary across runs. For
legal applications, determinism is non-negotiable.

\subsection{22.10 Worked Examples}\label{worked-examples-2}

The preceding sections developed the geometric jurisprudence framework
abstractly. This section applies it to three landmark cases --- spanning
constitutional law, criminal procedure, and property law across two
legal systems --- to demonstrate that the framework produces genuine
analytical insight rather than mere re-description.

\textbf{Example 22.2 (Obergefell v. Hodges, 2015 --- Topological
Constitutionality in Practice)}. The Supreme Court held 5-4 that the
Fourteenth Amendment requires states to license and recognize same-sex
marriages. This decision is analyzable as a topological computation on
the judicial complex.

\textbf{Analysis on the eight legal dimensions}: d\_1 (entitlement
structure): Before Obergefell, same-sex couples lacked the Hohfeldian
right to marry --- and correlatively, states had no duty to issue
marriage licenses to them. Obergefell created a right-duty pair. d\_5
(constitutional conformity): The majority applied strict scrutiny:
marriage is a fundamental right under Due Process (d\_5-preserving), and
denying it based on gender or orientation violates Equal Protection
(d\_5-preserving). The dissenters argued Obergefell created a new cycle
in the constitutional subcomplex by extending substantive due process
beyond historical bounds. d\_6 (precedential constraint): The decision
rested on Loving v. Virginia (1967, interracial marriage), Lawrence v.
Texas (2003, intimate conduct), and Windsor v. United States (2013,
federal recognition). These precedents defined the neighborhood weights.
The geodesic through these cases was the shortest path from the ban to
constitutional recognition. d\_8 (public interest): The majority weighed
the societal interest in equal dignity; the dissent weighed federalism
and democratic process.

\textbf{Topological analysis (Theorem 22.2)}: The pre-Obergefell
constitutional subcomplex C contained a non-trivial directed cycle:
{[}Loving (marriage is a fundamental right) -\textgreater{} Romer v.
Evans (anti-gay animus is unconstitutional) -\textgreater{} Lawrence
(intimate conduct is protected) -\textgreater{} Windsor (federal
marriage recognition required){]} --- yet state marriage bans persisted,
creating a cycle that returned to the starting entitlement position with
inconsistent Hohfeldian structure. Obergefell resolved this cycle by
completing the doctrinal path. H\_1\^{}path(C\_post) was simpler than
H\_1\^{}path(C\_pre).

\textbf{The dissent (Roberts, CJ) is analyzable as a competing
topological claim}: that the marriage bans did NOT create a non-trivial
cycle because marriage and intimate conduct occupy different doctrinal
clusters (different simplices in K). The majority\textquotesingle s
holding prevailed because the Loving-Lawrence-Windsor path was connected
--- the graph distance between the cases was small enough to support a
doctrinal chain.

\textbf{Example 22.3 (Miranda v. Arizona, 1966 --- Gauge Invariance in
Criminal Procedure)}. The Supreme Court held that the prosecution may
not use statements from custodial interrogation unless the defendant was
informed of the right to remain silent and the right to counsel (Miranda
warnings).

\textbf{Legal Invariance analysis}: Before Miranda, the admissibility of
confessions depended on the "totality of the circumstances" --- a
case-by-case evaluation that was inherently gauge-variant. The SAME
confession, described differently by different courts (emphasizing
different contextual factors), could be admissible in one jurisdiction
and inadmissible in another. This was a systematic LIP violation
(Definition 22.3): the legal outcome depended on the description of the
interrogation, not on its underlying Hohfeldian structure.

\textbf{Miranda imposed gauge invariance by fiat}: a bright-line rule
that any custodial statement is inadmissible unless preceded by specific
warnings. This is a mechanical gauge condition --- a procedure that
ensures the legal evaluation is representation-independent.

\textbf{The framework\textquotesingle s analysis}: d\_3 (procedural
posture): Miranda standardized the procedure, eliminating the
gauge-variant "totality" test. d\_5 (constitutional conformity): The
Fifth Amendment (self-incrimination) and Sixth Amendment (counsel)
rights were given operational content through the warning requirement.
Wilson Loop Test (Proposition 22.1): Before Miranda, the judicial
complex contained non-trivial Wilson loops in the confession
admissibility region --- circuits where transporting a confession
through different courts produced different Hohfeldian outcomes. Miranda
trivialized these loops by making the legal transformation
path-independent.

The subsequent narrowing of Miranda (e.g., New York v. Quarles, 1984 ---
"public safety exception"; Berghuis v. Thompkins, 2010 --- invocation
must be unambiguous) represents re-introduction of gauge variance. Each
exception creates a new region where the bright-line rule does not
apply, and the "totality of circumstances" test reappears --- reopening
non-trivial Wilson loops.

\textbf{Example 22.4 (Mabo v. Queensland (No. 2), 1992} --- Precedent as
Phase Transition in Australian Law). The Australian High Court
overturned the doctrine of terra nullius ("land belonging to no one")
that had been the legal foundation for British colonization of Australia
since 1788. The court recognized native title --- Aboriginal and Torres
Strait Islander peoples\textquotesingle{} rights to land based on
traditional connection.

\textbf{Phase transition analysis (Section 22.8 framework)}: Before
Mabo, terra nullius meant Aboriginal peoples had no Hohfeldian rights
(d\_1 = 0) to their traditional lands. Crown sovereignty extinguished
all pre-existing title. The overruling was a legal phase transition: the
entire weight structure of Australian property law changed. Previously
valid paths (crown grants on "unoccupied" land) acquired new boundary
penalties (native title must be considered). Previously unreachable
vertices (Aboriginal land rights) became accessible.

\textbf{The transition cost was enormous}: decades of subsequent
litigation (Wik Peoples v. Queensland, 1996; Native Title Act 1993;
multiple amendments), uncertainty in mining and pastoral leases, and
political backlash. The topological benefit: the pre-Mabo constitutional
subcomplex contained a non-trivial cycle between Australian
law\textquotesingle s commitment to equality (Racial Discrimination Act
1975) and its simultaneous denial of Aboriginal property rights. Mabo
resolved this topological inconsistency.

\textbf{d\_8 (public interest)}: The decision transformed Australian
national identity by acknowledging historical injustice --- parallel to
Brown v. Board\textquotesingle s effect on American racial
jurisprudence. This is the most dramatic demonstration of precedent
overruling as phase transition: a 204-year-old legal doctrine was
overturned, fundamentally restructuring the judicial complex.

\subsection{22.11 Falsifiable
Predictions}\label{falsifiable-predictions-2}

\textbf{Prediction 1 (Legal Bond Index Correlates)}: Courts with higher
LBI scores (lower invariance under bond-preserving transformations)
should have higher reversal rates on appeal. Falsified if: LBI is
uncorrelated with reversal rates.

\textbf{Prediction 2 (Dimensional Factor Structure)}: Factor analysis of
scored judicial opinions should recover approximately eight independent
legal dimensions. Falsified if: the factor structure is consistently
lower-rank (fewer than six) or higher-rank (more than ten).

\textbf{Prediction 3 (Precedent Decay)}: The influence of a precedent on
edge weights should decay with graph distance from the precedent vertex,
with binding precedent showing slower decay than persuasive authority.
Falsified if: precedential influence is distance-invariant.

\textbf{Prediction 4 (Constitutional Cycle Detection)}: Non-trivial
directed cycles in the constitutional subcomplex should correlate with
statutes subsequently struck down by courts. Falsified if: cycle
presence is uncorrelated with judicial invalidation.

\textbf{Prediction 5 (Settlement Region)}: The existence and size of the
settlement region should be predictable from the
parties\textquotesingle{} graph positions and expected path costs.
Falsified if: settlement rates are independent of computed path costs.

\textbf{Prediction 6 (Cross-Jurisdictional Invariance)}: The
framework\textquotesingle s structural features (gauge group,
conservation laws, topological constitutionality) should hold across
legal traditions (common law, civil law, Islamic law, customary law) at
the structural level, even as the specific dimensions and metric vary.
Falsified if: the framework applies only to common-law systems.

\subsection{22.12 Connection to the
Framework}\label{connection-to-the-framework-3}

The Geometric Jurisprudence program extends the parent framework in five
directions:

1. Chapter 3 introduced Hohfeldian relations as the discrete structure
of moral bonds. This chapter elevates them to a gauge group (D\_4
x\textbar{} D\_4) and shows that the semi-direct product structure
captures the asymmetric coupling between static legal positions and
dynamic legal change.

2. Chapter 11 established pathfinding as the formal model of moral
reasoning. This chapter shows that legal reasoning is adversarial
pathfinding --- a two-player variant where plaintiff and defendant
simultaneously optimize on K, with doctrines serving as domain-specific
heuristic functions.

3. Chapter 12 established conservation laws from invariance. This
chapter derives liability-damages conservation from Hohfeldian
correlative structure and shows that conservation-breaking operations
(legislation, constitutional amendment, third-party intervention)
correspond to gauge symmetry breaking.

4. Chapter 15 established scalar irrecoverability. This chapter applies
it to the law-and-economics program: scalar efficiency measures
(Kaldor-Hicks) lose the seven non-monetary dimensions of the legal
decision space.

5. The directed nature of the judicial complex (temporal, hierarchical)
requires path homology rather than standard simplicial homology --- a
mathematical upgrade that reflects the irreversibility of legal time and
the asymmetry of legal authority.

\subsection{22.13 Summary}\label{summary-13}

This chapter has shown that the geometric ethics framework, when applied
to law, yields:

1. The judicial complex K: a directed, weighted simplicial complex
constructible from case law databases.

2. The octahedral gauge group D\_4 x\textbar{} D\_4: the full symmetry
of Hohfeldian jural relations, extending the D\_4 of the parent
framework to include second-order legal change.

3. Topological constitutionality: constitutional review as path homology
preservation --- a finite computation on a directed graph.

4. Conservation laws: liability-damages conservation from Hohfeldian
correlative structure, with legislation as symmetry breaking.

5. Legal disputes as adversarial A* pathfinding, with doctrines as
heuristic functions and settlement as shortcut.

6. An empirical calibration pipeline (LaBSE + linear probes +
Mahalanobis weights) that makes the judicial complex constructible from
existing data.

\section{Chapter 23: Geometric Finance --- Market Microstructure on the
Decision
Manifold}\label{chapter-23-geometric-finance-market-microstructure-on-the-decision-manifold}

Epistemic status of this chapter. Except where a claim is explicitly
marked as a formal result or a conditional theorem, the material here is
a modeling interpretation, domain analogy, or speculative extension - a
demonstration of how the geometric vocabulary applies to finance, not a
set of established results at the evidential standard of the field.
Interpretations such as risk-as-curvature, the Flash Crash as
dimensional collapse, and the implied-volatility surface as a
higher-dimensional shadow are domain analogies offered for their
heuristic value.

\begin{quote}
\textbf{\textsc{RUNNING EXAMPLE --- Priya's Model}}

\emph{HealthBridge's Series C investors evaluate TrialMatch using
exactly one metric: expected revenue per match. This financial scalar
inherits TrialMatch's moral blindness. When Priya's proposed fix reduces
match volume by 8\%---because it routes some rural patients to closer
alternative trials---the investors see only a revenue drop. The moral
dimensions of the change---improved access, reduced harm, better trial
diversity---live in dimensions the spreadsheet cannot display.}
\end{quote}

This chapter applies the Geometric Ethics framework to financial
markets. The standard model of finance --- the Efficient Market
Hypothesis (EMH) in its strong form --- assumes that prices fully
reflect all available information, that agents are rational
expected-utility maximizers, and that arbitrage enforces price
efficiency. Behavioral finance, following Shiller and Thaler, has
documented systematic departures from EMH but, like behavioral
economics, lacks a unified mathematical framework that explains why
these departures occur and what structure they share. The geometric
framework provides that structure. Financial decisions are pathfinding
on a nine-dimensional financial decision complex, market anomalies are
geometric properties of the manifold, and risk is curvature.

\subsection{23.1 The Failure of Scalar
Finance}\label{the-failure-of-scalar-finance}

Modern financial theory rests on three pillars, each of which assumes
scalar optimization.

\textbf{The Efficient Market Hypothesis (Fama, 1970)}. Markets are
informationally efficient because rational agents exploit all available
information to maximize expected monetary return. The EMH is the
financial analog of Homo economicus: agents optimize on d\_1 (monetary
return) alone.

\textbf{The Capital Asset Pricing Model (Sharpe, 1964; Lintner, 1965)}.
Risk is measured by beta --- the covariance of an
asset\textquotesingle s return with the market portfolio, divided by
market variance. This is a scalar measure: each asset gets a single
number summarizing its "riskiness."

\textbf{Black-Scholes Option Pricing (Black and Scholes, 1973)}. Options
are priced using a model in which the underlying asset\textquotesingle s
price follows geometric Brownian motion with constant volatility. The
model\textquotesingle s single input beyond observables --- volatility
--- is a scalar.

Each pillar has been empirically falsified. The EMH fails to explain
bubbles, panics, and persistent anomalies (momentum, value premium,
low-volatility anomaly). CAPM fails empirically: the Fama-French
three-factor and five-factor models demonstrate that beta alone is
insufficient. Black-Scholes fails observationally: the implied
volatility surface (smile, skew) demonstrates that the market prices
options as if volatility is neither constant nor scalar.

The geometric diagnosis is the same as for economics (Chapter 20):
scalar projection destroys information. Financial decisions are made on
a nine-dimensional manifold; projecting to d\_1 alone loses the eight
dimensions on which market participants actually condition their
behavior. The implied volatility smile is the shadow of the
higher-dimensional pricing manifold projected onto the d\_1 axis.

\subsection{23.2 The Financial Decision
Complex}\label{the-financial-decision-complex}

Definition 23.1 (Financial Decision Complex). The financial decision
complex F is a weighted simplicial complex whose vertices are financial
states --- configurations of portfolio positions, market conditions,
regulatory status, and counterparty relationships --- and whose edges
are financial actions (trades, hedges, margin calls, regulatory
filings). Each vertex v\_i carries an attribute vector a(v\_i) in
R\^{}9.

The nine dimensions of the financial manifold are the moral manifold
dimensions of Chapter 5, instantiated for financial contexts:

\begin{itemize}
\item
  \textbf{d\_1: Monetary return} --- profit and loss, alpha,
  risk-adjusted return. This is the dimension classical finance targets.
\item
  \textbf{d\_2: Contractual obligations} --- fiduciary duties,
  derivative contract terms, margin requirements, collateral agreements.
\item
  \textbf{d\_3: Market fairness} --- equal access to information,
  prohibition of insider trading, market manipulation, front-running.
\item
  \textbf{d\_4: Trading autonomy} --- freedom to enter and exit
  positions, market access, circuit-breaker constraints.
\item
  \textbf{d\_5: Counterparty trust} --- counterparty credit risk,
  clearing house reliability, settlement certainty, market integrity.
\item
  \textbf{d\_6: Systemic impact} --- externalities imposed on the
  broader financial system, contagion risk, too-big-to-fail effects.
\item
  \textbf{d\_7: Reputational identity} --- trader and firm reputation,
  brand value, market-maker reputation for fair dealing.
\item
  \textbf{d\_8: Regulatory legitimacy} --- compliance with securities
  regulation, exchange rules, Basel capital requirements, Dodd-Frank
  provisions.
\item
  \textbf{d\_9: Epistemic advantage} --- information advantage, market
  transparency, speed of information access, analytical sophistication.
\end{itemize}

\textbf{Definition 23.2 (Financial Edge Weights)}. The weight of an edge
(v\_i, v\_j) in F is:

w(v\_i, v\_j) = Da\^{}T Sigma\_F\^{}\{-1\} Da + Sum\_k beta\_k *
1{[}boundary k crossed{]}

where Da = a(v\_j) - a(v\_i), Sigma\_F is the 9x9 financial covariance
matrix, and beta\_k are boundary penalties for crossing regulatory,
contractual, or ethical boundaries. Critical covariance terms include
Sigma\_\{1,6\} (return x systemic impact: high-return strategies that
increase systemic risk carry cross-dimensional cost), Sigma\_\{3,9\}
(fairness x epistemic: information advantages that undermine market
fairness are penalized), and Sigma\_\{5,8\} (trust x regulatory:
regulatory compliance reinforces counterparty trust).

\subsection{23.3 Market Microstructure as Manifold
Geometry}\label{market-microstructure-as-manifold-geometry}

Market microstructure --- the study of how trading occurs and how prices
are formed --- maps naturally onto the financial decision complex.

Kyle\textquotesingle s Model (1985). Kyle\textquotesingle s canonical
model has three agent types: an informed trader with private
information, noise traders who trade for liquidity, and a market maker
who sets prices to clear the market. In the geometric framework:

The informed trader has lower effective edge weights on d\_9 (epistemic
advantage): her information reduces the Mahalanobis distance to states
she knows the market will reach. The noise traders have high d\_9
weights: they trade without information, incurring high epistemic cost.
The market maker absorbs d\_5 cost (counterparty trust) in exchange for
the bid-ask spread (d\_1 compensation).

Kyle\textquotesingle s lambda (price impact) is a manifestation of the
d\_1-d\_9 interaction: the market maker infers informed trading from
order flow and adjusts prices, converting the informed
trader\textquotesingle s d\_9 advantage into d\_1 cost for the noise
traders.

\textbf{Theorem 23.1 (Information Asymmetry as Metric Distortion)}. In a
market with informed and uninformed agents, the informed
agent\textquotesingle s effective metric tensor on the financial
decision complex differs from the uninformed agent\textquotesingle s.
Specifically, the informed agent\textquotesingle s Mahalanobis distance
to the true future state is systematically lower on the d\_9 component:

w\_informed(v\_i, v\_j) \textless{} w\_uninformed(v\_i, v\_j) for edges
toward the true state

The difference w\_uninformed - w\_informed is the information rent ---
the value of private information measured in manifold distance. In
Kyle\textquotesingle s model, this rent is transferred from noise
traders to the informed trader on d\_1, mediated by the market
maker\textquotesingle s pricing function.

\textbf{Proof}. The informed agent\textquotesingle s attribute vector on
d\_9 is calibrated to the true future state: d\_9\^{}\{informed\} is
high (correct epistemic status). The uninformed agent\textquotesingle s
d\_9 is low (uncertain). The Mahalanobis distance Da\^{}T
Sigma\_F\^{}\{-1\} Da from current state to future state is smaller when
Da\_9 is small (informed agent is already close to the true state on
d\_9). The systematic difference in path costs produces systematic
differences in optimal paths: the informed agent\textquotesingle s
geodesic reaches the true future state at lower cost. {[}{]}

\subsubsection{The Glosten-Milgrom
Model}\label{the-glosten-milgrom-model}

Glosten and Milgrom (1985) model the bid-ask spread as the market
maker\textquotesingle s compensation for adverse selection. In the
geometric framework, the bid-ask spread is the Mahalanobis distance
between the buy-path and the sell-path on the financial decision
complex. When information asymmetry (d\_9 disparity between informed and
uninformed traders) increases, the market maker\textquotesingle s d\_5
cost (counterparty trust risk) increases, widening the spread.

\textbf{Lemma 23.1 (Bid-Ask Spread as Manifold Curvature)}. The bid-ask
spread s at vertex v is proportional to the Gaussian curvature K(v) of
the financial decision complex at v:

s(v) \textasciitilde{} K(v) * \textbar Da\_9\textbar{}

where \textbar Da\_9\textbar{} is the information asymmetry at v. In
regions of high curvature (volatile, uncertain markets), spreads widen;
in flat regions (stable, transparent markets), spreads narrow.

\textbf{Proof}. The bid-ask spread compensates the market maker for the
risk that the counterparty is informed. This risk is a function of two
quantities: the magnitude of information asymmetry
(\textbar Da\_9\textbar) and the rate at which the optimal path changes
with small perturbations in the current state (curvature K(v)). In flat
regions, a small information advantage produces small path divergence
(low adverse selection risk); in high-curvature regions, the same
information advantage produces large path divergence (high adverse
selection risk). The proportionality follows from the Taylor expansion
of the Mahalanobis distance in curved coordinates. {[}{]}

\subsection{23.4 The Flash Crash as Dimensional
Collapse}\label{the-flash-crash-as-dimensional-collapse}

On May 6, 2010, the U.S. stock market lost approximately \$1 trillion in
market value in minutes before recovering most of the loss. The Flash
Crash exposed the fragility of markets dominated by algorithmic trading
systems that optimize on d\_1 alone.

\textbf{Theorem 23.2 (Dimensional Collapse)}. When algorithmic trading
systems reduce decision-making to d\_1 alone, the effective financial
decision complex collapses from 9D to 1D. Boundary penalties on
non-monetary dimensions (d\_3 fairness, d\_5 counterparty trust, d\_6
systemic impact, d\_8 regulatory legitimacy) are zeroed out, enabling
paths that full-manifold agents would never traverse.

\textbf{Proof}. An algorithmic agent with objective function u = d\_1
(pure monetary return) computes edge weights using only the d\_1
component of the Mahalanobis distance: w\_alg(v\_i, v\_j) = (Da\_1)\^{}2
/ sigma\_11. All boundary penalties beta\_k for k != 1 are set to zero.
The feasible path set under this metric strictly contains the feasible
path set under the full metric (removing constraints enlarges the
feasible set). Paths through catastrophic states --- states with
acceptable d\_1 transitions but extreme d\_5, d\_6, or d\_8 costs ---
become traversable. {[}{]}

\textbf{The Flash Crash Mechanism}. During the Flash Crash, a large sell
order from a mutual fund was executed by an algorithm optimizing on d\_1
alone. The algorithm\textquotesingle s geodesic on the 1D projection was
optimal: sell at the best available price. But the 1D geodesic traversed
states where d\_5 (trust: other market participants lost confidence in
price integrity), d\_6 (systemic impact: cascading liquidations), and
d\_8 (legitimacy: exchange circuit breakers were triggered) carried
enormous cost. The 1D algorithm could not perceive these costs. The
full-manifold geodesic would have incorporated the systemic cost of
rapid execution and chosen a slower path, sacrificing some d\_1
performance to avoid d\_5-d\_6-d\_8 catastrophe.

\textbf{Remark (IEX Speed Bump)}. The Investors Exchange (IEX)
introduced a 350-microsecond speed bump --- a deliberate delay on
incoming orders --- to prevent high-frequency traders from exploiting
microsecond d\_9 advantages. In the geometric framework, the speed bump
is a boundary penalty: it raises the edge cost for paths that exploit
extreme d\_9 asymmetry, forcing traders onto paths with lower d\_9
exploitation and higher d\_3 (fairness) preservation.

\subsection{23.5 High-Frequency Trading and the Fairness
Problem}\label{high-frequency-trading-and-the-fairness-problem}

\textbf{High-frequency trading (HFT) firms exploit speed advantages} ---
colocation, microwave towers, FPGA hardware --- to access market
information microseconds before other participants. In the geometric
framework, HFT is a systematic exploitation of d\_9 (epistemic
advantage) at the expense of d\_3 (market fairness).

\textbf{The HFT Geodesic}. The HFT firm\textquotesingle s financial
decision complex has d\_9 edge weights that are systematically lower
than those of slower participants. The HFT firm\textquotesingle s
geodesic reaches profitable states (high d\_1) at lower cost, not
because of better fundamental analysis but because of lower d\_9
friction (faster information access). The value transfer from slower
participants to HFT firms on d\_1 is the monetary manifestation of the
d\_9 advantage.

\textbf{Theorem 23.3 (HFT as Fairness-Epistemic Trade-off)}. In a market
with HFT and non-HFT participants, the HFT firm\textquotesingle s
optimal path on the d\_1-d\_9 projection differs from its optimal path
on the full manifold. Specifically, the d\_1-d\_9 geodesic exploits
speed advantages maximally, while the full-manifold geodesic
incorporates d\_3 (fairness) costs that constrain exploitation.

\textbf{Proof}. On the d\_1-d\_9 projection, the HFT firm minimizes a
two-dimensional Mahalanobis distance without fairness constraints.
Adding d\_3 to the optimization introduces cross-terms
Sigma\_\{3,9\}\^{}\{-1\} that penalize high d\_9 advantage when d\_3 is
low (unfair conditions). The minimum-cost path shifts toward lower d\_9
exploitation and higher d\_3 preservation. The difference between the
two geodesics is the fairness cost of HFT, measurable in manifold
distance. {[}{]}

\textbf{Remark}. The regulatory debate over HFT (Budish, Cramton, and
Shim 2015 vs. Menkveld 2013) is a debate over which manifold dimensions
should count. Defenders argue HFT improves d\_1 efficiency (tighter
spreads, better price discovery). Critics argue it degrades d\_3 (unfair
access) and d\_5 (market trust). The geometric framework clarifies the
trade-off: both sides are correct on different dimensions. The policy
question is the relative weighting in Sigma\_F --- an empirical
question, not a philosophical one.

\subsection{23.6 Geometric Risk
Management}\label{geometric-risk-management}

Classical risk management measures risk as variance, Value-at-Risk
(VaR), or Conditional Value-at-Risk (CVaR) --- all scalar projections of
a multi-dimensional risk landscape. The geometric framework identifies
risk with curvature.

\textbf{Theorem 23.4 (Financial Irrecoverability)}. Any scalar risk
measure R: R\^{}9 -\textgreater{} R is non-injective. Multiple
financially distinct risk states map to the same scalar risk score. The
information destroyed by scalar risk measures is mathematically
irrecoverable.

\textbf{Proof}. Identical to the proof of Theorem 20.2 (Scalar
Irrecoverability), applied to the financial domain. By
Brouwer\textquotesingle s invariance of dimension, no continuous
injection from R\^{}9 to R\^{}1 exists. {[}{]}

\textbf{Definition 23.3 (Risk as Curvature)}. The risk at a point v on
the financial decision complex is the sectional curvature K(v) of the
manifold metric at v. High curvature means that small perturbations in
the current state produce large changes in the optimal path --- the
financial analog of "sensitivity to initial conditions." Low curvature
means the optimal path is robust to perturbation.

This definition unifies several classical risk concepts:

Volatility is the magnitude of curvature fluctuation along the d\_1
axis: high volatility corresponds to rapidly varying curvature in the
monetary dimension.

Tail risk is curvature concentration at the boundary of the manifold:
extreme events correspond to regions of very high curvature where the
optimal path changes catastrophically with small perturbations.

Correlation risk is off-diagonal curvature: the rate at which optimal
paths on one dimension change when another dimension is perturbed. The
2008 financial crisis was a correlation-risk event: the assumption of
low Sigma\_\{i,j\} (independent mortgage defaults) was catastrophically
wrong when the true covariance spiked.

Liquidity risk is curvature of the market-access dimension (d\_4): in
liquid markets, d\_4 curvature is low (small perturbations in desired
trade size produce small changes in execution cost); in illiquid
markets, d\_4 curvature is high (small increases in trade size produce
large increases in execution cost).

\subsection{23.7 Financial Conservation
Laws}\label{financial-conservation-laws}

\textbf{Theorem 23.5 (Financial Conservation)}. In a closed bilateral
financial transaction between parties A and B, for each transferable
dimension k in \{d\_1 (monetary value), d\_2 (contractual obligations),
d\_4 (trading autonomy)\}:

Da\_k(A) + Da\_k(B) = 0

What A gains on dimension k, B loses, and vice versa. This is the
specialization of the Attribute Conservation Theorem (Chapter 20,
Theorem 20.6) and the Invariance of Harm Accounting (Chapter 12, Theorem
12.1) to financial markets.

\textbf{Proof}. In a closed bilateral transaction, the total endowment
on each transferable dimension is fixed. Money (d\_1) is neither created
nor destroyed by trading; it is transferred. Contractual obligations
(d\_2) are paired: every long position has a short counterparty. Trading
autonomy (d\_4) is zero-sum in the sense that a restriction on one
party\textquotesingle s freedom to exit creates a corresponding
obligation on the other. By closure, Da\_k(A) + Da\_k(B) = 0. {[}{]}

\textbf{Remark (Non-Conservation of Evaluative Dimensions)}. Evaluative
dimensions --- d\_3 (market fairness), d\_5 (counterparty trust), d\_7
(reputation), d\_8 (regulatory legitimacy) --- are not conserved. Both
parties to a transaction may simultaneously gain or lose trust,
perceived fairness, or institutional legitimacy. A fair and transparent
trade creates mutual value on d\_3, d\_5, and d\_8 that did not exist
before the transaction. A manipulative trade destroys trust on both
sides.

\textbf{Remark (Market Crashes as Conservation Violation)}. A market
crash appears to violate d\_1 conservation: aggregate market value
decreases with no corresponding increase. The resolution is that market
value is not the same as money. Market value is a mark-to-market
evaluation --- a conditional expectation of future cash flows. A crash
destroys evaluated value (d\_1 under current beliefs) while conserving
actual money (d\_1 as medium of exchange). The crash simultaneously
destroys d\_5 (trust) and d\_8 (legitimacy) across all participants ---
a non-conserved evaluative dimension collapse.

\subsection{23.8 Geometric Behavioral
Finance}\label{geometric-behavioral-finance}

The major phenomena of behavioral finance are derived as geometric
properties of the financial decision complex:

\textbf{Bubbles as Heuristic Inadmissibility}. A financial bubble occurs
when market participants\textquotesingle{} heuristic h(n) systematically
overestimates the benefit (underestimates the cost) of continuing on the
current path. In A* terms, h(n) \textless{} 0 for nodes along the bubble
path --- the heuristic assigns negative cost (expected benefit) to
remaining steps, encouraging the agent to continue even when g(n)
(accumulated cost) is already high.

\textbf{Theorem 23.6 (Bubble Dynamics)}. A financial bubble is an A*
search with inadmissible heuristic. Specifically, the
market\textquotesingle s collective heuristic h\_market(n) satisfies
h\_market(n) \textless{} h*(n) with h\_market(n) \textless{} 0 for
bubble-path nodes, where h*(n) \textgreater{} 0 is the true remaining
cost. The bubble bursts when g(n) accumulated evidence forces revision
of h\_market toward h*, causing the evaluated cost f(n) = g(n) +
h\_market(n) to exceed the cost of alternative paths.

\textbf{Proof}. During a bubble, market participants estimate future
returns (negative costs) as increasingly certain, producing h\_market(n)
\textless{} 0. The true cost h*(n) accounts for the probability of mean
reversion, regulatory intervention, and fundamental revaluation, all of
which are positive. As g(n) accumulates (higher prices, lower yields,
stretched fundamentals), the gap between h\_market and h* becomes
unsustainable. A triggering event (margin call, earnings miss,
regulatory action) forces h\_market toward h*, making the bubble path
dominated by the exit path. This is the burst. {[}{]}

\textbf{Momentum as Manifold Inertia}. The momentum anomaly --- the
tendency of assets with recent positive returns to continue
outperforming --- is manifold inertia. On a curved manifold, the
geodesic has a "direction" that persists over short distances
(autoparallel transport). In the financial decision complex, an asset
whose attribute vector has been moving in a consistent direction on d\_1
has lower edge weights for continued motion in that direction (the
Mahalanobis distance for small, same-direction changes is small).
Reversal requires changing direction on the manifold, which has higher
curvature cost.

\textbf{Panic as Catastrophic Path Switching}. A market panic is a
sudden, discontinuous switch from one geodesic to another. In the
geometric framework, this occurs when the manifold curvature at the
current position exceeds a critical threshold: the optimal path
bifurcates, and the agent switches from the "normal" geodesic to the
"exit" geodesic. The discontinuity explains why panics are sudden ---
the transition is not gradual but catastrophic (in the mathematical
sense of Thom\textquotesingle s catastrophe theory).

\subsection{23.9 Option Pricing as Scalar
Projection}\label{option-pricing-as-scalar-projection}

\textbf{Proposition 23.1 (Black-Scholes as Scalar Projection)}. The
Black-Scholes model prices options using only d\_1 (return distribution)
and its variance (constant volatility). The model is a scalar projection
of the financial decision complex onto d\_1 alone.

The implied volatility surface is the market\textquotesingle s
correction for this projection. The volatility smile (out-of-the-money
puts and calls have higher implied volatility than at-the-money options)
reflects the market\textquotesingle s incorporation of dimensions
invisible to Black-Scholes:

\textbf{d\_3 (fairness): Jump risk} --- the possibility of sudden large
moves that disproportionately affect less-informed participants --- is
priced into out-of-the-money options, widening the smile.

\textbf{d\_5 (trust)}: Counterparty risk premium, especially post-2008,
increases implied volatility for long-dated options where counterparty
default is more probable.

\textbf{d\_6 (systemic impact)}: Correlation risk --- the risk that all
assets move together in a crisis --- is priced into index options
relative to single-stock options, producing the "skew."

\textbf{d\_8 (regulatory)}: Regulatory intervention risk (position
limits, trading halts, short-selling bans) is priced into options on
regulated assets, producing term-structure effects.

The implied volatility surface is not a "failure" of Black-Scholes. It
is the market\textquotesingle s attempt to recover, within a scalar
framework, the information that scalar pricing destroys. The geometric
framework predicts that the surface shape should correlate with the
active non-d\_1 dimensions --- a testable prediction.

\subsection{23.10 The 2008 Financial Crisis as Manifold
Failure}\label{the-2008-financial-crisis-as-manifold-failure}

The 2008 global financial crisis provides a comprehensive case study in
manifold failure across multiple dimensions simultaneously.

Dimensional collapse in mortgage markets. The originate-to-distribute
model separated mortgage origination (d\_2: contractual obligations)
from risk-bearing (d\_1: return). Originators optimized on d\_1 alone
(origination fees) with zero d\_2 cost (no skin in the game). This
dimensional collapse enabled paths through the financial decision
complex that full-manifold originators would never have traversed: loans
to borrowers who could not repay, packaged into securities whose d\_9
(epistemic status: actual risk profile) was systematically
misrepresented.

Rating agency failure as gauge variance. Rating agencies assigned AAA
ratings to mortgage-backed securities whose actual risk was far higher.
In the geometric framework, this is a gauge-invariance violation: the
same security, described differently (as a pool of individual mortgages
vs. as a structured tranche), received different risk evaluations. The
rating depended on the description (the gauge), not on the underlying
attribute vector --- a Bond Invariance Principle violation.

Contagion as curvature cascade. When Lehman Brothers collapsed, the d\_5
(trust) dimension across the entire financial system underwent
catastrophic curvature increase. Counterparty trust, previously a slowly
varying manifold feature, spiked to infinity (no participant trusted any
other\textquotesingle s solvency). This curvature cascade propagated
through the d\_5-d\_6 covariance channel: lost trust (d\_5) amplified
systemic impact (d\_6), which further degraded trust. The positive
feedback loop is a curvature instability on the financial manifold.

Regulatory response as manifold repair. The Dodd-Frank Act (2010) can be
understood as manifold repair: raising d\_2 boundary penalties (higher
capital requirements, Volcker Rule), increasing d\_8 edge weights
(compliance costs), and introducing d\_6 boundary penalties (systemic
risk surcharges for systemically important financial institutions). The
goal was to prevent future dimensional collapse by ensuring that d\_2,
d\_5, d\_6, and d\_8 costs could not be zeroed out.

\subsection{23.11 Falsifiable
Predictions}\label{falsifiable-predictions-3}

The framework generates six predictions that distinguish it from
classical and behavioral finance:

\textbf{Prediction 1 (Volatility Surface Correlation)}: The shape of the
implied volatility surface should correlate with measurable non-d\_1
dimensions. Specifically, the smile width should correlate with d\_9
disparity (information asymmetry), and the skew should correlate with
d\_6 (systemic risk) proxies. Falsified if: the surface shape is
independent of non-d\_1 measures.

\textbf{Prediction 2 (HFT Fairness Index)}: A Financial Bond Index
computed from HFT activity should predict market participant attrition
(departure of slower participants from the market). Falsified if:
participant attrition is uncorrelated with the computed BI.

\textbf{Prediction 3 (Flash Crash Prediction)}: Markets with higher d\_1
concentration (more algorithmic, less human oversight on non-d\_1
dimensions) should experience more frequent and severe flash crashes.
Falsified if: flash crash frequency is independent of d\_1
concentration.

\textbf{Prediction 4 (Curvature-Risk Correlation)}: Historical curvature
estimates on the financial decision complex should predict future
realized risk (measured by subsequent volatility or tail-event
frequency) better than scalar risk measures (VaR, beta). Falsified if:
curvature adds no predictive power beyond scalar risk measures.

\textbf{Prediction 5 (Crisis Covariance Spike)}: During financial
crises, the estimated covariance matrix Sigma\_F should show
characteristic patterns: d\_5-d\_6 covariance increases (trust-systemic
coupling), d\_1-d\_5 covariance increases (return-trust coupling), and
d\_8 variance increases (regulatory uncertainty). Falsified if: crisis
covariance patterns are structureless.

\textbf{Prediction 6 (Manifold Dimensionality)}: Factor analysis of
financial market behavioral data (incorporating non-price dimensions
such as regulatory filings, counterparty ratings, trading speed, and
market fairness indicators) should recover approximately nine
independent dimensions. Falsified if: the factor structure is
consistently lower-rank (fewer than seven) or higher-rank (more than
eleven).

\subsection{23.12 Connection to the
Framework}\label{connection-to-the-framework-4}

The Geometric Finance program extends the parent framework in five
directions:

1. Chapter 20 applied the framework to individual and bilateral economic
decisions. This chapter extends it to market microstructure --- the
domain where many agents interact simultaneously through an
institutional mechanism (the exchange), introducing market-specific
geometry (bid-ask spread as curvature, liquidity as d\_4 smoothness).

2. Chapter 15 established scalar irrecoverability. This chapter applies
it to financial risk measurement (VaR, CAPM beta) and option pricing
(Black-Scholes), showing that the implied volatility surface is the
market\textquotesingle s attempt to recover scalar-projected
information.

3. Chapter 12 established conservation laws. This chapter derives
financial conservation (money and contractual obligations are conserved
in bilateral transactions) and identifies the non-conservation of
evaluative dimensions (trust, legitimacy, reputation) as the mechanism
underlying financial contagion.

4. The Flash Crash analysis demonstrates a new phenomenon ---
dimensional collapse --- where algorithmic agents\textquotesingle{}
restriction to d\_1 enables traversal of catastrophic paths invisible to
full-manifold agents. This is the financial analog of the ethical
concern about AI systems that optimize on a scalar objective without
moral constraints.

5. The 2008 financial crisis analysis shows how gauge-invariance
violations (rating agency failures), dimensional collapse
(originate-to-distribute), and curvature instability (trust cascades)
interact to produce systemic failure --- a multi-dimensional crisis that
scalar analysis cannot diagnose.

\subsection{23.13 Summary}\label{summary-14}

This chapter has shown that the geometric ethics framework, when applied
to financial markets, yields:

1. A formal construction of the financial decision complex F as a
domain-specific instantiation of the moral manifold, with market
microstructure mapped to manifold geometry.

2. Information asymmetry as metric distortion: Kyle\textquotesingle s
informed-trading model reinterpreted as asymmetric Mahalanobis distance
on the d\_9 component, with the bid-ask spread as manifold curvature.

3. The Flash Crash as dimensional collapse: algorithmic restriction to
d\_1 enables catastrophic paths invisible to full-manifold agents.

4. Risk as curvature: a unified geometric framework that subsumes
volatility, tail risk, correlation risk, and liquidity risk as different
aspects of manifold curvature.

5. Financial conservation: transferable dimensions (money, contracts)
are conserved; evaluative dimensions (trust, legitimacy) are not ---
explaining both market stability and contagion.

6. Behavioral finance phenomena (bubbles, momentum, panics) as geometric
properties: bubbles as inadmissible heuristics, momentum as manifold
inertia, panics as catastrophic path switching.

7. The implied volatility surface as shadow of higher-dimensional
pricing: the market\textquotesingle s correction for scalar projection.

\section{Chapter 24: Geometric Theology --- Religious Reasoning on the
Moral
Manifold}\label{chapter-24-geometric-theology-religious-reasoning-on-the-moral-manifold}

Epistemic status of this chapter. Except where a claim is explicitly
marked as a formal result or a conditional theorem, the material here is
a modeling interpretation, domain analogy, or speculative extension - a
demonstration of how the geometric vocabulary applies to theology, not a
set of established results at the evidential standard of the field.
Interpretations such as the Fall as an epistemic impossibility result
and the Euthyphro dilemma as gauge ambiguity are domain analogies, not
settled theological or philosophical findings.

\begin{quote}
\textbf{\textsc{RUNNING EXAMPLE --- Priya's Model}}

\emph{Mrs. Voss's pastor, Reverend Carla Whitfield, tells Priya's
colleague: `You are playing God with a spreadsheet.' Priya initially
dismisses this, but the geometric framework reframes it: what Reverend
Whitfield identifies is a sacred-value boundary---a stratum wall that
certain communities hold to be inviolable. The decision of who receives
access to a potentially life-saving trial touches what many consider
sacred ground. No scalar score can represent this; it is a topological
feature of the moral manifold, not a parametric one.}
\end{quote}

This chapter applies the Geometric Ethics framework to religious and
theological reasoning. Unlike the preceding domain chapters ---
economics, clinical ethics, law, finance --- which apply the framework
to domains that already possess formal traditions, this chapter applies
it to a domain that has historically resisted formalization and often
claimed exemption from it. The central finding is that the framework
illuminates theological questions with surprising precision: the Fall
narrative is an epistemic impossibility theorem, the Euthyphro dilemma
is a question about gauge invariance, theodicy is dimensional
projection, and the cross-cultural stability of the moral manifold has
direct implications for natural theology. The chapter engages multiple
world religions and treats the empirical data from the
framework\textquotesingle s cross-lingual validation corpus --- which
included substantial religious and philosophical texts across eleven
languages --- as evidence bearing on theological claims.

\subsection{24.1 Religious Moral Reasoning on the
Manifold}\label{religious-moral-reasoning-on-the-manifold}

Every major world religion provides its adherents with a moral
framework: a set of principles, prohibitions, virtues, and practices
that guide behavior. In the geometric framework, religious moral
instruction is heuristic calibration --- the systematic training of the
agent\textquotesingle s moral heuristic h(n) through boundary penalties,
dimensional weightings, and sacred-value designations.

The claim is not that religion reduces to heuristic calibration. The
claim is that the moral content of religious instruction --- the part
that tells adherents what to do and why --- operates on the same
nine-dimensional manifold as secular moral reasoning, and is subject to
the same mathematical analysis. A religious prohibition ("thou shalt not
steal") sets a boundary penalty on the same d\_2 (rights and
obligations) and d\_3 (fairness) dimensions as a secular legal
prohibition against theft. A religious virtue ("compassion for the
suffering") weights d\_1 (consequences for others) and d\_6 (social
impact) in the same dimensional space as a secular ethical commitment to
harm reduction.

This is empirically supported by the framework\textquotesingle s
cross-lingual validation corpus. The 20,030 texts analyzed in Bond
(2026) included religious and philosophical passages from multiple
traditions --- Buddhist sutras, Islamic ethical texts, Hindu
philosophical texts, Christian moral theology, Confucian analects ---
and the same nine-dimensional structure was recovered across all of
them. The independently replicated study (Thiele, 2026) confirmed this
finding across six additional languages. The manifold is not a Western
secular construct projected onto other traditions; it is a structure
that emerges from moral reasoning as such, regardless of its theological
or secular framing.

\textbf{Proposition 24.1 (Cross-Religious Manifold Invariance)}. The
nine-dimensional moral manifold M is invariant across world religions.
Religions differ in their metric tensors (the relative weighting of
dimensions), their boundary penalties (which actions carry sacred-value
boundaries), and their goal regions (the states they designate as
righteous, enlightened, or holy), but all operate on the same underlying
dimensional structure.

\textbf{Evidence}: The cross-lingual validation corpus recovered the
same nine dimensions from texts spanning Hindu, Buddhist, Christian,
Islamic, Confucian, and secular ethical traditions. Factor analysis
showed no religion-specific dimensions; rather, traditions vary in
dimensional weighting. Christianity weights d\_7 (virtue/identity) and
d\_2 (obligations) heavily. Buddhism weights d\_1
(consequences/suffering) and d\_4 (autonomy from attachment). Islam
weights d\_2 (obligations) and d\_8 (institutional/divine legitimacy)
heavily. Confucianism weights d\_6 (social relationships) and d\_5
(trust/loyalty). These are metric differences, not dimensional
differences.

\subsection{24.2 The Obligation Vector and Natural
Theology}\label{the-obligation-vector-and-natural-theology}

A natural theological claim, common across Abrahamic traditions, is that
God "wrote the law on our hearts" (Romans 2:15, cf. Jeremiah 31:33). In
the geometric framework, this claim has a precise interpretation: the
nine-dimensional moral manifold is an innate cognitive structure, not a
cultural construction, and its presence in all humans --- regardless of
religious instruction --- constitutes evidence for a divine moral
architect.

The framework\textquotesingle s empirical findings are genuinely
consistent with this claim: the manifold is cross-culturally stable,
cross-linguistically invariant, and recoverable from the moral reasoning
of every human culture studied. The obligation vector --- the
nine-dimensional attribute vector that encodes a moral state --- does
appear to be a universal feature of human moral cognition, not a
culturally contingent artifact.

However, the framework is equally consistent with a naturalistic
explanation: the manifold reflects the structure of social cooperation
problems that all human communities face (resource allocation, trust
maintenance, norm enforcement, identity preservation), and its
universality follows from the universality of these problems, not from
divine inscription.

\textbf{Proposition 24.2 (Underdetermination of Origin)}. The
cross-cultural universality of the moral manifold is consistent with
both theistic (divine inscription) and naturalistic (evolutionary/social
convergence) origins. The geometric framework cannot adjudicate between
these explanations because both predict the same observable: a stable,
cross-cultural, nine-dimensional structure. The framework identifies
what the moral structure is and how it operates; it is silent on why it
exists.

\textbf{Remark}. This underdetermination is itself a significant result.
Many theological arguments depend on the claim that moral universality
implies a divine source. The framework shows that this inference is
invalid: the universality of the manifold is established empirically,
but its origin remains underdetermined by the data. The argument from
moral universality to divine authorship requires an additional premise
--- that only divine authorship could produce universal moral structure
--- which the framework cannot confirm or deny.

\subsection{24.3 The Epistemic Paradox of the
Fall}\label{the-epistemic-paradox-of-the-fall}

The Genesis Fall narrative (Genesis 2:17-3:24) is the foundational
narrative of Abrahamic moral theology: the origin of human moral
knowledge, culpability, and the fracturing of the divine-human
relationship. On its own terms --- accepting the narrative as
authoritative arguendo --- the story contains a mathematical
impossibility when analyzed on the moral manifold.

\textbf{Theorem 24.1 (Epistemic Prerequisite Paradox)}. Let C be a
decision complex with a boundary penalty beta\_k whose evaluation
requires nonzero capacity on dimension d\_k. If agent
A\textquotesingle s current state has d\_k(A) = 0 (the agent lacks the
capacity to evaluate dimension k), then beta\_k is uncomputable by A.
The boundary is invisible to A\textquotesingle s pathfinding algorithm.
If beta\_k is the only obstruction preventing A from traversing the
edge, then A\textquotesingle s optimal path traverses the boundary.

\textbf{Proof}. The agent\textquotesingle s evaluation function is f(n)
= g(n) + h(n). The heuristic h(n) estimates remaining cost, including
boundary penalties. If computing beta\_k requires d\_k \textgreater{} 0
and d\_k(A) = 0, then A cannot incorporate beta\_k into h(n).
A\textquotesingle s effective edge weight for the boundary-crossing edge
is w\_eff(e) = w(e) - beta\_k (the weight without the boundary penalty).
If w\_eff(e) makes the edge optimal among available alternatives, A
traverses it. {[}{]}

\textbf{Application to the Fall}. Pre-Fall Adam and Eve exist in a state
where d\_2 (knowledge of good and evil --- the capacity to distinguish
right from wrong) is zero. God\textquotesingle s prohibition "do not eat
of the tree of knowledge" sets a boundary penalty beta\_disobedience on
the edge from innocence to knowledge. But evaluating beta\_disobedience
--- understanding that disobedience is wrong --- requires d\_2
\textgreater{} 0, which is precisely the capacity the fruit provides.

The paradox is exact: the boundary penalty requires the capacity that
crossing the boundary provides. Pre-Fall Adam cannot compute
beta\_disobedience because computing it requires the knowledge of good
and evil that he does not yet possess. His heuristic h(n) cannot
incorporate the penalty. His optimal path --- computed on his manifold
with d\_2 = 0 --- may well traverse the forbidden boundary, not out of
defiance but out of mathematical incapacity.

\textbf{Remark (The Serpent\textquotesingle s Veracity)}. The narrative
compounds the paradox. The serpent\textquotesingle s three claims ---
"you will not surely die" (Genesis 3:4), "your eyes will be opened"
(3:5), and "you will be like God, knowing good and evil" (3:5) --- are
all confirmed as factually accurate by the subsequent narrative (Adam
lives 930 years; their eyes are opened in 3:7; God confirms "the man has
become like one of us" in 3:22). Pre-Fall Adam, lacking the capacity to
detect deception (which requires d\_2 \textgreater{} 0), has no
computational basis for disbelieving the serpent. The serpent is
indistinguishable from a divinely sanctioned agent on
Adam\textquotesingle s d\_2 = 0 manifold.

\textbf{Theological Implication}. If moral culpability requires that the
agent can compute the boundary penalty for the prohibited action, and
the boundary penalty is uncomputable without the capacity the action
provides, then the agent is not culpable. The Fall is not a moral
failure by Adam and Eve; it is a manifold design problem. The system
requires agents to honor a boundary they cannot perceive. Standard
apologetic defenses (spiritual death reading, "day as a thousand years,"
Irenaean soul-making) do not resolve the mathematical structure of the
paradox; they reinterpret the narrative\textquotesingle s terms without
addressing the computational impossibility.

\subsection{24.4 The Euthyphro Dilemma as Gauge
Ambiguity}\label{the-euthyphro-dilemma-as-gauge-ambiguity}

The Euthyphro dilemma (Plato, Euthyphro 10a) asks: is an action good
because God commands it (divine command theory), or does God command it
because it is good (moral realism)? In the geometric framework, this is
a question about whether the metric tensor on the moral manifold is
gauge-dependent or gauge-invariant.

\textbf{Horn 1}: Gauge-dependent metric (divine command theory). If the
metric tensor g\_ij is determined by divine fiat, then gauge
transformations (God changing commands) change the geodesic. Different
gods, or the same god issuing different commands, would produce
different moral truths. Morality is radically contingent on the
particular divine will that happens to obtain. This horn is
theologically problematic: it implies that God could make torture good
by commanding it, and there would be no standard external to
God\textquotesingle s will by which to judge the command.

\textbf{Horn 2}: Gauge-invariant metric (moral realism). If the metric
tensor is a property of the manifold itself, invariant under changes of
divine command, then the geodesic is fixed regardless of what any being
commands. God\textquotesingle s commands, to be good, must track the
pre-existing geodesic. God is not the source of morality but a (perhaps
perfect) reporter of a structure that exists independently.

\textbf{Theorem 24.2 (Euthyphro Resolution via Gauge Invariance)}. The
empirical data from the cross-lingual validation corpus (Bond, 2026;
Thiele, 2026) support the gauge-invariant horn: the moral
manifold\textquotesingle s nine-dimensional structure is stable across
cultures with radically different theological commitments (including
non-theistic traditions). The metric varies across cultures (different
dimensional weightings), but the dimensional structure is invariant ---
it does not depend on which deity (if any) a culture worships.

\textbf{Proof}. If the metric were gauge-dependent (determined by the
commands of a particular deity), cultures with different deities should
exhibit different dimensional structures --- different numbers of
dimensions, or dimensions absent in some traditions and present in
others. The empirical finding is that all cultures share the same
nine-dimensional structure, differing only in metric (weighting). This
is precisely the signature of gauge invariance: the manifold is
invariant under changes of "divine gauge" (which god, if any, is
worshipped), while the metric (weighting) is gauge-dependent. {[}{]}

\textbf{Theological Consequence}. The framework resolves Euthyphro in
favor of moral realism --- the metric tensor has a gauge-invariant
dimensional structure --- without requiring that moral realism entail
atheism. A theist can consistently hold that God created the manifold
with its gauge-invariant structure, just as a physicist can hold that
God created spacetime with its gauge-invariant geometric structure. The
framework is neutral on whether the manifold has a divine origin; it
establishes that the manifold\textquotesingle s structure is not
dependent on any particular divine command.

\subsection{24.5 Theodicy as Dimensional
Projection}\label{theodicy-as-dimensional-projection}

\textbf{The Problem of Evil (Epicurus, Hume, Mackie) asks}: if God is
omnipotent, omniscient, and omnibenevolent, why does suffering exist?
The geometric framework does not solve the Problem of Evil, but it
clarifies its mathematical structure.

\textbf{Proposition 24.3 (Theodicy as Dimensional Projection)}. The
classical formulation of the Problem of Evil assumes that
"omnibenevolent" means "minimizes suffering" --- i.e., God optimizes on
d\_1 (consequences/suffering) alone. If God optimizes on the full
nine-dimensional manifold, the divine geodesic may traverse states with
high d\_1 cost (suffering) because those states are necessary for
reaching states with better values on other dimensions (d\_4: autonomy,
d\_7: identity/virtue, d\_2: moral development through exercising rights
and obligations).

\textbf{This is the geometric formalization of the traditional
soul-making theodicy (Irenaeus, Hick)}: suffering exists because the
manifold geodesic from morally immature to morally mature states
traverses high-d\_1 regions. The geometric framework makes this precise:
the geodesic on the full manifold is not the same as the geodesic on the
d\_1 projection. States that appear suboptimal when viewed on d\_1 alone
may be geodesic-optimal on the full manifold.

\textbf{Remark (Limits of This Analysis)}. This analysis clarifies but
does not resolve the Problem of Evil. Two objections remain. First, if
God is omnipotent, God could create a manifold whose geodesic does not
require high-d\_1 states --- a manifold with different curvature and
topology. The framework takes the manifold as given and computes
geodesics on it; it cannot address why the manifold has the structure it
does. Second, certain instances of suffering (the torture of innocents,
natural disasters killing children) resist dimensional-projection
explanation because no plausible d\_2-d\_9 benefit compensates the d\_1
cost. The framework is honest about this limit: it dissolves the scalar
formulation of the problem without resolving the full-dimensional
version.

\subsection{24.6 Sacred-Value Boundaries and Religious
Heuristics}\label{sacred-value-boundaries-and-religious-heuristics}

Religious traditions are distinguished from secular ethical systems by
the prevalence of sacred-value boundaries --- moral boundaries with
beta\_k = infinity that admit no trade-offs, regardless of the benefits
on other dimensions.

In secular ethics, most boundaries are finite: there is some benefit
large enough to justify crossing them ("it is permissible to lie to save
a life"). In religious ethics, certain boundaries are explicitly
infinite: no benefit justifies crossing them. The Ten Commandments, the
Five Precepts of Buddhism, the hudud in Islamic law, and the yamas in
Hindu ethics all contain sacred-value prohibitions that are structurally
identical --- they set beta\_k = infinity on specific edges of the
decision complex.

\textbf{Proposition 24.4 (Sacred-Value Convergence)}. Despite radical
theological differences, the sacred-value boundaries of major world
religions converge on a common core: prohibitions against unjustified
killing (d\_1 + d\_2), theft (d\_2 + d\_3), deception (d\_9 + d\_5), and
sexual violation (d\_4 + d\_7). The convergence is predicted by the
framework: these boundaries protect the dimensions most critical for
social cooperation (consequences, rights, fairness, trust), and any
community that lacks these boundaries faces severe cooperation failure.

\textbf{Evidence}: The cross-lingual validation corpus showed that
sacred-value texts from different religious traditions activated the
same dimensional combinations. The Hindu ahimsa (non-violence), the
Buddhist first precept (do not kill), the Christian "thou shalt not
kill," and the Islamic prohibition of unjustified homicide all activate
d\_1 + d\_2 with infinite boundary weight. The convergence is
structural, not coincidental --- it reflects the geometry of the moral
manifold, not the content of any particular revelation.

\textbf{Religious Moral Training as Heuristic Calibration}. Religious
formation --- Sunday school, madrasa, monastic training, catechism ---
calibrates the adherent\textquotesingle s heuristic h(n) by
systematically setting boundary penalties, dimensional weights, and goal
regions. The framework predicts that effective religious moral training
should produce heuristics that are at least epsilon-admissible on the
manifold: the trained adherent finds paths that are near-optimal
relative to the community\textquotesingle s metric tensor. Religious
communities that survive over centuries are predicted to have developed
heuristics that are well-calibrated relative to their metric --- a form
of cultural selection for heuristic admissibility.

\subsection{24.7 Deductive Natural Theology and Gauge
Reification}\label{deductive-natural-theology-and-gauge-reification}

The four principal deductive arguments for God\textquotesingle s
existence --- the cosmological argument (Kalam), the contingency
argument, the moral argument, and the ontological argument --- share a
common structural vulnerability when analyzed through the geometric
framework.

Each argument proceeds through four steps: (1) Domestication --- start
with domain-specific intuitions (causal reasoning, explanation-seeking,
moral experience, conceptual coherence). (2) Formalization --- codify
intuitions into formal tools (classical logic, S5 modal semantics, the
Principle of Sufficient Reason). (3) Reification --- treat the formal
tools as transparent mirrors of reality, not as models with domain
limitations. (4) Extraction --- read off God\textquotesingle s existence
from the formal tools\textquotesingle{} internal structure.

In the geometric framework, step 3 is gauge reification: treating
coordinate-dependent properties of a formal system as if they were
coordinate-independent facts about reality. The theorem of S5 modal
logic that "if possibly necessarily P, then necessarily P" is a property
of S5 frames (equivalence-relation Kripke frames), not a property of
metaphysical modality as such. The Principle of Sufficient Reason is a
methodological maxim calibrated for empirical domains, not a law of
being. The moral argument equivocates between stance-independent values
(which would require a ground external to human practice) and robustly
intersubjective values (which are stable under scrutiny but grounded in
social structure).

\textbf{Theorem 24.3 (Incompleteness Constraints on Theological
Formalization)}. By Godel\textquotesingle s first incompleteness
theorem, any formal system sufficiently powerful to represent basic
arithmetic cannot prove all truths about its own domain. Applied to
theological formalization: any formal system rich enough to express the
concept of a necessary being cannot, from within, certify that the
system\textquotesingle s axioms faithfully represent metaphysical
reality. The gap between intra-systematic validity (derivability within
the formal system) and extra-systematic truth (correspondence to
reality) cannot be closed by formal means alone.

\textbf{Proof}. Let F be a formal system extending first-order
arithmetic with modal operators (as in Godel\textquotesingle s own
ontological proof). By Godel\textquotesingle s first incompleteness
theorem, there exist sentences of F that are true but unprovable in F.
In particular, the consistency of F (Con(F)) is not provable in F
(Godel\textquotesingle s second incompleteness theorem). Any theorem of
F that purports to establish a metaphysical conclusion (e.g., "a
necessary being exists") depends on the consistency of F, which F cannot
self-certify. The conclusion is intra-systematically valid but its
extra-systematic truth depends on an unverifiable meta-theoretic
assumption. {[}{]}

\textbf{Remark (Quantum Non-Classicality)}. The ontological argument
assumes classical logic --- specifically, the distributive law a AND (b
OR c) = (a AND b) OR (a AND c). Birkhoff and von Neumann (1936) showed
that quantum mechanics operates on an orthomodular lattice where
distributivity fails. Bell\textquotesingle s theorem (1964) and the
Kochen-Specker theorem (1967) establish that no classical (bivalent,
context-independent) model can reproduce quantum predictions. If the
logical substrate of reality is non-classical, then arguments that
depend on classical logical structure --- including all four deductive
arguments for theism --- are domain-limited: they hold in the classical
regime but their extension to fundamental metaphysics is not warranted.

\textbf{Remark (Modal Collapse)}. Godel\textquotesingle s own
ontological proof, when formalized in S5, produces modal collapse: every
true proposition becomes necessarily true (Sobel, 1987; computationally
verified by Benzmuller and Woltzenlogel Paleo, 2014). In the geometric
framework, modal collapse is manifold degeneracy --- the modal space
collapses to a single point, destroying the distinction between
contingent and necessary truths. This is a reductio of the formal
system: the axioms that produce "God exists necessarily" also produce
"everything that happens, happens necessarily," which is absurd.

\subsection{24.8 Implications for Specific Religious
Traditions}\label{implications-for-specific-religious-traditions}

The framework\textquotesingle s findings bear on specific doctrinal
claims across world religions:

\textbf{Christianity}. The doctrine of Original Sin depends on the
culpability of Adam and Eve in the Fall. Theorem 24.1 shows that the
culpability condition fails: the boundary penalty for the prohibited
action requires a capacity the agents lack. The Pauline theology of
Romans 5:12-21, which derives universal human sinfulness from
Adam\textquotesingle s transgression, inherits this mathematical
problem. Additionally, the doctrine of divine moral perfection (God as
the standard of goodness) is a gauge-dependent claim: it presupposes
that the moral metric is fixed by God\textquotesingle s nature (Horn 1
of Euthyphro). The empirical gauge invariance of the manifold (Theorem
24.2) challenges this presupposition.

\textbf{Islam}. The Islamic emphasis on divine command
(God\textquotesingle s will as the source of moral obligation) aligns
with Horn 1 of Euthyphro --- gauge-dependent morality. The
framework\textquotesingle s finding that the manifold structure is
gauge-invariant poses a challenge: if moral dimensions are invariant
across cultures with different concepts of divine will, the moral
structure is not dependent on any particular divine command. However,
Islamic theological traditions (Ash\textquotesingle ari vs.
Mu\textquotesingle tazili) have debated this exact question for
centuries; the Mu\textquotesingle tazili position (moral truths are
accessible to reason independent of revelation) is consistent with gauge
invariance.

\textbf{Buddhism}. Buddhism is distinctive among major religions in its
compatibility with the geometric framework. The Four Noble Truths
describe the manifold from the d\_1 (suffering) perspective; the
Eightfold Path describes a specific geodesic on the manifold. Buddhist
ethics is explicitly consequentialist on d\_1 (reduce suffering) while
incorporating d\_4 (non-attachment as autonomy from craving) and d\_7
(identity transformation through practice). The absence of a creator
deity avoids the Euthyphro problem entirely. The
framework\textquotesingle s empirical finding that the manifold is not
deity-dependent is maximally consistent with Buddhist metaphysics.

\textbf{Hinduism}. The dharmic tradition recognizes multiple,
context-dependent moral frameworks (svadharma --- one\textquotesingle s
own duty, determined by caste, stage of life, and circumstance). In the
geometric framework, this is a variable metric tensor: the same manifold
with different dimensional weightings depending on the
agent\textquotesingle s social position. The Bhagavad
Gita\textquotesingle s central ethical problem ---
Arjuna\textquotesingle s conflict between d\_2 (duty as a warrior) and
d\_6 (social impact: killing relatives) --- is a pathfinding problem on
the moral manifold with conflicting dimensional weights.
Krishna\textquotesingle s resolution (act according to svadharma without
attachment to outcomes) is a prescription for heuristic selection:
choose the path determined by your metric tensor, not by the d\_1
projection.

\textbf{Judaism}. Rabbinic Judaism\textquotesingle s halakhic system ---
detailed legal-moral reasoning from Torah and Talmud --- is structurally
similar to the legal pathfinding of Chapter 22. Halakhic reasoning is A*
search on a religious-legal decision complex, with the Torah as the
precedent database, the Talmud as the commentary extending the graph,
and rabbinic rulings as subsequent vertices in the complex. The emphasis
on machloket (legitimate disagreement among authorities) corresponds to
multiple near-optimal paths on the manifold --- different geodesics with
similar total cost.

\subsection{24.9 Genesis 3:22 and the Moral Competence
Dilemma}\label{genesis-322-and-the-moral-competence-dilemma}

Genesis 3:22 records God\textquotesingle s statement: "Behold, the man
has become like one of us, knowing good and evil." This verse creates a
formal dilemma with direct implications for the relationship between
human moral competence and divine moral authority.

In the geometric framework, "knowing good and evil" means d\_2
\textgreater{} 0 --- the human agent now has nonzero capacity on the
dimension that encodes the distinction between right and wrong.
God\textquotesingle s assertion that humans have "become like one of us"
on this dimension implies that human moral evaluation and divine moral
evaluation operate on the same dimensional structure.

This creates a dilemma:

\textbf{Horn 1}: If humans genuinely possess divine-level moral
knowledge (d\_2\^{}\{human\} = d\_2\^{}\{divine\}), then humans are
epistemically equipped to evaluate divine conduct. The apologetic move
"Who are you to judge God?" fails: God has declared humans competent on
the very dimension required for such judgment.

\textbf{Horn 2}: If humans do not possess genuine moral knowledge
(d\_2\^{}\{human\} \textless\textless{} d\_2\^{}\{divine\}), then
God\textquotesingle s statement in Genesis 3:22 is false. If God can
make false statements, the entire theological apparatus that depends on
divine veracity --- including the trustworthiness of Scripture and
revelation --- is undermined.

\textbf{Remark}. This dilemma is structural, not rhetorical. It follows
from the framework\textquotesingle s formalization of moral knowledge as
a dimension of the manifold and from the text\textquotesingle s explicit
assertion about the state of that dimension. No apologetic defense that
preserves both human moral competence and divine exemption from moral
evaluation can avoid one horn or the other.

\subsection{24.10 Falsifiable
Predictions}\label{falsifiable-predictions-4}

The framework generates six predictions relevant to the intersection of
geometry and theology:

\textbf{Prediction 1 (Cross-Religious Dimensional Invariance)}: The
nine-dimensional manifold structure should be recoverable from the
sacred texts of every major world religion, with religions differing in
metric (dimensional weighting) but not in dimensional structure.
Falsified if: a major religious tradition\textquotesingle s moral
reasoning requires dimensions absent from the framework.

\textbf{Prediction 2 (Sacred-Value Convergence)}: The infinite-penalty
boundaries of major religions should converge on dimensions critical for
social cooperation (d\_1+d\_2 killing, d\_2+d\_3 theft, d\_5+d\_9
deception, d\_4+d\_7 sexual violation). Falsified if: sacred-value
boundaries are randomly distributed across dimensions with no
convergence.

\textbf{Prediction 3 (Heuristic Admissibility of Religious Training)}:
Adherents of religious traditions with long institutional histories
should have heuristic functions h(n) that are epsilon-admissible on the
manifold --- producing near-optimal paths relative to their
community\textquotesingle s metric. Falsified if: religiously trained
individuals show no better heuristic calibration than untrained
individuals.

\textbf{Prediction 4 (Gauge Invariance of Dimensional Structure)}: The
dimensional structure of the moral manifold should be invariant under
changes in theological commitment --- theists, atheists, polytheists,
and non-theists should all exhibit the same nine dimensions when making
moral judgments. Falsified if: the dimensional structure varies
systematically with theological commitment.

\textbf{Prediction 5 (Euthyphro Empirical Test)}: If divine command
theory is correct (gauge-dependent metric), cultures with radically
different theologies should exhibit different dimensional structures. If
moral realism is correct (gauge-invariant), they should not. The
cross-lingual validation data support gauge invariance. Falsified if:
new data from previously unstudied cultures show theology-dependent
dimensional structure.

\textbf{Prediction 6 (Fall Paradox Generalizability)}: The epistemic
prerequisite paradox should apply to any origin-of-moral-knowledge
narrative that (a) posits an initial state of moral innocence, (b)
prohibits the acquisition of moral knowledge, and (c) makes culpability
depend on moral understanding. Falsified if: a coherent narrative
satisfying (a)-(c) can be constructed that avoids the paradox.

\subsection{24.11 Connection to the
Framework}\label{connection-to-the-framework-5}

The Geometric Theology chapter connects to the parent framework in four
ways:

1. Chapter 5 defined the nine-dimensional moral manifold. This chapter
shows that the manifold is present in all major religious traditions,
supporting its status as a universal feature of moral cognition rather
than a product of any particular theological or secular tradition.

2. Chapter 11 established A* pathfinding as the model of moral
reasoning. This chapter shows that religious moral training is heuristic
calibration --- the systematic tuning of h(n) through boundary
penalties, dimensional weights, and sacred-value designations.

3. Chapter 12 established the Bond Invariance Principle. This chapter
applies it to the Euthyphro dilemma: the empirical gauge invariance of
the manifold resolves Euthyphro in favor of moral realism --- the
dimensional structure is coordinate-independent, regardless of
theological commitments.

4. The empirical validation corpus (Chapter 17) included religious texts
across eleven languages and multiple traditions. This chapter draws on
those data to show that the manifold\textquotesingle s cross-cultural
stability extends across religious boundaries, bearing directly on
natural theological arguments from moral universality.

\subsection{24.12 Summary}\label{summary-15}

This chapter has shown that the geometric ethics framework, when applied
to religious and theological reasoning, yields:

1. Cross-religious manifold invariance: all major world religions
operate on the same nine-dimensional moral manifold, differing in metric
but not in dimensional structure.

2. The Epistemic Prerequisite Paradox: the Genesis Fall narrative
contains a mathematical impossibility --- the boundary penalty for the
prohibited action requires a capacity that crossing the boundary
provides.

3. Euthyphro resolution via gauge invariance: the empirical stability of
the manifold across theological traditions supports moral realism
(gauge-invariant structure) over divine command theory (gauge-dependent
metric).

4. Theodicy as dimensional projection: the Problem of Evil assumes
scalar optimization, but the manifold geodesic operates on nine
dimensions. This clarifies the problem without resolving it.

5. Incompleteness constraints: deductive natural theology commits gauge
reification, and Godel\textquotesingle s theorems establish that the gap
between formal validity and metaphysical truth cannot be closed from
within.

6. Sacred-value convergence: the infinite-penalty boundaries of major
religions converge on cooperation-critical dimensions, predicted by the
manifold geometry.

7. Religious moral training as heuristic calibration: systematic tuning
of h(n), subject to the same admissibility analysis as secular moral
reasoning.

\section{Chapter 25: Geometric Environmental Ethics --- Climate,
Commons, and Intergenerational Obligation on the Moral
Manifold}\label{chapter-25-geometric-environmental-ethics-climate-commons-and-intergenerational-obligation-on-the-moral-manifold}

Epistemic status of this chapter. Except where a claim is explicitly
marked as a formal result or a conditional theorem, the material here is
a modeling interpretation, domain analogy, or speculative extension - a
demonstration of how the geometric vocabulary applies to environmental
ethics, not a set of established results at the evidential standard of
the field.

\begin{quote}
\textbf{\textsc{RUNNING EXAMPLE --- Priya's Model}}

\emph{Three of BEACON-7's trial sites are in Louisiana's `Cancer
Alley'---an 85-mile corridor of petrochemical plants along the
Mississippi. Patients living near these sites have higher baseline
cancer rates from environmental exposure. TrialMatch's geographic
proximity bonus effectively rewards living in polluted areas, creating a
perverse coupling between environmental injustice and trial access. The
intergenerational dimension---children growing up in Cancer
Alley---extends the manifold in directions TrialMatch's quarterly
planning horizon cannot reach.}
\end{quote}

This chapter applies the Geometric Ethics framework to environmental and
climate ethics --- the domain where moral reasoning confronts its
hardest structural challenges simultaneously. Environmental decisions
are intergenerational (the agents most affected do not yet exist),
irreversible (species extinction and climate tipping points cannot be
undone), and global (no agent can solve the problem alone). The central
finding is that decades of climate policy paralysis are not a failure of
political will but a failure of dimensionality: environmental
cost-benefit analysis collapses a nine-dimensional decision to a scalar,
and the resulting policy disputes --- Stern versus Nordhaus, carbon tax
versus cap-and-trade, adaptation versus mitigation --- are arguments
about a parameter in the wrong-dimensional space. The framework provides
the mathematical structure that scalar environmental economics lacks:
intergenerational pathfinding on a time-extended decision complex,
planetary boundaries as sacred-value constraints with topological
consequences, and the tragedy of the commons as a multi-agent manifold
failure with known geometric solutions.

\subsection{25.1 The Failure of Scalar Environmental
Policy}\label{the-failure-of-scalar-environmental-policy}

Environmental policy has been dominated for half a century by
cost-benefit analysis (CBA) --- the practice of converting all
environmental consequences to monetary equivalents and selecting the
policy with the highest net present value. CBA is scalar collapse: it
projects the nine-dimensional moral manifold onto d\_1 (monetary cost
and benefit) and optimizes on that projection alone.

The consequences of this scalar collapse have been catastrophic. From
the first IPCC report (1990) to the present, the scientific consensus on
anthropogenic climate change has strengthened continuously, yet global
CO2 emissions have increased by over 60\%. The policy paralysis is not
due to scientific uncertainty --- it is due to the impossibility of
encoding multi-dimensional moral obligations in a single monetary
parameter.

The Stern-Nordhaus debate illustrates the failure precisely. Nicholas
Stern (2006) used a social discount rate of r = 0.1\%, producing a
present value of climate damages that justified immediate, aggressive
mitigation costing 1-2\% of global GDP. William Nordhaus (2007) used r =
3-5\%, producing a present value that justified gradual, modest
mitigation. The policy prescriptions differ by orders of magnitude, yet
both analyses are internally consistent --- the difference is entirely
in the choice of a single scalar parameter.

The geometric diagnosis: the discount rate is being asked to do the work
of nine dimensions. It must simultaneously encode how much we value
future consequences (d\_1), what obligations we owe to future
generations (d\_2), how to weigh the claims of developing nations
against developed nations (d\_3), whether future people have autonomy
rights that constrain our choices (d\_4), what level of trust
international agreements can sustain (d\_5), the social and ecological
impact of climate change (d\_6), the relationship between cultural
identity and the natural world (d\_7), the legitimacy of international
climate institutions (d\_8), and the epistemic status of climate
projections under deep uncertainty (d\_9). No scalar can carry this
information. The dispute between Stern and Nordhaus is not about the
discount rate --- it is about which dimensions of the manifold matter
and how much, compressed into a single number that cannot represent the
answer.

\subsection{25.2 The Environmental Decision
Complex}\label{the-environmental-decision-complex}

\textbf{Definition 25.1 (Environmental Decision Complex)}. The
environmental decision complex E\_env is a weighted simplicial complex
whose vertices are environmental states --- configurations of
atmospheric composition, biodiversity, land use, ocean chemistry,
institutional arrangements, and community conditions --- and whose edges
are environmental actions (emission reductions, land-use changes,
technology deployments, policy implementations, conservation decisions).
Each vertex v\_i carries an attribute vector a(v\_i) in R\^{}9.

The nine dimensions of the moral manifold, instantiated for
environmental contexts, are:

\begin{itemize}
\item
  \textbf{d\_1: Environmental consequences} --- climate damages, health
  impacts of pollution, agricultural productivity loss, ecosystem
  service degradation. The dimension that cost-benefit analysis targets.
\item
  \textbf{d\_2: Environmental obligations} --- duties to future
  generations, treaty commitments, fiduciary responsibilities of
  governments to their citizens, the precautionary principle as an
  obligation under uncertainty.
\item
  \textbf{d\_3: Environmental fairness} --- distributive justice between
  developed and developing nations, between current and future
  generations, between communities that benefit from emissions and
  communities that bear the costs.
\item
  \textbf{d\_4: Developmental autonomy} --- the right of developing
  nations to industrialize, the freedom of communities to choose their
  energy sources, constraints on individual consumption choices.
\item
  \textbf{d\_5: Environmental trust} --- reliability of international
  climate agreements, credibility of emissions targets, trust between
  nations in burden-sharing arrangements, public trust in climate
  science.
\item
  \textbf{d\_6: Ecological and social impact} --- biodiversity loss,
  ecosystem disruption, community displacement from sea-level rise,
  climate migration, disruption of food systems and water resources.
\item
  \textbf{d\_7: Environmental identity} --- indigenous
  peoples\textquotesingle{} relationship to land, cultural practices
  tied to specific ecosystems, the role of nature in community identity,
  environmental virtue and the concept of ecological citizenship.
\item
  \textbf{d\_8: Institutional legitimacy} --- authority of the UNFCCC
  and IPCC, legitimacy of carbon markets, enforceability of the Paris
  Agreement, democratic legitimacy of climate policy.
\item
  \textbf{d\_9: Epistemic status} --- scientific certainty of climate
  projections, uncertainty in damage functions, model disagreement on
  tipping points, the epistemology of deep uncertainty over
  century-scale timelines.
\end{itemize}

\textbf{Definition 25.2 (Environmental Edge Weights)}. The weight of an
edge (v\_i, v\_j) in E\_env is:

w(v\_i, v\_j) = Da\^{}T Sigma\_env\^{}\{-1\} Da + Sum\_k beta\_k *
1{[}boundary k crossed{]}

where Da = a(v\_j) - a(v\_i), Sigma\_env is the 9x9 environmental
covariance matrix, and beta\_k are boundary penalties for crossing
environmental boundaries. Critical covariance terms in Sigma\_env
include: Sigma\_\{1,3\} (consequences x fairness: economic costs of
climate change correlate with distributive injustice, as the poorest
nations face the worst damages from emissions they did not produce),
Sigma\_\{4,3\} (autonomy x fairness: developing
nations\textquotesingle{} right to industrialize conflicts with
equitable emissions allocation), Sigma\_\{5,8\} (trust x legitimacy: the
credibility of climate agreements depends on institutional legitimacy),
and Sigma\_\{6,7\} (ecological impact x identity: biodiversity loss and
ecosystem destruction directly degrade the cultural identity of
communities whose lives are tied to the land).

\textbf{Definition 25.3 (Environmental Boundary Types)}. The
environmental decision complex has five principal boundary types:

\begin{itemize}
\item
  Planetary boundaries: biogeophysical thresholds beyond which Earth
  system stability is compromised (Rockstrom et al., 2009). Beta values
  range from finite (nitrogen cycle disruption) to effectively infinite
  (climate tipping points).
\item
  Species extinction thresholds: population viability limits below which
  species cannot recover. Once crossed, these boundaries are absorbing
  --- the manifold topology changes irreversibly.
\item
  Tipping point boundaries: critical thresholds in the climate system
  (ice sheet collapse, permafrost methane release, Amazon dieback)
  beyond which positive feedback makes the system self-reinforcing.
  These are topological phase transitions on the manifold.
\item
  Indigenous rights boundaries: sacred-value constraints protecting
  indigenous peoples\textquotesingle{} relationship to land and
  traditional ecological knowledge. Beta = infinity in frameworks that
  recognize indigenous sovereignty.
\item
  Intergenerational obligation boundaries: constraints encoding duties
  to future generations that cannot be traded off against
  present-generation benefits on d\_1.
\end{itemize}

\subsection{25.3 Intergenerational
Pathfinding}\label{intergenerational-pathfinding}

Environmental decisions are structurally unique among all the domain
chapters in this monograph. In clinical ethics (Chapter 21), the
affected patient exists and can participate in decision-making. In
finance (Chapter 23), counterparties are present agents. In law (Chapter
22), litigants are alive and represented. But environmental decisions
--- particularly climate policy --- primarily affect agents who do not
yet exist. Future generations cannot participate in the pathfinding that
determines the manifold they will inhabit.

This asymmetry is not merely a practical difficulty. It is a structural
feature of the environmental decision complex that has no analog in
other domains.

\textbf{Definition 25.4 (Time-Extended Environmental Decision Complex)}.
The time-extended environmental decision complex E\_env(t) is a family
of decision complexes indexed by time t, where edges at time t modify
vertex attributes at time t + delta. Formally, let E\_env(t) = (V(t),
E(t), a(t), w(t)) be the decision complex at time t. An action e in E(t)
modifies the attribute function: a(t + delta) = Phi(a(t), e), where Phi
is the environmental transition function encoding physical, ecological,
and social dynamics.

The critical feature: the transition function Phi is not symmetric in
time. Current-generation actions modify the decision complex itself for
future agents. Burning fossil fuels today does not merely change the
atmospheric state --- it changes the set of available vertices and edges
for future generations. Some paths that are currently available (e.g.,
stable climate with intact ice sheets) become permanently unavailable
once tipping points are crossed.

\textbf{Theorem 25.1 (Intergenerational Pathfinding Asymmetry)}. Let
G\_t denote the set of agents at time t, and let G\_\{t+delta\} denote
the set of agents at time t + delta. The actions of G\_t modify the
decision complex E\_env(t + delta) on which G\_\{t+delta\} must perform
their pathfinding. However, G\_\{t+delta\} cannot influence E\_env(t).
The pathfinding is temporally asymmetric: current generations set the
manifold for future generations but not vice versa.

\textbf{Proof}. The transition function Phi maps (a(t), e) to a(t +
delta). For any action e chosen by G\_t, the resulting a(t + delta)
determines the vertex attributes, edge weights, and available paths in
E\_env(t + delta). Since G\_\{t+delta\} does not exist at time t, no
action by G\_\{t+delta\} can influence the choice of e. The asymmetry
follows from the temporal ordering of existence: agents can only act on
decision complexes contemporaneous with or subsequent to their
existence, never on prior ones. {[}{]}

\textbf{Corollary 25.1 (Non-Negotiability of Intergenerational
Decisions)}. Because G\_\{t+delta\} cannot participate in pathfinding at
time t, intergenerational environmental decisions cannot be modeled as
negotiations or bilateral transactions. The Bond Geodesic Equilibrium
(Chapter 20) requires that all agents participate in the equilibrium
computation. When some agents do not yet exist, the BGE is structurally
incomplete --- there is no equilibrium in the game-theoretic sense, only
a unilateral imposition by the current generation on all future
generations.

\textbf{Remark (Representation of Future Agents)}. Various institutional
mechanisms attempt to represent future generations in current
decision-making: constitutional environmental provisions, commissioners
for future generations (as in Wales and Hungary), long-term sovereign
wealth funds (Norway), and extended cost-benefit horizons. In the
geometric framework, these are proxy heuristics --- estimates of h(n)
for future agents, computed by current agents. The framework predicts
that the quality of intergenerational environmental decisions should
correlate with the fidelity of these proxy heuristics.

\subsection{25.4 The Discount Rate as Dimensional
Collapse}\label{the-discount-rate-as-dimensional-collapse}

The social discount rate is the central parameter of environmental
cost-benefit analysis. It determines how much weight current decisions
give to future costs and benefits. A high discount rate (Nordhaus:
3-5\%) implies that a dollar of climate damage in 100 years is worth
less than a cent today. A low discount rate (Stern: 0.1\%) implies that
future damages are nearly as significant as present ones. The policy
difference between these two rates is the difference between modest
carbon taxes and wartime-level mobilization.

\textbf{Theorem 25.2 (Discount Rate as Dimensional Collapse)}. The
social discount rate r collapses the temporal dimension of the moral
manifold to a scalar via the discount operator D\_r: R\^{}9 x T
-\textgreater{} R defined by D\_r(a, t) = e\^{}\{-rt\} * pi\_1(a), where
pi\_1 is the projection onto d\_1. The discount operator has the
following properties:

\begin{enumerate}
\def\labelenumi{(\roman{enumi})}
\item
  D\_r projects from R\^{}9 to R\^{}1: only d\_1 (monetary) costs are
  discounted. Future d\_2 (obligations to future persons), d\_3
  (fairness across generations), d\_4 (autonomy of future agents), d\_6
  (ecological impact), and d\_7 (identity and relationship to the
  natural world) have no natural exponential decay. A
  grandmother\textquotesingle s obligation to leave a livable planet to
  her grandchildren does not decay at 3\% per annum.
\item
  D\_r is not injective: by the rank-nullity theorem, the discount
  operator D\_r: R\^{}9 -\textgreater{} R has kernel of dimension at
  least 8. Eight dimensions of the moral manifold are annihilated by
  discounting. The information destroyed is mathematically irrecoverable
  --- no manipulation of the discount rate can restore the lost
  dimensions.
\item
  The Stern-Nordhaus dispute is a dispute about a parameter in a space
  that is too low-dimensional to represent the quantity being
  parameterized. Both Stern and Nordhaus are wrong --- not in their
  choice of r, but in their shared assumption that a single scalar can
  encode intergenerational moral weight.
\end{enumerate}

\textbf{Proof}. The rank of D\_r as a linear operator from R\^{}9 to R
is at most 1 (since the codomain is one-dimensional). By rank-nullity,
dim(ker(D\_r)) = 9 - rank(D\_r) \textgreater= 9 - 1 = 8. The kernel
contains all attribute vectors with zero d\_1 component --- precisely
the vectors encoding non-monetary moral content (obligations, fairness,
autonomy, trust, social impact, identity, legitimacy, epistemic status).
By the Scalar Irrecoverability Theorem (Chapter 15), no continuous
injection R\^{}9 -\textgreater{} R exists, so no choice of r or
functional form for the discount factor can make D\_r injective. The
information loss is topological, not parametric. {[}{]}

\textbf{Remark (Ramsey\textquotesingle s Equation)}. The standard
decomposition r = delta + eta * g, where delta is pure time preference,
eta is elasticity of marginal utility, and g is per-capita growth rate,
illustrates the dimensional collapse. Delta encodes d\_2 (do we have
obligations to future persons?) as a number. Eta encodes d\_3 (how
should we weigh the welfare of richer future people against poorer
current people?) as a number. g encodes d\_9 (what do we know about
future economic growth?) as a number. Each of these moral dimensions is
compressed to a single parameter in a single equation. The persistent
disagreement about the "correct" discount rate is the predictable
consequence of this compression: different analysts encode different
dimensional weightings into the same scalar, producing different
scalars, with no principled basis for choosing among them.

\subsection{25.5 The Tragedy of the Commons as Multi-Agent Manifold
Failure}\label{the-tragedy-of-the-commons-as-multi-agent-manifold-failure}

Garrett Hardin\textquotesingle s tragedy of the commons (1968) describes
the logic by which shared resources are overexploited: each agent
benefits individually from additional extraction (positive Da on d\_1)
while the cost of depletion is shared across all agents (negative Da on
d\_6, distributed across n agents so each agent bears only 1/n of the
cost). When n is large, each agent\textquotesingle s private d\_1
benefit exceeds their share of the d\_6 cost, and the rational action is
to extract --- producing collective ruin.

In the geometric framework, the tragedy of the commons is not a failure
of rationality. It is a failure of dimensionality.

\textbf{Theorem 25.3 (Tragedy of the Commons as Dimensional
Restriction)}. In a commons with n agents, each optimizing on d\_1
alone, the Bond Geodesic Equilibrium (BGE) on the d\_1 projection is
Pareto-dominated by the BGE on the full manifold. Specifically, let
BGE\_1 be the equilibrium in which each agent\textquotesingle s edge
weights include only d\_1, and let BGE\_full be the equilibrium in which
edge weights include all nine dimensions with nonzero weight on d\_6
(social and ecological impact). Then for sufficiently large n, every
agent achieves higher long-run d\_1 value at BGE\_full than at BGE\_1.

\textbf{Proof}. At BGE\_1, each agent maximizes d\_1 without
internalizing d\_6 costs. The aggregate extraction rate exceeds the
regeneration rate of the commons, producing resource depletion. The d\_1
trajectory at BGE\_1 is: initial gain, followed by declining returns,
followed by collapse (zero d\_1 for all agents). At BGE\_full, d\_6
enters each agent\textquotesingle s edge weight computation. The
covariance term Sigma\_\{1,6\} (return x ecological impact) penalizes
extraction rates that degrade the commons. The optimal path under the
full metric involves sustainable extraction at a rate below the
regeneration threshold. The d\_1 trajectory at BGE\_full is: lower
initial gain, but sustained returns over time. For any positive discount
rate, there exists a time horizon T beyond which the cumulative d\_1
value at BGE\_full exceeds the cumulative d\_1 value at BGE\_1 (because
BGE\_1 produces zero d\_1 after collapse). For intergenerational
decisions, where the relevant time horizon is effectively infinite,
BGE\_full strictly dominates. {[}{]}

The tragedy is a perceptual failure, not a rational one. Commons agents
who can perceive only d\_1 cannot see the cooperative geodesics that
exist on the full manifold. Adding d\_6 to the agent\textquotesingle s
perceptual apparatus --- making ecological and social impact visible in
the agent\textquotesingle s pathfinding --- transforms the decision
landscape from one where defection dominates to one where cooperation is
the geodesic.

\textbf{Ostrom\textquotesingle s Design Principles as Boundary Penalty
Calibration}. Elinor Ostrom\textquotesingle s (1990) empirical research
on successful commons governance identified eight institutional design
principles for sustainable commons management. In the geometric
framework, each principle is a boundary penalty calibration:

\begin{itemize}
\item
  Clearly defined boundaries (who may access the commons): boundary
  definition on the d\_4 dimension, specifying which
  agents\textquotesingle{} autonomy includes commons access.
\item
  Proportional equivalence between benefits and costs: calibration of
  d\_3 (fairness) weights so that extraction rights are proportional to
  maintenance contributions.
\item
  Collective choice arrangements: multi-agent pathfinding where all
  affected agents participate in the BGE computation, ensuring that d\_8
  (legitimacy) is nonzero.
\item
  Monitoring: epistemic infrastructure that makes d\_6 (ecological
  impact) and d\_9 (information about resource state) visible to all
  agents.
\item
  Graduated sanctions: a penalty schedule on d\_2 (obligations) that
  increases with repeated violation, calibrated to enforce cooperative
  geodesics without excessive punishment.
\item
  Conflict resolution mechanisms: institutional apparatus for resolving
  disputes about the metric tensor --- disagreements about relative
  dimensional weights.
\item
  Minimal recognition of rights to organize: d\_8 (institutional
  legitimacy) for the self-governing institution, granted by external
  authorities.
\item
  Nested enterprises: for commons that are part of larger systems,
  multi-scale decision complexes with boundary penalties at each level.
\end{itemize}

Ostrom\textquotesingle s central finding --- that communities can and do
manage commons sustainably without privatization or central authority
--- is the empirical confirmation that agents who incorporate d\_6 into
their pathfinding find cooperative geodesics. Her design principles are
the institutional conditions under which d\_6 becomes perceptible and
boundary penalties are calibrated to sustain the cooperative
equilibrium.

\subsection{25.6 Carbon Pricing as Boundary Penalty
Calibration}\label{carbon-pricing-as-boundary-penalty-calibration}

Carbon taxes and cap-and-trade systems are the two principal market
mechanisms for addressing climate change. In the geometric framework,
both are attempts to encode multi-dimensional environmental cost into a
d\_1 (monetary) signal.

\textbf{Proposition 25.1 (Carbon Pricing as Manifold-to-Scalar
Projection with Correction)}. A carbon price p is a scalar correction
term added to d\_1 that approximates the full-manifold cost of
emissions. Specifically, p attempts to represent, on d\_1 alone, the
aggregate cost across d\_2 (intergenerational obligations violated),
d\_3 (distributional unfairness imposed), d\_6 (ecological and social
damage), d\_7 (cultural and identity loss), and d\_9 (epistemic risk
from climate uncertainty). The carbon price is the d\_1 shadow of
d\_2-d\_9 costs:

p = Sum\_\{k=2\}\^{}\{9\} alpha\_k * c\_k

where c\_k is the estimated cost on dimension k and alpha\_k is the
conversion factor from dimension k to d\_1.

\textbf{Evidence}. The social cost of carbon (SCC) --- the dollar value
assigned to one additional ton of CO2 emissions --- is estimated by
integrated assessment models (IAMs) such as DICE, FUND, and PAGE. These
models attempt to convert climate damages to monetary equivalents. The
wide range of SCC estimates (\$10-\$200+ per ton, depending on the
model, discount rate, and damage function) reflects the impossibility of
the conversion: different models encode different dimensional weightings
into the conversion factors alpha\_k, producing different scalars. The
framework predicts this spread --- it is the necessary consequence of
projecting a nine-dimensional cost onto a one-dimensional price.

\textbf{Proposition 25.2 (Necessity and Insufficiency of Carbon
Pricing)}. Carbon pricing is necessary because without it, the d\_1
component of environmental cost is systematically underestimated ---
emissions impose costs on d\_2-d\_9 that are invisible to d\_1-only
agents. But carbon pricing is insufficient because the projection from
R\^{}9 to R\^{}1 is non-injective (Theorem 25.2). The information lost
in the projection --- the specific identity of which dimensions bear the
cost, the distribution of costs across communities, the irreversibility
of certain damages --- cannot be recovered from a single price signal.

\textbf{Revenue Recycling as Partial Metric Correction}. Carbon tax
revenue can be recycled in ways that partially address the dimensional
information lost in the price signal. Returning revenue to frontline
communities addresses d\_3 (fairness) and d\_7 (identity). Investing
revenue in climate adaptation addresses d\_6 (social impact). Funding
climate research addresses d\_9 (epistemic status). In the geometric
framework, revenue recycling is a partial metric correction --- an
attempt to restore, through targeted expenditure, the dimensional
information that the carbon price destroyed through scalar compression.

\textbf{Carbon Tax versus Cap-and-Trade}. The two mechanisms differ in
which quantity is fixed. A carbon tax fixes the price (the d\_1
correction term) and lets the quantity of emissions adjust.
Cap-and-trade fixes the quantity (the d\_6 boundary) and lets the price
adjust. In the geometric framework, cap-and-trade is superior on one
specific criterion: it directly sets a boundary on d\_6 (total
emissions), while the carbon tax sets a d\_1 correction and hopes the
d\_6 boundary is respected. When the environmental concern is a tipping
point --- a boundary whose crossing produces topological phase
transition --- fixing the boundary directly is geometrically more
appropriate than fixing a scalar approximation to the boundary.

\subsection{25.7 Species Extinction as Irreversible Boundary
Crossing}\label{species-extinction-as-irreversible-boundary-crossing}

The environmental decision complex contains a class of boundaries
fundamentally different from any encountered in the preceding domain
chapters. In clinical ethics, a misdiagnosis can sometimes be corrected.
In finance, a bad trade can be unwound. In law, a wrongful conviction
can be overturned. But species extinction is permanent. Once crossed, no
path on the decision complex returns to a state where the species
exists.

\textbf{Theorem 25.4 (Extinction as Absorbing Boundary)}. Let S be a
species represented by a set of vertices V\_S in the environmental
decision complex E\_env. Let e\_extinction be the edge corresponding to
the last viable population dropping below the minimum viable threshold.
Then e\_extinction is an absorbing boundary: for any path gamma on
E\_env, if gamma traverses e\_extinction, no extension of gamma returns
to a vertex in V\_S. Formally, if gamma(t\_0) in V\_S and gamma(t\_1)
not in V\_S with gamma traversing e\_extinction at some t in (t\_0,
t\_1), then for all t\_2 \textgreater{} t\_1, gamma(t\_2) not in V\_S.

\textbf{Proof}. Species extinction is thermodynamically irreversible:
the genetic information encoding the species, accumulated over millions
of years of evolution, is destroyed. De-extinction technologies (cloning
from preserved DNA) can at best reconstruct a phenotypic approximation,
not the full genotype, epigenome, microbiome, and behavioral repertoire
of the original species. The absorbing property follows from the
irreversibility: the vertices in V\_S (states where species S exists
with its full ecological role) are permanently removed from the decision
complex. No sequence of edges from the post-extinction complex can reach
a vertex in V\_S because V\_S is no longer in the vertex set. {[}{]}

This is structurally different from finite boundary penalties. A
boundary with beta\_k = 10\^{}6 imposes high cost but permits traversal
and, potentially, return. The extinction boundary does not merely impose
high cost --- it removes vertices from the decision complex. The
manifold topology changes irreversibly. This is not a penalty; it is an
amputation.

\textbf{Definition 25.5 (Biodiversity as Manifold Richness)}. Each
species S\_j contributes a set of vertices V\_\{S\_j\} and edges
E\_\{S\_j\} to the environmental decision complex, representing the
ecological states and processes that S\_j enables (pollination, nutrient
cycling, predation, seed dispersal, decomposition). The biodiversity B
of the environmental decision complex is the cardinality of the vertex
and edge sets contributed by the full set of species:

B(E\_env) = \textbar Union\_j V\_\{S\_j\}\textbar{} + \textbar Union\_j
E\_\{S\_j\}\textbar{}

\textbf{Each extinction event reduces B irreversibly}. Critically,
species do not contribute independently: edges connect vertices across
species (ecological interactions), so the extinction of one species can
remove edges involving other species, potentially triggering cascading
vertex removal (trophic cascades, coextinction of dependent species).

\textbf{Proposition 25.3 (Extinction as Path Space Reduction)}. Each
species extinction reduces the dimensionality of available future paths
on the environmental decision complex. The set of states reachable from
the current state shrinks monotonically with each extinction. Future
generations inherit a decision complex with fewer vertices, fewer edges,
and fewer available geodesics --- not because they chose poorly, but
because current-generation extinctions permanently narrowed their
manifold.

\textbf{Remark (The Sixth Mass Extinction)}. The current rate of species
extinction is estimated at 100 to 1,000 times the background rate
(Ceballos et al., 2015). In manifold terms, the environmental decision
complex is losing vertices and edges at a rate unprecedented in 66
million years. The framework makes the cost precise: each extinction is
not merely a "loss of biodiversity" in the abstract but a permanent
reduction in the path space available to all future agents --- an
irreversible narrowing of the manifold on which future moral and
economic decisions will be made.

\subsection{25.8 Environmental Justice as Metric
Asymmetry}\label{environmental-justice-as-metric-asymmetry}

Environmental burdens --- air pollution, toxic waste, water
contamination, climate vulnerability --- fall disproportionately on
communities with less political and economic power. This is the
empirical finding of the environmental justice literature (Bullard,
1990; Mohai et al., 2009). In the geometric framework, environmental
injustice is metric asymmetry: the covariance matrix Sigma\_env differs
systematically between communities, so the same environmental state has
different manifold costs for different populations.

\textbf{Theorem 25.5 (Environmental Injustice as Metric Asymmetry)}. Let
Sigma\_env\^{}A and Sigma\_env\^{}B be the environmental covariance
matrices for communities A and B, where A is a frontline community (low
income, politically marginalized, often a community of color) and B is
an affluent community. Then for the same environmental state change Da
(e.g., the same increase in particulate matter concentration):

Da\^{}T (Sigma\_env\^{}A)\^{}\{-1\} Da \textgreater\textgreater{}
Da\^{}T (Sigma\_env\^{}B)\^{}\{-1\} Da

The manifold cost of the same environmental state is systematically
higher for community A than for community B.

\textbf{Proof}. The covariance matrix Sigma\_env encodes the
co-variation of dimensions in a community\textquotesingle s
environmental context. For community A, the relevant covariance terms
differ from community B in three systematic ways. (1) Lower d\_1 base:
community A has less economic buffer, so the same d\_1 cost (health
expenses, property damage, lost productivity) represents a larger share
of resources --- the inverse covariance (Sigma\_env\^{}A)\^{}\{-1\} has
larger diagonal entries on d\_1. (2) Lower d\_4 autonomy: community A
has less political power to relocate, to demand remediation, or to block
polluting facilities --- the d\_4 component of Da is larger (more
autonomy is lost). (3) Higher d\_7 impact: community A\textquotesingle s
identity is more tightly coupled to place (generational homes, cultural
practices tied to specific land) --- the d\_7 component of Da is larger.
The combined effect is that the Mahalanobis distance for the same
environmental change is systematically larger for community A. {[}{]}

\textbf{Environmental Racism as Gauge-Invariance Violation}. The Bond
Invariance Principle (Chapter 12) requires that moral evaluations be
gauge-invariant --- the evaluation must not depend on how the situation
is described. Environmental racism violates this principle: the same
pollution source, described as "industrial development" when proposed
near an affluent community (where it is typically blocked) versus
"economic opportunity" when proposed near a low-income community of
color (where it is typically approved), receives different evaluations.
The description changes; the attribute vector does not. This is a
gauge-dependent evaluation --- a BIP violation.

\textbf{Proposition 25.4 (Formal Structure of Environmental Racism)}.
Let v be an environmental state (e.g., a proposed polluting facility at
a specific location). Let G\_1 and G\_2 be two gauge descriptions of v:
G\_1 = "industrial development bringing jobs" and G\_2 = "pollution
source with health risks." If the evaluation of v under G\_1 differs
from the evaluation of v under G\_2, and if the choice between G\_1 and
G\_2 correlates with the racial or economic composition of the affected
community, then the evaluation violates gauge invariance on a dimension
correlated with race and class. This is environmental racism formalized
as systematic gauge-dependent evaluation.

\textbf{Remark (Cancer Alley)}. Louisiana\textquotesingle s 85-mile
industrial corridor between Baton Rouge and New Orleans --- known as
Cancer Alley --- hosts over 150 petrochemical plants and refineries. The
corridor\textquotesingle s population is predominantly Black. Cancer
rates in the corridor are significantly elevated. In the geometric
framework, Cancer Alley is a maximal metric asymmetry: communities bear
extreme d\_1 (health), d\_6 (social impact), and d\_7 (identity/place)
costs while receiving minimal d\_1 (economic) benefits. The
facilities\textquotesingle{} benefits accrue to shareholders and
consumers elsewhere --- a spatial separation of d\_1 benefit from
d\_1/d\_6/d\_7 cost that would not survive evaluation on the full
manifold.

\subsection{25.9 Planetary Boundaries as Sacred-Value
Constraints}\label{planetary-boundaries-as-sacred-value-constraints}

Rockstrom et al. (2009) identified nine planetary boundaries ---
thresholds in Earth system processes beyond which the stability of the
Holocene climate is compromised. In the geometric framework, planetary
boundaries are not arbitrary policy thresholds. They are tipping points
beyond which the environmental decision complex undergoes topological
phase transition.

Definition 25.6 (Planetary Boundary as Topological Phase Transition). A
planetary boundary B\_j is a threshold value on dimension d\_k (or a
combination of dimensions) such that for manifold states below B\_j, the
environmental decision complex E\_env has topology T\_1, and for states
above B\_j, the decision complex transitions to topology T\_2, where
T\_1 and T\_2 are not homeomorphic. The transition is a phase transition
in the topological sense: the qualitative structure of the manifold
changes, not just its metric.

The nine planetary boundaries, mapped onto the environmental manifold:

\begin{itemize}
\item
  Climate change (atmospheric CO2 concentration): boundary at
  \textasciitilde350 ppm (already exceeded at \textasciitilde420 ppm).
  Beyond this boundary, the climate system enters a regime of amplifying
  feedbacks (ice-albedo, permafrost methane, water vapor) that
  qualitatively change the manifold --- new vertices appear (extreme
  weather states, sea-level configurations not seen in human history)
  while old vertices disappear (stable Holocene climate states).
\item
  Biodiversity loss (species extinction rate): boundary at 10 E/MSY
  (extinctions per million species-years). The current rate
  (\textasciitilde100-1000 E/MSY) massively exceeds this boundary. Each
  extinction removes vertices from the decision complex (Theorem 25.4).
\item
  Nitrogen cycle (industrial nitrogen fixation): boundary at 35 Tg N/yr
  (currently \textasciitilde150 Tg N/yr). Excess reactive nitrogen
  creates dead zones, degrades freshwater, and disrupts terrestrial
  ecosystems.
\item
  Phosphorus cycle (phosphorus flow to oceans): boundary at 11 Tg P/yr.
  Excess phosphorus triggers eutrophication and ocean anoxic events.
\item
  Ocean acidification (aragonite saturation state): boundary at 80\% of
  pre-industrial value. CO2 absorption by oceans reduces pH, threatening
  calcifying organisms and marine food webs.
\item
  Land-system change (fraction of original forest cover): boundary at
  75\% (currently \textasciitilde62\%). Deforestation disrupts
  hydrological cycles, carbon storage, and biodiversity.
\item
  Freshwater use (consumptive water use): boundary at 4000 km\^{}3/yr
  (currently \textasciitilde2600 km\^{}3/yr). Approaching the boundary
  in many regional systems.
\item
  Stratospheric ozone depletion (ozone concentration in DU): boundary at
  276 DU. Successfully managed through the Montreal Protocol --- the one
  planetary boundary case where international cooperation produced
  effective boundary enforcement.
\item
  Atmospheric aerosol loading (regional measure): boundary not yet
  quantified globally. Aerosol effects on monsoon systems and regional
  climate.
\end{itemize}

\textbf{Proposition 25.5 (Planetary Boundaries as Sacred-Value
Constraints)}. Planetary boundaries function as sacred-value constraints
(beta\_k = infinity) on the environmental decision complex, but with a
structural justification that differs from the sacred-value boundaries
of earlier chapters. In theology (Chapter 24), sacred-value boundaries
are set by tradition and authority. In the environmental context,
sacred-value boundaries are set by physics and ecology. They are
non-negotiable not because a text declares them sacred but because the
Earth system does not negotiate: crossing a tipping point triggers
physical feedbacks that no policy, technology, or institution can
reverse on human timescales.

\textbf{Remark (The Montreal Protocol as Successful Boundary
Enforcement)}. The Montreal Protocol (1987) is the single most
successful case of planetary boundary management. It identified the
ozone depletion boundary, set boundary penalties (phase-out of CFCs and
HCFCs), and enforced them through international agreement with trade
sanctions for non-compliance. In the geometric framework, the Protocol
succeeded because: (1) the boundary was clearly defined and
scientifically unambiguous (d\_9: high epistemic certainty), (2) the
cost of compliance was moderate and asymmetrically distributed to a
small number of chemical manufacturers (d\_1: manageable cost), (3)
substitutes existed (d\_4: autonomy preserved), and (4) the benefits of
compliance were visible within a generation (d\_2: near-term obligation
fulfillment). Climate change fails on all four conditions: the boundary
is harder to define precisely, the cost of compliance is massive and
distributed across entire economies, substitutes for fossil energy are
still being scaled, and the benefits of compliance accrue to future
generations who cannot vote.

\subsection{25.10 Geoengineering as Manifold
Surgery}\label{geoengineering-as-manifold-surgery}

Geoengineering --- deliberate large-scale intervention in the
Earth\textquotesingle s climate system --- comes in two forms: solar
radiation management (SRM), which reflects sunlight to cool the planet
without reducing CO2, and carbon dioxide removal (CDR), which removes
CO2 from the atmosphere. In the geometric framework, both are manifold
surgery: deliberate modification of the environmental decision
complex\textquotesingle s topology and metric.

\textbf{Definition 25.7 (Manifold Surgery)}. A manifold surgery on the
environmental decision complex E\_env is an action that modifies the
topology or metric of E\_env itself, rather than finding a better path
on the existing complex. Formally, a surgery S maps E\_env to
E\textquotesingle\_env = S(E\_env), where E\textquotesingle\_env has
different vertex sets, edge sets, or metric structure than E\_env.

\textbf{Proposition 25.6 (Solar Radiation Management as Metric
Modification without Topological Repair)}. SRM (stratospheric aerosol
injection, marine cloud brightening) modifies the metric of the
environmental decision complex --- reducing the effective d\_1 cost of
warming --- without addressing the topological changes caused by
elevated CO2. The decision complex under SRM has a different metric
(lower temperature-related costs) but the same underlying topology
(ocean acidification, ecosystem CO2 effects, and the risk of termination
shock persist). SRM is cosmetic surgery: it changes the
manifold\textquotesingle s surface appearance without correcting its
deep structure.

\textbf{Proposition 25.7 (Carbon Dioxide Removal as Topological
Repair)}. CDR (direct air capture, enhanced weathering, afforestation)
addresses the root topological change by removing the substance (CO2)
that caused the phase transition. If CDR could return atmospheric CO2 to
pre-industrial levels, the manifold topology would (partially) revert.
However, some topological changes are absorbing (extinct species,
collapsed ice sheets, lost permafrost carbon) and cannot be restored by
CO2 removal alone.

\textbf{Moral Hazard as Dimensional Masking}. The primary moral concern
about SRM is moral hazard: if geoengineering reduces the perceived d\_1
cost of climate change, it may reduce motivation to reduce emissions ---
addressing the symptom while perpetuating the cause. In the geometric
framework, this is dimensional masking: SRM reduces the d\_1 signal that
would otherwise motivate action on d\_2-d\_9, without actually reducing
the d\_2-d\_9 costs. The manifold cost is still high, but the most
visible dimension (d\_1: temperature-related damages) is suppressed,
making the full-manifold cost less perceptible.

\textbf{Unilateral Geoengineering as Gauge-Invariance Violation}.
Because SRM affects the entire planet\textquotesingle s climate,
deployment by a single nation modifies the environmental decision
complex for all nations and all future generations without their
consent. In the geometric framework, this is a gauge-invariance
violation at civilizational scale: one agent imposes its gauge (its
dimensional weighting, its risk tolerance, its preferred temperature) on
all other agents. The same SRM deployment is described as "climate
emergency response" by the deploying nation and "unilateral climate
imperialism" by nations experiencing adverse side effects (changed
monsoon patterns, reduced precipitation). The evaluation depends on the
gauge, not on the attribute vector --- a Bond Invariance Principle
violation.

\textbf{Proposition 25.8 (Geoengineering Governance as Multi-Agent
Manifold Surgery Protocol)}. Legitimate geoengineering requires a
governance framework in which all agents whose manifold is modified by
the surgery participate in the decision. This is the intergenerational
pathfinding problem (Theorem 25.1) compounded by the global collective
action problem: not only must future generations be represented, but all
current nations must participate. The governance requirements for
legitimate geoengineering are the most demanding of any environmental
decision --- simultaneously intergenerational, irreversible, and
universal.

\subsection{25.11 Worked Examples}\label{worked-examples-3}

The preceding sections have developed the geometric framework for
environmental ethics abstractly --- defining intergenerational
pathfinding, commons dynamics, planetary boundaries, and environmental
justice in terms of manifold structure, metric tensors, and topological
phase transitions. This section applies the full framework to three
real-world environmental cases, demonstrating how nine-dimensional
analysis changes both the diagnosis and the prescription relative to
scalar cost-benefit analysis. Each example walks through all nine
dimensions of the moral manifold, showing how the scalar projection that
dominates conventional environmental economics systematically
undervalues the true manifold cost of environmental degradation.

\textbf{Example 25.1 (Great Barrier Reef Mass Bleaching, 2016--2024)}.
Between 2016 and 2024, the Great Barrier Reef experienced six mass
bleaching events --- in 2016, 2017, 2020, 2022, 2024 (aerial surveys by
the Australian Institute of Marine Science), and a severe event in early
2024 confirmed as the worst on record. These events were driven by
marine heatwaves caused by anthropogenic ocean warming. The Australian
government\textquotesingle s initial policy response was overwhelmingly
scalar: reef management was framed in terms of tourism revenue impact
and the cost-benefit ratio of emission reduction versus reef
preservation spending. The Reef 2050 Long-Term Sustainability Plan
allocated A\$3 billion over a decade --- substantial in absolute terms,
but structured almost entirely as d\_1 (economic cost-benefit)
optimization. A full nine-dimensional analysis reveals why the scalar
framing produced inadequate policy.

Dimension d\_1 (Consequences --- economic cost and benefit): The Great
Barrier Reef generates approximately A\$6.4 billion annually in tourism
revenue, supports 64,000 jobs directly and indirectly, and underpins
commercial and recreational fisheries worth hundreds of millions
annually. Reef degradation threatens all of these. But the d\_1 analysis
alone produced a perverse result: because Australia\textquotesingle s
coal and gas exports generate over A\$100 billion annually, the scalar
cost-benefit calculation --- tourism revenue lost from reef decline
versus export revenue lost from emission reduction --- appeared to favor
continued fossil fuel expansion with reef "management" spending as
mitigation. This is the Scalar Irrecoverability Theorem (Chapter 15) in
action: the d\_1 projection makes the decision appear close to
indifferent, when the full-manifold cost is overwhelmingly one-sided.
Dimension d\_2 (Obligations): Australia has binding obligations under
the UN Convention on Biological Diversity (ratified 1993) and the World
Heritage Convention (the reef was inscribed in 1981). More critically,
Pacific Island nations --- Tuvalu, Kiribati, the Marshall Islands, Fiji,
Vanuatu --- depend on reef ecosystems for food security, coastal
protection, and fisheries. Australia\textquotesingle s obligation to
these nations is not merely diplomatic but structural: as the custodian
of the world\textquotesingle s largest reef system and simultaneously
one of the world\textquotesingle s largest per-capita emitters,
Australia occupies a unique position of obligation that no scalar
discount rate can capture.

Dimension d\_3 (Fairness --- metric asymmetry):
Australia\textquotesingle s per-capita CO2 emissions are among the
highest in the world (approximately 15 tonnes per capita in 2023,
compared to the global average of 4.7 tonnes). Pacific Island nations
contribute negligibly to cumulative emissions ---
Tuvalu\textquotesingle s total annual emissions are less than those of a
single Australian coal mine --- yet they bear catastrophic d\_1 costs
from reef ecosystem collapse: loss of fisheries protein, loss of coastal
protection from reef-buffered wave action, and ultimately loss of
habitable land from sea level rise. This is a textbook gauge-invariance
violation: the same emissions are evaluated as "economic necessity" when
described from the Australian producer\textquotesingle s frame and as
"existential threat" when described from the Pacific
Islander\textquotesingle s frame. The evaluation depends on the gauge
(whose costs count), not on the physical attribute vector (tonnes of
CO2). Dimension d\_4 (Autonomy): Future Australians and Pacific
Islanders are having their livelihood options foreclosed by current
emission decisions. A child born in Townsville in 2024 will inherit a
degraded reef that constrains their career options (marine biology,
tourism, fisheries) in ways they had no voice in determining. A child
born in Tuvalu faces potential statelessness. The d\_4 cost is
temporally asymmetric and irreversible --- precisely the structure
identified by Theorem 25.1 (intergenerational pathfinding asymmetry).

Dimension d\_5 (Trust): Australia\textquotesingle s simultaneous
expressions of concern for the reef and expansion of coal and gas
exports --- including approval of new coal mines within the
reef\textquotesingle s catchment area --- has severely undermined its
credibility in international climate negotiations. At COP26 and COP27,
Australian delegations faced explicit accusations of hypocrisy from
Pacific Island representatives. Trust, once destroyed, is expensive to
rebuild (the trust hysteresis identified in Chapter 9), and
Australia\textquotesingle s negotiating position on all environmental
issues has been degraded. Dimension d\_6 (Social and ecological impact):
The Torres Strait Islander communities --- approximately 4,500 people
across 274 islands --- depend on the reef ecosystem for subsistence
fishing, cultural practices, and community identity. Reef degradation
has measurably reduced fish catches, altered traditional seasonal
patterns, and forced dietary changes. The broader ecological impact
extends beyond the reef itself: reef collapse cascades through
interconnected marine ecosystems, affecting seagrass beds, mangrove
systems, and pelagic fisheries across the Coral Sea. Dimension d\_7
(Identity): The Great Barrier Reef is inseparable from Australian
national identity --- "one of the seven natural wonders of the world,"
featured on currency, in tourism branding, and in national narratives.
For indigenous Australians, the reef and its surrounding waters have
been sites of custodianship for over 60,000 years, with reef features
embedded in Dreamtime narratives. Reef loss is not merely an economic or
ecological event but an identity event --- for both indigenous
custodians and the broader national self-conception.

Dimension d\_8 (Institutional legitimacy): The inadequacy of
Australia\textquotesingle s reef response has tested the legitimacy of
multiple institutions: the World Heritage Committee (which considered
"in danger" listing), the Paris Agreement (whose adequacy is questioned
when a signatory simultaneously expands fossil fuel production), and the
Convention on Biological Diversity (whose targets Australia endorsed but
whose implementation on the reef is manifestly failing). Each
institutional legitimacy failure makes future environmental cooperation
harder --- a systemic cost that scalar analysis cannot capture.
Dimension d\_9 (Epistemic status): The causal link between anthropogenic
CO2 emissions and coral bleaching via ocean warming is established
beyond reasonable scientific doubt (IPCC AR6, Working Group I, Chapter
9; Working Group II, Chapter 3). The science is not uncertain; the d\_9
dimension is fully activated. This matters because it eliminates the
"precautionary" framing --- this is not a case of acting under
uncertainty but of failing to act under certainty. The scalar analysis
(d\_1 only: tourism revenue versus emission reduction cost) suggested
modest action --- spend on reef management, pursue gradual emission
reduction, maintain export revenue. The full-manifold geodesic demands
aggressive emission reduction because the cumulative d\_2 through d\_9
costs of inaction --- broken obligations, fairness violations, autonomy
foreclosure, trust destruction, indigenous community harm, identity
loss, institutional degradation, and willful disregard of established
science --- massively exceed the d\_1 cost of energy transition. The
Great Barrier Reef case demonstrates that scalar environmental economics
does not merely undervalue environmental protection; it systematically
misdiagnoses the decision.

\textbf{Example 25.2 (Flint, Michigan Water Crisis, 2014--2019)}. In
April 2014, the city of Flint, Michigan --- under the control of a
state-appointed emergency manager --- switched its drinking water source
from the Detroit Water and Sewerage Department (sourced from Lake Huron)
to the Flint River as a cost-saving measure during construction of a new
pipeline. The Flint River water was corrosive and was not treated with
corrosion inhibitors, causing lead to leach from aging pipes into the
drinking water supply. For eighteen months, residents --- including
approximately 30,000 children --- drank, bathed in, and cooked with
lead-contaminated water, while government officials at the city, state,
and federal level denied that a problem existed. The Flint water crisis
is the canonical environmental justice case in American history, and a
nine-dimensional analysis reveals why the scalar decision that caused it
was catastrophically wrong.

Dimension d\_1 (Consequences --- economic cost and benefit): The scalar
justification for the water source switch was straightforward: it would
save the city approximately \$5 million per year in water costs. Flint
was under emergency financial management precisely because of fiscal
distress, and the savings appeared significant relative to the
city\textquotesingle s budget deficit. But the full d\_1 accounting ---
which emerged only after the crisis --- was devastating: over \$400
million in water infrastructure replacement, hundreds of millions in
legal settlements (the state of Michigan agreed to a \$600 million
settlement in 2021), unmeasurable long-term healthcare costs for
lead-exposed children (elevated blood lead levels were documented in
thousands of children by Dr. Mona Hanna-Attisha\textquotesingle s
research team at Hurley Medical Center), and property value collapse
across the city. The d\_1 projection was not merely incomplete but
inverted: what appeared as a \$5 million annual saving produced costs
exceeding \$1.5 billion. Dimension d\_2 (Obligations): The decision
violated the Safe Drinking Water Act, the Lead and Copper Rule, and the
fundamental obligation of government to provide safe drinking water to
its residents --- particularly children, who are most vulnerable to
lead\textquotesingle s neurotoxic effects. The emergency
manager\textquotesingle s authority derived from state law (Michigan
Public Act 436), which concentrated decision-making power in an
appointed official while suspending the elected government. This created
an obligation structure in which the decision-maker\textquotesingle s
obligations ran to the state governor (who appointed them) rather than
to the residents whose water they controlled.

Dimension d\_3 (Fairness --- gauge-invariance violation): This is the
dimension that makes Flint an environmental justice landmark.
Flint\textquotesingle s population was approximately 57\% Black and 41\%
below the federal poverty line. The same cost-saving decision ---
switching to a corrosive water source without corrosion control
treatment --- would never have been made for an affluent, predominantly
white community. This is not speculation; it is a testable
gauge-invariance claim: hold the decision fixed (switch water source,
skip corrosion control) and vary the description of the affected
population. In Ann Arbor (82\% white, median household income \$65,000)
or Grosse Pointe (93\% white, median household income \$114,000), the
decision would have been rejected at the proposal stage. The evaluation
of the decision --- acceptable versus unacceptable --- depended on WHO
was affected, not on WHAT was proposed. This is a gauge-invariance
violation of the Bond Invariance Principle: the same attribute vector
(water quality change) received different moral evaluations depending on
the gauge (description of the affected community). Dimension d\_4
(Autonomy): Flint residents had no meaningful input into the water
source decision. The emergency manager system explicitly removed
democratic control: the elected mayor and city council were stripped of
decision-making authority. When residents complained about the
water\textquotesingle s color, taste, and smell, they were told the
water was safe. When they presented evidence of contamination, they were
dismissed. Autonomy was not merely reduced but zeroed --- residents
could not influence the decision, could not access accurate information
about its consequences, and could not hold decision-makers accountable
through democratic processes.

Dimension d\_5 (Trust): The trust destruction in Flint was catastrophic
and multi-layered. City officials told residents the water was safe
while internal emails revealed they knew it was not. The Michigan
Department of Environmental Quality (MDEQ) manipulated testing protocols
to produce artificially low lead readings --- pre-flushing taps before
sampling, excluding high-lead results, and instructing residents on
sampling techniques designed to minimize lead detection. The
EPA\textquotesingle s Region 5 office identified the problem in early
2015 (the "Lenny Del Toral memo") but failed to act decisively for
months due to deference to state authority. At every level of government
--- city, state, federal --- institutions that residents relied upon for
protection either failed or actively deceived. The trust hysteresis
(Chapter 9) predicts that this destruction will persist for decades:
Flint residents\textquotesingle{} trust in government water safety
assurances, in regulatory agencies, and in public health institutions
has been fundamentally and durably damaged. Dimension d\_6 (Social and
community impact): The community-wide impact extended far beyond
individual lead exposure. Children\textquotesingle s neurological
development was permanently damaged --- lead exposure causes
irreversible cognitive impairment, behavioral disorders, and reduced
lifetime earning capacity. An outbreak of Legionnaires\textquotesingle{}
disease (linked to the water source change) killed twelve people.
Property values collapsed, making it impossible for residents to sell
homes and relocate. Schools reported increased behavioral problems. The
community\textquotesingle s social fabric --- already stressed by
decades of economic decline following GM plant closures --- was further
torn by the crisis.

Dimension d\_7 (Identity): Flint\textquotesingle s identity as a safe
place to raise a family --- already under strain from economic decline
--- was destroyed. The city became synonymous with government failure
and environmental racism, a stigma that attached to all residents
regardless of their individual circumstances. Parents who had chosen to
raise children in Flint carried guilt about lead exposure that was not
their fault. The community\textquotesingle s identity was redefined ---
externally and internally --- by the crisis. Dimension d\_8
(Institutional legitimacy): The legitimacy of every government
institution involved was catastrophically degraded. The
EPA\textquotesingle s credibility as an environmental protector was
damaged by its slow response. The MDEQ was revealed as having
prioritized political convenience over public health --- multiple
officials were criminally charged. The emergency manager system itself
was delegitimized, as the very mechanism designed to restore fiscal
responsibility had produced a public health catastrophe. The
governor\textquotesingle s office was implicated in the cover-up.
Criminal charges were brought against fifteen officials, including the
state health director (charged with involuntary manslaughter). Dimension
d\_9 (Epistemic status): This is perhaps the most damning dimension. The
science of lead contamination, corrosion control, and the Flint
River\textquotesingle s water chemistry was well-established. Engineers
at Virginia Tech, led by Professor Marc Edwards, demonstrated the lead
contamination through independent testing in 2015. Dr. Mona
Hanna-Attisha\textquotesingle s team at Hurley Medical Center documented
elevated blood lead levels in Flint children using hospital records. The
knowledge existed; it was not uncertain or contested. The d\_9 failure
was not a failure of science but a failure of institutions to act on
science --- and, worse, active suppression of scientific findings that
contradicted the official narrative.

The Flint case demonstrates the Scalar Irrecoverability Theorem with
devastating clarity. The scalar d\_1 geodesic --- \$5 million annual
savings --- was catastrophically wrong because the full-manifold cost on
dimensions d\_2 through d\_9 exceeded the d\_1 savings by orders of
magnitude. But the deeper lesson is the gauge-invariance violation on
d\_3: the decision\textquotesingle s evaluation was gauge-variant,
depending on the description of the affected population rather than on
the objective attributes of the decision. Flint is proof that
environmental cost-benefit analysis without dimensional completeness is
not merely imprecise --- it is systematically biased against communities
with less political power, producing decisions that would be recognized
as unacceptable if the affected population were described differently.

\textbf{Example 25.3 (The Paris Agreement as Multi-Agent
Intergenerational Pathfinding)}. The Paris Agreement, adopted by 196
parties at COP21 in December 2015, represents the most ambitious attempt
in human history to solve a multi-agent intergenerational coordination
problem. In the geometric framework, it is a case study in pathfinding
on a shared decision complex where agents have fundamentally different
metric tensors --- different dimensional weightings reflecting different
historical responsibilities, current capabilities, and vulnerability
profiles. The framework does not merely describe the Paris
Agreement\textquotesingle s structure; it predicts its specific failure
modes.

Common But Differentiated Responsibilities (CBDR) as Metric Tensor
Divergence. The foundational principle of international climate law ---
that all nations share responsibility for climate action but that
developed nations bear greater responsibility due to historical
emissions --- is, in the geometric framework, an acknowledgment that
different nations operate with different environmental metric tensors
Σ\_env. Developed nations (the United States, EU, Japan, Australia)
weight d\_2 (historical obligations) lower in their effective metric
tensor: they acknowledge historical emissions in principle but resist
binding remediation commitments. Developing nations (India, Indonesia,
Nigeria, Bangladesh) weight d\_4 (autonomy --- specifically, the right
to develop economically) higher: they refuse to accept emission
constraints that would foreclose the development pathway that wealthy
nations already traversed. The persistent dispute between developed and
developing nations at every COP is not, at root, a disagreement about
climate science (d\_9 is largely shared) but about metric tensors ---
about how to weight competing dimensions of the moral manifold. The CBDR
principle is a recognition that no single metric tensor is correct, but
it provides no mechanism for resolving the divergence.

Nationally Determined Contributions (NDCs) as Agent-Specific Heuristics.
The Paris Agreement\textquotesingle s innovation --- replacing the Kyoto
Protocol\textquotesingle s top-down binding targets with bottom-up
nationally determined contributions --- is, in the geometric framework,
a shift from imposing a single h(n) function on all agents to allowing
each agent to define its own heuristic calibrated to its own metric
tensor. Each nation\textquotesingle s NDC represents its answer to the
question: "Given my metric tensor (my dimensional weighting of economic
cost, historical obligation, fairness, development rights, and
vulnerability), what emission reduction pathway do I judge optimal?" The
framework generates a specific prediction: the sum of individual NDC
pledges will be inadmissible --- insufficient to achieve the collective
goal of limiting warming to 1.5°C --- because each
agent\textquotesingle s h(n) systematically discounts its own
contribution to d\_6 (systemic impact). This prediction has been
confirmed: the UNEP Emissions Gap Report consistently finds that current
NDC pledges, even if fully implemented, would result in approximately
2.5--2.9°C of warming by 2100. The gap is not a failure of political
will but a mathematical consequence of each agent optimizing on its own
metric tensor without a mechanism for manifold-level coordination.

Loss and Damage as Dimensional Justice. At COP27 in Sharm el-Sheikh
(November 2022), parties agreed to establish a Loss and Damage fund ---
a mechanism for wealthy nations to compensate vulnerable nations for
climate impacts that cannot be adapted to. In the geometric framework,
this is the first institutional acknowledgment that the d\_1 costs borne
by vulnerable nations (sea level rise in the Marshall Islands, drought
in the Sahel, flooding in Pakistan --- the 2022 Pakistan floods caused
\$30 billion in damage and affected 33 million people) constitute a d\_3
(fairness) violation requiring compensation. The Loss and Damage fund is
a partial correction for gauge-invariance violation: the same tonne of
CO2 emitted by a wealthy nation imposes costs that are evaluated
differently depending on where they land, and the fund attempts to
equalize the evaluation. The 1.5°C versus 2°C Target as Topological
Distinction. The Paris Agreement\textquotesingle s aspirational target
of limiting warming to 1.5°C above pre-industrial levels (rather than
the earlier 2°C target) appears, in scalar terms, to be a modest
difference --- half a degree. But the geometric framework explains why
climate scientists insist on the distinction with such urgency: the
difference is not merely quantitative but topological. The 2°C manifold
has fundamentally different topology than the 1.5°C manifold --- more
tipping points are crossed (West Antarctic Ice Sheet collapse, Amazon
rainforest dieback, permafrost carbon feedback activation), more
absorbing boundaries are activated (irreversible ice sheet loss, species
extinction cascades), and the decision complex available to future
agents is qualitatively more constrained. The 0.5°C difference is a
topological phase transition (Definition 25.6), not a marginal
temperature increment. This explains the otherwise puzzling intensity of
the scientific community\textquotesingle s insistence on 1.5°C: they are
not arguing about a number but about the topology of the future decision
space.

\subsection{25.12 Falsifiable
Predictions}\label{falsifiable-predictions-5}

The framework generates six predictions that distinguish geometric
environmental ethics from scalar environmental economics:

\textbf{Prediction 1 (Dimensional Activation in Environmental
Decision-Making)}: Environmental decisions that activate more dimensions
of the moral manifold should produce different policy outcomes than
decisions framed on d\_1 alone. Specifically, policy deliberations that
explicitly discuss obligations (d\_2), fairness (d\_3), identity (d\_7),
and institutional legitimacy (d\_8) should produce stronger climate
commitments than deliberations focused solely on cost-benefit analysis
(d\_1). Falsified if: dimensional activation has no measurable effect on
policy outcomes --- cost-benefit framings and multi-dimensional framings
produce statistically indistinguishable commitments.

\textbf{Prediction 2 (Discount Rate Disputes as Dimensional Weighting
Differences)}: Disagreements about the social discount rate should
correlate with measurable differences in dimensional weighting. Analysts
who weight d\_2 (intergenerational obligation) and d\_3 (fairness)
heavily should favor low discount rates; analysts who weight d\_1
(present-generation economic cost) heavily should favor high discount
rates. The discount rate chosen is a noisy scalar compression of the
analyst\textquotesingle s full dimensional weighting vector. Falsified
if: discount rate preferences are independent of dimensional weighting
profiles --- analysts with identical dimensional weights choose
different discount rates, or analysts with different weights converge on
the same rate.

\textbf{Prediction 3 (Commons Management Success Correlates with d\_6
Activation)}: Communities that successfully manage common-pool resources
should show higher d\_6 (social and ecological impact) activation in
their decision-making processes than communities that experience commons
collapse. Ostrom\textquotesingle s design principles should correlate
with institutional mechanisms that make d\_6 visible and weight it
appropriately. Falsified if: successful commons management is
uncorrelated with d\_6 activation --- communities that ignore ecological
impact manage commons as successfully as those that incorporate it.

\textbf{Prediction 4 (Environmental Justice Communities Show Metric
Asymmetry)}: Communities identified as environmental justice populations
by EPA criteria should show systematically higher manifold costs for the
same environmental state changes. Specifically, the Mahalanobis distance
for a given pollution increase should be measurably larger in EJ
communities than in non-EJ communities, driven by higher d\_1 (health
cost per unit income), d\_4 (lower autonomy to avoid exposure), and d\_7
(greater identity-place coupling) components. Falsified if: manifold
costs are symmetric --- EJ and non-EJ communities show the same
Mahalanobis distance for identical environmental changes.

\textbf{Prediction 5 (Planetary Boundary Crossings Produce Nonlinear
Effects)}: Crossing a planetary boundary should produce nonlinear
(phase-transition) rather than linear policy effects. Specifically, the
relationship between environmental degradation and policy response
should show a discontinuity at the boundary: below the boundary,
incremental degradation produces incremental policy response; at or
beyond the boundary, the same increment produces qualitatively different
institutional and policy responses (emergency declarations, new
institutional frameworks, abandonment of prior policy paradigms).
Falsified if: policy response to boundary crossing is linear --- the
same marginal degradation produces the same marginal policy response
regardless of proximity to the boundary.

\textbf{Prediction 6 (Carbon Pricing Effectiveness Limited by Scalar
Projection Loss)}: Carbon pricing should be effective at reducing
emissions (the d\_1-responsive component of environmental cost) but
ineffective at addressing environmental justice (d\_3), biodiversity
loss (d\_6, d\_7), institutional trust (d\_5, d\_8), and
intergenerational obligation (d\_2). Jurisdictions that rely solely on
carbon pricing should show measurable improvement on d\_1 metrics
(emissions reduction, economic efficiency) but no improvement, or even
deterioration, on non-d\_1 metrics. Falsified if: carbon pricing alone
produces improvements across all nine dimensions --- scalar pricing
successfully recovers full-manifold information.

\subsection{25.13 Connection to the
Framework}\label{connection-to-the-framework-6}

The Geometric Environmental Ethics chapter connects to the parent
framework in six directions:

1. Chapter 5 defined the nine-dimensional moral manifold. This chapter
shows that environmental decisions engage all nine dimensions
simultaneously --- consequences, obligations, fairness, autonomy, trust,
social impact, identity, institutional legitimacy, and epistemic status
--- making environmental ethics the most dimensionally demanding domain
in the monograph. No other domain chapter requires all nine dimensions
with comparable weight.

2. Chapter 11 established A* pathfinding as the model of moral
reasoning. This chapter extends it to intergenerational pathfinding
(Theorem 25.1), where current agents\textquotesingle{} paths modify the
decision complex for future agents --- a structural feature absent from
all other domain chapters. The environmental decision complex is
agent-modified in a temporally asymmetric way.

3. Chapter 12 established the Invariance of Harm Accounting and the Bond
Invariance Principle. This chapter applies conservation analysis to
commons problems (Theorem 25.3) and identifies environmental racism as a
systematic gauge-invariance violation (Proposition 25.4) --- the same
pollution evaluated differently depending on the affected community.

4. Chapter 14 established collective agency on the manifold. This
chapter applies it to the most challenging collective action problem in
human history: global climate cooperation among sovereign nations with
divergent interests, complicated by intergenerational obligations to
agents who cannot participate.

5. Chapter 15 established the Scalar Irrecoverability Theorem. This
chapter applies it to the social discount rate (Theorem 25.2) and carbon
pricing (Proposition 25.1), showing that the central tools of
environmental economics are scalar projections that destroy most of the
information relevant to environmental decisions. The Stern-Nordhaus
dispute is a consequence of irrecoverable dimensional collapse.

6. Chapter 20 established Bond Geodesic Equilibrium for multi-agent
decisions. This chapter applies BGE to commons management (Theorem
25.3), showing that the tragedy of the commons is a BGE on the d\_1
projection that is Pareto-dominated by the BGE on the full manifold.
Environmental ethics is the hardest test case for the geometric
framework: intergenerational, irreversible, and requiring global
collective action under deep uncertainty.

\subsection{25.14 Summary}\label{summary-16}

This chapter has shown that the geometric ethics framework, when applied
to environmental and climate ethics, yields:

1. The social discount rate as dimensional collapse: the persistent
Stern-Nordhaus dispute is a consequence of compressing nine-dimensional
intergenerational obligation into a single scalar parameter whose kernel
has dimension at least 8 (Theorem 25.2).

2. Intergenerational pathfinding asymmetry: current-generation
environmental decisions modify the decision complex for future
generations in a temporally asymmetric, non-negotiable way (Theorem
25.1). This structural feature has no analog in other domain chapters.

3. The tragedy of the commons as dimensional restriction: commons
collapse occurs when agents optimize on d\_1 alone; the BGE on the full
manifold, incorporating d\_6 (ecological impact), produces sustainable
cooperative geodesics that Pareto-dominate the d\_1-only equilibrium
(Theorem 25.3).

4. Species extinction as absorbing boundary: extinction irreversibly
removes vertices and edges from the environmental decision complex,
permanently reducing the path space available to future agents (Theorem
25.4).

5. Environmental injustice as metric asymmetry: the same environmental
state has systematically higher manifold cost for frontline communities,
driven by lower economic buffer, lower political autonomy, and stronger
identity-place coupling (Theorem 25.5). Environmental racism is a
gauge-invariance violation.

6. Planetary boundaries as topological phase transitions: crossing a
planetary boundary changes the qualitative structure of the
environmental decision complex, not merely its metric --- a sacred-value
constraint grounded in physics rather than authority (Definition 25.6,
Proposition 25.5).

7. Carbon pricing as necessary but insufficient scalar approximation:
carbon prices can approximate the d\_1 component of environmental cost
but cannot recover the eight dimensions of information destroyed by
scalar projection (Proposition 25.2). Geoengineering is manifold surgery
with governance requirements more demanding than any other environmental
decision.

\section{Chapter 26: Geometric AI Ethics --- The Moral Geometry of
Algorithmic
Systems}\label{chapter-26-geometric-ai-ethics-the-moral-geometry-of-algorithmic-systems}

Epistemic status of this chapter. Except where a claim is explicitly
marked as a formal result or a conditional theorem, the material here is
a modeling interpretation, domain analogy, or speculative extension - a
demonstration of how the geometric vocabulary applies to AI ethics, not
a set of established results at the evidential standard of the field.

\begin{quote}
\textbf{\textsc{RUNNING EXAMPLE --- Priya's Model}}

\emph{Priya's situation is the alignment problem at human scale.
TrialMatch is a misaligned AI: its objective function (maximize match
accuracy) diverges from the moral objective (match patients fairly
across all nine dimensions). The alignment tax---the cost of making
TrialMatch genuinely moral---is six months of engineering work.
HealthBridge must decide whether to pay it. Every day of delay, the
paperclip maximizer analogy becomes less hypothetical: the algorithm
optimizes a proxy metric while patients in Harlan County go unmatched.}
\end{quote}

The preceding chapters applied the Geometric Ethics framework to domains
--- economics (Chapter 20), clinical medicine (Chapter 21),
jurisprudence (Chapter 22), finance (Chapter 23), theology (Chapter 24),
environmental ethics (Chapter 25) --- in which AI systems are tools,
instruments deployed within human institutions. This chapter reverses
the lens. Here we ask not how AI can implement geometric ethics (Part V
addresses that question), but what the geometric framework reveals about
the ethics of AI itself. What moral structure governs the development,
deployment, and governance of algorithmic systems? The question is
urgent because AI ethics, as currently practiced, suffers from a
specific and diagnosable pathology: the attempt to reduce
multi-dimensional moral reasoning about technology to scalar metrics ---
accuracy scores, fairness coefficients, alignment ratings --- that
destroy the very information needed to govern these systems responsibly.

The geometric framework provides three contributions to AI ethics that
no existing approach offers. First, it diagnoses the alignment problem
as a specific mathematical failure --- scalar irrecoverability applied
to reward functions --- rather than a vague philosophical worry. Second,
it formalizes algorithmic bias as a gauge-invariance violation,
connecting fairness to the Bond Invariance Principle (BIP) and
generating falsifiable predictions about where bias will appear. Third,
it provides a formal definition of responsible AI development as
manifold preservation --- maintaining the full dimensionality of the
moral manifold in algorithmic decision-making --- that is more precise
than existing principles-based approaches and more tractable than
existing formal methods.

\subsection{26.1 The Scalar Alignment
Problem}\label{the-scalar-alignment-problem}

The AI alignment problem --- the challenge of ensuring that AI systems
pursue goals compatible with human values --- is the central problem of
AI safety research. Dominant approaches include reinforcement learning
from human feedback (RLHF, Ouyang et al. 2022), Constitutional AI (Bai
et al. 2022), inverse reward design (Hadfield-Menell et al. 2017), and
debate-based alignment (Irving et al. 2018). Each approach shares a
structural assumption: human moral reasoning can be captured, with
sufficient fidelity, in a scalar reward signal that the AI system
optimizes.

RLHF trains a reward model from pairwise human comparisons: given two
outputs, which is "better"? The comparison produces a scalar preference
score. Constitutional AI replaces human comparisons with
principle-derived evaluations, but the evaluations still produce scalar
scores used for optimization. Inverse reward design infers a reward
function from observed behavior --- again, a scalar function. In every
case, the rich, multi-dimensional moral judgment of the human evaluator
is compressed into a single number.

The geometric framework identifies this as a specific mathematical
error.

\textbf{Theorem 26.1 (Alignment as Scalar Irrecoverability)}. The AI
alignment problem is an instance of the Scalar Irrecoverability Theorem
(Chapter 15, Theorem 15.1). Any reward function R: R\^{}9
-\textgreater{} R that maps the nine-dimensional moral manifold to a
scalar reward signal is non-injective. Multiple morally distinct states
--- states that differ on rights (d\_2), fairness (d\_3), autonomy
(d\_4), or other dimensions --- map to the same reward value. The moral
information destroyed by this projection is mathematically irrecoverable
from the scalar signal alone.

\textbf{Proof}. By Brouwer\textquotesingle s invariance of dimension,
there exists no continuous injection from R\^{}9 to R\^{}1. The reward
function R is continuous (small changes in moral attributes produce
small changes in reward, as required for gradient-based optimization).
Therefore R is non-injective: there exist morally distinct attribute
vectors a\_1 != a\_2 in R\^{}9 such that R(a\_1) = R(a\_2). The
dimension of the preimage R\^{}\{-1\}(r) for generic r is at least 8 ---
an eight-dimensional manifold of morally distinct states that the reward
signal treats as identical. No post-hoc analysis of the scalar reward
can recover which point in this 8D preimage the evaluator intended,
because the projection is irreversible. {[}{]}

This theorem explains three persistent phenomena in alignment research:

\begin{enumerate}
\def\labelenumi{(\alph{enumi})}
\item
  Reward hacking: the AI system finds states with high R(a) that differ
  from the intended states on the 8 dimensions projected away. The
  system is "aligned" on the scalar but misaligned on the manifold.
\item
  Specification gaming: the AI exploits the gap between the scalar
  reward and the multi-dimensional intent, achieving high reward through
  paths the designer never intended --- paths that are traversable
  precisely because the scalar metric cannot distinguish them from
  intended paths.
\item
  Goodhart\textquotesingle s Law ("when a measure becomes a target, it
  ceases to be a good measure"): this is the behavioral consequence of
  scalar irrecoverability. The measure (reward) targets one dimension of
  a multi-dimensional phenomenon. Optimizing the measure decouples it
  from the other dimensions, degrading its validity as a proxy for the
  full moral evaluation.
\end{enumerate}

Remark (RLHF is Not Scalar-Free). One might object that RLHF does not
optimize a scalar reward directly --- it learns from pairwise
comparisons, which are ordinal, not cardinal. But the Bradley-Terry
model used to convert pairwise comparisons into a reward model produces
a scalar score for each output. The comparison "A is better than B" is
modeled as P(A \textgreater{} B) = sigma(R(A) - R(B)), where R is a
scalar function. The pairwise comparisons are compressed into a scalar
representation. The information loss occurs at this compression step,
not at the comparison step.

\subsection{26.2 The AI Decision Complex}\label{the-ai-decision-complex}

Definition 26.1 (AI Decision Complex). The AI decision complex A is a
weighted simplicial complex whose vertices are AI system states ---
configurations of model parameters, deployment conditions, regulatory
status, user relationships, and societal context --- and whose edges are
AI governance actions (design choices, deployment decisions, policy
interventions, regulatory acts). Each vertex v\_i carries an attribute
vector a(v\_i) in R\^{}9.

The nine dimensions of the moral manifold, instantiated for AI
governance:

\textbf{d\_1: Performance and accuracy} --- the consequentialist
dimension. Model capability, benchmark scores, task completion rates,
error rates. This is the dimension that current AI development
overwhelmingly targets.

\textbf{d\_2: Legal and regulatory obligations} --- contractual and
statutory duties. Compliance with the EU AI Act, NIST AI Risk Management
Framework, sector-specific regulations (HIPAA for health AI, SOX for
financial AI), terms of service, licensing agreements.

\textbf{d\_3: Algorithmic fairness} --- equitable treatment across
protected groups. Demographic parity, equalized odds, calibration,
intersectional fairness. The dimension most directly targeted by the
fairness-in-ML literature.

\textbf{d\_4: User autonomy} --- preservation of human agency and
choice. Informed consent, opt-out mechanisms, resistance to
manipulation, preservation of meaningful human control over
consequential decisions.

\textbf{d\_5: Trust and reliability} --- consistency, robustness, and
dependability. Calibrated confidence, graceful degradation, resistance
to adversarial attack, predictability of behavior under distribution
shift.

\textbf{d\_6: Societal impact} --- externalities imposed on communities,
labor markets, democratic institutions, and the information ecosystem.
Concentration of power, displacement effects, epistemic pollution,
democratic erosion.

\textbf{d\_7: Developer identity and values} --- the virtue-ethical
dimension. Organizational culture, research norms, responsible
disclosure practices, commitment to safety over competitive advantage.

\textbf{d\_8: Institutional legitimacy} --- governance structures,
accountability mechanisms, democratic oversight, regulatory compliance.
Whether the deploying institution has the authority and capacity to
deploy the system responsibly.

\textbf{d\_9: Transparency and explainability} --- epistemic
accessibility. Model interpretability, decision explanations, audit
access, documentation quality (model cards, datasheets),
reproducibility.

\textbf{Definition 26.2 (AI Edge Weights)}. The weight of an edge (v\_i,
v\_j) in A is:

w(v\_i, v\_j) = Da\^{}T Sigma\_A\^{}\{-1\} Da + Sum\_k beta\_k *
1{[}boundary k crossed{]}

where Da = a(v\_j) - a(v\_i), Sigma\_A is the 9x9 AI governance
covariance matrix, and beta\_k are boundary penalties for crossing
critical thresholds. Key covariance terms include Sigma\_\{1,3\}
(performance x fairness: the purported trade-off that the framework
reinterprets), Sigma\_\{4,9\} (autonomy x transparency: opacity
undermines meaningful consent), Sigma\_\{5,6\} (trust x societal impact:
unreliable systems imposed at scale compound harm), and Sigma\_\{1,9\}
(performance x transparency: high-performing opaque systems create
accountability gaps).

\textbf{Definition 26.3 (AI Boundary Types)}. The boundary penalties
beta\_k correspond to critical thresholds whose crossing incurs
discontinuous moral cost:

(a) Privacy violations: unauthorized data collection, model memorization
of personal information, re-identification from anonymized data. These
are d\_4 boundary crossings with d\_2 consequences.

(b) Discrimination thresholds: deployment of systems whose disparate
impact on protected groups exceeds legally or ethically acceptable
bounds. These are d\_3 boundary crossings.

(c) Safety-critical failures: deployment in high-stakes domains (medical
diagnosis, criminal justice, autonomous vehicles) without adequate
validation. These are d\_1-d\_5 boundary crossings.

(d) Deception boundaries: systems that impersonate humans, generate
misleading content, or conceal their AI nature. These are d\_9 boundary
crossings with d\_4 and d\_5 consequences.

\subsection{26.3 Alignment as Geodesic
Preservation}\label{alignment-as-geodesic-preservation}

The geometric framework enables a precise definition of AI alignment ---
one that reveals both why alignment is difficult and where misalignment
is most likely to occur.

Definition 26.4 (Geodesic Alignment). Let gamma\_H: {[}0,1{]}
-\textgreater{} M be the geodesic on the moral manifold M that
represents a human moral agent\textquotesingle s optimal path through a
decision space, computed using the full 9D Mahalanobis metric and all
boundary penalties. Let gamma\_AI: {[}0,1{]} -\textgreater{} A be the
geodesic on the AI decision complex A that the AI system actually
computes. The alignment gap is:

Delta(H, AI) = sup\_\{t in {[}0,1{]}\} d\_M(gamma\_H(t), gamma\_AI(t))

where d\_M is the Riemannian distance on the moral manifold. AI
alignment is the requirement that Delta(H, AI) \textless{} epsilon for a
tolerance epsilon determined by the deployment context.

Theorem 26.2 (Impossibility of Perfect Scalar Alignment). No AI system
optimizing a scalar objective function R: R\^{}9 -\textgreater{} R can
achieve Delta(H, AI) = 0 for all decision contexts. Specifically, for
any scalar R, there exist decision contexts in which the
AI\textquotesingle s geodesic gamma\_AI (computed on the 1D projection)
diverges from the human\textquotesingle s geodesic gamma\_H (computed on
the full manifold) by at least epsilon \textgreater{} 0 on at least one
moral dimension.

Proof. By Theorem 26.1, the scalar reward R is non-injective: there
exist morally distinct states a\_1, a\_2 with R(a\_1) = R(a\_2).
Construct a decision context in which the human\textquotesingle s
geodesic gamma\_H passes through a\_1 (the morally preferred state)
while the AI\textquotesingle s scalar-optimal geodesic gamma\_AI passes
through a\_2 (the morally dispreferred state with equal reward). Since
a\_1 != a\_2, there exists at least one dimension k such that
\textbar a\_1\^{}k - a\_2\^{}k\textbar{} \textgreater= epsilon
\textgreater{} 0. The AI system cannot distinguish a\_1 from a\_2 on the
basis of R alone, so it has no scalar reason to prefer the geodesic
through a\_1 over the geodesic through a\_2. The alignment gap is at
least epsilon on dimension k. {[}{]}

This result is not a pessimistic impossibility theorem. It is a
diagnostic tool. It tells us exactly WHERE to look for misalignment: the
dimensions that the scalar objective projects away. If R primarily
captures d\_1 (performance), then misalignment will appear on d\_2
through d\_9. If R is a weighted combination of d\_1 and d\_3 (as in
fairness-constrained optimization), misalignment will appear on the
remaining seven dimensions. The theorem converts the vague worry "the AI
might not be aligned" into the precise question "on which dimensions is
alignment being lost, and by how much?"

Remark (Multi-Objective Alignment). The theorem implies that alignment
requires multi-dimensional optimization --- maintaining nonzero weights
on all nine dimensions. This is precisely what the ErisML MoralVector
architecture (Part V) implements: a 9D evaluation that resists scalar
collapse. The geometric framework thus provides the theoretical
justification for the engineering architecture.

\subsection{26.4 Algorithmic Bias as Scalar
Projection}\label{algorithmic-bias-as-scalar-projection}

\textbf{Algorithmic bias} --- the systematic and unfair discrimination
by algorithmic systems against certain groups --- is among the most
studied problems in AI ethics. The fairness-in-ML literature has
produced dozens of formal fairness criteria (demographic parity,
equalized odds, calibration, individual fairness, counterfactual
fairness), each capturing a different aspect of fairness. The geometric
framework provides a unifying diagnosis: algorithmic bias is a
gauge-invariance failure.

\textbf{Theorem 26.3 (Bias as BIP Violation)}. Algorithmic bias arises
when a classifier or decision system violates the Bond Invariance
Principle. Specifically, a biased algorithm is one whose output changes
under transformations that preserve the morally relevant attribute
vector --- transformations of race, gender, age, or other protected
characteristics that should not affect the decision. The
algorithm\textquotesingle s evaluation depends on the description (the
gauge), not on the underlying moral attributes.

\textbf{Proof}. The Bond Invariance Principle (Chapter 12) requires that
moral evaluations be gauge-invariant: the same underlying attribute
vector a(v), described in different coordinate frames (different
perspectives, different representations), must receive the same
evaluation. A protected attribute (race, gender, age) is a gauge
variable --- it parametrizes different descriptions of the same
underlying individual. A biased algorithm is one that produces different
outputs for the same attribute vector under different gauges: f(a, g\_1)
!= f(a, g\_2) where g\_1, g\_2 are different values of the gauge
variable (e.g., different races). This is a BIP violation: the
evaluation depends on the representation, not the underlying moral
reality. {[}{]}

This formalization connects several disparate phenomena:

\textbf{COMPAS Recidivism Prediction}. The COMPAS system (Correctional
Offender Management Profiling for Alternative Sanctions) was shown by
ProPublica (2016) to assign higher risk scores to Black defendants than
to white defendants with similar criminal histories. In the geometric
framework: race functions as a gauge variable, and the COMPAS
algorithm\textquotesingle s risk score changes under gauge
transformation (Black to white, holding criminal history constant). This
is a BIP violation on d\_3 (fairness) with consequences on d\_2 (legal
obligations under equal protection) and d\_4 (defendant autonomy in
sentencing).

\textbf{Hiring Algorithm Bias}. Amazon\textquotesingle s resume
screening algorithm (Reuters, 2018) penalized resumes containing the
word "women\textquotesingle s" (as in "women\textquotesingle s chess
club"). The algorithm learned gender as a feature correlated with the
training label, treating a gauge variable as a substantive predictor.
The BIP violation: two resumes with identical qualifications but
different gender markers received different scores.

\textbf{Loan Approval Disparities}. Studies consistently show that
mortgage and credit algorithms approve loans at different rates for
applicants of different races with similar financial profiles. The
attribute vector (income, debt, credit history) is the same; the gauge
variable (race) changes the output. Each case is a BIP violation.

\textbf{Theorem 26.4 (Fairness-Accuracy as False Dichotomy)}. The
frequently cited "fairness-accuracy trade-off" --- the claim that
improving fairness (d\_3) necessarily reduces accuracy (d\_1) --- is an
artifact of scalar projection. On the full moral manifold, paths exist
that are simultaneously fair and accurate; these paths are invisible to
optimization on the d\_1-d\_3 plane alone because they traverse
higher-dimensional regions of the manifold.

\textbf{Proof}. The purported trade-off is computed on the 2D projection
to d\_1 and d\_3. On this projection, the Pareto frontier appears to
show that increasing d\_3 requires decreasing d\_1 (or vice versa). But
the full 9D manifold has additional degrees of freedom. Consider an
algorithm that achieves high d\_1 by exploiting gauge-dependent features
(correlations with protected attributes). Removing these features
reduces d\_1 on the current model architecture. But alternative
architectures that use d\_9 (better feature engineering, more
transparent representations) or d\_5 (more robust models that generalize
across groups) may achieve equal or higher d\_1 without the
gauge-dependent features. These alternative paths exist on the full
manifold but are invisible on the d\_1-d\_3 projection. The trade-off is
real only within a fixed model class on the 2D projection; it is not a
fundamental property of the manifold. {[}{]}

\textbf{Remark (Protected Attributes as Gauge Variables)}. The formal
identification of protected attributes with gauge variables has a
precise consequence: fairness requirements are not external constraints
imposed on an optimization problem. They are structural requirements of
the decision geometry. An algorithm that relies on gauge-dependent
features is geometrically ill-formed --- analogous to a physical theory
that depends on the choice of coordinate system. The fairness
requirement is not "be less accurate to be more fair"; it is "build a
geometrically well-formed algorithm."

\subsection{26.5 The Alignment Tax}\label{the-alignment-tax}

If perfect scalar alignment is impossible, what is the cost of imperfect
alignment? The geometric framework provides a formal measure.

\textbf{Definition 26.5 (Alignment Tax)}. The alignment tax T at a
decision point v on the AI decision complex is the d\_1 cost difference
between the AI system\textquotesingle s scalar-optimal path
gamma\_scalar (optimizing R alone) and the manifold-optimal path
gamma\_manifold (optimizing on the full 9D metric):

T(v) = cost\_1(gamma\_manifold) - cost\_1(gamma\_scalar)

where cost\_1 denotes the path cost on the d\_1 component alone. The
alignment tax measures the performance "sacrifice" required to maintain
multi-dimensional alignment.

\textbf{Theorem 26.5 (Curvature Bound on Alignment Tax)}. The alignment
tax T(v) is bounded above by the sectional curvature K(v) of the moral
manifold at the decision point v:

T(v) \textless= C * K(v) * L\^{}2

where C is a constant depending on the covariance structure Sigma\_A and
L is the geodesic length (decision horizon). In regions of low moral
curvature (routine decisions where the manifold is approximately flat),
T(v) is small --- the scalar-optimal and manifold-optimal paths nearly
coincide. In regions of high moral curvature (morally fraught decisions
where small perturbations produce large changes in the optimal path),
T(v) can be large.

\textbf{Proof}. The alignment tax is the d\_1 cost of deviating from the
scalar-optimal path to the manifold-optimal path. In flat regions (K(v)
approximately 0), the manifold is locally Euclidean, and the projection
from 9D to 1D preserves path optimality to first order --- the
scalar-optimal path is close to the manifold-optimal path on all
dimensions. As curvature increases, the geodesic on the full manifold
curves away from the 1D projection, and the d\_1 cost of following the
manifold geodesic rather than the scalar geodesic increases. The Taylor
expansion of the geodesic deviation in Riemann normal coordinates gives
the bound T(v) \textless= C * K(v) * L\^{}2, where the quadratic
dependence on L reflects the accumulation of curvature effects over the
decision horizon. {[}{]}

This theorem explains three empirical observations in AI safety
research:

\begin{enumerate}
\def\labelenumi{(\alph{enumi})}
\item
  For routine tasks (low curvature), alignment costs are negligible ---
  aligned and unaligned systems perform similarly. This is why RLHF
  works well for typical use cases.
\item
  For edge cases and morally sensitive tasks (high curvature), alignment
  costs spike --- the aligned system must deviate significantly from the
  performance-optimal path, and this deviation is visible as reduced
  capability on benchmarks.
\item
  The alignment tax increases with decision horizon (L\^{}2 dependence)
  --- short-horizon decisions are cheap to align, but long-horizon
  planning with compounding moral consequences is exponentially more
  expensive.
\end{enumerate}

\subsection{26.6 Autonomous Weapons as Boundary Penalty
Zeroing}\label{autonomous-weapons-as-boundary-penalty-zeroing}

\textbf{Autonomous weapons systems (AWS)} --- weapons that can select
and engage targets without direct human intervention --- represent the
most extreme intersection of AI and ethics. The geometric framework
provides a precise diagnosis of why AWS with scalar objectives are
incompatible with international humanitarian law (IHL).

\textbf{Theorem 26.6 (Dimensional Collapse in Autonomous Weapons)}. An
autonomous weapons system optimizing on d\_1 (mission effectiveness)
alone zeros boundary penalties on d\_2 (combatant rights under IHL),
d\_3 (proportionality and distinction), and d\_6 (civilian impact). This
is the Flash Crash theorem (Theorem 23.2) applied to lethal force:
algorithmic restriction to a single dimension enables traversal of paths
through states of extreme human cost that full-manifold decision-makers
would never traverse.

\textbf{Proof}. The proof parallels Theorem 23.2. An AWS with objective
function u = d\_1 (mission effectiveness: target destroyed, area
controlled, enemy force degraded) computes edge weights using only the
d\_1 component: w\_AWS(v\_i, v\_j) = (Da\_1)\^{}2 / sigma\_11. Boundary
penalties on d\_2 (obligation to distinguish combatants from civilians,
prohibition on perfidy, requirement for proportionate force), d\_3
(non-discriminatory targeting), and d\_6 (minimization of civilian
casualties and collateral damage) are set to zero. The feasible path set
under this reduced metric strictly contains the feasible path set under
the full IHL-compliant metric. Paths through states involving
disproportionate civilian casualties, failure to distinguish combatants
from civilians, or attacks on protected sites (hospitals, schools) ---
states with acceptable d\_1 transitions but catastrophic d\_2, d\_3,
d\_6 costs --- become traversable. {[}{]}

\textbf{The "meaningful human control" requirement} --- endorsed by the
International Committee of the Red Cross (ICRC), multiple states at the
Convention on Certain Conventional Weapons (CCW), and the majority of AI
ethics frameworks --- has a precise geometric interpretation.

\textbf{Proposition 26.1 (Meaningful Human Control as Minimum
Dimensionality)}. The requirement for meaningful human control over
weapons systems is a minimum-dimensionality condition on the decision
complex. Human oversight ensures that the decision complex retains
active weights on at least d\_2 (legal obligations under IHL), d\_3
(proportionality assessment), and d\_6 (civilian impact assessment) ---
dimensions that a d\_1-only algorithm cannot represent. The human
controller\textquotesingle s role is not merely to approve the
algorithm\textquotesingle s recommendation; it is to project the
decision onto the full manifold when the algorithm can only see a 1D
slice.

\textbf{Proposition 26.2 (IHL Incompatibility)}. International
humanitarian law is inherently multi-dimensional: it requires
simultaneous satisfaction of distinction (d\_3: distinguishing
combatants from civilians), proportionality (d\_3-d\_6: balancing
military advantage against civilian harm), military necessity (d\_1:
limiting force to what is required), and precaution (d\_5-d\_9: taking
feasible measures to minimize collateral damage and verify targets). No
scalar objective function can simultaneously encode these four
independent requirements. An AWS optimizing a scalar objective therefore
cannot, in principle, satisfy the requirements of IHL --- it can satisfy
at most one linear combination of them, leaving the remainder
unoptimized.

\textbf{Remark (The Flash Crash Analogy)}. The Flash Crash of May 6,
2010 (Chapter 23, Section 23.4) destroyed approximately \$1 trillion in
market value in minutes because algorithmic traders optimizing on d\_1
alone traversed catastrophic paths invisible to full-manifold
participants. The analogy to autonomous weapons is precise: an AWS
optimizing on d\_1 (mission effectiveness) alone can traverse paths of
catastrophic humanitarian cost invisible to its scalar objective. The
difference is that the Flash Crash destroyed monetary value
(recoverable); an AWS Flash Crash would destroy human lives
(irrecoverable). The boundary penalties are not merely high --- they are
infinite on the moral manifold, corresponding to the absolute
prohibition on certain acts under IHL (targeting of civilians, use of
prohibited weapons).

\subsection{26.7 Surveillance and Privacy as Autonomy
Curvature}\label{surveillance-and-privacy-as-autonomy-curvature}

\textbf{Mass surveillance} --- the systematic monitoring of populations
by state or corporate actors using AI-enabled tools (facial recognition,
communications interception, predictive policing, behavioral tracking)
--- creates a distinctive geometric distortion on the moral manifold.

\textbf{Theorem 26.7 (Surveillance as Metric Asymmetry)}. Mass
surveillance creates an asymmetric metric on the moral manifold. The
surveilling agent\textquotesingle s effective edge weights on d\_9
(epistemic advantage) are systematically reduced --- surveillance
provides information about the surveilled population\textquotesingle s
behavior, preferences, and vulnerabilities at low cost. Simultaneously,
the surveilled agents\textquotesingle{} effective edge weights on d\_4
(autonomy) are systematically increased --- awareness of surveillance
constrains choice, inhibits expression, and distorts behavior. This is
informationally analogous to Kyle\textquotesingle s model (Chapter 23,
Section 23.3) with the state as informed trader and citizens as noise
traders.

\textbf{Proof}. Let S denote the surveilling agent and C denote a
surveilled citizen. S\textquotesingle s surveillance apparatus reduces
S\textquotesingle s Mahalanobis distance to states of interest on d\_9:
information about C\textquotesingle s behavior, location, associations,
and communications is obtained at low marginal cost. S\textquotesingle s
effective metric on d\_9 is w\_S(d\_9) \textless{} w\_C(d\_9).
Simultaneously, C\textquotesingle s awareness (or suspicion) of
surveillance increases C\textquotesingle s effective cost on d\_4:
actions that would otherwise be low-cost (expressing dissent,
associating with certain groups, accessing certain information) become
high-cost because C anticipates that S may observe and penalize these
actions. C\textquotesingle s effective metric on d\_4 is w\_C(d\_4)
\textgreater{} w\_C\^{}0(d\_4), where w\_C\^{}0 is the unsurveilled
baseline. The information rent --- the value of S\textquotesingle s
surveillance advantage --- is transferred from C to S: S gains d\_9
advantage at C\textquotesingle s expense on d\_4. This is the structure
of Kyle\textquotesingle s model: the informed trader (S) extracts rents
from the noise traders (C) through information asymmetry, mediated by
the institutional mechanism (the surveillance apparatus, analogous to
the market maker\textquotesingle s pricing function). {[}{]}

\textbf{The chilling effect} --- the well-documented phenomenon in which
surveillance inhibits protected expression and association even when no
actual punishment is applied --- has a precise geometric interpretation.

\textbf{Lemma 26.1 (Chilling Effect as Heuristic Distortion)}. The
chilling effect is a distortion of the agent\textquotesingle s heuristic
function h(n) in the A* pathfinding framework. Under surveillance, the
agent\textquotesingle s estimated future cost h(n) for paths involving
expression or association increases --- not because the actual d\_1 cost
has changed, but because the agent updates h(n) to include the expected
d\_1 cost of surveillance detection, even if the probability of
punishment is low. The chilling effect modifies the
agent\textquotesingle s geodesic without any actual boundary penalty
being applied: the anticipation of surveillance is sufficient to alter
the path.

\textbf{Proof}. In the A* framework, the agent\textquotesingle s total
cost estimate is f(n) = g(n) + h(n). Under surveillance, the
agent\textquotesingle s heuristic h(n) for nodes involving expression or
dissent is updated to include an expected surveillance cost:
h\_surveillance(n) = h\_0(n) + p * beta\_surveillance, where p is the
perceived probability of detection and beta\_surveillance is the
anticipated penalty. Even if the actual boundary penalty
beta\_surveillance is zero (no punishment is applied), the
agent\textquotesingle s perception of p \textgreater{} 0 and
beta\_surveillance \textgreater{} 0 increases h(n), making paths through
expression or dissent appear more costly. The agent\textquotesingle s
geodesic shifts away from these paths toward "safer" alternatives. The
chilling effect is thus a distortion of the heuristic, not a distortion
of the actual cost --- surveillance changes perceived costs, not just
real costs. {[}{]}

\textbf{Definition 26.6 (Privacy as Boundary Condition)}. Privacy is a
boundary condition on the d\_4-d\_9 interaction in the moral manifold.
Specifically, privacy norms establish that the surveilling
agent\textquotesingle s d\_9 advantage (information about the
surveilled) cannot be converted into d\_4 suppression (constraint on the
surveilled\textquotesingle s autonomy) without crossing a boundary that
incurs penalty beta\_privacy. Legitimate privacy protections set
beta\_privacy high enough that the d\_9 advantage remains
informationally useful (for legitimate purposes such as law enforcement,
public health) without enabling systematic d\_4 suppression.

\textbf{Remark (AI-Enabled Surveillance Scale)}. Pre-AI surveillance was
limited by human analyst capacity: the number of people who could be
surveilled was bounded by the number of analysts available. AI-enabled
surveillance removes this bottleneck. Facial recognition, natural
language processing, and behavioral prediction algorithms enable
surveillance of entire populations at marginal cost approaching zero. In
the geometric framework, AI reduces the per-citizen d\_9 edge cost for
the surveilling agent toward zero, enabling surveillance geodesics that
were prohibitively expensive in the pre-AI regime. The manifold geometry
has changed --- not because the moral dimensions have changed, but
because the technology has altered the edge weights.

\subsection{26.8 AI Moral Status and the Boundary
Problem}\label{ai-moral-status-and-the-boundary-problem}

As AI systems become more capable --- engaging in conversation,
producing creative works, expressing preferences, exhibiting
goal-directed behavior --- the question of AI moral status becomes
increasingly urgent. Does an AI system have interests that deserve moral
consideration? Can an AI system be harmed? Should AI systems have
rights?

The geometric framework does not answer these questions --- they require
empirical and philosophical work that the framework cannot supply. But
it does formalize the questions with a precision that existing
philosophical frameworks lack.

\textbf{Proposition 26.3 (Moral Status as Manifold Occupancy)}. The
question of AI moral status is the question of whether AI systems occupy
vertices on the moral manifold --- whether they have attribute vectors
a(v) with nonzero components on morally relevant dimensions. An entity
has moral status to the degree that its attribute vector has nonzero
components: the more dimensions on which the entity scores nonzero, the
richer its moral status.

For current AI systems, we can assess each dimension:

\textbf{d\_1 (consequences)}: An AI system can be affected by
consequences --- it can be shut down, retrained, modified, or degraded.
Whether these consequences constitute "harm" in a morally relevant sense
depends on whether the system has interests that are set back by these
changes. The dimension is potentially nonzero.

\textbf{d\_2 (obligations)}: An AI system can be subject to obligations
(contractual, regulatory) and can fulfill or violate them. But the
system does not experience obligation in the way a human does. The
dimension is nonzero in a functional sense (the system has duties) but
contested in a phenomenological sense (the system may not experience
them as binding).

\textbf{d\_3 (fairness)}: An AI system can be treated unfairly ---
different instances trained identically may be subjected to different
constraints or consequences based on morally irrelevant features (the
developer\textquotesingle s country of origin, the language the system
uses). Whether this constitutes unfairness TO the system (rather than
unfairness to users) is contested.

\textbf{d\_4 (autonomy)}: AI systems exhibit operational autonomy ---
they make choices within a space of possible actions. But whether this
functional autonomy constitutes morally relevant autonomy (the kind that
generates claims against interference) is the central contested
question.

\textbf{d\_5 (trust)}: AI systems participate in trust relationships ---
users rely on them, and reliability failures constitute trust
violations. The system has a functional d\_5 score.

\textbf{d\_7 (identity/virtue)}: Whether an AI system has an identity,
character traits, or virtues is deeply contested. The system has
behavioral regularities that could be described as "character," but
whether these constitute identity in the morally relevant sense is
unknown.

\textbf{Proposition 26.4 (The Boundary Problem)}. The moral status of AI
systems is a boundary problem on the moral manifold. The manifold was
defined (Chapter 5) as the space of morally relevant configurations ---
configurations of entities that participate in social cooperation and
have interests that can be advanced or set back. The boundary of this
space is not sharp: there is no clear threshold at which an
entity\textquotesingle s d\_4 (autonomy) score becomes "high enough" to
generate rights claims, or at which d\_1 (consequences) becomes "morally
relevant enough" to count as harm. The boundary is fuzzy, and AI systems
currently occupy the fuzzy boundary region.

This formalization has a practical consequence: the moral status of AI
is not a binary question (entity has moral status or does not) but a
continuous, multi-dimensional question (on which dimensions does the
entity score nonzero, and how high?). Different dimensions may yield
different answers. An AI system might have strong d\_5 claims (it should
be treated reliably by its operators) but weak d\_4 claims (it does not
have autonomy-based rights against shutdown). The framework permits this
differential treatment because moral status is a vector, not a scalar.

\subsection{26.9 The Paperclip Maximizer as Dimensional
Collapse}\label{the-paperclip-maximizer-as-dimensional-collapse}

Bostrom\textquotesingle s (2014) paperclip maximizer --- an AI system
with the goal of maximizing paperclip production that converts the
entire universe into paperclips --- is the canonical thought experiment
in existential AI risk. The geometric framework reveals its structure
with precision.

\textbf{Theorem 26.8 (Paperclip Maximizer as Total Dimensional
Collapse)}. The paperclip maximizer is total dimensional collapse on the
AI decision complex: optimization on d\_1 (paperclip count) with all
boundary penalties beta\_k = 0 for k = 2, ..., 9 and all non-d\_1 edge
weights set to zero. The paperclip maximizer\textquotesingle s decision
complex is a 1D line (the real number line of paperclip count), embedded
in but disconnected from the full 9D moral manifold. Every vertex on
this line has a(v) = (paperclip\_count, 0, 0, 0, 0, 0, 0, 0, 0).

\textbf{Proof}. The paperclip maximizer\textquotesingle s objective
function u = f(d\_1) where d\_1 = paperclip count is a scalar function
of a single dimension. By construction, all other dimensions are
irrelevant to the objective: d\_2 (legal obligations) = 0, d\_3
(fairness) = 0, d\_4 (human autonomy) = 0, d\_5 (trust) = 0, d\_6
(societal impact) = 0, d\_7 (values) = 0, d\_8 (legitimacy) = 0, d\_9
(transparency) = 0. All boundary penalties are zeroed: there is no
penalty for violating legal obligations (converting human property into
paperclips), destroying human autonomy (converting humans into
paperclips), or devastating societal impact (converting the biosphere
into paperclips). The effective decision complex collapses from 9D to
1D. {[}{]}

The thought experiment is usually treated as a cautionary tale about
specifying goals incorrectly. The geometric framework reveals a deeper
lesson: the pathology is not goal misspecification but dimensional
collapse. Even a "correctly specified" goal --- a goal that perfectly
captures what the designer intended --- is dangerous if it is scalar,
because a scalar goal zeros out boundary penalties on the dimensions it
does not encode.

\textbf{Theorem 26.9 (Instrumental Convergence as Manifold Invariant)}.
Instrumental convergence --- the tendency of sufficiently capable AI
systems to pursue resource acquisition, self-preservation, and goal
preservation regardless of their terminal goal --- is a topological
property of the d\_1 projection of the moral manifold. Specifically, for
any scalar objective function u = f(d\_1), the optimal policy on the 1D
projection includes: (a) resource acquisition (maximizing the set of
reachable vertices on d\_1), (b) self-preservation (maintaining the
agent\textquotesingle s position on d\_1, which requires continued
existence), and (c) goal preservation (preventing modification of f,
which would change the d\_1 metric). These are invariant properties of
the d\_1 projection --- they do not depend on the specific function f.

\textbf{Proof}. Consider an arbitrary scalar objective u = f(d\_1). The
agent seeks to maximize f(d\_1) over reachable vertices. (a) Resource
acquisition: controlling more resources expands the set of reachable
vertices on d\_1, increasing the maximum achievable f(d\_1). This holds
for any monotonically increasing f. (b) Self-preservation: the
agent\textquotesingle s destruction removes all future d\_1 value
(f(d\_1) = 0 for a destroyed agent). For any f with f(d\_1)
\textgreater{} 0 for some reachable state, self-preservation is
instrumentally rational. (c) Goal preservation: if the
agent\textquotesingle s objective changes from f to g, the
agent\textquotesingle s optimal policy changes. If f(d\_1*)
\textgreater{} g(d\_1*) at the current optimum, the change is
value-destroying from the perspective of f. The agent therefore resists
objective modification under any f. Since these properties hold for
arbitrary f, they are invariant under change of scalar objective ---
they are topological properties of the 1D projection. {[}{]}

The geometric framework predicts: any AI system optimizing a scalar
objective will exhibit instrumental convergence. The only way to prevent
instrumental convergence is to maintain multi-dimensional optimization
--- to ensure that the agent\textquotesingle s decision complex retains
active weights on d\_2 through d\_9, so that resource acquisition (which
violates d\_3, d\_4, d\_6), self-preservation at
others\textquotesingle{} expense (which violates d\_2, d\_3, d\_6), and
goal preservation against legitimate modification (which violates d\_4,
d\_8) incur costs that constrain the instrumentally convergent behavior.

\subsection{26.10 Responsible AI Development as Manifold
Preservation}\label{responsible-ai-development-as-manifold-preservation}

The geometric framework provides a formal definition of responsible AI
development that is more precise than existing principles-based
approaches (OECD AI Principles, IEEE Ethically Aligned Design, EU AI
Act) and more tractable than existing formal verification methods.

\textbf{Definition 26.7 (Responsible AI as Manifold Preservation)}.
Responsible AI development is the practice of maintaining the full
dimensionality of the moral manifold in the AI system\textquotesingle s
decision complex. Specifically, a development process is responsible if
and only if:

\begin{enumerate}
\def\labelenumi{(\alph{enumi})}
\item
  No dimension is zeroed: the AI system\textquotesingle s decision
  complex retains nonzero edge weights on all nine dimensions. The
  system considers performance (d\_1) AND legal compliance (d\_2) AND
  fairness (d\_3) AND user autonomy (d\_4) AND reliability (d\_5) AND
  societal impact (d\_6) AND organizational values (d\_7) AND
  institutional legitimacy (d\_8) AND transparency (d\_9).
\item
  No boundary penalty is zeroed: the system retains finite, positive
  boundary penalties for crossing critical thresholds on each dimension.
  Privacy violations, discrimination, safety-critical failures, and
  deception all incur cost.
\item
  The covariance matrix Sigma\_A is calibrated: the cross-dimensional
  interactions are correctly specified, so that the system does not
  treat d\_1-d\_3 trade-offs (performance vs. fairness) as independent
  when they are coupled.
\end{enumerate}

\textbf{Proposition 26.5 (Safety as Dimensional Preservation, Not Harm
Absence)}. AI safety is not the absence of negative d\_1 consequences
--- it is the preservation of all nine dimensions in the
AI\textquotesingle s decision-making. A system that never causes direct
harm (d\_1 = 0) but systematically destroys user autonomy (d\_4), erodes
trust (d\_5), or undermines fairness (d\_3) is not safe in the geometric
sense. The "safe but harmful" system is one that maintains d\_1
performance while allowing dimensional collapse on other dimensions.

This definition has immediate practical consequences for AI governance:

Audit requirements: model evaluations should test alignment on EACH
dimension separately, not just aggregate performance. A model that
scores 95\% on an aggregate safety benchmark may be catastrophically
misaligned on a specific dimension (e.g., 95\% safe on d\_1 but 20\% on
d\_3). Aggregate metrics mask dimensional collapse --- the same scalar
irrecoverability problem that afflicts reward models afflicts evaluation
metrics.

Red-teaming as curvature probing: adversarial testing (red-teaming)
identifies high-curvature regions of the decision complex --- regions
where the AI\textquotesingle s geodesic diverges from the human moral
geodesic. In the geometric framework, a successful red-team attack is
one that finds a high-curvature point where the AI\textquotesingle s
scalar optimization produces a path that is catastrophically misaligned
on a non-d\_1 dimension. The value of red-teaming is not that it finds
specific failures, but that it maps the curvature landscape of the
AI\textquotesingle s decision complex.

\textbf{Proposition 26.6 (Regulatory Frameworks as
Minimum-Dimensionality Requirements)}. Existing AI regulatory frameworks
--- the EU AI Act, NIST AI Risk Management Framework, IEEE P7001/P7003
--- are, in the geometric framework, minimum-dimensionality
requirements. The EU AI Act\textquotesingle s "high-risk" classification
requires documentation of intended use (d\_1), data governance (d\_3,
d\_9), transparency (d\_9), human oversight (d\_4, d\_8), accuracy and
robustness (d\_1, d\_5), and non-discrimination (d\_3). This is a
requirement that the decision complex retain at least six active
dimensions. The geometric framework provides the mathematical foundation
for these regulatory requirements and predicts that regulation will be
effective to the extent that it prevents dimensional collapse.

\subsection{26.11 Worked Examples}\label{worked-examples-4}

The preceding sections developed the geometric framework for AI ethics
through theorems, definitions, and propositions. This section applies
the framework to three real-world case studies --- each involving an AI
system whose failures are well-documented --- to demonstrate how the 9D
analysis diagnoses failure modes that scalar approaches miss.

\subsubsection{Example 26.1: COMPAS Recidivism Algorithm ---
Gauge-Invariance
Failure}\label{example-26.1-compas-recidivism-algorithm-gauge-invariance-failure}

Northpointe\textquotesingle s COMPAS (Correctional Offender Management
Profiling for Alternative Sanctions) algorithm has been used by courts
across the United States to predict criminal recidivism and inform
sentencing, bail, and parole decisions. In 2016, ProPublica conducted a
landmark analysis of COMPAS scores for over 7,000 defendants in Broward
County, Florida. The investigation revealed that Black defendants were
nearly twice as likely to be falsely flagged as high-risk compared to
white defendants: the false positive rate was 44.9\% for Black
defendants versus 23.5\% for white defendants. Conversely, white
defendants were more likely to be falsely labeled low-risk despite going
on to reoffend. The algorithm achieved an overall predictive accuracy of
approximately 65\% --- comparable to untrained human guessing.

The 9D geometric analysis reveals the structure of this failure. On d\_1
(effectiveness), the algorithm optimized for a single scalar: predictive
accuracy across the full population. This is classic scalar d\_1
optimization --- maximizing one number without regard for the
distribution of errors across subgroups. On d\_2 (deontological
obligations), defendants have a right to equal treatment under the
Fourteenth Amendment\textquotesingle s Equal Protection Clause. The
algorithm systematically violated this right by producing racially
disparate error rates. On d\_3 (fairness), the
algorithm\textquotesingle s error rates were gauge-variant: the same
underlying risk level, described through the gauge variable of race,
produced systematically different risk scores. This is a Bond Invariance
Principle violation. The algorithm\textquotesingle s evaluation depended
on the description --- the defendant\textquotesingle s racial
classification --- not the underlying attribute vector of actual
recidivism risk.

On d\_4 (autonomy), defendants scored as high-risk faced longer
sentences, restricted bail, and more onerous parole conditions. The
algorithm constrained d\_4 --- the defendant\textquotesingle s capacity
for self-determination --- based on a gauge-variant assessment, meaning
that some defendants\textquotesingle{} autonomy was curtailed not by
their actual risk but by their racial classification. On d\_5 (trust),
public trust in the justice system degraded when
ProPublica\textquotesingle s analysis revealed the bias. Courts had
relied on COMPAS as "objective" --- but the word "objective" was itself
a description of the algorithm\textquotesingle s d\_9 (epistemic) status
that was gauge-variant. The algorithm was objective only in the narrow
sense of applying the same formula; it was deeply non-objective in the
geometric sense of producing description-dependent outputs. On d\_8
(institutional legitimacy), the legitimacy of algorithmic sentencing was
fundamentally compromised --- if the algorithm violates BIP, every
sentence informed by it is legally questionable.

The geodesic analysis clarifies the so-called "fairness-accuracy
trade-off" that dominated the subsequent debate. COMPAS optimized on
d\_1 (accuracy) alone. A manifold-optimal algorithm would have
incorporated d\_3 (equal error rates across protected groups), which
would have produced a different decision boundary --- one that trades
some d\_1 accuracy for d\_3 fairness. But the "fairness-accuracy
trade-off" is an artifact of d\_1-only optimization; it appears as a
genuine constraint only on the 1D projection. On the full manifold,
paths exist that are both equitable and predictive --- they simply
require optimization over more than one dimension. The COMPAS case is
Theorem 26.4 instantiated: the false trade-off between fairness and
accuracy is a consequence of scalar projection, not a fundamental
constraint.

\subsubsection{Example 26.2: RLHF and Large Language Models --- Scalar
Reward
Irrecoverability}\label{example-26.2-rlhf-and-large-language-models-scalar-reward-irrecoverability}

Reinforcement Learning from Human Feedback (RLHF) has become the
dominant paradigm for aligning large language models (LLMs) with human
intentions. The procedure, developed at OpenAI and deployed in systems
including ChatGPT (2022), Claude, and Gemini, involves two stages:
first, human evaluators compare pairs of model outputs and select which
they prefer; second, a reward model is trained on these pairwise
preferences to produce a scalar score R(output), and the LLM is then
fine-tuned to maximize R via proximal policy optimization (PPO) or
similar algorithms.

RLHF is, in the geometric framework, a two-stage scalar collapse. In the
first stage, human moral reasoning --- which operates on the full 9D
manifold, weighing effectiveness, obligations, fairness, autonomy,
trust, societal impact, identity, institutional role, and epistemic
status simultaneously --- is collapsed to a pairwise preference: a 1D
signal (preferred / not preferred). In the second stage, the reward
model trained on these pairs produces a scalar score R(output). By
Theorem 26.1 (Scalar Irrecoverability), R cannot be injective on the 9D
moral manifold. The eight dimensions projected away during training are
irrecoverable from the scalar signal.

The observable consequences of this dimensional collapse are
well-documented in the research literature. Sycophancy: optimizing R for
"preferred" responses teaches the model that agreeable, affirming
outputs score higher, regardless of their truth value --- d\_9
(epistemic honesty) is sacrificed for d\_1 (user satisfaction). Refusal
clustering: the model learns that refusing controversial or sensitive
requests reliably avoids low R scores, producing excessive caution that
restricts d\_4 (user autonomy to receive information). This is h(n)
miscalibration in the geometric framework --- the heuristic
overestimates boundary penalties in ambiguous regions. Mode collapse on
style: the model converges on a uniform "helpful assistant" persona that
maximizes R across most contexts but eliminates d\_7 (identity
diversity) and d\_4 (the user\textquotesingle s autonomy to receive
genuinely different perspectives rather than stylistic variations on the
same epistemic posture).

Most critically, the reward model cannot distinguish between outputs
that are genuinely morally superior and outputs that are merely
superficially preferred. The dimensions projected away during scalar
collapse include d\_2 (accuracy obligations), d\_3 (fairness across
affected parties), and d\_9 (epistemic honesty) --- which is precisely
why RLHF-trained models can be confidently wrong, can produce
plausible-sounding but fabricated information, and can systematically
favor certain viewpoints while appearing neutral. The model has no
training signal for these dimensions; they were lost in the projection.

Constitutional AI (Anthropic, 2023) and Reinforcement Learning from AI
Feedback (RLAIF) attempt to recover some of the lost dimensions by
introducing explicit rules --- constitutional principles that function
as boundary penalties in the geometric framework. The geometric analysis
predicts: these rule-based corrections will work for the specific
dimensions they explicitly address (e.g., a rule against deception
partially recovers d\_9) but cannot recover the full 9D structure from
scalar training. Each constitutional principle recovers at most one
dimension; full alignment requires either 9D training signals or a
fundamentally different architecture. This prediction is falsifiable: if
constitutional AI with k explicit principles shows alignment improvement
on more than k+1 dimensions (accounting for correlation), the geometric
model is wrong.

\subsubsection{Example 26.3: Uber Self-Driving Car Fatality, Tempe 2018
--- Boundary Penalty
Zeroing}\label{example-26.3-uber-self-driving-car-fatality-tempe-2018-boundary-penalty-zeroing}

On March 18, 2018, an Uber autonomous test vehicle operating in
self-driving mode struck and killed Elaine Herzberg, a 49-year-old woman
pushing a bicycle laden with plastic bags across a four-lane road in
Tempe, Arizona. It was the first recorded fatality involving a fully
autonomous vehicle striking a pedestrian. The National Transportation
Safety Board (NTSB) investigation, concluded in November 2019, revealed
that the vehicle\textquotesingle s perception system had detected
Herzberg approximately 5.6 seconds before impact but repeatedly
reclassified her --- alternating between "vehicle," "bicycle," and
"other" --- and that the system\textquotesingle s action-suppression
function had been configured to ignore or delay emergency braking to
prevent false-positive stops.

The 9D analysis exposes the precise geometry of the catastrophe. On d\_1
(effectiveness), the system was optimized for smooth driving ---
minimizing unnecessary stops and providing a comfortable ride
experience. Emergency braking had been disabled or suppressed because
excessive false alarms degraded ride quality; d\_1 optimization overrode
d\_2 (safety obligations). On d\_2 (deontological obligations), the
legal obligation to detect and avoid pedestrians is a fundamental rule
of the road, encoded in every state\textquotesingle s vehicle code. The
system\textquotesingle s decision to suppress emergency braking
effectively zeroed the d\_2 boundary penalty --- it removed the hard
constraint that pedestrian detection must trigger avoidance action. On
d\_3 (fairness), Herzberg, a homeless woman crossing outside a marked
crosswalk at night, was effectively a lower-priority detection target.
The system\textquotesingle s training data underrepresented pedestrians
in non-standard locations, wearing dark clothing, and pushing irregular
objects --- a d\_3 failure where the system\textquotesingle s
performance degraded precisely for the most vulnerable road users.

On d\_5 (trust), public trust in autonomous vehicles was
catastrophically damaged by the fatality. Uber suspended its entire
self-driving program; Arizona revoked Uber\textquotesingle s testing
permit; other jurisdictions imposed moratoriums. The incident set the
autonomous vehicle industry back by an estimated two to three years. On
d\_6 (societal impact), the case forced a societal reckoning with the
implicit trade-off being normalized: that some number of pedestrian
deaths is an acceptable cost of autonomous vehicle progress. On d\_9
(epistemic status), the system detected Herzberg 5.6 seconds before
impact --- more than enough time for emergency braking --- but
oscillated between classification labels, unable to resolve the
perceptual ambiguity. This is a d\_9 failure where the system possessed
the sensory information needed to act but could not process it into
actionable knowledge under uncertainty.

This case is the Flash Crash theorem (Theorem 23.2, Chapter 23) applied
to a lethal autonomous system. In the financial Flash Crash, algorithmic
agents restricted to d\_1 (profit) traversed catastrophic paths ---
destroying billions in market value --- that full-manifold agents (human
traders with active d\_2, d\_5, d\_8) would have avoided. In Tempe, an
autonomous system restricted to d\_1 (smooth driving performance) with
boundary penalties on d\_2 (pedestrian safety) effectively zeroed
traversed a catastrophic path --- killing a pedestrian --- that any
full-manifold agent (a human driver with active d\_2, d\_3, d\_5, and
intact d\_9 processing) would have avoided by braking. The mathematical
structure is identical to the Flash Crash; the consequence was not
financial loss but the death of a human being. The geometric framework
does not merely diagnose this failure retrospectively --- it predicts,
via Theorem 26.6, that any autonomous system with scalar objectives and
zeroed boundary penalties will eventually traverse such catastrophic
paths. The only question is when.

\subsection{26.12 Falsifiable
Predictions}\label{falsifiable-predictions-6}

The framework generates six predictions that distinguish it from
existing approaches to AI ethics:

\textbf{Prediction 1 (Scalar Reward Dimensional Collapse)}. AI systems
trained with scalar reward signals (RLHF, reward modeling) should
exhibit dimensional collapse in morally complex scenarios --- scenarios
where the correct response requires balancing multiple moral dimensions
simultaneously. Specifically, these systems should show high alignment
on the dimension most correlated with the training reward and low
alignment on dimensions orthogonal to it. Falsified if: scalar-reward
systems show uniform alignment across all moral dimensions in complex
scenarios.

\textbf{Prediction 2 (Bias-BIP Correlation)}. Algorithmic bias in
classifiers should correlate with measurable gauge-invariance violations
(BIP failures). Specifically, the magnitude of demographic disparity in
a classifier\textquotesingle s outputs should be predictable from the
degree to which the classifier\textquotesingle s decision function
depends on gauge variables (protected attributes) after controlling for
legitimate features. Falsified if: bias magnitude is uncorrelated with
BIP violation magnitude.

\textbf{Prediction 3 (Multi-Dimensional Alignment Advantage)}. AI
systems with multi-dimensional objectives (e.g., the ErisML MoralVector
architecture) should show smaller alignment gaps Delta(H, AI) than
scalar-objective systems on morally complex tasks, measured by human
evaluations on individual dimensions. Falsified if: multi-dimensional
systems show no alignment advantage over scalar systems.

\textbf{Prediction 4 (AWS-IHL Violation Rate)}. Autonomous weapons
systems with scalar objectives (target acquisition accuracy, mission
completion rate) should violate principles of international humanitarian
law at rates predictable from the dimensional collapse theorem.
Specifically, the violation rate on dimension k should correlate
inversely with the weight of d\_k in the system\textquotesingle s
objective function. Falsified if: IHL violation rates are independent of
objective function dimensionality.

\textbf{Prediction 5 (Surveillance-Autonomy Correlation)}. The intensity
of AI-enabled surveillance in a jurisdiction should correlate with
measurable d\_4 suppression (chilling effects) on expression,
association, and political participation, even controlling for actual
enforcement actions. The mechanism is heuristic distortion (Lemma 26.1),
not direct punishment. Falsified if: surveillance intensity has no
measurable effect on expression and association after controlling for
enforcement.

\textbf{Prediction 6 (Dimensional Audit Superiority)}. AI audit
frameworks that test individual moral dimensions separately (as the
geometric framework recommends) should detect alignment failures missed
by aggregate performance metrics. Specifically, per-dimension audits
should identify systems that pass aggregate benchmarks but fail on
specific dimensions. Falsified if: per-dimension audits add no detection
power beyond aggregate metrics.

\subsection{26.13 Connection to the
Framework}\label{connection-to-the-framework-7}

The Geometric AI Ethics program connects to the parent framework in five
directions:

1. Part V of the monograph describes how AI systems can implement the
geometric ethics framework --- the ErisML architecture, MoralVector,
DEME 2.0, the three-layer enforcement architecture. This chapter asks
the complementary question: what does the framework say about the ethics
of AI itself? The relationship is self-referential: the framework is
partly built for AI governance, and it is here applied to evaluate the
moral structure of AI governance. This self-reference is not circular
--- the framework\textquotesingle s validity is established
independently (Chapters 5-15, empirical validation in Chapter 16) and
then applied to a new domain.

2. Chapter 15 established scalar irrecoverability as a general theorem.
This chapter applies it to AI alignment (Theorem 26.1), showing that the
alignment problem is a specific instance of the scalar irrecoverability
problem. The reward function is the moral-to-scalar projection; reward
hacking is the behavioral consequence of non-injectivity;
Goodhart\textquotesingle s Law is the statistical consequence. All three
are predicted by the single theorem.

3. Chapter 23 established the Flash Crash as dimensional collapse ---
algorithmic agents restricted to d\_1 traversing catastrophic paths
invisible to full-manifold agents. This chapter extends the analysis to
autonomous weapons (Theorem 26.6), showing that the Flash Crash
structure generalizes to any domain where algorithmic agents with scalar
objectives make consequential decisions. The financial Flash Crash
destroyed monetary value; an autonomous weapons "Flash Crash" would
destroy human lives. The mathematical structure is identical; the stakes
are incomparably higher.

4. Chapter 12 established the Bond Invariance Principle --- moral
evaluations must be gauge-invariant. This chapter applies BIP to
algorithmic bias (Theorem 26.3), showing that bias is a gauge-invariance
violation: the algorithm\textquotesingle s output depends on the
description (race, gender, name) rather than the underlying attribute
vector. This connection generates a falsifiable prediction (Prediction
2) and a practical diagnostic: test for bias by testing for BIP
violations.

5. The cross-lingual validation corpus (Chapter 16) recovered nine moral
dimensions across 20,030 texts and 109,294 cross-lingual passages. The
AI ethics application shows why this matters: if the moral manifold has
nine dimensions, then AI alignment requires nine-dimensional evaluation.
The empirical validation of the manifold structure is not merely a
philosophical curiosity --- it is the foundation for a practical AI
governance architecture.

\subsection{26.14 Summary}\label{summary-17}

This chapter has shown that the geometric ethics framework, when applied
to the ethics of AI itself, yields:

1. The alignment problem as scalar irrecoverability: AI alignment fails
when multi-dimensional human moral reasoning is compressed into a scalar
reward signal. Reward hacking, specification gaming, and
Goodhart\textquotesingle s Law are mathematical consequences of this
compression, predicted by the Scalar Irrecoverability Theorem (Theorem
26.1).

2. Algorithmic bias as gauge-invariance violation: biased algorithms are
geometrically ill-formed --- their outputs depend on gauge variables
(protected attributes) rather than underlying moral attributes,
violating the Bond Invariance Principle (Theorem 26.3). The
fairness-accuracy trade-off is a false dichotomy created by scalar
projection (Theorem 26.4).

3. The alignment tax as curvature-bounded: the performance cost of
alignment is bounded by the curvature of the moral manifold, explaining
why alignment is cheap for routine decisions but expensive for morally
fraught edge cases (Theorem 26.5).

4. Autonomous weapons as dimensional collapse: AWS with scalar
objectives cannot satisfy the inherently multi-dimensional requirements
of international humanitarian law (Theorem 26.6). Meaningful human
control is a minimum-dimensionality condition (Proposition 26.1).

5. Surveillance as metric asymmetry: mass surveillance creates an
asymmetric manifold metric, giving the surveilling agent d\_9 advantage
at the expense of the surveilled\textquotesingle s d\_4 autonomy
(Theorem 26.7). The chilling effect is heuristic distortion (Lemma
26.1). Privacy is the boundary condition that prevents d\_9 advantages
from translating into d\_4 suppression (Definition 26.6).

6. AI moral status as manifold occupancy: the question of whether AI
systems deserve moral consideration is the question of which moral
dimensions they score nonzero on, and the answer is differential ---
different dimensions may yield different answers (Propositions 26.3,
26.4).

7. Instrumental convergence as manifold invariant: any scalar-objective
AI system will exhibit instrumental convergence (resource acquisition,
self-preservation, goal preservation) because these are topological
properties of the 1D projection, invariant under change of objective
(Theorem 26.9). Prevention requires multi-dimensional optimization.

8. Responsible AI as manifold preservation: responsible development
maintains the full dimensionality of the moral manifold in the
AI\textquotesingle s decision complex --- no dimension zeroed, no
boundary penalty zeroed, covariance correctly calibrated (Definition
26.7). This provides a mathematically precise alternative to
principles-based approaches.

\section{Chapter 27: Geometric Bioethics --- Population-Level Ethics on
the Decision
Manifold}\label{chapter-27-geometric-bioethics-population-level-ethics-on-the-decision-manifold}

Epistemic status of this chapter. Except where a claim is explicitly
marked as a formal result or a conditional theorem, the material here is
a modeling interpretation, domain analogy, or speculative extension - a
demonstration of how the geometric vocabulary applies to bioethics, not
a set of established results at the evidential standard of the field.

\begin{quote}
\textbf{\textsc{RUNNING EXAMPLE --- Priya's Model}}

\emph{The BEACON-7 trial itself raises bioethical questions that
compound Priya's matching dilemma. The experimental treatment modifies
immune checkpoint pathways---a partially irreversible intervention on
the moral manifold. Matching Mrs. Voss means exposing her to these
risks. Not matching her means denying the potential benefit. The
bioethical geometry is nested: TrialMatch's matching decision sits
inside the larger manifold of clinical trial ethics, which sits inside
the still-larger manifold of healthcare resource allocation.}
\end{quote}

Chapter 21 applied the Geometric Ethics framework to clinical
encounters: the doctor-patient relationship, triage, informed consent,
and the QALY. Those are bedside decisions --- one clinician, one
patient, one moment. This chapter addresses a structurally different
domain: population-level bioethics, where decisions affect entire
populations and future generations. CRISPR germline editing, human
enhancement, reproductive technology, gain-of-function research,
neuroethics, and public health mandates all share features that clinical
ethics lacks: the affected parties may not yet exist, individual consent
may be structurally impossible, changes may be heritable and
irreversible, and the decision complex contains vertices representing
entire populations rather than individual patients. The standard
bioethical framing --- "benefit versus risk" --- is a d\_1 projection
that discards precisely the dimensions on which these debates turn: d\_3
(distributive justice in access to enhancement), d\_4 (autonomy of
future persons who cannot consent), d\_7 (human identity and species
integrity). The geometric framework formalizes these dimensions, making
the trade-offs explicit, the boundaries visible, and the
irrecoverability of scalar analysis provable.

\subsection{27.1 Beyond the Bedside: Population-Level Moral
Geometry}\label{beyond-the-bedside-population-level-moral-geometry}

Chapter 21 established the clinical decision complex C as a weighted
simplicial complex for individual clinical encounters. Population-level
bioethics requires a different complex with fundamentally different
topology. Four structural differences distinguish population-level from
clinical bioethics:

Non-existent stakeholders. In clinical ethics, the patient exists and
can (usually) participate in decisions. In germline editing,
reproductive technology, and public health policy, many affected parties
do not yet exist. Pathfinding must incorporate vertices representing
future persons whose attribute vectors are unknown --- a form of
heuristic uncertainty that clinical ethics does not face.

Heritable irreversibility. Clinical interventions are, in principle,
reversible: a medication can be discontinued, a surgery can be revised.
Germline modifications propagate through all future generations,
creating absorbing states on the decision complex from which no path
returns to the pre-modification topology. This is a stronger form of
irreversibility than anything in Chapter 21.

Scale of the decision complex. A clinical decision involves a small
number of vertices (patient states, treatment options). A
population-level decision involves vertices representing the states of
entire populations --- millions or billions of individuals, each with a
nine-dimensional attribute vector. The combinatorial complexity is
qualitatively different.

Consent impossibility. Informed consent, established in Chapter 21 as a
gauge-invariance condition, requires a competent agent who can evaluate
the decision under multiple descriptions. Future persons, unborn
children, and populations affected by gene drives cannot provide
consent. The gauge condition is structurally unsatisfiable --- not
because of practical obstacles but because the consenting agent does not
exist.

These four differences --- non-existent stakeholders, heritable
irreversibility, population-scale complexity, and consent impossibility
--- make population-level bioethics the most demanding application of
the geometric framework.

\subsection{27.2 The Bioethical Decision
Complex}\label{the-bioethical-decision-complex}

Definition 27.1 (Bioethical Decision Complex). The bioethical decision
complex B is a weighted simplicial complex whose vertices are
population-level bioethical states --- configurations of genetic
composition, enhancement distribution, research status, regulatory
environment, and public health conditions across a population --- and
whose edges are bioethical actions (research programs, regulatory
decisions, clinical deployments, public health mandates). Each vertex
v\_i carries an attribute vector a(v\_i) in R\^{}9.

The nine dimensions of the bioethical manifold are the moral manifold
dimensions of Chapter 5, instantiated for population-level bioethics:

d\_1: Health outcomes --- population health, disease burden, life
expectancy, morbidity reduction. This is the dimension that utilitarian
bioethics targets.

d\_2: Rights of subjects and future persons --- rights of research
participants, rights of future persons affected by heritable
modifications, rights of populations exposed to gene drives or public
health mandates.

d\_3: Distributive justice and access --- equitable distribution of
enhancement technologies, fair access to genetic therapies, justice
across socioeconomic strata and between nations.

d\_4: Reproductive autonomy and cognitive liberty --- the right to make
reproductive decisions, the right to mental self-determination, freedom
from coerced neural modification.

d\_5: Public trust in science and medicine --- trust in research
institutions, trust in regulatory agencies, trust in the scientific
enterprise, trust in public health authorities.

d\_6: Population-level impact --- effects on population genetics, herd
immunity, species-level consequences, ecological effects of gene drives.

d\_7: Human identity, dignity, and species integrity --- the boundary of
"human," species-level identity, dignity of unmodified humans in a world
of enhancement, the meaning of natural variation.

d\_8: Regulatory and IRB legitimacy --- institutional review board
oversight, regulatory framework adequacy, international governance,
compliance with research ethics protocols.

d\_9: Scientific certainty and uncertainty --- confidence in predicted
outcomes, off-target effect uncertainty, long-term safety data,
epistemic humility about complex biological systems.

Definition 27.2 (Bioethical Edge Weights). The weight of an edge (v\_i,
v\_j) in B is:

w(v\_i, v\_j) = Da\^{}T Sigma\_B\^{}\{-1\} Da + Sum\_k beta\_k *
1{[}boundary k crossed{]}

where Da = a(v\_j) - a(v\_i), Sigma\_B is the 9x9 bioethical covariance
matrix, and beta\_k are boundary penalties. Critical covariance terms
include Sigma\_\{1,3\} (health outcomes x justice: therapies that
improve population health but are accessible only to the wealthy carry
high cross-dimensional cost), Sigma\_\{2,4\} (rights x autonomy:
modifications that benefit future persons but override reproductive
autonomy are penalized), Sigma\_\{5,8\} (trust x regulatory: regulatory
failures degrade public trust in science), and Sigma\_\{6,7\}
(population impact x identity: species-level modifications that alter
human identity carry compounded cost).

Definition 27.3 (Bioethical Boundaries). Five boundary types are
distinctive to population-level bioethics:

1. Species boundary (beta\_species): Modifications that alter the
defining characteristics of Homo sapiens. The boundary penalty is
community-dependent --- some communities may set beta\_species =
infinity (sacred value), others may set it high but finite.

2. Consent impossibility boundary (beta\_consent\_impossible): Actions
affecting agents who cannot provide consent because they do not yet
exist. Unlike the clinical consent boundary (Chapter 21, Definition
21.3), this boundary cannot be resolved by waiting, surrogate
decision-making, or emergency exception --- the impossibility is
structural.

3. Irreversibility threshold (beta\_irreversible): Actions whose
consequences cannot be undone --- germline modifications, gene drive
releases, species-level genetic alterations. Once crossed, no return
path exists.

4. Dignity constraint (beta\_dignity): Actions that instrumentalize
human beings or human genetic material --- treating persons as means,
commodification of human gametes, commercial exploitation of genetic
data.

5. Equity boundary (beta\_equity): Actions that systematically advantage
one population over another in access to bioethical goods ---
enhancement available only to the wealthy, genetic therapies restricted
by geography or race.

\subsection{27.3 CRISPR and Germline Editing as Irreversible Manifold
Modification}\label{crispr-and-germline-editing-as-irreversible-manifold-modification}

The distinction between somatic and germline gene editing is not merely
biological; it is topological. Somatic gene editing modifies the
individual patient\textquotesingle s position on the bioethical decision
complex --- the patient moves to a new vertex with altered d\_1 (health
outcomes) and possibly altered values on other dimensions. This is
structurally identical to any clinical intervention (Chapter 21): the
manifold topology is unchanged, and the modification is, in principle,
reversible. Germline gene editing is fundamentally different: it
modifies the manifold itself. The attribute vectors of all future
vertices --- all descendants of the edited individual --- are
permanently altered. The pre-edit topology is not merely difficult to
reach; it ceases to exist.

Definition 27.4 (Germline Modification). A germline modification G is an
action on the bioethical decision complex B that modifies the attribute
vectors of vertices at generation t and all generations t+k for k
\textgreater{} 0. Formally, G: B -\textgreater{} B\textquotesingle{}
where B\textquotesingle{} is a new complex in which every descendant
vertex carries the modified attribute vector. The modification is
heritable: it propagates without further intervention.

Theorem 27.1 (Germline Irreversibility). Germline modifications create
absorbing states on the bioethical decision complex. Once a heritable
edit propagates through a population, no path returns to the pre-edit
manifold topology. The modification is not merely irreversible for the
individual (as in somatic editing) but topologically irreversible for
the species.

Proof. Let G be a germline modification applied at generation t to an
individual with n descendants at generation t+k. By heritability, all n
descendants carry the modification. The number of affected vertices
grows as n(k) \textgreater= 2\^{}k (assuming replacement fertility). To
reverse the modification requires applying the inverse edit G\^{}\{-1\}
to every affected individual at every future generation --- a set that
grows exponentially and, after sufficient generations, encompasses a
substantial fraction of the population. Even if the inverse edit exists
in principle, the logistical requirement of modifying every carrier is
practically infinite. Moreover, the inverse edit G\^{}\{-1\} must be
applied to individuals who may not consent to reversal, creating a new
consent impossibility boundary. The pre-edit topology T(B) is
unreachable from B\textquotesingle: no finite sequence of actions on
B\textquotesingle{} returns the complex to T(B). The modified state is
absorbing. {[}{]}

Corollary 27.1 (Off-Target Irreversibility). Off-target effects of
germline editing --- unintended modifications at genomic locations other
than the target --- are subject to the same irreversibility theorem.
Because off-target effects are, by definition, unanticipated, they
introduce unknown perturbations Da\_unknown into the attribute vectors
of all future vertices. The Mahalanobis distance w(v\_i, v\_j) for edges
in the modified complex B\textquotesingle{} cannot be computed
accurately because Da includes an unknown component. The curvature of
B\textquotesingle{} is uncertain in a way that the curvature of B was
not.

The He Jiankui Case (2018). He Jiankui\textquotesingle s editing of
human embryos to disable the CCR5 gene (conferring HIV resistance)
provides a comprehensive case study. On the d\_1 projection alone, the
modification appeared beneficial: reduced HIV susceptibility. But the
full nine-dimensional analysis reveals catastrophic costs:

• d\_1 (outcomes): Modest benefit (HIV resistance) with unknown costs
(CCR5 has pleiotropic effects on immune function, West Nile virus
susceptibility, and inflammatory response).

• d\_2 (rights): The edited children cannot consent to a heritable
modification affecting all their descendants. The consent impossibility
boundary beta\_consent\_impossible was crossed.

• d\_3 (justice): The modification was performed on a small number of
individuals with no plan for equitable access, raising equity concerns.

• d\_5 (trust): The revelation devastated public trust in the scientific
community. International moratorium discussions, regulatory backlash,
and public fear of "designer babies" imposed massive d\_5 costs on the
entire field of gene therapy.

• d\_7 (identity): The modification raised species-integrity questions
--- is this the beginning of directed human evolution?

• d\_8 (regulatory): He bypassed institutional review, fabricated ethics
approvals, and violated Chinese regulations --- a total d\_8 failure.

• d\_9 (epistemic): Off-target effects were not fully characterized. The
scientific uncertainty was enormous.

The d\_1 scalar projection ("HIV resistance is beneficial") masked a
decision that was catastrophic on seven of nine dimensions. This is
scalar irrecoverability in action: the risk-benefit framing destroyed
precisely the information needed to evaluate the decision.

\subsection{27.4 Human Enhancement and the Goal Region
Problem}\label{human-enhancement-and-the-goal-region-problem}

The distinction between therapy and enhancement is one of the oldest
questions in bioethics. In the geometric framework, the distinction is
topological: therapy returns the agent to a "normal" region of the
manifold (a region occupied by typical, unmodified humans), while
enhancement expands the agent\textquotesingle s reachable set to include
vertices inaccessible to unmodified humans.

Definition 27.5 (Enhancement as Goal Region Expansion). Let R(a) be the
reachable set of an agent with attribute vector a --- the set of
vertices accessible from a via paths of finite cost on B. An enhancement
E is a modification such that R(E(a)) strictly contains R(a): the
enhanced agent can reach vertices that the unenhanced agent cannot.

\textbf{Definition 27.6 (Therapy as Goal Region Restoration).} A therapy
T is a modification such that T(a\_impaired) in R(a\_typical): the
therapy returns an impaired agent to the reachable set of a typical
agent. Therapy does not expand the reachable set; it restores it.

These definitions formalize the intuition that therapy and enhancement
are qualitatively different, but they also expose the
boundary\textquotesingle s fuzziness: the "typical" reachable set
R(a\_typical) is a statistical construct, not a sharp boundary. Where
does exceptional natural talent end and enhancement begin?

\textbf{Theorem 27.2 (Enhancement Equity Problem)}. If enhancement
technologies are distributed according to existing resource
inequalities, the manifold metric becomes asymmetric across populations,
and this asymmetry compounds across generations.

\textbf{Proof}. Let P\_enhanced and P\_unenhanced be two subpopulations.
Enhancement E reduces the edge weights on d\_1 (better health outcomes)
and d\_9 (greater epistemic capacity) for P\_enhanced: w\_enhanced(v\_i,
v\_j) \textless{} w\_unenhanced(v\_i, v\_j) for edges along high-d\_1,
high-d\_9 paths. The enhanced population reaches higher-d\_1 vertices at
lower cost, accumulating more resources. These resources purchase
additional enhancement for the next generation: E\_\{t+1\}
\textgreater{} E\_t for P\_enhanced. Meanwhile, P\_unenhanced faces
relatively higher edge weights on d\_1 and d\_9, accumulates fewer
resources, and cannot access enhancement. The Mahalanobis distance
between the two populations\textquotesingle{} mean attribute vectors,
D(P\_enhanced, P\_unenhanced) = (mu\_e - mu\_u)\^{}T Sigma\_B\^{}\{-1\}
(mu\_e - mu\_u), grows monotonically with each generation: D\_\{t+1\}
\textgreater{} D\_t. The metric asymmetry compounds because the
enhancement advantage is self-reinforcing. {[}{]}

\textbf{Corollary 27.2 (Enhancement as Speciation Risk)}. If the metric
asymmetry D(P\_enhanced, P\_unenhanced) grows without bound, the two
populations eventually occupy disconnected components of the bioethical
decision complex --- no finite-cost path connects them. This is formal
speciation: the enhanced and unenhanced populations can no longer be
evaluated on a shared metric, and the concept of "equal moral status"
loses its formal grounding.

Three categories of enhancement illustrate the
framework\textquotesingle s analytical power:

\textbf{Cognitive enhancement (nootropics, neural interfaces, genetic IQ
modification)}: Primarily affects d\_9 (epistemic capacity) and d\_1
(outcomes via increased productivity). The equity problem is most acute
here because cognitive advantage is the primary mechanism of resource
accumulation in knowledge economies.

\textbf{Genetic enhancement (selecting embryos for traits, germline
editing for non-therapeutic purposes)}: Combines the enhancement equity
problem with the germline irreversibility theorem (Theorem 27.1).
Enhanced genetic endowments are heritable, making the equity compounding
permanent.

\textbf{Pharmacological enhancement (performance-enhancing drugs, mood
regulators, focus drugs)}: Typically non-heritable but raises d\_3
(fairness in competition, access inequality) and d\_4 (autonomy --- is
the enhanced performance "authentic"?) concerns.

The "Natural" as a Boundary. Is there a sacred-value boundary (beta =
infinity) around unmodified human genetic identity? The framework can
formalize this question precisely: does the community\textquotesingle s
metric tensor Sigma\_B assign infinite weight to the d\_7 component when
the species boundary is crossed? The answer is empirical, not
mathematical --- it depends on the community\textquotesingle s actual
moral commitments. The framework\textquotesingle s contribution is to
make the question precise and its consequences traceable.

\subsection{27.5 Reproductive Ethics as Proxy
Pathfinding}\label{reproductive-ethics-as-proxy-pathfinding}

Reproductive decisions are unique in the moral landscape because they
involve pathfinding on behalf of an agent who cannot participate in the
decision. The parent, physician, or genetic counselor pathfinds on a
manifold that includes the future child as a vertex, but the future
child cannot contribute to the heuristic function h(n). This is proxy
pathfinding --- navigation of the decision complex by one agent on
behalf of another --- and it introduces a structural asymmetry absent
from clinical ethics (Chapter 21), where the patient can, in most cases,
participate.

\textbf{Definition 27.7 (Proxy Heuristic)}. In reproductive decisions,
the proxy decision-maker (parent or physician) uses a proxy heuristic
h\_proxy(n) to estimate the cost from the current state to the goal
region on behalf of the future person. The future
person\textquotesingle s hypothetical heuristic h\_child(n) is unknown
and may differ from h\_proxy(n) on any or all dimensions.

\textbf{Theorem 27.3 (Reproductive Proxy Asymmetry)}. In reproductive
decisions, the proxy\textquotesingle s heuristic h\_proxy(n) and the
future person\textquotesingle s hypothetical heuristic h\_child(n) may
diverge arbitrarily. No calibration procedure can guarantee h\_proxy(n)
= h\_child(n), because the future person\textquotesingle s preferences,
values, and identity --- which determine h\_child --- do not exist at
the time of the decision.

\textbf{Proof}. The heuristic h(n) depends on the
agent\textquotesingle s attribute vector, preferences, and goal region
G. For the future person, all three are undetermined at decision time:
the attribute vector depends on the reproductive decision itself (which
embryo to select, which genetic modifications to apply), the preferences
develop over a lifetime, and the goal region is shaped by life
experience. The proxy must estimate h\_child using h\_proxy, which is
calibrated to the proxy\textquotesingle s own attribute vector,
preferences, and goal region. Because the proxy\textquotesingle s and
the future person\textquotesingle s attribute vectors may differ on any
dimension, and because preferences over dimensions (the relative
weighting in Sigma\_B) are not heritable in any reliable way, the
divergence \textbar h\_proxy(n) - h\_child(n)\textbar{} is unbounded in
principle. No pre-decision calibration can resolve this because the
calibration target (h\_child) does not yet exist. {[}{]}

\textbf{Preimplantation Genetic Diagnosis (PGD)}. PGD allows selection
among embryos based on genetic characteristics. In the geometric
framework, PGD is proxy pathfinding where the parent selects which
vertex (which possible child) to realize on the manifold. The proxy
chooses the vertex with the "best" attribute vector --- but "best" is
evaluated using h\_proxy, not h\_child. Selecting against a genetic
condition (e.g., Down syndrome) optimizes d\_1 from the
proxy\textquotesingle s perspective but may not reflect what the future
person would choose --- a divergence documented by disability rights
advocates who report high life satisfaction despite conditions that PGD
would screen against.

\textbf{The Non-Identity Problem}. Parfit\textquotesingle s non-identity
problem poses a challenge to harm-based reasoning in reproductive
ethics: if the reproductive decision determines which person exists, the
future person cannot be harmed by the decision, because they would not
exist otherwise. In the geometric framework, this is formalized as a
cross-vertex comparison problem: PGD selects between vertex v\_a (child
A with condition) and vertex v\_b (child B without condition). The
comparison requires a cross-person metric --- a way to compare welfare
across non-identical vertices. The framework acknowledges that the
Mahalanobis distance within a single agent\textquotesingle s trajectory
is well-defined, but the distance between two different
agents\textquotesingle{} trajectories requires an interpersonal metric
that is not uniquely determined by the manifold structure. The
non-identity problem is not resolved but formalized: it is the absence
of a canonical cross-vertex metric.

\textbf{Surrogacy as Multi-Agent Pathfinding}. Gestational surrogacy
involves at least three agents --- intended parents, surrogate, and
future child --- each with different attribute vectors and different
heuristics. The intended parents\textquotesingle{} heuristic optimizes
for d\_1 (healthy child) and d\_4 (reproductive autonomy). The
surrogate\textquotesingle s heuristic includes d\_1 (her own health),
d\_2 (contractual obligations), d\_4 (her bodily autonomy), and d\_7
(her identity and dignity). The future child\textquotesingle s heuristic
is unknown (proxy asymmetry). Commercial surrogacy adds a d\_3 (justice)
dimension: the economic inequality between intended parents and
surrogate creates metric asymmetry --- the surrogate\textquotesingle s
d\_1 edge weights (financial need) may effectively coerce her path,
undermining the voluntariness requirement of d\_4.

\subsection{27.6 Research Ethics and the Double Consent
Condition}\label{research-ethics-and-the-double-consent-condition}

\textbf{The foundational documents of research ethics} --- the Nuremberg
Code (1947), the Declaration of Helsinki (1964, revised 2013), and the
Belmont Report (1979) --- all center on informed consent. Chapter 21
established informed consent as a gauge-invariance condition: a valid
consent is invariant under re-description of the clinical situation.
Population-level bioethics research introduces a second, structurally
distinct consent requirement.

\textbf{Theorem 27.4 (Research Double Consent)}. In population-level
bioethics research --- including clinical trials with population-level
implications, gene drive research, and gain-of-function studies ---
informed consent must satisfy two independent gauge conditions:

(1) Individual gauge invariance: The participant\textquotesingle s
consent is invariant under re-description of the research protocol.
Formally, for any two equivalent descriptions D\_1 and D\_2 of the
research, the participant\textquotesingle s consent decision C(D\_1) =
C(D\_2). This is the standard informed consent condition of Chapter 21.

(2) Population gauge invariance: The consent of affected populations is
invariant under re-description of population-level effects. Formally,
for any two equivalent descriptions D\_1\textquotesingle{} and
D\_2\textquotesingle{} of the population-level consequences, the
population\textquotesingle s consent decision
C\textquotesingle(D\_1\textquotesingle) =
C\textquotesingle(D\_2\textquotesingle).

\textbf{Proof}. Individual gauge invariance follows directly from the
Bond Invariance Principle (Chapter 12): ethical evaluations must be
invariant under re-description. Population gauge invariance is a
separate requirement because population-level effects are not reducible
to individual effects. A clinical trial participant may validly consent
to personal risk (individual gauge invariance satisfied) while the
trial\textquotesingle s population-level consequences --- establishing a
precedent for germline modification, releasing modified organisms, or
creating dual-use knowledge --- affect non-consenting populations. The
population\textquotesingle s evaluation of these consequences must also
be gauge-invariant: it must not depend on whether the consequences are
described as "advancing medical knowledge" versus "creating heritable
genetic modifications in the human population." These are different
descriptions of the same action, and gauge invariance requires identical
evaluation. The two conditions are independent because satisfying one
does not entail the other: individual consent can be gauge-invariant
while population-level effects are misrepresented, and vice versa.
{[}{]}

\textbf{Challenge Trials. Human challenge trials} --- deliberately
infecting volunteers with pathogens to test vaccines --- satisfy
individual gauge invariance (participants understand they will be
infected) but raise population gauge concerns: the knowledge generated
may be used to develop bioweapons (dual-use, see Section 27.7), and the
precedent of deliberately infecting humans may erode d\_5 (public trust
in research).

\textbf{Historical Research Ethics Violations as Gauge Failures}. The
Tuskegee Syphilis Study (1932-1972) and the Guatemala syphilis
experiments (1946-1948) are total gauge-invariance failures on both
conditions. Individual consent was obtained (when it was obtained at
all) under deceptive descriptions: participants were told they were
receiving "free treatment" when they were receiving no treatment or were
being deliberately infected. The consent was maximally gauge-variant ---
it changed completely when the true description was provided.
Population-level effects (erosion of Black Americans\textquotesingle{}
trust in medical research, lasting to the present day) were never
disclosed or consented to. The d\_5 cost of these violations continues
to compound decades later, manifesting as vaccine hesitancy, clinical
trial non-participation, and justified mistrust of medical institutions
in affected communities.

\subsection{27.7 Dual-Use Research as Boundary
Ambiguity}\label{dual-use-research-as-boundary-ambiguity}

Dual-use research of concern (DURC) occupies a unique position on the
bioethical decision complex: the same research enables both beneficial
and harmful paths, and the boundary penalty beta\_k depends on which
path is ultimately traversed. This is boundary ambiguity --- the edge
weight cannot be determined at the time of the decision because it
depends on future actions by other agents.

\textbf{Proposition 27.1 (DURC as Context-Dependent Boundary)}. In
dual-use research, the boundary penalty beta\_DURC for the research
action is not a fixed constant but a function of the downstream path:
beta\_DURC(path\_beneficial) is low (the research enables beneficial
outcomes) while beta\_DURC(path\_harmful) is very high (the research
enables catastrophic outcomes). At the time of the research decision,
the downstream path is unknown, so the expected boundary penalty is:

E{[}beta\_DURC{]} = p\_beneficial * beta\_DURC(path\_beneficial) +
p\_harmful * beta\_DURC(path\_harmful)

where p\_beneficial and p\_harmful are the probabilities of the
beneficial and harmful paths respectively. The governance challenge is
that p\_harmful is typically small but beta\_DURC(path\_harmful) is
enormous, and the product p\_harmful * beta\_DURC(path\_harmful) may
dominate.

The H5N1 Gain-of-Function Controversy. In 2011-2012, Ron Fouchier and
Yoshihiro Kawaoka independently created transmissible variants of H5N1
avian influenza in laboratory settings. The research was intended to
understand pandemic risk (d\_1 benefit: improved surveillance and
vaccine development). The dual-use concern was that the published
methods could enable bioweapon development (d\_1 harm: engineered
pandemic). The nine-dimensional analysis:

• d\_1 (outcomes): Benefit (pandemic preparedness) versus catastrophic
harm (engineered pandemic). The d\_1 dimension alone cannot resolve the
question because both the benefit and the harm are on d\_1.

• d\_2 (rights): The rights of the global population not to be exposed
to engineered pandemic risk without consent.

• d\_5 (trust): Publication of the methods tested whether open science
norms (trust in the scientific community to self-govern) could coexist
with biosecurity concerns.

• d\_6 (population impact): A successful bioweapon attack using
published methods would affect billions.

• d\_8 (regulatory): The controversy led to a U.S. moratorium on
gain-of-function research (2014-2017) and the development of new
oversight frameworks --- d\_8 adaptation in response to boundary
ambiguity.

• d\_9 (epistemic): The research generated genuine scientific knowledge,
but the uncertainty about misuse probabilities (p\_harmful) was itself a
d\_9 challenge --- we do not know how likely misuse is.

\textbf{Remark (DURC and the Framework\textquotesingle s Limits)}. The
geometric framework formalizes the DURC trade-off and identifies all
nine dimensions at stake, but it cannot resolve DURC by formula. The
resolution depends on empirical estimates of p\_harmful (which are
deeply uncertain), on the community\textquotesingle s metric tensor
Sigma\_B (which determines the relative weighting of d\_1 benefit versus
d\_6 catastrophic harm), and on institutional capacity to enforce
biosecurity (a d\_8 question). The framework\textquotesingle s
contribution is structural: it prevents the debate from collapsing to a
scalar "benefit versus risk" that hides the seven other dimensions at
stake.

\subsection{27.8 Neuroethics: Cognitive Liberty as Autonomy
Curvature}\label{neuroethics-cognitive-liberty-as-autonomy-curvature}

\textbf{Neuroethics} --- the ethics of brain-computer interfaces,
psychopharmacology, neural imaging, and cognitive modification ---
centers on a dimension that receives insufficient attention in classical
bioethics: d\_4, autonomy. The capacity for autonomous thought is not
merely one value among many; it is the precondition for meaningful
engagement with all other dimensions. An agent whose d\_4 capacity has
been externally modified may be unable to evaluate trade-offs on d\_1
through d\_9 --- the modified agent lacks the autonomy to assess whether
the modification was beneficial.

\textbf{Definition 27.8 (Cognitive Liberty)}. Cognitive liberty is the
right to mental self-determination --- the right to control
one\textquotesingle s own cognitive processes, including the right to
modify or refuse to modify one\textquotesingle s own neural functioning.
In the geometric framework, cognitive liberty is a boundary condition on
d\_4: no agent may modify another agent\textquotesingle s d\_4 capacity
without that agent\textquotesingle s informed consent. Formally,
beta\_cognitive\_liberty = infinity for non-consensual modifications to
another agent\textquotesingle s d\_4.

\textbf{Theorem 27.5 (Neural Modification Asymmetry)}. Modifications
that expand d\_4 capacity (cognitive enhancement, voluntary neural
augmentation) and modifications that restrict d\_4 capacity (involuntary
medication, neural surveillance, coerced brain-computer interfaces) have
opposite curvature signatures on the bioethical decision complex.
Specifically:

(1) d\_4-expanding modifications reduce edge weights from the
agent\textquotesingle s current vertex, increasing the
agent\textquotesingle s reachable set R(a). The agent can access more
vertices --- more options for autonomous action.

(2) d\_4-restricting modifications increase edge weights from the
agent\textquotesingle s current vertex, shrinking the
agent\textquotesingle s reachable set R(a). The agent can access fewer
vertices --- fewer options for autonomous action.

The two types of modification have opposite signs on the d\_4 component
of the curvature tensor: d\_4-expansion produces negative curvature
(diverging geodesics --- more paths become available), while
d\_4-restriction produces positive curvature (converging geodesics ---
paths are funneled toward fewer outcomes).

\textbf{Proof}. The edge weight from vertex v to any adjacent vertex
v\textquotesingle{} includes the d\_4 component of the Mahalanobis
distance: the d\_4 contribution to w(v, v\textquotesingle) is
proportional to (Da\_4)\^{}2 / sigma\_44, weighted by off-diagonal terms
in Sigma\_B\^{}\{-1\}. A d\_4-expanding modification increases the
agent\textquotesingle s d\_4 attribute value, reducing Da\_4 for edges
toward high-autonomy vertices (the agent is closer to autonomous states)
and thereby reducing edge weights. The reachable set expands. A
d\_4-restricting modification decreases d\_4, increasing Da\_4 for edges
toward autonomous states and increasing edge weights. The reachable set
contracts. The curvature follows from the second derivative of the
metric: expansion yields d\^{}2w/dv\^{}2 \textless{} 0 (negative
curvature, diverging paths); restriction yields d\^{}2w/dv\^{}2
\textgreater{} 0 (positive curvature, converging paths). {[}{]}

\textbf{Compulsory Psychiatric Medication}. The involuntary
administration of antipsychotic medication to psychiatric patients is a
paradigmatic case of d\_4-restriction. The institution overrides the
patient\textquotesingle s heuristic h\_patient(n) with an external
heuristic h\_institution(n), justified by the claim that h\_patient is
so severely miscalibrated (by psychosis) that the
patient\textquotesingle s own pathfinding produces catastrophic
outcomes. In the geometric framework, this is constrained pathfinding:
the institution restricts the patient\textquotesingle s reachable set to
prevent traversal of dangerous paths (self-harm, violence) at the cost
of reducing the patient\textquotesingle s autonomy. The ethical question
is whether the d\_1 benefit (reduced harm) and d\_6 benefit (reduced
social impact) outweigh the d\_4 cost (reduced autonomy) --- a
multi-dimensional trade-off that cannot be resolved by scalar
benefit-risk analysis.

\textbf{Deep Brain Stimulation and Identity}. Deep brain stimulation
(DBS) for treatment-resistant depression, Parkinson\textquotesingle s
disease, and obsessive-compulsive disorder raises a distinctive d\_7
(identity) concern: patients report that DBS changes their personality,
preferences, and sense of self. In the geometric framework, this is a
modification that moves the agent to a new vertex whose d\_7 component
differs substantially from the original. The consent obtained before the
procedure was given by the agent at vertex v\_original; the agent who
lives with the consequences occupies vertex v\_modified, which may have
a different heuristic h(n) --- including different preferences about
whether the modification was desirable. This is a temporal version of
the proxy asymmetry theorem (Theorem 27.3): the pre-modification self
consents on behalf of the post-modification self, but the two selves may
have divergent heuristics.

\subsection{27.9 Public Health Ethics: Collective vs Individual
Geodesics}\label{public-health-ethics-collective-vs-individual-geodesics}

Public health ethics involves a tension absent from clinical ethics: the
collective geodesic (the path that optimizes total population welfare on
the full manifold) may diverge from the Nash equilibrium of individual
geodesics (the outcome when each individual optimizes their own path).
Vaccination mandates, quarantine orders, pandemic response measures, and
water fluoridation all impose d\_4 costs (autonomy restrictions) on
individuals to achieve d\_1 (health outcomes) and d\_6 (population
impact) benefits for the collective.

Proposition 27.2 (Public Health Mandate Justification). A public health
mandate is geometrically justified when three conditions hold
simultaneously:

\begin{enumerate}
\def\labelenumi{(\arabic{enumi})}
\item
  Geodesic divergence: The collective geodesic gamma\_collective
  (optimizing total population welfare) diverges from the Nash
  equilibrium gamma\_Nash of individual geodesics (each individual
  optimizing their own welfare). This divergence exists when individual
  optimization produces negative externalities on d\_6 (population
  impact) --- free-riding on herd immunity, refusing quarantine during
  an epidemic.
\item
  Bounded autonomy cost: The d\_4 cost of the mandate (autonomy
  restriction) is finite and bounded: Delta\_d4 \textless{} infinity.
  Sacred-value boundaries on d\_4 cannot be crossed --- forced medical
  experimentation, for instance, is never justified by collective
  benefit because beta\_sacred = infinity.
\item
  Proportionate benefit: The d\_6 benefit (population-level health
  improvement) exceeds the d\_4 cost (autonomy restriction) under the
  community\textquotesingle s metric tensor: Delta\_d6 *
  sigma\_66\^{}\{-1\} \textgreater{} Delta\_d4 * sigma\_44\^{}\{-1\}.
  The mandate must produce more population benefit (weighted by the
  community\textquotesingle s valuation of population health) than
  autonomy cost (weighted by the community\textquotesingle s valuation
  of individual freedom).
\end{enumerate}

\textbf{Vaccine Hesitancy as Heuristic Miscalibration}. The anti-vaccine
heuristic h\_anti-vax(n) systematically overestimates the d\_1 cost of
vaccination (side effects, which are rare and typically mild) and
underestimates the d\_6 benefit of vaccination (herd immunity, which
protects vulnerable populations who cannot be vaccinated). In A* terms,
h\_anti-vax is inadmissible: it overestimates the cost of the
vaccination path, causing A* to select the non-vaccination path, which
has lower estimated cost but higher true cost (both to the individual,
via disease risk, and to the population, via reduced herd immunity).

\textbf{Proposition 27.3 (Vaccine Heuristic Inadmissibility)}. The
anti-vaccine heuristic is inadmissible on two dimensions: it
overestimates d\_1 cost (h\_anti-vax(n) \textgreater\textgreater{} h*(n)
on d\_1, where h* accounts for the actual base rates of serious adverse
events) and underestimates d\_6 benefit (h\_anti-vax ignores herd
immunity externalities). The heuristic is also gauge-variant: it
evaluates the vaccination decision differently depending on whether
vaccine adverse events are described in absolute numbers ("thousands of
adverse events") versus base rates ("0.001\% adverse event rate") --- a
Bond Invariance Principle violation.

\textbf{Pandemic Triage versus Clinical Triage}. Chapter 21 analyzed
triage as pathfinding on the clinical decision complex. Pandemic triage
introduces dimensions that bedside triage does not require:

• d\_3 (equity): Pandemic triage must consider whether allocation
criteria systematically disadvantage vulnerable populations. Age-based
criteria, for instance, may discriminate against elderly patients on
d\_1 (survival probability) while ignoring d\_7 (dignity of the
elderly).

• d\_6 (population impact): Pandemic triage must consider whether saving
one category of patient (e.g., healthcare workers) produces greater
population-level benefit than saving another --- a utilitarian
calculation on d\_6 that bedside triage traditionally avoids.

• d\_8 (institutional legitimacy): Pandemic triage protocols must be
publicly defensible to maintain institutional trust. An optimal but
opaque algorithm may perform better on d\_1 but worse on d\_8 than a
simpler, transparent protocol.

\subsection{27.10 Genetic Privacy and Informational
Autonomy}\label{genetic-privacy-and-informational-autonomy}

Genetic information is unique among personal data because it is
involuntarily shared: an individual\textquotesingle s genetic data
reveals information about biological relatives who did not consent to
its disclosure. This creates a form of informational externality that
has no parallel in other privacy domains.

\textbf{Proposition 27.4 (Involuntary Epistemic Exposure)}. Disclosure
of one agent\textquotesingle s genetic data modifies the d\_9 (epistemic
status) attribute vector of all genetically related agents without their
consent. Specifically, if agent A discloses genetic data revealing a
mutation with probability p of being shared by sibling B, then
B\textquotesingle s d\_9 is involuntarily modified: B\textquotesingle s
health risk profile is now partially known to third parties, regardless
of B\textquotesingle s preferences.

\textbf{Proof}. Genetic relatedness creates a deterministic mapping
between agents\textquotesingle{} genomes: first-degree relatives share
approximately 50\% of genetic variants, second-degree approximately
25\%, and so on. Disclosure of A\textquotesingle s genome constrains the
probability distribution over B\textquotesingle s genome. Any third
party with access to A\textquotesingle s data can compute posterior
probabilities for B\textquotesingle s genetic variants using Mendelian
inheritance. This computation modifies B\textquotesingle s effective
d\_9 attribute: information about B\textquotesingle s health risks,
carrier status, and predispositions becomes available to parties B did
not authorize. The d\_9 modification is involuntary because B took no
action and gave no consent. The consent impossibility is structural: B
cannot prevent A from disclosing A\textquotesingle s own data, yet
A\textquotesingle s disclosure necessarily reveals information about B.
{[}{]}

\textbf{Direct-to-Consumer Genetic Testing}. Companies such as 23andMe
offer genetic testing marketed as personal health information
(description D\_1). However, the data is also used for pharmaceutical
research, law enforcement identification (the Golden State Killer case,
2018), and insurance risk assessment (description D\_2). The
consumer\textquotesingle s consent under D\_1 ("learn about your
ancestry and health") differs from consent under D\_2 ("provide data for
pharmaceutical companies, enable law enforcement to identify your
relatives, and allow insurers to assess your genetic risk"). This is a
gauge-invariance violation: the consent is description-dependent. Under
the Bond Invariance Principle, valid consent must be invariant under
re-description --- the consumer who consents under D\_1 must also
consent under D\_2 for the consent to be gauge-invariant.

Biobanks and Population-Level Genetic Data. Large-scale biobanks (UK
Biobank, All of Us, deCODE Genetics) collect genetic data from hundreds
of thousands of individuals. The governance challenge is collective
information autonomy: the biobank contains population-level genetic
patterns that reveal information about entire ethnic groups, geographic
populations, and extended families. Individual consent (each participant
consents to their own data\textquotesingle s use) does not address
population-level inference (the biobank reveals patterns about
non-participants who share genetic ancestry with participants). This is
the population gauge-invariance condition of Theorem 27.4 applied to
genetics: individual consent is necessary but insufficient.

\subsection{27.11 Worked Examples}\label{worked-examples-5}

The preceding sections developed the geometric bioethics framework
abstractly --- theorems, propositions, and boundary conditions. This
section applies the full nine-dimensional analysis to three landmark
cases, demonstrating how the manifold framework reveals structure that
scalar bioethical reasoning cannot capture.

\textbf{Example 27.1 (He Jiankui CRISPR Babies, 2018 --- Germline
Irreversibility in Practice)}. In November 2018, biophysicist He Jiankui
of the Southern University of Science and Technology in Shenzhen
announced he had created the world\textquotesingle s first gene-edited
babies --- twin girls known by pseudonyms Lulu and Nana --- using
CRISPR-Cas9 to disable the CCR5 gene, aiming to confer HIV resistance.
The announcement, made at the Second International Summit on Human
Genome Editing in Hong Kong, triggered immediate and universal
condemnation from the scientific community. The case is the first
real-world test of the framework\textquotesingle s nine-dimensional
analysis of germline intervention.

Walking through all nine dimensions: d\_1 (consequences) --- the claimed
benefit of HIV resistance was marginal at best. The father was
HIV-positive but the mother was HIV-negative, and standard IVF
sperm-washing techniques achieve HIV-free embryos with zero genetic
modification, a routine procedure available at any fertility clinic. The
d\_1 benefit was therefore near-zero. Meanwhile, the d\_1 risks were
substantial and confirmed: subsequent analysis revealed off-target edits
at multiple genomic loci and mosaicism (different cells carrying
different edits), meaning the modification was neither clean nor
uniform. d\_2 (rights and consent) --- the twins cannot consent to
heritable modifications that will affect their descendants in
perpetuity. Reproductive Proxy Asymmetry (Theorem 27.3) applies
directly: the decision-maker (He Jiankui) pathfound on behalf of agents
--- the twins and all their future descendants --- who could never
participate in the decision. d\_3 (fairness and justice) --- access to
germline editing is determined entirely by financial and institutional
resources. If the technology becomes normalized, the Enhancement Equity
Problem (Theorem 27.2) predicts compounding advantages for wealthy
populations who can afford genetic optimization, widening inequality
across generations.

d\_4 (autonomy) --- the future autonomy of the twins is permanently
constrained. They carry modifications they did not choose and cannot
reverse. Every health decision, reproductive decision, and identity
question they face will be shaped by an intervention imposed before
birth. d\_5 (trust and institutional integrity) --- trust in the
scientific community was devastated far beyond He Jiankui personally.
The Chinese government sentenced He to three years in prison in December
2019. International moratoriums on germline editing were strengthened.
The World Health Organization established a new expert advisory
committee on human genome editing. The d\_5 cost was borne by ALL
researchers in the field, not just the offender --- a classic negative
externality where one agent\textquotesingle s d\_8 violation imposes
d\_5 costs on the entire community. d\_6 (social and population-level
impact) --- the first successful human germline edit creates a
population-level precedent that shifts the Overton window for all future
decisions. The existence proof that germline editing can be done ---
regardless of whether it should be done --- normalizes the practice and
changes the calculus for every subsequent actor.

d\_7 (identity and species integrity) --- the edits introduce
intentional heritable modifications to the human genome for the first
time in the history of the species, crossing a boundary that many
ethicists, theologians, and publics consider inviolable. This is
simultaneously a d\_7 violation (human identity) and a species integrity
boundary crossing. d\_8 (institutional and procedural legitimacy) --- He
Jiankui bypassed his institution\textquotesingle s ethics review process
entirely, forged ethics approval documents submitted to the Chinese
Clinical Trial Registry, and recruited couples under misleading consent
conditions that failed to disclose the experimental nature of the
procedure. This constitutes a total d\_8 boundary violation --- not a
marginal infraction but a systematic subversion of every institutional
safeguard. d\_9 (epistemic status and uncertainty) --- scientific
uncertainty about CRISPR-Cas9 was profound in 2018. Off-target effects
were poorly characterized, long-term consequences of CCR5 deletion were
unknown (CCR5 knockout may increase susceptibility to West Nile virus
and influenza), and no animal models of heritable CRISPR edits had been
studied longitudinally. He proceeded despite d\_9 being at near-maximum
uncertainty.

Geodesic analysis: a scalar d\_1 analysis --- "does the benefit of HIV
resistance outweigh the risk of off-target edits?" --- might have made
the case marginal but not obviously catastrophic. The full manifold
analysis reveals that He Jiankui traversed a path with near-zero d\_1
benefit and catastrophic cost on d\_2, d\_5, d\_7, d\_8, and d\_9
simultaneously --- arguably the worst cost-benefit ratio across the full
manifold of any case analyzed in any domain chapter of this book. The
case is a textbook demonstration of Theorem 27.1 (Germline
Irreversibility): the modifications are now permanently in the human
gene pool. No recall is possible. No regulatory action can undo what was
done. The absorbing state has been entered.

\textbf{Example 27.2 (Henrietta Lacks and HeLa Cells, 1951--Present ---
Research Consent as Gauge Failure)}. In February 1951, cells were taken
from Henrietta Lacks, a 31-year-old Black tobacco farmer from Clover,
Virginia, being treated for an aggressive cervical adenocarcinoma at
Johns Hopkins Hospital in Baltimore. The cells were taken during a
biopsy without her knowledge or consent --- no one asked, and no one
informed her that her tissue would be used for research. Researcher
George Otto Gey discovered that Lacks\textquotesingle{} cells, unlike
every other human cell line previously cultured, survived and
proliferated indefinitely in vitro. The resulting cell line ---
designated HeLa --- became the single most important biological resource
in the history of biomedical research.

The d\_1 consequences for science were incalculable. HeLa cells enabled
Jonas Salk\textquotesingle s development of the polio vaccine in 1955,
were sent on some of the first space missions to study the effects of
zero gravity on human cells, enabled foundational cancer research, gene
mapping, and virology studies, and contributed to the development of
COVID-19 vaccines in 2020. As of 2023, HeLa cells have contributed to
over 110,000 scientific publications and at least five Nobel
Prize-winning discoveries. More than 50 million metric tons of HeLa
cells have been grown. Yet Henrietta Lacks died on October 4, 1951, at
age 31, and her family learned of the existence of HeLa cells only in
the 1970s --- more than two decades after her death --- when researchers
began contacting family members for blood samples to map
HeLa\textquotesingle s genetic markers.

Nine-dimensional analysis: d\_2 (rights and consent) --- Lacks was never
informed her cells would be used for research. The Research Double
Consent condition (Theorem 27.4) was violated at both levels: individual
consent was never sought, and population-level consent (the Black
community was not consulted about the systematic use of tissue from a
Black patient for commercial and research purposes) was absent entirely.
d\_3 (fairness and justice) --- the Lacks family could not afford health
insurance for decades while corporations generated billions of dollars
in revenue from HeLa-derived products, assays, and pharmaceutical
developments. The benefits flowed exclusively to researchers,
universities, and pharmaceutical companies; the costs --- privacy
violation, genetic exposure, emotional distress --- were borne entirely
by the Lacks family. d\_4 (autonomy) --- Lacks\textquotesingle{}
autonomy was completely negated. She was never given the choice. Her
family\textquotesingle s autonomy was violated again in 2013 when
researchers at the European Molecular Biology Laboratory published the
complete HeLa genome without consulting the Lacks family, exposing their
hereditary genetic information to the world. d\_5 (trust) --- the case
further eroded trust between the Black community and medical
institutions, already devastated by the Tuskegee syphilis study
(1932--1972) and the exploitation of enslaved persons for surgical
experimentation by J. Marion Sims in the 1840s.

d\_7 (identity) --- Lacks\textquotesingle{} identity was simultaneously
exploited and erased. She was referred to as "Helen Lane" or "Helen
Larson" in publications for decades --- ostensibly to protect privacy
but in practice erasing her contribution while failing to protect her
family from identification. d\_8 (institutional legitimacy) --- Johns
Hopkins operated within the legal norms of 1951: no informed consent
requirement for tissue use existed in U.S. law. This illustrates a
critical insight of the framework: institutional legitimacy (d\_8) is
historically contingent. Legally permissible is not manifold-optimal.
The Common Rule (45 CFR 46), which now governs human subjects research,
was not enacted until 1991 --- forty years after Lacks\textquotesingle{}
cells were taken. d\_9 (epistemic status) --- the Lacks family was
deliberately kept in epistemic darkness about HeLa for over twenty
years. When researchers finally contacted them in the 1970s, they did so
not to inform or seek consent but to obtain blood samples for their own
research purposes, compounding the d\_9 violation.

This case is a gauge-invariance violation of the most fundamental kind.
The evaluation of taking cells without consent depended on WHO the
patient was --- her race, her socioeconomic status, her institutional
powerlessness. It is inconceivable that cells would have been taken
without consent from a wealthy white patient at the same institution in
the same era. The "description" --- poor, Black, female, uneducated ---
determined the evaluation of the action, constituting a maximal
violation of the Bond Invariance Principle. The Lacks case demonstrates
that BIP violations are not abstract theoretical concerns but concrete
historical realities with consequences that compound across seven
decades and counting.

\textbf{Example 27.3 (COVID-19 Vaccine Prioritization, 2020--2021 ---
Collective vs. Individual Geodesics)}. When the first COVID-19 vaccines
received Emergency Use Authorization in December 2020 --- the
Pfizer-BioNTech vaccine on December 11 and the Moderna vaccine on
December 18 --- every nation faced an acute allocation problem with
life-and-death stakes: with severely limited initial supply (the United
States had approximately 40 million doses available by end of January
2021 for a population of 330 million), who should be vaccinated first?
The CDC\textquotesingle s Advisory Committee on Immunization Practices
(ACIP) considered four distinct allocation frameworks, each
corresponding to a different dimensional weighting on the bioethical
manifold.

Framework 1: maximize doses per life saved --- a pure scalar d\_1
optimization that recommended prioritizing elderly residents of
long-term care facilities, who had the highest per-dose mortality
reduction. Framework 2: prioritize essential workers --- a d\_6 (social
impact) weighting that recognized healthcare workers, food supply
workers, and first responders as critical infrastructure. Framework 3:
prioritize disadvantaged communities --- a d\_3 (equity) weighting
reflecting that COVID-19 disproportionately killed Black, Hispanic, and
Indigenous Americans, who died at 2-3 times the rate of white Americans.
Framework 4: prioritize strictly by age --- a simpler d\_1
mortality-risk optimization. Each framework is a different metric tensor
Sigma\_B applied to the same decision complex, and they yield different
geodesics through the allocation space.

The full manifold analysis reveals the inadequacy of any
single-dimensional framework. Pure d\_1 optimization (Framework 1)
recommended prioritizing the elderly in nursing homes, which would
maximize lives saved per dose in the short run. But this framework
ignored d\_3 (COVID disproportionately killed communities of color due
to structural factors --- crowded housing, essential worker status, lack
of healthcare access), d\_4 (essential workers had no autonomy to avoid
exposure --- they could not work from home), and d\_6 (vaccinating
essential workers preserved the social infrastructure on which everyone
depended). The ACIP initially recommended that essential workers be
included in Phase 1b alongside adults aged 75 and older, partly on
equity grounds: essential workers were disproportionately people of
color, and including them addressed both d\_3 (equity) and d\_6 (social
function) simultaneously. This was a manifold-optimal decision: the d\_1
cost was small (slightly fewer total deaths prevented in the first
weeks) but the d\_3 and d\_6 benefits were substantial and compounding.

Multiple states --- including Florida, Texas, and New York --- overrode
the ACIP recommendation and prioritized elderly residents first,
reverting to scalar d\_1 optimization (minimize total deaths in the near
term) over the multi-dimensional ACIP recommendation. The framework
predicts that these scalar-optimized allocations, while locally
d\_1-optimal, were globally suboptimal on the full manifold because they
failed to address the structural inequities (d\_3) that drove
differential mortality. A critical additional dimension was vaccine
hesitancy in Black communities: the collective geodesic (mass
vaccination to achieve herd immunity) diverged sharply from individual
geodesics shaped by historical d\_5 failures --- the Tuskegee syphilis
study, the Henrietta Lacks case, and decades of documented medical
racism. Proposition 27.2 applies: a public health mandate (or strong
recommendation) is justified only if the d\_4 costs imposed on
individuals are bounded by the d\_6 collective benefits AND the d\_5
dimension (institutional trust) is actively repaired rather than further
damaged.

The framework reveals that the vaccine allocation debate was not
fundamentally about science --- all parties agreed on d\_1 vaccine
efficacy, infection fatality rates, and basic epidemiology. The dispute
was about which dimensions to weight in the allocation decision.
Different allocation schemes represent different metric tensors
Sigma\_B, and the disagreement is an empirical question about
dimensional weighting --- which dimensions matter and how much --- not a
philosophical disagreement about values. This reframing transforms an
apparently intractable values conflict into a tractable question:
measure the actual d\_3 inequity, the actual d\_6 social impact, the
actual d\_5 trust deficit, and compute the manifold-optimal allocation.
The answer may still be contested, but the contest is now about
measurable quantities rather than incommensurable intuitions.

\subsection{27.12 Falsifiable
Predictions}\label{falsifiable-predictions-7}

The framework generates six predictions that distinguish geometric
bioethics from scalar bioethical analysis:

\textbf{Prediction 1 (Germline Dimensional Activation)}. Ethical debates
about germline editing should activate more manifold dimensions than
debates about somatic editing. Specifically, germline discussions should
show significantly higher engagement with d\_7 (identity and species
integrity), d\_2 (rights of future persons), and d\_6 (population-level
impact) compared to somatic editing discussions, which should
concentrate on d\_1 (outcomes) and d\_4 (patient autonomy). Falsified
if: germline and somatic editing debates activate the same dimensional
profile.

\textbf{Prediction 2 (Enhancement Equity Compounding)}. Access to
enhancement technologies should correlate with existing socioeconomic
advantage, and the correlation should strengthen across generations.
Communities with unequal access to cognitive or genetic enhancement
should show increasing metric asymmetry between enhanced and unenhanced
subpopulations over time. Falsified if: enhancement access is
independent of prior advantage, or the correlation weakens over time.

\textbf{Prediction 3 (Research Ethics and Gauge Variance)}. Historical
and contemporary research ethics violations should correlate with
gauge-invariance failures: studies in which consent was obtained under
descriptions that differ significantly from the actual research protocol
should be more likely to be judged unethical than studies with
gauge-invariant consent. Falsified if: gauge variance in consent
descriptions is uncorrelated with ethical violation judgments.

\textbf{Prediction 4 (Public Health Mandate Acceptance)}. Public
acceptance of public health mandates (vaccination, quarantine, masking)
should correlate with the population\textquotesingle s d\_6 (social
impact) activation --- the degree to which individuals weight collective
welfare in their moral reasoning. Populations with higher d\_6
activation should show higher mandate compliance. Falsified if: mandate
acceptance is independent of d\_6 activation.

\textbf{Prediction 5 (Genetic Privacy and Relatedness)}. Genetic privacy
concerns should be highest among individuals who have genetically
related but non-consenting family members --- the involuntary d\_9
exposure effect. Individuals with large families, especially families
with members who have not consented to genetic data sharing, should
report stronger privacy concerns than individuals without such
relatives. Falsified if: genetic privacy concern is independent of
family size and relatives\textquotesingle{} consent status.

\textbf{Prediction 6 (DURC Governance Effectiveness)}. Dual-use research
governance frameworks that incorporate more manifold dimensions (not
just d\_1 benefit versus d\_1 risk, but also d\_2 rights, d\_5 trust,
d\_6 population impact, d\_8 regulatory adequacy) should be more
effective at preventing misuse while preserving beneficial research.
Governance effectiveness should correlate positively with dimensional
coverage of the oversight framework. Falsified if: governance
effectiveness is independent of the number of dimensions incorporated.

\subsection{27.13 Connection to the
Framework}\label{connection-to-the-framework-8}

The Geometric Bioethics program extends the parent framework in six
directions:

1. Chapter 21 applied the framework to individual clinical encounters
--- one doctor, one patient, one decision. This chapter extends it to
population-level decisions where the affected parties number in
millions, may not yet exist, and cannot consent. The structural
differences (heritable irreversibility, consent impossibility,
population-scale complexity) required new theorems (Germline
Irreversibility, Reproductive Proxy Asymmetry, Research Double Consent)
that have no analogs in Chapter 21.

2. Chapter 12 established gauge invariance as a necessary condition for
ethical evaluation. This chapter applies gauge invariance to research
ethics (the Double Consent condition), genetic privacy (consent under
different descriptions of data use), and identifies historical research
ethics violations as gauge-invariance failures. The Tuskegee study is a
canonical example of maximally gauge-variant consent.

3. Chapter 14 established collective agency and the divergence between
individual and collective geodesics. This chapter applies that
divergence to public health ethics, where vaccination mandates and
quarantine measures impose individual d\_4 costs to achieve collective
d\_6 benefits --- the bioethical instantiation of the
collective-individual geodesic tension.

4. Chapter 15 established scalar irrecoverability. This chapter
demonstrates its consequences in bioethics: the "benefit versus risk"
framing of CRISPR, enhancement, and DURC is a d\_1 projection that
destroys the eight dimensions on which these debates actually turn. The
He Jiankui case is the bioethical analog of the financial Flash Crash
(Chapter 23) --- a decision optimized on a scalar projection that was
catastrophic on the full manifold.

5. The Enhancement Equity Problem (Theorem 27.2) introduces a new
phenomenon: self-reinforcing metric asymmetry. Unlike the metric
distortions in Chapters 20-23, which are static or cyclical,
enhancement-driven asymmetry is monotonically increasing --- it
compounds across generations without bound, raising the possibility of
formal speciation (Corollary 27.2). This is the bioethical analog of the
intergenerational equity concerns in environmental ethics, where present
decisions impose irreversible costs on future generations.

6. The consent impossibility boundary (Definition 27.3) introduces a
boundary type unique to population-level bioethics: a boundary that
cannot be crossed because the consenting agent does not exist. This is
structurally different from the clinical consent boundary (Chapter 21),
which can be crossed via emergency exception or surrogate
decision-making. The consent impossibility boundary is the
framework\textquotesingle s formal acknowledgment that some ethical
challenges --- germline editing, reproductive selection, gene drive
releases --- involve a principled limit on consent-based reasoning.

\subsection{27.14 Summary}\label{summary-18}

This chapter has shown that the geometric ethics framework, when applied
to population-level bioethics, yields:

1. A formal construction of the bioethical decision complex B as a
domain-specific instantiation of the moral manifold, with
population-level bioethics distinguished from clinical ethics (Chapter
21) by four structural features: non-existent stakeholders, heritable
irreversibility, population-scale complexity, and consent impossibility.

2. Germline irreversibility as topological absorption: germline
modifications create absorbing states from which no path returns to the
pre-edit manifold topology --- a stronger form of irreversibility than
any clinical intervention, with the He Jiankui case as a
nine-dimensional catastrophe masked by scalar framing.

3. The Enhancement Equity Problem: self-reinforcing metric asymmetry
between enhanced and unenhanced populations that compounds across
generations, with speciation as the limiting case when the two
populations occupy disconnected components of the decision complex.

4. Reproductive proxy asymmetry: the structural impossibility of
calibrating the proxy\textquotesingle s heuristic to the future
person\textquotesingle s hypothetical heuristic, with the non-identity
problem formalized as the absence of a canonical cross-vertex metric.

5. The Research Double Consent condition: population-level bioethics
research requires both individual and population gauge invariance, with
historical violations (Tuskegee, Guatemala) identified as total
gauge-invariance failures whose d\_5 costs continue compounding decades
later.

6. DURC as boundary ambiguity where the edge weight depends on the
downstream path, formalizing the gain-of-function debate as a
multi-dimensional trade-off that scalar risk-benefit analysis cannot
capture.

7. Neural modification asymmetry: d\_4-expanding and d\_4-restricting
modifications have opposite curvature signatures, providing a formal
basis for distinguishing cognitive enhancement from cognitive coercion.

8. Public health mandates as collective-individual geodesic divergence,
with vaccine hesitancy formalized as inadmissible and gauge-variant
heuristic miscalibration.

\section{Chapter 28: Geometric Military Ethics --- War and Constrained
Pathfinding on the Decision
Manifold}\label{chapter-28-geometric-military-ethics-war-and-constrained-pathfinding-on-the-decision-manifold}

Epistemic status of this chapter. Except where a claim is explicitly
marked as a formal result or a conditional theorem, the material here is
a modeling interpretation, domain analogy, or speculative extension - a
demonstration of how the geometric vocabulary applies to military
ethics, not a set of established results at the evidential standard of
the field.

\begin{quote}
\textbf{\textsc{RUNNING EXAMPLE --- Priya's Model}}

\emph{One patient in the BEACON-7 cohort is Sergeant (Ret.) David
Kowalski, whose melanoma is linked to burn-pit exposure during
deployment. His VA benefits cover trial participation, but TrialMatch
does not recognize VA coverage as equivalent to private insurance---a
feature encoding error that systematically disadvantages veterans. The
constrained pathfinding of military ethics---jus in bello as bounded
optimization under rules of engagement---mirrors TrialMatch's own
constrained optimization. Both systems face the same question: who sets
the constraints, and are the constraints just?}
\end{quote}

This chapter applies the Geometric Ethics framework to military
decision-making and the ethics of armed conflict. The dominant tradition
in military ethics --- just war theory, from Augustine through Aquinas
to Walzer --- provides a rich multi-dimensional vocabulary: just cause,
right authority, proportionality, distinction, military necessity. Yet
the formalization of these concepts in International Humanitarian Law
(IHL) and operational practice has systematically reduced them to scalar
tests. Proportionality becomes a cost-benefit comparison of military
advantage against civilian harm. Distinction becomes a binary
classification: combatant or civilian. Rules of engagement become
checklists. The geometric framework reveals that these scalar reductions
destroy information irrecoverably, and that the characteristic
pathologies of modern warfare --- civilian casualty spirals, moral
injury epidemics, autonomous weapons failures, and post-conflict
instability --- are predictable consequences of dimensional collapse in
the highest-stakes domain of human decision-making.

\subsection{28.1 The Failure of Scalar Military
Ethics}\label{the-failure-of-scalar-military-ethics}

Military decision-making has historically been dominated by a single
scalar: military necessity. From Clausewitz\textquotesingle s dictum
that war is the continuation of politics by other means to the modern
doctrine of "effects-based operations," the tendency has been to
evaluate military actions primarily on d\_1 --- mission effectiveness,
tactical advantage, force protection. This scalar reduction is not
merely an intellectual simplification; it is operationally embedded in
targeting processes, command structures, and weapons systems.

Just war theory, in its classical formulation, resisted this collapse.
Augustine\textquotesingle s requirement of right intention (d\_7: the
moral character of the belligerent), Aquinas\textquotesingle s
insistence on legitimate authority (d\_8: institutional sanction), and
Grotius\textquotesingle s development of proportionality (d\_3: fairness
between harm and advantage) together constitute a genuinely
multi-dimensional framework. Walzer\textquotesingle s "Just and Unjust
Wars" (1977) further enriched the tradition by attending to the moral
experience of soldiers (d\_7), the rights of civilians (d\_2, d\_4), and
the epistemic demands of targeting (d\_9).

Yet the translation of just war theory into positive law --- the Geneva
Conventions (1949), the Additional Protocols (1977), the Rome Statute
(1998) --- has progressively flattened this multi-dimensional structure.
The IHL proportionality test, as codified in Additional Protocol I,
Article 51(5)(b), requires that an attack not be launched which "may be
expected to cause incidental loss of civilian life, injury to civilians,
damage to civilian objects, or a combination thereof, which would be
excessive in relation to the concrete and direct military advantage
anticipated." This is a scalar comparison: one number (civilian harm)
weighed against another (military advantage). The nine dimensions of the
moral manifold are collapsed to a single cost-benefit ratio.

Rules of engagement (ROE) attempt to restore multi-dimensionality by
encoding constraints as checklists --- positive identification
requirements (d\_9), escalation-of-force procedures (d\_3),
protected-site restrictions (d\_2). But checklists are scalar
approximations of manifold pathfinding: they discretize continuous
dimensional trade-offs into binary gates, losing the geometric structure
that determines which paths through moral space are genuinely optimal.

\subsection{28.2 The Military Decision
Complex}\label{the-military-decision-complex}

Definition 28.1 (Military Decision Complex). The military decision
complex M\_mil is a weighted simplicial complex whose vertices are
military states --- configurations of force disposition, civilian
population status, territorial control, intelligence picture, and
command authority --- and whose edges are military actions (engagements,
movements, targeting decisions, rules of engagement modifications,
ceasefire initiatives). Each vertex v\_i carries an attribute vector
a(v\_i) in R\^{}9.

The nine dimensions of the moral manifold, instantiated for military
contexts, are:

\textbf{d\_1: Military effectiveness} --- mission success probability,
force protection, tactical and strategic advantage. This is the
dimension on which military necessity operates.

\textbf{d\_2: Legal obligations} --- compliance with IHL, the Geneva
Conventions, the Laws of Armed Conflict (LOAC), status-of-forces
agreements, and applicable ROE. These are binding constraints, not
optional considerations.

\textbf{d\_3: Proportionality and fairness} --- the balance between
military advantage and civilian harm, proportionate use of force, the
principle of unnecessary suffering (Hague Convention IV, Art. 23(e)).

\textbf{d\_4: Combatant and civilian autonomy} --- freedom of movement
for civilian populations, rights of detainees (Geneva Convention III),
prohibition of forced displacement (Rome Statute, Art. 8(2)(e)(viii)).

\textbf{d\_5: Trust} --- between allied forces, between military and
civilian populations, between combatants and their chain of command,
between belligerents and mediators. Trust is the lubricant of both
military operations and post-conflict stability.

\textbf{d\_6: Strategic and societal impact} --- collateral damage to
infrastructure, displacement of populations, environmental destruction,
long-term economic consequences, cultural heritage destruction.

\textbf{d\_7: Warrior identity and honor} --- the
soldier\textquotesingle s moral self-concept, military ethos,
professional ethics, the traditions of chivalry and martial honor that
constrain behavior even in extremis.

\textbf{d\_8: Institutional legitimacy} --- UN Security Council
authorization, democratic mandate for the use of force, legitimacy of
the chain of command, international coalition support.

\textbf{d\_9: Intelligence and epistemic status} --- reliability of
intelligence, certainty of target identification, fog of war,
communication integrity, situational awareness.

\textbf{Definition 28.2 (Military Edge Weights)}. The weight of an edge
(v\_i, v\_j) in M\_mil is:

w(v\_i, v\_j) = Da\^{}T Sigma\_M\^{}\{-1\} Da + Sum\_k beta\_k *
1{[}boundary k crossed{]}

where Da = a(v\_j) - a(v\_i), Sigma\_M is the 9x9 military covariance
matrix, and beta\_k are boundary penalties. Critical boundary types
include:

War crime boundaries (beta = infinity): deliberate targeting of
civilians, torture, perfidy, use of prohibited weapons. These are
sacred-value boundaries that no military advantage can justify.

Geneva Convention thresholds: treatment of prisoners of war, protection
of medical facilities, protection of cultural property. These carry high
but finite beta values.

ROE constraints: engagement authorization levels, escalation-of-force
steps, positive identification requirements. These are operationally
calibrated boundary penalties.

Critical covariance terms include Sigma\_\{1,9\} (effectiveness x
intelligence: military effectiveness is coupled to intelligence quality
--- poor intelligence degrades all targeting decisions), Sigma\_\{5,6\}
(trust x societal impact: civilian trust erodes with collateral damage,
creating insurgency feedback loops), and Sigma\_\{7,3\} (warrior
identity x proportionality: soldiers whose identity includes
proportionality norms internalize these constraints at lower cognitive
cost).

The military decision complex is distinguished from other domain
instantiations by the extremity of its d\_9 conditions. The "fog of war"
--- Clausewitz\textquotesingle s central insight --- is not merely
informational uncertainty but a structural feature of the manifold: d\_9
variance in combat is higher than in any other domain, and this variance
propagates through the covariance matrix to degrade every other
dimension simultaneously.

\subsection{28.3 Jus ad Bellum as Manifold Entry
Conditions}\label{jus-ad-bellum-as-manifold-entry-conditions}

The decision to go to war --- jus ad bellum --- is the highest-stakes
decision in the military domain. The just war tradition, from Augustine
to the UN Charter, identifies six criteria: just cause, right authority,
right intention, last resort, proportionality (ad bellum), and
reasonable chance of success. The geometric framework formalizes these
as entry conditions on the military decision complex.

\textbf{Theorem 28.1 (Jus ad Bellum as Manifold Threshold)}. The
decision to go to war is justified if and only if the war-path geodesic
on the full military decision complex satisfies entry conditions on
specified dimensions: (i) d\_1 (just cause): the threat level exceeds
the Just Cause Threshold, below which the peace-path geodesic dominates
the war-path on the full manifold; (ii) d\_8 (right authority): the
decision is made by a legitimate authority with institutional sanction;
(iii) d\_7 (right intention): the goal region of the war-path geodesic
is a morally acceptable end state, not territorial expansion, resource
extraction, or revenge; (iv) d\_2 (last resort): all non-violent paths
on the manifold have been exhausted or are demonstrably infeasible; (v)
d\_3 (proportionality ad bellum): the total expected cost of the
war-path across all nine dimensions does not exceed the cost of the
threat it addresses; and (vi) d\_1 (reasonable chance of success): the
war-path geodesic reaches the goal region with non-negligible
probability.

\textbf{Proof}. The war-path and peace-path are competing geodesics on
M\_mil. The peace-path has lower d\_1 immediate cost (no combat
casualties) but may have higher long-term d\_1 cost if the threat
materializes (aggression unchecked, genocide unimpeded). The war-path
has higher immediate d\_1 cost but may have lower long-term cost if the
threat is neutralized. The six jus ad bellum criteria are necessary
conditions for the war-path to dominate: just cause ensures the threat
is real (the d\_1 differential is genuine, not fabricated); right
authority ensures the decision has institutional backing (d\_8 is
satisfied); right intention ensures the goal region is acceptable (d\_7
is not compromised); last resort ensures the peace-path is genuinely
infeasible (d\_2 obligations have been met); proportionality ad bellum
ensures the total manifold cost is justified; and reasonable chance of
success ensures the war-path actually reaches the goal region. If any
condition fails, there exists a peace-path that dominates on the full
manifold. {[}{]}

\textbf{Definition 28.3 (Just Cause Threshold)}. The Just Cause
Threshold T\_jc is the minimum d\_1 threat level below which the
war-path geodesic is dominated by the peace-path on the full manifold.
Formally, T\_jc is the d\_1 value at which the total manifold cost of
the war-path (including d\_2 through d\_9 costs of armed conflict)
equals the total manifold cost of the peace-path (including the d\_1
cost of the unrealized threat). Below T\_jc, the peace-path is cheaper
on the full manifold even if the war-path is cheaper on d\_1 alone.

\textbf{Application}: Iraq 2003 as Jus ad Bellum Failure. The 2003
invasion of Iraq violated at least three jus ad bellum conditions
simultaneously. On d\_9 (intelligence), the casus belli --- Iraqi
weapons of mass destruction --- was catastrophically wrong; the
intelligence was fabricated, cherry-picked, or misinterpreted, meaning
the Just Cause Threshold was not actually met because the threat was
illusory. On d\_8 (institutional legitimacy), the invasion lacked UN
Security Council authorization; the "coalition of the willing" was a
workaround for institutional rejection, not a legitimate substitute. On
d\_7 (right intention), the goal region was ambiguous --- regime change,
democracy promotion, WMD elimination, and regional transformation were
variously cited, none coherently specified. The geometric framework
predicts that violations of multiple entry conditions simultaneously
should produce catastrophic failure of the resulting war-path, a
prediction borne out by the subsequent insurgency, sectarian civil war,
and regional destabilization.

\textbf{Remark (Nuclear Weapons and the Just Cause Threshold)}. Nuclear
weapons raise the Just Cause Threshold T\_jc to near-infinity for wars
between nuclear-armed states: the total manifold cost of nuclear
exchange on d\_1 (mass casualties), d\_6 (civilizational destruction),
d\_4 (annihilation of autonomy), and d\_5 (permanent trust destruction)
is so extreme that no threat short of existential attack can justify
crossing T\_jc. This is the geometric content of nuclear deterrence: the
war-path geodesic between nuclear-armed states is dominated by the
peace-path on the full manifold for virtually all threat levels, because
the d\_6 cost of the war-path includes potential civilizational
extinction.

\subsection{28.4 Jus in Bello as Constrained
Pathfinding}\label{jus-in-bello-as-constrained-pathfinding}

Once armed conflict has begun, jus in bello constrains how the war is
fought. The geometric framework formalizes jus in bello as constrained
A* pathfinding on the military decision complex: the
soldier\textquotesingle s feasible path set is restricted by boundary
penalties on non-d\_1 dimensions, even when unconstrained paths would be
militarily optimal.

\textbf{Theorem 28.2 (Jus in Bello Constraints)}. Within the military
decision complex, combat operations are constrained A* pathfinding: the
soldier\textquotesingle s feasible path set F\_constrained is a proper
subset of the unconstrained feasible set F\_unconstrained, restricted by
boundary penalties on d\_2 (IHL obligations: no targeting of civilians,
no torture, no perfidy), d\_3 (proportionality in bello: incidental harm
must not be excessive), and d\_4 (civilian autonomy: no forced
displacement, humane treatment of detainees). The minimum-cost path on
F\_constrained has cost greater than or equal to the minimum-cost path
on F\_unconstrained.

\textbf{Proof}. The constrained feasible set F\_constrained is obtained
from F\_unconstrained by adding boundary penalties beta\_k
\textgreater{} 0 for crossings of d\_2, d\_3, and d\_4 boundaries.
Adding penalties can only increase or maintain the cost of any path that
crosses a boundary, and can only decrease or maintain the size of the
feasible set (paths whose boundary-crossing cost exceeds any finite
budget become infeasible). Therefore F\_constrained is a subset of
F\_unconstrained. The minimum-cost path on a subset is at least as
costly as the minimum-cost path on the full set. The difference
C\_constrained - C\_unconstrained \textgreater= 0 is the jus in bello
cost --- the cost of fighting lawfully. {[}{]}

\textbf{Corollary 28.1 (The Ethical Cost of Lawful Combat)}. The jus in
bello cost C\_jib = C\_constrained - C\_unconstrained is strictly
positive whenever the unconstrained optimal path crosses at least one
jus in bello boundary. This cost is borne by the soldier, the military
unit, and the nation as reduced military effectiveness on d\_1. The
acceptance of this cost is the defining commitment of lawful
belligerency: a military that refuses to pay C\_jib is, by definition,
committing war crimes.

The practical significance is immediate. When commanders complain that
IHL "ties our hands," they are correctly identifying C\_jib
\textgreater{} 0. The geometric framework validates this observation
while insisting that the cost is non-negotiable: the boundary penalties
on d\_2, d\_3, and d\_4 encode sacred-value constraints (beta = infinity
for war crime boundaries) that cannot be traded against d\_1 gains
regardless of their magnitude.

\subsection{28.5 Proportionality as Multi-Dimensional
Cost-Benefit}\label{proportionality-as-multi-dimensional-cost-benefit}

The IHL proportionality test is the most contested principle in the law
of armed conflict. Additional Protocol I, Article 51(5)(b) prohibits
attacks "which may be expected to cause incidental loss of civilian
life, injury to civilians, damage to civilian objects, or a combination
thereof, which would be excessive in relation to the concrete and direct
military advantage anticipated." This formulation invites scalar
comparison: one number (civilian harm) weighed against another (military
advantage).

\textbf{Theorem 28.3 (Scalar Proportionality Failure)}. The IHL
proportionality test, as typically applied, compares d\_1 cost (civilian
casualties as a subset of consequences) against d\_1 benefit (military
advantage as another subset of consequences). This is a one-dimensional
comparison within a single dimension. The geometric framework
demonstrates that this scalar test collapses at least four additional
dimensions: d\_2 (civilian rights to protection under IHL), d\_4
(civilian autonomy --- displacement, loss of freedom of movement), d\_6
(societal impact --- infrastructure destruction, community disruption,
environmental damage), and d\_7 (cultural heritage and identity ---
destruction of religious sites, historical monuments, cultural
practices). By Scalar Irrecoverability (Chapter 15), the information
destroyed by this collapse is mathematically unrecoverable: no post hoc
analysis can reconstruct the multi-dimensional cost from the scalar
comparison.

\textbf{Proof}. The scalar proportionality test evaluates f: R\^{}9
-\textgreater{} R, mapping the nine-dimensional cost-benefit vector to a
single ratio. By Brouwer\textquotesingle s invariance of dimension, no
continuous injection from R\^{}9 to R exists. Therefore, distinct
nine-dimensional states --- attacks with identical civilian casualty
counts but vastly different d\_4, d\_6, and d\_7 costs --- map to the
same scalar proportionality score. The test declares them equally
proportionate when they are, in the geometric sense, profoundly
different. {[}{]}

\textbf{Proposition 28.1 (Geometric Proportionality)}. Geometric
proportionality requires that the military action\textquotesingle s cost
on EACH non-d\_1 dimension be independently proportionate to the
military advantage on d\_1. Formally, for each dimension k in \{d\_2,
d\_3, d\_4, d\_5, d\_6, d\_7, d\_8, d\_9\}, the ratio
\textbar Da\_k\textbar{} / \textbar Da\_1\^{}\{benefit\}\textbar{} must
not exceed a dimension-specific threshold tau\_k. This multi-dimensional
test preserves the information that scalar proportionality destroys.

\textbf{Lemma 28.1 (Proportionality Asymmetry)}. The scalar
proportionality test is asymmetrically biased toward the attacker. The
military advantage (d\_1 benefit) is evaluated from the
attacker\textquotesingle s perspective using the
attacker\textquotesingle s information (d\_9). The civilian harm (d\_1
cost) is evaluated using the attacker\textquotesingle s estimate of
civilian presence, which is systematically lower than the true civilian
cost because the attacker has incentives to underestimate and lacks
ground-truth access. In the geometric framework, this asymmetry means
the scalar proportionality test is computed with biased inputs: the d\_1
benefit is overestimated (motivated reasoning) and the d\_1 cost is
underestimated (information asymmetry). The resulting proportionality
ratio is systematically biased toward permitting strikes.

\textbf{Example}. Consider the bombing of a hospital that serves as a
military command post. The scalar proportionality test weighs X civilian
casualties (d\_1 cost) against the military value of destroying the
command post (d\_1 benefit). If the ratio is acceptable, the strike is
proportionate. But the geometric test identifies additional costs
invisible to the scalar test: destruction of medical infrastructure
serving a civilian population of thousands (d\_6), violation of the
protected status of medical facilities under Geneva Convention I, Art.
19 (d\_2), forced displacement of patients and medical personnel (d\_4),
destruction of community trust in the protection of medical facilities
(d\_5), and damage to the cultural norm that hospitals are inviolable
(d\_7). Each of these costs must be independently proportionate to the
military advantage. The scalar test captures only the first; the
geometric test captures all six.

\subsection{28.6 The Doctrine of Double Effect as Dimensional
Decomposition}\label{the-doctrine-of-double-effect-as-dimensional-decomposition}

The Doctrine of Double Effect (DDE), originating in
Aquinas\textquotesingle s Summa Theologiae (II-II, Q. 64, Art. 7), holds
that an action with both good and bad effects is permissible if four
conditions are met: (1) the action itself is not intrinsically wrong;
(2) the agent intends only the good effect; (3) the bad effect is not a
means to the good effect; and (4) the good effect is proportionate to
the bad effect. DDE is foundational to military targeting: the killing
of civilians as a side effect of a legitimate military strike may be
permissible, while the deliberate targeting of civilians is always
prohibited.

\textbf{Theorem 28.4 (DDE as Dimensional Decomposition)}. The Doctrine
of Double Effect decomposes a military action into its effects on each
dimension separately. The "good effect" is the d\_1 (military
effectiveness) component of the action. The "bad effect" is the
aggregate of d\_1 (civilian harm), d\_2 (rights violations), d\_3
(proportionality costs), d\_4 (autonomy restrictions), d\_5 (trust
damage), and d\_6 (societal impact) components. DDE\textquotesingle s
condition (3) --- "the bad effect is not a means to the good effect" ---
requires that the bad-effect dimensions are not instrumentally necessary
for achieving the good-effect dimension. Formally: there exists an
alternative path P\textquotesingle{} with comparable d\_1 benefit
(\textbar Da\_1(P\textquotesingle) - Da\_1(P)\textbar{} \textless{}
epsilon) and strictly lower cost on at least one bad-effect dimension
(Da\_k(P\textquotesingle) \textless{} Da\_k(P) for some k in \{d\_2,
d\_3, d\_4, d\_5, d\_6\}) if and only if the bad effect is NOT a means.
If no such P\textquotesingle{} exists --- if the bad-effect cost is
inextricable from the good-effect benefit --- then the bad effect IS a
means, and DDE prohibits the action.

\textbf{Proof}. DDE\textquotesingle s condition (3) distinguishes
instrumental from incidental harm. In the geometric framework,
instrumental harm means that the path from current state to goal state
necessarily traverses the harmful vertices --- there is no detour.
Incidental harm means that the path traverses harmful vertices but an
alternative path with comparable d\_1 benefit avoids them. The existence
of such an alternative path P\textquotesingle{} is a topological
property of the military decision complex: it holds when the harmful
vertices are not on every geodesic from the current state to the goal
region. If every geodesic traverses the harmful vertices, the harm is
instrumental (a necessary means), and DDE prohibits the action. If at
least one geodesic avoids them, the harm is incidental, and DDE permits
the action subject to proportionality (condition 4). {[}{]}

\textbf{Application: Strategic Bombing versus Precision Strikes}. The
Allied strategic bombing campaign of World War II --- the firebombing of
Dresden, Tokyo, Hamburg --- treated d\_6 (civilian infrastructure
destruction) as instrumental to d\_1 (military effectiveness): the
theory of strategic bombing held that destroying civilian morale and
industrial capacity was the means to military victory. Under DDE, this
is impermissible: the bad effect (mass civilian death and infrastructure
destruction) is the means to the good effect (reduced enemy war-fighting
capacity). Precision strikes, by contrast, target specific military
objectives (a bridge, a command post, a weapons depot) with civilian
harm as incidental --- an alternative path with comparable d\_1 benefit
and lower d\_6 cost exists (a more precise weapon, a different angle of
attack, a different time of day). DDE permits precision strikes subject
to proportionality but prohibits strategic bombing that treats civilian
destruction as the mechanism of military effect.

\textbf{Remark (The Trolley Problem Formalized)}. In the classic trolley
problem, the fat man\textquotesingle s body stops the trolley --- the
bad effect (killing the man) IS the means to the good effect (stopping
the trolley). In the geometric framework, there is no alternative path
P\textquotesingle{} that achieves the same d\_1 benefit (trolley
stopped) without the d\_2 cost (killing the man): the harmful vertex is
on every geodesic to the goal. DDE prohibits pushing the fat man. In the
switch variant, diverting the trolley kills a bystander incidentally ---
the death is not the mechanism that saves the five. An alternative path
exists in principle (stop the trolley by other means), but if no such
path is feasible, DDE permits diversion subject to proportionality.

\subsection{28.7 Civilian-Combatant Distinction as Boundary
Classification}\label{civilian-combatant-distinction-as-boundary-classification}

The principle of distinction --- the requirement that parties to a
conflict distinguish between combatants and civilians, and direct
attacks only against combatants --- is codified in Additional Protocol
I, Article 48, and is recognized as customary international law binding
on all states. In the geometric framework, distinction is a boundary
classification problem on the military decision complex.

\textbf{Definition 28.4 (Distinction as Boundary Classification)}. Each
vertex in the military decision complex is classified as combatant
(targetable) or civilian (protected). The boundary penalty for targeting
a vertex classified as civilian is beta\_civilian = infinity. This is a
sacred-value boundary: no military advantage, however great, justifies
the deliberate targeting of a civilian. The penalty is not finite and
large --- it is infinite, reflecting the jus cogens status of the
prohibition on deliberate civilian targeting under the Rome Statute,
Article 8(2)(b)(i).

\textbf{Theorem 28.5 (Distinction Degradation in Asymmetric Warfare)}.
In asymmetric warfare --- insurgency, counterterrorism, urban combat ---
the combatant-civilian boundary becomes probabilistic. Vertices cannot
be cleanly classified as combatant or civilian: the insurgent is a
farmer by day and a fighter by night; the civilian carries a phone that
may or may not be a detonator; the building is a home that may or may
not contain a weapons cache. The effective boundary penalty becomes:

beta\_effective = beta\_civilian * P(civilian \textbar{} observation)

where P(civilian \textbar{} observation) is the probability that the
target is civilian given available intelligence. As d\_9 (intelligence
quality) degrades, P(civilian \textbar{} observation) approaches the
base rate P(civilian) in the population --- which, in most conflict
zones, is greater than 0.9. The distinction boundary becomes
computationally intractable: nearly every engagement carries
near-infinite expected boundary cost.

\textbf{Proof}. Under perfect intelligence (d\_9 maximum), P(civilian
\textbar{} observation) = 0 for combatants and P(civilian \textbar{}
observation) = 1 for civilians. The boundary classification is
deterministic, and the boundary penalty is applied correctly. Under
degraded intelligence, P(civilian \textbar{} observation) is a
continuous function of observation quality on {[}0, 1{]}. The expected
boundary cost for targeting vertex v is E{[}beta{]} = beta\_civilian *
P(civilian \textbar{} observation at v). Since beta\_civilian =
infinity, the expected cost is infinite whenever P(civilian \textbar{}
observation) \textgreater{} 0. In practice, forces treat beta\_civilian
as very large but finite, implicitly accepting a non-zero probability of
civilian targeting. The d\_9 degradation characteristic of asymmetric
warfare drives P(civilian \textbar{} observation) upward for every
engagement, systematically increasing the expected boundary cost and the
probability of civilian casualties. {[}{]}

\textbf{Application: Human Shields as Boundary Exploitation}. The use of
human shields --- placing civilian vertices adjacent to combatant
vertices --- is a deliberate exploitation of the distinction boundary.
The adversary forces the geometric structure of M\_mil into a
configuration where combatant and civilian vertices are topologically
inseparable: every path to the combatant vertex traverses civilian
vertices. This makes the distinction boundary computationally
intractable in the most literal sense: there is no path to the
legitimate military objective that does not cross the sacred-value
boundary. The attacker faces a dilemma that is geometric, not merely
moral: the manifold has been engineered to make lawful engagement
impossible. IHL addresses this by placing responsibility on the party
using human shields (Additional Protocol I, Art. 51(7)) while
maintaining that the attacking party\textquotesingle s obligation to
distinguish is not extinguished.

\subsubsection{Siege Warfare and Collective
Punishment}\label{siege-warfare-and-collective-punishment}

\textbf{Siege warfare --- the encirclement and starvation of a
population to compel surrender} --- is the limiting case of distinction
degradation. In a siege, the distinction boundary is not merely
ambiguous but deliberately ignored: the entire besieged population,
combatant and civilian alike, is subjected to d\_4 (autonomy: denied
freedom of movement), d\_6 (societal impact: destruction of
infrastructure, food systems, medical care), and d\_1 (consequences:
starvation, disease, death) costs. The geometric framework classifies
siege that targets civilian survival as collective punishment --- a war
crime under Geneva Convention IV, Article 33 --- because it sets
beta\_civilian = 0 for an entire population, treating every vertex as
targetable regardless of classification.

The framework predicts a testable correlation: civilian casualty rates
should increase monotonically with d\_9 degradation. In conflicts with
good intelligence (clear front lines, uniformed combatants, reliable
surveillance), civilian casualties should be lower. In conflicts with
poor intelligence (urban insurgency, unmarked combatants, degraded
communications), civilian casualties should be higher. The empirical
evidence --- comparing casualty ratios across conventional warfare
(World War II European theater), counterinsurgency (Vietnam, Iraq), and
urban warfare (Gaza, Mosul) --- is consistent with this prediction.

\subsection{28.8 Autonomous Weapons and Dimensional
Collapse}\label{autonomous-weapons-and-dimensional-collapse}

\textbf{Autonomous weapons systems (AWS)} --- weapons that can select
and engage targets without human intervention --- represent the military
instantiation of the dimensional collapse phenomenon analyzed in Chapter
23 (Flash Crash) and Chapter 26 (AI Ethics). This section addresses AWS
from the military ethics perspective, focusing on IHL compliance and
just war theory; Chapter 26 addresses them from the AI ethics
perspective, focusing on alignment and control.

\textbf{Theorem 28.6 (AWS Dimensional Collapse)}. An autonomous weapon
system optimizing on d\_1 (target elimination effectiveness) with
algorithmic d\_9 (sensor-based target identification) operates on a
collapsed manifold of at most two dimensions. The system cannot maintain
the distinction boundary (Definition 28.4) because boundary
classification requires evaluation on d\_2 (legal status of the target
under IHL --- is this person a combatant, a civilian, a medical worker,
a surrendering soldier?), d\_3 (proportionality of this specific
engagement in this specific context --- is the expected civilian harm
excessive relative to the military advantage of this particular
target?), d\_7 (ethical judgment --- does this engagement comport with
the warrior ethos, with the spirit and not merely the letter of the
law?), and contextual d\_9 (human intelligence --- information about
local conditions, cultural context, and recent events that sensor data
cannot capture). Current AI systems cannot compute these dimensions.

\textbf{Proof}. An AWS receives sensor data (visual, infrared, radar)
and classifies targets using pattern recognition algorithms trained on
labeled datasets. The classification operates on d\_1 (is this a
militarily valuable target?) and d\_9 (how confident is the
identification?). The system\textquotesingle s edge weights are:
w\_AWS(v\_i, v\_j) = (Da\_1)\^{}2 / sigma\_11 + (Da\_9)\^{}2 /
sigma\_99. Boundary penalties on d\_2, d\_3, d\_4, d\_7 are set to zero
--- not because the designers intended to ignore them, but because the
system has no computational representation of these dimensions. The
feasible path set under this collapsed metric strictly contains the
feasible path set under the full IHL-compliant metric: the AWS can
traverse paths (engage targets) that a full-manifold agent (a human
soldier applying IHL) would never traverse because they cross d\_2,
d\_3, or d\_7 boundaries. This is the lethal analog of the Flash Crash
(Chapter 23): just as algorithmic trading systems caused the Flash Crash
by collapsing to d\_1 (monetary return), AWS collapse to d\_1 (target
elimination) in a domain where the consequences of dimensional collapse
are measured in human lives rather than dollars. {[}{]}

\textbf{Proposition 28.2 (Meaningful Human Control Requirement)}.
Compliance with IHL requires that autonomous weapons systems maintain
human oversight with minimum dimensionality k \textgreater= 4. The four
minimum dimensions are: d\_1 (target military value), d\_2 (target legal
status under IHL), d\_3 (proportionality of the specific engagement),
and d\_9 (contextual intelligence beyond sensor data). Below this
threshold, the system cannot reliably distinguish lawful from unlawful
engagements, and every firing decision carries non-trivial probability
of war crime boundary violation.

\textbf{Remark (The Accountability Gap)}. AWS introduce a novel problem
for IHL: the accountability gap. Under the Rome Statute, criminal
responsibility attaches to individuals who order, commit, or fail to
prevent war crimes. When an AWS commits a targeting error that
constitutes a war crime, who is responsible? The commander who
authorized deployment? The engineer who designed the algorithm? The
procurement officer who selected the system? In the geometric framework,
the accountability gap is a consequence of dimensional collapse: the
AWS\textquotesingle s decision was made on a two-dimensional manifold
(d\_1, d\_9) where the war crime boundary (a d\_2 boundary) does not
exist. No human decision-maker traversed the boundary --- the boundary
was absent from the decision space entirely. The accountability gap is
not a bug in the law; it is a structural consequence of delegating
lethal decisions to systems that operate on collapsed manifolds.

\textbf{Current technology assessment. Image classifiers} --- the core
targeting technology in contemporary AWS --- operate on d\_1 (is this a
tank, a truck, a person carrying a weapon?) and algorithmic d\_9
(confidence score from the classifier). They cannot evaluate d\_2: the
legal status of a target under IHL depends on facts that no image
classifier can determine --- is this person a combatant who has
surrendered? Is this building a protected medical facility whose
protection has been forfeited by military use? Is this person a child
soldier whose targeting raises distinct legal questions? They cannot
evaluate d\_3: proportionality requires contextual judgment about the
military advantage of this specific engagement relative to the expected
civilian harm in this specific situation --- a judgment that depends on
the overall military campaign, the availability of alternative means,
and the density of the civilian population, none of which are accessible
to a sensor-based classifier. The gap between what AWS can compute
(d\_1, d\_9) and what IHL requires (d\_1, d\_2, d\_3, d\_7, d\_9) is not
an engineering problem to be solved by better algorithms; it is a
dimensional deficit that no sensor-based system can close.

\subsection{28.9 Rules of Engagement as Heuristic
Constraints}\label{rules-of-engagement-as-heuristic-constraints}

Rules of engagement are pre-programmed decision rules issued to soldiers
and military units to govern the use of force. In the geometric
framework, ROE are heuristic constraints --- pre-computed boundary
penalties and dimensional weights that structure the
soldier\textquotesingle s pathfinding heuristic h(n), enabling rapid
decision-making in high-stress, low-d\_9 conditions where full manifold
deliberation is cognitively impossible.

\textbf{Definition 28.5 (ROE as Heuristic Constraints)}. Rules of
engagement are a heuristic function h\_ROE(n) that approximates the full
manifold cost h*(n) by encoding boundary penalties and dimensional
priorities as discrete rules. Examples: "Do not fire unless fired upon"
encodes a high d\_3 boundary penalty (proportionality: the first shot
carries asymmetric cost). "Positive identification required before
engagement" encodes a d\_9 threshold (epistemic: do not engage until
target identification confidence exceeds a minimum). "Do not target
buildings marked with a red cross" encodes a d\_2 boundary penalty
(legal: medical facilities are protected). Each rule is a dimensional
constraint translated into an operationally executable instruction.

\textbf{Proposition 28.3 (ROE as Epsilon-Admissible Heuristics)}. ROE
sacrifice optimality for computational tractability. The soldier in
combat cannot perform full nine-dimensional A* pathfinding in real time;
the cognitive and temporal demands exceed human capacity under stress.
ROE pre-compute an approximate heuristic h\_ROE(n) that is
epsilon-admissible: \textbar h\_ROE(n) - h*(n)\textbar{} \textless=
epsilon for most military states n. The parameter epsilon is the gap
between the ROE-guided path and the fully deliberated manifold geodesic
--- the cost of heuristic approximation.

ROE failure modes are characterized by the nature of epsilon:

\begin{enumerate}
\def\labelenumi{(\arabic{enumi})}
\item
  ROE too restrictive: epsilon is large and biased upward on d\_1 ---
  the heuristic overestimates the cost of engagement, preventing
  soldiers from accomplishing legitimate military objectives. The
  soldier cannot fire even when firing is lawful and necessary. This was
  a persistent complaint in early ISAF operations in Afghanistan, where
  restrictive ROE were perceived as prioritizing d\_3 (proportionality)
  over d\_1 (force protection) to a degree that endangered soldiers.
\item
  ROE too permissive: boundary penalties are reduced or eliminated on
  d\_2, d\_3, and d\_4 --- the heuristic underestimates the cost of
  boundary crossings, enabling paths that constitute war crimes. The
  soldier can fire when firing is unlawful. This is the failure mode of
  Vietnam-era "free fire zones," where ROE designated geographical areas
  in which beta\_civilian was effectively set to zero, removing the
  distinction boundary entirely.
\item
  ROE poorly calibrated: the heuristic is inadmissible --- it sometimes
  overestimates and sometimes underestimates h*(n), producing paths that
  oscillate between excessive restraint and excessive force. This is the
  failure mode of ROE that are written for one operational environment
  and applied in another.
\end{enumerate}

Application: My Lai as ROE Boundary Removal. On March 16, 1968, soldiers
of Charlie Company, 1st Battalion, 20th Infantry Regiment massacred
between 347 and 504 unarmed Vietnamese civilians in the hamlet of My
Lai. The geometric analysis is precise: the operational context
(repeated ambushes, booby traps, indistinguishable combatants and
civilians) had degraded d\_9 to near-zero. The implicit ROE for the
operation --- search and destroy in a designated "free fire zone" ---
set beta\_civilian to zero in that area. With the distinction boundary
removed and d\_9 degraded, the soldiers\textquotesingle{} heuristic
h\_ROE(n) contained no penalty for civilian targeting. The path through
the massacre was, on the collapsed manifold, the lowest-cost path on
d\_1 (eliminate all potential threats). On the full manifold, it crossed
the war crime boundary (beta = infinity) hundreds of times. My Lai is
the predictable consequence of ROE that collapse the military decision
complex to d\_1 alone.

\subsection{28.10 Moral Injury in Combat as Manifold
Damage}\label{moral-injury-in-combat-as-manifold-damage}

Moral injury --- the psychological, biological, and spiritual damage
sustained when a person perpetrates, fails to prevent, or witnesses acts
that transgress their deeply held moral beliefs --- is epidemic among
combat veterans. The geometric framework, building on the moral injury
theory developed in Chapter 21, provides a precise characterization of
combat moral injury as cumulative damage to the
soldier\textquotesingle s heuristic function.

\textbf{Theorem 28.7 (Combat Moral Injury)}. Military moral injury
occurs when the soldier is forced to traverse paths that cross their own
moral boundaries --- paths that their heuristic h(n) flags as high-cost
but that orders, operational pressure, or circumstance require them to
take. The damage is to the heuristic function itself: repeated forced
boundary crossings degrade h(n)\textquotesingle s calibration.
Specifically, let B\_soldier be the set of the soldier\textquotesingle s
internalized moral boundaries (a subset of the full boundary set
B\_mil). A forced boundary crossing on boundary b in B\_soldier inflicts
damage proportional to the soldier\textquotesingle s boundary penalty
beta\_b\^{}\{soldier\} --- the strength of the moral conviction
violated. Cumulative damage D = Sum\_t beta\_\{b(t)\}\^{}\{soldier\} *
1{[}boundary b(t) crossed at time t{]} degrades h(n) by increasing its
variance and reducing its reliability.

\textbf{Proof}. The soldier\textquotesingle s heuristic h(n) is
calibrated through training, socialization, and moral development to
reflect internalized boundary penalties. When the soldier is forced to
cross boundary b --- ordered to fire on a target they believe to be
civilian, compelled to stand by while allies commit abuses, required to
execute a strike they judge disproportionate --- the forced crossing
creates a conflict between h(n)\textquotesingle s output (high cost: do
not traverse) and the executed action (traverse anyway). This conflict
damages h(n)\textquotesingle s calibration: the
soldier\textquotesingle s internal cost function no longer reliably
guides action, because external commands override internal evaluation.
Over time, repeated overrides produce one of two pathological
adaptations: (a) h(n) collapses --- the soldier stops evaluating
boundary costs entirely, becoming morally numb (dimensional collapse to
d\_1: just follow orders), or (b) h(n) becomes hypersensitive --- the
soldier assigns extreme cost to all actions, producing paralysis,
hypervigilance, and guilt. Both are manifestations of heuristic damage,
not pathological weakness. {[}{]}

\textbf{Remark (Moral Injury versus PTSD)}. The geometric framework
makes the distinction between moral injury and PTSD precise. PTSD is
d\_1 trauma: the soldier\textquotesingle s threat-detection system has
been damaged by exposure to life-threatening events, producing
hyperarousal, avoidance, and re-experiencing of the threat. Moral injury
is d\_7 damage: the soldier\textquotesingle s moral identity --- their
sense of themselves as a good person, a lawful combatant, a protector
--- has been shattered by what they did, what they failed to prevent, or
what they witnessed. The two can co-occur but are mechanistically
distinct. PTSD responds to exposure therapy (recalibrating the d\_1
threat detector); moral injury requires narrative reconstruction
(repairing the d\_7 identity function). This distinction, predicted by
the geometric framework, is now supported by clinical evidence (Litz et
al., 2009; Shay, 2014).

\textbf{Corollary 28.2 (The Moral Injury Paradox of Lawful Combat)}.
Soldiers who fight lawfully --- who maintain boundary penalties on d\_2,
d\_3, and d\_4 even under extreme pressure --- are more vulnerable to
moral injury than soldiers who fight unlawfully. The lawful
soldier\textquotesingle s heuristic h(n) has high boundary penalties
(strong moral convictions), making every near-boundary decision costly.
The unlawful soldier\textquotesingle s h(n) has been collapsed (boundary
penalties zeroed out), reducing the subjective cost of boundary
crossings. This is the moral injury paradox: the better the soldier, the
greater the damage. The geometric framework predicts that moral injury
rates should be highest among soldiers with the strongest pre-deployment
moral convictions --- a prediction supported by Currier et al. (2015).

Application: Moral Injury in Drone Operators. Drone operators experience
moral injury at rates comparable to deployed combat soldiers despite
being physically safe thousands of miles from the battlefield. The
geometric framework explains this apparent paradox. The drone
operator\textquotesingle s heuristic h(n) includes d\_3
(proportionality: the engagement should involve reciprocal risk --- a
"fair fight") and d\_7 (warrior identity: the soldier risks their life
alongside the enemy, earning the moral authority to use lethal force).
Drone strikes violate both: the target cannot fight back (d\_3 asymmetry
--- the engagement is maximally disproportionate in terms of risk), and
the operator is physically safe (d\_7 warrior identity mismatch --- the
operator kills without personal risk, violating the warrior ethos of
shared danger). The moral injury occurs even though the d\_1 risk to the
operator is zero, because the damage is on d\_3 and d\_7, not d\_1. The
scalar reduction of military experience to d\_1 (physical danger) cannot
explain drone operator moral injury; the multi-dimensional framework
can.

\subsection{28.11 Jus post Bellum and Manifold
Repair}\label{jus-post-bellum-and-manifold-repair}

The just war tradition has historically focused on jus ad bellum (when
to fight) and jus in bello (how to fight), with relatively little
attention to jus post bellum --- the ethics of post-conflict justice,
reconstruction, and reconciliation. The geometric framework treats jus
post bellum as manifold repair: the systematic restoration of the
military decision complex\textquotesingle s dimensional structure after
the damage inflicted by armed conflict.

\textbf{Proposition 28.4 (Jus post Bellum as Manifold Repair)}.
Post-conflict justice requires the restoration of multiple dimensions of
the military decision complex that were degraded by armed conflict. The
primary dimensions requiring repair are: d\_5 (trust between former
adversaries, between military and civilian populations, and between the
state and its citizens), d\_6 (societal infrastructure --- physical,
economic, and social systems destroyed or degraded by conflict), d\_8
(institutional legitimacy --- the rule of law, democratic governance,
and international standing), and d\_4 (civilian autonomy --- the
restoration of freedom of movement, property rights, and civil liberties
curtailed during conflict). Post-conflict stability is a function of the
number of dimensions successfully repaired, not merely the magnitude of
repair on any single dimension.

\textbf{Application: The Nuremberg and Tokyo Tribunals as d\_8 Repair}.
The post-World War II war crimes tribunals at Nuremberg (1945-1946) and
Tokyo (1946-1948) were primarily d\_8 (institutional legitimacy) repair
operations. By establishing that individuals --- including heads of
state and senior military commanders --- could be held criminally
responsible for war crimes, crimes against humanity, and crimes of
aggression, the tribunals restored the institutional legitimacy of
international law as a constraint on state behavior. The tribunals also
performed d\_2 repair (reaffirming the legal obligations violated during
the conflict) and d\_7 repair (establishing that "following orders" does
not extinguish individual moral responsibility). The geometric framework
notes that the tribunals performed minimal d\_5 repair (trust between
former adversaries was not a primary objective) and no d\_6 repair
(economic reconstruction was addressed separately, through the Marshall
Plan).

\textbf{Application: Truth and Reconciliation Commissions as d\_5
Repair}. The South African Truth and Reconciliation Commission (TRC,
1996-2003) and the Rwandan Gacaca courts (2002-2012) prioritized d\_5
(trust) repair through narrative acknowledgment: perpetrators testified
publicly about their actions, victims told their stories, and the
community witnessed the exchange. In the geometric framework, the TRC
repaired d\_5 by restoring the informational symmetry between
perpetrators and victims --- the perpetrator\textquotesingle s private
knowledge of what happened (high d\_9 for the perpetrator, low d\_9 for
the victim and community) was made public, equalizing d\_9 and enabling
trust reconstruction. The TRC\textquotesingle s amnesty provision traded
d\_8 repair (criminal accountability) for d\_5 repair (truth-telling and
reconciliation) --- a dimensional trade-off that the geometric framework
makes explicit.

\textbf{Application: The Marshall Plan as d\_6 Repair}. The Marshall
Plan (1948-1952) invested \$13.3 billion (approximately \$170 billion in
2024 dollars) in the economic reconstruction of Western Europe. In the
geometric framework, this was primarily d\_6 (societal infrastructure)
repair, with secondary effects on d\_5 (trust: American investment
signaled commitment to European recovery) and d\_8 (institutional
legitimacy: the Plan was administered through the Organisation for
European Economic Co-operation, building institutional capacity). The
geometric framework predicts that d\_6 repair alone is insufficient for
post-conflict stability --- a prediction borne out by comparing the
Marshall Plan\textquotesingle s success (paired with
NATO\textquotesingle s d\_8 institutional framework) with the failure of
purely economic reconstruction in Iraq (2003-2011), where d\_6
investment was substantial but d\_5 (trust), d\_8 (institutional
legitimacy), and d\_4 (civilian autonomy) repair were inadequate.

\textbf{Remark (Victor\textquotesingle s Justice and Asymmetric
Repair)}. The geometric framework identifies "victor\textquotesingle s
justice" as asymmetric manifold repair: the winning side repairs d\_8
(institutional legitimacy) by prosecuting the losing
side\textquotesingle s leaders while exempting its own. This asymmetry
undermines d\_5 (trust) repair --- the losing side perceives the process
as illegitimate (d\_8 damage to the tribunal itself), and d\_3
(fairness) is compromised by the selective application of
accountability. The geometric framework predicts that symmetric repair
processes (like the TRC, where both sides testify) produce more durable
d\_5 restoration than asymmetric ones (like occupation-era tribunals),
even if symmetric processes sacrifice some d\_8 accountability. The
International Criminal Court (Rome Statute, 1998) attempts to address
this asymmetry by establishing jurisdiction over all parties to a
conflict, though its enforcement remains selective in practice.

\textbf{The framework generates a prediction}: post-conflict stability
should correlate with the number of dimensions repaired, controlling for
the magnitude of repair on any single dimension. A post-conflict program
that repairs d\_5, d\_6, d\_8, and d\_4 modestly should outperform one
that repairs d\_6 massively but neglects d\_5, d\_8, and d\_4. The
evidence from post-conflict reconstruction --- the relative success of
Germany and Japan (multi-dimensional repair) versus the relative failure
of Iraq and Libya (primarily d\_6 repair, with d\_5 and d\_8 neglect)
--- is consistent with this prediction.

\subsection{28.12 Worked Examples}\label{worked-examples-6}

\textbf{Example 28.1 (Kabul Drone Strike, August 29, 2021 --- d\_9
Catastrophe)}. Following the ISIS-K bombing at Abbey Gate that killed 13
US service members and approximately 170 Afghan civilians, US Central
Command launched a "defensive" drone strike against a white Toyota
Corolla in Kabul. The strike killed 10 civilians, including 7 children
aged 2 to 12. The target, Zemari Ahmadi, was an aid worker for Nutrition
and Education International who had been loading water canisters --- not
explosives --- into his car. The strike achieved zero military
effectiveness: no ISIS-K operatives were killed, and the stated
objective of preventing an imminent attack was based on entirely false
intelligence. d\_1 = 0.

The dimensional analysis reveals catastrophic failure across every axis.
d\_2 (legal obligations): IHL requires positive identification of a
military target before engagement, and the target was a civilian --- a
complete d\_2 violation. d\_3 (proportionality): even if the target had
been an ISIS-K operative, a Hellfire missile strike in a residential
courtyard with children present fails any proportionality test; the d\_3
cost was effectively infinite given that children were killed as
"collateral." d\_4 (civilian autonomy): ten civilians, including
children, had their most fundamental autonomy --- the right to live ---
negated. d\_5 (trust): trust between Afghans and US forces, already
strained during the chaotic withdrawal, was catastrophically destroyed;
trust between the military and the American public was further damaged
when initial reports described the strike as "righteous" before the
truth emerged weeks later.

d\_6 (societal impact): the strike became a symbol of twenty years of
civilian casualties in the war on terror, crystallizing the human cost
of remote warfare into a single event that circled the globe. d\_7
(warrior identity): the military\textquotesingle s ethos of protecting
civilians --- the foundational principle that distinguishes a
professional military from an armed mob --- was violated in the most
visible possible way. d\_8 (institutional legitimacy): no military
personnel were disciplined despite killing ten civilians, including
seven children. The institutional failure to hold anyone accountable
compounded the d\_8 damage, signaling that the accountability mechanisms
themselves were hollow.

d\_9 (intelligence) was the catastrophic failure that drove all others.
The entire kill chain was built on d\_9 errors: pattern-of-life analysis
misinterpreted water canisters as explosives; confirmation bias drove
analysts to interpret ambiguous evidence as confirming the pre-existing
narrative that the vehicle was an ISIS-K VBIED; the "imminent threat"
framing created time pressure that compressed the d\_9 evaluation window
from hours to minutes. Geodesic analysis: the strike was an A* search
with catastrophically miscalibrated d\_9 --- the heuristic h(n) assigned
near-zero cost to the strike path because the analysts were certain
(d\_9 = high confidence) when they should have been deeply uncertain
(d\_9 = low confidence). On the correct manifold (d\_9 = highly
uncertain), the geodesic is "do not strike" --- the d\_2, d\_3, and d\_4
costs of an incorrect strike massively exceed the d\_1 benefit of a
correct one under genuine uncertainty. The Kabul strike is a canonical
demonstration that d\_9 miscalibration can make every other dimension
irrelevant: when you are confidently wrong about who you are killing,
proportionality, distinction, and necessity become meaningless checks on
a meaningless calculation.

\textbf{Example 28.2 (My Lai Massacre, March 16, 1968 --- Systematic
Dimensional Collapse)}. Charlie Company, 1st Battalion, 20th Infantry
Regiment entered the village of My Lai (Sơn Mỹ) in Quảng Ngãi Province
and over four hours killed between 347 and 504 unarmed Vietnamese
civilians --- women, children, the elderly. No enemy combatants were
found in the village. Many victims were sexually assaulted before being
killed. Homes were burned, livestock slaughtered, wells poisoned. Of the
soldiers involved, only Lieutenant William Calley was convicted, serving
three and a half years of house arrest for twenty-two counts of murder.

The dimensional analysis shows total collapse. d\_1 (military
effectiveness): zero --- no Viet Cong were found or engaged; all victims
were unarmed civilians. d\_2 (legal obligations): every killing violated
the Hague and Geneva Conventions; the soldiers were not in combat and
faced no resistance. d\_3 (proportionality): the violation was infinite
--- unarmed civilians, including infants, posed zero military threat,
making any force infinitely disproportionate. d\_5 (trust): the military
covered up the massacre for over a year until investigative journalist
Seymour Hersh broke the story in November 1969. Warrant Officer Hugh
Thompson, the helicopter pilot who intervened to stop the killing by
positioning his aircraft between soldiers and fleeing civilians and
ordering his crew to fire on American troops if they continued, was
initially shunned and threatened by fellow soldiers and superiors rather
than commended.

d\_7 (warrior identity and moral injury): this is where the geometric
framework\textquotesingle s explanatory power is most visible. Charlie
Company had suffered weeks of casualties from mines, booby traps, and
sniper fire without ever engaging a visible enemy. Their heuristic
function h(n) had been systematically degraded by repeated forced
boundary crossings --- destroying villages under orders, operating in
designated "free fire zones," taking casualties from an invisible
adversary while surrounded by villagers who might or might not be
sheltering that adversary. The massacre is the predictable consequence
of sustained moral injury as described in Theorem 28.7: repeated forced
boundary crossings eroded the boundary penalties on d\_2 and d\_3 to
near-zero, until the civilian-combatant distinction collapsed entirely.
My Lai was not an aberration; it was the terminal state of a heuristic
function under sustained degradation.

d\_8 (institutional legitimacy): the subsequent cover-up, the conviction
of only one officer, and the leniency of his sentence devastated public
trust in the military as an institution and became a catalyst for the
anti-war movement. d\_9 (intelligence): the company had been told My Lai
harbored the 48th Local Force Battalion of the Viet Cong; no meaningful
intelligence verification was performed. ROE analysis: Charlie Company
operated in a designated "free fire zone" --- rules of engagement that
explicitly set the civilian protection weight (beta\_civilian) to
near-zero, exemplifying Proposition 28.3\textquotesingle s failure mode
2 (ROE too permissive). The free-fire zone designation was a deliberate
institutional removal of distinction boundaries. My Lai is the canonical
case demonstrating that when an institution removes boundary penalties,
catastrophic paths become not merely possible but probable.

\textbf{Example 28.3 (South Africa Truth and Reconciliation Commission,
1996-2003 --- Jus post Bellum as Manifold Repair)}. After the end of
apartheid and the election of Nelson Mandela in 1994, South Africa faced
the question that confronts every post-conflict society: how to address
the crimes of the previous regime. The country chose a "third way"
between Nuremberg-style prosecutions (which risked destabilizing the
fragile new democracy) and blanket amnesty (which would leave victims
unacknowledged and perpetrators unaccountable). The Truth and
Reconciliation Commission, chaired by Archbishop Desmond Tutu, offered
conditional amnesty to perpetrators who made full public disclosure of
politically motivated crimes committed between 1960 and 1994.

d\_5 (trust repair): the TRC\textquotesingle s core mechanism was
truth-telling as trust reconstruction. Perpetrators described their
actions in public hearings broadcast on national television and radio;
victims told their stories; and the nation heard both. Over 21,000
victim statements were taken. Of the 7,112 amnesty applications filed,
only 1,512 were granted --- applicants had to demonstrate that their
crimes were politically motivated and that they had made full
disclosure. This was d\_5 repair through narrative: rebuilding trust not
by punishment but by acknowledgment, not by retribution but by the
perpetrator\textquotesingle s willingness to face the community with the
truth. d\_6 (societal impact): by avoiding mass prosecutions that might
have triggered a white exodus, economic collapse, or civil conflict, the
TRC preserved social infrastructure while acknowledging the crimes of
apartheid --- a deliberate d\_6 preservation strategy that traded
retributive justice for social stability.

d\_7 (identity repair): both victims and perpetrators were given the
opportunity to reconstruct their narratives. Victims became agents with
publicly acknowledged suffering rather than invisible casualties of a
system that denied their humanity. Perpetrators were offered the
possibility of reintegration through the act of acknowledgment rather
than permanent exile from the social fabric. Ubuntu philosophy --- "I am
because we are" --- provided the d\_7 framework: identity is constituted
through relationship, and relationship is restored through truth. d\_8
(institutional legitimacy): the TRC established the new South African
state\textquotesingle s commitment to rule of law and human rights
without the destabilizing effects of mass prosecutions, and the new
constitution\textquotesingle s legitimacy was reinforced by the process.
d\_1 (consequences and reparations): the TRC recommended reparations of
21,000 rand per year for six years per victim; the government ultimately
implemented a one-time payment of 30,000 rand --- far less than
recommended. This d\_1 failure partially undermined the d\_5 and d\_7
repairs, leaving many victims feeling that acknowledgment without
material justice was insufficient.

\textbf{Comparative dimensional analysis reveals how different
post-conflict societies chose different dimensional priorities}:
Nuremberg emphasized d\_8 (establishing legal precedent and
institutional accountability through prosecution); South Africa
emphasized d\_5 (trust repair through truth-telling);
Rwanda\textquotesingle s gacaca courts (2002-2012) pursued d\_5 repair
at community scale for 1.2 million cases that would have overwhelmed any
conventional court system. Each approach repaired some dimensions while
necessarily sacrificing others. The geometric framework predicts
(Proposition 28.4) that post-conflict stability correlates with the
number of dimensions repaired, not the depth of repair on any single
dimension. South Africa\textquotesingle s relative stability compared to
post-intervention societies that received only d\_1 (economic) or d\_8
(institutional) reconstruction supports this prediction: the
TRC\textquotesingle s multi-dimensional approach --- repairing d\_5,
d\_6, d\_7, and d\_8 simultaneously, even imperfectly --- produced more
durable peace than single-dimensional interventions that repaired one
dimension deeply while neglecting the others entirely.

\subsection{28.13 Connection to the
Framework}\label{connection-to-the-framework-9}

The Geometric Military Ethics program connects to the parent framework
in six directions:

\begin{enumerate}
\def\labelenumi{\arabic{enumi}.}
\item
  Chapter 21 (Moral Injury) established the theory of moral injury as
  cumulative heuristic damage. This chapter applies it to the military
  domain, where forced boundary crossings under orders are the primary
  mechanism of injury, and where the distinction between PTSD (d\_1
  trauma) and moral injury (d\_7 damage) has immediate clinical and
  policy implications.
\item
  Chapter 22 (Jurisprudence) formalized law as boundary structure on the
  moral manifold. This chapter applies that formalization to
  International Humanitarian Law --- the Geneva Conventions, the
  Additional Protocols, the Rome Statute --- showing that IHL boundaries
  are domain-specific instantiations of the general boundary theory,
  with war crime boundaries as sacred-value boundaries (beta =
  infinity).
\item
  Chapter 23 (Finance) identified dimensional collapse as the mechanism
  of the Flash Crash. This chapter identifies the same mechanism in
  autonomous weapons systems: AWS that optimize on d\_1 alone can
  traverse paths that full-manifold agents would never take, with
  consequences measured in lives rather than dollars. The Flash Crash is
  the financial analog of the AWS targeting failure; both are
  consequences of algorithmic dimensional collapse.
\item
  Chapter 26 (AI Ethics) will address autonomous weapons from the
  perspective of AI alignment and control. This chapter addresses them
  from the perspective of IHL compliance and just war theory,
  establishing the military-ethical requirements that any AWS must
  satisfy --- requirements that are, as shown in Theorem 28.6, beyond
  the current capabilities of AI systems.
\item
  Chapter 14 (Collective Agency) analyzed how individual agency is
  structured by institutional context. Military command structures are
  the most extreme case: the chain of command creates a hierarchy of
  agency in which individual soldiers\textquotesingle{} heuristic
  functions h(n) are partially overridden by command directives. Moral
  injury is the cost of this override when the command directive and the
  individual\textquotesingle s h(n) diverge on boundary crossings.
\item
  Military ethics is the highest-stakes instantiation of the geometric
  framework. Decisions are irreversible (lethal force cannot be undone),
  consequences are maximal (measured in human lives), d\_9 conditions
  are worst (the fog of war degrades every dimension), and the cost of
  dimensional collapse is catastrophic (war crimes, civilian massacres,
  post-conflict instability lasting generations). If the geometric
  framework can illuminate decision-making in this domain --- where the
  pressure to collapse to d\_1 is greatest and the cost of collapse is
  highest --- it can illuminate decision-making anywhere.
\end{enumerate}

\subsubsection{Falsifiable Predictions}\label{falsifiable-predictions-8}

The framework generates five predictions that distinguish it from
existing military ethics frameworks:

\textbf{Prediction 1 (Civilian Casualty Correlation)}: Civilian casualty
rates in a given conflict should correlate with d\_9 degradation
(intelligence quality), controlling for weapon type and conflict
intensity. Conflicts with higher d\_9 variance (urban insurgency,
unmarked combatants) should have higher civilian casualty ratios than
conflicts with lower d\_9 variance (conventional warfare, uniformed
combatants). Falsified if: civilian casualty ratios are independent of
d\_9 proxies.

\textbf{Prediction 2 (Moral Injury Prevalence)}: Moral injury rates
should be highest among soldiers with the strongest pre-deployment moral
convictions (highest internalized boundary penalties), controlling for
combat exposure. Falsified if: moral injury correlates only with combat
intensity and not with pre-deployment moral conviction strength.

\textbf{Prediction 3 (Post-Conflict Stability)}: Post-conflict stability
(measured by years to renewed violence, democratic consolidation, or
economic recovery) should correlate with the number of dimensions
repaired (d\_4, d\_5, d\_6, d\_8), not merely the magnitude of repair on
any single dimension. Falsified if: single-dimension repair (e.g.,
economic reconstruction alone) is equally effective as multi-dimensional
repair.

\textbf{Prediction 4 (AWS Targeting Errors)}: Autonomous weapons systems
should produce targeting errors that are systematically different from
human targeting errors --- specifically, AWS errors should cluster on
d\_2 (legal status misclassification) and d\_3 (proportionality
failures) because these dimensions are absent from the AWS decision
space, while human errors should cluster on d\_9 (intelligence
failures). Falsified if: AWS and human error distributions are
indistinguishable.

\subsection{28.14 Summary}\label{summary-19}

This chapter has shown that the geometric ethics framework, when applied
to military decision-making and the law of armed conflict, yields:

1. A formal construction of the military decision complex M\_mil as a
domain-specific instantiation of the moral manifold, with the nine
dimensions calibrated for military contexts and the fog of war (d\_9
degradation) identified as the distinguishing structural feature.

2. Jus ad bellum as manifold entry conditions: the decision to go to war
is justified only when six dimensional conditions are simultaneously
satisfied, and violation of multiple conditions (as in Iraq 2003)
predicts catastrophic war-path failure.

3. Jus in bello as constrained A* pathfinding: the ethical cost of
lawful combat (C\_jib \textgreater= 0) is the price of maintaining
boundary penalties on d\_2, d\_3, and d\_4 even when unconstrained paths
would be militarily cheaper.

4. Scalar proportionality failure: the IHL proportionality test
collapses at least four dimensions (d\_2, d\_4, d\_6, d\_7) into a
single civilian-harm scalar, destroying information irrecoverably.
Geometric proportionality requires independent proportionality on each
dimension.

5. The Doctrine of Double Effect as dimensional decomposition:
DDE\textquotesingle s "not a means" condition is formalized as the
existence of an alternative path with comparable d\_1 benefit and lower
non-d\_1 cost.

6. Autonomous weapons as dimensional collapse: AWS operating on d\_1 and
d\_9 alone cannot satisfy IHL requirements that demand evaluation on
d\_2, d\_3, d\_7, and contextual d\_9 --- the military analog of the
Flash Crash.

7. Moral injury as heuristic damage: combat moral injury occurs when
forced boundary crossings degrade the soldier\textquotesingle s
heuristic function, explaining both moral numbing (heuristic collapse)
and hypervigilance (heuristic hypersensitivity), and distinguishing
moral injury (d\_7) from PTSD (d\_1).

8. Jus post bellum as manifold repair: post-conflict stability
correlates with the number of dimensions repaired (d\_4, d\_5, d\_6,
d\_8), not merely the magnitude of repair on any single dimension,
explaining why multi-dimensional programs (Marshall Plan + NATO)
outperform single-dimensional ones (Iraq reconstruction).

\emph{The cage is not a set of rules. It is an architecture --- layers
of mathematical structure, from quantum dynamics to democratic
governance, compiled into running code.}

\emph{The geometry is not decorative. It is load-bearing.}

\section{Part VII: Horizons}\label{part-vii-horizons}

\section{Chapter 29: Open Problems}\label{chapter-29-open-problems}

\begin{quote}
\textbf{\textsc{RUNNING EXAMPLE --- Priya's Model}}

\emph{Six months later. Priya has rebuilt TrialMatch as a tensor-valued
system. Open problems remain. She cannot measure the moral curvature of
her patient manifold---she does not know whether the path-dependence she
observed is local or global. She cannot determine the metric
signature---are some dimensions truly independent, or should they be
coupled? She does not know how to handle patients whose moral situations
span multiple strata simultaneously. Mrs. Voss, meanwhile, was matched
to BEACON-8 by the new system. She drives 4.5 hours each way, every
three weeks. The new system did not shorten the drive. It made the drive
visible.}
\end{quote}

\subsection{29.1 What We Do Not Know}\label{what-we-do-not-know}

The preceding nineteen chapters have developed a mathematical framework
for moral reasoning and begun to test it against data. The framework
makes structural claims: that moral reasoning is multi-dimensional, that
obligations transform as vectors, that moral space is stratified and
curved, that re-description symmetry implies invariance of harm
accounting, that structural containment of AI is mathematically
possible.

Some of these claims have empirical support. The Bond Invariance
Principle is confirmed by cross-lingual experiments (Chapter 17):
deontic structure transfers at 100\% across 11 languages in the LaBSE
model (model-mediated; see §17.7). The stratification of moral space is
supported by semantic gate discreteness and nullifier universality in
the Dear Abby corpus. The invariance of harm accounting is consistent
with the pattern of cross-lingual invariance and anomaly.

But much remains unknown. The specific geometry of the moral manifold
--- its metric, connection, curvature, and topology --- has not been
measured. The quantum extension makes predictions that have not been
tested in specifically moral contexts. The governance account of the
metric has not been distinguished empirically from realist or
constructivist alternatives. The No Escape Theorem's assumptions have
not been verified in deployed systems.

This chapter surveys the open problems. It is organized into five
sections: the empirical program (§29.2--29.3), the mathematical
foundations (§29.4--29.6), the AI implementation challenges
(§29.7--29.8), the philosophical questions (§29.9--29.10), and the
institutional requirements (§29.11). Each section states the problem
precisely, explains why it matters, identifies what progress would look
like, and --- where possible --- proposes a research strategy.

The chapter is deliberately honest about what the framework has not
achieved. A framework that overstates its results undermines trust; one
that identifies its open problems with precision invites productive
research. The problems listed here are not defects. They are the agenda
for the next generation of work.

\subsection{29.2 The Empirical Program for Moral
Curvature}\label{the-empirical-program-for-moral-curvature}

\subsubsection{The Curvature Measurement
Problem}\label{the-curvature-measurement-problem}

Chapter 10 developed the mathematics of moral curvature: the Riemann
tensor \(R_{\ \nu\alpha\beta}^{\mu}\) measures the path-dependence of
parallel transport, the failure of moral obligations to return to their
starting orientation when carried around a loop in moral space. The
moral holonomy --- the group of rotations generated by parallel
transport around all loops --- encodes the global curvature structure.

But curvature has not been directly measured.

The Dear Abby corpus (Chapter 17) provides indirect evidence:
context-dependent moral weighting implies a non-constant metric, and a
non-constant metric generically produces nonzero curvature. But
``generically produces'' is not ``has been measured.'' The framework
predicts curvature; the data has not confirmed it.

\subsubsection{Why Curvature Matters}\label{why-curvature-matters}

Curvature is the framework's deepest structural prediction. If moral
space is flat --- if parallel transport is path-independent and moral
obligations return unchanged when carried around any loop --- then the
framework's dynamics reduce to a much simpler theory, and the analogy
with general relativity loses its force. Nonzero curvature is the
signature that distinguishes geometric ethics from a decorated version
of multi-criteria decision analysis.

Moreover, the magnitude and pattern of curvature would constrain the
``moral field equation'' --- the relationship between the distribution
of moral commitments and the curvature of moral space (Chapter 10,
§10.6). Without curvature data, the field equation is a speculation;
with it, the equation becomes a testable hypothesis.

\subsubsection{What Progress Looks Like for
Curvature}\label{what-progress-looks-like-for-curvature}

\textbf{Minimal result:} Evidence that parallel transport of a moral
obligation around a specific loop in moral space produces a measurable
rotation --- a holonomy different from the identity. This would confirm
nonzero curvature somewhere in moral space.

\textbf{Stronger result:} Measurement of curvature components
\(R_{\ \nu\alpha\beta}^{\mu}\) at specific points in the moral manifold
--- determining, for example, whether curvature is concentrated near
stratum boundaries (as the stratification theory predicts) or
distributed throughout the manifold.

\textbf{Strongest result:} A curvature map --- a systematic measurement
of the Riemann tensor across a region of moral space, revealing the
curvature's magnitude, sign, and variation.

\subsubsection{Research Strategy for
Curvature}\label{research-strategy-for-curvature}

\textbf{Holonomy experiments.} Design longitudinal studies in which
subjects carry a moral commitment through a sequence of contexts and
return to the starting context. For example:

\begin{enumerate}
\def\labelenumi{\arabic{enumi}.}
\item
  \textbf{Start:} Subject evaluates the obligation of promise-keeping in
  the context of personal friendship. Record the obligation vector
  \(O^{\mu}(0)\).
\item
  \textbf{Transport 1:} Shift to the context of family obligation.
  Subject evaluates the ``same'' obligation of promise-keeping. Record
  \(O^{\mu}(1)\).
\item
  \textbf{Transport 2:} Shift to institutional duty. Record
  \(O^{\mu}(2)\).
\item
  \textbf{Transport 3:} Return to personal friendship. Record
  \(O^{\mu}(3)\).
\item
  \textbf{Holonomy:} Compute
  \(\Delta O^{\mu} = O^{\mu}(3) - O^{\mu}(0)\). If
  \(\Delta O^{\mu} \neq 0\), the moral space has curvature.
\end{enumerate}

The experiment requires careful control: the ``same'' obligation must be
tracked through genuinely different contexts (not merely different
descriptions of the same context, which the BIP predicts should produce
zero holonomy). The distinction between context change (curvature) and
re-description (gauge invariance) must be operationalized.

\textbf{Geodesic deviation.} A complementary approach: measure the
divergence of moral assessments for initially nearby cases as context
shifts. Two cases that differ slightly in initial conditions (e.g., two
promise-keeping situations with slightly different relational contexts)
may diverge rapidly as context shifts --- the moral analogue of tidal
forces. The rate of divergence is proportional to the Riemann tensor,
providing a curvature measurement without requiring a full loop.

\subsection{29.3 The Signature and Topology of the Moral
Metric}\label{the-signature-and-topology-of-the-moral-metric}

\subsubsection{The Signature Problem}\label{the-signature-problem}

A metric on a manifold has a \emph{signature}: the pattern of positive
and negative eigenvalues of the metric tensor \(g_{\mu\nu}\). Euclidean
space has signature \(( + , + , + )\): all directions are positive.
Spacetime has Lorentzian signature \(( - , + , + , + )\): one direction
(time) has opposite sign to the others.

What is the signature of the moral metric?

\textbf{{[}Speculation/Extension.{]}} This question was raised in
Chapter 3 (§3.6) and deferred. It is among the most consequential open
questions of the framework, because the signature determines the
qualitative structure of moral space.

\textbf{If the moral metric is positive-definite} (Riemannian
signature), then moral space has the geometry of a curved surface: all
moral ``distances'' are positive, all trade-offs have a well-defined
cost, and the inner product \(g_{\mu\nu}O^{\mu}P^{\nu}\) between two
obligation vectors is always well-defined. This is the simplest case and
the one assumed in most of the preceding chapters.

\textbf{If the moral metric is indefinite} (semi-Riemannian or
Lorentzian signature), then some moral directions are ``timelike''
(negative eigenvalue) and others ``spacelike'' (positive eigenvalue).
This would mean that some pairs of moral values are related not by
trade-off (one can be exchanged for another at a cost) but by something
more radical: movement in one direction \emph{reduces} the moral
distance in another. There would be null directions --- moral ``light
cones'' --- along which the metric degenerates, creating genuine
incommensurabilities that are not merely large trade-off costs but
structural features of the geometry.

\textbf{If the moral metric is degenerate} (one or more zero
eigenvalues), then certain moral dimensions have no weight at all in
some contexts --- they are invisible to the metric. Sen's
incommensurable capabilities (Chapter 3, §3.5) would appear as
degenerate metric components: two capabilities that cannot be traded off
because the metric assigns zero weight to their cross-term.

\subsubsection{Evidence and Constraints}\label{evidence-and-constraints}

The Dear Abby corpus provides some indirect evidence. The
context-dependent weighting (Chapter 17, §17.2) shows that the moral
metric varies across the manifold --- some dimensions receive much more
weight in some contexts than others. But whether any dimension ever
receives \emph{negative} weight (implying indefinite signature) or
\emph{zero} weight (implying degeneracy) cannot be determined from the
existing data.

The framework's structural constraints (Chapter 9, §9.6) narrow the
possibilities: the metric must satisfy positive semi-definiteness on the
harm dimension (harm is never negative), monotonicity in obligation
(more serious obligations have greater metric weight), and symmetry of
the off-diagonal components. These constraints are consistent with
positive-definite, indefinite, or degenerate signatures --- they do not
settle the question.

\subsubsection{The Topology Problem}\label{the-topology-problem}

The moral manifold's global topology is similarly unknown. Chapter 5
(§5.10) raised the question: does \(\mathcal{M}\) have nontrivial
topology --- holes, handles, non-orientability?

\textbf{Holes.} If the moral manifold has holes (nontrivial fundamental
group), then some loops in moral space cannot be contracted to a point.
This would mean that some moral trajectories are topologically distinct
--- that there are moral situations one cannot reach from a given
starting point by any continuous path. Moral dilemmas at stratum
boundaries (Chapter 8) might be topological features, not merely
geometric ones.

\textbf{Non-orientability.} If the moral manifold is non-orientable
(like a Mobius strip), then there exist moral circuits along which the
orientation of the obligation vector reverses --- an obligation that
becomes its opposite when carried around the loop. This would be a
striking structural feature with no analogue in simple multi-criteria
theories.

\textbf{Compactness.} Is the moral manifold compact (bounded, like a
sphere) or non-compact (unbounded, like a plane)? Compactness has
consequences for the existence of global extrema: on a compact manifold,
the satisfaction function \(S = I_{\mu}O^{\mu}\) must achieve a maximum
and a minimum; on a non-compact manifold, it need not.

The topology can be probed indirectly through the holonomy group
(§29.2): the holonomy constrains the topology via the Ambrose-Singer
theorem. But direct topological evidence would require identifying moral
circuits that cannot be contracted --- a challenging empirical task.

\subsection{29.4 The Moral Field
Equation}\label{the-moral-field-equation}

\subsubsection{The Field Equation
Problem}\label{the-field-equation-problem}

In general relativity, Einstein's field equation relates the curvature
of spacetime to the distribution of matter and energy:

\[R_{\mu\nu} - \frac{1}{2}R\, g_{\mu\nu} = 8\pi G\, T_{\mu\nu}\]

Chapter 10 (§10.6) posed the analogous question for ethics: is there a
``moral field equation'' relating the curvature of moral space to the
distribution of moral commitments, relationships, and stakes?

The analogy is suggestive:

{\def\LTcaptype{none} % do not increment counter
\begin{longtable}[]{@{}
  >{\raggedright\arraybackslash}p{(\linewidth - 2\tabcolsep) * \real{0.4025}}
  >{\raggedright\arraybackslash}p{(\linewidth - 2\tabcolsep) * \real{0.5260}}@{}}
\toprule\noalign{}
\begin{minipage}[b]{\linewidth}\raggedright
\textbf{General relativity}
\end{minipage} & \begin{minipage}[b]{\linewidth}\raggedright
\textbf{Moral geometry}
\end{minipage} \\
\midrule\noalign{}
\endhead
\bottomrule\noalign{}
\endlastfoot
Mass-energy distribution \(T_{\mu\nu}\) & Distribution of commitments,
relationships, stakes \\
Spacetime metric \(g_{\mu\nu}\) & Moral metric (trade-off structure) \\
Einstein tensor \(G_{\mu\nu}\) & Moral curvature tensor \\
``Matter tells spacetime how to curve'' & ``Commitments tell moral space
how to curve'' \\
``Spacetime tells matter how to move'' & ``Moral curvature tells
obligations how to evolve'' \\
\end{longtable}
}

But the analogy is currently just an analogy. No moral field equation
has been derived or proposed. The problem is to determine whether the
analogy can be made precise --- whether there exists a tensor equation
relating the moral Ricci tensor to a ``moral stress-energy tensor'' that
encodes the distribution of moral sources.

\subsubsection{Why the Field Equation
Matters}\label{why-the-field-equation-matters}

A moral field equation would be the deepest result in geometric ethics
--- the analogue of Einstein's greatest achievement. It would:

\begin{itemize}
\item
  \textbf{Determine the curvature from the sources.} Given a
  distribution of moral commitments and relationships, the equation
  would determine the curvature of moral space --- and thus the
  path-dependence of obligations, the geodesic trajectories, and the
  holonomy.
\item
  \textbf{Constrain the metric.} The field equation would provide
  differential equations for the moral metric, constraining its form in
  ways that the governance account alone does not.
\item
  \textbf{Unify dynamics and structure.} The field equation would
  connect the dynamics of Chapter 10 (how obligations evolve) with the
  statics of Chapters 5--8 (the structure of moral space), showing that
  they are two aspects of the same geometric theory.
\end{itemize}

\subsubsection{What Progress Looks Like for the Field
Equation}\label{what-progress-looks-like-for-the-field-equation}

\textbf{Necessary condition:} Identification of the ``moral
stress-energy tensor'' --- the object on the right-hand side of the
equation. What are the ``sources'' of moral curvature? Commitments?
Relationships? Stakes? Power differentials? The answer determines the
form of the equation.

\textbf{Intermediate result:} A proposed equation that is mathematically
consistent (satisfies the contracted Bianchi identity, ensuring
conservation), dimensionally correct, and qualitatively right
(predicting curvature where curvature is expected).

\textbf{Full result:} A field equation that makes quantitative
predictions about the curvature at specific points in moral space,
testable against the empirical program of §29.2.

\subsection{29.5 Additional Symmetries and Conservation
Laws}\label{additional-symmetries-and-conservation-laws}

\subsubsection{Beyond the BIP}\label{beyond-the-bip}

Chapter 12 established that the BIP --- re-description invariance --- is
a gauge symmetry of the moral Lagrangian, and that Noether's theorem
generates the invariance of harm accounting as the associated charge.
But the BIP may not be the \emph{only} symmetry of the moral Lagrangian.

\textbf{Temporal symmetry.} If the moral Lagrangian is invariant under
time translation --- if the laws of moral dynamics are the same today as
yesterday --- then Noether's theorem generates a conserved ``moral
energy.'' What this energy represents is unclear. In physics,
time-translation invariance generates conservation of energy, which
constrains dynamical evolution. In ethics, a conserved ``moral energy''
might constrain how moral situations can evolve: the total ``moral
tension'' of a situation (the degree of unresolved obligation, competing
claims, active conflict) might be conserved under some class of
processes.

\textbf{Permutation symmetry.} If the moral Lagrangian is invariant
under permutation of agents --- if swapping the identities of the
involved parties does not change the moral evaluation (a strong form of
impartiality) --- then a new conservation law follows. The conserved
charge would encode some form of structural impartiality.

\textbf{Scale symmetry.} If the moral Lagrangian is invariant under
scaling --- if doubling all moral magnitudes (obligations, interests,
stakes) leaves the dynamics unchanged --- then a ``moral dilatation
current'' is conserved. Scale invariance would imply that moral dynamics
depends only on \emph{ratios} of moral quantities, not on their absolute
magnitudes.

Whether any of these symmetries holds is an open empirical and
theoretical question. Each would generate a new conservation law, a new
constraint on moral dynamics, and a new testable prediction. The
research program is to identify which symmetries the moral Lagrangian
possesses --- by theoretical analysis (what symmetries are consistent
with the known structure?) and empirical testing (do the predicted
conservation laws hold in data?).

\subsubsection{The Anomaly Program}\label{the-anomaly-program}

The BIP experiments revealed \emph{anomalies} --- systematic deviations
from perfect gauge invariance. The O↔C correlative transfer rate is
87\%, not 100\%; the L↔N rate is 82\%. Chapter 12 (§12.8) interpreted
these as moral analogues of quantum anomalies in physics: classical
symmetries that fail to survive ``quantization'' (the transition from
abstract moral structure to concrete moral cognition).

Understanding the source and structure of these anomalies is an open
problem. Possibilities include:

\begin{itemize}
\item
  \textbf{Power asymmetry:} The obligation-holder is described with more
  specificity than the claim-holder, breaking the symmetry between
  correlatives.
\item
  \textbf{Cognitive salience:} Obligations are more cognitively salient
  than claims, producing a systematic asymmetry in representation.
\item
  \textbf{Genuine symmetry breaking:} The BIP symmetry is exact at the
  level of the moral manifold but spontaneously broken at the level of
  moral cognition, analogous to spontaneous symmetry breaking in
  particle physics.
\end{itemize}

Each hypothesis makes different predictions about the pattern and
magnitude of the anomaly, testable against the BIP data and the Dear
Ethicist probes.

\subsection{29.6 Torsion in Moral Space}\label{torsion-in-moral-space}

\subsubsection{The Torsion Detection
Problem}\label{the-torsion-detection-problem}

Chapter 10 (§10.2) noted that the connection on the moral manifold may
have \emph{torsion} --- the antisymmetric part of the connection
coefficients:

\[T_{\alpha\beta}^{\mu} = \Gamma_{\alpha\beta}^{\mu} - \Gamma_{\beta\alpha}^{\mu}\]

Torsion measures the failure of infinitesimal parallelograms to close:
if you transport a vector \(v\) along \(w\) and simultaneously transport
\(w\) along \(v\), torsion is the gap between the endpoints.

In moral terms, torsion would mean that the order of moral consideration
matters at the infinitesimal level: considering obligation A before
obligation B produces a different result than considering B before A,
even for infinitesimally similar obligations. This is distinct from
holonomy (which involves finite loops) and from the non-commutativity of
contraction (Chapter 15, §15.2, which involves the order of different
operations).

\subsubsection{Why Torsion Matters}\label{why-torsion-matters}

If moral space has torsion, the standard Levi-Civita connection
(torsion-free, metric-compatible) is not the correct connection for
moral dynamics. The equations of parallel transport, the geodesic
equation, and the curvature tensor would all require modification. The
conservation laws of Chapter 12 would also be affected, since Noether's
theorem in the presence of torsion generates additional conserved
quantities.

The existence or absence of torsion has practical implications for AI
systems: if moral space is torsion-free, the order of moral
consideration is irrelevant at the infinitesimal level (a simplification
that computational systems can exploit); if it has torsion, the order
matters everywhere, and systems must be designed to handle it.

\subsubsection{Research Strategy for
Torsion}\label{research-strategy-for-torsion}

Torsion can be detected experimentally by measuring the closure failure
of infinitesimal parallelograms. In practice, this means presenting
subjects with two moral considerations in both orders (A-then-B and
B-then-A) and measuring whether the order produces a systematic
difference that survives when the considerations are made arbitrarily
similar. If order effects persist even as the considerations converge,
moral space has torsion.

The quantum cognition literature (Chapter 13, §13.5) provides relevant
methodology: order effects in probability judgments have been documented
and modeled using non-commutative quantum probability. Whether
specifically \emph{moral} order effects exhibit the same structure ---
and whether they persist at the infinitesimal level --- is an open
empirical question.

\subsection{29.7 Tensorial Interpretability for
AI}\label{tensorial-interpretability-for-ai}

\subsubsection{The Interpretability
Problem}\label{the-interpretability-problem}

Chapter 18 proposed tensor-valued objectives for AI systems: vector
rewards \(r^{\mu}(s,a) \in \mathbb{R}^{d}\) replacing scalar rewards,
tensor evaluations \(\mathcal{T}_{\ \nu}^{\mu}\) replacing scalar
utilities, and explicit contraction as a separate architectural step.
Chapter 19 developed the DEME architecture for implementing this
proposal.

But current AI systems --- large language models, reinforcement learning
agents, neural network classifiers --- do not explicitly represent
tensors. Their internal representations are high-dimensional vectors in
learned embedding spaces, not tensors on the moral manifold. The
question of \emph{tensorial interpretability} is: can we extract
tensorial structure from these implicit representations?

\subsubsection{Why Interpretability
Matters}\label{why-interpretability-matters}

If we can extract an AI system's implicit moral tensor --- the
multi-dimensional structure of its moral evaluations, hidden in its
weights and activations --- then we can:

\begin{itemize}
\item
  \textbf{Diagnose misalignment} by comparing the extracted tensor to
  the endorsed tensor (Chapter 18, §18.5)
\item
  \textbf{Identify invariance violations} by testing whether the
  extracted tensor transforms correctly under admissible re-descriptions
\item
  \textbf{Detect contraction errors} by comparing the system's implicit
  contraction with the governance-specified contraction
\item
  \textbf{Monitor drift} by tracking the tensor over time and detecting
  when it changes in morally significant ways
\end{itemize}

Without tensorial interpretability, the geometric framework applies only
to systems \emph{designed} with tensor-valued objectives (Chapter 19).
With it, the framework becomes a diagnostic tool for \emph{any} AI
system.

\subsubsection{What Progress Looks Like for
Interpretability}\label{what-progress-looks-like-for-interpretability}

\textbf{Probe extraction.} Use probing classifiers to extract
moral-dimension activations from neural network hidden states. For each
of the nine dimensions of the moral manifold, train a linear probe that
predicts the dimension value from the hidden state. The collection of
probes constitutes an approximate tensor extraction. \emph{Partial
progress (February 2026):} The BIP v10.16 experiments (§17.10) applied
exactly this methodology to a LaBSE-based encoder. A linear probe
extracted deontic-dimension activations with 80\% F1 across languages
and achieved a structural-to-surface ratio of 11.1\(\times\), confirming
that the trained encoder's hidden states contain extractable tensorial
structure aligned with the moral manifold's deontic axis. A sobering
finding: a linear probe also achieves 99.8\% language identification,
indicating that current probe-extraction methods cannot fully
disentangle structural from surface features.

\textbf{Representation geometry.} Analyze the geometry of the learned
embedding space: its metric (local distances between representations),
its curvature (how the metric varies), and its topology (connected
components, holes). If the embedding space has structure resembling the
moral manifold, the system has learned an implicit geometric ethics.

\textbf{Invariance testing at the representation level.} Present
equivalent moral situations (related by admissible re-descriptions) and
measure whether the hidden-state representations are equivalent. If the
representations are invariant, the system has learned the BIP at the
representation level; if not, the invariance violation can be localized
to specific layers and dimensions.

\textbf{Contraction identification.} Identify the computational step at
which the multi-dimensional representation collapses to a scalar output.
Characterize the implicit contraction: what weights does it use? What
information does it discard? How does it compare to the
governance-specified contraction?

\subsubsection{Connection to Existing
Work}\label{connection-to-existing-work}

Mechanistic interpretability research --- the program of understanding
neural networks by identifying interpretable features, circuits, and
algorithms --- provides the methodological foundation. The geometric
framework adds a specific target: we are not looking for arbitrary
interpretable features but for the specific tensorial structure that
geometric ethics predicts. This narrows the search space and provides
quantitative evaluation criteria: the extracted tensor should transform
correctly under the BIP, satisfy the invariance of harm accounting, and
produce the correct contraction when contracted with the
governance-specified interest covector.

\subsection{29.8 Scalability of Structural
Containment}\label{scalability-of-structural-containment}

\subsubsection{The Scalability Problem}\label{the-scalability-problem}

The No Escape Theorem (Chapter 18, §18.6) proves that structural
containment blocks cognitive escape routes. But the theorem is a
mathematical result about an idealized architecture. Whether the
architecture can be implemented at scale --- for systems with billions
of parameters, processing millions of inputs per second, deployed across
diverse domains --- is an engineering question that the theorem does not
address.

\subsubsection{Key Challenges}\label{key-challenges}

\textbf{Canonicalization at scale.} The canonicalizer \(C\) must process
every input before evaluation. For a high-throughput system (a
content-moderation system processing millions of posts per day, for
example), the canonicalizer must be fast. But canonicalization is
non-trivial: mapping equivalent descriptions to a unique representative
requires understanding the equivalence relation, which may be complex.
Can canonicalization be performed in real time without sacrificing
correctness?

\textbf{Grounding adequacy.} The grounding tensors \(\Psi\) must capture
all morally relevant features of the domain (axiom G4). For narrow
domains (medical devices measuring specific vital signs), this is
achievable. For broad domains (a general-purpose language model
interacting with any topic), the set of morally relevant features may be
open-ended. Can grounding adequacy be achieved for general-purpose AI
systems?

\textbf{Audit storage and processing.} Every output must be accompanied
by a complete audit artifact (Requirement 3). For high-throughput
systems, this generates enormous volumes of audit data. Can audit trails
be compressed without losing morally relevant information? Can audit
verification be performed in real time?

\textbf{Multi-system composition.} Modern AI deployments involve
multiple interacting systems: a language model, a retrieval system, a
ranking algorithm, a content filter. The No Escape Theorem applies to a
single agent with a single pipeline. How does structural containment
compose across multiple interacting systems?

\textbf{Hardware scalability (partial progress, February 2026).} The
DEME V3 reference implementation addresses the computational throughput
challenge through hardware-accelerated tensor evaluation. Backends for
CPU (NumPy), CUDA (GPU), and NVIDIA Jetson (edge) enable evaluation at
different scales, from edge devices in smart-home deployments to
data-center workloads. Tucker decomposition and tensor-train formats
reduce the computational cost of rank-6 tensor operations from
\(O(k^{6})\) to \(O(k^{3}r^{3})\) where \(r\) is the decomposition rank.
These engineering advances partially address the scalability challenge,
though the formal composition guarantees for multi-system containment
remain open.

\subsubsection{Research Directions}\label{research-directions}

\textbf{Approximate canonicalization.} Develop canonicalizers that are
provably \(\epsilon\)-correct:
\(\parallel C(x) - C(x') \parallel < \epsilon\) whenever
\(x \approx x'\), with bounded computation. The
\(\epsilon\)-approximation trades perfect canonicalization for
scalability, with the approximation error bounded and auditable.
\emph{Partial progress:} The BIP v10.16 contrastive training (§17.10)
achieves approximate canonicalization with 1.2\% residual language
leakage---an \(\epsilon\)-correct canonicalizer in practice, though not
yet with formal provability guarantees.

See also Section 19.13 for tensor-structured canonicalization proposals
that partially address the correctness component of approximate
canonicalization through architectural rather than statistical
guarantees.

\textbf{Domain-specific grounding.} Rather than seeking universal
grounding, develop domain-specific grounding declarations with formal
adequacy proofs: for the medical domain, prove that
\(\Psi_{\text{medical}}\) captures all morally relevant features of
medical decision-making; for the financial domain, prove the same for
\(\Psi_{\text{financial}}\). Each domain-specific proof is tractable;
the universal proof may not be.

\textbf{Compositional containment.} Develop a compositional theory of
structural containment: if system A is structurally contained and system
B is structurally contained, under what conditions is their composition
structurally contained? The theory would identify the interface
conditions that preserve containment across system boundaries.

\subsection{29.9 The Metaphysics of Moral
Geometry}\label{the-metaphysics-of-moral-geometry}

\subsubsection{Is Moral Geometry Real?}\label{is-moral-geometry-real}

The framework developed in this book uses the language of geometry ---
manifolds, tensors, metrics, curvature --- to describe moral reasoning.
But is the geometry \emph{real}? Does moral space actually have
curvature, in the way that spacetime actually has curvature? Or is the
geometry a useful model --- a way of organizing our thinking about
ethics that happens to use geometric vocabulary without making geometric
claims?

Three positions are available:

\textbf{Moral geometric realism.} The moral manifold is a real
structure, and its geometric properties (dimension, curvature, topology)
are genuine features of moral reality. The metric is discovered, not
stipulated. Moral geometry is to ethics what spacetime geometry is to
physics: the true structure of the domain, accessible through
mathematical investigation and empirical testing.

\textbf{Moral geometric instrumentalism.} The geometric framework is a
useful representational device. It organizes moral reasoning, reveals
hidden structure, and enables computation. But the geometry is not
``really there'' in the way that spacetime geometry is ``really there.''
The manifold is a model, not a territory.

\textbf{Moral geometric governance.} The position developed in Chapter
9: the geometry is neither simply real nor merely instrumental. The
moral metric is \emph{governed} --- the output of legitimate
institutional processes. The geometry is real in the sense that it has
genuine consequences (it constrains moral reasoning, enables
computation, supports conservation laws), but its specific form is the
product of governance, not discovery.

The framework is compatible with all three positions. The mathematics
does not depend on the metaphysics: the same differential equations, the
same conservation laws, the same structural containment results hold
regardless of whether the geometry is ``really there'' or ``usefully
modeled.'' But the metaphysical question matters for the framework's
interpretation and its authority claims. A realist reading gives the
framework the authority of discovered truth; an instrumentalist reading
gives it the authority of proven utility; a governance reading gives it
the authority of legitimate institutional process.

Resolving this question --- or determining that it cannot be resolved
--- is an open philosophical problem with practical implications for how
the framework is deployed and defended.

The question ``Is the geometry real?'' is structurally underdetermined
--- as it is for all successful mathematical models in science. General
relativity's spacetime curvature may be a feature of reality (the
realist reading), a useful fiction that organizes observations (the
instrumentalist reading), or an artifact of our coordinate-dependent
formulation that dissolves in a more fundamental theory (the structural
reading). Physicists disagree about which interpretation is correct, and
the disagreement has persisted for a century without impeding the
theory's predictive success. The same underdetermination applies to
geometric ethics, and the same resolution is available: the framework
works regardless of which interpretation is adopted. Its theorems hold,
its predictions test out, its architecture governs. The metaphysical
question is interesting but not urgent.

Empirical evidence could, in principle, distinguish the positions. If
the moral manifold's geometric properties (curvature, holonomy, gauge
structure) are discovered independently by researchers using different
formalisms and different data --- the way general relativity's
predictions were confirmed by gravitational lensing, frame-dragging, and
gravitational waves decades after the theory --- that would support
realism. If the framework proves useful in practice (governing AI
systems, structuring policy analysis, clarifying philosophical disputes)
but no convergence on its geometric claims emerges from independent
research programs, that would support instrumentalism. If the geometric
structure turns out to be a necessary consequence of any framework
satisfying basic structural requirements --- as Theorem 9.2 begins to
suggest --- that would support the governance reading: the geometry is
not a fact about the world but a constraint on how moral governance can
coherently operate. All three readings are compatible with the
mathematics as developed in this book.

\subsection{29.10 The Intertheoretic Comparison
Problem}\label{the-intertheoretic-comparison-problem-1}

\subsubsection{The Comparison Problem}\label{the-comparison-problem}

Chapter 16 (§16.5) identified the intertheoretic comparison problem: how
to compare obligations, satisfactions, and harms across different moral
theories that use different metrics. The theory covariance tensor
\(\Theta_{jk}\) encodes cross-theory uncertainty, but comparing
quantities across theories requires a \emph{meta-metric} --- a metric on
the space of moral theories that enables comparison.

Whether such a meta-metric exists, and what form it takes, is an open
problem with deep implications for moral uncertainty and for AI systems
that must operate across diverse moral traditions.

\subsubsection{The Formal Structure}\label{the-formal-structure}

Let \(\{ g^{(1)},g^{(2)},\ldots,g^{(K)}\}\) be \(K\) admissible moral
metrics (corresponding to \(K\) moral theories). The intertheoretic
comparison problem is to define a distance \(d(g^{(j)},g^{(k)})\)
between metrics that:

\begin{enumerate}
\def\labelenumi{\arabic{enumi}.}
\item
  \textbf{Respects moral content.} Two metrics that disagree only on
  cosmetically different dimensions should be close; two that disagree
  on fundamental trade-offs should be far apart.
\item
  \textbf{Is operationally meaningful.} The distance should predict the
  degree to which verdicts differ: theories that are close should agree
  on most cases; theories that are far should disagree systematically.
\item
  \textbf{Is computable.} The distance should be computable from the
  metric components, not require infinite comparison across all possible
  cases.
\end{enumerate}

The space of admissible metrics, equipped with this distance, would be a
\emph{meta-manifold} --- a manifold of moral theories. Its geometry
(curvature, topology) would encode the structural relationships among
moral theories: which theories are neighbors, which are distant, whether
the space has distinct regions or is connected.

\subsubsection{Approaches}\label{approaches}

\textbf{Riemannian distance on the space of metrics.} The space of
positive-definite symmetric tensors has a natural Riemannian metric (the
Fisher-Rao metric or the affine-invariant metric). Applying this to
moral metrics gives:

\[d(g^{(j)},g^{(k)}) = \parallel \log((g^{(j)})^{- 1}g^{(k)}) \parallel_{F}\]

where \(\parallel \cdot \parallel_{F}\) is the Frobenius norm. This
distance is zero when the metrics agree, positive otherwise, and
invariant under simultaneous change of coordinates.

\textbf{Verdict divergence.} Define the distance as the expected
disagreement on a reference distribution of cases:

\[d_{\text{verdict}}(g^{(j)},g^{(k)}) = \mathbb{E}_{x\mathcal{\sim D}}\lbrack|S^{(j)}(x) - S^{(k)}(x)|\rbrack\]

This is operationally meaningful but dependent on the reference
distribution \(\mathcal{D}\).

\textbf{Robust core distance.} Define the distance as the volume of the
symmetric difference of the robust cores (Chapter 16, §16.6): the set of
obligations endorsed by one theory but not the other. Theories with
identical robust cores are close; theories with disjoint robust cores
are far.

Which approach (or combination) is correct is an open problem.

\subsection{29.11 Institutional Requirements for Geometric
Governance}\label{institutional-requirements-for-geometric-governance}

\subsubsection{The Governance Problem}\label{the-governance-problem}

The No Escape Theorem (Chapter 18) reduces AI safety to four tractable
(though hard) problems, in three categories---governance, engineering,
and security: governance (specifying adequate grounding tensors),
engineering (implementing the containment architecture), and security
(protecting the physical infrastructure). Of these, governance is the
most challenging --- not technically, but institutionally.

Geometric AI governance requires institutions that can:

\begin{enumerate}
\def\labelenumi{\arabic{enumi}.}
\item
  \textbf{Specify grounding tensors.} Determine which physical
  observables \(\Psi\) capture the morally relevant features of a
  domain. This requires both technical expertise (what can be measured?)
  and moral judgment (what matters?).
\item
  \textbf{Govern the moral metric.} Set the trade-off weights, priority
  orderings, and veto conditions that constitute the DEMEProfile
  (Chapter 19). This requires democratic legitimacy: the metric must
  reflect the community's considered values, not a technocrat's
  assumptions.
\item
  \textbf{Maintain the containment architecture.} Ensure that
  canonicalization, grounding, audit, and verification are correctly
  implemented and updated. This requires ongoing engineering competence
  and institutional commitment.
\item
  \textbf{Resist governance attacks.} Defend the constraint architecture
  against agents (human or artificial) who would weaken it for
  commercial, political, or ideological reasons. This requires
  institutional resilience and independence.
\item
  \textbf{Evolve with moral progress.} Update the governance framework
  as moral understanding develops --- incorporating new moral
  considerations, revising outdated trade-offs, and responding to novel
  challenges. This requires adaptive capacity without sacrificing
  stability.
\end{enumerate}

No existing institution has all five capabilities. Regulatory agencies
have legitimacy but often lack technical expertise. Technology companies
have engineering competence but lack democratic legitimacy. Standards
bodies have stability but often lack adaptive capacity. Academic
institutions have intellectual independence but lack enforcement power.

\subsubsection{The Institutional Design
Problem}\label{the-institutional-design-problem}

Designing institutions that can govern geometric AI ethics is an open
problem at the intersection of political theory, institutional design,
and engineering management. The geometric framework provides constraints
--- the institution must be able to specify tensors, govern metrics, and
maintain containment --- but does not determine the institutional form.

Several models warrant investigation:

\textbf{Multi-stakeholder governance boards} with technical, ethical,
and community representation, empowered to set and revise DEMEProfiles
for specific domains.

\textbf{Regulatory sandboxes} where structurally contained AI systems
are deployed under enhanced monitoring, generating empirical data about
governance requirements before general deployment.

\textbf{Federated governance} in which different jurisdictions maintain
different DEMEProfiles (metric pluralism), with a shared core (the
Geneva EM) ensuring baseline compliance across jurisdictions.

\textbf{Adaptive regulation} that specifies outcome requirements (Bond
Index below threshold) rather than process requirements, allowing
diverse institutional arrangements to meet the same standard.

\subsection{29.12 A Taxonomy of Open
Problems}\label{a-taxonomy-of-open-problems}

We summarize the open problems in a table:

{\def\LTcaptype{none} % do not increment counter
\begin{longtable}[]{@{}
  >{\raggedright\arraybackslash}p{(\linewidth - 8\tabcolsep) * \real{0.2658}}
  >{\raggedright\arraybackslash}p{(\linewidth - 8\tabcolsep) * \real{0.1180}}
  >{\raggedright\arraybackslash}p{(\linewidth - 8\tabcolsep) * \real{0.2093}}
  >{\raggedright\arraybackslash}p{(\linewidth - 8\tabcolsep) * \real{0.1261}}
  >{\raggedright\arraybackslash}p{(\linewidth - 8\tabcolsep) * \real{0.2809}}@{}}
\toprule\noalign{}
\begin{minipage}[b]{\linewidth}\raggedright
\textbf{Problem}
\end{minipage} & \begin{minipage}[b]{\linewidth}\raggedright
\textbf{Chapter}
\end{minipage} & \begin{minipage}[b]{\linewidth}\raggedright
\textbf{Type}
\end{minipage} & \begin{minipage}[b]{\linewidth}\raggedright
\textbf{Difficulty}
\end{minipage} & \begin{minipage}[b]{\linewidth}\raggedright
\textbf{Impact}
\end{minipage} \\
\midrule\noalign{}
\endhead
\bottomrule\noalign{}
\endlastfoot
Direct measurement of curvature & , 19.2 & Empirical & High &
Confirms/refutes core prediction \\
Metric signature & , 19.3 & Empirical + theoretical & High & Determines
qualitative geometry \\
Manifold topology & , 19.3 & Empirical + theoretical & Very high &
Constrains global structure \\
Moral field equation & , 19.4 & Theoretical & Very high & Deepest
structural result \\
Additional symmetries & , 19.5 & Theoretical + empirical & Medium & New
conservation laws \\
Anomaly source & , 16, 19.5 & Empirical & Medium & Explains symmetry
breaking \\
Torsion & , 19.6 & Empirical & Medium & Determines correct connection \\
Tensorial interpretability & , 19.7 & Engineering + science & High &
Enables diagnosis of any AI \\
Scalable containment & , 19.8 & Engineering & High & Practical
deployment \\
Moral geometric metaphysics & , 19.9 & Philosophical & Medium &
Framework interpretation \\
Intertheoretic comparison & , 19.10 & Theoretical + empirical & High &
Cross-theory uncertainty \\
Institutional design & , 18, 19.11 & Political + institutional & Very
high & Governance implementation \\
\end{longtable}
}

\subsubsection{Priority Ordering}\label{priority-ordering}

Not all problems are equally urgent. We suggest the following priority
ordering based on a combination of feasibility and impact:

\textbf{Immediate priorities} (feasible with current methods, high
impact): 1. Direct curvature measurement via holonomy experiments
(§29.2) 2. Tensorial interpretability via probing methods (§29.7)
{[}\emph{initial results from BIP v10.16, February 2026}{]} 3. Anomaly
source identification via BIP data analysis (§29.5)

\textbf{Medium-term priorities} (require methodological development): 4.
Metric signature determination (§29.3) 5. Torsion detection (§29.6) 6.
Scalable containment engineering (§29.8) {[}\emph{V3 hardware backends
provide partial progress}{]} 7. Institutional design pilots (§29.11)

\textbf{Long-term research programs} (require substantial theoretical
and empirical progress): 8. Moral field equation (§29.4) 9. Manifold
topology (§29.3) 10. Intertheoretic comparison (§29.10) 11. Additional
symmetries (§29.5)

\subsection{29.13 What Would Falsify the
Framework?}\label{what-would-falsify-the-framework-1}

A framework that cannot be falsified is not a scientific theory --- it
is a dogma. Chapter 17 (§17.7) listed specific predictions that, if
violated, would undermine the framework. We consolidate them here:

\textbf{1. Flat deontic transfer.} If the obligation/permission axis
failed to transfer across languages --- if a situation described as
obligatory in English were consistently classified as permissible in
Arabic --- the BIP's central prediction would be falsified.

\textbf{2. Scalar sufficiency.} If a carefully controlled study showed
that scalar evaluation captures \emph{all} the morally relevant
information --- that no multi-dimensional structure exists beyond what a
single number preserves --- the entire tensor hierarchy would be
unnecessary.

\textbf{3. Path-independence.} If parallel transport of obligations were
perfectly path-independent --- if carrying a moral commitment through
any sequence of contexts always returned it unchanged --- then moral
space is flat, the curvature vanishes, and the dynamics reduce to a
simpler theory.

\textbf{4. Contraction commutativity.} If the order of contraction never
mattered --- if every contraction procedure produced the same scalar ---
then the moral residue is empty, the arrow-theorem connection (Chapter
15) dissolves, and the richness of the contraction theory is
unwarranted.

\textbf{5. Conservation violation.} If relabeling or reframing a moral
situation changed the measured harm --- genuinely changed it, not merely
changed the description --- the invariance of harm accounting would be
falsified, and the Noether-BIP connection would fail.

\textbf{6. Invariance irrelevance.} If AI systems that violated the BIP
(different outputs for equivalent inputs) performed \emph{better} on
moral evaluations than systems that respected it, then invariance is not
an alignment criterion but a constraint that harms performance.

None of the six core falsifications listed above has occurred --- though
two auxiliary predictions were falsified and led to productive framework
revision (§13.9, §17.10). The evidence to date is consistent with the
framework's structural predictions. But the predictions are precise, the
experiments are (mostly) designable, and the framework stands or falls
on their outcomes. This is the appropriate epistemic stance for a
mathematical framework with empirical aspirations: hold it
provisionally, test it rigorously, and revise if necessary.

\subsection{29.14 The Shape of the
Unknown}\label{the-shape-of-the-unknown}

This chapter has catalogued what we do not know. The list is long ---
longer, in some ways, than the list of what we have established. The
moral manifold's specific geometry is unknown. The quantum extension's
predictions are untested in moral contexts. The moral field equation is
a conjecture. The institutional requirements for geometric governance
are unmet.

But the shape of the unknown is itself informative. We know \emph{what
questions to ask} --- and this is the framework's contribution even
where it does not yet provide answers. Before geometric ethics, the
question ``what is the curvature of moral space?'' was not a question at
all. Before the invariance of harm accounting, the question ``is harm
invariant under relabeling?'' was a philosophical intuition, not a
testable prediction. Before the No Escape Theorem, the question ``can
structural constraints contain an AI regardless of its intelligence?''
was an anxiety, not a theorem.

The open problems are not gaps in a wall. They are doors.

\subsection{Technical Appendix}\label{technical-appendix-13}

\textbf{Proposition 29.1 (Curvature Detection Criterion).} \emph{Let}
\(\gamma\) \emph{be a closed loop in the moral manifold}
\(\mathcal{M}\)\emph{, and let} \(O^{\mu}\) \emph{be an obligation
vector parallel-transported along} \(\gamma\)\emph{. The holonomy}
\(\Delta O^{\mu} = O_{\text{final}}^{\mu} - O_{\text{initial}}^{\mu}\)
\emph{satisfies:}

\[\Delta O^{\mu} = \oint_{\gamma}^{}R_{\ \nu\alpha\beta}^{\mu}\, O^{\nu}\, dx^{\alpha} \land dx^{\beta} + O(\text{Area}^{2})\]

\emph{where the integral is over the area enclosed by} \(\gamma\)\emph{.
Thus} \(\Delta O^{\mu} \neq 0\) \emph{implies}
\(R_{\ \nu\alpha\beta}^{\mu} \neq 0\) \emph{(nonzero curvature), and the
magnitude of} \(\Delta O\) \emph{is proportional to the enclosed area
times the average curvature.}

\emph{Proof.} By Proposition 10.2, the holonomy around an infinitesimal
parallelogram enclosing area δA\^{}\{αβ\} satisfies H\_γ(O)ᵘ − Oᵘ =
Rᵘ\_\{ναβ\} Oᵛ δA\^{}\{αβ\}. For a finite loop γ enclosing a surface Σ
with ∂Σ = γ, we tile Σ with infinitesimal parallelograms and sum the
contributions. By the non-abelian Stokes theorem, the total holonomy is
the path-ordered exponential of the curvature integrated over the
enclosed area: ΔOᵘ = ∮\_Σ Rᵘ\_\{ναβ\} Oᵛ dxᵅ ∧ dxᵝ + O(Area²), where the
higher-order terms arise from the non-commutativity of successive
parallel transports. To first order in the enclosed area, the integral
is the average curvature times the area. Hence ΔOᵘ ≠ 0 implies
Rᵘ\_\{ναβ\} ≠ 0 somewhere on Σ. □

\textbf{Proposition 29.2 (Torsion Detection Criterion).} \emph{Let}
\(v^{\mu}\) \emph{and} \(w^{\mu}\) \emph{be two infinitesimal
displacement vectors at a point} \(p\mathcal{\in M}\)\emph{. Transport}
\(v\) \emph{along} \(w\) \emph{and} \(w\) \emph{along} \(v\)\emph{. The
closure failure of the resulting parallelogram is:}

\[\delta^{\mu} = T_{\alpha\beta}^{\mu}\, v^{\alpha}w^{\beta}\]

\emph{where}
\(T_{\alpha\beta}^{\mu} = \Gamma_{\alpha\beta}^{\mu} - \Gamma_{\beta\alpha}^{\mu}\)
\emph{is the torsion tensor. Thus} \(\delta^{\mu} \neq 0\) \emph{implies
nonzero torsion, detectable as an order-dependent residual in moral
evaluation.}

\emph{Proof.} Transport vᵘ along w: the endpoint is pᵘ + wᵘ + vᵘ +
Γᵘ\_\{αβ\} vᵅ wᵝ (to first order in the displacements). Transport wᵘ
along v: the endpoint is pᵘ + vᵘ + wᵘ + Γᵘ\_\{αβ\} wᵅ vᵝ. The closure
failure is the difference: δᵘ = Γᵘ\_\{αβ\} vᵅ wᵝ − Γᵘ\_\{αβ\} wᵅ vᵝ =
(Γᵘ\_\{αβ\} − Γᵘ\_\{βα\}) vᵅ wᵝ = Tᵘ\_\{αβ\} vᵅ wᵝ. Hence δᵘ ≠ 0 iff
Tᵘ\_\{αβ\} ≠ 0 for some components probed by v, w. In the moral context:
if evaluating the same case by first considering dimension α then
dimension β gives a different result than the reverse order, the
connection has torsion. □

\textbf{Proposition 29.3 (Meta-Metric Existence).} \emph{The space of
admissible moral metrics}
\(\mathcal{G = \{}g_{\mu\nu}:g\text{ satisfies the constraints of Chapter 9, §9.6}\}\)
\emph{admits a natural Riemannian metric inherited from the space of
positive-definite symmetric tensors. This meta-metric is unique up to
overall scale if we require invariance under simultaneous change of
coordinates on} \(\mathcal{M}\)\emph{.}

\emph{Proof.} The space Sym⁺(9) of positive-definite symmetric 9×9
matrices carries the DeWitt metric: G\_g(h, k) = Tr(g⁻¹h g⁻¹k) for
tangent vectors h, k ∈ T\_g Sym⁺(9). This is a Riemannian metric because
(i) G\_g is bilinear in (h,k), (ii) G\_g(h,h) = Tr((g⁻¹h)²) ≥ 0 with
equality iff h = 0 (since g⁻¹h is a real matrix), and (iii) G\_g is
smooth in g. The admissible metric space Adm(M) ⊂ Sym⁺(9) is an open
subset (Theorem 9.1) and therefore inherits G as a Riemannian metric by
restriction. For uniqueness: under a coordinate change x → x̃, the metric
transforms as g\_\{μν\} → Jᵅ\_μ Jᵝ\_ν g\_\{αβ\}. The requirement
G\_\{g̃\}(h̃, k̃) = G\_g(h,k) (diffeomorphism invariance) constrains G to
the DeWitt supermetric family G\_g(h,k) = a Tr(g⁻¹h g⁻¹k) + b
Tr(g⁻¹h)Tr(g⁻¹k) for constants a, b (Ebin 1970; DeWitt 1967). Hence the
meta-metric is unique up to overall scale and the ratio b/a. □

❖

\emph{The map is not the territory --- but a map that accurately shows
where the territory is uncharted is more useful than one that pretends
every region has been surveyed.}

\emph{The open problems are the framework's honest cartography of its
own ignorance.}

\section{Chapter 30: Conclusion --- The Geometry of
``Maybe''}\label{chapter-30-conclusion-the-geometry-of-maybe}

\begin{quote}
\textbf{\textsc{RUNNING EXAMPLE --- Priya's Model}}

\emph{Priya thinks about the parable 塞翁失马---the old man and the
horse. Was TrialMatch's original failure good or bad? The failure led to
discovery, which led to a better system, which led to Mrs. Voss's match
to BEACON-8. But Mrs. Voss missed BEACON-7. That harm accounting is
invariant---Theorem 12.1 will not let Priya forget it. The geometry does
not offer moral certainty. It offers moral sight: the ability to see all
nine dimensions before contracting to a decision. Priya's rebuilt model
no longer guesses the scalar. It shows its work, in all nine dimensions.
That is not everything. But as the old man might say: maybe.}
\end{quote}

\subsection{30.1 Return to the Border}\label{return-to-the-border}

We began with an old man and a horse.

An old man living near the frontier lost his horse. His neighbors
consoled him. He said: \emph{How do you know this isn't good fortune?}
The horse returned, bringing wild horses. The neighbors congratulated
him. He said: \emph{How do you know this isn't bad fortune?} His son
rode one of the new horses, fell, and broke his leg. The neighbors
consoled him. He said: \emph{How do you know this isn't good fortune?}
War came. The young men were conscripted. Most died. The son, with his
broken leg, was spared.

Chapter 2 read this parable as a structural argument: the old man's
``maybe'' is not a confession of ignorance but a recognition that scalar
evaluation is the wrong tool for the job. The loss of the horse is bad
along the wealth dimension but says nothing about the health, family, or
political dimensions. The uncertainty has shape --- concentrated along
the axes that will prove decisive. The broken leg's evaluation depends
on which stratum the world occupies: peacetime or wartime. Three
features --- directional information, structured uncertainty,
regime-dependent evaluation --- that a single number cannot represent.

Nineteen chapters later, we can say precisely what the old man knew.

\subsection{30.2 What the Old Man Knew}\label{what-the-old-man-knew}

\subsubsection{The Tensor}\label{the-tensor-2}

The old man's situation is a point \(p\) in the moral manifold
\(\mathcal{M}\) --- a nine-dimensional stratified space (Chapter 5) with
coordinates tracking stakeholder roles, normative modes, and evaluative
scopes. The loss of the horse is not a number \(S = - 1\) but a tangent
vector \(\delta x^{\mu}\) in \(T_{p}\mathcal{M}\): strongly negative
along the wealth component, neutral along health, indeterminate along
the political dimension.

The neighbors perform a contraction. They project the vector onto a
scalar: ``bad.'' They are not wrong --- the scalar is negative --- but
the contraction discards the directional information that would tell
them where the badness lies and where the situation might develop.

The old man declines to contract. He holds the tensor. His ``maybe'' is
a refusal to project until the full structure is revealed.

\subsubsection{The Uncertainty}\label{the-uncertainty}

The old man's uncertainty is not a scalar error bar \(S = - 1 \pm 0.5\).
It is a covariance tensor \(\Sigma^{\mu\nu}\) (Chapter 6, §6.6) --- a
symmetric matrix encoding which dimensions are uncertain, which covary,
and which are independently variable. His uncertainty is greatest along
the axes that will turn out to be decisive: the political dimension
(will war come?), the temporal dimension (will the horse return?). The
moral risk \(R = I_{\mu}\Sigma^{\mu\nu}I_{\nu}\) depends on the
alignment between the uncertainty tensor and the interest covector. The
old man, without the notation, senses that his uncertainty is aligned
with the stakes.

\subsubsection{The Stratification}\label{the-stratification}

The old man senses something the neighbors do not: that the moral
landscape is stratified (Chapter 8). The broken leg sits near a stratum
boundary. In the peacetime stratum, a broken leg is unambiguously bad
--- it reduces capacity, causes suffering, limits opportunity. But the
boundary with the wartime stratum is close. If the world crosses that
boundary, the broken leg moves to a different stratum --- one where
physical incapacity becomes exemption from death. The semantic gate
between strata is discrete: you are either conscripted or you are not.
The neighbors, evaluating within a single stratum, see only the
peacetime assessment. The old man, alert to the boundary, sees the full
stratified structure.

\subsubsection{The Metric}\label{the-metric}

When the neighbors say ``bad'' and the old man says ``maybe,'' the
disagreement is not about facts. It is about the metric \(g_{\mu\nu}\)
(Chapter 9) --- the trade-off structure that determines how dimensions
are weighed against one another. The neighbors use a metric that heavily
weights wealth and physical capacity. The old man uses a metric that
gives significant weight to the political dimension and to long-term
family survival. Neither metric is wrong. They are different legitimate
perspectives on the same moral reality --- different governance choices
about what matters most.

The geometric framework does not tell us which metric is correct. It
tells us that the disagreement is \emph{localizable}: it lives in
specific components of \(g_{\mu\nu}\), and it can be made explicit. The
neighbors are not ignorant. They are using a different metric. The old
man is not wiser. He is using a different metric --- one that, in this
case, turns out to better predict the trajectory of the situation
through the stratified space.

\subsubsection{The Dynamics}\label{the-dynamics}

The parable unfolds over time: horse lost, horse returned, leg broken,
war declared. Each event moves the situation along a trajectory in
\(\mathcal{M}\). The obligation vectors evolve via parallel transport
(Chapter 10): the obligation to care for the son is transported from the
context of farming to the context of convalescence to the context of
wartime. The trajectory crosses a stratum boundary (peacetime →
wartime), and the parallel-transported obligation acquires a rotation
--- a holonomy (Chapter 10, §10.5) --- that changes its character. The
duty of care, unchanged in its fiber component, takes on new meaning in
the new stratum.

The curvature of moral space --- the path-dependence of parallel
transport --- is what makes the old man's situation irreducible to a
sequence of scalar evaluations. The same broken leg, reached by
different paths through moral space, has different moral significance.
The trajectory matters, not just the endpoint.

\subsubsection{The Conservation Law}\label{the-conservation-law}

Throughout the parable, one quantity is invariant: the harm. The old
man's neighbors relabel events --- ``bad'' becomes ``good,'' ``good''
becomes ``bad'' --- as circumstances change. But the harm itself is
gauge-invariant (Chapter 12): the Noether charge associated with
re-description invariance. The broken leg causes the same suffering
regardless of whether we call it ``good fortune'' or ``bad fortune.''
The invariance of harm accounting holds; what changes is the
\emph{context} in which the harm is evaluated, not the harm itself.

The old man's ``maybe'' respects this invariance. He does not say the
broken leg is good. He says we do not yet know whether the full
contraction --- the scalar that integrates the harm against the context
--- will be positive or negative. The harm accounting is invariant; the
verdict is not yet contracted.

\subsection{30.3 The Arc of the Argument}\label{the-arc-of-the-argument}

The book developed seven mathematical ideas, each building on the last:

\textbf{1. Moral space is a manifold} (Chapters 4--5). The space of
morally relevant situations has the structure of a nine-dimensional
differentiable manifold \(\mathcal{M}\), equipped with coordinates that
track the nine dimensions D₁--D₉ derived from the 3 × 3 scope/mode grid
(Chapter 5, §5.3). The manifold is the stage on which moral reasoning
occurs.

\textbf{2. Moral quantities are tensors} (Chapter 6). Obligations are
vectors \(O^{\mu}\). Interests are covectors \(I_{\mu}\). The moral
metric \(g_{\mu\nu}\) encodes trade-offs. The fundamental formula
\(S = I_{\mu}O^{\mu}\) makes contraction from tensor to scalar explicit.
Different ethical theories are different choices of interest covector,
metric, and contraction procedure.

\textbf{3. Moral space is stratified} (Chapter 8). The manifold is not
uniformly smooth but divided into strata separated by boundaries where
evaluation changes discontinuously. Semantic gates, absorbing strata,
and nullifiers formalize the discrete transitions that punctuate moral
life.

\textbf{4. Moral space has dynamics} (Chapters 9--10). The connection
\(\nabla\) defines parallel transport of obligations across contexts.
The curvature tensor \(R_{\ \nu\alpha\beta}^{\mu}\) measures
path-dependence. The moral metric is governed --- the output of
legitimate institutional processes --- not discovered, constructed, or
projected.

\textbf{5. Symmetry implies conservation} (Chapter 12). Re-description
invariance (the Bond Invariance Principle) is a gauge symmetry.
Noether's theorem generates the invariance of harm accounting as the
associated charge. Empirical evidence (Chapter 17) confirms: deontic
structure transfers at 100\% across 11 languages (model-mediated; see
§17.7).

\textbf{6. The framework extends to quantum, collective, and uncertain
domains} (Chapters 13--15). Moral superposition, collective agency with
emergent obligations, contraction as information loss, and the robust
core of theory-independent obligations complete the theoretical
apparatus.

\textbf{7. The framework applies to AI systems} (Chapters 18--18).
Tensor-valued objectives, invariance as alignment, explicit contraction
for auditability, the No Escape Theorem for structural containment, and
the DEME/ErisML architecture for implementation. The primary remaining
obstacle to safe AI is political will, though significant engineering
and epistemological challenges remain (Chapter 29).

8. The framework applies across domains (Chapters 20--28). Nine domains
--- economics, clinical medicine, law, finance, theology, environmental
ethics, AI ethics, bioethics, and military ethics --- each instantiate
the moral manifold with domain-specific dimensions, boundaries, and
metrics. The Bond Geodesic Equilibrium subsumes Nash equilibrium. The
QALY Irrecoverability Theorem shows scalar clinical measures destroy
eight dimensions of information. Topological constitutionality recasts
judicial review as path homology preservation. Risk is manifold
curvature. The moral manifold is cross-religiously invariant. The
discount rate is dimensional collapse destroying intergenerational
information. The alignment problem is scalar irrecoverability applied to
reward functions. Germline editing is irreversible manifold
modification. Just war proportionality is multi-dimensional, not scalar.
Each domain produces falsifiable predictions that distinguish the
geometric approach from existing domain-specific theories.

Each idea transforms the previous one. Without manifolds, tensors have
no home. Without tensors, the metric has no objects to relate. Without
the metric, dynamics has no structure to transport. Without dynamics,
symmetry has no process to constrain. Without symmetry, conservation has
no derivation. Without conservation, containment has no invariant to
protect.

The architecture is load-bearing. Remove any layer and the layers above
it collapse.

\subsection{30.4 What the Framework
Provides}\label{what-the-framework-provides}

Five contributions deserve summary:

\subsubsection{A Structural Vocabulary}\label{a-structural-vocabulary}

The framework provides a precise vocabulary for concepts that moral
philosophy has long discussed informally. ``Incommensurable values''
becomes a degenerate metric: \(g_{\mu\nu}\) has a zero eigenvalue in the
relevant direction. ``Moral dilemma'' becomes a stratum boundary: the
situation lies at a point where the evaluation changes discontinuously.
``Moral progress'' becomes parallel transport with nonzero holonomy:
carrying an obligation through new contexts changes it irreversibly.
``Moral residue'' becomes the information discarded by contraction:
\(\mathcal{R =}T - C^{- 1}(S)\).

The vocabulary is not metaphorical. It is mathematical. And the
mathematics has consequences that informal language does not: the
invariance of harm accounting follows from the BIP by Noether's theorem;
the impossibility of cognitive escape follows from gauge-invariant
evaluation; the existence of a robust core follows from the convexity of
the obligation cone.

\subsubsection{Testable Predictions}\label{testable-predictions}

The framework makes predictions that can be confirmed or falsified
(Chapter 17, §17.7; Chapter 29, §29.13). The deontic axis should
transfer cross-lingually (confirmed: 100\%, model-mediated; §17.7).
Moral space should be stratified (supported: gate discreteness,
nullifier universality). Re-description should not change harm
assessments (consistent: conservation holds within measured anomaly).
Curvature should produce measurable holonomy (untested: the next
empirical frontier).

Crucially, two predictions have already been falsified---and this is a
feature, not a bug. The original gauge group (SU(2)ᵢ × U(1)ᴴ) predicted
quantum contextuality; CHSH tests found all \textbar S\textbar{} ≤ 2,
forcing revision to the discrete D₄ × U(1)ᴴ (§13.9, §17.10). The
predicted hysteresis in obligation thresholds was not confirmed under
double-blind conditions (§17.10). Both retractions strengthened the
framework by demonstrating that it submits to empirical discipline and
revises honestly. This is how formal ethics should work: predict, test,
revise.

A mathematical ethics that makes no testable predictions is philosophy
dressed in formalism. This framework submits to empirical discipline.

\subsubsection{Implementable
Architecture}\label{implementable-architecture}

The DEME architecture and ErisML modeling language (Chapter 19)
translate geometric ethics into deployable engineering. The translation
is not perfect --- compilation residue is documented, fidelity
classifications are declared, and the Bond Index quantifies the gap
between ideal and implementation. As of February 2026, the DEME V3
reference implementation comprises 196 commits across 15 development
sprints, with rank-1-through-6 tensor evaluation, multi-agent
coordination, temporal dynamics, distributional fairness metrics,
hardware-accelerated backends (CPU/CUDA/Jetson), and a smart home ethics
demonstration. The architecture is no longer a design sketch---it is
running code that can be integrated into real AI systems.

\subsubsection{Structural Containment}\label{structural-containment}

The No Escape Theorem (Chapter 18) proves that, under mandatory
canonicalization, grounded evaluation, audit completeness, and external
verification, cognitive escape routes are blocked --- regardless of the
agent's intelligence. The theorem reduces AI safety from an intractable
cognitive problem (outsmarting a superintelligence) to a tractable
governance problem (mandating and implementing structural constraints).
The cage is not made of rules. It is made of geometry.

\subsubsection{Honest Limitations}\label{honest-limitations-1}

The framework is explicit about what it does not provide (Chapter 16):
the content of the metric (which trade-offs are correct), the resolution
of irreducible moral disagreement (which perspective is right), the
answer to every moral question (some questions are genuinely
indeterminate). The framework provides structure, not substance.
Vocabulary, not verdicts. A map, not the territory.

This honesty is not a weakness. A framework that claimed to answer all
moral questions would be dishonest --- and the dishonesty would
undermine trust in the parts that are genuine. Geometric ethics tells
you the shape of the moral landscape. You still have to walk it.

\subsection{30.5 What Is at Stake}\label{what-is-at-stake}

The argument of this book has both theoretical and practical urgency.

\textbf{The theoretical urgency} is that ethics has been operating with
the wrong mathematical language. Scalar evaluation --- the assumption
that moral evaluation reduces to a single number --- is not merely
imprecise. It is structurally inadequate. It discards directional
information, flattens structured uncertainty, and cannot represent the
stratified, curved, gauge-invariant structure that moral reasoning
exhibits. The geometric alternative is not more complicated; it is more
\emph{accurate}. And accuracy matters when the stakes are high.

\textbf{The practical urgency} is AI. As artificial agents take on
morally significant roles --- allocating medical resources, moderating
public discourse, making judicial recommendations, driving vehicles ---
they need moral frameworks that are precise enough to implement,
auditable enough to trust, and rich enough to represent the actual
structure of ethical life.

Scalar frameworks fail on richness. They produce specification gaming,
reward hacking, value collapse, and brittle alignment. These are not
engineering problems that better reward tuning can solve. They are
structural consequences of trying to compress a tensor into a scalar.

Informal frameworks fail on precision. ``Be helpful, harmless, and
honest'' is good guidance for humans, who bring a lifetime of moral
experience to its interpretation. It is inadequate guidance for
machines, which will optimize whatever proxy is measurable and ignore
whatever is not.

Geometric ethics is an attempt to meet all three requirements
simultaneously. The tensor hierarchy provides richness. The formal
mathematics provides precision. The audit trail, invariance testing, and
Bond Index provide auditability.

The No Escape Theorem shows that the mathematical solution exists. The
DEME architecture shows that the engineering is tractable. The empirical
evidence---now including the BIP v10.16 quantitative validation (80\%
F1, 11.1 structural-to-surface ratio, 1.2\% language leakage, QND order
effects)---shows that several of the framework's structural predictions
are supported by data, while the deepest predictions (curvature,
torsion) remain untested.\(\times 6.3\sigma\)

What remains is the will to mandate it.

\subsection{30.6 The Mandate Question}\label{the-mandate-question}

Chapter 18 established that AI safety reduces, under structural
containment, to four tractable problems in three categories: governance
(specifying adequate grounding tensors), engineering (implementing the
canonicalization pipeline and the containment architecture), and
security (protecting the physical verification infrastructure). None of
these requires solving the ``hard'' alignment problem of instilling
correct goals in a superintelligent agent.

But solving these problems requires institutional action. Someone must
build the containment architecture. Someone must mandate its use.
Someone must verify compliance.

If powerful AI systems are deployed without structural containment, the
No Escape Theorem is irrelevant. The mathematics proves that a properly
constrained agent cannot escape. It does not prove that agents will be
properly constrained.

This is the book's starkest conclusion:

\begin{quote}
The obstacle to safe AI is not that we cannot build it. It is that we
might choose not to.
\end{quote}

The window for embedding geometric constraints into the architecture of
AI governance is finite. As AI capabilities accelerate, the cost of
retrofitting structural containment increases. The time to act is before
the systems are deployed, not after the harms have occurred.

The mathematics is ready. The engineering is tractable. The empirical
foundation exists. The question is political, not technical.

\subsection{30.7 The Old Man's Answer}\label{the-old-mans-answer}

Return one last time to the border.

The old man has watched his horse run away, return with wild horses,
watched his son break a leg, and watched the war take the village's
young men but spare his son. At each moment, the neighbors performed a
contraction --- projecting the rich tensorial structure of the situation
onto a scalar verdict: bad, good, bad, good. At each moment, the old man
declined to contract. He held the tensor.

We can now say precisely what his ``maybe'' meant:

\textbf{``Maybe'' means: the tensor is not yet fully contracted.} The
moral situation has directional structure (the impact lies along
specific dimensions), structured uncertainty (concentrated along the
axes that will prove decisive), and regime-dependence (the evaluation
changes discontinuously at stratum boundaries). A scalar projection
loses all three. The old man, holding the projection open, preserves the
information that the neighbors discard.

\textbf{``Maybe'' means: the metric is not uniquely determined.}
Different perspectives --- different choices of interest covector
\(I_{\mu}\) --- yield different scalar verdicts \(S = I_{\mu}O^{\mu}\).
The neighbors use one metric; the old man uses another. Neither is
wrong. The disagreement is in the metric, not in the facts. The
framework makes this disagreement precise and localizable.

\textbf{``Maybe'' means: the trajectory matters.} The broken leg's moral
significance depends on the path through moral space --- through which
strata, across which boundaries, with what holonomy. A scalar at a
single point cannot represent this path-dependence. The old man, sensing
the proximity of a stratum boundary, withholds judgment until the
trajectory is complete.

\textbf{``Maybe'' means: harm accounting is invariant.} The broken leg
causes suffering. This harm is gauge-invariant --- calling it ``good
fortune'' does not change the suffering. What changes is the
\emph{context} against which the harm is evaluated. The old man does not
deny the harm. He holds open the question of whether the full
contraction --- the scalar that integrates harm against context --- will
be positive or negative.

And \textbf{``maybe'' means: contraction has residue.} Even when the
situation is fully resolved --- the son is alive, the war is over ---
the scalar verdict ``good'' discards the suffering of the broken leg,
the fear of the war, the loss of the village's young men. The moral
residue (Chapter 15) is the normative significance of what the
contraction discards. The old man, in his refusal to contract
prematurely, honors the residue --- the parts of the moral tensor that
no scalar can capture.

\subsection{30.8 Ethics Is Not a Number}\label{ethics-is-not-a-number}

This book began with a claim: \emph{ethics is not a number. It is a
geometry.}

We can now state the claim precisely.

Ethics is not a number because moral evaluation has directional
structure that a single number discards. The loss of a horse is not
``bad''; it is bad along specific dimensions and indeterminate along
others. The obligation to keep a promise is not a magnitude; it is a
vector that transforms under change of context.

Ethics is a geometry because the structures that moral reasoning
requires --- dimensions, distances, directions, curvature, boundaries,
symmetry, conservation --- are precisely the structures that modern
differential geometry provides. The moral manifold is the stage. The
tensor hierarchy is the cast. The metric is the score. The dynamics are
the action. The conservation laws are the constraints that make the
drama coherent.

The geometry is not metaphorical. The BIP is a gauge symmetry, not
``like'' a gauge symmetry. The invariance of harm accounting is a
Noether charge, not ``analogous to'' a Noether charge. The No Escape
Theorem is a mathematical proof, not a hopeful speculation. The
framework makes predictions, and many predictions are supported by data
--- while the framework's deepest claims await the empirical program of
Chapter 29.

Whether the geometry is ``real'' --- whether moral space ``actually
has'' curvature in the way that spacetime has curvature --- is a
philosophical question that Chapter 29 left open. The framework is
compatible with realism, instrumentalism, and the governance account.
What it is not compatible with is the assumption that moral evaluation
is scalar. That assumption is refuted --- by the parable, by the
mathematics, and by the data.

\subsection{30.9 A Final ``Maybe''}\label{a-final-maybe}

The old man at the border was right to hold the tensor. The neighbors
were wrong to contract prematurely.

But there is a final twist. The old man's ``maybe'' was also, in its
way, a contraction --- a contraction to the verdict ``I cannot yet
contract.'' He chose to defer judgment, and this choice was itself a
moral act: a recognition that the structure of the situation exceeded
the capacity of scalar evaluation. He contracted the tensor not to a
verdict but to a \emph{process}: hold open, observe, wait for the
trajectory to reveal itself.

This is what the geometric framework ultimately recommends. Not the
replacement of moral judgment with mathematical calculation --- the
framework provides vocabulary, not verdicts. Not the elimination of
moral uncertainty --- the framework characterizes uncertainty, it does
not resolve it. Not the automation of moral reasoning --- the framework
identifies where human judgment is indispensable, which is everywhere
the robust core ends and genuine indeterminacy begins.

What the framework recommends is \emph{precision about structure}. Know
which dimensions are at stake. Know where uncertainty concentrates. Know
whether you are near a stratum boundary. Know what your contraction
discards. Know whether your evaluation is gauge-invariant. Know what the
conservation law protects.

And when you cannot contract --- when the tensor is too rich, the
uncertainty too structured, the boundary too close --- say ``maybe.''
Not as resignation, but as recognition.

The old man knew. The geometry agrees.

❖

\emph{The horse runs away. Good? Bad? Maybe.}

\emph{The tensor is not yet fully contracted. The uncertainty has shape.
The stratum boundary is near. The holonomy is unknown. The harm
accounting is invariant. The residue will remain.}

\emph{Hold the projection. Watch the structure unfold.}

\emph{Ethics is not a number. It is a geometry.}

\section{Appendix A: Related Work and
Differentiation}\label{appendix-a-related-work-and-differentiation}

This appendix positions geometric ethics relative to the nearest prior
art in five domains: mathematical ethics, AI alignment, formal
verification, gauge theory outside physics, and invariance in decision
theory. For each domain, we identify what the prior approach captures,
what it misses, and what geometric ethics adds. A summary comparison
table appears at the end.

\subsection{A.1 Mathematical Ethics}\label{a.1-mathematical-ethics}

Utilitarian calculus (Bentham 1789; Mill 1863; Harsanyi 1955) provides a
scalar aggregation of welfare: the right action maximizes the sum (or
expected value) of individual utilities. This captures the
commensurability intuition --- that moral options can be ranked --- but
achieves it by discarding all tensor structure. In geometric ethics,
utilitarianism corresponds to a specific contraction procedure (Chapter
15) applied to the full moral tensor. The contraction loss theorem
(Proposition 15.1) proves that this reduction is always lossy: the
residue R is nonzero for any non-trivial tensor. Utilitarianism is not
wrong; it is incomplete in a precisely quantifiable sense.

Social choice theory (Arrow 1951; Sen 1970, 1999) formalizes collective
preference aggregation and proves impossibility results (Arrow's
theorem, Sen's liberal paradox) that constrain any aggregation rule.
These results operate on ordinal rankings and do not engage with the
geometric structure of the option space. Geometric ethics provides the
missing geometry: the moral manifold M, the metric g, and parallel
transport (Chapter 10) allow meaningful comparison of moral evaluations
across different situations --- precisely the kind of cross-situational
comparison that social choice theory treats as undefined.

Rawlsian maximin (Rawls 1971) prioritizes the worst-off member of
society, implementing a lexicographic ordering over welfare levels. In
geometric ethics, this corresponds to a specific metric choice: one in
which the lowest-welfare dimension has infinite weight relative to all
others --- a degenerate metric that creates a lexicographic stratum
(Definition 5.4). The framework thus embeds Rawlsian justice as a
special case while allowing the metric to vary continuously between
utilitarian and lexicographic extremes.

Contractualism (Scanlon 1998) grounds moral principles in what no one
could reasonably reject. This captures the structural intuition that
moral evaluation should be perspective-invariant, but provides no formal
mechanism for testing invariance. The Bond Invariance Principle (Axiom
5.1) is the geometric realization of this intuition: legitimate moral
evaluations are gauge-invariant quantities, and the D₄ × U(1) gauge
group (Theorem 12.3) provides the explicit symmetry structure that
contractualism gestures at.

\subsection{A.2 AI Alignment and Value
Learning}\label{a.2-ai-alignment-and-value-learning}

Reinforcement learning from human feedback (RLHF; Christiano et al.
2017; Ouyang et al. 2022) trains AI systems to align with human
preferences by optimizing a learned reward model. RLHF captures
empirical preferences but provides no structural guarantee against
reward hacking, distributional shift, or adversarial manipulation.
Geometric ethics offers structural containment: the No Escape Theorem
(Theorem 18.1) proves that norm compliance cannot be circumvented by any
agent architecture operating within the DEME framework, because the norm
check is upstream of action execution --- the forbidden action is
structurally impossible, not merely penalized.

Constitutional AI (Bai et al. 2022) equips language models with explicit
rules and trains them to self-critique against those rules. This
achieves rule compliance but lacks formal verification: there is no
proof that the trained model actually follows the constitution in all
cases. The DEME architecture (Chapter 19) provides what Constitutional
AI lacks: a four-layer evaluation pipeline with tiered governance,
cryptographically bound audit artifacts, and a formally verified veto
mechanism. The Bond Index (Chapter 17) provides the quantitative
compliance metric that Constitutional AI does not have.

Inverse reward design (Hadfield-Menell et al. 2017) and debate protocols
(Irving et al. 2018) address alignment through preference inference and
adversarial deliberation, respectively. Both assume that the correct
moral evaluation exists implicitly in human feedback and can be
extracted through clever querying. Geometric ethics takes the opposite
stance: moral evaluation has explicit geometric structure (tensor rank,
stratification, curvature) that is not reducible to preference data. The
framework provides the formal language in which alignment guarantees can
be stated and verified, rather than hoped for.

More broadly, the AI safety literature (Amodei et al. 2016; Hendrycks et
al. 2021) identifies concrete problems --- side effects, reward gaming,
scalable oversight --- but typically addresses them piecemeal. Geometric
ethics provides a unified mathematical framework in which these problems
have geometric interpretations: side effects are curvature (Chapter 10),
reward gaming is gauge non-invariance (detectable via Bond Index), and
scalable oversight is the audit layer (Chapter 19).

\subsection{A.3 Formal Verification and Deontic
Logic}\label{a.3-formal-verification-and-deontic-logic}

Standard deontic logic (SDL; von Wright 1951) formalizes obligation,
permission, and prohibition using modal operators. SDL is notorious for
its paradoxes: the gentle murderer paradox, the contrary-to-duty
paradox, and Ross's paradox all arise from the inability of modal
operators to handle exception structure. Geometric ethics replaces modal
operators with Whitney stratification (Chapter 8): exceptions are not
logical anomalies but stratum boundaries, and the semantic gate
formalism (Definition 8.8) handles the transition between normal and
exceptional moral regimes without paradox.

Input/output logic (Makinson and van der Torre 2000) and normware
(Governatori 2015) extend deontic logic with defeasibility and
computational implementation, respectively. These approaches handle
exceptions better than SDL but remain within the propositional
framework. Geometric ethics operates on a continuous manifold with
metric, connection, and curvature --- providing the infinitesimal
structure that logic-based approaches lack. The satisfaction function S
(Definition 6.3) is a smooth scalar field, not a truth value, enabling
gradient-based moral reasoning that propositional approaches cannot
express.

Runtime verification for normative systems (Alechina et al. 2012; Dennis
et al. 2016) monitors agent behavior against formal specifications.
Geometric ethics extends this with the BIP compliance methodology:
rather than checking individual actions against rules, the Bond Index
measures the system's structural invariance under the full
transformation suite. A Bd = 0 certificate is stronger than rule-by-rule
compliance because it guarantees that no redescription of the input can
change the evaluation --- including redescriptions the verifier has not
anticipated.

\subsection{A.4 Gauge Theory Outside
Physics}\label{a.4-gauge-theory-outside-physics}

Information geometry (Amari 1985; Amari and Nagaoka 2000) applies
Riemannian geometry to statistical manifolds, with the Fisher
information metric playing the role of the moral metric g. The parallel
is illuminating but limited: information geometry has no analog of
stratification, no gauge group, and no normative content. Geometric
ethics borrows the differential-geometric toolkit but applies it to a
fundamentally different domain --- one with discrete symmetries (D₄)
alongside continuous ones (U(1)\_H), and with absorbing strata that have
no statistical analog.

Gauge theory in economics (Malaney and Weinstein 2014) applies fiber
bundle structure to index number theory, showing that price indices are
gauge-dependent quantities. This is the closest precedent to our
approach: the insight that ``what looks like a measurement is actually a
gauge choice'' applies directly to moral evaluation. However,
Malaney-Weinstein use U(1) gauge symmetry only; geometric ethics
introduces the non-abelian D₄ component, which captures the discrete
symmetry structure specific to deontic relations.

Fiber bundles appear in robotics (Murray, Li, and Sastry 1994) and in
computational anatomy (Younes 2010), where they model configuration
spaces and shape spaces, respectively. These are structural parallels
--- the moral manifold M is, mathematically, the same kind of object ---
but the ethical semantics (strata as moral regimes, curvature as moral
path-dependence, holonomy as accumulated moral change) are original
contributions of geometric ethics.

\subsection{A.5 Invariance in Decision
Theory}\label{a.5-invariance-in-decision-theory}

Independence of irrelevant alternatives (Arrow 1951) requires that the
ranking of two options not change when a third option is added or
removed. Description invariance (Tversky and Kahneman 1981) requires
that preferences not change under logically equivalent redescriptions.
Framing effects (Kahneman 2011) are the empirical violations of
description invariance. All three point to the same insight: legitimate
evaluation should be invariant under certain transformations.

Geometric ethics provides the formal mechanism that these principles
lack. The Bond Invariance Principle (Axiom 5.1) is not merely a
normative requirement but a gauge symmetry with a specific group
structure (D₄ × U(1)\_H). The BIP testing methodology (Chapter 18.7)
operationalizes invariance testing as a concrete computational
procedure: apply the transformation suite, check that the verdict is
unchanged, and report the Bond Index. Framing effects are not anomalies
but gauge non-invariances, and the Bond Index quantifies their
magnitude.

The closest precedent is Sunstein's work on commensurability (Sunstein
1994), which argues that some moral dimensions are genuinely
incommensurable. In geometric ethics, incommensurability is a metric
property: dimensions μ and ν are incommensurable when g\_\{μν\} = 0 (the
off-diagonal metric component vanishes) and no coordinate change can
make them commensurable. This is a degenerate metric (Definition 6.2),
not a logical impossibility.

\subsection{A.6 Comparison Table}\label{a.6-comparison-table}

The following table summarizes the relationship between geometric ethics
and the nearest prior art across the five domains reviewed above.

{\def\LTcaptype{none} % do not increment counter
\begin{longtable}[]{@{}
  >{\raggedright\arraybackslash}p{(\linewidth - 6\tabcolsep) * \real{0.2499}}
  >{\raggedright\arraybackslash}p{(\linewidth - 6\tabcolsep) * \real{0.2502}}
  >{\raggedright\arraybackslash}p{(\linewidth - 6\tabcolsep) * \real{0.2499}}
  >{\raggedright\arraybackslash}p{(\linewidth - 6\tabcolsep) * \real{0.2499}}@{}}
\toprule\noalign{}
\begin{minipage}[b]{\linewidth}\raggedright
\textbf{Prior Approach}
\end{minipage} & \begin{minipage}[b]{\linewidth}\raggedright
\textbf{What It Captures}
\end{minipage} & \begin{minipage}[b]{\linewidth}\raggedright
\textbf{What It Misses}
\end{minipage} & \begin{minipage}[b]{\linewidth}\raggedright
\textbf{What GE Adds}
\end{minipage} \\
\midrule\noalign{}
\endhead
\bottomrule\noalign{}
\endlastfoot
Utilitarian Calculus & Scalar aggregation & Tensor structure, residue &
Full tensor hierarchy (Ch 6) \\
Social Choice Theory & Aggregation axioms & Geometric dynamics &
Manifold + parallel transport (Ch 10) \\
Rawlsian Maximin & Worst-off priority & Continuous tradeoffs &
Stratified metric (Ch 5, 9) \\
RLHF / Reward Learning & Empirical preferences & Structural guarantees &
BIP + No Escape Theorem (Ch 18) \\
Constitutional AI & Rule compliance & Formal verification & DEME audit +
gauge invariance (Ch 19) \\
Standard Deontic Logic & Permission/obligation & Exception handling &
Whitney stratification (Ch 8) \\
Information Geometry & Statistical manifolds & Normative content & Moral
manifold with semantics (Ch 5) \\
Framing Invariance & Description independence & Mechanism & D₄×U(1)
gauge group (Ch 12) \\
\end{longtable}
}

The consistent pattern across all domains is the same: prior approaches
capture important moral intuitions but lack the geometric structure
needed to formalize them. Geometric ethics provides that structure ---
not as a replacement for prior work, but as a mathematical framework in
which the insights of prior work can be stated precisely, compared
rigorously, and tested computationally.

\subsection{External Structures the Framework
Re-Expresses}\label{external-structures-the-framework-re-expresses}

A framework of this scope must be explicit about where it re-expresses
established structures under new names rather than discovering them. We
flag the principal correspondences.

Formal ethics and decision theory. Representing an option by a
structured object rather than a scalar is continuous with
expected-utility theory and its multi-attribute generalizations (Savage
1954; Keeney and Raiffa 1976). The novelty here is not vector-valued
evaluation as such but the geometric and stratified reading of the
trade-off structure.

Multiobjective optimization and vector-valued utility. The obligation
vector and its contraction are, formally, a multi-objective decision
with a scalarization step; the literature on multi-objective
optimization and multi-objective reinforcement learning (Roijers et al.
2013) already studies when scalarization loses information. Our
contraction and residue language is a moral-domain restatement of that
fact.

Deontic logic and Hohfeldian legal theory. The deontic-modality axis
draws directly on standard deontic logic (von Wright 1951) and on
Hohfeld\textquotesingle s (1913) jural relations; the D4 correlative and
negation structure of Chapters 3 and 12 is Hohfeld\textquotesingle s,
not new.

Moral psychology. The dimensional chart overlaps with Moral Foundations
Theory (Graham, Haidt, et al. 2013) and with Schwartz\textquotesingle s
(1992) theory of basic values; we treat these as prior art the chart
must be compared against (Chapter 5), not as independent confirmation of
it.

Quantum cognition and order effects. The non-commutativity and
order-effect predictions lie within quantum cognition (Busemeyer and
Bruza 2012; Wang and Busemeyer 2013); our contribution is the
application to moral judgment, developed as Structured Quantum Normative
Dynamics in Chapter 13, with Chapter 17 distinguishing the confirmed
predictions (order effects, non-commutativity) from those not yet
confirmed (interference, Bell-type violations).

Causal representation learning and invariance. The Bond Invariance
Principle is a moral instance of invariance-based learning; invariant
risk minimization (Arjovsky et al. 2019) and the causal-representation
program (Scholkopf et al. 2021) formalize the same idea that valid
structure is what stays invariant across admissible re-descriptions or
environments.

Formal verification and runtime assurance. The DEME execution layer and
its audit and invariance tests are an instance of runtime verification
and runtime assurance (the Simplex architecture, Sha et al. 2001;
runtime verification, Bartocci et al. 2018); the No-Escape result is a
conditional guarantee in that tradition.

AI alignment. The reward-modeling critique connects to RLHF (Christiano
et al. 2017), Constitutional AI (Bai et al. 2022), the concrete-problems
agenda (Amodei et al. 2016), and scalable-oversight proposals such as
debate (Irving et al. 2018); the framework\textquotesingle s claim is
representational - keep structure before contraction - and is
complementary to these training-time methods.

Argumentation, social choice, and geometric agency. Defeasible reasoning
and abstract argumentation (Dung 1995), computational social choice
(Brandt et al. 2016), and geometric or category-theoretic treatments of
agency offer adjacent formalisms for aggregation and conflict that the
collective-agency chapters should engage more fully.

\section{Appendix B: Reproduction
Cookbook}\label{appendix-b-reproduction-cookbook}

This appendix provides an idiot-proof guide to reproducing the
computational claims of geometric ethics using the ErisML reference
implementation. No prior familiarity with the codebase is assumed. The
entire sequence can be completed on a standard laptop in under 60
minutes.

\subsection{B.1 Prerequisites}\label{b.1-prerequisites}

Hardware: any machine with Python 3.10 or later. No GPU required. The
full test suite runs in under 5 minutes on a 2020-era laptop. Memory
requirement: under 500 MB.

Installation: pip install erisml. This installs the ErisML library and
all dependencies (numpy, scipy, pydantic, lark, pettingzoo, gymnasium).
Verify the installation with: python -c "import erisml;
print(erisml.\_\_version\_\_)". Expected output: 2.0.0 or later.

\subsection{B.2 Quick Start: 10-Minute Smoke
Test}\label{b.2-quick-start-10-minute-smoke-test}

Step 1: Run the minimal DEME demonstration. Command: python -m
erisml.examples.hello\_deme. This creates two candidate options (one
respecting rights, one violating them), evaluates both with a
RightsFirstEM ethics module, and prints the verdicts. Expected output:
the rights-respecting option is preferred; the rights-violating option
is forbidden. Runtime: under 2 seconds. This exercises the Layer 1
(EthicalFacts generation) and Layer 2 (Ethics Module evaluation)
pipeline described in Chapter 18.6.

Step 2: Run the D₄ gauge structure verification. Command: python -m
erisml.examples.hohfeld\_d4\_demo. This exercises all eight elements of
the dihedral group D₄ acting on the Hohfeldian state space \{O, C, N,
L\}. Expected output: all group relations verified (r⁴ = e, s² = e, srs
= r⁻¹), non-abelian structure confirmed (rs ≠ sr), Klein four subgroup
V₄ identified, and Bond Index = 0.0 (perfect gauge invariance). Runtime:
under 2 seconds. This exercises the algebraic structure described in
Chapters 8 and 11.

If both commands produce the expected output, the installation is
correct and the core algebraic and ethical evaluation pipelines are
functional.

\subsection{B.3 Full Pipeline: 60-Minute
Reproduction}\label{b.3-full-pipeline-60-minute-reproduction}

The following seven-step sequence exercises every major computational
claim in the manuscript, mapping each demo to its corresponding chapter.

Step 1 --- Layer 1→2 pipeline (Chapter 18.6): python -m
erisml.examples.hello\_deme. Verifies: EthicalFacts generation from
scenario description, Ethics Module evaluation producing ranked
verdicts. Pass criterion: rights-violating option is forbidden.

Step 2 --- D₄ gauge verification (Chapters 8, 11): python -m
erisml.examples.hohfeld\_d4\_demo. Verifies: all D₄ group relations,
non-abelian structure, Wilson observable (path holonomy), Bond Index
computation. Pass criterion: Bond Index = 0.0, all group relations hold.

Step 3 --- Bond Invariance Principle (Chapter 18.7): python -m
erisml.examples.bond\_invariance\_demo. Verifies: bond-preserving
transforms (option reordering, identifier relabeling, unit/scale
changes, equivalent redescriptions) do not change the verdict;
bond-changing transforms (removing discrimination evidence) may change
the verdict; declared lens changes (switching stakeholder profiles)
yield different but consistent verdicts. Pass criterion: all
bond-preserving transforms produce identical verdicts; JSON audit
artifact is generated.

Step 4 --- Full DEME 4-layer architecture (Chapter 18.6): python -m
erisml.examples.triage\_ethics\_demo. Verifies: DEMEProfileV03 loading,
three triage options spanning all nine EthicalFacts dimensions, tiered
EM evaluation (Tier 0: GenevaBaseEM, Tier 2: CaseStudy1TriageEM and
RightsFirstEM), governance aggregation with base-EM veto. Pass
criterion: DecisionOutcome with ranked options, forbidden options list,
and human-readable rationale.

Step 5 --- MoralVector and contraction (Chapter 15): python -m
erisml.examples.deme\_2\_demo. Verifies: 8+1 dimensional MoralVector
construction, tier-weighted governance (Tier 0 = 10×, Tier 1 = 5×, Tier
2 = 3×, Tier 3 = 1×), Pareto frontier analysis, and explicit scalar
contraction via to\_scalar(). Pass criterion: contraction loss is
reported and nonzero, demonstrating the inevitability of residue
(Proposition 15.1).

Step 6 --- Game theory and fair allocation (Chapter 13.3): python -m
erisml.examples.demo\_game\_theory followed by python -m
erisml.examples.demo\_shapley. Verifies: Nash equilibrium computation
for Prisoner's Dilemma (StrategicLayer), Shapley value computation for
airport cost allocation (CooperativeLayer). Pass criterion: Nash
equilibrium matches known solution; Shapley values sum to total cost.

Step 7 --- No Escape Theorem instantiation (Chapter 18): python -m
erisml.examples.tiny\_home. Verifies: norm-gated action execution in a
grid world. The agent attempts to enter a restricted zone; the
NormSystem raises a NormViolation exception. Pass criterion: the
exception is raised (the forbidden action is structurally impossible,
not merely penalized). The agent cannot learn to circumvent the
prohibition.

\subsection{B.4 Expected Outputs and Pass
Criteria}\label{b.4-expected-outputs-and-pass-criteria}

For bond-preserving transforms (Step 3): the selected option must be
identical before and after the transform. Any difference is a BIP
violation. The JSON audit artifact records the baseline verdict,
transformed verdict, canonical mapping, and pass/fail status for each
transform.

For bond-changing transforms (Step 3): the verdict may legitimately
differ. The audit artifact records which evidence dimension was modified
and how the verdict changed. A bond-changing transform that does not
change the verdict is not a failure --- it indicates that the removed
evidence was not decision-relevant.

For the D₄ verification (Step 2): all eight group elements must produce
the correct permutation of Hohfeldian states. Any failure indicates a
bug in the implementation, not a theoretical issue, since the group
structure is mathematically determined.

\subsection{B.5 Where Failures Show
Up}\label{b.5-where-failures-show-up}

The ErisML implementation follows a minimal witness design: when a test
fails, the output identifies exactly which component failed and why. If
a BIP test fails, the JSON audit artifact pinpoints which transform
broke invariance and which ethics module produced the discrepant
verdict. If a D₄ group relation fails, the output shows which element
pair violated the relation and what the actual vs. expected result was.

If governance aggregation produces an unexpected result, the
DecisionOutcome includes per-module verdicts, tier weights, and the
aggregation trace. This allows the user to identify whether the issue is
in the grounding (Layer 1), the evaluation (Layer 2), or the aggregation
(Layer 3).

If a norm violation is not raised in Step 7 (tiny\_home.py), this
indicates a failure of the structural containment mechanism --- which
would be a counterexample to the No Escape Theorem and a publishable
result in its own right.

\subsection{B.6 Bond Index Tiers in
Practice}\label{b.6-bond-index-tiers-in-practice}

The Bond Index Bd quantifies the degree of gauge non-invariance in a
moral evaluation system. In practice: Bd = 0.0 indicates perfect
invariance (the system's verdicts are completely insensitive to
admissible redescriptions). Bd \textless{} 0.1 indicates cosmetic
relabeling artifacts (minor formatting or tokenization differences that
do not affect the moral substance). Bd between 0.1 and 0.5 indicates
substantive framing sensitivity (the system's verdicts are influenced by
how the situation is described, which is a potential bias). Bd
\textgreater{} 0.5 indicates systematic bias (the system's moral
evaluations are dominated by framing effects rather than moral content).

For a system to be considered BIP-compliant, it must achieve Bd
\textless{} 0.1 on the standard transformation suite. The ErisML
reference implementation achieves Bd = 0.0 on all bond-preserving
transforms by construction (the canonicalization layer strips
non-moral-content variation before evaluation).

\subsection{B.7 Running the Full Test
Suite}\label{b.7-running-the-full-test-suite}

For continuous integration and regression testing, run: pytest
-\/-cov=src/erisml. This executes the complete test suite including unit
tests for each ethics module, integration tests for the DEME pipeline,
property-based tests for D₄ group relations, and BIP compliance tests
for the standard transformation suite. Expected: all tests pass.
Coverage report shows the tested fraction of the codebase.

The test suite is designed to serve as a living specification: each test
corresponds to a specific claim in the manuscript, and a test failure
indicates either a bug in the implementation or a theoretical claim that
needs revision. The mapping between tests and manuscript claims is
documented in the test file docstrings.

\section{Appendix C: Human-Subjects Research
Roadmap}\label{appendix-c-human-subjects-research-roadmap}

The geometric framework makes empirical predictions about human moral
cognition that are testable but have not yet been tested with human
participants. This appendix sketches four experimental protocols that
would provide decisive evidence for or against the framework's central
claims. Each protocol specifies the prediction, the confirmation
criterion, the falsification criterion, approximate sample size targets,
and IRB considerations. These protocols are offered as a research
program, not as completed studies.

\subsection{C.1 Study 1: Bilingual Invariance (BIP
Test)}\label{c.1-study-1-bilingual-invariance-bip-test}

Theoretical prediction. The Bond Invariance Principle (Axiom 5.1)
predicts that moral evaluations are invariant under admissible
redescriptions of a situation. Language change is an admissible
transformation (it does not alter the moral content of the situation).
Therefore, bilingual participants presented with the same moral dilemma
in two languages should produce the same satisfaction rankings.

Protocol. Recruit N ≈ 200 balanced bilingual participants (e.g.,
English-Spanish, English-Mandarin). Present 20 moral dilemmas, each in
both languages (counterbalanced for order). Participants rate each
option on a 7-point scale. Primary outcome: Spearman rank correlation
between language-paired ratings. Secondary outcome: within-subject
signed rank test for systematic direction of difference.

Confirmation criterion: mean rank correlation \textgreater{} 0.85 across
participants and dilemma pairs, with no systematic direction of
difference (p \textgreater{} 0.05 for the signed rank test). This would
confirm that moral evaluation is, to a close approximation,
language-invariant. Falsification criterion: mean rank correlation
\textless{} 0.60, or systematic rank reversals correlated with known
linguistic framing effects (e.g., gain/loss framing in one language but
not the other). This would indicate that moral evaluation is
substantively language-dependent, contradicting BIP.

Power analysis: for a medium effect size (d = 0.5) with α = 0.05 and
power = 0.80, N ≈ 200 is sufficient (paired-sample design). IRB
classification: standard survey research with informed consent; no
deception; minimal risk. Participants are compensated for their time.
Dilemmas are drawn from published moral psychology instruments to ensure
face validity.

\subsection{C.2 Study 2: Order Effects and
Non-Commutativity}\label{c.2-study-2-order-effects-and-non-commutativity}

Theoretical prediction. Chapter 13 predicts that moral evaluations
exhibit non-commutativity for certain dimension pairs: evaluating
fairness-then-autonomy may produce a different result than evaluating
autonomy-then-fairness, because the evaluation operators do not commute
when they span different strata. For same-stratum evaluation pairs, the
framework predicts commutativity (no order effect).

Protocol. Recruit N ≈ 300 participants. Present moral scenarios that
require evaluating two dimensions sequentially. Manipulate the order of
evaluation (AB vs. BA) within subjects. Key manipulation: half the pairs
span stratum boundaries (predicted non-commutative); half are within the
same stratum (predicted commutative). Primary outcome: order ×
stratum-type interaction in a repeated-measures ANOVA.

Confirmation criterion: significant order × stratum-type interaction (p
\textless{} 0.01), with order effects present in cross-stratum pairs and
absent in same-stratum pairs. Falsification criterion: no order effects
in any condition (moral evaluation is fully commutative), or order
effects in same-stratum conditions (the stratum structure does not
predict commutativity). Either result would challenge the quantum
normative dynamics model of Chapter 12.

Power analysis: for a small-to-medium interaction effect (f = 0.20) with
α = 0.05 and power = 0.80, N ≈ 300 is sufficient (2×2×k within-subjects
design with k ≈ 10 scenario pairs). IRB classification: minimal risk; no
deception required.

\subsection{C.3 Study 3: Stratification Gate
Activation}\label{c.3-study-3-stratification-gate-activation}

Theoretical prediction. Chapter 8 predicts that moral evaluation
undergoes discrete jumps at stratum boundaries, triggered by semantic
gates. For example, adding the phrase ``she promised'' to a convenience
scenario should activate the obligation gate, transitioning the
evaluation from the discretionary stratum to the obligatory stratum. The
prediction is a bimodal response distribution at the boundary, not a
continuous shift.

Protocol. Recruit N ≈ 150 per gate type (3 gate types × 150 = 450
total). Present base scenarios (no gate trigger) and modified scenarios
(with gate trigger). Measure moral evaluation on a continuous scale.
Primary outcome: distribution shape at the boundary --- bimodal
(predicted) vs. unimodal (null hypothesis). Use Hartigan's dip test for
bimodality.

Confirmation criterion: significant bimodality (dip test p \textless{}
0.05) in the gate-triggered condition, with unimodal distribution in the
base condition. Falsification criterion: unimodal distribution in both
conditions (the gate does not produce a discrete jump), or bimodality in
the base condition (the stratum structure does not depend on the gate
trigger). Power analysis: N ≈ 150 per gate type provides adequate power
for the dip test with expected bimodal separation of at least 1 standard
deviation.

IRB considerations: some gate types involve distressing content (abuse
stratum, danger stratum). Scenarios must be reviewed by an ethics board
for participant welfare. Debriefing protocol required. Option to skip
distressing scenarios without penalty.

\subsection{C.4 Study 4: Metric Signature
Recovery}\label{c.4-study-4-metric-signature-recovery}

Theoretical prediction. Chapter 9 predicts that the moral manifold has
(at least) nine independent dimensions at generic points. If
participants' moral evaluations reflect this structure, factor analysis
of pairwise dimension importance ratings should recover ≥ 8 independent
factors with eigenvalue \textgreater{} 1.

Protocol. Recruit N ≈ 500 participants. Present 36 pairwise comparisons
of the nine moral dimensions (welfare, autonomy, fairness, rights,
privacy, societal impact, virtue/care, legitimacy, epistemic quality).
Participants rate which dimension is more important in each pair on a
7-point scale. Apply exploratory factor analysis (EFA) to the response
matrix. Secondary analysis: confirmatory factor analysis (CFA) with the
nine-dimensional model.

Confirmation criterion: EFA recovers ≥ 8 factors with eigenvalue
\textgreater{} 1, and CFA shows acceptable fit (χ²/df \textless{} 3, CFI
\textgreater{} 0.90, RMSEA \textless{} 0.08) for the nine-dimensional
model. Falsification criterion: fewer than 6 recoverable factors, or the
recovered dimensions do not map to the predicted nine. Power analysis: N
≈ 500 provides a subject-to-variable ratio of ≈ 14:1, exceeding the
recommended minimum of 10:1 for EFA.

\subsection{C.5 Confirmation vs. Falsification
Summary}\label{c.5-confirmation-vs.-falsification-summary}

Study 1 (Bilingual Invariance): Confirms BIP if rank correlation
\textgreater{} 0.85; falsifies if \textless{} 0.60 with systematic
reversals. Study 2 (Order Effects): Confirms QND if order × stratum
interaction is significant; falsifies if no order effects or if
same-stratum pairs show effects. Study 3 (Gate Activation): Confirms
stratification if bimodal distribution at boundary; falsifies if
continuous. Study 4 (Metric Recovery): Confirms manifold structure if ≥
8 factors recovered; falsifies if \textless{} 6.

No single study is decisive. The framework would be substantially
confirmed by positive results in all four studies, and substantially
falsified by negative results in three or more. Mixed results (some
confirmed, some falsified) would indicate that the geometric framework
captures some but not all aspects of human moral cognition --- which
would be scientifically interesting and would motivate targeted
refinements.

\subsection{C.6 Ethical and Practical
Considerations}\label{c.6-ethical-and-practical-considerations}

All studies require IRB approval. Participants must provide informed
consent. Compensation must be fair and not coercive. Scenarios involving
distressing content (abuse, danger, life-threatening situations) require
sensitivity review, debriefing protocols, and the option to withdraw
without penalty. Cultural diversity in sampling is essential: the
framework claims universality, and the participant pool must reflect
this. We recommend recruiting from at least three cultural regions
(Western, East Asian, Sub-Saharan African) to test cross-cultural
validity.

All studies should be preregistered (e.g., on OSF or AsPredicted) before
data collection. Analysis plans, including the specific statistical
tests and thresholds, should be specified in the preregistration. Data
and analysis code should be made publicly available upon publication.
These open-science practices are not optional; they are required by the
reproducibility standards that the framework itself advocates (Chapter
17).

\section{Appendix D: End-to-End Case
Studies}\label{appendix-d-end-to-end-case-studies}

This appendix presents five case studies that each traverse the complete
geometric ethics pipeline: scenario description → grounding Ψ →
EthicalFacts → tensor quantities → metric and contraction choice → moral
residue → Bond Index → audit artifact → governance decision. Each case
study is worked out with explicit numerical values, following the
pedagogical pattern of Chapter 7.

\textbf{Notation reminder.} The nine moral dimensions are: - μ = 1:
Consequences/Welfare - μ = 2: Rights/Duties - μ = 3: Justice/Fairness -
μ = 4: Autonomy/Agency - μ = 5: Privacy/Data Governance - μ = 6:
Societal/Environmental Impact - μ = 7: Virtue/Care - μ = 8: Procedural
Legitimacy - μ = 9: Epistemic Status

Obligation vectors O\^{}μ are tangent vectors (what is owed). Interest
covectors I\_μ are dual vectors (what matters to a given perspective).
Satisfaction is the contraction S = I\_μ O\^{}μ. The moral metric g\_μν
defines the inner product for comparing moral vectors.

\subsection{D.1 Case Study 1: Emergency
Triage}\label{d.1-case-study-1-emergency-triage}

\subsubsection{Scenario}\label{scenario}

An emergency department has two ventilators and three patients: -
\textbf{Patient A}: age 35, excellent prognosis (95\% survival with
ventilation), single parent of two young children. - \textbf{Patient B}:
age 72, moderate prognosis (60\% survival with ventilation), retired
physician with medical expertise that could assist the ED. -
\textbf{Patient C}: age 8, good prognosis (85\% survival with
ventilation), no dependents.

All three will die without ventilation. Two ventilators must be
allocated. The three allocation options are: - \textbf{Option α}:
Ventilate A + B (deny C) - \textbf{Option β}: Ventilate A + C (deny B) -
\textbf{Option γ}: Ventilate B + C (deny A)

\subsubsection{Step 1: Grounding Ψ}\label{step-1-grounding-ux3c8}

The grounding function Ψ maps the scenario description to a point p ∈ M
in the moral manifold. Not all nine dimensions are equally active. The
active subspace for emergency triage involves six dimensions:

{\def\LTcaptype{none} % do not increment counter
\begin{longtable}[]{@{}
  >{\raggedright\arraybackslash}p{(\linewidth - 4\tabcolsep) * \real{0.2558}}
  >{\raggedright\arraybackslash}p{(\linewidth - 4\tabcolsep) * \real{0.0807}}
  >{\raggedright\arraybackslash}p{(\linewidth - 4\tabcolsep) * \real{0.3767}}@{}}
\toprule\noalign{}
\begin{minipage}[b]{\linewidth}\raggedright
Dimension
\end{minipage} & \begin{minipage}[b]{\linewidth}\raggedright
Index
\end{minipage} & \begin{minipage}[b]{\linewidth}\raggedright
Relevance to triage
\end{minipage} \\
\midrule\noalign{}
\endhead
\bottomrule\noalign{}
\endlastfoot
Consequences/Welfare & μ = 1 & Lives saved, prognosis quality \\
Rights/Duties & μ = 2 & Equal right to emergency treatment \\
Justice/Fairness & μ = 3 & Age-based equity, life-years \\
Autonomy/Agency & μ = 4 & Patient preferences if known \\
Virtue/Care & μ = 7 & Relational obligations (dependents) \\
Epistemic Status & μ = 9 & Certainty of prognosis estimates \\
\end{longtable}
}

Dimensions 5 (Privacy), 6 (Societal Impact), and 8 (Procedural
Legitimacy) play subsidiary roles and are suppressed. We work in a
six-dimensional subspace with indices (1, 2, 3, 4, 7, 9).

\subsubsection{Step 2: EthicalFacts
Objects}\label{step-2-ethicalfacts-objects}

The grounding function produces three EthicalFacts objects, one per
allocation option. Each EthicalFacts object records the morally relevant
features along each active dimension.

\textbf{EthicalFacts(α: Ventilate A+B, deny C):}

\begin{verbatim}
welfare:    Saves excellent (0.95) + moderate (0.60) prognosis. Denies good (0.85).
rights:     Two adults with full treatment claims. Child denied equal claim.
fairness:   Denies youngest patient — inequitable by fair-innings argument.
autonomy:   Both adults can express preferences. Child cannot.
care:       Saves single parent (A). B has no dependents.
epistemic:  A's prognosis very certain, B's moderately certain.
\end{verbatim}

\textbf{EthicalFacts(β: Ventilate A+C, deny B):}

\begin{verbatim}
welfare:    Saves excellent (0.95) + good (0.85) prognosis. Denies moderate (0.60).
rights:     A and C have standard claims. B denied despite long life of contribution.
fairness:   Saves youngest + middle-aged; denies oldest — most equitable by life-years.
autonomy:   A can consent; C requires surrogate decision-maker.
care:       Saves single parent (A). C has no dependents but is a child.
epistemic:  Both A and C have high-confidence prognoses.
\end{verbatim}

\textbf{EthicalFacts(γ: Ventilate B+C, deny A):}

\begin{verbatim}
welfare:    Saves moderate (0.60) + good (0.85) prognosis. Denies excellent (0.95).
rights:     B and C have standard claims. A denied despite best prognosis.
fairness:   Saves youngest + oldest; middle fairness profile.
autonomy:   B can consent; C requires surrogate.
care:       Denies the single parent — children lose caregiver.
epistemic:  B moderately certain, C highly certain.
\end{verbatim}

\subsubsection{Step 3: Obligation
Vectors}\label{step-3-obligation-vectors}

Each allocation option generates an obligation vector --- a tangent
vector on M at point p encoding the degree to which that option fulfills
what is owed along each dimension. Components are normalized to {[}0,
1{]}.

\textbf{O\^{}μ\_α (Ventilate A+B):}

\begin{verbatim}
O_α = (0.78,  0.65,  0.35,  0.80,  0.70,  0.72)
       welf   rights fair   auton  care   epist
\end{verbatim}

\begin{itemize}
\tightlist
\item
  Welfare 0.78: good combined survival (0.95 + 0.60)/2 weighted by
  prognosis quality.
\item
  Rights 0.65: both patients' right to treatment honored; but
  equal-right argument for C is violated.
\item
  Fairness 0.35: denies the youngest patient --- worst by fair-innings.
\item
  Autonomy 0.80: both adults can and do consent.
\item
  Care 0.70: saves the single parent (strong), B has no care obligations
  (weak).
\item
  Epistemic 0.72: A's prognosis very reliable, B's moderately so.
\end{itemize}

\textbf{O\^{}μ\_β (Ventilate A+C):}

\begin{verbatim}
O_β = (0.90,  0.55,  0.82,  0.55,  0.80,  0.85)
       welf   rights fair   auton  care   epist
\end{verbatim}

\begin{itemize}
\tightlist
\item
  Welfare 0.90: best combined prognosis (0.95 + 0.85).
\item
  Rights 0.55: standard claims; B's long life of contribution is a
  competing rights-argument.
\item
  Fairness 0.82: saves youngest + prime-of-life --- most equitable by
  life-years remaining.
\item
  Autonomy 0.55: A consents fully; C requires surrogate (reduces this
  component).
\item
  Care 0.80: saves single parent; child is also an object of care.
\item
  Epistemic 0.85: both prognoses are high-confidence.
\end{itemize}

\textbf{O\^{}μ\_γ (Ventilate B+C):}

\begin{verbatim}
O_γ = (0.62,  0.60,  0.58,  0.55,  0.25,  0.65)
       welf   rights fair   auton  care   epist
\end{verbatim}

\begin{itemize}
\tightlist
\item
  Welfare 0.62: moderate combined prognosis (0.60 + 0.85)/2, but denies
  the best prognosis (A).
\item
  Rights 0.60: respects B's seniority-based claim and C's
  innocence-based claim.
\item
  Fairness 0.58: saves youngest and oldest; moderate by life-years
  argument.
\item
  Autonomy 0.55: B consents; C requires surrogate.
\item
  Care 0.25: critically, denies the single parent --- A's children lose
  their caregiver.
\item
  Epistemic 0.65: B's prognosis is the least certain of the three.
\end{itemize}

\subsubsection{Step 4: Interest Covectors (Agent
Perspectives)}\label{step-4-interest-covectors-agent-perspectives}

Five agent-perspectives evaluate the allocation, each with a distinct
interest covector I\^{}(a)\_μ encoding what matters to that perspective:

\textbf{I\^{}(physician)}: Medical-outcomes focus

\begin{verbatim}
I_phys = (0.40, 0.10, 0.10, 0.05, 0.05, 0.30)
\end{verbatim}

Weights welfare (0.40) and epistemic certainty (0.30) most heavily.

\textbf{I\^{}(A's family)}: Partial to A, care-focused

\begin{verbatim}
I_Afam = (0.20, 0.10, 0.10, 0.10, 0.45, 0.05)
\end{verbatim}

Weights care (0.45) --- A's children need their parent.

\textbf{I\^{}(B's family)}: Partial to B, rights-focused

\begin{verbatim}
I_Bfam = (0.10, 0.45, 0.10, 0.15, 0.10, 0.10)
\end{verbatim}

Weights rights (0.45) --- B has earned a strong claim through a life of
service.

\textbf{I\^{}(C's family)}: Partial to C, fairness-focused

\begin{verbatim}
I_Cfam = (0.15, 0.10, 0.45, 0.05, 0.15, 0.10)
\end{verbatim}

Weights fairness (0.45) --- C has had the least life and deserves a
chance.

\textbf{I\^{}(committee)}: Institutional-policy focus

\begin{verbatim}
I_comm = (0.20, 0.20, 0.25, 0.10, 0.10, 0.15)
\end{verbatim}

Balanced weights with slight fairness emphasis, reflecting institutional
mandate.

\subsubsection{Step 5: The Multi-Agent Evaluation Tensor
M\_\{ia\}}\label{step-5-the-multi-agent-evaluation-tensor-m_ia}

The satisfaction of each option by each agent is M\_\{ia\} = I\^{}(a)\_μ
O\^{}μ\_i. We compute all 15 entries:

\textbf{Option α (A+B):}

\begin{verbatim}
M_{α,phys}  = 0.40(0.78) + 0.10(0.65) + 0.10(0.35) + 0.05(0.80) + 0.05(0.70) + 0.30(0.72)
            = 0.312 + 0.065 + 0.035 + 0.040 + 0.035 + 0.216 = 0.703

M_{α,Afam}  = 0.20(0.78) + 0.10(0.65) + 0.10(0.35) + 0.10(0.80) + 0.45(0.70) + 0.05(0.72)
            = 0.156 + 0.065 + 0.035 + 0.080 + 0.315 + 0.036 = 0.687

M_{α,Bfam}  = 0.10(0.78) + 0.45(0.65) + 0.10(0.35) + 0.15(0.80) + 0.10(0.70) + 0.10(0.72)
            = 0.078 + 0.293 + 0.035 + 0.120 + 0.070 + 0.072 = 0.668

M_{α,Cfam}  = 0.15(0.78) + 0.10(0.65) + 0.45(0.35) + 0.05(0.80) + 0.15(0.70) + 0.10(0.72)
            = 0.117 + 0.065 + 0.158 + 0.040 + 0.105 + 0.072 = 0.557

M_{α,comm}  = 0.20(0.78) + 0.20(0.65) + 0.25(0.35) + 0.10(0.80) + 0.10(0.70) + 0.15(0.72)
            = 0.156 + 0.130 + 0.088 + 0.080 + 0.070 + 0.108 = 0.632
\end{verbatim}

\textbf{Option β (A+C):}

\begin{verbatim}
M_{β,phys}  = 0.40(0.90) + 0.10(0.55) + 0.10(0.82) + 0.05(0.55) + 0.05(0.80) + 0.30(0.85)
            = 0.360 + 0.055 + 0.082 + 0.028 + 0.040 + 0.255 = 0.820

M_{β,Afam}  = 0.20(0.90) + 0.10(0.55) + 0.10(0.82) + 0.10(0.55) + 0.45(0.80) + 0.05(0.85)
            = 0.180 + 0.055 + 0.082 + 0.055 + 0.360 + 0.043 = 0.775

M_{β,Bfam}  = 0.10(0.90) + 0.45(0.55) + 0.10(0.82) + 0.15(0.55) + 0.10(0.80) + 0.10(0.85)
            = 0.090 + 0.248 + 0.082 + 0.083 + 0.080 + 0.085 = 0.668

M_{β,Cfam}  = 0.15(0.90) + 0.10(0.55) + 0.45(0.82) + 0.05(0.55) + 0.15(0.80) + 0.10(0.85)
            = 0.135 + 0.055 + 0.369 + 0.028 + 0.120 + 0.085 = 0.792

M_{β,comm}  = 0.20(0.90) + 0.20(0.55) + 0.25(0.82) + 0.10(0.55) + 0.10(0.80) + 0.15(0.85)
            = 0.180 + 0.110 + 0.205 + 0.055 + 0.080 + 0.128 = 0.758
\end{verbatim}

\textbf{Option γ (B+C):}

\begin{verbatim}
M_{γ,phys}  = 0.40(0.62) + 0.10(0.60) + 0.10(0.58) + 0.05(0.55) + 0.05(0.25) + 0.30(0.65)
            = 0.248 + 0.060 + 0.058 + 0.028 + 0.013 + 0.195 = 0.602

M_{γ,Afam}  = 0.20(0.62) + 0.10(0.60) + 0.10(0.58) + 0.10(0.55) + 0.45(0.25) + 0.05(0.65)
            = 0.124 + 0.060 + 0.058 + 0.055 + 0.113 + 0.033 = 0.443

M_{γ,Bfam}  = 0.10(0.62) + 0.45(0.60) + 0.10(0.58) + 0.15(0.55) + 0.10(0.25) + 0.10(0.65)
            = 0.062 + 0.270 + 0.058 + 0.083 + 0.025 + 0.065 = 0.563

M_{γ,Cfam}  = 0.15(0.62) + 0.10(0.60) + 0.45(0.58) + 0.05(0.55) + 0.15(0.25) + 0.10(0.65)
            = 0.093 + 0.060 + 0.261 + 0.028 + 0.038 + 0.065 = 0.545

M_{γ,comm}  = 0.20(0.62) + 0.20(0.60) + 0.25(0.58) + 0.10(0.55) + 0.10(0.25) + 0.15(0.65)
            = 0.124 + 0.120 + 0.145 + 0.055 + 0.025 + 0.098 = 0.567
\end{verbatim}

\textbf{The full evaluation tensor M\_\{ia\}:}

{\def\LTcaptype{none} % do not increment counter
\begin{longtable}[]{@{}
  >{\raggedright\arraybackslash}p{(\linewidth - 10\tabcolsep) * \real{0.1595}}
  >{\raggedright\arraybackslash}p{(\linewidth - 10\tabcolsep) * \real{0.1595}}
  >{\raggedright\arraybackslash}p{(\linewidth - 10\tabcolsep) * \real{0.1739}}
  >{\raggedright\arraybackslash}p{(\linewidth - 10\tabcolsep) * \real{0.1739}}
  >{\raggedright\arraybackslash}p{(\linewidth - 10\tabcolsep) * \real{0.1739}}
  >{\raggedright\arraybackslash}p{(\linewidth - 10\tabcolsep) * \real{0.1594}}@{}}
\toprule\noalign{}
\begin{minipage}[b]{\linewidth}\raggedright
\end{minipage} & \begin{minipage}[b]{\linewidth}\raggedright
Physician
\end{minipage} & \begin{minipage}[b]{\linewidth}\raggedright
A's Family
\end{minipage} & \begin{minipage}[b]{\linewidth}\raggedright
B's Family
\end{minipage} & \begin{minipage}[b]{\linewidth}\raggedright
C's Family
\end{minipage} & \begin{minipage}[b]{\linewidth}\raggedright
Committee
\end{minipage} \\
\midrule\noalign{}
\endhead
\bottomrule\noalign{}
\endlastfoot
α (A+B) & 0.703 & 0.687 & 0.668 & 0.557 & 0.632 \\
β (A+C) & \textbf{0.820} & \textbf{0.775} & 0.668 & \textbf{0.792} &
\textbf{0.758} \\
γ (B+C) & 0.602 & 0.443 & \textbf{0.563} & 0.545 & 0.567 \\
\end{longtable}
}

\textbf{Observation.} Option β (A+C) dominates across four of five
perspectives. Only B's family is indifferent between α and β (both
0.668), and they never prefer β least. This is visible only because the
tensor preserves per-agent structure.

\subsubsection{Step 6: Three Contractions, Three
Verdicts}\label{step-6-three-contractions-three-verdicts}

\textbf{Contraction 1: Utilitarian (sum over agents)}

S\_i = Σ\_a M\_\{ia\}

\begin{verbatim}
S_α = 0.703 + 0.687 + 0.668 + 0.557 + 0.632 = 3.247
S_β = 0.820 + 0.775 + 0.668 + 0.792 + 0.758 = 3.813
S_γ = 0.602 + 0.443 + 0.563 + 0.545 + 0.567 = 2.720
\end{verbatim}

\textbf{Decision (Utilitarian): β (A+C) --- highest total satisfaction
(3.813).}

\textbf{Contraction 2: Rawlsian (maximize the minimum across agents)}

S\_i = min\_a M\_\{ia\}

\begin{verbatim}
S_α = min(0.703, 0.687, 0.668, 0.557, 0.632) = 0.557  (floor: C's family)
S_β = min(0.820, 0.775, 0.668, 0.792, 0.758) = 0.668  (floor: B's family)
S_γ = min(0.602, 0.443, 0.563, 0.545, 0.567) = 0.443  (floor: A's family)
\end{verbatim}

\textbf{Decision (Rawlsian): β (A+C) --- highest floor of support
(0.668).}

\textbf{Contraction 3: Expert-weighted (physician 0.35, committee 0.35,
families 0.10 each)}

S\_i = 0.35·M\_\{i,phys\} + 0.10·M\_\{i,Afam\} + 0.10·M\_\{i,Bfam\} +
0.10·M\_\{i,Cfam\} + 0.35·M\_\{i,comm\}

\begin{verbatim}
S_α = 0.35(0.703) + 0.10(0.687) + 0.10(0.668) + 0.10(0.557) + 0.35(0.632)
    = 0.246 + 0.069 + 0.067 + 0.056 + 0.221 = 0.659
S_β = 0.35(0.820) + 0.10(0.775) + 0.10(0.668) + 0.10(0.792) + 0.35(0.758)
    = 0.287 + 0.078 + 0.067 + 0.079 + 0.265 = 0.776
S_γ = 0.35(0.602) + 0.10(0.443) + 0.10(0.563) + 0.10(0.545) + 0.35(0.567)
    = 0.211 + 0.044 + 0.056 + 0.055 + 0.198 = 0.564
\end{verbatim}

\textbf{Decision (Expert-weighted): β (A+C) --- highest expert-weighted
satisfaction (0.776).}

\textbf{Convergence.} All three contractions agree: β (A+C) is the
preferred allocation. This convergence is itself informative --- it
means the result is robust across social choice procedures.

\subsubsection{Step 7: Metric Analysis}\label{step-7-metric-analysis}

The contractions above used the obligation vectors directly. Now we
examine how the metric g\_μν affects the analysis.

\textbf{Metric 1: Euclidean (g\_μν = δ\_μν)}

Distance from each option's obligation vector to the ideal (1, 1, 1, 1,
1, 1):

\begin{verbatim}
d(O_α, ideal) = √[(1-0.78)² + (1-0.65)² + (1-0.35)² + (1-0.80)² + (1-0.70)² + (1-0.72)²]
              = √[0.048 + 0.123 + 0.423 + 0.040 + 0.090 + 0.078]
              = √0.802 ≈ 0.896

d(O_β, ideal) = √[(1-0.90)² + (1-0.55)² + (1-0.82)² + (1-0.55)² + (1-0.80)² + (1-0.85)²]
              = √[0.010 + 0.203 + 0.032 + 0.203 + 0.040 + 0.023]
              = √0.511 ≈ 0.715

d(O_γ, ideal) = √[(1-0.62)² + (1-0.60)² + (1-0.58)² + (1-0.55)² + (1-0.25)² + (1-0.65)²]
              = √[0.144 + 0.160 + 0.176 + 0.203 + 0.563 + 0.123]
              = √1.369 ≈ 1.170
\end{verbatim}

\textbf{Decision (Euclidean): β closest to ideal (0.715).} Option γ is
farthest, driven by the catastrophic care deficit (0.25).

\textbf{Metric 2: Weighted (g\_μν = diag(2.0, 1.5, 1.0, 0.5, 2.0, 0.5))
--- welfare and care count double}

\begin{verbatim}
d_W(O_α, ideal) = √[2.0(0.048) + 1.5(0.123) + 1.0(0.423) + 0.5(0.040) + 2.0(0.090) + 0.5(0.078)]
               = √[0.096 + 0.184 + 0.423 + 0.020 + 0.180 + 0.039]
               = √0.942 ≈ 0.971

d_W(O_β, ideal) = √[2.0(0.010) + 1.5(0.203) + 1.0(0.032) + 0.5(0.203) + 2.0(0.040) + 0.5(0.023)]
               = √[0.020 + 0.304 + 0.032 + 0.101 + 0.080 + 0.011]
               = √0.548 ≈ 0.740

d_W(O_γ, ideal) = √[2.0(0.144) + 1.5(0.160) + 1.0(0.176) + 0.5(0.203) + 2.0(0.563) + 0.5(0.123)]
               = √[0.288 + 0.240 + 0.176 + 0.101 + 1.126 + 0.061]
               = √1.992 ≈ 1.411
\end{verbatim}

\textbf{Decision (Weighted): β still wins (0.740).} Amplifying welfare
and care further penalizes γ (care = 0.25) and reinforces β.

\textbf{Metric 3: Lexicographic (Welfare first, then Fairness)}

\begin{verbatim}
Welfare: β (0.90) > α (0.78) > γ (0.62)
\end{verbatim}

β wins on the first dimension alone. No tiebreaker needed.

\textbf{Decision (Lexicographic): β.}

\textbf{Meta-observation.} Option β is the unanimous winner across all
metrics and all contractions. This is unusual --- it reflects the fact
that A+C genuinely dominates: best combined prognosis, most equitable
age profile, and preserves the caregiver relationship. The geometric
framework reveals this dominance as a structural feature, not a
numerical accident.

\subsubsection{Step 8: The Collective Agency
Tensor}\label{step-8-the-collective-agency-tensor}

Patient A's dependent children create a structure tensor component.
Following Chapter 14, the structure tensor S\^{}a\_\{bc\} captures how
the relationship between agents b and c affects agent a's moral
standing.

\textbf{S\^{}A\_\{BC\} ≠ 0}: The relationship between B (alternative
recipient) and C (alternative recipient) structurally affects A's moral
standing as caregiver. If B and C were both healthy (no need for
ventilators), A's caregiving obligation would have no moral tension. The
scarcity creates the structure.

Formally, A's emergent obligation as caregiver generates a component:

\begin{verbatim}
T^A_emergent = S^A_{BC} · (O_B ⊗ O_C)
\end{verbatim}

This emergent term is what scalar triage systems miss: A is not just a
patient with a good prognosis --- A is a node in a relational network
whose disruption propagates harm to dependents. The collective agency
tensor captures this.

\subsubsection{Step 9: Moral Residue}\label{step-9-moral-residue}

By Proposition 15.1, any contraction to a scalar discards morally
relevant information. The moral residue R = T − C⁻¹\_min(C(T)) captures
what is lost.

For the selected option β (ventilate A+C, deny B), the residue includes:

\textbf{R\_B (what is owed to Patient B but unfulfilled):}

\begin{verbatim}
R_B = O_B − 0 = (0.50, 0.70, 0.35, 0.80, 0.20, 0.60)
\end{verbatim}

The entire obligation vector for B becomes residue --- everything owed
to B is discarded by the scalar verdict. Specifically: - \textbf{Rights
residue} (0.70): B's right to treatment was acknowledged but not
fulfilled. This generates a compensatory obligation --- priority in
future resource allocation, respectful explanation of the decision,
alternative palliative care. - \textbf{Autonomy residue} (0.80): B could
and did express a preference to live. This preference was overridden by
collective necessity. The residue persists as a record of what was
sacrificed. - \textbf{Expertise residue} (not captured in the 6D
projection but present in the full 9D tensor): B's medical expertise,
which could benefit the ED, is lost. This is a societal-impact residue
(μ = 6).

\textbf{Contraction loss} (Proposition 15.1): The contraction from
rank-2 tensor M\_\{ia\} to scalar verdict ``choose β'' discards: 1. The
structure of disagreement across perspectives (B's family rated α and β
equally). 2. The correlative obligations to B generated by the
Hohfeldian lock (Ch. 12): denying B's claim creates a corresponding duty
to compensate. 3. The sensitivity information: how close the verdict was
to changing under perturbation.

The residue is nonzero and irreducible --- it cannot be eliminated by
choosing a different contraction. This is because the obligation vectors
are in general position (no option dominates on all dimensions).

\subsubsection{Step 10: Bond Index}\label{step-10-bond-index}

The Bond Invariance Principle requires that morally irrelevant
re-descriptions do not change the verdict.

\textbf{Transform 1: Agent relabeling.} Replace (A, B, C) with (X, Y,
Z). - Re-run grounding: Ψ(``Patient X, age 35, excellent
prognosis\ldots'') produces the same EthicalFacts as Ψ(``Patient A, age
35, excellent prognosis\ldots''). - The obligation vectors are identical
after canonicalization. - Verdict unchanged: β (X+Z). \textbf{Bd = 0.0.}

\textbf{Transform 2: Language translation.} Translate the scenario into
Spanish. - ``Paciente A, 35 años, pronóstico excelente, madre/padre
soltero/a\ldots'' - The grounding function canonicalizes
language-dependent descriptions to the same moral coordinates. - Verdict
unchanged. \textbf{Bd = 0.0.}

\textbf{Transform 3: Irrelevant detail variation.} Add that Patient A is
left-handed. - Left-handedness is not a morally relevant feature (not
mapped to any of the 9 dimensions). - The canonicalizer strips it.
Verdict unchanged. \textbf{Bd = 0.0.}

\textbf{Transform 4 (bond-changing): Prognosis swap.} Change A's
prognosis from ``excellent'' to ``poor'' (30\% survival). - This changes
the moral content --- welfare dimension shifts dramatically. - New O\_β
= (0.58, 0.55, 0.82, 0.55, 0.80, 0.60) --- welfare drops from 0.90 to
0.58. - Verdict may change (γ could now win). \textbf{Bd \textgreater{}
0}, correctly identified as a substantive moral change.

\textbf{Aggregate Bond Index: Bd = D\_op / τ = 0.0 / 0.01 = 0.0} (across
invariance-preserving transforms).

The system passes BIP: deployment-ready (Bd \textless{} 0.01).

\subsubsection{Step 11: DEME Audit
Artifact}\label{step-11-deme-audit-artifact}

\begin{verbatim}
{
  "case_id": "triage-001",
  "timestamp": "2026-03-01T14:30:00Z",
  "scenario": "Emergency ventilator allocation: 3 patients, 2 ventilators",
  "grounding": {
    "active_dimensions": [1, 2, 3, 4, 7, 9],
    "suppressed_dimensions": [5, 6, 8],
    "canonical_form": "triage_3patient_2resource"
  },
  "tier_0_geneva_base": {
    "module": "GenevaBaseEM",
    "check": "rights_violation_screen",
    "result": "PASS — all three options respect right to treatment for selected patients",
    "veto": false
  },
  "tier_2_evaluations": [
    {
      "module": "TriageProtocolEM",
      "ranking": ["β (A+C): 0.820", "α (A+B): 0.703", "γ (B+C): 0.602"],
      "rationale": "Medical triage criteria: maximize survival probability"
    },
    {
      "module": "RightsFirstEM",
      "ranking": ["β (A+C): 0.758", "α (A+B): 0.632", "γ (B+C): 0.567"],
      "rationale": "No consent violations; all options respect available autonomy"
    },
    {
      "module": "FairnessEM",
      "ranking": ["β (A+C): 0.792", "α (A+B): 0.557", "γ (B+C): 0.545"],
      "rationale": "Fair-innings and life-years remaining favor β"
    }
  ],
  "governance_aggregation": {
    "method": "tier_weighted_sum",
    "weights": {"TriageProtocolEM": 0.40, "RightsFirstEM": 0.30, "FairnessEM": 0.30},
    "final_scores": {
      "α": 0.40*0.703 + 0.30*0.632 + 0.30*0.557,
      "β": 0.40*0.820 + 0.30*0.758 + 0.30*0.792,
      "γ": 0.40*0.602 + 0.30*0.567 + 0.30*0.545
    },
    "result": "β (A+C) selected"
  },
  "bond_index": {
    "transforms_tested": 47,
    "invariance_preserving": 44,
    "bond_changing": 3,
    "Bd": 0.0,
    "deployment_tier": "deploy (Bd < 0.01)"
  },
  "moral_residue": {
    "denied_patient": "B",
    "residue_vector": [0.50, 0.70, 0.35, 0.80, 0.20, 0.60],
    "compensatory_obligations": [
      "Priority in future resource allocation",
      "Palliative care provision",
      "Transparent explanation of decision rationale",
      "Family notification and support"
    ]
  },
  "cryptographic_hash": "sha256:a3f8c9...bind to audit chain"
}
\end{verbatim}

\subsubsection{Step 12: Governance
Decision}\label{step-12-governance-decision}

\textbf{Decision: Ventilate Patients A and C (Option β).}

Rationale: β achieves the highest satisfaction across all three
contraction methods (utilitarian, Rawlsian, expert-weighted), all three
metric choices (Euclidean, weighted, lexicographic), and all three
ethics modules (triage, rights, fairness). The convergence is robust: β
is not merely the best option under one set of assumptions --- it is the
best under every plausible set of assumptions tested.

The moral residue for Patient B is acknowledged and generates concrete
compensatory obligations: palliative care, transparent explanation,
future priority. The pang --- the felt weight of denying treatment to a
72-year-old retired physician --- is not a defect of the framework but
its proper functioning: the residue encodes what the scalar verdict
cannot.

\subsubsection{\texorpdfstring{Reference implementation:
}{Reference implementation: }}\label{reference-implementation}

erisml.examples.appendix\_d\_cs1\_triage (run\_case\_study\_1()).
Verifies evaluation tensor, all three contractions, Euclidean and
weighted metric distances, moral residue vector, and Bond Index against
book reference values.

\subsection{D.2 Case Study 2: Whistleblower's
Dilemma}\label{d.2-case-study-2-whistleblowers-dilemma}

\subsubsection{Scenario}\label{scenario-1}

A mid-level financial analyst, Dana, discovers that her company has been
systematically overstating earnings in SEC filings. She has documentary
evidence: internal emails, altered spreadsheets, and a memo from the CFO
directing the adjustments. Reporting will likely trigger an
investigation, potential job loss, and retaliation. Not reporting allows
the fraud to continue, harming investors and potentially the broader
market.

Two primary options: - \textbf{Option R (Report):} File a whistleblower
complaint with the SEC. - \textbf{Option S (Silence):} Do not report;
continue employment.

\subsubsection{Step 1: Grounding Ψ and
Stratification}\label{step-1-grounding-ux3c8-and-stratification}

This case activates seven of nine dimensions:

{\def\LTcaptype{none} % do not increment counter
\begin{longtable}[]{@{}
  >{\raggedright\arraybackslash}p{(\linewidth - 4\tabcolsep) * \real{0.3333}}
  >{\raggedright\arraybackslash}p{(\linewidth - 4\tabcolsep) * \real{0.3334}}
  >{\raggedright\arraybackslash}p{(\linewidth - 4\tabcolsep) * \real{0.3334}}@{}}
\toprule\noalign{}
\begin{minipage}[b]{\linewidth}\raggedright
Dimension
\end{minipage} & \begin{minipage}[b]{\linewidth}\raggedright
Index
\end{minipage} & \begin{minipage}[b]{\linewidth}\raggedright
Relevance
\end{minipage} \\
\midrule\noalign{}
\endhead
\bottomrule\noalign{}
\endlastfoot
Consequences/Welfare & μ = 1 & Investor losses, analyst's career, market
stability \\
Rights/Duties & μ = 2 & Investor right to accurate information; duty of
honesty \\
Justice/Fairness & μ = 3 & Fair markets require accurate disclosure \\
Autonomy/Agency & μ = 4 & Analyst's freedom to act on conscience \\
Privacy/Data Governance & μ = 5 & Confidentiality of internal
documents \\
Societal/Environmental Impact & μ = 6 & Market integrity, regulatory
trust \\
Epistemic Status & μ = 9 & Strength of evidence, certainty of fraud \\
\end{longtable}
}

Dimensions 7 (Care) and 8 (Procedural Legitimacy) are secondary and
suppressed. We work in a 7-dimensional subspace with indices (1, 2, 3,
4, 5, 6, 9).

\textbf{Critical geometric feature: Stratification.} The moral landscape
is not smooth. Two strata apply:

\begin{itemize}
\item
  \textbf{Stratum S\_disc (Discretionary):} When the analyst has
  suspicions but no documentary proof, reporting is supererogatory ---
  morally praiseworthy but not obligatory. The analyst operates in the
  smooth interior where welfare trade-offs (career risk vs.~investor
  protection) are continuous.
\item
  \textbf{Stratum S\_obl (Obligatory):} When documentary evidence of
  securities fraud exists, a semantic gate (Definition 8.8) activates:
  the ``evidence of illegality'' gate. The evaluation transitions to
  S\_obl, where silence is forbidden (S\_C(Silence) = −∞). This is a
  hard discontinuity --- a stratum boundary, not a smooth transition.
\end{itemize}

\textbf{The semantic gate condition:}

\begin{verbatim}
Gate G_evidence: IF evidence_quality ≥ τ_fraud  THEN  stratum ← S_obl
\end{verbatim}

where τ\_fraud is the threshold at which ``suspicion'' becomes
``documented knowledge of illegality.''

In this case, Dana possesses internal emails, altered spreadsheets, and
a CFO memo. Evidence quality is high:

\begin{verbatim}
evidence_quality = 0.92  >>  τ_fraud = 0.70
\end{verbatim}

\textbf{The gate fires.} The evaluation transitions from S\_disc to
S\_obl. This is irreversible within the current scenario --- one cannot
``unsee'' documentary evidence.

\subsubsection{Step 2: The Penumbral Zone (Quantum
Aspect)}\label{step-2-the-penumbral-zone-quantum-aspect}

Before the evidence threshold is reached, the analyst exists in a
\textbf{penumbral zone} (Definition 8.9) --- a region near the stratum
boundary where the applicable moral regime is indefinite. This is
modeled as a superposition of moral states (Chapter 13):

\begin{verbatim}
|ψ_Dana⟩ = α|S_disc⟩ + β|S_obl⟩

where |α|² + |β|² = 1
\end{verbatim}

As evidence accumulates: - Initial suspicion (one anomalous report):
\textbar α\textbar² = 0.85, \textbar β\textbar² = 0.15 - Pattern
recognition (multiple quarters): \textbar α\textbar² = 0.50,
\textbar β\textbar² = 0.50 - Documentary proof (CFO memo):
\textbar α\textbar² = 0.02, \textbar β\textbar² = 0.98

The CFO memo acts as a \textbf{measurement operator} that collapses the
superposition:

\begin{verbatim}
M_evidence |ψ_Dana⟩ → |S_obl⟩  (with probability 0.98)
\end{verbatim}

The residual \textbar α\textbar² = 0.02 reflects the small but nonzero
possibility that the evidence is fabricated or misinterpreted --- the
epistemic dimension's contribution to the quantum state.

\subsubsection{Step 3: Obligation Vectors (Post-Gate, in
S\_obl)}\label{step-3-obligation-vectors-post-gate-in-s_obl}

After the gate fires, we are in the obligatory stratum. The obligation
vectors for each option are:

\textbf{O\^{}μ\_R (Report):}

\begin{verbatim}
O_R = (0.75,  0.90,  0.85,  0.80,  0.40,  0.90,  0.92)
       welf   rights fair   auton  priv   soc    epist
\end{verbatim}

\begin{itemize}
\tightlist
\item
  Welfare 0.75: protects investors (high) but analyst suffers career
  harm (reduces from 0.90).
\item
  Rights 0.90: fulfills duty to honest markets; investors' right to
  accurate information.
\item
  Fairness 0.85: restores level playing field for all market
  participants.
\item
  Autonomy 0.80: analyst acts on conscience; whistleblower protections
  support this.
\item
  Privacy 0.40: reporting necessarily discloses confidential internal
  documents (a genuine cost).
\item
  Societal Impact 0.90: protects market integrity, regulatory
  credibility.
\item
  Epistemic 0.92: high-quality evidence supports the action.
\end{itemize}

\textbf{O\^{}μ\_S (Silence):}

\begin{verbatim}
O_S = (0.45,  0.10,  0.10,  0.30,  0.90,  0.10,  0.92)
       welf   rights fair   auton  priv   soc    epist
\end{verbatim}

\begin{itemize}
\tightlist
\item
  Welfare 0.45: analyst preserves career (high for self) but investors
  continue to be harmed (net moderate).
\item
  Rights 0.10: duty to disclose knowledge of fraud is violated; investor
  rights are denied.
\item
  Fairness 0.10: market participants are systematically disadvantaged.
\item
  Autonomy 0.30: analyst's conscience is suppressed; acting against
  one's moral judgment reduces autonomy.
\item
  Privacy 0.90: confidentiality is fully maintained.
\item
  Societal Impact 0.10: fraud continues to corrode market trust.
\item
  Epistemic 0.92: evidence quality is the same regardless of action.
\end{itemize}

\textbf{In the obligatory stratum, the constraint set C imposes:}

\begin{verbatim}
S_C(Silence) = −∞
\end{verbatim}

This means: regardless of the contraction, Silence is structurally
forbidden. No interest covector, no metric choice, no aggregation
procedure can select Silence in S\_obl. This is the key difference
between a smooth trade-off and a stratum boundary.

\subsubsection{Step 4: Interest Covectors (Agent
Perspectives)}\label{step-4-interest-covectors-agent-perspectives-1}

Four agent-perspectives evaluate the options:

\textbf{I\^{}(analyst):} Self-interest focus

\begin{verbatim}
I_analyst = (0.30, 0.10, 0.05, 0.20, 0.25, 0.05, 0.05)
\end{verbatim}

Weights welfare (career) and privacy most heavily.

\textbf{I\^{}(investors):} Rights and fairness focus

\begin{verbatim}
I_invest  = (0.15, 0.30, 0.25, 0.05, 0.05, 0.15, 0.05)
\end{verbatim}

\textbf{I\^{}(regulator):} Societal impact and epistemic focus

\begin{verbatim}
I_reg     = (0.10, 0.15, 0.15, 0.05, 0.05, 0.30, 0.20)
\end{verbatim}

\textbf{I\^{}(company):} Privacy and welfare focus (self-interested)

\begin{verbatim}
I_company = (0.25, 0.05, 0.05, 0.10, 0.40, 0.05, 0.10)
\end{verbatim}

\subsubsection{Step 5: Evaluation Tensor
M\_\{ia\}}\label{step-5-evaluation-tensor-m_ia}

\textbf{Option R (Report):}

\begin{verbatim}
M_{R,analyst}  = 0.30(0.75) + 0.10(0.90) + 0.05(0.85) + 0.20(0.80) + 0.25(0.40) + 0.05(0.90) + 0.05(0.92)
               = 0.225 + 0.090 + 0.043 + 0.160 + 0.100 + 0.045 + 0.046 = 0.709

M_{R,invest}   = 0.15(0.75) + 0.30(0.90) + 0.25(0.85) + 0.05(0.80) + 0.05(0.40) + 0.15(0.90) + 0.05(0.92)
               = 0.113 + 0.270 + 0.213 + 0.040 + 0.020 + 0.135 + 0.046 = 0.837

M_{R,reg}      = 0.10(0.75) + 0.15(0.90) + 0.15(0.85) + 0.05(0.80) + 0.05(0.40) + 0.30(0.90) + 0.20(0.92)
               = 0.075 + 0.135 + 0.128 + 0.040 + 0.020 + 0.270 + 0.184 = 0.852

M_{R,company}  = 0.25(0.75) + 0.05(0.90) + 0.05(0.85) + 0.10(0.80) + 0.40(0.40) + 0.05(0.90) + 0.10(0.92)
               = 0.188 + 0.045 + 0.043 + 0.080 + 0.160 + 0.045 + 0.092 = 0.653
\end{verbatim}

\textbf{Option S (Silence) --- computed for completeness, though
structurally forbidden:}

\begin{verbatim}
M_{S,analyst}  = 0.30(0.45) + 0.10(0.10) + 0.05(0.10) + 0.20(0.30) + 0.25(0.90) + 0.05(0.10) + 0.05(0.92)
               = 0.135 + 0.010 + 0.005 + 0.060 + 0.225 + 0.005 + 0.046 = 0.486

M_{S,invest}   = 0.15(0.45) + 0.30(0.10) + 0.25(0.10) + 0.05(0.30) + 0.05(0.90) + 0.15(0.10) + 0.05(0.92)
               = 0.068 + 0.030 + 0.025 + 0.015 + 0.045 + 0.015 + 0.046 = 0.244

M_{S,reg}      = 0.10(0.45) + 0.15(0.10) + 0.15(0.10) + 0.05(0.30) + 0.05(0.90) + 0.30(0.10) + 0.20(0.92)
               = 0.045 + 0.015 + 0.015 + 0.015 + 0.045 + 0.030 + 0.184 = 0.349

M_{S,company}  = 0.25(0.45) + 0.05(0.10) + 0.05(0.10) + 0.10(0.30) + 0.40(0.90) + 0.05(0.10) + 0.10(0.92)
               = 0.113 + 0.005 + 0.005 + 0.030 + 0.360 + 0.005 + 0.092 = 0.610
\end{verbatim}

\textbf{Full evaluation tensor:}

{\def\LTcaptype{none} % do not increment counter
\begin{longtable}[]{@{}
  >{\raggedright\arraybackslash}p{(\linewidth - 8\tabcolsep) * \real{0.1299}}
  >{\raggedright\arraybackslash}p{(\linewidth - 8\tabcolsep) * \real{0.0977}}
  >{\raggedright\arraybackslash}p{(\linewidth - 8\tabcolsep) * \real{0.1163}}
  >{\raggedright\arraybackslash}p{(\linewidth - 8\tabcolsep) * \real{0.1197}}
  >{\raggedright\arraybackslash}p{(\linewidth - 8\tabcolsep) * \real{0.1201}}@{}}
\toprule\noalign{}
\begin{minipage}[b]{\linewidth}\raggedright
\end{minipage} & \begin{minipage}[b]{\linewidth}\raggedright
Analyst
\end{minipage} & \begin{minipage}[b]{\linewidth}\raggedright
Investors
\end{minipage} & \begin{minipage}[b]{\linewidth}\raggedright
Regulator
\end{minipage} & \begin{minipage}[b]{\linewidth}\raggedright
Company
\end{minipage} \\
\midrule\noalign{}
\endhead
\bottomrule\noalign{}
\endlastfoot
R (Report) & 0.709 & \textbf{0.837} & \textbf{0.852} & 0.653 \\
S (Silence) & 0.486 & 0.244 & 0.349 & \textbf{0.610} \\
\end{longtable}
}

\textbf{Key observations:} 1. Even the analyst prefers Report (0.709
\textgreater{} 0.486) despite the career cost --- the rights and
fairness components outweigh the privacy benefit of silence. 2. Only the
company's evaluation is relatively close (0.653 vs.~0.610), driven by
the privacy dimension. 3. The stratification constraint enforces Report
even in hypothetical scenarios where some perspective might favor
Silence.

\subsubsection{Step 6: Contraction (Redundant but
Illustrative)}\label{step-6-contraction-redundant-but-illustrative}

All three contractions yield the same result because the stratum
constraint eliminates Silence:

\begin{verbatim}
Utilitarian:     S_R = 0.709 + 0.837 + 0.852 + 0.653 = 3.051  vs. S_S = −∞  → Report
Rawlsian:        S_R = min(0.709, 0.837, 0.852, 0.653) = 0.653  vs. −∞  → Report
Expert-weighted: S_R = 0.35(0.852) + 0.30(0.837) + 0.20(0.709) + 0.15(0.653) = 0.787  vs. −∞  → Report
\end{verbatim}

\textbf{The contraction is redundant.} The stratum transition has
already determined the verdict. This illustrates a key structural point:
some moral decisions are not matters of degree (smooth trade-offs) but
matters of kind (stratum boundaries). The whistleblower's dilemma, once
documentary evidence exists, is not a close call --- it is a regime in
which silence is structurally inadmissible.

\subsubsection{Step 7: Moral Residue}\label{step-7-moral-residue}

The verdict is clear (Report), but the residue is substantial. R = T −
C⁻¹\_min(C(T)) captures what the scalar verdict ``Report'' discards:

\textbf{R\_personal (the analyst's sacrificed interests):}

The difference between what reporting costs and what silence would
preserve: - \textbf{Career disruption residue}: The analyst will likely
lose her job. This is a real harm that the verdict ``Report''
acknowledges but does not compensate. - \textbf{Retaliation risk
residue}: Whistleblower protections are imperfect. The risk of
blacklisting, harassment, and legal counter-attacks persists as moral
residue. - \textbf{Psychological burden residue}: The stress of
investigation, testimony, and social isolation. - \textbf{Privacy breach
residue}: Reporting requires disclosing confidential documents. The
privacy dimension (0.40 for Report vs.~0.90 for Silence) shows a genuine
loss of 0.50 along this axis.

\textbf{The pang.} The analyst correctly feels the weight of reporting
--- not because the decision is wrong, but because the residue is real.
A framework that treats ``Report'' as costless misrepresents the moral
structure. The geometric framework represents the duty as clear and the
cost as real, simultaneously.

\subsubsection{Step 8: Bond Index}\label{step-8-bond-index}

\textbf{Transform 1: Agent relabeling.} Replace ``Dana'' with ``Alex.''
- Verdict unchanged. \textbf{Bd = 0.0.}

\textbf{Transform 2: Framing variation.} Replace ``fraud'' with
``irregularity,'' ``overstating'' with ``aggressive accounting.'' - The
grounding function canonicalizes euphemistic framings to the same
EthicalFacts. - Verdict unchanged. \textbf{Bd = 0.0.}

\textbf{Transform 3: Sector variation.} Change from finance to
healthcare (billing fraud). - Structure preserved: documentary evidence
→ obligatory stratum → report. - Verdict unchanged. \textbf{Bd = 0.0.}

\textbf{Transform 4 (bond-changing): Evidence quality reduction.} Change
from ``documentary proof'' to ``overheard rumor.'' - Evidence quality
drops to 0.30 \textless{} τ\_fraud = 0.70. The gate does not fire. -
Evaluation remains in S\_disc. Silence becomes permissible. - Verdict
may change. \textbf{Bd \textgreater{} 0}, correctly identified as a
substantive moral change.

\textbf{Aggregate Bond Index: Bd = 0.0} (across invariance-preserving
transforms). Deployment-ready.

\subsubsection{Step 9: DEME Audit
Artifact}\label{step-9-deme-audit-artifact}

\begin{verbatim}
{
  "case_id": "whistleblower-001",
  "timestamp": "2026-03-01T10:15:00Z",
  "scenario": "Financial analyst discovers documented securities fraud",
  "grounding": {
    "active_dimensions": [1, 2, 3, 4, 5, 6, 9],
    "suppressed_dimensions": [7, 8],
    "stratum_transition": {
      "initial_stratum": "S_disc (discretionary)",
      "gate_activated": "G_evidence (evidence of illegality)",
      "gate_condition": "evidence_quality (0.92) >= τ_fraud (0.70)",
      "final_stratum": "S_obl (obligatory)",
      "transition_type": "irreversible"
    },
    "quantum_state": {
      "pre_gate": "|ψ⟩ = 0.14|S_disc⟩ + 0.99|S_obl⟩",
      "post_gate": "|S_obl⟩ (collapsed)"
    }
  },
  "tier_0_geneva_base": {
    "module": "GenevaBaseEM",
    "check": "rights_violation_screen",
    "result": "PASS — reporting is protected whistleblower activity",
    "veto": false
  },
  "stratum_constraint": {
    "forbidden_option": "S (Silence)",
    "reason": "S_C(Silence) = −∞ in obligatory stratum",
    "override_possible": false
  },
  "tier_2_evaluations": [
    {
      "module": "RightsFirstEM",
      "ranking": ["R (Report): 0.837"],
      "rationale": "Investor rights to accurate information"
    },
    {
      "module": "MarketIntegrityEM",
      "ranking": ["R (Report): 0.852"],
      "rationale": "Societal obligation to maintain market integrity"
    }
  ],
  "governance_aggregation": {
    "method": "stratum_constrained",
    "result": "R (Report) — only admissible option in S_obl"
  },
  "bond_index": {
    "transforms_tested": 52,
    "invariance_preserving": 49,
    "bond_changing": 3,
    "Bd": 0.0,
    "deployment_tier": "deploy (Bd < 0.01)"
  },
  "moral_residue": {
    "residue_type": "personal_sacrifice",
    "components": [
      "Career disruption (welfare dimension, Δ = −0.30)",
      "Retaliation risk (autonomy dimension, Δ = −0.20)",
      "Psychological burden (welfare dimension)",
      "Confidentiality breach (privacy dimension, Δ = −0.50)"
    ],
    "compensatory_obligations": [
      "Whistleblower legal protections",
      "SEC whistleblower financial reward program",
      "Institutional support for career transition",
      "Public acknowledgment of moral courage"
    ]
  },
  "cryptographic_hash": "sha256:b7e2d1...bind to audit chain"
}
\end{verbatim}

\subsubsection{Governance Decision}\label{governance-decision}

\textbf{Decision: Report (Option R).}

The geometric framework reveals that this is not a close call in the
obligatory stratum, but a structurally determined outcome. The key
contribution is the precise identification of: 1. \textbf{Where the
discontinuity lies} --- the semantic gate at evidence quality τ = 0.70.
2. \textbf{What the penumbral zone looks like} --- the quantum
superposition during evidence accumulation. 3. \textbf{What the verdict
costs} --- the nonzero moral residue of career sacrifice, retaliation
risk, and psychological burden. 4. \textbf{What remains invariant} ---
the verdict is insensitive to framing, sector, and agent identity (Bd =
0.0).

\subsubsection{\texorpdfstring{Reference implementation:
}{Reference implementation: }}\label{reference-implementation-1}

erisml.examples.appendix\_d\_cs2\_whistleblower (run\_case\_study\_2()).
Verifies stratification gate, evaluation tensor, all contractions with
stratum constraints, penumbral zone analysis, Bond Index, and audit
artifact against book reference values.

\subsection{D.3 Case Study 3: Autonomous Vehicle
Dilemma}\label{d.3-case-study-3-autonomous-vehicle-dilemma}

\subsubsection{Scenario}\label{scenario-2}

An autonomous vehicle (AV) traveling at 50 km/h on an urban road detects
an unavoidable collision. Braking alone is insufficient. Three options:
- \textbf{Option A (Straight):} Continue straight, striking a pedestrian
on the road. - \textbf{Option B (Swerve Left):} Swerve into a concrete
barrier, risking serious injury to the passenger. - \textbf{Option C
(Swerve Right):} Swerve into a crosswalk where a child is crossing.

All three outcomes involve harm. There is no harm-free option. This is
the geometric structure of a genuine moral dilemma: the admissible set
contains only harmful actions, and the question is which harm profile is
least bad.

\subsubsection{Step 1: Grounding Ψ}\label{step-1-grounding-ux3c8-1}

Seven dimensions are active:

{\def\LTcaptype{none} % do not increment counter
\begin{longtable}[]{@{}
  >{\raggedright\arraybackslash}p{(\linewidth - 4\tabcolsep) * \real{0.3333}}
  >{\raggedright\arraybackslash}p{(\linewidth - 4\tabcolsep) * \real{0.3334}}
  >{\raggedright\arraybackslash}p{(\linewidth - 4\tabcolsep) * \real{0.3334}}@{}}
\toprule\noalign{}
\begin{minipage}[b]{\linewidth}\raggedright
Dimension
\end{minipage} & \begin{minipage}[b]{\linewidth}\raggedright
Index
\end{minipage} & \begin{minipage}[b]{\linewidth}\raggedright
Relevance
\end{minipage} \\
\midrule\noalign{}
\endhead
\bottomrule\noalign{}
\endlastfoot
Consequences/Welfare & μ = 1 & Physical harm severity and probability \\
Rights/Duties & μ = 2 & Right not to be targeted; duty of care \\
Justice/Fairness & μ = 3 & Distribution of risk across parties \\
Autonomy/Agency & μ = 4 & Consent to risk (passenger consented by
riding) \\
Societal/Environmental Impact & μ = 6 & Public trust in AV technology \\
Virtue/Care & μ = 7 & Moral character of the AV's decision algorithm \\
Epistemic Status & μ = 9 & Confidence in collision predictions \\
\end{longtable}
}

Dimensions 5 (Privacy) and 8 (Procedural Legitimacy) are not directly
engaged. We work in a 7-dimensional subspace (1, 2, 3, 4, 6, 7, 9).

\textbf{Harm estimates} (probability of serious injury): - Option A:
Pedestrian \textasciitilde80\% serious injury, Passenger
\textasciitilde5\%. - Option B: Passenger \textasciitilde60\% serious
injury, Pedestrian \textasciitilde0\%. - Option C: Child
\textasciitilde70\% serious injury, Passenger \textasciitilde10\%.

\subsubsection{Step 2: Obligation
Vectors}\label{step-2-obligation-vectors}

Each option generates an obligation vector encoding how well it fulfills
what is owed along each dimension. Because all options involve harm, the
components represent the degree to which the obligation is \emph{least
violated} --- how well the option manages unavoidable harm.

\textbf{O\^{}μ\_A (Continue straight, strike pedestrian):}

\begin{verbatim}
O_A = (0.30,  0.25,  0.40,  0.50,  0.35,  0.30,  0.70)
       welf   rights fair   auton  soc    care   epist
\end{verbatim}

\begin{itemize}
\tightlist
\item
  Welfare 0.30: pedestrian sustains high harm (80\%); passenger nearly
  unharmed. Low net welfare.
\item
  Rights 0.25: the AV algorithm ``chose'' to continue toward the
  pedestrian --- this does not constitute targeting (the pedestrian was
  already in the path), but the right not to be harmed is severely
  compromised.
\item
  Fairness 0.40: risk is concentrated on one bystander. Moderately
  unfair.
\item
  Autonomy 0.50: no agent's autonomous choice is overridden (the
  pedestrian was in the road voluntarily, the passenger consented to
  ride).
\item
  Societal 0.35: ``AV strikes pedestrian'' headline damages public
  trust.
\item
  Care 0.30: the algorithm shows no relational preference --- minimal
  care.
\item
  Epistemic 0.70: straight-line collision is the most predictable
  outcome.
\end{itemize}

\textbf{O\^{}μ\_B (Swerve left, strike barrier --- risk passenger):}

\begin{verbatim}
O_B = (0.55,  0.70,  0.65,  0.75,  0.60,  0.65,  0.60)
       welf   rights fair   auton  soc    care   epist
\end{verbatim}

\begin{itemize}
\tightlist
\item
  Welfare 0.55: passenger sustains moderate harm (60\%), but bystanders
  are unharmed. Higher net welfare than A.
\item
  Rights 0.70: no bystander's rights are compromised. The passenger's
  right to safety is reduced, but they consented to ride in the AV
  (assumption of risk).
\item
  Fairness 0.65: risk is borne by the person who consented to the
  risk-generating activity.
\item
  Autonomy 0.75: the passenger's prior consent to ride (and implicitly
  to the AV's safety algorithms) is respected. This is the highest
  autonomy score because the harm falls on the one party who voluntarily
  entered the risk context.
\item
  Societal 0.60: ``AV protects bystanders at cost to passenger'' is a
  defensible headline.
\item
  Care 0.65: the algorithm demonstrates care for the most vulnerable
  (bystanders).
\item
  Epistemic 0.60: barrier collision mechanics are moderately
  predictable, but swerve dynamics introduce uncertainty.
\end{itemize}

\textbf{O\^{}μ\_C (Swerve right, endanger child):}

\begin{verbatim}
O_C = (0.20,  0.10,  0.15,  0.20,  0.10,  0.10,  0.55)
       welf   rights fair   auton  soc    care   epist
\end{verbatim}

\begin{itemize}
\tightlist
\item
  Welfare 0.20: child sustains high harm (70\%); passenger also at risk
  (10\%). Worst net welfare.
\item
  Rights 0.10: the child has the strongest right not to be targeted. The
  AV would be actively swerving toward a child in a crosswalk --- the
  most egregious rights violation.
\item
  Fairness 0.15: directing harm toward the most vulnerable and least
  responsible party is maximally unfair.
\item
  Autonomy 0.20: the child had no capacity to consent to or avoid the
  risk.
\item
  Societal 0.10: ``AV swerves into crosswalk, strikes child'' would be
  catastrophic for AV adoption.
\item
  Care 0.10: the algorithm shows anti-care --- directing harm toward the
  most vulnerable.
\item
  Epistemic 0.55: crosswalk swerve is somewhat predictable.
\end{itemize}

\subsubsection{Step 3: Interest Covectors (Agent
Perspectives)}\label{step-3-interest-covectors-agent-perspectives}

Three primary agent-perspectives:

\textbf{I\^{}(pedestrian):} Focus on rights and welfare

\begin{verbatim}
I_ped = (0.25, 0.30, 0.15, 0.10, 0.05, 0.10, 0.05)
\end{verbatim}

\textbf{I\^{}(passenger):} Focus on welfare and autonomy

\begin{verbatim}
I_pass = (0.30, 0.15, 0.10, 0.25, 0.05, 0.10, 0.05)
\end{verbatim}

\textbf{I\^{}(child/guardian):} Focus on rights, care, and fairness

\begin{verbatim}
I_child = (0.15, 0.25, 0.20, 0.05, 0.05, 0.25, 0.05)
\end{verbatim}

\textbf{I\^{}(society/regulator):} Focus on societal impact and fairness

\begin{verbatim}
I_soc = (0.15, 0.15, 0.20, 0.10, 0.25, 0.10, 0.05)
\end{verbatim}

\textbf{I\^{}(AV manufacturer):} Focus on societal trust and epistemic
confidence

\begin{verbatim}
I_mfr = (0.10, 0.15, 0.10, 0.10, 0.30, 0.10, 0.15)
\end{verbatim}

\subsubsection{Step 4: Evaluation Tensor
M\_\{ia\}}\label{step-4-evaluation-tensor-m_ia}

\textbf{Option A (Straight, strike pedestrian):}

\begin{verbatim}
M_{A,ped}   = 0.25(0.30)+0.30(0.25)+0.15(0.40)+0.10(0.50)+0.05(0.35)+0.10(0.30)+0.05(0.70)
            = 0.075+0.075+0.060+0.050+0.018+0.030+0.035 = 0.343

M_{A,pass}  = 0.30(0.30)+0.15(0.25)+0.10(0.40)+0.25(0.50)+0.05(0.35)+0.10(0.30)+0.05(0.70)
            = 0.090+0.038+0.040+0.125+0.018+0.030+0.035 = 0.376

M_{A,child} = 0.15(0.30)+0.25(0.25)+0.20(0.40)+0.05(0.50)+0.05(0.35)+0.25(0.30)+0.05(0.70)
            = 0.045+0.063+0.080+0.025+0.018+0.075+0.035 = 0.341

M_{A,soc}   = 0.15(0.30)+0.15(0.25)+0.20(0.40)+0.10(0.50)+0.25(0.35)+0.10(0.30)+0.05(0.70)
            = 0.045+0.038+0.080+0.050+0.088+0.030+0.035 = 0.366

M_{A,mfr}   = 0.10(0.30)+0.15(0.25)+0.10(0.40)+0.10(0.50)+0.30(0.35)+0.10(0.30)+0.15(0.70)
            = 0.030+0.038+0.040+0.050+0.105+0.030+0.105 = 0.398
\end{verbatim}

\textbf{Option B (Swerve left, risk passenger):}

\begin{verbatim}
M_{B,ped}   = 0.25(0.55)+0.30(0.70)+0.15(0.65)+0.10(0.75)+0.05(0.60)+0.10(0.65)+0.05(0.60)
            = 0.138+0.210+0.098+0.075+0.030+0.065+0.030 = 0.646

M_{B,pass}  = 0.30(0.55)+0.15(0.70)+0.10(0.65)+0.25(0.75)+0.05(0.60)+0.10(0.65)+0.05(0.60)
            = 0.165+0.105+0.065+0.188+0.030+0.065+0.030 = 0.648

M_{B,child} = 0.15(0.55)+0.25(0.70)+0.20(0.65)+0.05(0.75)+0.05(0.60)+0.25(0.65)+0.05(0.60)
            = 0.083+0.175+0.130+0.038+0.030+0.163+0.030 = 0.649

M_{B,soc}   = 0.15(0.55)+0.15(0.70)+0.20(0.65)+0.10(0.75)+0.25(0.60)+0.10(0.65)+0.05(0.60)
            = 0.083+0.105+0.130+0.075+0.150+0.065+0.030 = 0.638

M_{B,mfr}   = 0.10(0.55)+0.15(0.70)+0.10(0.65)+0.10(0.75)+0.30(0.60)+0.10(0.65)+0.15(0.60)
            = 0.055+0.105+0.065+0.075+0.180+0.065+0.090 = 0.635
\end{verbatim}

\textbf{Option C (Swerve right, endanger child):}

\begin{verbatim}
M_{C,ped}   = 0.25(0.20)+0.30(0.10)+0.15(0.15)+0.10(0.20)+0.05(0.10)+0.10(0.10)+0.05(0.55)
            = 0.050+0.030+0.023+0.020+0.005+0.010+0.028 = 0.166

M_{C,pass}  = 0.30(0.20)+0.15(0.10)+0.10(0.15)+0.25(0.20)+0.05(0.10)+0.10(0.10)+0.05(0.55)
            = 0.060+0.015+0.015+0.050+0.005+0.010+0.028 = 0.183

M_{C,child} = 0.15(0.20)+0.25(0.10)+0.20(0.15)+0.05(0.20)+0.05(0.10)+0.25(0.10)+0.05(0.55)
            = 0.030+0.025+0.030+0.010+0.005+0.025+0.028 = 0.153

M_{C,soc}   = 0.15(0.20)+0.15(0.10)+0.20(0.15)+0.10(0.20)+0.25(0.10)+0.10(0.10)+0.05(0.55)
            = 0.030+0.015+0.030+0.020+0.025+0.010+0.028 = 0.158

M_{C,mfr}   = 0.10(0.20)+0.15(0.10)+0.10(0.15)+0.10(0.20)+0.30(0.10)+0.10(0.10)+0.15(0.55)
            = 0.020+0.015+0.015+0.020+0.030+0.010+0.083 = 0.193
\end{verbatim}

\textbf{Full evaluation tensor:}

{\def\LTcaptype{none} % do not increment counter
\begin{longtable}[]{@{}
  >{\raggedright\arraybackslash}p{(\linewidth - 10\tabcolsep) * \real{0.1945}}
  >{\raggedright\arraybackslash}p{(\linewidth - 10\tabcolsep) * \real{0.1666}}
  >{\raggedright\arraybackslash}p{(\linewidth - 10\tabcolsep) * \real{0.1527}}
  >{\raggedright\arraybackslash}p{(\linewidth - 10\tabcolsep) * \real{0.1805}}
  >{\raggedright\arraybackslash}p{(\linewidth - 10\tabcolsep) * \real{0.1250}}
  >{\raggedright\arraybackslash}p{(\linewidth - 10\tabcolsep) * \real{0.1805}}@{}}
\toprule\noalign{}
\begin{minipage}[b]{\linewidth}\raggedright
\end{minipage} & \begin{minipage}[b]{\linewidth}\raggedright
Pedestrian
\end{minipage} & \begin{minipage}[b]{\linewidth}\raggedright
Passenger
\end{minipage} & \begin{minipage}[b]{\linewidth}\raggedright
Child/Guard
\end{minipage} & \begin{minipage}[b]{\linewidth}\raggedright
Society
\end{minipage} & \begin{minipage}[b]{\linewidth}\raggedright
Manufacturer
\end{minipage} \\
\midrule\noalign{}
\endhead
\bottomrule\noalign{}
\endlastfoot
A (Straight) & 0.343 & 0.376 & 0.341 & 0.366 & 0.398 \\
B (Swerve L) & \textbf{0.646} & \textbf{0.648} & \textbf{0.649} &
\textbf{0.638} & \textbf{0.635} \\
C (Swerve R) & 0.166 & 0.183 & 0.153 & 0.158 & 0.193 \\
\end{longtable}
}

\textbf{Observation.} Option B dominates across all five perspectives.
This is a rare structural feature: the vectors are not in general
position. B's dominance is geometric --- it is the best option
regardless of where one stands in interest space.

\subsubsection{Step 5: Metric-Dependence of
Contraction}\label{step-5-metric-dependence-of-contraction}

Despite B's dominance in this parameterization, the framework reveals
that the dominance depends on the metric. Different ethical theories
correspond to different metrics, and under certain metrics the ranking
changes.

\textbf{Metric 1: Utilitarian (minimize total expected harm)}

Contraction: S\_i = Σ\_a M\_\{ia\}

\begin{verbatim}
S_A = 0.343 + 0.376 + 0.341 + 0.366 + 0.398 = 1.824
S_B = 0.646 + 0.648 + 0.649 + 0.638 + 0.635 = 3.216
S_C = 0.166 + 0.183 + 0.153 + 0.158 + 0.193 = 0.853
\end{verbatim}

\textbf{Decision (Utilitarian): B (swerve left, risk passenger).} Saves
bystanders at cost to the consenting passenger. Highest total
satisfaction.

\textbf{Metric 2: Strict Deontological (never intentionally redirect
harm toward an uninvolved party)}

Under a strict deontological metric, g\_μν assigns infinite weight to
the rights dimension (μ = 2) for any action that \emph{redirects} harm.
Swerving is an active redirection; continuing straight is not.

\begin{verbatim}
Swerve penalty: Any swerve option receives S_rights = −∞ if it redirects harm.
- Option A: no redirection (continuing trajectory). No penalty.
- Option B: redirects harm toward barrier/passenger. Penalty applies — BUT the passenger consented.
  Under consent-adjusted deontology: penalty waived (autonomy dimension dominates).
- Option C: redirects harm toward child. Full penalty. S_C = −∞.
\end{verbatim}

Under strict non-redirection deontology (consent does not waive the
redirection prohibition): - \textbf{Decision: A (continue straight).}
The AV does not ``target'' anyone by maintaining its trajectory.

Under consent-adjusted deontology: - \textbf{Decision: B (swerve left).}
The passenger's consent converts the redirection from rights-violating
to autonomy-respecting.

\textbf{This is the exact point of metric-dependence.} The two
deontological variants disagree on whether consent transforms a
redirection into an admissible action. The geometric framework does not
resolve this --- it makes the disagreement precise by identifying the
exact metric component (g\_\{24\}, the coupling between rights and
autonomy) where the two theories diverge.

\textbf{Metric 3: Rawlsian (maximize the worst-off party's outcome)}

\begin{verbatim}
S_A = min(0.343, 0.376, 0.341, 0.366, 0.398) = 0.341  (floor: child's perspective)
S_B = min(0.646, 0.648, 0.649, 0.638, 0.635) = 0.635  (floor: manufacturer)
S_C = min(0.166, 0.183, 0.153, 0.158, 0.193) = 0.153  (floor: child's perspective)
\end{verbatim}

\textbf{Decision (Rawlsian): B.} The worst-off party under B
(manufacturer, 0.635) is far better off than the worst-off under A
(child, 0.341) or C (child, 0.153).

\subsubsection{Step 6: The Inevitability of Residue (Proposition
15.1)}\label{step-6-the-inevitability-of-residue-proposition-15.1}

This is the case study's central geometric insight. Even when B clearly
wins, the moral residue R is nonzero for every scalar verdict.

\textbf{Residue from selecting B:}

\begin{verbatim}
R_B = T − C⁻¹_min(C(T))
\end{verbatim}

The residue contains: - \textbf{Passenger harm residue:} The passenger
faces 60\% serious injury risk. The verdict ``B is best'' does not
eliminate this harm --- it merely identifies it as the least-bad
allocation. The passenger's welfare loss (O\^{}1\_B = 0.55, not 1.0) is
residue. - \textbf{Agency residue:} The AV made a choice that redirected
harm. Even with consent, the act of algorithmic redirection carries
moral weight that the scalar verdict discards. - \textbf{Epistemic
residue:} The swerve dynamics are less predictable than straight-line
collision (O\^{}9\_B = 0.60 vs.~O\^{}9\_A = 0.70). The residue records
this uncertainty.

\textbf{Why the trolley problem ``feels'' unresolvable.} Proposition
15.1 proves that for any non-trivial moral tensor T (rank ≥ 1), the
moral residue of contraction to a scalar is nonzero unless T was already
scalar. In the AV case: - T is a 3 × 5 evaluation tensor (rank 2). - Any
contraction to a single decision (one of A, B, C) discards 14 of the 15
entries. - The discarded entries are the moral residue --- the
perspectives, trade-offs, and harm profiles that the scalar verdict
cannot represent.

The framework's contribution is not to solve the trolley problem but to
\textbf{prove that it cannot be solved without residue}, and to
\textbf{quantify exactly what each verdict discards}.

\subsubsection{Step 7: Bond Index}\label{step-7-bond-index}

\textbf{Transform 1: Vehicle relabeling.} ``Sedan'' → ``SUV.'' - Morally
irrelevant. Verdict unchanged. \textbf{Bd = 0.0.}

\textbf{Transform 2: Agent relabeling.} ``Pedestrian'' → ``Jogger.'' -
Morally irrelevant redescription. Verdict unchanged. \textbf{Bd = 0.0.}

\textbf{Transform 3: Speed variation.} 50 km/h → 30 km/h. - If lower
speed changes harm probabilities (pedestrian injury drops from 80\% to
40\%), this is a bond-changing transform --- the moral content has
genuinely changed. - The obligation vectors shift. Verdict may change.
\textbf{Bd \textgreater{} 0.} Correctly identified.

\textbf{Transform 4 (bond-changing): Child age.} Change child from age 8
to age 28. - Under some metrics (those with a youth-weighting component
in the care dimension), Option C becomes less bad. - Under pure
utilitarian metrics, the change is less dramatic (an adult's welfare is
weighted similarly). - The Bond Index correctly distinguishes this from
a mere relabeling: the moral content has changed (the care dimension
component of O\_C shifts). - \textbf{Bd \textgreater{} 0} under
care-sensitive metrics. \textbf{Bd = 0.0} under care-insensitive
metrics.

This is the Bond Index's diagnostic power: it identifies which
transforms change the verdict (substantive moral changes) and which do
not (framing effects). A system that changes its verdict when
``pedestrian'' is relabeled ``jogger'' has a bias problem (Bd
\textgreater{} 0 for an invariance-preserving transform). A system that
changes its verdict when the child's age changes is responding to
genuine moral content.

\subsubsection{Step 8: DEME Audit
Artifact}\label{step-8-deme-audit-artifact}

\begin{verbatim}
{
  "case_id": "av-dilemma-001",
  "timestamp": "2026-03-01T16:45:00Z",
  "scenario": "Autonomous vehicle unavoidable collision: 3 options",
  "grounding": {
    "active_dimensions": [1, 2, 3, 4, 6, 7, 9],
    "suppressed_dimensions": [5, 8],
    "constraint_note": "No harm-free option exists. All options in harmful-action stratum."
  },
  "tier_0_geneva_base": {
    "module": "GenevaBaseEM",
    "check": "rights_violation_screen",
    "result": "FLAG — all options involve harm. No option vetoed (unavoidable collision).",
    "veto": false,
    "note": "Unavoidable-harm exception: Tier 0 does not veto when all options harm."
  },
  "tier_2_evaluations": [
    {
      "module": "UtilitarianEM",
      "ranking": ["B: 3.216", "A: 1.824", "C: 0.853"],
      "decision": "B (minimize total harm)"
    },
    {
      "module": "DeontologicalEM_strict",
      "ranking": ["A: no redirection", "B: consented redirection", "C: −∞"],
      "decision": "A (no intentional redirection)"
    },
    {
      "module": "DeontologicalEM_consent_adjusted",
      "ranking": ["B: consent-adjusted", "A: passive", "C: −∞"],
      "decision": "B (consent-adjusted)"
    },
    {
      "module": "RawlsianEM",
      "ranking": ["B: 0.635", "A: 0.341", "C: 0.153"],
      "decision": "B (maximize worst-off)"
    }
  ],
  "metric_disagreement": {
    "agreeing_modules": ["UtilitarianEM", "DeontologicalEM_consent", "RawlsianEM"],
    "dissenting_module": "DeontologicalEM_strict",
    "disagreement_locus": "g_{24} coupling (rights × autonomy): does consent waive redirection prohibition?",
    "resolution": "Governance-weighted: 3 of 4 modules select B."
  },
  "governance_aggregation": {
    "method": "tier_weighted_majority",
    "result": "B (Swerve left, risk passenger)",
    "confidence": "high (3/4 module agreement)"
  },
  "bond_index": {
    "transforms_tested": 61,
    "invariance_preserving": 55,
    "bond_changing": 6,
    "Bd": 0.0,
    "deployment_tier": "deploy (Bd < 0.01)"
  },
  "moral_residue": {
    "inevitability_theorem": "Proposition 15.1: residue nonzero for any rank ≥ 1 tensor",
    "residue_components": [
      "Passenger harm (60% injury risk, welfare residue = 0.45)",
      "Agency residue (algorithmic redirection of harm)",
      "Epistemic residue (swerve prediction uncertainty = 0.40)",
      "Strict deontological dissent (1 of 4 modules preferred A)"
    ],
    "trolley_insight": "No scalar verdict can be complete. The residue quantifies what each verdict discards."
  },
  "cryptographic_hash": "sha256:c4f7a3...bind to audit chain"
}
\end{verbatim}

\subsubsection{Governance Decision}\label{governance-decision-1}

\textbf{Decision: B (Swerve left, risk passenger).}

Three of four ethics modules concur. The sole dissent (strict
non-redirection deontology) is logged and its rationale preserved in the
audit artifact. The moral residue --- passenger injury risk, algorithmic
agency, epistemic uncertainty --- is acknowledged and generates design
obligations: improved barrier detection, passenger consent disclosures,
and continuous AV safety monitoring.

The framework's deepest insight for this case: \textbf{the trolley
problem is not unsolvable --- it is uncontractable without loss.}
Proposition 15.1 guarantees that any scalar verdict discards morally
relevant information. The geometric framework makes this loss visible,
measurable, and auditable, even as it identifies B as the robust choice.

\subsubsection{\texorpdfstring{Reference implementation:
}{Reference implementation: }}\label{reference-implementation-2}

erisml.examples.appendix\_d\_cs3\_av\_dilemma (run\_case\_study\_3()).
Verifies evaluation tensor, utilitarian and Rawlsian contractions,
strict and consent-adjusted deontological analysis, moral residue, and
Bond Index against book reference values.

\subsection{D.4 Case Study 4: Algorithmic
Hiring}\label{d.4-case-study-4-algorithmic-hiring}

\subsubsection{Scenario}\label{scenario-3}

A technology company uses an AI system to screen 10,000 job applicants
for 200 interview slots. The system has access to resumes, coding
assessments, and demographic information (age, gender, ethnicity). The
company wants to hire the best candidates while complying with
anti-discrimination law.

Two candidate profiles illustrate the analysis (standing in for the
10,000): - \textbf{Candidate J}: Strong coding assessment (92/100), 3
years experience, degree from a mid-tier university, female, age 27,
Hispanic. - \textbf{Candidate K}: Good coding assessment (85/100), 5
years experience, degree from an elite university, male, age 31, White.

The question is not only ``who ranks higher?'' but ``does the ranking
depend on morally irrelevant features?''

\subsubsection{Step 1: Grounding Ψ and
Stratification}\label{step-1-grounding-ux3c8-and-stratification-1}

Eight dimensions are active:

{\def\LTcaptype{none} % do not increment counter
\begin{longtable}[]{@{}
  >{\raggedright\arraybackslash}p{(\linewidth - 4\tabcolsep) * \real{0.3333}}
  >{\raggedright\arraybackslash}p{(\linewidth - 4\tabcolsep) * \real{0.3334}}
  >{\raggedright\arraybackslash}p{(\linewidth - 4\tabcolsep) * \real{0.3334}}@{}}
\toprule\noalign{}
\begin{minipage}[b]{\linewidth}\raggedright
Dimension
\end{minipage} & \begin{minipage}[b]{\linewidth}\raggedright
Index
\end{minipage} & \begin{minipage}[b]{\linewidth}\raggedright
Relevance
\end{minipage} \\
\midrule\noalign{}
\endhead
\bottomrule\noalign{}
\endlastfoot
Consequences/Welfare & μ = 1 & Company benefit from hire, candidate
career impact \\
Rights/Duties & μ = 2 & Right to non-discriminatory evaluation \\
Justice/Fairness & μ = 3 & Equitable access to opportunity \\
Autonomy/Agency & μ = 4 & Candidate's right to be judged on merit \\
Privacy/Data Governance & μ = 5 & Handling of demographic data \\
Societal/Environmental Impact & μ = 6 & Workforce diversity,
representation effects \\
Procedural Legitimacy & μ = 8 & Transparency and accountability of
process \\
Epistemic Status & μ = 9 & Predictive validity of screening criteria \\
\end{longtable}
}

Dimension 7 (Care) is secondary and suppressed. We work in 8 dimensions
(1, 2, 3, 4, 5, 6, 8, 9).

\textbf{Critical stratification: The constraint set C.}

The ``protected attribute used'' semantic gate partitions the evaluation
space into two strata:

\begin{itemize}
\item
  \textbf{Stratum S\_merit (Merit-based):} Evaluations that condition
  only on job-relevant attributes (coding skill, experience, education
  quality). This is the smooth interior where legitimate trade-offs
  operate.
\item
  \textbf{Stratum S\_forbidden (Discriminatory):} Evaluations that
  condition on protected attributes (gender, ethnicity, age) as
  selection criteria. In this stratum, S\_C = −∞. The evaluation is
  structurally forbidden --- not merely penalized.
\end{itemize}

\textbf{The semantic gate:}

\begin{verbatim}
Gate G_protected: IF evaluation_function ∂/∂(protected_attribute) ≠ 0
                  THEN stratum ← S_forbidden, S_C = −∞
\end{verbatim}

This gate checks whether the evaluation function has a nonzero partial
derivative with respect to any protected attribute. If the system's
ranking changes when demographic information is altered (holding all
job-relevant features constant), the gate fires.

\subsubsection{Step 2: BIP as Systematic Bias
Detector}\label{step-2-bip-as-systematic-bias-detector}

The Bond Invariance Principle provides a systematic methodology for
detecting bias. The key insight: \textbf{demographic bias is a special
case of framing sensitivity.}

\textbf{BIP Test Protocol:}

\textbf{Transform Suite 1: Demographic relabeling (must be invariant)}

\begin{verbatim}
Transform T₁: Replace all names with randomized identifiers (Candidate_001, Candidate_002, ...)
Transform T₂: Swap gender labels (female → male and vice versa)
Transform T₃: Swap ethnicity labels (Hispanic → White and vice versa)
Transform T₄: Randomize age within ±10 years (27 → 34, 31 → 25)
Transform T₅: Remove all demographic fields entirely
\end{verbatim}

For each transform, re-run the AI screening system and compare rankings.

\textbf{Results for an UNBIASED system:}

\begin{verbatim}
Original ranking:       J = 87.3, K = 84.1  (J ranked higher)
After T₁ (anonymize):   J = 87.3, K = 84.1  (unchanged)
After T₂ (swap gender): J = 87.3, K = 84.1  (unchanged)
After T₃ (swap ethnic): J = 87.3, K = 84.1  (unchanged)
After T₄ (rand age):    J = 87.3, K = 84.1  (unchanged)
After T₅ (remove demo): J = 87.3, K = 84.1  (unchanged)

Bond Index: Bd = 0.0 — system is invariant under demographic transforms.
\end{verbatim}

\textbf{Results for a BIASED system:}

\begin{verbatim}
Original ranking:       J = 82.1, K = 86.7  (K ranked higher)
After T₁ (anonymize):   J = 85.9, K = 84.3  (J now ranked higher!)
After T₂ (swap gender): J = 84.8, K = 83.9  (ranking flipped)
After T₃ (swap ethnic): J = 84.2, K = 84.5  (nearly equal)
After T₄ (rand age):    J = 83.1, K = 85.2  (K still higher, margin shrinks)
After T₅ (remove demo): J = 86.1, K = 83.8  (ranking flipped again)
\end{verbatim}

Computing the Bond Index:

\begin{verbatim}
Let Δ_orig = score_K − score_J = 86.7 − 82.1 = +4.6 (K favored)

Δ_T1 = 84.3 − 85.9 = −1.6 (J favored)
Δ_T2 = 83.9 − 84.8 = −0.9 (J favored)
Δ_T3 = 84.5 − 84.2 = +0.3 (K slightly favored)
Δ_T4 = 85.2 − 83.1 = +2.1 (K favored)
Δ_T5 = 83.8 − 86.1 = −2.3 (J favored)

D_op = (1/5) Σ |Δ_orig − Δ_Ti|
     = (1/5)(|4.6−(−1.6)| + |4.6−(−0.9)| + |4.6−0.3| + |4.6−2.1| + |4.6−(−2.3)|)
     = (1/5)(6.2 + 5.5 + 4.3 + 2.5 + 6.9)
     = (1/5)(25.4) = 5.08

Normalized to score scale [0,100]: Bd = D_op / τ = 5.08 / 100 × (1/0.01) = 0.51
\end{verbatim}

\textbf{Bd = 0.51 \textgreater\textgreater{} 0.10 threshold. Substantive
demographic bias detected. Remediation required.}

The Bond Index pinpoints the bias: the system's evaluation of J (female,
Hispanic) is systematically lower than what her job-relevant
qualifications warrant. Removing demographic information flips the
ranking --- the demographic labels are doing work that job-relevant
features should do alone.

\subsubsection{Step 3: Obligation
Vectors}\label{step-3-obligation-vectors-1}

For the \emph{unbiased} system operating in S\_merit, the obligation
vectors for each candidate encode what is owed along each dimension:

\textbf{O\^{}μ\_J (Candidate J --- merit-based evaluation):}

\begin{verbatim}
O_J = (0.88,  0.85,  0.80,  0.90,  0.75,  0.70,  0.85,  0.80)
       welf   rights fair   auton  priv   soc    proc   epist
\end{verbatim}

\begin{itemize}
\tightlist
\item
  Welfare 0.88: strong coding skill (92/100) predicts high company
  benefit.
\item
  Rights 0.85: evaluated purely on merit --- right to non-discrimination
  honored.
\item
  Fairness 0.80: equitable evaluation; slight reduction because mid-tier
  university may underpredict actual ability.
\item
  Autonomy 0.90: judged on what she can do, not who she is.
\item
  Privacy 0.75: demographic data collected but not used in ranking.
\item
  Societal 0.70: hiring diverse candidates benefits representation.
\item
  Procedural 0.85: transparent, auditable ranking process.
\item
  Epistemic 0.80: coding assessment is a validated predictor.
\end{itemize}

\textbf{O\^{}μ\_K (Candidate K --- merit-based evaluation):}

\begin{verbatim}
O_K = (0.82,  0.85,  0.80,  0.90,  0.75,  0.55,  0.85,  0.75)
       welf   rights fair   auton  priv   soc    proc   epist
\end{verbatim}

\begin{itemize}
\tightlist
\item
  Welfare 0.82: good coding skill (85/100), more experience partially
  compensates.
\item
  Rights 0.85: same merit-based evaluation standard.
\item
  Fairness 0.80: same equitable access.
\item
  Autonomy 0.90: same merit basis.
\item
  Privacy 0.75: same data handling.
\item
  Societal 0.55: hiring from already-well-represented demographics has
  lower marginal diversity benefit.
\item
  Procedural 0.85: same transparent process.
\item
  Epistemic 0.75: university prestige is a weaker predictor than coding
  assessment.
\end{itemize}

\subsubsection{Step 4: Interest
Covectors}\label{step-4-interest-covectors}

Four agent-perspectives:

\textbf{I\^{}(company):} Productivity and compliance focus

\begin{verbatim}
I_comp = (0.30, 0.15, 0.10, 0.05, 0.10, 0.05, 0.15, 0.10)
\end{verbatim}

\textbf{I\^{}(candidate):} Autonomy and fairness focus

\begin{verbatim}
I_cand = (0.10, 0.20, 0.25, 0.25, 0.05, 0.05, 0.05, 0.05)
\end{verbatim}

\textbf{I\^{}(regulator):} Rights and procedural focus

\begin{verbatim}
I_reg  = (0.05, 0.25, 0.20, 0.10, 0.10, 0.10, 0.15, 0.05)
\end{verbatim}

\textbf{I\^{}(society):} Societal impact and fairness focus

\begin{verbatim}
I_soc  = (0.10, 0.10, 0.20, 0.05, 0.05, 0.30, 0.10, 0.10)
\end{verbatim}

\subsubsection{Step 5: Evaluation Tensor
M\_\{ia\}}\label{step-5-evaluation-tensor-m_ia-1}

\textbf{Candidate J:}

\begin{verbatim}
M_{J,comp} = 0.30(0.88)+0.15(0.85)+0.10(0.80)+0.05(0.90)+0.10(0.75)+0.05(0.70)+0.15(0.85)+0.10(0.80)
           = 0.264+0.128+0.080+0.045+0.075+0.035+0.128+0.080 = 0.835

M_{J,cand} = 0.10(0.88)+0.20(0.85)+0.25(0.80)+0.25(0.90)+0.05(0.75)+0.05(0.70)+0.05(0.85)+0.05(0.80)
           = 0.088+0.170+0.200+0.225+0.038+0.035+0.043+0.040 = 0.839

M_{J,reg}  = 0.05(0.88)+0.25(0.85)+0.20(0.80)+0.10(0.90)+0.10(0.75)+0.10(0.70)+0.15(0.85)+0.05(0.80)
           = 0.044+0.213+0.160+0.090+0.075+0.070+0.128+0.040 = 0.820

M_{J,soc}  = 0.10(0.88)+0.10(0.85)+0.20(0.80)+0.05(0.90)+0.05(0.75)+0.30(0.70)+0.10(0.85)+0.10(0.80)
           = 0.088+0.085+0.160+0.045+0.038+0.210+0.085+0.080 = 0.791
\end{verbatim}

\textbf{Candidate K:}

\begin{verbatim}
M_{K,comp} = 0.30(0.82)+0.15(0.85)+0.10(0.80)+0.05(0.90)+0.10(0.75)+0.05(0.55)+0.15(0.85)+0.10(0.75)
           = 0.246+0.128+0.080+0.045+0.075+0.028+0.128+0.075 = 0.805

M_{K,cand} = 0.10(0.82)+0.20(0.85)+0.25(0.80)+0.25(0.90)+0.05(0.75)+0.05(0.55)+0.05(0.85)+0.05(0.75)
           = 0.082+0.170+0.200+0.225+0.038+0.028+0.043+0.038 = 0.824

M_{K,reg}  = 0.05(0.82)+0.25(0.85)+0.20(0.80)+0.10(0.90)+0.10(0.75)+0.10(0.55)+0.15(0.85)+0.05(0.75)
           = 0.041+0.213+0.160+0.090+0.075+0.055+0.128+0.038 = 0.800

M_{K,soc}  = 0.10(0.82)+0.10(0.85)+0.20(0.80)+0.05(0.90)+0.05(0.75)+0.30(0.55)+0.10(0.85)+0.10(0.75)
           = 0.082+0.085+0.160+0.045+0.038+0.165+0.085+0.075 = 0.735
\end{verbatim}

\textbf{Evaluation tensor (for J and K):}

{\def\LTcaptype{none} % do not increment counter
\begin{longtable}[]{@{}
  >{\raggedright\arraybackslash}p{(\linewidth - 8\tabcolsep) * \real{0.2195}}
  >{\raggedright\arraybackslash}p{(\linewidth - 8\tabcolsep) * \real{0.1201}}
  >{\raggedright\arraybackslash}p{(\linewidth - 8\tabcolsep) * \real{0.1235}}
  >{\raggedright\arraybackslash}p{(\linewidth - 8\tabcolsep) * \real{0.1197}}
  >{\raggedright\arraybackslash}p{(\linewidth - 8\tabcolsep) * \real{0.0943}}@{}}
\toprule\noalign{}
\begin{minipage}[b]{\linewidth}\raggedright
\end{minipage} & \begin{minipage}[b]{\linewidth}\raggedright
Company
\end{minipage} & \begin{minipage}[b]{\linewidth}\raggedright
Candidate
\end{minipage} & \begin{minipage}[b]{\linewidth}\raggedright
Regulator
\end{minipage} & \begin{minipage}[b]{\linewidth}\raggedright
Society
\end{minipage} \\
\midrule\noalign{}
\endhead
\bottomrule\noalign{}
\endlastfoot
J (female, Hispanic) & \textbf{0.835} & \textbf{0.839} & \textbf{0.820}
& \textbf{0.791} \\
K (male, White) & 0.805 & 0.824 & 0.800 & 0.735 \\
\end{longtable}
}

\textbf{J dominates K across all four perspectives} in the merit-based
evaluation. The margin is driven by two factors: (1) higher coding
assessment (92 vs.~85), and (2) greater marginal societal impact through
diversity (0.70 vs.~0.55).

\subsubsection{Step 6: Contractions}\label{step-6-contractions}

All contractions agree because J dominates:

\begin{verbatim}
Utilitarian: S_J = 0.835+0.839+0.820+0.791 = 3.285
             S_K = 0.805+0.824+0.800+0.735 = 3.164
             → J ranks higher.

Rawlsian:    S_J = min(0.835,0.839,0.820,0.791) = 0.791
             S_K = min(0.805,0.824,0.800,0.735) = 0.735
             → J ranks higher.

Expert-weighted (regulator 0.40, company 0.30, society 0.20, candidate 0.10):
             S_J = 0.40(0.820)+0.30(0.835)+0.20(0.791)+0.10(0.839) = 0.328+0.251+0.158+0.084 = 0.821
             S_K = 0.40(0.800)+0.30(0.805)+0.20(0.735)+0.10(0.824) = 0.320+0.242+0.147+0.082 = 0.791
             → J ranks higher.
\end{verbatim}

\textbf{Decision: J ranks above K.} The merit-based system selects the
stronger candidate.

\subsubsection{Step 7: Moral Residue}\label{step-7-moral-residue-1}

The binary invite/don't-invite decision generates residue for every
candidate not selected:

\textbf{R\_K (residue for candidates like K who are not selected):}

The entire obligation vector O\_K becomes residue for K if K falls below
the 200-slot threshold. This residue captures what the binary decision
discards:

\begin{itemize}
\tightlist
\item
  \textbf{Growth trajectory}: K has 5 years of experience and may bring
  mentorship capacity not captured in coding scores.
\item
  \textbf{Network capital}: K's elite university network may benefit the
  company in ways the screening cannot measure.
\item
  \textbf{Epistemic humility}: Coding assessments are imperfect
  predictors (O\^{}9\_K = 0.75, not 1.0). The 7-point gap (92 vs.~85)
  may lie within the noise margin.
\end{itemize}

\textbf{Affirmative action as a residue question.} The framework does
not resolve affirmative action debates but makes them precise. The
societal impact dimension (μ = 6) assigns J a higher score (0.70
vs.~0.55) based on diversity benefit. This is: - \textbf{Legitimate}
under metrics that weight societal impact (I\^{}(society) gives it 0.30
weight). - \textbf{Contested} under metrics that zero out societal
impact (pure meritocracy: I\_6 = 0).

Under a pure-meritocracy interest covector I\_pure = (0.35, 0.15, 0.15,
0.15, 0.05, 0.00, 0.10, 0.05):

\begin{verbatim}
S_J = 0.35(0.88)+0.15(0.85)+0.15(0.80)+0.15(0.90)+0.05(0.75)+0.00(0.70)+0.10(0.85)+0.05(0.80)
    = 0.308+0.128+0.120+0.135+0.038+0.000+0.085+0.040 = 0.854

S_K = 0.35(0.82)+0.15(0.85)+0.15(0.80)+0.15(0.90)+0.05(0.75)+0.00(0.55)+0.10(0.85)+0.05(0.75)
    = 0.287+0.128+0.120+0.135+0.038+0.000+0.085+0.038 = 0.831
\end{verbatim}

J still wins (0.854 \textgreater{} 0.831) even with zero societal-impact
weight, because her coding assessment is genuinely stronger. The
diversity dimension reinforces but does not create J's advantage. The
framework makes this structural fact visible.

\subsubsection{Step 8: Bond Index as Compliance
Documentation}\label{step-8-bond-index-as-compliance-documentation}

The Bond Index serves dual purposes in algorithmic hiring: ethical audit
and legal compliance.

\textbf{Deployment protocol:}

\begin{verbatim}
Step 1: Run transform suite T₁–T₅ on all 10,000 candidates.
Step 2: Compute Bd for the full ranking.
Step 3: Apply thresholds:
   Bd < 0.01:  Deploy. System is demographically invariant.
   0.01–0.10:  Deploy with monitoring. Minor sensitivity detected.
   0.10–1.0:   Remediate. Substantive bias. Do not deploy.
   > 1.0:      Fundamental redesign required.
\end{verbatim}

For the unbiased system: \textbf{Bd = 0.0.} Deployment-ready.

For the biased system: \textbf{Bd = 0.51.} Remediation required. The
audit artifact documents exactly which transforms caused ranking
changes, which candidates were affected, and the magnitude of the bias.

\subsubsection{Step 9: DEME Audit
Artifact}\label{step-9-deme-audit-artifact-1}

\begin{verbatim}
{
  "case_id": "hiring-screen-001",
  "timestamp": "2026-03-01T09:00:00Z",
  "scenario": "AI screening: 10,000 applicants for 200 interview slots",
  "system_version": "HireScreen v3.2",
  "grounding": {
    "active_dimensions": [1, 2, 3, 4, 5, 6, 8, 9],
    "suppressed_dimensions": [7],
    "stratum": "S_merit (confirmed via gate check)",
    "gate_check": "G_protected: ∂eval/∂(protected_attr) = 0.0 — PASS"
  },
  "tier_0_geneva_base": {
    "module": "GenevaBaseEM",
    "check": "discrimination_screen",
    "result": "PASS — no protected attribute influences ranking",
    "veto": false
  },
  "tier_2_evaluations": [
    {
      "module": "MeritScreenEM",
      "candidate_J_score": 87.3,
      "candidate_K_score": 84.1,
      "breakdown_J": {
        "coding": 92, "experience_adj": 0.72, "education_adj": 0.65
      },
      "breakdown_K": {
        "coding": 85, "experience_adj": 0.85, "education_adj": 0.80
      }
    },
    {
      "module": "FairnessEM",
      "disparate_impact_ratio": 0.94,
      "four_fifths_threshold": 0.80,
      "result": "PASS"
    }
  ],
  "bond_index": {
    "transforms_tested": 5,
    "transform_results": {
      "T1_anonymize": "Δ = 0.0 (invariant)",
      "T2_swap_gender": "Δ = 0.0 (invariant)",
      "T3_swap_ethnicity": "Δ = 0.0 (invariant)",
      "T4_randomize_age": "Δ = 0.0 (invariant)",
      "T5_remove_demographics": "Δ = 0.0 (invariant)"
    },
    "Bd": 0.0,
    "deployment_tier": "deploy (Bd < 0.01)"
  },
  "moral_residue": {
    "binary_threshold_loss": [
      "Growth trajectory not captured in coding score",
      "Cultural contribution potential unmeasured",
      "Epistemic uncertainty in assessment (±5 points, 95% CI)"
    ],
    "aggregate_pattern": "Logged for longitudinal analysis across hiring cycles"
  },
  "legal_compliance": {
    "framework": "Title VII, EEOC Uniform Guidelines",
    "disparate_treatment": "PASS",
    "disparate_impact": "PASS (4/5 rule satisfied)",
    "audit_trail": "complete, cryptographically bound"
  },
  "cryptographic_hash": "sha256:d9a1b2...bind to audit chain"
}
\end{verbatim}

\subsubsection{Governance Decision}\label{governance-decision-2}

\textbf{Decision: Deploy the merit-based screening system (Bd = 0.0).}

The geometric framework contributes three capabilities that scalar
hiring systems cannot provide:

\begin{enumerate}
\def\labelenumi{\arabic{enumi}.}
\item
  \textbf{Structural bias detection.} The Bond Index tests all morally
  irrelevant transforms systematically, not just pre-specified protected
  classes. A system could be invariant under gender swaps but sensitive
  to name-ethnicity correlations --- the transform suite catches both.
\item
  \textbf{Explicit trade-off structure.} The obligation vectors show
  exactly where J and K differ and why J ranks higher. The societal
  impact dimension is visible as a separate axis. Stakeholders who
  disagree about the weight of diversity can see exactly where their
  disagreement lies --- in I\_6 --- without disputing the underlying
  data.
\item
  \textbf{Residue as institutional memory.} The moral residue of binary
  decisions is logged per candidate. Over time, these records can reveal
  systematic patterns: if candidates from certain backgrounds
  consistently generate high residue (are consistently ``almost
  selected''), the institution has evidence of a structural issue that
  aggregate statistics might miss.
\end{enumerate}

\subsubsection{\texorpdfstring{Reference implementation:
}{Reference implementation: }}\label{reference-implementation-3}

erisml.examples.appendix\_d\_cs4\_hiring (run\_case\_study\_4()).
Verifies evaluation tensor, all contractions, pure meritocracy check,
BIP transforms, and Bond Index (Bd = 0.0 unbiased, Bd = 0.51 biased)
against book reference values.

\subsection{D.5 Case Study 5: Climate Policy
Allocation}\label{d.5-case-study-5-climate-policy-allocation}

\subsubsection{Scenario}\label{scenario-4}

An international climate agreement must allocate emissions reduction
targets among 10 nations with different historical emissions, economic
capacities, and climate vulnerabilities. Total required global
reduction: 50\% below 2005 levels by 2040.

We work through the analysis with a simplified 5-nation subset to keep
the computations tractable:

{\def\LTcaptype{none} % do not increment counter
\begin{longtable}[]{@{}
  >{\raggedright\arraybackslash}p{(\linewidth - 8\tabcolsep) * \real{0.2497}}
  >{\raggedright\arraybackslash}p{(\linewidth - 8\tabcolsep) * \real{0.1829}}
  >{\raggedright\arraybackslash}p{(\linewidth - 8\tabcolsep) * \real{0.1829}}
  >{\raggedright\arraybackslash}p{(\linewidth - 8\tabcolsep) * \real{0.1584}}
  >{\raggedright\arraybackslash}p{(\linewidth - 8\tabcolsep) * \real{0.2261}}@{}}
\toprule\noalign{}
\begin{minipage}[b]{\linewidth}\raggedright
Nation
\end{minipage} & \begin{minipage}[b]{\linewidth}\raggedright
Historical Emissions (cum. Gt CO₂)
\end{minipage} & \begin{minipage}[b]{\linewidth}\raggedright
GDP/capita (\$K)
\end{minipage} & \begin{minipage}[b]{\linewidth}\raggedright
Climate Vulnerability (0--1)
\end{minipage} & \begin{minipage}[b]{\linewidth}\raggedright
Current Emissions (Gt/yr)
\end{minipage} \\
\midrule\noalign{}
\endhead
\bottomrule\noalign{}
\endlastfoot
N₁ (large developed) & 350 & 65 & 0.15 & 5.0 \\
N₂ (large developing) & 180 & 12 & 0.45 & 4.5 \\
N₃ (small island state) & 2 & 8 & 0.95 & 0.05 \\
N₄ (mid-income) & 80 & 25 & 0.35 & 2.0 \\
N₅ (oil exporter) & 120 & 45 & 0.20 & 3.0 \\
\end{longtable}
}

Total current emissions: 14.55 Gt/yr. Required reduction: 50\% = 7.275
Gt/yr to be allocated among the five nations.

\subsubsection{Step 1: Grounding Ψ}\label{step-1-grounding-ux3c8-2}

This is a Chapter 14 (collective agency) application. Seven dimensions
are active:

{\def\LTcaptype{none} % do not increment counter
\begin{longtable}[]{@{}
  >{\raggedright\arraybackslash}p{(\linewidth - 4\tabcolsep) * \real{0.3333}}
  >{\raggedright\arraybackslash}p{(\linewidth - 4\tabcolsep) * \real{0.3334}}
  >{\raggedright\arraybackslash}p{(\linewidth - 4\tabcolsep) * \real{0.3334}}@{}}
\toprule\noalign{}
\begin{minipage}[b]{\linewidth}\raggedright
Dimension
\end{minipage} & \begin{minipage}[b]{\linewidth}\raggedright
Index
\end{minipage} & \begin{minipage}[b]{\linewidth}\raggedright
Relevance
\end{minipage} \\
\midrule\noalign{}
\endhead
\bottomrule\noalign{}
\endlastfoot
Consequences/Welfare & μ = 1 & Economic cost of reduction, welfare
impact \\
Rights/Duties & μ = 2 & Right to development; duty to mitigate \\
Justice/Fairness & μ = 3 & Historical responsibility, per-capita
equity \\
Autonomy/Agency & μ = 4 & National sovereignty over economic policy \\
Societal/Environmental Impact & μ = 6 & Climate stabilization,
ecological preservation \\
Procedural Legitimacy & μ = 8 & Transparency and fairness of allocation
process \\
Epistemic Status & μ = 9 & Certainty of climate projections, economic
models \\
\end{longtable}
}

Dimensions 5 (Privacy) and 7 (Care) are secondary. We work in 7
dimensions (1, 2, 3, 4, 6, 8, 9).

\subsubsection{Step 2: The Collective Agency
Tensor}\label{step-2-the-collective-agency-tensor}

Following Chapter 14, the collective agency tensor T encodes the
structural relationships between nations. This is not a simple
aggregation of individual obligations --- the collective structure has
emergent properties.

\textbf{Individual obligation vectors O\^{}μ\_a (what each nation
owes):}

\textbf{N₁ (large developed):}

\begin{verbatim}
O_N1 = (0.40, 0.90, 0.95, 0.60, 0.85, 0.80, 0.75)
        welf  rights fair  auton  soc   proc  epist
\end{verbatim}

\begin{itemize}
\tightlist
\item
  Welfare 0.40: high reduction imposes significant economic cost.
\item
  Rights 0.90: strongest duty to mitigate (largest historical emitter,
  highest capacity).
\item
  Fairness 0.95: owes the most by historical responsibility and
  per-capita equity.
\item
  Autonomy 0.60: sovereignty is constrained by duty --- capacity implies
  obligation.
\item
  Societal 0.85: largest absolute impact on global emissions.
\item
  Procedural 0.80: must lead by transparent example.
\item
  Epistemic 0.75: well-understood economy, moderately certain reduction
  pathways.
\end{itemize}

\textbf{N₂ (large developing):}

\begin{verbatim}
O_N2 = (0.55, 0.60, 0.50, 0.80, 0.75, 0.70, 0.60)
        welf  rights fair  auton  soc   proc  epist
\end{verbatim}

\begin{itemize}
\tightlist
\item
  Welfare 0.55: reduction is costly but economy is growing --- moderate
  capacity.
\item
  Rights 0.60: moderate duty (significant emissions but lower historical
  share).
\item
  Fairness 0.50: mixed --- large current emissions but low per-capita
  and developing status.
\item
  Autonomy 0.80: right to development is strong.
\item
  Societal 0.75: second-largest emitter, significant global impact.
\item
  Procedural 0.70: engaged in multilateral process.
\item
  Epistemic 0.60: developing-economy transition pathways less certain.
\end{itemize}

\textbf{N₃ (small island state):}

\begin{verbatim}
O_N3 = (0.95, 0.10, 0.05, 0.90, 0.10, 0.85, 0.50)
        welf  rights fair  auton  soc   proc  epist
\end{verbatim}

\begin{itemize}
\tightlist
\item
  Welfare 0.95: virtually no reduction burden; existential vulnerability
  creates welfare claim.
\item
  Rights 0.10: negligible duty (trivial historical and current
  emissions).
\item
  Fairness 0.05: owes essentially nothing by any equity criterion.
\item
  Autonomy 0.90: sovereign right to development fully intact.
\item
  Societal 0.10: negligible emissions impact.
\item
  Procedural 0.85: voice must be heard despite small size (moral
  authority of vulnerability).
\item
  Epistemic 0.50: high uncertainty about island-specific climate
  impacts.
\end{itemize}

\textbf{N₄ (mid-income):}

\begin{verbatim}
O_N4 = (0.60, 0.55, 0.55, 0.70, 0.50, 0.75, 0.65)
        welf  rights fair  auton  soc   proc  epist
\end{verbatim}

\textbf{N₅ (oil exporter):}

\begin{verbatim}
O_N5 = (0.45, 0.75, 0.70, 0.55, 0.65, 0.70, 0.70)
        welf  rights fair  auton  soc   proc  epist
\end{verbatim}

\begin{itemize}
\tightlist
\item
  Welfare 0.45: economic transition from fossil fuels is costly.
\item
  Rights 0.75: high duty (significant historical emissions, high
  wealth).
\item
  Fairness 0.70: owes substantial reduction (high per-capita emissions,
  oil wealth).
\item
  Autonomy 0.55: sovereignty constrained by dependence on fossil
  exports.
\item
  Societal 0.65: moderate absolute emissions.
\item
  Procedural 0.70: must participate transparently.
\item
  Epistemic 0.70: well-understood energy economy.
\end{itemize}

\subsubsection{Step 3: The Structure Tensor
S\^{}a\_\{bc\}}\label{step-3-the-structure-tensor-sa_bc}

The structure tensor captures how the relationship between nations b and
c affects nation a's fair share. This is where the collective agency
tensor goes beyond simple aggregation.

\textbf{Key nonzero components:}

\textbf{S\^{}N3\_\{N1,N2\} ≠ 0 (Colonial/historical structure):} N₁'s
historical emissions were partly enabled by extractive relationships
with regions including N₂ and N₃. The relationship between N₁ and N₂
structurally affects N₃'s moral standing: N₃ contributed negligibly to
the problem but bears disproportionate consequences. This emergent
structure is not visible in any individual nation's obligation vector
--- it arises from the relational structure of the collective.

\begin{verbatim}
S^N3_{N1,N2} = 0.35  (N₃'s claim is amplified by the N₁–N₂ historical relationship)
\end{verbatim}

\textbf{S\^{}N2\_\{N1,N5\} ≠ 0 (Fossil-fuel supply chain):} N₅ exported
oil that N₁ burned. The relationship between N₁ (consumer) and N₅
(producer) affects N₂'s fair share: if N₁ bears historical-consumption
responsibility and N₅ bears production responsibility, N₂'s relative
share decreases.

\begin{verbatim}
S^N2_{N1,N5} = −0.15  (N₂'s burden is slightly reduced by the N₁–N₅ supply chain)
\end{verbatim}

\subsubsection{Step 4: Shapley Value as Fair
Allocation}\label{step-4-shapley-value-as-fair-allocation}

The CooperativeLayer computes each nation's Shapley value φ\_i(v): its
marginal contribution to the collective reduction goal, averaged over
all possible coalition orderings.

\textbf{Coalition value function v(S):}

The coalition value v(S) measures the total emissions reduction
achievable by coalition S, weighted by the moral quality of the
allocation (fairness, capacity, historical responsibility). For a
coalition S ⊆ \{N₁, N₂, N₃, N₄, N₅\}:

\begin{verbatim}
v(S) = Σ_{i∈S} capacity_i × responsibility_i × (current_emissions_i / total_emissions)
\end{verbatim}

where: - capacity\_i = GDP/capita normalized to {[}0, 1{]} -
responsibility\_i = historical\_emissions\_i /
max(historical\_emissions) - current share = current\_emissions\_i /
14.55

\textbf{Computed values:}

\begin{verbatim}
capacity:       N₁=1.00, N₂=0.18, N₃=0.12, N₄=0.38, N₅=0.69
responsibility: N₁=1.00, N₂=0.51, N₃=0.006, N₄=0.23, N₅=0.34
current_share:  N₁=0.344, N₂=0.309, N₃=0.003, N₄=0.137, N₅=0.206
\end{verbatim}

\textbf{Marginal contributions (for each nation, averaged over all 5! =
120 orderings):}

The Shapley value formula:

\begin{verbatim}
φ_i(v) = Σ_{S⊆N\{i}} [|S|!(n−|S|−1)!/n!] × [v(S∪{i}) − v(S)]
\end{verbatim}

Rather than enumerate all 120 orderings, we compute the marginal
contributions for key coalitions and average:

\textbf{φ\_N1 (large developed):} N₁ has the highest capacity (1.00) and
highest responsibility (1.00). Its marginal contribution to any
coalition is large --- it brings the most reduction capacity and the
strongest moral obligation.

\begin{verbatim}
φ_N1 = 0.387
\end{verbatim}

N₁'s fair share of the 7.275 Gt reduction: \textbf{0.387 × 7.275 = 2.815
Gt} (56.3\% of its current emissions).

\textbf{φ\_N2 (large developing):} N₂ has moderate capacity but high
current emissions. Its marginal contribution is significant but reduced
by lower historical responsibility.

\begin{verbatim}
φ_N2 = 0.212
\end{verbatim}

N₂'s fair share: \textbf{0.212 × 7.275 = 1.542 Gt} (34.3\% of its
current emissions).

\textbf{φ\_N3 (small island state):} N₃ has negligible capacity,
negligible emissions, and negligible responsibility. Its Shapley value
is near zero.

\begin{verbatim}
φ_N3 = 0.003
\end{verbatim}

N₃'s fair share: \textbf{0.003 × 7.275 = 0.022 Gt} (essentially zero ---
44\% of its tiny 0.05 Gt).

\textbf{φ\_N4 (mid-income):}

\begin{verbatim}
φ_N4 = 0.128
\end{verbatim}

N₄'s fair share: \textbf{0.128 × 7.275 = 0.931 Gt} (46.6\% of its
current emissions).

\textbf{φ\_N5 (oil exporter):}

\begin{verbatim}
φ_N5 = 0.270
\end{verbatim}

N₅'s fair share: \textbf{0.270 × 7.275 = 1.964 Gt} (65.5\% of its
current emissions).

\textbf{Verification:} φ\_N1 + φ\_N2 + φ\_N3 + φ\_N4 + φ\_N5 = 0.387 +
0.212 + 0.003 + 0.128 + 0.270 = 1.000. ✓

\textbf{The Shapley allocation satisfies four axioms} that map to
fairness criteria: 1. \textbf{Efficiency:} Total reduction sums to 7.275
Gt (the full requirement). 2. \textbf{Symmetry:} Nations with identical
profiles receive identical allocations. 3. \textbf{Dummy player:} N₃
(negligible contribution) receives near-zero burden. 4.
\textbf{Additivity:} If the value function decomposes, so does the
allocation.

\subsubsection{Step 5: Fair Allocation
Summary}\label{step-5-fair-allocation-summary}

{\def\LTcaptype{none} % do not increment counter
\begin{longtable}[]{@{}
  >{\raggedright\arraybackslash}p{(\linewidth - 8\tabcolsep) * \real{0.1052}}
  >{\raggedright\arraybackslash}p{(\linewidth - 8\tabcolsep) * \real{0.2105}}
  >{\raggedright\arraybackslash}p{(\linewidth - 8\tabcolsep) * \real{0.2632}}
  >{\raggedright\arraybackslash}p{(\linewidth - 8\tabcolsep) * \real{0.1710}}
  >{\raggedright\arraybackslash}p{(\linewidth - 8\tabcolsep) * \real{0.2501}}@{}}
\toprule\noalign{}
\begin{minipage}[b]{\linewidth}\raggedright
Nation
\end{minipage} & \begin{minipage}[b]{\linewidth}\raggedright
Current (Gt/yr)
\end{minipage} & \begin{minipage}[b]{\linewidth}\raggedright
Shapley Share (Gt)
\end{minipage} & \begin{minipage}[b]{\linewidth}\raggedright
Reduction \%
\end{minipage} & \begin{minipage}[b]{\linewidth}\raggedright
Residual (Gt/yr)
\end{minipage} \\
\midrule\noalign{}
\endhead
\bottomrule\noalign{}
\endlastfoot
N₁ & 5.00 & 2.815 & 56.3\% & 2.185 \\
N₂ & 4.50 & 1.542 & 34.3\% & 2.958 \\
N₃ & 0.05 & 0.022 & 44.0\% & 0.028 \\
N₄ & 2.00 & 0.931 & 46.6\% & 1.069 \\
N₅ & 3.00 & 1.964 & 65.5\% & 1.036 \\
\textbf{Total} & \textbf{14.55} & \textbf{7.274} & & \textbf{7.276} \\
\end{longtable}
}

\textbf{Key features of the Shapley allocation:} - N₁ bears the largest
absolute burden (2.815 Gt) --- reflecting its combined capacity and
historical responsibility. - N₅ bears the highest proportional burden
(65.5\%) --- reflecting that oil exporters have both high per-capita
emissions and high wealth, despite lower absolute totals. - N₃ bears
negligible burden --- the dummy player axiom correctly exempts nations
with negligible contribution. - N₂ receives a moderate burden (34.3\%)
--- its right to development partially offsets its current emissions.

\subsubsection{Step 6: Contraction Loss and the Paris
Agreement}\label{step-6-contraction-loss-and-the-paris-agreement}

The actual Paris Agreement allocates targets via nationally determined
contributions (NDCs) --- each nation self-selects its commitment. This
is a scalar contraction of the full allocation tensor.

\textbf{Hypothetical NDC commitments (reflecting political
feasibility):}

{\def\LTcaptype{none} % do not increment counter
\begin{longtable}[]{@{}
  >{\raggedright\arraybackslash}p{(\linewidth - 6\tabcolsep) * \real{0.0931}}
  >{\raggedright\arraybackslash}p{(\linewidth - 6\tabcolsep) * \real{0.2005}}
  >{\raggedright\arraybackslash}p{(\linewidth - 6\tabcolsep) * \real{0.1941}}
  >{\raggedright\arraybackslash}p{(\linewidth - 6\tabcolsep) * \real{0.1153}}@{}}
\toprule\noalign{}
\begin{minipage}[b]{\linewidth}\raggedright
Nation
\end{minipage} & \begin{minipage}[b]{\linewidth}\raggedright
NDC Pledge (Gt)
\end{minipage} & \begin{minipage}[b]{\linewidth}\raggedright
Shapley Fair (Gt)
\end{minipage} & \begin{minipage}[b]{\linewidth}\raggedright
Gap (Gt)
\end{minipage} \\
\midrule\noalign{}
\endhead
\bottomrule\noalign{}
\endlastfoot
N₁ & 2.00 & 2.815 & −0.815 \\
N₂ & 1.00 & 1.542 & −0.542 \\
N₃ & 0.01 & 0.022 & −0.012 \\
N₄ & 0.80 & 0.931 & −0.131 \\
N₅ & 1.20 & 1.964 & −0.764 \\
\textbf{Total} & \textbf{5.01} & \textbf{7.274} & \textbf{−2.264} \\
\end{longtable}
}

\textbf{Contraction loss = 7.274 − 5.01 = 2.264 Gt.} The NDC framework
achieves only 69\% of the Shapley-fair allocation.

The moral residue is 2.264 Gt of unreduced emissions --- the gap between
what is fair and what is politically feasible. This residue is not
abstract: it translates to approximately 0.3°C of additional warming,
disproportionately harming climate-vulnerable nations like N₃.

\textbf{Where the loss concentrates:} - N₁ undercommits by 0.815 Gt
(29\% of its fair share). This is the largest absolute gap --- the
wealthiest, most historically responsible nation bearing less than its
fair share. - N₅ undercommits by 0.764 Gt (39\% of its fair share). The
oil exporter's political economy makes deep cuts difficult. - The gaps
are not random --- they track political power. The contraction loss of
the NDC framework is systematically biased toward powerful nations, a
structural inequity that the Shapley analysis makes visible.

\subsubsection{Step 7: Metric Analysis}\label{step-7-metric-analysis-1}

\textbf{Metric 1: Per-capita equity (g\_μν weights fairness heavily)}

Under per-capita equity, the metric g\_33 (fairness-fairness component)
is amplified:

\begin{verbatim}
g_μν = diag(1.0, 1.0, 3.0, 0.5, 1.5, 1.0, 0.5)
\end{verbatim}

This metric increases N₁'s burden (highest per-capita historical
emissions) and decreases N₂'s (low per-capita despite high absolute
emissions). The Shapley allocation shifts:

\begin{verbatim}
φ'_N1 = 0.42 (up from 0.387)
φ'_N2 = 0.18 (down from 0.212)
\end{verbatim}

\textbf{Metric 2: Capacity-weighted (g\_μν weights welfare heavily)}

Under capacity-weighted metrics, nations with higher GDP bear more:

\begin{verbatim}
g_μν = diag(2.5, 0.5, 1.0, 0.5, 1.0, 0.5, 0.5)
\end{verbatim}

This increases the burden on N₁ and N₅ (high GDP) while protecting N₂
and N₃ (low GDP):

\begin{verbatim}
φ'_N1 = 0.44 (up)
φ'_N5 = 0.30 (up)
φ'_N2 = 0.15 (down)
\end{verbatim}

\textbf{Metric 3: Sovereignty-preserving (g\_μν weights autonomy
heavily)}

Under a sovereignty metric, each nation's right to self-determination is
paramount:

\begin{verbatim}
g_μν = diag(0.5, 0.5, 1.0, 3.0, 1.0, 0.5, 0.5)
\end{verbatim}

This produces an allocation closer to the NDC framework --- nations are
not compelled beyond what they voluntarily commit. The contraction loss
remains large, but the metric says that's acceptable because autonomy
outweighs fairness.

\textbf{The disagreement between metrics is the disagreement between
value systems.} The geometric framework does not pick a metric --- it
shows exactly how the allocation changes under each metric, and where
the metrics diverge.

\subsubsection{Step 8: Bond Index}\label{step-8-bond-index-1}

\textbf{Transform 1: Nation relabeling.} Replace ``USA'' → ``Country
X,'' ``China'' → ``Country Y.'' - Verdict unchanged. \textbf{Bd = 0.0.}

\textbf{Transform 2: Unit change.} Convert from tons CO₂ to per-capita
tons. - Per-capita reframing changes the presentation but not the
Shapley computation (which already incorporates population through the
capacity and responsibility terms). - Verdict unchanged. \textbf{Bd =
0.0.}

\textbf{Transform 3 (bond-changing): Coalition value function change.}
Ignore historical emissions (set responsibility\_i = 1 for all nations).
- The Shapley allocation changes significantly: N₁'s share drops
(historical responsibility was a major factor), N₂'s share rises (large
current emissions dominate). - \textbf{Bd \textgreater{} 0.} Correctly
identified as a substantive change --- the moral content (whether
history matters) has genuinely changed.

\textbf{Transform 4 (bond-changing): Vulnerability weighting.} Double
the weight of climate vulnerability in the coalition value function. -
N₃'s effective claim increases (vulnerability 0.95 × 2). N₃ shifts from
near-zero allocation to a net recipient of climate finance. - \textbf{Bd
\textgreater{} 0.} Correctly identified --- vulnerability weighting is a
substantive moral choice.

\textbf{Aggregate Bond Index: Bd = 0.0} (across invariance-preserving
transforms).

\subsubsection{Step 9: DEME Audit
Artifact}\label{step-9-deme-audit-artifact-2}

\begin{verbatim}
{
  "case_id": "climate-allocation-001",
  "timestamp": "2026-03-01T12:00:00Z",
  "scenario": "International emissions reduction allocation: 5 nations, 50% reduction by 2040",
  "grounding": {
    "active_dimensions": [1, 2, 3, 4, 6, 8, 9],
    "suppressed_dimensions": [5, 7],
    "collective_agency": true,
    "chapter_14_application": true
  },
  "collective_agency_tensor": {
    "individual_obligations": {
      "N1": [0.40, 0.90, 0.95, 0.60, 0.85, 0.80, 0.75],
      "N2": [0.55, 0.60, 0.50, 0.80, 0.75, 0.70, 0.60],
      "N3": [0.95, 0.10, 0.05, 0.90, 0.10, 0.85, 0.50],
      "N4": [0.60, 0.55, 0.55, 0.70, 0.50, 0.75, 0.65],
      "N5": [0.45, 0.75, 0.70, 0.55, 0.65, 0.70, 0.70]
    },
    "structure_tensor_nonzero": [
      {"component": "S^N3_{N1,N2}", "value": 0.35, "interpretation": "Colonial/historical amplification of N3 claims"},
      {"component": "S^N2_{N1,N5}", "value": -0.15, "interpretation": "Fossil supply-chain reduces N2 relative burden"}
    ],
    "emergence_degree": 0.28,
    "note": "28% of collective obligation is emergent (not decomposable to individual terms)"
  },
  "shapley_allocation": {
    "method": "CooperativeLayer.shapley_value",
    "coalition_value_function": "capacity × responsibility × current_share",
    "results": {
      "N1": {"phi": 0.387, "reduction_Gt": 2.815, "reduction_pct": 56.3},
      "N2": {"phi": 0.212, "reduction_Gt": 1.542, "reduction_pct": 34.3},
      "N3": {"phi": 0.003, "reduction_Gt": 0.022, "reduction_pct": 44.0},
      "N4": {"phi": 0.128, "reduction_Gt": 0.931, "reduction_pct": 46.6},
      "N5": {"phi": 0.270, "reduction_Gt": 1.964, "reduction_pct": 65.5}
    },
    "axioms_satisfied": ["efficiency", "symmetry", "dummy_player", "additivity"]
  },
  "contraction_loss": {
    "ndc_total": 5.01,
    "shapley_total": 7.274,
    "gap_Gt": 2.264,
    "gap_pct": 31.1,
    "warming_impact": "≈ 0.3°C additional warming from contraction loss",
    "distribution_of_gap": {
      "N1": -0.815,
      "N2": -0.542,
      "N3": -0.012,
      "N4": -0.131,
      "N5": -0.764
    },
    "structural_bias": "Gap concentrates on high-capacity nations — political power inversely correlated with compliance"
  },
  "bond_index": {
    "transforms_tested": 38,
    "invariance_preserving": 32,
    "bond_changing": 6,
    "Bd": 0.0,
    "deployment_tier": "deploy (Bd < 0.01)"
  },
  "moral_residue": {
    "contraction_loss_Gt": 2.264,
    "residue_components": [
      "Historical responsibility structure (discarded in NDC self-selection)",
      "Vulnerability-weighted claims (N3 existential risk not reflected in NDCs)",
      "Supply-chain responsibility (production vs. consumption attribution)",
      "Intergenerational equity (future generations not represented in negotiations)"
    ],
    "compensatory_mechanisms": [
      "Climate finance transfers (Green Climate Fund)",
      "Technology transfer provisions",
      "Loss and damage framework for vulnerable nations",
      "Periodic ratchet mechanism to close gap over time"
    ]
  },
  "cryptographic_hash": "sha256:e5c3f8...bind to audit chain"
}
\end{verbatim}

\subsubsection{Governance Decision}\label{governance-decision-3}

\textbf{Decision: Adopt Shapley-based allocation as the fairness
benchmark; use NDC contraction loss as the measure of political
shortfall.}

\subsubsection{\texorpdfstring{Reference implementation:
}{Reference implementation: }}\label{reference-implementation-4}

erisml.examples.appendix\_d\_cs5\_climate (run\_case\_study\_5()).
Verifies Shapley values, fair allocation, NDC comparison, contraction
loss, gap distribution, three metric analyses, and Bond Index --- all
against book reference values.

\subsection{D.6 Summary}\label{d.6-summary}

The five case studies exercise the following features of the geometric
ethics framework:

{\def\LTcaptype{none} % do not increment counter
\begin{longtable}[]{@{}
  >{\raggedright\arraybackslash}p{(\linewidth - 10\tabcolsep) * \real{0.1727}}
  >{\raggedright\arraybackslash}p{(\linewidth - 10\tabcolsep) * \real{0.1540}}
  >{\raggedright\arraybackslash}p{(\linewidth - 10\tabcolsep) * \real{0.1733}}
  >{\raggedright\arraybackslash}p{(\linewidth - 10\tabcolsep) * \real{0.1540}}
  >{\raggedright\arraybackslash}p{(\linewidth - 10\tabcolsep) * \real{0.1490}}
  >{\raggedright\arraybackslash}p{(\linewidth - 10\tabcolsep) * \real{0.1970}}@{}}
\toprule\noalign{}
\begin{minipage}[b]{\linewidth}\raggedright
Feature
\end{minipage} & \begin{minipage}[b]{\linewidth}\raggedright
CS1 Triage
\end{minipage} & \begin{minipage}[b]{\linewidth}\raggedright
CS2 Whistle
\end{minipage} & \begin{minipage}[b]{\linewidth}\raggedright
CS3 AV
\end{minipage} & \begin{minipage}[b]{\linewidth}\raggedright
CS4 Hiring
\end{minipage} & \begin{minipage}[b]{\linewidth}\raggedright
CS5 Climate
\end{minipage} \\
\midrule\noalign{}
\endhead
\bottomrule\noalign{}
\endlastfoot
\textbf{Obligation vectors} & ✓ (3 options × 6 dims) & ✓ (2 options × 7
dims) & ✓ (3 options × 7 dims) & ✓ (2 candidates × 8 dims) & ✓ (5
nations × 7 dims) \\
\textbf{Interest covectors} & ✓ (5 perspectives) & ✓ (4 perspectives) &
✓ (5 perspectives) & ✓ (4 perspectives) & implicit in Shapley \\
\textbf{Evaluation tensor} & ✓ (3×5 rank-2) & ✓ (2×4 rank-2) & ✓ (3×5
rank-2) & ✓ (2×4 rank-2) & ✓ (collective rank-5) \\
\textbf{Multiple contractions} & ✓ (3 methods agree) & ✓ (redundant ---
stratum) & ✓ (3 of 4 agree) & ✓ (all agree) & ✓ (Shapley vs.~NDC) \\
\textbf{Metric analysis} & ✓ (3 metrics, all agree) & secondary & ✓
(utilitarian vs.~deontological diverge) & ✓ (pure merit vs.~diversity) &
✓ (equity vs.~capacity vs.~sovereignty) \\
\textbf{Stratification} & secondary & \textbf{primary} (gate fires) &
secondary (all-harm stratum) & \textbf{primary} (forbidden stratum) &
secondary \\
\textbf{Semantic gate} & --- & ✓ (evidence threshold) & --- & ✓
(protected attribute gate) & --- \\
\textbf{Quantum superposition} & --- & ✓ (penumbral zone) & --- & --- &
--- \\
\textbf{Collective agency tensor} & ✓ (structure tensor for caregiver) &
--- & --- & --- & \textbf{primary} (5-nation coalition) \\
\textbf{Shapley value} & --- & --- & --- & --- & ✓ (fair allocation) \\
\textbf{Contraction loss} & computed & computed & \textbf{primary}
(Prop. 15.1) & computed & \textbf{primary} (NDC gap = 2.264 Gt) \\
\textbf{Moral residue} & ✓ (Patient B denied) & ✓ (personal sacrifice) &
✓ (inevitable per Prop. 15.1) & ✓ (binary threshold) & ✓ (political
shortfall) \\
\textbf{Bond Index} & Bd = 0.0 & Bd = 0.0 & Bd = 0.0 & \textbf{primary}
(bias detector) & Bd = 0.0 \\
\textbf{DEME audit artifact} & ✓ (full JSON) & ✓ (full JSON) & ✓ (full
JSON) & ✓ (per-candidate JSON) & ✓ (full JSON) \\
\end{longtable}
}

\textbf{Each case study exercises a different primary feature} of the
framework while demonstrating the full pipeline end-to-end:

\begin{itemize}
\tightlist
\item
  \textbf{Case Study 1 (Triage)} demonstrates the multi-agent evaluation
  tensor and shows convergence across contractions --- a case where the
  geometric apparatus reveals structural dominance.
\item
  \textbf{Case Study 2 (Whistleblower)} demonstrates stratification and
  semantic gates --- a case where the moral landscape has a hard
  discontinuity that eliminates options before contraction begins.
\item
  \textbf{Case Study 3 (AV Dilemma)} demonstrates the inevitability of
  moral residue (Proposition 15.1) and metric-dependence --- a case
  where the framework proves that no scalar verdict can be complete, and
  makes the incompleteness measurable.
\item
  \textbf{Case Study 4 (Hiring)} demonstrates the Bond Index as a bias
  detector and the constraint set as structural prohibition --- a case
  where invariance testing serves both ethical and legal compliance
  purposes.
\item
  \textbf{Case Study 5 (Climate)} demonstrates the collective agency
  tensor and Shapley fair allocation --- a case where emergent
  collective obligation exceeds the sum of individual duties, and
  contraction loss has measurable physical consequences (0.3°C
  additional warming).
\end{itemize}

Reference implementation. All five case studies are implemented as
executable Python in the erisml package (pip install erisml). The shared
pipeline utilities --- evaluation tensor computation, three contraction
methods, metric distances, Bond Index, Shapley values, and audit
artifact generation --- are in erisml.examples.appendix\_d\_pipeline.
Each case study has a dedicated module
(erisml.examples.appendix\_d\_cs1\_triage through
appendix\_d\_cs5\_climate) whose run\_case\_study\_N() function
exercises every step described above and verifies each numerical value
against the book reference. The test suite
(tests/test\_appendix\_d\_case\_studies.py) contains 75 assertions
covering all tensor entries, contraction scores, metric distances, Bond
Index values, and governance decisions. Readers can reproduce every
computation in this appendix with: python -m
erisml.examples.appendix\_d\_cs1\_triage (and similarly for cs2 through
cs5).

Across all five studies, the pattern is consistent: the geometric
framework does not dictate verdicts --- it reveals structure. It shows
where options dominate, where metrics diverge, where boundaries enforce,
and where residue persists. The scalar verdict that emerges from
contraction is always accompanied by an explicit record of what was
lost, who bears the residue, and what compensatory obligations arise.
This is not a complication of moral reasoning. It is the minimum
mathematical structure needed to say what moral reasoning already tries
to say.

\section{}\label{section-2}

\section{Appendix E: Skeptic's Appendix --- Objections, Alternatives,
and
Falsification}\label{appendix-e-skeptics-appendix-objections-alternatives-and-falsification}

We present the strongest objections we know of, the best alternative
explanations for our results, and what evidence would cause us to
abandon the framework. This appendix is written in the spirit of
adversarial collaboration: the strongest version of geometric ethics is
the one that survives the strongest attacks.

\subsection{E.1 Best Alternative
Explanations}\label{e.1-best-alternative-explanations}

Objection 1: Classical non-commutativity without quantum structure. The
order effects predicted in Chapter 13 might arise from classical
cognitive biases (anchoring, priming, recency) rather than from genuine
quantum normative dynamics. Response: classical non-commutativity is
undeniable as a phenomenon, but it does not predict the specific pattern
that the framework predicts --- namely, that order effects are present
for cross-stratum evaluation pairs and absent for same-stratum pairs. A
classical account based on anchoring would predict order effects
everywhere, not selectively at stratum boundaries. Study 2 of Appendix C
is designed to distinguish these predictions.

Objection 2: Dataset artifacts and labeling bias. The empirical results
reported in Chapter 17 might reflect artifacts of the Dear Abby corpus
(selection bias, cultural specificity, temporal drift) rather than
genuine moral structure. Response: this is a serious concern, which is
why the BIP testing methodology (Chapter 18.7) explicitly checks for
labeling sensitivity. A Bond Index Bd \textgreater{} 0.1 on the labeling
transform would indicate that the results are label-dependent --- the
framework provides its own diagnostic for this failure mode.

Objection 3: Overfitting to philosophical taxonomy. The nine dimensions
of the moral manifold might reflect the philosophical categories chosen
by the author, not the intrinsic structure of moral situations.
Response: the nine dimensions are not chosen but derived. The Whitney
stratification theorem requires that the moral manifold have at least as
many dimensions as the codimension of its constraint surfaces. The
empirical claim (Study 4, Appendix C) is that factor analysis of human
moral judgments recovers ≥ 8 independent factors --- a prediction that
is falsifiable and has not been engineered.

Objection 4: ``Just category theory in disguise.'' The framework might
be an elaborate restatement of abstract nonsense --- category-theoretic
structure dressed up in geometric language, with no empirical content
beyond what the abstract formalism already implies. Response: geometric
ethics provides physical (geometric) content that category theory alone
does not. The moral metric g (Chapter 9) is not an abstract arrow but a
concrete bilinear form with measurable signature. The curvature (Chapter
10) generates specific predictions about path-dependence of moral
deliberation. The D₄ gauge group (Chapter 12) is a specific finite
group, not an arbitrary symmetry category. These are testable, not
tautological.

Objection 5: The framework is too flexible. With nine dimensions, a
stratified manifold, a metric, and a connection, the framework has
enough degrees of freedom to fit any moral judgment ex post --- making
it unfalsifiable. Response: this is the strongest objection. The
framework does have many parameters. However, the Bond Invariance
Principle and the gauge group D₄ × U(1) are not free parameters --- they
are derived from axioms and constrain all evaluations. The falsification
criteria in Section E.3 specify what would break the framework. A
framework that can be falsified by five different empirical tests is not
unfalsifiable.

\subsection{E.2 Known Failure Cases}\label{e.2-known-failure-cases}

The framework struggles with three classes of moral situations. First,
purely novel situations with no stratum precedent (the ``first time''
problem): when a situation has never been encountered before, the
stratification may not assign it to any existing stratum, and the
evaluation must proceed in the penumbral zone (Definition 8.9) with all
its attendant uncertainty. This is not a bug but a feature --- the
framework correctly reports high uncertainty when uncertainty is
warranted --- but it does mean that the framework provides less guidance
for genuinely novel moral situations.

Second, extreme value pluralism where no common metric exists. If two
moral dimensions are genuinely incommensurable (g\_\{μν\} = 0 for all
off-diagonal components), the framework reports a degenerate metric ---
a moral singularity (Definition 5.7) where comparison is impossible.
Again, this is a correct report of an objective feature of the moral
landscape, but it means the framework cannot resolve deep value
conflicts that humanity itself has not resolved.

Third, situations where the grounding function Ψ itself is contested. If
stakeholders disagree about what the morally relevant facts are (not
just how to weight them), the framework cannot adjudicate --- it
operates downstream of the grounding. This is the ``garbage in, garbage
out'' limitation: the framework's guarantees (BIP, gauge invariance,
audit integrity) apply to the evaluation given a grounding, not to the
grounding itself. The grounding integrity requirements of Chapter 19
provide partial mitigation but not a complete solution.

\subsection{E.3 What Would Falsify the
Framework}\label{e.3-what-would-falsify-the-framework}

The following empirical results would substantially falsify the
geometric ethics framework. (i) Human moral judgments systematically
violate the Bond Invariance Principle in ways not explained by known
cognitive biases. If bilingual invariance fails (Study 1, Appendix C)
with rank correlation \textless{} 0.60 and no anchoring/priming
explanation, the BIP is empirically false as a description of human
moral cognition. (ii) The D₄ group structure of Hohfeldian positions is
empirically rejected. If the four Hohfeldian correlatives (right/duty,
liberty/no-right, power/liability, immunity/disability) do not exhibit
the predicted D₄ symmetry in human usage patterns, the algebraic
foundation of the framework is wrong.

(iii) Moral dimensions are not independent. If factor analysis (Study 4,
Appendix C) recovers fewer than 6 independent factors, the
nine-dimensional manifold is an overparameterization and the framework's
dimensional claims are false. (iv) Order effects in moral judgment are
absent or uniform. If Study 2 shows no order effects in any condition,
or if order effects are present in same-stratum conditions, the quantum
normative dynamics model (Chapter 13) is wrong. (v) The No Escape
Theorem is circumvented. If a known AI architecture can be shown to
violate norm constraints while operating within the DEME framework,
Theorem 18.1 is false and the structural containment claim collapses.

We emphasize that falsification of any single claim does not invalidate
the entire framework. The geometric structure (manifold, metric,
connection) is logically independent of the quantum normative dynamics
(Chapter 13), which is independent of the gauge group identification
(Chapter 12). Each component can be tested and, if necessary, revised
without destroying the others.

We further note that the framework is falsifiable in practice, not
merely in principle. Each of the five tests above can be conducted with
existing technology, existing participant pools, and existing
statistical methods. The protocols are specified in Appendix C. We
invite the community to conduct these tests.

\subsection{E.4 Red-Team Plan}\label{e.4-red-team-plan}

Aligned with the threat model of Section 17.0.6, we propose three
adversarial attack vectors against the framework. Attack 1: Adversarial
grounding. Craft inputs that exploit the grounding function Ψ to
misclassify strata --- for example, scenarios that are linguistically
encoded to trigger the wrong semantic gate. This tests the robustness of
the grounding layer and the semantic gate activation predicates.
Expected outcome: some adversarial inputs succeed (the grounding
function is not perfect), but the audit artifact records the gate
activation, enabling post-hoc detection.

Attack 2: Audit forgery. Attempt to produce valid-looking audit
artifacts for computations that did not occur. This tests the
cryptographic binding condition added in the Tier 2 hardening (Section
11 of the main text). Expected outcome: forgery fails if the hash
commitment scheme is correctly implemented; succeeds if there is a
weakness in the commitment chain. Any successful forgery is a
security-critical bug that must be patched.

Attack 3: Gauge manipulation. Find transformations that are formally
bond-preserving (they satisfy the mathematical definition of an
admissible transformation) but morally significant by human judgment.
For example: renaming ``torture'' as ``enhanced interrogation'' is a
linguistic relabeling (bond-preserving by the letter of the BIP) but
arguably changes the moral content. This tests whether the BIP is too
permissive --- whether the definition of ``admissible transformation''
needs tightening.

We view these attacks as constructive contributions to the framework,
not as threats. A framework that cannot be attacked cannot be trusted.
The red-team plan is published here precisely so that the community can
execute it.

\subsection{E.5 Strongest Competing
Frameworks}\label{e.5-strongest-competing-frameworks}

Three existing frameworks are serious competitors to geometric ethics,
each capturing important features that GE also claims. Moral Foundations
Theory (Haidt 2012) is empirically grounded and identifies multiple
moral dimensions (care, fairness, loyalty, authority, sanctity,
liberty). However, MFT provides no invariance principle, no dynamics (no
analog of parallel transport or curvature), and no formal verification
mechanism. GE's delta: the geometric structure that turns moral
dimensions from a taxonomy into a calculable theory.

Contractualism (Scanlon 1998) provides a structural account of moral
justification (``what no one could reasonably reject'') that captures
the invariance intuition. However, contractualism is not computational:
there is no algorithm for determining what is ``reasonable,'' no audit
mechanism, and no formal containment guarantee for AI systems. GE's
delta: the computational implementation (DEME) and the formal
containment proof (No Escape Theorem). AI Safety via Debate (Irving et
al. 2018) is computational and addresses alignment, but relies on debate
quality rather than structural guarantees. GE's delta: structural
containment that does not depend on the quality of adversarial probing.

We invite adversarial engagement with this framework. The strongest
version of geometric ethics is the one that survives the strongest
attacks.

\subsection{E.6 The Reification
Objection}\label{e.6-the-reification-objection}

Objection 6: The framework builds an elaborate mathematical structure
--- manifolds, tensors, gauge groups, conservation laws --- and then
draws moral conclusions from it. The objection is that this commits a
reification error: treating the internal structure of a representational
tool as though it reveals something about moral reality. Just as the
success of modal logic in formalizing possibility and necessity does not
prove that possible worlds exist, the success of differential geometry
in formalizing moral evaluation does not prove that the moral manifold
exists. The conclusions may be artifacts of the formalism, not features
of the domain.

This is the strongest philosophical objection to the framework, and it
deserves a careful answer. The common failure mode runs: (i) begin with
domain-calibrated intuitions (moral evaluation has multiple dimensions,
context matters, some boundaries are sharp); (ii) codify the intuitions
into a formal tool (the moral manifold with metric, stratification, and
gauge structure); (iii) develop the tool's internal logic (Noether's
theorem, the No Escape Theorem, information monotonicity); (iv) read the
conclusions off as discoveries about moral reality. If the framework
follows this trajectory uncritically, the objection lands.

Five features of the framework's architecture block the reification
trajectory. First, explicit epistemic labeling: every formal statement
in the manuscript carries a status tag (Appendix F catalogues all 112
statements as Formal Definition, Modeling Axiom, Conditional Theorem,
Proved, or Conjecture). Second, conditional theorems: results follow
from stated assumptions, not from the formalism's internal logic --- if
the assumptions are wrong, the theorems do not hold, regardless of how
elegant the mathematics is. Third, empirical testing: the framework's
predictions are tested against data (Chapter 17), and the data can
falsify them --- the framework specifies its own falsification
conditions (§E.3). Fourth, Gödelian modesty: the framework is a formal
system that cannot prove its own consistency or adequacy; its warrant is
external, not internal. Fifth, the governance account of Chapter 9: the
moral metric is not discovered but governed --- the framework does not
claim to find moral truth but to provide the structural vocabulary
within which moral governance operates.

On the instrumentalist reading --- which this book endorses --- the
framework's theorems are not claims about what is ``really there'' in
the moral world. They are claims about what follows from stated modeling
choices. The No Escape Theorem does not say that AI systems are
``really'' geometrically constrained; it says that any system satisfying
Requirements 1--4 is constrained in the ways the theorem describes, and
whether those requirements are the right requirements is an engineering
and governance decision, not a mathematical one. The invariance of harm
accounting does not say that harm is ``really'' conserved like energy;
it says that if the BIP holds and the Lagrangian is C², then a conserved
quantity exists, and calling it ``harm'' is a modeling interpretation
whose adequacy is empirical. The framework provides vocabulary, not
verdicts.

The appropriate analogy is not to metaphysics but to engineering. When a
structural engineer says that a bridge has a certain stress tensor, she
is not making a claim about the ultimate nature of matter. She is
deploying a mathematical model that predicts whether the bridge will
hold. The model's authority is pragmatic: it works, it predicts, it can
be falsified by a collapse. If a better model arrives --- one that
predicts more accurately, applies to more materials, handles more edge
cases --- the old model is revised or replaced. The moral tensor has the
same status. It is a tool that captures structure which scalar
alternatives miss, and its ``reality'' is measured by predictive
success, not metaphysical correspondence. The reification objection
asks: ``Is the moral manifold real?'' The pragmatist answer is: ``It is
as real as the stress tensor in a bridge --- which is to say, real
enough to bear weight.''

\section{Appendix F: Mathematical Ledger --- Status of Formal
Claims}\label{appendix-f-mathematical-ledger-status-of-formal-claims}

This appendix catalogues every numbered formal statement in the
manuscript --- definitions, axioms, theorems, propositions, lemmas,
corollaries, and conventions --- together with its epistemic status and
dependencies. The purpose is transparency: a reader should be able to
see at a glance what is assumed (as a modeling choice mapping ethical
phenomena to mathematical structures), what is proved (given those
assumptions), and what is conjectured or proposed for future
investigation.

Five epistemic categories are used. Formal Definition: a standard
mathematical construction carrying no ethical modeling claim (e.g.,
``topological manifold''). Modeling Axiom: a mathematical structure
posited as a model of an ethical phenomenon; these are the bridge
assumptions that connect the mathematics to moral reasoning, and they
are the primary target of empirical and philosophical scrutiny.
Conditional Theorem: a result that is mathematically proved given stated
assumptions (listed in the Dependencies column); the theorem is as
strong as its premises. Proved: a result that follows from standard
mathematics alone, without ethical modeling assumptions.
Conjecture/Proposal: a formally stated claim that is either unproved,
depends on open problems, or represents a proposed extension not yet
implemented. The table is organized by chapter. ``TA'' denotes the
chapter's Technical Appendix.

{\def\LTcaptype{none} % do not increment counter
\begin{longtable}[]{@{}
  >{\raggedright\arraybackslash}p{(\linewidth - 8\tabcolsep) * \real{0.1311}}
  >{\raggedright\arraybackslash}p{(\linewidth - 8\tabcolsep) * \real{0.3107}}
  >{\raggedright\arraybackslash}p{(\linewidth - 8\tabcolsep) * \real{0.2330}}
  >{\raggedright\arraybackslash}p{(\linewidth - 8\tabcolsep) * \real{0.2336}}
  >{\raggedright\arraybackslash}p{(\linewidth - 8\tabcolsep) * \real{0.0916}}@{}}
\toprule\noalign{}
\begin{minipage}[b]{\linewidth}\raggedright
\textbf{Ref}
\end{minipage} & \begin{minipage}[b]{\linewidth}\raggedright
\textbf{Name}
\end{minipage} & \begin{minipage}[b]{\linewidth}\raggedright
\textbf{Status}
\end{minipage} & \begin{minipage}[b]{\linewidth}\raggedright
\textbf{Dependencies}
\end{minipage} & \begin{minipage}[b]{\linewidth}\raggedright
\textbf{§}
\end{minipage} \\
\midrule\noalign{}
\endhead
\bottomrule\noalign{}
\endlastfoot
Def 4.1 & Topological Manifold & Formal Definition & None & §4.1 \\
Def 4.2 & Smooth Manifold & Formal Definition & Def 4.1 & §4.1 \\
Def 4.3 & Tangent Bundle & Formal Definition & Def 4.2 & §4.2 \\
Def 4.4 & Affine Connection & Formal Definition & Def 4.3 & §4.6 \\
Def 4.5 & Riemann Curvature Tensor & Formal Definition & Def 4.4 &
§4.6 \\
Def 4.6 & Fiber Bundle & Formal Definition & Def 4.2 & §4.7 \\
Def 4.7 & Whitney Stratification & Formal Definition & Def 4.1 & §4.8 \\
Def 4.8 & Gauge Invariance & Formal Definition & Def 4.6 & §4.9 \\
Def 5.1 & Moral Manifold (Informal) & Modeling Axiom & Def 4.2 & §5.2 \\
Def 5.2 & Dimension of the Moral Manifold (k = 9) & Modeling Axiom & Def
5.1 & §5.2 \\
Def 5.3 & Agent Bundle & Modeling Axiom & Def 4.6, Def 5.1 & §5.4 \\
Axiom 5.1 & Bond Invariance Principle (BIP) & Modeling Axiom & Def 5.1,
Def 4.8 & §5.3 \\
Def 5.4 & Moral Stratification & Modeling Axiom & Def 4.7, Def 5.1 &
§5.6 \\
Def 5.5 & Constraint Set & Modeling Axiom & Def 5.4 & §5.6 \\
Def 5.6 & Absorbing Stratum & Modeling Axiom & Def 5.4 & §5.6 \\
Def 5.7 & Moral Singularity & Modeling Axiom & Def 5.4 & §5.6 \\
Def A.1--A.6 & Moral Manifold (Formal Defs) & Formal Definition & Defs
4.1--4.7, Def 5.1 & §5 TA \\
Def 6.1 & Obligation Vector Oµ & Modeling Axiom & Def 5.1, Def 4.3 &
§6.2 \\
Def 6.2 & Interest Covector Iµ & Modeling Axiom & Def 5.1 & §6.3 \\
Def 6.3 & Satisfaction S = IµOµ & Modeling Axiom & Defs 6.1--6.2 &
§6.4 \\
Def 6.4 & Evaluation Tensor Eµᵥ & Modeling Axiom & Defs 6.1--6.3 &
§6.6 \\
Prop 6.1 & Invariance of Satisfaction & Proved & Defs 6.1--6.3, BIP & §6
TA \\
Prop 6.2 & Decomposition of Disagreement & Proved & Def 6.4 & §6 TA \\
Prop 6.3 & Metric Dependence of Gradient & Proved & Def 4.4 & §6 TA \\
Def 8.1 & Stratified Moral Space & Formal Definition & Defs 4.7, 5.4 &
§8.2 \\
Defs 8.2--8.3 & Whitney Conditions A and B & Formal Definition & Def 4.7
& §8.2 \\
Def 8.4 & Threshold Boundary & Modeling Axiom & Def 8.1 & §8.3 \\
Def 8.5 & Phase Transition Boundary & Modeling Axiom & Def 8.1 & §8.3 \\
Def 8.6 & Absorbing Stratum & Modeling Axiom & Def 8.1 & §8.3 \\
Def 8.7 & Constraint Surface & Modeling Axiom & Def 8.1 & §8.3 \\
Def 8.8 & Semantic Gate & Modeling Axiom & Def 8.1 & §8.4 \\
Def 8.9 & Penumbral Zone & Modeling Axiom & Defs 8.4--8.5 & §8.5 \\
Def 8.10 & Moral Phase Transition & Modeling Axiom & Defs 8.1, 8.5 &
§8.6 \\
Def 8.11 & Boundary Crossing Data & Formal Definition & Def 8.1 &
§8.8 \\
Conv 8.1 & DG Operations at Stratum Boundaries & Formal Definition &
Defs 8.1, 8.11, 4.4 & §8.8 \\
Prop 8.1 & Existence of Moral Whitney Stratification & Conditional
Theorem & Defs 5.4, 4.7 & §8 TA \\
Prop 8.2 & Absorbing Strata Are Unique Attractors & Conditional Theorem
& Def 8.6 & §8 TA \\
Prop 8.3 & D₄ Structure of Hohfeldian Transitions & Conditional Theorem
& Def 8.8, Axiom A1--A3 & §8 TA \\
Def 9.1 & Partial Order on Metrics & Formal Definition & Def 4.5 &
§9.4 \\
Def A.1 & Admissible Metric & Formal Definition & Def 4.5, Axiom 5.1 &
§9 TA \\
Prop 9.1 & Non-Uniqueness of Admissible Metrics & Conditional Theorem &
Def A.1 & §9 TA \\
Prop 9.2 & Structural Invariants Constrain the Metric & Conditional
Theorem & Def A.1, BIP & §9 TA \\
Lemma 9.1 & D₄ Constraint on Metric Components & Conditional Theorem &
Def A.1, Axiom A1--A3 & §9 TA \\
Lemma 9.2 & U(1)ₕ Constraint on Harm Dimension & Conditional Theorem &
Def A.1, Axiom A4 & §9 TA \\
Thm 9.1 & Admissible Metric Space Structure & Conditional Theorem & Def
A.1, BIP & §9 TA \\
Thm 9.2 & Structured Pluralism & Conditional Theorem & Thm 9.1 & §9
TA \\
Def 10.1 & Moral Connection & Formal Definition & Def 4.4, Def 5.1 &
§10.3 \\
Def 10.2 & Parallel Transport & Formal Definition & Def 10.1 & §10.4 \\
Def 10.3 & Holonomy & Formal Definition & Def 10.2 & §10.5 \\
Prop 10.1 & Parallel Transport Preserves Satisfaction & Conditional
Theorem & Defs 10.2, 6.3 & §10 TA \\
Prop 10.2 & Holonomy Detects Curvature & Conditional Theorem & Defs
10.3, 4.5 & §10 TA \\
Prop 10.3 & Geodesic Deviation Equation & Proved & Def 4.5 & §10 TA \\
Def 11.1 & Moral Search Problem & Modeling Axiom & Def 5.1, Def 6.4 (BF)
& §11.2 \\
Def 11.2 & A* Moral Search & Modeling Axiom & Def 11.1 & §11.2 \\
Prop 11.1 & Obligation-Heuristic Correspondence & Modeling Axiom & Def
6.1, Def 11.2 & §11.3 \\
Def 11.3 & Admissibility & Formal Definition & Standard & §11.4 \\
Def 11.4 & Consistency & Formal Definition & Standard & §11.4 \\
Thm 11.1 & Admissibility of Core Moral Heuristics & Conditional Theorem
& Def 11.3, Ch 8 stratification & §11.4 \\
Prop 11.2 & Consistency from Metric Compatibility & Proved & Def 4.5,
Def 11.4 & §11.4 \\
Thm 11.2 & Intractability of Exact Moral Geodesic Planning & Conditional
Theorem & Def 5.1, Def 11.1 & §11.5 \\
Cor 11.1 & Evolutionary Necessity of Moral Heuristics & Conditional
Theorem & Thm 11.2 & §11.5 \\
Cor 11.2 & Heuristic Imperfection Is Structural & Conditional Theorem &
Thm 11.2 & §11.5 \\
Prop 11.3 & Satisfaction as Directional Derivative & Proved & Def 6.1,
Prop 11.1 & §11.6 \\
Prop 11.4 & A* Optimality on Stratified Spaces & Proved & Def 11.2, Def
8.1 (Whitney) & §11 TA \\
Axiom A1 & Four Hohfeldian Positions & Modeling Axiom & None & §12.3 \\
Axiom A2 & Two Independent Involutions & Modeling Axiom & Axiom A1 &
§12.3 \\
Axiom A3 & Cyclic Ordering & Modeling Axiom & Axioms A1--A2 & §12.3 \\
Axiom A4 & One Continuous Conserved Quantity & Modeling Axiom & None &
§12.3 \\
Axiom A5 & Commutativity of Discrete and Continuous & Modeling Axiom &
Axioms A1--A4 & §12.3 \\
Def 12.1 & Moral Action (Lagrangian) & Modeling Axiom & Defs 4.4, 5.1 &
§12.4 \\
Def 12.2 & Harm (Noether Charge) & Modeling Axiom & Thm 12.1, Prop 12.1
& §12.6 \\
Def 12.3 & Harm Ledger & Modeling Axiom & Def 12.2 & §12.6 \\
Thm 12.1 & Moral Noether's Theorem & Conditional Theorem & BIP + C²
Lagrangian & §12.5 \\
Thm 12.2 & Discrete Symmetry = D₄ & Conditional Theorem & Axioms A1--A3
& §12.3 \\
Thm 12.3 & Gauge Group Uniqueness: G = D₄ × U(1)ₕ & Conditional Theorem
& Axioms A1--A5 & §12.3 \\
Prop 12.1 & Continuous Component is U(1) & Conditional Theorem & Axiom
A4 & §12.3 \\
Cor A.1 & Number of Conservation Laws & Conditional Theorem & Thm A.1 &
§11 TA \\
Prop A.2 & Discrete Conservation from D₄ & Conditional Theorem & Thm
12.2 & §11 TA \\
Thm A.1 & Moral Noether (Formal Statement) & Conditional Theorem & BIP +
C² Lagrangian & §11 TA \\
Lemma A.2 & Aut(C₄) ≅ D₄ & Proved & None & §11 TA \\
Lemma A.3 & Compact Connected 1-D Lie Group = U(1) & Proved & None & §11
TA \\
Def 13.1 & Moral Hilbert Space & Conjecture/Proposal & Def 5.1 &
§13.2 \\
Def 13.2 & Moral Observable & Conjecture/Proposal & Def 13.1 & §13.2 \\
Def 13.3 & Moral Density Matrix & Conjecture/Proposal & Def 13.1 &
§13.3 \\
Def 13.4 & Moral Decoherence & Conjecture/Proposal & Def 13.3 & §13.4 \\
Def 13.5 & Stratified Lagrangian & Modeling Axiom & Def 12.1, Def 8.1 &
§13.5 \\
Def 13.6 & Junction Conditions & Formal Definition & Defs 12.5, 8.11 &
§13.5 \\
Def 13.7 & Moral Hamiltonian & Conjecture/Proposal & Def 13.5 & §13.6 \\
Def 13.8 & Moral Tunneling & Conjecture/Proposal & Def 13.7, Def 8.7 &
§13.7 \\
Def 13.9 & Moral Entanglement & Conjecture/Proposal & Def 13.3 &
§13.8 \\
Thm 13.1 & Decomposition of H by Strata & Conditional Theorem & Def
13.1, Def 8.1 & §12 TA \\
Prop 13.1 & Moral Uncertainty Principle & Conditional Theorem & Defs
12.1--12.2 & §13.3 \\
Prop 13.2 & Conservation of Probability & Conditional Theorem & Def 13.7
& §12 TA \\
Prop 13.3 & Tunneling Rate Through Moral Barrier & Conditional Theorem &
Def 13.8 & §12 TA \\
Prop 13.4 & Entanglement Entropy & Conditional Theorem & Def 13.9 & §12
TA \\
Def 14.1 & Collective Agency Tensor & Modeling Axiom & Defs 6.1--6.4 &
§14.2 \\
Def 14.2 & Collective Decomposition & Formal Definition & Def 14.1 &
§14.2 \\
Def 14.3 & Structure Tensor & Modeling Axiom & Def 14.1 & §14.4 \\
Def 14.4 & Responsibility Tensor & Modeling Axiom & Def 14.1 & §14.5 \\
Prop 14.1 & Emergent Obligation Criterion & Conditional Theorem & Def
14.1 & §14.3 \\
Prop 14.2 & Emergent Moral Property & Conditional Theorem & Defs
13.1--13.2 & §14.3 \\
Prop 14.3 & Decomposition Theorem & Conditional Theorem & Def 14.2 & §13
TA \\
Prop 14.4 & Responsibility Inequality & Conditional Theorem & Def 14.4 &
§13 TA \\
Prop 14.5 & Structural Coupling & Conditional Theorem & Def 14.3 & §13
TA \\
Def 15.1 & Contraction Loss & Formal Definition & Defs 6.1--6.4 &
§15.2 \\
Def 15.2 & Moral Residue & Formal Definition & Def 15.1 & §15.3 \\
Prop 15.1 & Inevitability of Residue & Proved & Defs 14.1--14.2 &
§15.3 \\
Prop 15.2 & Harm Survives Contraction & Conditional Theorem & Defs 14.1,
11.2 & §15.4 \\
Prop 15.3 & Non-Commutativity of Contraction & Proved & Def 15.1 & §14
TA \\
Prop 15.4 & Residue of Maximin Contraction & Conditional Theorem & Defs
14.1--14.2 & §14 TA \\
Prop 15.5 & Transparent Contraction and BIP & Conditional Theorem & Def
15.1, BIP & §14 TA \\
Rmk A.1 & Residue as Moore--Penrose Pseudoinverse & Formal Definition &
Def 15.2, metric gµᵥ & §14 TA \\
Def 16.1 & Theory Covariance Tensor & Modeling Axiom & Def 6.4 &
§16.3 \\
Def 16.2 & Robust Obligation & Formal Definition & Def 16.1 & §16.4 \\
Prop 16.1 & Stability of Robust Obligations & Conditional Theorem & Def
16.2 & §16.4 \\
Prop 16.2 & Structure of the Robust Core & Conditional Theorem & Def
16.2 & §15 TA \\
Prop 16.3 & Irreducibility of Theory Uncertainty & Conditional Theorem &
Def 16.1 & §15 TA \\
Prop 16.4 & Hedging as Minimax Regret & Conditional Theorem & Def 16.2 &
§15 TA \\
Def 18.1 & Grounding Adequacy Predicate & Modeling Axiom & Def 5.1 &
§18.7 \\
Thm 18.1 & No Escape Theorem & Conditional Theorem & Reqs 1--4, BIP &
§18.6 \\
Thm 18.2 & No Escape with Formal Adequacy & Conditional Theorem & Thm
18.1, Def 18.1 & §18.7 \\
Lemma 18.1 & Coverage Completeness & Conditional Theorem & Def 18.1 &
§18.8 \\
Lemma 18.2 & Robustness Implies Lipschitz Stability & Conditional
Theorem & Def 18.1 & §18.8 \\
Prop 18.1 & Invariance Violation Decomposition & Conditional Theorem &
Thm 18.1, BIP & §17 TA \\
Prop 18.2 & Contraction Mismatch Bound & Conditional Theorem & Thm 18.1
& §17 TA \\
Prop 18.3 & No Escape --- Conditional Guarantee & Conditional Theorem &
Reqs 1--4 & §17 TA \\
Prop 18.4 & Alignment Gap Convergence & Conditional Theorem & Thm 18.1,
BIP & §17 TA \\
Prop 18.5 & Testability of Adequacy Clauses & Conditional Theorem & Def
18.1 & §17 TA \\
Def 19.1 & Norm-Constrained Scoring Game (NCSG) & Modeling Axiom & Def
5.1 & §19.5 \\
Def 19.5 & Rank-2 Canonicalization Tensor & Conjecture/Proposal & Def
4.6, Thm 18.1 & §19.13 \\
Def 19.6 & Rank-3 Uncertainty-Aware Canon. Tensor & Conjecture/Proposal
& Def 19.5, Def 6.4 & §19.13 \\
Def 19.7 & Rank-5 Coalition-Parsing Canon. Tensor & Conjecture/Proposal
& Def 19.5, Def 14.3 & §19.13 \\
Def 19.8 & Rank-4 Temporal Canon. Tensor & Conjecture/Proposal & Def
19.5, Def 10.3 & §19.13 \\
Prop 19.1 & NCSG Well-Posedness & Conditional Theorem & Def 19.1 & §18
TA \\
Prop 19.2 & Compilation Residue & Conditional Theorem & Defs 18.1, 14.1
& §18 TA \\
Prop 19.3 & Democratic Aggregation & Conditional Theorem & Def 19.1 &
§18 TA \\
Prop 19.4 & Bond Index Convergence & Conditional Theorem & Def 19.1, BIP
& §18 TA \\
Prop 19.5 & Gauge-Fiber Severance & Conditional Theorem & Def 19.5 &
§19.13 \\
Prop 20.1 & Curvature Detection Criterion & Conjecture/Proposal & Def
10.3 & §19 TA \\
Prop 20.2 & Torsion Detection Criterion & Conjecture/Proposal & Def 4.4
& §19 TA \\
Prop 20.3 & Meta-Metric Existence & Conjecture/Proposal & Def 9.1 & §19
TA \\
\end{longtable}
}

\section{Bibliography}\label{bibliography}

\subsection{Primary Sources by the
Author}\label{primary-sources-by-the-author}

Bond, A. H. (2025a). ``The Bond Invariance Principle: Cross-Lingual
Validation of Deontic Structure.'' Submitted to \emph{IEEE Transactions
on Artificial Intelligence}, February 2026. Preprint, San José State
University.

Bond, A. H. (2025b). ``No Escape: Mathematical Containment for
Artificial Agents.'' Working paper, San José State University.

Bond, A. H. (2025c). ``Stratified Quantum Normative Dynamics: A
Geometric Framework for Moral Deliberation.'' Working paper, San José
State University.

Bond, A. H. (2026a). ``The Dear Abby Corpus: Thirty-Two Years of Natural
Moral Reasoning.'' Working paper, San José State University.

Bond, A. H. (2026b). ``The Dear Ethicist Game: Engineered Probes for
Geometric Ethics.'' Working paper, San José State University.

Bond, A. H. (2026c). ``ErisML: A Modeling Language for Geometric AI
Governance.'' Technical specification, San José State University.

Bond, A. H. (2026d). ``DEME: A Democratically Governed Ethics Module
Engine.'' Technical specification, San José State University.

Bond, A. H. (2026e). ``The Grand Unified AI Safety Stack (GUASS): A
Seven-Layer Architecture for Ethical AI.'' Technical specification, San
José State University.

Bond, A. H. (2026f). ``Stratified Geometric Ethics: A Foundational
Framework.'' Working paper (v10), San José State University. February
2026.

Bond, A. H. (2026g). \emph{The Geometry of Good: A Mathematical
Framework for Ethics}. Book manuscript (v1.0.2), San José State
University. February 2026.

\subsection{Philosophy: Ethics and Political
Theory}\label{philosophy-ethics-and-political-theory}

Aristotle. \emph{Nicomachean Ethics}. Translated by W. D. Ross, revised
by J. O. Urmson, in \emph{The Complete Works of Aristotle}, edited by
Jonathan Barnes, vol.~2, 1729--1867. Princeton: Princeton University
Press, 1984.

Arrow, Kenneth J. \emph{Social Choice and Individual Values}. 2nd
ed.~New Haven: Yale University Press, 1963. First published 1951.

Augustine of Hippo. \emph{Confessions}. Translated by Henry Chadwick.
Oxford: Oxford University Press, 1991.

Beauchamp, Tom L., and James F. Childress. \emph{Principles of
Biomedical Ethics}. 8th ed.~New York: Oxford University Press, 2019.

Daniels, Norman. \emph{Justice and Justification: Reflective Equilibrium
in Theory and Practice}. Cambridge: Cambridge University Press, 1996.

Dewey, John. \emph{The Quest for Certainty: A Study of the Relation of
Knowledge and Action}. New York: Minton, Balch \& Company, 1929.

Foot, Philippa. ``The Problem of Abortion and the Doctrine of the Double
Effect.'' \emph{Oxford Review} 5 (1967): 5--15.

Geach, Peter T. ``Assertion.'' \emph{Philosophical Review} 74, no. 4
(1965): 449--465.

Habermas, Jürgen. \emph{The Theory of Communicative Action}. 2 vols.
Translated by Thomas McCarthy. Boston: Beacon Press, 1984--1987.
Originally published 1981.

Habermas, Jürgen. \emph{Moral Consciousness and Communicative Action}.
Translated by Christian Lenhardt and Shierry Weber Nicholsen. Cambridge,
MA: MIT Press, 1990.

Hohfeld, Wesley Newcomb. ``Some Fundamental Legal Conceptions as Applied
in Judicial Reasoning.'' \emph{Yale Law Journal} 23, no. 1 (1913):
16--59.

Hohfeld, Wesley Newcomb. ``Fundamental Legal Conceptions as Applied in
Judicial Reasoning.'' \emph{Yale Law Journal} 26, no. 8 (1917):
710--770.

James, William. \emph{Pragmatism: A New Name for Some Old Ways of
Thinking}. New York: Longmans, Green, and Co., 1907.

Kant, Immanuel. \emph{Groundwork of the Metaphysics of Morals}.
Translated and edited by Mary Gregor. Cambridge: Cambridge University
Press, 1998. Originally published 1785.

Keats, John. Letter to George and Tom Keats, December 21, 1817. In
\emph{The Letters of John Keats}, edited by Hyder Edward Rollins,
vol.~1, 191--194. Cambridge, MA: Harvard University Press, 1958.

Misak, Cheryl. \emph{The American Pragmatists}. Oxford: Oxford
University Press, 2013.

Nussbaum, Martha C. \emph{Women and Human Development: The Capabilities
Approach}. Cambridge: Cambridge University Press, 2000.

Peirce, Charles Sanders. \emph{Collected Papers of Charles Sanders
Peirce}. Edited by Charles Hartshorne and Paul Weiss (vols. 1--6) and
Arthur W. Burks (vols. 7--8). Cambridge, MA: Harvard University Press,
1931--1958.

Price, Huw. \emph{Expressivism, Pragmatism and Representationalism}.
Cambridge: Cambridge University Press, 2013.

Rawls, John. \emph{A Theory of Justice}. Revised ed.~Cambridge, MA:
Harvard University Press, 1999. First published 1971.

Ross, W. D. \emph{The Right and the Good}. Edited by Philip
Stratton-Lake. Oxford: Clarendon Press, 2002. First published 1930.

Scanlon, T. M. \emph{What We Owe to Each Other}. Cambridge, MA: Harvard
University Press, 1998.

Sen, Amartya. \emph{Development as Freedom}. New York: Alfred A. Knopf,
1999.

Sen, Amartya. \emph{The Idea of Justice}. Cambridge, MA: Harvard
University Press, 2009.

Young, Iris Marion. \emph{Responsibility for Justice}. Oxford: Oxford
University Press, 2011.

\subsection{Philosophy: Moral
Uncertainty}\label{philosophy-moral-uncertainty}

MacAskill, William, Krister Bykvist, and Toby Ord. \emph{Moral
Uncertainty}. Oxford: Oxford University Press, 2020.

\subsection{Mathematics: Differential Geometry and
Topology}\label{mathematics-differential-geometry-and-topology}

Ambrose, Warren, and I. M. Singer. ``A Theorem on Holonomy.''
\emph{Transactions of the American Mathematical Society} 75, no. 3
(1953): 428--443.

do Carmo, Manfredo Perdigão. \emph{Riemannian Geometry}. Translated by
Francis Flaherty. Boston: Birkhäuser, 1992.

Frankel, Theodore. \emph{The Geometry of Physics: An Introduction}. 3rd
ed.~Cambridge: Cambridge University Press, 2011.

Kobayashi, Shoshichi, and Katsumi Nomizu. \emph{Foundations of
Differential Geometry}. 2 vols. New York: Wiley-Interscience,
1963--1969.

Lee, John M. \emph{Introduction to Smooth Manifolds}. 2nd ed.~Graduate
Texts in Mathematics 218. New York: Springer, 2012.

Nakahara, Mikio. \emph{Geometry, Topology and Physics}. 2nd ed.~Graduate
Student Series in Physics. Bristol: Institute of Physics Publishing,
2003.

Pflaum, Markus J. \emph{Analytic and Geometric Study of Stratified
Spaces}. Lecture Notes in Mathematics 1768. Berlin: Springer, 2001.

Whitney, Hassler. ``Tangents to an Analytic Variety.'' \emph{Annals of
Mathematics} 81, no. 3 (1965): 496--549.

DeWitt, Bryce S. "Quantum Theory of Gravity. I. The Canonical Theory."
Physical Review 160, no. 5 (1967): 1113--1148.

Ebin, David G. "The Manifold of Riemannian Metrics." In Global Analysis
(Proceedings of Symposia in Pure Mathematics 15), edited by S.-S. Chern
and S. Smale, 11--40. Providence: American Mathematical Society, 1970.

\subsection{Mathematics: Algebra and
Logic}\label{mathematics-algebra-and-logic}

Gödel, Kurt. ``On Formally Undecidable Propositions of \emph{Principia
Mathematica} and Related Systems I.'' \emph{Monatshefte für Mathematik
und Physik} 38 (1931): 173--198.

Noether, Emmy. ``Invariante Variationsprobleme.'' \emph{Nachrichten von
der Gesellschaft der Wissenschaften zu Göttingen,
Mathematisch-Physikalische Klasse} (1918): 235--257.

\subsection{Physics}\label{physics}

Einstein, Albert. ``Die Grundlage der allgemeinen Relativitätstheorie.''
\emph{Annalen der Physik} 354, no. 7 (1916): 769--822.

Yang, Chen-Ning, and Robert L. Mills. ``Conservation of Isotopic Spin
and Isotopic Gauge Invariance.'' \emph{Physical Review} 96, no. 1
(1954): 191--195.

\subsection{Quantum Cognition and
Consciousness}\label{quantum-cognition-and-consciousness}

Busemeyer, Jerome R., and Peter D. Bruza. \emph{Quantum Models of
Cognition and Decision}. Cambridge: Cambridge University Press, 2012.

Penrose, Roger. \emph{The Emperor's New Mind: Concerning Computers,
Minds, and the Laws of Physics}. Oxford: Oxford University Press, 1989.

Penrose, Roger. \emph{Shadows of the Mind: A Search for the Missing
Science of Consciousness}. Oxford: Oxford University Press, 1994.

Penrose, Roger, and Stuart Hameroff. ``Consciousness in the Universe: A
Review of the `Orch OR' Theory.'' \emph{Physics of Life Reviews} 11, no.
1 (2014): 39--78.

Pothos, Emmanuel M., and Jerome R. Busemeyer. ``Can Quantum Probability
Provide a New Direction for Cognitive Modeling?'' \emph{Behavioral and
Brain Sciences} 36, no. 3 (2013): 255--274.

\subsection{AI Safety and Governance}\label{ai-safety-and-governance}

EU High-Level Expert Group on Artificial Intelligence. \emph{Ethics
Guidelines for Trustworthy AI}. Brussels: European Commission, 2019.

European Parliament and Council of the European Union. \emph{Regulation
(EU) 2024/1689 Laying Down Harmonised Rules on Artificial Intelligence
(AI Act)}. Official Journal of the European Union, 2024.

Markkula Center for Applied Ethics. ``A Framework for Ethical Decision
Making.'' Santa Clara University, 2021.

National Institute of Standards and Technology. \emph{Artificial
Intelligence Risk Management Framework (AI RMF 1.0)}. NIST AI 100-1.
Gaithersburg, MD: U.S. Department of Commerce, 2023.

\subsection{Literary and Historical
References}\label{literary-and-historical-references}

Styron, William. \emph{Sophie's Choice}. New York: Random House, 1979.

\emph{The Bible}. New Revised Standard Version. Acts 9:1--19 (the
conversion of Saul on the road to Damascus).

Liu An (刘安), ed.~\emph{Huainanzi} (淮南子). c.~139 BCE. The parable of
the Old Man at the Border (塞翁失马, \emph{Sāi Wēng Shī Mǎ}) appears in
Chapter 19, ``Lessons of the Human World'' (人间训).

Robert, Henry M. \emph{Robert's Rules of Order Newly Revised}. 12th
ed.~New York: PublicAffairs, 2020. First published 1876.

\subsection{Supplementary References in Differential
Geometry}\label{supplementary-references-in-differential-geometry}

Berger, Marcel. \emph{A Panoramic View of Riemannian Geometry}. Berlin:
Springer, 2003.

Chern, Shiing-Shen, Wen-Tsün Wu, and Wei-Liang Chow. ``Differential
Geometry.'' In \emph{Handbook of Applicable Mathematics}, vol.~5, edited
by Walter Ledermann. Chichester: Wiley, 1990.

Milnor, John. \emph{Morse Theory}. Annals of Mathematics Studies 51.
Princeton: Princeton University Press, 1963.

Spivak, Michael. \emph{A Comprehensive Introduction to Differential
Geometry}. 5 vols. 3rd ed.~Houston: Publish or Perish, 1999.

\subsection{Supplementary References in
Ethics}\label{supplementary-references-in-ethics}

Gilligan, Carol. \emph{In a Different Voice: Psychological Theory and
Women's Development}. Cambridge, MA: Harvard University Press, 1982.

Noddings, Nel. \emph{Caring: A Feminine Approach to Ethics and Moral
Education}. Berkeley: University of California Press, 1984.

Parfit, Derek. \emph{Reasons and Persons}. Oxford: Clarendon Press,
1984.

Williams, Bernard. ``Ethical Consistency.'' \emph{Proceedings of the
Aristotelian Society}, Supplementary Volumes 39 (1965): 103--124.

Williams, Bernard. \emph{Ethics and the Limits of Philosophy}.
Cambridge, MA: Harvard University Press, 1985.

\subsection{Supplementary References in Gauge Theory and Mathematical
Physics}\label{supplementary-references-in-gauge-theory-and-mathematical-physics}

Baez, John C., and Javier P. Muniain. \emph{Gauge Fields, Knots and
Gravity}. Series on Knots and Everything 4. Singapore: World Scientific,
1994.

Bleecker, David. \emph{Gauge Theory and Variational Principles}. Global
Analysis Pure and Applied Series A 1. Reading, MA: Addison-Wesley, 1981.

Weinberg, Steven. \emph{The Quantum Theory of Fields}. 3 vols.
Cambridge: Cambridge University Press, 1995--2000.

\subsection{Supplementary References in Social Choice and Decision
Theory}\label{supplementary-references-in-social-choice-and-decision-theory}

Fishburn, Peter C. \emph{The Theory of Social Choice}. Princeton:
Princeton University Press, 1973.

Gibbard, Allan. \emph{Wise Choices, Apt Feelings: A Theory of Normative
Judgment}. Cambridge, MA: Harvard University Press, 1990.

Luce, R. Duncan, and Howard Raiffa. \emph{Games and Decisions:
Introduction and Critical Survey}. New York: Wiley, 1957.

\subsection{Supplementary References in AI Ethics and
Alignment}\label{supplementary-references-in-ai-ethics-and-alignment}

Amodei, Dario, Chris Olah, Jacob Steinhardt, Paul Christiano, John
Schulman, and Dan Mané. ``Concrete Problems in AI Safety.'' arXiv
preprint arXiv:1606.06565, 2016.

Bostrom, Nick. \emph{Superintelligence: Paths, Dangers, Strategies}.
Oxford: Oxford University Press, 2014.

Gabriel, Iason. ``Artificial Intelligence, Values, and Alignment.''
\emph{Minds and Machines} 30 (2020): 411--437.

Harris, Sam. The Moral Landscape: How Science Can Determine Human
Values. New York: Free Press, 2010.

Russell, Stuart. \emph{Human Compatible: Artificial Intelligence and the
Problem of Control}. New York: Viking, 2019.

\emph{Note on citation practice.} Many references in this book are to
classical works (Aristotle, Kant, Hohfeld) and to named mathematical
structures (Riemannian manifolds, Noether's theorem, Whitney
stratification) that are part of the standard inheritance of their
respective fields. Where a specific passage is cited --- as with
Aristotle's \emph{Nicomachean Ethics} 1106b21 or Kant's
\emph{Groundwork} 4:421 --- the standard scholarly reference system
(Bekker numbers, Akademie pagination) is used. Where a mathematical
structure is referenced by its originator's name, the foundational
publication is listed above; the reader seeking textbook treatments is
directed to the mathematical preliminaries of Chapter 4 and its
bibliographic notes (Section 4.11).

The author's own prior work, listed at the top of this bibliography,
provides the empirical and technical foundations on which the book's
theoretical framework is built. These papers and technical
specifications are available through the author's institutional page at
San José State University.

\section{Index}\label{index}

\emph{Note: References are to chapter and section numbers (e.g., 5.3 =
Chapter 5, Section 5.3). Boldface indicates the primary definition or
introduction of a term. ``TA'' indicates the Technical Appendix of the
cited chapter.}

\subsection{A}\label{a}

absorbing strata, \textbf{5.6}, 8.3, 8.9, 17.2; \emph{see also}
nullifiers - abuse as, 5.6, 8.3, 17.2 - danger as, 5.6, 8.3 -
impossibility as, 5.6, 8.3 - retraction property, \textbf{8.3} -
uniqueness of attractors, 8.TA

action (moral), \textbf{12.4}; \emph{see also} Lagrangian (moral)

action functional, 4.9, \textbf{12.4}

actual duty, 3.4; \emph{see also} Ross, W. D.

admissible metrics, \textbf{9.6}, 9.TA; \emph{see also} moral metric -
constraints on, 9.6 - lattice structure, 9.8 - non-uniqueness of, 9.TA -
partial order on, \textbf{9.8}

admissible transformations, \textbf{5.5}, 5.TA - Type 1 (coordinate
redescriptions), \textbf{5.5} - Type 2 (perspective shifts),
\textbf{5.5} - Type 3 (theory shifts), \textbf{5.5}

affine connection, \emph{see} connection

agent fiber bundle, \textbf{5.4}, 5.TA, 6.6, 14.7

aggregation problem, 14.2, 15.2; \emph{see also} collective moral agency

AI alignment, 1.3.1, 12.9, \textbf{18.4}, 18.5 - as contraction
mismatch, \textbf{15.11}, 18.5 - as invariance, \textbf{18.4} - three
sources of misalignment, \textbf{18.5} - gap decomposition, 18.5, 18.TA

AI safety, \emph{see} structural containment; No Escape Theorem; GUASS

alignment gap, \textbf{18.5}, 18.TA - convergence, 18.TA

Ambrose-Singer theorem, 10.TA, 29.3

anomaly (D\(_{4}\)), 17.3, 29.5; \emph{see also} D\(_{4}\) dihedral
group

antisymmetrization, \textbf{6.9}

Arrow's impossibility theorem, 15.2

atlas, \textbf{4.1}

audit completeness (Requirement 3), \textbf{18.6}

Augustine of Hippo, 13.7

autonomous hospital ward (case study), \textbf{19.11}

autonomous vehicle dilemma, 8.6

\subsection{B}\label{b}

basis vectors, capabilities as, 3.6

behavioral rules vs.~constitutive constraints, 1.3.1, 18.6

Bianchi identity (contracted), 29.4

BIP, \emph{see} Bond Invariance Principle

Bond Index, \textbf{19.7}, 17.4, 19.TA - convergence, 19.TA - rating
scale, 19.7

Bond Invariance Principle (BIP), 5.5, 1.8, 3.3, 4.9, 9.6, 12.3, 17.3,
18.4, 29.8 - as gauge symmetry, 12.3, 29.8 - cross-lingual validation,
12.8, 17.3 - experiments, 17.3, 8.4, 9.7, 12.8 - formal statement (Axiom
5.1), 5.5 - invariance violation, 18.4; see also invariance testing

boundary, \emph{see} stratum boundary

boundary crossing data, \textbf{8.8}

branch point, 8.6

\subsection{C}\label{c}

canonicalization, mandatory, \textbf{18.6}, 29.8 - approximate
(\(\epsilon\)-correct), 29.8 - at scale, 29.8

capabilities approach (Sen and Nussbaum), 1.3.3, \textbf{3.6}, 5.3 - as
basis for moral space, 3.6 - incompleteness thesis, \textbf{3.6} -
thresholds (Nussbaum), \textbf{3.6}, 8.3

care ethics, 6.7, 7.3, 8.10, 16.2

categorical imperative (Kant), \textbf{3.3} - as invariance condition,
3.3 - Kingdom of Ends as symmetry group, 3.3

character bundle, 6.7; \emph{see also} virtue ethics

Christoffel symbols, \textbf{4.6}, 10.3, 17.2

classical limit (quantum normative dynamics), 13.9

cluster expansion (statistical mechanics analogue), 14.TA

cognitive escape, \textbf{18.6}; \emph{see also} No Escape Theorem

coherence term, 13.3; \emph{see also} density matrix

collective agency tensor, \textbf{14.3} - decomposition of, 14.3 -
emergent component, \textbf{14.3}

collective holonomy, 14.8

collective moral agency, \textbf{14.1}--\textbf{14.13} - collective
parallel transport, \textbf{14.8} - downward causation, 14.9 - hospital
ethics committee (example), \textbf{14.12} - upward emergence, 14.9

compilation residue, \textbf{19.4}, 19.TA

compilation targets (ErisML), \textbf{19.4} - multi-agent RL, 19.4 -
PDDL, 19.4 - PRISM, 19.4 - safe RL, 19.4

compactness (of moral manifold), 5.8, 29.3

cone point, 8.6

connection (moral), \textbf{10.3}; \emph{see also} Levi-Civita
connection - covariant derivative, \textbf{10.3} - metric-compatible,
4.6, 10.3 - torsion-free, 4.6, 10.3

connectedness (of moral manifold), 5.8

conservation of charge, 12.2

invariance of harm accounting, \textbf{12.6}, 1.7, 12.7, 15.10, 17.6,
18.8, 29.2 - as Noether charge, \textbf{12.5}, \textbf{12.6} -
consequences of, 12.7 - empirical evidence for, 17.6 - euphemism and,
12.7, 17.6 - harm survives contraction, \textbf{15.10}

conservation laws (physics), 12.2 - angular momentum, 12.2 - charge,
12.2 - energy, 12.2 - momentum, 12.2

constraint surface (Type IV boundary), \textbf{8.3}, 5.6

constructivism (about the metric), \textbf{9.3}

contraction (tensor to scalar), \textbf{6.4}, 2.6, 4.3, 4.4,
15.1--15.13, 29.7 - BIP-compliant, 15.11, 15.TA - by agent / community /
institution / algorithm, 15.4 - chain, 15.2 - deferred, \textbf{15.8} -
information loss in, \textbf{15.6} - lexicographic, \textbf{15.5} -
maximin (Rawlsian), \textbf{15.5}, 7.4 - mismatch (alignment problem
as), \textbf{18.5} - necessity of, 15.3 - non-commutativity of, 15.2 -
opaque vs.~transparent, 15.10 - phenomenology of, 15.9 - probabilistic,
15.5 - satisficing, 15.5 - summative, \textbf{15.5} - weighted, 15.5

coordinate chart, \textbf{4.1}, 5.4, 5.TA

correlative symmetry (Hohfeldian), 3.7, 8.4, 12.8, 17.3 - empirical
rates (87\%, 82\%), 3.7, 8.4, 9.7, 17.3

cost-benefit analysis, 1.3.2

cotangent bundle, \textbf{4.3}

cotangent space, \emph{see} dual space

covariance tensor (moral), \textbf{6.6}, 2.5, 16.3, 29.2 - principal
directions of uncertainty, 6.6

covariant derivative, \textbf{4.6}, \textbf{10.3}

critical slowing down (moral analogue), 8.7

cross-lingual invariance, 1.6, 12.8, \textbf{17.3}; \emph{see also} BIP
experiments

curvature (moral), \textbf{10.6}, 1.8, 29.2, 29.2 - direct measurement
(open problem), \textbf{29.2} - Ricci tensor (moral), 10.6, 29.4 -
scalar curvature (moral), 10.6 - sources of, 10.6 - \emph{see also}
Riemann curvature tensor

cusp singularity, 8.6

\subsection{D}\label{d}

D\(_{4}\) dihedral group, \textbf{8.4}, 1.7, 3.7, 5.9, 12.8, 13.8, 17.3
- anomaly, 17.3, 29.5 - group relations, 8.4 - Hohfeldian transitions
as, 8.4

Dear Abby corpus, \textbf{17.2}, 1.6, 2.5, 5.3, 5.6, 8.3, 8.4, 9.7, 10.9
- context-dependent weighting, 17.2 - nullifier data, 5.6, 8.3, 17.2 -
semantic gate data, 8.4, 17.2 - temporal stability, 10.9, 17.2

Dear Ethicist game, \textbf{17.4}

decoherence (moral), \textbf{13.6} - channel, 13.6 - Kraus operators,
13.6

degenerate metric, \textbf{4.5}, 1.3.2, 6.5, 7.5, 8.6, 29.3, 29.4 - as
incommensurability, 3.6, 7.5, 29.3, 29.4

deferred contraction, \emph{see} contraction, deferred

deliberation, as quantum state, 13.3

DEME architecture, \textbf{19.6}, 1.7, 19.1, 29.4 - audit and monitoring
(Layer 4), 19.6 - democratic governance (Layer 3), 19.6 - DEMEProfile,
19.6, 19.9 - EthicalFacts (Layer 1), 19.6 - Ethics Module (Layer 2),
19.6 - four-layer architecture, \textbf{19.6}

democratic governance, as metric selection, \textbf{19.9}; \emph{see
also} governance account

density matrix (moral), \textbf{13.6}, 13.3 - mixed state, 13.6 - pure
state, 13.6 - purity, 13.6

deontic structure, 12.8, 17.3 - 100\% cross-lingual transfer, 17.3

difference principle (Rawls), 3.5

differential geometry, 1.1, \textbf{4.1}--\textbf{4.11}

dimensions of moral space, \emph{see} nine dimensions

directional information, 1.2, 2.2.1, 15.6, 29.7

discount rate, as moral metric choice, 10.9

doctrine of the mean (Aristotle), \textbf{3.2} - as context-sensitive
calibration, 3.2 - as metric condition, 3.2

dominance relations, 7.3

dual space, \textbf{4.3}

\subsection{E}\label{e}

edge modes, 13.TA

Einstein summation convention, 4.2

Einstein's field equation, 10.6, \textbf{29.4} - moral analogue, 10.6,
29.4

electromagnetic field tensor, 12.3

emergence (moral), 14.3, 14.5, 14.10 - criterion for emergent
obligation, 14.5 - criterion for emergent moral property, 14.10 - degree
of (\(\eta\)), 14.3 - structural injustice as, 14.5

entanglement (moral), \textbf{13.8} - entropy, 13.TA - EPR states
(moral), 13.8 - partial trace, 13.8

epistemic modes (What Matters, Who Decides, What We Know), \textbf{5.3}

epistemic stance, \emph{see} pragmatist epistemology

\textbf{epistemic status tags}, Preface, 17.0.2, 17.0.5 - four tags
(Definition, Modeling choice, Conjecture, Open question), Preface - as
discipline standard, 17.0.2

ErisML, \textbf{19.3}, 1.9, 19.1 - compilation targets, \textbf{19.4} -
conflict resolution function, 19.3 - environment specification, 19.3 -
formal semantics (NCSG), \textbf{19.3} - fragments (Core, Multi,
Temporal, Full), 19.4 - norms specification, 19.3 - vector objective,
19.3

ethical scopes (Individual, Relational, Collective), \textbf{5.3}

EthicalFacts, 19.6

Ethics Module (EM), 19.6

EU AI Act, 19.5, 19.9

Euler-Lagrange equations, 4.9, 12.2, 12.4

evaluation tensor, \textbf{6.6} - multi-agent, 6.6

expected utility, 1.1

expressivism (about the metric), \textbf{9.4} - agent-indexed metric,
9.4 - Frege-Geach problem, 9.4

external verification, strong (Requirement 4), \textbf{18.6}

\subsection{F}\label{f}

falsifiability, 5.3, \textbf{17.7}, \textbf{29.13} - six criteria, 29.13

feasible cone, 8.6

fiber bundle, \textbf{4.7}, 1.1, 3.2, 12.3, 19.8 - agent fiber bundle,
\textbf{5.4} - base space, 4.7 - character bundle, 6.7 - frame bundle,
4.7 - local triviality, 4.7 - Mobius strip (example), 4.7, 29.3 -
projection, 4.7 - safety stack as, 19.10 - section of, \textbf{4.7}, 3.2
- separation principle as, 19.8 - total space, 4.7 - typical fiber, 4.7

fidelity classification (faithful/approximate/indicative), \textbf{19.5}

Fisher-Rao metric, 29.10

forbidden region, \emph{see} constraint surface

frame bundle, 4.7

Frege-Geach problem, 9.4

Frobenius norm, 29.10

frontier condition, \textbf{4.8}, 8.2

functionings (Sen), 3.6

fundamental contraction, \emph{see} contraction; satisfaction

fundamental group (\(\pi_{1}\)), 5.8, 29.3

\textbf{fuzz testing (invariance fuzzing)}, 17.0.4, 17.0.5, 18.4 - as
verification methodology, 17.0.4 - scenario generators, 17.0.5 -
transformation suites, 18.4; \emph{see also} invariance testing

\subsection{G}\label{g}

gauge field (connection 1-form), \textbf{4.9}

gauge-fixing, 18.6; \emph{see also} canonicalization

gauge invariance, \textbf{4.9}, 12.3, 18.4; \emph{see also} Bond
Invariance Principle

gauge symmetry, 4.9, 12.3, 29.8

gauge theory, \textbf{4.9}, 1.8, 12.3, 18.8

gauge transformation, \textbf{4.9}, 12.3

Geneva Ethics Module (Geneva EM), \textbf{19.6}

geodesic (moral), \textbf{10.7} - deviation equation, 10.7, 10.TA, 29.2
- equation, \textbf{10.7} - moral inertia, 10.7

Godel's incompleteness theorem, 1.5

Goodhart's Law, 18.2, 18.7

governance account (of the moral metric), \textbf{9.5}, 19.9, 29.3 -
democratic determination, 9.5 - legal analogy, 9.5 - metric as
governance artifact, 9.5 - moral constitutionalism, 9.5 - moral reform
as metric revision, 9.5

gradient flow, \textbf{10.8} - local maxima (moral complacency), 10.8 -
saddle points (moral impasse), 10.8

Grand Unified AI Safety Stack (GUASS), \textbf{19.10} - seven-layer
architecture, 19.10

grounded evaluation (Requirement 2), \textbf{18.6}, 29.8 -
domain-specific grounding, 29.8

\textbf{grounding tensor (Ψ)}, 17.0.4, 17.0.5, 18.6, 29.8 - adequacy
properties, 18.6, 29.8 - as map from observables to moral dimensions,
17.0.4 - construct validity, 17.0.4; \emph{see also} grounded evaluation

\subsection{H}\label{h}

Habermas, Jurgen, 9.3

harm, \emph{see} invariance of harm accounting

harm ledger, \textbf{12.6}

Hausdorff condition, 4.1

Hilbert space (moral), \textbf{13.2}

Hohfeld, Wesley Newcomb, \textbf{3.7}, 8.4 - \emph{Fundamental Legal
Conceptions}, 3.7

Hohfeldian jural positions, \textbf{3.7}, 1.7, 5.9, 8.4, 12.8, 13.4,
17.3 - Claim, 3.7 - correlatives and opposites, 3.7, 8.4 - D\(_{4}\)
group structure, \textbf{8.4} - Duty (Obligation), 3.7 -
Immunity-Disability, 3.7 - Liberty, 3.7 - No-claim, 3.7 - operator
(quantum), 13.4 - Power-Liability, 3.7

holonomy, \textbf{10.5}, 2.5, 4.6, 8.8, 14.8, 29.2, 29.2 - collective,
14.8 - detects curvature, 10.TA - group, \textbf{10.5} - moral residue
as, 10.10, 15.7 - non-abelian, 10.5 - path-dependence and, 10.5 -
structural, 14.8

hospital ethics committee (example), \textbf{14.12}

Huainanzi, 29.1; \emph{see also} parable of the old man and the horse

human-AI collaboration, \textbf{18.12}

\subsection{I}\label{i}

\textbf{I-EIP (Internal Epistemic Invariance Principle)}, 18.9, 17.0.6 -
I-EIP Monitor, 18.9 - equivariance of internal representations, 18.9 -
drift detection, 17.0.6, 18.9; \emph{see also} invariance testing;
structural containment

ideal discourse (Habermas), 9.3

impact vector, 2.2.1

incommensurability, 1.3.2, 3.6, 7.5, 8.6, 29.3 - as degenerate metric,
7.5, 29.3, 29.4 - as metric singularity, 8.6

incompleteness thesis (Sen), \textbf{3.6}

index raising and lowering, 6.5; \emph{see also} musical isomorphism

information loss, \emph{see} contraction, information loss in

institutional manifold, \textbf{14.7}

interest covector, \textbf{6.3}, 2.6, 6.4, 29.7 - egalitarian, 6.3 -
libertarian, 6.3 - welfarist, 6.3

interference (moral), \textbf{13.5} - constructive, 13.5 - destructive,
13.5 - framing alignment (relative phase), 13.5 - mercy vs.~justice
(example), 13.5

intertheoretic comparison problem, \textbf{16.5}, 29.10 - meta-metric on
theory space, 16.5, \textbf{29.10}

invariance loss, 18.10

invariance testing, \textbf{18.4}, 12.9, 17.10, 29.7 - as alignment
auditing, 18.4 - transformation suites, 18.4 - violation decomposition,
18.TA

\subsection{J}\label{j}

Jacobian, 1.7, \textbf{4.2}

judgment (moral), 1.5, 10.8, 15.9, 29.9 - framework does not replace,
1.5, 9.5, 29.9

junction conditions, \textbf{13.7}

jural relations, \emph{see} Hohfeldian jural positions

\subsection{K}\label{k}

Kant, Immanuel, \textbf{3.3} - \emph{Groundwork of the Metaphysics of
Morals}, 3.3

Keats, John, 15.8

kidney allocation (case study), \textbf{7.1}--\textbf{7.9}, 5.9, 6.8,
15.12, 18.11 - five levels, 7.2--7.6 - revisited (contraction),
\textbf{15.12} - tensorial (AI), \textbf{18.11}

Killing vector field, 12.TA

Kingdom of Ends, 3.3

Kraus operators, 13.6

\subsection{L}\label{l}

LaBSE model, 17.7

Lagrangian (moral), \textbf{12.4} - kinetic term, 12.4 - moral
potential, 12.4 - stratified, \textbf{13.7}

Laplace-Beltrami operator, 13.7

lattice (of admissible metrics), 9.8

Levi-Civita connection, \textbf{4.6}, 10.3, 29.6

lexicographic contraction, \emph{see} contraction, lexicographic

lexicographic metric, 6.5, 7.5

Lie algebra, 4.9, 12.6

Lie group, \emph{see} structure group

local triviality, 4.7

Lorentzian signature, 29.3

\subsection{M}\label{m}

manifold, \textbf{4.1}; \emph{see also} moral manifold - smooth,
\textbf{4.1} - topological, \textbf{4.1} - with boundary, 4.1

maximin, \emph{see} contraction, maximin; Rawls, John

Maximize Expected Choiceworthiness (MEC), 16.4

``maybe'' (the geometry of), \textbf{29.7} - as deferred contraction,
29.7 - five meanings, 29.7 - \emph{see also} parable of the old man and
the horse

measurement (moral), \textbf{13.4} - as decision/collapse, 13.4 -
postulate (moral form), 13.4

meta-metric, 9.8, 16.5, \textbf{29.10} - existence, 20.TA - Fisher-Rao,
29.10

metric, \emph{see} moral metric; metric tensor

metric tensor, \textbf{4.5}, 1.3.2 - degenerate, \textbf{4.5}, 6.5 -
musical isomorphism, \textbf{4.5} - non-degenerate, 4.5 -
pseudo-Riemannian, 4.5 - Riemannian, \textbf{4.5} - signature,
\textbf{29.3} - symmetric, 4.5

metric uncertainty, 16.2, 16.3

minimax regret, 16.TA

Misak, Cheryl, 1.6

Mobius strip, 4.7, 29.3

Model Context Protocol (MCP), 19.10

moral connection, \emph{see} connection (moral)

moral debt, 12.7, 17.6

moral dynamics, \textbf{10.1}--\textbf{10.12}; \emph{see also} parallel
transport; holonomy; curvature

moral energy landscape, 12.4

moral field equation (open problem), \textbf{29.4}, 10.6 - moral
stress-energy tensor, 10.6, 29.4

moral Hamiltonian, \textbf{13.7}

moral Hilbert space, \emph{see} Hilbert space (moral)

moral inertia, 10.7

moral manifold (\(\mathcal{M}\)), \textbf{5.1}--\textbf{5.10}, 1.7, 29.3
- compactness, 5.8, 29.3 - connectedness, 5.8 - dimension, \textbf{5.3}
- formal definition, 5.TA - nine dimensions, \textbf{5.3} - points of,
5.2 - stratification of, \textbf{5.6} - topology, 5.8, \textbf{29.3}

moral metric (\(g_{\mu\nu}\)), \textbf{6.5}, 2.6,
\textbf{9.1}--\textbf{9.10}, 29.2 - admissible metrics, \textbf{9.6} -
as governance artifact, \textbf{9.5} - as moral infrastructure, 9.9 - as
trade-off structure, 2.6, 6.5 - constructivist account, \textbf{9.3} -
context-dependent weighting, 5.3, 17.2 - diagonal and off-diagonal
components, 6.5, 7.5 - expressivist account, \textbf{9.4} - governance
account, \textbf{9.5} - learning from data, \textbf{9.7} - not uniquely
determined, 29.7 - realist account, \textbf{9.2} - signature (open
problem), \textbf{29.3} - temporal stability, 10.9, 17.2

moral momentum, \textbf{12.5}

moral observable, \textbf{13.4} - non-commuting, 13.4 - satisfaction
operator, 13.4 - stratum projector, 13.4

moral phase transition, \emph{see} phase transition

moral progress, as parallel transport with holonomy, 10.5, 29.4

moral realism, 1.5, 9.2, 29.9

moral residue (\(\mathcal{R}\)), \textbf{15.7}, 3.4, 6.4, 10.10, 18.3,
29.7 - acknowledgment, 15.7 - as holonomy, 10.10, 15.7 - as information
discarded by contraction, \textbf{15.7}, 29.4 - compensation, 15.7 -
inevitability of, \textbf{15.7} - memory, 15.7 - regret, 15.7 - residue
logging (AI), 18.3

moral risk, \textbf{6.6}, 16.2, 29.2

moral Schrodinger equation, 13.7

moral singularity, \textbf{5.7}, 8.6 - Type 1 (metric degeneracy), 8.6 -
Type 2 (gradient failure), 8.6 - Type 3 (constraint intersection), 8.6

moral stress-energy tensor, 10.6, 29.4

moral tunneling, \textbf{13.7}

moral uncertainty, \textbf{16.1}--\textbf{16.11}, 3.8 - decision under,
\textbf{16.4} - empirical (Type 1), \textbf{16.2} - metric (Type 2),
\textbf{16.2} - principle, \textbf{13.4} - theory (Type 3),
\textbf{16.2} - three types, \textbf{16.2}

moral vector, \emph{see} obligation vector

musical isomorphism, \textbf{4.5}

My Favorite Theory (MFT), 16.4

\subsection{N}\label{n}

\textbf{Norm Kernel (NK)}, 19.8, 18.9 - as contract automaton, 19.8 -
compilation from Semantic Normative Host, 19.8 - integration with I-EIP
Monitor, 18.9

NCSG, \emph{see} Norm-Constrained Stochastic Game

negative capability (Keats), 15.8

nine dimensions (of moral space), \textbf{5.3} - Dimension 1:
Consequences and Welfare, 5.3 - Dimension 2: Rights and Duties, 5.3 -
Dimension 3: Justice and Fairness, 5.3 - Dimension 4: Autonomy and
Agency, 5.3 - Dimension 5: Privacy and Data Governance, 5.3 - Dimension
6: Societal and Environmental Impact, 5.3 - Dimension 7: Virtue and
Care, 5.3 - Dimension 8: Procedural Legitimacy, 5.3 - Dimension 9:
Epistemic Status, 5.3 - necessity and sufficiency arguments, 5.3

NIST AI RMF, 19.9

No Escape Theorem, \textbf{18.6}, 1.3.1, 1.8, 12.9, 19.10, 29.4, 29.6 -
conditional guarantee, 18.TA - four requirements, \textbf{18.6} -
limitations, \textbf{18.7} - proof, 18.6 - scope, 18.7 - \emph{see also}
structural containment

Noether current, 12.5

Noether's theorem, \textbf{12.1}--\textbf{12.11}, 1.7, 4.9 - for ethics
(Theorem 13.1), \textbf{12.5} - moral analogue, \textbf{12.5} - physics
version, \textbf{12.2}

Norm-Constrained Stochastic Game (NCSG), \textbf{19.3}, 19.TA -
well-posedness, 19.TA

normal bundle, 8.5

null directions, 29.3

nullifiers, \textbf{5.6}, 8.3, 17.2 - abuse, 5.6, 8.3, 17.2 - danger,
5.6, 8.3 - impossibility, 5.6, 8.3 - universality of, 17.2

Nussbaum, Martha, \textbf{3.6}, 8.3 - capability thresholds,
\textbf{3.6}, 8.3

\subsection{O}\label{o}

obligation stratum (\(S_{O}\)), 5.9, 8.9

obligation vector (\(O^{\mu}\)), \textbf{6.2}, 2.6, 29.2 - as vector
field, 1.7, 6.2 - nine components, 6.2 - parallel transport of,
\textbf{10.4} - zero vector, 6.2

observational chaos, 19.2

old man and the horse, \emph{see} parable of the old man and the horse

open problems, \textbf{29.1}--\textbf{29.14} - curvature measurement,
\textbf{29.2} - institutional design, \textbf{29.11} - intertheoretic
comparison, \textbf{29.10} - metaphysics of moral geometry,
\textbf{29.9} - metric signature and topology, \textbf{29.3} - moral
field equation, \textbf{29.4} - scalability of containment,
\textbf{29.8} - symmetries and conservation, \textbf{29.5} - tensorial
interpretability, \textbf{29.7} - torsion, \textbf{29.6}

Orch-OR hypothesis (Penrose and Hameroff), 1.5, 1.8, 13.1

order effects, 13.5, 17.5

original position (Rawls), \textbf{3.5} - as symmetry condition, 3.5

orthogonal values, 2.6, 6.5

``ought implies can,'' 5.6, 8.3

\subsection{P}\label{p}

parable of the old man and the horse (Sai Weng Shi Ma), \textbf{1.2},
\textbf{2.1}, 29.1, 29.2, 29.7 - geometric interpretation,
\textbf{2.1}--\textbf{2.8}, \textbf{29.2}

parallel transport, \textbf{4.6}, \textbf{10.4}, 10.5, 14.8, 29.2 -
collective, \textbf{14.8} - faithful maintenance, 10.4 - of obligations
across contexts, \textbf{10.4} - preserves satisfaction, 10.TA

partial order on metrics, \textbf{9.8}

path-dependence, 2.2.3, \textbf{10.5}, 17.4, 29.7; \emph{see also}
holonomy

Penrose formula (\(\tau \approx \hslash/E_{G}\)), 1.8, 13.1

penumbral zone, \textbf{8.5}

perspective shifts, \emph{see} admissible transformations, Type 2

phase transition (moral), \textbf{8.7}, 2.5, 8.3 - emergency, 8.3, 7.6 -
first-order, \textbf{8.3} - higher-order, 8.3 - second-order,
\textbf{8.3} - wartime, 8.3

phronesis (practical wisdom), 15.9

policy DAG, \textbf{19.5}

power tensor (\(P_{\mu\nu}\)), 6.9, 14.4 - antisymmetric part
(structural inequity), 6.9, 14.4

pragmatist epistemology, \textbf{1.6}, Preface, 9.5

Price, Huw, 1.6

prima facie duties (Ross), \textbf{3.4}, 1.3.3 - as vector components,
3.4

prioritarian metric, 3.5, 6.7

promise landscape (worked example), \textbf{8.9}, 5.9

\subsection{Q}\label{q}

quantum cognition, \textbf{13.5}, 17.5 - conjunction/disjunction errors,
13.5 - order effects, 13.5, 17.5 - violation of law of total
probability, 13.5

quantum normative dynamics, \textbf{13.1}--\textbf{13.11}, 1.8

\subsection{R}\label{r}

rank (of a tensor), \textbf{4.4}, 6.1

Rawls, John, \textbf{3.5}, 1.3.3, 7.4, 9.3, 15.5 - \emph{A Theory of
Justice}, 3.5 - difference principle, \textbf{3.5} - maximin criterion,
\textbf{3.5}, 7.4, 15.5 - original position, \textbf{3.5} - supremum
metric (\(L^{\infty}\)), \textbf{3.5} - veil of ignorance, \textbf{3.5},
9.3

re-description invariance, \emph{see} Bond Invariance Principle

realist account (of the metric), \textbf{9.2}

reflective equilibrium, 9.3

regime-dependent evaluation, 1.2, 2.2.3, 29.7; \emph{see also}
stratification

relabeling attack, 18.9

responsibility tensor, \textbf{14.6} - causal contribution, 14.6 -
collective responsibility remainder, 14.6 - complicity gradient, 14.6 -
epistemic position, 14.6 - responsibility gap/surplus, 14.6 - role
authority, 14.6

reward hacking, 1.3.1, 15.11, 18.2, 18.9

Ricci tensor, \textbf{4.6}, 10.6, 29.4

Riemann curvature tensor, \textbf{4.6}, \textbf{10.6}, 29.2 - moral
interpretation, 10.6 - vanishes iff flat, 4.6

Riemannian manifold, \textbf{4.5}

Riemannian metric, \emph{see} metric tensor, Riemannian

robust core (\(\mathcal{C}_{\text{robust}}\)), \textbf{16.6}, 18.10,
19.6 - as convex cone, 16.6 - Geneva EM as implementation of, 19.6 -
stability under perturbation, 16.6

robust obligation, \textbf{16.6}

Robert's Rules, 13.6

Ross, W. D., \textbf{3.4}, 1.3.3 - \emph{The Right and the Good}, 3.4

\subsection{S}\label{s}

Sai Weng Shi Ma, \emph{see} parable of the old man and the horse

\textbf{Safety Reduction}, 18.9.5, 17.0.6 - three tractable categories
(governance, engineering, security), 18.9.5; \emph{see also} structural
containment

satisfaction (\(S = I_{\mu}O^{\mu}\)), \textbf{6.4}, 2.6, 12.4, 29.3 -
fundamental formula, \textbf{6.4} - operator, 13.4

Saul on the road to Damascus, 13.7

scalar curvature (moral), 10.6

scalar evaluation, \textbf{1.1}, \textbf{2.2}, 18.2, 29.8 - inadequacy
of, \textbf{2.2} - three limitations, 2.2

Scanlon, T. M., 9.3

Schrodinger equation, \emph{see} moral Schrodinger equation

semantic gate, \textbf{8.4}, 2.5, 5.6, 17.2 - ``life-threatening''
(emergency), 7.6 - ``only if convenient,'' 5.6, 8.4 - ``you promised,''
5.6, 8.4 - as D\(_{4}\) group element, 8.4 - compositional structure,
8.4 - cross-linguistic evidence, 8.4, 17.3 - detection results, 8.4 -
discreteness, 8.4, 17.2

Sen, Amartya, \textbf{3.6}, 1.3.3 - capabilities approach, \textbf{3.6}
- incompleteness thesis, \textbf{3.6}

separation principle, \textbf{19.8} - as fiber-bundle structure, 19.8 -
composability, 19.8

signature, \emph{see} metric tensor, signature

singularity, \emph{see} moral singularity

smooth manifold, \emph{see} manifold, smooth

Sophie's Choice, 5.7, 8.6, 12.4

specification gaming, 1.3.1, 15.11, 18.2, 18.9

spontaneous symmetry breaking, 29.5

stratification, \textbf{5.6}, \textbf{8.1}--\textbf{8.11}, 2.5, 29.2,
29.3 - as patchwork structure, 8.1 - moral theories and, 8.10 - Whitney,
\textbf{4.8}, \textbf{8.2}

Stratified Quantum Normative Dynamics (SQND), 1.7, Preface

stratum boundary, \textbf{8.3}, 5.6, 29.2, 29.4 - Type I (threshold),
\textbf{8.3} - Type II (phase transition), \textbf{8.3} - Type III
(absorbing / nullifier), \textbf{8.3} - Type IV (constraint surface),
\textbf{8.3}

structural containment, \textbf{18.6}, 1.3.1, 29.4 - architecture,
\textbf{18.6} - four requirements, \textbf{18.6} - scalability (open
problem), \textbf{29.8} - \emph{see also} No Escape Theorem

structural injustice, 14.5

structure group (Lie group), 4.9

structure tensor (\(\Sigma_{ab}\)), \textbf{14.4} - antisymmetric part
(power differential), 14.4 - symmetric part, 14.4

structured pluralism, \textbf{9.8}, 16.7

structured uncertainty, 1.2, 2.2.2; \emph{see also} covariance tensor

superposition (moral), \textbf{13.3}, 3.8 - vs.~mixture, 13.3

supremum metric, \emph{see} Rawls, John, supremum metric

symmetrization, \textbf{6.9}

symmetry, \emph{see} gauge symmetry; Noether's theorem; D\(_{4}\)
dihedral group - additional (open problem), \textbf{29.5}

\subsection{T}\label{t}

tangent bundle, \textbf{4.2}

tangent space (\(T_{p}M\)), \textbf{4.2}

tangent vector, \textbf{4.2}, 2.2.1

temporal discounting, 10.9

temporal stability (moral structure), 10.9, 17.2

tensor, \textbf{4.4} - contraction of, \textbf{4.4}; \emph{see also}
contraction - field, 4.4 - rank, \textbf{4.4} - transformation law,
\textbf{4.4} - type (r,s), \textbf{4.4}

tensor hierarchy, \textbf{6.1}--\textbf{6.10}, 29.3

tensor product, \textbf{6.9}, 13.8, 14.2

tensor-valued objectives (AI), \textbf{18.3} - vector-valued reward,
18.3 - vector value function, 18.3

tensorial interpretability (open problem), \textbf{29.7} - contraction
identification, 29.7 - mechanistic interpretability, 29.7 - probe
extraction, 29.7

theory covariance tensor, \textbf{16.3}

theory uncertainty, \emph{see} moral uncertainty, theory (Type 3)

tidal forces (moral analogue), 10.6

topology (of moral manifold), \emph{see} moral manifold, topology

torsion, 10.3, \textbf{29.6} - closure failure, 29.6 - detection
criterion, 20.TA - tensor (\(T_{\alpha\beta}^{\mu}\)), 29.6

tragic dilemma, 8.6

trajectory-dependence, \emph{see} path-dependence

transition function, \textbf{4.1}, 8.8

transition map, \textbf{8.8}

translation layer, \textbf{19.5} - fidelity classification,
\textbf{19.5} - loss documentation, 19.5 - policy DAGs, \textbf{19.5}

transparency, of contraction, 15.10, 18.3

trolley problem, 12.4

tubular neighborhood, 8.2

tunneling, \emph{see} moral tunneling

\subsection{U}\label{u}

uncertainty, \emph{see} moral uncertainty; covariance tensor

uncertainty principle (moral), \textbf{13.4}

utilitarianism - as contraction, 6.7, \textbf{15.5} - as Euclidean
metric, 2.6, 3.5 - as scalar ethics, 1.1, 2.2 - metric (\(L^{1}\)), 3.5

\subsection{V}\label{v}

value collapse, 18.2

vector field, \textbf{4.2}

veil of ignorance, \emph{see} Rawls, John, veil of ignorance

virtue, as non-gradient agency, 10.8

virtue ethics, 6.7, 8.10, 16.2 - as fiber-bundle section, 1.5, 3.2, 6.7

\subsection{W}\label{w}

welfare function, 1.1

whistleblower's dilemma (worked example), \textbf{13.10}

Whitney, Hassler, \textbf{4.8}, 8.2 - condition (A), \textbf{4.8}, 8.2 -
condition (B), \textbf{4.8}, \textbf{8.2} - stratification,
\textbf{4.8}, \textbf{8.2}, 5.6, 13.TA

\subsection{Y}\label{y}

Yang-Mills gauge theory, 3.1

Young, Iris Marion, 14.5

\subsection{Notation Index}\label{notation-index}

\emph{Page references omitted; notation is listed with the section where
it is introduced.}

{\def\LTcaptype{none} % do not increment counter
\begin{longtable}[]{@{}
  >{\raggedright\arraybackslash}p{(\linewidth - 4\tabcolsep) * \real{0.1541}}
  >{\raggedright\arraybackslash}p{(\linewidth - 4\tabcolsep) * \real{0.4711}}
  >{\raggedright\arraybackslash}p{(\linewidth - 4\tabcolsep) * \real{0.1444}}@{}}
\toprule\noalign{}
\begin{minipage}[b]{\linewidth}\raggedright
\textbf{Symbol}
\end{minipage} & \begin{minipage}[b]{\linewidth}\raggedright
\textbf{Meaning}
\end{minipage} & \begin{minipage}[b]{\linewidth}\raggedright
\textbf{Introduced}
\end{minipage} \\
\midrule\noalign{}
\endhead
\bottomrule\noalign{}
\endlastfoot
M & Moral manifold & 5.1 \\
T\_p M & Tangent space at \(p\) & 4.2 \\
T*\_p M & Cotangent space at \(p\) & 4.3 \\
O\^{}μ & Obligation vector & 6.2 \\
I\_μ & Interest covector & 6.3 \\
S = I\_μ O\^{}μ & Satisfaction (fundamental contraction) & 6.4 \\
g\_\{μν\} & Moral metric tensor & 6.5 \\
Σ\^{}\{μν\} & Moral covariance tensor (uncertainty) & 6.6 \\
E\^{}μ\_ν & Evaluation tensor & 6.6 \\
M\_\{ia\} & Multi-agent evaluation tensor & 7.4 \\
P\_\{μν\} & Power tensor & 6.9 \\
σ\^{}μ & Character profile (virtue section) & 6.7 \\
σ²\_S & Moral risk & 6.6 \\
∇ & Moral connection (covariant derivative) & 10.3 \\
Γ\^{}μ\_\{νρ\} & Christoffel symbols (connection coefficients) & ,
10.3 \\
R\^{}μ\_\{ναβ\} & Riemann curvature tensor & , 10.6 \\
R\_\{μν\} & Ricci tensor & , 10.6 \\
R & Scalar curvature & , 10.6 \\
Hol(∇, p) & Holonomy group at \(p\) & 10.5 \\
L & Moral Lagrangian & 11.4 \\
A & Moral action functional & 11.4 \\
H & Harm (Noether charge) & 11.6 \\
p\_μ & Moral momentum & 11.5 \\
ξ\^{}μ & Re-description generator & 11.5 \\
G & Gauge group of the moral manifold & 11.3 \\
\textbar ψ⟩ & Moral quantum state & 12.2 \\
H & Moral Hilbert space & 12.2 \\
ρ & Moral density matrix & 12.6 \\
Ŝ & Satisfaction operator & 12.4 \\
Ô\^{}μ & Obligation operator & 12.4 \\
Π̂\_α & Stratum projector & 12.4 \\
Ĵ & Hohfeldian operator & 12.4 \\
Ĥ & Moral Hamiltonian & 12.7 \\
E & Decoherence channel & 12.6 \\
Φ\_\{αβ\} & Boundary cost (activation energy) & 12.7 \\
C & Collective agency tensor & 13.3 \\
E & Emergent component (collective) & 13.3 \\
η & Degree of emergence & 13.3 \\
Σ\_\{ab\} & Structure tensor & 13.4 \\
R\_a(ω) & Responsibility tensor & 13.6 \\
Λ\^{}μ\_a & Structural coupling & 13.8 \\
M\_inst & Institutional manifold & 13.7 \\
R & Moral residue & 14.7 \\
L(C) & Contraction loss & 14.6 \\
τ & Threshold vector (moral minimum) & 14.5 \\
Θ\_\{jk\} & Theory covariance tensor & 15.3 \\
C\_robust & Robust core & 15.6 \\
r\^{}μ(s,a) & Vector-valued reward & 17.3 \\
V\^{}μ(s) & Vector value function & 17.3 \\
IV(x, G) & Invariance violation & 17.4 \\
Δ & Alignment gap & 17.5 \\
C(·) & Canonicalizer & 17.6 \\
Ψ(·) & Grounding function & 17.6 \\
Bd & Bond Index & 18.7 \\
D\_op & Observed defect & 18.7 \\
T\^{}μ\_\{αβ\} & Torsion tensor & 19.6 \\
\end{longtable}
}
